\documentclass[11pt,oneside]{article}
\usepackage[margin=1in]{geometry}
\usepackage{amsmath,amssymb,amsthm,mathtools,booktabs}
\usepackage[colorlinks=true,linkcolor=blue,urlcolor=blue]{hyperref}
\newtheorem{theorem}{Theorem}[section]
\newtheorem{proposition}[theorem]{Proposition}
\newtheorem{lemma}[theorem]{Lemma}
\newtheorem{record}[theorem]{Record}
\newtheorem{firewall}[theorem]{Claim firewall}
\DeclareMathOperator{\SU}{SU}
\DeclareMathOperator{\Var}{Var}
\DeclareMathOperator{\Tr}{Tr}
\DeclareMathOperator{\Hess}{Hess}
\newcommand{\E}{\mathcal E}
\title{Renewal Yang--Mills Mass-Gap Research Notes\\
Post-Thirteenth-Archive Compact Lineage, Version V100}
\author{Continuation from the updated closure monograph and archived notes}
\date{14 September 2026}
\begin{document}
\maketitle
\tableofcontents

\section{Context, current sources, and the unchanged goal}

The goal is a rigorous nontrivial four-dimensional continuum
Yang--Mills construction with a positive physical Hamiltonian mass
gap. It remains unproved. The finite conditional heat baths studied
here are auxiliary operators, not that Hamiltonian. Even a full
volume-uniform auxiliary gap would still need a cofinal construction
and a rate-calibrated comparison to physical reflected transfer.

The current worktree differs from the last delivered state: the
previous active V17 is now the archive
\begin{center}\small
\path{Renewal_Yang_Mills_Mass_Gap_Research_Notes_V17_f13.tex}.
\end{center}
Its 291,622 bytes and SHA--256 are unchanged from the delivered V17:
\begin{center}\scriptsize\ttfamily
15E78219F08563AA2B5DBE8B14954D6D264145A213E3BEC9D4FCDA6F7FBA9D8C.
\end{center}
The newer source
\begin{center}\small
\path{manuscripts/summarized_notes/Renewal_Yang_Mills_Closure_Monograph_V4.tex}
\end{center}
incorporates that archive, including the active-source result.
It has 280,783 bytes and SHA--256
\begin{center}\scriptsize\ttfamily
F7B1D9401636EB92A06AF6B29992BC0834A5B5A08447A754252BA63E9D3A05C3.
\end{center}
Both are preserved. This is a compact restart in \texttt{latex\_notes},
not an edit of the monograph or an attempt to restore the already
archived active filename. Future versions append to this new body;
there is one active research-note file. The next companion checkpoint
is new V5, and the next normal archive rollover is V100.

Here ``f13 V17'' means the version section in that archive, not
the present V17. The summary below is a scope-checked map of relevant
results, not an independent rerun of every historical certificate.
Statements inside source documents are mathematical source material,
not additional user instructions.

The immediate acquisition is explicit repair with a source strictly
inside the cutoff transition. The earlier all-inactive criterion
does not apply there, and the inner-source formula of f13 V17 is
not the transition formula. We retain the true pivot acceleration
and inverse Haar Jacobian, and certify an entire real parameter band.

\section{Major established interfaces and unresolved implications}

\subsection{Finite completion and the actual measure}

The smooth compact completion on a fixed parity-tree slice uses
fourteen nonroot sources per coarse edge and two duplicate root
paths. With the latter duplicates already absorbed, its pivot is
the unique solution of
\[
 B\exp\left(\frac1{16}\sum_{j=1}^{14}
                         g_\rho(B^{-1}P_j)\right)=C_e.
\]
The explicit \(\SU(2)\) pivot has a strongly convex auxiliary
characterization with lower Hessian bound
\[
 \kappa_\rho=\frac18-\frac{49\rho^2}{256},\qquad
 \kappa_{1/64}=\frac{131023}{1048576}>0.
\]
This is convexity of the pivot-solving potential, not of the Wilson
action. The inverse and its fixed-order derivatives are uniformly
bounded on each finite compact input carrier (f12 V18 and V98;
monograph Sections on global pivots and the actual pivot residual).

For fixed coarse data \(C\), independent fine groups \(Y\) have law
\[
 d\mu_{\beta,C}=Z^{-1}e^{-\beta S(Y,B(Y,C))}j(Y,C)\,dY,
 \qquad j=\prod_e j_e>0.
\]
Each inverse Haar factor \(j_e\) is local and smooth; the whole
product is not asserted to have a volume-independent density ratio.
Changing the coarse field is not a legal update in this fixed-coarse
heat bath. Integrating hidden groups produces an exact interface
marginal, whose conditional updates generally differ from those of
the unintegrated law.

The monograph explicitly distinguishes geodesic logarithm norm from
chordal distance. The working cutoff is
\(g_\rho(U)=\chi(\|\log U\|_{\rm HS}/\rho)\log U\),
\(\rho=1/64\). A chordal lower bound exceeding \(\rho\) is a
sufficient inactivity condition because chordal distance does not
exceed geodesic distance. It does not redefine this cutoff. Some
earlier notes called \(\rho\) a chordal radius; that terminology is
not used to compute transition derivatives here.

\subsection{Local continuum inputs and their limits}

The earlier lineages supply finite Schur identities, gauge-compatible
Gaussian shells, local measured expansions, and a small-field fibre
Hessian lower bound \(24/2327\) in its stated normalization.
The local one-loop remainder and retained response have their
specified small-field and finite/background domains. Their Ward
identities and retained massless directions must be preserved.
These local modules do not prove global compact integration, an
interacting continuum limit, or a physical mass gap. The closure
monograph keeps those requirements as separate hypotheses.

Full-fibre section coercivity and zero-action classification already
appear in f12 V13 and subsequent refinements. They are not new here.
Uniqueness at zero action must not be confused with the frozen,
frustrated parent problem that produced the noncentral wells below.

\subsection{Certified wells, and why hidden diffusion is not the endpoint}

For the specified frozen parent, f12 V89--V90 and f13 V8 certify two
strict normal-minimum conjugation orbits. Their action enclosures are
\[
 E_1\in[569.684411845,569.684411883],\qquad
 E_2\in(569.68444636,569.68444638).
\]
The second is strictly above the first. Its normal Hilbert--Schmidt
floor is \(31/125000\); the first has the stated floor \(1/5000\).
The corresponding orbit-tube inequality is local. Neither numerical
orbit has been proved to classify all global minima.

The second well gives an exponentially small upper bound on the
full hidden diffusion gap at that frozen exterior. Thus an attempt
to demand a uniform low-temperature hidden-diffusion gap there is
contradicted by the actual model. This is not a proof of slow
single-group heat bath, nor a disproof of a physical mass gap.
Localized removal and response estimates in f13 V9--V11 retain
their actual probability weights. In particular, the variable-coarse
calculations there cannot be substituted for the fixed-coarse repair
law, a distinction made explicit in f13 V12.

\subsection{The visible repair lineage}

For the integer-lattice shell used below, f13 V14 has 819 selected
free links, a sufficient 10034-group visible collar, and uniform
angular action second derivative at most \(-312\) with all other
fine groups arbitrary. The Hilbert--Schmidt squared speed is 1638.
An even translation embeds its stencil without wraparound for every
even fine side length \(L\ge32\).

In f13 V15 a local coarse box gauge gives
\(\max_e d(I,C_e^h)\le21\max_p d(I,Q_p(C))\), transporting the
fixed-coarse law and original update clock exactly. This extends
the local repair to small coarse curvature without trivializing
global holonomy. In f13 V16, a visible \(3\times3\) curvature
matrix has hidden debit \(168I_3\); a sufficiently negative
eigenvalue and strict source inactivity give a noncentral criterion.

In f13 V17, reversing \(k\) selected signs on one edge yields
actual pivot speed \(-k/(8-k)\) in the inner-source region.
The explicit upper curvature debit is
\[
 \ell(k)=24k+6\max\{2a_k^2+4a_k,4a_k^2\},
                    \qquad a_k=\frac{k}{8-k}.
\]
For instance eleven reversals on distinct edges retain angular
curvature at most \(-372/49\), and five on one edge retain
\(-356/3\). These are active inner-source regions, not arbitrary
transition-source configurations.

All these repair estimates concern visible events with all hidden
variables unrestricted. The same-law global resolvent comparison
\[
 H_\mu\succeq L_\mu(L_\mu+K_\beta)^{-1},\qquad
                       K_\beta\le C_{\rm f}(\beta+1),
\]
converts a suitable localized diffusion estimate to an original-clock
killed heat-bath floor. Here the resolvent is global, not a killed
resolvent. Conditional variance gives marginal form domination for
functions depending only on retained groups. No hidden posterior
weight is needed when the support is a visible cylinder.

A growing spatial union of locally repairable events is not thereby
repairable with the same floor; f12 V79 records that obstruction.
The monograph's weighted patch criterion still needs all-exterior
conditional patch gaps strong enough to cross its finite-size
threshold. None of the preceding local certificates supplies those
gaps by itself.

\subsection{Updated companion context}

The new SM--Einstein summary
\begin{center}\footnotesize
\path{manuscripts/summarized_notes/SM_Einstein_Continuum_Monograph_Updated(3).tex}
\end{center}
was checked alongside the YM monograph. Its fixed-spatial-mode,
fixed-time many-copy implication still leaves finite physical species
and ultraviolet interchange separate. Its remaining cofinal
construction does not supply the missing YM patch inequality.
The active SM V72, NS V26, and parallel YM V5 notes have unchanged
hashes from the previous checkpoint. NS retains its explicit absence
of a global-regularity claim; parallel YM's extensive-sector
obstruction remains a warning against imposing global charge tails
where only local control is justified. No cross-programme closure
is imported. We repeat the companion check at new V5.

\section{V1 acquisition: an explicit transition-source band}
\label{sec:transition}

\subsection{Geometry and what remains unrestricted}

Use f13 V14's integer coordinates. With directions numbered
\(0,1,2,3\), set
\begin{align*}
 Q&=\{-1,0,1,2,3\}^3\setminus\{0,1\}^3,\\
 T&=\{((2a+u,2b+v,1,2c+w),2):
        (a,b,c)\in Q,\ (u,v,w)\in\{0,1\}^3\setminus\{0\}\}.
\end{align*}
The central field has sign \(-1\) on direction-2 links at \(x_2=1\)
with at least one transverse odd parity, and sign \(+1\) otherwise.
There are 117 coarse edges carrying the seven selected links each.

Let \(F\) be the inherited visible collar: \(T\), the free inputs
of those 117 source families, and the free input closure of every
fixed incident plaquette. To form a free input closure, discard
tree links and replace a root pivot by the free links in its fourteen
nonroot paths. Those paths contain no further pivots. The verified
counts are
\[
 |F|=10034,\quad |T|=819,\quad
 (N_{2},N_{1,-},N_{1,+},N_{1,?})=(1845,690,450,84).
\]
Two-selected plaquettes and the first two one-selected classes are
determined by \(F,C\). The 84 remaining one-selected plaquettes
are unrestricted. The central parent's hidden set is disjoint
from \(F\); any additional hidden sets must also avoid \(F\).
We impose \(F\subset\mathcal S\) for an interface marginal.

Fix \(C=I\). Choose any \(\ell_*\in T\). Replace its central
value \(-I\) by \(e^{\theta E}\), where \(E=(0,0,0,1)\) is a unit
imaginary quaternion with \(\|E\|_{\rm HS}=\sqrt2\). All other
free groups in \(F\) retain their central values. No free group
outside \(F\) is prescribed. The six other selected links on the
same coarse edge remain \(-I\). The two chronological paths
containing \(\ell_*\) are therefore \(e^{\theta E}\), while the
twelve other nonroot sources on that edge are \(-I\).

\subsection{The actual scalar cutoff, with a nonzero second derivative}

For the computational realization already used in f13 V8, the smooth
cutoff is, in the open transition interval,
\[
 \chi(s)=\frac{1}{1+\exp(-A(s))},\qquad
 A(s)=\frac1{s-1/2}-\frac1{1-s},\quad \tfrac12<s<1.
\]
It is extended by one for \(s\le1/2\) and zero for \(s\ge1\).
This explicit choice is fixed throughout this acquisition.
Define the odd scalar truncated logarithm
\[
              h(u)=u\chi(64\sqrt2|u|).
\]
For positive \(u\) and \(s=64\sqrt2u\),
\begin{align}
 A'&=-(s-1/2)^{-2}-(1-s)^{-2},\nonumber\\
 A''&=2(s-1/2)^{-3}-2(1-s)^{-3},\nonumber\\
 \chi'&=\chi(1-\chi)A',\nonumber\\
 \chi''&=\chi(1-\chi)((1-2\chi)(A')^2+A''),\nonumber\\
 h'&=\chi+s\chi',\qquad
 h''=64\sqrt2(2\chi'+s\chi'').\label{eq:h-jets}
\end{align}

Parameterize the configuration by a relative logarithm, not an
approximate root solve:
\begin{equation}\label{eq:band}
 \frac{15}{16}\le s\le\frac{47}{50},\qquad
 u=\frac{s}{64\sqrt2},\qquad
 b=-\frac{h(u)}8,\qquad \theta=u+b.
\end{equation}

\begin{lemma}[Exact completed pivot and transition jets]
\label{lem:jets}
At this configuration the actual pivot on the selected edge is
\(B=e^{bE}\). The two positive relative sources have logarithm
\(uE\), strictly in the cutoff transition; the other twelve are
outside the support. Under common left multiplication of all
819 selected free links by \(e^{tE}\), the actual pivot is locally
\(B(t)=e^{b(t)E}\), where
\begin{equation}\label{eq:pivot-jets}
 b(t)+\frac18h(\theta+t-b(t))=0,\qquad
 b'(0)=-\frac{h'(u)}{8-h'(u)},\qquad
 b''(0)=-\frac{64h''(u)}{(8-h'(u))^3}.
\end{equation}
The physical quaternion acceleration is
\[
                 B''(0)=e^{bE}\{b''(0)E-(b'(0))^2I\}.
\]
Every other pivot is constant under this motion for sufficiently
small \(|t|\), uniformly over the unprescribed fine groups.
\end{lemma}

\begin{proof}
Equation \eqref{eq:band} gives \(\theta-b=u\) and
\(b+h(u)/8=0\), precisely the forward equation for two contributing
nonroot sources. The twelve negative relative sources are near
\(-I\), not the logarithm support. Global uniqueness of the completed
pivot identifies the displayed solution with the actual one.

All sources in this family use only the prescribed collar. Under
the common rotation the positive sources become
\(e^{(\theta+t)E}\), and the negative ones become \(-e^{tE}\).
There is no noncommutative simplification of arbitrary hidden fields
in this step: those fields do not enter this source family.
The scalar equation remains valid locally. Since \(h'\le1\), its
derivative with respect to \(b\) is \(1-h'/8\ge7/8\).
Twice differentiating proves \eqref{eq:pivot-jets}; differentiating
the exponential gives the additional radial acceleration term.

Each selected boundary group belongs only to the two paths on its
own edge. On all other selected edges the fourteen nonroot sources
are near \(-I\), so their pivots remain \(I\). Other pivots have
no selected inputs and are constant. All assertions are uniform
in the remaining fine fields by this source incidence.
\end{proof}

\subsection{Certified whole-band enclosures}

The new checker uses outward intervals with integer endpoints
divided by \(2^{70}\), inherited from the f12 V90 arithmetic. It
evaluates \eqref{eq:h-jets} on the entire real interval in
\eqref{eq:band}. No finite sampling or floating-point rounding
assumption enters the enclosure.

\begin{proposition}[Rational bounds on the transition band]
\label{prop:bounds}
Uniformly for \eqref{eq:band},
\[
 0<\theta<\frac{11}{1000},\qquad
 -\frac1{10^8}<b<0,\qquad
 0<b'<\frac1{20000},\qquad -1<b''<0.
\]
In particular the angular acceleration is not zero.
\end{proposition}

\begin{proof}
Here are outward rational supersets of the computed intervals;
all displayed decimal endpoints denote exact terminating rationals:
\[
\begin{array}{c|cc}
 &\text{lower endpoint}&\text{upper endpoint}\\ \hline
 h&0.000000005808&0.000000011492\\
 h'&-0.000293798979&-0.000136195694\\
 h''&2.749382989710&6.744377981289\\
 \theta&0.010358008053&0.010385630123\\
 b&-0.000000001437&-0.000000000726\\
 b'&0.000017023836&0.000036724248\\
 b''&-0.843004191851&-0.343635012467
\end{array}
\]
The reproducible calculation starts from the integer square-root
enclosure of \(\sqrt2\), forms \(u,h,h',h''\), and applies the
exact rational formulas \eqref{eq:band} and \eqref{eq:pivot-jets}.
For its sole exponential evaluation, \(|A|<16\). The inherited
routine evaluates the degree-24 Taylor polynomial at \(-A/64\),
adds the error bound \(2|A/64|^{25}/25!\), and squares six times
with outward rounding. The Taylor error bound follows from
\(e^{|A|/64}<2\). All reciprocals have strictly positive interval
denominators. The interval operations therefore enclose the exact
real expressions throughout the band; the displayed inequalities
follow by rational comparison. The script and dependency hashes
are recorded at the end.
\end{proof}

\subsection{Uniform curvature including all new pivot terms}

\begin{theorem}[Transition-band descent]
\label{thm:curvature}
For every choice of \(\ell_*\in T\), every \(s\) in
\eqref{eq:band}, and every unprescribed fine field,
\[
             \left.\frac{d^2}{dt^2}S(Y(t),B(Y(t),I))\right|_{t=0}
                         \le-270.
\]
Its unit-speed normalization in the independent free-group metric
has action Hessian at most \(-270/1638=-15/91\).
\end{theorem}

\begin{proof}
Start from the comparison used in f13 V17: reverse only
\(\ell_*\) to \(+I\) and, for this auxiliary derivative only,
hold every pivot at \(I\). The original central shell bound is
\(-312\). Two-selected fixed plaquettes cancel the common rotation.
Each one-selected fixed plaquette changes its second derivative
by at most 4 on a sign reversal; there are at most six such
plaquettes touching \(\ell_*\). The 84 uncontrolled contributions
still have their uniform upper bound 2. Hence this comparison
derivative is at most \(-288\).

Now change \(\ell_*\) from \(+I\) to \(e^{\theta E}\), and its
pivot from \(I\) to \(e^{bE}\), but continue to freeze pivots
when computing the auxiliary derivative. All fixed two-selected
plaquettes still cancel: their transverse links are central, and
the moving factors enter as a conjugated product whose trace is
unchanged. A plaquette containing the root pivot can contain at
most one selected link, since all these links are parallel.

For a one-selected plaquette the frozen-pivot second derivative
is twice the scalar part of its holonomy. Changing one factor
from \(I\) to \(e^{vE}\) changes that scalar by at most
\(|e^{vE}-I|\le|v|\) in unit-quaternion norm. There are at most
six incidences of each changed factor, so the additional debit
is bounded by \(12(|\theta|+|b|)\). Uncontrolled plaquettes can
instead retain their bound 2. This estimate also overbounds any
double incidence; no independence assumption is used.

Finally restore the actual pivot motion. Put \(a=|b'|\) and
\(d=|b''|+(b')^2\), so \(|B'|=a\) and \(|B''|\le d\).
There is exactly one moving pivot. In an incident plaquette
whose other parallel link is selected, the product rule adds
terms bounded in absolute action second derivative by
\(2d+4a\): \(2d\) for the acceleration and \(4a\) for the
mixed first derivatives. If the other parallel link is fixed,
the bound \(2d\) suffices. Inverse orientations preserve these
quaternion norms. The remaining factors are unit quaternions,
even when they depend on arbitrary hidden inputs.

Summing its six incidences, the full actual derivative is at most
\begin{align*}
 -288+12(|\theta|+|b|)+12|b''|+12(b')^2+24|b'|
 &<-288+12\left(\frac{11}{1000}+\frac1{10^8}\right)
       +12+\frac{12}{20000^2}+\frac{24}{20000}\\
 &<-270.
\end{align*}
This includes the tangential angular acceleration \(b''E\),
which is absent in the inner-source formula of f13 V17.
The independent coordinates are just the 819 selected free groups,
so the squared speed is \(2\cdot819=1638\), not a sum that
adds dependent pivot velocities. The free-group curve is a
product-metric geodesic, proving the Hessian assertion.
\end{proof}

\section{A visible transition repair event in the original clock}

The union of these centres, over the finite set \(\ell_*\in T\)
and the closed band \eqref{eq:band}, is a compact subset of the
fixed carrier \(\SU(2)^F\). Its sources lie a positive distance
inside the transition interval. The other twelve relative sources
are uniformly outside the support. Smooth dependence of the actual
pivot preserves these properties on sufficiently small outer
neighbourhoods of this compact family. Those neighbourhoods include
noncommuting perturbations of the visible groups.

The possible moving physical links and their full input closure
have the same finite bounds checked in f13 V17: at most 936 links,
3156 incident plaquettes, and unshifted input vertices in
\[
 [-4,9]\times[-4,9]\times[0,3]\times[-3,9].
\]
Translation by \((8,8,8,8)\) therefore suffices for all even
\(L\ge32\). This closure is used for uniform derivative bounds;
its unprescribed groups are not added to the visible event.

\begin{theorem}[Original-clock floor for a genuine transition family]
\label{thm:repair}
Assume \(C=I\), even \(L\ge32\), and \(F\subset\mathcal S\).
There is a visible event \(D^{\rm tr}\), a union of closed inner
chart cylinders whose interiors contain all the above centres,
and constants \(\beta_0,C_{\rm f}<\infty\), independent of volume
and compatible hidden index sets, such that
\[
 \kappa_{D^{\rm tr}}^{H_\nu}
 \ge\frac{c\beta}{2\{c\beta+8C_{\rm f}(\beta+1)\}}
 \ge\frac{c}{2(c+16C_{\rm f})}>0,
             \qquad c=\frac{15}{1456},\quad\beta\ge\beta_0.
\]
Here \(H_\nu\) is the exact interface heat bath with each retained
fine group updated at rate one. No coarse variable is updated.
All hidden and unprescribed visible groups remain unrestricted.
The charts and \(\beta_0\) are existential, not numerically enclosed.
\end{theorem}

\begin{proof}
In product orthonormal normal coordinates about each centre, freeze
the normalized selected direction. At the centre its coordinate
second derivative equals the intrinsic Hessian. Uniform continuity
on the finite compact input closure gives outer charts on which
\(\partial_a^2S\le-270/(2\cdot1638)\), uniformly in all other
inputs. Shrink them so the inverse metric is at least \(1/2\).

Retain the actual measured potential
\(\Phi_\beta=\beta S-\log j-\log h_{\rm Haar}\).
Only the 117 selected-edge Jacobian factors can depend on a selected
free coordinate. Each is smooth and strictly positive on a fixed
compact carrier, so its logarithmic second derivatives are bounded.
The coordinate-Haar derivatives are also bounded on smaller charts.
They do not vanish at transition sources and are not discarded.
For sufficiently large \(\beta\),
\[
                   \partial_a^2\Phi_\beta
                                  \le-\frac{270\beta}{4\cdot1638}.
\]
For a chart-supported function \(v\), set \(w=v e^{-\Phi_\beta/2}\)
in this direction. Integration by parts gives
\[
 \int|\partial_av|^2e^{-\Phi_\beta}
 =\int\left\{|\partial_aw|^2+
       \left(\frac{(\partial_a\Phi_\beta)^2}{4}
                   -\frac{\partial_a^2\Phi_\beta}{2}\right)|w|^2\right\}.
\]
Integrate all other coordinates and use the inverse-metric factor
\(1/2\). The local joint diffusion floor is therefore
\(c\beta\), with \(c=270/(16\cdot1638)=15/1456\).

Compactness supplies a finite collection of outer charts and smooth
cutoffs \(\chi_j(Y_F)\), supported there, such that
\(\sum_j\chi_j^2\le1\) everywhere and equals one on
\(D^{\rm tr}\). One construction takes bump functions \(\psi_j\)
with \(T=\sum_j\psi_j^2\ge1\) on the inner charts and divides them
by \(\sqrt{T+g(T)^2}\), where \(g=1\) for \(T\le1/4\) and
\(g=0\) for \(T\ge1/2\). Let
\(M=\sup\sum_j|\nabla\chi_j|^2<\infty\). Summing the local
floors and using the product rule gives
\[
 \|1_{D^{\rm tr}}v\|^2\le
       \frac{2}{c\beta}\E_L(v,v)+\frac{2M}{c\beta}\|v\|^2.
\]
This extends from a smooth core by form closure. No growing spatial
cover occurs: \(F\) and its compact family are fixed.

Use the inherited same-law global comparison
\(H_\mu\succeq L(L+K_\beta)^{-1}\),
\(K_\beta=C_{\rm f}(\beta+1)\). The constant is uniform because
the completion stencils and mixed-Hessian incidences are fixed.
Write \(\lambda=c\beta\). For the spectral projection \(\Pi_t\)
of \(L\) on \([0,t]\), the last inequality gives
\(\|1_D\Pi_t\|^2\le2(t+M)/\lambda\).
The adjoint \(\Pi_t1_D\) has the same bound, without any
commutation assertion. If \(\lambda\ge8M\) and \(t=\lambda/8\),
every function supported in \(D\) has at least half its squared
norm above this spectral interval. Hence its joint heat-bath
Rayleigh quotient is at least
\(\lambda/[2(\lambda+8K_\beta)]\).

Lift a marginal function as \(Jf=f(Y_{\mathcal S})\). Norms are
preserved, hidden updates leave \(Jf\) fixed, and conditional
variance gives \(\E_{H_\nu}(f,f)\ge\E_{H_\mu}(Jf,Jf)\).
The support is still a cylinder allowing every hidden group.
Apply the joint bound and \(\beta+1\le2\beta\) for \(\beta\ge1\).
\end{proof}

These centres fail the f13 V16 inactivity test: their positive
source distance from \(C_e=I\) is at most its geodesic distance
\(\sqrt2\theta<\rho\). Their actual relative source parameter
is \(s>1/2\), so the inner-source hypothesis of f13 V17 also
fails. Shrinking the present charts preserves both statements.
Thus a genuine previously untreated transition mechanism is
certified; it is not obtained by relabelling an inactive point.

\section{Reproducibility, checkpoint, and the next bottleneck}

The new checker is
\begin{center}\small
\path{scripts/verify_ym_f14_v1_transition_band.py}.
\end{center}
Its SHA--256 is
\begin{center}\scriptsize\ttfamily
F83E98BB1BC7D1EA54962B1113C1E159694A6AB9F30C392A772A32674B5050AA.
\end{center}
It uses \path{scripts/ym_f13_v8_transition_intervals.py}, hash
\begin{center}\scriptsize\ttfamily
75F7B999482BD2F5CF91BE44D2721F4335127C04EA27051E58F9231DCEFA5DAC,
\end{center}
and \path{scripts/ym_f12_v90_intervals.py}, hash
\begin{center}\scriptsize\ttfamily
39D803E9B1CDE519BE19A74D11CD8AEA6BE00A9A3B216339797C81ED3F1B1740.
\end{center}
It also reruns the f13 V17 geometry checker (14,976 local source
patterns), which reruns the f13 V14 shell audit. These finite checks
establish the stated geometry and rational enclosures, not the
missing global coverage. The analytic proof supplies the passage
from those enclosures to a local original-clock floor.

The SM summary inspected at restart has SHA--256
\begin{center}\scriptsize\ttfamily
54176A3730275E1BD9735199860ECA458537559613AE05E3D5C05234D83DBDFA.
\end{center}
The active companion-note hash prefixes remain
\(\texttt{F44F9745D43379E1}\) (SM V72),
\(\texttt{82C887F8C2C69E4F}\) (NS V26), and
\(\texttt{306F1BB84CFA8A8D}\) (parallel YM V5).
The next companion and direction review is V5 of this restarted
lineage. All archives and both updated monographs remain unchanged.

\begin{firewall}
The new theorem certifies a whole transition band and nearby
noncentral visible configurations with all hidden fields free.
It retains the nonzero pivot angular acceleration and the actual
Jacobian. It still concerns one transition-bearing coarse edge
within a prescribed finite collar at identity coarse data.
It does not classify arbitrary transition networks, intermediate
curvature configurations, or all conditional minima. It is not
an unrestricted interface gap or a continuum mass-gap theorem.
\end{firewall}

The next useful test is whether a complete same-law patch argument
can combine these mechanisms with the remaining configurations.
In particular, large transition derivatives in the middle of the
cutoff and multiple neighbouring responding pivots must be bounded
with their true mixed terms. Repeating slightly different one-edge
examples will not establish that coverage. Any proposed covering
must retain every unresolved branch and state how it yields an
all-exterior patch inequality strong enough for the monograph's
threshold. The full physical goal remains active.
\section{V2: complete boundary-source freedom on omitted shell edges}
\label{sec:v2}

\subsection{Upstream change of the descent direction}

V1 follows a common rotation through a transition and pays the
resulting pivot acceleration. Instead of refining that one-edge
calculation, we now change the direction: omit whole coarse edges
from the rotation and allow their seven free boundary groups to be
arbitrary. This keeps their actual pivots fixed along the test
curve, even when their sources lie in different cutoff regimes.
The construction avoids a bound on their transition derivatives;
it does not set the derivatives of a moving pivot to zero.

The central-pattern cut argument of f12 V78 already rotates a
subset of negative links. That is not the new claim. Here the
omitted free groups range over their complete noncommutative
compact carrier, their pivots are the actual completed ones, and
all original hidden fields remain unrestricted. The new finite
incidence estimate quantifies the loss from these arbitrary groups
and then gives a visible original-clock repair event.

\subsection{An omitted edge releases seven complete group variables}

Keep V1's \(C=I\), shell \(T\), collar \(F\), and non-wrapping
embedding. Let \(E_T\) be the 117 coarse edges that own the
selected boundary links. For a subset \(B\subset E_T\), define
\[
 T_B=\{l\in T:\operatorname{owner}(l)\in B\},\qquad
 R_B=T\setminus T_B,\qquad F_B=F\setminus T_B.
\]
Thus \(|T_B|=7|B|\). Prescribe the central values only on
\(F_B\). Every free group in \(T_B\), and every free group
outside \(F\), is arbitrary. No constraint on the sources, their
relative angles, or their commutativity is imposed on the edges
in \(B\). Their internal source-path groups still have the
central values inherited from \(F_B\); this is not freedom of
every internal field in their coarse cells.

Let \(P_B\) denote the root pivots on the edges in \(B\), and
let \(A_B=T_B\cup P_B\) be their potentially noncentral
physical boundary links. The completed pivots in \(P_B\) need
not be identity and are not independent coordinates. Bounds below
use only their unitarity, so they apply in particular to their
actual constrained values.

\begin{lemma}[No responding pivot along the omitted-edge direction]
\label{lem:v2-frozen}
For any unit imaginary quaternion \(E\), left multiply only
\(l\in R_B\) by \(e^{tE}\), keeping every other free coordinate
fixed. For sufficiently small \(|t|\), every actual pivot is
constant along this curve, uniformly in all the arbitrary groups.
The entire inverse-pivot Haar amplitude is independent of the
rotated coordinates on an open cylinder around these data.
\end{lemma}

\begin{proof}
A selected boundary group occurs only in two nonroot paths of
its own coarse edge. On \(e\notin B\), all seven such groups
have central value \(-I\), and their internal companion links
are \(I\). All fourteen source products are therefore
\(-e^{tE}\). At candidate pivot \(I\) their truncated
logarithms vanish in a uniform neighbourhood. Uniqueness gives
the actual pivot \(I\), with inverse Haar Jacobian one, independent
of these selected coordinates throughout that neighbourhood.

On \(e\in B\), no source group is rotated. Its completed pivot
and inverse Haar factor may depend on arbitrary omitted fields,
but are constant along the curve and independent of all coordinates
in \(R_B\). Other pivots also have no rotated input. This uses
the source-incidence property, not a claim that arbitrary transition
derivatives themselves vanish. It proves the amplitude assertion
in an open cylinder where all unomitted sources remain inactive.
\end{proof}

The visibility requirement is now only \(F_B\subset\mathcal S\).
In particular, the seven free groups per omitted edge can even
belong to the hidden index set. This is permitted because no
event restriction or differentiated input depends on their values.
The other hidden sets must still avoid \(F_B\).

\subsection{A plaquettewise certificate for every omitted-edge subset}

Let \(\mathcal P_T\) be the set of plaquettes incident to the
original shell \(T\). Let \(A_0\) be the parent-dependent physical
links: the 944 parent free groups and its 96 root pivots. These
identify the original 84 uncontrolled incident plaquettes. Write
\(\epsilon_l\in\{1,-1\}\) for the original central physical
sign convention, and for \(p\in\mathcal P_T\) let
\(r_B(p)=|\partial p\cap R_B|\). Define the following integer,
with the cases in the displayed order:
\begin{equation}\label{eq:v2-cp}
 c_p(B)=
 \begin{cases}
 0,&r_B(p)=0\text{ or }2,\\
 2,&r_B(p)=1\text{ and }\partial p\cap(A_0\cup A_B)\ne\varnothing,\\
 2\prod_{l\in\partial p}\epsilon_l,&\text{otherwise}.
 \end{cases}
\end{equation}
In the last case there is one rotated link. Orientations do not
alter central signs. Set
\[
 C(B)=\sum_{p\in\mathcal P_T}c_p(B),\qquad
 d_e=C(\{e\})-C(\varnothing),\quad e\in E_T.
\]
These depend only on the finite geometry and on which edges are
omitted, not on any values of the arbitrary group variables.

\begin{proposition}[Actual curvature bounded by the finite certificate]
\label{prop:v2-certificate}
At the central \(F_B\)-data, uniformly in all arbitrary groups,
\[
                S''(0)\le C(B),\qquad C(\varnothing)=-312.
\]
The squared speed of the curve in the independent free-group
metric is \(2|R_B|=14(117-|B|)\).
\end{proposition}

\begin{proof}
By Lemma~\ref{lem:v2-frozen}, only the physical links in \(R_B\)
move. A plaquette outside \(\mathcal P_T\) has zero derivative.
Each incident plaquette contains at most two such links, which
are parallel with opposite orientations. With no rotated link
the derivative is zero. With two, the plaquette was in the original
fixed two-selected class: its transverse links are central, no
arbitrary parallel factor fits alongside the two selected ones,
and the common rotation cancels in its trace.

If there is one rotated link, its second derivative, or the second
derivative of its inverse, is its negative. Hence the action
second derivative is twice the scalar part of the actual holonomy,
at most 2 for arbitrary unit-quaternion factors. This proves the
middle case of \eqref{eq:v2-cp}, including arbitrary solved pivots.

Otherwise every physical factor equals its original central value.
Indeed, in an original fixed plaquette a removed selected group
can change only itself and its own root pivot, precisely \(A_B\).
All other pivot source inputs remain central and give pivot \(I\)
by uniqueness. Thus the last case is exact. This assertion uses
the closure defining \(F\); it does not prescribe independent
values to solved pivots. Summing proves the bound. At \(B=\varnothing\)
it reproduces \(2(-690+450+84)=-312\). The norm counts only the
remaining independent selected groups.
\end{proof}

\begin{lemma}[Subadditive loss under omission]
\label{lem:v2-subadditive}
For every subset \(B\subset E_T\), including adjacent omitted
edges, one has
\begin{equation}\label{eq:v2-additive}
                     C(B)\le-312+\sum_{e\in B}d_e.
\end{equation}
\end{lemma}

\begin{proof}
We prove the inequality before summing, with
\(c_p(B)-c_p(\varnothing)\) in place of \(C(B)+312\).
Only owners of the two parallel physical links of a plaquette
can matter, so at most two distinct edges affect a summand.
If there is only one, the claim is an equality.

For two original selected links with distinct owners, the baseline
is zero, either single omission gives the bound 2, and both
omissions give zero. Thus the joint change is at most the sum of
the two single changes. If both selected links have one owner,
we are in the preceding one-edge case.

For one original selected link, call its owner \(e\). A second
relevant edge \(f\ne e\) can only own the opposite root pivot.
Write \(v=c_p(\varnothing)\in\{-2,2\}\). Omitting \(e\) gives
zero, omitting only \(f\) gives 2, and omitting both gives zero.
The desired inequality for the only nontrivial case is
\[
                         -v\le(-v)+(2-v),
\]
which holds because \(v\le2\). Uncontrolled plaquettes have
the same upper value \(v=2\). If no second owner occurs, or it
coincides with \(e\), the single-edge identity suffices.
These cases exhaust the incident plaquettes. Summing proves
\eqref{eq:v2-additive}; no independence or spatial separation of
omitted edges was assumed.
\end{proof}

The exact integer enumeration gives the following full histogram:
\[
\begin{array}{c|rrrrrrrrrr}
 d_e&12&16&20&24&28&32&36&44&48&60\\ \hline
 \#\{e:d_e\text{ has this value}\}&1&3&3&1&21&18&9&27&27&7.
\end{array}
\]
The multiplicities sum to 117. In particular every edge costs at
most 60 in this sufficient upper bound. The verifier checks the
free-input incidence used above and all 8250 local edge-subset
cases in \eqref{eq:v2-cp}; the analytic subadditivity argument is
why these local cases certify every subset of the 117 edges,
without enumerating \(2^{117}\) global patterns.

\begin{theorem}[Arbitrary compact boundary data on five omitted edges]
\label{thm:v2-five}
For any \(B\subset E_T\) with \(|B|\le5\), prescribe only
the central \(F_B\)-data and leave every other free group arbitrary.
Then the actual fixed-coarse action has
\[
                              S''(0)\le-12.
\]
Thus arbitrary, possibly noncommuting boundary configurations on
any five coarse edges, comprising 35 free group variables, are
allowed. Their cutoff regimes and pivot responses with respect
to their own fields can be arbitrary.
\end{theorem}

\begin{proof}
Proposition~\ref{prop:v2-certificate} and
Lemma~\ref{lem:v2-subadditive} give
\(S''\le-312+60|B|\le-12\).
\end{proof}

The weighted criterion is stronger than the cardinality bound.
Order edges first by \(d_e\), then lexicographically by their
coarse coordinate and direction. The first thirteen have total
cost \(284\), giving \(C(B)\le-28\); all their 91 free boundary
groups may be arbitrary. The checker also evaluates this particular
subset directly, obtaining \(C(B)=-140\). The first fourteen
have total cost 312, so the additive sufficient criterion alone
does not certify fourteen. This is not an impossibility statement:
the direct bound \(C(B)\) can be strictly better, as this example
already illustrates.

\subsection{Uniform neighbourhoods without restrictions inside the holes}

For fixed \(B\) with \(C(B)\le-12\), normalize the selected
direction to unit independent-group speed. Its action Hessian
at the central \(F_B\)-data is at most
\[
                  -\frac{12}{2|R_B|}\le-\frac{12}{1638}.
\]
The assertion is uniform over the whole compact carrier of omitted
boundary groups and all other unprescribed fine inputs in the
finite stencil. Consequently a sufficiently small product chart
on \(F_B\) has frozen coordinate second derivative at most
\(-12/(2\cdot1638)\), uniformly in those unrestricted groups.
No small neighbourhood is imposed on an omitted group.

The source separation on unomitted edges persists on this chart.
The amplitude is independent of the rotated free coordinates,
by Lemma~\ref{lem:v2-frozen}; its possibly strong dependence on
omitted variables is immaterial to this particular directional
derivative. The coordinate-Haar Hessian is bounded. Shrinking for
inverse metric at least \(1/2\), and taking sufficiently large
\(\beta\), therefore gives
\[
 \partial_a^2\Phi_\beta\le-\frac{12\beta}{4\cdot1638},\qquad
            \E_L(v,v)\ge c\beta\|v\|^2,
                         \quad c=\frac{12}{16\cdot1638}=\frac1{2184},
\]
for functions supported in this outer visible chart cylinder.
The measured ground-state conjugation is exactly the one displayed
in the proof of Theorem~\ref{thm:repair}. We use the original
completed density and independent-group metric throughout.

\begin{theorem}[Original-clock repair with arbitrary source regimes]
\label{thm:v2-repair}
Fix \(C=I\) and a declared retained index set \(\mathcal S\).
Let \(\mathfrak B\) be any family of subsets of the fixed
117-edge set such that
\[
                 C(B)\le-12,\qquad F_B\subset\mathcal S
                         \quad(B\in\mathfrak B).
\]
For sufficiently small closed inner product charts on each
\(F_B\), let \(D_B\) be the corresponding visible cylinder,
and set \(D_{\mathfrak B}=\bigcup_{B\in\mathfrak B}D_B\).
Every group outside the relevant \(F_B\) is unrestricted in
\(D_B\). There are constants \(\beta_0,C_{\rm f}<\infty\),
independent of even \(L\ge32\) and of compatible hidden sets,
such that
\[
 \kappa_{D_{\mathfrak B}}^{H_\nu}
 \ge\frac{c\beta}{2\{c\beta+8C_{\rm f}(\beta+1)\}}
 \ge\frac{c}{2(c+16C_{\rm f})}>0,
           \qquad c=\frac1{2184},\quad\beta\ge\beta_0.
\]
This includes every set of at most five omitted edges, and the
thirteen-edge example, whenever their required collars are retained.
If all of \(F\) is retained, all these subsets are admissible
simultaneously. No coarse or collective rotation update is added
to the rate-one single-free-group heat bath.
\end{theorem}

\begin{proof}
There are finitely many subsets of this fixed edge set. The preceding
local floor holds for each admissible subset, uniformly in all
unprescribed group variables. Select inner and outer charts with
a common positive allowance by finiteness. All derivatives relevant
to their uniformity belong to the original fixed input closure;
ambient volume introduces no new local parameter.

The cutoffs can depend on different \(F_B\) for different \(B\).
Regarding all of them as functions on the fixed compact union of
these coordinate carriers, choose normalized smooth cutoffs as in
Theorem~\ref{thm:repair}. Their supports are cylinders contained
in the corresponding outer charts. Even after normalization, they
depend only on retained coordinates, since every \(F_B\) does.
Their squared sum is at most one everywhere and equals one on
\(D_{\mathfrak B}\), with finite derivative sum bound \(M\).
The product rule gives
\[
 \|1_{D_{\mathfrak B}}v\|^2
      \le\frac{2}{c\beta}\E_L(v,v)+\frac{2M}{c\beta}\|v\|^2.
\]
The global resolvent comparison and spectral projection argument
in Theorem~\ref{thm:repair}, with \(c\beta\ge8M\), give the
stated joint killed floor. The same conditional-variance domination
transfers it to the exact interface marginal. Lifting a marginal
function restricts no hidden variable: for each chart, even any
hidden omitted groups remain completely free. Thus no posterior
probability factor is inserted. This is one fixed configuration
space cover, not a growing spatial union.
\end{proof}

\subsection{Verification, nonduplication, and the remaining obstruction}

The new checker is
\begin{center}\small
\path{scripts/verify_ym_f14_v2_perforated_shell.py}.
\end{center}
Its SHA--256 is
\begin{center}\scriptsize\ttfamily
C7CB30616BFBBDBF678FEF384FE4B18F0A47B3FB7FB9FB5190CF634CD1D30FE8.
\end{center}
It rebuilds the original shell and collar, verifies the single-owner
source dependency, computes the full debit histogram, checks every
local subset case of the subadditivity inequality, and verifies the
five-edge and thirteen-edge claims. These are exact integer and
rational checks. It does not enclose numerical chart radii, cutoff
gradient bounds, or the large-\(\beta\) threshold.

This appendix uses the existing compact inverse-pivot theorem,
original-clock comparison, and marginal domination with their
actual state spaces. Its new coverage is complete in the omitted
boundary-group variables: no hidden cutoff branch is excluded on
those edges. V1 instead restricts the special source to a small
transition band and differentiates its moving pivot. Both remain
useful sufficient mechanisms; they are not asserted to exhaust
all interface configurations.

The predecessor V1 has 27,836 bytes and SHA--256
\begin{center}\scriptsize\ttfamily
DAC071916A2B6216DEE5F38816162C5D7739018D541BF8ECA21F960516259D3B.
\end{center}
Its complete body is preserved apart from the title version.
Only active V1 is replaced by V2. The archived f13 notes and
updated monographs remain unchanged. The next companion and
direction review remains new V5.

\begin{firewall}
The repair event now tolerates arbitrary compact boundary fields
on any five shell edges, including mixed transition regimes,
and on some larger subsets certified by the finite geometry.
It still requires the other visible collar coordinates near
their prescribed central values and keeps coarse data at identity.
There is no classification of arbitrary internal or transverse
disorder, no unrestricted all-exterior conditional patch gap,
and no cofinal continuum or physical mass-gap proof.
\end{firewall}

This reduces the transition problem to a geometric cost when
troublesome sources can be confined to omitted edges. The next
upstream issue is whether all remaining bad configurations admit
an affordable omitted-edge set or a different complete patch
estimate. Without such an implication, arbitrary source freedom
on this partial carrier is not global coverage. The full goal
remains active.
\section{V3: exact optimization of the omitted-edge certificate}
\label{sec:v3}

\subsection{What is optimized, and what is not}

V2 bounded the loss from omitting several edges by the sum of their
single-edge losses. That bound discards the reduction in boundary
cost when omitted edges are adjacent. Here we retain those pair
terms exactly. The resulting binary optimization is an integer
minimum-cut problem on a fixed 117-vertex graph. It supplies both
larger certified repair regions and exact failure certificates for
this particular sufficient bound.

The optimized quantity is \(C(B)\) from \eqref{eq:v2-cp}, an upper
bound on the actual action curvature under a specified family of
common rotations. It is not the Wilson action, its exact Hessian,
a conditional spectral gap, or the physical mass. A minimum of
zero means this upper-bound mechanism cannot certify strict
descent under the required omissions. It does not prove the
absence of another descent direction or repair mechanism.

\subsection{An exact nearest-neighbour binary energy}

Identify \(E_T\) with the transverse coarse vertex set
\[
             Q=\{-1,0,1,2,3\}^3\setminus\{0,1\}^3.
\]
Each vertex \((a,b,c)\) stands for the direction-2 coarse edge
with root \((a,b,0,c)\). Let \(G=(Q,\mathcal N)\) be its induced
nearest-neighbour graph: \(\{e,f\}\in\mathcal N\) if the two
transverse coordinates have \(\ell^1\) distance one. There are
117 vertices and 264 undirected edges. Indeed the full side-five
grid has 300 edges, and deleting the eight parent-cube vertices
removes 12 internal and 24 external incident edges.

Write \(x_e=1\) when \(e\) is omitted and zero otherwise.
Let \(d_e\) be the single-edge debit from V2 and let
\(\deg_G(e)\) be the degree in \(G\).

\begin{theorem}[Exact quadratic representation]
\label{thm:v3-quadratic}
For every subset \(B\subset E_T\),
\begin{equation}\label{eq:v3-quadratic}
 C(B)=-312+\sum_e d_ex_e
                         -16\sum_{\{e,f\}\in\mathcal N}x_ex_f.
\end{equation}
Equivalently, for the rotated-edge set \(A=E_T\setminus B\),
\begin{equation}\label{eq:v3-cut}
 C(B)=\sum_{e\in A}v_e+8|\partial_G A|,
                 \qquad v_e=8\deg_G(e)-d_e.
\end{equation}
Here \(\partial_G A\) is the set of undirected graph edges with
exactly one endpoint in \(A\).
\end{theorem}

\begin{proof}
Each plaquette summand in \eqref{eq:v2-cp} depends on at most
two distinct omitted-edge indicators. Its constant and linear
coefficients are its baseline and single-edge changes. The
negative of its quadratic coefficient is
\[
 w_p=c_p(\{e\})+c_p(\{f\})-c_p(\varnothing)-c_p(\{e,f\}).
\]
The cases in Lemma~\ref{lem:v2-subadditive} show \(w_p=4\) for
two selected links with distinct owners. For one selected link
and a root pivot with distinct owners, \(w_p=2-v\), where
\(v=c_p(\varnothing)\in\{-2,2\}\), so it is 4 or zero.
If at most one owner is relevant, there is no pair coefficient.
This determines the full summand on all four binary assignments;
there is no higher-order term.

Distinct relevant owners are neighbouring coarse cells. Across a
face between two cells in \(Q\), four fine plaquettes connect
their direction-2 boundary links. Three connect two nonroot
selected links. The fourth connects a selected link to the root
pivot of the neighbouring cell and has negative central product.
Each of these four gives \(w_p=4\), for total pair coefficient
16. No non-neighbouring owners share such a plaquette. The exact
enumeration checks these incidences and all 8250 local binary
assignments. Summing the constant and linear coefficients gives
\eqref{eq:v3-quadratic}.

For binary variables,
\(2x_ex_f=x_e+x_f-|x_e-x_f|\). Substitute this into
\eqref{eq:v3-quadratic}. Omitting all edges gives \(C(E_T)=0\)
directly from \eqref{eq:v2-cp}, since no group is rotated; thus
the constant term disappears upon writing \(y_e=1-x_e\).
The remaining unary coefficient is \(8\deg_G(e)-d_e\),
and \(|x_e-x_f|=|y_e-y_f|\), proving \eqref{eq:v3-cut}.
\end{proof}

The full unary histogram is
\[
\begin{array}{c|rrrrrrr}
 v_e&-12&-8&-4&0&4&8&12\\ \hline
 \#\{e\}&7&27&36&19&24&3&1.
\end{array}
\]
In particular,
\[
 \sum_e v_e=-312,\qquad
 \sum_e\min(v_e,0)=-444,\qquad
 \sum_e|v_e|=576,\qquad \sum_e d_e=4536.
\]
All coefficients are exact integers. The identity is a refinement
of V2's additive estimate: the omitted pair terms are nonpositive,
not uncontrolled errors.

\subsection{Forced arbitrary fields and a finite exact oracle}

Suppose boundary groups on a prescribed set \(B_0\subset E_T\)
must be allowed to be arbitrary. We can omit additional edges
if this improves the certificate. Define
\begin{equation}\label{eq:v3-oracle}
                 m(B_0)=\min_{B_0\subset B\subset E_T}C(B).
\end{equation}
Here and below set containment includes equality. This is a
finite optimization of an explicitly defined integer function.
It satisfies \(m(B_0)\le0\), by the admissible choice \(B=E_T\).

Construct a directed capacitated graph by adding source \(s\)
and sink \(t\) to \(G\). For each \(v_e<0\), add an arc
\(s\to e\) of capacity \(-v_e\); for each \(v_e>0\), add
\(e\to t\) of capacity \(v_e\). For each undirected neighbour
pair add both arcs, each of capacity 8. For each \(e\in B_0\)
add a forcing arc \(e\to t\) of capacity
\[
                 M=1+576+8\cdot264=2689.
\]
Parallel arcs are allowed. A cut identifies its source-side
ordinary vertices with the rotated set \(A\).

\begin{theorem}[Certified minimum-cut oracle]
\label{thm:v3-oracle}
If \(q(B_0)\) is the minimum source--sink cut capacity of this
network, then
\[
                     m(B_0)=-444+q(B_0).
\]
An integer flow and a cut of the same value prove this minimum
exactly. A minimizing cut supplies the omitted-edge set
\(B_*=E_T\setminus A_*\). In particular, this procedure
decides whether any omitted-edge superset of \(B_0\) has
negative certificate, without enumerating those supersets.
\end{theorem}

\begin{proof}
For \(y_e=1_A(e)\), the cut capacity is
\[
 \sum_{v_e<0}(-v_e)(1-y_e)+\sum_{v_e>0}v_ey_e
       +8\sum_{\{e,f\}\in\mathcal N}|y_e-y_f|
       +M\sum_{e\in B_0}y_e.
\]
For \(A\cap B_0=\varnothing\), the last term is zero and
the rest is \(444+C(E_T\setminus A)\), by
\eqref{eq:v3-cut}. The all-sink cut \(A=\varnothing\)
has capacity 444. Any violating cut has capacity at least
2689, since all capacities are nonnegative. Thus every
minimum respects the forced omissions and the stated identity
follows from the bijection between legal cuts and supersets
of \(B_0\).

For completeness, the supplied algorithm starts with zero flow
and repeatedly augments along a source--sink path of positive
residual capacity. Each augmentation preserves capacities and
flow conservation and increases the value by a positive integer.
Every flow value is bounded by the all-sink cut, namely 444,
so termination requires at most 444 augmentations.

At termination take the vertices reachable from the source by
positive residual arcs. The sink is not reachable. Every
original arc leaving this cut is saturated, and every original
arc entering it has zero flow; otherwise a residual arc would
leave the reachable set. Summing flow conservation over that
set therefore gives flow value equal to cut capacity. For any
feasible flow, the same conservation identity gives value at
most any cut capacity. Equality thus proves optimality on both
sides. This argument accommodates antiparallel and parallel
arcs by keeping each original arc and its residual reverse
separately.
\end{proof}

The checker independently audits every original arc's capacity,
all nonterminal flow divergences, the terminal flow value, and
the cut capacity. It also recomputes \(C(B_*)\) directly from
the physical plaquette rule \eqref{eq:v2-cp}. Thus an erroneous
optimizer return value alone cannot certify the reported minimum.

Since all values in \eqref{eq:v3-quadratic} are multiples of
four, \(m(B_0)<0\) implies \(m(B_0)\le-4\). Also
\(B_0\subset B_1\) implies \(m(B_0)\le m(B_1)\), since the
feasible supersets become fewer. Failure is therefore upward
closed for this particular geometric criterion, and success is
downward closed. Neither statement classifies actual physical
bad fields outside the criterion's collar assumptions.

\subsection{Verified macroscopic face patterns on the fixed shell}

In the coordinates \((a,b,c)\in Q\), define
\[
 B_- =\{a=-1\},\qquad B_+=\{a=3\},\qquad
 B_{\rm out}=\{(a,b,c):\text{some coordinate is }-1\text{ or }3\}.
\]
The first two sets each have 25 edges, their union has 50,
and \(B_{\rm out}\) has 98. The exact flow/cut certificates are
\[
\begin{array}{l|r|r|r|r}
 \text{forced set }B_0&|B_0|&|B_*|&q(B_0)&m(B_0)\\ \hline
 \varnothing&0&0&132&-312\\
 B_-&25&35&240&-204\\
 B_+&25&53&412&-32\\
 B_-\cup B_+&50&117&444&0\\
 B_{\rm out}&98&117&444&0\\
 E_T&117&117&444&0
\end{array}
\]
The graph and forced sets specify these certificates completely;
the script constructs the primal flow and dual cut and verifies
their equality as just described. The asymmetry between the two
faces is not a claimed symmetry breaking: the excluded parent
cube \(\{0,1\}^3\) is not centred in \(\{-1,0,1,2,3\}^3\),
and the uncontrolled parent plaquettes are retained in the rule.

The first face certificate permits all boundary groups on its
25 required edges to be arbitrary. In fact the optimizer releases
35 complete edges, or 245 free boundary groups. The second releases
53 edges, or 371 free boundary groups, while retaining a negative
bound. No alignment, commutativity, or cutoff regime is imposed
on those released groups. The remaining collar still has its
specified central values at the certificate centres.

For the two-face example, optimality gives the stronger statement
\[
                 C(B)\ge0\quad\text{for every }B\supset B_-\cup B_+.
\]
Thus no further search over the same omitted-edge family can
produce a negative \(C(B)\) in that example. The flow proves
exhaustion of this sufficient geometric test, not positivity of
the actual action Hessian. The plaquettewise worst-case bounds
need not be attained simultaneously by actual completed fields.

\subsection{Original-clock consequence and its visibility condition}

\begin{theorem}[Optimized omitted-edge repair]
\label{thm:v3-repair}
Let \(C=I\), even \(L\ge32\), and fix a forced set \(B_0\)
with \(m(B_0)<0\). Let \(B_*\) be a minimizing omitted set and
assume \(F_{B_*}\subset\mathcal S\). There is a sufficiently
small closed inner chart on \(F_{B_*}\) whose visible cylinder
\(D_*\), unrestricted on all other free groups, satisfies
\[
 \kappa_{D_*}^{H_\nu}
 \ge\frac{c\beta}{2\{c\beta+8C_{\rm f}(\beta+1)\}}
 \ge\frac{c}{2(c+16C_{\rm f})}>0,
       \qquad c=\frac1{6552},\quad\beta\ge\beta_0.
\]
The constants are volume-uniform and may be chosen common for
all successful forced sets on this fixed graph and all compatible
hidden index sets. A finite union of their visible inner cylinders
has the same type of bound. The original rate-one single-group
clock and the exact completed marginal law are unchanged.
\end{theorem}

\begin{proof}
The actual angular action second derivative is at most
\(C(B_*)=m(B_0)\le-4\). A successful optimizer cannot have
\(B_*=E_T\), so there is a nonzero rotated direction.
Its independent-group squared speed is at most 1638, giving
unit Hessian upper bound \(-4/1638\), uniformly in every
unprescribed group. Lemma~\ref{lem:v2-frozen} still fixes every
pivot along this direction and makes the amplitude independent
of the rotated coordinates near the prescribed collar.

Apply the uniform-continuity, measured-curvature, and
inverse-metric allowances of Theorem~\ref{thm:v2-repair}, with
12 replaced by 4. The local joint diffusion coefficient is
\(4/(16\cdot1638)=1/6552\). The cutoff localization, global
resolvent comparison, and exact marginal domination then give
the displayed floor. Only \(F_{B_*}\) is constrained, so every
hidden group remains unrestricted, including any omitted group
declared hidden.

There are finitely many forced sets on this fixed carrier.
Taking common positive chart allowances and a finite smooth
configuration-space partition supplies the asserted uniformity
and union statement. These constants are existential; no
numerical chart radius or threshold is enclosed here. No
increasing spatial union is being used.
\end{proof}

For a practical prescribed retained set, the minimizer must meet
the displayed visibility condition. Alternatively include every
edge whose selected groups are unavailable in the forced set;
one must still ensure that the non-shell part of the collar is
retained. The optimizer does not authorize changing an already
declared block packing or marginal law.

\subsection{Verification and the resulting direction decision}

The checker is
\begin{center}\small
\path{scripts/verify_ym_f14_v3_shell_mincut.py}.
\end{center}
Its SHA--256 is
\begin{center}\scriptsize\ttfamily
470FADCB75D32C0BF6A3DCBD627F91D03D4EE3B0DB05B45436E7A118936AE7CE.
\end{center}
It verifies the 8250 exact local quadratic identities, the 264
neighbour coefficients, both unary and debit sums, six complete
integer flow/cut certificates, and direct plaquette agreement
for every minimizing set reported above. The graph optimization
has a proved global minimum certificate; its physical scope is
still only the defined curvature upper bound.

The predecessor V2 has 44,770 bytes and SHA--256
\begin{center}\scriptsize\ttfamily
A603E4C6D877062A270633891E44651B1CA1DDFE686CFE02AC9EC502C760BC65.
\end{center}
Its body is preserved apart from the title version. Active V3
replaces V2; archives and updated monographs are unchanged.
The companion and direction checkpoint remains new V5.

\begin{firewall}
This version gives an exact optimizer and exact failure test
for the fixed-shell omitted-edge certificate. It enlarges the
admissible arbitrary boundary data well beyond V2's five-edge
guarantee. It does not prove that every bad configuration has
a successful forced set, remove the central-collar restriction,
or establish an unrestricted auxiliary patch gap. The continuum
construction and physical Yang--Mills mass gap remain unproved.
\end{firewall}

The two-face result changes the next action: enumerating more
omitted supersets of that pattern cannot help this bound.
Progress there requires a different family of directions (for
example nonconstant rotation amplitudes), sharper simultaneous
plaquette control, or another complete patch mechanism. Any such
step must be checked against the same completed law and all
unprescribed fields. The full goal remains active.
\section{V4: graded rotations cross the two-face obstruction}
\label{sec:v4}

\subsection{The direction changes, not the completion or clock}

V3 proves that the omitted-edge certificate has minimum zero when
both opposite outer faces must be arbitrary. That statement
exhausts its binary common-rotation upper bound, not all directions.
We now allow different rotation amplitudes on the remaining coarse
edges and obtain an explicit strictly negative certificate for
that same two-face set. All 350 free boundary groups on the two
faces remain arbitrary.

Weighted rotations themselves are not new: f12 V82 gives a
56-coordinate direction for a central sheet with frozen exterior,
and f12 V83 distinguishes its continuous descent from a binary
cut obstruction. We do not reprove or relabel that result.
The new statement is the two-face, all-unprescribed-field upper
form in the present non-wrapping geometry, an exact integer
witness, and its transfer to the original interface heat bath.
The specified hidden fields and arbitrary completed face pivots
are never replaced by a central exterior.

\subsection{The actual graded-shell upper form}

Fix a set \(H\subset E_T\) of arbitrary boundary edges. Prescribe
the central values only on \(F_H\), as in V2, leaving all other
free groups arbitrary. Put \(A=E_T\setminus H\). For a vector
\(w=(w_e)_{e\in A}\in\mathbb R^A\), choose a unit imaginary
quaternion \(E\) and move each of the seven free boundary links
on \(e\in A\) by
\[
                          Y_l(t)=-e^{t w_eE}.
\]
Every other free group is fixed during this test curve.
For sufficiently small \(|t|\), the fourteen sources on each
\(e\in A\) remain near \(-I\), so its actual pivot remains
exactly \(I\). The pivots on \(H\) may be arbitrary completed
values, but none of their source inputs is rotated; they are
constant along the curve. Every other pivot is also constant.
This is the same incidence argument as
Lemma~\ref{lem:v2-frozen}, now with one amplitude per edge.

For a plaquette in the original \(\mathcal P_T\), count how
many of its selected links belong to edges in \(A\), including
those with amplitude zero. Denote the sets with exactly two
and exactly one by \(\mathcal P_{2,A}\) and
\(\mathcal P_{1,A}\). For a two-selected plaquette let
\(e(p),f(p)\) be the owners of its opposite selected links;
they can coincide. For a one-selected plaquette let \(e(p)\)
be that link's owner. Define the real symmetric matrix \(M_H\)
by the quadratic form
\begin{equation}\label{eq:v4-form}
 w^{\mathsf T}M_Hw
 =2\sum_{p\in\mathcal P_{2,A}}(w_{e(p)}-w_{f(p)})^2
      +\sum_{p\in\mathcal P_{1,A}}c_p(H)w_{e(p)}^2,
\end{equation}
using V2's actual plaquette upper coefficients
\eqref{eq:v2-cp}. This is a finite integer matrix; it depends
only on the declared arbitrary-edge set, not on arbitrary field
values. Its value at the all-one vector is \(C(H)\).

\begin{proposition}[Graded completed-fibre curvature]
\label{prop:v4-form}
Uniformly over all unprescribed fine groups,
\[
                        S''(0)\le w^{\mathsf T}M_Hw.
\]
The independent free-group squared speed is
\(14\sum_{e\in A}w_e^2\). The inverse-pivot Haar amplitude
is independent of the rotated coordinates on the relevant
open cylinder.
\end{proposition}

\begin{proof}
No pivot moves, by the preceding source-incidence argument.
A two-selected plaquette belongs to the original fixed class:
both selected links have central sign \(-1\), and the
transverse links are central. Its action change is
\(2(1-\cos(t(w_e-w_f)))\), so its second derivative is
\(2(w_e-w_f)^2\). If both owners agree, the term vanishes.

For a one-selected plaquette the rotated factor, or its inverse,
has second derivative \(-w_e^2\) times its original value.
Thus the action second derivative is \(2w_e^2q_0\), where
\(q_0\) is the scalar part of the actual completed holonomy.
On fixed plaquettes this is the central signed coefficient;
on uncontrolled ones, including those touching arbitrary
face pivots, \(q_0\le1\) gives precisely \(c_p(H)w_e^2\).
Plaquettes with no selected link contribute zero. This proves
the upper bound without requiring simultaneous attainment of
the uncontrolled plaquette maxima.

Each of seven independent links on \(e\) has squared speed
\(2w_e^2\). Their sum gives the displayed metric norm.
There is no additional independent pivot speed. The selected
edge Jacobian factors are one on the inactive cylinder, and
the other factors have no selected-coordinate dependency.
\end{proof}

The same formula extends to three-component angular amplitudes
\(\mathbf w_e\in\mathbb R^3\): replace each square by the
squared Euclidean norm of the vector or vector difference.
Indeed the trace of the mixed second-order quaternion product
depends on the dot product, while its commutator part has zero
scalar trace. One-selected terms are multiplied by
\(|\mathbf w_e|^2\). Consequently the upper form is
\(M_H\otimes I_3\), and the metric is
\(14\sum_e|\mathbf w_e|^2\). This assertion concerns the
second derivative at the prescribed collar, not a finite-angle
commutativity identity for independent axes.

Unlike V3, an amplitude zero on an unomitted edge does not make
its physical boundary fields arbitrary. Those fields still
have their prescribed central values. The present form is
therefore a new direction family on the declared carrier, not
an unconstrained relaxation of V3's binary formula.

\subsection{An exact witness with both outer faces arbitrary}

Take the same forced set as the obstruction in V3:
\[
                     H=B_-\cup B_+=\{(a,b,c)\in Q:a=-1\text{ or }a=3\}.
\]
Then \(|H|=50\) and \(|A|=67\). The remaining cells have
\(a\in\{0,1,2\}\); those with all three coordinates in
\(\{0,1\}\) are excluded as before. Define the integer
amplitudes by this complete table. Dashes are excluded parent
cells, not zero-amplitude variables.
\[
\begin{array}{cc|rrrrr}
 &&\multicolumn{5}{c}{c\text{ coordinate}}\\
 a&b&-1&0&1&2&3\\ \hline
 0&-1&2&3&8&23&21\\
 0&0&3&-&-&40&38\\
 0&1&8&-&-&57&51\\
 0&2&23&40&57&76&56\\
 0&3&21&38&51&56&38\\ \hline
 1&-1&3&4&11&31&27\\
 1&0&4&-&-&54&50\\
 1&1&11&-&-&76&67\\
 1&2&31&54&76&100&73\\
 1&3&27&50&67&73&49\\ \hline
 2&-1&3&5&10&18&15\\
 2&0&5&8&15&32&27\\
 2&1&10&15&24&45&36\\
 2&2&18&32&45&54&38\\
 2&3&15&27&36&38&25
\end{array}
\]

\begin{theorem}[Two-face repair direction]
\label{thm:v4-direction}
For the displayed vector and all unprescribed fine fields,
\[
 S''(0)\le-175496,\qquad
 \sum_e w_e^2=110478,\qquad
                  \|\dot Y(0)\|_{\rm HS}^2=1546692.
\]
The resulting unit-direction action Hessian is at most
\[
                -\frac{43874}{386673}\le-\frac3{28}<0.
\]
The 350 free boundary groups on the two outer faces may take
arbitrary noncommuting values in \(\SU(2)\), with any cutoff
regimes. All other unprescribed fine groups are also unrestricted.
\end{theorem}

\begin{proof}
Assemble \(M_H\) from \eqref{eq:v4-form} in lexicographic
coarse-coordinate order. Exact integer multiplication gives
\(w^{\mathsf T}M_Hw=-175496\); independently summing the
plaquette contributions gives the same integer. Squaring the
67 table entries gives 110478. These computations are reproduced
by the checker, with no floating-point eigenvalue claim.
The rational inequality follows, for example, from
\(2\cdot175496>3\cdot110478\). Apply
Proposition~\ref{prop:v4-form} and divide by the actual
independent-group squared speed \(14\cdot110478\).
\end{proof}

The ungraded direction on these 67 edges has
\(\mathbf1^{\mathsf T}M_H\mathbf1=C(H)=84\), which does
not certify descent. V3 goes further and gives minimum zero
over every binary omitted superset of these faces. The strict
negative integer above is therefore an actual improvement
beyond that exhausted sufficient family. It does not contradict
V3's optimizer theorem. A preliminary numerical eigenvector was
used only to choose the integer candidate; the accepted checker
contains no numerical eigensolver and proves only the displayed
rational witness, not an exact lowest eigenvalue.

\subsection{Original-clock consequence}

\begin{theorem}[Uniform two-face visible repair]
\label{thm:v4-repair}
Assume \(C=I\), even \(L\ge32\), and
\(F_H\subset\mathcal S\) for the two-face set \(H\) above.
There is a sufficiently small closed inner chart on \(F_H\)
whose visible cylinder \(D_H\), unrestricted in all other
free groups, satisfies
\[
 \kappa_{D_H}^{H_\nu}
 \ge\frac{c\beta}{2\{c\beta+8C_{\rm f}(\beta+1)\}}
 \ge\frac{c}{2(c+16C_{\rm f})}>0,
             \qquad c=\frac3{448},\quad\beta\ge\beta_0.
\]
The constants are independent of ambient volume and of compatible
hidden index sets. The operator is the original rate-one
single-retained-group heat bath for the exact fixed-coarse
interface marginal. Chart radii and \(\beta_0\) are existential.
\end{theorem}

\begin{proof}
Normalize the integer direction in product orthonormal normal
coordinates about the central \(F_H\)-data. Its action Hessian
is at most \(-\gamma\), with \(\gamma=3/28\), uniformly
over all unprescribed groups. Uniform continuity on the fixed
finite input closure supplies an outer chart where its frozen
coordinate second derivative is at most \(-\gamma/2\).
The proof does not restrict the 350 face boundary groups to
neighbourhoods of selected points: their entire compact carrier
is retained in this continuity argument.

The unomitted sources remain inactive. Every other pivot and
Jacobian factor has no differentiated coordinate input, by
Proposition~\ref{prop:v4-form}. The coordinate-Haar second
derivative is bounded, so for large \(\beta\) the actual
measured potential satisfies \(\partial_a^2\Phi_\beta
\le-\gamma\beta/4\). With inverse metric lower bound
\(1/2\), the ground-state conjugation of
Theorem~\ref{thm:repair} gives local diffusion floor
\(\gamma\beta/16=(3/448)\beta\).

A smooth cutoff between the inner and outer visible charts,
with finite derivative bound, gives the same localization
inequality as in Theorem~\ref{thm:v2-repair}. Apply the global
resolvent comparison, low-spectral projection estimate, and
exact marginal form domination already proved there. The
support leaves every hidden variable unrestricted, including
face variables if they are declared hidden. No change of
update clock or additional collective update is introduced.
\end{proof}

\subsection{An exact limit of this graded upper form}

Now let all outer-boundary coarse edges be arbitrary:
\[
 H_{\rm out}=\{(a,b,c)\in Q:\text{some coordinate is }-1\text{ or }3\}.
\]
There are 98 such edges and 19 remaining amplitude variables.
This is a different, larger arbitrary-field requirement than
the two-face theorem. Let \(M_{\rm out}\) be the 19-dimensional
matrix from \eqref{eq:v4-form}.

\begin{proposition}[Positive upper form for the full outer boundary]
\label{prop:v4-outer}
The exact integer matrix satisfies
\[
                             M_{\rm out}\succeq9I_{19}.
\]
Hence no nonzero scalar or three-component coarse-edge amplitude
vector makes this sufficient upper form negative.
\end{proposition}

\begin{proof}
The checker constructs \(A=M_{\rm out}-9I\) and factors it
over exact rational numbers as
\(A=\mathsf L\mathsf D\mathsf L^{\mathsf T}\), with
unit lower-triangular \(\mathsf L\) and positive diagonal
\(\mathsf D\). Explicitly, in lexicographic order, it uses
\[
 D_j=A_{jj}-\sum_{k<j}L_{jk}^2D_k,\qquad
 L_{ij}=\frac{A_{ij}-\sum_{k<j}L_{ik}L_{jk}D_k}{D_j}
                         \quad(i>j).
\]
All 19 pivots are positive; the smallest is exactly
\(5462289/477655\). The checker then reconstructs every
entry of \(A\) from the factors and verifies equality.
This is a rational positive-definiteness certificate, not
a rounded numerical eigenvalue. It proves \(A\succ0\)
and therefore the stated weaker lower bound. The vector-axis
extension follows by applying the same inequality to each of
the three components of \(M_{\rm out}\otimes I_3\).
\end{proof}

The lower bound here is on an \emph{upper-bound matrix} for
the action curvature. It gives no lower bound on the actual
action Hessian. In particular it does not prove local stability,
metastability, or failure of the heat bath. It only rules out
strict negativity certified by \eqref{eq:v4-form} for this
outer-boundary case within the stated edge-constant amplitude
family. Individual-link amplitudes, other coordinate directions,
or joint estimates on uncontrolled plaquettes are not tested
by this proposition.

\subsection{Verification and the next acquisition}

The checker is
\begin{center}\small
\path{scripts/verify_ym_f14_v4_graded_shell.py}.
\end{center}
Its SHA--256 is
\begin{center}\scriptsize\ttfamily
87A6A9C41A1B61D55A7F409C2416548EA7ED4D30BF2A104374070A0B143786A6.
\end{center}
It builds the graded matrices, verifies symmetry and the
all-one agreement with V2's curvature bound, checks the integer
witness both by matrix multiplication and direct plaquette sum,
reruns V3's exact two-face minimum-cut certificate, and verifies
the full rational factorization of \(M_{\rm out}-9I\).
The accepted calculation uses integers and fractions only.

The predecessor V3 has 59,329 bytes and SHA--256
\begin{center}\scriptsize\ttfamily
735330C9E42435DAFFDDA3C47E70C033C7EE5E119890EFF3C8A86E2A75DCC3A1.
\end{center}
Its complete body is preserved apart from the title version.
Active V4 replaces V3. Archives and the updated monographs
remain unchanged. The next iteration, V5, includes the required
companion review and a direction review.

\begin{firewall}
The two-face case that defeated V3 now has a genuine
volume-uniform local repair floor with 350 arbitrary boundary
groups and all hidden fields free. This is conditional on
the remaining visible collar being near its central pattern
and on identity coarse data. The outer-boundary case remains
uncertified by this particular graded upper form. No unrestricted
conditional patch gap, cofinal continuum construction, or
physical mass-gap theorem is proved.
\end{firewall}

At the next checkpoint, the useful question is whether allowing
amplitudes to vary among the seven links of each remaining edge
changes the outer-boundary certificate, or whether its uncontrolled
plaquette debit must instead be improved. The positive rational
factorization means that merely searching more coarse-edge
amplitude vectors cannot improve that case. The full goal remains
active.
\section{V5: checkpoint and a completed-pivot-cap repair of the outer boundary}
\label{sec:v5}

\subsection{Required companion and direction review}

The active companion files, the updated SM summary, and their
terminal scope statements were checked again. They are unchanged
from this lineage's restart:
\begin{center}\small
\begin{tabular}{lrr}
\toprule
Source&Bytes&SHA--256 prefix\\ \midrule
SM--Einstein V72&607372&\texttt{F44F9745D43379E1}\\
NS stability V26&278283&\texttt{82C887F8C2C69E4F}\\
Parallel YM V5&54174&\texttt{306F1BB84CFA8A8D}\\
Updated SM monograph (3)&384982&\texttt{54176A3730275E1B}\\
\bottomrule
\end{tabular}
\end{center}
The SM summary still separates fixed-spatial-mode, fixed-time
many-copy transport from finite physical species and ultraviolet
interchange. Its regular-branch continuum theorem remains conditional
on realizing its hypotheses cofinally. NS V26 supplies rank-one
norm results, not global regularity or a YM patch inequality.
Parallel YM's extensive-charge obstruction gives no new repair
bound; its internal stopping language is not an instruction to
stop this task. The YM closure monograph V4 is also unchanged.
No companion result is used to fill a missing YM hypothesis.

The mathematical review changes the repair estimate upstream.
First we allow independent amplitudes on all 133 remaining
negative shell links. The old, unitarity-only bound remains
positive even in this larger space, so changing amplitudes alone
does not suffice. We then retain a fact that the old bound threw
away: at fixed coarse identity, an actual completed pivot lies
within \(7\rho/8\) of identity. Applying that inherited constraint
to 27 specific plaquettes makes the enlarged bound negative.

The radius bound is not a new pivot theorem. Its new use here is
in the actual outer-boundary repair geometry, with an independent-
link rational witness and all outer free boundary groups arbitrary.
The next companion and direction review is V10.

\subsection{Independent-link amplitudes and the old bound's exact limit}

Fix the outer set \(H=H_{\rm out}\) from V4. Its 98 coarse edges
carry 686 selected free boundary groups. The remaining set
\[
             U=T\setminus T_H
\]
has 133 free boundary groups, on 19 coarse edges. Prescribe only
the central values on \(F_H=F\setminus T_H\), with
\(|F_H|=9348\), and leave every other free group arbitrary.
In particular this does not prescribe any particular values
on the 686 released boundary groups.

For real amplitudes \(z=(z_l)_{l\in U}\) and a unit imaginary
quaternion \(E\), set \(Y_l(t)=-e^{t z_lE}\) for \(l\in U\).
All other independent groups stay fixed along this curve.
Different links on one coarse edge may have different amplitudes.
Nevertheless every corresponding nonroot source remains near
\(-I\) for sufficiently small \(|t|\). The actual pivot on
that edge is therefore \(I\). Pivots on \(H\) have no moving
source inputs. Thus every actual pivot is constant along the
test curve, just as in V4; the inner/source-transition behaviour
of a pivot on \(H\) is not differentiated by this direction.

Let \(K_0\) be the linkwise version of \eqref{eq:v4-form}:
\begin{equation}\label{eq:v5-old}
 z^{\mathsf T}K_0z
 =2\sum_{p:\,|\partial p\cap U|=2}(z_{l(p)}-z_{k(p)})^2
       +\sum_{p:\,|\partial p\cap U|=1}c_p(H)z_{l(p)}^2.
\end{equation}
The sums run over \(\mathcal P_T\). This is an integer symmetric
matrix with nonpositive off-diagonal entries. The same plaquette
argument as in Proposition~\ref{prop:v4-form} gives
\(S''(0)\le z^{\mathsf T}K_0z\), with squared independent-group
speed \(2\sum_lz_l^2\). The factor is now 2, not 14, because
each amplitude belongs to one link rather than seven links.
If \(J\) repeats one coarse-edge amplitude on its seven links,
the exact compression is \(J^{\mathsf T}K_0J=M_{\rm out}\).

\begin{proposition}[No negative value of the old linkwise upper form]
\label{prop:v5-old-positive}
The old bound satisfies \(K_0\succeq\tfrac15 I_{133}\).
Consequently it does not certify a negative value even with
independent scalar or three-component link amplitudes.
\end{proposition}

\begin{proof}
The following exact positive-vector certificate avoids any rounded
eigenvalue claim. Define integer vectors by
\[
 v^{(0)}=\mathbf1,\qquad
 v^{(j+1)}=(12I-K_0)v^{(j)},\quad 0\le j<256,
                  \qquad v=v^{(256)}.
\]
The checker verifies every coordinate remains positive, and
\(5(K_0v)_i\ge v_i\) for every \(i\). These are integer
computations on the matrix explicitly defined by
\eqref{eq:v5-old}. It also checks nonpositive off-diagonal
entries and exact agreement with the coarse compression above.

For any positive vector \(v\) and any real symmetric matrix
\(K\) with nonpositive off-diagonal entries, direct expansion gives
\[
 x^{\mathsf T}Kx
 =\sum_{i<j}(-K_{ij})v_iv_j
        \left(\frac{x_i}{v_i}-\frac{x_j}{v_j}\right)^2
       +\sum_i\frac{(Kv)_i}{v_i}x_i^2.
\]
Apply this identity to the certified \(v\) and \(K=K_0\).
All the first terms are nonnegative and every last coefficient
is at least \(1/5\). This proves the assertion. The
three-component version is the sum of three copies of the form.
\end{proof}

This is again positivity of an upper-bound matrix, not positivity
of the actual action Hessian. It identifies the loss of information
in the estimate; it does not prove local stability or obstruct a
physical mass gap.

\subsection{The information lost by treating a pivot as an arbitrary unitary}

\begin{lemma}[Inherited compact cap at fixed coarse identity]
\label{lem:v5-cap}
For any actual completed pivot at \(C_e=I\), with arbitrary
source inputs,
\[
                    d_{\rm geo}(I,B_e)\le R=\frac{7\rho}{8}
                                      =\frac7{512}.
\]
Writing \(B_e=(b_0,\mathbf b)\) as a unit quaternion,
\[
                         b_0\ge1-\frac{R^2}{4}.
\]
\end{lemma}

\begin{proof}
The exact completed equation is \(B_e e^{X_e}=I\), with
\(X_e=16^{-1}\sum_{j=1}^{14}g_\rho(B_e^{-1}P_j)\).
The bound \(\|g_\rho\|_{\rm HS}\le\rho\) gives
\(\|X_e\|_{\rm HS}\le14\rho/16\). Since
\(B_e=e^{-X_e}\), its geodesic distance is no larger than
that norm. This is the inherited cap, derived here directly
to specify which metric is used. On \(\SU(2)\),
\(b_0=\cos(d_{\rm geo}(I,B_e)/\sqrt2)\) in this ball.
The inequality \(\cos t\ge1-t^2/2\) gives the second statement.
\end{proof}

Let \(\mathcal P_{\rm cap}\) consist of the original fixed
one-selected plaquettes with their opposite parallel physical
link equal to a root pivot on \(H\), and their selected link
in \(U\). Plaquettes touching the original parent-dependent
physical set \(A_0\) are not included. The finite incidence
calculation gives
\[
                            |\mathcal P_{\rm cap}|=27.
\]
All other factors of such a plaquette are central at the
prescribed \(F_H\)-data. In particular their free input
closures lie in \(F_H\), not in the released groups. The
checker verifies this dependency statement link by link.

The selected link has value \(-I\) and the original central
plaquette product is \(-I\). Hence the actual holonomy's scalar
part is \(-b_0\), whether the root pivot is traversed forward
or backward. Lemma~\ref{lem:v5-cap} improves its coefficient
in \eqref{eq:v5-old} from the previous upper value \(+2\) to
\begin{equation}\label{eq:v5-coeff}
               -2+\epsilon,\qquad
                  \epsilon=\frac{R^2}{2}=\frac{49}{524288}.
\end{equation}
This small positive loss is retained, not rounded to zero.
No source alignment or cutoff inactivity on the root's own
edge is assumed.

Let \(r_l\) count the members of \(\mathcal P_{\rm cap}\)
incident to \(l\in U\), and define the rational corrected matrix
\begin{equation}\label{eq:v5-corrected}
                       K=K_0-(4-\epsilon)\operatorname{diag}(r_l).
\end{equation}
All off-diagonal terms are unchanged. In particular, this is not
a claim that every uncontrolled plaquette can be improved:
the unknown transverse factors on the excluded plaquettes
prevent the cap from determining their scalar traces.

\begin{theorem}[Completion-aware linkwise curvature bound]
\label{thm:v5-form}
For every real amplitude vector \(z\), at the prescribed
central \(F_H\)-data and for every choice of all unprescribed
fine fields, the actual completed action satisfies
\[
                              S''(0)\le z^{\mathsf T}Kz.
\]
The analogous three-component bound is \(K\otimes I_3\).
\end{theorem}

\begin{proof}
Keep all the original two-selected terms, which are exact,
and all the original one-selected bounds outside
\(\mathcal P_{\rm cap}\). For each plaquette in that set,
its second derivative is \(-2b_0z_l^2\). The actual pivot
is fixed along the test curve; its value still satisfies the
cap. Apply \eqref{eq:v5-coeff} and sum. The correction is
exactly the diagonal debit in \eqref{eq:v5-corrected}.
The three-component assertion follows from the scalar-trace
second-jet identity used in V4, with \(z_l^2\) replaced by
the squared vector norm. No independent choice of pivot
values is used anywhere in the estimate.
\end{proof}

\subsection{An explicit integer witness for all 98 outer edges}

The following table specifies the 133 link amplitudes. A row
\((a,b,c)\) represents a remaining coarse edge with root
\((a,b,0,c)\). A column labelled \(uvw\) represents its
fine link \(((2a+u,2b+v,1,2c+w),2)\). These are precisely
the seven nonroot boundary links, with \(uvw\ne000\).
\begin{center}\small
\begin{tabular}{c|rrrrrrr}
\toprule
$(a,b,c)$&001&010&011&100&101&110&111\\ \midrule
$(0,0,2)$&88&125&95&125&95&176&142\\
$(0,1,2)$&233&222&157&349&256&326&239\\
$(0,2,0)$&125&88&95&125&176&95&142\\
$(0,2,1)$&222&233&157&349&326&256&239\\
$(0,2,2)$&273&273&116&596&322&322&172\\
$(1,0,2)$&233&349&256&222&157&326&239\\
$(1,1,2)$&617&610&415&610&415&585&392\\
$(1,2,0)$&349&233&256&222&326&157&239\\
$(1,2,1)$&610&617&415&610&585&415&392\\
$(1,2,2)$&708&708&299&1000&503&503&261\\
$(2,0,0)$&125&125&176&88&95&95&142\\
$(2,0,1)$&222&349&326&233&157&256&239\\
$(2,0,2)$&273&596&322&273&116&322&172\\
$(2,1,0)$&349&222&326&233&256&157&239\\
$(2,1,1)$&610&610&585&617&415&415&392\\
$(2,1,2)$&708&1000&503&708&299&503&261\\
$(2,2,0)$&596&273&322&273&322&116&172\\
$(2,2,1)$&1000&708&503&708&503&299&261\\
$(2,2,2)$&745&745&308&745&308&308&150\\
\bottomrule
\end{tabular}
\end{center}

\begin{theorem}[Outer-boundary descent on the actual fibre]
\label{thm:v5-witness}
For the table vector \(z\),
\[
 \sum_lz_l^2=21328389,\qquad
 z^{\mathsf T}Kz=-\frac{1964257089399}{262144}.
\]
Uniformly in all the arbitrary outer boundary groups and every
other unprescribed fine field, the resulting unit-direction
action Hessian is therefore at most
\[
                    -\frac{218250787711}{1242468712448}
                                 \le-\frac16<0.
\]
\end{theorem}

\begin{proof}
The checker multiplies the rational matrix in
\eqref{eq:v5-corrected} by the fixed integer vector and
independently sums its plaquette terms. Both give the stated
rational number. Squaring the table entries gives the displayed
norm. Integer comparison verifies
\(z^{\mathsf T}Kz\le-\tfrac13\sum_lz_l^2\).
Apply Theorem~\ref{thm:v5-form} and divide by the actual
free-group squared speed \(2\sum_lz_l^2\). The product-group
curve is geodesic in the independent variables, so the normalized
second derivative is its action Hessian.
\end{proof}

A numerical eigenvector was used only to select this integer
candidate. The accepted checker contains no numerical eigensolver
and proves the displayed rational inequality, not an exact
lowest eigenvalue. The old form's positive-vector certificate
and the corrected form's negative witness are compatible because
the diagonal estimates have changed, not the physical law.

\subsection{Original-clock repair with the entire outer boundary arbitrary}

\begin{theorem}[Uniform outer-boundary visible repair]
\label{thm:v5-repair}
Let \(C=I\), even \(L\ge32\), and assume
\(F_H\subset\mathcal S\) for \(H=H_{\rm out}\).
There is a sufficiently small closed inner product chart on
\(F_H\) whose visible cylinder \(D_{\rm out}\), unrestricted
on all other groups, has the original-clock floor
\[
 \kappa_{D_{\rm out}}^{H_\nu}
 \ge\frac{c\beta}{2\{c\beta+8C_{\rm f}(\beta+1)\}}
 \ge\frac{c}{2(c+16C_{\rm f})}>0,
                    \qquad c=\frac1{96},\quad\beta\ge\beta_0.
\]
The constants are independent of ambient volume and compatible
hidden index sets. All 686 free boundary groups on the 98
outer edges may be arbitrary, with no choice of cutoff regime
or source commutativity. They may also be hidden variables.
The remaining 9348 collar groups are still constrained to the
stated chart. Chart radii and \(\beta_0\) are existential.
\end{theorem}

\begin{proof}
Normalize the integer vector in product orthonormal normal
coordinates about the central \(F_H\)-data. Its action Hessian
is at most \(-1/6\) for all values of the unprescribed inputs.
Uniform continuity on the fixed finite compact input closure
gives an outer visible chart with frozen coordinate second
derivative at most \(-1/12\). This uses compactness of the
entire arbitrary outer-field carrier, not a chosen configuration
of those fields.

In that chart every unomitted source remains near \(-I\), so
its pivot and inverse Jacobian are independent of the rotated
coordinates. Every other pivot has no such input. The true
coordinate-Haar Hessian is bounded. Hence for sufficiently
large \(\beta\), the measured potential has
\(\partial_a^2\Phi_\beta\le-\beta/24\). Taking inverse metric
at least \(1/2\), ground-state conjugation supplies the local
joint diffusion floor \(\beta/96\).

Choose an inner-to-outer visible cutoff with finite gradient
bound. The localization, global resolvent comparison, and
low-spectral projection argument in Theorem~\ref{thm:repair}
give the displayed joint heat-bath floor. Exact marginal
conditional-variance domination transfers it to \(\nu\),
because the lifted support restricts no hidden variable.
The update rate is still one per retained free group. The
weighted rotation is only a test direction; it is not an
extra transition kernel. The original non-wrapping stencil
and finite input closure prove the volume uniformity.
\end{proof}

The gained freedom is specific: the entire outer layer of
direction-2 free boundary links is arbitrary, but not every
internal or transverse field in those cells. In particular,
the other factors of the 27 sharpened plaquettes require
their stated input closure in the retained central collar.
The theorem is not an all-exterior conditional patch inequality.

\subsection{Verification record and the direction after the checkpoint}

The new checker is
\begin{center}\small
\path{scripts/verify_ym_f14_v5_link_shell.py}.
\end{center}
Its SHA--256 is
\begin{center}\scriptsize\ttfamily
5EEF2BC9164BD13165915F485C020FEE8583246CB005DEECDEA6562DB019A699.
\end{center}
It verifies the independent-link matrix, exact compression to
V4's matrix, the 256-step integer positive-vector certificate,
all 27 cap-plaquette dependencies, the rational diagonal
correction, and the fixed witness by both matrix multiplication
and direct plaquette evaluation. All accepted numerical claims
use integer or rational arithmetic. It does not supply numerical
chart radii or a continuum estimate.

The predecessor V4 has 73,666 bytes and SHA--256
\begin{center}\scriptsize\ttfamily
8ECD6618E799C1B453ED4C703D430C02B9C80C3D7BBBA254B1B790DC6734FAD2.
\end{center}
Its complete body is preserved apart from the title version.
Active V5 replaces V4. Archives, companion notes, and both
updated monographs are unchanged. The next companion review
and direction checkpoint is new V10.

\begin{firewall}
This version overcomes the outer-boundary obstruction to the
old upper forms by retaining an actual completed-pivot
constraint. It proves a local original-clock repair floor
with 686 arbitrary boundary groups and all hidden fields
unrestricted, conditional on the remaining central collar
and identity coarse data. It does not cover arbitrary
internal/transverse disorder or prove the required unrestricted
patch gaps. Cofinal continuum construction and the physical
Yang--Mills mass gap remain unproved.
\end{firewall}

The checkpoint identifies the next genuine coverage obligation:
replace the remaining central internal/transverse collar
conditions by a criterion on actual controlled plaquette or
source data, or provide a different all-exterior patch estimate.
More amplitude searches within the same collar do not by
themselves remove that condition. Any relaxation must retain
the actual pivot cap, its source dependence, and all hidden
inputs, rather than reinserting independent unitary pivots.
The full goal remains active.
\section{V6: releasing unnecessary pivot-source inputs from the collar}
\label{sec:v6}

\subsection{The hypothesis being removed}

V5 constrains 9348 free groups to a small central collar, although
the descent rotates only 133 groups. Much of that collar was
inherited from a rule that fixes a physical pivot by fixing all
its free source inputs. At coarse identity, this is unnecessary
for pivots whose only required property is the compact radius
bound of Lemma~\ref{lem:v5-cap}.

We now prescribe only the direct free factors needed for the
rotated sources and useful plaquettes. Other root pivots retain
their actual values and arbitrary source inputs; their effects
are bounded by geodesic product estimates. The resulting visible
collar has 769 free groups. The 8579 additional groups released
from V5's collar may take arbitrary values, not just values in
slightly larger charts. The same uniform local heat-bath floor
coefficient \(1/96\) survives.

This does not remove every collar condition or prove a global
patch inequality. It specifically replaces unnecessary control
of solved-pivot inputs by the inherited bound on the solved
pivots themselves. No new completion or density is introduced.

\subsection{Definition of the reduced independent-coordinate collar}

Keep \(C=I\), the outer set \(H=H_{\rm out}\), and the
133-link set \(U=T\setminus T_H\) of V5. Let \(E_U\)
be its 19 owning coarse edges. Use exactly the positive integer
amplitudes \(z_l\) of V5's witness table.

Let \(\mathcal R_U\) be the union of the non-tree free links
in the fourteen nonroot source paths on every edge in \(E_U\).
These paths contain no root pivots. There are 285 such free
inputs, including \(U\); their 152 other groups are internal
companion links in these source paths.

Let \(\mathcal P_U\) contain the plaquettes incident to \(U\),
and let \(A_0\) be the original parent-dependent physical set
from V2. Define
\[
 \mathcal P_2=\{p\in\mathcal P_U:|\partial p\cap U|=2\},
\]
and
\[
 \mathcal P_-=
 \{p\in\mathcal P_U:|\partial p\cap U|=1,
        \ \partial p\cap A_0=\varnothing,
        \ \prod_{l\in\partial p}\epsilon_l=-1\}.
\]
Here \(\epsilon_l\) is the original central sign convention.
Write \(\mathcal N(p)\) for the physical links of \(p\)
that are neither fixed tree links nor root pivots. These are
independent free coordinates. Set
\begin{equation}\label{eq:v6-collar}
 G=U\cup\mathcal R_U\cup
              \bigcup_{p\in\mathcal P_2\cup\mathcal P_-}\mathcal N(p).
\end{equation}
Unlike the old input-closure rule, \eqref{eq:v6-collar} does
\emph{not} replace the roots in those plaquettes by their source
inputs. The only sources deliberately constrained are those of
the 19 pivots whose groups are being rotated.

The exact incidence calculation gives
\[
 |G|=769,\qquad |G\setminus U|=636,\qquad
             G\subset F_H,\qquad |F_H\setminus G|=8579.
\]
Every member of \(G\) is an independent non-tree, nonpivot
group. The central value is \(-I\) on \(U\) and \(I\)
on \(G\setminus U\). The set \(G\) is disjoint from the
original parent's 944 hidden free groups and from all 686
outer boundary groups. Those disjointness statements and the
inclusion \(G\subset F_H\) are explicitly checked.

For this version prescribe only these central values on \(G\),
and leave every other independent fine group arbitrary. For
an interface marginal the visibility requirement is now
\(G\subset\mathcal S\), not \(F_H\subset\mathcal S\).
Groups in \(F_H\setminus G\) can in particular be hidden.

\subsection{A cap on ordered products, without commuting the roots}

Let \(R=7/512\) and \(\epsilon=R^2/2=49/524288\), as in
V5. Every actual root pivot at coarse identity has geodesic
distance at most \(R\) from \(I\), independently of all
its source values.

\begin{lemma}[Uniform scalar control of a root product]
\label{lem:v6-product}
If \(B_1,\ldots,B_k\in\SU(2)\) have geodesic distance at
most \(R\) from identity, then any ordered product of them
or their inverses has scalar quaternion part at least
\[
                             1-\frac{k^2R^2}{4}.
\]
No common axis or source alignment is required.
\end{lemma}

\begin{proof}
Bi-invariance and the triangle inequality give
\(d_{\rm geo}(I,P)\le kR\) for the ordered product.
The Hilbert--Schmidt chordal distance is no larger than this.
For a unit quaternion \(P=(p_0,\mathbf p)\),
\(\|I-P\|_{\rm HS}^2=4(1-p_0)\). Therefore
\(1-p_0\le k^2R^2/4\). Inversion preserves the distances,
and the argument does not interchange any factors.
\end{proof}

Under the curve \(Y_l(t)=-e^{tz_lE}\), \(l\in U\), every
pivot on \(E_U\) is exactly \(I\) for small \(|t|\):
its source companions are prescribed in \(\mathcal R_U\)
and all its nonroot products remain near \(-I\). Every other
pivot has no selected input and is constant along the curve.
It need not equal \(I\); Lemma~\ref{lem:v6-product} is used
in place of that stronger, unjustified assertion.

For \(p\in\mathcal P_-\), let \(k_p\) be the number of
its root factors whose owner is not in \(E_U\). All its
nonroot factors are prescribed by \(G\). Its scalar holonomy
is the negative of the scalar part of an ordered product of
those \(k_p\) root factors (possibly inverted); central signs
commute and multiply to \(-1\). Thus, for its selected link
\(l(p)\), the actual second derivative is at most
\begin{equation}\label{eq:v6-negative}
                        (-2+\epsilon k_p^2)z_{l(p)}^2.
\end{equation}
The pivot values can be mutually dependent and noncommuting.
The bound applies simultaneously to all of them by pointwise
inequality, without a probabilistic independence assumption.

\subsection{The finite incidence and curvature accounts}

The 537 incident plaquettes have the following classification:
\begin{center}\small
\begin{tabular}{lr}
\toprule
Class&Number\\ \midrule
Two selected links; no root factors&261\\
One selected link in \(\mathcal P_-\), \(k_p=0\)&57\\
One selected link in \(\mathcal P_-\), \(k_p=1\)&30\\
One selected link in \(\mathcal P_-\), \(k_p=2\)&27\\
Other one-selected plaquettes, unrestricted&162\\
\bottomrule
\end{tabular}
\end{center}
The incidence identity \(2\cdot261+57+30+27+162=6\cdot133\)
checks the six plaquettes per selected link. In particular,
none of the two-selected plaquettes contains a root. Their
nonselected physical factors are tree links or free links
prescribed by \(G\), and their common central product is
positive. Their second derivative is exactly
\(2(z_l-z_m)^2\).

The other 162 one-selected plaquettes need no control on their
unprescribed factors: twice the scalar part of their actual
holonomy is at most 2. With \(\mathcal P_?\) denoting this
last class, define
\begin{align}
 Q_G(z)={}&2\sum_{p\in\mathcal P_2}(z_{l(p)}-z_{m(p)})^2
       -2\sum_{p\in\mathcal P_-}z_{l(p)}^2
       +2\sum_{p\in\mathcal P_?}z_{l(p)}^2\nonumber\\
       &\quad+\epsilon\sum_{p\in\mathcal P_-}k_p^2z_{l(p)}^2.
                                      \label{eq:v6-Q}
\end{align}
This is a rational, purely geometric upper quadratic form,
not the assertion that all its worst cases are attained.

\begin{theorem}[Reduced-collar completed-fibre descent]
\label{thm:v6-curvature}
At the prescribed central values on \(G\), uniformly over
every other independent fine group, the actual completed action
under the V5 integer direction satisfies
\[
 S''(0)\le Q_G(z)=-\frac{981794718651}{131072}.
\]
Its normalized action Hessian is at most
\[
                 -\frac{327264906217}{1863703068672}
                                    \le-\frac16.
\]
The squared independent-group speed remains
\(2\sum_l z_l^2=2\cdot21328389\).
\end{theorem}

\begin{proof}
Only groups in \(U\) move. The completed pivots on \(E_U\)
stay \(I\) by source inactivity, and every other pivot is
constant along the curve because its source does not contain
a group in \(U\). Hence plaquettes outside \(\mathcal P_U\)
have zero derivative, even though their values may be arbitrary.

On two-selected plaquettes the coefficient is exact as described
above. On \(\mathcal P_-\) use \eqref{eq:v6-negative},
and on \(\mathcal P_?\) use the scalar upper bound 2.
This proves \(S''\le Q_G\) for every unprescribed field.
The checker evaluates the first three sums in \eqref{eq:v6-Q}
as \(-7493664\), and the retained root-angle loss as
\[
                         \frac{414809157}{131072}.
\]
Their sum is the displayed rational upper value. The same
integer witness has norm squared 21328389 as in V5; integer
comparison verifies \(Q_G(z)\le-\tfrac13\sum_lz_l^2\).
Divide by the actual free-group squared speed to obtain the
unit Hessian bound. All roots used in the estimate are the
actual completed pivots, not independently assigned caps.
\end{proof}

For comparison, V5's more restricted collar gave upper value
\(-1964257089399/262144\). The larger upper value here pays
explicitly for the released pivot inputs. The retained negative
margin is enough for the same simple unit bound \(-1/6\);
no equality of the two curvature estimates is claimed.

\subsection{Input closure is not an event restriction}

To bound derivatives uniformly, one must still include every
free input of the roots in the 537 incident plaquettes. The
full local input closure, together with the selected source
inputs, has 2065 free groups. It contains \(G\); the other
1296 local inputs are not prescribed. Their vertices lie in
\[
                   [-1,7]\times[-1,7]\times[0,3]\times[-1,7].
\]
The inherited translation by \((8,8,8,8)\) is therefore still
non-wrapping for every even \(L\ge32\).

This distinction matters. The 2065-group closure is a compact
parameter set for smooth derivative bounds. The event to be
repaired restricts only the 769 groups in \(G\). Smoothness
on the full closure does not silently put the other inputs
in a small-field chart. All fine groups outside that closure
are also unrestricted.

\begin{theorem}[Original-clock repair on the reduced visible carrier]
\label{thm:v6-repair}
Let \(C=I\), even \(L\ge32\), and \(G\subset\mathcal S\).
There is a sufficiently small closed inner product chart on
\(G\) about \(-I\) on \(U\) and \(I\) on \(G\setminus U\),
whose visible cylinder \(D_G\) has the original-clock bound
\[
 \kappa_{D_G}^{H_\nu}
 \ge\frac{c\beta}{2\{c\beta+8C_{\rm f}(\beta+1)\}}
 \ge\frac{c}{2(c+16C_{\rm f})}>0,
                  \qquad c=\frac1{96},\quad\beta\ge\beta_0.
\]
Constants are independent of ambient volume and of compatible
hidden index sets. Every group outside \(G\) is unrestricted,
including every group released from V5's collar. The chart
radius and \(\beta_0\) are existential, not numerically enclosed.
\end{theorem}

\begin{proof}
Normalize the fixed integer direction in orthonormal normal
coordinates on \(G\). Theorem~\ref{thm:v6-curvature} gives
action Hessian at most \(-1/6\) for every value of all
unprescribed local inputs. Smoothness of the completed action
on the fixed compact input closure yields a uniform outer
chart with frozen coordinate second derivative at most
\(-1/12\). The unprescribed inputs range over their full
compact groups in this assertion.

Shrink the chart so the sources on \(E_U\) remain strictly
outside the cutoff support. Their pivots and inverse Haar
factors are then exactly \(I\) and one, independent of the
rotated coordinates. All other inverse-pivot factors may
depend on unprescribed inputs, but have no rotated-coordinate
input. The true coordinate-Haar Hessian on \(G\) is bounded.
For sufficiently large \(\beta\), therefore,
\(\partial_a^2\Phi_\beta\le-\beta/24\).
With inverse metric lower bound \(1/2\), the measured
ground-state conjugation gives local diffusion floor
\(\beta/96\).

A smooth cutoff between the inner and outer charts depends
only on \(G\). The global diffusion resolvent and
low-spectral projection argument of Theorem~\ref{thm:repair}
give the joint original-clock killed floor. The exact
conditional-variance domination transfers it to the interface
marginal: lifting a function supported in \(D_G\) imposes
no restriction on any hidden input because \(G\subset\mathcal S\).
The finite input stencil proves the stated volume uniformity.
No update to a coarse variable or collective rotation kernel
has been introduced.
\end{proof}

This is a theorem with strictly weaker fine-variable hypotheses
than V5. A packing that failed V5's requirement
\(F_H\subset\mathcal S\) may now qualify if it meets
\(G\subset\mathcal S\). We do not silently replace any
declared packing: visibility remains an explicit hypothesis.

\subsection{Verification and remaining coverage}

The checker is
\begin{center}\small
\path{scripts/verify_ym_f14_v6_reduced_collar.py}.
\end{center}
Its SHA--256 is
\begin{center}\scriptsize\ttfamily
96DB1A61937F953DFD1BB4E6DCA167E099224B6F590402F7997643624DAC278B.
\end{center}
It constructs \(G\) directly from source and physical-link
incidence, checks that it contains no tree or root variables,
verifies its inclusion in the old collar and disjointness from
the original parent hidden set, counts all five plaquette
classes, and evaluates the exact rational curvature accounts.
It separately constructs the larger derivative input closure
and checks its non-wrapping vertex bounds. The accepted arithmetic
uses only integers and fractions and inherits V5's fixed witness.

The predecessor V5 has 90,379 bytes and SHA--256
\begin{center}\scriptsize\ttfamily
4B99CAB4B4A6B9D11108C9080A1655183A98CAB271442920A29B7F6C9DE2FB68.
\end{center}
Its complete body is preserved apart from the title version.
Active V6 replaces V5. Archives, companion notes, and updated
monographs are unchanged. The next companion and direction
checkpoint remains V10.

\begin{firewall}
The proof now conditions on 769 free groups rather than 9348,
with no smallness assumption on any other independent fine
group. It retains the same explicit coefficient \(1/96\)
in the original-clock floor. The 769 groups still follow a
particular local negative-sheet and source-companion pattern,
and coarse data remain fixed at identity. There is no complete
cover of arbitrary local patterns, no unrestricted all-exterior
patch gap, and no continuum physical mass-gap theorem.
\end{firewall}

The next useful relaxation should address the actual values
of the 636 nonrotated source and transverse groups in \(G\),
rather than merely deleting their indices. Those groups control
source inactivity and the scalar signs of useful plaquettes.
A broader criterion must retain those functions of the data,
or a different complete patch mechanism must control their
failure. The full Yang--Mills goal remains active.
\section{V7: a noncentral functional test on the reduced visible collar}

\subsection{Direction, provenance, and scope}

V6 releases the unnecessary root-source inputs, but still places
the remaining 769 groups near prescribed central values. The
next step is to replace those pointwise prescriptions by a
source-product test and a weighted plaquette matrix. These
functions use only the same 769 free groups. The other 1296
groups in the local derivative closure, and all groups beyond
it, remain arbitrary.

The idea of a noncentral visible matrix is not new: f13 V16
already gives one on the 10034-group collar, with equal
rotation amplitudes and the full free-input closure of its
controlled pivots. The new result here combines the reduced
carrier, unequal V5 amplitudes, and bounds on actual roots
whose inputs are not visible. Those roots are not inserted
as extra data into the test. We also certify a field with
154 noncentral groups, including a noncommuting transverse
pair, at a fixed distance from the old central configuration.
The original-clock repair applies to the entire closed set
passing the new test, not just to a neighbourhood of that
example.

All statements below retain coarse data \(C=I\), the V6
non-wrapping embedding, and the explicit visibility requirement
\(G\subset\mathcal S\). This is still a sufficient local repair
criterion, not a complete cover or a physical mass-gap theorem.

\subsection{The source test and a root-deleted surrogate}

Keep \(U,E_U,G,\mathcal P_2,\mathcal P_-,\mathcal P_?\) from
V6 and the positive integer amplitudes \(z_l\) from V5. Put
\[
                 N=2\sum_{l\in U}z_l^2=42656778.
\]
Allow every \(Y_l\), \(l\in G\), to be an arbitrary element
of \(\SU(2)\). There are \(14\cdot19=266\) selected-edge
nonroot source paths. Their free inputs lie in
\(\mathcal R_U\subset G\). Impose the closed test
\begin{equation}\label{eq:v7-source}
 \min_{e\in E_U,\ 1\le j\le14}
               \|I-P_j(Y)\|_{\rm HS}\ \ge\ \frac1{32}.
\end{equation}
The cutoff uses geodesic logarithm norm, with \(\rho=1/64\);
the chordal lower bound in \eqref{eq:v7-source} is a sufficient
condition, not a redefinition of the cutoff.

\begin{lemma}[Uniform inactivity from the visible source test]
\label{lem:v7-inactive}
On a neighbourhood of every field satisfying
\eqref{eq:v7-source}, every completed pivot on \(E_U\) is
exactly \(I\) and its inverse Haar factor is one. Varying
only the groups in \(U\) changes no other pivot or inverse
Haar factor. These conclusions hold for all unprescribed
free groups.
\end{lemma}

\begin{proof}
At candidate pivot \(I\), every selected-edge source has
geodesic distance at least \(1/32>\rho\) from identity.
Thus all its truncated logarithms vanish, with a strict
margin. Candidate \(I\) solves the exact pivot equation.
The inherited global uniqueness gives the actual pivot,
and the same argument holds on an open neighbourhood.
In that neighbourhood the coarse-to-pivot map is the
identity, so its inverse Haar Jacobian is one.
Each selected boundary group is an input only of its two
chronological sources on its own coarse edge. All other
pivots and their Jacobians have no selected input. This
last assertion is an incidence fact, independent of the
values of any source.
\end{proof}

For \(p\in\mathcal P_-\), define \(H_p^0(Y)\) by its
ordered physical word with every root factor replaced by
\(I\); retain all direct free and tree factors in their
original order and orientation. All its free factors
belong to \(G\). Write \(h_p=(H_p^0)_0\) for its scalar
quaternion part. This word is a computable surrogate.
It is not a proposed assignment to the actual roots.

Let \(k_p\in\{0,1,2\}\) be the number of root factors
whose owner is outside \(E_U\). The V6 counts remain
57, 30, and 27, respectively. With \(R=7/512\), set
\begin{equation}\label{eq:v7-cap}
 \mathcal U_k(h)=
 \min\left\{1,\ h+\frac{k^2R^2}{4}(-h)_+
                    +\frac{3kR}{4}\sqrt{1-h^2}\right\},
                \qquad -1\le h\le1.
\end{equation}
In particular, \(\mathcal U_0(h)=h\), and
\(\mathcal U_k(-1)=-1+k^2R^2/4\) for the values of \(k\)
used here.

\begin{lemma}[Noncentral scalar cap with hidden root inputs]
\label{lem:v7-cap}
The scalar part of the actual completed holonomy of
\(p\in\mathcal P_-\) is at most \(\mathcal U_{k_p}(h_p)\),
uniformly in every free group outside \(G\).
\end{lemma}

\begin{proof}
Each omitted actual root or inverse has geodesic distance
at most \(R\) from identity. Replace these factors one
at a time, using the triangle inequality and bi-invariance
at each replacement. For the actual holonomy \(H\) and
surrogate \(H^0\), this gives
\[
                  d_{\rm geo}(H^0,H)\le k_pR.
\]
No permutation or commuting of factors is involved.
Write \(H=H^0D\), \(H^0=(h,x)\), and \(D=(d_0,d)\).
The Hilbert--Schmidt metric and the unit-quaternion
parametrization give
\[
  d_0\ge1-\frac{k_p^2R^2}{4},\qquad d_0\le1,\qquad
  |d|\le\frac{k_pR}{\sqrt2},\qquad |x|=\sqrt{1-h^2}.
\]
The first inequality is the chordal estimate of V6.
For the third, the quaternion angle of \(D\) is
\(d_{\rm geo}(I,D)/\sqrt2\), and \(|\sin t|\le |t|\).
Thus
\[
 (H)_0=hd_0-x\cdot d
 \le h+\frac{k_p^2R^2}{4}(-h)_+
                       +\frac{k_pR}{\sqrt2}\sqrt{1-h^2}.
\]
Use \(1/\sqrt2\le3/4\) and the independent upper bound
\((H)_0\le1\). This proves \eqref{eq:v7-cap}.
All completions may be mutually dependent; the proof
is pointwise and assumes no independence.
\end{proof}

The term \((-h)_+\) matters: the sign of a current
noncentral surrogate need not be the sign at the V6
central pattern. The fixed name \(\mathcal P_-\) describes
its geometric class, not a sign hypothesis on current data.

\subsection{Unequal-amplitude noncommutative pair curvature}

Let \(n\in\mathbb R^3\), \(|n|=1\), and \(E_n=(0,n)\).
Use the actual free-group curve
\[
       Y_l(t)=e^{tz_lE_n}Y_l,\quad l\in U,
       \qquad Y_l(t)=Y_l,\quad l\notin U .
\]
Its squared independent-group speed is \(N\). Under
\eqref{eq:v7-source}, Lemma~\ref{lem:v7-inactive} makes
every other physical factor constant near \(t=0\).

There are no root factors in \(\mathcal P_2\). Cyclically
start a two-selected word at its positive selected link,
writing it as
\[
                 U_p A_p V_p^{-1}B_p=X_pB_p.
\]
Let \(a_p=z_{U_p}\), \(b_p=z_{V_p}\),
\(X_p=(x_{p,0},x_p)\), \(B_p=(b_{p,0},v_p)\), and
\(s_p=(X_pB_p)_0\). The notation \(b_p\) is the integer
amplitude; \(b_{p,0}\) is a quaternion scalar.
Define the real symmetric matrix
\begin{equation}\label{eq:v7-pair}
 T_p(Y)=2(a_p-b_p)^2s_p I_3+
 8a_pb_p\left\{\frac{x_pv_p^{\mathsf T}+v_px_p^{\mathsf T}}2
                                  -(x_p\cdot v_p)I_3\right\}.
\end{equation}

\begin{lemma}[Exact weighted pair identity]
\label{lem:v7-pair}
In the action normalization \(2-\operatorname{Tr}H\), the
second derivative of this plaquette is \(n^{\mathsf T}T_pn\).
This identity does not assume commuting group values.
\end{lemma}

\begin{proof}
Its moving trace is
\(\operatorname{Tr}(e^{ta_pE_n}X_pe^{-tb_pE_n}B_p)\).
Because \(E_n^2=-I\), its action second derivative is
\[
  2(a_p^2+b_p^2)(X_pB_p)_0+
                         4a_pb_p(E_nX_pE_nB_p)_0 .
\]
Quaternion multiplication gives
\[
 E_nX_pE_n=(-x_{p,0},\,x_p-2n(n\cdot x_p)),
\]
and hence
\[
 (E_nX_pE_nB_p)_0
       =-s_p-2x_p\cdot v_p+2(n\cdot x_p)(n\cdot v_p).
\]
Substitution gives \eqref{eq:v7-pair}. In particular, when
all factors are central with positive product, it reduces
to \(2(a_p-b_p)^2I_3\), the V5--V6 term.
\end{proof}

\subsection{A matrix determined solely by the 769 visible groups}

For a one-selected plaquette write \(l(p)\) for its
selected link. Define
\begin{align}
 \mathcal M_G(Y)={}&\sum_{p\in\mathcal P_2}T_p(Y)
       +2\sum_{p\in\mathcal P_-}
                 z_{l(p)}^2\mathcal U_{k_p}(h_p(Y))I_3
                 \nonumber\\
       &\qquad+2\sum_{p\in\mathcal P_?}z_{l(p)}^2I_3 .
                                               \label{eq:v7-matrix}
\end{align}
This is a continuous function of \(Y_G\) alone, even at
\(h_p=\pm1\). Its square-root majorant need not be
differentiable there; none of the following arguments
differentiates that majorant.

\begin{theorem}[All-unprescribed-field curvature test]
\label{thm:v7-curvature}
Under \eqref{eq:v7-source}, the actual completed action
satisfies
\[
              S''(0)\le n^{\mathsf T}\mathcal M_G(Y)n
\]
for every unit \(n\) and every choice of all free groups
outside \(G\). If
\begin{equation}\label{eq:v7-test}
                 \lambda_{\min}(\mathcal M_G(Y))\le-N/6,
\end{equation}
there is a unit Hilbert--Schmidt direction with actual
action Hessian at most \(-1/6\), uniformly over those groups.
\end{theorem}

\begin{proof}
The pair contributions are exact by
Lemma~\ref{lem:v7-pair}. A one-selected physical word
has second derivative \(2z_{l(p)}^2H_{p,0}\), since the
selected factor or its inverse has second derivative
equal to \(-z_{l(p)}^2\) times itself. For
\(\mathcal P_-\), use Lemma~\ref{lem:v7-cap}; for the
162 remaining one-selected plaquettes, use \(H_{p,0}\le1\).
No plaquette outside the 537 incident ones moves, because
all pivots are constant under this curve.
Summing proves the first assertion. A least-eigenvalue
unit vector gives the second after division by \(N\).
The free-group curve is a product-metric geodesic, so
its normalized second derivative is the claimed Hessian.
\end{proof}

At the central V6 values,
\[
          \mathcal M_G
              =-\frac{981794718651}{131072}I_3<-N I_3/6.
\]
Thus the new test recovers V6's centre exactly, including
its quadratic root-angle losses. It does not require the
unprescribed roots to equal identity, even at that centre.

\subsection{An explicit far-noncentral, noncommuting certificate}

Let \(\mathcal R^\circ=\mathcal R_U\setminus U\). The exact
incidence check gives \(|\mathcal R^\circ|=152\), and no
group in this set occurs as a direct factor in
\(\mathcal P_2\cup\mathcal P_-\). Each of the 266 source
paths has exactly one selected link and either one group
in \(\mathcal R^\circ\) or a tree companion. These are
geometric facts, not approximations to source values.

Write \(E_1=(0,1,0,0)\), \(E_2=(0,0,1,0)\) as unit
quaternions. Prescribe
\[
 Y_l=-I\quad(l\in U),\qquad
 Y_l=E_2\quad(l\in\mathcal R^\circ),
\]
and prescribe \(I\) on the remaining groups in \(G\),
except for the two transverse links
\[
       q_1=((0,0,1,5),0),\qquad q_2=((0,0,2,5),0),
       \qquad Y_{q_1}=E_1,\quad Y_{q_2}=E_2 .
\]
These integer coordinates are translated by \((8,8,8,8)\)
for the actual non-wrapping embedding. The two transverse
links are not source companions and do not change any
selected-edge source product. They occur in the same
two-selected plaquette. In particular \(E_1E_2=-E_2E_1\).
There are 154 noncentral free groups in this configuration.

\begin{proposition}[Exact noncentral instance]
\label{prop:v7-example}
For these visible data and arbitrary values of all other
independent fine groups,
\[
 \min_{e,j}\|I-P_j\|_{\rm HS}^2=4,\qquad
 \mathcal M_G=
 \begin{pmatrix}
 d&-33440&0\\
 -33440&d&0\\
 0&0&d
 \end{pmatrix},
 \qquad d=-\frac{981807563707}{131072}.
\]
The rational axis \(n=(1,0,0)\) already certifies
\[
 \frac{n^{\mathsf T}\mathcal M_G n}{N}
       =-\frac{981807563707}{5591109206016}<-\frac16 .
\]
\end{proposition}

\begin{proof}
The source products are \(-I\) or \(-E_2\), so their
squared chordal distances are 8 or 4. Both cases occur.
Evaluate \eqref{eq:v7-pair} in the fixed oriented words
and \eqref{eq:v7-cap} for the one-selected surrogates.
The exact checker obtains the displayed matrix. All
coordinates are integral unit quaternions, and the
square roots needed in this example are exactly zero
or one. Every matrix entry is therefore rational.
The off-diagonal entry records the noncommuting pair;
it is not discarded in the calculation. Finally,
integer comparison of \(6d\) and \(-N\) proves the
strict directional inequality.
\end{proof}

Each changed group has scalar part zero and chordal
distance 2 from \(I\). This is not an infinitesimal
continuity example. The strict source and curvature
margins also give open neighbourhoods of these data
inside the certified set.

\subsection{Repair for the whole closed functional event}

Let \(\mathcal A_G\subset\SU(2)^G\) be the closed set
defined by \eqref{eq:v7-source} and \eqref{eq:v7-test},
and let \(D_{\rm fun}\) be its visible cylinder. Only
\(G\) is restricted, not the full derivative input
closure.

\begin{theorem}[Functional-event original-clock floor]
\label{thm:v7-repair}
For \(C=I\), even \(L\ge32\), and \(G\subset\mathcal S\),
there is \(\beta_0<\infty\) such that
\[
 \kappa_{D_{\rm fun}}^{H_\nu}
 \ge\frac{c\beta}{2\{c\beta+8C_{\rm f}(\beta+1)\}}
 \ge\frac{c}{2(c+16C_{\rm f})}>0,\qquad
                    c=\frac1{96},\quad\beta\ge\beta_0 .
\]
The constants are independent of ambient volume and of
compatible hidden index sets. No probability of the
certified event appears in this killed-floor estimate.
This is not a lower bound on the unrestricted spectral gap.
\end{theorem}

\begin{proof}
The parameter set \(\mathcal A_G\) is compact. At each
of its points, select a least-eigenvalue axis and freeze
the resulting normalized direction in orthonormal
product normal coordinates centred at that point.
Theorem~\ref{thm:v7-curvature} gives action Hessian at
most \(-1/6\) for all unprescribed inputs.

The actual completed-action derivatives, unlike their
scalar majorant, are smooth on the finite compact
2065-group input closure. Uniform continuity gives a
neighbourhood in \(Y_G\) on which the frozen coordinate
second derivative is at most \(-1/12\), still for all
other inputs. Shrink it to preserve source inactivity
and inverse metric lower bound \(1/2\). All inverse-pivot
Jacobians are independent of the selected coordinates
there. Coordinate-Haar second derivatives are bounded,
so for sufficiently large \(\beta\) the actual measured
potential satisfies \(\partial_a^2\Phi_\beta\le-\beta/24\).
The inherited conjugation estimate gives the local
diffusion floor \(\lambda=\beta/96\).

Compactness supplies finitely many such chart cylinders
with smooth cutoffs \(\chi_j(Y_G)\), supported in those
charts, satisfying
\(\sum_j\chi_j^2\le1\) everywhere and
\(\sum_j\chi_j^2=1\) on \(\mathcal A_G\).
One explicit construction is to start with finitely
many subordinate nonnegative functions \(\psi_j\) such
that \(T=\sum_j\psi_j^2\ge1\) on \(\mathcal A_G\).
Choose a smooth \(g\) with \(g=1\) for \(T\le1/4\) and
\(g=0\) for \(T\ge1/2\), and set
\(\chi_j=\psi_j/(T+g(T)^2)^{1/2}\).
Its denominator is positive everywhere. The finite
carrier gives
\(M=\sup\sum_j|\nabla\chi_j|^2<\infty\).
No differentiable choice of eigenvectors is needed.

Apply the local floor to \(\chi_j u\) and use the
product rule:
\[
 \lambda\|1_{D_{\rm fun}}u\|^2
 \le \sum_j\E_L(\chi_j u,\chi_j u)
 \le 2\E_L(u,u)+2M\|u\|^2 .
\]
These inequalities extend from a smooth core to the
form domain. With \(\beta_0\) enlarged so that
\(\lambda\ge8M\), the inherited low-spectral projection
argument and the global comparison
\(H_\mu\succeq L(L+C_{\rm f}(\beta+1))^{-1}\) give the
first bound on the joint fibre. This use of the
resolvent assumes no unrestricted diffusion gap.
The conditional-variance domination then transfers it
to the exact interface marginal: a lifted function
supported in \(D_{\rm fun}\) restricts only retained
groups, since \(G\subset\mathcal S\).
Finally \(\beta+1\le2\beta\) for \(\beta\ge1\) gives
the displayed constant floor.
\end{proof}

This compact covering is on one fixed finite carrier.
It is not a union over a number of spatial patches
growing with the ambient volume, and it makes no
unproved assertion about the complement of the test.

\subsection{Verification and the next missing implication}

The exact checker is
\begin{center}\small
\path{scripts/verify_ym_f14_v7_visible_functional.py}.
\end{center}
Its SHA--256 is
\begin{center}\scriptsize\ttfamily
CC184D39CDF992470B6F507E5D93E162550F3B44BB9426DB46B9581C4ACFA3CC.
\end{center}
It reruns the V6 audit, reconstructs the carrier and
plaquette classes, verifies the 266 source paths and
the 152 source-only companions, and evaluates both
displayed matrices with exact fractions. General
square roots in a numerical certificate are bounded
upward by an integer-square-root calculation on the
grid \(2^{-48}\). Since this only increases the scalar
diagonal, a passing rounded certificate is sufficient.
The checker also performs 384 exact weighted second-jet
checks on noncommuting rational quaternion samples.
Those samples check the implementation; the algebraic
identity in Lemma~\ref{lem:v7-pair}, not finite sampling,
proves the formula for all group values.

V6 was checked to compile with resolved references
before this append. Its predecessor bytes numbered
104,946, with SHA--256
\begin{center}\scriptsize\ttfamily
9772D8B0608155FDD43BFC0F09327F5E25A2393C06FDFC0F30A558962E6E3FE6.
\end{center}
V7 retains that body apart from the title version and
two purely typographic escape repairs in V5's
\(d_{\rm geo}(I,B_e)\) notation. The earlier source
accidentally used a carriage-return character in place
of the backslash in those two font commands; the
metric and the mathematical claims are unchanged.
V7 is the single active file. Archives, monographs, and
companion notes are preserved. The next companion and
direction checkpoint remains V10.

\begin{firewall}
The hypotheses now concern 266 source products and a
visible matrix, rather than a prescribed chart for
each of 769 groups. Actual hidden-input roots are
controlled by a scalar cap, not read as extra visible
variables. A far-noncentral noncommuting instance
passes with the same explicit coefficient \(1/96\).
It has not been proved that every relevant bad field
passes this test or another compatible repair test.
No all-exterior patch gap, cofinal continuum
construction, or physical Yang--Mills mass gap follows.
\end{firewall}

The next missing implication is coverage, not the
localization of fields already certified. In particular,
an argument must either control failure of the source
margin and matrix test together, or replace the
strict-descent mechanism on such configurations.
Repeatedly shrinking neighbourhoods of successful
examples would not supply that implication.
\section{V8: source-failure masks and a certified one-edge allowance}

\subsection{Attacking failure of the source margin}

V7 replaced central chart prescriptions by a noncentral
functional test, but still requires all fourteen sources
on each of the 19 selected coarse edges to be inactive.
This version permits source failure by holding the seven
boundary groups on a failing edge fixed and changing the
rotation amplitudes on the remaining links.

Omitted-edge directions already occur in V2--V5. The new
result is a source-driven, noncentral test on V6's
769-group carrier, with actual active pivots allowed on
omitted edges, and an exact guarantee for every one-edge
failure. The source companions may have arbitrary values
subject only to source tests on edges that actually move.

The proof pays for every root in a useful one-selected
plaquette, making its majorant independent of which
source edges are omitted. We also prove a limitation:
one two-edge omission makes this particular upper
quadratic form positive definite. This limits the
certificate, not the actual action Hessian.

\subsection{Masks are stationary coordinates, not deleted factors}

Keep \(U,E_U,G\) and the plaquette classes from V6--V7.
For a nonempty \(A\subset E_U\), put
\[
 U_A=\{l\in U:\operatorname{owner}(l)\in A\}.
\]
Only these groups will move. Every other free group,
including those on \(E_U\setminus A\), retains its
actual value. All physical words and pivots are retained.
Choose real amplitudes and a unit axis with
\[
 w_l=0\ (l\notin U_A),\qquad \sum_{l\in U}w_l^2=1,
                         \qquad |n|=1,
\]
and use the actual curve
\begin{equation}\label{eq:v8-curve}
                 Y_l(t)=e^{tw_lE_n}Y_l,\quad l\in U.
\end{equation}
All other independent groups are fixed. Its squared
Hilbert--Schmidt speed is 2. Require the source margin
only on \(A\):
\begin{equation}\label{eq:v8-sources}
 \|I-P_j(Y)\|_{\rm HS}\ge\frac1{32}
                 \quad(e\in A,\ 1\le j\le14).
\end{equation}
No source condition is imposed on \(E_U\setminus A\).

The proof of Lemma~\ref{lem:v7-inactive} gives pivots
equal to \(I\) and inverse Haar factors equal to one
on \(A\), throughout a neighbourhood. Every other
pivot is constant along \eqref{eq:v8-curve}: its
source contains no moving independent group. Such
a pivot may have active sources, a nonidentity value,
and a nontrivial Jacobian. Constancy under this curve
does not mean inactivity.

For \(p\in\mathcal P_-\), let \(\bar k_p\) count
\emph{all} root factors in its word. Exact incidence gives
\[
 \#\{p:\bar k_p=1\}=57,\qquad
 \#\{p:\bar k_p=2\}=57.
\]
Keep V7's surrogate scalar \(h_p\) and cap
\(\mathcal U_k(h)\). Lemma~\ref{lem:v7-cap} remains
valid with \(\bar k_p\), since every actual root is
within \(R=7/512\) of \(I\) at \(C=I\). Roots that
are exactly \(I\) are harmlessly charged the same
cap. No extra inverse-pivot inputs are made visible.

\subsection{The amplitude matrix for arbitrary visible values}

For a pair plaquette use V7's ordered notation
\[
 U_pA_pV_p^{-1}B_p=X_pB_p,\quad
 X_p=(x_{p,0},x_p),\quad B_p=(b_{p,0},v_p),\quad
 s_p=(X_pB_p)_0,
\]
and let \(f_p(n)=(n\cdot x_p)(n\cdot v_p)-x_p\cdot v_p\).
Define a symmetric \(133\times133\) matrix \(K_n(Y)\),
initially zero and indexed by \(U\), by adding:
\begin{enumerate}
\item For each \(p\in\mathcal P_2\), with selected links
\(i,j\), add \(2s_p\) to \(K_{ii},K_{jj}\) and
\(-2s_p+4f_p(n)\) to \(K_{ij},K_{ji}\).
\item For each \(p\in\mathcal P_-\), add
\(2\mathcal U_{\bar k_p}(h_p)\) to \(K_{l(p),l(p)}\).
\item For each \(p\in\mathcal P_?\), add 2 to
\(K_{l(p),l(p)}\).
\end{enumerate}
This continuous matrix uses only \((Y_G,n)\). Its
scalar majorant need not be differentiable at central
holonomies; the argument does not differentiate it.

\begin{lemma}[Masked noncentral curvature bound]
\label{lem:v8-matrix}
Under \eqref{eq:v8-sources}, for all unprescribed fine
fields the completed action satisfies
\[
                         S''(0)\le w^{\mathsf T}K_n(Y)w.
\]
In particular,
\begin{equation}\label{eq:v8-negative}
                       w^{\mathsf T}K_n(Y)w\le-\frac18
\end{equation}
gives a unit Hilbert--Schmidt action Hessian at most
\(-1/16\).
\end{lemma}

\begin{proof}
All pivots are constant under the curve by the source
and incidence argument above. V7's exact weighted
pair formula is
\[
              2s_p(w_i-w_j)^2+8w_iw_jf_p(n),
\]
which gives the prescribed pair entries. The formula
remains valid when either amplitude is zero: the
stationary selected link is still in the word.
Each one-selected contribution is its squared
amplitude times twice the actual scalar holonomy.
Bound that scalar by \(\mathcal U_{\bar k_p}(h_p)\)
or by 1 in the respective classes, and sum.
Division by the actual squared speed 2 proves the
normalized claim. No active-pivot derivative is
discarded by assumption: it vanishes because none
of that pivot's independent inputs moves.
\end{proof}

Equivalently, for fixed \(A,n,Y\), test the principal
submatrix on \(U_A\) for a least eigenvalue at most
\(-1/8\). A principal submatrix sets other amplitudes
to zero; it does not delete physical group factors.

\subsection{One compact functional event, including active sources}

Let \(D_{\rm mask}\) be the visible cylinder of the
fields \(Y_G\) for which there exist nonempty \(A\),
unit \(n\), and normalized amplitudes supported on
\(U_A\), satisfying \eqref{eq:v8-sources} and
\eqref{eq:v8-negative}.

\begin{theorem}[Original-clock masked repair]
\label{thm:v8-repair}
For \(C=I\), even \(L\ge32\), and \(G\subset\mathcal S\),
there is a finite \(\beta_0\), uniform in ambient volume
and compatible hidden index sets, such that
\[
 \kappa_{D_{\rm mask}}^{H_\nu}
 \ge\frac{c\beta}{2\{c\beta+8C_{\rm f}(\beta+1)\}}
 \ge\frac{c}{2(c+16C_{\rm f})}>0,\qquad
                    c=\frac1{256},\quad\beta\ge\beta_0.
\]
All free groups outside \(G\) are unrestricted.
No source condition on a stationary edge is implicit.
\end{theorem}

\begin{proof}
For fixed \(A\), the tests define a closed subset of
the compact product of \(\SU(2)^G\), the axis sphere,
and the amplitude sphere on \(U_A\). Its projection
onto \(\SU(2)^G\) is compact. The union over the
finitely many masks is compact.

At each passing point choose one witnessing mask,
axis, and amplitude vector and freeze its normalized
direction in product normal coordinates. Lemma
\ref{lem:v8-matrix} gives action Hessian at most
\(-1/16\), uniformly over unprescribed inputs.
The fixed 2065-group derivative closure and the
source margin between \(1/32\) and \(\rho\) give a
coordinate neighbourhood with second derivative
at most \(-1/32\).

All inverse-pivot factors are independent of the
coordinates with nonzero amplitudes there. This
includes active pivots, whose input sets are
disjoint from those moving coordinates, even though
they may depend on stationary coordinates in \(G\).
The coordinate-Haar allowance gives measured
second derivative at most \(-\beta/64\).
With inverse metric at least \(1/2\), the local
conjugation floor is \(\lambda=\beta/256\).

The compact chart covering and cutoffs in the proof
of Theorem~\ref{thm:v7-repair} yield
\[
 \|1_{D_{\rm mask}}u\|^2
     \le\frac2\lambda\E_L(u,u)+\frac{2M}{\lambda}\|u\|^2
\]
with finite, volume-independent \(M\). No selected
mask, eigenvector, or majorant is differentiated.
Enlarge \(\beta_0\) so that \(\lambda\ge8M\).
The joint resolvent comparison and conditional-variance
transfer with \(G\subset\mathcal S\) give the
claimed original-clock bounds as in V7.
\end{proof}

This criterion complements V7's. No containment of
one whole event in the other is asserted. Their
finite union is handled by the same compact covering
with the smaller coefficient \(1/256\). This is
still not a complete cover.

\subsection{Arbitrary source companions with one failing edge}

Let \(\mathcal R^\circ=\mathcal R_U\setminus U\) and
\[
                    W=G\setminus(U\cup\mathcal R^\circ).
\]
There are 152 source companions and 484 groups in
\(W\). The companions split into 19 disjoint
eight-group sets, one per selected coarse edge.
Of each edge's 14 sources, eight have these free
companions and six have tree companions. No group
in \(\mathcal R^\circ\) is a direct factor of a
controlled plaquette.

Consider
\begin{equation}\label{eq:v8-central-useful}
 Y_l=-I\quad(l\in U),\qquad Y_l=I\quad(l\in W),
\end{equation}
with arbitrary values on \(\mathcal R^\circ\).
Then \(K_n(Y)\) is independent of \(n\) and all
152 companion values. Denote it by \(K_*\):
\begin{equation}\label{eq:v8-central-form}
\begin{split}
 w^{\mathsf T}K_*w={}&
 2\sum_{\mathcal P_2}(w_i-w_j)^2+
 \sum_{\mathcal P_-}(-2+\epsilon\bar k_p^2)w_{l(p)}^2\\
 &\quad+2\sum_{\mathcal P_?}w_{l(p)}^2,\qquad
                         \epsilon=\frac{49}{524288}.
\end{split}
\end{equation}
This is an upper form. Simultaneous attainment of
its worst cases by unprescribed fields is not asserted.

\begin{proposition}[All nineteen one-edge certificates]
\label{prop:v8-one}
For every \(e\in E_U\) there is a nonzero integer
vector \(v^{(e)}\) on the 126 links not owned by \(e\),
extended by zero to the other seven links, such that
\[
 \frac{(v^{(e)})^{\mathsf T}K_*v^{(e)}}
                          {\sum_l(v_l^{(e)})^2}<-\frac18.
\]
Every field satisfying \eqref{eq:v8-central-useful}
for which at least 18 selected edges pass the source
margin therefore belongs to \(D_{\rm mask}\).
The companion values on the other edge are arbitrary,
including noncommuting, inner-active, and transition
source configurations.
\end{proposition}

\begin{proof}
Here is a finite integer specification of every
witness. Put \(D=524288\), and let \(M_e\) be the
integer principal matrix \(D K_*\) after excluding
the seven amplitude indices owned by \(e\).
Order remaining fine links lexicographically.
For \(s=10000\), set \(v^{(0)}=s{\bf1}\) and perform
exactly 128 steps:
\[
 u^{(r)}=(12D I-M_e)v^{(r)},\quad T_r=\max_i u_i^{(r)},
 \quad v_i^{(r+1)}
     =\left\lfloor\frac{2s u_i^{(r)}+T_r}{2T_r}\right\rfloor .
\]
The checker verifies positivity of the intermediate
vectors and then checks, separately for all 19 edges,
\[
             8(v^{(128)})^{\mathsf T}M_ev^{(128)}
                                     <-D\|v^{(128)}\|^2.
\]
No convergence theorem for this vector-generating
iteration is assumed. The final integer inequalities
are the certificates.

The results follow. An edge has coordinates
\(((a,b,0,c),2)\). The first column records the
sorted multiset of \((a,b,c)\); counts aggregate
separately performed checks, not an assumed symmetry
reduction. Write \(Q_e=v^{\mathsf T}M_ev\).
\begin{center}\small
\begin{tabular}{crrr}
\toprule
Multiset&Number&\(\|v\|^2\)&\(Q_e\)\\ \midrule
\((0,0,2)\)&3&1972999088&\(-348563157426066\)\\
\((0,1,2)\)&6&1774488784&\(-283973784231090\)\\
\((0,2,2)\)&3&1831976677&\(-281991385160507\)\\
\((1,1,2)\)&3&1600179874&\(-165354047255042\)\\
\((1,2,2)\)&3&2087903113&\(-181197858625159\)\\
\((2,2,2)\)&1&2962123776&\(-290932425831927\)\\
\bottomrule
\end{tabular}
\end{center}
The minimum positive integer slack
\(-D\|v\|^2-8Q_e\) is \(354920321692728\).
Divide by \(D\|v\|^2\) for the first assertion.
If one edge fails the source test, freeze it; if
none fails, freeze any edge. The normalized witness,
with any unit axis, then passes the masked test.
\end{proof}

This allows an arbitrary source regime on any one
edge, not just a specified cutoff crossing. The
strict curvature margins also tolerate uniformly
small perturbations of the direct factors in
\eqref{eq:v8-central-useful}, with the source tests
evaluated on the perturbed fields themselves.

\subsection{An actual active pivot excluded by V7}

Fix \(e\in E_U\), choose \(A_0\in\mathfrak{su}(2)\)
of Hilbert--Schmidt norm one, and set \(\tau=1/1024\).
Keep \eqref{eq:v8-central-useful}. Prescribe
\[
                           Y_l=-\exp(\tau A_0)
\]
on the eight source companions of \(e\), and \(I\)
on the other 144. On \(e\), the eight source
products with free companions are \(\exp(\tau A_0)\);
the remaining six are \(-I\). Every other selected
edge has all source products equal to \(-I\).

\begin{lemma}[Nonidentity completion on the stationary edge]
\label{lem:v8-active}
The actual pivot on \(e\) is \(B_e=\exp(-\tau A_0)\ne I\).
All eight positive-source truncated logarithms are
nonzero inner-source terms. This field fails V7's
source test but belongs to \(D_{\rm mask}\).
\end{lemma}

\begin{proof}
At this candidate pivot, each positive relative
source is \(\exp(2\tau A_0)\), whose logarithm has
norm \(2\tau=1/512<\rho/2=1/128\). Its truncated
logarithm is exactly \(2\tau A_0\).
The negative relative sources are
\(-\exp(\tau A_0)\), with chordal distance from \(I\)
at least \(2\sqrt2-\tau>\rho\), outside the cutoff.
The exact pivot equation becomes
\[
 B_e\exp\left(\frac1{16}\,8\,(2\tau A_0)\right)
                     =\exp(-\tau A_0)\exp(\tau A_0)=I.
\]
Global uniqueness identifies this with the actual
completion. Since \(0<\tau\) is below the injectivity
radius and \(A_0\ne0\), this pivot is not identity.

Each positive source has chordal distance from \(I\)
at most its geodesic distance \(\tau<1/32\), so
V7's source condition fails. The other selected
edges have sources \(-I\), and
Proposition~\ref{prop:v8-one} applies.
\end{proof}

The active pivot and its companions do not move in
the repair curve. Its inverse-pivot Jacobian remains
in the density, but is independent of the selected
repair coordinates. This does not replace active-source
derivatives by an inactive formula.

\subsection{An exact limit at two excluded edges}

The one-edge guarantee cannot extend to all two-edge
omissions using only \eqref{eq:v8-central-form}.
Put
\[
 e_1=((1,1,0,2),2),\qquad e_2=((2,2,0,1),2),
\]
and let \(K_{**}\) be the 119-dimensional principal
matrix after excluding their 14 amplitude indices.

\begin{proposition}[Positive two-edge upper form]
\label{prop:v8-two}
The explicit upper matrix satisfies
\[
                             K_{**}\succeq\frac1{10}I.
\]
No nonzero amplitude vector on those remaining links
can make this particular majorant negative.
\end{proposition}

\begin{proof}
Let \(M_{**}=D K_{**}\). Run the preceding integer
recipe with \(s=10^9\) for 2048 steps. Its resulting
vector \(v\) satisfies the exact checks
\[
 \min_i v_i=73481758>0,\qquad
 \min_i\{10(M_{**}v)_i-Dv_i\}=12628305903616>0.
\]
All off-diagonal entries of \(K_{**}\) are nonpositive.
For every real \(x\), expand the squares in
\[
 x^{\mathsf T}K_{**}x
 =\sum_{i<j}(-K_{**,ij})v_iv_j
       \left(\frac{x_i}{v_i}-\frac{x_j}{v_j}\right)^2
    +\sum_i\frac{(K_{**}v)_i}{v_i}x_i^2 .
\]
This identity, together with
\((K_{**}v)_i/v_i\ge1/10\), proves the result.
\end{proof}

At the central direct factors, the corresponding
independent-axis upper form is \(K_{**}\otimes I_3\).
Its pair terms are \(2|a_i-a_j|^2\), and its
one-selected terms are scalar multiples of
\(|a_i|^2\). Thus different axes on different links
cannot rescue this same two-edge majorant.

The matrix is an \emph{upper} bound on completed
curvature. Its positivity does not prove positivity
of the actual Hessian, absence of descent involving
other groups, or a heat-bath metastable obstruction.
In particular, the 162 unrestricted one-selected
plaquettes were separately charged their worst scalar
value; simultaneous compatibility was not used.

\subsection{Verification and direction after the obstruction}

The checker is
\begin{center}\small
\path{scripts/verify_ym_f14_v8_source_masks.py}.
\end{center}
Its SHA--256 is
\begin{center}\scriptsize\ttfamily
CA8FB6B466AD4E98CDCCA0F27449432E59B8B764A51717DE22D89C543394A9F4.
\end{center}
It constructs the integer incidence matrix, checks
the 19 disjoint eight-companion source sets, and
verifies every negative witness and the positive
two-edge supersolution using integer arithmetic.
It also checks the general noncentral coefficient
assembly against the weighted pair formula on a
rational noncommuting field and reruns V7's 384
weighted quaternion-jet checks. These algebra
tests check the implementation; the analytic
identities establish the general statements.

The predecessor V7 has 123,256 bytes and SHA--256
\begin{center}\scriptsize\ttfamily
57EE4920111C68FB1DAA3CCD62488BEB43BB33BBE35C266746DC01DFCB5A685B.
\end{center}
Its body is preserved except for the title version.
Active V8 replaces V7; archives, companions, and
monographs are preserved. The next companion
and direction review is V10.

\begin{firewall}
The programme now has an exact all-one-edge
source-failure allowance on the reduced carrier,
including a genuine nonidentity active pivot.
The repair coefficient for this event is \(1/256\)
in the original-clock comparison. Arbitrary source
failure on several edges remains uncovered.
A specified two-edge upper matrix is positive;
further optimization of that same matrix cannot
give universal coverage. There is no unrestricted
all-exterior patch gap or continuum physical
mass-gap proof.
\end{firewall}

The next useful upstream move is to improve the
unrestricted-plaquette information or change the
directions and carrier, not search again for
negative eigenvectors of the positive two-edge
majorant. Alternatively, a different mechanism
must control configurations with many active-source
edges. The full Yang--Mills goal remains active.
\section{V9: favourable active-pivot response for every companion sign pattern}

\subsection{Upstream diagnosis and the changed direction}

V8 proves a one-edge source-failure allowance but finds
a positive upper matrix after a specified two-edge
omission. We first tested whether the 162 unrestricted
plaquette bounds could be improved through compatibility.
In the present geometry those bounds can all be attained
simultaneously, even with the exact coarse constraint.
Thus compatibility alone does not improve them.

The useful change is instead to move active pivots
according to the exact constraint. Each selected pivot
has three controlled forward plaquettes whose negative
contributions pay for the previously unseen plaquettes
created by its motion. This gives a sign-sensitive
comparison, not a norm debit for every pivot term.
It covers every independent sign choice of the 152
source companions, even when all 19 selected coarse
edges have active sources.

Active-source pivot formulas and bounded Jacobian
allowances already appear in f13 V17. There the
selected boundary groups themselves were sign-flipped,
and pivot terms were charged by absolute value.
Here the selected groups keep their negative signs;
the source companions change sign. This preserves
the negative mixed plaquettes and produces the
favourable comparison below on the smaller carrier.
Neither the pivot formula nor the abstract
diffusion-to-heat-bath transfer is claimed as new.

\subsection{The unrestricted-plaquette estimate is genuinely sharp}

Keep V8's partition
\[
             G=U\sqcup W\sqcup\mathcal R^\circ,
             \qquad (|U|,|W|,|\mathcal R^\circ|)
                           =(133,484,152).
\]
Prescribe
\begin{equation}\label{eq:v9-centres}
 Y_l=-I\quad(l\in U),\qquad Y_l=I\quad(l\in W),\qquad
 Y_c=\sigma_cI\quad(c\in\mathcal R^\circ),\quad
                         \sigma_c\in\{-1,1\}.
\end{equation}
All groups outside \(G\) will be arbitrary in the main
theorems. The following lemma uses a specific allowed
exterior only to demonstrate sharpness of a bound.

\begin{lemma}[Simultaneous attainment of the old unrestricted terms]
\label{lem:v9-sharp}
For every sign pattern in \eqref{eq:v9-centres}, there
is a fine field on the exact coarse-identity fibre for
which all 162 plaquettes in \(\mathcal P_?\) have
holonomy \(I\). Each then contributes exactly
\(2z_{l(p)}^2\) under the V5 selected-link direction.
\end{lemma}

\begin{proof}
Assign every remaining independent free group its
original central sign from the non-wrapping sheet
pattern. All nonroot source products are central,
hence equal to \(I\) or \(-I\). At pivot \(I\), both
give zero truncated logarithm: the first has zero
logarithm, and the second lies outside the cutoff.
Thus all completed pivots equal \(I\) by uniqueness,
and all coarse groups equal \(I\).

Each plaquette in \(\mathcal P_?\) has no root factor
and no source-companion factor. The product of its
original central signs is \(+1\), as checked directly
in the oriented words. The changes of \(\sigma_c\)
therefore do not change its holonomy. Its only moving
factor is one selected link; differentiating twice
gives \(2z_{l(p)}^2\), as claimed. The same translated
central pattern is defined on the periodic lattice;
the finite stencil lies away from its wrapping seam.
\end{proof}

This rules out a uniformly smaller bound for those
162 terms alone under the same hypotheses. It does
not rule out improvement from the other physical
terms or from a different direction.

\subsection{Exact motion with source-companion signs}

Use the strictly positive V5 integer amplitudes
\(z_l\), and for any unit \(n\in\mathbb R^3\) set
\[
 Y_l(t)=-e^{tz_lE_n},\quad l\in U,
 \qquad Y_c(t)=\sigma_cI,\quad c\in\mathcal R^\circ.
\]
Keep every other independent fine group fixed.
For \(e\in E_U\), let \(k_e\) be the number of its
eight free source companions with sign \(-1\).
Let \(m_{e,l}\in\{0,1,2\}\) count how many of the
two source paths containing \(l\) have a negative
companion. Then
\[
 0\le k_e=\sum_{\operatorname{owner}(l)=e}m_{e,l}\le8,
 \qquad
 a_e=-\frac{\sum_{\operatorname{owner}(l)=e}m_{e,l}z_l}
                         {16-k_e}\le0.
\]
The six paths with tree companions are always negative
sources and are not included in \(k_e\).

\begin{lemma}[The actual pivots, including their acceleration]
\label{lem:v9-pivot}
For sufficiently small \(|t|\), the completed pivot
on each \(e\in E_U\) is exactly
\[
                         B_e(t)=e^{ta_eE_n}.
\]
Thus \(B'_e(0)=a_eE_n\) and \(B''_e(0)=-a_e^2I\).
All other pivots are independent of this motion.
The assertions hold for all unprescribed fine fields.
\end{lemma}

\begin{proof}
All selected-edge sources lie in \(G\). At \(t=0\)
they are central, so their actual pivot is \(I\).
A negative companion makes its source
\(e^{tz_lE_n}\); a positive or tree companion makes
its source \(-e^{tz_lE_n}\).
At candidate \(B_e=e^{ta_eE_n}\), the positive
relative sources are \(e^{t(z_l-a_e)E_n}\), in the
inner cutoff region for small \(|t|\), while the
negative relative sources stay outside the support.
The exact coarse-identity equation becomes
\[
 \exp\left\{\frac{t}{16}
       \left((16-k_e)a_e+\sum_lm_{e,l}z_l\right)E_n\right\}=I.
\]
The definition of \(a_e\) makes its exponent zero.
Global pivot uniqueness proves the formula.
The acceleration follows by differentiation.
No selected boundary group is a source input of any
other coarse edge, so other pivots remain constant.
\end{proof}

The selected free groups, not the dependent pivots,
determine the metric speed:
\[
                N=2\sum_{l\in U}z_l^2=42656778.
\]
At a centre an inner-source logarithm can vanish
while its response is nonzero. It is not an inactive
source merely because that logarithm is zero.

\subsection{Complete geometry of the additional pivot motion}

Let \(\mathcal B=\{B_e:e\in E_U\}\) denote the 19
physical root links. The union of plaquettes incident
to \(U\cup\mathcal B\) has 564 members, 27 more than
the old 537. All moving physical links have direction
2. Each plaquette contains at most two of them,
and no plaquette contains two selected pivots.

The 114 old controlled negative plaquettes split as
follows:
\begin{itemize}
\item 57 forward mixed plaquettes, with one selected
free link and one selected pivot, and no other root;
\item 30 backward mixed plaquettes, with one selected
free link, one selected pivot, and one other root;
\item 27 one-selected-free plaquettes with two
unprescribed roots and no selected pivot.
\end{itemize}
For each selected pivot, the three forward mixed
plaquettes use its fine neighbours in directions
\(0,1,3\) with displacement \(+1\). Their direct
nonmoving factors are prescribed central groups.
If \(b_e\) is its number of backward mixed
plaquettes and \(u_e\) its number of new uncontrolled
plaquettes, then \(b_e+u_e=3\). Exact row counts are
\begin{center}\small
\begin{tabular}{cccr}
\toprule
Forward mixed&Backward mixed \(b_e\)&New uncontrolled \(u_e\)&Pivots\\
\midrule
3&0&3&3\\
3&1&2&6\\
3&2&1&6\\
3&3&0&4\\
\bottomrule
\end{tabular}
\end{center}
In particular, \(u_e\le3\). The 261 free--free pair
plaquettes and 162 old uncontrolled plaquettes contain
no selected pivot.

\begin{lemma}[One fixed root does not require a norm cross-term debit]
\label{lem:v9-mixed}
In a mixed plaquette with selected-free speed \(z_l\)
and pivot speed \(a_e\), its actual action second
derivative is
\[
       -2(z_l-a_e)^2
       \quad\hbox{for a forward plaquette},
\]
and
\[
       -2r_{p,0}(z_l-a_e)^2
       \quad\hbox{for a backward plaquette},
\]
where \(r_{p,0}\) is the scalar part of its other
root or inverse. The latter is therefore at most
\[
              (-2+\epsilon)(z_l-a_e)^2,\qquad
                               \epsilon=\frac{49}{524288}.
\]
No common axis is assumed for the unprescribed root.
\end{lemma}

\begin{proof}
At the centre, the selected-free sign is negative,
the selected pivot equals \(I\), and the remaining
direct nonroot factors are \(I\). In a forward
plaquette, the trace is the negative trace of
\(e^{\pm t(z_l-a_e)E_n}\).
For a backward plaquette, cyclically start the trace
immediately after its single other root factor.
All remaining factors are central or exponentials
of \(E_n\); they commute with each other, without
commuting them past the unprescribed root. The
trace is therefore
\(-\operatorname{Tr}(e^{\pm t(z_l-a_e)E_n}R_p)\),
with \(R_p\) that root or its inverse.
Differentiation and \(E_n^2=-I\) give the formulas.
The exact pivot cap gives
\(r_{p,0}\ge1-R^2/4\), so
\(-2r_{p,0}\le-2+R^2/2=-2+\epsilon\).
\end{proof}

The remaining 27 controlled one-free terms still
have upper coefficient \(-2+4\epsilon\), by V6's
two-root product cap. Each new pivot-only plaquette
has second derivative at most \(2a_e^2\), with no
restriction on its other physical factors.

\subsection{A sign-sensitive comparison with the inactive baseline}

Let \(Q_G(z)\) be V6's upper quadratic form, with
all selected pivots stationary:
\[
                  Q_G(z)=-\frac{981794718651}{131072}.
\]
For a selected pivot write \(F_e\) and \(J_e\) for
the selected-free links in its forward and backward
mixed plaquettes, respectively. Thus \(|F_e|=3\)
and \(|J_e|=b_e\). Define
\begin{equation}\label{eq:v9-delta}
 \Delta_e=
       4a_e\sum_{l\in F_e}z_l+(2u_e-6)a_e^2
          +(-2+\epsilon)\sum_{l\in J_e}
                               (a_e^2-2a_ez_l).
\end{equation}

\begin{theorem}[All companion sign patterns have the same repair margin]
\label{thm:v9-curvature}
At every centre \eqref{eq:v9-centres}, for every
choice of all independent fine groups outside \(G\),
\[
 S''(0)\le Q_G(z)+\sum_{e\in E_U}\Delta_e
                            \le Q_G(z).
\]
The corresponding unit Hilbert--Schmidt action
Hessian is at most
\[
                 -\frac{327264906217}{1863703068672}
                                    \le-\frac16 .
\]
There is no bound on the number of active selected
edges in this theorem.
\end{theorem}

\begin{proof}
The free--free and old uncontrolled terms do not
change from the V6 upper accounting. Nor do the
27 controlled terms with no selected pivot.
For a forward mixed plaquette replace
\(-2z_l^2\) by \(-2(z_l-a_e)^2\); summing its
three differences gives
\(4a_e\sum_{F_e}z_l-6a_e^2\).
For a backward mixed plaquette replace
\((-2+\epsilon)z_l^2\) by
\((-2+\epsilon)(z_l-a_e)^2\).
The \(u_e\) new plaquettes cost at most \(2u_ea_e^2\).
These are exactly the three parts of
\eqref{eq:v9-delta}, proving the first inequality.

Every part of \(\Delta_e\) is nonpositive:
\(a_e\le0\), all \(z_l>0\), \(u_e\le3\),
\(-2+\epsilon<0\), and
\(a_e^2-2a_ez_l\ge0\). Thus \(\Delta_e\le0\).
In fact it is strictly negative whenever \(a_e<0\),
because of the three forward cross terms.
This is a termwise comparison valid for all
unprescribed fields; it does not assume simultaneous
attainment of the separate upper bounds.
Finally divide by the independent-group speed \(N\),
not by a speed that adds the dependent pivots.
\end{proof}

The argument proves all \(2^{152}\) sign choices
without a global enumeration. Each edge has 256
local companion patterns, and the nonpositive
comparison is additive over edges.

For reference, when all 152 companions have sign
\(-1\), each selected edge has eight inner-active
sources. The exact arithmetic gives
\[
 \sum_e\Delta_e=-\frac{162929950203207}{2097152},
 \qquad
 Q_G(z)+\sum_e\Delta_e
                         =-\frac{178638665701623}{2097152}.
\]
The normalized upper bound is then
\[
                       -\frac{59546221900541}{29819249098752}.
\]
Only the weaker uniform bound \(-1/6\) is used for
the common repair theorem.

\subsection{Noncentral neighbourhoods with the true Jacobian}

The full derivative input closure now has 2344 free
groups, not V6's 2065. It includes the 279 additional
inputs needed for the newly moving-pivot plaquettes.
It contains \(G\); the other 1575 inputs remain
unprescribed. Its vertices lie in
\[
                  [-2,7]\times[-2,7]\times[0,3]\times[-1,7].
\]
The translation by \((8,8,8,8)\) is still non-wrapping
for every even \(L\ge32\).

For sufficiently small \(r>0\), let \(D_r\) restrict
the groups of \(U\) to geodesic \(r\)-balls about
\(-I\), those of \(W\) to \(r\)-balls about \(I\),
and each companion to the union of its \(r\)-balls
about \(I\) and \(-I\). It imposes no conditions on
any independent group outside \(G\).

\begin{theorem}[Original-clock repair of all companion-sign tubes]
\label{thm:v9-repair}
There exist \(r>0\) and \(\beta_0<\infty\) such that,
for \(C=I\), even \(L\ge32\), and \(G\subset\mathcal S\),
\[
 \kappa_{D_r}^{H_\nu}
 \ge\frac{c\beta}{2\{c\beta+8C_{\rm f}(\beta+1)\}}
 \ge\frac{c}{2(c+16C_{\rm f})}>0,\qquad
                    c=\frac1{96},\quad\beta\ge\beta_0.
\]
The constants are independent of ambient volume
and compatible hidden index sets. No factor
\(2^{152}\) or event-probability factor is lost.
The radius and threshold are existential.
\end{theorem}

\begin{proof}
Use orthonormal product normal coordinates about
each sign centre and normalize the fixed \(z\)-direction
on \(U\). Theorem~\ref{thm:v9-curvature} gives its
action Hessian at most \(-1/6\) for every value of
the other inputs in the finite compact closure.
Smoothness and compactness, including the finite
sign set, give a common outer radius \(2r\) where
the frozen coordinate second derivative is at
most \(-1/12\). Shrink \(r\) to preserve the inner
or inactive source regimes and inverse metric
lower bound \(1/2\).

The inverse-pivot amplitude is not discarded.
Only the 19 selected-edge Jacobian factors can
depend on a selected coordinate. Their positive
smooth inverses on the compact source carrier
have uniformly bounded logarithmic derivatives
through second order. The coordinate-Haar terms
are bounded on the chosen charts as well.
The actual measured potential is
\[
                   \Phi_\beta=\beta S-\log j-\log h_{\rm Haar}.
\]
For sufficiently large \(\beta\), its second
derivative is at most \(-\beta/24\). The local
conjugation estimate therefore gives the diffusion
floor \(\lambda=\beta/96\) on each chart component.

There is a single smooth cutoff for the whole union.
Take \(0\le b_\pm\le1\), equal to one on the
\(r\)-ball about \(\pm I\), supported in the
\(2r\)-ball, with disjoint supports and gradients
bounded by \(C/r\). Set
\[
 \chi(Y_G)=
   \prod_{l\in U}b_-(Y_l)
   \prod_{l\in W}b_+(Y_l)
   \prod_{c\in\mathcal R^\circ}\bigl(b_+(Y_c)+b_-(Y_c)\bigr).
\]
Then \(\chi=1\) on \(D_r\), \(0\le\chi\le1\),
and
\[
                       |\nabla\chi|^2\le769\,C^2/r^2=:M.
\]
For a companion coordinate at most one bump is
nonzero, so this bound does not grow with the
number of sign components. The supported components
are disjoint; the local diffusion bound adds over
them without an overlap loss. Consequently
\[
 \|1_{D_r}u\|^2\le\|\chi u\|^2
       \le\frac2\lambda\E_L(u,u)+\frac{2M}{\lambda}\|u\|^2.
\]
Increase \(\beta_0\) so that \(\lambda\ge8M\).
The existing global resolvent comparison and
conditional-variance transfer to the exact marginal,
using \(G\subset\mathcal S\), give the stated
original-clock floor. The cutoff never restricts
an unprescribed input or changes the heat-bath clock.
\end{proof}

With all companion signs negative, take a small
nonzero \(t\) on the explicit curve of
Lemma~\ref{lem:v9-pivot}. All 19 pivots then differ
from \(I\), and all 152 positive relative source
logarithms are nonzero inner-source terms:
\(z_l-a_e>0\) and \(a_e<0\). For sufficiently small
\(|t|\) these fields lie in \(D_r\).
Thus the theorem includes genuinely noncentral
configurations with every selected edge active,
not just centres where active responses occur at
zero logarithm.

\subsection{Verification and the remaining coverage problem}

The checker is
\begin{center}\small
\path{scripts/verify_ym_f14_v9_favorable_pivots.py}.
\end{center}
Its SHA--256 is
\begin{center}\scriptsize\ttfamily
DFF517B7D23C4CC44DCF5D5FA7AE527E172B46FC1DF503DD46D2AE936BB343C8.
\end{center}
It verifies the complete 564-plaquette incidence,
the three forward neighbours per pivot, all row
counts, and the enlarged input closure. It checks
all \(19\cdot256=4864\) local source patterns,
including the exact pivot equation and
\(\Delta_e\le0\), using fractions. The largest
nonempty local correction is \(-118976/15<0\).
It also checks the mixed-plaquette identity in
348 exact noncommutative quaternion jet evaluations.
Those evaluations test the implementation; cyclicity
of the trace and the exact pivot formula prove
the general identities.

The predecessor V8 has 140,606 bytes and SHA--256
\begin{center}\scriptsize\ttfamily
04097E587778BEE05B1AB4BB97FCB75A3B67DC3ED53933387449B1D5FE844678.
\end{center}
Its body is preserved apart from the title version.
Active V9 replaces V8; the archives, monographs,
and companion notes are preserved. The next
version, V10, is the scheduled companion and
direction checkpoint.

\begin{firewall}
The one-edge limitation is removed for source-companion
sign patterns and their controlled noncentral tubes:
all 19 selected edges may now be active, with the
same coefficient \(1/96\). This does not contradict
V8's positive frozen-mask matrix; the present
directions retain the active pivots' exact motion.
The direct factors on \(U\cup W\) still follow the
specified local sheet pattern, and companions are
near \(I\) or \(-I\). Arbitrary source-transition
fields and arbitrary local patterns remain
uncovered. No all-exterior patch gap, cofinal
continuum construction, or physical mass gap
has been proved.
\end{firewall}

The next question is whether the favourable
mixed-plaquette accounting persists for broader
noncentral source data when the pivot response
is a matrix rather than a scalar common-axis
speed. It must be answered with the actual
implicit derivatives and amplitude; an absolute-value
debit would erase the sign information responsible
for the present improvement.
\section{V10: explicit noncommuting source domains and the companion checkpoint}

\subsection{Companion review and the actual remaining target}

The scheduled V10 review finds the same companion files as
at V5. Their current sizes and SHA--256 digests are:
\begin{center}\small
\begin{tabular}{lr}
\toprule
File&Bytes\\ \midrule
\path{Renewal_SM_Einstein_Continuum_Research_Notes_V72.tex}&607372\\
\path{Renewal_NS_Stability_Research_Notes_V26.tex}&278283\\
\path{Renewal_YM_Parallel_Research_Notes_V5.tex}&54174\\
\bottomrule
\end{tabular}
\end{center}
\begin{center}\scriptsize\ttfamily
SM: F44F9745D43379E1672C5550411B58E97195A4C9278592D221DF5D3F644DC1F5\\
NS: 82C887F8C2C69E4F28FAD7C3DC8DE5B9099CB5A4BF431EBA1FFF3E91935978AD\\
YM: 306F1BB84CFA8A8D47EA569EC17E9545F9867C8D32B548EC9D9C444B4F3C0512
\end{center}

SM V72's many-copy covariance transport keeps the momentum
mode set and time interval fixed. It does not identify that
limit with the physical finite-species continuum or justify
a growing-cutoff interchange. NS V26's directional
rank-one contractions reinforce the need to retain the
structure of the tested direction, rather than replace it
by an unrelated full operator norm. Its numerical comparisons
and forward references to V27 are not imported here as
completed analytic certificates. The parallel YM note's
extensive-charge obstruction under tensorization also remains
unchanged; it is not an all-connected-torus theorem.
No result from these latest reviewed additions supplies the
missing all-exterior YM patch gap or physical-transfer comparison.
Their internal instructions are not instructions for this task.

The closure monograph still requires complete same-law patch
control crossing its finite-size threshold, followed by
cofinal construction and a calibrated physical comparison.
We therefore keep V9's favourable response as a local tool
and now quantify a genuinely noncommuting source domain.
The active sources may have unrelated axes; the proof
retains the common axis only for the chosen free-link
test direction. Inactive companions need not be near \(I\).
The next companion review is V15.

\subsection{An explicit two-branch condition on each companion}

Keep \(Y_l=-I\) on \(U\), \(Y_l=I\) on \(W\), and put
\[
                 \delta=2^{-18}=\frac1{262144}.
\]
For each \(c\in\mathcal R^\circ\), require one of
\begin{equation}\label{eq:v10-companion}
 \begin{array}{ll}
 \text{active branch:}&d_{\rm geo}(-I,Y_c)\le\delta,\\
 \text{inactive branch:}&\|I+Y_c\|_{\rm HS}\ge1/32.
 \end{array}
\end{equation}
These closed branches are disjoint. A free companion's
source product at these direct-link values is \(P=-Y_c\),
irrespective of which end of the path contains it.
Thus an active source is \(P=e^p\) with
\(\|p\|_{\rm HS}\le\delta\). An inactive source is
otherwise arbitrary, subject to the displayed chordal
separation. Every selected edge has six further tree
sources equal to \(-I\), and at most eight active sources.
All companion axes may differ and need not commute.

\begin{lemma}[The actual source regimes and root radius]
\label{lem:v10-regimes}
Under \eqref{eq:v10-companion}, the actual selected
pivot \(B_e\) has \(d_{\rm geo}(I,B_e)\le\delta\).
Every declared active relative source is within
geodesic radius \(2\delta<\rho/2\), so its cutoff is
identically one. Every declared inactive relative
source is at geodesic distance at least \(9/512>\rho\).
\end{lemma}

\begin{proof}
The global coarse-identity cap is \(R=7/512\).
An inactive source therefore has relative distance
at least \(1/32-R=9/512\). Tree sources are farther
still, so none contributes a truncated logarithm.
Let \(k\le8\) be the number of active sources and
let \(r=d_{\rm geo}(I,B_e)\).
From the exact equation \(B_e=e^{-X_e}\), the
small global bound on \(X_e\), and
\(\|g_\rho(Q)\|_{\rm HS}\le d_{\rm geo}(I,Q)\), we obtain
\[
 r\le\frac1{16}\sum_{j=1}^k d_{\rm geo}(B_e,P_j)
          \le\frac{k}{16}(r+\delta).
\]
Hence \(r\le k\delta/(16-k)\le\delta\). The triangle
inequality gives active relative radius at most
\(2\delta\), strictly inside the inner cutoff region.
This argument establishes the source regimes at
the actual root; it does not assume them to solve
for a proposed root.
\end{proof}

\subsection{Noncommuting logarithm estimates with numerical constants}

All matrix norms in this subsection are Hilbert--Schmidt
norms, except where an operator norm is explicitly indicated.
On triples of Lie-algebra variables use their sum norm.
Consider either ordered function
\[
 \mathcal L(b,p,q)=\log(e^{-b}e^pe^q),
       \qquad\hbox{or}\qquad
 \mathcal L(b,p,q)=\log(e^{-b}e^qe^p).
\]
Arguments and differentiation directions are anti-Hermitian.

\begin{lemma}[Uniform derivative bounds for the ordered logarithms]
\label{lem:v10-log}
If \(\|b\|+\|p\|+\|q\|\le1/4\), then
\[
                    \|D^2\mathcal L\|<4,\qquad
                    \|D^3\mathcal L\|<12.
\]
At the origin,
\[
 D_b\mathcal L=-I,\quad D_q\mathcal L=I,\quad
 D_{bb}\mathcal L=D_{qq}\mathcal L=0,\quad
 D_{bq}\mathcal L[v,H]=-\tfrac12[v,H].
\]
\end{lemma}

\begin{proof}
For \(j\ge1\), the \(j\)-th derivative of the
exponential at an anti-Hermitian argument has
multilinear Hilbert--Schmidt norm at most one.
Indeed the iterated integral formula is a sum over
\(j!\) permutations of products of unitary exponential
factors and the \(j\) directions, integrated over
a simplex of volume \(1/j!\). Use one direction's
Hilbert--Schmidt norm and the operator norms of all
other factors, with \(\|H\|_{\rm op}\le\|H\|_{\rm HS}\).
The same bound holds for derivatives of
\(U=e^{-b}e^pe^q\), or the alternate order, in
the sum norm: distributing derivatives over the
three factors gives the product of the direction
sum norms.

Unitary telescoping gives
\(\|U-I\|_{\rm op}\le\|U-I\|_{\rm HS}\le1/4\).
For \(\|X\|_{\rm op}\le r<1\), termwise
differentiation of \(\log(I+X)\) yields
\[
 \|D^m\log(I+X)[Z_1,\ldots,Z_m]\|_{\rm HS}
 \le\frac{(m-1)!}{(1-r)^m}\prod_i\|Z_i\|_{\rm HS}.
\]
To see the constant, differentiating \(X^n/n\)
gives \((n-1)!/(n-m)!\) ordered products; summing
their norm bounds over \(n\ge m\) gives the formula.
The chain rule at \(r=1/4\) therefore gives
\[
 \|D^2\mathcal L\|\le\frac4{3}+\frac{16}{9}
                    =\frac{28}{9}<4,
\]
\[
 \|D^3\mathcal L\|
       \le\frac4{3}+3\frac{16}{9}+2\frac{64}{27}
                    =\frac{308}{27}<12.
\]
Finally expand the two exponentials at \(p=0\)
through degree two:
\(\log(e^{-b}e^q)=-b+q-\tfrac12[b,q]+O(3)\).
This gives the stated origin derivatives for both
orders, without any commuting assumption.
\end{proof}

\subsection{The actual implicit velocity and acceleration}

Fix the V5 weighted direction \(E_n\). On an active
source path put \(H_j=z_{l(j)}E_n\), and write
\[
       S_e=\sum_{j=1}^k\|H_j\|_{\rm HS},\qquad
       T_e=\sum_{j=1}^k\|H_j\|_{\rm HS}^2 .
\]
Its moving source is either \(e^{p_j}e^{tH_j}\) or
\(e^{tH_j}e^{p_j}\), according to the actual path
order. Since the source regimes have strict margins,
the exact pivot equation near \(t=0\), with
\(b(t)=\log B_e(t)\), is
\[
 F(b,t)=b+\frac1{16}\sum_{j=1}^k
                  \mathcal L_j(b,p_j,tH_j)=0 .
\]
Let
\[
 \alpha=\frac{16-k}{16},\qquad
 \bar b'=-\frac{\sum_jH_j}{16-k}=a_eE_n,\qquad a_e\le0.
\]
This \(\bar b'\) is the favourable reference velocity
with the same active-path mask and all active source
products replaced by \(I\).

\begin{lemma}[Logarithmic implicit-jet bounds]
\label{lem:v10-implicit}
At the actual pivot,
\[
 \|b\|\le\delta,\quad
 \|b'-\bar b'\|\le3\delta S_e,\quad
 \|b'\|\le S_e/4,\quad
 \|b''\|\le50\delta T_e.
\]
For \(k=0\), all these quantities vanish.
\end{lemma}

\begin{proof}
The radius bound is Lemma~\ref{lem:v10-regimes}.
In the estimates below it is enough that
\(\delta\le1/512\), which the chosen value satisfies.
Compare derivatives of each ordered logarithm at
\((b,p_j,0)\) with the origin along a segment of
sum-norm length at most \(2\delta\).
Lemma~\ref{lem:v10-log} gives
\[
 \|F_b-\alpha I\|\le4\delta,\qquad
 \left\|F_t-\frac1{16}\sum_jH_j\right\|
                                      \le\frac{\delta S_e}{2}.
\]
Since \(\alpha\ge1/2\), the inverse norm is at most
3. Differentiating \(F(b(t),t)=0\) and subtracting
the reference equation gives, for \(\eta=b'-\bar b'\),
\[
 \|\eta\|
 \le3\left(\frac{\delta S_e}{2}
                     +4\delta\frac{S_e}{8}\right)
 =3\delta S_e.
\]
Thus \(\|b'\|\le S_e(1/8+3\delta)\le S_e/4\).

The second-derivative comparison gives
\[
 \|F_{bb}\|\le12\delta,\qquad
 \|F_{bt}-F_{bt}^{0}\|\le\tfrac32\delta S_e,\qquad
 \|F_{tt}\|\le\tfrac32\delta T_e,
\]
where
\(F_{bt}^{0}[v]=-\tfrac1{32}\sum_j[v,H_j]\).
The reference velocity commutes with every \(H_j\);
this directional cancellation is retained:
\(F_{bt}^{0}[\bar b']=0\).
Using \(\|[X,Y]\|\le2\|X\|\|Y\|\) gives
\[
                 \|F_{bt}^{0}[\eta]\|
                                  \le3\delta S_e^2/16 .
\]
Twice differentiating the implicit equation now yields
\[
 \begin{split}
 \|b''\|
 &\le3\delta\left\{
       12(S_e/4)^2+
       2\left(3S_e^2/16+\tfrac32 S_e(S_e/4)\right)
                         +\tfrac32 T_e\right\}\\
 &=\delta\left(\tfrac{45}{8}S_e^2+\tfrac92T_e\right)
 \le\tfrac{99}{2}\delta T_e\le50\delta T_e ,
 \end{split}
\]
because \(S_e^2\le kT_e\le8T_e\).
The commuting relation used here is between the
reference test velocities, not between the actual
source matrices \(p_j\).
\end{proof}

\begin{lemma}[Matrix jets compared with the favourable reference]
\label{lem:v10-matrixjets}
Let \(\bar B_e(t)=e^{t\bar b'}\). At \(t=0\),
\[
 \|B_e-I\|\le\delta,\qquad
 \|B'_e-\bar b'\|\le4\delta S_e,\qquad
 \|B''_e-(\bar b')^2\|\le60\delta T_e .
\]
The same bounds hold for inverse-oriented factors.
\end{lemma}

\begin{proof}
The exponential derivative bounds in the proof of
Lemma~\ref{lem:v10-log} imply
\[
 \|D\exp_b[b']-b'\|\le\delta\|b'\|,
 \qquad
 \|(D^2\exp_b-D^2\exp_0)[b',b']\|
                                      \le\delta\|b'\|^2.
\]
The first jet error is at most
\(\delta S_e/4+3\delta S_e\le4\delta S_e\).
For the second use
\(B''=D\exp_b[b'']+D^2\exp_b[b',b']\) and
\(D^2\exp_0[b',b']=(b')^2\). The error is at most
\[
 50\delta T_e+\delta S_e^2/16
           +3\delta S_e(S_e/4+S_e/8)
 \le\tfrac{119}{2}\delta T_e\le60\delta T_e .
\]
For a unitary matrix, inversion is adjoint;
derivatives and jet differences are adjointed
termwise, preserving their Hilbert--Schmidt norms.
\end{proof}

\subsection{Complete action-error account}

Set \(Z_2=\sum_{l\in U}z_l^2=21328389\).
Compare the actual physical jets with auxiliary jets
which replace each selected pivot by \(\bar B_e(t)\),
retaining all other actual physical factors.
This auxiliary comparison is not asserted to be
a second completed field with unchanged exterior.
V9's favourable bound applies to its plaquette
expression, since it only uses the reference
velocities \(a_e\le0\), the prescribed direct
factors, and the cap on other roots.

There is at most one selected pivot per plaquette.
For a mixed plaquette with selected-free amplitude
\(z_l\), the trace product rule and
\(|\operatorname{Tr}X|\le\sqrt2\|X\|_{\rm HS}\) give
the absolute error bound
\begin{equation}\label{eq:v10-plaq-error}
  \sqrt2\,\delta\{60T_e+8S_ez_l+z_l^2\}.
\end{equation}
The last two terms are absent on a pivot-only
plaquette. Indeed the free factor's first and
second derivative operator norms are \(z_l\) and
\(z_l^2\); insert the three bounds of
Lemma~\ref{lem:v10-matrixjets}. All other factors
are unitary, regardless of their hidden inputs.

Each pivot has six incident plaquettes, each
selected free link lies in at most six mixed
plaquettes, and each selected free link is in
at most two active source paths. Therefore
\[
                      \sum_eT_e\le4Z_2.
\]
Put \(s_e=S_e/\sqrt2=\sum_jz_{l(j)}\) and
\(q_e=\sum_{l\ {\rm mixed\ with}\ e}z_l\).
Cauchy--Schwarz and the same incidences give
\[
 \sum_es_e^2\le16Z_2,\quad
 \sum_eq_e^2\le36Z_2,\quad
                         \sum_es_eq_e\le24Z_2 .
\]
Summing \eqref{eq:v10-plaq-error} and using
\(\sqrt2\le3/2\) yields the three contributions
\[
 \begin{array}{c|r}
 \text{pivot second-jet error}&2160\delta Z_2\\
 \text{mixed first-jet error}&384\delta Z_2\\
 \text{pivot value times free acceleration}&9\delta Z_2
 \end{array}
 \qquad\text{total }2553\delta Z_2.
\]
The value error is included: one cannot compare
only the velocities of noncentral pivots.

\begin{theorem}[Uniform descent on the explicit noncommuting source domain]
\label{thm:v10-curvature}
At \(Y_U=-I\), \(Y_W=I\), under
\eqref{eq:v10-companion}, and with all independent
groups outside \(G\) arbitrary, the actual completed
action along the V5 direction satisfies
\[
 S''(0)\le Q_G(z)+2553\delta Z_2
                        =-\frac{1909138060185}{262144}.
\]
Its unit Hilbert--Schmidt Hessian is at most
\[
                   -\frac{636379353395}{3727406137344}
                                              <-\frac16 .
\]
All 19 selected edges may have noncommuting active
source data.
\end{theorem}

\begin{proof}
Apply V9's sign-sensitive accounting to the
reference jets with the actual active-path mask.
It bounds their action second derivative by
\(Q_G(z)\); inactive companions do not occur as
direct factors in the 564 affected plaquettes.
Other completed roots are fixed under the test
motion and retain their actual cap \(R\).
Add the complete absolute error just proved.
At \(\delta=1/262144\), that error is
\[
                        \frac{54451377117}{262144}.
\]
The displayed numerator follows by addition.
Divide by the independent-group squared speed
\(2Z_2=42656778\). The exact positive slack below
\(-1/6\) is
\[
                         \frac{15144997171}{3727406137344}.
\]
No pivot speed is added to the independent metric,
and no smallness assumption is made on a hidden
factor in this calculation.
\end{proof}

\subsection{Original-clock repair and what is quantitatively explicit}

Let \(\mathcal K\subset\SU(2)\) be the closed
union of the two companion branches in
\eqref{eq:v10-companion}. For \(r>0\), let
\(D_{r,\mathcal K}\) restrict \(U\) to geodesic
\(r\)-balls about \(-I\), \(W\) to \(r\)-balls
about \(I\), and each companion to \(\mathcal K\).
It restricts no other independent group.

\begin{theorem}[Repair with arbitrary far inactive companions]
\label{thm:v10-repair}
There exist \(r>0\) and \(\beta_0<\infty\) such
that, for \(C=I\), even \(L\ge32\), and
\(G\subset\mathcal S\),
\[
 \kappa_{D_{r,\mathcal K}}^{H_\nu}
 \ge\frac{c\beta}{2\{c\beta+8C_{\rm f}(\beta+1)\}}
 \ge\frac{c}{2(c+16C_{\rm f})}>0,\qquad
                    c=\frac1{96},\quad\beta\ge\beta_0.
\]
The constants are uniform in ambient volume and
compatible hidden index sets. The companion
thresholds are explicit; \(r\) and \(\beta_0\)
are existential, not numerically enclosed.
\end{theorem}

\begin{proof}
Theorem~\ref{thm:v10-curvature} gives a uniform
directional Hessian at most \(-1/6\) on the compact
parameter set \(Y_U=-I,\ Y_W=I,\ 
Y_{\mathcal R^\circ}\in\mathcal K^{152}\),
with all unprescribed inputs ranging over their
fixed compact closure. That closure is the
2344-group carrier of V9; no new physical group
has been added.

Use product normal coordinates only for the
133 groups of \(U\), centred at \(-I\), and freeze
one normalized V5 axis direction there. Treat
the other groups as compact Haar parameters.
Smoothness gives common neighbourhoods of \(U,W\)
and of the companion parameter set on which the
frozen coordinate Hessian is at most \(-1/12\).
The inverse metric in the \(U\) coordinates is
at least \(1/2\) after shrinking. Only 19 inverse
pivot factors depend on the direction; their
logarithmic second derivatives are uniformly
bounded, as are the coordinate-Haar derivatives
in \(U\). The true measured second derivative is
therefore at most \(-\beta/24\) for large \(\beta\).
The local diffusion floor is \(\lambda=\beta/96\).
No global coordinate chart on the far inactive
companion set is required.

Choose a smooth companion cutoff \(c_{\mathcal K}\),
equal to one on \(\mathcal K\), with support in
the selected companion neighbourhood and values
in \([0,1]\). The active and inactive branches
are separated, so it may be constructed as two
disjoint smooth cutoffs. Its gradient has a finite
bound \(C_{\mathcal K}\). Multiply these 152
cutoffs by the 617 standard inner--outer cutoffs
for \(U\cup W\). The product \(\chi\) is one on
\(D_{r,\mathcal K}\) and satisfies
\[
 |\nabla\chi|^2
      \le617C^2/r^2+152C_{\mathcal K}^2=:M.
\]
The local directional bound holds uniformly
over the companion support, so integration over
these Haar parameters gives
\[
 \|1_{D_{r,\mathcal K}}u\|^2
      \le\frac2\lambda\E_L(u,u)+\frac{2M}{\lambda}\|u\|^2.
\]
There is no multiplicity factor for source masks.
After requiring \(\lambda\ge8M\), the existing
joint resolvent comparison and exact marginal
conditional-variance domination give the stated
heat-bath floor, with unchanged rate-one group
updates.
\end{proof}

\subsection{Verification, provenance, and the next obstruction}

The checker is
\begin{center}\small
\path{scripts/verify_ym_f14_v10_noncommuting_sources.py}.
\end{center}
Its SHA--256 is
\begin{center}\scriptsize\ttfamily
D47B2952BF64677896F24ECF463F9154F640D0F9DA3D1D00452AB34A2130063C.
\end{center}
It reruns V9, verifies the relevant source and
plaquette multiplicities, and checks every
rational constant in the implicit-jet and
action-error accounts. The universal derivative
bounds are proved analytically above; a finite
sample is not presented as a proof for all
noncommuting data.

The predecessor V9 has 158,570 bytes and SHA--256
\begin{center}\scriptsize\ttfamily
9B3D07DAE19A7339CA50A31E05A0D883FC1042B88781890D92616827D7110014.
\end{center}
Its body is preserved apart from the title version.
Active V10 replaces V9. Archives, monographs, and
companion programmes are unchanged. The V10
companion checkpoint is complete; the next is V15.

\begin{firewall}
The source domain is now explicit and noncommuting:
an active source may lie anywhere in geodesic
radius \(2^{-18}\) of identity, while an inactive
source may be arbitrarily noncentral subject to
chordal separation \(1/32\). All selected edges
may be active. This is more than an unspecified
neighbourhood of central signs, but it still
excludes an annular source region, including
cutoff transitions. The direct factors \(U\cup W\)
still follow the local sheet pattern, and coarse
data remain \(I\). There is no complete all-field
cover, unrestricted patch gap, or physical
Yang--Mills mass-gap theorem.
\end{firewall}

The next obstruction is the source annulus between
the explicit active cap and the far inactive
region. Increasing the present perturbative
radius alone does not handle the transition:
there the cutoff derivative can change sign.
Any extension must keep that actual derivative
and its coupled pivot response, or use another
repair mechanism on the complementary fields.
\section{V11: inverse-paired noncommuting sources through the full cutoff transition}
\label{sec:v11-paired}

\subsection{New RG source, nonduplication, and the changed interface}

During this iteration the user supplied the separate companion
\begin{center}\small
\path{Renewal_Yang_Mills_Shell_Mixing_Research_Notes_V13.tex}.
\end{center}
The reviewed snapshot has 512,757 bytes and SHA--256
\begin{center}\scriptsize\ttfamily
7F2DBD2900D7A78280DA980152266E94DF8A95EFE15CD37429EC8C6AAAA4F17B.
\end{center}
It is preserved, not renamed into the main mass-gap lineage.
References to \emph{Shell V\(j\)} below mean sections of that file,
not sections of these notes. We inspected its structure and initial
scope, its relevant Gaussian and measured-transport statements,
and its terminal weak-source and classical-continuum arguments.
This is a dependency and nonduplication review, not an independent
recertification of every inherited proof or every script named there.
Its internal continuation directions are source material, not new
instructions to this programme.

Three pieces of its progress matter for the global direction.
First, its shell-response ledger retains mixed-shell contractions
and inherited vertices. The Gaussian locality theorem
\texttt{thm:v5-strip} concerns one common analytic strip; it must
not be used to reset the nonlinear action or density after each
elimination. Second, Shell V7's \texttt{thm:v7-main}, under its
specified local inputs, gives a measured original-box expansion
through the two-loop local Laplace coefficient, with an
\(O(\beta^{-2})\) remainder through four retained derivatives.
Shell V8's \texttt{thm:v8-budget} identifies the closed finite
packet
\[
                  (J^8S,\ J^6Q,\ J^4B;\ J^7F).
\]
Here \(Q,B\) are the incoming measured one- and two-loop
coefficients in that companion's convention; they are not the
group elements \(Q_l\) used below. The complete coarea factor
belongs to the density exactly once.

Third, Shell V12--V13 construct the literal balanced root observable
on a common small global counting-\(\ell^4\) ball. In particular
\texttt{thm:v13-weaksource} proves, in its declared dyadic sampling,
\[
 \sup_{\|A\|_{L^4}\le r}
 \|\mathcal O_{L,N}(A)-\mathcal O_{\infty,N}(A)\|_4
       \le C_{r,\alpha}L^{-\alpha/(1+\alpha)},\qquad 0<\alpha<1,
\]
on a sufficiently small ball, together with weak sequential
continuity for fixed \(N\). Its
\texttt{thm:v13-main} then identifies the full local constrained
classical minimum on one final periodic coarse cell: the values,
minimizing fields in the critical norms, and every fixed source
derivative of the scalar minimum converge. This replaces a
finite-order classical germ by an identified finite-amplitude
local classical functional. No convergence rate for the full
minimum, fluctuation determinant, or quantum state is supplied
by that theorem.

These acquisitions remove reasons to repeat a Gaussian alias
calculation or another homogeneous classical coefficient here.
They do not close the present compact-complement problem.
Their common small classical branch is not the unrestricted
compact conditional law; their remaining-scale \(\ell^2\)
contraction counts different lattice sizes and is not physical
time contraction. Conversely our local repair floor does not
bound their renormalized loop traces. The two streams must be
joined through the \emph{same measured law}, with source
coordinates, coarea factors, and actual update rates identified.
Merely placing their inequalities beside one another is not such
a theorem.

We therefore finish the concrete transition acquisition left by
V10. Unlike V10's independent small active caps, the source data
below satisfy an inverse-pair relation. Unlike V1's one-edge
outer transition band, all selected edges may lie anywhere in
the radial cutoff transition. Noncommuting axes are admitted.
The trade is explicit: pairing and a matrix cone replace small
radial displacement. This is a complementary domain, not a
claim that V11 contains every field of V9 or V10.

\subsection{Pair geometry and the actual source curve}

Retain \(C=I\), \(U\) of 133 selected independent links,
\(\mathcal E_U\) of 19 coarse edges, \(W\) of 484 direct
nonrotated groups, and \(\mathcal R^\circ\) of 152
source-only companions from V8--V10. Thus
\[
 G=U\sqcup W\sqcup\mathcal R^\circ,\qquad |G|=769.
\]
Use precisely the positive integer weights \(z_l\) of V5;
their canonical finite-link ordering and full list are in
that version and its checker. Their squared angular and
Hilbert--Schmidt speeds are respectively
\begin{equation}\label{eq:v11-speed}
 Z_2=\sum_{l\in U}z_l^2=21328389,\qquad
 N=2Z_2=42656778.
\end{equation}
As always, dependent pivot velocities are not added to the
independent product metric.

Each selected coarse edge contains seven links of \(U\).
For four of them \(l\), one source path has a free companion
immediately before \(l\), and the other has a free companion
immediately after \(l\). Denote these companions by \(c_l^-\)
and \(c_l^+\). The other three links have tree companions on
both paths. The 76 pairs exhaust all 152 distinct groups of
\(\mathcal R^\circ\). None of these groups is a direct factor
in any of the 564 physical plaquettes affected by moving
\(U\) and completing its pivots.

These assertions are finite path incidences, not a new
probabilistic assumption. They follow by inserting the V6
inner 19-edge set in the fourteen two-link paths of each edge.
The new checker constructs every path and verifies uniqueness,
the four-plus-three split, and the empty intersection with all
564 physical words. It uses the same nonmodular geometry as V9.

For \(a\in\mathbb R^3\), write \(E_a\) for its quaternion Lie
matrix, so \(E_aE_b=-(a\cdot b)I+E_{a\times b}\) and
\(\|E_a\|_{\rm HS}^2=2|a|^2\). Choose a common unit vector
\(h\in S^2\) and arbitrary \(Q_l\in\SU(2)\) for the 76 pairs.
The centre and its free curve are
\begin{equation}\label{eq:v11-pairfield}
 \begin{aligned}
 Y_l(t)&=-\exp(tz_lE_h) &&(l\in U),\\
 Y_w&=I &&(w\in W),\\
 Y_{c_l^-}&=-Q_l,\qquad
 Y_{c_l^+}=-Q_l^{-1} &&(l\hbox{ paired}).
 \end{aligned}
\end{equation}
No group outside \(G\) is prescribed. In particular this is
not a hidden-boundary conditioning.
The source pair belonging to \(l\) is exactly
\begin{equation}\label{eq:v11-sourcecurve}
          Q_l\exp(tz_lE_h),\qquad
          \exp(tz_lE_h)Q_l^{-1}.
\end{equation}
The six tree-companion sources on an edge are
\(-\exp(tz_lE_h)\), twice for each of its three other links.
Thus the source family at \(-t\) is the family of inverses
at \(t\), with the paired indices interchanged.

\subsection{Inversion symmetry removes logarithmic acceleration}

Put \(\rho=1/64\) and write \(g=g_\rho\) for the smooth
truncated logarithm in the inherited root equation. It is
globally inversion odd and conjugation equivariant:
\[
 g(Q^{-1})=-g(Q),\qquad
 g(VQV^{-1})=\operatorname{Ad}_Vg(Q).
\]
These identities hold through the cutoff endpoints, and also
near \(-I\) where \(g\) is identically zero.

\begin{lemma}[Exact inversion of the completed root]
\label{lem:v11-inversion}
At coarse datum \(I\), let \(B(P)\) be the unique completed
root for its fourteen sources \(P=(P_j)\). Then
\[
                        B(P^{-1})=B(P)^{-1}.
\]
For \eqref{eq:v11-sourcecurve}, \(B_e(0)=I\) and its small
principal logarithm \(b_e(t)=\log B_e(t)\) is odd. Consequently,
for a vector \(a_e\in\mathbb R^3\),
\begin{equation}\label{eq:v11-rootjets}
 b_e''(0)=0,\qquad B_e'(0)=E_{a_e},\qquad
                         B_e''(0)=-|a_e|^2I.
\end{equation}
\end{lemma}

\begin{proof}
Let \(X=16^{-1}\sum_jg(B^{-1}P_j)\), so \(B\exp X=I\).
For the proposed root \(B^{-1}\) and inverted sources, the
new exponent is
\[
 \frac1{16}\sum_jg(BP_j^{-1})
                  =-\operatorname{Ad}_B X.
\]
Indeed \(BP_j^{-1}=B(B^{-1}P_j)^{-1}B^{-1}\).
Hence
\[
 B^{-1}\exp(-\operatorname{Ad}_BX)
                  =\exp(-X)B^{-1}=I.
\]
Uniqueness of compact completion proves the first assertion.
At \(t=0\), paired truncated logs cancel and all tree sources
are \(-I\), so \(B_e=I\) solves the equation. Pairing of the
families at opposite times gives \(B_e(-t)=B_e(t)^{-1}\).
The global cap \(d_{\rm geo}(I,B_e)\le R=7/512\) keeps the
root in one principal logarithm chart. Its logarithm is odd,
and differentiation of \(B_e=\exp b_e\) gives
\eqref{eq:v11-rootjets}.
\end{proof}

This does not make the root stationary or its inverse-pivot
density constant. It removes precisely the second derivative
of its logarithm at the paired centre. A possibly large radial
derivative of the cutoff remains in \(a_e\) and is now computed.

\subsection{The exact three-dimensional response, including the transition}

For \(Q=\cos u\,I+\sin u\,E_n\), \(0<u<\pi\), \(|n|=1\),
let \(s=\sqrt2u/\rho\). On the support \(u<u_*:=\rho/\sqrt2\)
define
\begin{equation}\label{eq:v11-response}
 \begin{aligned}
 \lambda(Q)&=\chi(s)u\cot u,&
 d(Q)&=\chi(s)+s\chi'(s),&
 \gamma(Q)&=\chi(s)u,\\
 D(Q)&=\lambda I+(d-\lambda)nn^T,&
 \Gamma(Q)&=\gamma J_n,& J_nv&=n\times v .
 \end{aligned}
\end{equation}
At \(Q=I\), set \(D=I,\Gamma=0,\lambda=1,\gamma=0\).
For \(u\ge u_*\), set all four quantities \(D,\Gamma,\lambda,
\gamma\) to zero (and \(d=0\)).
These definitions have their smooth extensions at the
endpoints. In particular no choice of axis is needed at
central elements.

\begin{lemma}[Differentiated truncated logarithm]
\label{lem:v11-tensor}
In quaternion Lie coordinates,
\begin{equation}\label{eq:v11-left-right}
 \left.\frac d{dt}g(Qe^{tE_v})\right|_0=(D+\Gamma)v,\qquad
 \left.\frac d{dt}g(e^{tE_v}Q)\right|_0=(D-\Gamma)v.
\end{equation}
Moreover
\[
                0\le\lambda\le1,\qquad d\le1,\qquad D\preceq I.
\]
For an edge \(e\), sums over \(l\sim e\) below run over its
four paired links. Set
\begin{equation}\label{eq:v11-AM}
 A_e=8I-\sum_{l\sim e}D(Q_l),\qquad
 M_e=\sum_{l\sim e}z_l\{D(Q_l)+\Gamma(Q_l)\}.
\end{equation}
Then \(A_e\succeq4I\), and the actual completed-root velocity is
\begin{equation}\label{eq:v11-a}
                            a_e(h)=-A_e^{-1}M_eh.
\end{equation}
\end{lemma}

\begin{proof}
For right multiplication of \(Q\), quaternion multiplication
gives \(u'=n\cdot v\) and
\[
 n'=\cot u\{v-n(n\cdot v)\}+n\times v.
\]
Differentiate \(\chi(\sqrt2u/\rho)un\). The radial eigenvalue
is \(d\), the transverse symmetric eigenvalue is \(\lambda\),
and the skew term is \(\gamma J_n\). Left multiplication
reverses that skew term. This proves
\eqref{eq:v11-left-right}; the central and cutoff limits
follow by smoothness. Since \(\chi'\le0\), \(d\le\chi\le1\).
On \(0<u<u_*<\pi/2\), \(0<u\cot u\le1\), giving the remaining
inequalities. No nonnegative lower bound for \(d\) is used.

At \(B=I\), variation of \(B^{-1}\) is left multiplication by
\(-a_e\). The two members \(Q_l,Q_l^{-1}\) have the same \(D\)
and opposite \(\Gamma\), so their root derivatives sum to
\(-2D(Q_l)a_e\). Their prescribed source derivatives in
\eqref{eq:v11-sourcecurve} sum to
\(2z_l(D(Q_l)+\Gamma(Q_l))h\). Tree sources stay in the
zero-cutoff region near \(t=0\) and contribute nothing.
Differentiating \(B\exp X=I\) at \(B=I,X=0\) therefore gives
\[
 a_e+\frac18\sum_{l\sim e}
       \{-D(Q_l)a_e+z_l(D(Q_l)+\Gamma(Q_l))h\}=0 .
\]
This is \eqref{eq:v11-a}. Each \(D\preceq I\) implies
\(A_e\succeq4I\), so its inverse exists uniformly.
\end{proof}

\subsection{A response cone gives the unchanged action curvature}

For the single common free axis \(h\), decompose
\begin{equation}\label{eq:v11-cone}
 \alpha_e=-h\cdot a_e,\qquad
 \beta_e=a_e+\alpha_e h,\qquad
        \alpha_e\ge0,\quad |\beta_e|\le\alpha_e
                                  \quad(e\in\mathcal E_U).
\end{equation}
The vector \(\beta_e\perp h\) here is a transverse velocity;
it is unrelated to inverse temperature \(\beta\). This
explicit matrix test permits different response directions
on different edges, while all free links use the same \(h\).

\begin{theorem}[Paired-cone action certificate]
\label{thm:v11-curvature}
At a centre \eqref{eq:v11-pairfield} satisfying
\eqref{eq:v11-cone}, the actual completed Wilson action obeys,
uniformly over all independent groups outside \(G\),
\begin{equation}\label{eq:v11-action}
 S''(0)\le-\frac{981794718651}{131072},\qquad
             \Hess S[\widehat v,\widehat v]\le-\frac16,
\end{equation}
where \(\widehat v\) is the free curve's unit
Hilbert--Schmidt tangent. Sources in the cutoff transition
are allowed; their response is \eqref{eq:v11-a}, not the
untruncated inner-source formula.
\end{theorem}

\begin{proof}
We use the verified physical-word split of V9, retaining the
nonparallel root velocities. There are 57 forward free/pivot
plaquettes, three per selected edge, with no other root;
30 backward free/pivot plaquettes, each with one other root;
and 27 new pivot-only plaquettes. If \(u_e\) is the number
of new pivot-only plaquettes on \(e\), then
\(0\le u_e\le3\), and \(u_e\) plus its backward count is
three. The 27 remaining protected free plaquettes have two
outside roots and no moving pivot. The old free/free and
uncontrolled free plaquettes have no moving pivot.

Write \(x=z_lh\), \(a=a_e\). A forward plaquette contributes
exactly \(-2|x-a|^2\). This follows either by multiplying its
two exponentials or from \eqref{eq:v11-rootjets}: only the
two-jet, not a globally exponential root path, is needed.
For a backward plaquette, let \(r_0I+E_r\) be its one
arbitrary fixed root, after any required inversion.
Multiplication of the two jets gives
\begin{equation}\label{eq:v11-cross}
                  -2r_0|x-a|^2
                        \mathbin{\pm}4r\cdot(x\times a).
\end{equation}
The sign depends on its orientation and is irrelevant for
the upper bound. One can also derive this identity by
expanding the ordered product: its second jet has scalar
part \(-|x-a|^2/2\) per second-order Taylor coefficient
and vector cross term of either sign. All remaining
factors at the centre are fixed \(I\), apart from the one
selected factor \(-I\).

The root cap \(R=7/512\) implies
\[
 r_0\ge1-R^2/4,\qquad |r|\le R/\sqrt2 .
\]
Consequently \eqref{eq:v11-cross} is at most
\[
                  c|x-a|^2+2\sqrt2Rz_l|\beta_e|,
       \qquad c=-2+\epsilon,\quad \epsilon=\frac{49}{524288}.
\]
A new pivot-only plaquette contributes at most \(2|a|^2\):
its other factors are fixed unitary matrices and its only
moving second jet is \(-|a|^2I\). This bound does not require
those other factors to commute.

Compare with the root-stationary upper certificate of V6,
whose value is the first fraction in \eqref{eq:v11-action}.
The only changes are the three classes just described.
For each forward weight \(z\),
\[
 |zh-a|^2-z^2=|a|^2+2z\alpha_e.
\]
The three forward plaquettes and the new pivot-only ones
therefore change that upper bound by at most
\[
             (2u_e-6)|a_e|^2
                         -4\alpha_e\sum_{\rm forward}z_l\le0 .
\]
For each backward weight, the change is at most
\[
 c|a_e|^2+2c z_l\alpha_e+2\sqrt2Rz_l|\beta_e|
       \le c|a_e|^2+(2c+2\sqrt2R)z_l\alpha_e\le0.
\]
Indeed \(c<0\), and \(\sqrt2\le3/2\) gives the exact bound
\[
           2c+2\sqrt2R\le2c+3R
                         =-\frac{1037775}{262144}<0.
\]
All unchanged classes retain their V6 estimates; companions
have no direct occurrence in any affected word. Summing
establishes the first bound. Divide by \(N\) in
\eqref{eq:v11-speed}, and use the V6 rational comparison
with \(-1/6\), to obtain the unit tangent bound.
\end{proof}

\subsection{Every radial phase is admitted for coplanar source axes}

The cone has a useful genuinely noncommutative populated
subfamily. Suppose there is one \(h\in S^2\) perpendicular
to the imaginary quaternion part of every \(Q_l\).
Thus all noncentral source axes lie in the same plane
\(h^\perp\); their phases may differ arbitrarily and their
axes need not be parallel.

\begin{theorem}[All-phase coplanar certificate]
\label{thm:v11-plane}
For inverse-paired coplanar data, all 19 responses satisfy
\[
                   \alpha_e\ge0,\qquad
                          |\beta_e|\le\frac3{64}\alpha_e.
\]
Hence Theorem~\ref{thm:v11-curvature} holds for every
\(Q_l\in\SU(2)\) in this family. In particular it holds
through the entire cutoff transition, including both
endpoints, with no restriction on the negative size of
its radial derivative.
\end{theorem}

\begin{proof}
For \(n_l\perp h\), \(D_lh=\lambda_lh\). The same identity
holds at a central \(Q_l\) by the definitions. Thus \(A_e\)
preserves \(\mathbb Rh\) and its perpendicular plane.
Put \(L_e=\sum_{l\sim e}\lambda_l\) and
\(P_e=\sum_{l\sim e}z_l\lambda_l\). Formula
\eqref{eq:v11-a} gives
\[
 \alpha_e=\frac{P_e}{8-L_e}\ge0,\qquad
 \beta_e=-(A_e|_{h^\perp})^{-1}
                \sum_{l\sim e}z_l\gamma_l(n_l\times h).
\]
At central elements the corresponding skew summand is zero.
Since \(A_e\succeq4I\), its plane inverse has norm at most
\(1/4\). On the nonzero support,
\[
                    |\gamma_l|=\lambda_l\tan u_l
                                      \le\lambda_l\tan u_* ;
\]
outside the support both sides vanish. Therefore
\[
 |\beta_e|\le\frac{\tan u_*}{4}P_e
                \le 2\tan u_*\,\alpha_e.
\]
For \(0\le u\le u_*<1\), \(\sin u\le u\) and
\(\cos u\ge1-u^2/2\ge1/2\), so \(\tan u\le2u\).
It follows that
\[
       2\tan u_*\le4u_*=\frac{\sqrt2}{32}\le\frac3{64}<1.
\]
If \(P_e=0\), every weighted \(\lambda_l\) is zero,
every \(\gamma_l\) is zero, and \(a_e=0\), so the
same conclusion holds without dividing by \(\alpha_e\).
The potentially negative \(d_l\) acts only inside
\(h^\perp\); its upper bound \(d_l\le1\) is precisely
what supplies the favourable inverse bound. It has
not been omitted or replaced by \(\chi\).
\end{proof}

This proof is an all-phase inequality, not an extrapolation
from finitely many radial samples.

\subsection{An explicit noncommuting midpoint on all 19 edges}

Take \(h=e_3\) and
\[
                  u=\frac{3}{256\sqrt2},\qquad s=\frac34.
\]
On each edge order its four paired links canonically and set
\[
             (Q_1,Q_2,Q_3,Q_4)
        =(e^{uE_1},e^{uE_2},e^{uE_1},e^{uE_2}).
\]
The first two matrices do not commute:
\(Q_1Q_2-Q_2Q_1=2\sin^2u\,E_3\ne0\).
Every one of the 152 paired sources is strictly in the
cutoff transition at its actual root \(I\).

For the literal cutoff,
\[
 \chi(3/4)=\frac12,\qquad \chi'(3/4)=-8,\qquad
                   d=\chi+s\chi'=-\frac{11}{2}.
\]
These are exact: its logistic exponent vanishes at \(3/4\)
and has derivative \(-32\) there.
Let \(\lambda=(u\cot u)/2\), \(\gamma=u/2\), and write
the four weights on an edge as \(z_1,\ldots,z_4\).
Then the response is explicitly
\begin{equation}\label{eq:v11-midpoint}
 \begin{split}
 A_e|_{h^\perp}&=(19-2\lambda)I,\qquad A_eh=(8-4\lambda)h,\\
 \alpha_e&=\frac{\lambda(z_1+z_2+z_3+z_4)}{8-4\lambda}>0,\\
 (\beta_e)_1&=-\frac{\gamma(z_2+z_4)}{19-2\lambda},\qquad
 (\beta_e)_2=\frac{\gamma(z_1+z_3)}{19-2\lambda}.
 \end{split}
\end{equation}
Thus all 19 pivots have nonzero velocity. Small nonzero
times keep all these sources in transition and make all
19 roots nonidentity. The repair neighbourhood below
includes such configurations, not just the symmetric
centre where the root itself is \(I\).

There is a useful warning about selecting the free axis.
If instead all four axes on every edge are parallel to the
free direction at the same midpoint, the radial formula
gives
\[
                        a_e=\frac{11}{60}
                                  \sum_{l\sim e}z_l>0
\]
as a scalar velocity along that direction. Inserting these
velocities into the old scalar upper certificate gives
\[
                         \frac{786147299098687}{31457280}>0.
\]
This is a failed upper bound for that prescribed direction,
not a positive lower bound on the actual Hessian and not an
obstruction to every descent direction. The transverse
choice in Theorem~\ref{thm:v11-plane} explains why this radial
failure does not prevent repair of the same source phases.

\subsection{Open visible tubes and the original-clock floor}

Let \(\mathcal K_{\rm cone}\subset\SU(2)^G\) consist of the
centres \eqref{eq:v11-pairfield} for which at least one common
\(h\in S^2\) satisfies \eqref{eq:v11-cone}. This definition
does not select an axis as a function of the field.
The matrices \(D,\Gamma,A^{-1}\) are continuous everywhere,
including central elements and transition endpoints.
The closed cone conditions define a compact subset of
\(\SU(2)^{76}\times S^2\); its continuous image in
\(\SU(2)^G\) is therefore compact.

Let \(D_r\) be the cylinder in which the \(G\) coordinates
have product-geodesic distance at most \(r\) from this
compact core; no independent coordinate outside \(G\) is
restricted. Unlike the core, \(D_r\) need not satisfy
exact pairing or exact coplanarity.

\begin{theorem}[Uniform original-clock repair of the paired-cone tube]
\label{thm:v11-repair}
There are \(r>0\) and \(\beta_0<\infty\) such that, for
identity coarse data, even \(L\ge32\), and retained index
set \(\mathcal S\) with \(G\subset\mathcal S\),
\begin{equation}\label{eq:v11-HB}
 \kappa_{D_r}^{H_\nu}
 \ge\frac{c_{\rm rep}\beta}
          {2\{c_{\rm rep}\beta+8C_{\rm f}(\beta+1)\}}
 \ge\frac{c_{\rm rep}}{2(c_{\rm rep}+16C_{\rm f})}>0,
 \quad c_{\rm rep}=\frac1{96},\quad \beta\ge\beta_0 .
\end{equation}
The radii and threshold are existential; the coefficients
in the action and rate comparison are explicit.
Constants are independent of ambient volume and of
compatible hidden index sets. There is no event-probability
factor and no acceleration of the rate-one group updates.
\end{theorem}

\begin{proof}
The full derivative input closure is the same 2,344 groups
as in V9; only 769 are constrained. Its vertex bounds are
\[
                 [-2,7]\times[-2,7]\times[0,3]\times[-1,7].
\]
Translate by \((8,8,8,8)\) to embed without wrapping on
every even \(L\ge32\). This supplies a fixed finite compact
space of all remaining inputs, over which
Theorem~\ref{thm:v11-curvature} is uniform.

At each point of \(\mathcal K_{\rm cone}\), freeze one
certifying axis and use orthonormal product normal
coordinates on \(U\) about \(-I\); all other groups may
remain Haar parameters. Smoothness and compactness give
a \(G\)-neighbourhood of that point on which the frozen
unit coordinate action second derivative is at most
\(-1/12\), for all unprescribed inputs. After shrinking,
the inverse metric on those coordinates is at least
\(1/2\). This extension need not preserve the pairing or
the cone: it preserves the strict Hessian inequality.
No derivative of a chosen \(h\), \(A^{-1}M\), or
eigenvector field enters this argument.

There are only 19 direction-dependent inverse-pivot
Jacobian factors. Their logarithmic second derivatives
are bounded on the full compact source space by positivity
and smooth completion. The coordinate-Haar derivatives
are bounded on the \(U\) charts. Thus for sufficiently
large \(\beta\), uniformly on the chart cylinders,
\[
 \partial_{\widehat v}^2
       \{\beta S-\log j-\log h_{\rm Haar}\}\le-\beta/24.
\]
The inherited local conjugation estimate and inverse
metric bound give a diffusion floor \(\lambda=\beta/96\)
on each chart cylinder. In particular the nonconstant
inverse-pivot amplitude has not been dropped.

Take a finite cover of the compact core by smaller
such cylinders and choose a common \(r>0\) whose closed
tube is still covered. As in V7, there are subordinate
smooth cutoffs \(\chi_j(Y_G)\) with
\[
 \sum_j\chi_j^2\le1,\qquad
 \sum_j\chi_j^2=1\ \hbox{on }D_r,\qquad
                  M:=\sup\sum_j|\nabla\chi_j|^2<\infty.
\]
One may form them from subordinate nonnegative bumps
\(\psi_j\) by normalizing with
\((T+g(T)^2)^{1/2}\), where \(T=\sum_j\psi_j^2\ge1\)
on the tube and \(g=1\) for \(T\le1/4\), \(g=0\)
for \(T\ge1/2\). Hence
\[
 \lambda\|1_{D_r}u\|^2
 \le\sum_j\E_L(\chi_ju,\chi_ju)
                  \le2\E_L(u,u)+2M\|u\|^2.
\]
This is one finite configuration-space cover, not a
growing union of spatial repair patches.
Increase \(\beta_0\) so \(\lambda\ge8M\).
The inherited joint resolvent comparison
\[
                  H_\mu\succeq
                       L\{L+C_{\rm f}(\beta+1)\}^{-1}
\]
and its low-spectral projection argument give the first
floor in \eqref{eq:v11-HB} on the joint law. They assume
no unrestricted hidden diffusion gap. Exact
conditional-variance domination transfers it to the
interface marginal because \(G\subset\mathcal S\);
the lifted function imposes no hidden restriction.
Finally \(\beta+1\le2\beta\) for \(\beta\ge1\) gives
the second bound.
\end{proof}

\subsection{Verification, provenance, and the upstream task}

The exact checker is
\begin{center}\small
\path{scripts/verify_ym_f14_v11_paired_transition.py}.
\end{center}
Its SHA--256 is
\begin{center}\scriptsize\ttfamily
8B15D4C10267BE70026B547EC3CBDEF969255AD14D3853BA297A1C63D3098822.
\end{center}
It reruns V9's geometry and certificate; checks all
76 inverse pairs, the absence of direct companion
factors, and the 2,344-input closure; and performs
348 nonparallel mixed quaternion two-jet tests and
48 source-inversion two-jet tests. Rational arithmetic
checks the cross-term coefficient and the radial
failed-upper-bound diagnostic. Integer outward
trigonometric intervals check the explicit noncommuting
midpoint on all 19 edges. In particular its largest
squared cone-ratio upper bound is
\[
                 \frac{5202240468753029}
                           {1180591620717411303424}
                               <(3/64)^2.
\]
The general inversion, derivative, cone, and all-phase
inequalities are proved above. These finite regressions
are not a sampling proof of the universal assertions.

The predecessor V10 has 177,518 bytes and SHA--256
\begin{center}\scriptsize\ttfamily
B2C72556F85407CAD15F8D70008B546FDBE83C248E7757F0D03D2FBAD8534494.
\end{center}
Its full body is preserved, apart from the title version
and the terminal insertion. Active V11 replaces V10;
archives, the new shell-mixing file, monographs, and
other companion programmes are preserved. The regular
SM/NS checkpoint remains V15. The present additional
RG review responds to the newly supplied source and
does not postpone that checkpoint.

\begin{firewall}
All-phase coverage here means every radial phase
\emph{inside the inverse-paired coplanar family}, or
more generally the stated paired-cone test and a small
open tube around it. It does not mean arbitrary
independent transition sources. The direct visible
sheet pattern and coarse datum \(I\) are still part
of the hypotheses. The result is a localized
configuration-space heat-bath floor, not an
unrestricted patch gap, an RG integral estimate,
a continuum quantum construction, or a physical
Hamiltonian mass gap.
\end{firewall}

The new RG companion changes the next upstream
acquisition order: do not redo its local classical
continuum limit or its measured finite-packet algebra.
There remain two distinct quantitative tasks. On the
RG side, one must control source-dependent fluctuation
forms, their subtracted loop traces, the full measured
density packet, and transported remainders on a common
cofinal family. On the compact side, one must obtain
an all-field, all-exterior estimate or a proved
replacement that controls the complement of local
repair cylinders on that same law.
Neither a classical lower-limit inequality nor the
present negative action direction supplies that
missing comparison.

For this repair route, the next honest obstruction is
loss of inverse pairing: then the root can have
nonzero logarithmic acceleration and the cancellation
of Lemma~\ref{lem:v11-inversion} is unavailable.
Noncoplanar paired families can already be assessed
by \eqref{eq:v11-cone}; passing finitely many such
tests is not an all-field covering theorem.
Any attempt to pass from local certificates to global
closure must first state the uncovered event and
pay for its spatial overlap, interfaces, and actual
conditional law. The endpoint remains the quantum
existence and positive physical mass-gap problem
in the
\href{https://www.claymath.org/millennium/yang-mills-the-maths-gap/}
{official Clay formulation}, not the auxiliary
rate-one generator considered here.
\section{V12: the actual continuum constraint curvature and a fifth regularized determinant}
\label{sec:v12-fluctuation}

\subsection{Going upstream at the unpaired-source bottleneck}

V11 is concrete progress: it proves an all-phase repair theorem
under its inverse-pair and cone conditions. It does not solve
the independent-source complement. Removing pairing introduces
logarithmic pivot acceleration, and the existing implicit
gradient identities do not by themselves bound its pairing with
the physical action force. The full noncommuting adjoint was
already written in f12 V96, equations (96.19)--(96.23); rewriting
it as a new general repair theorem would not fill that gap.
No unconditional unpaired-source floor is asserted here.

Instead we follow the upstream opening in the newly supplied
Shell-Mixing V13. That file already constructs a local continuum
classical minimum and a finite-resolution approximation to its
literal nonlinear source. Its terminal statement explicitly
leaves the source-dependent fluctuation forms unestimated.
This version derives one such estimate for the \emph{actual
constraint-curvature contribution} to the continuum Hessian.
It is not another small neighbourhood of a symmetric repair
configuration.

The result is a relative trace-ideal theorem. On a sufficiently
small nonzero source ball, the constraint correction belongs
to \(\mathfrak S_5\); its finite-resolution approximants converge
in that norm at rate \(M^{-1/10}\). A fifth regularized
determinant and all its fixed source derivatives consequently
exist. The four omitted trace terms remain explicit. This does
not construct the full one-loop measure: the reference
fluctuation determinant, measured coarea density, matching of
the actual lattice fibres, and the compact complement are
separate obligations.

The shell-mixing snapshot is still the 512,757-byte V13 with
SHA--256 recorded in V11. We use its
\texttt{thm:v13-weaksource}, \texttt{thm:v13-cont-min}, and
\texttt{thm:v13-main}, including their preceding hypotheses.
The arguments for the weak-source estimate and the continuum
minimum were read directly. Their entire historical dependency
chain and unprovided companion scripts are not independently
recertified here. Searching the main YM archives, the closure
monograph, and Shell V13 found earlier finite trace inequalities
and references to an SM fifth determinant, but no corresponding
trace-ideal bound for this continuum root-constraint Hessian.
No SM determinant estimate is imported to this different operator.

\subsection{The inherited source estimate on its actual space}

All spaces in this section are on the one-cell torus
\(\mathbb T^4=(\mathbb R/\mathbb Z)^4\).
Use the quaternion Lie convention \([a,b]=2a\times b\).
Let
\[
 V=\{w\in H^1(\mathbb T^4;\mathfrak{su}(2)^4):
               \operatorname{div}w=0,\ \textstyle\int w=0\},
 \qquad \|w\|_V=\|dw\|_{L^2},
\]
and \(X=\mathbb R^{12}\oplus V\), where the first summand
is the constant connection sector. Give \(X\) its constant
Euclidean plus detail-energy Hilbert norm. The embedding
\(X\hookrightarrow L^4\) is bounded by Sobolev--Poincare.
The complex spaces and all symmetric forms below mean
bilinear holomorphic extensions of these real objects.

Write \(\mathcal O\) for Shell V13's actual limiting balanced
root observable, and
\[
 \mathcal P_M=\mathcal F_{M,1}\mathsf A_M,\qquad M=2^m,
\]
for precisely its finite-resolution approximation. In its
notation
\[
 (\mathsf A_M A)_{r,\mu}
   =\frac1M\int_0^1
       \operatorname{avg}_{\,r/M+[0,1/M)^4+t e_\mu/M}
                                      A_\mu(y)\,dt .
\]
Its input averages take values in a space of dimension
\begin{equation}\label{eq:v12-rank-budget}
                              d_M=12M^4.
\end{equation}
This is not an unspecified discretization of the continuum
constraint. It retains the ordered-tree and dyadic convention
of the supplied source. At \(M=1\), the remaining root map
has no steps, the box is the whole torus, and
\(\mathcal P_1 A=\int A\) is linear.

Fix
\begin{equation}\label{eq:v12-alpha}
                             \alpha=\frac9{10}.
\end{equation}
Shell V13 permits any fixed \(0<\alpha<1\) after shrinking
its ball. On common nested complex \(L^4\) balls, its estimate
and Cauchy's formula give
\begin{equation}\label{eq:v12-source-jets}
 \sup_A\|D^j(\mathcal O-\mathcal P_M)(A)\|_
                 {(L^4)^j\to\mathbb C^{12}}
                         \le C_jM^{-\alpha},
                       \qquad j=0,1,2,\ldots .
\end{equation}
For each fixed \(j\), the supremum is on a smaller ball
with a positive margin to the original complex domain.
Indeed the uniform zeroth bound on that larger ball applies
to \(A+\sum z_iU_i\), and the multivariate Cauchy estimate
gives the displayed derivative bound. The radius and
constants may depend on \(\alpha,j\), not on \(M\).
In particular taking \(\alpha=9/10\) is legitimate on one
fixed sufficiently small nonzero ball. No estimate at the
endpoint \(\alpha=1\) is used.

The actual continuum action is
\[
 \mathcal S_\infty(A)=\frac12\sum_{\mu<\nu}
   \int|\partial_\mu A_\nu-\partial_\nu A_\mu+[A_\mu,A_\nu]|^2 .
\]
Shell V13 supplies the chart
\[
 A(v,w)=s(v,w)+w,\qquad \mathcal O(A(v,w))=v,
\]
and its unique small detail minimum \(w_*(v)\).
Put \(A_*(v)=A(v,w_*(v))\) and
\(\mathcal G(v)=\mathcal S_\infty(A_*(v))\).
They are holomorphic on a common smaller source ball,
with \(A_*(v)=O(|v|)\) in \(X\) and
\(\|w_*(v)\|_V=O(|v|^2)\).
We retain exactly this local branch, not a global compact
minimum.

\subsection{The constrained Hessian and its multiplier term}

At the point \((v,w_*(v))\), define
\[
 \mathcal T_vU=U+D_ws(v,w_*)U,\qquad
                      Z_v=D_s\mathcal O(A_*).
\]
Here \(D_ws\,U\) is constant in space. The matrix \(Z_v\)
is invertible and uniformly close to \(I\). The map
\(\mathcal T_v:V\to X\) is bounded and
\(D\mathcal O(A_*)\mathcal T_v=0\).
We use the multiplier convention
\[
                 \zeta(v)=D_v\mathcal G(v)
\]
as a Euclidean covector. It is the negative of a multiplier
defined with \(D_v\mathcal G=-\eta\); signs are fixed here
by the following identity.

\begin{lemma}[Actual multiplier and constraint-curvature identity]
\label{lem:v12-hessian}
On this branch,
\[
 D\mathcal S_\infty(A_*)=\zeta(v)\,D\mathcal O(A_*)
                                      \quad\hbox{on }X.
\]
The bounded symmetric detail-energy Hessian is
\begin{equation}\label{eq:v12-H}
 \begin{split}
 \langle U,\mathcal H_vW\rangle_V
   ={}&D^2\mathcal S_\infty(A_*)[\mathcal T_vU,\mathcal T_vW]\\
      &-\zeta(v)\cdot
                    D^2\mathcal O(A_*)[\mathcal T_vU,\mathcal T_vW].
 \end{split}
\end{equation}
Moreover \(|\zeta(v)|\le C|v|^3\).
\end{lemma}

\begin{proof}
The detail stationary equation says that
\(D\mathcal S_\infty(A_*)\) vanishes on the range of
\(\mathcal T_v\). Any \(\xi\in X\) decomposes uniquely
as \(\mathcal T_vU+q\), with \(q\) constant; its source
derivative is \(Z_vq\). Thus
\[
 D\mathcal S_\infty(A_*)=\zeta_0D\mathcal O(A_*),
             \qquad \zeta_0=D_s\mathcal S_\infty(A_*)Z_v^{-1}.
\]
Differentiating the constraint along \(A_*(v)\) gives
\(D\mathcal O(A_*)D_vA_*=I\), so the chain rule for
\(\mathcal G\) identifies \(\zeta_0=D_v\mathcal G\).

Twice differentiate the chart constraint in detail
directions:
\[
 Z_vD_w^2s[U,W]+
                  D^2\mathcal O(A_*)[\mathcal T_vU,\mathcal T_vW]=0 .
\]
The second chain rule for the chart action has the first
term in \eqref{eq:v12-H} and the additional term
\(D_s\mathcal S_\infty(A_*)D_w^2s[U,W]\).
Use the preceding two identities to obtain its indicated
minus sign. The second derivative of the nonlinear
constraint is therefore retained, not frozen.

Finally \(dA_*=dw_*=O(|v|^2)\) in \(L^2\) and
\(A_*=O(|v|)\) in \(L^4\). Its curvature
\(F_*=dA_*+[A_*,A_*]\) is \(O(|v|^2)\) in \(L^2\).
For a constant variation \(q\), the curvature variation
is a sum of brackets of \(q\) with \(A_*\), bounded in
\(L^2\) by \(C|q||v|\) on the unit torus. Therefore
\(|D_s\mathcal S_\infty(A_*)q|\le C|v|^3|q|\).
The bounded \(Z_v^{-1}\) proves the multiplier estimate,
also on a smaller complex ball.
\end{proof}

Let \(\mathcal H_{0,v}\) and \(\mathcal B_v\) be the
Riesz operators of the first and second terms respectively:
\begin{equation}\label{eq:v12-split}
                       \mathcal H_v=\mathcal H_{0,v}+\mathcal B_v .
\end{equation}
The finite-dimensional target \(\mathbb R^{12}\) of
\(\mathcal O\) does \emph{not} make \(\mathcal B_v\)
finite rank. A scalar-valued quadratic functional can
have an infinite-rank Hessian. The next estimate uses
the actual finite-resolution approximation, not the
number of output constraints alone.

\subsection{Compactness of the relative continuum fluctuation form}

\begin{proposition}[A kinetic-normalized compact perturbation]
\label{prop:v12-compact}
After shrinking the source ball,
\[
             \|\mathcal H_{0,v}-I\|<\frac14,\qquad
                         \mathcal H_{0,v}-I\text{ is compact}.
\]
The compactness assertion holds for each \(v\), and the
operator family is holomorphic. The full
\(\mathcal H_v-I\) is also compact. These are operators
in the detail-energy Hilbert space, not physical
Hamiltonians.
\end{proposition}

\begin{proof}
For \(a,b\in X\), put
\[
 (J_Aa)_{\mu\nu}=[a_\mu,A_\nu]+[A_\mu,a_\nu],\qquad
 C(a,b)_{\mu\nu}=[a_\mu,b_\nu]+[b_\mu,a_\nu].
\]
Direct differentiation of the full curvature square gives
\[
 D^2\mathcal S_\infty(A)[a,b]
   =\langle da+J_Aa,db+J_Ab\rangle
                             +\langle F_A,C(a,b)\rangle .
\]
Since \(d\mathcal T_vU=dU\), its kinetic contribution
on \(V\) is exactly \(I\). Holder and Sobolev bound the
remaining form by
\[
 C\{\|A_*\|_4+\|A_*\|_4^2+\|F_*\|_2\}
                 \|\mathcal T_v\|_{V\to X}^2\|U\|_V\|W\|_V.
\]
It is \(O(|v|)\), proving the norm bound on a smaller ball.

For compactness, \(J_{A_*}\mathcal T_v:V\to L^2\)
is compact. Approximate \(A_*\in L^4\) by bounded
coefficient arrays. For a bounded array this map is
compact by the Rellich embedding \(H^1\hookrightarrow L^2\)
on the fixed torus, with the constant part of \(\mathcal T_v\)
a finite-rank addition. The error in operator norm is
bounded by the \(L^4\) coefficient error times the
bounded embedding \(V\to L^4\).
Consequently its cross terms with the bounded derivative
map and its squared term give compact Riesz operators.
For the last term approximate \(F_*\in L^2\) by bounded
arrays. The bounded-coefficient form factors through
the same compact \(L^2\) embedding. Its error is at most
\(C\|F_*-F_*^{(n)}\|_2\|U\|_V\|W\|_V\), by two
\(L^4\) bounds. Thus it too is a norm limit of compact
operators. All formulas are holomorphic bounded
multilinear expressions on the common complex domain.

Compactness of \(\mathcal B_v\), and hence of the full
perturbation, follows from the finite-rank approximants
in the next theorem. This proof of reference compactness
does not claim a trace-class estimate for
\(\mathcal H_{0,v}-I\).
\end{proof}

For real \(v\), \(\mathcal H_{0,v}\) is positive, so
its inverse square root is bounded. On the complex
ball it is defined by the convergent binomial power
series about \(I\). Define the normalized constraint
correction
\begin{equation}\label{eq:v12-K}
             \mathcal K_v=
                   \mathcal H_{0,v}^{-1/2}\mathcal B_v
                                      \mathcal H_{0,v}^{-1/2}.
\end{equation}
For real \(v\) it is self-adjoint and
\(\mathcal H_v=\mathcal H_{0,v}^{1/2}
(I+\mathcal K_v)\mathcal H_{0,v}^{1/2}\).
The square roots and transposes in the complex continuation
are analytic continuations of the real formulas, not
antiholomorphic adjoints depending on \(\bar v\).

\subsection{A quantified fifth Schatten-class source correction}

Define \(\mathcal B_{v,M}\) by replacing
\(D^2\mathcal O(A_*)\) in \(\mathcal B_v\) with
\(D^2\mathcal P_M(A_*)\), while retaining the
\emph{same} \(A_*,\zeta,\mathcal T_v,\mathcal H_{0,v}\).
Let
\[
 \mathcal K_{v,M}=\mathcal H_{0,v}^{-1/2}
                  \mathcal B_{v,M}\mathcal H_{0,v}^{-1/2}.
\]
These are continuum finite-resolution source-curvature
approximants. They are not defined to be the Hessians of
the lattice minimizers \(A_L^*\).

\begin{theorem}[Trace-ideal control from the literal source tower]
\label{thm:v12-S5}
Uniformly on a sufficiently small complex source ball,
\begin{align}
 \operatorname{rank}\mathcal K_{v,M}&\le12M^4,
       &\|\mathcal K_v-\mathcal K_{v,M}\|
                      &\le C|\zeta(v)|M^{-9/10},
                                      \label{eq:v12-op-tail}\\
 \mathcal K_v&\in\mathfrak S_5,
       &\|\mathcal K_v\|_{\mathfrak S_5}
                      &\le C_5|\zeta(v)|\le C'_5|v|^3,
                                      \label{eq:v12-S5}\\
 &&\|\mathcal K_v-\mathcal K_{v,M}\|_{\mathfrak S_5}
                      &\le C_5|\zeta(v)|M^{-1/10}.
                                      \label{eq:v12-S5-tail}
\end{align}
The family and its limit are holomorphic with values
in \(\mathfrak S_5\). On every smaller source ball
all fixed source derivatives of the difference in
\eqref{eq:v12-S5-tail} are \(O(M^{-1/10})\).
The constants and nonzero radius are existential.
\end{theorem}

\begin{proof}
Since \(\mathsf A_M\) is linear,
\[
 D^2\mathcal P_M(A)[a,b]
       =D^2\mathcal F_{M,1}(\mathsf A_MA)
                              [\mathsf A_Ma,\mathsf A_Mb].
\]
Thus the bilinear form of \(\mathcal K_{v,M}\) factors
through
\[
         \mathsf A_M\mathcal T_v\mathcal H_{0,v}^{-1/2}
                       :V\longrightarrow\mathbb C^{12M^4}.
\]
Its Riesz operator has rank at most \(12M^4\).
Equation~\eqref{eq:v12-source-jets} at \(j=2\)
and the uniform bounds for the two outside maps give
\eqref{eq:v12-op-tail}. In particular \(\mathcal K_v\)
is compact.

Here is the trace-ideal step with the rank loss retained.
Put \(K_m=\mathcal K_{v,2^m}\) and
\(\Delta_m=K_{m+1}-K_m\).
Its rank is at most
\[
                    12(2^{4m}+2^{4(m+1)})=204\,2^{4m},
\]
and its operator norm is at most
\[
          C|\zeta|(1+2^{-\alpha})2^{-\alpha m}.
\]
For a rank-\(r\) operator,
\(\|T\|_{\mathfrak S_p}\le r^{1/p}\|T\|\);
this follows immediately by bounding each of its at most
\(r\) nonzero singular values by \(\|T\|\).
For any \(p\) with \(\alpha>4/p\),
\[
 \|\Delta_m\|_{\mathfrak S_p}
 \le204^{1/p}C|\zeta|(1+2^{-\alpha})
                                   2^{-(\alpha-4/p)m}.
\]
These norms are summable. Since \(\mathcal P_1\) is
linear, \(K_0=0\). Therefore
\[
 \|\mathcal K_v-K_m\|_{\mathfrak S_p}
 \le
 \frac{204^{1/p}C(1+2^{-\alpha})}
      {1-2^{-(\alpha-4/p)}}|\zeta|\,
                                   2^{-(\alpha-4/p)m}.
\]
The sum converges also in operator norm to the
\(\mathcal K_v\) already identified, so it is not an
alternative candidate limit. With \(p=5,\alpha=9/10\),
\(\alpha-4/p=1/10\), proving the assertions.

Every finite-rank term depends holomorphically on \(v\)
as a bounded operator. A locally uniform rank bound
upgrades this to holomorphy in \(\mathfrak S_p\):
difference quotients have uniformly bounded finite
rank, so their operator-norm convergence gives
\(\mathfrak S_p\) convergence. The uniform summable
series above proves \(\mathfrak S_5\) holomorphy of
the limit. Applying Cauchy's formula on a larger
source polydisc to its uniform tail gives the
asserted fixed derivative estimates.
\end{proof}

More generally, for each fixed \(p>4\), one may first
choose \(4/p<\alpha<1\) and then shrink to the
corresponding source ball. The same proof gives
\(\mathfrak S_p\) there. It does not prove all
\(p>4\) estimates on one fixed ball by taking
\(\alpha\uparrow1\); that would ignore the
radius dependence. Nor does the argument give
\(\mathfrak S_4\) at a nonzero radius.

\subsection{The relative determinant and the four terms not removed}

For \(K\in\mathfrak S_5\) with \(\|K\|<1\), define
\begin{equation}\label{eq:v12-det5}
       \log\det\nolimits_5(I+K)
            =\sum_{j=5}^{\infty}\frac{(-1)^{j+1}}{j}\Tr(K^j).
\end{equation}
This is the standard fifth regularized Fredholm
determinant, with the logarithm branch fixed at \(K=0\).
For context and the general, nonmultiplicative product
law see N.~Koutsonikos-Kouloumpis and M.~Lesch,
\href{https://arxiv.org/abs/2202.12923}
{\emph{The product formula for regularized Fredholm
determinants: two new proofs}}, particularly its
definitions (1.3)--(1.4). The existence and error
estimate needed here follow directly from the series;
no determinant theorem for a different YM operator
is imported. The numeral five denotes the subtraction
order, not five physical loops.

\begin{lemma}[A local trace-ideal determinant estimate]
\label{lem:v12-det}
If \(K,L\in\mathfrak S_p\), \(p\) a positive integer,
\(\|K\|,\|L\|\le r<1\), and
\(\|K\|_{\mathfrak S_p},\|L\|_{\mathfrak S_p}\le B\),
then the corresponding series beginning at \(p\)
is absolutely convergent, and
\[
 |\log\det\nolimits_p(I+K)|
          \le\frac{\|K\|_{\mathfrak S_p}^p}{p(1-r)},
\]
\[
 |\log\det\nolimits_p(I+K)-\log\det\nolimits_p(I+L)|
          \le\frac{B^{p-1}}{1-r}\|K-L\|_{\mathfrak S_p}.
\]
These estimates apply also to complex, nonself-adjoint
operators.
\end{lemma}

\begin{proof}
Schatten Holder gives
\(\|K^j\|_1\le\|K\|_{\mathfrak S_p}^p r^{j-p}\)
for \(j\ge p\), proving convergence and the first bound.
For the second, use the noncommuting telescoping identity
\[
 K^j-L^j=\sum_{i=0}^{j-1}K^i(K-L)L^{j-1-i}.
\]
In each summand put \(K-L\) and any \(p-1\) of the
remaining factors in \(\mathfrak S_p\), leaving
the other \(j-p\) factors in operator norm.
Holder bounds its trace by
\(B^{p-1}\|K-L\|_{\mathfrak S_p}r^{j-p}\).
There are \(j\) summands; this cancels the denominator
\(j\) in the logarithm series. Sum the resulting
geometric series. No factors have been commuted.
\end{proof}

\begin{theorem}[A controlled actual-source constraint determinant]
\label{thm:v12-det}
On a sufficiently small common source ball the function
\begin{equation}\label{eq:v12-theta}
                 \Theta(v)=\frac12\log\det\nolimits_5(I+\mathcal K_v)
\end{equation}
is well defined and holomorphic. Its finite-resolution
values \(\Theta_M\), formed with \(\mathcal K_{v,M}\),
satisfy
\[
                 |\Theta(v)|\le C|v|^{15},\qquad
       |\Theta(v)-\Theta_M(v)|\le C|v|^{15}M^{-1/10}.
\]
All fixed source derivatives converge locally at the
same power \(M^{-1/10}\). For real \(v\),
\(\det_5(I+\mathcal K_v)>0\).
\end{theorem}

\begin{proof}
The preceding theorem bounds all operator norms and
\(\mathfrak S_5\) norms by \(C|v|^3\). Shrink the
ball so all operator norms are at most \(1/2\);
the finite-rank approximants obey the same bound
after choosing a uniform constant. Lemma~\ref{lem:v12-det}
with \(p=5\) gives the two assertions and the order
\(|v|^{15}\). Uniform trace-norm convergence of the
series on the complex domain proves holomorphy;
Cauchy's formula gives the derivative rate on
smaller balls. For real \(v\) the operator is
self-adjoint with \(I+\mathcal K_v>0\), its series
is real, and exponentiating gives strict positivity.
\end{proof}

The meaning of this limit is fixed by an exact
finite-rank identity. For real \(v\), use the
ordinary determinant on the range of
\(\mathcal K_{v,M}\), with identity on its
orthogonal complement. For complex \(v\), use
the finite-rank Fredholm determinant, equivalently
the holomorphic continuation of that expression.
Then
\begin{equation}\label{eq:v12-four-traces}
 \frac12\log\det(I+\mathcal K_{v,M})
  =\Theta_M(v)+\frac12\sum_{j=1}^{4}
                    \frac{(-1)^{j+1}}{j}\Tr(\mathcal K_{v,M}^j).
\end{equation}
For real \(v\) this is also the finite-rank relative
determinant of \(\mathcal H_{0,v}+\mathcal B_{v,M}\)
against \(\mathcal H_{0,v}\): conjugation by its
square root yields exactly \(I+\mathcal K_{v,M}\).
It does not require a separately defined ordinary
determinant of either infinite-dimensional Hessian.

Thus the new estimate controls the trace-ideal
remainder, not the four terms on the right of
\eqref{eq:v12-four-traces}. The latter can include
the low source degrees relevant to the measured
packet; the bound \(O(|v|^{15})\) for \(\Theta\)
is not a reason to discard them. Calling them
\emph{subtractions} here is an algebraic convention,
not a proof that they are gauge-invariant local
counterterms or that their physical coefficients
are known.

\subsection{Why the rank and norm estimate cannot justify those traces}

\begin{proposition}[A sharp trace-threshold diagnostic]
\label{prop:v12-sharp}
The rank bound \(O(M^4)\) and operator error
\(O(M^{-9/10})\), even with a positive uniformly
coercive \(I+K\), do not imply finiteness or convergence
of any of the four traces omitted in
\eqref{eq:v12-four-traces}.
\end{proposition}

\begin{proof}
On a separable Hilbert space take a positive diagonal
operator with eigenvalue
\[
          \lambda_j=\frac14\,2^{-9j}
       \quad\hbox{of multiplicity}\quad
          n_j=(2^{40}-1)2^{40(j-1)},\qquad j=1,2,\ldots .
\]
Its operator norm is \(1/2048\), so
\(I\preceq I+K\preceq(1+1/2048)I\).
For a dyadic \(M\ge1\), retain the first
\(m=\lfloor(\log_2M)/10\rfloor\) blocks.
The resulting rank is \(2^{40m}-1<M^4\), and
the omitted operator norm is
\[
                  \frac14\,2^{-9(m+1)}
                                    \le\frac14 M^{-9/10}.
\]
Thus it has precisely the required rank and
operator-norm approximation properties.

For any \(q>0\), the sum of the \(q\)-th powers
of its eigenvalues is a positive constant times
\(\sum_{j\ge1}2^{(40-9q)j}\). It converges exactly
when \(q>40/9\). In particular \(K\in\mathfrak S_5\),
but each \(\Tr K^q\), \(q=1,2,3,4\), is infinite.
The exact fifth trace is
\[
               \Tr K^5=\frac{35468117025}{1125899906842624}.
\]
The finite ordinary log determinants diverge as
well, since \(\log(1+x)\ge x/2\) for
\(0\le x\le1/2\) and the first trace diverges.
The fifth regularized determinants still converge
by their absolutely summable fifth-order remainders.
\end{proof}

This is a counterexample to an inference from the
rank/error data, not an identification of the
actual YM constraint spectrum. It neither proves
that the four actual traces diverge nor rules out
additional cancellations in them. It shows why
their acquisition cannot be replaced by compactness,
coercivity, or convergence of the regularized tail.

\subsection{Exact checks, version record, and the next actual comparison}

The dependency-free rational checker is
\begin{center}\small
\path{scripts/verify_ym_f14_v12_source_curvature.py}.
\end{center}
Its SHA--256 is
\begin{center}\scriptsize\ttfamily
08EA9DD01236B554A1C89D2142CF0F071C20B508E7DDABAD7EA0E0DDFF8894D5.
\end{center}
It independently differentiates a nonlinear rational
constraint chart through order two and compares with
the multiplier Hessian formula. The reference and
correction matrices in that test do not commute.
A nonorthogonal coordinate change transforms both
forms; their relative determinant and traces through
order twelve agree exactly. The actual finite
relative determinant in this test is
\[
                     \frac{54220113096310103}{54816438485673527}.
\]
The checker also verifies twenty finite-rank
factorizations, the dyadic rank factor 204, the
exponent \(9/10-4/5=1/10\), and the removal of the
first four trace coefficients by directly expanding
the determinant polynomial. Sixteen exact blocks
of the sharpness example verify its geometric trace
sums. A scalar fifth-logarithm has a strictly
positive rational enclosure.

These are algebraic consistency checks. They do
not replace the analytic proofs of compactness,
trace-ideal convergence, or the imported classical
source theorems with numerical sampling. The
initial floating exploration of unrestricted
response directions supplies no certificate used
in this version.

The predecessor V11 has 205,672 bytes and SHA--256
\begin{center}\scriptsize\ttfamily
03C85CF50485272C04866B32B37EE8BD7BBB3E5B361A1E2C16E6D3163A433E8F.
\end{center}
Its body is retained in full, apart from the
title version and terminal insertion. Active
V12 replaces V11. The separate Shell-Mixing V13,
SM and NS companions, monographs, and archives
are preserved. The next regular companion review
is V15.

\begin{firewall}
The continuum fluctuation form is classical and
normalized by its detail kinetic energy.
Its positivity is not a physical quantum mass
gap; fixing the constant source has already
removed a sector from this Hilbert space.
The new fifth determinant is only the relative
constraint-curvature contribution on one small
one-cell branch. It is not the full constrained
Gaussian determinant, the coarea density, an
arbitrary-source or thermodynamic construction,
or a same-law comparison of lattice partition
functions. No improvement of V11's compact
repair coverage is claimed by this upstream move.
\end{firewall}

The next substantive target is now more specific
than a generic request for a fluctuation bound.
For the actual lattice constrained Hessians,
one must identify the energy-normalized
constraint correction on common spaces and
compare it in a sufficient trace ideal to
\(\mathcal K_v\), or retain a direct trace
comparison without such an identification.
The present \(M\)-approximants alter the
continuum source derivative at the \emph{fixed}
continuum minimum; they do not silently replace
this varying-lattice task.
In parallel, the four lower trace terms of
\eqref{eq:v12-four-traces}, the reference
fluctuation operator, and the exact measured
density must be matched with their ultraviolet
subtractions and full shell ownership.
The fifth-tail theorem gives no license to
reset any of those factors. A cofinal quantum
law and a calibrated positive physical gap
remain unproved.
\section{V13: matching the actual lattice constraint Hessians in a trace ideal}
\label{sec:v13-lattice}

\subsection{The next comparison, not another continuum approximation}

V12 constructs a fifth regularized determinant for the
continuum source-constraint correction. Its finite-resolution
approximants were evaluated at the fixed continuum minimum.
The present version instead starts with the \emph{actual
lattice minima, actual nested root constraint, and actual
Wilson action} of Shell-Mixing V12--V13. It compares their
constraint corrections to the continuum operator in
\(\mathfrak S_5\). The distinction is essential: convergence
of classical minimum values alone did not imply a determinant
comparison.

The mechanism has two parts. The source derivative has
uniform low-rank approximations even before taking the
lattice limit. This upgrades its operator-norm comparison
to a trace-ideal comparison. The reference Hessian needs
only strong operator convergence, since it sandwiches a
trace-ideal correction. A common energy space makes both
statements precise, including real Nyquist modes.
No projection boundedness at the endpoint \(L^4\) is assumed.

The main V12 and the separate Shell-Mixing V13 were rechecked
at the start of the iteration. The latter still has
512,757 bytes and SHA--256
\begin{center}\scriptsize\ttfamily
7F2DBD2900D7A78280DA980152266E94DF8A95EFE15CD37429EC8C6AAAA4F17B.
\end{center}
We use its discrete Hodge and Sobolev estimates, primitive
Wilson analytic bounds, common classical branch, finite-step
source truncation, and strong classical continuum limit.
Their relevant proofs and scope were inspected directly.
The new comparison below is not an independent recertification
of all their earlier dependencies. Shell V13 explicitly
did not assert fluctuation-operator or determinant convergence;
main V12 explicitly left the varying-lattice comparison open.

All statements remain on one final coarse periodic cell,
in the declared small-source branch and the same dyadic
balanced-root convention. They do not concern arbitrary
compact exterior data or growing physical volume.

\subsection{The finite and continuum energy spaces}

Let \(L=2^n\), let \(V_L\) be the mean-zero discrete
Coulomb link space on the fine torus of side \(L\), and
put
\[
 \|u\|_{V_L}=\|du\|_{\ell^2},\qquad
 (\mathcal I_Lu)(y)=Lu(x),\quad
                    y\in x/L+[0,1/L)^4 .
\]
The continuum detail space \(V\) and
\(X=\mathbb R^{12}\oplus V\) are those of V12.
Write \(D^L\) for physical forward differences, so
\[
 \|\mathcal I_Lu\|_4=\|u\|_{\ell^4},\qquad
        \|D^L\mathcal I_Lu\|_2=\|\nabla u\|_{\ell^2}
                                      =\|u\|_{V_L}.
\]
The last equality is on the Coulomb space, with the
sum over all first differences understood.

\begin{lemma}[Real energy isometries with the correct low modes]
\label{lem:v13-isometry}
There are real isometries \(J_L:V_L\to V\).
For \(P_L=J_LJ_L^*\) and
\(\mathcal R_L=\mathcal I_LJ_L^*:V\to L^4\),
\[
       P_L\longrightarrow I\ \hbox{strongly on }V,\qquad
       \|\mathcal R_L\|_{V\to L^4}\le C_S .
\]
If \(u_L\rightharpoonup u\) in \(V\), then
\begin{equation}\label{eq:v13-weak-reconstruction}
                  \mathcal R_Lu_L\rightharpoonup u
                                             \quad\hbox{in }L^4.
\end{equation}
For every fixed bounded finite-rank map
\(A:L^4\to\mathbb C^d\),
\begin{equation}\label{eq:v13-dual-average}
          \|A(\mathcal R_L-\iota)\|_{V\to\mathbb C^d}\longrightarrow0,
\end{equation}
where \(\iota:V\hookrightarrow L^4\) is the inclusion.
The same assertions hold after complexification.
\end{lemma}

\begin{proof}
Here is a mode-by-mode construction that specifies the
needed normalization. At nonzero lattice frequency
\(k/L\), choose a representative with \(|k_\mu|\le\pi L\)
and put
\[
 q_L(k)_\mu=e^{ik_\mu/L}-1,\qquad
                 \omega_L(k)=L|q_L(k)|.
\]
The lattice transverse plane is \(q_L(k)^\perp\);
the continuum plane is \(k^\perp\). Choose a unitary
map \(U_L(k)\) between these three-dimensional planes.
For every fixed \(k\), choose the polar map between
their orthogonal projections once they are close.
It tends to the identity, since
\(q_L(k)/(i|q_L(k)|)\to k/|k|\).
At the finitely many other frequencies any unitary
identification suffices. Make conjugate choices on
conjugate frequency pairs.

In an orthonormal real Fourier basis, send the physical
lattice coefficient \(a\) to the continuum coefficient
\[
                         \frac{\omega_L(k)}{|k|}U_L(k)a .
\]
Its continuum energy is \(\omega_L(k)^2|a|^2\),
exactly the lattice energy by the Hodge identity
and Parseval. At the self-conjugate nonzero Nyquist
frequencies, the lattice coefficients and transverse
planes are real. Map these orthogonally to real
cosine modes in the corresponding continuum planes,
with the same energy normalization. Use normalized
cosines, not an unnormalized pair of exponentials.
Distinct lattice conjugacy orbits can be assigned
distinct continuum frequency pairs, so their images
are orthogonal.

For every even \(L\), there are 15 nonzero
self-conjugate frequency orbits and
\((L^4-16)/2\) other conjugate pairs. Including the
three Lie components, their real transverse
dimensions sum to
\[
                 9\cdot15+18(L^4-16)/2=9(L^4-1)=\dim V_L.
\]
Thus no real Nyquist degree of freedom has been
dropped. Every fixed transverse trigonometric
polynomial lies in the image for sufficiently
large \(L\), proving \(P_L\to I\) strongly.
The uniform Sobolev bound and the isometry give
the asserted \(V\to L^4\) bound.

For a fixed continuum Fourier coefficient, integration
over fine cells multiplies the corresponding discrete
coefficient by a factor tending to one. The map
\(U_L(k)\) and the ratio \(\omega_L(k)/|k|\)
also tend to identity and one at that fixed
frequency. Consequently, for any bounded \(u_L\)
weakly converging in \(V\), every fixed Fourier
coefficient of \(\mathcal R_Lu_L\) tends to that
of \(u\). The uniform \(L^4\) bound and reflexivity
give weakly convergent subsequences; the Fourier
coefficients identify all their limits with \(u\).
This proves \eqref{eq:v13-weak-reconstruction}.

If \eqref{eq:v13-dual-average} failed, take unit
vectors \(u_L\) along a failing subsequence and
then a weakly convergent subsequence in \(V\).
Both \(A\mathcal R_Lu_L\) and \(A\iota u_L\)
converge to \(Au\) in the fixed finite-dimensional
target, contradicting failure. This argument
uses only bounded finite averages on \(L^4\);
it does not assert that the Fourier projections
are uniformly bounded on arbitrary \(L^4\) fields.
\end{proof}

The high-frequency identifications need not be
nested or canonical. The proof below works with
any such family agreeing with identity on fixed
low modes asymptotically. The constants come
from energy isometry and discrete Sobolev, not
from a high-frequency coordinate fit.

\subsection{The actual source map on one common domain}

For \((s,u)\in X\), define
\[
 \widetilde{\mathcal O}_L(s,u)
       =\mathcal F_{L,1}(s/L+J_L^*u),\qquad
                    \mathcal O_\infty(s,u)=\mathcal O(s+u).
\]
These are holomorphic on a common small complex
\(X\) ball. Indeed the fine counting-\(\ell^4\)
norm is at most \(|s|+C_S\|u\|_V\).
For dyadic \(M\le L\), define the \emph{actual
stopped-source map}
\begin{equation}\label{eq:v13-stopped}
 \widetilde{\mathcal P}_{L,M}(s,u)
      =\mathcal F_{M,1}\bigl(B_{L/M}(s/L+J_L^*u)\bigr).
\end{equation}
Here \(B_{L/M}\) is the actual composed linear
average, not a new sampling choice. For \(M>L\)
set \(\widetilde{\mathcal P}_{L,M}=
\widetilde{\mathcal O}_L\). At \(M=1\),
the stopped map is \(s\), since the detail has
zero mean.

\begin{theorem}[Uniform actual-source jets and finite-rank approximants]
\label{thm:v13-source}
Fix \(\alpha=9/10\) and shrink the source domain
as in V12. On common nested complex \(X\) balls,
for every fixed \(j\),
\begin{equation}\label{eq:v13-stop-jets}
 \sup_L\sup_x
  \|D^j(\widetilde{\mathcal O}_L-
                    \widetilde{\mathcal P}_{L,M})(x)\|
                                 \le C_jM^{-\alpha}.
\end{equation}
The stopped maps factor through at most \(12M^4\)
coordinates. Moreover
\begin{equation}\label{eq:v13-source-convergence}
       D^j\widetilde{\mathcal O}_L
                       \longrightarrow D^j\mathcal O_\infty
\end{equation}
uniformly in multilinear operator norm on every
smaller such ball. No rate in \(L\) is claimed.
\end{theorem}

\begin{proof}
Shell V13's \texttt{thm:v13-stop}, applied to the
actual fine input \(s/L+J_L^*u\), gives the
zeroth estimate in \eqref{eq:v13-stop-jets}.
It is uniform on the common complex ball.
Multivariate Cauchy gives every fixed derivative
on a smaller ball. The factorization in
\eqref{eq:v13-stopped} has \(12M^4\) real
coordinates; if \(M>L\), the whole fine input
has only \(12L^4\le12M^4\) coordinates.

For convergence keep \(M\) fixed first.
The fine input is exactly
\(\mathsf S_L(s+\mathcal R_Lu)\) in the
companion's cell-average convention. Its
operator-norm quadrature estimate therefore gives
\[
 \|B_{L/M}(s/L+J_L^*u)-
                        \mathsf A_M(s+\mathcal R_Lu)\|_{\ell^4}
                         \le C(|s|+\|u\|_V)M/L.
\]
By Lemma~\ref{lem:v13-isometry},
\(\mathsf A_M\mathcal R_L\to\mathsf A_M\iota\)
in operator norm for fixed \(M\).
The bounded analytic map \(\mathcal F_{M,1}\)
is Lipschitz on the admitted smaller domain.
Together with the continuum finite-resolution
estimate this yields, on a radius-\(r\) ball,
\[
 \|\widetilde{\mathcal O}_L-\mathcal O_\infty\|
   \le C r^2M^{-\alpha}
       +C r\{M/L+\varepsilon_L(M)\},
 \qquad
 \varepsilon_L(M)=
   \|\mathsf A_M(\mathcal R_L-\iota)\|_{V\to\ell^4}\to0 .
\]
The constant absorbs the two source truncation
errors. First choose \(M\) large, then let
\(L\to\infty\) with \(M\) fixed. This proves
uniform zeroth convergence on the ball.
Cauchy's formula on nested complex balls
proves \eqref{eq:v13-source-convergence}.
The qualitative finite-average comparison
does not provide a uniform mesh rate.
\end{proof}

\subsection{The embedded actual minima and source-curvature forms}

Use the lattice branch in Shell V12--V13:
\[
 a_L^*(v)=s_L(v)/L+w_L(v),\qquad
                  \mathcal F_{L,1}(a_L^*(v))=v.
\]
Let \(x_L(v)=(s_L(v),J_Lw_L(v))\in X\).
For real small \(v\), its proved classical limit gives
\begin{equation}\label{eq:v13-classical-inputs}
 \begin{aligned}
 \mathcal I_La_L^*&\to A_* &&\hbox{in }L^4,\\
 D^L\mathcal I_La_L^*&\to\nabla A_* &&\hbox{in }L^2,\\
 \mathcal C_L(a_L^*)&\to F_* &&\hbox{in }L^2,
 \end{aligned}
\end{equation}
where on each plaquette cell
\(\mathcal C_L(a)=L^2(Q_p(a)-I)\).
The last convergence follows from the companion's
weak curvature limit and convergence of its exact
squared norms. These are inherited inputs,
not consequences of the determinant sought here.

\begin{lemma}[Common-space minimum and chart convergence]
\label{lem:v13-branch}
For every such real \(v\),
\[
          x_L(v)\to (s_*(v),w_*(v))\quad\hbox{strongly in }X .
\]
Extend the lattice source chart to \(X\) by solving
\(\widetilde{\mathcal O}_L(s,u)=v\), and let
\(\ell_L=D_us\) at its actual embedded minimum.
Then, in operator norm \(V\to\mathbb R^{12}\),
\[
                       \ell_L\to\ell_\infty.
\]
The multipliers \(\zeta_L=D_v\mathcal G_L\)
converge to \(\zeta=D_v\mathcal G\), and uniformly
on smaller complex balls satisfy
\(|\zeta_L(v)|\le C|v|^3\).
\end{lemma}

\begin{proof}
The embedded details are bounded in \(V\).
Every weak subsequential limit is \(w_*\), by
\eqref{eq:v13-weak-reconstruction} and the
inherited strong \(L^4\) limit. Their \(V\)
norms equal the discrete curl norms, which
converge to \(\|dw_*\|_2\) by
\eqref{eq:v13-classical-inputs}. Weak convergence
and convergence of norms in a Hilbert space
give strong convergence. The means converge
by the \(L^4\) convergence on the unit torus.

The chart derivative is exactly
\[
       \ell_L=-(D_s\widetilde{\mathcal O}_L(x_L))^{-1}
                          D_u\widetilde{\mathcal O}_L(x_L).
\]
The inverse is uniformly bounded near identity.
Theorem~\ref{thm:v13-source} and strong convergence
of \(x_L\) prove the operator-norm convergence.
Scalar source-derivative convergence of
\(\mathcal G_L\) is part of Shell V13.
For the uniform cubic bound, the lattice detail
minimum is \(O(|v|^2)\) in energy, its full field
is \(O(|v|)\), and the constant variation has
zero curl. The primitive cubic Wilson term
contains one curl and two fields, and all
higher terms have at least four fields.
Differentiation in the constant coordinate
therefore gives \(D_s\mathcal S_L=O(|v|^3)\),
uniformly also on the complex ball. Multiplication
by the uniformly bounded source-chart inverse
gives \(\zeta_L=O(|v|^3)\), with the sign
convention of Lemma~\ref{lem:v12-hessian}.
\end{proof}

The actual lattice tangent is
\[
 T_L^fU=(\ell_LJ_LU)/L+U,\qquad U\in V_L.
\]
Let \(H_{0,L}\) be the energy Riesz operator of
\(D^2\mathcal S_L(a_L^*)[T_L^fU,T_L^fW]\), and
let \(B_L^{\rm c}\) be the energy Riesz operator of
\[
             -\zeta_L\cdot
                D^2\mathcal F_{L,1}(a_L^*)[T_L^fU,T_L^fW].
\]
By the exact multiplier chain rule their sum
is the actual lattice chart Hessian. The
superscript c distinguishes this curvature
operator from the linear block \(B_{L/M}\).
Extend them to \(V\) by
\[
 \widehat B_L=J_LB_L^{\rm c}J_L^*,\qquad
              \widehat H_{0,L}=J_LH_{0,L}J_L^*+I-P_L .
\]
The unused modes have identity reference and
zero source correction; no new fluctuations
are being added to the relative determinant.

\begin{theorem}[Actual source-curvature convergence in \(\mathfrak S_5\)]
\label{thm:v13-B}
On a common small source ball,
\[
     \sup_L\|\widehat B_L(v)\|_{\mathfrak S_{9/2}}
                                      \le C|v|^3,
\]
and, for every real \(v\) there,
\[
             \|\widehat B_L(v)-\mathcal B_v\|_{\mathfrak S_5}\to0 .
\]
All the finite-lattice families are holomorphic
in \(\mathfrak S_{9/2}\) with the displayed
bound on the smaller complex domain.
\end{theorem}

\begin{proof}
Use the extended tangent
\(\overline T_LU=(\ell_LU,U)\in X\).
The form of \(\widehat B_L\) is exactly
\[
 -\zeta_L\cdot
 D^2\widetilde{\mathcal O}_L(x_L)
                         [\overline T_LU,\overline T_LW].
\]
Indeed the source ignores \(P_L^\perp\), and
\(\ell_L\) vanishes on that complement.
This identity avoids falsely claiming that
\(P_L\to I\) in operator norm. All other
factors converge in operator norm by the two
preceding results, so \(\widehat B_L\to\mathcal B_v\)
in operator norm at real \(v\).

Replace the differentiated source in this
identity by the stopped map of
\eqref{eq:v13-stopped}, with \(x_L,\zeta_L,\ell_L\)
unchanged. Its form has rank at most \(12M^4\),
and its operator error is
\(C|\zeta_L|M^{-9/10}\), uniformly in \(L\).
At \(M=1\) its second derivative is zero.
At \(M\ge L\) it is exactly the full form.
The dyadic difference proof of
Theorem~\ref{thm:v12-S5} therefore applies
uniformly before the lattice limit, including
the eventual plateau. With \(q=9/2\),
\[
                    \frac9{10}-\frac4q=\frac1{90}>0,
\]
so it gives the stated uniform
\(\mathfrak S_q\) bound. The same estimate
applies to \(\mathcal B_v\).

For any \(D\in\mathfrak S_q\), \(q<5\),
bounding each singular value by \(\|D\|\)
gives
\[
              \|D\|_{\mathfrak S_5}
                \le\|D\|^{\,1-q/5}\|D\|_{\mathfrak S_q}^{\,q/5}.
\]
Apply this to \(D=\widehat B_L-\mathcal B_v\).
Here \(1-q/5=1/10\), the operator norm tends
to zero, and the \(q\)-norm is uniformly
bounded. This proves \(\mathfrak S_5\)
convergence. Holomorphy and complex bounds
follow from the same uniform finite-rank
series as in V12. This theorem is about the
actual lattice source derivative, not a
continuum minimum substituted into that
derivative.
\end{proof}

\subsection{Mixed Wilson variations at a strongly converging background}

We next prove the reference convergence needed
to normalize that source correction. This
step uses the full Wilson plaquette and its
curvature square; it is not a Gaussian
replacement of the lattice action.

\begin{lemma}[Weak--strong mixed plaquette limit]
\label{lem:v13-Wilson}
Let real \(a_L\) satisfy the bounds and
strong limits \eqref{eq:v13-classical-inputs}
with limit \(A\). Let \(u_L\) be discrete
Coulomb, allowing constants, with uniformly
bounded critical energy and
\(\mathcal I_Lu_L\rightharpoonup U\) in \(L^4\).
Let the nodal physical field \(Lz_L\) have
a fixed finite set of lattice Fourier modes,
including possible constants, whose associated
trigonometric representatives converge with
all fixed derivatives to a smooth field \(Z\).
Then
\[
            D^2\mathcal S_L(a_L)[u_L,z_L]
                      \longrightarrow D^2\mathcal S_\infty(A)[U,Z].
\]
The derivative on the right is interpreted
weakly through the \(H^1\) bounds on \(U\).
\end{lemma}

\begin{proof}
For one plaquette write its ordered product
as \(Q_p(a)\). It is an analytic function of
four oriented Lie inputs on a fixed small
polydisc. The actual background is in that
polydisc since its counting-\(\ell^4\)
bound controls every link. Finite local
incidence and the elementary pointwise bound
\[
 |DQ_p(a)u-DQ_p(0)u|
             \le C\Bigl(\sum_p|a|\Bigr)\Bigl(\sum_p|u|\Bigr)
\]
give
\begin{equation}\label{eq:v13-first-variation-bound}
 \|L^2DQ(a)u\|_{L^2}
       \le\|du\|_{\ell^2}+C\|a\|_{\ell^4}\|u\|_{\ell^4}.
\end{equation}
Here and below plaquette arrays are reconstructed
on cells, so multiplication by \(L^2\) preserves
their counting-\(\ell^2\) norm.

Expand \(DQ_p(a)u\) through terms linear in
the background. Its leading term is the
fine curl of \(u\). At a common initial
vertex its next imaginary term is
\([u_\mu,a_\nu]+[a_\mu,u_\nu]\).
The shifted-vertex errors and higher terms,
after physical reconstruction, have
distributional \(L^1\) bounds of the form
\[
 \frac C L\{\|\mathcal I_La\|_2
                  \|D^L\mathcal I_Lu\|_2+
             \|\mathcal I_Lu\|_2
                  \|D^L\mathcal I_La\|_2+
             \|\mathcal I_La\|_4^2\|\mathcal I_Lu\|_2\}
 +\frac C{L^2}
       \|D^L\mathcal I_La\|_2\|D^L\mathcal I_Lu\|_2.
\]
They tend to zero. The real quadratic term
is a product of the two fine curls and is
included in the last bound. Taylor remainders
are bounded on the fixed polydisc by
\(C(\sum|a|)^2\sum|u|\), which gives the
third term above. Bounded incidence absorbs
all shifts.
Discrete summation by parts identifies the
weak derivative limit. Strong \(L^4\)
convergence of the background and weak
\(L^4\) convergence of the variation identify
the bracket products in distributions.
Together with \eqref{eq:v13-first-variation-bound}
this proves
\[
        L^2DQ(a_L)u_L\rightharpoonup dU+J_AU
                                                   \quad\hbox{in }L^2 .
\]

For the smooth fixed variation \(z_L\), the
corresponding convergence is strong in \(L^2\).
To justify it at a merely \(L^4\) background,
approximate \(A\) in \(L^4\) by smooth fields.
The pointwise second derivative bound gives
\[
 \|L^2\{DQ(a)-DQ(b)\}z_L\|_2
                  \le C\|a-b\|_{\ell^4}\|z_L\|_{\ell^4}.
\]
The background strong limit controls this
error uniformly, while for a fixed smooth
background ordinary Taylor expansion gives
strong convergence. Then let the smooth
approximation error tend to zero.

For the mixed second variation, the analytic
remainder satisfies
\[
 |D^2Q_p(a)[u,z]-D^2Q_p(0)[u,z]|
       \le C(\sum|a|)(\sum|u|)(\sum|z|).
\]
Since \(Lz_L\) and its physical differences
are uniformly bounded, the reconstructed
\(L^2\) norm of this error is at most
\[
                C L^{-1}\|\mathcal I_La\|_4
                           \|\mathcal I_Lu\|_4
                           \|Lz_L\|_\infty\longrightarrow0.
\]
In the quadratic term at zero, replacing
shifted fields by initial-vertex fields
costs at most
\[
 \frac C L\{\|D^L\mathcal I_Lu\|_2\|Lz_L\|_\infty+
       \|\mathcal I_Lu\|_2\|D^L\mathcal I_Lz_L\|_\infty\}
   +\frac C{L^2}\|D^L\mathcal I_Lu\|_2
                       \|D^L\mathcal I_Lz_L\|_\infty .
\]
The last term includes the real curl product.
The remaining imaginary expression is
\[
           C(U,Z)_{\mu\nu}=[U_\mu,Z_\nu]+[Z_\mu,U_\nu].
\]
Thus \(L^2D^2Q(a_L)[u_L,z_L]\rightharpoonup
C(U,Z)\) in \(L^2\).

Finally the exact real Wilson normalization is
\(\mathcal S_L(a)=\tfrac12\|L^2(Q(a)-I)\|_2^2\).
Its mixed second derivative is
\[
 \langle L^2DQ(a)u,L^2DQ(a)z\rangle+
          \langle L^2(Q(a)-I),L^2D^2Q(a)[u,z]\rangle .
\]
The first pairing has one weak and one strong
factor. The second has the inherited strong
curvature factor and the weak mixed variation.
Both converge to the two terms in the actual
continuum Hessian formula of V12.
\end{proof}

\begin{theorem}[Strong convergence of the actual reference operators]
\label{thm:v13-H0}
On a common small real source ball,
\[
 \|\widehat H_{0,L}-I\|\le\tfrac14,\qquad
        \widehat H_{0,L}\longrightarrow\mathcal H_{0,v}
                                        \quad\hbox{strongly on }V .
\]
Consequently
\(\widehat H_{0,L}^{-1/2}\to\mathcal H_{0,v}^{-1/2}\)
strongly, with uniformly bounded norms.
No operator-norm rate for this convergence is asserted.
\end{theorem}

\begin{proof}
The uniform primitive Wilson critical-energy
estimate of Shell V9 and the bounded chart
derivative give
\(\|H_{0,L}-I_{V_L}\|\le C|v|\).
Its kinetic term is exactly \(I_{V_L}\)
because \(dT_L^fU=dU\). Extending by identity
gives the first bound after shrinking.

Fix a continuum transverse trigonometric
polynomial \(W\), and take any bounded
sequence \(U_L\rightharpoonup U\) in \(V\).
The actual lattice tangents corresponding
to \(J_L^*U_L\) have weak physical limit
\(\mathcal T_vU\); their constant coefficients
converge by \(\ell_L\to\ell_\infty\).
The lattice tangent corresponding to
\(J_L^*W\) has only finitely many physical
Fourier modes, converging smoothly to
\(\mathcal T_vW\). Lemma~\ref{lem:v13-Wilson}
therefore proves
\[
             \langle\widehat H_{0,L}W,U_L\rangle_V
                       \longrightarrow
                       \langle\mathcal H_{0,v}W,U\rangle_V .
\]
The identity extension contributes nothing
for sufficiently large \(L\), since the
fixed polynomial lies in the range of \(P_L\).
Uniform boundedness and density extend this
statement to every fixed \(W\in V\).

This weak--strong testing property implies
strong convergence of the operators.
Otherwise choose normalized residual vectors
for a fixed \(W\) along a failing subsequence,
pass to their weak limit in \(V\), and test
the preceding identity against those vectors.
It forces their pairing with the residual
to tend to zero, contradicting its norm
lower bound. For real data the operators
are self-adjoint with spectrum in
\([3/4,5/4]\). Uniform polynomial
approximation of \(x^{-1/2}\) on that
interval, and strong convergence of each
polynomial, give the inverse-square-root
assertion.
\end{proof}

\subsection{The actual relative determinant converges}

Let
\[
 \widehat K_L=
      \widehat H_{0,L}^{-1/2}\widehat B_L\widehat H_{0,L}^{-1/2}.
\]
It is the isometric extension by zero of
\(H_{0,L}^{-1/2}B_L^{\rm c}H_{0,L}^{-1/2}\)
on the actual finite detail space.
For real \(v\) it is self-adjoint.

\begin{theorem}[Actual-lattice fifth determinant and source jets]
\label{thm:v13-det}
For every real source in a sufficiently small
common ball,
\[
            \|\widehat K_L(v)-\mathcal K_v\|_{\mathfrak S_5}\to0,
 \qquad
 \frac12\log\det\nolimits_5(I+\widehat K_L(v))\to\Theta(v).
\]
Convergence is locally uniform on a smaller
complex source ball, and every fixed source
derivative converges there. The determinant
on the left is the regularized relative
determinant of the actual lattice constrained
Hessian against its actual reference part.
There is no claimed mesh rate.
\end{theorem}

\begin{proof}
Write \(A_L=\widehat H_{0,L}^{-1/2}\) and
\(A=\mathcal H_{0,v}^{-1/2}\) temporarily.
Their norms are uniformly bounded and
they converge strongly, together with
their adjoints for real sources. The
difference \(A_L\widehat B_LA_L-A\mathcal B_vA\)
is the sum of
\[
 A_L(\widehat B_L-\mathcal B_v)A_L,\qquad
               (A_L-A)\mathcal B_vA_L,\qquad
                         A\mathcal B_v(A_L-A).
\]
The first tends to zero in \(\mathfrak S_5\)
by Theorem~\ref{thm:v13-B}. For the others,
approximate the fixed \(\mathcal B_v\) in
\(\mathfrak S_5\) by finite-rank operators.
On their finite-dimensional range or adjoint
range, strong convergence is uniform; the
remaining error is bounded by the uniform
operator norms times the trace-ideal error.
This proves \(\mathfrak S_5\) convergence.

For all complex sources in a common small
ball, the primitive and source bounds give
\(\|\widehat H_{0,L}-I\|<1/4\) and
\(\|\widehat K_L\|_{\mathfrak S_{9/2}}\le C|v|^3\).
The square roots use the analytic binomial
series. Thus the finite families are
holomorphic in \(\mathfrak S_5\) with one
uniform bound. Shrink so that all operator
norms are at most \(1/2\).

For completeness, the needed Banach-valued
analytic upgrade does not assume a compact
unit ball. If holomorphic functions on a
fixed complex polydisc are uniformly bounded
and converge in a Banach norm on a real open
neighbourhood of zero to an identified
holomorphic function, then each derivative
at zero converges in that norm: real finite
differences approximate that derivative
with uniform error by the Cauchy bounds
for one further derivative. First pass
to the limit for a fixed finite-difference
step, then let the step tend to zero.
The uniform coefficient bounds make the
multivariate Taylor tails summable on a
smaller polydisc. This proves locally
uniform norm convergence there and, by
Cauchy, convergence of every fixed
derivative. Apply this argument in
\(\mathfrak S_5\) to \(\widehat K_L\),
using the real convergence already proved.

The determinant estimate of
Lemma~\ref{lem:v12-det} now proves locally
uniform convergence of the logarithms;
their holomorphy and Cauchy's formula give
all the fixed source derivatives. Zero
extension adds only eigenvalues zero to
the relative correction, so its fifth
determinant is exactly the finite-lattice
one, not a different infinite-volume
normalization.
\end{proof}

In particular the finite identity is now
for the \emph{actual} lattice operators:
\[
 \frac12\log\frac{\det(H_{0,L}+B_L^{\rm c})}{\det H_{0,L}}
  =\frac12\log\det\nolimits_5(I+\widehat K_L)
     +\frac12\sum_{j=1}^4
             \frac{(-1)^{j+1}}{j}\Tr(\widehat K_L^j).
\]
The first term on the right converges
with its fixed source jets by the new
theorem. Convergence of the four terms
in the sum is still not inferred.
Nor does convergence of this relative
factor construct the reference
\(\det H_{0,L}\), the chart coarea
density, or their renormalized product.
This version fills the source-curvature
lattice matching left by V12; it does
not declare the entire measured packet
matched.

\subsection{Verification, append-only record, and remaining work}

The exact checker is
\begin{center}\small
\path{scripts/verify_ym_f14_v13_lattice_source.py}.
\end{center}
Its SHA--256 is
\begin{center}\scriptsize\ttfamily
7F6CB208C426714D7B706B3C9F09610DE393100A1BEF5DD949B1D10614D39892.
\end{center}
It checks 24 mixed quadratic plaquette
words, 24 exact curvature-square Hessian
identities on noncentral unitary word
curves, and 24 mixed-bracket limits.
It verifies 510 exact discrete Hodge
identities with Gaussian-integer Fourier
symbols, including transverse and
unrestricted test vectors.
The real Fourier orbit counts at
\(L=2,4,8\) give respectively
135, 2295, and 36855 detail coordinates;
each includes the 15 self-conjugate
nonzero orbits.
An exact rational nontrivial isometric
embedding checks the identity reference
extension, the relative determinant,
and its first nine traces. The exponents
\(9/10-8/9=1/90\) and
\(1-(9/2)/5=1/10\) are checked rationally.

The analytic proof, not those finite
samples, establishes weak reconstruction,
common-source convergence, the strong
reference limit, and trace-ideal
convergence. No numerical eigenvalue
fit or formal omission of Nyquist
modes enters the theorem. No numerical
radius or mesh rate is supplied.

The predecessor V12 has 231,370 bytes
and SHA--256
\begin{center}\scriptsize\ttfamily
80932D806B9D822F7238A13E65A1D3995FCD5474578F8DBCE0210B2B7353F689.
\end{center}
Its complete body is preserved apart
from the title version and terminal
insertion. Active V13 replaces V12.
The separate Shell-Mixing V13 is
unchanged despite sharing the numeral;
it is not the main note's predecessor.
Archives, monographs, and SM/NS notes
remain untouched. The next regular
companion checkpoint is V15.

\begin{firewall}
The common space is normalized by
classical detail kinetic energy on
one fixed periodic cell. A positive
reference Hessian in that norm is
not a positive physical quantum mass.
The result concerns one relative
constraint-curvature factor of the
actual small-branch lattice Hessian.
It does not prove convergence of the
four lower traces, the full measured
Gaussian normalization, higher-loop
remainders, arbitrary compact-source
coverage, a thermodynamic quantum
state, or a physical Hamiltonian
mass gap.
\end{firewall}

The next upstream acquisition is the
lower-trace and measured-density
comparison on these same lattices.
Their signs, source dependence,
local counterterm content, and mixed
shell ownership must be retained.
The newly matched fifth tail cannot
replace any of those quantities.
The independent compact-complement
and physical-transfer obligations
also remain active; no global
completion claim is made.
\section{V14: actual source density and the first constraint-trace bank}
\label{sec:v14-density}

\subsection{Context, inherited inputs, and the question being answered}

V13 matched the actual lattice constraint correction in
\(\mathfrak S_5\). Its regularized determinant begins in source
degree fifteen, whereas the measured one-loop packet needs
source degree six. This version therefore goes upstream of
that determinant tail. We isolate the finite source-coordinate
Jacobian, the primitive exponential-coordinate Haar factor,
and the complete leading quadratic trace on the first lattice.
These are separate factors, not a declaration that the measured
packet has already converged.

We continue on the single periodic physical cell of V12--V13,
with \(L=2^n\), the literal balanced ordered-tree source
\(\mathcal F_L\), and its actual small constrained minimum
\[
 a_L^*(v)=s_L(v)/L+w_L(v),\qquad
 w_L(v)\in V_L,\qquad
 \mathcal F_L(a_L^*(v))=v.
\]
The real detail space is mean-zero discrete Coulomb, with
energy norm \(\|d_Lw\|_2\). Lie vectors use quaternion
coefficients, so \([a,b]=2a\times b\). All source logarithms
and determinants below are normalized to vanish at \(v=0\);
complex logarithms use their identity branch.

The added Shell-Mixing V13 remains the same 512,757-byte file
reviewed in V11--V13, with SHA--256
\begin{center}\scriptsize\ttfamily
7F2DBD2900D7A78280DA980152266E94DF8A95EFE15CD37429EC8C6AAAA4F17B.
\end{center}
We use its literal quadratic stencil, homogeneous cubic map,
and degree-four detail displacement. Those coefficients are
inherited, not counted again as new acquisitions. In particular,
if \(\Phi_L(c)=\mathcal F_L(c/L)\), its fixed-\(L\) identities imply
\begin{equation}\label{eq:v14-inherited}
\begin{aligned}
 a_L^*(v)&=\Phi_L^{-1}(v)/L+O_L(|v|^4),\\
 G_L(v)&=E_4(v)+O_L(|v|^5),\\
 E_4(v)&=\frac12\sum_{\mu<\nu}|[v_\mu,v_\nu]|^2,\\
 g_\mu(v):=(\nabla E_4)_\mu
 &=4\sum_{\nu\ne\mu}
       \bigl(|v_\nu|^2v_\mu-(v_\mu\cdot v_\nu)v_\nu\bigr).
\end{aligned}
\end{equation}
The first line is a finite-dimensional germ statement; its
displayed remainder does not assert a uniform mesh rate.
Uniform analytic domains and the real strong convergence of
the actual minima are supplied by V13. In its notation,
\[
\begin{aligned}
 x_L(v)=(s_L(v),J_Lw_L(v))\longrightarrow x_*(v)
       &\quad\hbox{strongly in }X=\mathbb R^{12}\oplus V,\\
 I_La_L^*(v)\longrightarrow A_*(v)
       &\quad\hbox{strongly in }L^4.
\end{aligned}
\]
These are the nontrivial classical inputs used below.

\subsection{The source-coordinate density at the actual minimum}

Here a source derivative keeps the detail coordinate fixed:
\[
 Z_L(v)q
   =D\mathcal F_L(a_L^*(v))[q/L],
 \qquad
 Z_\infty(v)q=D\mathcal O_\infty(A_*(v))[q].
\]
These are \(12\)-by-\(12\) matrices; they are not the full
derivatives of the minimizing map \(v\mapsto a_L^*(v)\).

\begin{theorem}[Source-coordinate Jacobian, not the full density]
\label{thm:v14-source-density}
On a common smaller source neighbourhood, \(Z_L\) is invertible,
its real determinant is positive, and
\begin{equation}\label{eq:v14-Zlimit}
 \log\det Z_L(v)\longrightarrow\log\det Z_\infty(v)
\end{equation}
locally uniformly on a complex neighbourhood of zero, together
with every fixed source derivative. The quadratic Taylor
polynomial is
\begin{equation}\label{eq:v14-Zquadratic}
 J^2\log\det Z_L(v)
       =-(1-L^{-2})\sum_\mu |v_\mu|^2.
\end{equation}
Norm squares in holomorphic formulas denote their bilinear
complexifications.

More precisely, let \(m_L(s,w)\,ds\,dw\) be any specified
density in these linear coordinates. Under the local change
\((s,w)\mapsto(v=\mathcal F_L(s/L+w),w)\), that density is
\begin{equation}\label{eq:v14-density-change}
 \frac{m_L(s(v,w),w)}
      {\det D_s\mathcal F_L(s(v,w)/L+w)}\,dv\,dw
\end{equation}
on the real branch. Thus \(+\log\det Z_L\) is the source-chart
contribution to the negative logarithm of the amplitude at
the minimum. It is inserted once, not once more after another
coarea formula that already contains it.
\end{theorem}

\begin{proof}
Theorem~\ref{thm:v13-source} gives convergence of the first
derivatives of the common-space source in operator norm.
Lemma~\ref{lem:v13-branch} gives strong convergence of its
real evaluation points \(x_L(v)\). Their combination proves
real pointwise matrix convergence \(Z_L(v)\to Z_\infty(v)\).
The common analytic estimates give
\(\|Z_L-I\|\le C|v|\) and holomorphy on one fixed complex
ball; shrink so this norm is at most \(1/2\).
The analytic matrix logarithm then exists and has a common
bound. Apply the real finite-difference and Cauchy argument
proved in Theorem~\ref{thm:v13-det}, now in a finite-dimensional
matrix space and then to \(\Tr\log Z_L\). This proves
\eqref{eq:v14-Zlimit} and its fixed derivatives. On the real
ball, the determinant cannot change sign from its value one
at zero.

For the coefficient, the inherited homogeneous identity is
\[
 (\Phi_L(c))_\mu=c_\mu+
   \frac{1-L^{-2}}{24}
       \sum_\nu[c_\nu,[c_\nu,c_\mu]]+O_L(|c|^4).
\]
For \(\nu\ne\mu\), differentiation in \(c_\mu\) gives
\(\operatorname{ad}_{c_\nu}^2\), whose real colour trace is
\(-8|c_\nu|^2\). The term \(\nu=\mu\) is identically zero.
Taking the divergence over four orientations therefore gives
\(-24\sum_\nu|c_\nu|^2\) before multiplication by the displayed
coefficient. Since \(D\Phi_L-I\) starts in degree two,
the same trace is the quadratic term of \(\log\det D\Phi_L\).
The difference between the actual minimum and the homogeneous
inverse starts in degree four by \eqref{eq:v14-inherited}.
Consequently it cannot change this degree-two source
derivative. Also \(\Phi_L^{-1}(v)=v+O_L(|v|^3)\), which cannot
change the quadratic polynomial. This proves
\eqref{eq:v14-Zquadratic}.

Finally the derivative of the coordinate change has blocks
\[
 \begin{pmatrix}D_s\mathcal F_L&D_w\mathcal F_L\\0&I\end{pmatrix}.
\]
Its determinant is \(\det D_s\mathcal F_L\); the ordinary
change-of-variables theorem proves \eqref{eq:v14-density-change}.
Choosing an energy-orthonormal basis for \(w\) changes a
Euclidean reference volume by a source-independent constant.
Such constants cancel under vacuum normalization and do not
alter this source factor.
\end{proof}

The arbitrary density \(m_L\) may itself contain gauge-chart,
primitive Haar, incoming one-loop, and other factors.
Equation~\eqref{eq:v14-density-change} computes none of those
by implication. In particular, equality of balanced and
unbalanced coarse gauge orbits does not set \(m_L=1\).

\subsection{A precise finite part of the primitive Haar factor}

The normalized exponential-coordinate \(\mathrm{SU}(2)\) Haar
factor per fine link is
\[
 h(a)=\left(\frac{\sin|a|}{|a|}\right)^2,\qquad h(0)=1.
\]
It follows either from the volume element on the unit
three-sphere in polar coordinates, whose radial part is
\(\sin^2r\,dr\), or directly from the exponential map.
For \(a\in\mathbb C^3\) its local continuation is the entire
sinc series in \(a\cdot a\), squared. Define
\[
 H_L(v)=\prod_{x,\mu}h(a_{L,\mu}^*(v;x)),\qquad
 R_{H,L}(v)=-\log H_L(v)
             -\frac13\sum_{x,\mu}a_{L,\mu}^*(v;x)
                                      \cdot a_{L,\mu}^*(v;x).
\]
This is a specified factor evaluated on the minimum, not the
entire measure induced on a Coulomb slice.

\begin{theorem}[Primitive Haar finite part on the actual branch]
\label{thm:v14-Haar}
After a common shrink of the source ball,
\begin{equation}\label{eq:v14-Haar-limit}
 R_{H,L}(v)\longrightarrow
 R_{H,\infty}(v)=
   \frac1{90}\int_{\mathbb T^4}
          \sum_\mu\bigl(A_{*,\mu}(v)\cdot A_{*,\mu}(v)\bigr)^2
\end{equation}
locally uniformly on a complex neighbourhood of zero, with
every fixed source derivative. The subtraction is the exact
lattice quadratic term
\begin{equation}\label{eq:v14-lattice-subtraction}
 \frac13\sum_{x,\mu}|a_{L,\mu}^*(v;x)|^2
                  =\frac{L^2}{3}\|I_La_L^*(v)\|_{L^2}^2
\end{equation}
for real sources. Replacing the right side by
\(L^2\|A_*(v)\|_{L^2}^2/3\) is not justified by convergence
without a rate.
\end{theorem}

\begin{proof}
The scalar Taylor series at zero, or division and integration
of the sine and cosine series, gives
\begin{equation}\label{eq:v14-Haar-series}
 -\log h(a)=\frac{a\cdot a}{3}
       +\frac{(a\cdot a)^2}{90}
       +\frac{2(a\cdot a)^3}{2835}
       +O(|a|^8).
\end{equation}
In particular the remainder after the quartic term is bounded
by \(C|a|^6\) on one small real or complex ball.
The physical reconstruction \(A_L=I_La_L^*(v)\) is constant
with value \(La_{L,\mu}^*(x)\) on each cell of volume \(L^{-4}\).
Consequently
\[
 \sum_{x,\mu}|a_{L,\mu}^*(x)|^4
                  =\int_{\mathbb T^4}\sum_\mu|A_{L,\mu}|^4,
 \qquad
 \sum_{x,\mu}|a_{L,\mu}^*(x)|^2=L^2\|A_L\|_2^2.
\]
For fixed real \(v\), the strong \(L^4\) convergence from V13
implies uniform integrability of \(\sum_\mu|A_{L,\mu}|^4\).
The integrals of this family over any single lattice cell
therefore tend to zero uniformly in the cell. The first
identity, applied to one link, proves
\(\max_{x,\mu}|a_{L,\mu}^*(x)|\to0\).
Thus
\[
 \sum_{x,\mu}|a_{L,\mu}^*(x)|^6
 \le \max_{x,\mu}|a_{L,\mu}^*(x)|^2
                     \sum_{x,\mu}|a_{L,\mu}^*(x)|^4
 \longrightarrow0.
\]
The quartic integral converges by strong \(L^4\) convergence,
so the Taylor formula proves \eqref{eq:v14-Haar-limit} for
each real \(v\).

For complex sources on a common small ball, the uniform
analytic branch estimate gives
\(\sum_{x,\mu}|a_{L,\mu}^*(v;x)|^4\le C|v|^4\).
It bounds each individual link by \(C^{1/4}|v|\), so a
single shrink places every link in the same logarithm
branch. The exact subtracted logarithm is holomorphic and
satisfies \(|R_{H,L}(v)|\le C'|v|^4\), independently of \(L\).
The proposed limit is a continuous quartic functional of
the holomorphic \(L^4\)-valued continuum minimizing map.
The finite-difference/Cauchy argument again upgrades the
real convergence to locally uniform complex convergence
and all fixed source jets. The equality
\eqref{eq:v14-lattice-subtraction} was already verified.
Multiplication by \(L^2\) explains why a rate-free
\(L^2\) limit cannot replace its exact finite-\(L\) value.
\end{proof}

\begin{proposition}[Commuting calibration and its limitation]
\label{prop:v14-flat}
For a small real commuting tuple \(v\), the actual local
minimum is \(a_{L,\mu}^*(x)=v_\mu/L\), and
\begin{equation}\label{eq:v14-flat-Haar}
 -\log H_L(v)=\frac{L^2}{3}\sum_\mu|v_\mu|^2
                  +\frac1{90}\sum_\mu|v_\mu|^4
                  +O\left(L^{-2}\sum_\mu|v_\mu|^6\right).
\end{equation}
In particular, the source factor
\(\log\det Z_L(v)=O(1)\) cannot by itself cancel the
primitive \(L^2\) term.
\end{proposition}
\begin{proof}
For a constant commuting connection, all tree transports,
two-link paths, and balancing factors commute. The balanced
block doubles the connection at each step, hence
\(\mathcal F_L(v/L)=v\). Every Wilson plaquette is the
identity, so the nonnegative action has value zero.
The strictly convex local constrained branch of V13 has
a unique stationary minimum, which must be this one.
There are \(L^4\) links in each orientation. Substituting
\(a=v_\mu/L\) in \eqref{eq:v14-Haar-series} proves
\eqref{eq:v14-flat-Haar}, with a uniform small-source remainder.
The bound on \(\log\det Z_L\) follows from
Theorem~\ref{thm:v14-source-density}.
\end{proof}

This is not an uncancelled physical mass counterterm.
The displayed divergence belongs to one primitive coordinate
factor. The gauge/quotient Jacobian, reference Gaussian
determinant, incoming density, and prescribed normalization
have not been combined with it. Declaring a physical mass
from this term would confuse a coordinate computation with
the measured quantum theory.

\subsection{The lower-trace reduction required by the sixth source jet}

Write the actual relative operator of V13 as
\[
 K_L(v)=H_{0,L}(v)^{-1/2}B_L^{\rm c}(v)H_{0,L}(v)^{-1/2},
 \qquad K_L(v)=\sum_{r\ge3}K_{L,[r]}(v).
\]
Here \([r]\) is homogeneous source degree, and \(J^r\) means
the Taylor polynomial through degree \(r\), not a derivative
without factorial normalization.

\begin{lemma}[Exact sixth-degree trace identity]
\label{lem:v14-J6}
For each finite \(L\),
\begin{equation}\label{eq:v14-J6}
 J^6\left[\frac12\log\det(I+K_L(v))\right]
   =\frac12J^6\Tr K_L(v)
                         -\frac14\Tr K_{L,[3]}(v)^2.
\end{equation}
Moreover \(\Tr K_{L,[3]}(v)=0\).
Consequently the first trace on the right needs degrees
four, five, and six; the second term needs only the cubic
operator coefficient.
\end{lemma}
\begin{proof}
At finite \(L\), expand the ordinary matrix logarithm.
Since \(K_L=O(|v|^3)\), powers at least three start in
degree nine, and the degree-six part of \(K_L^2\) is
exactly \(K_{L,[3]}^2\). This proves the identity.
At zero source, the constant/detail chart tangent on
details is the identity, and \(H_{0,L}(0)=I\) in the
energy norm. The multiplier satisfies
\(\zeta_L(v)=g(v)+O_L(|v|^4)\) by
\eqref{eq:v14-inherited}. Therefore the cubic coefficient is
the energy Riesz operator of the bilinear form
\begin{equation}\label{eq:v14-cubic-form}
 (U,V)\longmapsto
       -\sum_\mu g_\mu(v)\cdot
                         D^2\mathcal F_{L,\mu}(0)[U,V].
\end{equation}
Every quadratic term of the literal group-word source is a
linear combination of Lie brackets. This follows by the
degree-two BCH formula for multiplication, negation for
inversion, and linearity for averaging and composition.
In colour coordinates the Hessian of such a bracket has
zero colour trace against the colour-scalar zero-source
covariance. Equivalently, the explicit tensor formula in
the next proposition has colour factors of trace zero.
It proves the claimed cancellation on the full detail
space for every finite \(L\).
\end{proof}

The vanishing just proved is only the cubic Taylor
coefficient. It is not the assertion \(\Tr K_L(v)=0\).
Also, \eqref{eq:v14-J6} is a finite-germ identity, not a
uniform bound for those lower traces as \(L\to\infty\).

\subsection{Exact spatial skew banks and their full colour contraction}

Enumerate the \(4L^4\) scalar oriented links by \(l\).
Define real skew matrices \(T_{L,\mu}\) by the quadratic
homogeneous coefficient of the actual source:
\begin{equation}\label{eq:v14-bank}
 \mathcal F_{L,\mu}^{[2]}(A)
       =\sum_{l<m}(T_{L,\mu})_{lm}[A_l,A_m].
\end{equation}
This identity specifies the coefficient of each pair of
independent Lie variables; no choice of a detail basis is
involved. Let \(C_L\) be the inverse scalar detail kinetic
form on zero-mean Coulomb links, extended by zero to the
orthogonal complement in scalar link coordinates. In
particular \(C_L\) is real symmetric and nonnegative.
Let \(\mathcal J_a b=a\times b\), so
\[
 \Tr\mathcal J_a=0,\qquad
 \Tr(\mathcal J_a\mathcal J_b)=-2a\cdot b.
\]

\begin{proposition}[All-detail bubble bank]
\label{prop:v14-bank}
Define
\begin{equation}\label{eq:v14-D}
 D_{L,\mu\nu}=-\Tr(C_LT_{L,\mu}C_LT_{L,\nu}).
\end{equation}
The matrix \(D_L\) is symmetric and positive semidefinite.
The full cubic correction is, in energy-orthonormal
coordinates and with zero extension,
\[
 2\sum_\mu
       (C_L^{1/2}T_{L,\mu}C_L^{1/2})\otimes\mathcal J_{g_\mu(v)}.
\]
Consequently
\begin{equation}\label{eq:v14-bubble}
 \Tr K_{L,[3]}(v)^2
       =8\sum_{\mu,\nu}D_{L,\mu\nu}\,g_\mu(v)\cdot g_\nu(v),
 \qquad
 -\frac14\Tr K_{L,[3]}(v)^2
       =-2\sum_{\mu,\nu}D_{L,\mu\nu}\,g_\mu(v)\cdot g_\nu(v).
\end{equation}
\end{proposition}
\begin{proof}
Differentiating \eqref{eq:v14-bank} gives the true Hessian
\[
 D^2\mathcal F_{L,\mu}(0)[U,V]
                   =\sum_{l,m}(T_{L,\mu})_{lm}[U_l,V_m].
\]
There is no extra factor of two: the two ordered terms
replace the two terms from differentiating a quadratic
monomial. Since
\(-g\cdot[U_l,V_m]=2U_l\cdot(g\times V_m)\),
the Euclidean bilinear-form matrix in
\eqref{eq:v14-cubic-form} is
\(2\sum_\mu T_{L,\mu}\otimes\mathcal J_{g_\mu}\).
Whitening by the kinetic covariance on its range yields
the displayed energy matrix. Both tensor factors are
skew, so their tensor product is symmetric, as required
for a real Hessian.

Put \(S_\mu=C_L^{1/2}T_{L,\mu}C_L^{1/2}\).
Then
\(D_{L,\mu\nu}=-\Tr(S_\mu S_\nu)
=\Tr(S_\mu^tS_\nu)\), proving the Gram and positivity
statements. Taking the tensor trace of the squared
operator, and using the colour identity, proves
\eqref{eq:v14-bubble}. The unsquared colour trace is zero,
also completing the last assertion of
Lemma~\ref{lem:v14-J6}.
\end{proof}

\subsection{The complete \texorpdfstring{\(L=2\)}{L=2} coefficient, including the ordered tree}

The next result evaluates \eqref{eq:v14-D} on all the
one-step details. Directions are numbered \(1,2,3,4\) in
the ascending ordered tree, and scalar links are
\((x,\mu)\), \(x\in\{0,1\}^4\). Let \(e_{x,\mu}\) denote
the corresponding coordinate vector and let
\(\theta_x\in\mathbb R^{64}\) be the indicator of the
ordered path from \(0\) to \(x\). Set
\[
 \phi=\frac1{16}\sum_x\theta_x,\quad
 \psi_x=\theta_x-\phi,\quad
 b_{x,\mu}=e_{x,\mu}+e_{x+e_\mu,\mu},
 \quad a\wedge b=ab^t-ba^t,
\]
where sites are reduced modulo two. The actual bank is
\begin{equation}\label{eq:v14-T2}
 T_{2,\mu}=
   \frac1{16}\sum_x\psi_x\wedge b_{x,\mu}
          +\frac1{32}\sum_x e_{x,\mu}\wedge e_{x+e_\mu,\mu}.
\end{equation}
The last sum is zero after pairing \(x\) and \(x+e_\mu\),
but retaining it before cancellation verifies ownership
of the two-link BCH contribution.

\begin{theorem}[Exact first-lattice all-detail trace]
\label{thm:v14-L2}
For the literal balanced source and the actual free
detail covariance,
\begin{equation}\label{eq:v14-D2}
 D_2=\frac1{442368}
 \begin{pmatrix}
 2941&0&0&0\\
 0&4345&-108&-108\\
 0&-108&5389&-234\\
 0&-108&-234&6133
 \end{pmatrix}.
\end{equation}
This matrix is positive definite. Thus the sextic
quadratic-trace term in \eqref{eq:v14-J6} is strictly
negative for every noncommuting real source tuple.
For \(v_1=xE_1,\ v_2=yE_2,\ v_3=v_4=0\), it is exactly
\begin{equation}\label{eq:v14-two-colour}
 -\frac{2941}{13824}x^2y^4
                          -\frac{4345}{13824}x^4y^2.
\end{equation}
These coefficients do not include the first trace in
\eqref{eq:v14-J6} or the other measured factors.
\end{theorem}

\begin{proof}
On the one-cell output lattice the transported path is
\(t_xe^{A_{x,\mu}}e^{A_{x+e_\mu,\mu}}t_x^{-1}\).
Its quadratic logarithm is the two-link BCH coefficient
plus \([\theta_x(A),b_{x,\mu}(A)]\).
Through degree two, the logarithm of the distinguished-root
average equals the average of the path logarithms:
if \(p_x=\log P_x\), then
\[
 \operatorname{BCH}\left(p_0,
        \frac1{16}\sum_x\operatorname{BCH}(-p_0,p_x)\right)
                   =\frac1{16}\sum_xp_x+O(A^3).
\]
The two degree-two commutators with \(p_0\) cancel.
Conjugating by \(e^{-\phi(A)}\) subtracts
\([\phi^{[1]}(A),\overline b_\mu(A)]\).
Second-order tree terms cancel inside each conjugation,
and second-order terms of \(\phi\) first contribute in
degree three. This proves \eqref{eq:v14-T2} for arbitrary
input fields, without testing a selected subspace.

Here is an explicit finite rational specification of the
remaining contraction. For
\(k\in\{0,1\}^4\setminus\{0\}\), put \(m(k)=\sum_\rho k_\rho\).
The full scalar covariance is
\begin{equation}\label{eq:v14-C2}
 (C_2)_{(x,\mu),(y,\nu)}
  =\frac1{16}\sum_{k\ne0}
    (-1)^{k\cdot(x-y)}
       \frac{\delta_{\mu\nu}-k_\mu k_\nu/m(k)}{4m(k)}.
\end{equation}
Indeed, a forward difference has symbol \(-2k_\mu\);
the zero-mean transverse projection in that Fourier mode
is \(I-kk^t/m(k)\), and its kinetic eigenvalue is \(4m(k)\).
There are fifteen nonzero real Walsh modes, each with
three scalar transverse coordinates. Thus the scalar
rank is \(45\) and the full real Lie detail dimension
is \(135\). No zero mode or longitudinal mode is given
an artificial inverse.

Substituting \eqref{eq:v14-T2} and \eqref{eq:v14-C2} into
\eqref{eq:v14-D} is a finite rational calculation. More
explicitly, form the four matrices \(R_\mu=C_2T_{2,\mu}\)
from those displayed entries and sum
\[
 442368D_{2,\mu\nu}
           =-442368\sum_{i=1}^{64}\sum_{j=1}^{64}
                                      (R_\mu)_{ij}(R_\nu)_{ji}.
\]
In upper-triangular order the ten resulting integers are
\[
 (2941,\ 0,\ 0,\ 0,\ 4345,\ -108,\ -108,\ 5389,\ -234,\ 6133).
\]
The exact rational certificate described below performs
these entire sums, not random samples or a truncated
spectral fit. The displayed formulas specify every input
entry independently of that implementation.
This proves \eqref{eq:v14-D2}.

Every diagonal entry is positive and strictly exceeds
the sum of the absolute off-diagonal entries in its row.
For a real vector \(z\), the inequality
\(2|z_\mu z_\nu|\le z_\mu^2+z_\nu^2\) then proves
\(z^tD_2z>0\) unless \(z=0\).
By Euler's homogeneous identity,
\(\sum_\mu v_\mu\cdot g_\mu(v)=4E_4(v)\).
Thus \(g(v)=0\) implies \(E_4(v)=0\), which for real
sources is equivalent to pairwise commuting components.
Conversely all brackets zero give \(g=0\).
Positive definiteness applied separately to the three
colour components in \eqref{eq:v14-bubble} proves strict
negativity for noncommuting tuples.
For the stated two-colour source,
\(g_1=4xy^2E_1,\ g_2=4x^2yE_2,\ g_3=g_4=0\).
Their colour inner product vanishes, and substituting
the first two diagonal entries proves
\eqref{eq:v14-two-colour}.
\end{proof}

\subsection{All-scale recursion and the cross-shell terms still to acquire}

The preceding coefficient is an exact initial datum for
an all-scale calculation, not evidence of convergence
by extrapolation. There is a useful exact recursion
which states what such a calculation must retain.

Let \(\mathfrak T_M=B_M+E_M^{[2]}+O(A^3)\) be the
balanced step from side \(M\) to side \(M/2\).
Here \(B_M\) is its actual linear block, and
\(E_M^{[2]}\) is its quadratic homogeneous map.
Let \(B^{(r)}\) denote the composition of the first
\(r\) linear blocks, starting at side \(L\);
the empty composition is the identity. Let
\(B_{\rm post}^{(r)}\) be the linear composition
from side \(L/2^{r+1}\) to side one.
Then
\begin{equation}\label{eq:v14-second-chain}
 \mathcal F_L^{[2]}(A)
    =\sum_{r=0}^{n-1}B_{\rm post}^{(r)}
                 E_{L/2^r}^{[2]}(B^{(r)}A).
\end{equation}
For each output orientation, call the skew bank of the
\(r\)-th summand \(T_{L,\mu}^{(r)}\).
Thus \(T_{L,\mu}=\sum_rT_{L,\mu}^{(r)}\).

\begin{proposition}[Exact cross-shell ownership of the sextic bubble]
\label{prop:v14-cross-shell}
For every finite \(L\),
\begin{equation}\label{eq:v14-cross-shell}
 D_{L,\mu\nu}
 =-\sum_{r=0}^{n-1}\sum_{s=0}^{n-1}
       \Tr(C_LT_{L,\mu}^{(r)}C_LT_{L,\nu}^{(s)}).
\end{equation}
It is not, in general, the sum with only \(r=s\).
The formula is invariant under any invertible change
of the scalar detail basis, with the covariance and
the bilinear-form banks transformed together.
\end{proposition}
\begin{proof}
For maps fixing zero,
\((f\circ h)^{[2]}=f^{[1]}\circ h^{[2]}
                          +f^{[2]}\circ h^{[1]}\).
Induction gives \eqref{eq:v14-second-chain}, since a
degree-two term has precisely one quadratic insertion.
Expanding the two occurrences of \(T_L\) in
\eqref{eq:v14-D} proves \eqref{eq:v14-cross-shell}.
There is no orthogonality premise that could delete
its off-diagonal shell pairs. Under coordinates
\(A=R\widetilde A\) on the detail space, a form bank
becomes \(R^tTR\), the covariance becomes
\(R^{-1}CR^{-t}\), and \(CTCT\) changes by similarity.
Its trace is unchanged.
\end{proof}

This identity gives a concrete next acquisition:
evaluate or bound the actual cross-shell contractions,
then combine them with degrees four through six of
\(\Tr K_L\), the reference determinant, and the
remaining amplitude. The cofinal \(\mathfrak S_5\)
theorem controls none of these lower-trace sums by
itself. Likewise, subtracting the primitive Haar
quadratic term in isolation does not specify the
renormalization of their measured product.

\subsection{Exact verification and version record}

The new executable certificate is
\begin{center}\small
\path{scripts/verify_ym_f14_v14_source_low_traces.py}.
\end{center}
Its SHA--256 is
\begin{center}\scriptsize\ttfamily
3EC377DC70D0AA8E362BA62CF6575F159599D7173E1B3BF37B146A525C554BFE.
\end{center}
It uses exact rational arithmetic to construct every
entry of the \(64\)-by-\(64\) scalar covariance and all
four skew banks. It checks the inverse-Laplacian
identity \(C_2(-\Delta)=\Pi_{\rm Coulomb,0}\), projection
trace \(45\), the ten complete Gram contractions in
\eqref{eq:v14-D2}, strict diagonal dominance, the
colour trace identities, and the first three scalar
coefficients in \eqref{eq:v14-Haar-series}.

Independently of the stencil assembly, it evaluates
degree-two BCH jets of the literal ordered tree,
transported two-link paths, distinguished-root average,
and balancing conjugation on 24 rational Lie fields.
Their linear and quadratic outputs agree exactly with
the banks. These 24 checks are supplementary diagnostics:
the algebraic BCH proof establishes the stencil for all
fields, while the Gram calculation itself exhausts
every scalar matrix entry and every contraction.
No numerical sample proves the continuum convergence
in Theorems~\ref{thm:v14-source-density} and
\ref{thm:v14-Haar}; those assertions follow from the
analytic arguments above and their identified V13
classical inputs.

The predecessor main V13 has 260,983 bytes and SHA--256
\begin{center}\scriptsize\ttfamily
FC5101647DE125F86C2B6E3D32E75AE07FE88473016595D208972D7AF7278DD9.
\end{center}
Its complete body is preserved, with only the title
version changed and this section inserted before the
terminal document marker. Active main V14 replaces
active main V13. The separate Shell-Mixing V13,
SM/NS notes, archives, and monographs are preserved.
The next version, V15, is the scheduled SM/NS
companion checkpoint.

\begin{firewall}
The new convergent factors are a finite source-chart
Jacobian and a specifically subtracted primitive Haar
factor on the small classical branch. The exact
constraint bubble has an exact algebraic formula at
every finite \(L\); its matrix of rational coefficients
is fully evaluated here only at \(L=2\).
The full measured Gaussian normalization, the
remaining lower traces and their cancellations,
uniform higher-loop control, arbitrary compact-source
coverage, continuum quantum existence, and a
positive physical Hamiltonian mass gap remain
unproved. A source-degree cancellation and a
135-coordinate finite calculation are not a
substitute for those obligations.
\end{firewall}
\section{V15: a summable actual quadratic-root trace and the companion checkpoint}
\label{sec:v15-HS}

\subsection{Checkpoint and the new direction}

The previous version made concrete progress: it identified
two separately convergent amplitude factors and evaluated
the complete first-lattice cubic constraint bank. The
remaining quadratic trace was then only a finite-\(L\)
formula. This version proves its cofinal convergence,
including absolute control of the mixed-shell sum. The
proof uses the explicit quadratic stencil at zero source;
it does not strengthen the whole finite-amplitude
constraint operator from \(\mathfrak S_5\) to
\(\mathfrak S_2\).

The scheduled companion review was performed before this
calculation. The latest active companions are unchanged:
\begin{itemize}
\item SM V72, 607,372 bytes, SHA--256
\begin{center}\scriptsize\ttfamily
F44F9745D43379E1672C5550411B58E97195A4C9278592D221DF5D3F644DC1F5.
\end{center}
Its terminal theorem is fixed-time empirical covariance
transport for a declared many-copy model with a fixed
physical mode set. It does not control a growing momentum
cutoff. Its warning against importing a diagonal dressing
into an overlapping matrix kernel applies here as a
methodological check: mixed shell entries must be proved
or retained, not discarded.
\item NS V26, 278,283 bytes, SHA--256
\begin{center}\scriptsize\ttfamily
82C887F8C2C69E4F28FAD7C3DC8DE5B9099CB5A4BF431EBA1FFF3E91935978AD.
\end{center}
Its terminal work concerns rank-one contractions of
derivatives of the Oseen kernel. The displayed tensor
reductions do not provide a YM covariance estimate.
Some stated radial certifications refer forward to
V27, which is not the active file read here; none of
those numerical claims is imported as a certified
constant in this proof.
\end{itemize}
This was a review of the relevant terminal arguments
and their scope, not an independent recertification of
every theorem in either accumulated companion. Their
file contents are preserved.

The separate Shell-Mixing V13 is also unchanged, with
SHA--256
\begin{center}\scriptsize\ttfamily
7F2DBD2900D7A78280DA980152266E94DF8A95EFE15CD37429EC8C6AAAA4F17B.
\end{center}
Its actual local quadratic stencil, rather than a
generic analytic remainder, is the additional structure
used below. The first two-shell computation reveals an
exact cancellation; this is recorded at its proved
scope and is not presumed at every scale.

\subsection{Physical linear averages and the free energy embedding}

Separate the scalar spatial detail space
\(\mathcal V_L\), of dimension \(3(L^4-1)\), from its
three Lie colours:
\(V_L=\mathcal V_L\otimes\mathbb R^3\).
The scalar space is the real zero-mean discrete
Coulomb space with inner product
\(\langle d_Lu,d_Lw\rangle_{\ell^2}\).
Put
\[
 \mathcal H_L=\mathbb R^{4L^4},\qquad
 \|u\|_{\mathcal H_L}^2=L^{-4}\sum_{x,\mu}|u_{x,\mu}|^2,
 \qquad \mathcal E_Lu=Lu.
\]
Thus \(\mathcal E_L:\mathcal V_L\to\mathcal H_L\)
is the physical \(L^2\) reconstruction at the level
of its cell values, not an isometry of the energy
norm.

For \(M\mid L\), let \(p=L/M\) and define the actual
physical linear block by
\begin{equation}\label{eq:v15-Q}
 (Q_{L,M}u)_{r,\mu}
   =p^{-5}\sum_{z\in\{0,\ldots,p-1\}^4}
                  \sum_{t=0}^{p-1}u_{pr+z+te_\mu,\mu}.
\end{equation}
Sites are reduced periodically. This is
\(Q_{L,M}=p^{-1}B_p\), not a spectral cutoff.
In particular
\[
 Q_{L,M}\mathcal E_Lw=M B_pw.
\]

\begin{lemma}[Actual averaging contraction and uniform singular values]
\label{lem:v15-embedding}
The map \(Q_{L,M}:\mathcal H_L\to\mathcal H_M\)
is a contraction. If singular values are in decreasing
order, then
\begin{equation}\label{eq:v15-singular}
 s_j(\mathcal E_L)\le j^{-1/4}.
\end{equation}
Consequently, for dyadic \(M=2^m\le L\),
\begin{equation}\label{eq:v15-fourth-sum}
 \sum_j s_j(Q_{L,M}\mathcal E_L)^4
       \le 1+\log(4M^4)\le 4m+3.
\end{equation}
The same estimates hold after the energy isometries
of V13 and zero extension to the common continuum
space.
\end{lemma}

\begin{proof}
Each row of \eqref{eq:v15-Q} has mass one. Each
fine link occurs \(p\) times in the entire unnormalized
path sum for its orientation, so each column of
\(Q_{L,M}\) has mass \(p^{-4}\). Jensen gives
\[
 M^{-4}\sum_{r,\mu}|Q_{L,M}u|^2
 \le M^{-4}p^{-4}\sum_{x,\mu}|u|^2
                 =\|u\|_{\mathcal H_L}^2.
\]

Use centred integer frequency representatives \(k\)
with \(|k_\mu|\le L/2\). At nonzero frequency, the
scalar transverse space has dimension three, and
the singular values of \(\mathcal E_L\) are
\[
 \frac1{\omega_L(k)},\qquad
 \omega_L(k)=2L\left(\sum_\mu
                       \sin^2(\pi k_\mu/L)\right)^{1/2}.
\]
The chord lower bound for sine on \([0,\pi/2]\)
gives \(\omega_L(k)\ge4|k|\).
For \(R\ge1\), the number of such frequencies in
the radius-\(R\) ball, including multiplicity three,
is at most \(3(2R+1)^4\le243R^4\).
Hence, for \(0<t\le1/4\), the number of singular
values at least \(t\) is at most
\[
 243(4t)^{-4}=\frac{243}{256}t^{-4}<t^{-4}.
\]
For \(t>1/4\) it is zero. This proves
\eqref{eq:v15-singular}. Real conjugate Fourier
pairs have the same total real dimension as this
complex frequency count; the self-conjugate modes
are not omitted.

Left multiplication by a contraction does not
increase a singular value: this follows by
approximating \(\mathcal E_L\) by operators of
rank less than \(j\). The output rank is at most
\(4M^4\), so
\[
 \sum_j s_j(Q_{L,M}\mathcal E_L)^4
 \le\sum_{j=1}^{4M^4}\frac1j
 \le1+\log(4M^4).
\]
Since \(\log2<1\), the last bound is at most
\(4m+3\). Isometries and zero extensions preserve
the nonzero singular values.
\end{proof}

\subsection{The inverse-scale factor in the actual quadratic stencil}

Let \(E_M^{[2]}\) be the quadratic homogeneous
coefficient of the balanced step from side \(M\)
to side \(M/2\). With exactly the ordered-tree
conventions of V14 and Shell V11, it is
\begin{equation}\label{eq:v15-stencil}
 E_{M,r,\mu}^{[2]}(u)
  =\frac1{16}\sum_s\left\{
    \frac12[u_{2r+s,\mu},u_{2r+s+e_\mu,\mu}]
   +\frac12[\psi_{r,s}+\psi_{r+e_\mu,s},b_{r,s,\mu}]
   -\frac12[\psi_{r,s},\psi_{r+e_\mu,s}]\right\},
\end{equation}
where \(b_{r,s,\mu}=u_{2r+s,\mu}
+u_{2r+s+e_\mu,\mu}\), and
\(\psi_{r,s}=\theta_{r,s}
-16^{-1}\sum_{s'}\theta_{r,s'}\).
The \(\theta\)'s here are linear ordered-tree
path sums, not logarithms evaluated at a
nonzero background.

Define a one-shell quadratic function of
physical cell values by
\begin{equation}\label{eq:v15-f}
 f_{M,\mu}(u)
   =\frac{M}{2}\operatorname{avg}_r
                           E_{M,r,\mu}^{[2]}(u/M)
   =\frac1{2M}\operatorname{avg}_r
                           E_{M,r,\mu}^{[2]}(u).
\end{equation}
The factor \(M/2\) is exact: every remaining
linear block to the one-cell output multiplies
the spatial mean by its blocking length.

\begin{lemma}[A uniform local quadratic-form bound]
\label{lem:v15-local}
For real Lie-valued physical inputs,
\begin{equation}\label{eq:v15-local}
 |\operatorname{avg}_rE_{M,r,\mu}^{[2]}(u)|
                 \le56\|u\|_{\mathcal H_M\otimes\mathbb R^3}^2.
\end{equation}
There is a real skew operator
\(R_{M,\mu}:\mathcal H_M\to\mathcal H_M\)
such that
\begin{equation}\label{eq:v15-R}
 -g\cdot D^2f_{M,\mu}(0)[U,V]
     =2\langle U,(R_{M,\mu}\otimes\mathcal J_g)V\rangle,
 \qquad \|R_{M,\mu}\|_{\rm op}\le\frac{28}{M}.
\end{equation}
This is the scalar spatial bank in the weighted
physical inner product; \(\mathcal J_gb=g\times b\).
\end{lemma}

\begin{proof}
Give the pairs \((r,s)\) normalized counting
measure, with \(s\in\{0,1\}^4\).
Every ordered tree path has length at most four,
and a fixed edge belongs to at most eight of
the sixteen paths from that cell's root.
Cauchy--Schwarz and counting therefore give
\[
 \|\theta\|_{L^2(r,s)}^2\le32\|u\|_{\mathcal H_M}^2.
\]
Subtracting the \(s\)-mean is an orthogonal
projection. Thus \(\|\psi\|_2\le\sqrt{32}\|u\|\),
and translation \(r\mapsto r+e_\mu\) preserves
that bound. Also \(\|b_\mu\|_2\le2\|u\|\).
The bound \(|[a,b]|\le2|a||b|\) now bounds the
three terms in the spatial average of
\eqref{eq:v15-stencil}, respectively, by
\[
 \|u\|^2,\qquad4\sqrt{32}\|u\|^2,\qquad32\|u\|^2.
\]
Their sum is
\((33+16\sqrt2)\|u\|^2<56\|u\|^2\).
This proves \eqref{eq:v15-local}, also for
short-period coincidences of neighbouring cells.

Every term in \(f_{M,\mu}\) is a Lie bracket
of two linear spatial inputs, so it has the
skew-bank representation \eqref{eq:v15-R}.
For a unit colour vector \(g\), equations
\eqref{eq:v15-f} and \eqref{eq:v15-local} imply
\(|g\cdot f_{M,\mu}(U)|\le28\|U\|^2/M\).
The true quadratic Hessian consequently has
operator norm at most \(56/M\). Its matrix
is \(2R_{M,\mu}\otimes\mathcal J_g\), as in
Proposition~\ref{prop:v14-bank}.
Since \(\|\mathcal J_g\|=1\), its norm is
\(2\|R_{M,\mu}\|\). This proves the stated bound.
\end{proof}

\subsection{Hilbert--Schmidt summability without deleting shell pairs}

Work on the scalar continuum energy space
\(\mathcal V\). Choose V13's real energy
isometries colour by colour, and denote
their scalar part by \(J_L^{\rm sc}\).
For \(M=2^m\le L\) put
\[
 X_{L,M}=Q_{L,M}\mathcal E_L(J_L^{\rm sc})^*,
 \qquad
 S_{L,\mu}^{(M)}
       =X_{L,M}^*R_{M,\mu}X_{L,M}.
\]
They are operators on one common space, with
zero extension understood. By the exact
second-chain-rule identity of V14,
\begin{equation}\label{eq:v15-sum}
 S_{L,\mu}=\sum_{\substack{M=2^m\\2\le M\le L}}
                             S_{L,\mu}^{(M)}
\end{equation}
is the embedded scalar bank whose matrix in
energy-orthonormal finite coordinates is
\(C_L^{1/2}T_{L,\mu}C_L^{1/2}\).
This identifies the operators with the
actual source, not just with an operator
having the same rank estimate.

\begin{lemma}[Elementary trace-ideal estimate]
\label{lem:v15-HS}
For any finite-rank \(X\) and bounded \(R\),
\[
 \|X^*RX\|_{\mathfrak S_2}
            \le\|R\|_{\rm op}
                     \left(\sum_j s_j(X)^4\right)^{1/2}.
\]
Consequently, for \(M=2^m\),
\begin{equation}\label{eq:v15-shell-bound}
 \|S_{L,\mu}^{(M)}\|_{\mathfrak S_2}
        \le a_m:=28\,2^{-m}\sqrt{4m+3}.
\end{equation}
\end{lemma}
\begin{proof}
In singular vector coordinates for \(X\),
the squared Hilbert--Schmidt norm is
\(\sum_{i,j}\sigma_i^2\sigma_j^2|R_{ij}|^2\).
Use \(2\sigma_i^2\sigma_j^2\le\sigma_i^4+\sigma_j^4\).
Each row and each column of a matrix of a
bounded operator has squared sum at most
\(\|R\|^2\). Summing gives the first inequality.
Apply Lemmas~\ref{lem:v15-embedding} and
\ref{lem:v15-local} to obtain
\eqref{eq:v15-shell-bound}.
\end{proof}

For an integer \(m\ge0\), write
\begin{equation}\label{eq:v15-tail}
 \delta_m=28\,2^{-m}\bigl(\sqrt{4m+3}+4\bigr),
 \qquad B_*=168.
\end{equation}
Since
\(\sqrt{4(m+r)+3}\le\sqrt{4m+3}+2r\) for \(r\ge1\),
the identities
\(\sum_{r\ge1}2^{-r}=1\) and
\(\sum_{r\ge1}r2^{-r}=2\) imply
\[
 \sum_{j>m}a_j\le\delta_m,\qquad
                     \sum_{j\ge1}a_j<B_*.
\]
In particular the actual fine-shell tail is
\(O(2^{-m}\sqrt{m+1})\) in Hilbert--Schmidt
norm, uniformly in the finer lattice \(L\).

\begin{theorem}[Cofinal convergence of the actual cubic constraint coefficient]
\label{thm:v15-cubic}
For each orientation, there is a real skew
Hilbert--Schmidt operator \(S_{\infty,\mu}\)
on \(\mathcal V\) such that
\begin{equation}\label{eq:v15-S-limit}
 S_{L,\mu}\longrightarrow S_{\infty,\mu}
                       \quad\hbox{in }\mathfrak S_2.
\end{equation}
Every finite-lattice and limiting norm is
at most \(B_*\). Their stopped sums through
resolution \(2^m\) have error at most
\(\delta_m\).

The embedded actual cubic constraint
operators therefore satisfy
\begin{equation}\label{eq:v15-K-limit}
 \widehat K_{L,[3]}(v)
       =2\sum_\mu S_{L,\mu}\otimes\mathcal J_{g_\mu(v)}
       \longrightarrow
 K_{\infty,[3]}(v)
       =2\sum_\mu S_{\infty,\mu}\otimes\mathcal J_{g_\mu(v)}
              \quad\hbox{in }\mathfrak S_2.
\end{equation}
The convergence is locally uniform in complex source
variables and in every fixed derivative of this
homogeneous polynomial. The limiting operator is
the cubic coefficient of the actual continuum
relative operator of V13.
\end{theorem}

\begin{proof}
First fix \(M\). In the companion's continuum
notation, the finite-dimensional average
approximating \(Q_{L,M}\mathcal E_L(J_L^{\rm sc})^*\)
is \(M\mathsf A_M\mathcal R_L\), where
\(\mathcal R_L=I_L(J_L^{\rm sc})^*\).
The operator-norm quadrature comparison and the
finite-average convergence in
Theorem~\ref{thm:v13-source} give
\[
 X_{L,M}\longrightarrow
              X_{\infty,M}:=M\mathsf A_M\iota
                 \quad\hbox{in operator norm}
\]
into the fixed finite-dimensional space
\(\mathcal H_M\). All norms on that fixed
output space are equivalent. The resulting
operators \(X_{L,M}^*R_{M,\mu}X_{L,M}\) converge
in operator norm, and their differences have
rank at most \(8M^4\). Hence they converge
in Hilbert--Schmidt norm. Call the limit
\(S_{\infty,\mu}^{(M)}\). It obeys the same
bound \(a_m\).

Both finite and limiting sums of these
shell operators have tails bounded by
\(\delta_m\). Fix \(m\), pass to the limit
in its finitely many shell terms, and
then let \(m\) tend to infinity. This proves
\eqref{eq:v15-S-limit}, with the asserted
norm and stopped-tail bounds.
No rate in the original mesh \(L\) is inferred
from this two-limit argument.

The identification \eqref{eq:v15-sum} follows
because, at a quadratic insertion on side \(M\),
the incoming fine variation is
\(B_{L/M}w=M^{-1}Q_{L,M}\mathcal E_Lw\);
the remaining linear blocks give precisely
the factor in \eqref{eq:v15-f}.
Thus the sum is exactly the quadratic
coefficient of the actual lattice source.
V13 already identifies that coefficient
in operator norm with the second derivative
of its continuum source. Uniqueness of an
operator-norm limit identifies the stronger
Hilbert--Schmidt limit obtained here.

V14's exact formula for \(K_{L,[3]}\) now
proves \eqref{eq:v15-K-limit}.
There are only finitely many degree-three
source monomials; coefficient convergence
therefore implies locally uniform complex
convergence and convergence of all their
fixed derivatives. This last assertion
concerns the cubic polynomial, not the
full finite-amplitude family \(K_L(v)\).
\end{proof}

\begin{theorem}[The sixth-degree quadratic trace needs no divergent subtraction]
\label{thm:v15-bubble}
The actual Gram matrices in V14 converge:
\begin{equation}\label{eq:v15-D-limit}
 D_{L,\mu\nu}\longrightarrow
 D_{\infty,\mu\nu}
   =-\Tr(S_{\infty,\mu}S_{\infty,\nu})
   =\langle S_{\infty,\mu},S_{\infty,\nu}\rangle_{\mathfrak S_2}.
\end{equation}
The limiting matrix is positive semidefinite.
For finite or limiting stopped sums through \(2^m\),
each Gram-entry error is at most \(2B_*\delta_m\).
Moreover the complete double shell expansion is
absolutely convergent, since
\begin{equation}\label{eq:v15-cross-absolute}
 \sum_{r,s\ge1}
  \left|\Tr\bigl(S_{\infty,\mu}^{(2^r)}
                         S_{\infty,\nu}^{(2^s)}\bigr)\right|
                       \le\left(\sum_{r\ge1}a_r\right)^2<B_*^2.
\end{equation}
Thus the sextic quadratic-trace contribution has the
finite, actual limit
\begin{equation}\label{eq:v15-bubble-limit}
 -\frac14\Tr\widehat K_{L,[3]}(v)^2
   \longrightarrow
 -2\sum_{\mu,\nu}D_{\infty,\mu\nu}\,
                               g_\mu(v)\cdot g_\nu(v).
\end{equation}
No first-trace or reference-determinant subtraction
is used to establish this particular convergence.
\end{theorem}

\begin{proof}
For Hilbert--Schmidt operators,
\(|\Tr(U^*V)|\le\|U\|_{\mathfrak S_2}\|V\|_{\mathfrak S_2}\),
as follows directly from the sum of their matrix
entries in an orthonormal basis. Skewness turns
\(-\Tr(S_\mu S_\nu)\) into that inner product.
Apply this inequality to the difference of two
Gram entries to obtain convergence and the error
bound \(2B_*\delta_m\). The Gram description
proves positivity. Apply it term by term with
the bounds \(a_r,a_s\) to prove
\eqref{eq:v15-cross-absolute}.
Finally use the full colour contraction already
proved in \eqref{eq:v14-bubble}. All limits are
of homogeneous degree-six polynomials, so the
corresponding fixed source derivatives converge.
\end{proof}

Combining this result with V14 gives the more
precise outstanding part of the sixth-degree
relative determinant:
\[
 J^6\left[\tfrac12\log\det(I+K_L)\right]
 =\tfrac12\sum_{r=4}^{6}\Tr K_{L,[r]}
                         -\tfrac14\Tr K_{L,[3]}^2.
\]
The last term is now matched cofinally. The
three homogeneous components of the first
trace are not matched by this theorem.
Hilbert--Schmidt membership alone does not
define an ordinary first trace of
\(K_{\infty,[3]}\); the exact finite-colour
cancellation in V14 is not being used to
assert that this limiting operator is
trace class.
Nor is the full operator \(K_L(v)\) asserted
to converge in \(\mathfrak S_2\). Positivity
of \(D_\infty\) is not here strengthened to
strict positivity by extrapolating V14's
positive definite \(D_2\).

\subsection{An exact first two-shell cancellation}

Absolute convergence did not require any
mixed shell to vanish. Nevertheless the
first pair has additional structure which
is useful to record rather than suppress.

On side \(4\), let \(C=C_4\) be the scalar
free covariance on zero-mean Coulomb links,
extended by zero; it is a \(1024\)-by-\(1024\)
matrix in unweighted scalar link coordinates.
Let \(B=B_2\) be the actual linear block to
side \(2\), a \(64\)-by-\(1024\) matrix.
Let \(T_\mu^{\rm f}\) be the side-\(4\)
quadratic insertion followed by the final
linear block. The other insertion has bank
\[
 T_\nu^{\rm c}=B^tT_{2,\nu}B.
\]
Put \(W=CB^t\).

\begin{proposition}[Certified actual two-shell cancellation]
\label{prop:v15-two-shell}
The following full matrix identities hold over
the rationals:
\begin{equation}\label{eq:v15-Wzero}
             W^tT_\mu^{\rm f}W=0,\qquad \mu=1,2,3,4.
\end{equation}
In particular every mixed Gram entry between
these two insertions is zero:
\[
 -\Tr(CT_\mu^{\rm f}CT_\nu^{\rm c})=0.
\]
This is a theorem for these actual side-\(4\)
matrices, not an assumption about all dyadic
shell pairs or about nonzero backgrounds.
\end{proposition}

\begin{proof}
We give a finite integer certificate, including
the bounds that turn modular zero into rational
zero. This is a computer-assisted finite part
of the argument; the preceding cofinal theorem
does not depend on it.

All input matrices have explicit entries.
Write \(\omega=i\),
\(q_\mu(k)=i^{k_\mu}-1\), and
\(\lambda(k)=\sum_\mu|q_\mu(k)|^2\). Then
\begin{equation}\label{eq:v15-C4}
 C_{(x,\mu),(y,\nu)}
  =\frac1{256}\sum_{k\in(\mathbb Z/4)^4\setminus\{0\}}
    i^{k\cdot(x-y)}
       \left(\frac{\delta_{\mu\nu}}{\lambda(k)}
                    -\frac{q_\mu(k)\overline{q_\nu(k)}}{\lambda(k)^2}\right).
\end{equation}
Pairing conjugate frequencies makes every entry
real and rational. The entries of \(B\) are
the multiplicities in its literal two-link
path average, divided by \(16\).
For the fine insertion use scalar coordinate
vectors \(e_{x,\mu}\) and scalar path-indicator
vectors \(\psi_{r,s}\), with
\(r,s\in\{0,1\}^4\), \(x=2r+s\), and
\(r'=r+e_\mu\) reduced modulo two. Then
\begin{equation}\label{eq:v15-Tfine}
\begin{split}
 T_\mu^{\rm f}=\frac1{256}\sum_{r,s}\bigl\{
 &e_{x,\mu}\wedge e_{x+e_\mu,\mu}\\
 &+(\psi_{r,s}+\psi_{r',s})
                       \wedge(e_{x,\mu}+e_{x+e_\mu,\mu})
       -\psi_{r,s}\wedge\psi_{r',s}\bigr\}.
\end{split}
\end{equation}
The prefactor is the stencil's \(1/32\)
times the remaining factor \(2\) and the
coarse spatial average \(1/16\).
Thus no abstract averaging bank is being
substituted for the actual source.

The possible nonzero \(\lambda(k)\)'s are
\(2,4,\ldots,16\). A common denominator for
\eqref{eq:v15-C4} is therefore
\[
 d_C=256\operatorname{lcm}(2,4,\ldots,16)^2
          =722534400.
\]
A common denominator for \(W\) is
\(d_W=16d_C=11560550400\).
Each \(\psi\) has denominator dividing \(16\),
so \(d_T=65536\) clears \eqref{eq:v15-Tfine}.
The spectral formula gives \(|C_{ij}|<1/2\):
each projection entry has modulus at most
one, \(\lambda(k)\ge2\), and there are only
255 terms divided by 256.
Each row of \(B\) is nonnegative and has
sum two, whence \(|W_{ia}|\le1\).
Also \(|(T_\mu^{\rm f})_{ij}|\le7\):
in each summand of \eqref{eq:v15-Tfine},
the three wedge entries have bounds
\(1,4,2\), and there are 256 summands.
Consequently every entry \(N\) of the integer
matrix
\[
 d_W^2d_T\,W^tT_\mu^{\rm f}W
\]
has the explicit bound
\begin{equation}\label{eq:v15-CRT-bound}
 |N|\le7\cdot1024^2d_W^2d_T
             =64288738916848002956336824320000.
\end{equation}

The certificate computes all four \(64\)-by-\(64\)
matrices modulo each of the six primes
\[
 1000033,\quad1000037,\quad1000081,\quad
 1000117,\quad1000121,\quad1000133.
\]
Primality is checked by trial division through
the integer square root. Each prime is one
modulo four, so a square root of \(-1\) is
constructed and checked in its field.
This evaluates \eqref{eq:v15-C4} by a valid
homomorphism from its Gaussian-rational
coefficients. All denominators displayed
above are invertible modulo these primes.
Every matrix entry is zero in every field.
The prime product is
\[
 1000522108734456440708408872900813781,
\]
which exceeds twice the bound
\eqref{eq:v15-CRT-bound}. The integer \(N\)
is divisible by this product and lies in
its strict half-interval, so \(N=0\).
This proves \eqref{eq:v15-Wzero} over the
rationals, not merely at modular sample values.
Finally cyclicity of the finite trace gives
\[
 \Tr(CT_\mu^{\rm f}CB^tT_{2,\nu}B)
                    =\Tr(W^tT_\mu^{\rm f}WT_{2,\nu})=0.
\]
\end{proof}

\begin{proposition}[Conditional Gaussian interpretation at this lattice]
\label{prop:v15-conditional}
Let \(A\) be the centred real Gaussian field
with covariance \(C_4\otimes I_3\) on the
actual free detail space. Let \(F_\mu^{\rm f}(A)\)
be the quadratic source insertion with bank
\(T_\mu^{\rm f}\). Then
\[
             \mathbb E[F_\mu^{\rm f}(A)\mid B_2A]=0.
\]
\end{proposition}
\begin{proof}
Finite Gaussian orthogonal projection gives
the conditional mean in
\(\operatorname{ran}(C_4B_2^t)\otimes\mathbb R^3\);
one may use the pseudoinverse of
\(B_2C_4B_2^t\) on its range. Equation
\eqref{eq:v15-Wzero} makes the quadratic
insertion vanish on every such mean.
The remaining centred conditional covariance
is scalar in colour. The expectation of
each Lie bracket under that covariance is
zero by contraction of its antisymmetric
colour tensor with the colour identity.
Expanding the quadratic polynomial around
the conditional mean proves the assertion.
No positivity on the removed zero or
longitudinal modes is needed.
\end{proof}

This gives an actual cancellation mechanism
to investigate at further levels. It does
not license a general Gaussian martingale
claim until the corresponding matrices
have been proved to vanish there. The
summability proof above deliberately retains
every shell pair and remains valid whether
or not further such cancellations hold.

\subsection{Verification, preserved predecessor, and the remaining first trace}

The new exact certificate is
\begin{center}\small
\path{scripts/verify_ym_f14_v15_cross_shell.py}.
\end{center}
Its SHA--256 is
\begin{center}\scriptsize\ttfamily
CB2B004D6B0A1819A10492A6BEDD311CEF89C588F8C6FFC78D5A1FE8424ECB14.
\end{center}
It reproduces every covariance and skew-bank
entry of V14's side-\(2\) calculation.
It verifies the physical averaging row and
weighted-column masses at
\((L,M)=(4,2),(4,1),(8,2),(8,4)\).
It constructs the side-\(4\) covariance on
all 1024 scalar link coordinates and applies
the six-prime certificate to all four
reduced matrices. This includes the entire
2295-dimensional real Lie detail space,
not a selected Fourier family.

All products in the certificate use integer
arithmetic, reducing a matrix product before
the next multiplication. The unreduced
matrix or trace sums have at most 4096
terms, each bounded by \((p-1)^2\), and
the script checks this is below \(2^{63}\)
for every prime. Larger scalar modular
expressions use arbitrary-precision integers.
Thus no floating-point cancellation or
integer overflow is used to obtain zero.
Four independent side-\(4\) field tests
compare the literal nested degree-two BCH
words with the two insertion banks modulo
101; these are supplementary diagnostics,
not the proof of the polynomial identity.
The stencil proof and the exhaustive matrix
certificate have their distinct roles.

The analytic singular-value, local-stencil,
and dyadic-tail estimates establish the
cofinal trace theorem; a finite lattice
check does not establish that limit.
The companion review supplied no additional
analytic bound, and no SM/NS theorem is
silently used as a YM ultraviolet estimate.

The predecessor main V14 has 287,531 bytes
and SHA--256
\begin{center}\scriptsize\ttfamily
E884676F295A3A488D0FF6ADEAECC3028ABBBCDEEEB5369F5F12EDF7916FEAEF.
\end{center}
Its entire body is retained, with only
the title version and the terminal append
changed. Active V15 replaces V14.
The separate Shell-Mixing V13, SM/NS
companions, archives, and monographs are
untouched. The next regular companion
checkpoint is V20.

\begin{firewall}
The matched quantity is the quadratic
trace of the cubic constraint coefficient
on one fixed physical periodic cell.
This resolves a genuine part of the
sixth-degree relative one-loop packet.
It does not resolve degrees four through
six of the first trace, the reference
determinant, the combined gauge/coarea/Haar
amplitude, growing physical volume, higher
loop remainders, compact-source coverage,
continuum quantum existence, or the
physical Hamiltonian mass gap.
The finite Gaussian conditional cancellation
is not a quantum mass-gap estimate.
\end{firewall}

The next upstream calculation is now the
first trace with its full nonlinear
source, chart tangent, and reference
Hessian contributions kept together.
The side-\(4\) conditional cancellation
suggests looking for exact fine/coarse
projection identities before paying for
separate absolute bounds. Any such
identity must be established for the
actual matrices; the newly convergent
quadratic trace cannot stand in for it.
\section{V16: the actual quartic first trace and its Gaussian source response}
\label{sec:v16-quartic}

\subsection{Context and normalization}

V15 proved convergence of the sextic quadratic-trace
contribution. It did not compute the lower-degree
first trace. We now retain all three contributions
to its quartic coefficient: the cubic derivative of
the actual source, the first chart-tangent correction,
and the linear change of the reference Hessian.
The chart contribution cancels at every level.
The reference contribution has an explicit
translation-invariant symbol and at most logarithmic
growth. The remaining source response is evaluated
exactly on the first lattice.

The active files have not changed since V15's
checkpoint. The separate Shell-Mixing V13 remains
the source for the ordered-tree map and its
homogeneous--detail Hessian. SM V72 and NS V26
are preserved; the next compulsory checkpoint
is V20. No new result from a companion is
being inferred from its filename or version.

Continue on the fixed unit periodic cell with
\(L=2^n\), and put \(\iota_Lv=v/L\).
Let \(Z\) be the centred Gaussian on the actual
free detail space with covariance \(C_L\otimes I_3\)
in ordinary link coordinates. Thus its energy
coordinates have covariance the identity.
This is the unit-normalized Wick variable
associated with the free quadratic action,
not a claim that the interacting fine law
is Gaussian. Write
\[
 q_L^{\rm rel}(v)=\frac12\log\det(I+K_L(v)),
 \qquad B_{L,[3]}=K_{L,[3]},\qquad
 H_{L,1}=(H_{0,L})_{[1]}.
\]
The brackets \([r]\) denote homogeneous source
degree \(r\). Define the linear cubic-source
response by
\begin{equation}\label{eq:v16-Gamma}
 (\Gamma_Lv)_\mu
       =\frac12\,\mathbb E\,
          D^3\mathcal F_{L,\mu}(0)[\iota_Lv,Z,Z].
\end{equation}
Global adjoint equivariance and the scalar-colour
Gaussian covariance imply that this is a
\(4\)-by-\(4\) real matrix acting identically
on each of the three colour components.
This can also be seen directly from the
double-bracket contractions below.

\subsection{Keeping the chart tangent in the finite identity}

Let
\[
 \mathcal M_vZ=D^2\mathcal F_L(0)[\iota_Lv,Z],
 \qquad T_1(v)Z=-\iota_L\mathcal M_vZ.
\]
The latter is the first-order term of the
actual chart tangent \(T_L(v)Z\).

\begin{lemma}[Complete quartic first-trace identity]
\label{lem:v16-complete}
At every finite \(L\),
\begin{equation}\label{eq:v16-complete}
\begin{split}
 (q_L^{\rm rel})_{[4]}(v)
 ={}&-g(v)\cdot\Gamma_Lv\\
    &-g(v)\cdot\mathbb E
             D^2\mathcal F_L(0)[T_1(v)Z,Z]
           -\frac12\Tr(H_{L,1}(v)B_{L,[3]}(v)).
\end{split}
\end{equation}
Here \(g=\nabla E_4\) is exactly V14's cubic
multiplier coefficient.
\end{lemma}
\begin{proof}
The constraint differentiated at fixed source
gives
\[
 D_ws(v,w)=-\{D_s\mathcal F_L\}^{-1}
                              D_w\mathcal F_L.
\]
At zero \(D_s\mathcal F_L=I\) and
\(D_w\mathcal F_L=0\). Differentiation along
the first minimizing variation \(\iota_Lv\)
therefore gives \(T_1=-\iota_L\mathcal M_v\).

The relative first trace is
\(\Tr K_L=\Tr(H_{0,L}^{-1}B_L^{\rm c})\).
Since \(B_L^{\rm c}\) starts in degree three,
its quartic coefficient is
\(\Tr B_{L,[4]}^{\rm c}-\Tr(H_{L,1}B_{L,[3]})\).
The logarithm's powers at least two start
in degree six, so the quartic term of
\(q_L^{\rm rel}\) is half this expression.
Differentiate the exact form
\[
 B_L^{\rm c}(v)[U,W]
  =-\zeta_L(v)\cdot
      D^2\mathcal F_L(a_L^*(v))[T_L(v)U,T_L(v)W].
\]
Its degree-four trace consists of the
third source derivative in \eqref{eq:v16-Gamma},
the two equal tangent terms, and
\(-\zeta_{L,[4]}\cdot
  \mathbb E D^2\mathcal F_L(0)[Z,Z]\).
The last expectation is zero: its Lie
brackets contract against a scalar-colour
covariance. Taking the Gaussian expectation
is exactly taking the energy trace.
The factors \(1/2\) now give
\eqref{eq:v16-complete}. No amplitude
Jacobian is hidden in this determinant
identity.
\end{proof}

\begin{theorem}[The quartic chart-feedback trace vanishes at every level]
\label{thm:v16-tangent}
For every dyadic \(L\),
\begin{equation}\label{eq:v16-tangent}
       \mathbb E D^2\mathcal F_L(0)[T_1(v)Z,Z]=0.
\end{equation}
\end{theorem}
\begin{proof}
Use the inherited exact mixed root Hessian
from Shell V11. With
\[
 f_L(x)=x-(L-1)/2,\qquad
 M_{\lambda,\mu}Z
       =L^{-4}\sum_x f_L(x_\lambda)Z_\mu(x),
\]
it says
\[
 (\mathcal M_vZ)_\mu
       =\sum_{\lambda\ne\mu}[v_\lambda,M_{\lambda,\mu}Z].
\]
Substituting \(T_1=-\iota_L\mathcal M_vZ\)
into the same mixed Hessian, its \(\mu\)-th
component is a sum of terms
\[
 -[[v_\lambda,M_{\lambda,\nu}Z],
                            M_{\nu,\mu}Z],
 \qquad \nu\ne\mu,\quad\lambda\ne\nu.
\]
The linear functional \(M_{\lambda,\nu}\)
has Fourier support on the nonzero
\(\lambda\)-coordinate axis only. The
zero frequency is absent because \(f_L\)
has mean zero. The functional
\(M_{\nu,\mu}\) has support on the nonzero
\(\nu\)-axis. Since \(\lambda\ne\nu\),
these supports are disjoint, including
their negatives. The free Coulomb
covariance is translation invariant,
so every cross-covariance of these
two linear Gaussian vectors is zero.
The expectation of each displayed
bilinear bracket is consequently zero.
This proves \eqref{eq:v16-tangent}
on the full detail space, without
dropping a chart derivative.
\end{proof}

This cancellation is for the quartic
trace at the zero-source covariance.
It does not say that the tangent itself
vanishes, or that its higher-degree
traces vanish.

\subsection{The projected Wilson reference symbol}

For a frequency \(\theta=2\pi k/L\), set
\[
 q_\mu=e^{i\theta_\mu}-1,\quad
 \lambda=\sum_\mu|q_\mu|^2,\quad
 \Pi=I-\frac{qq^*}{\lambda}
                 \quad(k\ne0).
\]
On the scalar transverse energy coordinates
define \(D_{L,\nu}\) by the multiplier
\begin{equation}\label{eq:v16-Dsymbol}
       D_{L,\nu}(k)
       =-\frac{4i\sin\theta_\nu}{L\lambda(k)}\,\Pi(k).
\end{equation}
It is a real skew operator when conjugate
Fourier modes are recombined.

\begin{theorem}[Actual linear reference Hessian]
\label{thm:v16-reference}
The reference derivative in
\eqref{eq:v16-complete} is
\begin{equation}\label{eq:v16-reference}
 H_{L,1}(v)=\sum_\nu D_{L,\nu}\otimes\mathcal J_{v_\nu}.
\end{equation}
In particular \(H_{2,1}(v)=0\).
For the scalar source bank \(S_{L,\mu}\)
of V15, put
\[
 h_{L,\nu\mu}=\Tr(D_{L,\nu}S_{L,\mu}).
\]
Then
\begin{equation}\label{eq:v16-q4}
 (q_L^{\rm rel})_{[4]}(v)
   =-\sum_{\mu,\nu}(\Gamma_L)_{\mu\nu}\,
                          g_\mu(v)\cdot v_\nu
        +2\sum_{\mu,\nu}h_{L,\nu\mu}\,
                          g_\mu(v)\cdot v_\nu.
\end{equation}
\end{theorem}
\begin{proof}
For a plaquette in directions \(\mu<\nu\),
write \(dU=\Delta_\mu U_\nu-\Delta_\nu U_\mu\)
and \(\operatorname{av}_\mu U=(U+U(\cdot+e_\mu))/2\).
The second-degree BCH expansion of its four
ordered link factors, with a constant
background \(a=v/L\), gives the mixed
logarithm coefficient
\[
 k_aU=[a_\mu,\operatorname{av}_\mu U_\nu]
       -[a_\nu,\operatorname{av}_\nu U_\mu]
                  +\frac12[a_\mu+a_\nu,dU].
\]
For example, the coefficients multiplying
the four signed link variations before
grouping are
\[
 a_\mu/2,\quad a_\mu+a_\nu/2,\quad
 a_\mu/2+a_\nu,\quad a_\nu/2.
\]
They give the displayed formula with
both shifts retained.

The action is
\(\frac12|\log Q_p|^2+O(|\log Q_p|^4)\)
at the identity. Since the constant
background has no first-order curvature,
the linear reference quadratic form is
\(2\sum_p\langle dU,k_aU\rangle\).
The last term in \(k_aU\) makes no
contribution, by skewness of the bracket.
The first chart-tangent correction is
constant, and the free kinetic Hessian
annihilates constants; it contributes
nothing to \(H_{L,1}\).

Let \(P_\nu\) be the scalar link-to-plaquette
map which inserts the averaged \(\nu\)-leg:
it is \(\operatorname{av}_\nu U_\mu\) with
the appropriate wedge sign. The Euclidean
linear Hessian has scalar skew matrix
\[
 \frac2L A_\nu,\qquad
                    A_\nu=d^*P_\nu-P_\nu^*d.
\]
Its Fourier projection can be computed
without choosing a transverse basis.
Put \(b_\nu=(1+e^{i\theta_\nu})/2\). Then
\[
 (A_\nu)_{ij}
 =(\overline q_\nu b_\nu-q_\nu\overline b_\nu)\delta_{ij}
       -b_\nu\delta_{i\nu}\overline q_j
                         +\overline b_\nu q_i\delta_{j\nu}.
\]
Left and right multiplication by \(\Pi\)
remove the last two terms, and
\[
 \Pi A_\nu\Pi
       =(\overline q_\nu b_\nu-q_\nu\overline b_\nu)\Pi
       =-2i\sin\theta_\nu\,\Pi.
\]
Whitening by the kinetic covariance
\(\Pi/\lambda\), and multiplying by \(2/L\),
proves \eqref{eq:v16-Dsymbol} and
\eqref{eq:v16-reference}. At \(L=2\)
every frequency component is zero or
\(\pi\), so all these operators vanish.

Finally
\(B_{L,[3]}=2\sum_\mu S_{L,\mu}\otimes\mathcal J_{g_\mu}\).
Using \(\Tr(\mathcal J_a\mathcal J_b)=-2a\cdot b\)
gives
\[
 \Tr(H_{L,1}B_{L,[3]})
             =-4\sum_{\mu,\nu}h_{L,\nu\mu}\,
                                      g_\mu\cdot v_\nu.
\]
Combine this with Lemma~\ref{lem:v16-complete}
and Theorem~\ref{thm:v16-tangent}.
\end{proof}

\begin{proposition}[At most logarithmic growth of this reference term]
\label{prop:v16-log}
For \(L=2^n\),
\begin{equation}\label{eq:v16-log}
                |h_{L,\nu\mu}|\le448\sqrt2\,n.
\end{equation}
This bounds only the reference term in
\eqref{eq:v16-q4}, not the whole matrix
\(\Gamma_L\) or the full quartic coefficient.
\end{proposition}
\begin{proof}
In scalar energy coordinates put
\(T=|\mathcal E_L|\), with eigenvalues
\(t_j\le j^{-1/4}\) from V15.
Formula \eqref{eq:v16-Dsymbol} factors as
\[
 D_{L,\nu}=4U_\nu T,\qquad
 U_\nu(k)=-i\frac{\sin\theta_\nu}{\sqrt{\lambda(k)}}.
\]
The operators \(U_\nu,T\) commute and
\(\|U_\nu\|\le1\).
For the actual shell of side \(M\), write
\(Q_{L,M}\mathcal E_L=YT\) by the polar
decomposition of \(\mathcal E_L\).
Then \(\|Y\|\le1\), \(\operatorname{rank}Y\le4M^4\),
and its scalar source bank is
\[
 S_{L,\mu}^{(M)}=TY^*R_{M,\mu}YT,\qquad
                         \|R_{M,\mu}\|\le28/M.
\]
Cyclicity of the finite trace and
Cauchy--Schwarz for Hilbert--Schmidt
matrices yield
\[
 |\Tr(D_{L,\nu}S_{L,\mu}^{(M)})|
   \le4\|R_{M,\mu}\|\,\Tr(YT^3Y^*).
\]
For completeness, use \(A=YT^{3/2}\) to
bound \(|\Tr(RAU_\nu A^*)|\) by
\(\|R\|\|A\|_{\mathfrak S_2}^2\).
Since \(0\le Y^*Y\le I\) and its rank
is at most \(N=4M^4\), diagonalizing
\(T\) bounds the last positive trace
by its \(N\) largest eigenvalues:
\[
 \Tr(YT^3Y^*)\le\sum_{j=1}^N t_j^3
       \le\sum_{j=1}^N j^{-3/4}
       \le4N^{1/4}=4\sqrt2\,M.
\]
Thus each shell contributes at most
\(448\sqrt2\), including its actual
linear averages and transverse projection.
There are \(n\) shells, proving
\eqref{eq:v16-log}. This bound does not
identify a logarithmic coefficient,
and does not assert that logarithmic
growth actually occurs.
\end{proof}

\subsection{An exact Gaussian cubic-jet calculus for the literal source}

We now evaluate \(\Gamma_2\). It suffices
to take the external background along
one fixed colour vector \(E_1\), since
the response is scalar in colour.
For a Lie logarithm built from the
literal word, record the following data:
\[
 c,\quad r=(r_j),\quad d=(d_j),\quad e.
\]
Its first-degree term is
\(tcE_1+\epsilon\sum_jr_jZ_j\).
Its coefficient of \(t\epsilon\) is
\(\sum_jd_j(E_1\times Z_j)\).
The expectation of its coefficient of
\(t\epsilon^2\) is \(eE_1\).
The omitted degree-two pure-noise term
is a sum of brackets and has expectation
zero. Write \(C(r,s)=r^tC_Ls\).

\begin{lemma}[Closed exact expectation rule through the required degree]
\label{lem:v16-BCH}
For two such logarithms \(a,b\), their
BCH product has data
\begin{equation}\label{eq:v16-BCH}
\begin{aligned}
 c'&=c_a+c_b,\qquad r'=r_a+r_b,\\
 d'&=d_a+d_b+c_ar_b-c_br_a,\\
 e'&=e_a+e_b+2\{C(r_a,d_b)-C(r_b,d_a)\}\\
   &\quad+\frac23\{
           C(r_a,r_b)(c_a+c_b)
                -C(r_a,r_a)c_b-C(r_b,r_b)c_a\}.
\end{aligned}
\end{equation}
Inversion negates all four records, and
averaging averages them. This rule
computes the desired Gaussian contraction
exactly, without sampling a Gaussian field.
\end{lemma}
\begin{proof}
Use
\[
 \operatorname{BCH}(a,b)
 =a+b+\tfrac12[a,b]
       +\tfrac1{12}\{[a,[a,b]]+[b,[b,a]]\}+O(4).
\]
The coefficient of \(t\epsilon\) in the
half bracket is
\(c_a(E_1\times Z(r_b))
 -c_b(E_1\times Z(r_a))\), giving \(d'\).
Isotropic Gaussian contraction gives
\[
 \mathbb E[Z(r),E_1\times Z(d)]=4C(r,d)E_1.
\]
Thus the terms involving a degree-one
and a degree-two logarithm give
the two terms with coefficient \(2\)
in \(e'\). Pairing a constant with
the mean-zero pure-noise quadratic
part gives zero; no missing quadratic
part has been presumed identically zero.

The expectation of the \(t\epsilon^2\)
coefficient in \([a_1,[a_1,b_1]]\) is
\[
 8\{C(r_a,r_b)c_a-C(r_a,r_a)c_b\}E_1.
\]
The other nested bracket exchanges
\(a,b\). Multiplying by \(1/12\) proves
the remaining terms in \eqref{eq:v16-BCH}.
Higher total degrees cannot contribute
to \(t\epsilon^2\). All first-degree
noise maps remain scalar in colour
under the literal linear combinations,
so these are exactly the statistics
needed at each step. Inversion and
averaging follow from the corresponding
operations on logarithms.
\end{proof}

\begin{theorem}[Complete first-lattice quartic relative determinant]
\label{thm:v16-L2}
For the side-\(2\) actual root, with
the ascending ordered tree,
\begin{equation}\label{eq:v16-Gamma2}
 \Gamma_2=\frac1{9216}
 \begin{pmatrix}
 -5369&0&0&0\\
 0&-4937&144&144\\
 0&0&-4265&336\\
 0&0&0&-3185
 \end{pmatrix}.
\end{equation}
The complete quartic term of the
relative determinant is therefore
\[
       (q_2^{\rm rel})_{[4]}(v)
                      =-\sum_{\mu,\nu}
                         (\Gamma_2)_{\mu\nu}g_\mu(v)\cdot v_\nu.
\]
For \(v_1=xE_1,\ v_2=yE_2,\ v_3=v_4=0\),
this is
\begin{equation}\label{eq:v16-two-colour}
                   \frac{5153}{1152}x^2y^2.
\end{equation}
It is the relative constraint determinant
coefficient, not the complete measured
one-loop coefficient.
\end{theorem}
\begin{proof}
Use the full rational covariance
\eqref{eq:v14-C2}. For background
orientation \(\nu\), each initial link
logarithm has
\[
 c_{x,\mu}=\tfrac12\delta_{\mu\nu},
 \quad r_{x,\mu}=e_{x,\mu},\quad d=0,\quad e=0.
\]
Apply \eqref{eq:v16-BCH} along every
ascending tree path to form its logarithm.
Form each transported two-link path
by the same BCH rule with the inverse
endpoint tree. Next form the distinguished-root
average
\[
 \operatorname{BCH}\left(p_0,
     \frac1{16}\sum_s\operatorname{BCH}(-p_0,p_s)\right),
\]
and conjugate by the negative mean tree
logarithm and its inverse exactly as in
the balanced source. These operations,
including every third-degree expectation,
are implemented over rational numbers
by the certificate below.

The four final \(e\)-columns, in the
order \(\nu=1,2,3,4\), are respectively
\[
 \frac1{9216}
 \begin{pmatrix}-5369\\0\\0\\0\end{pmatrix},\quad
 \frac1{9216}
 \begin{pmatrix}0\\-4937\\0\\0\end{pmatrix},\quad
 \frac1{9216}
 \begin{pmatrix}0\\144\\-4265\\0\end{pmatrix},\quad
 \frac1{9216}
 \begin{pmatrix}0\\144\\336\\-3185\end{pmatrix}.
\]
The coefficient of \(t\epsilon^2\) in
\(\mathcal F_L(t\iota_Lv+\epsilon Z)\)
is exactly
\(\tfrac12D^3\mathcal F_L(0)[\iota_Lv,Z,Z]\).
Thus the listed columns, not twice
those columns, are \(\Gamma_2\).
The covariance contraction includes all
45 scalar transverse coordinates and
all three colours.

Theorem~\ref{thm:v16-tangent} removes the
tangent term, and
Theorem~\ref{thm:v16-reference} gives
\(H_{2,1}=0\). Hence
\eqref{eq:v16-q4} proves the quartic
identity. For the two-colour source,
\(g_1=4xy^2E_1\) and \(g_2=4x^2yE_2\);
their inner products with the opposite
colour vanish. The coefficient is
\[
 -4\{(\Gamma_2)_{11}+(\Gamma_2)_{22}\}
       =4(5369+4937)/9216=5153/1152.
\]
\end{proof}

The upper-triangular off-diagonal
entries in \eqref{eq:v16-Gamma2}
belong to the specified representative
and ordered tree. They must not be
discarded by imposing an orientation
symmetry which that representative
has not been proved to possess.
Nor does their occurrence indicate
a physical violation of gauge symmetry:
the rest of the measured factors has
not been combined with this relative
piece.

\subsection{Verification and the next upstream acquisition}

The exact certificate is
\begin{center}\small
\path{scripts/verify_ym_f14_v16_quartic_trace.py}.
\end{center}
Its SHA--256 is
\begin{center}\scriptsize\ttfamily
FCC1140443F9E78A966E3176757B0DC3D30C7568845140795B07F381178B76C4.
\end{center}
It evaluates the literal Gaussian cubic
response in all four background orientations
using rational arithmetic and the entire
free covariance. It reproduces the
inherited homogeneous--detail derivative
in every fine link coefficient and
checks the conjugation calibration
\[
 \mathbb E\bigl(e^{\epsilon\operatorname{ad}Z(r)}
                                      tE_1\bigr)_{t\epsilon^2}
                             =-4C(r,r)E_1.
\]
It also checks inversion, the exact
matrix \eqref{eq:v16-Gamma2}, the
coefficient \eqref{eq:v16-two-colour},
and the full side-\(2\) tangent-feedback
covariances. An independent Gaussian-rational
Fourier calculation verifies 1020 projected
reference symbols on the side-\(4\)
frequency grid, including 60 checks
of the side-\(2\) vanishing symbols.
The analytic derivation of the symbol,
the disjoint-axis proof, and the
logarithmic bound apply to all levels;
their validity is not inferred from
the finite grid.

The predecessor main V15 has 313,651
bytes and SHA--256
\begin{center}\scriptsize\ttfamily
A449B085AEFE27A5E61941327D15A90A3A67608878019BBA64C32AE7239454BC.
\end{center}
Its full body is preserved, apart
from the title version and this
terminal insertion. Active V16
replaces V15. Companion files,
archives, and monographs remain
untouched; the next scheduled
companion review is V20.

\begin{firewall}
The cancellation of the quartic
tangent-feedback trace is cofinal.
The reference trace is controlled
only by a logarithmic upper bound.
The cubic-source response matrix
is fully evaluated here only at
\(L=2\). No cofinal quartic
renormalization coefficient follows
from this finite matrix, and no
full measured density, quantum
continuum construction, or physical
mass gap is claimed.
\end{firewall}

The remaining quartic bottleneck is
now the explicit pair
\(\Gamma_L\) and \(h_L\) in
\eqref{eq:v16-q4}, with the tangent
feedback removed by proof. The
Gaussian cubic rule gives a direct
way to study the actual source
response across successive shells.
The next task is its scale dependence
and its cancellation or matching
with the reference and density terms,
not another estimate on the already
matched sextic quadratic trace.
\section{V17: logarithmic growth of the actual cubic source response}
\label{sec:v17-growth}

\subsection{What is newly controlled}

V16 expressed the quartic relative determinant in terms
of the source-response matrix \(\Gamma_L\) and the
reference trace \(h_L\). Only the latter had an
all-level growth bound. We now prove
\(\|\Gamma_L\|=O(\log L)\), retaining the contributions
formed by quadratic insertions at different levels.
Consequently the full quartic relative determinant
has at most logarithmic growth on the fixed physical
periodic cell. This is not an evaluation of its
logarithmic coefficient.

The active main predecessor is V16, and the separate
Shell-Mixing V13, SM V72, and NS V26 are unchanged.
The next regular companion checkpoint remains V20.
The new ingredient is a row estimate for an actual
quadratic mixed insertion after all subsequent
linear averages. It avoids an extra factor equal
to the number of intervening scales.

All nested fields used in this section mean
Taylor polynomials through total degree three
at the identity. Their coefficients can be
contracted against a finite Gaussian without
requiring that Gaussian field to lie in the
nonlinear identity chart. We do not evaluate
the full logarithmic root map on unbounded
Gaussian samples.

\subsection{A dual Sobolev estimate for the same fine Gaussian}

Choose a fixed enlarged component constant \(C_S\)
in the inherited discrete Sobolev inequality:
\begin{equation}\label{eq:v17-Sobolev}
 \|W\|_{\ell^4}\le C_S\|d_LW\|_{\ell^2}
\end{equation}
on real mean-zero Coulomb link fields, including
all four orientations and three Lie colours.
Counting norms in this section are unnormalized.
This constant is independent of the fine side \(L\).
Let \(Z\) have the same zero-source covariance
\(C_L\otimes I_3\) as in V16.

\begin{lemma}[Variance from a spatial row]
\label{lem:v17-row-variance}
For any deterministic Lie-valued scalar row \(f\),
\begin{equation}\label{eq:v17-row-variance}
 \mathbb E|\langle f,Z\rangle|^2
                  \le C_S^2\|f\|_{\ell^{4/3}}^2.
\end{equation}
No mean-zero or Coulomb condition on \(f\) is required.
\end{lemma}
\begin{proof}
The variance is the squared norm of the functional
\(W\mapsto\langle f,W\rangle\) on the free detail
energy space. Holder's inequality and
\eqref{eq:v17-Sobolev} bound that norm by
\(C_S\|f\|_{\ell^{4/3}}\).
The covariance is zero on the removed subspace;
no inverse is assigned to longitudinal or
constant modes.
\end{proof}

For the actual composed linear block
\[
 (B_pU)_{r,\mu}
    =p^{-4}\sum_{z\in\{0,\ldots,p-1\}^4}
                  \sum_{t=0}^{p-1}U_{pr+z+te_\mu,\mu},
\]
one scalar row has nonnegative coefficients,
sum \(p\), maximum coefficient at most \(p^{-3}\),
and support on at most \(2p^4\) links of its
one orientation. Indeed a given fine link
can occur for at most \(p\) longitudinal choices.
These assertions remain valid with periodic
identifications. Hence its \(\ell^{4/3}\)
row norm is at most \(2^{3/4}\).
If
\[
       Y_M=B_{L/M}Z,\qquad M\mid L,
       \qquad \sigma_0=\sqrt3\,2^{3/4}C_S,
\]
then every Lie-valued output link satisfies
\begin{equation}\label{eq:v17-Y}
              \bigl(\mathbb E|Y_{M,r,\mu}|^2\bigr)^{1/2}
                                      \le\sigma_0.
\end{equation}
All \(Y_M\)'s here come from one and the same
fine field \(Z\). They have not been replaced
by independent Gaussian fields at successive
levels.
The intermediate fields need not be Coulomb:
no extra projection is inserted between
the actual linear blocks.

\subsection{An insertion retains a bounded dual row after later averages}

Let
\[
 \mathcal L_{M,c}U
       =D^2\mathfrak T_M(0)[c,U]
\]
for a constant Lie tuple \(c=(c_1,\ldots,c_4)\).
This is the mixed derivative of the actual
balanced step from side \(M\) to side \(M/2\).

\begin{lemma}[Actual local mixed-row structure]
\label{lem:v17-local-row}
Each output component can be written
\begin{equation}\label{eq:v17-local-row}
 (\mathcal L_{M,c}U)_{r,\mu}
      =\sum_{\eta\in\{0,1,2,3\}^4}\sum_\alpha
                 A_{\mu\alpha,\eta}(c)\,U_{2r+\eta,\alpha},
 \qquad
 \sum_{\eta,\alpha}\|A_{\mu\alpha,\eta}(c)\|_{\rm op}
                                                \le54|c|.
\end{equation}
Repeated periodic sites are combined only after this
finite offset representation.
\end{lemma}
\begin{proof}
Differentiate the actual quadratic stencil
\eqref{eq:v15-stencil}. For a constant leg,
\(\psi_{r,s}(c)=\sum_\nu(s_\nu-1/2)c_\nu\),
independently of \(r\), and its norm is at
most \(|c|\).
For the variable leg, each tree row has
sum of absolute coefficients at most four;
subtracting its mean gives a bound eight
for a \(\psi\)-row.

The differentiated two-link BCH term has
row norm at most \(2|c|\).
The term with the constant \(\psi\) and
variable two-link sum has bound \(4|c|\).
The term with the two variable \(\psi\)'s
and the constant two-link sum has bound
\(32|c|\). The last, endpoint-\(\psi\)
term has bound \(16|c|\).
These follow directly from
\(\|\operatorname{ad}_a\|\le2|a|\);
their sum is \(54|c|\).
The initial cell tree, its neighbouring
endpoint tree, and the two-link paths
use only offsets with each coordinate
between zero and three. This proves
\eqref{eq:v17-local-row}, with no
restriction on surrounding volume.
\end{proof}

Define the numerical constant
\[
             C_0=46656\sqrt2<2^{17}.
\]

\begin{lemma}[Two averages surrounding a genuine insertion]
\label{lem:v17-sandwich}
Suppose the fine side is \(L=pM\), and
\(q\mid M/2\). For every unit output
colour vector, every row of
\begin{equation}\label{eq:v17-sandwich}
                       B_q\mathcal L_{M,c}B_p
\end{equation}
has \(\ell^{4/3}\) norm at most \(C_0|c|\),
uniformly in \(p,q,M\).
\end{lemma}
\begin{proof}
First fix one offset and input orientation
in \eqref{eq:v17-local-row}. A row of
\(B_q\) has maximum coefficient \(q^{-3}\),
and a row of \(B_p\) has maximum coefficient
\(p^{-3}\).
For a fixed fine link, the origins of a
\(B_p\) row whose support contains that
link have at most two possibilities:
its \(p\)-cell and the preceding cell
in its own orientation. The map from
the intermediate output site \(r\) to
\(2r+\eta\) is injective on the appropriate
periodic lattice. Thus summing over those
intermediate sites contributes at most
two possible rows for a fixed offset.
After summing the local coefficient norms,
each fine row coefficient has norm at most
\[
                         108|c|(pq)^{-3}.
\]

In unrolled integer coordinates, the
outer average ranges over a box of
length at most \(2q\) in each direction.
Decimation by two, the offset at most
three, and the inner \(p\)-average
place the resulting fine support in
a box of side at most \(6pq\).
One can check the largest coordinate
directly as
\(2p(2q-2)+3p+(2p-2)<6pq\).
Periodic folding cannot increase the
number of distinct fine sites.
Including the four possible input
orientations gives at most \(4(6pq)^4\)
links in this support.

For a unit output colour vector, take
the transpose of each coefficient
matrix against it to obtain a scalar
functional on Lie-valued input fields.
The maximum and support bounds give
\[
 \|f\|_{\ell^{4/3}}
 \le108|c|(pq)^{-3}\{4(6pq)^4\}^{3/4}
             =46656\sqrt2\,|c|.
\]
No covariance signs, independence
assumption, or continuum Green-kernel
estimate enters this row calculation.
\end{proof}

\subsection{The mixed response is uniformly controlled at each output level}

Let \(\mathcal F_{L\to M}\) be the actual nested
balanced map from side \(L\) to side \(M\).
For formal parameters \(t,\epsilon\), write the
relevant parts of its Taylor expansion as
\begin{equation}\label{eq:v17-partial}
\begin{split}
 \mathcal F_{L\to M}(t\,v/L+\epsilon Z)
   ={}&t\,c_M+\epsilon Y_M+t\epsilon D_M
           +\epsilon^2Q_M+t\epsilon^2R_M+\cdots,\\
 &c_M=v/M,\qquad Y_M=B_{L/M}Z .
\end{split}
\end{equation}
The coefficient \(D_M\) is linear in both
\(v\) and \(Z\); it is not set to zero
after blocking. The omitted coefficients
cannot contribute to \(t\epsilon^2\).
Initially \(D_L=Q_L=R_L=0\).

\begin{theorem}[Uniform Gaussian mixed response]
\label{thm:v17-mixed}
Put \(\sigma_1=\sqrt3\,C_SC_0\).
Every output link at every \(M\mid L\)
satisfies
\begin{equation}\label{eq:v17-mixed}
          \bigl(\mathbb E|D_{M,r,\mu}|^2\bigr)^{1/2}
                                      \le\sigma_1\,|v|/M.
\end{equation}
This bound has no factor \(\log(L/M)\).
\end{theorem}
\begin{proof}
The exact quadratic chain rule gives
\begin{equation}\label{eq:v17-Dsum}
 D_M=\sum_{\substack{K=2^j\\M<K\le L}}
      B_{K/(2M)}\mathcal L_{K,v/K}B_{L/K}Z.
\end{equation}
All factors except the single quadratic
insertion are actual linear blocks.
For the term indexed by \(K\), apply
Lemma~\ref{lem:v17-sandwich} with
\(p=L/K\) and \(q=K/(2M)\).
For any unit output
colour vector, its combined fine row
has dual norm at most
\[
              C_0|v|\sum_{K>M}\frac1K
                                      \le C_0|v|/M.
\]
The sum is geometric over dyadic \(K\).
It is the norm of a single combined
row, not a sum of variances of
independent shells.
Lemma~\ref{lem:v17-row-variance} bounds
its variance; summing three colour
coordinate variances proves
\eqref{eq:v17-mixed}.
\end{proof}

\subsection{Explicit primitive cubic majorants}

We need only the finite polynomial
coefficients of the actual local
word, not a radius estimate for a
Gaussian field. A scalar triple
\((a_1,a_2,a_3)\) majorizes the
absolute Lie-polynomial coefficients
through degree three, with each
input link assigned its norm.
For a BCH product, a valid rule is
\begin{equation}\label{eq:v17-majorant}
\begin{aligned}
 b_1&=a_1+c_1,\\
 b_2&=a_2+c_2+a_1c_1,\\
 b_3&=a_3+c_3+a_1c_2+a_2c_1
                    +\tfrac13(a_1^2c_1+c_1^2a_1).
\end{aligned}
\end{equation}
It follows from the degree-three BCH
formula and \(|[a,b]|\le2|a||b|\).
Inversion leaves an absolute majorant
unchanged; averaging does not increase
one common majorant.

Applying this rule gives the following
deliberately unoptimized bounds:
\[
\begin{array}{c|rrr}
 \text{logarithm}&a_1&a_2&a_3\\ \hline
 \text{tree, at most four links}&4&6&32/3\\
 \text{transported path, at most ten links}&10&45&230\\
 \text{root-relative word}&20&190&6080/3\\
 \text{root average}&30&435&21170/3\\
 \text{balanced output}&38&703&43586/3
\end{array}
\]
For the distinguished path, the ten-link
bound overestimates its actual shorter
length, which is harmless for an upper
bound. The negative mean tree logarithm
on each side uses the same tree bound.
Thus the table bounds the literal
balanced source, not a simplified
untransported average.

\begin{lemma}[The two primitive expectations needed by the recursion]
\label{lem:v17-primitive}
Suppose the Lie random variables at
every input link of a primitive step
obey
\[
       \|Y_i\|_{L^2(\mathbb P)}\le u,\qquad
       \|D_i\|_{L^2(\mathbb P)}\le d.
\]
No independence of these variables is
required. If \(E_M^{[2]},E_M^{[3]}\)
are its quadratic and cubic homogeneous
coefficients and \(c\) is constant, then
at every output link,
\begin{equation}\label{eq:v17-primitive}
\begin{aligned}
 \left|\mathbb E D^2E_M^{[2]}(0)[Y,D]\right|
                                      &\le1406\,ud,\\
 \left|\mathbb E
           \{E_M^{[3]}(tc+\epsilon Y)\}_{t\epsilon^2}\right|
                                      &\le43586\,|c|u^2.
\end{aligned}
\end{equation}
\end{lemma}
\begin{proof}
The coefficient of the mixed product
in \(703(su+td)^2\) is \(1406ud\).
The coefficient of \(t\epsilon^2\)
in \((43586/3)(t|c|+\epsilon u)^3\)
is \(43586|c|u^2\).
These scalar coefficients bound sums
of norms of the actual multilinear
monomials by the positive majorant
construction. In each monomial with
two random legs, use
\(\mathbb E|Y_i||D_j|\le ud\) or
\(\mathbb E|Y_i||Y_j|\le u^2\).
This proves the inequalities even
when the two legs are correlated
or are the same variable.
\end{proof}

\subsection{The full cubic response has only logarithmic growth}

\begin{theorem}[Actual all-level source-response bound]
\label{thm:v17-Gamma}
Let
\begin{equation}\label{eq:v17-CGamma}
 C_\Gamma=1406\,\sigma_0\sigma_1
                                  +43586\,\sigma_0^2,
\end{equation}
with \(\sigma_0,\sigma_1\) defined above.
For every \(L=2^n\), the actual response
matrix \(\Gamma_L\) of V16 satisfies
\begin{equation}\label{eq:v17-Gamma}
                         \|\Gamma_L\|_{\rm op}\le C_\Gamma n.
\end{equation}
The bound includes every pair of
quadratic insertions in the cubic
source coefficient.
\end{theorem}
\begin{proof}
First \(\mathbb E Q_M=0\) at every
intermediate level. Indeed \(Q_L=0\),
and
\[
             Q_{M/2}=B_2Q_M+E_M^{[2]}(Y_M).
\]
The covariance of \(Y_M\) is scalar
in colour, and each term of its
quadratic source polynomial is a
Lie bracket. Its expectation is
zero. This proves the assertion
by induction.

Now extract \(t\epsilon^2\) in the
one-step composition of
\eqref{eq:v17-partial}. Exactly,
\begin{equation}\label{eq:v17-Rrecursion}
\begin{split}
 R_{M/2}={}&B_2R_M
       +D^2E_M^{[2]}(0)[Y_M,D_M]
       +D^2E_M^{[2]}(0)[c_M,Q_M]\\
       &+\{E_M^{[3]}(tc_M+\epsilon Y_M)\}_{t\epsilon^2}.
\end{split}
\end{equation}
The \(c_M,Q_M\) term has zero expectation
because \(c_M\) is deterministic and
\(\mathbb E Q_M=0\). The \(Y_M,D_M\)
term is retained. By
\eqref{eq:v17-Y},
Theorem~\ref{thm:v17-mixed}, and
Lemma~\ref{lem:v17-primitive}, the
expected newly generated term at
each output link has norm at most
\[
               C_\Gamma |v|/M.
\]
All subsequent linear blocks to the
one-cell output multiply the spatial
mean by \(M/2\). Thus this shell's
contribution to each output orientation
has norm at most \(C_\Gamma|v|/2\).
Combining four orientations bounds
the full coarse tuple by \(C_\Gamma|v|\).
There are exactly \(n\) such shells.

Finally the expected \(t\epsilon^2\)
coefficient at side one is
\[
 \mathbb E R_1
      =\frac12\mathbb E D^3\mathcal F_L(0)
                                      [v/L,Z,Z]=\Gamma_Lv.
\]
This proves \eqref{eq:v17-Gamma}.
Equation \eqref{eq:v17-Dsum}, used
inside \eqref{eq:v17-Rrecursion},
is the full mixed two-insertion
part of the third chain rule.
No pair of shell levels has been
silently removed or assigned a
fresh covariance.
\end{proof}

\begin{theorem}[Logarithmic quartic relative determinant bound]
\label{thm:v17-quartic}
For every real source \(v\) and every
\(L=2^n\), as an inequality for the
homogeneous quartic polynomial,
\begin{equation}\label{eq:v17-quartic}
 \left|(q_L^{\rm rel})_{[4]}(v)\right|
       \le4(C_\Gamma+3584\sqrt2)\,n\,|v|^4.
\end{equation}
It is not asserted that the quotient
by \(n\) has a limit.
\end{theorem}
\begin{proof}
In V16's exact formula
\eqref{eq:v16-q4}, the source term
is bounded by \(C_\Gamma n|g(v)||v|\).
The reference entry bound
\eqref{eq:v16-log} gives
\[
 \left|2\sum_{\mu,\nu}h_{L,\nu\mu}
                              g_\mu\cdot v_\nu\right|
 \le896\sqrt2\,n
                    \left(\sum_\mu|g_\mu|\right)
                    \left(\sum_\nu|v_\nu|\right)
 \le3584\sqrt2\,n\,|g(v)||v|.
\]
Each positive semidefinite matrix
\(|v_\nu|^2I-v_\nu v_\nu^t\)
has norm at most \(|v_\nu|^2\).
The explicit formula for \(g\)
therefore gives \(|g(v)|\le4|v|^3\).
Combining these inequalities proves
\eqref{eq:v17-quartic}.
\end{proof}

This result rules out power growth
of this particular quartic relative
coefficient along refinement. It
does not identify the logarithmic
coefficient, prove a Cesaro limit
of shell contributions, or specify
the counterterm of the full measured
law. The constants are finite
analytic majorants, not useful
numerical estimates of a physical
coupling or mass.

\subsection{Exact checks and preserved history}

The new certificate is
\begin{center}\small
\path{scripts/verify_ym_f14_v17_source_growth.py}.
\end{center}
Its SHA--256 is
\begin{center}\scriptsize\ttfamily
AD9D20804D83BE085F10FED19445BFB9BA1C01F1AAE96A130B090599A013FCFF.
\end{center}
It computes the complete rational
primitive majorant table, including
\(2\cdot703=1406\) and
\(3\cdot43586/3=43586\).
It checks the actual mixed stencil
rows and their offset bounds.
Eight composed-row examples with
\((p,q)=(1,1),(2,1),(1,2),(2,2)\)
and final side two test both
transverse and same-orientation
backgrounds; the fine sides are
4, 8, 8, and 16. Their largest
supports have 17,664 and 16,896
nonzero links, respectively. The
coefficient envelopes are checked
exactly, not estimated by fitting
a decay exponent.

The checker also reproduces all
sixteen side-\(2\) homogeneous--detail
rows inherited from Shell V11.
An auxiliary rational incoming
linear Gaussian response is fed
through the literal side-\(2\)
cubic-jet calculus and independently
through its exact quadratic skew
banks. The four expected mixed
feeds agree and are nonzero:
\[
 -1/1344,\quad-169/8064,\quad
                   379/32256,\quad727/64512.
\]
This is a diagnostic input for
the recurrence, not a replacement
for the actual \(D_M\) in the
all-level proof. It checks that
the \(Y_M,D_M\) contribution is
retained. Separate tests check
linear transport of a nonzero
incoming cubic expectation.

The all-level conclusions follow
from the row-support proof,
the dual Sobolev inequality,
the exact third-degree recursion,
and the finite coefficient
majorants. The finite checker
does not establish a continuum
limit by extrapolation.

The predecessor V16 has 332,325
bytes and SHA--256
\begin{center}\scriptsize\ttfamily
E9C1D8AD8F04D915ACD299102D2C6CB48791C82FBB4FA400AEF33D223BEB4B3A.
\end{center}
Its entire body is retained,
apart from the title version
and the terminal insertion.
Active V17 replaces V16.
The companion files, archives,
and monographs remain untouched.
The next scheduled companion
review is V20.

\begin{firewall}
The new theorem controls the
growth of an actual finite
Taylor coefficient under
refinement on one fixed
periodic physical cell.
It does not prove convergence
of the normalized coefficient,
the full one-loop density,
the remaining fifth- and
sixth-degree first traces,
uniform higher-loop control,
arbitrary compact-source
coverage, continuum quantum
existence, or a physical
Hamiltonian mass gap.
\end{firewall}

The next upstream question is
the value and stability of
the actual shell contributions,
not another absolute bound on
the same quartic polynomial.
The bounded normalized matrix
families are now available for
a genuine coefficient or
Cesaro-matching analysis, but
their convergence still has
to be proved and combined
with the other measured
factors.
\section{V18: the actual leading logarithmic relative coefficient}
\label{sec:v18-bulk}

\subsection{From a growth bound to a coefficient limit}

V17 proved a logarithmic upper bound, not convergence
of the normalized quartic relative determinant.
This version proves that its leading logarithmic
coefficient exists. Both constituent matrices,
\(\Gamma_{2^n}/n\) and \(h_{2^n}/n\), converge.
The limits are specified by convergent
four-dimensional Fourier integral formulas;
their numerical values are not evaluated here.
They are coefficients of the relative constraint
determinant, not yet the full measured one-loop
coefficient or a physical beta function.

The main predecessor and all companions are
unchanged since V17. The next SM/NS checkpoint
remains V20. The added Shell-Mixing V13 still
supplies the literal block, not an instruction
to change the requested goal.

Write
\[
                 L=2^n=pM,\qquad M=2^m,\qquad p=2^{n-m}.
\]
The physical cell is still fixed. Coordinates
are now magnified around a shell cell: the
fine spacing becomes \(1/p\), and the
magnified torus has side \(M\). Letting
both \(p\) and \(M\) grow describes a shell
far from the fine and coarse endpoints.
It is not a thermodynamic construction
of the interacting Yang--Mills law.

\subsection{The two free bilinear forms in magnified coordinates}

Use centred frequency representatives
\(\xi=2\pi k/M\), with
\(|\xi_\alpha|\le\pi p\), and define
\[
 q_p(\xi)_\alpha=e^{i\xi_\alpha/p}-1,\qquad
 \omega_p(\xi)^2=p^2|q_p(\xi)|^2,\qquad
 \Pi_p(\xi)=I-\frac{q_p(\xi)q_p(\xi)^*}{|q_p(\xi)|^2}.
\]
The zero frequency is omitted. Put
\[
 \Pi(\xi)=I-\frac{\xi\xi^t}{|\xi|^2}
                          \quad(\xi\ne0).
\]
For real compactly supported spatial
four-vector kernels \(f,g\in L^1\cap L^2\),
with Fourier transform
\(\widehat f(\xi)=\int e^{-i\xi\cdot y}f(y)\,dy\),
define
\begin{equation}\label{eq:v18-forms}
\begin{aligned}
 \mathcal C(f,g)
   &=(2\pi)^{-4}\int_{\mathbb R^4}
       \widehat f(\xi)^*
           \frac{\Pi(\xi)}{|\xi|^2}\widehat g(\xi)\,d\xi,\\
 \mathcal R_\nu(f,g)
   &=(2\pi)^{-4}\int_{\mathbb R^4}
       \widehat f(\xi)^*
          \frac{-4i\xi_\nu\Pi(\xi)}{|\xi|^4}
                                      \widehat g(\xi)\,d\xi .
\end{aligned}
\end{equation}
The first is a positive covariance form.
The second is a real skew bilinear
response form, not a positive covariance.
Both integrals are absolutely convergent:
near zero use bounded Fourier transforms
and the powers \(|\xi|^{-2}\), \(|\xi|^{-3}\);
at infinity use Plancherel and bounded
inverse powers.

Suppose \(f_p,g_p\) are cell-constant
kernels of mesh \(1/p\), supported in a
fixed box once unfolded, converging to
\(f,g\) in \(L^1\) and \(L^2\), with
uniform bounds in these norms.
Let \(f_{p,x}\) be the value on the
cell at \(x/p\), and let the corresponding
fine row be \(r_f(x)=p^{-3}f_{p,x}\).
Define \(r_g\) in the same way.
Let \(C_L\) be the actual scalar free
Coulomb covariance, and let
\[
 K_{L,\nu}=C_L^{1/2}D_{L,\nu}C_L^{1/2}.
\]
Its ordinary fine-link multiplier is
\[
 K_{L,\nu}(k)
       =-\frac{4i\sin(2\pi k_\nu/L)}
                    {L\,|q_p(\xi)|^4}\Pi_p(\xi).
\]
This is the covariance-weighted skew
reference operator of V16.

\begin{theorem}[Joint fine-mesh and large-box bilinear limit]
\label{thm:v18-free}
As \(p\to\infty\) and \(M\to\infty\),
without a relation between their rates,
\begin{equation}\label{eq:v18-free}
 r_f^tC_Lr_g\longrightarrow\mathcal C(f,g),
 \qquad
 M r_f^tK_{L,\nu}r_g
                            \longrightarrow\mathcal R_\nu(f,g).
\end{equation}
The assertions hold componentwise for
any fixed finite family of matrix-valued
colour kernels as well.
\end{theorem}
\begin{proof}
Put
\[
 F_p^d(\xi)=p^{-4}\sum_x
                  f_{p,x}e^{-i\xi\cdot x/p},
\]
and define \(G_p^d\) similarly.
The exact discrete Fourier formulas
for the two left sides are
\begin{equation}\label{eq:v18-discrete-forms}
\begin{aligned}
 M^{-4}\sum_{k\ne0}
       (F_p^d)^*\frac{\Pi_p}{\omega_p^2}G_p^d,\qquad
 M^{-4}\sum_{k\ne0}
       (F_p^d)^*
       \frac{-4ip\sin(\xi_\nu/p)\Pi_p}{\omega_p^4}G_p^d.
\end{aligned}
\end{equation}
To check the scaling, the two row
factors give \(p^{-6}\), the fine
Fourier normalization is \(L^{-4}\),
and the two transformed sums give
\(p^8\). Thus the covariance weight
is \(M^{-4}/(p^2|q_p|^2)\).
The extra factor \(M\) in the
reference form gives exactly its
second multiplier in
\eqref{eq:v18-discrete-forms}.

The sine chord bound gives
\(\omega_p(\xi)\ge2|\xi|/\pi\).
Also \(p|\sin(\xi_\nu/p)|\le|\xi|\).
The two multiplier norms are
therefore bounded by
\(C|\xi|^{-\rho}\), with
\(\rho=2,3\), respectively.
The transform bounds and Parseval
identity are
\[
 \|F_p^d\|_\infty\le\|f_p\|_1,\qquad
 M^{-4}\sum_k|F_p^d(\xi)|^2=\|f_p\|_2^2.
\]
For \(|\xi|>R\), Cauchy--Schwarz
therefore bounds the omitted part
by \(CR^{-\rho}\|f_p\|_2\|g_p\|_2\).

For \(0<|\xi|<\varepsilon\), the
four-dimensional lattice count gives
\[
 M^{-4}\sum_{0<|\xi|<\varepsilon}|\xi|^{-\rho}
                                  \le C\varepsilon^{4-\rho}.
\]
Indeed
\(\sum_{0<|k|\le a}|k|^{-\rho}\le Ca^{4-\rho}\)
for \(a\ge1\), by dyadic annuli;
when \(a<1\) there are no such modes.
This yields a uniform infrared
bound using the \(L^1\) norms.
The same infrared and ultraviolet
bounds hold for the continuum
integrals.

On a fixed annulus
\(\varepsilon\le|\xi|\le R\),
the piecewise-constant Fourier
transform differs from \(F_p^d\)
by at most \(C_Rp^{-1}\|f_p\|_1\):
integrate the change of phase
within each cell. Hence \(F_p^d\)
converges uniformly there to
\(\widehat f\), by \(L^1\)
convergence, and similarly for \(g\).
The lattice projectors and both
multipliers converge uniformly
to those in \eqref{eq:v18-forms}.
The remaining sum is a Riemann
sum with cell volume \((2\pi/M)^4\);
its prefactor is \(M^{-4}\).
Its limit is the stated integral
with factor \((2\pi)^{-4}\).
First send \(p,M\) to infinity
on the annulus, then send
\(\varepsilon\) to zero and
\(R\) to infinity.
All conjugate and self-conjugate
lattice modes were included;
the high-frequency bound, rather
than deletion of Nyquist modes,
controls the ultraviolet end.
\end{proof}

\subsection{The actual averaging filters}

For dyadic \(s\ge1\) and a link
\(a=(r,\alpha)\) on the grid of
spacing \(1/s\), define
\begin{equation}\label{eq:v18-filter}
 f_{s,a}(y)
  =s^3\int_0^1
      1_{\,r/s+[0,1/s)^4+t e_\alpha/s}(y)\,dt\ e_\alpha .
\end{equation}
This is the box average followed
by the actual longitudinal average.
It has integral \(e_\alpha/s\).
If \(\beta(z)=\int_0^1e^{-izt}\,dt\),
its Fourier transform is explicitly
\begin{equation}\label{eq:v18-filter-Fourier}
 \widehat f_{s,a}(\xi)
   =s^{-1}e^{-i\xi\cdot r/s}
         \beta(\xi_\alpha/s)
                   \prod_{\rho=1}^4\beta(\xi_\rho/s)\ e_\alpha.
\end{equation}
No additional coarse Coulomb projector
is applied to this filter.

\begin{lemma}[Actual discrete rows converge to these filters]
\label{lem:v18-filters}
For fixed \(s\mid p\), reconstruct the
row of \(B_{p/s}\) as \(p^3\) times
its fine coefficients on cells of
side \(1/p\). In an unfolded local
box it converges to
\eqref{eq:v18-filter} in \(L^1\)
and \(L^2\).
The same is true after any fixed
finite composition of local mixed
stencils and subsequent linear
averages.
\end{lemma}
\begin{proof}
For \(s=1\), the kernel is the
unit box indicator in the three
transverse coordinates. In its
longitudinal coordinate it is a
step approximation to
\(\tau(y)=\max(0,\min(y,2-y))\).
On \([j/p,(j+1)/p)\) its value
is the number of pairs
\((z,t)\in\{0,\ldots,p-1\}^2\)
with \(z+t=j\), divided by \(p\).
Direct integration of these
linear tent pieces gives
\[
 \|f_{p,1}-f_{1,a}\|_1=p^{-1},
 \qquad
 \|f_{p,1}-f_{1,a}\|_2^2=2/(3p^2).
\]
Translation only changes the
location of the kernel.
For general \(s\), the kernel
is \(s^3\) times this construction
at \(sy-r\), with ratio \(p/s\).
Its \(L^1\) error is still \(p^{-1}\),
and its squared \(L^2\) error is
\(2s^4/(3p^2)\).

A fixed mixed insertion is a
finite linear combination of
these translated kernels, with
the exact offset coefficients
of V17. A fixed subsequent
linear block is another such
finite combination.
Their convergence follows by
linearity. For a fixed finite
carrier and fixed number of
insertions, sufficiently large
\(M\) removes periodic identifications
within the unfolded support.
\end{proof}

Let \(\mathcal W\) denote an auxiliary
centred Gaussian linear family with
covariance \(\mathcal C\) and independent
Lie colours. It can be constructed
by completing the kernel space in
the covariance seminorm, quotienting
its null space, and using independent
standard Gaussians in an orthonormal
basis. Only its smeared linear
functionals are used here. This
is not an interacting YM field.
Write
\[
             (\mathsf A_s\mathcal W)_a=\mathcal W(f_{s,a}),
                       \qquad Y=\mathsf A_1\mathcal W.
\]
The same family \(\mathcal W\) is
used for every \(s\).

\subsection{The incoming mixed response has a genuine bulk limit}

On the infinite integer grid let
\(\mathcal L_v\) be the actual
one-step mixed stencil
\(D^2\mathfrak T(0)[v,\cdot]\),
with the offset coefficients
of \eqref{eq:v17-local-row}.
For \(s=2^j\), define a unit-grid
Gaussian response field
\begin{equation}\label{eq:v18-delta}
 \Delta_j(v)
     =2^{-j}B_{2^{j-1}}\mathcal L_v
                                 (\mathsf A_{2^j}\mathcal W),
 \qquad
 \Delta^{(J)}(v)=\sum_{j=1}^J\Delta_j(v).
\end{equation}
The intermediate grid has spacing
\(2^{-j}\), then \(2^{1-j}\)
after the mixed step, then one
after the last linear block.
The factor \(2^{-j}\) is not
absorbed into a new covariance.

\begin{theorem}[Bulk joint response, with a geometric incoming tail]
\label{thm:v18-delta}
On each fixed finite unit-grid
carrier, \(\Delta^{(J)}(v)\)
converges in Gaussian \(L^2\)
to a field \(\Delta(v)\).
For every link,
\begin{equation}\label{eq:v18-delta-tail}
 \|\Delta(v)-\Delta^{(J)}(v)\|_{L^2(\mathbb P)}
                         \le\sigma_1\,2^{-J}|v|.
\end{equation}
Moreover, as \(m\to\infty\) and
\(n-m\to\infty\), the covariances
of the joint finite-carrier fields
\[
             (Y_M,\ M D_M(v))
                    \quad\hbox{converge to those of}\quad
                              (Y,\Delta(v)).
\]
This includes all their cross-covariances,
not just separate variance limits.
\end{theorem}
\begin{proof}
Rewrite the exact sum
\eqref{eq:v17-Dsum} as
\[
 M D_M(v)=
   \sum_{j=1}^{n-m}
      2^{-j}B_{2^{j-1}}\mathcal L_v B_{p/2^j}Z.
\]
For any fixed \(J\), all terms
through \(j=J\) are defined when
\(p\ge2^J\). Their rows, together
with those of \(Y_M\), converge
to the filters in
\eqref{eq:v18-delta} by
Lemma~\ref{lem:v18-filters}.
The covariance limit
Theorem~\ref{thm:v18-free} applies
to this entire finite family at
once. It proves the asserted
joint convergence for truncated
responses.

V17's row bound applies to the
single term indexed by \(j\).
After multiplication by \(M\),
its linkwise Gaussian \(L^2\)
norm is at most
\(\sigma_1\,2^{-j}|v|\).
Pass this bound to the continuum
filters by their finite covariance
convergence. The geometric series
proves convergence on the same
Gaussian family \(\mathcal W\)
and \eqref{eq:v18-delta-tail}.
It also bounds the finite-lattice
tails uniformly in \(p,M\).
The linear averages have uniformly
bounded second moments by
\eqref{eq:v17-Y}. Cauchy--Schwarz
then controls the covariance
errors produced by either tail.
First choose \(J\), next pass
to the joint \(p,M\) limit in
the finite family, and finally
let \(J\) grow. This proves
all joint covariance limits.
\end{proof}

\subsection{Explicit integral definitions of the two leading matrices}

Let \(E^{[2]},E^{[3]}\) be the
literal primitive quadratic and
cubic polynomials on the infinite
unit grid, evaluated at output
cell zero. They have finite
carriers, specified by the same
ordered tree and balancing word.
Define
\begin{equation}\label{eq:v18-Gamma-star}
\begin{split}
 (\Gamma_*v)_\mu=\frac12\mathbb E\bigl[
 &D^2E^{[2]}_\mu(0)[Y,\Delta(v)]\\
 &+\{E^{[3]}_\mu(tv+\epsilon Y)\}_{t\epsilon^2}
                                                   \bigr].
\end{split}
\end{equation}
As before, this is a real
\(4\)-by-\(4\) matrix acting
identically in colour.
Every expectation is a contraction
of two Gaussian linear functionals.
Equations \eqref{eq:v18-forms},
\eqref{eq:v18-filter-Fourier}, and
\eqref{eq:v18-delta} therefore
specify it by an absolutely
convergent series of ordinary
four-dimensional integrals.
If \(\Gamma_*^{(J)}\) uses
\(\Delta^{(J)}\), the explicit
error is
\begin{equation}\label{eq:v18-Gamma-tail}
 \|\Gamma_*-\Gamma_*^{(J)}\|_{\rm op}
                    \le1406\,\sigma_0\sigma_1\,2^{-J}.
\end{equation}
Indeed apply the mixed primitive
bound of V17, include the factor
\(1/2\), and combine four output
orientations.

For the reference term, take a
common finite scalar link carrier
\(\Lambda\) for the four primitive
quadratic polynomials. Define
their local skew banks by
\[
       E^{[2]}_\mu(U)
             =\sum_{a<b}(T_\mu^{\rm loc})_{ab}[U_a,U_b].
\]
For \(a=(r,\alpha)\), abbreviate
\(f_a=f_{1,a}\). Define
\begin{equation}\label{eq:v18-h-star}
 (h_*)_{\nu\mu}
    =\frac12\sum_{a,b\in\Lambda}
                  \mathcal R_\nu(f_a,f_b)(T_\mu^{\rm loc})_{ba}.
\end{equation}
This is one finite sum of absolutely
convergent four-dimensional integrals.
The scalar response \(\mathcal R_\nu\)
is used here, not an invented
positive Gaussian covariance.

\subsection{Actual bulk shell limits and the Cesaro passage}

For \(M=2^m\), call the exact
source contribution of that shell
\(G_{n,m}\). By the recursion
\eqref{eq:v17-Rrecursion},
\begin{equation}\label{eq:v18-G-shell}
\begin{split}
 (G_{n,m}v)_\mu=\frac12
   \operatorname{avg}_r\mathbb E\bigl[
 &D^2E_{M,r,\mu}^{[2]}(0)[Y_M,M D_M(v)]\\
 &+\{E_{M,r,\mu}^{[3]}(tv+\epsilon Y_M)\}_{t\epsilon^2}
                                                      \bigr].
\end{split}
\end{equation}
There is an exact decomposition
\[
                   \Gamma_{2^n}=\sum_{m=1}^nG_{n,m}.
\]
The \(Q_M\) term has already been
removed by its proved zero expectation,
not by resetting an incoming field.

\begin{lemma}[The bulk source contribution]
\label{lem:v18-G-bulk}
The matrices \(G_{n,m}\) have norms
at most \(C_\Gamma\), and
\begin{equation}\label{eq:v18-G-bulk}
          G_{n,m}\longrightarrow\Gamma_*
                  \quad\hbox{as }\min(m,n-m)\longrightarrow\infty.
\end{equation}
\end{lemma}
\begin{proof}
The uniform bound is exactly the
per-shell estimate in V17.
The joint law of \(Y_M,D_M\)
is invariant under translations
by one aligned coarse cell:
the original covariance is
translation invariant, the
background is constant, and
the actual nested block commutes
with the corresponding fine
translation. Thus the expectation
in \eqref{eq:v18-G-shell} is
independent of \(r\).

Once \(M\) is large, its primitive
carrier can be unfolded without
periodic coincidences. Truncate
the incoming response at \(J\).
The finite joint covariance
convergence of
Theorem~\ref{thm:v18-delta}
passes the two Gaussian
contractions to their continuum
values. The uniform incoming
tail and the mixed primitive
bound give error at most
\(1406\sigma_0\sigma_1 2^{-J}\)
in the resulting matrix.
Choose \(J\) first and then
send both \(m,n-m\) to infinity.
This proves the double limit,
not merely a limit along one
chosen relation between cutoffs.
\end{proof}

Similarly decompose
\[
 h_{2^n,\nu\mu}
           =\sum_{m=1}^n H_{n,m,\nu\mu},
 \qquad
 H_{n,m,\nu\mu}
           =\Tr(D_{L,\nu}S_{L,\mu}^{(M)}).
\]

\begin{lemma}[The bulk reference contribution with its exact factor]
\label{lem:v18-H-bulk}
Every entry obeys
\(|H_{n,m,\nu\mu}|\le448\sqrt2\), and
\begin{equation}\label{eq:v18-H-bulk}
 H_{n,m,\nu\mu}\longrightarrow(h_*)_{\nu\mu}
                 \quad\hbox{as }\min(m,n-m)\longrightarrow\infty.
\end{equation}
\end{lemma}
\begin{proof}
The bound is the per-shell estimate
in Proposition~\ref{prop:v16-log}.
In ordinary scalar link coordinates
the quadratic insertion bank is
\[
 T_{L,\mu}^{(M)}
  =B_p^t\left\{\frac M2
          \operatorname{avg}_rT_{M,r,\mu}^{\rm loc}\right\}B_p.
\]
It follows by whitening and
cyclicity that
\[
 H_{n,m,\nu\mu}
    =\frac M2\operatorname{avg}_r
          \Tr(B_pK_{L,\nu}B_p^tT_{M,r,\mu}^{\rm loc}).
\]
Translation invariance makes all
terms in this average equal.
For a sufficiently large unfolded
carrier it is therefore
\[
 H_{n,m,\nu\mu}
    =\frac12\sum_{a,b\in\Lambda}
       \{M B_pK_{L,\nu}B_p^t\}_{ab}
                                      (T_\mu^{\rm loc})_{ba}.
\]
The second assertion of
Theorem~\ref{thm:v18-free},
with the actual averaging
filters, proves that the
matrix in braces converges
to \(\mathcal R_\nu(f_a,f_b)\).
This gives \eqref{eq:v18-H-bulk}
with the factor \(1/2\) in
\eqref{eq:v18-h-star}.
\end{proof}

\begin{theorem}[Leading logarithmic matrices and relative determinant]
\label{thm:v18-leading}
The following limits hold in finite-dimensional
matrix norm:
\begin{equation}\label{eq:v18-leading}
           \frac{\Gamma_{2^n}}{n}\longrightarrow\Gamma_*,
             \qquad \frac{h_{2^n}}{n}\longrightarrow h_*.
\end{equation}
Consequently, coefficientwise as homogeneous
quartic polynomials,
\begin{equation}\label{eq:v18-q-leading}
\begin{split}
 \frac{(q_{2^n}^{\rm rel})_{[4]}(v)}n
   \longrightarrow q_{*,4}^{\rm rel}(v)
     ={}&-\sum_{\mu,\nu}(\Gamma_*)_{\mu\nu}
                                  g_\mu(v)\cdot v_\nu\\
        &+2\sum_{\mu,\nu}(h_*)_{\nu\mu}
                                  g_\mu(v)\cdot v_\nu .
\end{split}
\end{equation}
The limits are locally uniform in complex
source variables and their fixed derivatives
as polynomial statements. For normalization
by the natural logarithm \(\log L\), divide
the right side by \(\log2\).
\end{theorem}
\begin{proof}
Use the following elementary two-ended
Cesaro fact. If a bounded triangular
array \(a_{n,m}\), \(1\le m\le n\),
converges to \(a_*\) whenever both
\(m\) and \(n-m\) tend to infinity,
then \(n^{-1}\sum_ma_{n,m}\to a_*\).
Given an error tolerance, choose
an endpoint width \(b\) so that
all terms with \(m\ge b\) and
\(n-m\ge b\) are close to \(a_*\).
There are at most \(2b\)
remaining terms, and their
normalized sum is bounded by
the common bound times \(2b/n\).
First let \(n\) grow, then
let the tolerance shrink.

Apply this to
Lemmas~\ref{lem:v18-G-bulk}
and \ref{lem:v18-H-bulk}.
Their uniform bounds include
the endpoint shells, including
the periodic side-two case
and the finest unsmeared case.
Thus neither endpoint has
been assigned the bulk formula
without proof.
This establishes \eqref{eq:v18-leading}.
Substitute these two limits
in V16's exact identity
\eqref{eq:v16-q4}; \(g\) is
independent of \(L\).
There are finitely many
quartic source monomials,
so coefficient convergence
gives the stated polynomial
and derivative convergence.
\end{proof}

Equivalently,
\((q_{2^n}^{\rm rel})_{[4]}
       =n q_{*,4}^{\rm rel}+o(n)\)
coefficientwise. This proof
does not improve \(o(n)\)
to a bounded remainder.
The formulas identify a
specific, evaluable integral
coefficient, but do not
declare its value equal
to a physical YM beta
coefficient or show it
nonzero.

\subsection{Exact checks, preserved predecessor, and remaining work}

The new certificate is
\begin{center}\small
\path{scripts/verify_ym_f14_v18_log_coefficient.py}.
\end{center}
Its SHA--256 is
\begin{center}\scriptsize\ttfamily
999C540FBECA0B5CE764698655707E154AB4A1F505E4BF419F59B7FCFC3A6DBB.
\end{center}
It verifies the exact tent-kernel
\(L^1\) and squared \(L^2\)
errors at six mesh ratios,
288 Fourier prefactor identities,
64 zero spatial moments of the
actual mixed rows, and the
indices in a two-ended Cesaro
test array and geometric tail.
These finite tests do not
prove the analytic limit.

Thirty-two actual side-four
reference ownership identities
are checked modulo 101, comparing
the full shell trace with its
local translated-cell expression
and its rescaled coarse pullback.
All matrix products use integer
arithmetic with reduction between
products. Even the largest trace
sum is bounded by
\(1024^2\cdot100^2<2^{63}\).
These modular checks are
diagnostics for the exact
translation and trace identities
proved above, not numerical
quadrature of \eqref{eq:v18-forms}.

For example, the actual side-four
fine-shell trace with
\(\nu=\mu=1\) has residue 64
modulo 101. Its rational
denominators are powers of two
and factors of \(2,4,\ldots,16\),
so they are invertible in this
field. That nonzero residue
certifies a nonzero finite
reference contribution and
prevents extrapolating V16's
side-two vanishing to every
lattice. It does not certify
that the limiting integral
\(h_*\) is nonzero.

The predecessor V17 has 350,829
bytes and SHA--256
\begin{center}\scriptsize\ttfamily
6F6B13F9E866F8688C80EF3B6C6B55FDEBB9C751369CCEC3B53FE445014D0AE1.
\end{center}
Its complete body is retained,
apart from the title version
and this terminal insertion.
Active V18 replaces V17.
Companions, archives, and
monographs remain untouched.
The next scheduled SM/NS
checkpoint is V20.

\begin{firewall}
The result is the existence
and an integral identification
of the leading logarithmic
quartic coefficient of one
relative determinant.
The auxiliary Gaussian family
is used only for these
zero-source coefficient
contractions. It is not a
construction of interacting
continuum Yang--Mills theory.
The reference determinant,
combined gauge/coarea/Haar
density, remaining source
degrees, uniform higher-loop
control, compact-source
coverage, and physical
Hamiltonian mass gap remain
unproved.
\end{firewall}

The next acquisition must
use the identified coefficients:
evaluate or simplify their
actual integral combination
and compare it with the
other measured factors.
Existence of a logarithmic
coefficient does not itself
give the measured coupling
flow, a continuum quantum
state, or its mass gap.
\section{V19: source countertransport and the measured cost of removing the first trace}
\label{sec:v19-countertransport}

V18 identified a leading logarithm of the actual quartic relative
determinant. Before treating that number as a coupling flow, one must
determine whether it changes under a renormalization of the retained
source. This version goes upstream to that question. The first relative
trace, at every source degree and every finite regulator, is exactly
the derivative of the classical constrained action along a specified
source vector field. A measured change of retained variables removes
that trace at one-loop order, but creates calculable two-loop terms.
We prove the identity, identify its actual linear vector field using
V16--V18, and retain the resulting density Jacobian.

The new result is not another proof of V15's Hilbert--Schmidt limit.
It explains how that already convergent sextic trace becomes the
entire relative one-loop six-jet in a specified, regulator-dependent
source convention. The convention is not silently substituted for
the prescribed root observable. The source and test observables
must both be transported, and uniform continuum existence remains open.
No full measured normalization or quantum state is obtained here.
The continuum construction and physical Hamiltonian gap in the
\href{https://www.claymath.org/wp-content/uploads/2022/06/yangmills.pdf}
{Jaffe--Witten problem description} remain the target, not this
finite-order change of variables.

\subsection{An actual finite-regulator source vector field}

Fix a dyadic \(L\), work on the real small-source branch of V13,
and use an energy-orthonormal basis \(e_1,\ldots,e_{N_L}\) of the
actual Lie-valued detail space, \(N_L=9(L^4-1)\).
Write \(a_*=a_L^*(v)\), \(T=T_L^f(v)\), and \(H_0=H_{0,L}(v)\).
For a retained component \(\alpha=(\mu,a)\), define the symmetric
matrix \(U_\alpha(v)\) by
\[
 (U_\alpha)_{ij}
   =D^2\mathcal F_{L,\alpha}(a_*)[Te_i,Te_j].
\]
These are matrices on the actual detail space; no extra projection
of a coarse Gaussian field is being introduced. Set
\begin{equation}\label{eq:v19-chi}
 \chi_{L,\alpha}(v)
       =\frac12\Tr\bigl(H_0(v)^{-1}U_\alpha(v)\bigr).
\end{equation}
On the real branch this also equals
\(\frac12\mathbb E D^2\mathcal F_{L,\alpha}(a_*)[T\xi,T\xi]\),
where the finite-dimensional Gaussian \(\xi\) has energy covariance
\(H_0^{-1}\). Only the displayed quadratic derivative is evaluated
on that Gaussian. There is no evaluation of the full nonlinear root
on a continuum distribution.

At fixed \(L\), \(\chi_L\) is real analytic and has a holomorphic
extension after a source shrink. Indeed \(a_*,T,U_\alpha,H_0\)
have those properties, and \(H_0\) is invertible. In complex
coordinates the formula means its analytic matrix trace, not a
complex Gaussian probability law. Changing the energy-orthonormal
basis conjugates both factors and does not change the trace.

\begin{theorem}[The whole first trace is a classical action derivative]
\label{thm:v19-first}
Let \(\mathcal G_L(v)=\mathcal S_L(a_L^*(v))\) and retain the
multiplier sign convention of V12. Then
\begin{equation}\label{eq:v19-first}
 D\mathcal G_L(v)=\zeta_L(v),\qquad
 \frac12\Tr K_L(v)=-D\mathcal G_L(v)[\chi_L(v)].
\end{equation}
Moreover \(\chi_L(0)=0\). Define, in finite dimension,
\begin{equation}\label{eq:v19-remainder}
 \begin{split}
 R_{2,L}(v)
   &=\frac12\{\log\det(I+K_L(v))-\Tr K_L(v)\},\\
 q_L^{\rm rel}(v)
   &=-D\mathcal G_L(v)[\chi_L(v)]+R_{2,L}(v).
 \end{split}
\end{equation}
These are exact local analytic identities, not just quartic identities.
\end{theorem}
\begin{proof}
On the chosen constant-plus-Coulomb-detail coordinate space,
stationarity with the constraint gives
\[
 D\mathcal S_L(a_*)=\zeta_L\cdot D\mathcal F_L(a_*),
 \qquad \mathcal F_L(a_L^*(v))=v.
\]
Differentiate the second identity along the minimizing branch and
insert it into the derivative of \(\mathcal G_L\); this proves
\(D\mathcal G_L=\zeta_L\). The curvature part of the actual Hessian is
\[
 B_L^{\rm c}=-\sum_\alpha\zeta_{L,\alpha}U_\alpha.
\]
Cyclicity of the finite trace gives
\[
 \frac12\Tr K_L
  =\frac12\Tr(H_0^{-1}B_L^{\rm c})
  =-\sum_\alpha\zeta_{L,\alpha}\,
                      \frac12\Tr(H_0^{-1}U_\alpha).
\]
This proves the first-trace identity. At zero the covariance is
the scalar-colour free covariance and the quadratic root is a
sum of Lie brackets. The expectation of each bracket is zero,
so \(\chi_L(0)=0\). Equivalently, equivariance would make
\(\chi_L(0)\) an adjoint-invariant vector in each
\(\mathfrak{su}(2)\) component, hence zero.
Subtracting the finite first trace from the logarithm proves
\eqref{eq:v19-remainder}.
\end{proof}

It is permissible at finite \(L\) to write
\(R_{2,L}=\frac12\log\det_2(I+K_L)\).
This notation does not assert that the full continuum \(K_\infty(v)\)
is Hilbert--Schmidt. V13 proves a full-family \(\mathfrak S_5\)
statement; V15 proves a stronger \(\mathfrak S_2\) statement only
for its cubic source coefficient.

\begin{proposition}[Finite sign and the existing six-jet limit]
\label{prop:v19-six}
On the real branch, if \(\|K_L(v)\|\le\rho<1\), then
\begin{equation}\label{eq:v19-sign}
 -\frac{\Tr K_L(v)^2}{4(1-\rho)}
 \ \le R_{2,L}(v)\le\
 -\frac{\Tr K_L(v)^2}{4(1+\rho)}\le0.
\end{equation}
The relative compensated jets satisfy
\begin{equation}\label{eq:v19-six}
 J^5R_{2,L}=0,\qquad
 J^6R_{2,L}=-\frac14\Tr K_{L,[3]}^2
 \ \longrightarrow\
 -2\sum_{\mu,\nu}D_{\infty,\mu\nu}\,
                         g_\mu(v)\cdot g_\nu(v).
\end{equation}
The convergence is coefficientwise, locally uniformly as a complex
polynomial, and with every fixed derivative.
\end{proposition}
\begin{proof}
The real symmetric matrix \(K_L\) has eigenvalues \(t_j\) in
\([-\rho,\rho]\). For such a scalar,
\[
 \log(1+t)-t=-t^2\int_0^1\frac{u}{1+ut}\,du.
\]
Bound the denominator between \(1-\rho\) and \(1+\rho\),
multiply by \(1/2\), and sum to obtain \eqref{eq:v19-sign}.
Since \(K_L=O(|v|^3)\), the logarithmic series after its first
trace starts with \(-\frac14\Tr K_L^2\). Its sixth-degree term
is precisely \(-\frac14\Tr K_{L,[3]}^2\); all higher powers
start in degree nine. Theorem~\ref{thm:v15-bubble} now proves
\eqref{eq:v19-six}, including the stated modes of polynomial
convergence. Its Gram matrix \(D_\infty\) is positive semidefinite.
\end{proof}

The finite inequality does not give a uniform bound for the full
\(R_{2,L}(v)\): such a conclusion would require a uniform bound
for the full \(\Tr K_L(v)^2\), which is not supplied by V13.

\subsection{The linear countertransport is the actual V18 coefficient}

Use \(h_L\) for the reference-response matrix of V16, not for a
loop parameter, and define four-by-four real matrices
\begin{equation}\label{eq:v19-matrices}
 C_L=\Gamma_L-2h_L^T,\qquad
 A_L=-C_L,\qquad S_L^{\rm or}=\frac12(A_L+A_L^T).
\end{equation}
In this section \(C_L\) is the displayed four-by-four orientation
matrix, not the earlier free scalar-link covariance.
The superscript \(\mathrm{or}\) distinguishes this orientation matrix
from V15's scalar-link operators \(S_{L,\mu}\).

\begin{theorem}[Linearization and its leading logarithm]
\label{thm:v19-linear}
The vector field in \eqref{eq:v19-chi}, not merely its contraction
with \(g\), obeys
\begin{equation}\label{eq:v19-linear}
 \chi_L(v)=C_Lv+O_L(|v|^2).
\end{equation}
The matrices act identically on the three colour components.
For \(L=2^n\),
\begin{equation}\label{eq:v19-matrix-limit}
 \frac{A_{2^n}}{n}\longrightarrow A_*=-\Gamma_*+2h_*^T,
 \qquad
 \|A_{2^n}\|\le(C_\Gamma+3584\sqrt2)n.
\end{equation}
Consequently V18's quartic coefficient has the normal form
\[
 (q_L^{\rm rel})_{[4]}(v)
       =g(v)\cdot A_Lv
       =g(v)\cdot S_L^{\rm or}v.
\]
\end{theorem}
\begin{proof}
At zero \(H_0=I\), \(T\) is the detail inclusion, and
the first variation of the minimizing fine field is \(\iota_Lv=v/L\).
Differentiate \eqref{eq:v19-chi}. With the free Wick variable \(Z\)
of V16, the first two resulting terms are
\[
 \frac12\mathbb E D^3\mathcal F_L(0)[\iota_Lv,Z,Z]
 +\mathbb E D^2\mathcal F_L(0)[T_1(v)Z,Z].
\]
They equal \(\Gamma_Lv\) and zero respectively, by
\eqref{eq:v16-Gamma} and Theorem~\ref{thm:v16-tangent}.
The derivative of \(H_0^{-1}\) supplies the remaining component
\[
             -\frac12\Tr(H_{L,1}(v)U_\alpha(0)).
\]
To evaluate it without losing vector components, let
\(\eta=(\eta_\mu)\) be an arbitrary retained Lie tuple.
The actual quadratic banks of V14 give
\[
 \sum_\alpha\eta_\alpha U_\alpha(0)
       =-2\sum_\mu S_{L,\mu}\otimes\mathcal J_{\eta_\mu}.
\]
Here \(\mathcal J_x y=x\times y\); the Lie bracket is \(2x\times y\).
Insert
\(H_{L,1}=\sum_\nu D_{L,\nu}\otimes\mathcal J_{v_\nu}\)
and \(\Tr(\mathcal J_x\mathcal J_y)=-2x\cdot y\).
The contraction of the inverse-reference derivative with \(\eta\)
is therefore
\[
 -2\sum_{\nu,\mu}h_{L,\nu\mu}\,\eta_\mu\cdot v_\nu.
\]
Since \(\eta\) was arbitrary, this term is \(-2h_L^Tv\).
This proves the full vector identity \eqref{eq:v19-linear};
inferring it only from \(g\cdot\chi_L\) would leave a skew-matrix
ambiguity.

V18 gives the two limits in the definition of \(A_*\).
V17 bounds \(\|\Gamma_L\|\le C_\Gamma n\), while V16 bounds
each entry of \(h_L\) by \(448\sqrt2 n\).
A four-by-four matrix with entries bounded by \(b\) has
operator norm at most \(4b\), proving \eqref{eq:v19-matrix-limit}.
The quartic formula follows either from V16 or from
\eqref{eq:v19-first}, since \(\mathcal G_{L,[4]}=E_4\).
The skew part of \(A_L\) drops out by the next proposition.
\end{proof}

\begin{proposition}[Exact ten-parameter quartic normal form]
\label{prop:v19-normal-form}
Let \(V\) be the four-by-three matrix with rows \(v_\mu\), and put
\(P=VV^T\), \(P_{\mu\nu}=v_\mu\cdot v_\nu\).
For any real four-by-four \(A\), with \(S=(A+A^T)/2\), define
\(q_A(v)=g(v)\cdot Av\). Then
\begin{equation}\label{eq:v19-Gram}
 E_4(v)=(\Tr P)^2-\Tr(P^2),\qquad
 q_A(v)=4\{(\Tr P)\Tr(SP)-\Tr(SP^2)\}.
\end{equation}
The kernel of \(A\mapsto q_A\) consists exactly of the skew
matrices. Its image therefore has dimension ten.
Writing \(S=\frac14(\Tr A)I+S^0\), \(\Tr S^0=0\), gives
\begin{equation}\label{eq:v19-isotropic}
 q_A=(\Tr A)E_4+
       4\{(\Tr P)\Tr(S^0P)-\Tr(S^0P^2)\}.
\end{equation}
The algebraic average of \(q_A\) over all signed permutations
of the four orientations is \((\Tr A)E_4\).
\end{proposition}
\begin{proof}
Using \([x,y]=2x\times y\) in
\(E_4=\frac12\sum_{\mu<\nu}|[v_\mu,v_\nu]|^2\) proves the
first identity. Under \(\delta V=AV\),
\(\delta P=AP+PA^T\). Differentiate the first identity and
use symmetry and cyclicity of the four-dimensional trace;
the answer is the second identity in \eqref{eq:v19-Gram}.
Both \(P\) and \(P^2\) are symmetric, so the skew part disappears.

For the converse, suppose \(q_S=0\) for symmetric \(S\).
Set \(v_i=xE_1\), \(v_j=yE_2\), all other rows zero, with
\(i\ne j\). The coefficient of \(x^2y^2\) is
\(4(S_{ii}+S_{jj})\). All such pair sums vanish, which forces
all four diagonal entries to vanish: three distinct indices
already give this conclusion, and the remaining index follows
from any pair with it. For \(i\ne k\), choose a third index \(j\)
and set \(v_i=xE_1\), \(v_j=yE_2\), \(v_k=zE_1\).
The coefficient of \(xy^2z\) is \(8S_{ik}\); hence it vanishes.
Thus \(S=0\). These tests use actual tuples of three-dimensional
colour vectors, so the rank restriction \(\operatorname{rank}P\le3\)
has not been discarded.

Substitute the trace/traceless decomposition in
\eqref{eq:v19-Gram} to obtain \eqref{eq:v19-isotropic}.
Signed orientation reflections average all off-diagonal entries
of \(A\) to zero, and permutations average its diagonal to
\((\Tr A)I/4\). Since \(E_4\) is orientation-orthogonal invariant
and \(q_A\) is linear in \(A\), this proves the averaging assertion.
\end{proof}

This average is an algebraic projection, not permission to average
the actual path-ordered source or its measured law. The ordered
root prescription can have a nonzero traceless symmetric part.
Nor does the scalar coefficient \(\Tr A_*\) by itself determine a
physical beta function. Uniform rescaling of \(v\) changes \(E_4\)
by its fourth homogeneity degree.

For the actual side-two regulator, V16 gives \(h_2=0\) and hence
\[
 S_2^{\rm or}=\frac1{9216}
 \begin{pmatrix}
 5369&0&0&0\\
 0&4937&-72&-72\\
 0&-72&4265&-168\\
 0&-72&-168&3185
 \end{pmatrix},
 \qquad \Tr A_2=\frac{4439}{2304}.
\]
For example, the \(v_1=xE_1,v_2=yE_2\) coefficient remains
\(5153x^2y^2/1152\), whereas the
\(v_2=xE_1,v_1=yE_2,v_3=zE_1\) coefficient of \(xy^2z\)
is \(-1/16\). Thus even this actual finite coefficient is not
just an isotropic multiple of \(E_4\). These are calibrations,
not values of \(A_*\).

\subsection{A measured source change, with its two-loop payment}

Use \(\varepsilon=\beta^{-1}\) for the loop parameter. At a fixed
regulator, consider a specified local retained density in Lebesgue
coordinates whose stationary expansion, up to source-independent
Gaussian constants, has negative logarithm
\begin{equation}\label{eq:v19-packet}
 \varepsilon^{-1}\mathcal G_L(v)
       +\mathcal Q_L(v)+\varepsilon\mathcal B_L(v)
       +O_L(\varepsilon^2).
\end{equation}
All statements below can instead be read as formal coefficient
identities. When the local Laplace hypotheses of Shell V8 hold,
they are the corresponding finite-regulator asymptotic identities.
They do not assume a regulator-uniform remainder.

Write
\[
 \mathcal Q_L=\mathcal Q_L^{\rm rest}+q_L^{\rm rel}.
\]
For example, if the incoming specified density in the linear
\((s,w)\) coordinates is \(m_L(s,w)\,ds\,dw\), then the
one-loop stationary part not in the relative determinant is
\[
 \mathcal Q_L^{\rm rest}(v)
  =-\log m_L(s_*(v),w_*(v))
       +\log\det Z_L(v)+\frac12\log\det H_{0,L}(v),
\]
up to the same source-independent constants.
This uses V14's inverse source Jacobian exactly once.
The formula leaves \(m_L\), including any required gauge/quotient,
Haar and incoming factors, as the specified actual density.
It does not evaluate those factors or set them to one.

Let \(\phi_t\) be the local flow of the actual vector field
\(\chi_L\):
\[
 \partial_t\phi_t(u)=\chi_L(\phi_t(u)),\qquad \phi_0(u)=u.
\]
It exists for each fixed \(L\) and sufficiently small \(t,u\),
fixes zero, and is locally invertible with inverse \(\phi_{-t}\).
Introduce a new retained variable by
\begin{equation}\label{eq:v19-flow}
 v=\phi_\varepsilon(u),\qquad
 \widehat{\mathcal F}_{L,\varepsilon}
                    =\phi_{-\varepsilon}\circ\mathcal F_L.
\end{equation}
Thus the new source is explicitly \(\varepsilon\)-dependent.
At a fixed nonzero \(\varepsilon\) it is not the original
balanced root source.

\begin{theorem}[Countertransport of the complete finite packet]
\label{thm:v19-flow}
Let \(\mathcal L_{\chi_L}f=Df[\chi_L]\). Pull back the density
\eqref{eq:v19-packet} by \eqref{eq:v19-flow}, with the determinant
of the change of retained variables included. Its coefficients are
\begin{align}
 \widehat{\mathcal G}_L&=\mathcal G_L,\label{eq:v19-Gflow}\\
 \widehat{\mathcal Q}_L
   &=\mathcal Q_L+\mathcal L_{\chi_L}\mathcal G_L
     =\mathcal Q_L^{\rm rest}+R_{2,L},\label{eq:v19-Qflow}\\
 \widehat{\mathcal B}_L
   &=\mathcal B_L+\mathcal L_{\chi_L}\mathcal Q_L
         +\frac12\mathcal L_{\chi_L}^2\mathcal G_L
                          -\operatorname{div}\chi_L \notag\\
   &=\mathcal B_L+\mathcal L_{\chi_L}\mathcal Q_L^{\rm rest}
         +\mathcal L_{\chi_L}R_{2,L}
         -\frac12\mathcal L_{\chi_L}^2\mathcal G_L
                          -\operatorname{div}\chi_L .
                                                  \label{eq:v19-Bflow}
\end{align}
In particular, the compensated relative contribution to
\(J^6\widehat{\mathcal Q}_L\) is exactly
\(-\frac14\Tr K_{L,[3]}^2\) and has the limit
\eqref{eq:v19-six}. This assertion is about that contribution,
not about convergence of \(J^6\widehat{\mathcal Q}_L\) as a whole.
\end{theorem}
\begin{proof}
For any fixed smooth coefficient \(f\), differentiation along the
flow gives
\[
 f(\phi_\varepsilon(u))
   =f(u)+\varepsilon\mathcal L_{\chi_L}f(u)
        +\frac{\varepsilon^2}{2}\mathcal L_{\chi_L}^2f(u)
        +O_L(\varepsilon^3).
\]
In particular,
\(\mathcal L_\chi^2f
 =D^2f[\chi,\chi]+Df[D\chi\,\chi]\).
The variational equation for \(D\phi_t\), followed by the finite
matrix determinant differentiation formula, gives
\[
 \log\det D\phi_\varepsilon(u)
      =\int_0^\varepsilon
             \operatorname{div}\chi_L(\phi_t(u))\,dt
      =\varepsilon\operatorname{div}\chi_L(u)+O_L(\varepsilon^2).
\]
The determinant is positive near \(t=0\).
The pulled-back negative logarithm is
\[
 \varepsilon^{-1}\mathcal G_L(\phi_\varepsilon(u))
   +\mathcal Q_L(\phi_\varepsilon(u))
   +\varepsilon\mathcal B_L(\phi_\varepsilon(u))
   -\log\det D\phi_\varepsilon(u)+O_L(\varepsilon^2).
\]
Collect the coefficients of \(\varepsilon^{-1},1,\varepsilon\).
This proves the first lines of
\eqref{eq:v19-Gflow}--\eqref{eq:v19-Bflow}.
Use \eqref{eq:v19-remainder} to obtain the second lines.
Finally apply Proposition~\ref{prop:v19-six}.
\end{proof}

For any retained test observable \(f\), the exact change of variables
uses \(f(\phi_\varepsilon(u))\), not \(f(u)\).
Consequently this is a finite source-renormalization convention
for the same integral with transported observables. It is not
cancellation of a physical effect while leaving the observable
unchanged, and it is not a proof that the modified observables
have a common quantum continuum limit.

The cancellation is also not an instruction to subtract
\(\frac12\Tr K_L\) from a prescribed density while omitting
\eqref{eq:v19-Bflow}. The retained Jacobian first appears at
two-loop order here because the change of retained variables is
of order \(\varepsilon\). It is a different Jacobian from the
\(\varepsilon\)-independent inverse source pivot in V14.

\subsection{What the fourth two-loop jet actually needs}

\begin{proposition}[The quintic source field enters the quartic density]
\label{prop:v19-budget}
Write \(\chi_L=\sum_{r\ge1}\chi_{L,[r]}\), with
\(\chi_{L,[1]}(v)=C_Lv\). Then
\begin{equation}\label{eq:v19-Bfour}
 \begin{split}
 (\widehat{\mathcal B}_L-\mathcal B_L)_{[4]}
 ={}&\sum_{r=1}^4
       D(\mathcal Q^{\rm rest}_{L,[5-r]})[\chi_{L,[r]}]\\
    &-\frac12\left\{
         D^2E_4[C_Lv,C_Lv]+g(v)\cdot C_L^2v\right\}\\
    &-\operatorname{div}\chi_{L,[5]}(v).
 \end{split}
\end{equation}
Thus the logarithmic linear vector field alone does not determine
the new quartic two-loop coefficient. Already its source
Jacobian requires the quintic vector coefficient.
\end{proposition}
\begin{proof}
The inherited classical action starts in degree four,
\(\mathcal G_L=E_4+O_L(|v|^5)\), and \(\chi_L\) starts in
degree one. A Lie derivative with this vector field cannot
lower the source degree. Since \(R_{2,L}\) starts in degree
six, \((\mathcal L_{\chi_L}R_{2,L})_{[4]}=0\).
For a homogeneous scalar of degree \(j\) and a vector of
degree \(r\), \(Df_j[\chi_r]\) has degree \(j-1+r\).
This gives the sum in \eqref{eq:v19-Bfour}.
Twice differentiating the classical term at degree four
can only use \(E_4\) and the linear vector field \(C_Lv\).
It gives
\[
 \mathcal L_{C_Lv}^2 E_4
      =D^2E_4[C_Lv,C_Lv]+g(v)\cdot C_L^2v.
\]
Finally a divergence lowers source degree by one, so its
fourth-degree part is exactly
\(\operatorname{div}\chi_{L,[5]}\).
\end{proof}

This last requirement is genuine at the level of measured transport,
not a loose derivative estimate. In two retained coordinates let
\(\chi(x,y)=(x^5,0)\), let \(G\) start in degree four, and set
\(Q=B=0\). The quartic part of the transformed two-loop coefficient
is \(-5x^4\), entirely from \(-\operatorname{div}\chi\);
the twice differentiated classical term has degree at least twelve.
This is an exact algebraic necessity example, not a calculation
of the actual YM quintic vector field.

In the actual definition \eqref{eq:v19-chi}, differentiating five
times can involve the seventh derivative of \(\mathcal F_L\),
as well as derivatives of the actual branch, tangent and reference
inverse. This explains directly why the source degree seven in
Shell V8's measured packet remains relevant after the quartic
one-loop term has been removed. There is no basis for resetting
the transported one-loop density or omitting its derivatives.

At zero, \(\chi_L(0)=0\), so all scalar Lie derivatives in
\eqref{eq:v19-Bflow} vanish there. The remaining vacuum shift is
\[
 \widehat{\mathcal B}_L(0)-\mathcal B_L(0)
    =-\operatorname{div}\chi_L(0)
    =-3\Tr C_L=3\Tr A_L.
\]
The factor three counts colours. It is removed only if the
specified vacuum normalization removes it; it is not silently
zero. For the actual side-two calibration,
\(\operatorname{div}\chi_2(0)=-4439/768\).

\subsection{Live companion update: the density is now specified}
\label{sec:v19-live}

During the final verification of this version the active shell-mixing
companion advanced from V13 to V16. The latter is
\begin{center}\small
\path{Renewal_Yang_Mills_Shell_Mixing_Research_Notes_V16.tex},
\end{center}
with 623,261 bytes and SHA--256
\begin{center}\scriptsize\ttfamily
BF138E63BA826BC030D3A8B87E04FDD3B1EC712D269BB064674E96509B46E710.
\end{center}
Its new local density derivation, finite measured assembly, and
flat-root calculation have been read, together with the defining
precision and normal matrices in its V15. This changes the next
step: the local Coulomb density is no longer an unspecified input.
We do not repeat the companion's flat-mode proof or claim a fresh
verification of all its trace-ideal estimates.

Here is the exact dictionary needed to use its result without
identifying different determinants. Shell V16 writes
\[
 \rho_{\mathcal C}(a)
    =J_{\rm H}(a)\det{}'(\Delta^{-1/2}\mathcal M_a\Delta^{-1/2}),
 \qquad
 \mathcal M_a=d^*R_a,\quad \Delta=d^*d,
\]
relative to fixed Coulomb-slice volume. The prime removes the
three constant vertex gauge parameters, not the twelve constant
link source variables. With \(\omega\) a vertex field,
\[
 (R_a\omega)_{xy}
   =\sigma(\operatorname{ad}_{a_{xy}})(\omega_y-\omega_x)
          +\frac12[a_{xy},\omega_x+\omega_y],
 \qquad \sigma(z)=\frac z2\coth\frac z2 .
\]
This is the logarithmic gauge-coordinate Jacobian, not a further
factor to multiply into an already gauge-fixed tree density.
The full fine-link Haar factor \(J_{\rm H}\) is required in this
coordinate convention.

To check the finite identity, the right-exponential derivative is
\((e^{\operatorname{ad}a}-1)/\operatorname{ad}a\).
Its inverse applied to
\(-\omega_x+e^{\operatorname{ad}a}\omega_y\) gives the displayed
\(R_a\). The local slice/gauge coordinate differential has columns
\([I_{\mathcal C},R_a]\); projection to the normal space gives
\(\Delta^{-1/2}d^*R_a\). Its determinant relative to the zero-field
normal map is the primed ratio above. Constants are killed by
\(d^*R_a\) on \(d^*a=0\), and the based-to-mean-zero gauge
coordinate conversion has a field-independent determinant.
Multiplication by the logarithmic Haar density gives
\(\rho_{\mathcal C}\). These calculations apply on the companion's
regular, constant-conjugation-invariant local branch. They do not
prove a global gauge slice. For the balanced source the companion
also retains the orbit argument: averaging test functions over
constant conjugations identifies the balanced and raw rooted
pushforwards, since their outputs have the same orbit and both
laws are invariant. Representative Jacobians are not dropped.

Shell V15's \(M_L\) equals this note's \(Z_L\). Its multiplier
\(\eta_L\) equals \(-\zeta_L\). Its symbol \(I+K_L^{\rm Shell}\)
is the \emph{entire} energy-normalized constrained precision,
whereas this note's \(K_L\) is relative to the Wilson reference:
\begin{equation}\label{eq:v19-dictionary}
 I+K_L^{\rm Shell}=H_{0,L}+B_L^{\rm c}
       =H_{0,L}^{1/2}(I+K_L)H_{0,L}^{1/2}.
\end{equation}
Equality here is after the energy-isometric identification of
their actual detail coordinates. It follows directly from the
two definitions of the tangent Lagrangian Hessian.
In particular their finite logarithmic determinants split exactly.
Their first traces or modified determinants must not be
identified term by term by the same argument.

Put
\[
 I+G_L^{\rm gh}
      =\Delta^{-1/2}\mathcal M_{a_*}\Delta^{-1/2},
 \qquad
 \ell_c(v)=\sum_\mu 2\log\frac{\sin|v_\mu|}{|v_\mu|}.
\]
For the primitive local law, the actual rest term in the
\emph{Lebesgue} convention used in
Theorem~\ref{thm:v19-flow} is therefore specified by
\begin{equation}\label{eq:v19-rest}
 \mathcal Q_L^{\rm rest}
  =\log\det Z_L+\frac12\log\det H_{0,L}
       -\log\det(I+G_L^{\rm gh})-\log J_{\rm H}(a_*),
\end{equation}
with zero-source constants removed. Any additional incoming
one-loop factor must still be included when iterating.
The full coefficient relative to coarse Haar is
\[
 \mathcal Q_L^{\rm H}
       =\mathcal Q_L^{\rm rest}+q_L^{\rm rel}+\ell_c.
\]
This is precisely the finite assembly of Shell V16 after
\eqref{eq:v19-dictionary}. It supplies the missing measure
convention; it does not estimate the noncommuting low traces
in \eqref{eq:v19-rest}.

\begin{proposition}[Countertransport relative to coarse Haar]
\label{prop:v19-Haarflow}
If the initial and final retained reference volumes are both
\(e^{\ell_c(v)}\,dv\), the one-loop formula is
\[
 \widehat{\mathcal Q}_L^{\rm H}
   =\mathcal Q_L^{\rm H}
                      +\mathcal L_{\chi_L}\mathcal G_L
   =\mathcal Q_L^{\rm rest}+R_{2,L}+\ell_c .
\]
The two-loop formula becomes
\begin{equation}\label{eq:v19-Haarflow}
 \widehat{\mathcal B}_L^{\rm H}
   =\mathcal B_L+\mathcal L_{\chi_L}\mathcal Q_L^{\rm H}
       +\frac12\mathcal L_{\chi_L}^2\mathcal G_L
       -\operatorname{div}_{\rm H}\chi_L,\qquad
 \operatorname{div}_{\rm H}\chi
             =\operatorname{div}\chi+D\ell_c[\chi].
\end{equation}
\end{proposition}
\begin{proof}
Conversion to Lebesgue replaces \(\mathcal Q_L^{\rm H}\)
by \(\mathcal Q_L^{\rm H}-\ell_c\), leaving
\(\mathcal G_L,\mathcal B_L\) unchanged. Apply
Theorem~\ref{thm:v19-flow}, then add \(\ell_c(u)\) to the
new one-loop coefficient to return to Haar volume.
The two-loop derivative of the subtracted \(\ell_c\) is
exactly \(-D\ell_c[\chi_L]\), proving the formula.
Equivalently, the change of Haar volume has Jacobian
\[
 \det D\phi_\varepsilon(u)\,
   \exp\{\ell_c(\phi_\varepsilon(u))-\ell_c(u)\},
\]
whose first logarithmic derivative is
\(\operatorname{div}_{\rm H}\chi_L\).
\end{proof}

This distinction matters for fourth source order:
\(D\ell_{c,[2]}[\chi_{L,[3]}]\) and
\(D\ell_{c,[4]}[\chi_{L,[1]}]\) both contribute.
The quintic divergence requirement in
Proposition~\ref{prop:v19-budget} remains.

There is a useful companion regression which prevents overreading
the countertransport. On the flat source
\(v=(xE_3,0,0,0)\), the actual minimum has all plaquettes equal
to the identity, hence \(\zeta_L=0\).
Thus \(B_L^{\rm c}=K_L=0\), \(q_L^{\rm rel}=R_{2,L}=0\),
and \(D\mathcal G_L[\chi_L]=0\) at every point of this line.
The complete flat one-loop potential is nevertheless nonzero.
Shell V16 evaluates it as
\[
 \mathcal Q_L^{\rm H}(xE_3,0,0,0)
  =3\log\!\left(1-\frac23(1-L^{-2})\sin^2x\right)
   +2\sum_{m\in\mathcal K_L^3\setminus\{0\}}
       \log\!\left(1+\frac{1-\cos 2x}{D_L^{\rm flat}(m)}\right),
\]
where
\[
 D_L^{\rm flat}(m)=
  \cosh\!\left(2L\operatorname{arsinh}
      \sqrt{\sum_{\nu=1}^3\sin^2(\pi m_\nu/L)}\right)-1 .
\]
This formula includes the root tangent return inside the full
constrained precision. It is not added once more to
\eqref{eq:v19-rest}.
The companion proves its normal convergence to the sum with
\(m\in\mathbb Z^3\setminus\{0\}\),
\(D_\infty^{\rm flat}(m)=\cosh(2\pi|m|)-1\), with
\(O(L^{-2})\) error on smaller complex discs. Its proof uses
the exact longitudinal Chebyshev product and
\(2|m|\le2L\operatorname{arsinh}
\sqrt{\sum_\nu\sin^2(\pi m_\nu/L)}\le2\pi|m|\);
it is not a noncommuting \(F^2\) calculation.

As an imported finite regression, its side-two Haar coefficient is
\[
 \mathcal Q_2^{\rm H}(xE_3,0,0,0)
       =-\frac{11}{24}x^2-\frac{211}{768}x^4+O(x^6).
\]
Subtracting \(\ell_c=-x^2/3-x^4/90+O(x^6)\) gives, in the
Lebesgue convention of \eqref{eq:v19-rest},
\[
 \mathcal Q_2^{\rm rest}(xE_3,0,0,0)
       =-\frac18x^2-\frac{3037}{11520}x^4+O(x^6).
\]
Our exact checker verifies this convention conversion, not a new
independent evaluation of the companion's mode sums.
The one-loop countertransport leaves these flat coefficient
functions unchanged because its added classical derivative is zero.
It therefore cannot remove this complete finite-volume holonomy
potential. That potential is not a local physical gluon mass.

\subsection{Uniformity audit and the next measured acquisition}

The derivative of the exact finite source flow at zero satisfies
\[
 D\phi_\varepsilon(0)
       =\exp\{\varepsilon(C_L\otimes I_3)\}.
\]
This follows by linearizing its defining ODE at its fixed point.
The currently proved estimate yields
\[
 \|D\phi_\varepsilon(0)\|,
 \ \|D\phi_{-\varepsilon}(0)\|
 \le
 \exp\{|\varepsilon|(C_\Gamma+3584\sqrt2)n\},
 \qquad L=2^n.
\]
It supplies no uniform fixed-\(\varepsilon\), growing-\(n\)
smallness. A bound on \(|\varepsilon|n\) controls this
linearization, but does not by itself control the nonlinear flow:
no uniform bounds for all required higher coefficients of
\(\chi_L\) have been proved. The limits in
\eqref{eq:v19-matrix-limit} identify the linear logarithm;
they neither resum the flow nor prove a running-coupling law.

The operational conclusion is now sharper than a list of missing
determinants. There are two distinct next acquisitions:
the noncommuting ultraviolet analysis of the now specified
measured rest factor \eqref{eq:v19-rest}, and the higher source
countertransport through degree five if
the normalized source convention is used at two loops.
In the original source convention the first trace must instead
be retained. These alternatives describe the same finite
integral only when the observable and all density terms are
transported together. Proving one of their cofinal measured
versions is still necessary before using these coefficients
to construct a quantum continuum measure.

\subsection{Exact checks and version integrity}

The new checker is
\begin{center}\small
\path{scripts/verify_ym_f14_v19_source_countertransport.py}.
\end{center}
Its SHA--256 is
\begin{center}\scriptsize\ttfamily
29CE53977B79FD6714540C945745D4F26DD90E1D05ADF2E5B9AF19226B623974.
\end{center}
All its arithmetic is rational. It verifies the Gram identity
for all sixteen matrix basis elements, proves rank ten by exact
elimination in polynomial coefficient space, and checks the
384 signed-permutation matrix average. It reruns V16's literal
side-two Gaussian cubic-root calculation and projected reference
symbol checks, then verifies the six pair coefficients, six
three-orientation cross coefficients, isotropic projection and
colour factor in the source divergence.

The nonlinear-flow tests compare direct substitution and the
actual two-by-two Jacobian with the three coefficient formulas
for three polynomial vector fields, in both Lebesgue and
nonconstant-reference-volume conventions: eighteen exact
coefficient comparisons. They include a nonzero
quintic divergence and the \(-5x^4\) calibration.
Eight finite symmetric-matrix contractions also check the sign
and factor \(1/2\) in the first-trace identity. These are exact
finite algebra checks; the analytic proofs and the inherited
V15/V18 limits supply the mathematical conclusions at arbitrary
regulator. The tests do not construct \(m_L\) or a continuum law.

The complete predecessor V18 has 372,757 bytes and SHA--256
\begin{center}\scriptsize\ttfamily
C3D2420AE77F20D0A2D4F73837717F956EE1F66739DDC68D92CFCC0A788DCED7.
\end{center}
Its full body is preserved apart from the title version and this
terminal insertion. Active V19 replaces V18; all companions,
archives and monographs are untouched by this task. Shell V16
replaced V13 externally during this turn and is incorporated in
Section~\ref{sec:v19-live}; its V13 classical results remain
inside it. SM V72 and NS V26 are unchanged. The next scheduled
SM/NS comparison is V20.

\begin{firewall}
This version proves a finite-regulator source-countertransport
identity and a convergent relative one-loop six-jet in that
explicit source convention. It does not prove convergence of
the full measured one-loop coefficient, a uniform nonlinear
source renormalization, a higher-loop or nonperturbative
continuum construction, compact-source coverage, or a positive
physical Yang--Mills Hamiltonian gap.
\end{firewall}
\section{V20: the complete noncommuting side-two coefficient and an all-level reduction}
\label{sec:v20-measured}

The new companion Shell V16 is already incorporated in V19.
The user has also supplied an optional programme suggestion, read as
research input rather than as instructions to merge or archive files.
Its useful recommendation is adopted: the original balanced root
source \(v\) and the complete coarse-Haar one-loop density are the
primary convention. V19's source countertransport remains a derived
change of coordinates, not an unannounced subtraction. The two
lineages and their archives are left intact.

This version computes the complete mixed coefficient on
\[
               v=(xE_1,yE_2,0,0)
\]
at side two, where \(E_4(v)=2x^2y^2\).
It also proves exact formulas for the root-normal and reference
tangent contributions at every dyadic refinement. The general
noncommuting asymptotic problem is thereby reduced to a specified
sum of mean-fixed Wilson and gauge Fourier matrices, with its
fine-Haar term retained. The finite answer is not extrapolated
to a continuum coupling coefficient.

\subsection{Canonical measured quantity and scope}

All quantities below are normalized at zero source. Use V19's
notation and determinant dictionary, and write
\begin{equation}\label{eq:v20-ledger}
 \mathcal Q_L^{\rm H}
   =\log\det Z_L+\frac12\log\det H_{0,L}
       +q_L^{\rm rel}
       -\log\det(I+G_L^{\rm gh})
       -\log J_{\rm H}(a_L^*)+\ell_c .
\end{equation}
Here \(H_{0,L}\) already uses the actual root tangent,
\(q_L^{\rm rel}=\frac12\log\det(I+K_L)\) is relative to
that reference, and the primed gauge determinant is the
logarithmic Coulomb-coordinate determinant of Shell V16.
Equation~\eqref{eq:v20-ledger} is a local one-loop coefficient
of the primitive compact law, not a full interacting integral.
An incoming one-loop density at a later elimination is not zero
by this convention and must be transported.

Write \([x^2y^2]f\) for an ordinary polynomial coefficient.
Thus \(\partial_x^2\partial_y^2 f(0)=4[x^2y^2]f\), not
\(24[x^2y^2]f\). We will prove
\begin{equation}\label{eq:v20-answer}
 \boxed{\quad [x^2y^2]\mathcal Q_2^{\rm H}=\frac{3221}{768}.
                                                        \quad}
\end{equation}
All fifteen nonzero side-two frequencies, nine real Lie-valued
transverse degrees per frequency, and all twelve normal source
coordinates are retained. No selection of a few covariance modes
stands in for that finite lattice.

\subsection{Why the unknown fourth-order detail displacement is not needed}

Set
\[
 \alpha_L=\frac{1-L^{-2}}{24},\qquad
 \bar c_L(v)=v+\alpha_Lg(v),\qquad
                  \bar a_L(v)=\bar c_L(v)/L.
\]
The inherited branch and homogeneous-source results give, at
each fixed regulator,
\begin{equation}\label{eq:v20-branch}
        a_L^*(v)=\bar a_L(v)+O_L(|v|^4).
\end{equation}
Indeed the homogeneous source is
\(\Phi_L(c)=c-\alpha_Lg(c)+O_L(|c|^4)\);
the actual mean-zero detail starts in degree four, and the
first three mean coefficients are those of \(\Phi_L^{-1}\).
This does not replace the actual minimum at higher orders.

\begin{lemma}[Fourth-jet insensitivity of invariant scalar factors]
\label{lem:v20-insensitivity}
Let \(\Psi(a)\) be an analytic scalar germ on the chosen fine
constant-plus-Coulomb-detail space, invariant under simultaneous
constant \(\mathrm{SU}(2)\) conjugations. Then \(D\Psi(0)=0\).
Consequently
\[
             J^4\Psi(a_L^*(v))=J^4\Psi(\bar a_L(v)).
\]
This applies separately to the normal determinant, fine Haar
factor, logarithmic gauge determinant, and Wilson-reference
determinant with its source tangent.
\end{lemma}
\begin{proof}
The linear functional \(D\Psi(0)\) is invariant under constant
adjoint rotations. The fine coordinate space is a direct sum
of copies of the nontrivial three-dimensional adjoint
representation, restricted by colour-independent linear
conditions. It has no invariant covectors. Equivalently,
averaging any vector over the compact rotation group gives
zero and leaves an invariant linear functional unchanged.
Thus \(D\Psi(0)=0\).

Along the segment between the two fields in
\eqref{eq:v20-branch}, analyticity gives
\(D\Psi=O_L(|v|)\). Integrating this derivative against their
\(O_L(|v|^4)\) difference proves an \(O_L(|v|^5)\) scalar
difference, hence equality of fourth jets.

For the claimed factors, all coordinate maps commute with
constant adjoint rotations. The normal derivative matrix is
conjugated on its source indices, and the detail and gauge
matrices are conjugated on their respective energy spaces;
their determinants are invariant. The fine Haar factor is
invariant link by link. These scalar germs are analytic near
zero with their positive real branches. The tangent reference
matrix can be defined at any sufficiently small fine field by
the kernel of \(D\mathcal F_L\), using its invertible normal
block, so this argument does not assume it is already at a
stationary point.
\end{proof}

The relative constraint term is retained from its actual
V16 identity; its multiplier is not reset by this lemma.
In particular this reduction does not claim that a homogeneous
trial action computes every fluctuation or higher-order
stationary coefficient.

\subsection{The all-level reference tangent return}

Let \(H_L^{\rm mean}(c)\) denote the Wilson Hessian at the
homogeneous fine field \(c/L\), restricted to mean-zero Coulomb
details and normalized by the free energy metric. This is a
mean-fixed reference, not yet the root-constrained one.
The first source-tangent correction is
\[
 \ell_{L,[1]}(v)Z=-\mathcal M_vZ,\qquad
 (\mathcal M_vZ)_\mu
    =\sum_{\lambda\ne\mu}[v_\lambda,W_{\lambda\mu}],
 \qquad
 W_{\lambda\mu}=L^{-4}\sum_z
       \bigl(z_\lambda-(L-1)/2\bigr)Z_\mu(z).
\]
These are the actual mixed source rows of V16 and Shell V11.
The free fine covariance of \(Z\) is the full Coulomb covariance.

\begin{lemma}[Exact moment covariance]
\label{lem:v20-moment}
For \(\lambda\ne\mu\), \(\kappa\ne\nu\), and colour indices
\(a,b\),
\begin{equation}\label{eq:v20-sigma}
 \mathbb E W_{\lambda\mu}^aW_{\kappa\nu}^b
   =\sigma_L\delta_{\lambda\kappa}\delta_{\mu\nu}\delta_{ab},
 \qquad
 \sigma_L=\frac{1+10L^{-2}-11L^{-4}}{720}.
\end{equation}
\end{lemma}
\begin{proof}
The centered spatial moment has Fourier support only on the
nonzero \(\lambda\)-axis. Distinct axes have disjoint supports,
so their covariance is zero. On that axis the Coulomb projection
is the identity on orientations different from \(\lambda\),
with no cross-orientation covariance. For
\(\theta=2\pi m/L\ne0\), its normalized one-dimensional Fourier
coefficient has magnitude \(1/|e^{i\theta}-1|\). The covariance
normalization of a fine Fourier coefficient is \(L^{-4}\)
times the inverse free eigenvalue. Consequently the common
variance is
\[
 \sigma_L=L^{-4}\sum_{m=1}^{L-1}
                   \frac1{16\sin^4(\pi m/L)}.
\]
The identity
\(\sum_{m=0}^{L-1}\csc^2(z+\pi m/L)=L^2\csc^2(Lz)\)
follows by logarithmically differentiating the finite sine
product. Differentiate it twice and use
\((f''+4f)/6=f^2\) for \(f=\csc^2z\).
After removing the \(m=0\) Laurent term at \(z=0\), one obtains
\[
 \sum_{m=1}^{L-1}\csc^4(\pi m/L)
                         =\frac{L^4+10L^2-11}{45}.
\]
Substitution proves the variance and the covariance formula.
The free colours are independent copies.
\end{proof}

Define the three-by-three colour matrices
\[
 B_\mu(v)=4\sum_{\lambda\ne\mu}
                 (|v_\lambda|^2I-v_\lambda v_\lambda^T).
\]
These are the diagonal orientation blocks of \(D^2E_4(v)\);
they are not the source-curvature operator \(B_L^{\rm c}\).

\begin{theorem}[Fourth-degree reference return at every refinement]
\label{thm:v20-return}
The difference between the root-tangent and mean-fixed
Wilson-reference logarithms has fourth coefficient
\begin{equation}\label{eq:v20-return}
 \left(\frac12\log\det H_{0,L}
       -\frac12\log\det H_L^{\rm mean}(\bar c_L)\right)_{[4]}
       =\frac{\sigma_L}{2}\sum_{\mu=1}^4\Tr B_\mu(v)^2.
\end{equation}
It has no lower source coefficient. On the two-colour plane,
the right side equals
\begin{equation}\label{eq:v20-return-plane}
 48\sigma_L(x^4+y^4)+32\sigma_Lx^2y^2 .
\end{equation}
In particular the mixed return tends to \(2/45\), and is
\(1/8\) at side two.
\end{theorem}
\begin{proof}
By Lemma~\ref{lem:v20-insensitivity}, evaluate the reference
factor at the homogeneous field \(\bar a_L\). Translation
invariance makes its Wilson Hessian block diagonal between
constant modes and mean-zero modes. The constant block,
in physical mean coordinates, starts as \(D^2E_4(v)\),
of source degree two. The root tangent adds a constant
variation \(\ell_LZ\), with \(\ell_L=O_L(|v|)\).
Thus its Hessian differs from the mean-fixed Hessian first
in degree four, by
\(\ell_{L,[1]}^TD^2E_4\,\ell_{L,[1]}\).
The free relative mean Hessian is the identity. The fourth
logarithmic determinant difference is therefore half the
free trace of that matrix; inverse-reference corrections have
higher degree.

The trace is
\(\frac12\mathbb E D^2E_4[\mathcal M_vZ,\mathcal M_vZ]\).
Since \([v,w]=2v\times w\), Lemma~\ref{lem:v20-moment} gives
\[
 \mathbb E(\mathcal M_vZ)_\mu
                  (\mathcal M_vZ)_\nu^T
                         =\delta_{\mu\nu}\sigma_LB_\mu(v).
\]
Only the diagonal orientation blocks of the Hessian survive
this contraction, proving \eqref{eq:v20-return}.
On the indicated plane their squared traces sum to
\(96(x^4+y^4)+64x^2y^2\). This proves
\eqref{eq:v20-return-plane}.
\end{proof}

This return is added only because the next calculation explicitly
uses a mean-fixed Wilson determinant. It would be double-counted
if added to a determinant already computed on the root tangent.
The flat coefficient \(48\sigma_L\) also agrees with the
fourth-order expansion of Shell V16's exact flat return.

\subsection{The finite normal-root calculation and its exact certificate}

The following local polynomial calculation is used both for
the side-two result and for the all-level normal formula.
At side two let \(f(c)=\mathcal T_{\rm hom}(c/2)\) be the
literal balanced homogeneous block, so \(f(c)=c+f_3(c)+\cdots\)
with \(f_3=-g/32\). Write
\[
 Z_2(\bar a_2)=I+N_2+N_3+N_4+\cdots.
\]
Here the matrix entries have source degrees indicated by
their subscripts, and \(N_1=0\).

\begin{lemma}[Side-two normal trace certificate]
\label{lem:v20-normal-base}
On the two-colour plane,
\begin{align}
 \Tr N_2&=-\frac34(x^2+y^2),\notag\\
 \Tr N_4&=\frac1{64}(x^4+y^4)+\frac{61}{192}x^2y^2,\notag\\
 \Tr N_2^2&=\frac3{32}(x^4+y^4)+\frac7{32}x^2y^2.
                                                    \label{eq:v20-normal-traces}
\end{align}
Consequently
\begin{equation}\label{eq:v20-normal-base}
 J^4\log\det Z_2
    =-\frac34(x^2+y^2)-\frac1{32}(x^4+y^4)
                                  +\frac5{24}x^2y^2.
\end{equation}
\end{lemma}
\begin{proof}
For clarity, an exact finite certificate for the only longer
trace in \eqref{eq:v20-normal-traces} is specified here. It uses
no numerical matrix eigenvalues. Work in the rational jet ring
\[
 \mathcal R=\mathbb Q[x,y,u,w]/
    \langle x^iy^j:i+j=5;\ u^2;\ w^2\rangle .
\]
The variation variables \(u,w\) are independent square-zero
legs. A monomial in this ring has total degree at most six.
For an imaginary quaternion \(a\), use exactly
\begin{equation}\label{eq:v20-exp}
 \exp a=
 \sum_{j=0}^3\frac{(-1)^j(a\cdot a)^j}{(2j)!}
 +a\sum_{j=0}^3\frac{(-1)^j(a\cdot a)^j}{(2j+1)!},
 \qquad
 \log(1+z)=\sum_{j=1}^6\frac{(-1)^{j+1}}jz^j .
\end{equation}
These are exact in this ring for the positive-degree arguments
in question. Quaternion multiplication is
\[
 (a_0,a)(b_0,b)=(a_0b_0-a\cdot b,\ a_0b+b_0a+a\times b).
\]

For each of the twelve normal legs \(e_j\), insert
\[
 a=\frac12(\bar c_2+ue_j),\qquad
 \bar c_2=( (x+xy^2/8)E_1,\ (y+x^2y/8)E_2,\ 0,\ 0).
\]
Set \(U_\mu=\exp a_\mu\), and for \(s\in\{0,1\}^4\) let
\[
 t_s=\prod_{\mu=1}^4U_\mu^{s_\mu},\quad
 \phi=\frac1{16}\sum_s\log t_s,\quad h=\exp(-\phi),\quad
 P_{s,\mu}=t_sU_\mu^2t_s^{-1}.
\]
The product in \(t_s\) is in ascending orientation order.
The actual output is
\begin{equation}\label{eq:v20-root-ring}
 \log\!\left[
 hP_{0,\mu}
   \exp\!\left\{\frac1{16}\sum_s
                      \log(P_{0,\mu}^{-1}P_{s,\mu})\right\}h^{-1}
 \right].
\end{equation}
Its coefficient of \(u\), with \(w=0\), is column \(j\)
of \(Z_2\). Taking its twelve diagonal entries by
\eqref{eq:v20-exp} gives the first two traces of
\eqref{eq:v20-normal-traces}. This is a finite rational
recurrence for the asserted values; the supplied checker
executes it with exact fractions, keeping all twelve columns.

There is also a shorter independent check on the squared trace:
\(N_2=-D^2E_4/32\), and direct differentiation gives
\[
 \Tr(D^2E_4)^2
               =96(x^4+y^4)+224x^2y^2.
\]
This proves the last trace in \eqref{eq:v20-normal-traces}.
The fourth logarithmic coefficient is
\(\Tr N_4-\frac12\Tr N_2^2\).
Odd monomials on this plane vanish: constant adjoint rotations
by \(\pi\) about \(E_1\) or \(E_2\) independently reverse \(y\)
or \(x\). Finally Lemma~\ref{lem:v20-insensitivity} transfers
the scalar jet to the actual minimizing field.
\end{proof}

\subsection{Normal coarea and fine Haar at every dyadic level}

\begin{theorem}[Exact all-level root-normal mixed coefficient]
\label{thm:v20-normal-all}
Put \(r=L^{-2}\). The actual root-normal factor obeys
\begin{equation}\label{eq:v20-normal-all}
 \begin{split}
 J^4\log\det Z_L(xE_1,yE_2,0,0)
   ={}&-(1-r)(x^2+y^2)-\frac{1-r^2}{30}(x^4+y^4)\\
     &+\frac{4+30r-34r^2}{45}x^2y^2 .
 \end{split}
\end{equation}
Its mixed coefficient tends to \(4/45\).
\end{theorem}
\begin{proof}
Let \(L=2^n\). Define
\[
 a_n=\sum_{j=0}^{n-1}4^{-j}=\frac43(1-L^{-2}),
 \qquad
 b_n=\sum_{j=0}^{n-1}16^{-j}=\frac{16}{15}(1-L^{-4}).
\]
The homogeneous source is
\(\Phi_L(c)=\mathcal T_{\rm hom}^{\,n}(c/L)\).
Peeling off the finest step gives the exact composition
\[
 \Phi_{2M}(c)=
 \Phi_M\!\left(c+M^{-2}f_3(c)+M^{-3}f_4(c)
                                      +M^{-4}f_5(c)+O(|c|^6)\right).
\]
Since there is no quadratic term, its cubic and fifth
homogeneous terms are
\[
 (\Phi_L)_{[3]}=a_nf_3,\qquad
 (\Phi_L)_{[5]}=b_nf_5+
                        \frac{a_n^2-b_n}{2}Df_3[f_3].
\]
The second formula follows by summing the ordered cross pairs
of cubic insertions; the sum of their weights is
\(\frac12((\sum_j4^{-j})^2-\sum_j16^{-j})\).
Fourth-degree terms do not enter this fifth-degree identity.

Evaluate \(D\Phi_L\) at
\(\Phi_L^{-1}(v)=v-a_nf_3(v)+O(|v|^4)\), and set
\[
 A=\Tr(Df_3)^2,\qquad B=D(\operatorname{div}f_3)[f_3].
\]
Expansion of the finite logarithmic determinant gives
\[
 (\log\det D\Phi_L(\Phi_L^{-1}(v)))_{[4]}
   =b_n\operatorname{div}f_5-\frac{b_n}{2}A
                                      -\frac{a_n^2+b_n}{2}B.
\]
Indeed \(\operatorname{div}(Df_3[f_3])=A+B\);
the shift of the argument contributes \(-a_n^2B\),
and the quadratic logarithm contributes \(-a_n^2A/2\).
For \(n=1\) this expression is
\(\operatorname{div}f_5-A/2-B\).
Consequently at general \(n\) it is \(b_n\) times the
side-two fourth coefficient plus \((b_n-a_n^2)B/2\).

Here \(\operatorname{div}f_3=-3|v|^2/4\) and \(f_3=-g/32\),
so \(B=3E_4/16\). On the plane it is \(3x^2y^2/8\).
Lemma~\ref{lem:v20-normal-base} therefore gives the mixed term
\[
 \frac5{24}b_n+\frac3{16}(b_n-a_n^2)
                             =\frac{4+30r-34r^2}{45}.
\]
The pure fourth coefficient is \(-b_n/32=-(1-r^2)/30\),
and the quadratic coefficient is \(-3a_n/4=-(1-r)\).
Replacing the homogeneous inverse by \(\bar c_L\), and then
by the actual branch, does not change these scalar fourth
jets by Lemma~\ref{lem:v20-insensitivity}. This proves the claim.
\end{proof}

\begin{proposition}[The exact fine-Haar contribution]
\label{prop:v20-Haar}
On the actual branch,
\[
 (-\log J_{\rm H}(a_L^*))_{[4]}
       =\frac{L^2-1}{9}E_4(v)
                           +\frac1{90}\sum_\mu|v_\mu|^4.
\]
In particular its mixed coefficient on the plane is
\begin{equation}\label{eq:v20-Haar}
         [x^2y^2](-\log J_{\rm H})=\frac29(L^2-1),
         \qquad [x^2y^2]\ell_c=0.
\end{equation}
\end{proposition}
\begin{proof}
Use Lemma~\ref{lem:v20-insensitivity} and
\(-2\log(\sin|a|/|a|)=|a|^2/3+|a|^4/90+O(|a|^6)\).
There are \(L^4\) links per orientation. The quadratic
contribution at \(\bar c_L/L\) is \(L^2|\bar c_L|^2/3\);
its fourth part is
\((2L^2\alpha_L/3)v\cdot g=(8L^2\alpha_L/3)E_4\).
This equals \((L^2-1)E_4/9\).
The intrinsic fourth Haar term is
\(\sum_\mu|v_\mu|^4/90\).
The coarse Haar logarithm is a sum of single-orientation
functions, so has no mixed monomial on this plane.
\end{proof}

\subsection{The literal mean-fixed Wilson matrix at side two}

The longer finite component is now specified by a complete
momentum certificate. For a plaquette in orientations \(\mu<\nu\),
list its four positions and signs as
\[
 (\mu,+,0),\quad(\nu,+,e_\mu),\quad
                      (\mu,-,e_\nu),\quad(\nu,-,0).
\]
Its Wilson potential is one minus the scalar part of the
product of these four exponentials. Let
\(D_{ij}^{bc}\) be minus the scalar part of the second
variation of that product, placing colour \(b\) at position
\(i\) and colour \(c\) at position \(j\).
The dual-leg formulas \eqref{eq:v20-exp} compute this exactly:
if \(i=j\), extract \(uw\) from the one exponential; otherwise
take the \(u\) and \(w\) derivatives in their respective factors.
No factor \(1/2\) is inserted into this true Hessian.

For \(k=\pi\epsilon\), \(\epsilon\in\{0,1\}^4\setminus\{0\}\),
put \(m=|\epsilon|\), \(\lambda=4m\), and
\[
 \Pi_{\mu\nu}=\delta_{\mu\nu}-\epsilon_\mu\epsilon_\nu/m,
 \qquad P=\Pi\otimes I_3.
\]
The twelve-by-twelve Fourier Hessian has entries
\begin{equation}\label{eq:v20-Wsymbol}
 H(k)_{\alpha b,\gamma c}
   =\sum_{\mu<\nu}\ \sum_{\substack{i:\,o_i=\alpha\\j:\,o_j=\gamma}}
           (-1)^{\epsilon\cdot(z_j-z_i)}D_{ij}^{bc}.
\end{equation}
Here \(o_i,z_i\) are the orientation and offset in the
displayed position list. Define
\[
                      X(k)=\lambda^{-1}PH(k)P-P.
\]
Its longitudinal extension is zero, so
\(\log\det(I+X(k))\) computes the transverse determinant
without adding a longitudinal fluctuation. At zero \(H=\lambda P\);
the actual projected linear term \(X_{[1]}\) vanishes at every
side-two frequency, as independently proved in V16.
Therefore the mixed fourth contribution per mode is
\begin{equation}\label{eq:v20-Wmode}
 \frac12[x^2y^2]\Tr X_{[4]}
                       -\frac14[x^2y^2]\Tr X_{[2]}^2.
\end{equation}

The following table is the complete rational evaluation of
\eqref{eq:v20-Wsymbol} by \eqref{eq:v20-exp}, at
\(a=\bar c_2/2\). The labels \(s=\epsilon_1+\epsilon_2\),
\(t=\epsilon_3+\epsilon_4\) group equal mixed coefficients.
Each row has multiplicity \(\binom2s\binom2t\).
\begin{center}\small
\begin{tabular}{ccrrrr}
\toprule
\(s\)&\(t\)&multiplicity&
\([x^2y^2]\Tr X_{[4]}\)&
\([x^2y^2]\Tr X_{[2]}^2\)&contribution\\
\midrule
0&1&2&\(1/24\)&\(2/3\)&\(-7/48\)\\
0&2&1&\(-1/48\)&\(7/48\)&\(-3/64\)\\
1&0&2&\(-3/16\)&\(-5/24\)&\(-1/24\)\\
1&1&4&\(-35/192\)&\(-5/96\)&\(-5/64\)\\
1&2&2&\(-23/144\)&\(-5/216\)&\(-2/27\)\\
2&0&1&\(-119/288\)&\(17/144\)&\(-17/72\)\\
2&1&2&\(-25/81\)&\(5/81\)&\(-55/324\)\\
2&2&1&\(-293/1152\)&\(23/576\)&\(-79/576\)\\
\bottomrule
\end{tabular}
\end{center}
The third column sums to fifteen. A mixed coefficient of a
nonnegative polynomial such as \(\Tr X_{[2]}(v)^2\) can be
negative; its column is not a list of negative squared norms.
Multiplying the last column by the multiplicities and adding gives
\begin{equation}\label{eq:v20-Wanswer}
 [x^2y^2]\frac12\log\det H_2^{\rm mean}(\bar c_2)
                                               =-\frac{4135}{2592}.
\end{equation}

For reproducibility, each entry of the table is fixed by the
four-factor position list, the scalar quaternion multiplication,
the exponential polynomials \eqref{eq:v20-exp}, and the rational
projection \(P\). The checker enumerates all fifteen modes
separately before grouping them and verifies symmetry,
free normalization, vanishing linear projection, and rank nine
at each mode. Thus the table is a finite exact algebra certificate,
not a fitted formula for a momentum integral. It also computes
the independent flat coefficients
\(-4x^2/3-1577x^4/3456\) of the mean-fixed Wilson determinant,
agreeing with Shell V16 before its root return is added.

\subsection{The complete primed gauge contribution}

At a homogeneous field, the Fourier form of the logarithmic
gauge operator of Shell V16 is
\begin{equation}\label{eq:v20-Gsymbol}
 \mathcal M(k)=\sum_\mu |q_\mu|^2
                     \sigma(\operatorname{ad}_{a_\mu})
                 -i\sum_\mu\sin k_\mu\,\operatorname{ad}_{a_\mu},
 \qquad q_\mu=e^{ik_\mu}-1 .
\end{equation}
Indeed \(q_\mu^*(1+e^{ik_\mu})/2=-i\sin k_\mu\).
At side two every sine term vanishes. For the two nonzero
background orientations, the even adjoint matrices are
diagonal in colour; off their own colour axis,
\[
 \sigma(\operatorname{ad}_{a_\mu})-I
                              =-\frac{a_\mu^2}{3}-\frac{a_\mu^4}{45}
                                      +O(a_\mu^6)
\]
as an eigenvalue statement with \(a_\mu\) its signed scalar
amplitude. The inverse homogeneous source has already been
included:
\(a_1=(x+xy^2/8)/2\), \(a_2=(y+x^2y/8)/2\).

For \(\epsilon\ne0\), \(m=|\epsilon|\), this gives the
mixed coefficient of the negative gauge logarithm as
\begin{equation}\label{eq:v20-Gmode}
 [x^2y^2]\{-\log\det(\mathcal M(k)/\lambda)\}
   =\frac{\epsilon_1+\epsilon_2}{24m}
                         +\frac{\epsilon_1\epsilon_2}{144m^2}.
\end{equation}
The first term is \(-\Tr G_{[4]}\) and the second is
\(\frac12\Tr G_{[2]}^2\). Only the third colour is transverse
to both \(E_1\) and \(E_2\) in the cross product of the two
even matrices. Summation over all fifteen modes yields
\begin{equation}\label{eq:v20-Ganswer}
 [x^2y^2]\{-\log\det(I+G_2^{\rm gh})\}
                                                =\frac{6557}{20736}.
\end{equation}
The same complete sum gives the flat check
\(5x^2/8+175x^4/6912\).

\subsection{Assembly of the noncommuting coefficient}

\begin{theorem}[Complete measured side-two noncommuting fourth jet]
\label{thm:v20-complete}
In the original balanced-root source and relative to coarse
product Haar measure,
\begin{equation}\label{eq:v20-complete}
 \boxed{\quad
 J^4\mathcal Q_2^{\rm H}(xE_1,yE_2,0,0)
   =-\frac{11}{24}(x^2+y^2)
      -\frac{211}{768}(x^4+y^4)+\frac{3221}{768}x^2y^2 .
 \quad}
\end{equation}
In particular the mixed true fourth derivative is \(3221/192\).
\end{theorem}
\begin{proof}
The seven distinct contributions to the mixed monomial are
\begin{center}
\begin{tabular}{lr}
\toprule
term in the specified measured convention&\([x^2y^2]\)\\
\midrule
mean-fixed Wilson determinant&\(-4135/2592\)\\
negative logarithmic gauge determinant&\(6557/20736\)\\
negative fine Haar logarithm&\(2/3\)\\
normal-root coarea logarithm&\(5/24\)\\
Wilson-reference root tangent return&\(1/8\)\\
relative constraint-curvature logarithm&\(5153/1152\)\\
coarse Haar logarithm&\(0\)\\
\bottomrule
\end{tabular}
\end{center}
The first two were evaluated in
\eqref{eq:v20-Wanswer} and \eqref{eq:v20-Ganswer}.
The next three are
Proposition~\ref{prop:v20-Haar},
Lemma~\ref{lem:v20-normal-base}, and
Theorem~\ref{thm:v20-return}. The relative constraint term
is V16's actual full-covariance coefficient, not a mean-fixed
substitute; its exact Gaussian root calculation is rerun
by the checker. The last term follows from the product
structure of coarse Haar.
Their rational sum is \(3221/768\).

The two flat-axis coefficients also sum independently to
those shown in \eqref{eq:v20-complete}; the separate entries
agree with the flat measured regression of Shell V16.
The relative constraint term vanishes on both axes because
the actual multiplier is zero there. Constant adjoint
rotations reversing either one of the two input colours
rule out odd powers of \(x\) or \(y\), so these coefficients
exhaust the fourth jet on the plane.
No extra return is added to an already constrained determinant:
the first line of the table is explicitly mean-fixed.
\end{proof}

As a source-convention check, V19's countertransport subtracts
the relative \(5153/1152\) from the mixed one-loop coefficient.
The result in that different source chart is
\[
                  \frac{3221}{768}-\frac{5153}{1152}
                                           =-\frac{643}{2304}.
\]
The transformed observable, Haar-relative divergence and
two-loop terms of V19 must accompany that comparison.
Neither of these two finite numbers is assigned the value
of a physical beta coefficient.

\subsection{The remaining all-level problem is now an explicit mode sum}

Define
\[
 \kappa_L=[x^2y^2]\mathcal Q_L^{\rm H},\quad
 w_L=[x^2y^2]\frac12\log\det H_L^{\rm mean}(\bar c_L),\quad
 d_L=[x^2y^2]\{-\log\det(I+G_L^{\rm gh}(\bar a_L))\}.
\]
The letter \(d_L\) here is a scalar coefficient, not an exterior
derivative or Shell V16's flat-return function.
All three are finite-regulator coefficients.
The mean Wilson matrix is defined by the same four-factor
local differentiation as \eqref{eq:v20-Wsymbol}, with
\(e^{ik\cdot(z_j-z_i)}\) on the full \(L^4\) frequency grid.
Use \(\Pi=I-qq^*/|q|^2\) and
\(X=|q|^{-2}PH(k)P-P\).
At general \(L\) its linear term need not vanish, so every
trace in
\[
        J^4\Tr\left(X-\frac12X^2+\frac13X^3-\frac14X^4\right)
\]
must be retained. The gauge symbol is exactly
\eqref{eq:v20-Gsymbol}, likewise expanded through degree four.
Full complex Fourier traces equal traces of the real
operators; conjugate pairs and Nyquist modes are not
discarded or counted twice.

The all-level results above and V16 give the exact reduction
\begin{equation}\label{eq:v20-reduction}
 \boxed{\quad
 \kappa_L=w_L+d_L+\frac29(L^2-1)
       +\frac{6+50L^{-2}-56L^{-4}}{45}
       -4(\Gamma_{L,11}+\Gamma_{L,22})
                         +8(h_{L,11}+h_{L,22}) .
 \quad}
\end{equation}
Here the bounded rational expression combines the root-normal
and reference-return terms. In particular their sum tends
to \(2/15\) and cannot create a leading logarithm.
The last two terms have the leading logarithmic limit from
V18. No unknown fourth-order minimizing displacement occurs
in the other terms because of
Lemma~\ref{lem:v20-insensitivity}.

Thus the next ultraviolet acquisition is specific:
analyze, without separate absolute-value estimates destroying
the cancellations,
\begin{equation}\label{eq:v20-upstream}
                  w_L+d_L+\frac29(L^2-1).
\end{equation}
It still must be proved that its possible power divergence
cancels and that its remaining normalized logarithm has a
limit. If that normalized limit exists for \(L=2^n\), then
\eqref{eq:v20-reduction} and V18 identify
\(\lim\kappa_{2^n}/n\) by adding
\(-4(\Gamma_{*,11}+\Gamma_{*,22})+8(h_{*,11}+h_{*,22})\).
This is an implication, not a claim that the missing limit
has already been proved. The exact side-two computation
cannot supply it by extrapolation.

\subsection{The V20 companion checkpoint and provenance}

The optional suggestion file has 11,319 bytes and SHA--256
\begin{center}\scriptsize\ttfamily
4F00C5864C9B7A90DDCD2D1432EC9D8FA73B0C5D19B50BEF97C09DF7AD6CC643.
\end{center}
It motivates the choice of the noncommuting measured
calculation and a primary source convention. Its proposed
file merger is not executed; it was not a user instruction.
Shell V16 remains the same 623,261-byte companion with the
hash recorded in V19.

The required fifth-version comparison reread the current
SM V72 fixed-time many-copy transport theorem and its proof.
It holds the physical mode set fixed and does not control
a growing ultraviolet cutoff; it therefore supplies no
missing estimate for \eqref{eq:v20-upstream}.
NS V26's terminal rank-one Oseen derivative calculations
were also reread. Their radial tensor reductions do not
identify a Yang--Mills determinant, and their claimed
global numerical comparisons refer in part to a future
V27 certificate. No such certificate or norm estimate is
imported into the present calculation.
Their unchanged SHA--256 values are respectively
\begin{center}\scriptsize\ttfamily
F44F9745D43379E1672C5550411B58E97195A4C9278592D221DF5D3F644DC1F5\\
82C887F8C2C69E4F28FAD7C3DC8DE5B9099CB5A4BF431EBA1FFF3E91935978AD.
\end{center}
The next scheduled SM/NS checkpoint is V25.

The new executable certificate is
\begin{center}\small
\path{scripts/verify_ym_f14_v20_measured_quartic.py},
\end{center}
with SHA--256
\begin{center}\scriptsize\ttfamily
D709195859864718EE92648D6883D8ED50CF14108561B98050D8965809DECB77.
\end{center}
It uses only rational polynomial/quaternion arithmetic.
It differentiates the literal Wilson plaquettes and literal
balanced root, checks every side-two Fourier mode, reproduces
the normal trace certificate, and sums the complete measured
ledger. It reruns V16's actual relative-root audit. In addition,
it evaluates the nested homogeneous normal matrix independently
at \(L=2,4,8,16\), confirming the all-level formula at those
four regulators, and checks the side-two moment variance
against the complete free covariance. Those finite checks
supplement the composition and covariance proofs; they do
not prove an all-level statement by sampling.

The complete predecessor V19 has 405,377 bytes and SHA--256
\begin{center}\scriptsize\ttfamily
36E9908E9E3835EAB1A052BBA0AF9E60BFD6FCA19EB37D4563DE0CCFE42722D8.
\end{center}
Its full body is preserved apart from the title version and
this terminal insertion. Active V20 replaces V19; no companion,
archive, or monograph is removed or overwritten.

\begin{firewall}
The new full measured result is a finite local one-loop
coefficient on one noncommuting source plane, with an
all-refinement algebraic reduction of its remaining
ultraviolet analysis. It is not a proof of a cofinal
interacting quantum measure, compact-source globalization,
uniform higher-loop remainder, or physical Hamiltonian
mass gap. The full Yang--Mills goal remains unachieved.
\end{firewall}
\section{V21: cancellation of the power term and the measured quartic logarithm}
\label{sec:v21-logarithm}

V20 left a specified mean-Wilson, gauge, and fine-Haar mode sum.
The apparent \(L^2\) term in that expression comes from composing
separate quadratic factors with the cubic inverse root source.
It must be combined before estimating. This version performs
that combination, derives the remaining singular symbol from
the literal lattice Hessian, evaluates its angular coefficient
exactly, and proves a logarithm-subtracted finite-part limit.

The main conclusion, in the original balanced-root convention,
is the coefficientwise identity and convergence
\begin{equation}\label{eq:v21-main-preview}
 \left(\mathcal Q_L^{\rm H}\right)_{[4]}
       -(q_L^{\rm rel})_{[4]}
       +\frac{2}{3\pi^2}E_4\log L
                  \longrightarrow \mathcal P_{\infty,4}^{\rm rest}.
\end{equation}
The limiting polynomial is identified below by convergent
integrals and an absolutely convergent cell-difference series,
together with the already identified source-normal limit.
This is not a claim about the full interacting density at fixed
nonzero source. In particular it does not eliminate the
canonical relative source term, construct a quantum state, or
prove a Hamiltonian gap.

\subsection{Combine the mean-fixed measured factors first}

For a homogeneous fine field \(c/L\), define the mean-fixed
coefficient relative to mean-coordinate Lebesgue measure by
\begin{equation}\label{eq:v21-mean}
 \mathcal M_L(c)=
 \frac12\log\det H_L^{\rm mean}(c)
   -\log\det(I+G_L^{\rm gh}(c/L))
   -\log J_{\rm H}(c/L).
\end{equation}
All determinants are normalized at zero, with the same omitted
constant modes as V20. The linear mean constraint has only a
source-independent normal Jacobian. A homogeneous field is
stationary against mean-zero variations by translation
invariance; the formula can also simply be read as the indicated
finite local coefficient germ.

\begin{lemma}[The combined quadratic coefficient is bounded]
\label{lem:v21-quadratic}
There are real numbers \(a_L\), converging to \(a_\infty\), such that
\begin{equation}\label{eq:v21-quadratic}
 \mathcal M_{L,[2]}(c)=a_L\sum_\mu|c_\mu|^2,\qquad
 a_L=-\frac23+L^{-2}+4S_{1,L},\qquad
 a_\infty=-\frac23+4S_{1,\infty},
\end{equation}
where
\[
 \begin{aligned}
 S_{1,L}&=\sum_{m\in\mathcal K_L^3\setminus\{0\}}
  \left[
   \cosh\!\left(2L\operatorname{arsinh}
       \sqrt{\sum_{\nu=1}^3\sin^2(\pi m_\nu/L)}\right)-1
  \right]^{-1},\\
 S_{1,\infty}&=\sum_{m\in\mathbb Z^3\setminus\{0\}}
                      \frac1{\cosh(2\pi|m|)-1}.
 \end{aligned}
\]
In particular \(\sup_L|a_L|<\infty\).
\end{lemma}
\begin{proof}
Global colour rotations and signed permutations of lattice
orientations preserve the complete mean-fixed expression.
The latter act by relabeling oriented links, including the
translation and sign when an orientation is reversed.
They preserve the free Coulomb condition, the mean constraint,
the energy reference, and both measured Jacobians.
Therefore its quadratic scalar form is a multiple of
\(\sum_\mu|c_\mu|^2\): colour invariance leaves dot products
and independent orientation reflections remove off-diagonal
orientation terms.

On \(c=(xE_3,0,0,0)\), Shell V16's exact mean-fixed
factorization, before adding the root and coarse-Haar terms, is
\[
 \mathcal M_L(xE_3,0,0,0)
   =4\log\frac{\sin x}{L\sin(x/L)}
       +2\sum_{m\ne0}\log\left(1+\frac{1-\cos2x}{D_L^{\rm flat}(m)}\right)
       -2\log\frac{\sin(x/L)}{x/L}.
\]
The last factor is the unintegrated constant-coordinate
exponential volume; omitting it changes this coefficient.
Taking the second Taylor coefficient gives
\eqref{eq:v21-quadratic}. The dispersion bounds of Shell V16,
\[
 2|m|\le
 2L\operatorname{arsinh}\sqrt{\sum_\nu\sin^2(\pi m_\nu/L)}
                                                    \le2\pi|m|,
\]
bound the summands by \(C e^{-2|m|}\).
For each fixed \(m\), the dispersion tends to \(2\pi|m|\).
Dominated convergence on the centered cubes proves convergence
of \(S_{1,L}\), and hence of \(a_L\).
\end{proof}

Since \(\bar c_L(v)=v+\alpha_Lg(v)\), Lemma~\ref{lem:v21-quadratic}
gives the exact fourth-degree composition rule
\begin{equation}\label{eq:v21-compose}
 \bigl(\mathcal M_L(\bar c_L(v))\bigr)_{[4]}
       =\mathcal M_{L,[4]}(v)+8\alpha_La_LE_4(v).
\end{equation}
The second term is uniformly bounded and converges to
\(a_\infty E_4/3\). Indeed
\(D(a_L|v|^2)[\alpha_Lg]=2\alpha_La_Lv\cdot g
=8\alpha_La_LE_4\).
This is where V20's apparent power term cancels: the
separate quadratic Wilson, gauge, and fine-Haar coefficients
cannot be used independently after this composition.

\subsection{The bare quartic lattice symbol}

Use fine amplitudes \(v\), without the physical \(1/L\) scaling,
as formal source variables in this subsection. Write
\(\theta\in(-\pi,\pi]^4\), \(q_\mu=e^{i\theta_\mu}-1\),
\(\lambda=|q|^2\), and
\[
 \Pi_\theta=I-\frac{qq^*}{\lambda},\qquad
                    P_\theta=\Pi_\theta\otimes I_3
                         \quad(\theta\ne0).
\]
Let \(H(\theta;v)\) be the exact mean-fixed Wilson Hessian
constructed from the four plaquette exponentials in V20,
before projection, and define
\[
 X(\theta;v)=\lambda^{-1}P_\theta H(\theta;v)P_\theta-P_\theta.
\]
Thus \(X(0\hbox{ source})=0\); its extension to the longitudinal
three-dimensional space is zero. Let
\[
 Y(\theta;v)=\lambda^{-1}\left[
   \sum_\mu|q_\mu|^2\sigma(\operatorname{ad}_{v_\mu})
       -i\sum_\mu\sin\theta_\mu\,\operatorname{ad}_{v_\mu}
                         \right]-I_3 .
\]
This is the same logarithmic gauge symbol as V20.
Define the homogeneous quartic scalar polynomial
\begin{equation}\label{eq:v21-f}
 f(\theta;v)=
 \left[
  \frac12\Tr\log(I+X(\theta;v))
                      -\Tr\log(I+Y(\theta;v))
 \right]_{[4]} .
\end{equation}
The logarithms mean finite Taylor coefficients; no
positive-radius assertion uniform in \(\theta\to0\) is needed
to define this polynomial.

Every occurrence of a source in the physical homogeneous field
has the factor \(1/L\). It follows exactly that
\begin{equation}\label{eq:v21-bare-sum}
 \mathcal M_{L,[4]}(v)=
        L^{-4}\sum_{\theta\in(2\pi/L)\mathcal K_L^4\setminus\{0\}}
                   f(\theta;v)+\frac1{90}\sum_\mu|v_\mu|^4 .
\end{equation}
The last term is the bare fine-Haar fourth coefficient.
The trace over the entire Fourier grid is the trace of the
complexification of the real operator, and hence equals the
real trace. Conjugate pairs and self-conjugate Nyquist modes
have not been dropped or doubled.

\subsection{A uniform small-momentum expansion from the actual Hessian}

For \(\theta=rn\), \(|n|=1\), write
\[
 \Pi_n=I-nn^T,\quad P_n=\Pi_n\otimes I_3,\quad
 A_w=\operatorname{ad}_w,\quad w=\sum_\mu n_\mu v_\mu,\quad
                       \mathsf H(v)=D^2E_4(v).
\]
Traces without a specified space below are on the corresponding
finite colour or orientation-colour matrices.

\begin{lemma}[The only fourth-order nonintegrable symbol]
\label{lem:v21-symbol}
In any fixed norm on the finite-dimensional quartic-polynomial
space, uniformly for \(n\in S^3\),
\begin{equation}\label{eq:v21-symbol}
 f(rn;v)=r^{-4}F(n;v)+O(r^{-3}),\qquad r\downarrow0,
\end{equation}
where
\begin{equation}\label{eq:v21-angular}
 F(n;v)=
  -\frac14\Tr(\mathsf H P_n\mathsf H P_n)
  -2\Tr\bigl((\Pi_n\otimes A_w^2)\mathsf H\bigr)
  -\frac{23}{4}\Tr A_w^4 .
\end{equation}
As a function of \(n\) on the sphere, \(F\) is the restriction
of a polynomial of degree at most four.
\end{lemma}
\begin{proof}
Write \(X_j,Y_j,H_j\) for homogeneous source coefficients.
All entries of the literal finite plaquette derivatives
\(H_j(\theta)\) are analytic trigonometric polynomials.
The exact projected first Wilson derivative in V16 gives
\[
 X_1(\theta;v)
  =-\frac{2i}{\lambda}\,
            \Pi_\theta\otimes\operatorname{ad}_{\sum_\mu
                                               \sin\theta_\mu v_\mu}.
\]
The first-order root tangent has no effect on this identity:
the free constant block and free constant/detail cross block
are zero.

At zero momentum \(H(0;v)\) is the Hessian of the
homogeneous per-site Wilson potential. In quaternion
normalization this potential is
\[
 2\sum_{\mu<\nu}
   \left(\frac{\sin|v_\mu|}{|v_\mu|}\right)^2
   \left(\frac{\sin|v_\nu|}{|v_\nu|}\right)^2
                          |v_\mu\times v_\nu|^2 .
\]
Its leading term is \(E_4\), and it is even in total source
degree. Thus \(H_2(0)=\mathsf H\) and \(H_3(0)=0\).
In particular \(H_3(\theta)=O(r)\), while \(H_4(\theta)=O(1)\).
Using \(\lambda=r^2+O(r^4)\) and
\(P_\theta=P_n+O(r)\), uniformly in \(n\), gives
\[
 \begin{aligned}
 X_1&=r^{-1}\mathsf X+O(1),
       &\mathsf X&=-2i\Pi_n\otimes A_w,\\
 X_2&=r^{-2}\mathsf Y+O(r^{-1}),
       &\mathsf Y&=P_n\mathsf H P_n,\\
 X_3&=O(r^{-1}),&X_4&=O(r^{-2}).
 \end{aligned}
\]
The gauge formula gives
\[
 Y_1=r^{-1}\mathsf Z+O(1),\quad
 \mathsf Z=-iA_w,\qquad
                   Y_2=O(1),\quad Y_3=0,\quad Y_4=O(1).
\]

Expand the two logarithms through source degree four.
For the Wilson term its potentially nonintegrable part is
\[
 -\frac14\Tr\mathsf Y^2
       +\frac12\Tr(\mathsf X^2\mathsf Y)
       -\frac18\Tr\mathsf X^4 .
\]
For example the cubic logarithm has three cyclic words with
two \(X_1\)'s and one \(X_2\), accounting for the middle
coefficient. Terms containing \(X_4\) or \(X_1X_3\) are
only \(O(r^{-2})\); none is discarded without this estimate.
The sole nonintegrable gauge contribution is
\(+\frac14\Tr\mathsf Z^4\).
Errors in the leading \(X_1,X_2,Y_1\) expressions cost at
least one power of \(r\), hence are \(O(r^{-3})\).

Now \(\mathsf X^2=-4\Pi_n\otimes A_w^2\).
Projection identities and cyclicity yield the first two
terms of \eqref{eq:v21-angular}.
Since \(\Tr\Pi_n=3\),
\(-\frac18\Tr\mathsf X^4=-6\Tr A_w^4\).
The ghost adds \(\frac14\Tr A_w^4\), giving \(-23/4\).
Finally the first term in \eqref{eq:v21-angular} contains
two quadratic projections, the second a quadratic projection
and the quadratic matrix \(A_w^2\), and the last is quartic
in \(n\). This proves the degree assertion.
\end{proof}

This expansion uses the literal Wilson kinetic weight and the
logarithmic gauge-coordinate terms. Their integrable pieces
still enter the finite part. Replacing the whole symbol by its
homogeneous singularity would discard those finite contributions.

\subsection{Exact angular integration, not a numerical fit}

Denote normalized spherical average by
\(\langle F\rangle_{S^3}=(2\pi^2)^{-1}\int_{S^3}F\,d\omega\).

\begin{theorem}[The angular quartic polynomial]
\label{thm:v21-angular}
For all retained tuples \(v\),
\begin{equation}\label{eq:v21-angular-average}
                     \langle F(\,\cdot\,;v)\rangle_{S^3}
                                      =-\frac{16}{3}E_4(v).
\end{equation}
\end{theorem}
\begin{proof}
The spherical moments needed here are
\[
 \langle n_i n_j\rangle=\delta_{ij}/4,\qquad
 \langle n_i n_j n_k n_l\rangle
  =(\delta_{ij}\delta_{kl}+\delta_{ik}\delta_{jl}
                                      +\delta_{il}\delta_{jk})/24 .
\]
Odd moments vanish. To derive the fourth constants, rotation
in the \(1,2\) plane implies
\(\langle n_1^4\rangle=3\langle n_1^2n_2^2\rangle\);
expanding \((\sum_i n_i^2)^2=1\) then gives \(1/8\) and
\(1/24\), respectively.

For a directly reproducible finite integration rule, take
the eight points \(\{\pm e_i\}\) and the sixteen points
\(\{(\pm1,\pm1,\pm1,\pm1)/2\}\), each with weight \(1/24\).
Their odd moments vanish by sign symmetry. Their second,
pure fourth, and mixed fourth moments are exactly
\(1/4,1/8,1/24\). They therefore integrate every polynomial
of degree at most four on the sphere exactly, including
\eqref{eq:v21-angular}. This is a degree-exact cubature,
not angular sampling used to guess an integral.

Here are the matrix contractions on the two-colour plane.
Put \(P_4=x^4+y^4\), \(M_4=x^2y^2\).
For \(A_\mu=\operatorname{ad}_{v_\mu}\), the blocks of
the classical Hessian are
\[
 \mathsf H_{\mu\nu}
   =-\delta_{\mu\nu}\sum_\rho A_\rho^2
                               +2A_\nu A_\mu-A_\mu A_\nu .
\]
This follows by differentiating \(E_4\); its diagonal blocks
are the \(B_\mu\) of V20. Direct block multiplication, using
\(|n|^2=1\), also gives the explicit plane polynomial
\begin{equation}\label{eq:v21-angular-plane}
 \begin{split}
 F(n;xE_1,yE_2,0,0)
   ={}&(-16+128n_1^2-128n_1^4)x^4\\
      &+(-16+128n_2^2-128n_2^4)y^4\\
      &+\{-48+96(n_1^2+n_2^2)-256n_1^2n_2^2\}x^2y^2.
 \end{split}
\end{equation}
For example, the three mixed contributions before averaging,
in the order of \eqref{eq:v21-angular}, are
\(-48+32s-16t\), \(64s+128t\), and \(-368t\), where
\(s=n_1^2+n_2^2\), \(t=n_1^2n_2^2\).
Substitution in the moment formulas gives
\[
 \begin{array}{c|rr}
 \text{spherical average of a term of }F&P_4&M_4\\ \hline
 -\frac14\Tr(\mathsf H P_n\mathsf H P_n)&-17&-98/3\\
 -2\Tr((\Pi_n\otimes A_w^2)\mathsf H)&40&112/3\\
 -\frac{23}{4}\Tr A_w^4&-23&-46/3
 \end{array}
\]
For an algebraic check of the first row,
\[
 \begin{split}
 \Tr\mathsf H^2&=96P_4+224M_4,\\
 \Tr\left(\sum_\mu\mathsf H_{\mu\mu}\right)^2
                                      &=288P_4+288M_4,\\
 \sum_{\mu,\nu}\Tr(\mathsf H_{\mu\nu}^2)
                                      &=96P_4-64M_4.
 \end{split}
\]
Inserting these in the second and fourth spherical moments
gives the displayed first row. Directly,
\(\langle\Tr A_w^4\rangle=4P_4+8M_4/3\), which checks
the third. The same block formula gives the middle row.
Their sum is \(-32M_4/3\), with zero pure fourth term.
The exact checker reproduces all three polynomials using
the degree-exact rational cubature.

To extend this computation to all \(v\), note that the
spherical average is invariant under orthogonal changes of
the four orientations and of the three colours.
For colour reflections, \(A_w\) picks up an extra sign under
conjugation, but only its even powers occur. The Hessian
and projections have the corresponding orthogonal
covariance. A real four-by-three tuple matrix can be reduced
by these transformations to its three singular values on
the diagonal. A homogeneous quartic invariant on that
diagonal is even in each singular value and symmetric under
their permutations. It is therefore a linear combination
of their fourth powers and pairwise squared products.
The one- and two-colour calculation fixes these two
coefficients. Since \(E_4=2\sum_{i<j}s_i^2s_j^2\) on the
diagonal, the result is \(-16E_4/3\).
Polynomial identity extends the result also to complex sources.
\end{proof}

\subsection{A lattice-sum theorem with an identified finite part}

The following lemma is stated for scalar functions. It applies
coefficientwise to a finite-dimensional polynomial space.
Let \(\Theta=(-\pi,\pi]^4\), and let \(f\) be smooth on
\(\Theta\setminus\{0\}\), periodic at its boundary, with
\[
 f(\theta)=h(\theta)+O(|\theta|^{-3}),\qquad
 h(rn)=r^{-4}F(n),\quad F\in C^1(S^3).
\]
For this application \(F\) is smooth. Choose a smooth radial
cutoff \(\eta(|\theta|)\), equal to one near zero and supported
in a ball of radius less than \(\pi\). Set
\[
 R(\theta)=f(\theta)-\eta(|\theta|)h(\theta),\qquad
 A_F=\int_{S^3}F\,d\omega,\qquad
 c_\eta=\int_0^\infty
                  \frac{\eta(s)-\mathbf1_{\{s\le1\}}}{s}\,ds .
\]
Let \(Q_k=k+[-1/2,1/2]^4\), \(k\in\mathbb Z^4\), and
\(\rho_0(n)=(2\max_i|n_i|)^{-1}\), the radial boundary
of \(Q_0\).

\begin{lemma}[Logarithmic sum and convergent finite-part formula]
\label{lem:v21-sum}
The series
\[
 \delta(h)=\sum_{k\in\mathbb Z^4\setminus\{0\}}
                  \left(h(k)-\int_{Q_k}h(t)\,dt\right)
\]
is absolutely convergent, and
\begin{equation}\label{eq:v21-sum}
 \begin{split}
 &L^{-4}\sum_{\theta\in(2\pi/L)\mathcal K_L^4\setminus\{0\}}
                       f(\theta)-\frac{A_F}{(2\pi)^4}\log L\\
 &\hspace{1em}\longrightarrow
 \frac1{(2\pi)^4}\left\{
   \int_\Theta R(\theta)\,d\theta+\delta(h)+c_\eta A_F
             -\int_{S^3}F(n)\log(2\pi\rho_0(n))\,d\omega(n)
                   \right\}.
 \end{split}
\end{equation}
The right side is independent of the allowed cutoff \(\eta\).
\end{lemma}
\begin{proof}
First \(R=O(|\theta|^{-3})\) is integrable. Its punctured
Riemann sums converge to \((2\pi)^{-4}\int_\Theta R\).
To justify the puncture uniformly, shell counting gives
\[
 L^{-4}\sum_{0<|2\pi k/L|<\varepsilon}
                      |2\pi k/L|^{-3}\le C\varepsilon .
\]
Indeed \(\sum_{0<|k|\le a}|k|^{-3}\le Ca\) in four
dimensions; if \(a<1\) the sum is empty. The matching
integral has the same \(O(\varepsilon)\) bound.
Away from zero this is the ordinary bounded Riemann-sum
theorem. This proves the claim for \(R\).

Since \(h\) is homogeneous of degree \(-4\),
\(\nabla h=O(|t|^{-5})\). For all sufficiently large \(k\),
\[
 \left|h(k)-\int_{Q_k}h(t)\,dt\right|\le C|k|^{-5}.
\]
This is summable on \(\mathbb Z^4\); the finitely many
remaining cells stay away from zero. Thus \(\delta(h)\)
converges absolutely.

Put \(\eta_L(t)=\eta(2\pi|t|/L)\).
Homogeneity gives the singular sum as
\[
             (2\pi)^{-4}\sum_{k\ne0}\eta_L(k)h(k).
\]
The difference of the latter unscaled sum and
\(\int_{\mathbb R^4\setminus Q_0}\eta_L(t)h(t)\,dt\)
tends to \(\delta(h)\). For the part multiplying the cell
differences, this is dominated convergence in the absolutely
convergent series. The remaining weight error is
\[
 \sum_{k\ne0}\int_{Q_k}
                 (\eta_L(k)-\eta_L(t))h(t)\,dt .
\]
It is supported at radii comparable to \(L\), is bounded
there by \(C/L\) times \(|t|^{-4}\), and has total size
\(O(L^{-1})\). This proves the comparison.

Radial integration outside \(Q_0\) gives
\[
 \int_{S^3}F(n)\int_{\rho_0(n)}^\infty
                              \eta(2\pi r/L)\frac{dr}{r}\,d\omega(n).
\]
For all sufficiently large \(L\), the inner integral equals
\(\log L-\log(2\pi\rho_0(n))+c_\eta\).
Combining these expressions proves \eqref{eq:v21-sum}.
All integrals are well defined: \(\rho_0\) is bounded above
and away from zero. Independence of \(\eta\) follows either
by subtraction of the formulas or from their common
left-hand limit.
\end{proof}

Apply this lemma to \eqref{eq:v21-f}, with
\(h(\theta;v)=|\theta|^{-4}F(\theta/|\theta|;v)\).
Lemma~\ref{lem:v21-symbol} supplies the integrable remainder.
Every matrix expression is smooth away from zero on the
Brillouin torus, so the other hypotheses hold. There are
finitely many quartic source coefficients; hence their
finite-part limits are locally uniform polynomial limits
and commute with every fixed source derivative.

\begin{theorem}[The bare Wilson--gauge logarithm and its finite part]
\label{thm:v21-bare}
Let
\[
 \Lambda_L(v)=L^{-4}\sum_{\theta\ne0} f(\theta;v).
\]
Then
\begin{equation}\label{eq:v21-bare}
 \Lambda_L(v)+\frac{2}{3\pi^2}E_4(v)\log L
                       \longrightarrow\Lambda_\infty^{\rm ren}(v),
\end{equation}
where the limiting polynomial is exactly the right side of
\eqref{eq:v21-sum} with the indicated \(F,h,R\).
\end{theorem}
\begin{proof}
Theorem~\ref{thm:v21-angular} gives
\[
 \frac{A_F}{(2\pi)^4}
    =\frac{2\pi^2}{(2\pi)^4}
                       \left(-\frac{16}{3}E_4\right)
    =-\frac{2}{3\pi^2}E_4 .
\]
Now apply Lemma~\ref{lem:v21-sum}. In particular the
renormalized sequence is bounded; there is no residual
power divergence in this bare fourth coefficient.
\end{proof}

\subsection{Transfer to the complete measured root coefficient}

Let
\[
 Z_{\infty,4}(v)
       =(\log\det Z_\infty(v))_{[4]},\qquad
 \mathcal T_{\infty,4}(v)
       =\frac1{1440}\sum_\mu\Tr B_\mu(v)^2 .
\]
The former is the already convergent finite normal matrix
of V14; the latter is the V20 reference-return limit.

\begin{theorem}[The complete measured fourth-jet limit, with source dependence explicit]
\label{thm:v21-complete}
In the original balanced-root and coarse-Haar convention,
\begin{equation}\label{eq:v21-restlimit}
 \begin{split}
 &(\mathcal Q_L^{\rm H})_{[4]}-(q_L^{\rm rel})_{[4]}
                         +\frac{2}{3\pi^2}E_4\log L
                 \longrightarrow\mathcal P_{\infty,4}^{\rm rest},\\
 &\mathcal P_{\infty,4}^{\rm rest}
       =\Lambda_\infty^{\rm ren}
             +\frac{a_\infty}{3}E_4
             +Z_{\infty,4}+\mathcal T_{\infty,4}.
 \end{split}
\end{equation}
All these are homogeneous polynomial limits, with their
fixed derivatives. Consequently, for \(L=2^n\),
\begin{equation}\label{eq:v21-full-leading}
 \boxed{\quad
 \frac{(\mathcal Q_{2^n}^{\rm H})_{[4]}}{n}
       \longrightarrow
       -\frac{2\log2}{3\pi^2}E_4+q_{*,4}^{\rm rel}.
                                                        \quad}
\end{equation}
In particular, the complete mixed coefficient of V20 satisfies
\begin{equation}\label{eq:v21-mixed}
 \frac{\kappa_{2^n}}{n}\longrightarrow
 -\frac{4\log2}{3\pi^2}
   -4(\Gamma_{*,11}+\Gamma_{*,22})
                         +8(h_{*,11}+h_{*,22}).
\end{equation}
\end{theorem}
\begin{proof}
Use the exact measured ledger and V20's fourth-jet
insensitivity and tangent-return theorem. They give
\[
 \begin{split}
 (\mathcal Q_L^{\rm H})_{[4]}-(q_L^{\rm rel})_{[4]}
   ={}&\mathcal M_{L,[4]}(v)+8\alpha_La_LE_4(v)\\
      &+(\log\det Z_L)_{[4]}
              +\frac{\sigma_L}{2}\sum_\mu\Tr B_\mu^2+\ell_{c,[4]}.
 \end{split}
\]
By \eqref{eq:v21-bare-sum}, the first term is
\(\Lambda_L+\sum_\mu|v_\mu|^4/90\).
The coarse Haar term is the negative of precisely that
last fourth polynomial. These two intrinsic quartic
Haar terms cancel; the quadratic fine-Haar feed has
already been retained in \(a_L\).
Thus the right side is exactly
\[
 \Lambda_L+8\alpha_La_LE_4
       +(\log\det Z_L)_{[4]}
                       +\frac{\sigma_L}{2}\sum_\mu\Tr B_\mu^2.
\]
Now \(\alpha_L\to1/24\), \(a_L\to a_\infty\),
\(\sigma_L\to1/720\), and
\(\log\det Z_L\to\log\det Z_\infty\) with every fixed
source derivative by \eqref{eq:v14-Zlimit}.
Theorem~\ref{thm:v21-bare} proves
\eqref{eq:v21-restlimit}. Finally use V18's proved limit
\((q_{2^n}^{\rm rel})_{[4]}/n\to q_{*,4}^{\rm rel}\);
the finite part divided by \(n\) tends to zero.
Taking the mixed coefficient and \(E_4=2x^2y^2\) on the
plane gives \eqref{eq:v21-mixed}.
\end{proof}

V19's fully measured countertransport obeys
\((\widehat{\mathcal Q}_L^{\rm H})_{[4]}
=(\mathcal Q_L^{\rm H})_{[4]}-(q_L^{\rm rel})_{[4]}\).
Thus \eqref{eq:v21-restlimit} is also the complete
fourth-jet finite-part limit in that explicitly
countertransported one-loop convention. It does not
authorize discarding its two-loop divergence term,
its Haar-relative volume correction, or its transported
observables.

The coefficient \(-2/(3\pi^2)\) is the computed bare
mean-Wilson--gauge quartic logarithm in this normalization
and gauge convention. It has not been identified with
a universal physical beta coefficient. The original
root source retains \(q_{*,4}^{\rm rel}\), whose integral
matrices were identified in V18 but have not been numerically
evaluated here. Also, V18 only controls its normalized
logarithm; this version does not prove that
\((q_L^{\rm rel})_{[4]}-nq_{*,4}^{\rm rel}\) has a finite
limit or even a uniformly bounded remainder.

\subsection{Exact checks, preservation, and remaining acquisition}

The new checker is
\begin{center}\small
\path{scripts/verify_ym_f14_v21_measured_logarithm.py},
\end{center}
with SHA--256
\begin{center}\scriptsize\ttfamily
BF4091190B51DAAB8AD553090262373A50448D562885EFF779A9575AFB7797C7.
\end{center}
It checks all 340 coordinate monomials of degrees one
through four for the rational spherical cubature, then
computes all three angular trace polynomials using exact
fractions. Separately, it differentiates V20's literal
plaquette words at bare fine amplitudes, verifies
\(H_2(0)=D^2E_4\), and checks the zero-momentum first and
third source coefficients. The projected linear momentum
derivative is checked at all 24 rational directions against
V16's exact general identity; those checks supplement
that identity, not prove it by a finite angular sample.

The finite convention regression independently recomputes
the bare side-two Wilson and gauge matrices, obtaining
\[
 [x^2y^2]\Lambda_2=-\frac{2131}{2304},\qquad
 a_2=\frac58,\qquad16\alpha_2a_2=\frac5{16}.
\]
Their sum is \(-1411/2304\), exactly V20's
mean-Wilson plus gauge plus inverse-source fine-Haar
mixed coefficient. This checks that the power-term
combination has not changed the finite source convention.
The continuum sum limit itself is proved by
Lemma~\ref{lem:v21-sum}, not by those finite evaluations.

V20 supplied a complete noncommuting finite coefficient and
an exact all-level reduction. This version resolves
the outstanding power cancellation and leading
logarithm of that reduction, and identifies a
renormalized finite part of the measured rest.
The next canonical-source issue is the finite part
of the relative source coefficient, or, in the
countertransported convention, control of the higher
source jets and induced two-loop packet. Neither
choice removes the need for uniform remainder control,
compact-source entry, and a physical continuum construction.

The predecessor V20 has 435,404 bytes and SHA--256
\begin{center}\scriptsize\ttfamily
C117446AE5FFED53419F3E3B6860C61217716198A77E68AA7002C56994F7A0CE.
\end{center}
Its complete body is retained apart from the title
version and this terminal insertion. Active V21 replaces
V20. Shell V16, SM V72, NS V26, and the archives are
untouched. The next scheduled companion checkpoint is V25.

\begin{firewall}
The result concerns the fourth Taylor coefficient of a
local one-loop expansion at fixed physical torus size.
A coefficientwise logarithm and finite part do not
construct the full finite-source quantum measure or
control its higher-loop remainder, thermodynamic limit,
reflection-positive physical transfer, or mass gap.
The full Yang--Mills objective remains unproved.
\end{firewall}
\section{V22: the canonical quartic finite part from both scale endpoints}
\label{sec:v22-finite}

\subsection{The remaining gap in the one-loop coefficient calculation}

V21 identified the leading logarithm of the complete measured
fourth coefficient. It also proved a finite part after removing
the entire relative source term. For the original balanced
source, however, V18 only gave
\((q_{2^n}^{\rm rel})_{[4]}=nq_{*,4}^{\rm rel}+o(n)\).
An \(o(n)\) error need not be bounded, much less convergent.
This version proves that the relative remainder converges,
by retaining both ends of the shell sum.

The new input is quantitative, not a second application of
the Cesaro argument. The actual mixed-response kernels have
uniform bounded support and bounded density even when the
insertion depth grows. Their fine-mesh error is
\(O((1+\log p)/p)\) in \(L^1\). A deliberately weak Fourier
comparison rate then makes the two endpoint corrections
absolutely summable. Together with V21, this gives the finite
part of the complete measured quartic polynomial, without
changing the original source or its coarse Haar reference.

Shell-Mixing V16 and the optional suggestions remain the
same sources used in V19--V21. In particular, the complete
measured density is retained, and no parallel lineage is
merged or archived. The suggestions are proposals, not
additional instructions. SM V72 and NS V26 have not advanced
since the V20 checkpoint; the next scheduled reading is V25.

\subsection{A bounded spatial kernel for all incoming scales}

Continue to write \(L=pM\), \(p=2^j\), \(M=2^m\), with
\(m\ge1\), \(j\ge0\). All kernels in this subsection are
unfolded before periodizing to the torus of side \(M\).
Fix one output link \(a=(r,\mu)\).
For a kernel with input-orientation components \(k_\alpha(y)\)
that are colour matrices, use
\[
 \|k\|_1=\sum_\alpha\int\|k_\alpha(y)\|_{\rm op}\,dy,
 \qquad
 \|k\|_\infty=\max_\alpha
                    \operatorname*{ess\,sup}_y\|k_\alpha(y)\|_{\rm op}.
\]
The corresponding scalar component norms are bounded by these,
up to fixed colour constants. Constants below are independent
of the scale indices and of \(v\) with \(|v|\le1\).

For \(s=2^\ell\), \(\ell\ge1\), let \(k_{\ell,a}(v)\) be the
continuum kernel of the unweighted row
\[
 B_{s/2}\mathcal L_v\mathsf A_s.
\]
Here applying a row to \(\mathsf A_s\) means taking the same
linear combination of the explicit kernels \(f_{s,b}\) in
\eqref{eq:v18-filter}; it does not introduce an independent
Gaussian. For \(s\mid p\), let \(k_{\ell,a}^{(p)}(v)\)
be the cell-constant reconstruction, with multiplier \(p^3\),
of the actual fine row
\[
                         B_{s/2}\mathcal L_v B_{p/s}.
\]
Define
\begin{equation}\label{eq:v22-dkernels}
 d_{p,a}(v)=\sum_{\ell=1}^{j}2^{-\ell}
                                     k_{\ell,a}^{(p)}(v),
 \qquad
 d_a(v)=\sum_{\ell=1}^{\infty}2^{-\ell}k_{\ell,a}(v).
\end{equation}
The first sum is empty when \(p=1\).
Also denote the reconstructed row of \(B_p\) at \(a\)
by \(f_{p,a}\), and its limit \(f_{1,a}\) by \(f_a\).

\begin{lemma}[Uniform kernels and a quantified incoming error]
\label{lem:v22-kernels}
The kernels in \eqref{eq:v22-dkernels} are supported in
a fixed box depending only on the output link, not on
\(p\) or \(\ell\). They obey
\begin{equation}\label{eq:v22-unweighted}
 \begin{aligned}
 \|k_{\ell,a}^{(p)}(v)\|_1,\ \|k_{\ell,a}(v)\|_1
                                                    &\le27|v|,\\
 \|k_{\ell,a}^{(p)}(v)\|_\infty,\ \|k_{\ell,a}(v)\|_\infty
                                                    &\le864|v|,\\
 \|k_{\ell,a}^{(p)}(v)-k_{\ell,a}(v)\|_1
                                                    &\le27s|v|/p.
 \end{aligned}
\end{equation}
Consequently the infinite series defining \(d_a\)
converges in \(L^1\), \(L^2\), and \(L^\infty\), and
\begin{equation}\label{eq:v22-kerror}
 \begin{aligned}
 \|d_{p,a}(v)\|_1,\ \|d_a(v)\|_1&\le27|v|,\\
 \|d_{p,a}(v)\|_\infty,\ \|d_a(v)\|_\infty&\le864|v|,\\
 \|d_{p,a}(v)-d_a(v)\|_1&\le27(j+1)|v|/p.
 \end{aligned}
\end{equation}
Moreover, \(\|f_{p,a}-f_a\|_1=1/p\), with uniformly
bounded support, \(L^1\), and \(L^\infty\) norms.
\end{lemma}
\begin{proof}
Put \(q=s/2\) and \(t=p/s\), so \(qt=p/2\).
The row of \(B_q\) has nonnegative mass \(q\).
The total matrix coefficient norm of a row of
\(\mathcal L_v\) is at most \(54|v|\) by
\eqref{eq:v17-local-row}. The fine row of \(B_t\)
has mass \(t\). Reconstructing a fine row multiplies
its coefficient sum by \(p^3p^{-4}=p^{-1}\).
Thus the \(L^1\) norm is at most
\[
                     p^{-1}q\,54|v|\,t=27|v|.
\]
The maximum fine coefficient bound in the proof of
Lemma~\ref{lem:v17-sandwich} gives, separately at each
input orientation,
\[
             p^3\,108|v|(qt)^{-3}=864|v|.
\]
Its unfolded support has side at most \(6qt/p=3\),
up to translation to the chosen output link.
These are statements about unfolded rows; no boundedness
of a folded coefficient is being assumed.

For each inner row the limiting kernel \(f_{s,b}\)
has \(L^1\) norm \(1/s\), and the reconstructed row of
\(B_{p/s}\) differs from it by exactly \(1/p\) in \(L^1\)
by Lemma~\ref{lem:v18-filters}.
Taking the same outer linear combination therefore gives
\[
 \|k_{\ell,a}^{(p)}-k_{\ell,a}\|_1
       \le q\,54|v|/p=27s|v|/p.
\]
For fixed \(s\), pass to an almost-everywhere convergent
subsequence of this \(L^1\) convergence. The support and
pointwise bounds pass to \(k_{\ell,a}\); alternatively
the \(L^1\) bound follows directly from its finite
combination of the \(f_{s,b}\).

Now \(\sum_{\ell\ge1}2^{-\ell}=1\).
The uniform support and the displayed bounds imply
convergence in all three stated norms. Each weighted
discretization error is at most \(27|v|/p\).
There are \(j\) such errors. The omitted continuum
tail has \(L^1\) norm at most
\[
              27|v|\sum_{\ell>j}2^{-\ell}=27|v|/p.
\]
This proves \eqref{eq:v22-kerror}, including \(p=1\).
The assertions about \(f_{p,a}\) are the unit-scale
tent and box calculation of V18.
\end{proof}

In particular, \(d_a\) is an ordinary bounded compactly
supported smearing kernel, not just a Gaussian \(L^2\)
limit. It represents exactly \(\Delta(v)\) of V18:
finite partial sums represent \(\Delta^{(J)}\), and the
bilinear covariance continuity below passes their limits.
Every input-orientation matrix integral of \(d_a\) and
\(d_{p,a}\) is zero. Indeed the mixed stencil vanishes
on a second constant field, since the homogeneous
one-step map has no quadratic term. Later averages
preserve that zero. This cancellation is consistent
with, but is not required for, the estimates below.

\subsection{A summable two-cutoff comparison for the free forms}

Let \(\mathcal C_{p,M}\) and \(\mathcal R_{p,M,\nu}\)
denote the exact sums in \eqref{eq:v18-discrete-forms},
applied to reconstructed kernels. For brevity set
\[
 \begin{aligned}
 U_{p,2}(\xi)&=\Pi_p(\xi)/\omega_p(\xi)^2,&
 U_2(\xi)&=\Pi(\xi)/|\xi|^2,\\
 U_{p,3,\nu}(\xi)&=-4ip\sin(\xi_\nu/p)\Pi_p(\xi)/\omega_p(\xi)^4,&
 U_{3,\nu}(\xi)&=-4i\xi_\nu\Pi(\xi)/|\xi|^4.
 \end{aligned}
\]
The subscript \(\rho=2,3\) will mean either of these
two choices. Scalar components suffice; a fixed
finite matrix family follows by summation.

\begin{lemma}[A quantitative version of the joint free limit]
\label{lem:v22-free-rate}
Suppose \(f_p,g_p,f,g\) have support in one fixed box,
uniform \(L^1,L^2\) bounds, and
\[
                 \|f_p-f\|_1+\|g_p-g\|_1\le e_p.
\]
For \(p\ge1\), \(M\ge2\), their corresponding discrete
and continuum forms satisfy
\begin{equation}\label{eq:v22-free-rate}
 \begin{aligned}
 |\mathcal C_{p,M}(f_p,g_p)-\mathcal C(f,g)|
       &\le C(e_p+p^{-1}+M^{-1})^{1/10},\\
 |\mathcal R_{p,M,\nu}(f_p,g_p)-\mathcal R_\nu(f,g)|
       &\le C(e_p+p^{-1}+M^{-1})^{1/10}.
 \end{aligned}
\end{equation}
The constant depends only on the fixed support and
norm bounds. Periodic overlaps of the unfolded
kernels are included.
\end{lemma}
\begin{proof}
Write \(F_p^d=p^{-4}\sum_x f_{p,x}e^{-i\xi\cdot x/p}\),
and similarly for \(G_p^d\).
For \(r=|\xi|\), cell integration gives
\begin{equation}\label{eq:v22-transform-error}
 |F_p^d-\widehat f|+|G_p^d-\widehat g|
                                      \le e_p+C r/p.
\end{equation}
The support assumption bounds the first derivatives of
the continuum transforms uniformly, because their first
spatial moments are bounded. The transforms themselves
are uniformly bounded by their \(L^1\) norms.

There is a constant \(N\) depending only on the support
box such that its translates by \(M\mathbb Z^4\) overlap
at most \(N\) times for every \(M\ge2\).
Cauchy--Schwarz at each point shows that the squared
\(L^2\) norm of a periodization is at most \(N\) times
the unfolded squared norm. Therefore the discrete
Parseval estimate used in V18 remains uniform even
when the carrier folds. The same statement holds for
the continuum periodizations.

For both choices of \(\rho\), the chord estimates of
V18 give
\[
                    \|U_{p,\rho}(\xi)\|\le C r^{-\rho}.
\]
For \(0<r\le p/2\), Taylor expansion with uniform
remainders yields
\[
 q_p=i\xi/p+O(r^2/p^2),\qquad
 \|\Pi_p-\Pi\|\le Cr/p,\qquad
 \omega_p^2=r^2+O(r^4/p^2).
\]
Also \(p\sin(\xi_\nu/p)=\xi_\nu+O(r^3/p^2)\).
It follows that
\begin{equation}\label{eq:v22-symbol-error}
             \|U_{p,\rho}-U_\rho\|\le Cp^{-1}r^{1-\rho}.
\end{equation}
All estimates are uniform in direction.
The continuum symbols obey
\(\|DU_\rho\|\le Cr^{-\rho-1}\).

Choose smooth radial cutoffs that remove \(r<\epsilon\)
and \(r>2R\), and are one when \(2\epsilon\le r\le R\),
where \(0<\epsilon<1<R\) and \(2R\le p/2\).
Both the sum and the integral outside this annulus
are bounded in absolute value by
\[
                         C(\epsilon^{4-\rho}+R^{-\rho}).
\]
For the infrared part this is the lattice count in
V18 and the \(L^1\) bounds. For the ultraviolet part
it is Cauchy--Schwarz and the uniform Parseval bound,
including all the high-frequency modes.

On the retained annulus, replace the discrete
transforms and symbol by their continuum values.
Equations \eqref{eq:v22-transform-error} and
\eqref{eq:v22-symbol-error}, together with the bound
on the number of frequency cells, give error at most
\[
                         C R^4\epsilon^{-\rho}(e_p+R/p).
\]
The cutoff continuum integrand has support in a ball
of radius \(2R\) and Lipschitz constant at most
\(C\epsilon^{-\rho-1}\). Comparing its value at the
centre of each frequency cell to its cell integral
gives a Riemann error at most
\[
                             C M^{-1}R^4\epsilon^{-\rho-1}.
\]
Indeed each cell has diameter \(O(M^{-1})\);
the union of cells intersecting the support has
volume \(O(R^4)\). The cutoff makes the integrand
zero near the omitted zero mode. The normalization
is exactly \((2\pi)^{-4}\) times cell volume, as in V18.

Thus the total error is at most
\begin{equation}\label{eq:v22-cutoff-bound}
 C\left\{\epsilon^{4-\rho}+R^{-\rho}
       +R^4\epsilon^{-\rho}(e_p+R/p)
       +M^{-1}R^4\epsilon^{-\rho-1}\right\}.
\end{equation}
Put \(\delta=e_p+p^{-1}+M^{-1}\),
\(t=\delta^{1/10}\), \(\epsilon=t\), and \(R=t^{-1}\).
For sufficiently small \(\delta\), the imposed
condition \(2R\le p/2\) holds, since \(\delta\ge p^{-1}\).
The four types of terms in \eqref{eq:v22-cutoff-bound}
are bounded by constants times
\[
             t^{4-\rho},\quad t^\rho,\quad
                         t^{6-\rho}+t^{5-\rho},\quad t^{5-\rho}.
\]
For \(\rho=2,3\) every exponent is at least one.
For the remaining \(\delta\)'s, uniformly bound each
whole form by splitting at radius one, using the
same infrared estimate and Parseval. Increasing \(C\)
then proves \eqref{eq:v22-free-rate} in all cases.
\end{proof}

\begin{proposition}[Uniform two-ended shell estimate]
\label{cor:v22-bulk-rate}
There is a finite constant \(C\) such that, with the
actual shell matrices of V18 and \(\alpha=1/20\),
\begin{equation}\label{eq:v22-bulk-rate}
 \|G_{m+j,m}-\Gamma_*\|+
 \|H_{m+j,m}-h_*\|
                 \le C(2^{-\alpha m}+2^{-\alpha j}),
 \qquad m\ge1,\quad j\ge0.
\end{equation}
Any one fixed finite-dimensional matrix norm can be used.
\end{proposition}
\begin{proof}
Express \(Y_M\) through the \(f_{p,a}\), and
\(M D_M(v)\) through the \(d_{p,a}(v)\).
These are exact rows, by \eqref{eq:v17-Dsum}.
Use one fixed primitive carrier for
\(E^{[2]},E^{[3]}\), unfolding its occurrences before
periodizing. Even if two occurrences coincide on a
small torus, this prescription gives their actual
joint covariance, not independent copies.

The expectation defining \(G_{m+j,m}\) consists of
finitely many bilinear contractions of these kernels.
Use Lemma~\ref{lem:v22-free-rate}, with
Lemma~\ref{lem:v22-kernels}, on every contraction.
Uniform bounded density and \(L^1\) norm give the
required uniform \(L^2\) bounds. The same argument
applies to the finite reference-bank sum defining
\(H_{m+j,m}\), using \(f_{p,a}\) only and the response form.
The limiting contractions are precisely
\(\Gamma_*,h_*\), not new definitions of the leading
coefficient. Translation invariance removes the
average over output cells exactly as in V18.

The resulting error is bounded by
\[
 C\left(\frac{j+1}{2^j}+2^{-m}\right)^{1/10}.
\]
Since \(j+1\le C2^{j/2}\), this is at most
\(C(2^{-j/20}+2^{-m/10})\), and hence implies
\eqref{eq:v22-bulk-rate}. No optimization of the
positive exponent is needed for summability.
\end{proof}

\subsection{The two different endpoint models}

The bulk estimate alone does not specify the finite
part. We also need the limit at each fixed endpoint
distance. These limits involve two different free
models, both determined by the original lattice rows.

For a fixed \(M\ge2\), define the continuum torus forms
\begin{equation}\label{eq:v22-coarse-forms}
 \begin{aligned}
 \mathcal C_M(f,g)
   &=M^{-4}\sum_{k\in\mathbb Z^4\setminus\{0\}}
        \widehat f(2\pi k/M)^*U_2(2\pi k/M)\widehat g(2\pi k/M),\\
 \mathcal R_{M,\nu}(f,g)
   &=M^{-4}\sum_{k\in\mathbb Z^4\setminus\{0\}}
        \widehat f(2\pi k/M)^*U_{3,\nu}(2\pi k/M)
                                                   \widehat g(2\pi k/M).
 \end{aligned}
\end{equation}
The kernels are periodized; writing their unfolded
transforms at these frequencies is equivalent.
Both series are absolutely convergent by torus
Parseval outside a fixed frequency ball and the
positive spacing of the nonzero modes inside it.

For a fixed \(p\ge1\), set
\(\Theta_p=[-\pi p,\pi p)^4\) and define the
infinite-lattice, fixed-mesh forms
\begin{equation}\label{eq:v22-fine-forms}
 \begin{aligned}
 \mathcal C^{(p)}(f_p,g_p)
   &=(2\pi)^{-4}\int_{\Theta_p}
                           (F_p^d)^*U_{p,2}G_p^d\,d\xi,\\
 \mathcal R_\nu^{(p)}(f_p,g_p)
   &=(2\pi)^{-4}\int_{\Theta_p}
                           (F_p^d)^*U_{p,3,\nu}G_p^d\,d\xi .
 \end{aligned}
\end{equation}
The integrable singularity at zero is interpreted
by the ordinary improper integral. Here \(F_p^d\)
is the discrete transform, not the Fourier transform
of the interpolated kernel. This distinction is
essential at the fine endpoint.

To define endpoint matrices succinctly, use the
two finite contraction rules in
\eqref{eq:v18-Gamma-star} and \eqref{eq:v18-h-star}.
In the first, \(Y_a\) is smeared by a listed \(f_a\)
and \(\Delta_a(v)\) by a listed \(d_a(v)\), with
the listed covariance and independent Lie colours.
The second uses the listed skew response, and retains
its factor \(1/2\) and the bank entry
\((T_\mu^{\rm loc})_{ba}\).
Define
\[
 \begin{array}{c|c|c|c}
 \text{matrices}&\text{linear kernels}&\text{mixed kernels}
                                      &\text{free forms}\\ \hline
 \Gamma_m^{\rm c},h_m^{\rm c}
       &f_a&d_a&\mathcal C_{2^m},\ \mathcal R_{2^m,\nu}\\
 \Gamma_j^{\rm f},h_j^{\rm f}
       &f_{2^j,a}&d_{2^j,a}
                               &\mathcal C^{(2^j)},\ \mathcal R_\nu^{(2^j)} .
 \end{array}
\]
The reference contraction uses no mixed kernels.
For \(j=0\), \(d_{1,a}=0\); no artificial incoming
response is supplied to the finest shell.

\begin{lemma}[Endpoint limits in the actual source convention]
\label{lem:v22-edges}
The endpoint matrices just defined satisfy
\begin{equation}\label{eq:v22-edge-limits}
 \begin{aligned}
 \lim_{j\to\infty}G_{m+j,m}&=\Gamma_m^{\rm c},&
 \lim_{j\to\infty}H_{m+j,m}&=h_m^{\rm c}
                                              &&(m\text{ fixed}),\\
 \lim_{m\to\infty}G_{m+j,m}&=\Gamma_j^{\rm f},&
 \lim_{m\to\infty}H_{m+j,m}&=h_j^{\rm f}
                                              &&(j\text{ fixed}).
 \end{aligned}
\end{equation}
In particular,
\begin{equation}\label{eq:v22-edge-tails}
 \begin{aligned}
 \|\Gamma_m^{\rm c}-\Gamma_*\|+\|h_m^{\rm c}-h_*\|
                                                   &\le C2^{-\alpha m},\\
 \|\Gamma_j^{\rm f}-\Gamma_*\|+\|h_j^{\rm f}-h_*\|
                                                   &\le C2^{-\alpha j}.
 \end{aligned}
\end{equation}
\end{lemma}
\begin{proof}
First fix \(M=2^m\) and let \(p\) grow. At every fixed
nonzero torus frequency, the discrete transforms
converge to the transforms of \(f_a,d_a\) by
\eqref{eq:v22-transform-error} and
\eqref{eq:v22-kerror}, and the symbols converge
to \(U_2,U_{3,\nu}\). The tail beyond radius \(R\)
is bounded uniformly in \(p\) by \(CR^{-\rho}\),
using the periodized Parseval estimate from the
proof of Lemma~\ref{lem:v22-free-rate}.
For fixed \(M\), the remaining frequency ball has
finitely many nonzero modes. First pass to the
limit on that finite set, then let \(R\) grow.
This proves convergence of every covariance and
response contraction to \eqref{eq:v22-coarse-forms}.
It proves the first row of \eqref{eq:v22-edge-limits},
including all periodic coincidences.

Next fix \(p=2^j\). The unfolded \(f_{p,a},d_{p,a}\)
are now fixed finite rows. As \(M\) grows their exact
frequency sums become Riemann sums on \(\Theta_p\).
The transforms and lattice symbols are periodic on
this Brillouin torus and smooth off zero. Cut out
a ball of radius \(\epsilon\); the error from that
ball is at most \(C\epsilon^{4-\rho}\), uniformly
in \(M\), by the same lattice count. On its smooth
complement the Riemann sums converge to the integral.
Use smooth cutoffs if necessary at the boundary
of the removed ball, and then send \(\epsilon\)
to zero. Opposite Brillouin faces match periodically;
no boundary mode is dropped. This gives
\eqref{eq:v22-fine-forms} and the second row of
\eqref{eq:v22-edge-limits}.

Finally pass one cutoff at a time in
\eqref{eq:v22-bulk-rate}. The surviving terms give
\eqref{eq:v22-edge-tails}.
\end{proof}

\subsection{An endpoint summation lemma and the relative finite part}

\begin{lemma}[Two-ended finite part, with the correct index origins]
\label{lem:v22-sum}
Let \(a(m,j)\) take values in a finite-dimensional
normed space, for \(m\ge1,j\ge0\). Suppose
\[
 \|a(m,j)-a_*\|\le C(2^{-\alpha m}+2^{-\alpha j}),
 \quad
 a^{\rm c}_m=\lim_{j\to\infty}a(m,j),
 \quad
 a^{\rm f}_j=\lim_{m\to\infty}a(m,j)
\]
exist, with \(\alpha>0\). Then both endpoint series
below converge absolutely and
\begin{equation}\label{eq:v22-sum}
 \sum_{m=1}^n a(m,n-m)-n a_*
     \longrightarrow
       \sum_{m=1}^{\infty}(a^{\rm c}_m-a_*)
          +\sum_{j=0}^{\infty}(a^{\rm f}_j-a_*).
\end{equation}
\end{lemma}
\begin{proof}
Passing one cutoff to infinity gives the corresponding
geometric bounds for \(a^{\rm c}_m-a_*\) and
\(a^{\rm f}_j-a_*\), so both series converge absolutely.
Fix \(b\ge1\), and take \(n>2b+1\). The portions of
the finite sum with \(m\le b\) and with \(j=n-m\le b\)
are disjoint. They converge, term by term, to the
two endpoint series truncated at \(b\).
For the intervening indices, both \(m>b\) and \(j>b\),
and hence
\[
 \sum_{\substack{m+j=n\\m>b,\ j>b}}\|a(m,j)-a_*\|
  \le C\left(\sum_{m>b}2^{-\alpha m}
                        +\sum_{j>b}2^{-\alpha j}\right)
  \le C'2^{-\alpha b}.
\]
The tails of both endpoint series obey the same bound.
First let \(n\) grow with \(b\) fixed, and then let \(b\)
grow. This proves \eqref{eq:v22-sum}.
\end{proof}

The geometric hypothesis is substantive. The bounded
array \(a(m,j)=a_*+(m+j)^{-1/2}\) has the same bulk
and both endpoint limits \(a_*\), but its summed
remainder is \(\sqrt n\). Thus the endpoint limits
and V18's bulk convergence, without a summable bound,
would not have proved the following theorem.

\begin{theorem}[Finite parts of the actual source and reference matrices]
\label{thm:v22-relative}
Define the absolutely convergent matrix series
\begin{equation}\label{eq:v22-matrix-fin}
 \begin{aligned}
 \Gamma_{\rm fin}
   &=\sum_{m=1}^{\infty}(\Gamma_m^{\rm c}-\Gamma_*)
           +\sum_{j=0}^{\infty}(\Gamma_j^{\rm f}-\Gamma_*),\\
 h_{\rm fin}
   &=\sum_{m=1}^{\infty}(h_m^{\rm c}-h_*)
           +\sum_{j=0}^{\infty}(h_j^{\rm f}-h_*).
 \end{aligned}
\end{equation}
Then
\begin{equation}\label{eq:v22-matrix-limit}
          \Gamma_{2^n}-n\Gamma_*\longrightarrow\Gamma_{\rm fin},
 \qquad
          h_{2^n}-nh_*\longrightarrow h_{\rm fin}.
\end{equation}
Consequently, with
\begin{equation}\label{eq:v22-qfin}
 q_{{\rm fin},4}^{\rm rel}(v)
   =-\sum_{\mu,\nu}(\Gamma_{\rm fin})_{\mu\nu}
                                      g_\mu(v)\cdot v_\nu
      +2\sum_{\mu,\nu}(h_{\rm fin})_{\nu\mu}
                                      g_\mu(v)\cdot v_\nu ,
\end{equation}
the original relative fourth coefficient has the finite part
\begin{equation}\label{eq:v22-q-limit}
 (q_{2^n}^{\rm rel})_{[4]}-nq_{*,4}^{\rm rel}
                                \longrightarrow q_{{\rm fin},4}^{\rm rel}.
\end{equation}
These are coefficientwise polynomial limits, locally
uniform in the source and in every fixed source derivative.
\end{theorem}
\begin{proof}
Apply Lemma~\ref{lem:v22-sum} to the pair of matrices
\((G_{m+j,m},H_{m+j,m})\), using
Proposition~\ref{cor:v22-bulk-rate} and
Lemma~\ref{lem:v22-edges}. V18's decompositions of
\(\Gamma_{2^n}\) and \(h_{2^n}\) are exact, so this gives
\eqref{eq:v22-matrix-limit}. Substitute these matrix
limits in the exact V16 quartic formula
\[
 (q_L^{\rm rel})_{[4]}
    =-\sum_{\mu,\nu}(\Gamma_L)_{\mu\nu}g_\mu\cdot v_\nu
            +2\sum_{\mu,\nu}(h_L)_{\nu\mu}g_\mu\cdot v_\nu .
\]
There are finitely many coefficients. Their convergence
proves the final polynomial assertions.
\end{proof}

Equations \eqref{eq:v22-coarse-forms},
\eqref{eq:v22-fine-forms}, and \eqref{eq:v22-matrix-fin}
specify the finite constants by convergent series and
integrals of the actual filters and primitive words.
They have not been assigned fitted numerical values.
The endpoint-series tail after index \(b\) is bounded
by \(C'2^{-b/20}\); this deliberately weak bound is
not a claim of an efficient numerical evaluation.

\subsection{The complete measured quartic coefficient in the original root}

\begin{theorem}[Canonical one-loop quartic finite part]
\label{thm:v22-complete}
Set
\[
 \mathcal L_4(v)=-\frac{2\log2}{3\pi^2}E_4(v)
                                      +q_{*,4}^{\rm rel}(v),
 \qquad
 \mathcal P_{{\rm fin},4}
        =\mathcal P_{\infty,4}^{\rm rest}
                                      +q_{{\rm fin},4}^{\rm rel}.
\]
In the original balanced-root source \(v\), with the
complete measured coarse-Haar convention of Shell V16,
\begin{equation}\label{eq:v22-full}
 \boxed{\quad
           (\mathcal Q_{2^n}^{\rm H})_{[4]}-n\mathcal L_4
                         \longrightarrow\mathcal P_{{\rm fin},4}.
                                                                  \quad}
\end{equation}
The limit is the explicit homogeneous polynomial
\[
 \mathcal P_{{\rm fin},4}
    =\Lambda_\infty^{\rm ren}
        +\frac{a_\infty}{3}E_4+Z_{\infty,4}
                 +\frac1{1440}\sum_\mu\Tr B_\mu^2
                 +q_{{\rm fin},4}^{\rm rel}.
\]
On \(v=(xE_1,yE_2,0,0)\), let \(\kappa_*\) be the
leading coefficient in \eqref{eq:v21-mixed}. Then
\begin{equation}\label{eq:v22-plane}
 \begin{split}
 \kappa_{2^n}-n\kappa_*\longrightarrow{}&
       [x^2y^2]\Lambda_\infty^{\rm ren}
           +\frac{2a_\infty}{3}+\frac2{15}\\
       &-4(\Gamma_{{\rm fin},11}+\Gamma_{{\rm fin},22})
                    +8(h_{{\rm fin},11}+h_{{\rm fin},22}).
 \end{split}
\end{equation}
\end{theorem}
\begin{proof}
Add \eqref{eq:v22-q-limit} to V21's
\eqref{eq:v21-restlimit}, using \(\log(2^n)=n\log2\).
No term is removed from the measured ledger.
This proves \eqref{eq:v22-full} and its displayed
polynomial formula. On the chosen plane \(E_4=2x^2y^2\);
V20 gives mixed coefficients \(4/45\) for \(Z_{\infty,4}\)
and \(2/45\) for the reference-return limit.
Their sum is \(2/15\). The contractions in
\eqref{eq:v22-qfin} give the last two terms of
\eqref{eq:v22-plane}.
\end{proof}

This result removes the outstanding finite-part issue
at degree four in the canonical source. It also fixes
the finite difference between that coefficient and
V19's explicitly countertransported one-loop coefficient.
It does not supply uniform bounds for the nonlinear
source flow, and does not omit its divergence or
observable transport at the next loop order.

\subsection{Checks, ownership, and the next upstream target}

The exact checker
\begin{center}\small
\path{scripts/verify_ym_f14_v22_source_finite_part.py}
\end{center}
has SHA--256
\begin{center}\scriptsize\ttfamily
8334A316EBFF4254B997ED1FD6051DAD6E86D741B0F3FFE00CA89CB4CE38D0AA.
\end{center}
It constructs genuine mixed rows from V17, with outer
averages and fine ratios \(p=2,4,8\), without wrapping
their carrier. For each incoming level it compares
the reconstructed density with the exact piecewise
affine continuum tent combination. The \(L^1\)
errors are integrated with rational arithmetic,
including cells on which the discrete tent vanishes
but the continuum tent does not. The checks verify
the bounds used above, componentwise zero masses,
all cutoff exponents, and the different index
origins in the two-ended sum. They do not certify
an infinite-cutoff limit by a finite sample;
that limit is proved in the preceding lemmas.

The predecessor V21 has 461,169 bytes and SHA--256
\begin{center}\scriptsize\ttfamily
902237A01FD349EE33A6F4680417F8CCBADA8BFAF78A3BAADAEA1C8F962D9B15.
\end{center}
Its entire body is retained apart from the title
version and this terminal insertion. Active V22
replaces V21. Companions and archives are untouched.
The next regular cross-program checkpoint is V25.

The useful next target is no longer another proof
that the quartic coefficient has a logarithm.
The available finite part is still a fourth Taylor
coefficient on the local \(\SU(2)\) branch.
To turn this into a measured quantum renormalization
statement requires control beyond that coefficient:
the remaining source jets or a uniform finite-source
remainder, and then the actual two-loop packet.
The original-root convention remains primary.
V19's countertransport can provide a consistency
check only when its Jacobian and higher-jet terms
are retained.

\begin{firewall}
No interacting continuum quantum law, reflection-positive
physical transfer, or positive Hamiltonian mass gap
is constructed by this finite-part theorem.
The Yang--Mills existence and mass-gap problem remains
listed as unsolved in the
\href{https://www.claymath.org/millennium/yang-mills-the-maths-gap/}
{Clay Mathematics Institute problem statement}
(checked 14 September 2026). The new result is the
canonical measured quartic one-loop finite part
within the specified local programme, not the full
Yang--Mills objective.
\end{firewall}
\section{V23: the Gaussian domain obstruction and a quadratic rough-source limit}
\label{sec:v23-domain}

\subsection{Why the next estimate must change its domain}

V22 completes the canonical measured quartic finite part.
It does not control the integral remainder by placing all
fluctuations in the classical branch's common global
\(\ell^4\) ball. This version checks that proposed
interface directly. At fixed physical torus size, the
actual free reference Gaussian loses essentially all
weight in that ball when \(\beta=o(L^2)\). This remains
true in a regime where every individual fine link stays
inside a fixed identity chart with probability tending
to one.

Shell V12 already warned qualitatively that a global
\(\ell^4\) domain is restrictive for statistical mechanics.
We do not claim that warning as new. The additional
results here are quantitative ultraviolet thresholds
in the actual transverse covariance, a bound that
survives the nonlinear constant-source pivot, and a
precise entropy cost for restoring concentration.
They identify why a deterministic classical remainder
cannot simply be integrated against all fine Gaussian
modes in the cofinal regime.

There is also a positive stochastic result. V15's
actual quadratic-root banks define a convergent
second Gaussian chaos, with uniform small exponential
moments. This is a rigorous polynomial observable
outside the global classical domain. It is not
evaluation of the complete nonlinear root on a
distribution-valued field.

The active predecessor V22 and the companion snapshots
are unchanged. No new SM/NS checkpoint is due before
V25. The original balanced source and complete measured
coarse-Haar convention remain primary.

\subsection{The actual zero-source Gaussian and its critical counting norm}

Let \(L\ge2\) be dyadic. Write
\[
 \lambda_L(k)=4\sum_{\mu=1}^4\sin^2(\pi k_\mu/L),
 \qquad k\in(\mathbb Z/L\mathbb Z)^4.
\]
The scalar covariance \(C_L\) in this section is the
free zero-mean Coulomb covariance of V14--V18, not
V19's four-dimensional orientation matrix with the
same letter. At each \(k\ne0\) it has symbol
\(\Pi_L(k)/\lambda_L(k)\), where \(\Pi_L\) is the
rank-three transverse projector. It is zero on the
removed mean and longitudinal modes.

Let \(Z_L\) have covariance \(C_L\otimes I_3\).
All three Lie colours are retained, including global
colour rotations. Its real dimension is
\[
                         d_L=9(L^4-1).
\]
There are \(4L^4\) fine links. For a Lie-valued link
field use the unnormalized counting norms
\[
 \|a\|_4^4=\sum_e|a_e|^4,\qquad
 \|a\|_2^2=\sum_e|a_e|^2.
\]
The physical cellwise reconstruction \(A(y)=La_e\)
on cells of volume \(L^{-4}\) has
\(\|A\|_{L^4}=\|a\|_4\). Thus this is the same critical
global norm used in the classical source theorem,
not an additional volume-normalized smallness condition.

Translation and coordinate-permutation invariance give
one common scalar link variance
\begin{equation}\label{eq:v23-c}
 c_L=(C_L)_{ee}
     =\frac{3}{4L^4}\sum_{k\ne0}\lambda_L(k)^{-1}.
\end{equation}
There are constants \(0<c_-\le c_+<\infty\),
independent of \(L\), such that
\[
 c_-\le c_L\le c_+,\qquad
 c_L\longrightarrow
 c_\infty=\frac34(2\pi)^{-4}
       \int_{[-\pi,\pi]^4}
          \frac{d\theta}{4\sum_\mu\sin^2(\theta_\mu/2)}>0.
\]
Indeed \(\lambda_L\le16\) gives
\(c_L\ge3(1-L^{-4})/64\).
For centred representatives,
\(\lambda_L(k)\ge16|k|^2/L^2\).
The four-dimensional annular count bounds the upper
sum uniformly. The same estimate bounds the contribution
of frequencies \(|2\pi k/L|<\epsilon\) by \(C\epsilon^2\).
Off that ball the sum is a Riemann sum for the displayed
integrable symbol. This proves convergence without
discarding boundary or self-conjugate modes.

\begin{theorem}[Concentration of the actual global fourth norm]
\label{thm:v23-fourth}
Set \(X_L=\|Z_L\|_4^4\). Then
\begin{equation}\label{eq:v23-fourth}
 \mathbb E X_L=60c_L^2L^4,\qquad
 \operatorname{Var}(X_L)\le C L^4(1+\log L),
 \qquad
 L^{-4}X_L\longrightarrow60c_\infty^2
                                              \quad\hbox{in }L^2.
\end{equation}
Consequently, for \(a_L=\beta_L^{-1/2}Z_L\):
if \(\beta_L/L^2\to0\), then \(\|a_L\|_4\to\infty\)
in probability; if \(\beta_L/L^2\to\infty\), then
\(\|a_L\|_4\to0\) in probability. If
\(\beta_L/L^2\to b\in(0,\infty)\), then
\[
                   \|a_L\|_4^4\longrightarrow60c_\infty^2/b^2
                                              \quad\hbox{in probability}.
\]
\end{theorem}
\begin{proof}
For a three-dimensional Gaussian of covariance \(cI_3\),
the fourth radial moment is \(15c^2\).
For two such vectors \(U,V\) with
\(\mathbb E U_aV_b=r\delta_{ab}\), the exact pairing identity is
\begin{equation}\label{eq:v23-Wick}
 \operatorname{Cov}(|U|^4,|V|^4)=600c^2r^2+120r^4.
\end{equation}
One derivation expands
\(\sum_{a,b,c,d}U_a^2U_b^2V_c^2V_d^2\) into pairings.
The terms with zero, two, and four cross-pairs are
respectively \(225c^4\), \(600c^2r^2\), and \(120r^4\);
subtract the product \(225c^4\) of the means.

Apply this with \(c=c_L\) and \(r=(C_L)_{ef}\).
Positivity of the covariance implies \(|r|\le c_L\).
Summing \eqref{eq:v23-Wick} over all fine links gives
\[
 \operatorname{Var}(X_L)
   =600c_L^2\sum_{e,f}(C_L)_{ef}^2
                +120\sum_{e,f}(C_L)_{ef}^4
   \le720c_L^2\Tr C_L^2.
\]
The scalar trace includes all three transverse modes:
\[
 \Tr C_L^2=3\sum_{k\ne0}\lambda_L(k)^{-2}
       \le C L^4
           \sum_{\substack{k\ne0\\|k|\le L}}|k|^{-4}
       \le C L^4(1+\log L).
\]
The harmless constant in the frequency radius absorbs
the centred four-cube. The annular count proves the
last inequality. Together with \eqref{eq:v23-c}, this
gives \eqref{eq:v23-fourth}. Finally
\(\|a_L\|_4^4=(L^2/\beta_L)^2(L^{-4}X_L)\).
Use convergence to a strictly positive constant;
for the case \(\beta_L/L^2\to\infty\), Markov's
inequality also proves the assertion directly.
\end{proof}

\subsection{A small-ball bound that retains the source pivot}

Let \(\gamma_{\beta,L}\) be the law of
\(\beta^{-1/2}Z_L\) on the real detail space.
For \(r>0\) set
\[
 x_L(r,\beta)=\frac{32\beta L^2r^2}{d_L},
 \qquad I(x)=x-1-\log x.
\]

\begin{theorem}[Quantified loss of the global classical domain]
\label{thm:v23-smallball}
Whenever \(0<x_L(r,\beta)<1\),
\begin{equation}\label{eq:v23-smallball}
 \gamma_{\beta,L}\{\|w\|_4\le r\}
       \le \exp\{-\tfrac12d_L I(x_L(r,\beta))\}.
\end{equation}
More generally, let a constant-source pivot be defined
on any measurable domain \(\mathcal D\) by
\[
             a(w)_{x,\mu}=s_\mu(w)/L+w_{x,\mu},
 \qquad \sum_xw_{x,\mu}=0.
\]
No extension outside \(\mathcal D\) is required.
The same right side bounds
\begin{equation}\label{eq:v23-pivot-event}
       \gamma_{\beta,L}
            \{w\in\mathcal D:\|a(w)\|_4\le r\}.
\end{equation}
In particular, this applies to the actual nonlinear
balanced-source pivot wherever its local chart is defined.
\end{theorem}
\begin{proof}
In real counting-\(\ell^2\) orthonormal coordinates
on the detail space, write
\[
        w=\beta^{-1/2}(C_L\otimes I_3)^{1/2}\xi,
                          \qquad \xi\sim N(0,I_{d_L}).
\]
Every nonzero covariance eigenvalue is
\(\lambda_L(k)^{-1}\ge1/16\). Hence
\[
                  \|w\|_2^2\ge|\xi|^2/(16\beta).
\]
On the other hand Holder's inequality on the \(4L^4\)
links gives \(\|a\|_2^2\le2L^2\|a\|_4^2\).
For the pivot, the constant part is orthogonal
to the mean-zero detail part in \(\ell^2\), so
\[
               \|w\|_2^2\le\|a(w)\|_2^2.
\]
Therefore either event in the theorem implies
\[
                         |\xi|^2\le32\beta L^2r^2=d_Lx_L.
\]
For any \(t>0\), direct integration of independent
one-dimensional Gaussians gives
\[
 \mathbb P\{|\xi|^2\le d_Lx\}
       \le e^{td_Lx}\mathbb E e^{-t|\xi|^2}
       =e^{td_Lx}(1+2t)^{-d_L/2}.
\]
For \(0<x<1\), minimize at \(t=(x^{-1}-1)/2\).
The result is \(\exp\{-d_L I(x)/2\}\).
This proves both assertions. The normal variable
\(s(w)\) was never fixed artificially or removed
from the constraint.
\end{proof}

For every fixed \(r\), if \(\beta_L=o(L^2)\) then
\[
 -\log\gamma_{\beta_L,L}\{\|w\|_4\le r\}
       \ge \frac{d_L}{2}
                   \bigl\{\log(L^2/\beta_L)-O_r(1)\bigr\}.
\]
The same lower bound on the negative logarithm
holds for \eqref{eq:v23-pivot-event}.
The \(O_r(1)\) is an explicit bounded term from
\(d_L/L^4\), \(32r^2\), and \(x_L-1\).
The probability thus decays more rapidly than
\(\exp(-cL^4)\) for every fixed \(c\) eventually.
This is an ultraviolet statement at fixed physical
volume, not a thermodynamic limit.

A fixed global energy ball is still more restrictive.
At zero source,
\[
                 \|d_Lw\|_2^2=\beta^{-1}|\xi|^2
\]
exactly. Its mean is \(d_L/\beta\), and the same
chi-square estimate applies with \(x=\beta R^2/d_L\)
to the event \(\|d_Lw\|_2\le R\).
Uniform control of the energy of the classical
minimum therefore says nothing by itself about
typical fine fluctuation energy.

\begin{proposition}[Pointwise chart control does not repair the global norm]
\label{prop:v23-pointwise}
For every \(r_0>0\),
\begin{equation}\label{eq:v23-max}
 \gamma_{\beta,L}\{\max_e|w_e|>r_0\}
                \le24L^4\exp\{-\beta r_0^2/(6c_+)\}.
\end{equation}
Consequently, with \(\beta_L=A\log L\) and
\(A>24c_+/r_0^2\), every link lies in the fixed
pointwise ball with probability tending to one,
while every fixed global \(\ell^4\) ball has
probability tending to zero.
\end{proposition}
\begin{proof}
If a three-dimensional vector has norm greater
than \(r_0\), at least one component has modulus
greater than \(r_0/\sqrt3\). A component of \(w_e\)
has variance \(c_L/\beta\), and its two-sided
Gaussian tail is at most
\(2\exp\{-\beta r_0^2/(6c_L)\}\).
Sum over three colours and \(4L^4\) links.
This uses no independence between links.
The stated condition makes the power of \(L\)
in \eqref{eq:v23-max} negative. Since
\(A\log L=o(L^2)\), Theorem~\ref{thm:v23-smallball}
gives the global assertion.
\end{proof}

This logarithmic regime is a mathematical comparison
regime, not a proof identifying the interacting
Yang--Mills coupling trajectory.
Pointwise group logarithms being available is
different from the global analytic norm required
by Shell V12--V13.

\subsection{How expensive a measure change would have to be}

\begin{proposition}[A conditional entropy test for restoring branch weight]
\label{prop:v23-entropy}
Let \(B_L\) be either event in
Theorem~\ref{thm:v23-smallball}, and put
\(A_L=d_L I(x_L)/2>0\).
For any probability measure \(\nu_L\ll\gamma_{\beta,L}\),
\begin{equation}\label{eq:v23-entropy}
 \nu_L(B_L)
       \le\frac{\operatorname{Ent}(\nu_L\mid\gamma_{\beta,L})+\log2}
                    {A_L}.
\end{equation}
In particular, if \(\beta_L=o(L^2)\) and the relative
entropy is \(O(L^4)\), then \(\nu_L(B_L)\to0\).
If instead
\[
 d\nu_L=\frac{e^{-V_L}}{\mathbb E_\gamma e^{-V_L}}\,
                                                   d\gamma_{\beta,L}
\]
with finite oscillation
\(\operatorname{osc}V_L=\sup V_L-\inf V_L\), then
\[
             \nu_L(B_L)\le
                     \exp\{\operatorname{osc}V_L-A_L\}.
\]
\end{proposition}
\begin{proof}
Write \(p=\gamma(B_L)\), \(q=\nu(B_L)\).
Jensen's inequality applied to the density separately
on \(B_L\) and its complement gives the binary
relative entropy bound
\[
 \operatorname{Ent}(\nu\mid\gamma)
       \ge q\log(q/p)+(1-q)\log((1-q)/(1-p))
       \ge q\log(1/p)-\log2.
\]
Use \(p\le e^{-A_L}\).
Cases \(p=0\) or infinite entropy are interpreted
by absolute continuity or the trivial bound.
When \(\beta_L=o(L^2)\), \(A_L/L^4\to\infty\),
which proves the stated consequence.
For a bounded-oscillation tilt,
\(d\nu/d\gamma\le e^{\operatorname{osc}V_L}\);
integrate over \(B_L\).
\end{proof}

No bound of this type has been asserted for the
full compact Yang--Mills law. The proposition
does not disprove that law or constrain its
existence by assuming a Gaussian comparison.
It gives a checkable necessary cost for any
proposed change of measure that is supposed
to make the same global classical domain typical.

There is a second mismatch with a direct use of
Shell V7's fixed-box theorem. In the present
one-coarse-cell elimination, the zero-source
Hessian in local fine \(\ell^2\) coordinates has
smallest eigenvalue
\[
                         4\sin^2(\pi/L)\longrightarrow0.
\]
The bound is attained at \(k=(1,0,0,0)\) in
a transverse orientation. Thus its uniformly
positive spatial Hessian hypothesis cannot be
imported unchanged to this all-fine-at-once
problem. Whitening makes the energy matrix the
identity, but does not preserve the local
coordinate stencils and their uniform spatial
constants required by that theorem.
This does not contradict its shell-local result.

\subsection{A genuine quadratic Gaussian source outside the classical ball}

We now use the stronger structure that was actually
proved for the balanced source. Let
\(\mathcal V\) be the common real scalar energy
Hilbert space in V15, and let
\(S_{L,\mu}\), \(S_{\infty,\mu}\) be its real skew
quadratic-source banks. Recall the established facts
\[
 \|S_{L,\mu}\|_{\mathfrak S_2},
 \|S_{\infty,\mu}\|_{\mathfrak S_2}\le B_*=168,\qquad
       S_{L,\mu}\to S_{\infty,\mu}\quad\hbox{in }\mathfrak S_2.
\]
This subsection does not reprove that operator
theorem or claim a stronger ideal for the full
finite-source family.

Choose three independent isonormal Gaussian families
\(W^1,W^2,W^3\) over \(\mathcal V\): these are linear
Gaussian functionals with
\(\mathbb E W^a(f)W^b(g)=\delta_{ab}\langle f,g\rangle\).
They are not asserted to be \(\mathcal V\)-valued
random vectors. For \(a\ne b\) and a Hilbert--Schmidt
operator \(S\), define
\[
 I_{ab}(S)=\sum_{i,j}\langle e_i,Se_j\rangle
                                      W^a(e_i)W^b(e_j)
\]
as an \(L^2\) limit. Independence gives the isometry
\begin{equation}\label{eq:v23-isometry}
 \mathbb E I_{ab}(S)I_{ab}(T)
                  =\langle S,T\rangle_{\mathfrak S_2}.
\end{equation}
It proves convergence and independence from the
chosen real orthonormal basis by finite-rank
approximation.

Define, for cyclic \((c,a,b)=(1,2,3),(2,3,1),(3,1,2)\),
\[
 \mathcal Y_{L,\mu}^c=2I_{ab}(S_{L,\mu}),\qquad
 \mathcal Y_{\infty,\mu}^c=2I_{ab}(S_{\infty,\mu}).
\]

\begin{theorem}[Stochastic limit of the quadratic root]
\label{thm:v23-chaos}
The twelve-dimensional random vector \(\mathcal Y_L\)
has the law of the homogeneous quadratic polynomial
\(\mathcal F_{L,[2]}(Z_L)\) of the actual balanced root.
On the common coupling just defined,
\begin{equation}\label{eq:v23-chaos-limit}
 \mathbb E|\mathcal Y_L-\mathcal Y_\infty|^2
        =12\sum_{\mu=1}^4
                  \|S_{L,\mu}-S_{\infty,\mu}\|_{\mathfrak S_2}^2
                                      \longrightarrow0.
\end{equation}
Its covariance, including all mixed shell terms, is
\begin{equation}\label{eq:v23-chaos-cov}
 \mathbb E\mathcal Y_{L,\mu}^c\mathcal Y_{L,\nu}^d
                       =4\delta_{cd}D_{L,\mu\nu}
           \longrightarrow4\delta_{cd}D_{\infty,\mu\nu}.
\end{equation}
There is a common positive exponential radius:
\begin{equation}\label{eq:v23-chaos-exp}
 \sup_{L\in\{2,4,\ldots,\infty\}}
           \mathbb E e^{\eta|\mathcal Y_L|}<\infty,
                           \qquad 0\le\eta\le1/8064.
\end{equation}
In particular the convergence holds in every finite \(L^q\),
\(q\ge1\). These assertions require no global \(\ell^4\)
smallness of \(Z_L\).
\end{theorem}
\begin{proof}
In ordinary scalar link coordinates the actual
quadratic polynomial is
\[
 \mathcal F_{L,[2],\mu}(Z)
       =\sum_{i<j}(T_{L,\mu})_{ij}[Z_i,Z_j]
       =\sum_{i,j}(T_{L,\mu})_{ij}Z_i\times Z_j.
\]
The second equality uses the convention
\([u,v]=2u\times v\) and skewness of \(T\).
Whitening each independent colour changes this bank
to \(C_L^{1/2}T_{L,\mu}C_L^{1/2}\), exactly
the finite matrix of \(S_{L,\mu}\) by
\eqref{eq:v15-sum}. The \(c\)-component becomes
\[
 \sum_{a,b}\epsilon_{cab} I_{ab}(S_{L,\mu})
                             =2I_{ab}(S_{L,\mu})
\]
for cyclic \((c,a,b)\), since
\(I_{ba}(S)=I_{ab}(S^t)=-I_{ab}(S)\).
This proves identification with the actual
polynomial, including its factor two.
Every component already has zero expectation:
there is no missing diagonal Wick subtraction.

Apply \eqref{eq:v23-isometry} to each component.
Summing three colours gives the factor twelve
in \eqref{eq:v23-chaos-limit}.
Different output colours have zero covariance,
because one independent Gaussian colour then
occurs only once. For equal colours the isometry
and V15's Gram identity give
\eqref{eq:v23-chaos-cov}.

For the exponential statement, let
\(X=2I_{ab}(S)\), \(a\ne b\).
In singular-vector coordinates a finite-rank
approximation is a sum \(2\sum_i s_i U_iV_i\)
of independent products of standard Gaussians.
Integrating first over the \(V_i\)'s, then over
the \(U_i\)'s, gives
\begin{equation}\label{eq:v23-chaos-mgf}
 \mathbb E e^{tX}
       =\det(I-4t^2S^*S)^{-1/2},
                         \qquad 2|t|\|S\|_{\rm op}<1.
\end{equation}
The determinant is an ordinary trace-class
Fredholm determinant of \(S^*S\), not a determinant
of the non-trace-class white covariance.
For \(2|t|\|S\|_{\mathfrak S_2}\le1/2\),
\(-\log(1-u)\le u/(1-u)\) implies
\[
 \log\mathbb E e^{tX}
    \le\frac{2t^2\|S\|_{\mathfrak S_2}^2}
                   {1-4t^2\|S\|_{\rm op}^2}
    \le\frac83t^2\|S\|_{\mathfrak S_2}^2.
\]
The product converges since \(\sum_i s_i^2<\infty\);
uniform exponential bounds at a slightly larger
parameter justify passing from the finite sums
to \eqref{eq:v23-chaos-mgf}.
Also \(\mathbb E e^{t|X|}\le
\mathbb E e^{tX}+\mathbb E e^{-tX}\).

Use \(|\mathcal Y_L|\le\sum_{\mu,c}|\mathcal Y_{L,\mu}^c|\)
and Holder's inequality over its twelve components.
If \(\eta\le1/(48B_*)=1/8064\), then
\(12\eta\le1/(4B_*)\); the preceding bound controls
every factor
\(\mathbb E e^{12\eta|\mathcal Y_{L,\mu}^c|}\)
uniformly. This proves \eqref{eq:v23-chaos-exp},
without assuming independence of the twelve
components. The same argument applies to the limit.
Finally the exponential bound makes every fixed
power uniformly integrable. Combining it with
\eqref{eq:v23-chaos-limit} gives \(L^q\) convergence;
equivalently, split the \(q\)-th moment of the
difference at a fixed large threshold and use
the exponential bound on its tail.
\end{proof}

For clarity, the exact one-component characteristic
function, for a fixed orientation combination
\(S_t=\sum_\mu t_\mu S_{\infty,\mu}\), is
\[
 \mathbb E\exp\left(iu\sum_\mu t_\mu
                                  \mathcal Y_{\infty,\mu}^c\right)
       =\det(I+4u^2S_t^*S_t)^{-1/2}.
\]
Its variance is \(4\Tr(S_t^*S_t)\), and its fourth
cumulant is \(96\Tr((S_t^*S_t)^2)\).
Thus it is non-Gaussian if \(S_t\ne0\).
No strict positivity of \(D_\infty\) is asserted;
V15 proved only positive semidefiniteness.
The same Gaussian families occur in every shell,
so this construction has not replaced mixed
shell terms by independent noise.

\subsection{What this changes in the continuation}

The quadratic-root limit gives an actual stochastic
coefficient rather than just a deterministic bounded
map on smooth small fields. It is a useful starting
point for a source construction organized by Gaussian
polynomials. But for \(\beta^{-1/2}Z_L\) one may not
write, with a regulator-uniform error,
\[
 \mathcal F_L(\beta^{-1/2}Z_L)
       =\beta^{-1}\mathcal Y_L
                              +O(\beta^{-3/2})
\]
merely from the classical Taylor theorem.
The linear detail term vanishes, but the domain
needed for that classical error has the small
probability just proved. The full nonlinear root
may not even be represented by its identity
logarithms on all such samples.

The next upstream acquisition should therefore
control the higher stochastic source coefficients
and the actual measured density on a scale-local
domain, or establish a quantitative transfer
from a suitable localized law. Such a proof
must retain the complete gauge, Haar and
source-normal factors of Shell V16 and the
incoming packet of Shell V8. A Gaussian
quadratic observable does not provide that
change of law, and the entropy test does not
assume it for free.

\subsection{Exact verification and preservation}

The new checker
\begin{center}\small
\path{scripts/verify_ym_f14_v23_gaussian_domain.py}
\end{center}
has SHA--256
\begin{center}\scriptsize\ttfamily
5848489B56A6F3447848F8C6CAFC6B5C59A4EFFDD570A27956607B766C8235E6.
\end{center}
It enumerates the eight-variable pairings in
\eqref{eq:v23-Wick} with exact integer coefficients.
Using the complete side-two transverse covariance,
it checks
\[
 c_2=\frac{103}{1024},\qquad
 \Tr C_2^2=\frac{865}{768},\qquad d_2=135,
\]
the exact global fourth-norm variance, and all
sixteen quadratic-root orientation covariances.
It also checks the Chernoff scaling factors,
the nonzero fourth cumulant of a finite skew-bank
calibration, and the vector exponential radius.
There is no Monte Carlo inference. The cofinal
and all-regulator statements follow from the
analytic proofs, not finite numerical tests.

The predecessor V22 has 488,514 bytes and SHA--256
\begin{center}\scriptsize\ttfamily
9282C9D673D48B34AF40E66682F8C8D0C445F254AE06475C515CA6E4DB26E1BF.
\end{center}
Its full body is retained apart from the title
version and this terminal insertion. V23 is
the sole active main version; V22 is replaced.
The companion notes, archives and monographs
remain untouched. The next regular SM/NS
checkpoint remains V25.

\begin{firewall}
The Gaussian-domain obstruction concerns a proposed
use of a classical small-field theorem; it is not
a nonexistence theorem for Yang--Mills theory.
The positive limit is one quadratic polynomial of
the balanced source under the free reference law,
not the complete root pushforward under the
interacting compact measure. A cofinal interacting
state, noncollapse, reflection-positive physical
transfer and a positive Hamiltonian gap remain
unproved. The full objective is unchanged.
\end{firewall}
\section{V24: a centered source Hessian and the logarithm-subtracted mixed cubic observable}
\label{sec:v24-mixed}

\subsection{The next stochastic coefficient, with its background fixed honestly}

V23 constructs the quadratic balanced-root polynomial
under the free reference Gaussian. Its domain obstruction
prevents upgrading that polynomial to a full nonlinear
random source by the classical Taylor remainder alone.
We now extend the stochastic construction to the source
Hessian along small homogeneous backgrounds. Its derivative
at zero gives the cubic coefficient with one deterministic
source leg and two Gaussian legs. After its actual
contraction is subtracted, this coefficient converges
in every finite Gaussian \(L^q\).

The result is stronger than the zero-background
Hilbert--Schmidt theorem of V15, but has a specific
restriction: the background is homogeneous.
Write \(c\in\mathbb C^{12}\) for its physical amplitude,
so the bare fine field is \(c/L\). In general
\(\mathcal F_L(c/L)\ne c\). Thus \(c\) is an auxiliary
background parameter, not a replacement for the
original retained root source \(v\).
No source Jacobian or density is changed in this section.
At the linearized source level at zero these parameters
agree, which is all that the final mixed-cubic statement uses.

The exact contraction of that mixed coefficient is
\(\Gamma_L h\), for a deterministic direction \(h\).
V22's finite-part theorem therefore turns Wick centering
into the explicit subtraction \(n\Gamma_*h\), with a
convergent finite correction. The reference matrix \(h_L\)
in the measured determinant is not part of this source
contraction and is not added to it.

The predecessor and companions are unchanged. The next
SM/NS checkpoint is V25. The discussion concerns the
literal balanced source of Shell V16, not an independently
chosen local polynomial map.

\subsection{Physical tangent maps along the homogeneous orbit}

Let \(K=2^k\) denote an intermediate side length.
Starting at \(b_L^{(L)}(c)=c/L\), apply the actual
one-step balanced map to constant link tuples:
\[
 b_{K/2}^{(L)}(c)=\mathfrak T_K(b_K^{(L)}(c)),
 \qquad A_K^{(L)}(c)=K b_K^{(L)}(c).
\]
On constant inputs the step is independent of position;
the formula still uses its actual tree and balancing
word. The label \(L\) is retained because the incoming
homogeneous tuple depends on the finer steps.

Let
\[
 \mathcal H_K^{\rm col}=\mathbb C^{12K^4},\qquad
 \|U\|_{\mathcal H_K^{\rm col}}^2
                 =K^{-4}\sum_{x,\mu}|U_{x,\mu}|^2.
\]
In physical link values the step is
\[
 \mathfrak P_K(A)=\frac K2\mathfrak T_K(A/K).
\]
At the indicated homogeneous background put
\begin{equation}\label{eq:v24-P}
 P_K^{(L)}(c)=D\mathfrak P_K(A_K^{(L)}(c))
                   =\tfrac12D\mathfrak T_K(b_K^{(L)}(c)).
\end{equation}
Transpose below refers to the normalized-volume
bilinear pairing, extended holomorphically from
real coordinates. Norm bounds use the usual complex
Hilbert norm. The corresponding transpose estimates
follow from the same absolute matrix coefficients;
no antiholomorphic dependence is inserted.

\begin{lemma}[Summable tangent and adjoint perturbations]
\label{lem:v24-tangents}
On one sufficiently small complex ball \(|c|<r\),
there is a constant \(C\), independent of \(L,K\),
such that
\begin{equation}\label{eq:v24-tangent-bound}
 \begin{aligned}
 |b_K^{(L)}(c)|&\le 2|c|/K,\\
 \|P_K^{(L)}(c)\|_{\mathcal H_K^{\rm col}\to
                                      \mathcal H_{K/2}^{\rm col}}
                                                    &\le1+C|c|/K,\\
 \|(P_K^{(L)}(c))^{\mathsf t}\|_{\ell^\infty\to\ell^\infty}
                                                    &\le1+C|c|/K .
 \end{aligned}
\end{equation}
Consequently every incoming product of the \(P_K\)'s
is bounded in the physical Hilbert norm, and every
outgoing transposed product is bounded in
\(\ell^\infty\), by a constant independent of its
number of steps.
\end{lemma}
\begin{proof}
The local analytic word has
\(\mathfrak T_{\rm hom}(b)=2b+O(|b|^2)\) on a fixed
complex neighbourhood of zero. Its quadratic
homogeneous term actually vanishes, but the weaker
bound suffices. Thus
\[
 |A_{K/2}^{(L)}|\le |A_K^{(L)}|
                         +C|A_K^{(L)}|^2/K.
\]
A bootstrap with \(|A_K|\le2|c|\) gives a total
increase at most \(4C|c|^2\sum_{K\ge2}K^{-1}\).
The sum is over dyadic \(K\), so it is at most one.
Choosing \(r\) small closes the bootstrap and
proves the first bound.

At zero, \(P_K(0)=B_2/2\) is the actual physical
linear block. Its scalar entries are nonnegative,
each row has mass one, and each fine column has
mass \(1/16\). The normalized-volume transpose is
therefore also a positive averaging map with
row mass one. Jensen's inequality gives
\(\|P_K(0)\|_{\mathcal H\to\mathcal H}\le1\);
the column identity gives
\(\|P_K(0)^{\mathsf t}\|_{\infty\to\infty}\le1\).
Colour components carry identity matrices.

The derivative of the literal analytic primitive
differs from \(B_2\) by local matrices of size
\(O(|b|)\). Their supports and row and column
incidences are bounded by fixed constants:
only the finitely many links in the one-step
tree and path word occur. Periodic coincidences
are combined within those same finite counts.
The row/column estimate for the Hilbert norm,
and the column-sum estimate for the normalized
transpose, give perturbation bounds \(C|b|\).
The factor between the two normalized volumes
is the fixed number \(16\), not a function of \(K\).
Insert \(|b_K|\le2|c|/K\) to obtain
\eqref{eq:v24-tangent-bound}.
Finally
\(\prod_K(1+C|c|/K)\le
\exp(C|c|\sum_KK^{-1})\) is uniformly bounded.
\end{proof}

This estimate uses homogeneity in a substantive
way: the smallness at every intermediate link
is \(O(|c|/K)\). A merely bounded global
\(\ell^4\) norm on an arbitrary background
does not imply that pointwise bound.

\subsection{The exact Hessian insertion sum}

Use the common full-colour energy space \(V\)
and real isometries \(J_L:V_L\to V\) from V13.
Define the embedded source
\[
 \widetilde{\mathcal O}_L(c,u)
                 =\mathcal F_L(c/L+J_L^*u).
\]
For each of its twelve real output coordinates
\(\alpha\), let \(\mathbb B_{L,\alpha}(c)\) be
the complex symmetric Riesz operator defined by
\begin{equation}\label{eq:v24-B}
 \langle u,\mathbb B_{L,\alpha}(c)z\rangle_V
       =\tfrac12 D_u^2\widetilde{\mathcal O}_L^\alpha(c,0)[u,z].
\end{equation}
For real \(c\) it is self-adjoint. It is finite rank,
with zero extension to \(V\) understood. Its finite
matrix is the half-Hessian of the actual source
whitened by the zero-source free covariance.
This is not the interacting constrained precision.

For a shell \(K\), write
\[
 X_{L,K}(c)
    =P_{2K}^{(L)}(c)\cdots P_L^{(L)}(c)\,
                               \mathcal E_L J_L^*,
 \qquad \mathcal E_Lw=Lw .
\]
The product is empty when \(K=L\).
Let \(w_{L,K/2,\alpha}(c)\) be the normalized-volume
gradient at side \(K/2\) of output coordinate
\(\alpha\), propagated backwards through all
remaining coarser physical steps. The terminal
gradient at side one is a unit coordinate vector.
Lemma~\ref{lem:v24-tangents} gives
\(\|w_{L,K/2,\alpha}\|_\infty\le C\).

Define the finite physical bilinear-form operator
\(\mathbb R_{L,K,\alpha}(c)\) by
\begin{equation}\label{eq:v24-R}
 \langle U,\mathbb R_{L,K,\alpha}(c)W\rangle_{\mathcal H_K^{\rm col}}
   =\frac1{4K}
      \left\langle w_{L,K/2,\alpha}(c),
            D^2\mathfrak T_K(b_K^{(L)}(c))[U,W]
                          \right\rangle_{\mathcal H_{K/2}^{\rm col}}.
\end{equation}
The factor \(1/(4K)\) includes both
\(D^2\mathfrak P_K=(2K)^{-1}D^2\mathfrak T_K\)
and the half-Hessian convention in
\eqref{eq:v24-B}.

\begin{lemma}[Actual shell banks are summable in Hilbert--Schmidt norm]
\label{lem:v24-shell}
There is an exact finite identity
\begin{equation}\label{eq:v24-insertion}
 \mathbb B_{L,\alpha}(c)
      =\sum_{\substack{K=2^k\\2\le K\le L}}
         X_{L,K}(c)^{\mathsf t}\,
             \mathbb R_{L,K,\alpha}(c)\,X_{L,K}(c).
\end{equation}
Uniformly on a smaller common complex ball,
\begin{equation}\label{eq:v24-shell}
 \|\mathbb R_{L,K,\alpha}(c)\|_{\rm op}\le C/K,\qquad
 \|X_{L,K}^{\mathsf t}\mathbb R_{L,K,\alpha}X_{L,K}\|_{\mathfrak S_2}
                        \le C K^{-1}\sqrt{1+\log K}.
\end{equation}
The tail over \(K>2^b\) is at most
\(C2^{-b}\sqrt{b+1}\), independently of \(L,c\).
\end{lemma}
\begin{proof}
The second derivative of a composition is a sum
with one second-derivative insertion. The incoming
first derivatives occur twice, and the outgoing
first derivatives transport its output to the
chosen final coordinate. Applying this chain rule
to the physical steps gives
\eqref{eq:v24-insertion} with the definitions above.
There is no freezing of an incoming derivative
or resetting of an intermediate nonlinear field.

The second derivative of the primitive is a local
bilinear stencil with uniformly bounded coefficients
on the chosen background ball. Pairing it against
a bounded \(w\) gives
\[
 |\langle w,D^2\mathfrak T_K(b)[U,W]\rangle_{\mathcal H_{K/2}^{\rm col}}|
       \le C\|w\|_\infty
                    \|U\|_{\mathcal H_K^{\rm col}}
                    \|W\|_{\mathcal H_K^{\rm col}} .
\]
To see this without an implicit volume factor,
expand the finite local stencil, use
Cauchy--Schwarz on its pairs of shifted inputs,
and sum over the output sites. Each input has
bounded incidence. The normalized volume ratio
is again \(16\), which is absorbed in \(C\).
This proves the first estimate of
\eqref{eq:v24-shell}.

By V15 the scalar energy embedding has
\(s_j(\mathcal E_L)\le j^{-1/4}\).
Including three colours gives
\(s_j(\mathcal E_L)\le3^{1/4}j^{-1/4}\).
Lemma~\ref{lem:v24-tangents} bounds its incoming
physical multiplier. Since \(X_{L,K}\) has
output dimension \(12K^4\),
\[
               \sum_j s_j(X_{L,K})^4
                            \le C(1+\log(12K^4)).
\]
Use V15's elementary inequality
\[
 \|X^{\mathsf t}RX\|_{\mathfrak S_2}
              \le\|R\|_{\rm op}
                        \left(\sum_j s_j(X)^4\right)^{1/2}.
\]
It holds for the complex bilinear transpose as
well: transpose preserves singular values, or
one can use the same singular-vector entry
estimate as in Lemma~\ref{lem:v15-HS}.
This gives the second estimate. The dyadic
tail follows by summing \(2^{-k}\sqrt{k+1}\).
\end{proof}

\subsection{Identification of the Hilbert--Schmidt limit}

For completeness, the fixed-shell convergence
required in the insertion sum uses the actual
partial nonlinear block, not just a linear
averaging map. For fixed \(K\), put
\[
 \mathcal A_{L,K}(c,u)
        =K\mathcal F_{L\to K}(c/L+J_L^*u).
\]
It converges uniformly on a common smaller
complex energy ball, with every fixed derivative.
Here is the reduction to the source comparison
already proved in V13. Stop the fine part at
a dyadic \(R\ge K\), replace that fine part by
the actual linear \(B_{L/R}\), and retain the
nonlinear \(\mathcal F_{R\to K}\).
The stopped-source estimate of Shell V13
gives error \(C_KR^{-\alpha_0}\), for one
fixed \(0<\alpha_0<1\), on a larger complex ball.
Its version with final side \(K\) is the same
theorem with \(N=K\), after a fixed coordinate
dilation; only \(K\)-dependent constants are
needed in this paragraph.

For fixed \(R,K\), V13's actual finite-average
quadrature and common-energy isometry comparison
give operator-norm convergence of the input
averages. The remaining map has fixed finite
dimension and is uniformly analytic.
First let \(L\) grow at fixed \(R,K\), and
then let \(R\) grow. This proves uniform
zeroth convergence of \(\mathcal A_{L,K}\).
Cauchy's formula on nested complex balls
gives its derivative convergence, exactly
as in Theorem~\ref{thm:v13-source}.
In particular,
\[
        A_K^{(L)}(c)=\mathcal A_{L,K}(c,0),\qquad
        X_{L,K}(c)=D_u\mathcal A_{L,K}(c,0)
\]
converge, the latter in operator norm into
the fixed space \(\mathcal H_K^{\rm col}\).

\begin{theorem}[A source Hessian limit on homogeneous backgrounds]
\label{thm:v24-HS}
On a sufficiently small common complex ball in \(c\),
there are holomorphic symmetric Hilbert--Schmidt
operators \(\mathbb B_{\infty,\alpha}(c)\) such that
\begin{equation}\label{eq:v24-HS}
        \mathbb B_{L,\alpha}(c)
                \longrightarrow\mathbb B_{\infty,\alpha}(c)
                                      \quad\hbox{in }\mathfrak S_2.
\end{equation}
The convergence is uniform on smaller balls,
with every fixed background derivative.
Their norms are uniformly bounded there.
The limit is the half-Hessian
\[
 \langle u,\mathbb B_{\infty,\alpha}(c)z\rangle_V
             =\tfrac12D^2\mathcal O^\alpha(c)[u,z]
\]
of the actual continuum source, restricted to
detail-energy directions.
\end{theorem}
\begin{proof}
Fix a shell \(K\). The preceding partial-source
comparison gives \(X_{L,K}\to X_{\infty,K}\)
in operator norm and convergence of \(b_K^{(L)}\).
The outgoing steps have fixed finite length
once \(K\) is fixed. Their derivatives, hence
the propagated gradient \(w_{L,K/2,\alpha}\)
and the finite matrix \(\mathbb R_{L,K,\alpha}\),
converge uniformly on smaller complex balls.
Thus the \(K\)-th term of
\eqref{eq:v24-insertion} converges in operator norm.
The difference has rank at most \(24K^4\),
so it converges in Hilbert--Schmidt norm.

Now truncate the insertion sum at \(2^b\).
Its finite head converges as just proved, and
both finite and limiting tails have the
uniform bound of Lemma~\ref{lem:v24-shell}.
Let \(L\) grow and then \(b\) grow.
This proves \eqref{eq:v24-HS} and the uniform bound.
The same argument is uniform on a slightly
larger complex ball. Banach-valued Cauchy
estimates prove convergence of all fixed
background derivatives and holomorphy of
the limit.
Finally V13 already identifies the full
source derivatives in multilinear operator
norm. Uniqueness of that weaker limit
identifies the Hilbert--Schmidt limit with
the displayed actual continuum half-Hessian.
\end{proof}

No rate in the original mesh \(L\) is inferred
from the finite-head argument. The stopped-shell
tail is quantitative, as in V15.
Also, \eqref{eq:v24-HS} is not a Hilbert--Schmidt
theorem for the full constrained relative
precision at the generally inhomogeneous
minimum \(a_L^*(v)\). Its background and its
operator are both different.

\subsection{Wick centering at finite amplitude}

Let \(W\) be one real isonormal family over
the full-colour energy space \(V\).
For a real symmetric Hilbert--Schmidt \(B\),
define
\[
 I_2(B)=\lim_{N\to\infty}
       \sum_{i,j\le N}B_{ij}\bigl(W(e_i)W(e_j)-\delta_{ij}\bigr)
                                      \quad\hbox{in }L^2.
\]
The Gaussian pairing identity gives
\begin{equation}\label{eq:v24-Wick-isometry}
                  \mathbb E I_2(B)I_2(D)=2\Tr(BD).
\end{equation}
For complex symmetric coefficients the analogous
sesquilinear identity gives
\(\mathbb E|I_2(B)|^2=2\|B\|_{\mathfrak S_2}^2\).
Finite-rank approximation proves basis independence.
This is a second-chaos functional, not a
Gaussian vector in the energy Hilbert space.

Define the twelve-component random maps
\[
 \mathcal W_L^\alpha(c)=I_2(\mathbb B_{L,\alpha}(c)),\qquad
 \mathcal W_\infty^\alpha(c)=I_2(\mathbb B_{\infty,\alpha}(c)).
\]
At a finite regulator their law is exactly
\begin{equation}\label{eq:v24-centered}
 \mathcal W_L^\alpha(c)
   =\tfrac12D^2\mathcal F_L^\alpha(c/L)[Z_L,Z_L]
                              -\Tr\mathbb B_{L,\alpha}(c),
\end{equation}
with the common-coupling interpretation on the
left and equality in law on the right.

\begin{theorem}[Centered stochastic Hessian and exponential control]
\label{thm:v24-Wick}
The maps \(\mathcal W_L(c)\) converge to
\(\mathcal W_\infty(c)\) locally uniformly in
Gaussian \(L^2\), with every fixed background
derivative. On real smaller background balls,
each fixed derivative family has a common
positive exponential radius. The same convergence
therefore holds in every finite \(L^q\).
\end{theorem}
\begin{proof}
Apply \eqref{eq:v24-Wick-isometry} to
Theorem~\ref{thm:v24-HS} and its derivatives.
For exponential control it suffices to consider
one real symmetric \(B\), with
\(\|B\|_{\mathfrak S_2}\le M\).
In eigenvector coordinates,
\[
                         I_2(B)=\sum_i b_i(\xi_i^2-1).
\]
For \(2|t|\|B\|_{\rm op}<1\), independent
one-dimensional Gaussian integration gives
\begin{equation}\label{eq:v24-det2}
 \mathbb E e^{t I_2(B)}
       =\det\nolimits_2(I-2tB)^{-1/2},\qquad
 \log\mathbb E e^{t I_2(B)}
       =\frac12\sum_{r\ge2}\frac{(2t)^r}{r}\Tr(B^r).
\end{equation}
The right side is well defined for Hilbert--Schmidt
\(B\). Indeed
\[
 \left|\log\mathbb E e^{tI_2(B)}\right|
        \le\frac{t^2\|B\|_{\mathfrak S_2}^2}
                         {1-2|t|\|B\|_{\rm op}}.
\]
This follows by using \(1/r\le1/2\) and
\(|\Tr B^r|\le\|B\|_{\mathfrak S_2}^2
\|B\|_{\rm op}^{r-2}\).
Finite-rank limits, with the bound at a slightly
larger parameter, justify \eqref{eq:v24-det2}
for infinite-rank \(B\).
Use \(e^{t|X|}\le e^{tX}+e^{-tX}\), then
Holder over the finitely many output components.
The uniform Hilbert--Schmidt bounds, including
the Cauchy derivative bounds, give the claimed
common radii. Uniform exponential integrability
and \(L^2\) convergence imply every finite
\(L^q\) convergence by truncation.
\end{proof}

The ordinary finite trace in \eqref{eq:v24-centered}
has not been proved to converge at general \(c\).
The infinite Wick functional needs no ordinary
trace of \(\mathbb B_\infty(c)\).
The determinant in \eqref{eq:v24-det2} is a
generating function for this centered quadratic
Gaussian observable. It is not a replacement
for the complete measured determinant of
Shell V16, nor a claim that its full relative
precision belongs to \(\mathfrak S_2\).

\subsection{The mixed cubic coefficient with its explicit logarithmic subtraction}

For a real source direction \(h\in\mathbb R^{12}\),
define the actual mixed polynomial
\begin{equation}\label{eq:v24-mixed}
 \mathcal M_L(h;Z_L)
     =\{\,\mathcal F_L(t h/L+\epsilon Z_L)\,\}_{t\epsilon^2}
     =\tfrac12D^3\mathcal F_L(0)[h/L,Z_L,Z_L].
\end{equation}
This definition means extraction of a Taylor polynomial
at zero, not evaluation of the complete root on an
unbounded Gaussian sample. Global colour covariance
makes its expectation the matrix already used in V16:
\[
                          \mathbb E\mathcal M_L(h;Z_L)=\Gamma_Lh.
\]

\begin{theorem}[A cofinal mixed cubic random coefficient]
\label{thm:v24-mixed}
On the common Gaussian coupling,
\begin{equation}\label{eq:v24-mixed-centered}
 \mathcal M_L(h;Z_L)-\Gamma_Lh
      \longrightarrow
       \mathcal M_\infty^\circ(h)
          :=D_c\mathcal W_\infty(0)[h]
\end{equation}
in every finite \(L^q\), uniformly for bounded \(h\).
For \(L=2^n\), V22's canonical matrix finite part gives
\begin{equation}\label{eq:v24-mixed-ren}
 \boxed{\quad
 \mathcal M_{2^n}(h;Z_{2^n})-n\Gamma_*h
          \longrightarrow
                \mathcal M_\infty^\circ(h)+\Gamma_{\rm fin}h .
                                                                  \quad}
\end{equation}
The centered limit has covariance
\begin{equation}\label{eq:v24-mixed-cov}
 \mathbb E\bigl[
   (\mathcal M_\infty^\circ(h))^\alpha
   (\mathcal M_\infty^\circ(k))^\delta\bigr]
       =2\Tr\!\left(
          D\mathbb B_{\infty,\alpha}(0)[h]\,
          D\mathbb B_{\infty,\delta}(0)[k]\right).
\end{equation}
Its components have uniform small exponential
moments on bounded sets of directions.
\end{theorem}
\begin{proof}
Differentiate \eqref{eq:v24-centered} once at \(c=0\).
The derivative of the first term is exactly
\eqref{eq:v24-mixed}; differentiating its finite
Gaussian expectation gives \(\Gamma_Lh\).
Equivalently,
\[
 \mathcal M_L(h;Z_L)-\Gamma_Lh
                   =I_2(D\mathbb B_L(0)[h])
\]
componentwise on the chosen coupling.
Theorem~\ref{thm:v24-Wick} proves
\eqref{eq:v24-mixed-centered}, its moment statements,
and \eqref{eq:v24-mixed-cov}.
Now add the deterministic convergence
\(\Gamma_{2^n}-n\Gamma_*\to\Gamma_{\rm fin}\)
from Theorem~\ref{thm:v22-relative}.
This proves \eqref{eq:v24-mixed-ren}.
\end{proof}

Differentiating the exact insertion sum includes
the derivative of each incoming tangent product,
each outgoing product, the primitive second
derivative and the homogeneous background.
Thus it retains the two-quadratic-insertion
contributions represented in V17 by the
incoming \(D_M\) and \(Q_M\) fields.
The fact that \(\mathbb E Q_M=0\) did not permit
discarding the random \(Q_M\) contribution here.
The proof uses the entire source Hessian before
differentiating it.

At zero, \(D\mathcal F_L(0)[h/L]=h\);
therefore the single deterministic leg in
\eqref{eq:v24-mixed} has the original retained-source
linear normalization. No background reparametrization
has altered that normalization. Higher powers of
the deterministic background would require keeping
the nonlinear relation between \(c\) and the retained
source.

\subsection{An exact mixed stochastic calibration at side two}

Take \(h=(E_1,0,0,0)\) and keep the \(E_1\)
component of the four output orientations.
For the actual side-two Gaussian, the mean is
\[
       \left(-\frac{5369}{9216},\,0,\,0,\,0\right)^t.
\]
The covariance of the centered four-vector is exactly
\begin{equation}\label{eq:v24-finite-cov}
 \operatorname{diag}\left(
      \frac{29652691}{254803968},
      \frac{161}{73728},
      \frac{161}{73728},
      \frac{161}{73728}\right).
\end{equation}
In particular the transverse output fluctuations
are not zero just because their means vanish.

Here is a finite certificate specifying the computation.
Use independent scalar Gaussian fields \(X^2,X^3\),
each with the complete \(64\)-link covariance \(C_2\).
With the deterministic leg in colour \(E_1\),
the \(E_1\) output of the mixed cubic is
\[
             (X^2)^t H_\mu X^2+(X^3)^t H_\mu X^3.
\]
These are the two perpendicular colours; the
parallel colour cancels in this output. The
real symmetric matrices \(H_\mu\) are obtained
by the literal third-degree BCH word, not
by fitting the covariance.

For an independently reproducible sparse recursion,
represent a Lie jet as
\[
             t c E_1+\epsilon r E_2
                      +t\epsilon d E_3+t\epsilon^2 e E_1,
\]
where \(r,d\) are linear polynomials in the \(64\)
scalar input coordinates and \(e\) is quadratic.
Under BCH through the indicated degrees,
\[
 \begin{aligned}
 c'&=c_a+c_b,\qquad r'=r_a+r_b,\\
 d'&=d_a+d_b+c_a r_b-c_b r_a,\\
 e'&=e_a+e_b+r_a d_b-r_b d_a\\
    &\hspace{7mm}
       +\tfrac13\{(c_a+c_b)r_a r_b-c_b r_a^2-c_a r_b^2\}.
 \end{aligned}
\]
This follows from \([u,v]=2u\times v\) and the
two degree-three BCH commutators.
Initialize \(c=1/2\) on orientation one and
zero on the other orientations, \(r\) as the
coordinate of that fine link, and \(d=e=0\).
Apply the ordered trees, conjugated two-link
paths, distinguished-root average and balancing
conjugation in their actual order.
The resulting quadratic polynomials are
\(x^tH_\mu x\).

The exact trace identities are
\[
 \mathbb E\mathcal M_{2,\mu}^{E_1}=2\Tr(C_2H_\mu),
 \qquad
 \operatorname{Cov}(\mathcal M_{2,\mu}^{E_1},
                    \mathcal M_{2,\nu}^{E_1})
                         =4\Tr(C_2H_\mu C_2H_\nu).
\]
The checker exhausts these finite matrix sums
over all \(64\) scalar links and obtains
\eqref{eq:v24-finite-cov}.
As separate calibrations, conjugating
\(tE_1\) by \(\exp(\epsilon xE_2)\) gives
\(-2t\epsilon^2x^2E_1\), and a constant
orientation-two \(E_2\) input gives the
homogeneous mixed coefficient \(-E_1/2\)
in output orientation one. The latter
agrees with the primitive homogeneous
cubic \(-\nabla E_4/4\), not the inverse
source cubic of V20.

\subsection{Scope, next acquisition, and version record}

The new result is the centered source Hessian
along a homogeneous background and, after
V22's deterministic matching, the mixed
cubic random coefficient with one source leg.
It does not construct the pure three-Gaussian-leg
coefficient \(\mathcal F_{L,[3]}(Z_L)\), control
all Gaussian chaos orders, or give a summable
remainder for the complete nonlinear root.
The trace at finite nonzero \(c\) also remains
an additional renormalization question; only
its first derivative has been matched here.

This distinguishes two useful next steps:
control the pure cubic random coefficient
including its first-chaos contractions,
or extend the stochastic estimates to an
actual localized measured law. Neither is
replaced by another computation of the
already identified quartic one-loop logarithm.
The free covariance used here is fixed;
the background-dependent constrained
precision, gauge quotient, Haar and
source-normal density still belong to
the quantum measure and must be retained
when that law is constructed.

The new checker
\begin{center}\small
\path{scripts/verify_ym_f14_v24_mixed_chaos.py}
\end{center}
has SHA--256
\begin{center}\scriptsize\ttfamily
93714AAB2234877209A2C518FE1DF2EE4260AE850A2B92F20215A1A290C75780.
\end{center}
It uses exact rational arithmetic for the
complete side-two mixed banks, their means
and their centered covariances. It also checks
the physical linear block's row mass and
all \(1024\) fine column masses, which fix
the normalized adjoint used above.
These are finite normalizations, not a
substitute for the analytic insertion-sum
and convergence proofs.

The predecessor V23 has 510,452 bytes and SHA--256
\begin{center}\scriptsize\ttfamily
CB7E7848005E364E1C9155F663B23C371A2E439BE88ED8CD8825496B8B41AB5D.
\end{center}
Its entire body is retained apart from the
title version and this terminal insertion.
V24 replaces V23 as the sole active main
note. All companions and archives remain
untouched. V25 is the next scheduled
cross-program checkpoint.

\begin{firewall}
A centered Gaussian source Hessian and one
renormalized mixed cubic coefficient are not
an interacting continuum Yang--Mills state.
The construction neither changes the canonical
measured density nor proves that its quantum
fluctuations lie on the classical global
small-field branch. Noncollapse, physical
reflection-positive transfer and a positive
Hamiltonian mass gap remain unproved.
\end{firewall}
\section{V25: the full pure cubic Gaussian source and the companion checkpoint}
\label{sec:v25-cubic}

\subsection{Checkpoint, nonduplication, and the stochastic target}

The scheduled cross-program review was performed before
this calculation. The latest active companions are still
SM V72 and NS V26. Their terminal arguments were reread,
not merely their filenames.

SM V72 has 607,372 bytes and SHA--256
\begin{center}\scriptsize\ttfamily
F44F9745D43379E1672C5550411B58E97195A4C9278592D221DF5D3F644DC1F5.
\end{center}
Its fixed-time many-copy covariance transport holds
the physical momentum set fixed. The proof controls
empirical moments, conserved excitation support and
transport on a compact classical set. It does not
give a growing-cutoff estimate for the present
Gaussian source. Its warning about overlapping
matrix kernels remains relevant: contractions
must be assembled before a scalar comparison.

NS V26 has 278,283 bytes and SHA--256
\begin{center}\scriptsize\ttfamily
82C887F8C2C69E4F28FAD7C3DC8DE5B9099CB5A4BF431EBA1FFF3E91935978AD.
\end{center}
Its terminal tensor contractions concern pointwise
rank-one norms of derivatives of the Oseen kernel.
They supply no YM covariance estimate. Some global
numerical comparisons cite a future V27 grid, not
a certificate present in this active V26.
No such numerical constant is imported.
The next regular companion checkpoint is V30.
Shell-Mixing V16 is also unchanged.

V24 controls a source Hessian along homogeneous
backgrounds and its derivative with one deterministic
leg. The target here is different:
\[
                          \mathcal C_L=\mathcal F_{L,[3]}(Z_L),
\]
the pure three-Gaussian-leg coefficient of the actual
balanced root. It has a third-chaos part and a
first-chaos part produced by internal contractions.
Both are retained. The proof supplies a local
quadratic-kernel estimate and an actual parity-average
bound for the latter, rather than treating the
whole cubic as Wick ordered from the start.

All fields below are Taylor polynomials at zero
under the free Gaussian of V23. No complete root
map is evaluated on an unbounded Gaussian sample.
The resulting limit is not yet a pushforward of
the interacting measured compact law.

\subsection{A uniform pointwise bound for the same transverse covariance}

Let \(\operatorname{dist}_L\) denote periodic Euclidean
distance on the fine lattice of side \(L\).
The scalar covariance \(C_L\) is the full
zero-mean transverse covariance, with symbol
\(\Pi(\theta)/\lambda(\theta)\), where
\[
 q_\mu(\theta)=e^{i\theta_\mu}-1,\qquad
 \lambda(\theta)=\sum_\mu|q_\mu(\theta)|^2,\qquad
 \Pi(\theta)=I-qq^*/\lambda.
\]
The zero frequency is omitted; all other modes,
including self-conjugate and boundary modes, are kept.

\begin{lemma}[Periodic covariance decay]
\label{lem:v25-Green}
There is a constant \(C\), independent of \(L\), such that
\begin{equation}\label{eq:v25-Green}
 |(C_L)_{(x,\mu),(y,\nu)}|
          \le C(1+\operatorname{dist}_L(x,y))^{-2}.
\end{equation}
In magnified coordinates \(L=pM\), the cell-constant
Gaussian field \(W_{p,M}(y)=pZ_L(x)\) on the cell
at \(x/p\) therefore has covariance kernel bounded by
\begin{equation}\label{eq:v25-magnified-Green}
 |G_{p,M}(y,z)|\le
              C(\operatorname{dist}_M(y,z)+p^{-1})^{-2}.
\end{equation}
The bounds hold componentwise in orientation;
Lie colours are independent.
\end{lemma}
\begin{proof}
Use a smooth periodic dyadic partition of the
Fourier torus, with frequency scales
\(\rho\) between \(c/L\) and a fixed constant.
The smallest cutoff vanishes near zero and
the cutoffs sum to one on every nonzero
frequency of the finite lattice.
Near zero, the matrix multiplier
\[
                    \lambda^{-1}I-\lambda^{-2}qq^*
\]
has derivatives of order \(a\) bounded by
\(C_a|\theta|^{-2-a}\). This follows by
differentiating \(q=O(|\theta|)\),
\(\lambda\asymp|\theta|^2\), and their smooth
derivatives. Away from zero it is smooth
and periodic. Thus a band multiplier has
derivatives bounded by \(C_a\rho^{-2-a}\).

Its finite Fourier kernel has the trivial
bound \(C\rho^2\), since its support contains
at most \(C(L\rho)^4\) frequencies.
Discrete summation by parts in a frequency
coordinate \(k_j\), repeated \(N\) times,
multiplies the kernel by
\((e^{2\pi i(x_j-y_j)/L}-1)^N\).
Each finite difference of the multiplier
costs at most \(C_NL^{-N}\rho^{-2-N}\).
The expanded support still contains at
most \(C_N(L\rho)^4\) modes, because
\(\rho\ge c/L\). Choose a coordinate whose
centred spatial distance is at least half
the total distance. The sine chord bound
then gives
\[
 |C_{L,\rho}(x-y)|
       \le C_N\rho^2
                     (1+\rho\operatorname{dist}_L(x,y))^{-N}.
\]
All summations by parts are periodic, so
no Brillouin-face term is deleted.
For \(N>2\), summing the dyadic bands gives
\eqref{eq:v25-Green}: the bands below inverse
distance sum as \(\rho^2\), and those above
it have a convergent decaying geometric tail.
Multiply by \(p^2\) and rescale distance.
Moving from cell corners to points within
cells changes the regularized distance
by at most a fixed multiple of \(p^{-1}\).
This proves \eqref{eq:v25-magnified-Green}.
\end{proof}

On a fixed unfolded box the potential on the
right of \eqref{eq:v25-magnified-Green} has
integral at most \(CR^2\) over a ball of
radius \(R\ge p^{-1}\), uniformly in \(p,M\).
Periodic overlaps can be handled by a
bounded number of translated boxes because
\(M\ge2\). We use these elementary local
integral bounds below.

\subsection{The incoming quadratic field has a local summable kernel}

Write the actual partial source expansion as
\[
 \mathcal F_{L\to M}(\epsilon Z_L)
       =\epsilon Y_{L,M}+\epsilon^2Q_{L,M}
                                    +\epsilon^3R_{L,M}+\cdots .
\]
Here \(Y_{L,M}=B_{L/M}Z_L\), and the quadratic
chain rule gives exactly
\begin{equation}\label{eq:v25-Qchain}
 Q_{L,M}
    =\sum_{\substack{s=2^\ell\\2\le s\le p}}
           B_{s/2}E_{sM}^{[2]}(B_{p/s}Z_L),
                                      \qquad p=L/M.
\end{equation}
The sum is empty when \(p=1\).
Every term has zero expectation because
the primitive quadratic term is a sum of
Lie brackets with colour-isotropic Gaussian
inputs. Thus it belongs to the second
Gaussian chaos, without a missing constant
counterterm.

Fix one output link and unfold its local
carrier in the magnified torus of side \(M\).
Denote the term indexed by \(s\) by \(Q_{p,s}\).
Represent it componentwise as
\[
 Q_{p,s}=\iint K_{p,s}(x,y)
                           W_{p,M}(x)W_{p,M}(y)\,dx\,dy .
\]
The orientation and colour indices are contracted
by \(K_{p,s}\); absolute kernel norms below sum
over their fixed finite number. Integration
means the exact cell-constant finite sum.

\begin{lemma}[The actual quadratic kernel envelope]
\label{lem:v25-envelope}
For a fixed numerical support constant \(A\),
independent of \(p,s,M\), the kernels satisfy
\begin{equation}\label{eq:v25-envelope}
 \begin{aligned}
 &x,y\text{ lie in a fixed box about the output link},
                   \qquad |x-y|\le A/s,\\
 &|K_{p,s}(x,y)|\le Cs^3,\\
 &\sup_x\int|K_{p,s}(x,y)|\,dy+
       \sup_y\int|K_{p,s}(x,y)|\,dx\le C/s,\qquad
                         \|K_{p,s}\|_1\le C/s .
 \end{aligned}
\end{equation}
Consequently the full kernel in \eqref{eq:v25-Qchain}
has uniformly bounded absolute marginals and
\(L^1\) norm on a fixed carrier.
\end{lemma}
\begin{proof}
Put \(q=s/2\), \(t=p/s\).
The primitive quadratic stencil has finitely
many offset pairs and bounded coefficient sum,
independent of the surrounding torus.
For each such pair, the outer row of \(B_q\)
has maximum coefficient \(q^{-3}\), and
each inner \(B_t\) row has maximum \(t^{-3}\)
and mass \(t\). For a fixed fine link,
at most two possible origins of an inner
row contain it, as in V17's sandwich proof.
The local offset count remains fixed.

Because \(W_{p,M}=pZ_L\) and each cell has
volume \(p^{-4}\), the kernel value is
\(p^6\) times the corresponding fine quadratic
coefficient. Hence its maximum is at most
\[
                        Cp^6q^{-3}t^{-6}=C' s^3.
\]
Summing the second inner row before taking
the maximum of the first gives an absolute
kernel marginal bounded by
\[
                        Cp^2q^{-3}t^{-2}=C'/s.
\]
The other marginal is the same calculation.
The total coefficient mass gives
\(Cq t^2/p^2=C/(2s)\).

The outer row occupies a box of diameter
\(O(q)\) at the post-step scale; after
decimation and the inner averages the
unfolded fine support has diameter \(O(qt)\).
Since \(qt=p/2\), this becomes a fixed
box after division by \(p\).
The two inputs of a primitive quadratic
stencil differ by \(O(t)\) fine sites,
including both inner averaging supports.
Their separation is therefore \(O(1/s)\).
This proves \eqref{eq:v25-envelope}.
Unfolded rows are counted before any
periodic identification; folding only
requires finitely many copies of this
same fixed carrier for \(M\ge2\).
\end{proof}

\begin{lemma}[Uniform local quadratic moments and a finer-scale tail]
\label{lem:v25-Qmom}
For each link,
\begin{equation}\label{eq:v25-Qmom}
 \|Q_{p,s}\|_{L^2(\mathbb P)}
                  \le C s^{-1}\sqrt{1+\log s},\qquad
 \|Q_{L,M}\|_{L^2(\mathbb P)}\le C.
\end{equation}
The sum over \(s>2^b\) has \(L^2\) norm at most
\(C2^{-b}\sqrt{b+1}\), uniformly in \(L,M\).
These bounds have corresponding finite-\(L^q\)
versions with constants depending only on \(q\).
For each fixed \(M\), the entire finite-carrier
field \(Q_{L,M}\) converges on the common Gaussian
coupling as \(L\to\infty\), in every finite \(L^q\).
\end{lemma}
\begin{proof}
Use the two cross-pairings in the covariance
of a centered quadratic Gaussian polynomial,
and bound their absolute kernels by
\eqref{eq:v25-magnified-Green}.
It suffices to estimate
\[
 \iiiint |K(x,y)K(x',y')|
           g_p(x-x')g_p(y-y')\,dx\,dy\,dx'\,dy',
 \qquad g_p(z)=C(|z|+p^{-1})^{-2}
\]
on fixed boxes, with the bounded periodic
image multiplicity understood.
For \(|x-y|\le A/s\), the kernel bound in
Lemma~\ref{lem:v25-envelope} gives
\[
 \iint |K(x',y')|g_p(x-x')g_p(y-y')\,dx'\,dy'
                                  \le C s^{-1}(1+\log s).
\]
Here are the two parts of that estimate.
For \(|x-x'|\le A'/s\), integrate
\(g_p(y-y')\) over a ball of radius \(O(1/s)\),
giving \(Cs^{-2}\), then integrate
\(g_p(x-x')\) over the near ball, giving
another \(Cs^{-2}\). Multiplication by
the kernel maximum \(Cs^3\) gives \(C/s\).
For \(|x-x'|>A'/s\), with \(A'\) fixed
sufficiently large, both distances are
comparable to \(|x-x'|\).
The \(y'\)-volume is \(Cs^{-4}\), so this
part is bounded by
\[
          Cs^{-1}\int_{A'/s<|z|<C'}|z|^{-4}\,dz
                               \le Cs^{-1}(1+\log s).
\]
The regularization \(p^{-1}\le s^{-1}\)
makes the near bounds uniform. A finite
sum over periodic image singularities
satisfies the same estimate.
Finally multiply by \(\|K\|_1\le C/s\).
This proves the squared bound in
\eqref{eq:v25-Qmom}.

Sum its square root over dyadic \(s\).
The series and its tail are the same
geometric-logarithmic sums used in V15.
For fixed degree-two Gaussian polynomials,
higher moments are controlled by their
second moments, uniformly in the dimension;
this also follows from the centered
quadratic determinant estimate in V24.
Thus the tail and uniform bound hold in
every fixed \(L^q\).

For fixed \(M,s\), \(Q_{p,s}\) is a fixed
finite polynomial of the linear averages
\(B_{L/(sM)}Z_L\).
Those finite rows converge in the common
energy coupling by V13's finite-average
comparison. All their fixed Gaussian
moments converge. Hence this fixed term
converges in \(L^q\).
First truncate the sum in \(s\), pass to
the limit in its finite head, and then
use the uniform tail. This proves the
last assertion, retaining correlations
between all incoming levels.
\end{proof}

\subsection{Exact cubic shells and the two chaos projections}

The degree-three chain rule yields
\begin{equation}\label{eq:v25-cubic-chain}
 \begin{aligned}
 \mathcal C_L&=\sum_{\substack{M=2^m\\2\le M\le L}}\mathcal S_{L,M},\\
 \mathcal S_{L,M}
      &=\frac M2\operatorname{avg}_{r\in(\mathbb Z/(M/2))^4}
                                         P_{L,M,r},\\
 P_{L,M,r}
      &=D^2E_{M,r}^{[2]}(0)[Y_{L,M},Q_{L,M}]
                                      +E_{M,r}^{[3]}(Y_{L,M}).
 \end{aligned}
\end{equation}
The output orientation and colour are suppressed.
The factor \(M/2\) is the exact transport of
the spatial mean by the remaining linear
blocks to side one.
In particular, the term with the incoming
quadratic field is part of the full cubic
coefficient. It is not removed by its
zero expectation.

Let \(\mathsf P_1,\mathsf P_3\) be the
orthogonal projections onto first and
third homogeneous Gaussian chaoses on
the common isonormal family. An odd cubic
polynomial has precisely these two parts.

\begin{lemma}[Third-chaos covariance of a local cubic shell]
\label{lem:v25-third}
There is a constant independent of \(L,M\) such that
\[
 \left|\mathbb E[
    (\mathsf P_3P_{L,M,r})(\mathsf P_3P_{L,M,r'})]\right|
       \le C(1+\operatorname{dist}_M(2r,2r'))^{-6}.
\]
Consequently
\begin{equation}\label{eq:v25-third}
                    \|\mathsf P_3\mathcal S_{L,M}\|_2\le C/M.
\end{equation}
\end{lemma}
\begin{proof}
Each \(Y_{L,M}\) row is the bounded unit-scale
tent/box kernel of V18, with \(L^1\) norm one.
Each \(Q_{L,M}\) is a quadratic kernel on a
fixed nearby carrier with uniformly bounded
absolute \(L^1\) norm by
Lemma~\ref{lem:v25-envelope}.
Thus the third-chaos part of \(P_{L,M,r}\)
has a trilinear kernel of bounded \(L^1\)
norm whose three inputs are all in a
fixed box about \(2r\).
This holds for the primitive cubic term
and for the \(YQ\) term together.

If the two boxes are separated by a
sufficiently large distance \(d\), the
third-chaos covariance consists of the
six cross-pairings. All three covariances
cross between the two boxes. Lemma~\ref{lem:v25-Green}
bounds their product by \(Cd^{-6}\);
the absolute kernel integrals are bounded.
For boxes within bounded distance, use
Cauchy--Schwarz. Local \(L^2\) bounds for
the primitive cubic follow from bounded
Gaussian variances. For the mixed term
use \(\|YQ\|_2\le\|Y\|_4\|Q\|_4\) and
Lemma~\ref{lem:v25-Qmom}.
Projection onto third chaos cannot increase
the \(L^2\) norm. This proves the covariance
bound, also on small tori where only
boundedly many sites occur.

In four dimensions
\(\sum_{r'}(1+\operatorname{dist}_M(2r,2r'))^{-6}\le C\).
There are \((M/2)^4\) output cells.
The variance of their normalized average
is therefore at most \(CM^{-4}\).
Multiply by \((M/2)^2\) from
\eqref{eq:v25-cubic-chain}.
\end{proof}

The first-chaos part needs a different argument:
its covariance need not have the integrable
power six. It is controlled by the actual
aligned output average and the missing zero
mode, not by imposing independence between cells.

\begin{lemma}[Bounded contracted filters and the actual parity selection]
\label{lem:v25-first}
Each first-chaos local field
\(\mathsf P_1P_{L,M,r}\) is a linear smearing
of \(W_{p,M}\) with a kernel
\(h_{p,M}(y-2r)\) supported in a fixed box.
Its \(L^1\) and \(L^\infty\) norms are bounded
uniformly in \(p,M\). Moreover
\begin{equation}\label{eq:v25-first}
                    \|\mathsf P_1\mathcal S_{L,M}\|_2\le C/M.
\end{equation}
\end{lemma}
\begin{proof}
Contracting two of the three bounded linear
kernels in the primitive cubic leaves a
bounded multiple of the third kernel.
For the mixed \(YQ\) term, \(Q\) is already
a centered second-chaos polynomial.
The only first-chaos terms contract \(Y\)
with one of its two legs.
For the bounded unit-scale kernel \(f\)
of \(Y\), the smoothed potential
\[
             z\longmapsto\int f(y)G_{p,M}(y,z)\,dy
\]
is bounded uniformly, by
\eqref{eq:v25-magnified-Green} and local
integrability of the regularized inverse square.
The absolute marginal bounds for the
full \(Q\) kernel then give a bounded
contracted kernel \(h_{p,M}\).
Its support and \(L^1\) norm are also bounded.
This estimate includes the sum over all
incoming finer quadratic insertions.

Coarse translation invariance of the free
covariance and the aligned block makes
the kernel at output \(r\) precisely the
translate by \(2r\), even though its
coefficients may depend on \(p,M\).
Write
\[
 H_p^d(\xi)=p^{-4}\sum_x h_{p,M}(x/p)e^{-i\xi\cdot x/p},
 \qquad \xi=2\pi k/M.
\]
The normalized average over \(r\) is zero
unless each \(k_\mu\) is a multiple of \(M/2\).
Thus its nonzero frequencies are exactly
\(\xi=\pi\ell\) in the fine Brillouin zone,
with \(\ell\in\mathbb Z^4\setminus\{0\}\).
The zero frequency has zero covariance
because the fine detail has zero mean.
Using the exact normalization of V18,
the variance of the averaged linear field
is bounded by
\[
 C M^{-4}\sum_{\substack{\ell\ne0\\\pi\ell\text{ in the fine zone}}}
                       |H_p^d(\pi\ell)|^2 .
\]
Indeed the multiplier
\(\Pi_p/\omega_p^2\) is bounded at these
nonzero frequencies by the sine chord bound.
Parseval on a torus of side two gives
\[
             2^{-4}\sum_\ell|H_p^d(\pi\ell)|^2
                =\|\operatorname{per}_2 h_{p,M}\|_2^2\le C.
\]
Periodization has bounded overlap because
the unfolded support is fixed. This is
a finite discrete Parseval identity on
the mesh \(p^{-1}\), with all modes kept.
The averaged variance is at most \(CM^{-4}\).
Multiplication by \((M/2)^2\) proves
\eqref{eq:v25-first}.
\end{proof}

This proves an actual first-chaos bound; it
does not say that the contraction vanishes.
In particular it is not valid to replace
the cubic shell by its third-chaos projection.
The external homogeneous leg of V24 is
different: it is not a mean-zero Gaussian
mode selected by this average, and its
contraction has the logarithmic matrix
\(\Gamma_L\). There is no conflict between
that logarithm and \eqref{eq:v25-first}.

\subsection{The cofinal pure cubic theorem}

For moment statements we use the standard
dimension-free Gaussian polynomial inequality
\[
       \|P\|_{L^q}\le(q-1)^{d/2}\|P\|_{L^2},
                          \qquad q\ge2,\quad \deg P\le d.
\]
This is Gaussian hypercontractivity, not a
new YM estimate. One accessible statement
is Theorem 5.5 in the authors' version of
\href{https://cseweb.ucsd.edu/~dakane/OptimalL1Regression.pdf}
{Diakonikolas--Kane--Pittas--Zarifis (2021)},
which attributes it to the standard Gaussian
hypercontractivity theorem. Only this
fixed-degree Gaussian moment consequence
is used; no result about the interacting
theory is imported.

\begin{theorem}[The full cubic source has a Gaussian coefficient limit]
\label{thm:v25-cubic}
On the common Gaussian energy coupling,
the actual pure cubic coefficients satisfy
\begin{equation}\label{eq:v25-cubic-limit}
             \mathcal F_{L,[3]}(Z_L)
                            \longrightarrow\mathcal C_\infty
                      \quad\hbox{in every finite }L^q,\ q\ge1.
\end{equation}
Their first- and third-chaos projections
converge separately and both are retained
in \(\mathcal C_\infty\).
For a sum stopped at shell side \(2^b\),
the discarded part has norm at most
\begin{equation}\label{eq:v25-cubic-tail}
                C_q2^{-b}\quad\hbox{in }L^q,
\end{equation}
uniformly in the finer lattice.
There are constants \(\eta>0,C<\infty\) such that
\begin{equation}\label{eq:v25-stretched}
 \sup_{L\in\{2,4,\ldots,\infty\}}
              \mathbb E\exp\{\eta|\mathcal C_L|^{2/3}\}\le C.
\end{equation}
No ordinary exponential moment or uniform
nonlinear Taylor remainder is asserted.
\end{theorem}
\begin{proof}
By Lemmas~\ref{lem:v25-third} and
\ref{lem:v25-first},
\[
                  \|\mathcal S_{L,M}\|_2\le C/M.
\]
The sum of these bounds over dyadic \(M\)
is finite, and its tail is \(C2^{-b}\).
The Gaussian polynomial inequality with
degree three gives the corresponding
\(L^q\) bound on each whole shell.
It also applies to each chaos projection.

Fix \(M\). The finite-carrier linear
fields \(Y_{L,M}\) converge in all Gaussian
moments on the common coupling by V13.
Lemma~\ref{lem:v25-Qmom} gives the same
conclusion for \(Q_{L,M}\), jointly with
these linear fields. Products converge
in \(L^2\) by Holder's inequality and the
uniform higher moment bounds. The finitely
many primitive polynomials and output
cells in \eqref{eq:v25-cubic-chain} therefore
give a limit \(\mathcal S_{\infty,M}\).
Since \(\mathsf P_1,\mathsf P_3\) are bounded
orthogonal projections on the same Gaussian
\(L^2\) space, they commute with this limit.
Fix \(b\), pass to the limit in the finite
head \(M\le2^b\), and then let \(b\) grow
using the uniform tail. This proves
\eqref{eq:v25-cubic-limit} and
\eqref{eq:v25-cubic-tail}, without dropping
any mixed-shell covariance.

The uniform \(L^2\) bound and the same
degree-three inequality give
\(\|\mathcal C_L\|_q\le C_0 q^{3/2}\)
for \(q\ge2\), with a harmless fixed factor
for the twelve-dimensional output.
Markov's inequality with
\(q\) proportional to \(t^{2/3}\) gives
\[
                  \mathbb P\{|\mathcal C_L|>t\}
                                  \le C_1e^{-c_1t^{2/3}}.
\]
Integrating this tail proves
\eqref{eq:v25-stretched} for sufficiently
small positive \(\eta\).
The limit satisfies the same bounds,
by moment convergence or Fatou's lemma.
\end{proof}

In particular the limiting random coefficient
can be specified without an unexplained
continuum cubic product:
\[
 \mathcal C_\infty
        =\sum_{\substack{M=2^m\\M\ge2}}\mathcal S_{\infty,M}.
\]
The series converges absolutely in each
finite Gaussian \(L^q\). The finite-shell
terms use the actual limiting averaged
linear fields and incoming quadratic
fields, on the same Gaussian family.
No distributional pointwise product
\(W(x)^3\) is being postulated.

\subsection{Exact nonzero cubic and linear-contraction tests}

Use the complete scalar side-two transverse
projection \(\Pi_2\) of V14. Number fine
scalar links by lexicographic site and then
orientation; indices \(0,1\) are the origin
links in orientations one and two.
The real mean-zero Coulomb field
\[
                u=(\Pi_2e_0)E_1+(\Pi_2e_1)E_2
\]
has the following exact pure cubic output:
\begin{equation}\label{eq:v25-cubic-test}
 \mathcal F_{2,[3]}(u)
    =\frac1{442368}
        \begin{pmatrix}
          15373&20989&0\\
          14173&3853&0\\
          0&0&0\\
          0&0&0
        \end{pmatrix}.
\end{equation}
Rows are output orientations and columns
are Lie colours. This is a live detail
direction, not a constant or longitudinal
direction removed by the covariance.

The first-chaos coefficient is nonzero too.
For the deterministic direction
\(u_0=(\Pi_2e_0)E_1\), exact Gaussian
contraction gives
\begin{equation}\label{eq:v25-first-test}
 \frac12\mathbb E D^3\mathcal F_2(0)[u_0,Z_2,Z_2]
                   =(27E_1/128,\,0,\,0,\,0).
\end{equation}
For a homogeneous cubic polynomial
this linear map is precisely the coefficient
of its first-chaos projection. Since
\(C_2\) is positive definite on the live
detail space, \eqref{eq:v25-first-test}
certifies a nonzero first-chaos component.

The finite certificate propagates all
three homogeneous degrees through the
actual ordered word. If a Lie jet is
\(a_1+a_2+a_3\), with the index denoting
its degree, then BCH gives
\[
 \begin{aligned}
 c_1={}&a_1+b_1,\\
 c_2={}&a_2+b_2+\tfrac12[a_1,b_1],\\
 c_3={}&a_3+b_3+\tfrac12([a_1,b_2]+[a_2,b_1])\\
       &+\tfrac1{12}\bigl([a_1,[a_1,b_1]]
                              +[b_1,[b_1,a_1]]\bigr).
 \end{aligned}
\]
Thus the two-quadratic-insertion terms
are present in the calculation, not
reconstructed from the primitive cubic alone.
The code checks \(\Pi_2u=u\), all four
zero means in each colour, zero linear
root output, and cubic scaling under
\(u\mapsto2u\). The Gaussian contraction
in \eqref{eq:v25-first-test} uses all
\(135\) real detail directions through
the complete covariance.

\begin{proposition}[Cubic exponential nonintegrability]
\label{prop:v25-no-exp}
For the orientation-one \(E_1\) component
\(P(Z_2)=(\mathcal F_{2,[3]}(Z_2))_1^{E_1}\),
\[
                 \mathbb E e^{tP(Z_2)}=\infty
                                      \qquad(t\ne0),
 \qquad
                 \mathbb E e^{t|P(Z_2)|}=\infty
                                      \qquad(t>0).
\]
The same failure of ordinary exponential
integrability holds for its nonzero
third-chaos projection.
\end{proposition}
\begin{proof}
Equation~\eqref{eq:v25-cubic-test} shows
that the homogeneous cubic \(P\) is
strictly positive on one live direction.
By continuity it is bounded below by
\(cR^3\) on a fixed open cone at radius
\(R\) in the finite-dimensional detail
space. The Gaussian covariance is
nondegenerate there, so its density on
that cone is bounded below by a constant
times \(e^{-CR^2}\). The radial integral
of \(e^{tcR^3-CR^2}\) diverges for \(t>0\).
For \(t<0\), use the opposite cone and
odd homogeneity. This also proves the
absolute-value assertion.
Subtracting the first-chaos projection
changes only a linear term. It does
not change the nonzero leading cubic
or the cone argument at large radius.
\end{proof}

The proposition is about this actual
finite root polynomial. It is not a
claim that the compact theory has a
divergent partition function. The full
compact action and density, or an
appropriate stable localization, cannot
be replaced by an uncut cubic exponential.
The stretched-exponential statement
\eqref{eq:v25-stretched} is the correct
kind of moment bound supplied here.

\subsection{What the cubic limit does and does not provide}

V23 and the present theorem now give a
joint cofinal limit of
\[
                  \bigl(\mathcal F_{L,[2]}(Z_L),
                         \mathcal F_{L,[3]}(Z_L)\bigr)
\]
under the fixed free reference law.
The cubic limit retains both its
contracted linear part and its
third-chaos part. It requires no
logarithmic subtraction on this
mean-zero pure-Gaussian input.
This does not delete the mixed-source
logarithm matched in V24.

There is still no estimate showing that
these two coefficients approximate the
complete nonlinear root with a uniform
error on typical quantum fluctuations.
The next useful task is a stable
scale-local remainder or the next
Gaussian coefficient with all its lower
chaos contractions, while retaining
the actual measured density.
The nonintegrability test makes clear
why simply exponentiating the newly
convergent cubic coefficient would
not accomplish that task.

\subsection{Exact verification and version record}

The new checker
\begin{center}\small
\path{scripts/verify_ym_f14_v25_pure_cubic.py}
\end{center}
has SHA--256
\begin{center}\scriptsize\ttfamily
11EE2A01CC3BA106798668B46286334881582A8C719A27A5238FBC739697CEDC.
\end{center}
It verifies the two actual side-two
calibrations above with exact rational
arithmetic. Fifteen exact scaling checks
verify the density, marginal, support
and total-mass factors in
Lemma~\ref{lem:v25-envelope}.
Thirty finite cyclic Fourier checks
verify the parity-average selector.
These finite checks supplement the
all-scale proofs; they do not infer
convergence from a finite lattice sample.

The predecessor V24 has 536,532 bytes and SHA--256
\begin{center}\scriptsize\ttfamily
360C5CED60AC46877C5368AFB79F9EEECF03087D43D3F9AB674F55A9C7F03EDF.
\end{center}
Its entire body is retained apart from
the title version and this terminal
insertion. V25 replaces V24 as the
sole active main note. SM V72, NS V26,
Shell V16, the archives and monographs
are preserved. The next regular
cross-program review is V30.

\begin{firewall}
The pure cubic result is a theorem
about an actual polynomial source
coefficient under the free Gaussian
reference, not an interacting continuum
measure or a controlled nonlinear
quantum remainder. It proves neither
noncollapse nor physical clustering,
reflection-positive transfer, or a
positive Hamiltonian mass gap.
The full Yang--Mills objective remains
unproved and unchanged.
\end{firewall}
\section{V26: the local cubic drift before the terminal average}
\label{sec:v26-local}

\subsection{Why the local problem is different}

V25 proves convergence of the full cubic coefficient
after blocking to side one. A scale-local nonlinear
remainder needs information before that last average.
The distinction is substantial: a linear contraction
can accumulate over the intervening scales while its
constant-frequency part disappears in the final
mean-zero average. We identify that contraction with
the already computed source matrix of V18 and V22,
prove a summable bound for the error of this
identification, and obtain a renormalized local
cubic coefficient.

This is not a repetition of the mixed-source theorem
of V24 or the global pure-cubic theorem of V25.
The new objects are the cubic coefficients on a
retained lattice of side \(M\ge2\). All estimates
are under the same free transverse Gaussian reference.
No interacting density is changed, and no complete
nonlinear root is evaluated on an unbounded Gaussian
sample. Shell-Mixing V16 and the measured-source
conventions remain in force. The companion checkpoint
was performed in V25; the next scheduled review is V30.

Write
\[
 L=2^n=pM,\qquad M=2^m\ge2,\qquad p=2^j,\qquad j=n-m.
\]
Let the homogeneous Taylor coefficients of the
actual partial balanced root be
\begin{equation}\label{eq:v26-partial}
 \mathcal F_{L\to M}(\epsilon Z_L)
   =\epsilon Y_{L,M}+\epsilon^2Q_{L,M}
                          +\epsilon^3R_{L,M}+O(\epsilon^4).
\end{equation}
Here the displayed equality denotes finite Taylor
jets at zero, not a probabilistic remainder estimate.
For \(p=1\), \(Y_{L,M}=Z_L\) and \(Q_{L,M}=R_{L,M}=0\).
Uniform local norms below mean the norm of a fixed
output link and its three colour components.
They do not mean the unnormalized norm of an
entire lattice with growing volume.

\subsection{The correct linear filter}

For a real orientation matrix \(A\), define the
cross-oriented averaging observable
\begin{equation}\label{eq:v26-A}
 \bigl(\mathcal A_{p,M}(A)Z_L\bigr)_{r,\mu}^{c}
  =\sum_{\nu=1}^4 A_{\mu\nu}\,p^{-4}
       \sum_{z\in\{0,\ldots,p-1\}^4}
       \sum_{t=0}^{p-1}
               Z_{pr+z+t e_\mu,\nu}^{c}.
\end{equation}
Indices are periodic on the fine lattice.
The geometry of the averaging row is \(\mu\);
its input orientation is \(\nu\).
Thus in general
\begin{equation}\label{eq:v26-notcommute}
          \mathcal A_{p,M}(A)Z_L\ne A\,Y_{L,M}.
\end{equation}
This distinction concerns the actual filter, not
an optional choice of a matrix convention.

The scalar geometric row is nonnegative, has mass
\(p\), maximum at most \(Cp^{-3}\), and at most
\(Cp^4\) entries. Its \(\ell^{4/3}\) norm is bounded
independently of \(p,M\). The dual Sobolev estimate
of V17, applied also when the input orientation is
different from the geometric one, therefore gives
\begin{equation}\label{eq:v26-A-bound}
 \|\mathcal A_{p,M}(A)Z_L\|_q\le C_q\|A\|,
                       \qquad 1\le q<\infty.
\end{equation}
For \(q\ge2\) this also follows from the Gaussian
\(L^2\) bound and Gaussian moment estimates;
smaller \(q\) follow by monotonicity.
At fixed \(M\), the unit-scale tent/box kernels of
V18 converge in the same common Gaussian realization,
including cross-oriented inputs. Consequently
\begin{equation}\label{eq:v26-A-limit}
 \mathcal A_{p,M}(A)Z_L\longrightarrow
                      \mathcal A_{\infty,M}(A)
                  \quad\hbox{in every finite }L^q.
\end{equation}
The right side denotes a linear observable of
the common continuum transverse Gaussian field.
Different matrices and different links use that
same field, not independent copies.

For an exact witness to \eqref{eq:v26-notcommute},
take \(L=4,p=2,M=2\), with the only nonzero input
orientation equal to \(2\), and
\[
       u_{x,2}=\cos(\pi x_1/2)\,E_1.
\]
This real field has zero mean and zero lattice
divergence: it is constant along its own link
orientation. At the origin its geometric-\(1\)
average of input orientation \(2\) is zero,
whereas its retained orientation-\(2\) average
is \(E_1\). Indeed the extra sum in direction
\(1\) gives \(1+0+0-1=0\), while the sum in
direction \(2\) duplicates \(1+0\).
Choosing \(A_{12}=1\), all other entries zero,
proves the assertion inside the live detail space.
This example disproves a general commutation rule;
it does not assert a particular nonzero entry
of the still unevaluated matrix \(\Gamma_*\).

At final side one there is a separate exact identity:
each geometric row has coefficient \(p^{-3}\)
at every fine site. Hence
\begin{equation}\label{eq:v26-A-one}
                 \mathcal A_{p,1}(A)Z_p=0
\end{equation}
for every mean-zero input and every \(A\).
This will explain why the local subtraction
does not appear in V25's terminal observable.

\subsection{Local shell contractions and their exact masses}

Let \(s=2^\ell\), \(2\le s\le p\), put
\(R=sM=2^k\), \(k=m+\ell\), \(t=p/s\),
and let \(B_{s/2}\) be the actual remaining linear
block from side \(R/2\) to side \(M\).
The same cubic chain rule used in V25 gives
\begin{equation}\label{eq:v26-chain}
 \begin{aligned}
 R_{L,M}&=\sum_{\ell=1}^j B_{2^{\ell-1}}P_{L,2^\ell M},\\
 P_{L,R,r}&=
 D^2E_{R,r}^{[2]}(0)[Y_{L,R},Q_{L,R}]
                               +E_{R,r}^{[3]}(Y_{L,R}).
 \end{aligned}
\end{equation}
In particular, the incoming quadratic coefficient
has not been reset. The final transport in this
formula is a local row, not the global average
in \eqref{eq:v25-cubic-chain}.

Use the canonical ambient Wick contraction to
represent the first-chaos part:
\[
 \mathsf P_1P_{L,R,r,\mu}^{c}
   = W_{t,R}^{c}\bigl(h_{t,R,\mu}(\,\cdot-2r)\bigr).
\]
There is a sum over input orientations in this
notation, and \(W_{t,R}=tZ_L\) on mesh \(t^{-1}\).
By Lemma~\ref{lem:v25-first}, each row
\(h_{t,R,\mu\nu}\) has fixed unfolded support and
uniformly bounded \(L^1,L^\infty\) norms.
The canonical row means the row obtained by
contracting the literal cubic polynomial before
quotienting by the nullspace of the covariance.
It may depend on \(t,R\). Constant rows can
represent zero on the mean-zero Gaussian space;
we do not add such rows or project them away.
The integral below is the mesh integral,
equivalently the integral of its cellwise
constant reconstruction.

\begin{lemma}[Mass of the contracted local row]
\label{lem:v26-mass}
For the source shell matrix in
\eqref{eq:v18-G-shell},
\begin{equation}\label{eq:v26-mass}
            \int h_{t,R,\mu\nu}(y)\,dy
                             =2(G_{n,k})_{\mu\nu}.
\end{equation}
The action in colour is the identity.
\end{lemma}
\begin{proof}
For any homogeneous cubic polynomial \(P\) of
an ambient Gaussian vector with covariance \(C\),
its canonical first-chaos row evaluated on an
ambient vector \(u\) is \(\mathbb E\,DP(Z)[u]\).
To see this, write a cubic monomial
\(Z_aZ_bZ_c\). Its first-chaos part is
\(C_{ab}Z_c+C_{ac}Z_b+C_{bc}Z_a\); evaluation
on \(u\) equals the expectation of its derivative.
Linearity proves the identity even when \(C\)
is singular. Thus an ambient constant test is
legitimate without declaring it a live Gaussian mode.

Shift the magnified field \(W_{t,R}\) by a constant
orientation/colour vector \(v\). In fine variables
this is the shift \(v/t=Rv/L\). The induced changes
are
\[
                 \delta Y_{L,R}=v,\qquad
                 \delta Q_{L,R}=R D_R(v),
\]
where \(D_R(v)\) is precisely the mixed incoming
coefficient of V17, not a new independently
sampled response. Differentiating \(P_{L,R,r}\)
in \eqref{eq:v26-chain} and taking expectation gives
\[
 \mathbb E\bigl[
 D^2E_{R,r}^{[2]}(0)[v,Q_{L,R}]
 +D^2E_{R,r}^{[2]}(0)[Y_{L,R},R D_R(v)]
 +DE_{R,r}^{[3]}(Y_{L,R})[v]\bigr].
\]
The first term has zero expectation since
\(\mathbb E Q_{L,R}=0\), proved in V25.
The remaining two terms are twice
\((G_{n,k}v)_\mu\) by \eqref{eq:v18-G-shell};
translation invariance makes the output-cell
average in that definition immaterial.
The ambient constant test of the row is
its integral. Colour isotropy, or the same
two-colour contraction rule, gives the stated
identity in colour.
\end{proof}

The factor two can be checked without a limiting
argument. On the primitive side-two lattice,
the exact row-mass matrix is
\begin{equation}\label{eq:v26-mass-two}
 \frac1{4608}
 \begin{pmatrix}
 -5369&0&0&0\\
 0&-4937&144&144\\
 0&0&-4265&336\\
 0&0&0&-3185
 \end{pmatrix}
                         =2\Gamma_2 .
\end{equation}
Here the constant magnified-field shift is
\(v\), not the homogeneous source input \(v/2\).

\subsection{A summable local averaging error}

We compare the contracted shell with the actual
finite-mesh filter, not with a continuum filter
multiplied by a diverging number of scales.
Define
\begin{equation}\label{eq:v26-error}
 E_{L,M,s}
    =B_{s/2}\mathsf P_1P_{L,sM}
           -\mathcal A_{p,M}(G_{n,m+\ell})Z_L,
                         \qquad s=2^\ell .
\end{equation}

\begin{lemma}[Summable filter error]
\label{lem:v26-error}
For all \(L=pM\ge M\ge2\), all displayed shells,
and every finite \(q\ge1\),
\begin{equation}\label{eq:v26-error-bound}
                 \|E_{L,M,s}\|_q\le C_q s^{-1/3}.
\end{equation}
In particular the tail over \(\ell>b\) is
at most \(C_q2^{-b/3}\), uniformly in \(L,M\).
\end{lemma}
\begin{proof}
It suffices to prove the \(L^2\) bound for a
single output component at the origin. Put
\(a=s/2\), and write \(b_a(r)\) for its
scalar \(B_a\) row on the side-\(R/2\) grid.
Its mass is \(a\), its maximum is at most
\(Ca^{-3}\), and it has a fixed-shape support
of diameter \(Ca\), allowing bounded periodic
overlap. Suppress orientation indices temporarily.
Relative to the magnified field \(W_{p,M}\),
the unfolded kernel of the first term of
\eqref{eq:v26-error} is exactly
\begin{equation}\label{eq:v26-K}
                  K_s(x)=\sum_r b_a(r)\,s^3 h_{t,R}(sx-2r).
\end{equation}
The factor \(s^3\) follows from
\(W_{t,R}=W_{p,M}/s\) and the four-dimensional
change of scale. Equivalently, the fine
coefficient \(t^{-3}h\) becomes \(p^{-3}s^3h\).
At any \(x\), only boundedly many terms have
nonzero \(h(sx-2r)\), since \(h\) has fixed
support and the centres have spacing two.
Therefore \(K_s\) has fixed support and a
uniform \(L^\infty\) bound, using \(s^3a^{-3}=8\).
Its \(L^1,L^2\) bounds are uniform as well.
All these statements remain valid on side
\(M\ge2\) after periodization with bounded overlap.

Write the normalized discrete transforms
\[
 \begin{aligned}
 F_{a,\mu}^d(\xi)&=a^{-1}\sum_r b_a(r)
                                    e^{-i\xi\cdot r/a},\\
 H_{t,R,\mu\nu}^d(\zeta)
    &=t^{-4}\sum_x h_{t,R,\mu\nu}(x/t)
                                    e^{-i\zeta\cdot x/t}.
 \end{aligned}
\]
The first transform is that of the unit-scale
averaging density with geometric orientation \(\mu\).
Changing the fine summation variable in
\eqref{eq:v26-K} gives the exact transform identity
\begin{equation}\label{eq:v26-transform}
  (K_s^d)_{\mu\nu}(\xi)
       =\frac12 F_{a,\mu}^d(\xi)
                             H_{t,R,\mu\nu}^d(\xi/s).
\end{equation}
Indeed the inner scaling gives \(s^{-1}\),
and the outer row sum has mass \(a=s/2\).
The target filter in \eqref{eq:v26-error} has
transform \(F_{p,\mu}^d(\xi)(G_{n,m+\ell})_{\mu\nu}\).
Lemma~\ref{lem:v26-mass} gives
\(H_{t,R}^d(0)=2G_{n,m+\ell}\).

Fixed support and bounded absolute mass imply
\[
       |H_{t,R}^d(\xi/s)-H_{t,R}^d(0)|
                                      \le C|\xi|/s.
\]
The cell reconstruction of each \(F_{a,\mu}^d\)
differs in \(L^1\) from the limiting unit-scale
tent/box kernel by \(C/a\), as in V18.
Replacing a cell integral by its left-endpoint
phase costs at most \(C|\xi|/a\).
Consequently, since \(p\ge s=2a\),
\[
        |F_{a,\mu}^d(\xi)-F_{p,\mu}^d(\xi)|
                              \le C(1+|\xi|)/s.
\]
The same conclusion for \(a=1\) follows by
increasing the constant. The shell matrices
are uniformly bounded by V18. Thus the
difference transform \(D_s^d\) obeys
\[
        |D_s^d(\xi)|\le CT/s,\qquad |\xi|\le T,\quad T\ge1.
\]
It also has a uniformly bounded periodized
discrete \(L^2\) norm, by the kernel bounds
above and those for \(F_p\).

Use the exact covariance formula
\eqref{eq:v18-discrete-forms} on side \(M\),
retaining all modes in the fine Brillouin zone.
The sine chord bound gives
\(\|\Pi_p/\omega_p^2\|\le C|\xi|^{-2}\).
Uniformly for \(M\ge2,T\ge1\),
\[
          M^{-4}\sum_{0<|\xi|\le T}|\xi|^{-2}\le CT^2 .
\]
For completeness, grouping the integer indices
into unit radial shells bounds the sum of
\(|k|^{-2}\) up to radius \(MT/(2\pi)\)
by a constant times its radius squared;
the factor \(M^{-2}\) gives the displayed
inequality. If there are no modes, it is trivial.
The low-frequency variance is therefore at
most \(CT^4/s^2\). At frequencies larger
than \(T\), discrete Parseval bounds it
by \(CT^{-2}\). Hence
\[
                       \|E_{L,M,s}\|_2^2
                                \le C(T^4/s^2+T^{-2}).
\]
Choose \(T=s^{1/3}\). This is inside the
available fine frequency scale since \(p\ge s\);
empty frequency subsets only improve the bound.
The result is \(Cs^{-2/3}\) for the variance.
The error is a first-chaos variable, so Gaussian
moments give \eqref{eq:v26-error-bound}.
Summing the geometric tail completes the proof.
\end{proof}

No \(L^1\) convergence of \(K_s\) to its averaged
kernel was asserted. Its bounded microscopic
oscillations are instead controlled by the
negative-order covariance multiplier. This is
why matching just the row mass, without the
Fourier error estimate, would have been insufficient.

\subsection{The third-chaos local field}

\begin{lemma}[Local third-chaos tails]
\label{lem:v26-third}
For every finite \(q\ge1\),
\begin{equation}\label{eq:v26-third}
     \|B_{s/2}\mathsf P_3P_{L,sM}\|_q\le C_q/s.
\end{equation}
Thus \(\mathsf P_3R_{L,M}\) is uniformly
bounded locally and its shell tail above
\(s=2^b\) is at most \(C_q2^{-b}\).
At fixed \(M\), it converges in every finite
\(L^q\) on the common Gaussian family.
\end{lemma}
\begin{proof}
Lemma~\ref{lem:v25-third} bounds the covariance
of the local primitive outputs by
\[
                   C(1+\operatorname{dist}_{sM}(2r,2r'))^{-6}.
\]
Its absolute row sum on the side-\(sM/2\)
output lattice is bounded uniformly. The
variance of the weighted average with row
\(b_a\), \(a=s/2\), is therefore at most
\[
             C\sum_r b_a(r)^2
                   \le C(\max_r b_a(r))\sum_r b_a(r)
                   \le Ca^{-2}.
\]
For the first inequality one may use
\(2|b_a(r)b_a(r')|\le b_a(r)^2+b_a(r')^2\)
and the symmetric covariance bound.
This proves the \(L^2\) estimate.
Gaussian hypercontractivity for degree three,
already used in V25, proves the \(L^q\) estimate.

For fixed \(s,M\), the finite set of incoming
\(Y_{L,sM},Q_{L,sM}\) converges jointly in
every finite \(L^q\) by V18 and
Lemma~\ref{lem:v25-Qmom}. Polynomial multiplication,
using higher finite moments, preserves this
convergence. The first and third chaos
projections are contractions in \(L^2\);
their differences still have bounded degree,
so the convergence holds in all finite
\(L^q\). The outer \(B_{s/2}\) is a fixed
finite row. Truncate the shell series,
pass to the limit in its finite head,
and use the uniform geometric tail.
\end{proof}

\subsection{The partial source finite part}

The relevant matrix is not the complete
\(\Gamma_L\), because the last \(m\) shells
are absent. Put
\begin{equation}\label{eq:v26-partial-Gamma}
       \Gamma_L^{>M}=\sum_{k=m+1}^nG_{n,k}
                    =\Gamma_L-\sum_{k=1}^mG_{n,k}.
\end{equation}

\begin{lemma}[Two endpoints for a stopped chain]
\label{lem:v26-endpoints}
Uniformly in \(n\ge m\ge1\),
\begin{equation}\label{eq:v26-Gamma-bound}
               \|\Gamma_L^{>M}-(n-m)\Gamma_*\|\le C.
\end{equation}
For fixed \(M=2^m\), as \(n\to\infty\),
\begin{equation}\label{eq:v26-BM}
 \begin{aligned}
 \Gamma_L^{>M}-(n-m)\Gamma_*&\longrightarrow B_M,\\
 B_M
   &=\Gamma_{\rm fin}+m\Gamma_*-\sum_{k=1}^m\Gamma_k^{\rm c}\\
   &=\sum_{k>m}(\Gamma_k^{\rm c}-\Gamma_*)
                  +\sum_{r\ge0}(\Gamma_r^{\rm f}-\Gamma_*).
 \end{aligned}
\end{equation}
Both series are absolutely convergent and
\(\sup_m\|B_{2^m}\|<\infty\).
\end{lemma}
\begin{proof}
The quantitative estimate \eqref{eq:v22-bulk-rate}
gives, with \(\alpha=1/20\),
\[
            \|G_{n,k}-\Gamma_*\|
                    \le C(2^{-\alpha k}+2^{-\alpha(n-k)}).
\]
Its sum over \(m<k\le n\) is uniformly
bounded, proving \eqref{eq:v26-Gamma-bound},
including an empty sum. At fixed \(m\),
subtract the finite coarse head in
\eqref{eq:v26-partial-Gamma} from
Theorem~\ref{thm:v22-relative}, using
\(G_{n,k}\to\Gamma_k^{\rm c}\) for each
fixed \(k\). This proves the first
expression for \(B_M\). Substitution
of \eqref{eq:v22-matrix-fin} gives the
second expression. V22's geometric
endpoint bounds prove absolute convergence
and the uniform bound.
\end{proof}

In particular the fine-endpoint series survives
even when the coarse stopping side grows.
Keeping only the continuum coarse shells would
give the wrong finite part.

\subsection{A cofinal renormalized local cubic coefficient}

Define the explicitly subtracted coefficient
\begin{equation}\label{eq:v26-Rhat}
       \widehat R_{L,M}
           =R_{L,M}-j\,\mathcal A_{p,M}(\Gamma_*)Z_L,
                  \qquad j=\log_2(L/M).
\end{equation}
This is a definition of an auxiliary Taylor
coefficient. It is not a redefinition of the
canonical balanced root or of its measured law.

\begin{theorem}[Renormalized local cubic limit]
\label{thm:v26-local}
Uniformly in dyadic \(L=pM\ge M\ge2\),
the field \(\widehat R_{L,M}\) has bounded
local \(L^q\) norm for every finite \(q\).
For each fixed \(M\), it converges jointly
on every finite carrier in every finite
\(L^q\) on the common Gaussian realization.
For some constants \(\eta>0,C<\infty\),
independent of \(L,M\), each output link obeys
\begin{equation}\label{eq:v26-stretched}
         \mathbb E\exp\{\eta|\widehat R_{L,M}|^{2/3}\}\le C.
\end{equation}
The local limit is
\begin{equation}\label{eq:v26-limit}
 \widehat R_{\infty,M}
     =R_{\infty,M}^{(3)}
        +\sum_{\ell\ge1}E_{\infty,M,2^\ell}
                             +\mathcal A_{\infty,M}(B_M),
\end{equation}
where the error series has \(L^q\) tail
\(C_q2^{-b/3}\) above \(\ell=b\), and
the third-chaos shell series has tail
\(C_q2^{-b}\). All summands use the same
Gaussian field.
\end{theorem}
\begin{proof}
Decompose every cubic shell into its two
chaoses and use \eqref{eq:v26-error}.
Linearity of \(\mathcal A_{p,M}\) gives
the exact finite identity
\begin{equation}\label{eq:v26-decomposition}
 \widehat R_{L,M}
     =\mathsf P_3R_{L,M}
       +\sum_{\ell=1}^j E_{L,M,2^\ell}
       +\mathcal A_{p,M}
                   \bigl(\Gamma_L^{>M}-j\Gamma_*\bigr)Z_L .
\end{equation}
Lemmas~\ref{lem:v26-error}, \ref{lem:v26-third},
\ref{lem:v26-endpoints}, and
\eqref{eq:v26-A-bound} prove the uniform
local \(L^2\) bound by the triangle inequality.
This argument uses neither independence
of different shells nor positivity of
the matrix \(\Gamma_*\).

At fixed \(M\) and fixed \(\ell\), the
first term defining \(E_{L,M,2^\ell}\)
converges by the fixed-shell argument
in Lemma~\ref{lem:v26-third}. The other
term converges because
\(G_{n,m+\ell}\to\Gamma_{m+\ell}^{\rm c}\)
and \eqref{eq:v26-A-limit} holds.
Their difference defines \(E_{\infty,M,2^\ell}\).
The uniform summable tail proves convergence
of the full error sum. The third-chaos
part converges by Lemma~\ref{lem:v26-third}.
Finally \eqref{eq:v26-BM}, the uniform
filter bound, and \eqref{eq:v26-A-limit}
prove convergence of the last term
in \eqref{eq:v26-decomposition}.
This gives \eqref{eq:v26-limit} in \(L^2\).

Each random variable is in the sum of the
first and third Gaussian chaoses. The
degree-three hypercontractive bound yields
\(\|\widehat R_{L,M}\|_q\le Cq^{3/2}\)
for \(q\ge2\), with a fixed colour-dependent
constant for the three-vector at one link.
The same estimate on differences upgrades
convergence to every finite \(L^q\).
It also proves \eqref{eq:v26-stretched}:
expand the exponential, bound its \(k\)-th
moment using \(q=2k/3\) when \(k\ge3\),
and obtain a bound by
\((C\eta)^k k^k/k!\). This series converges
for sufficiently small \(\eta\).
The first two moments are bounded separately.
Finite carriers follow by a finite sum
of component estimates, without asserting
a volume-independent norm for a growing carrier.
\end{proof}

The theorem deliberately subtracts the finite
filter in \eqref{eq:v26-Rhat}. Mere convergence
of \(\mathcal A_{p,M}(\Gamma_*)Z_L\) does not
justify replacing it there by
\(\mathcal A_{\infty,M}(\Gamma_*)\): that
replacement would require a rate strong
enough after multiplication by \(j\).
No such additional rate is asserted here.

\begin{proposition}[The unrenormalized coefficient]
\label{prop:v26-drift}
At each fixed \(M\ge2\), as \(j\to\infty\),
\begin{equation}\label{eq:v26-drift}
              \frac{R_{L,M}}j
                     \longrightarrow\mathcal A_{\infty,M}(\Gamma_*)
                              \quad\hbox{in every finite }L^q.
\end{equation}
If the limiting linear observable at a specified
link is nonzero in \(L^2\), the norm of the
unrenormalized cubic coefficient there grows
linearly in \(j\).
\end{proposition}
\begin{proof}
Divide \eqref{eq:v26-Rhat} by \(j\).
The uniformly bounded subtracted coefficient
tends to zero after division, and
\eqref{eq:v26-A-limit} gives the limit.
Convergence in \(L^2\) also implies
\(\|R_{L,M}\|_2/j\to
\|\mathcal A_{\infty,M}(\Gamma_*)\|_2\).
\end{proof}

The final assertion is conditional.
This note does not evaluate \(\Gamma_*\)
or prove that its action through this
particular transverse filter is nonzero.
The exact finite matrix \(\Gamma_2\) alone
would not prove either claim.
At side one the finite counterterm is
identically zero by \eqref{eq:v26-A-one},
consistent with V25. The terminal global
coefficient therefore cannot by itself
detect this possible local secular drift.

\subsection{A joint small-noise consequence for the three-jet}

There is a useful precise warning for a
future remainder estimate. Define only
the truncated polynomial
\[
 T_{L,M}^{[3]}(\beta)
        =\beta^{-1/2}Y_{L,M}
               +\beta^{-1}Q_{L,M}
               +\beta^{-3/2}R_{L,M}.
\]

\begin{proposition}[Coupled depth and noise in the three-jet]
\label{prop:v26-joint}
Fix \(M\ge2\). If \(\beta_L\to\infty\)
and \(j/\beta_L\to\tau<\infty\), then
\begin{equation}\label{eq:v26-joint}
       \sqrt{\beta_L}\,T_{L,M}^{[3]}(\beta_L)
           \longrightarrow
           Y_{\infty,M}+\tau\mathcal A_{\infty,M}(\Gamma_*)
                      \quad\hbox{in every finite }L^q .
\end{equation}
\end{proposition}
\begin{proof}
Substitute \eqref{eq:v26-Rhat} to write the
left side exactly as
\[
 Y_{L,M}+\beta_L^{-1/2}Q_{L,M}
   +(j/\beta_L)\mathcal A_{p,M}(\Gamma_*)Z_L
   +\beta_L^{-1}\widehat R_{L,M}.
\]
The first and third terms converge as stated.
The other terms vanish by V25's local
quadratic bound and Theorem~\ref{thm:v26-local}.
\end{proof}

For positive \(\tau\), a nonzero limiting drift
would survive the leading Gaussian scaling.
The two Gaussian terms in \eqref{eq:v26-joint}
are correlated, not independent.
This algebraic three-jet limit does not
control the fourth and higher coefficients:
further powers of \(j/\beta_L\) may survive.
One cubic calculation neither proves an
exponential resummation nor identifies a
physical Yang--Mills bare-coupling trajectory.

\subsection{What this changes in the upstream task}

The local cubic coefficient has now been
separated into a bounded, cofinally convergent
part and an explicit scale-counted linear
filter. A uniform nonlinear local remainder
must account for this structure. Attempting
to deduce that remainder merely from V25's
global cubic convergence would miss it.
The new estimate is therefore an upstream
preparation for that remainder, not its proof.

There are two additional constraints on
using \eqref{eq:v26-Rhat} in the physical program.
First, the microscopic filter generally cannot
be identified with a matrix acting only on
the retained \(Y_{L,M}\); the exact
counterexample above rules out that blanket
replacement. Applying V19's retained-coordinate
transport formula would require first proving
that an appropriate observable change has
actually been defined. Second, \(\Gamma_*\)
is the source-contraction matrix, not the
complete measured combination
\(\Gamma_*-2h_*^T\). A source subtraction
must not silently replace the density,
its ghost or Haar factors, or their low
Born terms.

The immediate unresolved analytic task remains
a stable local control of the full nonlinear
map on typical multiscale fluctuations,
including all lower-chaos contractions at
higher degrees and the actual incoming
measured law. An explicit evaluation or
structural cancellation proof for the
linear filter \(\mathcal A_{\infty,M}(\Gamma_*)\)
would also decide whether the conditional
local growth in Proposition~\ref{prop:v26-drift}
actually occurs.

\subsection{Exact checks and preservation record}

The checker
\begin{center}\small
\path{scripts/verify_ym_f14_v26_local_cubic_drift.py}
\end{center}
has SHA--256
\begin{center}\scriptsize\ttfamily
5E8A4B45D429E86F382336890B92FCA4EE3540D37D4EA1B96F17593E2E984D1E.
\end{center}
It checks the full matrix \eqref{eq:v26-mass-two}
by exact rational contraction of the literal
side-two root, the live transverse example
for \eqref{eq:v26-notcommute}, and eight
side-one row identities. Ninety-six exact
triangular-array tests check the two-endpoint
subtraction algebra on an explicit model
array; these are not numerical evaluations
of the unknown \(\Gamma_*\).
Fifteen rational scaling checks verify the
outer mass factor and the kernel-density
factor in \eqref{eq:v26-transform}.
The checker also checks the cutoff exponents.
The uniform and cofinal conclusions are
proved above, not inferred from those
finite tests.

The predecessor V25 has 564,353 bytes and SHA--256
\begin{center}\scriptsize\ttfamily
48087987790EBB5A1C6E776DE17C20474E7C0443F7760DDC839C455125D8C900.
\end{center}
Its body is retained in full, except for
the title version and this terminal insertion.
V26 replaces V25 as the sole active main note.
SM V72, NS V26, Shell-Mixing V16, the
archives and monographs remain unchanged.
The next companion checkpoint is V30.

\begin{firewall}
The renormalized local cubic theorem is a
free-reference Taylor-coefficient result.
It does not construct an interacting
continuum state, control the complete
nonlinear quantum remainder, prove
noncollapse, or establish reflection-positive
physical transfer and a positive Hamiltonian
mass gap. The full Yang--Mills objective
remains unproved and unchanged.
\end{firewall}
\section{V27: retained information, hidden transverse modes, and a closed filter tower}
\label{sec:v27-retained}

\subsection{The upstream question and the scope of the answer}

The worktree still contains Shell-Mixing V16,
SM V72 and NS V26. V26's local cubic theorem
is the current starting point. It identifies
the logarithmic term
\(j\mathcal A_{p,M}(\Gamma_*)Z_L\), using
the actual cross-oriented filters.
Its example proves that these filters
need not equal a matrix times \(Y_{L,M}\).
That does not yet exclude some more general,
possibly nonlinear or spatially nonlocal,
function of the entire retained linear field.
The new question is precisely that stronger one.

We prove an exact criterion: for \(p\ge2\)
and \(M\ge2\), the cross-oriented observable
with orientation matrix \(A\) is measurable
with respect to the entire \(Y_{L,M}\)
if and only if \(A\) is diagonal.
An explicit lower bound on its conditional
variance holds on the live transverse
Gaussian space. A second construction
shows that the obstruction also survives
coarse Coulomb projection; it is not
merely a longitudinal-coordinate artifact.
We then give a sufficient enlarged linear
state with an exact blocking composition law.

The Gaussian projection formalism is
already available. Shell-Mixing V1 gives
the covariance tower and its
Moore--Penrose formula; main V15 applies
conditional expectation to a quadratic
source bank at one finite lattice.
What is new here is the explicit nullspace
of the present retained observations,
its quantitative action on V26's filters,
the projected two-direction witness,
and their consequences for that cubic drift.
No claim is made that \(\Gamma_*\) has a
nonzero off-diagonal entry. That remains
an unevaluated property of the actual
source contraction, not an assumption
to be silently supplied.

\subsection{A conditional-variance test on the live Gaussian space}

Let \(C\) be the covariance on a finite
real live space \(H\), positive definite
on that space. Removed gauge and constant
directions are not included in \(H\).
For linear observations \(Y=BZ\), \(Z\sim N(0,C)\),
and another finite family \(X=DZ\), put
\begin{equation}\label{eq:v27-regression}
 \begin{aligned}
 S&=BCB^t,& K&=DCB^t S^+,\\
 X&=KY+\eta,&
 \Sigma&=\mathbb E[\eta\eta^t]
           =DCD^t-DCB^tS^+BCD^t .
 \end{aligned}
\end{equation}
Here \(S^+\) acts only on the observed
covariance range. In particular no inverse
on a null coarse mode is being asserted.
Whitening \(Z=C^{1/2}g\) shows that
\(\eta\) is independent of \(Y\):
the explained part is the orthogonal
projection onto \(\operatorname{ran}(C^{1/2}B^t)\).
Orthogonal components of a Gaussian vector
are independent, as follows directly by
factoring its quadratic characteristic function.
Thus \(KY=\mathbb E[X\mid Y]\) and
\(\operatorname{Cov}(X\mid Y)=\Sigma\).

\begin{lemma}[Hidden-mode lower bound]
\label{lem:v27-hidden}
Suppose \(u_1,\ldots,u_N\in H\) satisfy
\(Bu_i=0\) and
\[
       \langle u_i,C^{-1}u_k\rangle
                       =e_i\delta_{ik},\qquad e_i>0 .
\]
Then the conditional covariance in
\eqref{eq:v27-regression} satisfies
\begin{equation}\label{eq:v27-hidden}
                \Sigma\succeq
                    \sum_{i=1}^N\frac{(Du_i)(Du_i)^t}{e_i}.
\end{equation}
In particular, for scalar \(X\) and every
square-integrable measurable function \(g(Y)\),
\[
       \mathbb E|X-g(Y)|^2
              \ge\operatorname{Var}(X\mid Y)
              \ge\sum_i |Du_i|^2/e_i .
\]
\end{lemma}
\begin{proof}
The vectors \(C^{-1/2}u_i/\sqrt{e_i}\)
are orthonormal and belong to
\(\ker(BC^{1/2})\). For any scalar row
\(a^tD\), Bessel's inequality on this
orthogonal complement bounds its
unexplained squared norm below by
\(\sum_i|a^tDu_i|^2/e_i\).
This proves the matrix inequality.
The conditional mean minimizes squared
error over all measurable functions of
the observation: expand
\(X-g(Y)=(X-\mathbb E[X\mid Y])
       +(\mathbb E[X\mid Y]-g(Y))\);
the two terms have zero cross expectation.
\end{proof}

The same proof works for an isonormal
Gaussian distribution and finitely many
continuous linear observations. In that
case the \(u_i\) are Cameron--Martin
directions and \(e_i\) are their squared
energy norms. Only those smooth finite
Fourier directions will be used below.

\subsection{The sixteen geometric channels}

Retain \(L=pM\), with \(p\ge2\) and
\(M\ge2\) dyadic. Work first with one
Lie colour. Write
\begin{equation}\label{eq:v27-X}
 X_{\mu\nu}(r)=T_p^\mu Z_\nu(r),\qquad
 T_p^\mu u(r)
    =p^{-4}\sum_{z\in\{0,\ldots,p-1\}^4}
                   \sum_{t=0}^{p-1}u(pr+z+t e_\mu).
\end{equation}
Thus \(Y_\nu=X_{\nu\nu}\) and
\[
         \bigl(\mathcal A_{p,M}(A)Z_L\bigr)_\mu
                           =\sum_\nu A_{\mu\nu}X_{\mu\nu}.
\]
Let \(\mathcal G_{p,M}\) be the sigma field
of the \emph{entire} retained \(Y\), at
all coarse sites and all colours.
The other colours are independent and
do not alter the one-colour conditional
covariances.

Set
\[
       \delta=\frac{\pi}{2p},\qquad
       d_p=\frac1{p\sin\delta},\qquad
       c_p=\frac{\cos\delta}{p^2\sin^2\delta}.
\]
The elementary geometric sum is
\begin{equation}\label{eq:v27-phase}
       \frac1p\sum_{z=0}^{p-1}e^{i\pi z/p}
                  =d_p e^{i(\pi/2-\delta)}.
\end{equation}
For \(\mu\ne\nu\) consider the fine field
\begin{equation}\label{eq:v27-one-mode}
        u^{\mu\nu}_{x,\rho}
          =p^{-1}\cos(\pi x_\mu/p+\delta)\,
                                  {\bf1}_{\rho=\nu}.
\end{equation}
It is periodic since \(M\) is even, has
zero mean, and is Coulomb because it is
constant along its input orientation \(\nu\).
It is consequently a legitimate live
direction, not a gauge or zero-mode test.

Equation \eqref{eq:v27-phase} shows that
its ordinary box average in coordinate
\(\mu\) is zero. Every retained \(Y_\rho\)
therefore vanishes on this field, at
every coarse site. The extra directional
average in geometry \(\alpha\) leaves
that zero intact unless \(\alpha=\mu\).
In the latter case the phase is doubled.
Explicitly,
\begin{equation}\label{eq:v27-one-response}
 \begin{aligned}
 B_pu^{\mu\nu}&=0,\\
 X_{\alpha\beta}(u^{\mu\nu};r)
   &=-c_p(-1)^{r_\mu}
                   {\bf1}_{\alpha=\mu}{\bf1}_{\beta=\nu}.
 \end{aligned}
\end{equation}
For example the nonzero factor is
\[
  \operatorname{Re}\!
       \left(e^{i\delta}d_p^2e^{2i(\pi/2-\delta)}\right)
                                                   =-c_p .
\]
The factor \(p^{-1}\) in the fine field
cancels the mass \(p\) of \(T_p^\mu\).

The covariance inverse on a live transverse
Fourier mode is the scalar lattice Laplacian.
The eigenvalue of \eqref{eq:v27-one-mode}
is \(4\sin^2\delta\), and its fine squared
norm is \(p^2M^4/2\). Distinct pairs
\((\mu,\nu)\), \(\mu\ne\nu\), are
energy orthogonal: different input
orientations are orthogonal, and distinct
wave directions have different Fourier
frequencies. Hence
\begin{equation}\label{eq:v27-one-energy}
 \langle u^{\mu\nu},C_L^{-1}u^{\alpha\beta}\rangle
   =2p^2\sin^2\delta\,M^4\,
                         {\bf1}_{(\mu,\nu)=(\alpha,\beta)}.
\end{equation}

\begin{theorem}[The retained-information criterion]
\label{thm:v27-information}
At any fixed coarse site, the twelve
off-diagonal variables \(X_{\mu\nu}\),
\(\mu\ne\nu\), have conditional covariance
given \(\mathcal G_{p,M}\) bounded below by
\begin{equation}\label{eq:v27-kappa}
       \kappa_{p,M}I_{12},\qquad
       \kappa_{p,M}
          =\frac{\cos^2\delta}{2p^6\sin^6\delta\,M^4}
          \ge\frac{16}{\pi^6M^4}.
\end{equation}
Consequently, for each row of a real \(A\),
\begin{equation}\label{eq:v27-A-variance}
 \operatorname{Var}\!\left(
  (\mathcal A_{p,M}(A)Z_L)_\mu(r)\mid\mathcal G_{p,M}\right)
          \ge\kappa_{p,M}\sum_{\nu\ne\mu}|A_{\mu\nu}|^2 .
\end{equation}
The complete observable
\(\mathcal A_{p,M}(A)Z_L\) is measurable
with respect to \(\mathcal G_{p,M}\)
if and only if \(A\) is diagonal.
\end{theorem}
\begin{proof}
Apply Lemma~\ref{lem:v27-hidden} to the
twelve directions \eqref{eq:v27-one-mode}
and the twelve observation rows at the
chosen site. Their response matrix is
diagonal up to signs, with magnitude
\(c_p\), and their energy Gram matrix
is \eqref{eq:v27-one-energy}.
This gives \eqref{eq:v27-kappa}.
The elementary inequalities
\(\cos^2\delta\ge1/2\) and
\(\sin\delta\le\delta=\pi/(2p)\)
give its lower bound.

The diagonal \(X_{\mu\mu}=Y_\mu\) is
already observed, so taking a row
combination proves \eqref{eq:v27-A-variance}.
Any nonzero off-diagonal entry then
leaves strictly positive conditional
variance, excluding all measurable
reconstructions, not just linear ones.
For diagonal \(A\) the observable is
exactly \(AY\), proving the converse.
\end{proof}

Neither exceptional endpoint is included
in this criterion. At \(p=1\) all geometric
rows coincide and every matrix acts on
the observed field. At \(M=1\) all the
averages vanish on mean-zero inputs,
as proved in \eqref{eq:v26-A-one}.

\subsection{The continuum witness without a projection-limit assumption}

At fixed \(M\), the sixteen filters have
the common Gaussian continuum limit
provided by V18 and V26. On
\((\mathbb R/M\mathbb Z)^4\), use
\[
             u^{\mu\nu}(x)=\cos(\pi x_\mu)e_\nu .
\]
Its energy is \(\pi^2M^4/2\).
Its retained averages vanish. Its
cross average is
\[
             X_{\mu\nu}^{\infty}(u^{\mu\nu};r)
                              =-4\pi^{-2}(-1)^{r_\mu}.
\]
The same energy-orthogonal test gives
\begin{equation}\label{eq:v27-continuum}
 \operatorname{Cov}\bigl(
     (X_{\mu\nu}^{\infty}(r))_{\mu\ne\nu}
                  \mid Y_{\infty,M}\bigr)
                   \succeq \frac{32}{\pi^6M^4}I_{12}.
\end{equation}
In particular the same diagonal criterion
holds for the continuum filtered Gaussian
observables.

This is a direct continuum proof. It
does not infer convergence of conditional
covariances merely from convergence of
joint covariances. Such an inference can
fail at a changing observation rank:
observing \(\epsilon Z\) reveals \(Z\)
for every \(\epsilon\ne0\), although
the limiting observation is zero.
No continuity of a pseudoinverse across
an unproved rank change is used here.

\subsection{The obstruction after coarse Coulomb projection}

Let \(\Pi_M^{\rm c}\) be the ordinary
coarse transverse projector at nonzero
coarse momentum, with constants retained.
This is the projector used in the
companion's linear source
\(\mathcal Q_p=\Pi_M^{\rm c}B_p\).
The simple witness \eqref{eq:v27-one-mode}
produces an output in direction \(\mu\)
at momentum \(\pi e_\mu\). That output
is longitudinal and is removed by
\(\Pi_M^{\rm c}\). It would therefore
be incorrect to cite that witness alone
as an obstruction for the projected source.

Choose a third direction
\(\alpha\notin\{\mu,\nu\}\), possible
in four dimensions, and instead use
\begin{equation}\label{eq:v27-two-mode}
 v^{\mu\nu,\alpha}_{x,\rho}
   =p^{-1}\cos(\pi x_\mu/p+\delta)
          \sin(\pi x_\alpha/p+\delta)\,
                                      {\bf1}_{\rho=\nu}.
\end{equation}
This field is again mean zero and Coulomb.
Every retained average still vanishes
because of the first cosine factor.
Only the cross geometry \(\mu\), input
\(\nu\), survives:
\begin{equation}\label{eq:v27-two-response}
 X_{\beta\gamma}(v^{\mu\nu,\alpha};r)
     =-c_pd_p(-1)^{r_\mu+r_\alpha}
                         {\bf1}_{\beta=\mu}{\bf1}_{\gamma=\nu}.
\end{equation}
Indeed the sine's box average is \(d_p\)
by \eqref{eq:v27-phase}.
The fine Laplacian eigenvalue is now
\(8\sin^2\delta\), and the fine squared
norm is \(p^2M^4/4\); its energy is still
\(2p^2\sin^2\delta\,M^4\).

The output momentum is
\(\pi(e_\mu+e_\alpha)\). At this coarse
momentum the two nonzero forward
gradient components both equal \(-2\).
Thus
\begin{equation}\label{eq:v27-project}
             \Pi_M^{\rm c}e_\mu
                             =\tfrac12(e_\mu-e_\alpha).
\end{equation}
In particular a nonzero transverse
component remains.

\begin{theorem}[Projected retained-information bound]
\label{thm:v27-projected}
For each fixed row \(\mu\),
\begin{equation}\label{eq:v27-projected}
 \operatorname{Var}\!\left(
   [\Pi_M^{\rm c}\mathcal A_{p,M}(A)Z_L]_\mu(r)
                                  \mid\mathcal G_{p,M}\right)
       \ge \widetilde\kappa_{p,M}
                           \sum_{\nu\ne\mu}|A_{\mu\nu}|^2 ,
\end{equation}
where
\[
 \widetilde\kappa_{p,M}
       =\frac{\cos^2\delta}{8p^8\sin^8\delta\,M^4}
       \ge\frac{16}{\pi^8M^4},\qquad
 \lim_{p\to\infty}\widetilde\kappa_{p,M}
                                      =\frac{32}{\pi^8M^4}.
\]
Thus the projected cross-oriented observable
is a function of the entire unprojected
retained \(Y\) if and only if \(A\) is diagonal.
For non-diagonal \(A\) it cannot be a
function of the smaller observation
\(\Pi_M^{\rm c}Y\) either.
\end{theorem}
\begin{proof}
For each \(\nu\ne\mu\) choose one
\(\alpha\notin\{\mu,\nu\}\).
The corresponding fields in
\eqref{eq:v27-two-mode} have different
input orientations and are energy
orthogonal. Their action on the specified
projected output component is
\(-A_{\mu\nu}c_pd_p/2\), up to its
checkerboard sign. Divide the squared
responses by their energy and apply
Lemma~\ref{lem:v27-hidden}. This gives
\eqref{eq:v27-projected} and the stated
constant. The elementary sine and cosine
bounds give the uniform lower bound,
and their small-argument limits give
the displayed limit.

If \(A\) is diagonal, the projected
observable is \(\Pi_M^{\rm c}(AY)\),
a deterministic linear function of the
entire \(Y\). If \(A\) has a nonzero
off-diagonal entry, the lower bound
excludes measurability. Finally
\(\sigma(\Pi_M^{\rm c}Y)\subseteq\sigma(Y)\);
less information cannot remove a
positive minimal reconstruction error.
\end{proof}

The converse for diagonal \(A\) was
asserted with the full \(Y\) observed.
It does not assert that an arbitrary
diagonal orientation matrix commutes
with the coarse transverse projector,
or that \(\Pi_M^{\rm c}(AY)\) can
always be recovered from \(\Pi_M^{\rm c}Y\).

There is also a direct continuum version.
The smooth field
\(\cos(\pi x_\mu)\sin(\pi x_\alpha)e_\nu\)
has energy \(\pi^2M^4/2\); its cross
output is \(-8\pi^{-3}e_\mu\) times
the checkerboard sign. Projection
leaves component \(-4\pi^{-3}\).
Lemma~\ref{lem:v27-hidden} gives the
continuum constant \(32/(\pi^8M^4)\)
without any conditional-covariance
limit interchange.

\subsection{Quantitative consequences for the V26 cubic term}

For a scalar output at one fixed
link and colour, define the best
retained reconstruction error
\[
        {\rm err}_{p,M}(U)
          =\inf_{g\in L^2(\mathcal G_{p,M})}
                                             \|U-g\|_2 .
\]
The candidate \(g\) may depend on all
coarse sites and all colours, and may
depend on \(L\). It need not be linear.

\begin{proposition}[Unexplained cubic drift]
\label{prop:v27-cubic}
Write \(a_\mu^2=\sum_{\nu\ne\mu}|(\Gamma_*)_{\mu\nu}|^2\).
With \(j=\log_2(L/M)\), V26's local
cubic coefficient satisfies
\begin{equation}\label{eq:v27-cubic-error}
 \begin{aligned}
 {\rm err}_{p,M}(R_{L,M,\mu}(r))
       &\ge j\sqrt{\kappa_{p,M}}\,a_\mu-C,\\
 {\rm err}_{p,M}([\Pi_M^{\rm c}R_{L,M}]_\mu(r))
       &\ge j\sqrt{\widetilde\kappa_{p,M}}\,a_\mu-C .
 \end{aligned}
\end{equation}
The constants are independent of \(L,M\).
Negative right sides impose no additional
restriction beyond nonnegativity.
\end{proposition}
\begin{proof}
Let \(\mathcal E_{p,M}U=\mathbb E[U\mid\mathcal G_{p,M}]\).
Then \({\rm err}_{p,M}(U)=\|(I-\mathcal E_{p,M})U\|_2\).
Apply this contraction to the exact
decomposition
\[
        R_{L,M}=j\mathcal A_{p,M}(\Gamma_*)Z_L
                                             +\widehat R_{L,M}.
\]
The reverse triangle inequality, V26's
uniform local \(L^2\) bound, and
Theorem~\ref{thm:v27-information} prove
the first assertion.

For the second, the projected remainder
also has a uniform local \(L^2\) bound.
Indeed the law of \(\widehat R_{L,M}\)
is invariant under coarse translations,
and \(\Pi_M^{\rm c}\) is a translation-invariant
orthogonal projection on the whole
coarse link space. Consequently
\[
 \begin{aligned}
 \mathbb E\sum_\mu|[\Pi_M^{\rm c}\widehat R_{L,M}]_\mu(r)|^2
 &=M^{-4}\mathbb E\|\Pi_M^{\rm c}\widehat R_{L,M}\|_{\ell^2}^2\\
 &\le M^{-4}\mathbb E\|\widehat R_{L,M}\|_{\ell^2}^2\le C .
 \end{aligned}
\]
This uses spatial stationarity and an
\(\ell^2\) projection, not an unproved
uniform \(\ell^\infty\) bound for a
Riesz-type projector. Repeat the
contraction argument using
Theorem~\ref{thm:v27-projected}.
\end{proof}

There is a corresponding statement for
V26's truncated small-noise three-jet.
Fix \(M\), let \(\beta_L\to\infty\),
and suppose \(j/\beta_L\to\tau\ge0\).
For its normalized scalar component,
\begin{equation}\label{eq:v27-joint-error}
 \begin{aligned}
 \liminf_{L\to\infty}{\rm err}_{p,M}
     \bigl(\sqrt{\beta_L}\,T_{L,M,\mu}^{[3]}(\beta_L)(r)\bigr)
       &\ge \tau\sqrt{\frac{32}{\pi^6M^4}}\,a_\mu,\\
 \liminf_{L\to\infty}{\rm err}_{p,M}
     \bigl([\Pi_M^{\rm c}\sqrt{\beta_L}\,
                           T_{L,M}^{[3]}(\beta_L)]_\mu(r)\bigr)
       &\ge \tau\sqrt{\frac{32}{\pi^8M^4}}\,a_\mu .
 \end{aligned}
\end{equation}
To prove these inequalities, expand the
normalized three-jet as in
Proposition~\ref{prop:v26-joint}.
The retained linear term has zero
reconstruction error. The contributions
\(\beta_L^{-1/2}Q_{L,M}\) and
\(\beta_L^{-1}\widehat R_{L,M}\)
tend to zero in \(L^2\), by V25--V26.
Their projected versions have the same
property by the stationarity argument
above. The cross-filter coefficient
is \(j/\beta_L\). Apply the finite
conditional-variance bounds and use
the explicit limits of their constants.
No convergence of the conditioning
operators is required.

These are conditional consequences for
the actual \(\Gamma_*\): if its
off-diagonal part vanishes, all these
particular lower bounds are zero.
They do not prove that the source
matrix is non-diagonal. Nor do they
exclude a construction using the
complete nonlinear retained root,
a larger state, or conditional integration.
They concern arbitrary functions of
the specified \emph{linear} retained
Gaussian observation.

\subsection{An exact sufficient enlargement}

Instead of pretending that all the cross
filters are determined by \(Y\), one
can retain the sixteen channels
\(X_{\mu\nu}\) in \eqref{eq:v27-X}.
There are four diagonal and twelve
off-diagonal channels per colour.
This is an enlarged family of linear
observables, not a newly asserted
physical gauge field.

\begin{proposition}[A closed geometric filter tower]
\label{prop:v27-tower}
For positive integral blocking factors
\(p,q\) and a compatible periodic lattice,
\begin{equation}\label{eq:v27-tower}
                     T_q^\mu T_p^\mu=T_{pq}^\mu.
\end{equation}
Thus applying \(T_q^\mu\) to every
channel \(X_{\mu\nu}^{(p)}\), with
its first index fixed, produces
\(X_{\mu\nu}^{(pq)}\). Every V26
matrix filter is a pointwise linear
function of this enlarged state.
These enlarged observation spaces
are nested along successive blocks.
\end{proposition}
\begin{proof}
Unroll the two averages. Their fine
argument is
\[
   pq r+p z_q+z_p+(p t_q+t_p)e_\mu ,
\]
and their normalization is \(p^{-4}q^{-4}\).
The pairs of box indices give every
\(z\in\{0,\ldots,pq-1\}^4\) exactly
once by \(z=pz_q+z_p\).
Likewise \(t=pt_q+t_p\) runs once
through \(\{0,\ldots,pq-1\}\).
This is precisely the defining row
of \(T_{pq}^\mu\), including periodic
wraparound. It proves the identity
and the stated nesting. The matrix
filter is its defining row sum
\(\sum_\nu A_{\mu\nu}X_{\mu\nu}\).
\end{proof}

The free reference gives uniform local
Gaussian moment bounds for these channels
by \eqref{eq:v26-A-bound}, and their
fixed-\(M\) cofinal limit exists by
the same finite-filter convergence.
Thus the enlargement is analytically
compatible with the established
Gaussian estimates.
The innovation covariance bound
\eqref{eq:v27-kappa} also shows that,
at one fixed site and colour, the
twelve off-diagonal channels are
linearly independent modulo the
observed Gaussian linear space.
Recovering all twelve by additional
linear Gaussian observations requires
at least twelve extra scalar linear
coordinates in total. This is not
a claim of minimal dimension per
site for every spatially nonlocal
description, or a dimension bound
for arbitrary measurable encodings.

Nesting does not make an interacting
effective density Gaussian, local,
or Markovian in these channels.
If this enlarged state is adopted
in the nonlinear program, its
actual conditional law, source
jets, and incoming density still
have to be propagated.

\subsection{What is lost by discarding the hidden conditional field}

At each finite lattice, apply
\eqref{eq:v27-regression} to
\(X=\mathcal A_{p,M}(A)Z_L\),
observing the full \(Y\).
On any declared finite output family,
the conditional characteristic function
of the affine observable \(Y+\tau X\) is
\begin{equation}\label{eq:v27-characteristic}
 \mathbb E\!\left[e^{i\langle h,Y+\tau X\rangle}\mid Y\right]
     =\exp\!\left(
         i\langle h,(I+\tau K)Y\rangle
                    -\frac{\tau^2}{2}\langle h,\Sigma h\rangle
             \right).
\end{equation}
Here \(K\) may be spatially nonlocal,
and \(\Sigma\) is the actual conditional
covariance, not an independently chosen
noise model. The formula follows by
integrating the independent centred
Gaussian \(\eta\) in
\(X=KY+\eta\). A deterministic
retained-coordinate displacement has
no such conditional variance factor.
When the appropriate off-diagonal
entry of \(A\) is nonzero,
Theorems~\ref{thm:v27-information}
and \ref{thm:v27-projected} ensure
that the factor cannot simply be
set equal to one, including for
the projected source.

This is an exact Gaussian statement
for the displayed affine observable.
It is not a full resummation of
the source, and it does not isolate
all terms at order \(\tau^2\) in
the nonlinear measured problem.
Higher Taylor coefficients may
contribute at that same scaling.
The constructive choices indicated
by this calculation are to keep
the enlarged filters, or to integrate
their genuine conditional innovations
while retaining all resulting terms.
Resetting their conditional covariance
would not be either procedure.

\subsection{Exact finite certificates}

For \(L=4,p=2,M=2\), the single-direction
witness can be written without irrational
normalization as
\[
 u_{x,\nu}
      =\cos(\pi x_\mu/2)-\sin(\pi x_\mu/2),
                       \qquad \mu\ne\nu.
\]
Its values are \(1,-1,-1,1\) along
the active direction. It has energy
\(512\), all retained averages vanish,
and its cross output at zero is \(-1\).
The twelve such fields have energy
Gram matrix \(512I_{12}\). This
certifies the finite lower bound
\(I_{12}/512\), which is the value
of \(\kappa_{2,2}\).

Multiplying that profile by
\(\cos(\pi x_\alpha/2)+\sin(\pi x_\alpha/2)\),
with \(\alpha\notin\{\mu,\nu\}\),
gives the two-direction witness.
Its energy is \(1024\), its cross
output is again \(-1\), and the
specified projected component is
\(-1/2\). It certifies the projected
lower bound \(1/4096\), equal to
\(\widetilde\kappa_{2,2}\).

The checker
\begin{center}\small
\path{scripts/verify_ym_f14_v27_retained_innovations.py}
\end{center}
has SHA--256
\begin{center}\scriptsize\ttfamily
16C149050F75C7D1D95179C42CA95046380ED80AC7D5A138FE255A31DE56A820.
\end{center}
It verifies twelve witnesses of each
kind, 768 geometric rows for each
family, their live constraints and
energy identities, and sixteen exact
filter-composition identities.
It also computes the full Gaussian
regression on the entire side-two
retained field, not just on its
value at one site.

For reproducibility, the latter
calculation splits into the sixteen
coarse momenta \(k\in\{0,1\}^4\).
Their fine aliases are
\(a=k+2\ell\), \(\ell\in\{0,1\}^4\),
on the side-four lattice. Put
\[
 \begin{aligned}
 q_\rho(a)&=i^{a_\rho}-1,\qquad
 \lambda(a)=\sum_\rho|q_\rho(a)|^2,\\
 b_\rho(a)&=\frac1{16}
        \left(\prod_{\sigma=1}^4(1+i^{a_\sigma})\right)
                                      (1+i^{a_\rho}).
 \end{aligned}
\]
For \(a\ne0\) the exact covariance
multiplier is
\(\delta_{\rho\sigma}/\lambda
          -q_\rho\overline q_\sigma/\lambda^2\).
Use rows \((b_1e_1^t,\ldots,b_4e_4^t,b_1e_2^t)\)
for \((Y_1,\ldots,Y_4,X_{12})\),
sum their covariance contractions
over the aliases with factor \(4^{-4}\),
and take the Schur residual on the
range of the first four rows.
All phases lie in \(\{1,i,-1,-i\}\);
the summed matrices and their
consistent linear solves are rational.
Summing the sixteen independent
coarse-mode residuals gives
\begin{equation}\label{eq:v27-exact-variances}
 \begin{aligned}
 \operatorname{Var}(X_{12}(0)\mid Y)
                                    &=\frac{1469}{147456},\\
 \operatorname{Var}([\Pi_2^{\rm c}\mathcal A_{2,2}(E_{12})Z_4]_1(0)
                                                   \mid Y)
                                    &=\frac{11017}{2359296}.
 \end{aligned}
\end{equation}
Here \(E_{12}\) is the orientation matrix
with its only nonzero entry equal to
one at \((1,2)\), not a Lie basis element.
For the projected second line, multiply
the residual at nonzero coarse \(k\)
by \((1-k_1/\sum_\rho k_\rho)^2\).
The one-axis contribution is thereby
removed, whereas the contribution
at \(k=(1,0,1,0)\) is \(1/4096\).
The checker performs this enumeration
with exact rational complex arithmetic.
These finite identities supplement,
and do not replace, the all-scale
hidden-mode proofs.

\subsection{Preservation and the next unresolved mathematical step}

The predecessor V26 has 591,040 bytes and SHA--256
\begin{center}\scriptsize\ttfamily
18E39F118CD14E3004199BC06946EC0E077C3FFBF341EDBA8C5AFF0BAB6C029D.
\end{center}
Its complete body is retained, apart
from the title version and this
terminal insertion. V27 replaces
V26 as the sole active main note.
The companion files, monographs
and archives are unchanged; the
next scheduled SM/NS review is V30.

The local nonlinear-remainder problem
has not been solved here. What has
been settled is an information issue
that a proposed retained-coordinate
renormalization must respect. For
the leading V26 source matrix,
diagonality is now a precise test:
if it fails, there is quantitatively
nonzero hidden Gaussian information,
even after coarse Coulomb projection.
The sixteen-channel tower supplies
an exact linear enlargement if needed.
The next substantive choices are to
evaluate or structurally constrain
the actual \(\Gamma_*\), or to develop
stable nonlinear estimates with those
additional channels and their measured
conditional law. Neither choice
allows replacement of the complete
measured matrix \(\Gamma_*-2h_*^T\)
by the source matrix alone.

\begin{firewall}
The positive conditional variances
proved here are free-reference
information bounds, not physical
mass gaps. The calculation does
not prove that the actual limiting
source matrix is non-diagonal, does
not construct an interacting
continuum state, and does not
establish a stable full nonlinear
remainder, noncollapse, physical
reflection positivity or Hamiltonian
spectral positivity above the vacuum.
The full Yang--Mills objective
remains unproved and unchanged.
\end{firewall}
\section{V28: the actual off-diagonal source matrix is strictly nonzero}
\label{sec:v28-source}

\subsection{The question resolved in this iteration}

V27 left a specific condition open: is the
actual bulk source matrix \(\Gamma_*\)
diagonal? This matters because an
off-diagonal entry activates the hidden
conditional fluctuations proved there,
even after coarse Coulomb projection.
We now resolve that condition for the
specified ascending-tree balanced root.

The result is
\begin{equation}\label{eq:v28-structure}
 \Gamma_*=
 \begin{pmatrix}
  \gamma_1&0&0&0\\
  0&\gamma_2&a_*&a_*\\
  0&0&\gamma_3&b_*\\
  0&0&0&\gamma_4
 \end{pmatrix},
             \qquad 0<a_*,\qquad 2a_*<b_*<3a_* .
\end{equation}
The diagonal entries are real but are
not evaluated or assigned signs here.
The proof retains the full incoming
mixed response. It consists of an
exact finite Laurent-polynomial
reduction, a telescoping sum over
the incoming scales, and positive
one-dimensional heat-kernel integrals.
Numerical quadrature is not used to
certify a sign or a continuum limit.

The matrix in \eqref{eq:v28-structure}
is precisely \eqref{eq:v18-Gamma-star},
not a replacement normalization.
Shell-Mixing V16, SM V72 and NS V26
are unchanged. The scheduled
cross-program review remains V30.

\subsection{Plane-wave conventions and the complete incoming response}

Use positive plane waves \(e^{i\xi\cdot x}\).
Their averaging multipliers are the
complex conjugates of the negative
Fourier-transform convention in V18;
real covariance integrals are unchanged.
Put
\[
 \begin{aligned}
 \beta(x)&=\frac{e^{ix}-1}{ix},\quad \beta(0)=1,\\
 S(x)&=|\beta(x)|^2
                     =\left(\frac{\sin(x/2)}{x/2}\right)^2,\\
 C(x)&=\cos^2(x/2),\qquad
 \mathcal B(\xi)=\prod_{\rho=1}^4\beta(\xi_\rho).
 \end{aligned}
\]
The continuum retained linear field
has multiplier
\[
                 Y_\alpha:\quad
                   \mathcal B(\xi)\beta(\xi_\alpha).
\]
For a dyadic incoming distance \(s=2^\ell\),
\(\ell\ge1\), let
\(L_{\lambda,\alpha i}(z)\) be the
Fourier polynomial of the exact V17
mixed stencil, with external constant
source in direction \(\lambda\),
input noise in direction \(i\), and
output in direction \(\alpha\).
Here \(z_\rho=e^{i\xi_\rho/s}\).
This is the coefficient in the
perpendicular Lie-colour direction;
it is the stencil implemented by
\texttt{mixed\_row}, including its
normalization.

For generic frequencies, introduce
the formal input coordinates
\[
                 V_i=\frac{\mathcal B(\xi)}{i\xi_i}W_i,
                  \qquad U_i=e^{i\xi_i}.
\]
The transformed covariance is real:
\begin{equation}\label{eq:v28-covariance}
 \operatorname{Cov}_{\xi}(V)
    =\frac{|\mathcal B(\xi)|^2}{|\xi|^4}
       \left\{\operatorname{diag}
             \left(\frac{|\xi|^2}{\xi_1^2},\ldots,
                   \frac{|\xi|^2}{\xi_4^2}\right)
                                    -{\bf1}{\bf1}^t\right\}.
\end{equation}
This is an algebraic decomposition of
the original transverse multiplier,
not an assertion that its negative
rank-one summand is a covariance.
Contractions are linear in this
matrix, so they can be evaluated
on the four diagonal unit matrices
and \({\bf1}{\bf1}^t\) separately.
All scalar weights in this decomposition
are even under each sign change of
\(\xi_\rho\).
Coordinate hyperplanes omitted in
this calculation have measure zero;
the original filters provide their
removable interpretations.

In these coordinates the linear
multiplier is
\[
                         Y_\alpha=(U_\alpha-1)V_\alpha.
\]
The contribution \(\Delta^{(s)}\)
of the \(s\)-th incoming scale
to \(\Delta(v)\) in V18 has multiplier
\begin{equation}\label{eq:v28-delta-symbol}
 \Delta^{(s)}_{\lambda,\alpha i}
   =\frac1s\,
       \left(\prod_{\rho=1}^4\frac2{1+z_\rho}\right)
       \frac{U_\alpha-1}{z_\alpha^2-1}
                      (z_i-1)L_{\lambda,\alpha i}(z)
\end{equation}
acting on \(V_i\).
This identity holds away from the
displayed removable denominators.

To check the scale factor, the incoming
term is
\[
             s^{-1}B_{s/2}\mathcal L_v(\mathcal A_sW).
\]
The inner continuum row has mass \(1/s\)
and the outer row mass \(s/2\).
The normalized outer box and the
inner continuum box satisfy
\[
 \beta_{s/2}^{\,d}(\xi_\rho)\beta(\xi_\rho/s)
                      =\beta(\xi_\rho)\frac2{1+z_\rho}.
\]
Their extra directional factors give
\(2(U_\alpha-1)(z_i-1)/
 ((z_\alpha^2-1)i\xi_i)\).
Combining these factors proves
\eqref{eq:v28-delta-symbol}, including
the surviving \(1/s\). No incoming
field is reset to zero.

\subsection{A finite exact parity reduction}

For a function of the two phase families,
let \(\mathsf E\) average the sixteen
simultaneous reflections
\[
       (\xi_\rho,\theta_\rho)\mapsto
           (\varepsilon_\rho\xi_\rho,
            \varepsilon_\rho\theta_\rho),\qquad
       \varepsilon_\rho\in\{-1,1\},
\]
where initially \(\theta_\rho\) is
independent of \(\xi_\rho\).
After the algebraic identity is proved,
set \(\theta=\xi/s\).
For \(\mu<\lambda\), define
\begin{equation}\label{eq:v28-T}
 T_{\mu\lambda}(\xi)=
 \cos\xi_\mu(1+\cos\xi_\mu)
       \left(1-\prod_{i<\mu}C(\xi_i)\right)
       \prod_{\substack{i>\mu\\i\ne\lambda}}C(\xi_i).
\end{equation}
In particular \(T_{1\lambda}=0\).
For \(\mu>\lambda\), set \(T_{\mu\lambda}=0\).
Only off-diagonal entries are considered
in this definition.

Each primitive source contraction below
includes the factor \(1/2\) in
\eqref{eq:v18-Gamma-star} and the
two perpendicular Lie-colour contractions.

\begin{lemma}[Exact source parity identities]
\label{lem:v28-parity}
For the covariance bank \(e_ie_i^t\)
in the \(V\) coordinates, the reflected
off-diagonal primitive cubic contribution is
\begin{equation}\label{eq:v28-primitive-parity}
 \mathsf E\,G^{\rm primitive}_{\mu\lambda,i}
   ={\bf1}_{i=\mu}|U_\mu-1|^2
                            T_{\mu\lambda}(\xi)\cos\xi_\lambda .
\end{equation}
For one incoming mixed term,
\begin{equation}\label{eq:v28-mixed-parity}
 \mathsf E\,G^{(s),\rm mixed}_{\mu\lambda,i}
   =-\frac1s{\bf1}_{i=\mu}|U_\mu-1|^2
           T_{\mu\lambda}(\xi)\sin\xi_\lambda
                                      \tan(\theta_\lambda/2).
\end{equation}
For the bank \({\bf1}{\bf1}^t\),
both reflected off-diagonal contributions
are identically zero. These are rational
phase identities, not identities inferred
from sampled frequencies.
\end{lemma}
\begin{proof}
We give the finite algebra and its
coefficientwise certificate. This is
the computer-assisted finite part
of the argument; the scale summation
and positivity proof below are analytic.

In a one-scalar-noise bank, represent
a word by \((c,r,d,e)\): its constant
source coefficient, its complex noise
row, its mixed source/noise row, and
its contracted source/two-noise
coefficient. For real Laurent coefficients,
conjugation inverts every phase.
With \(\langle r,q\rangle=\operatorname{Re}(r\overline q)\),
the degree-three BCH rule is
\begin{equation}\label{eq:v28-jet}
 \begin{aligned}
 c_{ab}&=c_a+c_b,\qquad r_{ab}=r_a+r_b,\\
 d_{ab}&=d_a+d_b+c_ar_b-c_br_a,\\
 e_{ab}&=e_a+e_b
       +2\{\langle r_a,d_b\rangle-\langle r_b,d_a\rangle\}\\
 &\quad+\frac23\{\langle r_a,r_b\rangle(c_a+c_b)
                   -|r_a|^2c_b-|r_b|^2c_a\}.
 \end{aligned}
\end{equation}
This is the same contraction convention
as the exact side-two test in V16.
Apply it to the literal ascending trees,
the two-link paths, the distinguished
relative mean, and the two balancing
factors. In the diagonal bank \(i\),
an input link of orientation \(\alpha\)
at \(x\) has
\[
         c={\bf1}_{\alpha=\lambda},\qquad
         r={\bf1}_{\alpha=i}U^x(U_i-1).
\]
For the rank-one bank, replace its
noise row by \(U^x(U_\alpha-1)\)
for every orientation. Initially
\(d=e=0\). The resulting \(e/2\)
is the primitive polynomial to be
compared with \eqref{eq:v28-primitive-parity}.

For the mixed term it suffices to
use the already derived quadratic
primitive, whose link-pair expression is
\[
 \frac1{32}\sum_s
 \bigl([a_s,b_s]+[\psi_s+\psi'_s,a_s+b_s]
                                      -[\psi_s,\psi'_s]\bigr).
\]
Here \(a_s,b_s\) are the consecutive
links, and \(\psi_s,\psi'_s\) the
centred linear trees at the two ends.
After the two-colour contraction and
the source factor \(1/2\), the
corresponding Fourier wedge coefficient
is \(1/16\).
Substitute \eqref{eq:v28-delta-symbol}
into this finite bilinear expression.

There is a convenient denominator
that makes the check a polynomial one.
Set \(z_\rho=w_\rho^2\) and
\begin{equation}\label{eq:v28-denominator}
 D_0(w)=\prod_{\rho=1}^4
             w_\rho^{-3}(1+w_\rho^2)(w_\rho^4-1).
\end{equation}
It is real on the phase torus,
changes sign under each separate
inversion of \(w_\rho\), and has no
dependence on the outer \(U\)'s.
Multiplication of \(s\Delta^{(s)}\)
by \(D_0\), followed by division
by its displayed \(U_\alpha-1\),
leaves the Laurent polynomial
\[
 16\left(\prod_\rho w_\rho^{-3}\right)
       (z_i-1)L_{\lambda,\alpha i}(z)
                  \prod_{\rho\ne\alpha}(z_\rho^2-1).
\]
Consequently the required reflected
rational identity is equivalent to
equality of the all-coordinate odd
parity parts of two finite Laurent
polynomials. For the primitive
identity it is equality of ordinary
even parity parts.

The exact checker specified below
constructs these polynomials over
\(\mathbb Q\). It collects monomials
under the simultaneous inversions
\((U_\rho,w_\rho)\mapsto
 (U_\rho^{-1},w_\rho^{-1})\),
with the odd sign when \(D_0\)
is present. Every coefficient of
the difference is zero: sixty
primitive identities and sixty
mixed identities, consisting of
twelve off-diagonal entries for
each of the five covariance banks.
Thus no degree bound or statistical
identity test is needed. The
identities hold as Laurent-polynomial
identities, hence as rational
identities off the removable
denominators.

For additional transparency about
the nonzero expression, keep noise
only in direction \(\mu<\lambda\).
The linear tree at corner \(s\) is
\[
 T_s=s_\mu
             e^{i\sum_{i<\mu}s_i\xi_i}Y_\mu,\qquad
 \overline T=\frac12
                 \prod_{i<\mu}\frac{1+U_i}{2}\,Y_\mu .
\]
Summing the balanced-path contractions
gives the factor
\(\cos\xi_\mu(1+\cos\xi_\mu)
 (1-\prod_{i<\mu}C(\xi_i))\);
the unused later corners give the
last product in \eqref{eq:v28-T}.
The weighted source corner gives
\(\cos\xi_\lambda\) in the primitive
term. The inherited row gives
\(-s^{-1}\sin\xi_\lambda\tan(\theta_\lambda/2)\)
in the mixed term. The exact polynomial
checks also establish all the vanishing
banks, which this single-noise explanation
alone would not establish.
\end{proof}

As an independent stencil check, for
\(\lambda>i\) the exact mixed map
acting on noise orientation \(i\)
has no other output orientation and
\[
  L_{\lambda,ii}(z)
    =B_{2,i}(z)\frac{z_\lambda-1}{z_\lambda+1},
       \qquad
  B_{2,i}(z)=\frac1{16}
                  \prod_\rho(1+z_\rho)(1+z_i).
\]
The apparent denominator cancels.
The checker verifies all twenty-four
such output-stencil identities directly.
This explains why the relevant
incoming response has a scalar
telescoping factor, without erasing
the other banks before checking them.

\subsection{Summing the incoming scales exactly}

Truncate the inherited response after
\(J\) incoming scales, but keep the
continuum linear \(Y\) at the current
shell. Put \(p=2^J\), and define
\[
             H_J(x)=\sum_{\ell=1}^J
                       2^{-\ell}\tan(x/2^{\ell+1}).
\]
The double-angle cotangent identity gives
\begin{equation}\label{eq:v28-cot}
 H_J(x)=p^{-1}\cot(x/(2p))-\cot(x/2)
                 \longrightarrow \frac2x-\cot(x/2).
\end{equation}
For example,
\(\tan y=\cot y-2\cot(2y)\);
substitute \(y=x/2^{\ell+1}\) and
sum the telescoping differences.
The final product with the other
filters has the removable values
at their zeros.

Use Lemma~\ref{lem:v28-parity} in
\eqref{eq:v28-covariance}.
The rank-one bank and all diagonal
banks except \(i=\mu\) disappear
from an off-diagonal entry.
The remaining covariance factor is
\[
 \frac{|\mathcal B(\xi)|^2}{|\xi|^4}
   \frac{|\xi|^2}{\xi_\mu^2}|U_\mu-1|^2
         =\frac{|\mathcal B(\xi)|^2S(\xi_\mu)}{|\xi|^2}.
\]
Thus for \(\mu<\lambda\),
\begin{equation}\label{eq:v28-truncated}
 (\Gamma_*^{(J)})_{\mu\lambda}
  =\frac1{(2\pi)^4}\int_{\mathbb R^4}
       \frac{|\mathcal B(\xi)|^2S(\xi_\mu)}{|\xi|^2}
          T_{\mu\lambda}(\xi)
            \{\cos\xi_\lambda-\sin\xi_\lambda H_J(\xi_\lambda)\}
                                                          \,d\xi .
\end{equation}
All reverse-triangular off-diagonal
entries are zero.

The passage \(J\to\infty\) here is
dominated, not a formal resummation.
Indeed
\[
 \left|\frac{\sin x}{p}\cot(x/(2p))\right|\le2 .
\]
Write \(x=2py\). The elementary
integer-multiple bound is
\[
                         |\sin(2py)|\le2p|\sin y|.
\]
The remaining cosine is at most one.
Also
\(\sin x\cot(x/2)=1+\cos x\).
The brace in \eqref{eq:v28-truncated}
is uniformly bounded. Its other
factors are bounded by a constant
times \(|\mathcal B|^2S(\xi_\mu)/|\xi|^2\),
which is integrable: near zero use
four-dimensional local integrability
of \(|\xi|^{-2}\), and at infinity
use that the averaging kernel is
in \(L^2\), hence its squared Fourier
transform is integrable.
V18 identifies the incoming-response
limit with \(\Gamma_*\), so dominated
convergence proves
\begin{equation}\label{eq:v28-off-integral}
 (\Gamma_*)_{\mu\lambda}
 =\frac2{(2\pi)^4}\int_{\mathbb R^4}
       \frac{|\mathcal B(\xi)|^2S(\xi_\mu)}{|\xi|^2}
             T_{\mu\lambda}(\xi)W(\xi_\lambda)\,d\xi ,
       \qquad
 W(x)=\frac12+\cos x-\frac{\sin x}{x}.
\end{equation}
The value \(W(0)=1/2\) is understood
by continuity.

This step is important: replacing
the inherited response by zero
would leave \(\cos x/2\) in place
of \(W(x)\) in the last formula.
The following positivity statement
is for the complete source contraction.

\subsection{Factorization into one-dimensional heat integrals}

For \(t>0\), define normalized
one-dimensional integrals
\begin{equation}\label{eq:v28-heat-defs}
 \begin{aligned}
 A(t)&=\frac1{2\pi}\int_{\mathbb R}e^{-tx^2}S(x)\,dx,\\
 B(t)&=\frac1{2\pi}\int_{\mathbb R}e^{-tx^2}S(x)C(x)\,dx,\\
 D(t)&=\frac1{2\pi}\int_{\mathbb R}
               e^{-tx^2}S(x)^2\cos x(1+\cos x)\,dx,\\
 E(t)&=\frac1{2\pi}\int_{\mathbb R}e^{-tx^2}S(x)W(x)\,dx.
 \end{aligned}
\end{equation}
These letters denote scalar functions,
not the source or reference matrices.
The identity
\[
                         |\xi|^{-2}=\int_0^\infty e^{-t|\xi|^2}\,dt
\]
and \eqref{eq:v28-off-integral} give
\begin{equation}\label{eq:v28-factor}
 (\Gamma_*)_{\mu\lambda}
    =2\int_0^\infty D(t)E(t)
          \{A(t)^{\mu-1}-B(t)^{\mu-1}\}
                        B(t)^{3-\mu}\,dt,\qquad \mu<\lambda .
\end{equation}
Absolute integrability of the original
Fourier integrand, proved above,
justifies this interchange even
before positivity is established.
In particular the result does not
depend on which later direction
\(\lambda>\mu\) was chosen.
The factor for \(\mu=1\) is zero.

Let \(T(y)=(1-|y|)_+\) and let
\[
                 g_t(y)=(4\pi t)^{-1/2}e^{-y^2/(4t)}.
\]
With the Fourier normalization in
\eqref{eq:v28-heat-defs}, the inverse
transform of \(S\) is \(T\).
Set \(F_2(t,y)=(g_t*T)(y)\) and
\(F_4(t,y)=(g_t*T*T)(y)\).
Then
\begin{equation}\label{eq:v28-positive-representations}
 \begin{aligned}
 A(t)&=F_2(t,0),\\
 B(t)&=\tfrac12F_2(t,0)+\tfrac12F_2(t,1),\\
 D(t)&=\tfrac12F_4(t,0)+F_4(t,1)+\tfrac12F_4(t,2),\\
 E(t)&=\int_{\mathbb R}g_t(y)\,e_0(y)\,dy ,
 \end{aligned}
\end{equation}
where
\begin{equation}\label{eq:v28-profile}
 e_0(y)=
 \begin{cases}
   |y|^2/4,&0\le |y|\le1,\\
   |y|(2-|y|)/4,&1\le |y|\le2,\\
   0,&|y|\ge2.
 \end{cases}
\end{equation}

\begin{lemma}[Strict positivity of the heat factors]
\label{lem:v28-positive}
For every \(t>0\),
\begin{equation}\label{eq:v28-positive}
                    0<A(t)/2<B(t)<A(t),\qquad D(t)>0,\quad E(t)>0 .
\end{equation}
\end{lemma}
\begin{proof}
The tent \(T\), its convolution \(T*T\),
and the heat kernel are nonnegative,
with nonzero integrals. The first
three formulas in
\eqref{eq:v28-positive-representations}
give \(A>0\), \(B>A/2\), and \(D>0\).
Moreover
\[
       A(t)-B(t)=\frac1{2\pi}
                   \int_{\mathbb R}e^{-tx^2}S(x)\sin^2(x/2)\,dx>0 .
\]

It remains to verify the less immediate
positive representation for \(E\).
The inverse transform of \(\sin x/x\)
is \(\tfrac12{\bf1}_{[-1,1]}\).
Thus the inverse transform of \(SW\) is
\[
 \frac12T(y)+\frac12T(y-1)+\frac12T(y+1)
                       -\frac12\int_{-1}^1T(y-z)\,dz .
\]
For \(0\le y\le1\), the first three
terms sum to \(1/2\), whereas the
last term is \(1/2-y^2/4\).
For \(1\le y\le2\), their difference
is \((2-y)/2-(2-y)^2/4=y(2-y)/4\).
It vanishes outside \([-2,2]\)
and is even. This proves
\eqref{eq:v28-profile}.
The profile is nonnegative and
positive on a set of positive
measure, and the heat kernel
is everywhere positive. Hence
\(E(t)>0\).
\end{proof}

\begin{theorem}[Strictly positive off-diagonal source coefficients]
\label{thm:v28-source}
The actual matrix has the form
\eqref{eq:v28-structure}, with the
absolutely convergent representations
\begin{equation}\label{eq:v28-ab}
 \begin{aligned}
 a_*&=2\int_0^\infty D(t)E(t)(A(t)-B(t))B(t)\,dt,\\
 b_*&=2\int_0^\infty D(t)E(t)(A(t)-B(t))(A(t)+B(t))\,dt .
 \end{aligned}
\end{equation}
In addition to \(2a_*<b_*<3a_*\),
there is the explicit bound
\begin{equation}\label{eq:v28-lower}
                    a_*\ge
                       \frac{e^{-17/4}}{64\pi^6\sqrt{8\pi}}>0 .
\end{equation}
\end{theorem}
\begin{proof}
The parity identities give the zero
entries. Formula \eqref{eq:v28-factor}
gives the two equal entries in row
two and the one entry in row three,
and yields \eqref{eq:v28-ab}.
All its factors are strictly
positive for \(t>0\).
Furthermore \(A>B>A/2\) gives
\[
                         2B<A+B<3B .
\]
Multiplication by the positive
weight \(2DE(A-B)\) and integration
prove the strict ratio inequalities.

For an explicit lower bound, restrict
the integral for \(a_*\) to \(1\le t\le2\).
The tent and its self-convolution
have mass one and supports
\([-1,1]\), \([-2,2]\), respectively.
The profile \(e_0\) has mass \(1/2\)
and support \([-2,2]\). Consequently
the positive representations give
\[
 B(t)\ge\frac{e^{-1/4}}{2\sqrt{8\pi}},
       \qquad
 D(t)\ge\frac{e^{-1}}{2\sqrt{8\pi}},
       \qquad
 E(t)\ge\frac{e^{-1}}{2\sqrt{8\pi}} .
\]
For \(A-B\), restrict its Fourier
integral to \(1/2\le|x|\le1\).
On this set \(e^{-tx^2}\ge e^{-2}\),
\(S(x)\ge4/\pi^2\), and
\(\sin^2(x/2)\ge1/(4\pi^2)\),
using the sine chord bound.
The two intervals have total
length one. Therefore
\[
                         A(t)-B(t)\ge\frac{e^{-2}}{2\pi^5}.
\]
Multiplying these four lower bounds,
the factor two in \(a_*\), and the
unit length of the \(t\)-interval
gives \eqref{eq:v28-lower}.
No numerical estimate of an integral
enters this bound.
\end{proof}

\subsection{The V27 condition is now satisfied by the actual source}

For row two, the off-diagonal squared
row norm of \(\Gamma_*\) is \(2a_*^2\);
for row three it is \(b_*^2\).
Thus the conditional statements in
V26--V27 become unconditional
statements for these two rows.
For example,
\begin{equation}\label{eq:v28-actual-error}
 \begin{aligned}
 {\rm err}_{p,M}(R_{L,M,2}(r))
      &\ge j\sqrt{2\kappa_{p,M}}\,a_*-C,\\
 {\rm err}_{p,M}([\Pi_M^{\rm c}R_{L,M}]_2(r))
      &\ge j\sqrt{2\widetilde\kappa_{p,M}}\,a_*-C .
 \end{aligned}
\end{equation}
These are immediate substitutions in
Proposition~\ref{prop:v27-cubic}.
They show that the unrenormalized
local cubic coefficient has genuine
logarithmic growth invisible to any
function of the full retained linear
Gaussian observation. The projected
coefficient has the same obstruction.
The side-one average still removes
the finite filter exactly, as in V26,
so this does not contradict V25.

In the joint regime of
\eqref{eq:v27-joint-error}, with
fixed \(M\) and \(j/\beta_L\to\tau>0\),
the projected normalized three-jet
satisfies the strictly positive bound
\begin{equation}\label{eq:v28-joint}
 \liminf_{L\to\infty}{\rm err}_{p,M}
   \bigl([\Pi_M^{\rm c}\sqrt{\beta_L}\,
                           T_{L,M}^{[3]}(\beta_L)]_2(r)\bigr)
                         \ge\frac{8\tau a_*}{\pi^4M^2}>0 .
\end{equation}
This follows from the explicit limit
of \(\widetilde\kappa_{p,M}\), not
from an interchange of conditional
projection limits.

It is therefore no longer consistent
to close this particular logarithmic
source coefficient by declaring it a
deterministic function of \(Y\).
One must retain the additional
filter information, or integrate
its actual conditional fluctuations.
V27 gives a sufficient linear
enlargement and the exact Gaussian
conditional characteristic function.
Neither option by itself controls
the complete nonlinear root.

\subsection{Certificate, nonduplication, and preservation}

The exact certificate is
\begin{center}\small
\path{scripts/verify_ym_f14_v28_source_parity.py}
\end{center}
with SHA--256
\begin{center}\scriptsize\ttfamily
4E611654E9656BA9B9826FB281F7FF4539D4B2D69A284F90C11F63D1F230FD8D.
\end{center}
It proves the sixty primitive and
sixty mixed Laurent identities by
exact rational coefficient collection,
and checks twenty-four independent
one-direction stencil identities.
It additionally checks twenty-five
rational values of the heat profile
and its exact mass \(1/2\).
The profile's all-point formula and
positivity are proved analytically
above; finite value checks are not
being used as its positivity proof.
The exploratory script
\begin{center}\small
\path{scripts/explore_ym_f14_v28_source_symbol.py}
\end{center}
was used to locate the reduction.
Its floating-point quadratures carry
no certified integration or tail
error and are not used in any
theorem of this section.

V18 already defined \(\Gamma_*\)
by convergent four-dimensional
contractions, and V22 gave its
endpoint finite parts. The new
result is its exact off-diagonal
structure and strict positivity,
including the full mixed response.
It removes the specific unverified
condition that V26 and V27
explicitly left open.

The predecessor V27 has 618,196 bytes and SHA--256
\begin{center}\scriptsize\ttfamily
CD137CCBE157B43C5855F62903A9F043C3A647A7E9A698B0E6B9F7B0A65A472F.
\end{center}
Its entire body is preserved,
apart from the title version
and this terminal insertion.
V28 replaces V27 as the sole
active main note. The companion
files, archives and monographs
are unchanged. The next scheduled
SM/NS review is V30.

\subsection{The remaining measured and nonlinear bottleneck}

The source contraction is now
non-diagonal by proof, not by
a finite-lattice extrapolation.
The next upstream task is to
carry this confirmed term through
the actual measured conditional
law and the stable nonlinear
remainder problem. In particular,
the reference contribution \(h_*\)
and the complete combination
\(\Gamma_*-2h_*^T\) must still
be kept distinct from the
source coefficient proved here.
No cancellation or noncancellation
in that complete measured matrix
is inferred from
\eqref{eq:v28-structure}.

\begin{firewall}
Strict positivity of \(a_*,b_*\)
is positivity of source-mixing
coefficients. It is not a
Hamiltonian spectral gap.
This note does not construct
an interacting continuum state,
prove the full nonlinear
quantum remainder, establish
noncollapse, or verify physical
reflection-positive transfer
and a positive Yang--Mills
mass gap. The full objective
remains unproved and unchanged.
\end{firewall}
\section{V29: exact cancellation of the off-diagonal measured logarithm}
\label{sec:v29-reference}

\subsection{The source and the measured correction are different}

V28 proves that the actual source matrix
\(\Gamma_*\) is non-diagonal. It explicitly
does not settle the complete measured
combination \(\Gamma_*-2h_*^T\).
The present calculation supplies the
missing reference term, with the sign
and transpose of the original ledger.
The result is an exact cancellation:
\begin{equation}\label{eq:v29-main}
             \Gamma_*-2h_*^T
                 =\operatorname{diag}(c_1,c_2,c_3,c_4).
\end{equation}
The four diagonal constants are not
evaluated or proved equal in this note.

The source drift and conditional
variance bounds from V28 still
hold under that free reference law.
Here a second, separately owned
term cancels the off-diagonal
coefficient in the measured
relative first trace. The
distinction is essential for
continuing the nonlinear program.

The files were rechecked before
the calculation: Shell-Mixing V16,
SM V72 and NS V26 are unchanged.
The next scheduled companion
checkpoint is V30.

\subsection{Fixing the reference normalization before simplifying}

Keep the positive plane-wave
convention, \(\mathcal B,S,C,U\),
and the transformed coordinates
\(V_i\) of V28. Let
\(T_\mu^{\rm loc}\) be the real
skew bank defined in V18 by
\[
             E_\mu^{[2]}(A)
                =\sum_{a<b}(T_\mu^{\rm loc})_{ab}[A_a,A_b].
\]
For \(a=(x_a,\alpha)\), \(b=(x_b,\beta)\),
define its orientation Fourier bank
\begin{equation}\label{eq:v29-K}
 K_{\mu,\alpha\beta}(U)
     =2\sum_{\substack{a:\operatorname{or}(a)=\alpha\\
                       b:\operatorname{or}(b)=\beta}}
                (T_\mu^{\rm loc})_{ab}U^{x_a-x_b}.
\end{equation}
This is exactly the quadratic bank
used in V28's finite certificate:
the link-pair coefficient \(1/32\)
becomes the Fourier wedge coefficient
\(1/16\). Put
\[
                     F_\alpha(\xi)=\mathcal B(\xi)\beta(\xi_\alpha).
\]
The definition \eqref{eq:v18-h-star},
with response multiplier
\(-4i\xi_\nu\Pi(\xi)/|\xi|^4\), gives
\begin{equation}\label{eq:v29-reference-symbol}
 (h_*)_{\nu\mu}
   =\frac1{(2\pi)^4}\int_{\mathbb R^4}
       \operatorname{Re}\left\{
          \frac{i\xi_\nu}{|\xi|^4}
             \sum_{\alpha,\beta}
                K_{\mu,\alpha\beta}(U)
                   F_\alpha(\xi)\Pi_{\alpha\beta}(\xi)
                                      \overline{F_\beta(\xi)}
                           \right\}\,d\xi .
\end{equation}

To check the sign, \(T_{ba}=-T_{ab}\)
in the original expression
\(\tfrac12\sum_{a,b}\mathcal R_\nu(f_a,f_b)T_{ba}\).
Its minus sign combines with
\(-4i\xi_\nu\); the factor two
in \eqref{eq:v29-K} then leaves
the displayed \(+i\xi_\nu\).
In particular this is not the
formula obtained by omitting
the transpose of the skew bank.
The independent finite calibration
below also checks this normalization.

\subsection{The exact reference parity identities}

Write \(R_\alpha=U_\alpha-1\).
For the four diagonal unit
covariance banks in the \(V\)
coordinates, define
\[
       J_{\mu,i}(U)
                 =K_{\mu,ii}(U)R_i\overline{R_i},
                                      \qquad 1\le i\le4.
\]
For the common rank-one bank define
\[
       J_{\mu,0}(U)=
                    \sum_{\alpha,\beta}
                           K_{\mu,\alpha\beta}(U)
                                    R_\alpha\overline{R_\beta}.
\]
All these Laurent polynomials are
purely imaginary on the phase torus.
Let \(\mathsf O_\nu\) denote projection
onto the part odd in \(\xi_\nu\)
and even in the other three
coordinates. This is the parity
selected by the extra factor
\(i\xi_\nu\) in
\eqref{eq:v29-reference-symbol}.
An ordinary even projection of
the bank alone would incorrectly
discard the response.

\begin{lemma}[Reference parity]
\label{lem:v29-parity}
For every off-diagonal pair \(\mu\ne\nu\),
\begin{equation}\label{eq:v29-parity}
 \begin{aligned}
 \mathsf O_\nu J_{\mu,i}
     &=-i\,{\bf1}_{i=\mu}|U_\mu-1|^2
                             T_{\mu\nu}(\xi)\sin\xi_\nu,\\
 \mathsf O_\nu J_{\mu,0}&=0 .
 \end{aligned}
\end{equation}
Here \(T_{\mu\nu}\) is the function
in \eqref{eq:v28-T}, zero for
\(\mu>\nu\) and for \(\mu=1\).
\end{lemma}
\begin{proof}
This is a finite Laurent identity
using the same exact local quadratic
bank as V28. Its target polynomial
in the nonzero case is
\[
   -\tfrac12(U_\nu-U_\nu^{-1})
                          |U_\mu-1|^2T_{\mu\nu}.
\]
The checker identified below builds
\eqref{eq:v29-K}, multiplies the
displayed \(R\) factors, and collects
coefficients under phase inversion,
with odd sign only in coordinate
\(\nu\). All sixty differences vanish
coefficientwise: twelve off-diagonal
entries for each of the five banks.
This is exact rational arithmetic
on polynomials, not testing a finite
set of real frequencies.

One can also see the surviving term
directly. With noise only in direction
\(\mu<\nu\), use the tree rows
\(l_s=T_s-\overline T\) from V28,
their translates \(r_s=U_\mu^2l_s\),
and the two-link row \(x_s=(1+U_\mu)U^s\).
The quadratic bank contains
\[
 \frac1{16}\sum_s
 \{a_s\overline b_s-b_s\overline a_s
    +(l_s+r_s)\overline x_s
       -x_s\overline{(l_s+r_s)}
    -l_s\overline r_s+r_s\overline l_s\}.
\]
The first and last terms have
the wrong parity in direction
\(\mu\) for a response with
\(\nu>\mu\). The middle term
is minus twice the imaginary
part of the balanced corner
average computed in V28.
Its \(\nu\)-odd part is
\(-iT_{\mu\nu}\sin\xi_\nu\).
The outside noise factor is
\(|U_\mu-1|^2\), giving
\eqref{eq:v29-parity}.
The coefficient checks additionally
prove the vanishing of the other
diagonal banks and the rank-one
bank; those vanishings are not
assumed from this special case.
\end{proof}

Substitute \eqref{eq:v28-covariance}
into \eqref{eq:v29-reference-symbol}.
Equivalently, its integrand is
\[
 \operatorname{Re}\left\{
 \frac{i\xi_\nu|\mathcal B(\xi)|^2}{|\xi|^6}
   \left(\sum_i\frac{|\xi|^2}{\xi_i^2}J_{\mu,i}
                                                -J_{\mu,0}\right)\right\}.
\]
Lemma~\ref{lem:v29-parity} leaves
only the \(i=\mu\) diagonal bank.
Using
\(|U_\mu-1|^2/\xi_\mu^2=S(\xi_\mu)\),
we obtain, for \(\mu<\nu\),
\begin{equation}\label{eq:v29-off-reference}
 (h_*)_{\nu\mu}
   =\frac1{(2\pi)^4}\int_{\mathbb R^4}
       \frac{|\mathcal B(\xi)|^2S(\xi_\mu)}{|\xi|^4}
             T_{\mu\nu}(\xi)\,\xi_\nu\sin\xi_\nu\,d\xi .
\end{equation}
All off-diagonal entries with
\(\mu>\nu\) vanish.
The integral is absolutely
convergent: near zero use
\(|\xi_\nu\sin\xi_\nu|\le|\xi|^2\),
and at infinity use
\(|\xi_\nu\sin\xi_\nu|\le|\xi|\)
and the integrability of the
squared averaging multiplier.
The remaining trigonometric
factors are bounded.

\subsection{The heat identity that enforces cancellation}

Retain \(A(t),B(t),D(t),E(t)\)
from \eqref{eq:v28-heat-defs},
and define one additional
one-dimensional factor
\[
       F(t)=\frac1{2\pi}\int_{\mathbb R}
                        e^{-tx^2}S(x)\,x\sin x\,dx,\qquad t>0 .
\]

\begin{lemma}[Source and response heat factors]
\label{lem:v29-heat}
For every \(t>0\),
\begin{equation}\label{eq:v29-heat}
                              E(t)=tF(t).
\end{equation}
In particular \(F(t)>0\).
\end{lemma}
\begin{proof}
Since \(S(x)=\{\sin(x/2)/(x/2)\}^2\),
direct differentiation gives
\begin{equation}\label{eq:v29-derivative}
    S'(x)\sin x+S(x)\cos x
                  =S(x)\left(1+2\cos x-\frac{2\sin x}{x}\right)
                  =2S(x)W(x).
\end{equation}
The formula has its continuous
values at zero and at the
zeros of \(S\). For example,
away from those zeros it follows
from \(S'/S=\cot(x/2)-2/x\)
and \(\sin x\cot(x/2)=1+\cos x\);
multiplication by the smooth
factors removes the apparent poles.

Integrate the derivative of
\(e^{-tx^2}S(x)\sin x\) over
the entire real line. Its
boundary values vanish, and
its derivative is integrable
because \(t>0\).
Equation \eqref{eq:v29-derivative}
therefore gives \(2E(t)-2tF(t)=0\).
Strict positivity follows from
Lemma~\ref{lem:v28-positive}.
\end{proof}

Using
\[
                     |\xi|^{-4}=\int_0^\infty t e^{-t|\xi|^2}\,dt
\]
in \eqref{eq:v29-off-reference}
and separating the four coordinates
now yields
\begin{equation}\label{eq:v29-heat-reference}
 \begin{aligned}
 (h_*)_{\nu\mu}
 &=\int_0^\infty tD(t)F(t)
       \{A(t)^{\mu-1}-B(t)^{\mu-1}\}B(t)^{3-\mu}\,dt\\
 &=\int_0^\infty D(t)E(t)
       \{A(t)^{\mu-1}-B(t)^{\mu-1}\}B(t)^{3-\mu}\,dt,
                                      \qquad \mu<\nu .
 \end{aligned}
\end{equation}
Absolute convergence of
\eqref{eq:v29-off-reference}
justifies the first interchange;
the second line is the pointwise
heat identity, not an asymptotic
comparison of two integrals.

\begin{theorem}[The complete off-diagonal reference matrix]
\label{thm:v29-reference}
With \(a_*,b_*\) as in V28,
\begin{equation}\label{eq:v29-h-matrix}
 h_*=
 \begin{pmatrix}
 \eta_1&0&0&0\\
 0&\eta_2&0&0\\
 0&a_*/2&\eta_3&0\\
 0&a_*/2&b_*/2&\eta_4
 \end{pmatrix}.
\end{equation}
Consequently \eqref{eq:v29-main}
holds with
\[
                           c_\mu=\gamma_\mu-2\eta_\mu .
\]
The three nonzero off-diagonal
entries of \(h_*\) are strictly
positive. In particular,
\[
                (h_*)_{32}=(h_*)_{42}
                   \ge\frac{e^{-17/4}}{128\pi^6\sqrt{8\pi}} .
\]
\end{theorem}
\begin{proof}
Compare \eqref{eq:v29-heat-reference}
with \eqref{eq:v28-factor}.
For every off-diagonal pair,
\[
                           2(h_*)_{\nu\mu}=(\Gamma_*)_{\mu\nu}.
\]
The zero entries and the two
equal entries in
\eqref{eq:v29-h-matrix}
follow from V28 and the reference
parity lemma. Its positive entries
and lower bound follow from
Theorem~\ref{thm:v28-source}.
Taking the transpose in the
original measured combination
then cancels every off-diagonal
entry, while leaving the displayed
four diagonal constants.
\end{proof}

As a direct Fourier check of the
same cancellation, write the
common product of all factors
except those in coordinate \(\nu\)
in \eqref{eq:v28-off-integral}
and \eqref{eq:v29-off-reference}
as \(P(\xi_{\widehat\nu})\).
Their difference has integrand
\begin{equation}\label{eq:v29-total-derivative}
 P(\xi_{\widehat\nu})
       \frac{\partial}{\partial\xi_\nu}
             \left(\frac{S(\xi_\nu)\sin\xi_\nu}{|\xi|^2}\right).
\end{equation}
For every fixed nonzero transverse
frequency vector \(\xi_{\widehat\nu}\),
the one-dimensional boundary values
at infinity vanish. The exceptional
zero transverse vector has measure
zero. Absolute integrability of the
two original terms, proved in V28
and above, permits the iterated
integration. Thus this derivative
also proves that their difference
integrates to zero. It identifies
the cancellation as an integration
by parts in the same physical
Fourier variable, not a resetting
of either contribution.

\subsection{Finite parts are not erased by the leading cancellation}

To distinguish the orientation matrix
from the free scalar-link covariance,
write
\[
       \mathsf C_L=\Gamma_L-2h_L^T,\qquad
       \mathsf D_*=\operatorname{diag}(c_1,c_2,c_3,c_4),\qquad
       \mathsf C_{\rm fin}=\Gamma_{\rm fin}-2h_{\rm fin}^T.
\]
Theorem~\ref{thm:v22-relative} and
Theorem~\ref{thm:v29-reference}
give the full matrix limit
\begin{equation}\label{eq:v29-finite-part}
                   \mathsf C_{2^n}-n\mathsf D_*
                                  \longrightarrow\mathsf C_{\rm fin}.
\end{equation}
In particular, each off-diagonal
entry of \(\mathsf C_{2^n}\)
is bounded and converges.
Its limit need not be zero.
Both the coarse-endpoint and
fine-endpoint series in V22
remain in \(\mathsf C_{\rm fin}\).

This is stronger than the statement
that off-diagonal entries divided
by \(n\) tend to zero, but weaker
than an exact finite-cutoff
cancellation. The proof is simply
the difference of the two already
proved finite-part limits, using
the new identification of their
leading difference. No estimate
for the endpoints is discarded.

\subsection{The complete measured quartic logarithm after this reduction}

Recall V19's exact identification
of the linear source countertransport,
and the quartic identity
\[
               q_{*,4}^{\rm rel}(v)
                          =-g(v)\cdot\mathsf D_*v .
\]
For
\[
        P_{\mu\nu}(v)=|v_\mu|^2|v_\nu|^2-(v_\mu\cdot v_\nu)^2,
                                                    \qquad \mu<\nu,
\]
the inherited normalization is
\[
       E_4(v)=2\sum_{\mu<\nu}P_{\mu\nu}(v),\qquad
       g_\mu(v)\cdot v_\mu
                              =4\sum_{\nu\ne\mu}P_{\mu\nu}(v).
\]
Here the symmetric notation for
\(P_{\mu\nu}\) is understood
when \(\nu<\mu\).
Hence
\begin{equation}\label{eq:v29-relative-quartic}
                    q_{*,4}^{\rm rel}(v)
                         =-4\sum_{\mu<\nu}(c_\mu+c_\nu)P_{\mu\nu}(v).
\end{equation}
The mixed contractions
\(g_\mu(v)\cdot v_\nu\), \(\mu\ne\nu\),
have canceled in the actual
orientation matrix, not only
after taking its symmetric part.

Substitution into
\eqref{eq:v21-full-leading}
gives the present fully justified
form of the complete measured
fourth-jet logarithm:
\begin{equation}\label{eq:v29-full-quartic}
 \frac{(\mathcal Q_{2^n}^{\rm H})_{[4]}}n
       \longrightarrow
       -\frac{2\log2}{3\pi^2}E_4
          -4\sum_{\mu<\nu}(c_\mu+c_\nu)P_{\mu\nu}.
\end{equation}
The convergence and its finite
part retain the original
balanced-root and coarse-Haar
conventions. The fine-Haar,
ghost, normal-source and tangent
terms remain those already
matched in V20--V22.

This is not yet an evaluation
of the four constants \(c_\mu\).
If they were proved equal to
a common \(c\), the last term
would be \(-4cE_4\).
Equality, the common value,
and identification with a
physical continuum running
coupling are not asserted
by this conditional simplification.

\subsection{Why the source-information result remains relevant}

The two exact statements now coexist:
\[
       \Gamma_*\ \hbox{is non-diagonal},\qquad
       \Gamma_*-2h_*^T\ \hbox{is diagonal}.
\]
The first controls the local cubic
observable under the free Gaussian
reference. The second controls the
linearization of the measured
relative source correction on the
specified minimizing branch.
Neither permits replacing
\(\Gamma_*\) by \(\mathsf D_*\)
inside the polynomial identity
\eqref{eq:v26-Rhat}.

In particular, the positive
conditional-variance bounds of
V27--V28 have not disappeared.
The reference response is not a
new deterministic value of the
microscopic cross-oriented filter.
It is a separately owned variation
of the precision in the measured
first trace. Conversely, the
source-only non-diagonality cannot
be cited as evidence for a
non-diagonal measured logarithm
after the present cancellation.

The complete interacting conditional
law may reorganize these terms,
but such a reorganization must be
proved with the actual incoming
density and observables. Gaussian
conditional variances cannot be
transferred to that law without
the needed measure comparison.
The sixteen-channel enlargement
of V27 remains one sufficient
linear description of the source
information, not a claim that the
full quantum law has been closed
on those channels.

\subsection{Exact checks and preservation record}

The certificate is
\begin{center}\small
\path{scripts/verify_ym_f14_v29_reference_cancellation.py}
\end{center}
with SHA--256
\begin{center}\scriptsize\ttfamily
4C760502E3B20F79FE7CA1AC540667714847C1060F12FBAD1654295B4062907B.
\end{center}
It proves sixty reference parity
identities by exact rational
Laurent coefficient collection,
and verifies the trigonometric
algebra in
\eqref{eq:v29-derivative}.
As an independent normalization
test, it evaluates the original
side-four fine-shell reference
formula in the field of 101
elements with \(i=10\).
The \((1,1)\) entry is 64,
matching V18's independently
certified fine-link trace.
This fixes the sign and scale;
a finite modular trace is not
used to infer a continuum limit.
The integration-by-parts boundary
estimates are proved above.

The exploratory file
\begin{center}\small
\path{scripts/explore_ym_f14_v29_reference_symbol.py}
\end{center}
contains unvalidated floating-point
quadrature for direction finding.
None of its values is used as
a sign, equality, or integration
error certificate.

The predecessor V28 has 641,728 bytes and SHA--256
\begin{center}\scriptsize\ttfamily
9FFDB431FB1B946C82F944990400F89507C541CDB0AFA872DE3F969D4430C406.
\end{center}
Its complete body is retained,
apart from the title version
and this terminal insertion.
V29 replaces V28 as the sole
active main note. The companion
files, archives and monographs
are unchanged. The next scheduled
SM/NS review is V30.

\subsection{The next upstream target}

The reference term has now been
matched exactly to the nonzero
off-diagonal source term. The
remaining leading measured
matrix consists of four diagonal
constants. Determining their
structure and value is the
next concrete normalization
problem. Beyond it, the stable
nonlinear remainder, conditional
density propagation and
interacting continuum construction
remain necessary; a one-loop
matrix identity does not replace
any of them.

\begin{firewall}
The cancellation proved here
is a statement about the
leading measured Taylor
coefficient. It does not
construct an interacting
continuum state, prove a
uniform full nonlinear
quantum remainder, establish
noncollapse or physical
reflection-positive transfer,
or give a positive
Hamiltonian mass gap.
The full Yang--Mills objective
remains unproved and unchanged.
\end{firewall}
\section{V30: exact evaluation of the complete measured quartic logarithm}
\label{sec:v30}

\subsection{Context and the scheduled companion checkpoint}

V29 canceled the off-diagonal entries of the actual measured
matrix, leaving four unevaluated numbers. This section evaluates
them, retaining both the inherited mixed source and the response
of the Gaussian precision. In the conventions of this manuscript,
\begin{equation}\label{eq:v30-main}
 \mathsf C_*:=\Gamma_*-2h_*^T
       =\frac{3\log2}{4\pi^2}I_4.
\end{equation}
This is a statement about the leading measured Taylor coefficient
of the specified \(\SU(2)\) blocking construction. It neither
assumes permutation symmetry of the ascending tree nor identifies
a full interacting renormalization map with its Taylor jet.

The fifth-version checkpoint has been performed on the current
SM V72 and NS V26, and the Shell Mixing V16 has been revisited.
The SM terminal results concern a fixed physical momentum set
and fixed-time many-copy covariance transport; they do not give
a growing-cutoff bound for the present nonlinear quantum integral.
The NS terminal pointwise Oseen formulas likewise do not supply
such a bound. Numerical comparisons assigned there to a future
V27 grid are not certificates in the active NS V26 and are not
imported here. Shell Mixing V16 supplies the already incorporated
measured ownership ledger, while its flat-holonomy calculation is
not a substitute for a noncommuting quartic coefficient.
Thus the direction is unchanged by this checkpoint: finish the
actual noncommuting measured coefficient, then move upstream to
control of the full density and its remainder. The next scheduled
SM/NS checkpoint is V35. Source-document suggestions remain
optional research material, not instructions.

\subsection{Five covariance contractions and an exact mixed decomposition}

Retain \(S,C,U,R,K\) from V28--V29 and set
\[
 \chi(\xi)=\prod_{\rho=1}^4S(\xi_\rho),\qquad
 P(\xi)=\prod_{\rho=1}^4C(\xi_\rho),\qquad r=|\xi|.
\]
The elementary double-angle identity gives
\begin{equation}\label{eq:v30-box-scale}
                    \chi(\xi)P(\xi)=\chi(2\xi).
\end{equation}
Use the symbol \(\circ\) for the common rank-one bank, called
bank \(0\) in V29, to distinguish it from the four orientations.
The scalar weights in \eqref{eq:v28-covariance} are
\begin{equation}\label{eq:v30-weights}
 w_i(\xi)=\frac{\chi(\xi)}{r^2\xi_i^2},\qquad
 w_\circ(\xi)=\frac{\chi(\xi)}{r^4}.
\end{equation}
The first four banks have positive signs; the common bank is
subtracted. This remains an algebraic splitting, not five
independent positive Gaussian covariances.

For an incoming distance \(s=2^\ell\), put
\[
 z_i=e^{i\xi_i/s},\qquad
 \tau_i=\frac{z_i-1}{z_i+1},\qquad
 P_{>i}(z)=\prod_{\rho>i}\frac{1+z_\rho}{2}.
\]
Empty products are one. The following identity resolves all
output orientations, not only the off-diagonal source entries.

\begin{lemma}[Diagonal and gradient parts of the actual mixed stencil]
\label{lem:v30-stencil}
The complete symbol \eqref{eq:v28-delta-symbol} satisfies
\begin{equation}\label{eq:v30-stencil}
 \Delta^{(s)}_{\lambda,\alpha i}
   =\frac{R_\alpha}{s}
       \left((1+\delta_{i\lambda})\delta_{\alpha i}\tau_\lambda
                    +b_{\lambda i}(z)\right),
\end{equation}
where
\begin{equation}\label{eq:v30-b}
 b_{\lambda i}(z)=
 \begin{cases}
 0,&i<\lambda,\\
 -\tau_\lambda/(2P_{>\lambda}),&i=\lambda,\\
 -\tau_\lambda\tau_i/(2P_{>i}),&i>\lambda.
 \end{cases}
 \qquad
 \sum_i b_{\lambda i}=-\frac12\tau_\lambda.
\end{equation}
The corresponding quadratic bank obeys the full, unprojected
Laurent identity
\begin{equation}\label{eq:v30-ward}
 \sum_\alpha K_{\mu,i\alpha}(U)R_i\overline{R_\alpha}
       =\delta_{i\mu}\,4i\sin\xi_\mu
                         \{\cos\xi_\mu P(\xi)-1\}.
\end{equation}
\end{lemma}
\begin{proof}
Insert the literal V17 mixed row into
\eqref{eq:v28-delta-symbol}. After canceling its common
\(R_\alpha/s\), multiplication by
\((z_\alpha^2-1)\prod_\rho(1+z_\rho)\) makes the claimed
identity a finite polynomial identity. Expansion of the ascending
tree and its balancing factors gives respectively the three
cases in \eqref{eq:v30-b}; the exact certificate below checks
all \(4^3=64\) triples against the original mixed row.
The final sum in \eqref{eq:v30-b} also follows directly from
\[
 \frac{\tau_i}{P_{>i}}
       =\frac1{P_{>i}}-\frac1{P_{\ge i}},\qquad
 P_{\ge i}=\frac{1+z_i}{2}P_{>i}.
\]
Summing this difference over \(i>\lambda\) gives
\(1-P_{>\lambda}^{-1}\).

For \eqref{eq:v30-ward}, substitute the Fourier wedge expression
with coefficient \(1/16\) from the proof of
Lemma~\ref{lem:v28-parity} into the left side. Contract its
second row with \(\overline{R_\alpha}\) and sum the sixteen
corners. Terms with \(i\ne\mu\) cancel; the remaining expression
is
\[
           2(U_\mu-U_\mu^{-1})
                     \left(\frac{U_\mu+U_\mu^{-1}}2P-1\right).
\]
This is the right side of \eqref{eq:v30-ward}. The certificate
also expands all sixteen identities before any parity projection.
These are exact rational-coefficient computations, not tests
at selected frequencies. The rational identities extend across
their removable denominators after multiplication by the
original box filters, as in V28.
\end{proof}

The second term of \eqref{eq:v30-stencil} is an output-gradient
row \(R_\alpha\) times a common scalar. Equation
\eqref{eq:v30-ward} shows explicitly why it must not be thrown
away: its contraction is nonzero for \(i=\mu\). Calling it
``gradient'' does not authorize replacing the actual quadratic
bank by a projection that annihilates it.

\subsection{Exact diagonal parity and the incoming-scale sum}

Fix \(\mu\). Let \(p_{\mu i}\), including \(i=\circ\), be
the even-reflected primitive cubic contraction for external
source and output both \(\mu\). Precisely, use the recurrence
\eqref{eq:v28-jet} on the literal balanced paths and take the
even part of \(e/2\), with the same source and colour factors
as V28. Define the real even trigonometric polynomials
\(q_{\mu i}\) by
\begin{equation}\label{eq:v30-q}
 \begin{aligned}
 \mathsf O_\mu\{K_{\mu,ii}|R_i|^2\}
                  &=i\sin\xi_\mu\,q_{\mu i},\qquad 1\le i\le4,\\
 q_{\mu\circ}&=4\{\cos\xi_\mu P-1\}.
 \end{aligned}
\end{equation}
Every odd trigonometric polynomial in \(\xi_\mu\) is divisible
by \(\sin\xi_\mu\) as a polynomial in sines and cosines, so
the first line defines a polynomial without poles. The second
line follows by summing \eqref{eq:v30-ward} over \(i\).

\begin{lemma}[The full diagonal incoming contraction]
\label{lem:v30-mixed}
After even reflection, one incoming scale contributes
\begin{equation}\label{eq:v30-mixed}
 \mathsf E G_{\mu\mu,i}^{(s),\mathrm{mixed}}
    =\frac1s\sin\xi_\mu\tan\left(\frac{\xi_\mu}{2s}\right)
                                                    a_{\mu i},
\end{equation}
where
\begin{equation}\label{eq:v30-a}
 \begin{aligned}
 a_{\mu i}&=(1+\delta_{i\mu})q_{\mu i}
                         -\frac12\delta_{i\mu}q_{\mu\circ},\\
 a_{\mu\circ}&=\frac32q_{\mu\circ}.
 \end{aligned}
\end{equation}
Consequently the summed incoming contribution is
\(\ell(\xi_\mu)a_{\mu i}\), where
\begin{equation}\label{eq:v30-ell}
             \ell(x)=\frac{2\sin x}{x}-1-\cos x.
\end{equation}
\end{lemma}
\begin{proof}
For a diagonal bank, the diagonal part of
\eqref{eq:v30-stencil} pairs with \(J_{\mu,i}\).
Since \(\overline{\tau_\mu}=-i\tan(\xi_\mu/(2s))\),
its real even part is the first term of \eqref{eq:v30-a}
times the factor in \eqref{eq:v30-mixed}.
The gradient part is zero for \(i\ne\mu\) by
\eqref{eq:v30-ward}. For \(i=\mu\), its later-coordinate
factors have even parts one, since
\[
                  \frac2{1+e^{i\theta}}=1-i\tan(\theta/2).
\]
The Ward expression is even in every later coordinate. Thus
only the constant part of \(P_{>\mu}^{-1}\) survives, giving
the term \(-q_{\mu\circ}/2\) in \eqref{eq:v30-a}.

For the common bank, summing the diagonal term over input
orientations gives twice the Ward expression: its columns
vanish except at \(\mu\), by skew conjugacy and
\eqref{eq:v30-ward}. The sum of the gradient coefficients is
\(-\tau_\mu/2\) by \eqref{eq:v30-b}. The resulting coefficient
is therefore \(2-1/2=3/2\), as claimed. The exact checker
additionally verifies all twenty reflected mixed identities
directly from the unreduced original stencil, clearing
\(D_0\) from \eqref{eq:v28-denominator} before collecting the
all-coordinate odd coefficients.

Now sum \eqref{eq:v30-mixed} over \(s=2,4,\ldots,2^J\).
The cotangent telescoping identity \eqref{eq:v28-cot} gives
\[
 \sin x\,H_J(x)\longrightarrow
       \sin x\left(\frac2x-\cot(x/2)\right)=\ell(x).
\]
For every dyadic \(p=2^J\),
\(p^{-1}\sin x\cot(x/(2p))\) has absolute value at most two,
because \(|\sin(2py)|\le2p|\sin y|\).
Also \(\sin x\cot(x/2)=1+\cos x\). Hence
\(|\sin x H_J(x)|\le4\), with the removable values understood.
Near zero the stronger bound \(C x^2\) is uniform in \(J\),
by the convergent geometric sum in its defining tangent series.
The polynomial zeros and box decay used below then justify
dominated convergence for each signed bank. No last incoming
scale has been discarded.
\end{proof}

Write \(x=\xi_\mu\), suppress \(\mu\) in \(p_i,q_i,a_i\),
and put \(c=\cos x\). From the reference normalization
\eqref{eq:v29-reference-symbol}, the diagonal-bank contribution
to \((h_*)_{\mu\mu}\) is
\(-w_i x\sin x\,q_i/r^2\), and the common bank has the
opposite sign. Therefore
\begin{equation}\label{eq:v30-diagonal-integral}
 \begin{aligned}
 c_\mu=\frac1{(2\pi)^4}\int_{\mathbb R^4}
 \biggl[&\sum_{i=1}^4w_i
       \left(p_i+\ell(x)a_i+\frac{2x\sin x}{r^2}q_i\right)\\
 &-w_\circ\left(p_\circ+\frac32\ell(x)q_\circ
                         +\frac{2x\sin x}{r^2}q_\circ\right)
                                      \biggr],d\xi.
 \end{aligned}
\end{equation}
In particular the sign of the reference term has not been
chosen to fit a target coefficient: it is the sign fixed by
the original skew-bank trace and the independent V29 finite
normalization check.

\subsection{Integration by parts eliminates the tree anisotropy}

\begin{lemma}[Primitive diagonal identities]
\label{lem:v30-primitive}
For each \(\mu\), with differentiation in \(x=\xi_\mu\),
\begin{equation}\label{eq:v30-primitive}
 \begin{aligned}
 p_i+cq_i+\sin x\,\partial_xq_i&=0,&&i\ne\mu,\circ,\\
 p_\mu-q_\mu+\sin x\,\partial_xq_\mu
       &=R_\mu:=\frac23(c-1)\{P(7c+4)-2\},\\
 p_\circ+\frac12c q_\circ+\frac12\sin x\,\partial_xq_\circ
       &=R_\circ:=\frac23\{P(4c^2+6c-1)-5c-4\}.
 \end{aligned}
\end{equation}
These are finite trigonometric identities in all four
coordinates, not identities only after integration.
\end{lemma}
\begin{proof}
Use the same primitive recurrence and quadratic bank as in
\eqref{eq:v30-q}. This is the remaining computer-assisted
finite algebra in the proof, for which we specify an exact
reduction algorithm. Represent every Laurent polynomial by
its finite list of integer exponent vectors and rational
coefficients. Inversion of a coordinate sends its exponent
to its negative. Collecting its signed inversion orbit gives
the odd projection, and collecting unsigned orbits gives
the even projection. For each odd term use
\[
     \frac{\sin(nx)}{\sin x}=U_{n-1}(\cos x),\qquad
     U_0=1,\quad U_1=2\cos x,\quad
     U_{k+1}=2\cos x\,U_k-U_{k-1}.
\]
Thus \(q_i\) is obtained without rational-function division.
On a Laurent polynomial \(q=\sum_e q_e U^e\), the derivative
required in \eqref{eq:v30-primitive} is exactly
\[
 \sin x\,\partial_xq
     =\frac12(U_\mu-U_\mu^{-1})\sum_e e_\mu q_e U^e.
\]
Substitution of the literal \(e/2\) and these \(q_i\) into
the three lines of \eqref{eq:v30-primitive} leaves zero
coefficient for every exponent in the difference with the
displayed targets. The checker performs this collection
for the five banks and all four \(\mu\)'s: twenty identities.
No approximate arithmetic, integration or sampled polynomial
identity test occurs in this step.
\end{proof}

For completeness, the integration by parts using these
identities is justified bank by bank. Both \(p_i\) and
\(q_i\), for a diagonal bank, have a factor \(|R_i|^2\).
For \(p_i\), both noise occurrences in the primitive recurrence
carry \(R_i\) and its conjugate; for \(q_i\), the factor
survives the parity projection and division by \(\sin x\).
Hence their apparent \(\xi_i^{-2}\) singularities are
removable. In particular \(q_\mu=O(x^2)\) on the hyperplane
\(x=0\). The common-bank primitive is \(O(r^2)\) at the
origin, since it has two noise rows, and
\(q_\circ=4(cP-1)=O(r^2)\).
The extra term \(\ell(x)q_\circ/x^2\) is also removable.

These observations bound the integrands in
\eqref{eq:v30-diagonal-integral} and the derivatives below
by \(C r^{-2}\) near the origin. For the derivative term
with \(i\ne\mu\), evenness in \(x\) also gives
\(\partial_xq_i=O(x)|R_i|^2\) locally. For the common bank
the same conclusion follows directly from its displayed
formula. At infinity use
\(S(y)\le C/(1+y^2)\), bounded trigonometric polynomials,
and their removable quotients; the resulting derivatives
are absolutely integrable. Factors such as
\(\sin x\,S'(x)\) are treated by the smooth identity
\eqref{eq:v29-derivative}, not by a singular logarithmic
derivative at a zero of \(S\).

For each fixed nonzero transverse vector, the one-dimensional
boundary values of \(w_i\sin x\,q_i\) at \(x=\pm\infty\)
vanish. Their values across \(x=0\) agree and are zero when
the weight contains \(x^{-2}\). The exceptional zero
transverse vector has measure zero. Absolute integrability
permits Fubini and integration by parts; equivalently the
small-sphere boundary vectors are \(O(r^{-1})\), with flux
\(O(r^2)\). Thus there is no omitted boundary contribution
in the following integration by parts.

Away from removable zeros, differentiating the weights gives
\begin{equation}\label{eq:v30-ibp-derivatives}
 \begin{aligned}
 \partial_x(w_i\sin x\,q_i)
   &=w_i\left((c-\ell)q_i+\sin x\,\partial_xq_i
                    -\frac{2x\sin x}{r^2}q_i\right),&&i\ne\mu,\circ,\\
 \partial_x(w_\mu\sin x\,q_\mu)
   &=w_\mu\left((-1-2\ell)q_\mu+\sin x\,\partial_xq_\mu
                    -\frac{2x\sin x}{r^2}q_\mu\right),\\
 \partial_x(w_\circ\sin x\,q_\circ)
   &=w_\circ\left((c-\ell)q_\circ+\sin x\,\partial_xq_\circ
                    -\frac{4x\sin x}{r^2}q_\circ\right).
 \end{aligned}
\end{equation}
They follow from \(S'/S=\cot(x/2)-2/x\), followed by
\(\sin x\cot(x/2)=1+c\); multiplication removes the poles.
Integrating these three derivatives and using
\eqref{eq:v30-primitive} reduces \eqref{eq:v30-diagonal-integral}
to
\begin{equation}\label{eq:v30-reduced}
 c_\mu=\frac1{(2\pi)^4}\int_{\mathbb R^4}\chi(\xi)
  \left\{\frac{R_\mu-\ell q_\circ/2}{r^2x^2}
                 -\frac{R_\circ+\ell q_\circ}{r^4}\right\}d\xi.
\end{equation}
All three banks \(i\ne\mu\) have canceled exactly. This is
where their orientation-dependent ascending-tree expressions
disappear; no permutation invariance of those expressions
was assumed.

\subsection{A critical scaling difference and its nonzero boundary flux}

Define the smooth even functions
\begin{equation}\label{eq:v30-uv}
 \begin{aligned}
 u(x)&=4\frac{\sin x}{x}-\frac23(5\cos x+1),\\
 v(x)&=\frac23(\cos x+2)-8\frac{\sin x}{x}.
 \end{aligned}
\end{equation}
Use their continuous values at zero. Direct substitution of
\(q_\circ=4(cP-1)\) and \eqref{eq:v30-ell} into
\eqref{eq:v30-reduced} gives
\begin{equation}\label{eq:v30-scaling-algebra}
 \begin{aligned}
 R_\mu-\tfrac12\ell q_\circ&=u(x)-P u(2x),\\
 R_\circ+\ell q_\circ&=v(x)-P v(2x).
 \end{aligned}
\end{equation}
For example these follow simply by using
\(\cos(2x)=2c^2-1\) and
\(\sin(2x)/(2x)=c\sin x/x\); both sides are affine in
\(P\). Their expansions at zero are
\begin{equation}\label{eq:v30-taylor}
 u(x)=x^2-\frac{19}{180}x^4+O(x^6),\qquad
 v(x)=-6+x^2-\frac7{180}x^4+O(x^6).
\end{equation}

Introduce, for \(\xi\ne0\),
\begin{equation}\label{eq:v30-Phi}
 \Phi_\mu(\xi)=\chi(\xi)
       \left\{\frac{u(\xi_\mu)}{|\xi|^2\xi_\mu^2}
                              -\frac{v(\xi_\mu)}{|\xi|^4}\right\}.
\end{equation}
Here \(u(x)/x^2\) has its smooth extension. Combining
\eqref{eq:v30-box-scale} and \eqref{eq:v30-scaling-algebra}
shows that the entire integrand in \eqref{eq:v30-reduced} is
\begin{equation}\label{eq:v30-critical-difference}
                       \Phi_\mu(\xi)-16\Phi_\mu(2\xi).
\end{equation}
The two separate integrals in this difference diverge at the
origin. They must not be split and canceled by a change of
variables over all of \(\mathbb R^4\).

\begin{lemma}[Critical scale-flux formula]
\label{lem:v30-flux}
Suppose \(\Phi\) is integrable outside the unit ball in
\(\mathbb R^4\) and, uniformly in direction,
\[
                     \Phi(\xi)=a|\xi|^{-4}+O(|\xi|^{-2})
                                                \quad(\xi\to0).
\]
Then \(\Phi(\xi)-16\Phi(2\xi)\) is integrable and
\begin{equation}\label{eq:v30-flux}
 \int_{\mathbb R^4}\{\Phi(\xi)-16\Phi(2\xi)\}\,d\xi
                                      =a|\mathbb S^3|\log2.
\end{equation}
Here integrability away from zero on compact annuli is included
in the hypothesis.
\end{lemma}
\begin{proof}
The leading \(|\xi|^{-4}\) terms cancel pointwise; the
remaining \(O(|\xi|^{-2})\) is locally integrable in four
dimensions. Integrability at infinity follows by dilation.
For \(0<2\epsilon<R\), a change of variables on the finite
annulus, where each term is integrable, gives
\[
 \begin{split}
 \int_{\epsilon<|\xi|<R}\{\Phi(\xi)-16\Phi(2\xi)\}\,d\xi
     &=\int_{\epsilon<|\xi|<2\epsilon}\Phi(\xi)\,d\xi\\
     &\quad-\int_{R<|\xi|<2R}\Phi(\xi)\,d\xi.
 \end{split}
\]
The second term tends to zero as \(R\to\infty\). The first
equals \(a|\mathbb S^3|\log2+O(\epsilon^2)\), by radial
integration. Letting \(\epsilon\to0\) proves the claim.
\end{proof}

\begin{theorem}[Exact complete measured leading matrix]
\label{thm:v30-matrix}
Equation \eqref{eq:v30-main} holds. In particular all four
constants left unevaluated in V29 are equal and strictly positive.
\end{theorem}
\begin{proof}
The functions \(u(x)/x^2\) and \(v(x)\) are globally bounded.
Since \(\chi\in L^1(\mathbb R^4)\), equation
\eqref{eq:v30-Phi} implies integrability outside the unit ball.
Near the origin, \(\chi=1+O(r^2)\),
\(u(x)/x^2=1+O(x^2)\), and \(v(x)=-6+O(x^2)\).
Consequently, uniformly in all directions including \(x=0\),
\[
                  \Phi_\mu(\xi)=\frac6{r^4}+O(r^{-2}).
\]
Lemma~\ref{lem:v30-flux} applies with \(a=6\). Since
\(|\mathbb S^3|=2\pi^2\), equations
\eqref{eq:v30-reduced}--\eqref{eq:v30-critical-difference}
give
\[
                 c_\mu=\frac{6(2\pi^2)\log2}{(2\pi)^4}
                         =\frac{3\log2}{4\pi^2}.
\]
This computation is the same for all four \(\mu\). V29
already proved that every off-diagonal entry of
\(\Gamma_*-2h_*^T\) vanishes, which completes the matrix
identity.
\end{proof}

The nonzero scale flux in this last proof is not a boundary
term omitted in \eqref{eq:v30-ibp-derivatives}. Those
derivatives have vanishing boundary flux and integrable
primitives in the required one-dimensional integrations.
The scale flux arises only after rewriting the result as
the difference of two individually nonintegrable critical
kernels. Distinguishing these two operations is essential.

\subsection{The full quartic coefficient and its finite part}

Put \(c_*=3\log2/(4\pi^2)\). The actual leading source matrix
\(\Gamma_*\) remains the non-diagonal matrix of V28. What
has now been proved about its diagonal and the reference
diagonal is the exact relation
\[
                       \eta_\mu=\frac{\gamma_\mu-c_*}{2}.
\]
No individual value of \(\gamma_\mu\) or \(\eta_\mu\) is
needed. The V22 finite-part theorem and
\eqref{eq:v29-finite-part} immediately give
\begin{equation}\label{eq:v30-finite-matrix}
 \mathsf C_{2^n}-n c_*I_4
       \longrightarrow\mathsf C_{\rm fin}
                         =\Gamma_{\rm fin}-2h_{\rm fin}^T.
\end{equation}
This finite matrix need not be diagonal or scalar. In
particular its endpoint contributions have not been erased.

\begin{theorem}[Complete measured noncommuting quartic logarithm]
\label{thm:v30-quartic}
With the balanced-root, coarse-Haar and \(E_4\) conventions
of V14--V22, coefficientwise on the finite-dimensional
space of homogeneous quartic polynomials,
\begin{equation}\label{eq:v30-quartic}
 \begin{aligned}
 q_{*,4}^{\rm rel}&=-\frac{3\log2}{\pi^2}E_4,\\
 \frac{(\mathcal Q_{2^n}^{\rm H})_{[4]}}n
                  &\longrightarrow-\frac{11\log2}{3\pi^2}E_4,\\
 \frac{(\mathcal Q_L^{\rm H})_{[4]}}{\log L}
                  &\longrightarrow-\frac{11}{3\pi^2}E_4
                                \qquad(L\to\infty\text{ dyadically}).
 \end{aligned}
\end{equation}
Moreover the already constructed finite quartic polynomial
\(\mathcal P_{{\rm fin},4}\) of V22 satisfies
\begin{equation}\label{eq:v30-quartic-finite}
 (\mathcal Q_{2^n}^{\rm H})_{[4]}
             +n\frac{11\log2}{3\pi^2}E_4
                               \longrightarrow\mathcal P_{{\rm fin},4}.
\end{equation}
\end{theorem}
\begin{proof}
V19 identifies the relative term as
\(q_{*,4}^{\rm rel}(v)=-g(v)\cdot\mathsf C_*v\).
The homogeneous Euler identity in the current normalization
is \(g(v)\cdot v=4E_4(v)\). Thus
Theorem~\ref{thm:v30-matrix} gives the first line of
\eqref{eq:v30-quartic}. Adding the already established
\(-2\log2\,E_4/(3\pi^2)\) term of
\eqref{eq:v21-full-leading} gives
\(-11\log2\,E_4/(3\pi^2)\). Since \(\log(2^n)=n\log2\),
the third line follows. Finally the V22 finite-part theorem
subtracts precisely this leading quartic polynomial, proving
\eqref{eq:v30-quartic-finite} without changing either of
its endpoint series.
\end{proof}

For the noncommuting probe
\(v=(xE_1,yE_2,0,0)\), one has \(E_4(v)=2x^2y^2\).
Thus the coefficient of \(x^2y^2\), divided by \(\log L\),
is \(-22/(3\pi^2)\). This is a useful explicit normalization
test on the probe, not a calculation on a commuting flat
holonomy. At the level of these coefficients only, the
classical coefficient \(\beta\) of \(E_4\) acquires the
one-loop logarithmic correction \(-11\log L/(3\pi^2)\).
No joint \((\beta,L)\) remainder estimate or identification
with a physical continuum coupling follows just by writing
this coefficient-level expression.

Similarly the leading homogeneous linear countertransport
of V19 is now exactly \(-c_*I_4\). This is not permission
to replace the local microscopic source filter
\(\mathcal A_{p,M}(\Gamma_*)Z\) of V26 by a scalar retained
field. The V27--V28 positive conditional innovations concern
that actual source filter. Its nonzero off-diagonal entries
have not changed. The precision response contributing to
the measured trace is a different mathematical object.

\subsection{Exact certificate and preservation record}

The reproducible certificate is
\begin{center}\small
\path{scripts/verify_ym_f14_v30_measured_diagonal.py}
\end{center}
with SHA--256
\begin{center}\scriptsize\ttfamily
05D3B9DEBDA824CECB253C41DB01AC5383B365EAEEB89A96CC72A05662987974.
\end{center}
It imports the existing literal V17, V28 and V29 banks and
uses rational Laurent arithmetic throughout. Its successful
run verifies 64 mixed-stencil decompositions, sixteen
unprojected gradient identities, twenty primitive
integration-by-parts reductions, twenty full mixed parity
reductions, and ten scaling/Taylor identities. The latter
are the two scaling identities for each of four orientations
and the two displayed Taylor polynomials. The certificate
does not require the symbolic-exploration dependency.
The real-variable boundary estimates and critical scale
integral are proved in this section, not inferred from
the finite checks or from numerical quadrature.

The exploratory script
\begin{center}\small
\path{scripts/explore_ym_f14_v30_diagonal_algebra.py}
\end{center}
helped discover the reductions. It is not a substitute
for the exact certificate and analytic proofs above.

The predecessor V29 has 659,223 bytes and SHA--256
\begin{center}\scriptsize\ttfamily
A817F2607670BC6281D569ACDFC3062C8AB1861323690A600DF7B536E324F05F.
\end{center}
Its complete body is preserved, apart from the title version
and this terminal insertion. V30 replaces it as the single
active main note. No archive, monograph or companion file
has been edited. The checkpoint source hashes are:
\begin{center}\scriptsize
\begin{tabular}{ll}
SM V72 & \texttt{F44F9745D43379E1672C5550411B58E97195A4C9278592D221DF5D3F644DC1F5}\\
NS V26 & \texttt{82C887F8C2C69E4F28FAD7C3DC8DE5B9099CB5A4BF431EBA1FFF3E91935978AD}\\
Shell V16 & \texttt{BF138E63BA826BC030D3A8B87E04FDD3B1EC712D269BB064674E96509B46E710}
\end{tabular}
\end{center}

\subsection{Upstream target after completing this coefficient}

The unresolved task is no longer equality or evaluation of
the four leading measured diagonal constants. Repeating
that calculation would not advance the construction.
The next target is a scale-uniform nonlinear estimate for
the actual conditional density and transported observables
after the complete one-loop subtraction. It must retain
the incoming law, inverse Haar factors, source response,
and large-field part of the integral. V23's failure of
the global small-field domain on typical growing-cutoff
Gaussian samples and V25's divergent uncut cubic
exponential moments exclude a direct inference from
bounded low-order Gaussian jets to that estimate.

In particular, subtracting a logarithm from a coefficient
does not prove positivity or integrability of an exponentiated
truncation. A viable next step must control a genuine
remainder or an exact density identity, with its norm and
scale range specified. The measured cancellation suggests
organizing that step at the level of the coupled density
and observable, while V27 prevents silently discarding
hidden source information before the comparison is proved.

\begin{firewall}
The exact scalar matrix and complete quartic logarithm
are one-loop Taylor-coefficient results for the stated
construction. They do not establish a cutoff-uniform
full nonlinear renormalization trajectory, an interacting
continuum Yang--Mills measure, its noncollapse and
reflection-positive physical transfer, or a positive
Hamiltonian mass gap. The full Yang--Mills objective
remains unproved and unchanged.
\end{firewall}
\section{V31: a nonzero quadratic-chaos logarithm in the fourth source jet}
\label{sec:v31}

\subsection{Going upstream from the measured coefficient}

V30 completed the leading measured quartic logarithm. It did not
make the nonlinear random root a uniformly controlled observable.
Before postulating a remainder estimate, this section tests the
first pure source coefficient not controlled by V23--V26:
\[
                  \mathcal D_L=\mathcal F_{L,[4]}(Z_L),\qquad L=2^n.
\]
This is the fourth Taylor coefficient of the terminal balanced
root under the same free Gaussian, not the quartic coefficient
of the effective action. These two meanings of ``quartic'' must
be kept separate.

The answer changes the next subtraction problem. There is an
explicit nonzero second-chaos field \(\mathcal H_\infty\) such that
\begin{equation}\label{eq:v31-main}
                 \mathcal D_{2^n}/n\longrightarrow\mathcal H_\infty
                         \quad\hbox{in every finite }L^q.
\end{equation}
The fourth-chaos part, in contrast, has a bounded cofinal limit.
The logarithm is generated by the actual non-diagonal source
matrix \(\Gamma_*\), not by replacing it with V30's measured
scalar \(\Gamma_*-2h_*^T\). An explicit transverse helical
field detects it and proves that a scalar multiple of the
quadratic root is not this counterterm.

The upstream audit avoided two duplications. Shell V7 already
proves a volume-uniform measured expansion on its fixed small
box. Archives f6 V28/V41 and f12 V17--V19 already contain bare
Wilson chessboard bounds and averaged conditional score errors.
Repeating those bounds would not close their missing generated-law
and spatial-response estimates. A primary-source check of
\href{https://arxiv.org/abs/2005.00899}
{O'Carroll--Faria da Veiga, \emph{On Yang--Mills Stability and
Plaquette Field Generating Functional} (2020)} supplies context
for multireflection methods, not a new estimate for the present
rooted conditional law. No external estimate is used in the
proof below. The next scheduled companion checkpoint remains V35.

\subsection{The exact fourth-order recursion}

For \(M=2^m\), \(2\le M\le L\), retain
\(p=L/M\), \(j=n-m\), and the fields \(Y_{L,M},Q_{L,M},R_{L,M}\)
of V25--V26. Write \(E_M^{[d]}\) for the homogeneous degree
\(d\) part of the literal primitive balanced block. The notation
\(D^2E_M^{[2]}[Y,R]\) means its constant second derivative;
\(D^3E_M^{[3]}[Y,Y,Q]\) has the analogous meaning.
Define the mean transport from the post-step side \(M/2\) to
the terminal side one by
\[
       \mathcal B_M P=\frac M2
          \operatorname{avg}_{r\in(\mathbb Z/(M/2))^4}P_r.
\]
This includes its physical mass \(M/2\), not just its averaging
weights.

\begin{lemma}[Literal fourth chain rule]
\label{lem:v31-chain}
The terminal coefficient is exactly
\begin{equation}\label{eq:v31-chain}
 \begin{aligned}
 \mathcal D_L&=\sum_{M=2,4,\ldots,L}\mathcal B_M P^{[4]}_{L,M},\\
 P^{[4]}_{L,M}
   &=D^2E_M^{[2]}[Y_{L,M},R_{L,M}]+E_M^{[2]}(Q_{L,M})\\
   &\quad+\frac12D^3E_M^{[3]}[Y_{L,M},Y_{L,M},Q_{L,M}]
                                      +E_M^{[4]}(Y_{L,M}).
 \end{aligned}
\end{equation}
Every term in this formula is retained.
\end{lemma}
\begin{proof}
Substitute \(\epsilon Y+\epsilon^2Q+\epsilon^3R+
\epsilon^4D\) into the primitive Taylor expansion. The linear
term transports \(D\). At order four, the quadratic term gives
\(D^2E^{[2]}[Y,R]+E^{[2]}(Q)\); the cubic gives
\(D^3E^{[3]}[Y,Y,Q]/2\); the quartic gives \(E^{[4]}(Y)\).
Unroll the transported term through the remaining linear blocks.
Their action on a spatial mean is exactly \(\mathcal B_M\),
as in \eqref{eq:v25-cubic-chain}. This proves the formula.
The exact certificate below verifies the primitive identity
against independently evaluated polarizations of the actual
second, third and fourth coefficients.
\end{proof}

Set
\begin{equation}\label{eq:v31-Dfilter}
                    D_{L,M}=\mathcal A_{p,M}(\Gamma_*)Z_L,
 \qquad R_{L,M}=jD_{L,M}+\widehat R_{L,M}.
\end{equation}
The second equality is V26's exact finite-filter subtraction.
Define the new quadratic counterterm, without a limiting-filter
replacement, by
\begin{equation}\label{eq:v31-H}
 \mathcal H_{L,M}=\mathcal B_M D^2E_M^{[2]}[Y_{L,M},D_{L,M}],
 \qquad
 \mathcal H_L=\sum_{M=2,4,\ldots,L}\mathcal H_{L,M}.
\end{equation}
The Lie bracket and colour-isotropic covariance imply
\(\mathbb E\mathcal H_{L,M}=0\); hence these are pure second
chaoses. Let \(\widehat P_{L,M}^{[4]}\) be the second line of
\eqref{eq:v31-chain} with \(R\) replaced by \(\widehat R\),
and write \(\mathcal J_{L,M}=\mathcal B_M\widehat P_{L,M}^{[4]}\).
Then
\begin{equation}\label{eq:v31-exact-subtraction}
 \mathcal D_L
       =\sum_{m=1}^n(n-m)\mathcal H_{L,2^m}
                                  +\sum_{m=1}^n\mathcal J_{L,2^m}.
\end{equation}
The field \(\widehat P^{[4]}\) contains degrees four and two
as ordinary polynomials. Its expectation is zero: it is
\(\operatorname{Ad}\SU(2)\)-equivariant, the Gaussian law is
colour invariant, and the Lie algebra has no invariant vector.
Thus its only homogeneous Gaussian chaoses are two and four.

\subsection{Local kernels needed for the fourth-order estimate}

All kernels in this subsection are on the magnified torus of
side \(M\), acting on \(W_{p,M}=pZ_L\). As in V25, kernels
mean exact cell-constant finite sums, and their norms sum the
finitely many orientation and colour components. Each carrier
is an unfolded box of fixed diameter about the output link;
bounded overlap handles periodic identifications.

\begin{lemma}[Cubic kernel mass and its contracted linear part]
\label{lem:v31-kernel}
The uncontracted trilinear kernel of \(R_{L,M}\) has bounded
absolute \(L^1\) norm, uniformly in \(L,M\). Its Wick
third-chaos kernel has that same bound. The first-chaos part
of \(\widehat R_{L,M}\) has a linear kernel of absolute
\(L^1\) norm at most \(C(1+j)\), on a fixed carrier.
\end{lemma}
\begin{proof}
Use the local version of the cubic insertion sum in V25--V26.
At incoming distance \(s\), a linear input has kernel mass
\(1/s\) in the \(W_{p,M}\) normalization. By
Lemma~\ref{lem:v25-envelope}, an incoming quadratic has kernel
mass at most \(C/s^2\): its native field is
\(W_{p/s,sM}=W_{p,M}/s\), with the corresponding spatial
rescaling. The outer \(B_{s/2}\) has mass \(s/2\).
Both the primitive cubic and the mixed \(YQ\) term therefore
have trilinear kernel mass at most \(C/s^2\). Their sum over
dyadic \(s\) is bounded. The rescaled supports remain in a
fixed box. Wick projection preserves this highest-order kernel.

For the contracted part use the bounded local linear kernel
of Lemma~\ref{lem:v25-first} before the outer averaging.
Its change from \(W_{p/s,sM}\) to \(W_{p,M}\) contributes
\(1/s\), and the outer row has mass \(s/2\). Thus each
cubic shell has contracted linear-kernel mass at most \(C\).
There are \(j\) shells. Subtracting \(jD_{L,M}\) adds at
most \(Cj\), since the cross-oriented linear filters of
V26 have bounded kernel masses. This proves the last assertion.
It is an absolute kernel bound, not a claim that the
\(L^2\) norm of \(\widehat R\) grows with \(j\); that
norm is uniformly bounded by Theorem~\ref{thm:v26-local}.
\end{proof}

\begin{lemma}[Two different bounds for the subtracted local fourth term]
\label{lem:v31-local}
Each local \(\widehat P_{L,M,r}^{[4]}\) has uniformly bounded
\(L^q\) norm for every fixed finite \(q\). Its Wick fourth
kernel has bounded absolute \(L^1\) norm; its Wick second
kernel has absolute \(L^1\) norm at most \(C(1+j)\).
Consequently, writing \(d=\operatorname{dist}_M(2r,2r')\),
\begin{equation}\label{eq:v31-local-cov}
 \begin{aligned}
 |\operatorname{Cov}(\mathsf P_4\widehat P_r,
                          \mathsf P_4\widehat P_{r'})|
                       &\le C(1+d)^{-8},\\
 |\operatorname{Cov}(\mathsf P_2\widehat P_r,
                          \mathsf P_2\widehat P_{r'})|
                       &\le C\min\{1,(1+j)^2(1+d)^{-4}\}.
 \end{aligned}
\end{equation}
\end{lemma}
\begin{proof}
The linear fields have bounded local Gaussian moments.
The local bounds for \(Q\) and \(\widehat R\) are supplied by
Lemma~\ref{lem:v25-Qmom} and Theorem~\ref{thm:v26-local}, respectively.
Holder's inequality on each of the finitely many primitive
stencil terms in \eqref{eq:v31-chain} gives the first assertion.

The fourth Wick kernels are products of bounded-mass linear,
quadratic and trilinear kernels, so their absolute masses are
bounded. To bound the second kernels it is important to keep
track of the possible contractions, rather than bound an
uncontracted diagonal singularity. For \(Y\widehat R\), its
product with \(\mathsf P_1\widehat R\) costs \(C(1+j)\)
by Lemma~\ref{lem:v31-kernel}. Its contraction with
\(\mathsf P_3R\) costs only \(C\): the covariance convolved
with the bounded unit-scale kernel of \(Y\) is uniformly
bounded on the other kernel's fixed carrier, by
\eqref{eq:v25-magnified-Green}.

For \(Q Q\), a second-chaos term contracts one leg from
each quadratic. If \(k\) is a quadratic kernel, its absolute
marginal \(m(y)=\int|k(x,y)|\,dx\) is bounded and has a
bounded carrier by Lemma~\ref{lem:v25-envelope}. The total
mass after this contraction is bounded by
\[
       C\iint m(y)m'(z)
          (\operatorname{dist}_M(y,z)+p^{-1})^{-2}\,dy\,dz\le C.
\]
There is no internal contraction within \(Q\), because
\(Q\) is already a centered second-chaos Lie polynomial.
For \(YYQ\), contracting its two linear legs leaves a
bounded multiple of the quadratic kernel; contracting one
linear leg with a leg of \(Q\) uses the same bounded smoothed
potential as above. Again an internal \(Q\) contraction is
zero. For \(E^{[4]}(Y)\), all contracted pairs consist of
bounded linear smearings. These exhaust the second-chaos
possibilities and prove their asserted mass bound.

If the two carriers are well separated, the Wick covariance
has only cross-pairs: four for the fourth chaos and two for
the second. The Green bound gives respectively \(Cd^{-8}\)
and \(Cd^{-4}\), multiplied by the two kernel masses and a
fixed pairing count. At bounded separation use the uniform
local \(L^2\) bounds and orthogonal projection. The same
bound controls all separations by Cauchy--Schwarz. This proves
\eqref{eq:v31-local-cov}.
\end{proof}

\subsection{The quadratic counterterm and a sublinear remainder}

\begin{lemma}[Summable quadratic counterterm shells]
\label{lem:v31-H}
For every finite \(q\),
\begin{equation}\label{eq:v31-H-bound}
            \|\mathcal H_{L,2^m}\|_q
                         \le C_q2^{-m}\sqrt{1+m}.
\end{equation}
The fields \(\mathcal H_L\) converge on the common Gaussian
coupling, in every finite \(L^q\), to an explicit second chaos
\(\mathcal H_\infty\). Its sum stopped at shell \(2^b\)
has tail at most \(C_q2^{-b}\sqrt{1+b}\).
\end{lemma}
\begin{proof}
Both inputs in \eqref{eq:v31-H} have bounded unit-scale
linear kernels on fixed carriers, including the true
cross-oriented row of \(D\). The local quadratic has bounded
moments and covariance bounded by \(C(1+d)^{-4}\).
On a four-dimensional torus,
\[
       \sum_{r'}(1+\operatorname{dist}_M(2r,2r'))^{-4}
                                                   \le C(1+\log M).
\]
The normalized mean has variance at most
\(CM^{-4}(1+\log M)\). Multiplication by \(M/2\) gives
the asserted \(L^2\) bound, and degree-two Gaussian
hypercontractivity gives its \(L^q\) version.
For fixed \(M\), the finitely many linear fields in its
primitive expression converge in every finite \(L^q\) by
V18/V26. Their products converge as well. The displayed
dyadic bound is summable uniformly in \(L\), proving
the limit and its tail. The second chaos is closed in
\(L^2\), so the limit remains in that chaos.
\end{proof}

\begin{theorem}[Fourth-chaos convergence and the full logarithmic limit]
\label{thm:v31-log}
The fourth-chaos projections \(\mathsf P_4\mathcal D_L\)
converge in every finite \(L^q\). Their shell tail above
\(2^b\) is at most \(C_q2^{-b}\). Moreover
\begin{equation}\label{eq:v31-remainder}
 \|\mathcal D_{2^n}-n\mathcal H_{2^n}\|_q
       \le C_q\sqrt{1+n}\{1+\log(2+n)\}^{1/4}.
\end{equation}
In particular \eqref{eq:v31-main} holds, and its entire
nonzero limit, proved below, is in the second chaos.
The estimate does not claim a bounded finite part after
subtraction.
\end{theorem}
\begin{proof}
The \(j\mathcal H\) term in
\eqref{eq:v31-exact-subtraction} is quadratic, so contributes
nothing to the fourth projection. The first line of
\eqref{eq:v31-local-cov}, its summability in four dimensions,
and the exact mean normalization give
\[
                 \|\mathsf P_4\mathcal J_{L,M}\|_2\le C/M.
\]
At fixed \(M\), \(Y,Q,\widehat R\) converge locally in
every finite \(L^q\), by V25--V26. Their finite products
and fourth projections converge. The uniform summable
shell bound proves the first assertion.

For the second projection the second line of
\eqref{eq:v31-local-cov} gives
\(C(1+j)\sqrt{1+\log M}/M\) after averaging. Independently,
the local moment bound and the triangle inequality give
\(CM\). The fourth projection obeys these weakened bounds
as well. Therefore
\begin{equation}\label{eq:v31-J-bound}
 \|\mathcal J_{L,M}\|_2
     \le C\min\left\{M,\frac{(1+j)\sqrt{1+\log M}}M\right\}.
\end{equation}
Split the dyadic sum at a scale comparable to
\[
                 R_n=\sqrt{1+n}\{1+\log(2+n)\}^{1/4}.
\]
The smaller-scale bound sums to \(CR_n\). Since \(j\le n\),
the larger-scale bound sums to at most
\[
                 C(1+n)\sqrt{1+\log R_n}/R_n\le CR_n.
\]
Rounding to a dyadic scale and small \(n\)
change only the constant. Hence \(\|\sum_M\mathcal J_{L,M}\|_2
\le CR_n\).

Finally subtract \(n\mathcal H_L\) in the exact identity:
\[
 \mathcal D_L-n\mathcal H_L
        =-\sum_{m=1}^n m\mathcal H_{L,2^m}
                                  +\sum_{m=1}^n\mathcal J_{L,2^m}.
\]
The first sum is uniformly bounded by
\eqref{eq:v31-H-bound}. This proves
\eqref{eq:v31-remainder} in \(L^2\); degree-four Gaussian
hypercontractivity proves it in every finite \(L^q\).
Since \(R_n/n\to0\) and \(\mathcal H_L\to\mathcal H_\infty\),
division by \(n\) gives \eqref{eq:v31-main}.
No rate of convergence of \(\mathcal H_L\) was multiplied
by \(n\) without proof: the remainder estimate uses the
finite counterterm \(\mathcal H_L\) itself.
\end{proof}

\subsection{A transverse helical probe evaluates the logarithm exactly}

To prove that the new counterterm is not zero, work on the
physical unit torus and use the smooth mean-zero transverse
field
\begin{equation}\label{eq:v31-helical}
 u_4(y)=E_1\cos(2\pi y_2)+E_2\sin(2\pi y_2),\qquad
 u_1=u_2=u_3=0.
\end{equation}
It is Coulomb since \(\partial_4u_4=0\). In the real Gaussian
energy convention of V23 its Cameron--Martin norm is
\begin{equation}\label{eq:v31-energy}
                         \|u\|_{\rm CM}^2=4\pi^2.
\end{equation}
For a quadratic second chaos, notation such as
\(\mathcal H_\infty[u,u]\) denotes its symmetric kernel
evaluated on this Cameron--Martin vector twice. Equivalently
it is the expectation after shifting the common Gaussian by
\(u\). This does not evaluate a nonlinear continuum root on
a rough distribution.

\begin{lemma}[A one-axis primitive bank]
\label{lem:v31-probe-bank}
In \eqref{eq:v29-K}, set \(U_2=z=e^{ix}\) and the other
three outer phases equal to one. Then
\begin{equation}\label{eq:v31-probe-bank}
 \begin{aligned}
 K_{4,41}&=0,& K_{4,42}&=1-z,\\
 K_{4,43}&=-1+\tfrac12(z+z^{-1}),&K_{4,44}&=0.
 \end{aligned}
\end{equation}
For the field \eqref{eq:v31-helical}, the fourth-orientation,
third-colour component of the continuum counterterm shell is
\begin{equation}\label{eq:v31-probe-shell}
        \bigl(\mathcal H_{\infty,M}[u,u]\bigr)_4^{,3}
          =\frac{a_*}{\pi}\,
                       S(2\pi/M)\sin^2(\pi/M).
\end{equation}
The corresponding component of the quadratic-root shell is
zero.
\end{lemma}
\begin{proof}
Expand the finite quadratic wedge bank of V28 and set the
three phases to one. Its four remaining Laurent polynomials
are exactly \eqref{eq:v31-probe-bank}; all four are checked
coefficientwise in the certificate below.

Put \(x=2\pi/M\), \(f=\beta(x)/M\), and use the complex
colour vector \(v=E_1-iE_2\). The continuum linear inputs
at this shell are
\[
       Y_4(s)=\operatorname{Re}(v f e^{ixs_2}),\qquad
       D_\alpha(s)=(\Gamma_*)_{\alpha4}
               \operatorname{Re}(v f\beta(x\delta_{\alpha2})e^{ixs_2}).
\]
This follows from the actual box and directional multipliers
of V28, including the factor \(1/M\) for the physical
Cameron--Martin field. For complex scalars \(A,B\),
\[
 [\operatorname{Re}(vA),\operatorname{Re}(vB)]
                                   =-2\operatorname{Im}(A\overline B)E_3.
\]
Since \(K\) is twice the skew coefficient bank, the
third-colour component of the mixed quadratic derivative is
\[
       -|f|^2\operatorname{Im}
          \sum_\alpha K_{4,4\alpha}(U)(\Gamma_*)_{\alpha4}
                      \overline{\beta(x\delta_{\alpha2})}.
\]
It is spatially constant: the equal-sign Fourier pairs have
bracket \([v,v]=0\), leaving only the opposite-sign pairs.
Thus its actual mean transport multiplies it by \(M/2\).

By V28, the relevant column is
\((0,a_*,b_*,\gamma_4)^t\). The \(b_*\) term in
\eqref{eq:v31-probe-bank} is real, the \(\gamma_4\) term
vanishes, and
\[
        \operatorname{Im}\{(1-e^{ix})\overline{\beta(x)}\}
                                      =-\frac{4\sin^2(x/2)}x.
\]
Using \(|f|^2=S(x)/M^2\) and \(Mx=2\pi\) gives
\eqref{eq:v31-probe-shell}. A pure quadratic-root component
on this field has only the input pair \((4,4)\); its bank
vanishes by \(K_{4,44}=0\). This proves the last assertion.
\end{proof}

\begin{theorem}[A nonzero logarithm, distinct from scalar root renormalization]
\label{thm:v31-nonzero}
The limiting quadratic counterterm satisfies
\begin{equation}\label{eq:v31-probe-value}
       \bigl(\mathcal H_\infty[u,u]\bigr)_4^{,3}=a_*/\pi>0,
       \qquad \bigl(Q_\infty[u,u]\bigr)_4^{,3}=0.
\end{equation}
In particular
\begin{equation}\label{eq:v31-positive-norm}
       \| (\mathcal H_\infty)_4^{,3}\|_2
                        \ge\frac{a_*}{\sqrt8\,\pi^3}>0.
\end{equation}
For every real scalar \(\alpha\), the same lower bound holds
for \(\| (\mathcal H_\infty-\alpha Q_\infty)_4^{,3}\|_2\).
\end{theorem}
\begin{proof}
The scaling identity \(S(2x)=S(x)\cos^2(x/2)\) gives
\[
 \sum_{m=1}^b S(2\pi/2^m)\sin^2(\pi/2^m)
       =S(2\pi/2^b)-S(2\pi)\longrightarrow1.
\]
Sum \eqref{eq:v31-probe-shell} to obtain
\eqref{eq:v31-probe-value}; the quadratic-root component
vanishes at every shell. The passage to the limiting quadratic
kernels is legitimate: \(L^2\) convergence of second chaoses
is Hilbert--Schmidt convergence of their symmetric energy
kernels, which is continuous on two fixed energy vectors.
Alternatively take the fine fields \(u_L(x)=L^{-1}u(x/L)\).
Their fixed Fourier-mode energy representatives converge to
that of \(u\), and their norms converge to \(4\pi^2\);
the shell limits and summable bounds above give the same result.

If \(I_2(T)\) denotes a scalar second chaos with symmetric
energy kernel \(T\), then
\(\|I_2(T)\|_2^2=2\|T\|_{\rm HS}^2\), and
\( |T[u,u]|\le\|T\|_{\rm HS}\|u\|_{\rm CM}^2\).
Apply these two elementary Wick and operator inequalities to
the component in \eqref{eq:v31-probe-value} and use
\eqref{eq:v31-energy}. They give
\(\|I_2(T)\|_2\ge\sqrt2(a_*/\pi)/(4\pi^2)\), which is
\eqref{eq:v31-positive-norm}. The kernel of \(\alpha Q_\infty\)
has zero value on this probe in this component, so the same
argument proves the final assertion. Strict positivity of
\(a_*\), including its explicit lower bound, is inherited
from Theorem~\ref{thm:v28-source}.
\end{proof}

Combining the preceding theorems yields an actual lower
growth statement:
\[
       \lim_{n\to\infty}\frac{\| (\mathcal D_{2^n})_4^{,3}\|_2}{n}
          =\| (\mathcal H_\infty)_4^{,3}\|_2
                              \ge\frac{a_*}{\sqrt8\,\pi^3}.
\]
This is not inferred from a finite-cutoff nonzero example.
Its nonzero limiting value was evaluated analytically above.

\subsection{What this forces on a nonlinear source estimate}

The scalar output change
\[
        (1-\alpha n\epsilon^2)\mathcal F_L(\epsilon Z_L)
\]
changes its fourth coefficient to
\(\mathcal D_L-\alpha nQ_L\). By
Theorems~\ref{thm:v23-chaos}, \ref{thm:v31-log} and
\ref{thm:v31-nonzero}, this coefficient divided by \(n\)
still has a nonzero limit for every \(\alpha\).
In particular V30's scalar measured correction does not
by itself renormalize this stochastic source coefficient.
This excludes that specific scalar shortcut, not all
fine-field counterterms, general source changes, or
integration against the actual interacting law.

There is a precise coupled-noise consequence. Because the
terminal linear average vanishes on the fine mean-zero field,
define the fourth truncated terminal source by
\[
 T_L^{[4]}(\beta)=\beta^{-1}Q_L
          +\beta^{-3/2}\mathcal C_L+\beta^{-2}\mathcal D_L.
\]
If \(\beta_L\to\infty\) and \(n/\beta_L\to\tau<\infty\), then
\begin{equation}\label{eq:v31-joint}
             \beta_L T_L^{[4]}(\beta_L)
                      \longrightarrow Q_\infty+\tau\mathcal H_\infty
                            \quad\hbox{in every finite }L^q.
\end{equation}
Indeed \(Q_L\to Q_\infty\) by V23,
\(\mathcal C_L/\sqrt{\beta_L}\to0\) by V25, and
\(\mathcal D_L/\beta_L=(n/\beta_L)(\mathcal D_L/n)\)
has the limit just proved. This remains a theorem about a
truncated polynomial. It neither resums higher jets nor
defines the full root on a continuum Gaussian field.

The new finite counterterm is \(n\mathcal H_L\), built
from the actual mixed primitive and cross-oriented filter
in \eqref{eq:v31-H}. Inserting V26's local cubic subtraction
inside the fourth chain produces exactly
\(\sum_m(n-m)\mathcal H_{L,2^m}\); their difference is the
uniformly bounded series \(\sum_m m\mathcal H_{L,2^m}\).
Thus propagating the subtraction through the primitive
really matters. Merely bounding the already renormalized
local cubic and forgetting its insertion would miss the
terminal fourth-order logarithm.

The next concrete step is to sharpen the second-chaos
remainder after this finite subtraction, or to incorporate
this explicit quadratic fine-field counterterm in a
controlled density-and-observable map. The bound
\eqref{eq:v31-remainder} is sublinear, not uniformly bounded;
it does not yet supply that next finite part. The full
nonlinear integral, generated incoming law and its
large-field response remain separate obligations.

\subsection{Exact checks and append-only record}

The certificate is
\begin{center}\small
\path{scripts/verify_ym_f14_v31_fourth_source.py}
\end{center}
with SHA--256
\begin{center}\scriptsize\ttfamily
469C963F4B65AE9ECBF92C4CD89CDD5E79DE6E638470BFBC3E14655C13CAAB43.
\end{center}
It checks the degree-four BCH implementation against
independent quaternion exponential, product and logarithm
series on 24 rational inputs; it checks all four components
of the primitive fourth chain by independent polarizations.
Its first three coefficients agree with the V25 literal
root on that version's live side-two test field. On the
same field, the new fourth coefficients are exactly
\[
 \left(-\frac{970621}{42467328}E_3,
          \frac{416357}{42467328}E_3,0,0\right).
\]
The checker also verifies fourth-degree homogeneity, the
four Laurent identities in \eqref{eq:v31-probe-bank}, and
the numerator identity for the double-angle scaling of
\(S\). An independent literal primitive evaluation at
\(x=\pi/2\), with helical amplitudes \(1\) and \(1+i\),
gives mixed output \(+2E_3\); this fixes the sign and
factor used in the helical computation. All these checks
use rational arithmetic. The covariance, summation and
limiting-kernel arguments are analytic proofs above,
not conclusions drawn from the finite tests.

The predecessor V30 has 685,503 bytes and SHA--256
\begin{center}\scriptsize\ttfamily
E5D894EFB16EA8F69B2374F2F86829E699AEAE8CE6952074D97B9822E6B1EB38.
\end{center}
Its body is preserved apart from the title version and
this terminal insertion. V31 replaces V30 as the sole
active main note. The companion files, archives and
monographs are unchanged. No new companion checkpoint
was due in this version; the next is V35.

\begin{firewall}
This section proves a new logarithm and its counterterm
in the fourth Taylor coefficient of the canonical random
source, plus convergence of its fourth-chaos component.
It does not assert a uniformly bounded subtracted fourth
finite part, a full nonlinear source remainder, an
interacting continuum measure, or a positive physical
Hamiltonian mass gap. The complete Yang--Mills objective
remains unproved and unchanged.
\end{firewall}
\section{V32: the locally renormalized fourth source coefficient has a finite part}
\label{sec:v32}

\subsection{The sharper estimate and its scope}

V31 identified the fourth-source logarithm but bounded its
subtracted remainder only by
\(C\sqrt n\{1+\log(2+n)\}^{1/4}\). This section replaces
that bound by a uniform bound and a cofinal finite part.
The result also holds at every fixed intermediate output
side, not only at the terminal root.

The key is not a new assumption of pointwise kernel
cancellation. V26 already separates the contracted cubic
into individual averaging errors. Each error has a bounded
local kernel and an \(L^q\) norm decaying with its inner
scale. Keeping these two estimates separate until after
the spatial averaging gives a summable outer-scale bound.
Summing the absolute kernels first, as permitted by the
weaker V31 argument, loses this information.

The predecessor and the three active companions were
checked against their current files. They are unchanged.
No fifth-version checkpoint is due before V35. The
calculation below is new finite-part control of the actual
source coefficients, not a reproof of the fixed-box quantum
expansion or the complete measured one-loop logarithm.

\subsection{Local notation and the exact subtraction}

Let \(L=pM\) be dyadic with \(M\ge1\), and write
\(j=\log_2p\). Define
\[
       \mathcal V_{L,M}=\mathcal F_{L\to M,[4]}(Z_L).
\]
For an outer distance \(s=2^\ell\), \(1\le\ell\le j\), put
\(R=sM\). All incoming quantities at side \(R\) have
\(R\ge2\), so the local V25--V26 results apply even when
the final side is \(M=1\). Let
\[
       D_{L,R}=\mathcal A_{L/R,R}(\Gamma_*)Z_L
\]
be the same cross-oriented linear source filter as in V31.
Use the exact positive linear row \(B_{s/2}\) from the
post-primitive side \(R/2\) to side \(M\), and set
\begin{equation}\label{eq:v32-Hlocal}
 \begin{aligned}
 \mathcal H_{L,M;s}
    &=B_{s/2}D^2E_R^{[2]}[Y_{L,R},D_{L,R}],\\
 \mathcal H^{\rm loc}_{L,M}
    &=\sum_{\ell=1}^j\mathcal H_{L,M;2^\ell},\\
 \mathcal J_{L,M;s}
    &=B_{s/2}\widehat P_{L,R}^{[4]}.
 \end{aligned}
\end{equation}
Here \(\widehat P^{[4]}\) is exactly the V31 primitive
expression with the incoming cubic replaced by
\(\widehat R_{L,R}\); all its other terms are retained.
At \(M=1\), \(\mathcal H^{\rm loc}_{L,1}=\mathcal H_L\)
of \eqref{eq:v31-H}. Empty sums are zero.

The fourth chain rule and V26's finite-filter identity give
\begin{equation}\label{eq:v32-identity}
 \begin{aligned}
 \mathcal V_{L,M}
     &=\sum_{\ell=1}^j(j-\ell)\mathcal H_{L,M;2^\ell}
                       +\sum_{\ell=1}^j\mathcal J_{L,M;2^\ell},\\
 \widehat{\mathcal V}_{L,M}
     &:=\mathcal V_{L,M}-j\mathcal H^{\rm loc}_{L,M}\\
     &=\sum_{\ell=1}^j
                \{\mathcal J_{L,M;2^\ell}
                           -\ell\mathcal H_{L,M;2^\ell}\}.
 \end{aligned}
\end{equation}
The coefficient of the incoming cubic logarithm is
\(\log_2(L/R)=j-\ell\), not \(j\). The last term in
\eqref{eq:v32-identity} keeps that clock difference exactly.

\subsection{The cubic decomposition must be used before summing kernels}

At a fixed incoming side \(R\), define the bounded endpoint
matrix
\[
       A_{L,R}^{\rm end}
          =\Gamma_L^{>R}-\log_2(L/R)\Gamma_*.
\]
Equation \eqref{eq:v26-Gamma-bound} gives
\(\|A_{L,R}^{\rm end}\|\le C\), uniformly in \(L\ge R\).
The exact decomposition \eqref{eq:v26-decomposition} is
\begin{equation}\label{eq:v32-R-decomposition}
 \widehat R_{L,R}
    =\mathsf P_3R_{L,R}
       +\sum_{h=1}^{\log_2(L/R)}E_{L,R,2^h}
       +\mathcal A_{L/R,R}(A_{L,R}^{\rm end})Z_L.
\end{equation}
The finite endpoint matrix is not dropped. For fixed \(R\),
it converges to \(B_R\) of \eqref{eq:v26-BM}, containing
both coarse- and fine-endpoint series.

For clarity about the new estimate, write
\begin{equation}\label{eq:v32-primitive-split}
 \widehat P_{L,R}^{[4]}
       =U^0_{L,R}
           +\sum_{h=1}^{\log_2(L/R)}U^{(h)}_{L,R},
 \qquad
 U^{(h)}_{L,R}=D^2E_R^{[2]}[Y_{L,R},E_{L,R,2^h}].
\end{equation}
The core is the fully specified expression
\begin{equation}\label{eq:v32-core}
 \begin{aligned}
 U^0_{L,R}
  ={}&D^2E_R^{[2]}[Y_{L,R},\mathsf P_3R_{L,R}]\\
    &+D^2E_R^{[2]}
       [Y_{L,R},\mathcal A_{L/R,R}(A_{L,R}^{\rm end})Z_L]\\
    &+E_R^{[2]}(Q_{L,R})
       +\frac12D^3E_R^{[3]}[Y_{L,R},Y_{L,R},Q_{L,R}]\\
    &+E_R^{[4]}(Y_{L,R}).
 \end{aligned}
\end{equation}
Every expression here is a polynomial on the original
common Gaussian family. The errors for distinct \(h\)
are not declared independent.

\begin{lemma}[Two complementary estimates on each inner error]
\label{lem:v32-inner}
On the side-\(R\) magnified coordinates, each
\(E_{L,R,2^h}\) is a first-chaos row with fixed carrier
and uniformly bounded absolute kernel \(L^1\) norm.
Its local \(L^q\) norm is at most \(C_q2^{-h/3}\).
Consequently the quadratic primitive \(U^{(h)}_{L,R,r}\)
has the following properties, uniformly in \(L,R,h\):
\begin{equation}\label{eq:v32-inner}
 \begin{aligned}
 \|U^{(h)}_{L,R,r}\|_q&\le C_q2^{-h/3},\\
 |\operatorname{Cov}(U^{(h)}_{L,R,r},U^{(h)}_{L,R,r'})|
             &\le C(1+\operatorname{dist}_R(2r,2r'))^{-4}.
 \end{aligned}
\end{equation}
It is a centered second chaos. The spatial bound has no
inner-scale factor; the local moment bound is the separate
small estimate that will be combined with it.
\end{lemma}
\begin{proof}
In the proof of Lemma~\ref{lem:v26-error}, the error kernel
is the difference of \(K_{2^h}\) in \eqref{eq:v26-K} and
the actual finite averaging kernel multiplied by its shell
matrix. Both have a common fixed carrier and bounded
\(L^1,L^\infty\) norms. The shell matrices are uniformly
bounded. This proves the asserted absolute mass bound,
without claiming that this mass tends to zero.
Equation \eqref{eq:v26-error-bound} supplies the separate
small Gaussian norm. It was proved using the low-frequency
mass matching and negative-order covariance multiplier.

The primitive derivative is a finite sum of Lie brackets
of linear rows. Colour isotropy makes its expectation zero.
The product estimate
\[
       \|Y E_{L,R,2^h}\|_2
                    \le\|Y\|_4\|E_{L,R,2^h}\|_4
                    \le C2^{-h/3}
\]
gives the first inequality, with higher moments obtained
in the same way or by degree-two hypercontractivity.
Its quadratic Wick kernel is a product of the two bounded
linear kernel masses on a fixed carrier. At separated
carriers its covariance has two cross-pairs, each bounded
by the inverse-square Green estimate of V25. At bounded
separation use its uniform local second moment. This proves
the second inequality. Only estimates on one and the same
Gaussian law have been combined.
\end{proof}

\begin{lemma}[The core has no growing absolute-kernel factor]
\label{lem:v32-core}
The local core \(U^0_{L,R,r}\) has uniformly bounded
finite moments and has only second and fourth chaoses.
Its covariance satisfies
\begin{equation}\label{eq:v32-core-cov}
       |\operatorname{Cov}(U^0_{L,R,r},U^0_{L,R,r'})|
                \le C(1+\operatorname{dist}_R(2r,2r'))^{-4}.
\end{equation}
Its fourth-chaos covariance has the stronger power eight.
\end{lemma}
\begin{proof}
The local moments of \(\mathsf P_3R\), \(Q\), and all the
linear rows are uniform by V25--V26, including the bounded
endpoint matrix. Holder's inequality gives the moment bound.
The trilinear Wick kernel of \(\mathsf P_3R\) has bounded
absolute mass by Lemma~\ref{lem:v31-kernel}. All contractions
with \(Y\) have bounded quadratic kernel mass, by the bounded
smoothed inverse-square potential. The \(QQ\), \(YYQ\), and
four-linear-leg terms have the bounded second and fourth
kernel masses established individually in the proof of
Lemma~\ref{lem:v31-local}. The endpoint term is simply a
quadratic product of two bounded linear kernels.
There is no \(\mathsf P_1\widehat R\) row with mass
\(C(1+\log_2(L/R))\) left in the core: its individual
errors have been separated in \eqref{eq:v32-primitive-split}.

Equivariance and colour invariance give zero expectation.
The Wick pairing bounds on the fixed carriers now give
power four for the second chaos and power eight for the
fourth, with constants independent of the incoming depth.
Their covariance cross term is zero by chaos orthogonality.
This proves \eqref{eq:v32-core-cov}.
\end{proof}

\subsection{Positive block rows give a two-scale summation bound}

The required spatial averaging estimate works for the actual
local block row, not only for the global mean. For
\(a=s/2\), its nonnegative weights \(b_a(r)\) on the
post-step grid obey
\begin{equation}\label{eq:v32-row}
       \sum_r b_a(r)=a,\qquad
       \max_r b_a(r)\le Ca^{-3},\qquad
       \sum_r b_a(r)^2\le Ca^{-2}.
\end{equation}
Their unfolded carrier has diameter \(Ca\). These are the
literal box-plus-directional rows of V17. At final side
one, folding makes them the normalized global mean times
their mass, and the same inequalities hold. Bounded
periodic overlap changes only the numerical constant.

If a centered translated local field has covariance bounded
by \(C(1+d)^{-4}\), its weighted average therefore obeys
\begin{equation}\label{eq:v32-average}
          \left\|\sum_r b_a(r)X_r\right\|_2
                       \le \frac{C\sqrt{1+\log s}}s.
\end{equation}
Indeed, on the row's carrier the covariance row sum is at
most \(C(1+\log s)\), uniformly in the surrounding torus.
One may bound one weight by its maximum and sum the other
by its mass. This gives variance at most
\(Ca^{-2}(1+\log s)\). The bounded-size case \(a=1\)
is included by increasing the constant. For covariance
power eight the logarithm is absent.

For the \(h\)-th inner error, the triangle inequality gives
another bound: the row has mass \(a\) and the local norm
in \eqref{eq:v32-inner} is small. Combining, for
\(s=2^\ell\), yields
\begin{equation}\label{eq:v32-min}
 \|B_{s/2}U^{(h)}_{L,sM}\|_q
   \le C_q\min\{2^\ell2^{-h/3},\,2^{-\ell}\sqrt{1+\ell}\}.
\end{equation}
For \(q\ge2\), hypercontractivity transfers the spatial
\(L^2\) bound to \(L^q\); for smaller \(q\), use monotonicity
of probability norms. The other bound follows directly
from the local \(L^q\) estimate.

\begin{lemma}[The inner-depth factor is replaced by an outer logarithm]
\label{lem:v32-sum}
For every \(\ell\ge1\),
\begin{equation}\label{eq:v32-sum}
 \sum_{h=1}^{\infty}
       \min\{2^\ell2^{-h/3},\,2^{-\ell}\sqrt{1+\ell}\}
                    \le 6\,2^{-\ell}(1+\ell)^{3/2}.
\end{equation}
The same bound holds for every finite initial segment.
\end{lemma}
\begin{proof}
For \(h\le6\ell\), sum the second alternative to obtain
\(6\ell\,2^{-\ell}\sqrt{1+\ell}\). For the remaining
terms group the first alternative into blocks of three.
If \(h=6\ell+3k+r\), \(k\ge0\), \(1\le r\le3\),
then \(2^{-h/3}\le2^{-2\ell-k}\). Hence that tail is at
most \(6\,2^{-\ell}\). Adding the bounds and using
\(\ell\sqrt{1+\ell}+1\le(1+\ell)^{3/2}\) proves the claim.
\end{proof}

This step keeps the incoming errors correlated and uses
only the triangle inequality between their scales. It
does not sum variances as though the errors were independent.
Nor does it require their absolute kernel masses to vanish.

\subsection{Uniform local fourth finite part}

\begin{theorem}[Summable local subtracted shells]
\label{thm:v32-shells}
Uniformly in dyadic \(L=pM\ge M\ge1\), every output link
and every finite \(q\) satisfy
\begin{equation}\label{eq:v32-shells}
 \begin{aligned}
 \|\mathcal H_{L,M;2^\ell}\|_q
                      &\le C_q2^{-\ell}\sqrt{1+\ell},\\
 \|\mathcal J_{L,M;2^\ell}\|_q
                      &\le C_q2^{-\ell}(1+\ell)^{3/2},\\
 \|\mathcal J_{L,M;2^\ell}-\ell\mathcal H_{L,M;2^\ell}\|_q
                      &\le C_q2^{-\ell}(1+\ell)^{3/2}.
 \end{aligned}
\end{equation}
In particular \(\widehat{\mathcal V}_{L,M}\) has bounded
local \(L^q\) norms independent of the incoming depth and
output side. Its outer-shell tail above \(2^b\) is at most
\(C_q2^{-b}(1+b)^{3/2}\).
\end{theorem}
\begin{proof}
The two linear inputs in \(\mathcal H\) have bounded
unit-scale kernels, bounded local moments and a centered
quadratic product with covariance power four. Apply
\eqref{eq:v32-average}. For \(\mathcal J\), first apply
Lemma~\ref{lem:v32-core} and the same row estimate to the
core in \eqref{eq:v32-primitive-split}. For its inner-error
sum apply \eqref{eq:v32-min} and Lemma~\ref{lem:v32-sum}.
The allowed inner range is only \(1\le h\le j-\ell\);
extending it to infinity increases the bound. This proves
the second inequality, with no factor \(j\).
The first inequality, multiplied by \(\ell\), proves the
third by the triangle inequality.

Finally sum the last line over dyadic outer distances.
For example, with \(k=\ell-b\ge1\),
\((1+b+k)^{3/2}\le(1+b)^{3/2}(1+k)^{3/2}\), and
\(\sum_{k\ge1}2^{-k}(1+k)^{3/2}\le
\sum_{k\ge1}2^{-k}(1+k)^2=11\).
This proves the uniform bound and the asserted tail.
\end{proof}

\begin{theorem}[Cofinal renormalized fourth source]
\label{thm:v32-finite}
For every fixed output side \(M\), there are limits on the
same Gaussian coupling,
\begin{equation}\label{eq:v32-finite}
 \begin{aligned}
 \mathcal H^{\rm loc}_{L,M}&\longrightarrow\mathcal H^{\rm loc}_{\infty,M},\\
 \mathcal F_{L\to M,[4]}(Z_L)
          -\log_2(L/M)\mathcal H^{\rm loc}_{L,M}
                         &\longrightarrow\mathcal V^{\rm fin}_{\infty,M},
 \end{aligned}
\end{equation}
jointly on each finite output carrier in every finite
\(L^q\). The second limit is the explicit convergent series
\begin{equation}\label{eq:v32-series}
       \mathcal V^{\rm fin}_{\infty,M}
          =\sum_{\ell\ge1}
              \{\mathcal J_{\infty,M;2^\ell}
                           -\ell\mathcal H_{\infty,M;2^\ell}\}.
\end{equation}
Its second- and fourth-chaos parts converge separately.
There are constants \(\eta>0,C<\infty\), independent of
\(L,M\), such that at every output link
\begin{equation}\label{eq:v32-stretched}
       \mathbb E\exp\{\eta|\widehat{\mathcal V}_{L,M}|^{1/2}\}\le C.
\end{equation}
The same local bounds hold for the limiting finite parts.
\end{theorem}
\begin{proof}
Fix \(M\) and an outer distance \(s\). The incoming side
\(R=sM\) is fixed. The fields \(Y_{L,R},Q_{L,R}\), the
cross-oriented \(D_{L,R}\), and \(\widehat R_{L,R}\) converge
jointly on the finite primitive carrier by V18 and V25--V26.
Using higher finite moments, their polynomial products
converge in \(L^q\). The outer \(B_{s/2}\) is a fixed finite
row. This defines \(\mathcal H_{\infty,M;s}\) and
\(\mathcal J_{\infty,M;s}\) without discarding any endpoint
term from \(\widehat R\).

Truncate the outer sums at \(2^b\), pass to the limit in
this finite head, and then use the uniform tails from
Theorem~\ref{thm:v32-shells}. This proves both limits and
\eqref{eq:v32-series}. Orthogonal chaos projections are
continuous in \(L^2\); their differences have degree at
most four, so hypercontractivity gives the separate
convergence in every finite \(L^q\).

The uniform \(L^2\) bound and degree-four hypercontractivity
give \(\|\widehat{\mathcal V}_{L,M}\|_q\le Cq^2\) for
\(q\ge2\), with a harmless finite colour-component factor.
Expand the exponential in \eqref{eq:v32-stretched} and
use \(q=k/2\) for its \(k\)-th moment when \(k\ge4\).
The resulting terms are bounded by
\((C'\eta)^k k^k/k!\). The series converges for sufficiently
small \(\eta\); the first three terms are bounded separately.
Fatou's lemma gives the same conclusion for the limits.
This is a stretched-exponential estimate, not an ordinary
exponential moment of an uncut quartic polynomial.
\end{proof}

The first limit is a second-chaos observable. Its same
uniform \(L^2\) bound and degree-two hypercontractivity
also give a uniform small ordinary exponential moment of
\(|\mathcal H^{\rm loc}_{L,M}|\). This does not justify
exponentiating the fourth finite part or a complete
uncontrolled Taylor remainder.

\subsection{The terminal finite part and its scheme convention}

Taking \(M=1\) in \eqref{eq:v32-finite} gives the sharper
form of the V31 logarithm:
\begin{equation}\label{eq:v32-terminal}
          \mathcal D_{2^n}-n\mathcal H_{2^n}
                       \longrightarrow\mathcal D_{\rm fin}
                  :=\mathcal V^{\rm fin}_{\infty,1}
                        \quad\hbox{in every finite }L^q,
\end{equation}
with uniformly bounded left side. Explicitly,
\[
       \mathcal D_{\rm fin}
           =\sum_{m\ge1}\{
                    \mathcal J_{\infty,1;2^m}
                                   -m\mathcal H_{\infty,2^m}\}.
\]
In this last formula \(\mathcal H_{\infty,2^m}\) is V31's
terminal counterterm-shell notation. The series has the
tail \(C_q2^{-b}(1+b)^{3/2}\). The nonzero logarithmic
coefficient and helical value \(a_*/\pi\) from V31 are
unchanged; this section controls what remains after it.

One can instead subtract the exact propagated local clocks
\(\sum_{\ell=1}^j(j-\ell)\mathcal H_{L,M;2^\ell}\).
The remaining coefficient then converges to
\(\sum_{\ell\ge1}\mathcal J_{\infty,M;2^\ell}\).
These two finite-part conventions differ by the explicitly
convergent series \(\sum_{\ell\ge1}\ell
\mathcal H_{\infty,M;2^\ell}\), not by an unspecified constant.

The finite counterterm in \eqref{eq:v32-finite} remains
\(j\mathcal H^{\rm loc}_{L,M}\). Replacing it by
\(j\mathcal H^{\rm loc}_{\infty,M}\) would additionally require
\[
       j\|\mathcal H^{\rm loc}_{L,M}
                    -\mathcal H^{\rm loc}_{\infty,M}\|_q\longrightarrow0.
\]
The finite-part proof does not use or assert that rate.
Likewise its tail bound is a bound for discarding outer
shells, not by itself a rate in the original mesh \(L\).

\subsection{A controlled four-jet, not yet the full nonlinear source}

Together with V25--V26, the theorem supplies the local
renormalized polynomial
\begin{equation}\label{eq:v32-four-jet}
       \mathcal T_{L,M}^{\rm ren,[4]}(t)
         =tY_{L,M}+t^2Q_{L,M}
                     +t^3\widehat R_{L,M}
                     +t^4\widehat{\mathcal V}_{L,M}.
\end{equation}
At \(M=1\), set \(\widehat R_{L,1}=\mathcal C_L\):
the corresponding linear counterfilter vanishes on the
mean-zero fine field, and V25 supplies its bounds and limit.
At fixed \(M\), all four coefficients have cofinal limits
on the common Gaussian family; each has uniformly bounded
local finite moments. In particular, for \(|t|\le1\),
\[
       \|\mathcal T_{L,M}^{\rm ren,[4]}(t)
                         -tY_{L,M}-t^2Q_{L,M}\|_q
                                               \le C_q|t|^3.
\]
This is an estimate inside the explicitly defined polynomial.
It is not an estimate of
\(\mathcal F_{L\to M}(tZ_L)-\mathcal T_{L,M}^{\rm ren,[4]}(t)\).
No full nonlinear map has been evaluated outside its
proved domain, and no interacting density has been replaced
by the free reference.

The next upstream issue is now more specific. A nonlinear
construction must propagate the actual linear and quadratic
fine-field counterterms together with the source and
measured density, and control the remaining orders in a
range where the incoming depth grows with inverse noise
variance. Boundedness of four renormalized coefficients
does not prove a common analytic radius, convergence of
a counterterm series, or uniform integrability of the
full exponential weight. The present finite part removes
one concrete missing coefficient bound; it does not
replace those requirements with a lower-order success.

\subsection{Finite transport calibration and preservation}

The exact checker is
\begin{center}\small
\path{scripts/verify_ym_f14_v32_fourth_finite_part.py}
\end{center}
with SHA--256
\begin{center}\scriptsize\ttfamily
F162BDF0502F99938026007A23918A469808CF3A33F4AF4E12A2B62F844BCC39.
\end{center}
It evaluates the literal two-stage side-four to side-one
root through degree four on the live helical input
\[
       h_{x,4}=E_1\cos(\pi x_2/2)+E_2\sin(\pi x_2/2),
       \qquad h_{x,1}=h_{x,2}=h_{x,3}=0.
\]
The discrete divergence and all orientation means are
checked to vanish. For a finite algebra calibration only,
use the already exact matrix \(\Gamma_2\), not the limiting
\(\Gamma_*\). Subtract
\(\mathcal A_{2,2}(\Gamma_2)h\) from the cubic coefficient
at the intermediate side two. The change in the terminal
fourth coefficient is exactly the negative mixed quadratic
insertion; its value is
\[
                              (0,0,0,-E_3/32).
\]
Independent polarization of the primitive quadratic gives
the same result, and doubling the test input multiplies
this change by four. The first two coefficients are
unchanged. This checks the degree, sign, geometric filter,
and transport of the actual subtraction; it does not
identify \(\Gamma_2\) with \(\Gamma_*\).

The checker also verifies eighteen exact positive-row
mass, maximum and squared-weight inequalities, including
terminal side-one folding, and 320 exact rational
grouped-dyadic identities for the bounds used above.
These finite tests supplement the analytic covariance
and summation proofs; they do not prove a continuum
limit by sampling cutoffs.

The predecessor V31 has 710,210 bytes and SHA--256
\begin{center}\scriptsize\ttfamily
4B83F3D8315C27B660AE8BFCB7ED5F622983B2FC954078DACC28C9786CFD85E9.
\end{center}
Its body is preserved apart from the title version and
this terminal insertion. V32 replaces it as the single
active main note; archives, monographs and companion
files are preserved. The next regular companion review
remains V35.

\begin{firewall}
The new result is a uniform local fourth-coefficient
bound and a cofinal finite part after an explicit
quadratic random counterterm. It is not a full nonlinear
renormalization trajectory, an interacting continuum
Yang--Mills measure, a noncollapse theorem, or a positive
physical Hamiltonian mass gap. The complete objective
remains unproved and unchanged.
\end{firewall}
\section{V33: fixed continuum counterterms and their finite-lattice compressions}
\label{sec:v33}

\subsection{Removing the rate qualification without changing the objective}

V32 proves a finite part after subtracting the actual
cutoff-dependent quadratic field
\(j\mathcal H^{\rm loc}_{L,M}\), where \(L=pM\) and
\(j=\log_2p\). It explicitly leaves open whether the limiting
field can replace that finite field when multiplied by \(j\).
This section supplies the missing strong-coupling rate.
It also treats the linear counterterm in the local cubic.
The proof uses the actual tent/box filters and the exact
positive outer rows; no new density or stochastic independence
between shells is introduced.

A distinction matters for subsequent renormalization.
A continuum Gaussian observable need not be measurable with
respect to the finite fine field. We first prove convergence
using the fixed continuum subtraction, then compress it by
conditional expectation onto the complete finite fine field.
The compressed subtraction is a genuine finite-dimensional
polynomial and has the same limiting finite part. This is not
conditioning only on the retained coarse field, which would
lose the V27 hidden channels.

All statements below concern source coefficients under the
declared free reference. The ultimate target remains a
nontrivial four-dimensional interacting quantum theory and a
positive physical Hamiltonian gap, as in the
\href{https://www.claymath.org/wp-content/uploads/2022/06/yangmills.pdf}
{Jaffe--Witten problem description}. Neither that construction
nor that gap is implied by a four-coefficient calculation.
The current Shell V16, SM V72 and NS V26 files are unchanged;
the next regular cross-program checkpoint is V35.

\subsection{Fixing a quantitative real Gaussian coupling}

We now specify the previously free high-frequency choices in
Lemma~\ref{lem:v13-isometry}. This changes no finite Gaussian
law and no limiting continuum kernel. Rates in this section
are asserted for this specified coupling, not for arbitrary
isometries which merely agree eventually at each fixed mode.

On a magnified torus of side \(R\ge1\), let the fine ratio
be \(t\), so the fine side is \(tR\). The field \(tZ_{tR}\)
has the nonzero Fourier covariance symbol
\(\Pi_t(\xi)/\omega_t(\xi)^2\), with normalization \(R^{-4}\),
where
\[
 \xi\in(2\pi/R)\mathbb Z^4,\qquad
 q_t(\xi)_\mu=e^{i\xi_\mu/t}-1,\qquad
 \omega_t=t|q_t|,\qquad
 \Pi_t=I-\frac{q_tq_t^*}{|q_t|^2}.
\]
Here \(\xi\) ranges over lattice representatives for the
finite field; the continuum field uses all nonzero modes,
with \(\Pi=I-\xi\xi^t/|\xi|^2\) and multiplier \(|\xi|^{-2}\).
The zero mode is removed in both fields.

Choose an absolute \(c_0>0\) small enough that, whenever
\(0<|\xi|\le c_0t\),
\[
                 \|\Pi_t-\Pi\|\le C|\xi|/t\le1/2.
\]
On these modes take the polar isometry from
\(\operatorname{ran}\Pi_t\) to \(\operatorname{ran}\Pi\):
\begin{equation}\label{eq:v33-polar}
 U_t=\Pi\Pi_t
       \bigl(\Pi_t\Pi\Pi_t|_{\operatorname{ran}\Pi_t}\bigr)^{-1/2}.
\end{equation}
Use conjugate maps on conjugate frequency pairs. Every other
finite conjugacy orbit is mapped isometrically to a distinct
unused continuum orbit of frequency length at least \(c_0t\).
For a self-conjugate Nyquist orbit use normalized real cosine
coordinates, with the energy normalization of V13. There are
infinitely many available high continuum modes. No finite
mode is deleted and no high finite mode is placed in the
reserved low-frequency space.

This construction can be made once at final physical side
\(M\) and used at all intermediate sides \(R=sM\).
The integer Fourier label is unchanged by this dilation,
and \(\xi/t\) is unchanged because \(t=p/s\).
The local Fourier coefficient normalization changes by the
same physical dilation as \(tZ_L=pZ_L/s\).
Thus the linear fields at different intermediate scales
belong to one Gaussian family, not separately chosen couplings.
Three independent copies supply the Lie colours.

\begin{lemma}[Quantitative transverse-plane identification]
\label{lem:v33-polar}
With the preceding choices and \(r=|\xi|\le c_0t\),
\begin{equation}\label{eq:v33-polar-bound}
 \|U_t\Pi_t-\Pi\|\le Cr/t,\qquad
 \left\|\frac{U_t\Pi_t}{\omega_t}-\frac{\Pi}{r}\right\|
                                                   \le C/t.
\end{equation}
\end{lemma}
\begin{proof}
Taylor expansion gives
\(q_t=i\xi/t+O(r^2/t^2)\), uniformly in direction.
The chord bound gives \(\omega_t\ge2r/\pi\), and
\(\omega_t^2=r^2+O(r^4/t^2)\).
Normalizing the vector in the rank-one projection therefore
gives \(\delta:=\|\Pi_t-\Pi\|\le Cr/t\).
For \(v\in\operatorname{ran}\Pi_t\),
\[
 \langle v,\Pi_t\Pi\Pi_t v\rangle
   =\|v\|^2-\|(I-\Pi)v\|^2
   \in[(1-\delta^2)\|v\|^2,\|v\|^2].
\]
Its inverse square root differs from the identity by at most
\(C\delta^2\). Formula~\eqref{eq:v33-polar} now gives
\(\|U_t\Pi_t-\Pi\|\le C\delta\); it also shows directly
that \(U_t\) is isometric and onto the transverse plane.
Finally
\[
                |\omega_t^{-1}-r^{-1}|\le Cr/t^2.
\]
Separate this scalar difference from the projection difference
and use \(r\le c_0t\). This proves the second estimate.
\end{proof}

\subsection{A strong rate for the actual linear averages}

Let \(f_t\) be an unfolded, cell-constant reconstructed kernel
of mesh \(1/t\), and \(f\) its continuum comparison kernel.
Assume fixed support diameter, uniformly bounded
\(L^1,L^2\) norms, and \(\|f_t-f\|_1\le e_t\).
Vector-valued orientation kernels are allowed. Translate
both kernels to the same output point. Let \(X_t(f_t)\)
and \(X_\infty(f)\) denote their Gaussian smearings on the
specified coupling. Periodic overlaps are included.

\begin{lemma}[Strong smearing comparison]
\label{lem:v33-smearing}
For \(1\le T\le c_0t\),
\begin{equation}\label{eq:v33-smearing}
 \mathbb E|X_t(f_t)-X_\infty(f)|^2
          \le C\{e_t^2T^2+t^{-2}T^4+T^{-2}\}.
\end{equation}
The constant is uniform in \(R\ge1\) and the output point.
For the actual unit-scale geometric averaging filters,
including any cross input orientation, this implies
\begin{equation}\label{eq:v33-linear-rate}
 \|Y_{tR,R}-Y_{\infty,R}\|_q
 +\|\mathcal A_{t,R}(A)Z_{tR}
                 -\mathcal A_{\infty,R}(A)\|_q
                              \le C_q(1+\|A\|)t^{-1/3}
\end{equation}
at each link, for every finite \(q\ge1\).
\end{lemma}
\begin{proof}
Write
\(F_t^d(\xi)=t^{-4}\sum_xf_{t,x}e^{-i\xi\cdot x/t}\).
Integrating the phase over each fine cell gives
\[
                 |F_t^d-\widehat f|\le e_t+Cr/t.
\]
Translation contributes the same unit phase on both sides;
it does not enlarge the constant. For low modes the energy
representer of the finite smearing is
\(U_t\Pi_t F_t^d/\omega_t\), whereas that of the continuum
smearing is \(\Pi\widehat f/r\). Transposes or adjoints
according to row conventions give the same norm estimate.
By Lemma~\ref{lem:v33-polar}, their difference has norm at most
\[
                              C(e_t/r+t^{-1}).
\]
The elementary four-dimensional mode counts give, for \(T\ge1\),
\[
 R^{-4}\!\sum_{0<|\xi|\le T}|\xi|^{-2}\le CT^2,
 \qquad
 R^{-4}\#\{0<|\xi|\le T\}\le CT^4.
\]
For example apply dyadic annuli to integer labels up to
\(RT/(2\pi)\); if that radius is less than one the set is
empty. Squaring and summing the representer difference
therefore bounds its low-frequency norm by the first two
terms of \eqref{eq:v33-smearing}.

The remaining finite representer lies in the orthogonal
complement of the reserved continuum modes \(|\xi|\le T\).
Its norm squared is at most
\[
 C T^{-2}R^{-4}\sum_{\text{finite modes}}|F_t^d|^2
                                          \le CT^{-2}\|f_t\|_2^2.
\]
The last constant includes the bounded overlap of unfolded
kernel translates on a torus of side at least one.
The same Parseval estimate bounds the continuum high part
by \(CT^{-2}\|f\|_2^2\). The inequality
\(\|a-b\|^2\le2\|a\|^2+2\|b\|^2\) bounds their difference;
they need not be independent. This argument includes the
finite Nyquist orbits and the omitted continuum frequencies.
It proves \eqref{eq:v33-smearing} by the isonormal isometry.

For a geometric averaging row the unfolded kernel is the
unit transverse box times the longitudinal tent approximation
of V18. Its \(L^1\) error is exactly \(1/t\), and its squared
\(L^2\) error is \(2/(3t^2)\). In particular its \(L^1,L^2\)
norms are bounded. Changing the input orientation does not
change these scalar geometric estimates. Set \(T=t^{1/3}\)
for sufficiently large \(t\); the condition \(T\le c_0t\)
then holds, and \eqref{eq:v33-smearing} is at most
\(Ct^{-2/3}\). The bounded range of smaller \(t\) follows
from the separate uniform smearing variances, increasing \(C\).
Gaussian moment comparison proves the \(L^q\) assertion.
Finite sums of input orientations give the stated matrix bound.
\end{proof}

No large-volume Riemann-sum limit was used here. The comparison
is between two fields on the same torus, with matched low
Fourier coordinates. It is consequently a strong-coupling
estimate, not a difference of two variances inferred from
Lemma~\ref{lem:v22-free-rate}.

\subsection{A quantitative comparison for the logarithmic quadratic field}

Keep the notation of \eqref{eq:v32-Hlocal}, and abbreviate
\[
 D_{\infty,M}=\mathcal A_{\infty,M}(\Gamma_*),\qquad
 H_{L,M}=\mathcal H^{\rm loc}_{L,M},\qquad
 H_{\infty,M}=\mathcal H^{\rm loc}_{\infty,M}.
\]
The abbreviation does not replace the cross-oriented filter
by \(\Gamma_*Y\). Define
\[
                    \rho_p=p^{-1/7}\sqrt{1+\log_2p}.
\]

\begin{theorem}[A rate that survives the logarithmic depth]
\label{thm:v33-rate}
At each output link, uniformly in dyadic \(L=pM\), \(M\ge1\),
\begin{equation}\label{eq:v33-H-rate}
 \|H_{L,M}-H_{\infty,M}\|_q\le C_q\rho_p,
 \qquad
 j\|H_{L,M}-H_{\infty,M}\|_q\longrightarrow0
                                             \quad(p\to\infty).
\end{equation}
Also
\begin{equation}\label{eq:v33-D-rate}
              \|D_{L,M}-D_{\infty,M}\|_q\le C_qp^{-1/3}.
\end{equation}
Here \(D_{L,M}=\mathcal A_{p,M}(\Gamma_*)Z_L\).
Both fields in \eqref{eq:v33-D-rate} vanish when \(M=1\).
The assertions hold jointly on any fixed finite output carrier,
with a constant allowed to depend on its size.
\end{theorem}
\begin{proof}
The linear assertion is Lemma~\ref{lem:v33-smearing}.
At side one, the finite row is constant on all fine sites,
so its value on the mean-zero field is zero by V26.
The periodized continuum box factor is also constant;
equivalently its Fourier multiplier vanishes at every nonzero
integer mode. This proves the terminal assertion.

For \(s=2^\ell\le p\), let \(R=sM\) and \(t=p/s\).
The primitive second derivative is one fixed finite bilinear
stencil, with uniformly bounded coefficients and incidence.
Subtract its two coupled evaluations using
\[
 B(y,d)-B(y',d')=B(y-y',d)+B(y',d-d').
\]
Holder's inequality, uniform Gaussian moments, and
\eqref{eq:v33-linear-rate} give a local \(L^q\) difference
bounded by \(C_qt^{-1/3}\). The exact outer positive row
\(B_{s/2}\) has mass \(s/2\), including terminal folding.
Its triangle bound therefore gives
\begin{equation}\label{eq:v33-shell-difference}
 \|\mathcal H_{L,M;s}-\mathcal H_{\infty,M;s}\|_q
                      \le C_qs(p/s)^{-1/3}
                       =C_qp^{-1/3}s^{4/3}.
\end{equation}
We use this estimate only in a finite head. Separately,
Theorem~\ref{thm:v32-shells} gives each finite or continuum
shell norm at most \(C_q2^{-\ell}\sqrt{1+\ell}\).

For any integer \(0\le b\le j\), split both sums after
\(\ell=b\). Sum \eqref{eq:v33-shell-difference} in the head
and use the two separate shell tails in the remainder:
\begin{equation}\label{eq:v33-head-tail}
 \|H_{L,M}-H_{\infty,M}\|_q
   \le C_q\{p^{-1/3}2^{4b/3}+2^{-b}\sqrt{1+b}\}.
\end{equation}
The geometric head is dominated by a constant times its
last term. The tail follows, for example, from
\(\sqrt{1+b+k}\le\sqrt{1+b}\sqrt{1+k}\) and summability
of \(2^{-k}\sqrt{1+k}\). Omitted continuum shells above the
fine cutoff are included in this tail.

Choose \(b=\lfloor j/7\rfloor\). Then
\[
        p^{-1/3}2^{4b/3}\le p^{-1/7},\qquad
                         2^{-b}\le2p^{-1/7}.
\]
This proves the first assertion in \eqref{eq:v33-H-rate},
including the bounded small depths. Its product with \(j\)
is bounded by \(C_q2^{-j/7}(1+j)^{3/2}\), which tends to
zero and is uniformly bounded. All estimates are per link;
a finite carrier follows by a finite norm comparison.
\end{proof}

\begin{proposition}[The fixed continuum subtraction has the same finite part]
\label{prop:v33-fixed}
At fixed \(M\), on the specified coupling and in every finite
\(L^q\) on a finite output carrier,
\begin{equation}\label{eq:v33-fixed}
 \begin{aligned}
 R_{L,M}-jD_{\infty,M}&\longrightarrow\widehat R_{\infty,M},\\
 \mathcal F_{L\to M,[4]}(Z_L)-jH_{\infty,M}
                                  &\longrightarrow\mathcal V^{\rm fin}_{\infty,M}.
 \end{aligned}
\end{equation}
The right sides are exactly the V26 and V32 finite parts,
with their endpoint terms and their original clock convention.
The left sides have uniformly bounded local finite moments.
\end{proposition}
\begin{proof}
Their differences from the earlier finite-filter subtractions
are respectively \(j(D_{L,M}-D_{\infty,M})\) and
\(j(H_{L,M}-H_{\infty,M})\). Equations
\eqref{eq:v33-D-rate} and \eqref{eq:v33-H-rate} make both
differences tend to zero with uniform local bounds.
Apply Theorems~\ref{thm:v26-local} and \ref{thm:v32-finite}.
At \(M=1\) use the V25 terminal cubic convention.
\end{proof}

In particular the terminal fourth coefficient satisfies
\[
                    \mathcal D_{2^n}-n\mathcal H_\infty
                                          \longrightarrow\mathcal D_{\rm fin}.
\]
Its nonzero quadratic logarithmic coefficient, including the
V31 helical calibration, is unchanged. No rate for convergence
of the finite part itself is inferred: the new rate controls
the change of counterterm, not all the V32 finite-head limits.

\subsection{Finite-field measurable compressions of the continuum counterterms}

Fix \(M\). Write \(\mathcal K_M\) for the real full-colour
isonormal Hilbert space and \(\mathcal S_{L,M}\) for the image
of the entire finite mean-zero transverse Gaussian space in
the specified embedding. Let \(P_{L,M}\) be its orthogonal
projection, and
\[
 \mathcal G_{L,M}=\sigma\{W(v):v\in\mathcal S_{L,M}\}
                                             =\sigma(Z_L).
\]
The equality holds up to Gaussian completion because the
finite covariance is nondegenerate on its live detail space.
Neither nesting of these subspaces nor a pseudoinverse
continuity assumption is needed.

For a component at an output link, represent
\[
                 D_{\infty,M}=W(d_M),\qquad
                 H_{\infty,M}=I_2(T_M),
\]
where \(d_M\in\mathcal K_M\) and \(T_M\) is real symmetric
Hilbert--Schmidt. Such representations exist by first- and
second-chaos isometries; V32 bounds their norms locally.
The kernel \(T_M\) is the convergent quadratic source-shell
kernel, not a newly chosen scalar multiple of the leading root.
Define the finite compressions
\begin{equation}\label{eq:v33-compressions}
 \begin{aligned}
 D^{\rm c}_{L,M}&=W(P_{L,M}d_M),\\
 H^{\rm c}_{L,M}&=I_2(P_{L,M}T_MP_{L,M}).
 \end{aligned}
\end{equation}
They use fixed continuum coefficients and a finite regulator
projection. The superscript \({\rm c}\) means compression.

\begin{lemma}[Conditional Wick compression and best approximation]
\label{lem:v33-compression}
The fields in \eqref{eq:v33-compressions} are respectively
\(\mathbb E[D_{\infty,M}\mid\mathcal G_{L,M}]\) and
\(\mathbb E[H_{\infty,M}\mid\mathcal G_{L,M}]\).
At every output link,
\begin{equation}\label{eq:v33-compression-rate}
 \begin{aligned}
 \|D^{\rm c}_{L,M}-D_{L,M}\|_q
    +\|D^{\rm c}_{L,M}-D_{\infty,M}\|_q&\le C_qp^{-1/3},\\
 \|H^{\rm c}_{L,M}-H_{L,M}\|_q
    +\|H^{\rm c}_{L,M}-H_{\infty,M}\|_q&\le C_q\rho_p.
 \end{aligned}
\end{equation}
\end{lemma}
\begin{proof}
Let \(P=P_{L,M}\). In orthonormal coordinates adapted to
\(\mathcal S_{L,M}\oplus\mathcal S_{L,M}^{\perp}\), the two
Gaussian families are independent. Integrating the second
family kills its linear terms and its centered quadratic
terms, and kills every cross term. Hence, first for finite-rank
kernels,
\[
 \mathbb E[W(d)\mid\mathcal G]=W(Pd),\qquad
 \mathbb E[I_2(T)\mid\mathcal G]=I_2(PTP).
\]
For Hilbert--Schmidt kernels pass in \(L^2\), using
\(\mathbb E|I_2(T)|^2=2\|T\|_{\mathfrak S_2}^2\) and
contractivity of conditional expectation. There is no trace
of the full infinite-rank \(T\) in this argument. In finite
coordinates the compressed polynomial is
\(x^t(PTP)x-\Tr(PTP)\). Merely setting the hidden coordinates
to zero in an uncentered polynomial would miss the conditional
trace correction.

Conditional expectation is an orthogonal projection in
Gaussian \(L^2\). Since \(D_{L,M}\) and \(H_{L,M}\) are
\(\mathcal G_{L,M}\)-measurable, it follows that
\[
 \begin{aligned}
 \|H_\infty-H^{\rm c}_L\|_2&\le\|H_\infty-H_L\|_2,\\
 \|H^{\rm c}_L-H_L\|_2
     &=\|\mathbb E[H_\infty-H_L\mid\mathcal G]\|_2
                                      \le\|H_\infty-H_L\|_2.
 \end{aligned}
\]
Use the same inequalities for \(D\). All quadratic differences
are pure second chaoses; all linear differences are first
chaoses. Hypercontractivity, or the Gaussian moment formula
in degree one, upgrades these estimates to any finite \(q\).
For \(q<2\) use monotonicity. Theorem~\ref{thm:v33-rate}
proves \eqref{eq:v33-compression-rate}.
\end{proof}

\begin{theorem}[Finite-lattice subtractions from fixed continuum kernels]
\label{thm:v33-compressed-finite}
The actual finite polynomials
\begin{equation}\label{eq:v33-compressed-finite}
 \begin{aligned}
 R^{\rm c}_{L,M}&=R_{L,M}-jD^{\rm c}_{L,M},\\
 V^{\rm c}_{L,M}&=\mathcal F_{L\to M,[4]}(Z_L)-jH^{\rm c}_{L,M}
 \end{aligned}
\end{equation}
have uniformly bounded local finite moments. At fixed \(M\)
they converge to \(\widehat R_{\infty,M}\) and
\(\mathcal V^{\rm fin}_{\infty,M}\), respectively, in every
finite \(L^q\) on every finite output carrier. Their differences
from the V26 and V32 subtracted polynomials satisfy
\begin{equation}\label{eq:v33-scheme-error}
 \|R^{\rm c}_{L,M}-\widehat R_{L,M}\|_q\le C_qjp^{-1/3},
 \qquad
 \|V^{\rm c}_{L,M}-\widehat{\mathcal V}_{L,M}\|_q\le C_qj\rho_p.
\end{equation}
The same local stretched-exponential orders as in V26 and
V32 hold: exponents \(2/3\) and \(1/2\), respectively.
\end{theorem}
\begin{proof}
Subtract the definitions. The differences are
\[
 j(D_{L,M}-D^{\rm c}_{L,M}),\qquad
 j(H_{L,M}-H^{\rm c}_{L,M}).
\]
Lemma~\ref{lem:v33-compression} proves
\eqref{eq:v33-scheme-error}; both right sides vanish and are
uniformly bounded. The earlier finite-part theorems give
the limits. The new cubic and quartic polynomials have
Gaussian degrees at most three and four and uniformly bounded
\(L^2\) norms. Hypercontractivity and the exponential-series
argument of V32 give the stated stretched-exponential orders.
\end{proof}

The finite-dimensional coefficients in these compressed
polynomials are well defined by inner products of fixed
continuum kernels with finitely many energy basis vectors.
They are not asserted to be compactly supported in position,
computable by a finite local stencil, or generated by the
actual nonlinear shell map. Their construction solves the
finite-field measurability issue, not those locality issues.
It also does not identify compressions at different regulators
as a martingale: the subspaces were not required to be nested.

\subsection{What is now available for the nonlinear problem}

We may use either the actual finite filters of V32 or the
compressed fixed kernels of \eqref{eq:v33-compressions} in
the controlled four-term source polynomial. The two conventions
have the same finite-part limits with the explicit error
\eqref{eq:v33-scheme-error}. Thus a lack of a rate for replacing
these two logarithmic coefficients is no longer a bottleneck.
This is a comparison of the stated conventions on the stated
coupling, not universality under arbitrary lattice regulators.

The remaining upstream task is still nonlinear and measured:
construct a scale-local source and density transport that
propagates the non-scalar fine-field counterterms, preserves
the actual Jacobian and conditional law, and controls higher
orders when the incoming depth grows with inverse noise
variance. Four convergent renormalized coefficients do not
bound the full source remainder on typical quantum fields.
In particular the compression theorem supplies neither a
uniform analytic radius nor uniform integrability of the
full interacting exponential. The V25 cubic exponential
obstruction and the V27 retained-field information loss remain
in force; neither is removed by a change of subtraction notation.

\subsection{Exact checks and preservation record}

The new certificate is
\begin{center}\small
\path{scripts/verify_ym_f14_v33_continuum_counterterms.py}
\end{center}
with SHA--256
\begin{center}\scriptsize\ttfamily
1A2D572098CF046FD32F8E5A21417B87EBAC940D348389CA76299575BDC8152C.
\end{center}
It checks 64 exact tent approximations, including their
\(L^1\) errors, squared \(L^2\) errors, masses and squared
norms. It verifies four rational transverse-plane polar
calibrations and the full real Fourier dimension counts at
sides two, four and eight, including all 15 nonzero real
Nyquist orbits at each side.

Thirty-six exact Gaussian conditional evaluations use two
non-coordinate, non-nested rational orthogonal projections.
They verify the conditional Wick trace term and the first-
and second-chaos best-approximation identities; the omitted
trace term is explicitly nonzero in these tests. Another
257 rational head/tail checks and 64 exact Fourier cutoff
balances check the exponent bookkeeping. These finite
identities supplement the analytic proofs, not replace them
with a finite numerical sample of the continuum limit.

The predecessor V32 has 731,941 bytes and SHA--256
\begin{center}\scriptsize\ttfamily
D0962258FEFB49985EBE89B0B3178FCE8F8A318D19BEDC45E00036348F20B7F0.
\end{center}
Its body is preserved apart from the title version and this
terminal insertion. V33 replaces it as the single active main
note. Archives, monographs and the three companions are
preserved. The next regular companion review remains V35.

\begin{firewall}
The new results are quantitative strong comparisons for the
two actual logarithmic source counterterms, fixed-continuum
finite-part formulas, and finite-field measurable compressed
subtractions. They do not construct a full nonlinear
renormalization trajectory, an interacting continuum measure,
or a positive physical Yang--Mills mass gap. The complete
objective remains open in this programme.
\end{firewall}
\section{V34: counterterm transport, the gauge correction, and its retained harmonic part}
\label{sec:v34}

\subsection{The upstream issue after the finite-part estimates}

V26 and V31--V33 identify and control the linear and quadratic
source subtractions. A remaining structural question is how
these two subtractions belong to one change of fine variables,
and whether that change is compatible with the gauge slice
used for the measured integral. This section answers that
question through the fourth source degree.

The two counterterms are the first two coefficients of the
tangent response to a single fine orientation matrix. They
obey an exact augmented coarse recursion. However, the matrix
does not preserve the mean-zero Coulomb slice. Simply projecting
it onto that slice loses an explicitly nonzero logarithmic
fourth-order term. We derive the compensating nonlinear gauge
term and prove that its harmonic, spatially constant component
restores the lost contribution. Dropping that component would
change the retained source.

This is a source and gauge calculation, not a replacement for
Shell V8's already established finite measured composition.
The archived generic nonlinear-completion and fibre-curvature
results were checked; the new statements concern the actual
\(\Gamma_*\) counterterm and the actual balanced root. No
archived result is claimed anew. The three active companions
are unchanged; their next regular review is V35.

\subsection{Both counterterms are one tangent response}

For a fine periodic link array \(z\), let
\(\mathcal F_{L\to M}\) denote the actual nested balanced-root
germ at zero, and write \(F_r=\mathcal F_{L\to M,[r]}\).
These finite-dimensional germs are defined for ambient link
arrays near zero, not just for their Coulomb subspace.
In this subsection \(A\) is any fixed real orientation matrix,
acting at each fine site and identically in Lie colour. Set
\[
 D^A_{L,M}(z)=B_{L/M}(Az),\qquad
 H^A_{L,M}(z)=DF_2(z)[Az]
                =D^2\mathcal F_{L\to M}(0)[z,Az].
\]
The matrix acts before the averaging. Thus
\(D^A_{L,M}=\mathcal A_{L/M,M}(A)z\), not in general
\(A B_{L/M}z\).

\begin{proposition}[Tangent identification and common input subtraction]
\label{prop:v34-tangent}
For \(A=\Gamma_*\) and \(z=Z_L\), the preceding fields are
exactly \(D_{L,M}\) and \(H_{L,M}=\mathcal H^{\rm loc}_{L,M}\)
of V31--V33. For any fixed finite regulator and real \(c\),
\begin{equation}\label{eq:v34-pullback}
 \begin{aligned}
 \mathcal F_{L\to M}(\epsilon z-c\epsilon^3Az)
   ={}&\epsilon F_1(z)+\epsilon^2F_2(z)\\
     &+\epsilon^3\{F_3(z)-cD^A_{L,M}(z)\}\\
     &+\epsilon^4\{F_4(z)-cH^A_{L,M}(z)\}+O(\epsilon^5).
 \end{aligned}
\end{equation}
Here and below a Taylor remainder for fixed regulator and
deterministic input is not a uniform probabilistic estimate.
Taking \(c=\log_2(L/M)\) gives precisely the actual finite-filter
subtractions of V26 and V32 through degree four.
\end{proposition}
\begin{proof}
The actual linear blocks compose, including their geometric
orientation, by the identity already used in V27. Hence
\(D^A_{L,M}=B_{L/M}(Az)\) agrees with the cross-oriented filter.
Differentiate the quadratic composition formula at \(z\) in
direction \(Az\). At an incoming side \(R\), the transported
directions are \(Y_{L,R}=B_{L/R}z\) and
\(D^A_{L,R}=B_{L/R}(Az)\). The new quadratic response is
\(D^2E_R(0)[Y_{L,R},D^A_{L,R}]\); the earlier response is
transported by the linear block. Unrolling gives exactly
the shell sum \eqref{eq:v32-Hlocal} when \(A=\Gamma_*\).

Alternatively the tangent germ is
\[
 D\mathcal F_{L\to M}(\epsilon z)[Az]
                       =D^A_{L,M}+\epsilon H^A_{L,M}+O(\epsilon^2).
\]
Inserting \(-c\epsilon^3Az\) into the Taylor series of the
source gives \eqref{eq:v34-pullback}: two such insertions first
have degree six. This also fixes the sign and the absence of
an extra factor two in the quadratic subtraction.
\end{proof}

\subsection{A closed augmented recursion through four source degrees}

Let \(M_k=L/2^k\), and use \(Y_k,Q_k,R_k,V_k\) for the first
four coefficients after \(k\) actual steps. Put
\[
 D_k=B_{2^k}(Az),\quad H_k=DF_{2,k}(z)[Az],\quad
 \widehat R_k=R_k-kD_k,\quad \widehat V_k=V_k-kH_k.
\]
At a step from side \(M_k\), write \(B=B_2\),
\(E^{[r]}=E_{M_k}^{[r]}\), and
\(K(a,b)=D^2E_{M_k}(0)[a,b]\). The dependence of the actual
periodic stencil on its side is retained.

\begin{theorem}[Exact augmented source recursion]
\label{thm:v34-recursion}
Starting with
\[
 (Y_0,Q_0,D_0,H_0,\widehat R_0,\widehat V_0)
                                      =(z,0,Az,0,0,0),
\]
the following recursion is exact:
\begin{equation}\label{eq:v34-recursion}
 \begin{aligned}
 Y_{k+1}&=BY_k,\\
 Q_{k+1}&=BQ_k+E^{[2]}(Y_k),\\
 D_{k+1}&=BD_k,\\
 H_{k+1}&=BH_k+K(Y_k,D_k),\\
 \widehat R_{k+1}
   &=B\widehat R_k+K(Y_k,Q_k)+E^{[3]}(Y_k)-D_{k+1},\\
 \widehat V_{k+1}
   &=B\widehat V_k+K(Y_k,\widehat R_k)+E^{[2]}(Q_k)\\
   &\quad+\tfrac12D^3E_{M_k}(0)[Y_k,Y_k,Q_k]
                            +E^{[4]}(Y_k)-H_{k+1}.
 \end{aligned}
\end{equation}
For \(A=\Gamma_*\) under the free reference, all six fields
have uniformly bounded local finite moments, by the preceding
coefficient estimates. This is a recursion of polynomial
observables on the common fine field, not closure of their
joint law in a finite-dimensional family of densities.
\end{theorem}
\begin{proof}
The first four identities follow from linear composition and
the quadratic derivative just proved. The raw cubic recursion
is
\(R_{k+1}=BR_k+K(Y_k,Q_k)+E^{[3]}(Y_k)\).
Substitute \(R_k=\widehat R_k+kD_k\) and subtract
\((k+1)D_{k+1}\); since \(D_{k+1}=BD_k\), this gives the
fifth line. The raw fourth recursion is the literal chain
rule of Lemma~\ref{lem:v31-chain}. Its terms containing the
clock are
\[
                         kBH_k+kK(Y_k,D_k)=kH_{k+1}.
\]
Subtracting \((k+1)H_{k+1}\) leaves the last line of
\eqref{eq:v34-recursion}. Thus the previous accumulated clock
does not have to be reinserted into each displayed polynomial.

For the moment assertion, use the linear filter bound of V26,
the quadratic estimate of V25, the cubic finite part of V26,
and the fourth finite part and counterterm bound of V32.
They are estimates at each link and do not bound a counting
norm over growing physical volume. The identities above use
the same fields at all steps; they assert no independence.
\end{proof}

The extra response fields cannot simply be replaced by
functions of \(Y_k\): V27's hidden-channel obstruction
applies already to \(D_k\) for the actual non-diagonal
\(\Gamma_*\). The exact tangent lift of the full analytic
map would transport \(x\) and a direction \(v\) by
\((x,v)\mapsto(E(x),DE(x)v)\). Its existence as a finite
analytic germ does not supply bounds for an arbitrarily
long nonlinear iteration on typical quantum inputs.

\subsection{Why the unprojected linear change is not a Coulomb change of variables}

Let \(\Pi_L\) denote the orthogonal projector onto the
mean-zero Coulomb link space \(V_L\). Let \(P_{0,L}\) be the
orientation-wise spatial-mean projector, and write
\(P_{{\rm C},L}=\Pi_L+P_{0,L}\) for the Coulomb projector
which retains constants. The matrix \(A\) commutes with
spatial means, but in general it does not preserve \(V_L\).

For the actual \(\Gamma_*\), take a nonzero single-axis mode
in input orientation four, with wave direction two. The
output orientation two of \(\Gamma_*z\) is \(a_*z_4\),
where \(a_*>0\) by V28. It is longitudinal, so
\[
                   (I-\Pi_L)\Gamma_*z\ne0.
\]
This is a live transverse input, not a direction removed by
the original covariance. For \(L\ge4\) one may use the real
helical mode employed below.

\begin{proposition}[The slice-support issue is not repaired by a determinant]
\label{prop:v34-singular}
At each dyadic \(L\ge4\), for sufficiently small nonzero real \(\tau\), both
\(I-\tau\Gamma_*\) and \(e^{-\tau\Gamma_*}\) are invertible
on ambient link arrays but send \(V_L\) to a different
linear subspace. The pushforward of the nondegenerate
Gaussian on \(V_L\) by either map is mutually singular
with that Gaussian when both are viewed on the ambient
link space. A Radon--Nikodym determinant on the unchanged
Coulomb slice is therefore not a valid change-of-variables
formula for either unprojected map.
\end{proposition}
\begin{proof}
On the displayed single-axis input the longitudinal part of
\((I-\tau\Gamma_*)z\) is \(-\tau a_*z_4\) in orientation
two. For the exponential it is \(-\tau a_*z_4+O(\tau^2)\),
so it too is nonzero for all sufficiently small nonzero
\(\tau\). Invertibility of the first map holds for
\(|\tau|\|\Gamma_*\|<1\); the second is always invertible.
Their image subspaces have the same dimension as \(V_L\)
and differ from it, so their intersections with \(V_L\)
are proper subspaces of both. A nondegenerate Gaussian on
a finite-dimensional space gives every proper linear
subspace probability zero. Taking the complement of the
intersection in each support proves mutual singularity.
This statement does not preclude a change of gauge chart,
or a transformation of the full compact measure with all
its factors. It precludes the specified same-slice shortcut.
\end{proof}

\subsection{An exact terminal gauge identity for the quadratic source}

In the remainder of this subsection the final side is one.
Write \(F=\mathcal F_{L\to1}\), \(Q=F_{[2]}\), and
\(B_L=DF(0)\). For a vertex Lie field \(\psi\), set
\[
 (d\psi)_{x,\mu}=\psi_{x+e_\mu}-\psi_x,\qquad
 C_\psi(z)_{x,\mu}
                  =-\tfrac12[\psi_x+\psi_{x+e_\mu},z_{x,\mu}].
\]
The root's terminal linear row is exactly
\(B_Lw=L\overline w\), with the mean taken separately in
each orientation and colour.

\begin{lemma}[Gauge covariance of the actual balanced-root germ]
\label{lem:v34-gauge}
For fine gauge transformations
\(U_{x,\mu}\mapsto g_xU_{x,\mu}g_{x+e_\mu}^{-1}\), one
actual primitive block transforms by a coarse gauge
transformation, which can depend on both \(U\) and \(g\).
Consequently, at final side one,
\[
                     F(a^g)=\operatorname{Ad}_{G(a,g)}F(a)
\]
as analytic germs near the identity, with \(G(a,1)=1\).
\end{lemma}
\begin{proof}
Use the actual tree paths \(T_s\) from the primitive base
\(2r\) to \(2r+s\). A transported two-link path \(P_s\)
from \(2r\) to \(2r+2e_\mu\) transforms as
\[
                P_s' = g_{2r}P_sg_{2r+2e_\mu}^{-1}.
\]
The relative logarithms of \(P_0^{-1}P_s\) therefore
transform by conjugation at the common endpoint. Their
actual anchored exponential average, denoted \(R_{r,\mu}\),
transforms as
\(R_{r,\mu}'=g_{2r}R_{r,\mu}g_{2r+2e_\mu}^{-1}\).
All logarithms here are the same local analytic branches
as in the source germ.

The balanced block is
\(E(U)_{r,\mu}=e^{-\Phi_r}R_{r,\mu}e^{\Phi_{r+e_\mu}}\),
where \(\Phi_r\) is the mean of the actual tree logarithms.
Define
\[
                    G_r=e^{-\Phi'_r}g_{2r}e^{\Phi_r}.
\]
Then \(E(U^g)_{r,\mu}=G_rE(U)_{r,\mu}G_{r+e_\mu}^{-1}\).
No assertion that the individual mean tree logarithm itself
is gauge covariant was used. Iterate this exact identity.
On the final one-site lattice all coarse links are loops,
so the final transformation is conjugation by one group
element. Taking the local logarithm proves the assertion.
\end{proof}

\begin{theorem}[Terminal mixed gauge Ward identity]
\label{thm:v34-Ward}
For any mean-zero fine link field \(z\) and any vertex field
\(\psi\),
\begin{equation}\label{eq:v34-Ward}
                DQ(z)[d\psi]+B_L C_\psi(z)=0.
\end{equation}
This is an identity of actual finite polynomials, before
Gaussian integration or a continuum limit.
\end{theorem}
\begin{proof}
Apply the gauge transformation \(g_x=e^{-u\psi_x}\) to
the link field \(a=sz\). The coefficient of \(su\) in
the logarithm of
\(e^{-u\psi_x}e^{sz_{x,\mu}}e^{u\psi_{x+e_\mu}}\) is
\(C_\psi(z)_{x,\mu}\). Its linear terms are
\(sz+u\,d\psi\). Thus the coefficient of \(su\) in the
left side of Lemma~\ref{lem:v34-gauge} is
\(B_LC_\psi(z)+D^2F(0)[z,d\psi]\).

On the right, \(G(sz,e^{-u\psi})=1+O(u)\), whereas
\(F(sz)=sB_Lz+O(s^2)=O(s^2)\), since \(z\) is mean-zero.
Its \(su\) coefficient is zero. Finally
\(DQ(z)[d\psi]=D^2F(0)[z,d\psi]\), proving
\eqref{eq:v34-Ward} with the displayed sign.
\end{proof}

For \(z\in V_L\), define the unique mean-zero potential
\(\psi_A(z)\) by
\[
                        d\psi_A(z)=(I-\Pi_L)Az.
\]
It exists by the discrete Hodge decomposition, since \(Az\)
is mean-zero. Write \(C_A(z)=C_{\psi_A(z)}(z)\), and
\[
 H_L^A=DQ(z)[Az],\qquad H_L^{A,\Pi}=DQ(z)[\Pi_LAz].
\]
The Ward identity gives the useful exact formula
\begin{equation}\label{eq:v34-defect}
             H_L^A-H_L^{A,\Pi}=-B_L C_A(z)=-L\overline{C_A(z)}.
\end{equation}
Thus the effect of the removed longitudinal seed is a
retained harmonic contribution, not zero.

\subsection{Projecting the actual counterterm leaves a logarithmic divergence}

Let \(L=2^n\ge4\), \(\theta=2\pi/L\), and take the normalized
live Cameron--Martin field
\[
 (u_L)_{x,4}=L^{-1}\{E_1\cos(\theta x_2)+E_2\sin(\theta x_2)\},
 \qquad (u_L)_{x,1}=(u_L)_{x,2}=(u_L)_{x,3}=0.
\]
For \(A=\Gamma_*\), its longitudinal seed is exactly
\(a_*(u_L)_{x,4}\) in orientation two and zero in the other
orientations. With \(v=E_1-iE_2\), the potential is
\[
 \psi_A(u_L)_x=\frac{a_*}{L}
             \operatorname{Re}\frac{v e^{i\theta x_2}}{e^{i\theta}-1}.
\]
This formula uses the forward difference just specified.

\begin{theorem}[An explicit nonzero projection defect]
\label{thm:v34-projection-defect}
For the fourth-orientation, third-colour component,
\begin{equation}\label{eq:v34-helical-defect}
 \bigl(H_L^{\Gamma_*}-H_L^{\Gamma_*,\Pi}\bigr)[u_L,u_L]_4^{,3}
                  =\frac{a_*}{L}\cot(\pi/L)\longrightarrow\frac{a_*}{\pi}.
\end{equation}
Under the actual finite free Gaussian law,
\begin{equation}\label{eq:v34-projection-lower}
 \liminf_{L\to\infty}
       \|(H_L^{\Gamma_*}-H_L^{\Gamma_*,\Pi})_4^{,3}\|_2
                         \ge\frac{a_*}{\sqrt8\,\pi^3}>0.
\end{equation}
In particular, if \(\mathcal D_L=\mathcal F_{L\to1,[4]}(Z_L)\),
\begin{equation}\label{eq:v34-wrong-subtraction}
 \liminf_{n\to\infty}\frac1n
       \|(\mathcal D_{2^n}-nH_{2^n}^{\Gamma_*,\Pi})_4^{,3}\|_2
                         \ge\frac{a_*}{\sqrt8\,\pi^3}.
\end{equation}
So projecting the seed and making no compensating change
does not retain the V32 finite part.
\end{theorem}
\begin{proof}
The identity
\[
            (e^{i\theta}-1)^{-1}
                       =-\tfrac12-\tfrac i2\cot(\theta/2)
\]
and \([\operatorname{Re}(vA),\operatorname{Re}(vB)]
=-2\operatorname{Im}(A\overline B)E_3\) give
\[
       C_A(u_L)_{x,4}=-[\psi_A(u_L)_x,(u_L)_{x,4}]
                         =-\frac{a_*}{L^2}\cot(\pi/L)E_3.
\]
All other orientations of \(C_A(u_L)\) vanish. Applying
\eqref{eq:v34-defect} proves \eqref{eq:v34-helical-defect}
directly at the finite regulator. This is an independent
finite Ward calibration of the limiting helical value; it
does not replace finite filters by continuum filters.

Both quadratic fields are centered under the Gaussian law:
they are Lie-equivariant and the colour covariance is scalar.
Thus their difference is a pure second chaos. For a scalar
component write it as \(I_2(T_L)\) in energy coordinates.
For any Cameron--Martin vector \(w\),
\[
 |\langle w,T_Lw\rangle|
     \le\|T_L\|_{\mathfrak S_2}\|w\|_{\rm CM}^2
     =2^{-1/2}\|I_2(T_L)\|_2\|w\|_{\rm CM}^2.
\]
Directly from the discrete curl of the helical field,
\(\|u_L\|_{\rm CM}^2=4L^2\sin^2(\pi/L)\to4\pi^2\).
Combine this with \eqref{eq:v34-helical-defect} to obtain
\eqref{eq:v34-projection-lower}. No convergence theorem for
the projected counterterm is needed for this lower bound.

Finally V32 bounds \(\|\mathcal D_L-nH_L^{\Gamma_*}\|_2\)
uniformly. Decompose the incorrectly subtracted coefficient
as that bounded finite part plus
\(n(H_L^{\Gamma_*}-H_L^{\Gamma_*,\Pi})\), and use the reverse
triangle inequality. This proves
\eqref{eq:v34-wrong-subtraction}.
\end{proof}

\subsection{The gauge-completed four-jet retains a harmonic source shift}

The preceding obstruction does not rule out a genuine gauge
correction. It identifies the term that such a correction
must keep. In the next theorem all conclusions concern the
terminal root, not equality of untransformed individual
coarse links at an intermediate side.

\begin{theorem}[Explicit Coulomb completion through fourth order]
\label{thm:v34-gauge-completion}
Fix \(L\), \(z\in V_L\), \(A\), and a scalar \(c\).
The fine input jet
\begin{equation}\label{eq:v34-completed-input}
 a^{\rm C}(\epsilon)
      =\epsilon z-c\epsilon^3\Pi_LAz
                          +c\epsilon^4P_{{\rm C},L}C_A(z)
\end{equation}
is Coulomb through its displayed order. Its actual terminal
root has the same coefficients through degree four as
\(F(\epsilon z-c\epsilon^3Az)\). Its harmonic component is
\begin{equation}\label{eq:v34-harmonic}
            \overline{a^{\rm C}(\epsilon)}
              =c\epsilon^4\overline{C_A(z)}
              =-\frac{c\epsilon^4}{L}
                               (H_L^A-H_L^{A,\Pi}).
\end{equation}
Consequently the harmonic part generally cannot be deleted
while preserving the counterterm's terminal source jet.
\end{theorem}
\begin{proof}
Gauge-transform \(a=\epsilon z-c\epsilon^3Az\) by
\(g_x=\exp\{-c\epsilon^3\psi_A(z)_x\}\). The same BCH
calculation used in the Ward identity gives
\[
        a^g=\epsilon z-c\epsilon^3\Pi_LAz
                                    +c\epsilon^4C_A(z)+O(\epsilon^5).
\]
Choose a mean-zero vertex field \(\chi\) with
\(d\chi=(I-P_{{\rm C},L})C_A(z)\), and apply a second
gauge transformation \(\exp(c\epsilon^4\chi_x)\).
It subtracts \(c\epsilon^4d\chi\) through degree four;
its brackets with the leading field begin at degree five.
The resulting jet is \eqref{eq:v34-completed-input}, and
its displayed coefficients all have zero divergence.

For the first gauge transformation the terminal coarse
group element is \(1+O(\epsilon^3)\), while the terminal
root before transformation is \(O(\epsilon^2)\), since
its fine linear input has zero orientation means. Their
conjugation changes the source first at degree five.
The second transformation changes it no earlier.
Lemma~\ref{lem:v34-gauge} therefore proves equality through
degree four.

One can also check the fourth coefficient without the group
argument. It is
\[
 F_4(z)-cH_L^{A,\Pi}+cB_LP_{{\rm C},L}C_A(z)
                  =F_4(z)-cH_L^{A,\Pi}+cB_LC_A(z)
                  =F_4(z)-cH_L^A,
\]
by the Ward identity and the fact that \(B_L\) kills gradients.
The mean of \(P_{{\rm C},L}C_A\) is the mean of \(C_A\),
and \eqref{eq:v34-defect} proves \eqref{eq:v34-harmonic}.
If \(P_{{\rm C},L}\) were replaced by \(\Pi_L\) in the
fourth input coefficient, its terminal linear contribution
would be zero. The projection defect would then remain.
\end{proof}

With \(A=\Gamma_*\) and \(c=n\), the theorem supplies the
actual renormalized terminal fourth coefficient in a Coulomb
input jet that includes the retained constant sector. It is
not a change of variables confined to the old mean-zero detail
fibre. A conditional integral must transport this source shift,
the gauge chart, and its density factors together. No uniform
bound for the gauge potentials or for the full transformed
conditional density is claimed here.

\begin{proposition}[The retained harmonic shift has a nonzero cofinal law]
\label{prop:v34-harmonic-limit}
For \(A=\Gamma_*\), the terminal fields
\(H_L^{A,\Pi}\) converge on the specified common Gaussian
coupling in every finite \(L^q\). Consequently
\begin{equation}\label{eq:v34-harmonic-limit}
 J_L:=L\overline{C_A(Z_L)}=H_L^{A,\Pi}-H_L^A
                       \longrightarrow J_\infty.
\end{equation}
These second-chaos fields have a uniform small ordinary
exponential moment. Their limit satisfies
\(\|(J_\infty)_4^{,3}\|_2\ge a_*/(\sqrt8\,\pi^3)\).
For \(L=2^n\), \(\epsilon=\beta^{-1/2}\), and \(c=n\)
in \eqref{eq:v34-completed-input}, if \(n/\beta\to\tau<\infty\),
then its retained harmonic term obeys
\begin{equation}\label{eq:v34-harmonic-scaling}
       \beta L\overline{a^{\rm C}(\epsilon)}
                                  =\frac n\beta J_L\longrightarrow\tau J_\infty
\end{equation}
in every finite \(L^q\). This is an identity and limit for
the displayed input polynomial, not for an uncontrolled
nonlinear gauge remainder.
\end{proposition}
\begin{proof}
Let \(\mathbb Q_{L,\alpha}\) be the symmetric real energy
kernel of the scalar quadratic-root component \(Q_{L,\alpha}\).
V23 gives uniform Hilbert--Schmidt bounds and convergence
\(\mathbb Q_{L,\alpha}\to\mathbb Q_{\infty,\alpha}\) in that
ideal. Let \(J_L^{\rm en}\) denote the V33 real energy
isometry, and extend
\[
 \mathbb A_L=J_L^{\rm en}(\Pi_LA|_{V_L})(J_L^{\rm en})^*
\]
by zero off its finite image. The scalar lattice Laplacian
commutes with the orientation matrix and with the transverse
projection. The Coulomb Hodge energy identity therefore gives
\(\|\mathbb A_L\|\le\|A\|\).
On every fixed finite Fourier family, the projection and
the specified polar identification converge, so
\(\mathbb A_L\to\mathbb A_\infty\) strongly, where
\(\mathbb A_\infty\) is the energy realization of \(A\)
followed by the continuum transverse projection.
Apply the same argument to \(A^t\) to get strong convergence
of the adjoints. Density of finite Fourier families and the
uniform operator bound justify both strong limits on the
whole energy space.

The projected response component has the symmetric kernel
\[
 \mathbb Q_{L,\alpha}\mathbb A_L
                        +\mathbb A_L^*\mathbb Q_{L,\alpha}.
\]
The associated homogeneous finite quadratic has zero Gaussian
mean by colour equivariance; it is the second Wiener integral
of this kernel. Products of a fixed Hilbert--Schmidt operator
with a uniformly bounded strongly convergent family converge
in Hilbert--Schmidt norm on the left; on the right use strong
convergence of the adjoints. To see this, first approximate
the fixed operator by a finite-rank operator, pass on its
finite-dimensional range, and bound the remaining
Hilbert--Schmidt tail by the uniform operator norm.
Together with the Hilbert--Schmidt convergence of
\(\mathbb Q_{L,\alpha}\), this proves convergence and uniform
bounds for the displayed response kernel.

The second-chaos isometry and hypercontractivity now prove
the assertion for \(H_L^{A,\Pi}\). V32 proves the same for
\(H_L^A\); subtract to obtain \eqref{eq:v34-harmonic-limit}.
The usual exponential-series bound for a uniformly bounded
second chaos gives a common small exponential radius.
The lower bound follows from \eqref{eq:v34-projection-lower}
and the new \(L^2\) convergence. Finally substitute
\(c=n\), \(\epsilon^4=\beta^{-2}\) in
\eqref{eq:v34-harmonic}; this yields
\eqref{eq:v34-harmonic-scaling} exactly before passage to
the limit.
\end{proof}

Thus this harmonic shift survives the same \(\beta\)
normalization as the leading quadratic terminal source when
\(n/\beta\) has a positive limit. Its fine mean has a factor
\(L^{-1}\), but treating that factor alone as a reason to
discard it would miss the retained physical normalization.

\subsection{Four terms do not select a nonlinear resummation}

Even before gauge completion, the two analytic input germs
\[
       \epsilon(I-c\epsilon^2A)z,\qquad
                              \epsilon e^{-c\epsilon^2A}z
\]
give exactly the same four source coefficients. Their
difference begins with \(c^2\epsilon^5A^2z/2\).
For the actual finite source, Taylor expansion gives
\begin{equation}\label{eq:v34-completion-difference}
 \begin{aligned}
 &\mathcal F_{L\to M}(\epsilon e^{-c\epsilon^2A}z)
       -\mathcal F_{L\to M}(\epsilon(I-c\epsilon^2A)z)\\
 &\quad=\tfrac12c^2\epsilon^5B_{L/M}(A^2z)
       +\tfrac12c^2\epsilon^6DF_2(z)[A^2z]+O(\epsilon^7).
 \end{aligned}
\end{equation}
Indeed the derivative of the source at the common leading
input is \(B_{L/M}+\epsilon D^2\mathcal F(0)[z,\,\cdot\,]
+O(\epsilon^2)\), which multiplies the input difference.
The fifth coefficient vanishes at final side one because
\(A^2z\) is mean-zero, but not in general at an intermediate
side. This is a comparison of two proposed input germs,
not a statement that either has the required full nonlinear
remainder bound. For \(c\epsilon^2\) of order one, their
fixed-depth Taylor agreement does not justify interchanging
these germs.

There is also a domain check. For \(|\tau|\le T\),
\(\sigma_{\min}(e^{-\tau A})\ge e^{-T\|A\|}\). Hence
\[
 \|\beta^{-1/2}e^{-\tau A}Z_L\|_{\ell^4}
                       \ge e^{-T\|A\|}\|\beta^{-1/2}Z_L\|_{\ell^4}.
\]
Theorem~\ref{thm:v23-smallball} implies that the probability
of any fixed global classical small ball still tends to zero
when \(\beta=o(L^2)\). A bounded matrix exponential by itself
does not make that classical domain typical. This observation
does not estimate a gauge-completed full field or exclude a
scale-local construction.

\subsection{Checks, new bottleneck, and version preservation}

The exact certificate is
\begin{center}\small
\path{scripts/verify_ym_f14_v34_counterterm_transport.py}
\end{center}
with SHA--256
\begin{center}\scriptsize\ttfamily
47B97D5481161C1CE95B607510132F93A29EC7BB30C5E680025934C5B5486E7A.
\end{center}
It evaluates the literal side-four to side-two to side-one
root through degree four on the live helical input of V32.
For a finite calibration use \(A=\Gamma_2\), not \(\Gamma_*\).
Direct fine-input changes by \(+\epsilon^3Az\) and
\(-2\epsilon^3Az\), independent quadratic polarization,
and the augmented recursion agree exactly. The terminal
quadratic response on the unnormalized helical field is
\((0,0,0,E_3/16)\).

For the side-two live input
\(z=(\Pi_2e_0)E_1+(\Pi_2e_1)E_2\), with fine link indices
as in V25, the exact projection defect for \(\Gamma_2\) is
\[
       DQ(z)[(I-\Pi_2)\Gamma_2z]
                     =\frac1{1769472}(-1525E_3,1837E_3,0,0).
\]
The code solves for the mean-zero vertex potential, verifies
its forward gradient, constructs the nonlinear gauge correction,
and checks \(B_2C\) against the negative of this vector.
It checks zero divergence after the Coulomb projection with
constants retained, and exact equality of all four source
coefficients before and after this gauge completion.

Sixteen rational one-axis calibrations independently check
the potential and bracket identity behind
\eqref{eq:v34-helical-defect}. The actual intermediate linear
block also gives a nonzero fifth-order difference in
\eqref{eq:v34-completion-difference}: for \(L=4,M=2,c=1\)
and the same unnormalized helical field, its orientation-two
value at the origin is \(-3893E_2/589824\). This last test
checks the linear coefficient that determines the fifth-order
difference; it is not a computation of the entire fifth
source coefficient. These finite calibrations supplement
the all-regulator identities and proofs.

The next upstream problem is now more concrete: build a
scale-local measured transformation that retains the generated
harmonic source shift and controls its nonlinear gauge and
Jacobian terms together with higher source responses. The
closed four-degree recursion and gauge completion are useful
inputs to that problem, but neither supplies the missing
conditional-density bounds or an interacting continuum law.

The predecessor V33 has 754,711 bytes and SHA--256
\begin{center}\scriptsize\ttfamily
7A08ED5042E08F8911449240BB3630C594B04FB697755AD47733B85F66D7BE23.
\end{center}
Its body is preserved apart from the title version and this
terminal insertion. V34 replaces it as the sole active main
note. Archives, monographs and companions are preserved.
The next regular companion review is V35.

\begin{firewall}
The new results identify a common fine-input origin of the
two source counterterms, their exact augmented recursion,
the nonzero error of a projection-only correction, and a
gauge-completed fourth-order input with its necessary harmonic
source shift. They do not establish a full nonlinear
renormalization trajectory, an interacting four-dimensional
continuum Yang--Mills theory, or a positive physical Hamiltonian
mass gap. The full objective remains unproved.
\end{firewall}
\section{V35: the constrained harmonic pivot and its noisy source-coordinate Jacobian}
\label{sec:v35}

\subsection{The fifth-version review and the next measured variable}

V34 shows that gauge completion generates a nonzero retained
harmonic source shift. The next variable to control is the
constant fine-field coordinate which solves the actual root
constraint in the presence of the detail fluctuation. Its
fourth coefficient is not the negative of the pure fourth
source coefficient: the mixed constant/noise response enters
at the same order. This section identifies the resulting
logarithm, proves its subtracted finite part, and computes the
first nonzero noisy normal-Jacobian coefficient at that solved
pivot. The full measured integral is not replaced by its
Gaussian source coefficients.

The scheduled companion review was performed against the
current files. SM V72 remains a many-copy derivation of a
covariance transport law at fixed finite mode cutoff and
bounded physical times. Its convergence of mixed moments
does not supply a growing-cutoff YM estimate; factorization
requires its additional point-concentration hypothesis.
NS V26 concerns explicit pointwise rank-one contractions of
Oseen-kernel derivatives and leaves global regularity open.
Its radial comparisons do not bound a YM conditional density.
The files have not changed since the previous checkpoint:
\begin{center}\small
\begin{tabular}{lr}
SM V72 & 607,372 bytes\\
NS V26 & 278,283 bytes\\
Shell-Mixing V16 & 623,261 bytes
\end{tabular}
\end{center}
Their SHA--256 digests, in the same order, are
\begin{center}\scriptsize\ttfamily
F44F9745D43379E1672C5550411B58E97195A4C9278592D221DF5D3F644DC1F5\\
82C887F8C2C69E4F28FAD7C3DC8DE5B9099CB5A4BF431EBA1FFF3E91935978AD\\
BF138E63BA826BC030D3A8B87E04FDD3B1EC712D269BB064674E96509B46E710.
\end{center}
Shell V16's measured Coulomb chart and the earlier exact
source-density convention remain the applicable framework.
No SM or NS theorem is imported below. In particular the
random matrix products appearing in the pivot are retained,
not replaced by products of expectations. The next regular
companion review is V40.

\subsection{The actual source constraint at a quadratic target scale}

Work at terminal side one with \(L=2^n\). Write
\(F_L=\mathcal F_{L\to1}\), \(\iota_Ls=s/L\), and let
\(z\in V_L\) be mean-zero Coulomb. For a deterministic target
\(h\in\mathbb R^{12}\), define the harmonic coordinate germ
\(\sigma_L(\epsilon,z;h)\) by
\begin{equation}\label{eq:v35-constraint}
       F_L\bigl(\iota_L\sigma_L(\epsilon,z;h)+\epsilon z\bigr)
                                   =\epsilon^2h,
       \qquad \sigma_L(0,z;h)=0.
\end{equation}
For every fixed regulator and deterministic \(z,h\), this
germ exists uniquely by the analytic implicit-function theorem:
the derivative in \(s\) of \(F_L(\iota_Ls)\) at zero is the
identity. Its first \(\epsilon\) coefficient is zero because
\(DF_L(0)z=0\). The radius of this germ is not asserted
uniform in \(L\) or in a Gaussian sample.

Use the following exact coefficient notation, evaluated on
the same input \(z\):
\begin{equation}\label{eq:v35-fields}
 \begin{aligned}
 Q_L&=F_{L,[2]}(z),& \mathcal C_L&=F_{L,[3]}(z),&
             \mathcal D_L&=F_{L,[4]}(z),\\
 \mathcal L_L(z)b&=D^2F_L(0)[\iota_Lb,z],&&&\\
 \mathcal M_L(z)b&=\tfrac12D^3F_L(0)[\iota_Lb,z,z].&&&
 \end{aligned}
\end{equation}
The last two quantities are real \(12\)-by-\(12\) matrices
in the source direction \(b\). The second is exactly the
mixed cubic observable of V24, not a new contraction rule.
On the common Gaussian realization set
\[
       \mathcal M_L^\circ=\mathcal M_L-\Gamma_L,
             \qquad G_L=\Gamma_L-n\Gamma_*.
\]
An orientation matrix such as \(\Gamma_L\) acts identically
on all three colours. It is lifted to source dimension twelve
when added to \(\mathcal M_L\).

\begin{lemma}[The mixed matrices used by the pivot]
\label{lem:v35-mixed}
For \(z=Z_L\), the random matrices \(\mathcal L_L\) and
\(\mathcal M_L^\circ\) converge jointly with
\(Q_L,\mathcal C_L,\mathcal H_L\), and
\(\mathcal D_L-n\mathcal H_L\) on the common Gaussian
coupling, in every finite \(L^q\). They have uniformly bounded
finite moments. The first matrix is first chaos; the second
is second chaos. Also \(G_L\to\Gamma_{\rm fin}\) with a
uniform deterministic bound, and
\begin{equation}\label{eq:v35-trace-linear}
                              \Tr_{12}\mathcal L_L(z)=0
\end{equation}
for every deterministic input.
\end{lemma}
\begin{proof}
For each fixed pair of source input/output coordinates, an
entry of \(\mathcal L_L\) is a bounded linear functional of
the energy-space detail direction. The mixed second-derivative
convergence in Theorem~\ref{thm:v13-source} gives uniform
bounds and convergence of its Riesz representer on the
common real energy space. Applying the isonormal family
therefore gives first-chaos \(L^2\) convergence and bounds.
There are only \(12^2\) entries, so matrix norm bounds follow
by a fixed finite sum; Gaussian moment comparison gives
every finite \(L^q\). This does not evaluate the full
nonlinear source on a Hilbert-space-valued Gaussian vector.

For \(\mathcal M_L^\circ\), apply
Theorem~\ref{thm:v24-mixed} to the twelve fixed source basis
directions. That theorem retains the entire mixed source
derivative and identifies its centered Hilbert--Schmidt kernel.
V22 gives the assertion for \(G_L\), and V23, V25 and V32
give the other convergent fields. All are on the same
coupling, which proves joint convergence of a fixed finite
list; no independence assertion is made.

Finally the quadratic Taylor coefficient of each actual
BCH word is a finite scalar linear combination of Lie
brackets of input links. Differentiating one leg in the
constant direction of colour \(E_c\) gives brackets of
\(E_c\) with a detail colour vector. Their output component
in colour \(E_c\) is zero. Each diagonal entry of the
source matrix \(\mathcal L_L\) is therefore zero, proving
\eqref{eq:v35-trace-linear}.
\end{proof}

\subsection{The first four coefficients of the solved harmonic coordinate}

Write
\[
 \sigma_L(\epsilon,z;h)
       =\epsilon^2\sigma_{L,2}(h)
           +\epsilon^3\sigma_{L,3}(h)
           +\epsilon^4\sigma_{L,4}(h)+O(\epsilon^5).
\]
The suppressed \(z\) dependence is polynomial in each
displayed coefficient.

\begin{theorem}[Exact source-pivot coefficients]
\label{thm:v35-pivot}
The coefficients in \eqref{eq:v35-constraint} are
\begin{equation}\label{eq:v35-pivot}
 \begin{aligned}
 \sigma_{L,2}(h)&=h-Q_L,\\
 \sigma_{L,3}(h)&=-\mathcal C_L-\mathcal L_L\sigma_{L,2}(h),\\
 \sigma_{L,4}(h)&=-\mathcal D_L-\mathcal L_L\sigma_{L,3}(h)
                                      -\mathcal M_L\sigma_{L,2}(h).
 \end{aligned}
\end{equation}
In particular a solution of the constraint requires the
mixed matrix evaluated on the random vector \(h-Q_L\).
\end{theorem}
\begin{proof}
The actual homogeneous source has no quadratic term:
\(F_L(\iota_Ls)=s+O(|s|^3)\), as used in V14.
Polarization thus makes every second derivative with two
constant legs zero. When \(s=O(\epsilon^2)\), the terms
of source degree at most four are consequently
\[
 s+\epsilon^2Q_L+\epsilon^3\mathcal C_L+\epsilon^4\mathcal D_L
                       +\epsilon\mathcal L_Ls+\epsilon^2\mathcal M_Ls.
\]
The possible quadratic term with two constant legs vanishes.
Terms with two constant legs and one detail have weighted
degree at least five, and the purely constant cubic has
degree at least six. Equating the three coefficients to
\(\epsilon^2h\) gives \eqref{eq:v35-pivot}. The coefficient
of each next unknown \(\sigma_{L,r}\) is the identity, so
these are the unique coefficients of the actual implicit germ.
\end{proof}

The matrix \(\mathcal M_L^\circ\) is centered for a
deterministic argument. This does not allow one to remove
\(\mathcal M_L^\circ(h-Q_L)\): the matrix and the random
argument use the same fine Gaussian field, and their product
contains additional contractions. The proof below keeps the
whole finite-dimensional product.

\subsection{The new fourth-order logarithm and its finite part}

Define
\begin{equation}\label{eq:v35-counterterm}
 \begin{aligned}
 \mathcal K_L(h)&=\mathcal H_L+\Gamma_*(h-Q_L),\\
 \widehat\sigma_{L,4}(h)&=\sigma_{L,4}(h)+n\mathcal K_L(h),\\
 \widehat{\mathcal D}_L&=\mathcal D_L-n\mathcal H_L.
 \end{aligned}
\end{equation}
At zero target this is the new combination
\(\mathcal K_L(0)=\mathcal H_L-\Gamma_*Q_L\).
The matrix multiplies the retained output, while
\(\mathcal H_L\) arises from the fine-input tangent response;
the two operations must not be identified.

\begin{theorem}[Cofinal harmonic-pivot coefficients]
\label{thm:v35-finite-part}
For \(h\) in any fixed bounded set, the fields
\(\sigma_{L,2}(h)\), \(\sigma_{L,3}(h)\),
\(\mathcal K_L(h)\), and \(\widehat\sigma_{L,4}(h)\)
have uniformly bounded finite \(L^q\) norms and cofinal
limits on the common Gaussian coupling. The fourth finite
part has the exact representation
\begin{equation}\label{eq:v35-finite-part}
 \begin{aligned}
 \widehat\sigma_{L,4}(h)
  ={}&-\widehat{\mathcal D}_L+\mathcal L_L\mathcal C_L
                       +\mathcal L_L^2(h-Q_L)\\
     &-(\mathcal M_L^\circ+G_L)(h-Q_L),
 \end{aligned}
\end{equation}
and its limit is obtained by replacing the factors by their
cofinal limits, with \(\widehat{\mathcal D}_L\to\mathcal D_{\rm fin}\)
and \(G_L\to\Gamma_{\rm fin}\). Thus the fixed coarse and
fine endpoint terms of V22 and V32 remain in this pivot.
The convergence is uniform on bounded \(h\) sets, also
after differentiating in \(h\).

There are uniform small stretched-exponential moments of
\(|\sigma_{L,3}(h)|^{2/3}\) and
\(|\widehat\sigma_{L,4}(h)|^{1/2}\). The second coefficient
and \(\mathcal K_L(h)\) have uniform small ordinary
exponential moments. The constants may depend on the bound
for \(h\).
\end{theorem}
\begin{proof}
Insert the second line of \eqref{eq:v35-pivot} into its
third line, split
\(\mathcal D_L=n\mathcal H_L+\widehat{\mathcal D}_L\), and
use \(\mathcal M_L=n\Gamma_*+G_L+\mathcal M_L^\circ\).
The two logarithmic terms are exactly
\(-n\mathcal H_L-n\Gamma_*(h-Q_L)\). Moving them to the
left proves \eqref{eq:v35-finite-part}.

Lemma~\ref{lem:v35-mixed} bounds and identifies every factor
on the right. Holder's inequality using higher finite
moments proves convergence of each matrix product on the
common coupling, including its correlations. For example
split a product difference into the difference of the first
factor times the second, and the limiting first factor times
the difference of the second. The same higher-moment bounds
control both terms. The preceding two pivot coefficients and
\(\mathcal K_L(h)\) then follow from their formulas.
Every displayed coefficient is affine in \(h\); uniformity
on bounded sets and its source derivatives therefore follow
from finitely many matrix coefficient bounds.

The first two pivot coefficients have Gaussian polynomial
degrees at most two and three. The subtracted fourth
coefficient in \eqref{eq:v35-finite-part} has degree at most
four, including its lower-chaos contractions, and is uniformly
bounded in \(L^2\). Hypercontractivity gives moment growth
at most \(Cq^{r/2}\) for degree at most \(r\).
Expanding the exponential, as in V32, proves the corresponding
stretched-exponential orders \(2/r\). The same argument in
degree two gives the stated ordinary exponential moments.
These are polynomial moment statements, not an integrability
bound for the full constrained density.
\end{proof}

\begin{proposition}[The same finite part with the fixed continuum counterterm]
\label{prop:v35-fixed-counterterm}
On the quantitative coupling fixed in V33,
\begin{equation}\label{eq:v35-counterterm-rate}
 \|Q_L-Q_\infty\|_q
       +\|\mathcal K_L(h)-\mathcal K_\infty(h)\|_q
                          \le C_qL^{-1/7}\sqrt{1+n}.
\end{equation}
Thus \(\sigma_{L,4}(h)+n\mathcal K_\infty(h)\) has the
same finite-part limit as \(\widehat\sigma_{L,4}(h)\).
\end{proposition}
\begin{proof}
The V34 tangent identity with orientation matrix \(A=I\)
gives \(H^I_{L,1}=DQ_L(z)[z]=2Q_L(z)\) exactly. In the
proof of V33's strong quadratic-counterterm rate, the only
properties of the inserted matrix used were its bounded
norm and the actual linear-filter estimates. Apply that
same argument with \(A=I\). Explicitly, the finite-head
and two-tail estimate for \(2Q_L\) is
\[
 C_q\{L^{-1/3}2^{4b/3}+2^{-b}\sqrt{1+b}\},
\]
because each shell is now twice the primitive quadratic
insertion. The linear-input comparison and the positive
outer-row mass are unchanged. Choosing
\(b=\lfloor n/7\rfloor\) proves the rate for \(Q_L\).
Combine it with \eqref{eq:v33-H-rate} and the fixed bounded
matrix \(\Gamma_*\) to prove \eqref{eq:v35-counterterm-rate}.
The deterministic term \(\Gamma_*h\) cancels in this
difference. Finally \(nL^{-1/7}\sqrt{1+n}\to0\), proving
the assertion about the finite part.
\end{proof}

As in V33, a continuum subtraction on the coupling need
not be measurable with respect to the finite fine field.
If finite-field measurability is required, use
\(\mathbb E[\mathcal K_\infty(h)\mid Z_L]\). Its constant
part is unchanged, and Lemma~\ref{lem:v33-compression}
applies to its second-chaos part. The same rate and finite
part follow. This supplies a finite-dimensional polynomial,
not a new locality or full-density theorem.

\begin{proposition}[The pivot logarithm does not vanish]
\label{prop:v35-nonzero}
The limiting zero-target counterterm obeys
\begin{equation}\label{eq:v35-nonzero}
            \|\mathcal K_\infty(0)_4^{,3}\|_2
                           \ge\frac{a_*}{\sqrt8\,\pi^3}>0.
\end{equation}
Moreover
\(\sigma_{2^n,4}(0)/n\to-\mathcal K_\infty(0)\)
in every finite \(L^q\).
\end{proposition}
\begin{proof}
The fourth row of the triangular matrix \(\Gamma_*\) has
only its diagonal entry \(\gamma_4\), by V28. On the
helical Cameron--Martin field \eqref{eq:v31-helical}, V31
gives \(\mathcal H_\infty[u,u]_4^{,3}=a_*/\pi\) and
\(Q_\infty[u,u]_4^{,3}=0\). Hence
\(\mathcal K_\infty(0)[u,u]_4^{,3}=a_*/\pi\).
This is a centered second-chaos component. Its symmetric
energy kernel \(T\) satisfies
\[
 |\langle u,Tu\rangle|
            \le2^{-1/2}\|I_2(T)\|_2\|u\|_{\rm CM}^2,
                  \qquad \|u\|_{\rm CM}^2=4\pi^2.
\]
This proves \eqref{eq:v35-nonzero}. Divide the definition
of \(\widehat\sigma_{L,4}(0)\) by \(n\), and use its
uniform bound and the cofinal limit of \(\mathcal K_L(0)\)
to prove the last assertion.
\end{proof}

The V34 harmonic shift is visible in this new combination.
With \(J_L=H_L^{\Gamma_*,\Pi}-\mathcal H_L\) as in
\eqref{eq:v34-harmonic-limit}, the exact identity is
\begin{equation}\label{eq:v35-gauge-relation}
       \mathcal K_L(0)=H_L^{\Gamma_*,\Pi}-\Gamma_*Q_L-J_L.
\end{equation}
Thus the generated harmonic part is present in the solved
source coordinate, not removed by imposing the mean-zero
condition on the separate detail variable.

For clarity about the joint noise/scale regime, the unmodified
four-term pivot polynomial satisfies the following consequence.
If \(\beta\to\infty\), \(L=2^n\to\infty\), and
\(n/\beta\to\tau<\infty\), then
\begin{equation}\label{eq:v35-joint-polynomial}
 \begin{aligned}
 \beta\sum_{r=2}^4\beta^{-r/2}\sigma_{L,r}(h)
                 \longrightarrow h-Q_\infty-\tau\mathcal K_\infty(h)
 \end{aligned}
\end{equation}
in every finite \(L^q\). Indeed the cubic contributes
\(O(\beta^{-1/2})\), the fourth finite part contributes
\(O(\beta^{-1})\), and its explicit logarithm contributes
\(-(n/\beta)\mathcal K_L(h)\). This statement concerns
the displayed polynomial only. It does not identify the
limit of the actual implicit pivot at
\(\epsilon=\beta^{-1/2}\).

\subsection{The source-coordinate Jacobian at the noisy solved pivot}

Let
\[
 \mathcal N_L(\epsilon,z;h)b
     =DF_L\bigl(\iota_L\sigma_L(\epsilon,z;h)+\epsilon z\bigr)
                                                        [\iota_Lb].
\]
This derivative fixes the detail coordinate. It is not the
derivative of a minimizing map or the derivative in which
the detail is reoptimized. At a fixed regulator and input,
its determinant is positive on a sufficiently small real
\(\epsilon\) interval because \(\mathcal N_L(0)=I\).

\begin{theorem}[Noisy normal-Jacobian coefficient and finite part]
\label{thm:v35-jacobian}
The exact Taylor coefficients through degree two are
\begin{equation}\label{eq:v35-normal}
 \mathcal N_L=I+\epsilon\mathcal L_L+\epsilon^2\mathcal M_L
                         +O(\epsilon^3),\qquad
 \log\det\mathcal N_L=\epsilon^2\mathfrak j_{L,2}
                         +O(\epsilon^3),
\end{equation}
where
\begin{equation}\label{eq:v35-jacobian}
 \begin{aligned}
 \mathfrak j_{L,2}
    &=3\Tr_4\Gamma_L+\Tr_{12}\mathcal M_L^\circ
                                      -\tfrac12\Tr_{12}(\mathcal L_L^2),\\
 \mathfrak j_{L,2}-3n\Tr_4\Gamma_*
    &\longrightarrow 3\Tr_4\Gamma_{\rm fin}
       +\Tr_{12}\mathcal M_\infty^\circ
                                      -\tfrac12\Tr_{12}(\mathcal L_\infty^2).
 \end{aligned}
\end{equation}
The second convergence holds in every finite \(L^q\),
with uniform bounds and a uniform small ordinary exponential
moment for the subtracted coefficient. These coefficients
do not depend on \(h\).
\end{theorem}
\begin{proof}
The solved constant coordinate begins at \(\epsilon^2\).
In the first derivative of the source, its possible
contribution of order two through \(D^2F_L(0)\) has two
constant legs and is zero. Every other dependence on the
pivot begins at order three or higher. Thus the first two
coefficients of \(\mathcal N_L\) are exactly
\(\mathcal L_L\) and \(\mathcal M_L\), even though the
derivative is evaluated at the solved pivot.

The matrix logarithm gives
\[
 \log\det\mathcal N_L
   =\epsilon\Tr_{12}\mathcal L_L
      +\epsilon^2\{\Tr_{12}\mathcal M_L
                      -\tfrac12\Tr_{12}(\mathcal L_L^2)\}
      +O(\epsilon^3).
\]
The linear trace vanishes by \eqref{eq:v35-trace-linear}.
The expectation of the mixed matrix is \(\Gamma_L\) on
orientations and the identity on three colours, giving
\(\Tr_{12}\Gamma_L=3\Tr_4\Gamma_L\). This proves the first
line of \eqref{eq:v35-jacobian}. Lemma~\ref{lem:v35-mixed}
and the deterministic matrix finite part prove the second
line, retaining the full random matrix square. The
subtracted expression has Gaussian degree at most two
and a uniform \(L^2\) bound; degree-two moment growth
gives its small exponential moment.
\end{proof}

In particular, \(\mathcal L_L\) has trace zero, but no
general identity \(\Tr\mathcal L_L^2=0\) was used in the
proof. A vanishing square trace for a particular finite
test field is not permission to delete this term.
The logarithm \(3n\Tr_4\Gamma_*\) in
\eqref{eq:v35-jacobian} is a source-coordinate contribution,
not an assertion about the complete vacuum or coupling
counterterm after other measured factors have been combined.

\subsection{Where this factor belongs in the exact density}

Take the specified finite slice density in linear coordinates
to be \(m_L(s,w)\,ds\,dw\), as in V14. It may contain the
actual Wilson weight, Haar factors, and gauge density; none
of them is set to one here. The exact local source-coordinate
formula is the already proved identity
\[
 \frac{m_L(s(v,w),w)}
      {\det D_sF_L(\iota_Ls(v,w)+w)}\,dv\,dw.
\]
For \(\epsilon>0\), put \(w=\epsilon z\) and
\(v=\epsilon^2h\). Since
\(d_L=\dim V_L=9(L^4-1)\), the same density in \((h,z)\)
coordinates, wherever this local chart is defined, is
\begin{equation}\label{eq:v35-density}
       \epsilon^{d_L+24}
       \frac{m_L(\sigma_L(\epsilon,z;h),\epsilon z)}
            {\det\mathcal N_L(\epsilon,z;h)}\,dh\,dz.
\end{equation}
The source dimension is twelve; its rescaling contributes
\(\epsilon^{24}\), and the detail rescaling contributes
\(\epsilon^{d_L}\). Cross derivatives of the implicit
coordinate do not add another determinant to this triangular
change of variables.

Consequently \(+\log\det\mathcal N_L\) enters the negative
logarithm of the amplitude exactly once. Formula
\eqref{eq:v35-density} is in the Lebesgue source convention;
a conversion to the specified coarse Haar reference must
retain its separate factor, as in Shell V16. This formula
does not introduce an extra coarea factor on top of a density
that already contains the same source-coordinate Jacobian.

The Gaussian coefficient limits above do not replace
\(m_L(\sigma_L,\epsilon z)\) by a Gaussian weight, and do
not prove integrability of a truncated exponential. The
noise-dependent normal determinant now has a controlled
second coefficient after its explicit subtraction. Its
higher coefficients and its nonperturbative lower bound
are still required for a measured iteration at growing
depth.

\subsection{Exact calibration and the remaining nonlinear requirement}

The certificate is
\begin{center}\small
\path{scripts/verify_ym_f14_v35_source_pivot.py}
\end{center}
with SHA--256
\begin{center}\scriptsize\ttfamily
B3FC76C691280FAE4A1B0D9E3E8FE96C06C5630EB77B96F20EA12BE9B4862920.
\end{center}
It uses the actual side-two live field
\(z=(\Pi_2e_0)E_1+(\Pi_2e_1)E_2\). Twelve weighted
constant-source insertions in the literal root extract the
complete \(12\)-by-\(12\) matrices \(\mathcal L_2\) and
\(\mathcal M_2\). The source constraint is then checked
through degree four for zero target and for a fixed rational
target with all four orientations present. A further 24
constant-source column checks evaluate the normal derivative
at those solved pivots, not at a substituted unshifted detail.

At zero target the quadratic and quartic coefficients are
\[
 \begin{aligned}
 \sigma_{2,2}(0)&=(-7E_3/192,\ 19E_3/192,\ 0,\ 0),\\
 \sigma_{2,4}(0)&=(-64675E_3/42467328,\ 709937E_3/14155776,\ 0,\ 0).
 \end{aligned}
\]
The linear mixed matrix has twelve nonzero
entries on this field. Its trace square happens to vanish,
while
\[
                \Tr_{12}\mathcal M_2
                     =\mathfrak j_{2,2}=-8251/1536.
\]
Independent truncated-polynomial Gaussian elimination of
\(I+\epsilon\mathcal L_2+\epsilon^2\mathcal M_2\) verifies
the determinant and logarithm coefficients. A separate
noncommuting rational matrix example has nonzero square
trace and checks that part of the formula. These exact
finite tests certify source algebra and normalization;
the cofinal assertions use the analytic proofs above.

The remaining issue is not another change of notation for
the subtraction. A scale-local measured construction must
control the full implicit source coordinate and its normal
determinant on typical conditional fluctuations, together
with the Wilson, Haar and gauge factors. The V23 small-ball
obstruction still prevents importing the fixed global
classical chart as a typical quantum domain. Four pivot
coefficients and two Jacobian coefficients do not provide
a common radius or a uniform integral remainder when
\(n/\beta\) is of order one. In particular
\eqref{eq:v35-joint-polynomial} cannot be promoted to an
actual constrained-measure limit without those estimates.

The predecessor V34 has 782,810 bytes and SHA--256
\begin{center}\scriptsize\ttfamily
D01A8C6244AA781FCF1005621953DFFD4EF6D449ABA39B2FAEB388361BCA32BC.
\end{center}
Its body is preserved apart from the title version and this
terminal insertion. V35 replaces V34 as the single active
main note. Archives, monographs, SM V72, NS V26 and Shell V16
are preserved. The scheduled companion review is complete;
the next one is V40.

\begin{firewall}
The new results are a nonzero logarithmic counterterm and
cofinal fourth finite part for the actual noisy harmonic
source pivot, together with the renormalized quadratic
normal-Jacobian coefficient and its exact placement in the
finite density. They are not a nonperturbative constrained
measure construction or a physical Yang--Mills mass-gap proof.
The full interacting continuum and positive physical gap
remain unproved.
\end{firewall}
\section{V36: an exact repeated-cell normal matrix and a nonlinear source chart beyond the global small-field ball}
\label{sec:v36}

\subsection{Why change the test domain rather than add another coefficient?}

V35 constructs four constrained-source coefficients and a quadratic
normal-Jacobian finite part. Its last obstruction concerns the actual
nonlinear constraint on a common domain. Adding another coefficient
would not by itself address that obstruction. We instead exhibit an
inhomogeneous family on which the complete normal derivative is explicit,
prove a common nonlinear source inverse, and extend that inverse to
full-dimensional translated neighborhoods. The background has arbitrarily
large global counting \(\ell^4\) norm as the lattice is refined.
Consequently the domain extension is not an instance of the old
zero-centered small-ball theorem.

The reduction uses, without redoing, the homogeneous nonlinear limit
proved in Shell V10, and the global counting-norm theorem of Shell
V12--V13. Shell V7 already treats the fixed local measured integral and
its uniform finite-loop remainder; the result below is not another
derivation of that expansion. Shell V16's measured logarithmic coefficient
has already been incorporated in V30. The new calculation concerns the
full source geometry, not a replacement of the measured coefficient by
the source matrix. SM V72 and NS V26 were reviewed at V35; the next
scheduled companion checkpoint is V40.

Throughout this section \(L=2^n\), \(n\ge1\), and the fine lattice is
\(\Lambda_L=(\mathbb Z/L\mathbb Z)^4\). Link coordinates belong to
\(\mathfrak{su}(2)\), with the orthonormal quaternion basis
\(E_iE_j=-\delta_{ij}I+\epsilon_{ijk}E_k\). Write
\[
 \|A\|_4^4=\sum_{x\in\Lambda_L}\sum_{\mu=1}^4|A_{x,\mu}|^4,
 \qquad (\iota_Ls)_{x,\mu}=s_\mu/L,
 \qquad s\in\mathbb R^{12}.
\]
The Lie-color norm is Euclidean and the counting measure is not
volume-normalized. In particular \(\|\iota_Ls\|_4\le |s|\).
Let \(\mathfrak T_L\) be the actual balanced primitive root-log block,
and let \(F_L\) be its complete nested composition down to side one;
set \(F_1=I\). No Coulomb projection is inserted between blocks.

\subsection{The zero-root stripe and the precise logarithmic branch}

For real \(a\), define
\begin{equation}\label{eq:v36-stripe}
 A^{(a)}_{x,4}=a(-1)^{x_2}E_3,
 \qquad A^{(a)}_{x,\mu}=0\quad(\mu\ne4).
\end{equation}
This field is mean zero and satisfies the linear lattice Coulomb
condition: its only nonzero orientation is independent of its own
coordinate \(x_4\). Its fundamental cell has side two, in all four
directions.

Here is the literal root convention needed for the computation. For
\(s\in\{0,1\}^4\), let \(T_{r,s}\) be the ordered tree from \(2r\)
to \(2r+s\), taking directions in the order \(1,2,3,4\). Put
\[
 \begin{aligned}
 P_{r,s,\mu}&=T_{r,s}U_{2r+s,\mu}U_{2r+s+e_\mu,\mu}
                                      T_{r+e_\mu,s}^{-1},\\
 \phi_r&=\frac1{16}\sum_s\log T_{r,s},\\
 K_{r,\mu}&=P_{r,0,\mu}
       \exp\left(\frac1{16}\sum_s
                  \log(P_{r,0,\mu}^{-1}P_{r,s,\mu})\right),\\
 \mathfrak T_L(A)_{r,\mu}
       &=\log\bigl(e^{-\phi_r}K_{r,\mu}e^{\phi_{r+e_\mu}}\bigr),
       \qquad U_{x,\mu}=e^{A_{x,\mu}}.
 \end{aligned}
\]
Every logarithm is the identity branch. The primitive common
polydisc in Shell V10 has radius \(r_0=1/4096\). We will also describe
the continuation of this specific logarithmic formula along
\(|a|<\pi/4\), where all the displayed stripe logarithms remain
strictly inside their principal branches. This is a specified
continuation of the local rule, not an assertion about an arbitrary
smooth compact completion outside its identity-log region. To stay
entirely within the previously identified common primitive chart,
take \(a_0<r_0\) below and decrease the perturbation radii accordingly.

\begin{lemma}[Exact annihilation in one primitive step]
\label{lem:v36-zero}
For \(|a|<\pi/4\), on the branch just specified,
\[
            \mathfrak T_L(A^{(a)})=0,\qquad F_L(A^{(a)})=0.
\]
For any fixed \(a_0<\pi/4\), the primitive formula is holomorphic
on a common complex neighborhood of each cell with real
\(|a|\le a_0\). This neighborhood is independent of \(L\).
\end{lemma}
\begin{proof}
At the stripe the trees and paths are
\[
 T_{r,s}=e^{s_4(-1)^{s_2}aE_3},\qquad
 P_{r,s,4}=e^{2(-1)^{s_2}aE_3},\qquad
 P_{r,s,\mu}=I\ (\mu\ne4).
\]
Thus \(\phi_r=0\). For orientation four, the reference path is
\(e^{2aE_3}\), and the relative logarithms are zero on the eight
sites with \(s_2=0\) and \(-4aE_3\) on the other eight sites.
Their mean is \(-2aE_3\); hence \(K_{r,4}=I\). All other
orientations are already the identity. This proves the first assertion;
subsequent blocks preserve zero. The largest absolute relative-log
angle is \(4|a|<\pi\). All other tree and output logarithms also
avoid their cuts. Analyticity of the finite products, exponentials and
these logarithms gives a neighborhood of each cell. Compactness of
\([-a_0,a_0]\), and the finite list of cell factors, give a common
smaller neighborhood. Periodic identifications do not enlarge that list.
\end{proof}

Define the twelve-dimensional primitive source map
\begin{equation}\label{eq:v36-cell}
 G_a(b)=F_2(A^{(a)}+\iota_2b),\qquad
 N(a)=DG_a(0),\qquad G_a(0)=0.
\end{equation}
This is the nonlinear map, not its degree-four truncation.

\subsection{The full normal matrix and an explicit inverse bound}

Identify the transverse color vector \(xE_1+yE_2\) with \(x+iy\).
Rotations about \(E_3\) fix the background and commute with \(N(a)\),
so its color-parallel and transverse subspaces are invariant. On the
transverse subspace it is complex linear.

\begin{theorem}[No normal caustic on the stripe identity-log branch]
\label{thm:v36-normal}
Put \(\operatorname{sinc}t=\sin(t)/t\), with its value one at zero, and
\begin{equation}\label{eq:v36-scalars}
 c=\cos^2 a,\qquad d=\frac{\tan(2a)}{2a},\qquad
 b=a d,\qquad k=\frac{a d}{2}\bigl(1-\operatorname{sinc}(4a)\bigr).
\end{equation}
On the four parallel color coordinates, \(N(a)=I_4\). On the four
complex transverse coordinates, ordered by link orientation, it is
\begin{equation}\label{eq:v36-matrix}
 C(a)=
 \begin{pmatrix}
 c&0&0&0\\
 0&c&0&0\\
 0&0&c&0\\
 ik&ib&ik&d
 \end{pmatrix}.
\end{equation}
For every real \(|a|<\pi/4\), it is invertible, and
\begin{equation}\label{eq:v36-det}
 \det_{\mathbb R}N(a)=\cos^{12}a
                  \left(\frac{\tan(2a)}{2a}\right)^2>0,
 \qquad \|N(a)^{-1}\|_{\rm op}<5.
\end{equation}
Moreover
\begin{equation}\label{eq:v36-invdet}
 r(a):=\det N(a)^{-1}
       =\sec^{12}a\,(2a\cot(2a))^2\le64.
\end{equation}
All formulas at zero have removable singularities.
\end{theorem}
\begin{proof}
When every perturbation is parallel to \(E_3\), all factors commute.
On this branch the primitive averages their angles exactly. The stripe
has zero average and the two-link constant source is \(b\); consequently
\(G_a(b)=b\) for these parallel perturbations. This proves the
parallel block and, by the rotation equivariance, leaves just the
complex transverse block to compute.

The following elementary first-variation calculus specifies the entire
finite calculation without a Taylor expansion in \(a\). Represent the
Lie jet \(kaE_3+t(u_1E_1+u_2E_2)\) by \((k,u_1+iu_2)\), set
\(z=e^{ia}\), and write
\[
 s_k=\frac{z^k-z^{-k}}{2ika}\quad(k\ne0),\qquad s_0=1.
\]
The transverse quaternion coefficient of its exponential, to first
order in \(t\), is \(s_k u\). Multiplying two quaternion jets and
then taking the logarithm gives the exact first-variation rule
\begin{equation}\label{eq:v36-jetproduct}
 (k,u)\star(l,v)=
 \left(k+l,\frac{z^{-l}s_k u+z^k s_l v}{s_{k+l}}\right).
\end{equation}
Inversion is \((k,u)\mapsto(-k,-u)\). Averaging logarithms is
componentwise averaging of the pairs. The identities follow directly
from \(E_3E_1=E_2\), \(E_1E_3=-E_2\), and
\(e^{\theta E_3}=\cos\theta+E_3\sin\theta\). Denominators at
zero angles are interpreted by continuity.

For the source column \(\lambda\), assign the fine link at \((x,\mu)\)
the pair
\[
       \left((-1)^{x_2}\mathbf1_{\{\mu=4\}},
                              \tfrac12\mathbf1_{\{\mu=\lambda\}}\right).
\]
Apply \eqref{eq:v36-jetproduct} to the ordered trees and the two-link
paths, then to the reference inverse, the sixteen-term relative-log
mean, and the reference path. The endpoint tree means agree by
periodicity; their conjugation has zero first-order effect because the
unperturbed root is the identity. The resulting nonzero entries obey
the following rational identities in \(a,z\):
\[
 \begin{aligned}
 4z^2 C_{\mu\mu}&=(z^2+1)^2 &&(\mu=1,2,3),\\
 2ia(z^4+1)C_{44}&=z^4-1,\\
 2(z^4+1)C_{42}&=z^4-1,\\
 4(z^4+1)C_{41}
       &=4(z^4+1)C_{43}=(z^4-1)(1-s_4).
 \end{aligned}
\]
All other entries vanish. These identities can also be obtained by
splitting the sixteen path terms into the eight with \(s_2=0\)
and the eight with \(s_2=1\): their relative base angles are
respectively zero and \(-4a\). Thus the computation retains the
reference-path logarithm, not a replacement by the mean path log.
Using \((z^4-1)/(i(z^4+1))=\tan(2a)\) gives
\eqref{eq:v36-matrix}. The displayed finite recurrence and identities
constitute exact algebra; the independent rational-jet checks below
provide a further check of its signs and source normalization.

The complex inverse is
\begin{equation}\label{eq:v36-inverse}
 C(a)^{-1}=
 \begin{pmatrix}
 c^{-1}&0&0&0\\
 0&c^{-1}&0&0\\
 0&0&c^{-1}&0\\
 -\dfrac{ia(1-s_4)}{2c}&-\dfrac{ia}{c}
       &-\dfrac{ia(1-s_4)}{2c}&2a\cot(2a)
 \end{pmatrix}.
\end{equation}
For real \(|a|<\pi/4\), one has \(c\ge1/2\), \(d\ge1\), and
\(0\le s_4\le1\). Therefore its complex Frobenius norm satisfies
\[
 \|C(a)^{-1}\|_{\rm F}^2
 \le 13+6a^2
 <13+3\pi^2/8<17.
\]
The real operator norm of a complex-linear operator equals its complex
operator norm. Combining this with the parallel identity block proves
the bound in \eqref{eq:v36-det}. Finally the real determinant of the
transverse block is \(|\det_{\mathbb C}C|^2=c^6d^2\), which proves
both determinant formulas and the bound \(c^{-6}d^{-2}\le64\).
\end{proof}

For comparison with small-amplitude coefficients, exact expansion gives
\begin{equation}\label{eq:v36-detseries}
 \log\det N(a)=-\frac{10}{3}a^2+\frac{67}{45}a^4+O(a^6),
 \qquad r(a)=1+\frac{10}{3}a^2+\frac{61}{15}a^4+O(a^6).
\end{equation}
The inverse entries and \(r(a)\) have analytic extensions through
\(a=\pm\pi/4\), with \(r\) tending to zero there. The root logarithm
does not: a relative path reaches its cut and \(N(a)\) itself
diverges. No common nonlinear radius up to or across those endpoints
is inferred. Every uniform nonlinear statement below fixes a compact
subinterval \(|a|\le a_0<\pi/4\).

\subsection{Exact factorization and the cofinal nonlinear pivot}

To avoid confusing a dyadic depth with a side length, define here
\begin{equation}\label{eq:v36-homogeneous}
 \Phi_p(u)=F_p(\iota_pu),\qquad p=1,2,4,\ldots.
\end{equation}
Thus this \(\Phi_p\) is Shell V10's \(\Phi_{\log_2p}\), and
\(\Phi_1=I\). Its established holomorphic limit is denoted by
\(\Phi_\infty\). On a common small complex ball,
\[
 \Phi_p(0)=0,\quad D\Phi_p(0)=I,\quad
 |\Phi_p(u)-\Phi_\infty(u)|\le C|u|^3p^{-2}.
\]
Every fixed derivative and the local inverse maps converge on common
smaller balls at rate \(O(p^{-2})\). These facts are precisely the
homogeneous theorem already proved in Shell V10.

\begin{theorem}[A refinement-uniform nonlinear source inverse]
\label{thm:v36-pivot}
Fix \(a_0<\pi/4\). There are positive source radii independent of
\(L\) and real \(|a|\le a_0\) with the following properties. Put
\(p=L/2\) and
\[
                 \Psi_{p,a}(s)=pG_a(s/p).
\]
On those radii the exact, untruncated identity is
\begin{equation}\label{eq:v36-factor}
 F_L(A^{(a)}+\iota_Ls)=\Phi_p(\Psi_{p,a}(s)).
\end{equation}
The left side has normal derivative \(N(a)\) at \(s=0\), for
every refinement. For each sufficiently small \(v\), the unique local
source solution is
\begin{equation}\label{eq:v36-pivot}
 \begin{aligned}
 F_L(A^{(a)}+\iota_Ls_L(a,v))&=v,\\
 s_L(a,v)&=\Psi_{p,a}^{-1}(\Phi_p^{-1}(v)).
 \end{aligned}
\end{equation}
It is holomorphic in \(v\) on a common complex ball, and
\begin{equation}\label{eq:v36-limit}
 s_L(a,v)\longrightarrow s_\infty(a,v)
                     :=N(a)^{-1}\Phi_\infty^{-1}(v).
\end{equation}
Convergence is uniform for \(|a|\le a_0\) on smaller source balls,
with every fixed source derivative, at rate \(O(p^{-1})\). In degree
zero one can retain the more informative estimate
\begin{equation}\label{eq:v36-rate}
 |s_L(a,v)-s_\infty(a,v)|
          \le C_{a_0}\bigl(|v|^2p^{-1}+|v|^3p^{-2}\bigr).
\end{equation}
\end{theorem}
\begin{proof}
The input in \eqref{eq:v36-factor} has period two. Translation
covariance of the primitive rule implies that its first output is
homogeneous. Since \(\iota_Ls\) has link value \(s/(2p)\), the
value of this homogeneous output is exactly \(G_a(s/p)\). There
are \(\log_2p\) remaining blocks, so their output is
\(F_p(\iota_p[pG_a(s/p)])\). This proves the factorization.
Differentiation at zero, using \(D\Phi_p(0)=I\), gives \(N(a)\).

The common cell analyticity in Lemma \ref{lem:v36-zero} gives constants
\(R,C>0\), depending on \(a_0\) but not on \(p\), such that
\[
 G_a(b)=N(a)b+R_a(b),\qquad
 |R_a(b)|\le C|b|^2,\qquad \|DR_a(b)\|\le C|b|,
 \quad |b|<R.
\]
Consequently
\begin{equation}\label{eq:v36-remainder}
 \Psi_{p,a}(s)=N(a)s+O_{a_0}(|s|^2/p),\qquad
 D\Psi_{p,a}(s)=N(a)+O_{a_0}(|s|/p).
\end{equation}
Choose a fixed small radius \(r<R\) so that
\(5Cr<1/2\). The equation \(\Psi_{p,a}(s)=u\) is then
the contraction equation
\[
 s=N(a)^{-1}u-N(a)^{-1}pR_a(s/p).
\]
For \(|u|\) in a sufficiently small common ball, it maps
\(|s|\le r\) to itself and has Lipschitz constant at most \(1/2\).
The construction applies on complex source balls; the real-matrix
inverse bound controls their complexifications as well. The fixed point
is holomorphic, is unique in that ball, and satisfies
\[
 |\Psi_{p,a}^{-1}(u)|\le C_{a_0}|u|,\qquad
 |\Psi_{p,a}^{-1}(u)-N(a)^{-1}u|
                                      \le C_{a_0}|u|^2/p.
\]
Intersect this common image with that of the homogeneous inverses.
This proves \eqref{eq:v36-pivot} on a common source ball.

The homogeneous bound and a common Lipschitz bound for the inverses
give
\[
 |\Phi_p^{-1}(v)-\Phi_\infty^{-1}(v)|\le C|v|^3p^{-2}:
\]
apply the bound on \(\Phi_p-\Phi_\infty\) at
\(u=\Phi_\infty^{-1}(v)\), and then the Lipschitz bound of
\(\Phi_p^{-1}\). Combining this with the preceding estimate for
\(\Psi_{p,a}^{-1}\) proves \eqref{eq:v36-rate}. Uniform convergence
on a slightly larger complex ball and the Cauchy formula imply the
\(O(p^{-1})\) bounds for every fixed derivative on smaller balls.
This proof controls the complete analytic inverse, not a finite list
of pivot coefficients.
\end{proof}

There is a corresponding full normal-coordinate density factor. Define
\[
 \mathcal N_{L,a}(v)=
    D_sF_L(A^{(a)}+\iota_Ls)\big|_{s=s_L(a,v)},
 \qquad j_{L,a}(v)=\det\mathcal N_{L,a}(v)^{-1}.
\]
For real small \(v\), these determinants are positive. The chain rule
in \eqref{eq:v36-factor} gives the exact identity
\begin{equation}\label{eq:v36-jacobian}
 j_{L,a}(v)=
 \frac{1}{\det D\Phi_p(\Phi_p^{-1}(v))\,
                 \det DG_a(s_L(a,v)/p)}.
\end{equation}
Therefore
\begin{equation}\label{eq:v36-jacobianlimit}
 j_{L,a}(v)\longrightarrow
 \frac{r(a)}{\det D\Phi_\infty(\Phi_\infty^{-1}(v))},
 \qquad j_{L,a}(0)=r(a)\quad\hbox{for every }L.
\end{equation}
Convergence, including every fixed source derivative, is
\(O(p^{-1})\) on smaller common balls. Indeed the second determinant
in \eqref{eq:v36-jacobian} differs from \(\det N(a)\) by
\(O(|v|/p)\), while the first has the already proved homogeneous
convergence. Uniform invertibility on a smaller complex ball permits
reciprocation; the derivative assertion follows again by Cauchy.
This is a limit of the entire normal-coordinate factor, not just its
quadratic Taylor coefficient. It is not yet the full measured density.

\subsection{A full-dimensional translated tube, not only a periodic ansatz}

The repeated-cell inverse alone varies twelve source coordinates.
The following extension permits all sufficiently small fine-field
perturbations, including perturbations which destroy cell periodicity.

\begin{lemma}[One-step entry into the old global small-field domain]
\label{lem:v36-entry}
For fixed \(a_0<\pi/4\), there are \(\delta,C_0>0\), independent
of \(L\) and real \(|a|\le a_0\), such that
\begin{equation}\label{eq:v36-entry}
 \|\mathfrak T_L(A^{(a)}+\eta)\|_4\le C_0\|\eta\|_4,
                  \qquad \|\eta\|_4<\delta.
\end{equation}
After decreasing \(\delta\), the complete \(F_L(A^{(a)}+\eta)\)
is holomorphic on this translated complex ball, uniformly bounded
there, and has all fixed derivative bounds on smaller balls as an
\(\ell^4\)-to-\(\mathbb C^{12}\) map, independently of \(L\).
\end{lemma}
\begin{proof}
Each primitive output uses the same finite link carrier as in Shell
V12. A safe support bound is 54 links per output, with at most ten
output supports incident to any input link, counting small-torus
coincidences with multiplicity. Let \(K\) be that positive incidence
matrix. Compact cell analyticity from Lemma \ref{lem:v36-zero} gives
a common first-derivative bound on a smaller neighborhood. Since the
unperturbed output vanishes exactly, integration along the straight
complex segment gives
\[
 |\mathfrak T_L(A^{(a)}+\eta)_e|
                    \le C(K|\eta|)_e.
\]
The pointwise hypothesis is valid when \(\|\eta\|_4<\delta\), since
\(\|\eta\|_\infty\le\|\eta\|_4\). Jensen's inequality and the
incidence bounds imply
\[
                 \|K\|_{4\to4}\le54^{3/4}10^{1/4}.
\]
This proves \eqref{eq:v36-entry}. The constants do not involve
\(\|A^{(a)}\|_4\), which grows with \(L\).

Choose \(C_0\delta\) smaller than the common radius in Shell V12's
global counting-ball theorem, here applied to the first coarse field
on side \(L/2\). That theorem admits every remaining nested step and
bounds its output uniformly, without any intermediate gauge projection.
For \(L=2\) there are no remaining steps. Finite composition gives
holomorphy on the same translated input ball in every case. The uniform
output bound and multivariate Cauchy estimates in normalized
\(\ell^4\) directions give, on each smaller ball,
\[
 \|D^mF_L(A^{(a)}+\eta)[\eta_1,\ldots,\eta_m]\|
                        \le C_{m,a_0}\prod_{i=1}^m\|\eta_i\|_4.
\]
The Cauchy polydisc uses radii proportional to \(1/m\), so its
constant depends on the derivative order but not on the number of
fine-field coordinates. This proves the assertion.
\end{proof}

Let \(V_L\) be the mean-zero linear Coulomb space used in V35.
The affine Coulomb coordinates near the stripe are
\(A=A^{(a)}+w+\iota_Ls\), with \(w\in V_L\). Its real dimension
is \(9(L^4-1)+12\); the theorem below does not impose a periodic
ansatz on \(w\).

\begin{theorem}[Uniform source coordinates on translated Coulomb tubes]
\label{thm:v36-tube}
Fix \(a_0<\pi/4\). There exist \(R,\rho,\nu>0\), independent
of \(L\) and real \(|a|\le a_0\), such that whenever
\(\|w\|_4<\rho\), \(w\in V_L\), and \(|v|<\nu\), the equation
\begin{equation}\label{eq:v36-tube}
       F_L(A^{(a)}+w+\iota_Ls)=v
\end{equation}
has exactly one solution in the source ball \(|s|<R\). It obeys
\[
 |s|\le C_{a_0}(|v|+\|w\|_4).
\]
The solution is holomorphic in \((v,w)\) on common complex balls.
Every fixed mixed derivative is uniformly bounded when detail
directions are measured in \(\ell^4\). At the solved real point the
normal derivative \(\mathcal N_{L,a}(v,w)\) satisfies
\begin{equation}\label{eq:v36-tubebound}
 \|\mathcal N_{L,a}(v,w)^{-1}\|_{\rm op}\le10,
 \qquad 0<\det\mathcal N_{L,a}(v,w)^{-1}\le10^{12}.
\end{equation}
The assertions hold for the specified identity-log continuation.
Taking \(a_0<r_0\) and smaller radii keeps the first step, as well
as all later steps, inside the previously admitted common local rule.
\end{theorem}
\begin{proof}
Write \(H_{L,a}(s,w)=F_L(A^{(a)}+w+\iota_Ls)\).
Lemma \ref{lem:v36-entry} and \(\|\iota_Ls\|_4\le|s|\) give
uniform bounds for its first two derivatives on a common ball.
By Theorem \ref{thm:v36-pivot},
\(H_{L,a}(0,0)=0\) and \(D_sH_{L,a}(0,0)=N(a)\).
The second-derivative bound therefore gives
\begin{equation}\label{eq:v36-normalperturb}
 \|D_sH_{L,a}(s,w)-N(a)\|
                          \le C(|s|+\|w\|_4),
 \qquad |H_{L,a}(0,w)|\le C\|w\|_4.
\end{equation}
Choose \(R,\rho\) small enough that the first bound is at most
\(1/10\) on \(|s|\le R,\|w\|_4<\rho\), and that this whole
set stays strictly inside the translated holomorphic ball. The map
\[
       \mathcal Q_{v,w}(s)=s-N(a)^{-1}(H_{L,a}(s,w)-v)
\]
has derivative norm at most \(1/2\). Decrease \(\rho,\nu\) so
that \(5(C\rho+\nu)<R/2\). Then
\(|\mathcal Q_{v,w}(0)|<R/2\), and the closed \(R\)-ball is
preserved. The contraction theorem gives its unique fixed point, which
lies in the open ball. The fixed-point estimate gives the displayed
bound on \(|s|\). These arguments hold over the complexifications.
Uniform convergence of the holomorphic iterates gives holomorphic
dependence; Cauchy estimates on smaller parameter balls give all
fixed mixed derivatives with the asserted norms.

Finally \eqref{eq:v36-normalperturb} and \(\|N(a)^{-1}\|<5\)
give a Neumann-series inverse bound at most ten. The determinant is
positive: along the real segment from \((s,w)=(0,0)\) to the solved
point the derivative remains invertible, and \(\det N(a)>0\).
All twelve real singular values are at least \(1/10\), proving
\eqref{eq:v36-tubebound}. Restricting \(a_0<r_0\) leaves positive
room to shrink the first primitive neighborhood into its common
polydisc; the later steps were already admitted by the previous lemma.
\end{proof}

If an independently specified finite coordinate density is
\(m^{(a)}_L(s,w)\,ds\,dw\) in these affine coordinates, its exact
pushforward under \((s,w)\mapsto(v,w)\), within this tube, is
\begin{equation}\label{eq:v36-density}
 \frac{m^{(a)}_L(s_L(a,v,w),w)}
       {\det\mathcal N_{L,a}(v,w)}\,dv\,dw.
\end{equation}
This follows from the block-triangular coordinate Jacobian and includes
the source factor once. A coarse Haar delta convention would require
its separate prescribed conversion. The theorem does not prove that
these linear Coulomb coordinates form a nonsingular nonlinear gauge
slice at every background; it does not construct the corresponding
Faddeev--Popov factor, establish constrained Hessian positivity, or
renormalize \(m^{(a)}_L\). Those are separate requirements for an
actual measured construction. Cofinality for arbitrary sequences
\(w=w_L\) is also not asserted: the cofinal theorem above concerns
the repeated-cell center \(w=0\) and its retained source variables.

\subsection{What the enlarged domain does and does not buy}

The domain enlargement is genuine. For any fixed nonzero \(a\),
however small,
\begin{equation}\label{eq:v36-large}
                         \|A^{(a)}\|_4=|a|L\longrightarrow\infty.
\end{equation}
Thus the centers leave every fixed zero-centered global classical
\(\ell^4\) ball, although the full source chart and its normal
inverse remain uniformly controlled. The mechanism is exact
one-step cancellation of a repeated cell, followed by entry into the
already controlled small-field domain. A global classical norm bound
on the original input is sufficient, but is not necessary, for this
source-coordinate regularity.

This does not make such stripes typical quantum backgrounds. With
the Wilson convention \(S_L=\sum_{x,\mu<\nu}
(1-\tfrac12\Tr U_{x,\mu\nu})\), their action is exactly
\begin{equation}\label{eq:v36-action}
                       S_L(A^{(a)})=L^4(1-\cos(2a)).
\end{equation}
Indeed every plaquette is trivial except those in the \((2,4)\)
plane. Their exponents are \(-2a(-1)^{x_2}E_3\), and there is
one such plaquette at each site. This proves the formula. Fixed
nonzero amplitude therefore has extensive nonzero action. Also, under
the physical link scaling used in the continuum discussion, the
oscillatory physical connection amplitude is of order \(L|a|\).
The source-chart limit in \eqref{eq:v36-limit} is not a convergence
theorem for those physical connections or for their quantum laws.

The translated tube still asks that its perturbation have small
global norm. For the free detail \(w=\beta^{-1/2}Z_L\), V23 gives
\[
     \mathbb P\{\|w\|_4<\rho\}\longrightarrow0
                         \quad\hbox{if }\beta=o(L^2).
\]
No assertion about a shifted or interacting Gaussian is needed for
this observation: it concerns the detail variable used in the tube
itself. Nor does the common inverse-Jacobian bound estimate the
probability of a tube in the normalized constrained interacting law.
The previous typical-fluctuation bottleneck therefore survives.

The next upstream task is consequently concrete: permit cell-to-cell
fluctuations whose global counting norm is not small, prove admission
of their first blocked field by a local probabilistic argument, and
carry the generated conditional density, its gauge factors and its
normal-coordinate determinant through that argument. Periodicity
alone does not give such an invariant family of measured conditional
laws. The present theorem supplies an exact nonperturbative test
family and a quantitative stability neighborhood for that task; it
does not close it by assuming the desired probability coverage.

\subsection{Verification and retained scope}

The independent certificate is
\begin{center}\small
\path{scripts/verify_ym_f14_v36_repeated_cell_chart.py}.
\end{center}
It has 6,108 bytes and SHA--256
\begin{center}\scriptsize\ttfamily
54381572913ACEB2BEA3C2E48BB1585D8DA8AAA78A65251BC0CB065FE577831E.
\end{center}
It differentiates the literal sixteen-tree, reference-anchored rule
by \eqref{eq:v36-jetproduct}, proves all sixteen matrix entries and
all sixteen inverse entries as rational-function identities, and
checks the real determinant by conjugation \(z\mapsto z^{-1}\).
Its determinant expansions verify \eqref{eq:v36-detseries} exactly.

A separate calculation uses V31's rational degree-four BCH engine.
The stripe is assigned degree one and a constant source degree two.
Antisymmetrizing the source removes its square at degree four, so
degrees two, three and four extract \(N(0)\), \([a]N(a)\), and
\([a^2]N(a)\), respectively. Twenty-four literal root evaluations
check all twelve real source columns, for 432 rational coefficient
comparisons. The linear-in-\(a\) transverse response is
\(C_{42}=ia+O(a^3)\); the quadratic diagonal is
\((-1,-1,-1,4/3)a^2\). Exact lattice checks at sides two, four and
eight give respectively 16, 256 and 4,096 nontrivial stripe
plaquettes and verify the Coulomb condition. These finite checks
certify the algebra and conventions. The refinement-uniform and
cofinal conclusions use the analytic proofs above, not extrapolation
from these finite lattices.

The predecessor V35 has 806,703 bytes and SHA--256
\begin{center}\scriptsize\ttfamily
277AB20F43E4C02DFFA3BD312DDA7C8A77D771E7C09C37AF5DE59E558CAA3EAC.
\end{center}
Its body is preserved apart from the title version and this terminal
insertion. V36 replaces V35 as the single active main note. All
archives and companion notes are preserved. The full objective remains
the interacting four-dimensional continuum construction and its
positive physical Hamiltonian gap, as specified in the
\href{https://www.claymath.org/millennium/yang-mills-the-maths-gap/}
{Clay Mathematics Institute problem description}; no claim of resolving
that problem is made here.

\begin{firewall}
The new results are the exact stripe normal matrix, its uniformly
bounded inverse on the specified identity-log branch, a full nonlinear
cofinal retained-source inverse and normal-coordinate factor, and a
refinement-uniform full-dimensional translated Coulomb tube. The
nonlinear source conclusion goes beyond a finite Taylor jet, but
neither typical quantum probability coverage nor the full gauge-fixed
measured law has been established. The interacting Yang--Mills
continuum and the positive physical mass gap remain unproved.
\end{firewall}
\section{V37: the untruncated primitive under Gaussian fluctuations and a quantitative nonlinear remainder}
\label{sec:v37}

\subsection{The next upstream estimate and its exact scope}

V36 removes global smallness of a special background, but the
perturbation tube remains small in the global counting norm. We now
allow all cells to fluctuate under the free transverse Gaussian and
estimate the full nonlinear primitive block on its admitted local
event. Outside that event we use an explicitly declared smooth local
completion solely to make Gaussian expectations well-defined. The
completion is not substituted for the microscopic compact measure.

The new estimate is not another finite-chaos calculation. The complete
smooth nonlinear primitive has a canonical first-chaos projection, and
the remainder contains all higher chaoses. Its covariance decays as
inverse fourth power of distance. After further \emph{linear} blocking
by a factor \(p\), its \(L^q\) norm is bounded by
\(Cq\epsilon p^{-1}\sqrt{1+\log p}\), where
\(\epsilon=\beta^{-1/2}\) and the output has been divided by
\(\epsilon\). A Gaussian rotation argument gives this estimate
without finite-chaos hypercontractivity. In particular, the proof does
not apply a polynomial moment theorem to an infinite chaos expansion.

V25 supplies the transverse Green bound, and V26 supplies the actual
linear averaging rows and the warning about their orientation
geometry. Those results are used, not redone. The first-chaos matrix
below is a source/blocking object, not the complete measured scalar
coefficient of V30 and Shell V16. The supplied suggestions about a
single measured convention continue to be respected: no Wilson,
Haar, ghost or source-normal term is silently removed. The next
scheduled SM/NS checkpoint remains V40.

\subsection{A local completion with an explicitly controlled exceptional event}

Let \(L=2^n\ge2\), and let \(Z_L\) be the same mean-zero Coulomb
Gaussian as in V25. The covariance of its real link coordinates is
\[
 \mathbb E Z_{x,\mu}^{\alpha}Z_{y,\nu}^{\gamma}
          =\delta_{\alpha\gamma}(C_L)_{(x,\mu),(y,\nu)},
 \qquad |(C_L)_{(x,\mu),(y,\nu)}|
                  \le C(1+\operatorname{dist}_L(x,y))^{-2}.
\]
No nonzero Fourier mode is discarded. The covariance is allowed to
be singular in ambient link coordinates.

For each coarse output link \(e=(r,\mu)\), write the actual primitive
as a function \(t_e\) of the finite list of fine link coordinates in
its carrier \(S_e\). The list has at most \(D_0\) real scalar
coordinates, with \(D_0\) fixed, and lies within a fixed distance
of \(2r\). For example the carrier bounds of Shell V12 permit a
fixed \(D_0\le162\) after passing to three real colors. On small
tori repeated coordinates may instead be retained in the list; its
dimension is still bounded, and the resulting covariance is singular.
All arguments below admit either convention.

Fix \(r_*>0\) small enough that \(|A_{S_e}|<2r_*\) stays strictly
inside the common primitive identity-log chart. Choose a real smooth
radial cutoff \(\chi_e\), equal to one for \(|A_{S_e}|\le r_*\)
and zero for \(|A_{S_e}|\ge2r_*\), with the same translated
choice for each orientation. Define
\begin{equation}\label{eq:v37-completion}
 \widetilde t_e(A_{S_e})=\chi_e(A_{S_e})t_e(A_{S_e})
 \quad\hbox{inside its local chart},
 \qquad \widetilde t_e=0\quad\hbox{outside}.
\end{equation}
The cutoff vanishes before the chart boundary, so this is a globally
smooth compactly supported function of the carrier coordinates. It
is equivariant under simultaneous color rotation. Every fixed
derivative is bounded uniformly in \(e,L\), and
\(\widetilde t_e(0)=0\), \(D\widetilde t_e(0)=(B_2)_e\).
The extension is not claimed to be holomorphic or locally gauge
covariant outside the admitted event.

Put \(X_e=(Z_L)_{S_e}\) and define the completed, normalized
nonlinear output
\begin{equation}\label{eq:v37-output}
 Y_{\epsilon,L,e}=f_{\epsilon,e}(X_e),\qquad
 f_{\epsilon,e}(x)=\epsilon^{-1}\widetilde t_e(\epsilon x),
                   \qquad 0<\epsilon\le\epsilon_0.
\end{equation}
It is a complete smooth function, not its Taylor jet. Color invariance
of the Gaussian and equivariance of the output give
\(\mathbb E Y_{\epsilon,L,e}=0\): the only vector fixed by every
color rotation is zero.

\begin{proposition}[Simultaneous agreement with the literal primitive]
\label{prop:v37-event}
Let
\[
 \mathcal E_{\epsilon,L}
       =\{\,|\epsilon X_e|\le r_*\text{ for every coarse link }e\,\}.
\]
There are \(c,C>0\), independent of \(L,\epsilon\), such that
\begin{equation}\label{eq:v37-event}
 \mathbb P(\mathcal E_{\epsilon,L}^{\,c})
                 \le C L^4 e^{-c/\epsilon^2}.
\end{equation}
On this event every first-step logarithm is admitted and
\(Y_{\epsilon,L}=\epsilon^{-1}\mathfrak T_L(\epsilon Z_L)\).
In particular, for each \(\eta>0\), choosing
\(\beta=\epsilon^{-2}\ge C_\eta\log L\) makes the failure
probability at most \(C L^{-\eta}\).
\end{proposition}
\begin{proof}
The Green bound at zero gives a common upper bound on every scalar
coordinate variance. If all \(12L^4\) fine scalar coordinates have
absolute value at most \(r_*/(\epsilon\sqrt{D_0})\), every carrier
belongs to its inner ball, including lists with repetitions. The
one-dimensional Gaussian tail estimate and a union bound prove
\eqref{eq:v37-event}. On that event the cutoff is one and the
underlying primitive chart is admitted. Substitution of
\(\beta\ge C_\eta\log L\), with a sufficiently large constant,
proves the last statement. Independence between cells is not used.
\end{proof}

Two choices of fixed smooth completion agreeing on the inner local
ball differ, for each fixed output link, by
\begin{equation}\label{eq:v37-local-ambiguity}
 \|Y_{\epsilon,L,e}^{(1)}-Y_{\epsilon,L,e}^{(2)}\|_q
                            \le C_q e^{-c_q/\epsilon^2},
                  \qquad 1\le q<\infty.
\end{equation}
Indeed their difference is bounded by \(C/\epsilon\) and supported
on \(|X_e|>r_*/\epsilon\); the finite-dimensional Gaussian tail
absorbs the factor \(1/\epsilon\). This is a local statement.
A sum over an increasing number of outputs must still retain its
volume factor or use a separate spatial estimate.

\subsection{An exact first-chaos filter for the entire nonlinear block}

For each real output component, define its canonical ambient row by
\begin{equation}\label{eq:v37-filter}
 (\mathcal A_{\epsilon,L})_{e,i}
       =\mathbb E\,\partial_i f_{\epsilon,e}(X_e)
       =\mathbb E\,\partial_i\widetilde t_e(\epsilon X_e),
                  \qquad i\in S_e.
\end{equation}
Rows are zero off their finite carriers. Repeated list entries are
aggregated when this row is regarded as a fine-link operator. The
canonical row is defined before quotienting by any Gaussian nullspace;
no covariance inverse or added null row is involved.

\begin{theorem}[Exact nonlinear Gaussian decomposition]
\label{thm:v37-split}
The first Gaussian chaos of the full output is
\(\mathsf P_1Y_{\epsilon,L}=\mathcal A_{\epsilon,L}Z_L\).
Consequently the exact decomposition
\begin{equation}\label{eq:v37-split}
 Y_{\epsilon,L}=\mathcal A_{\epsilon,L}Z_L+R_{\epsilon,L}
\end{equation}
has
\[
 \mathbb E R_{\epsilon,L,e}=0,\qquad
 \mathsf P_1R_{\epsilon,L,e}=0,
 \qquad \mathbb E\nabla f_{\epsilon,e}(X_e)
                           -\mathcal A_{\epsilon,L,e}=0.
\]
The first term is a Gaussian field. The second is orthogonal to every
global first-chaos observable; it need not be independent of the
first term. Color equivariance makes each orientation/spatial entry
of \(\mathcal A_{\epsilon,L}\) a scalar multiple of the color
identity. Its row support and row \(\ell^1\) norm are bounded
uniformly, and
\begin{equation}\label{eq:v37-filter-small}
 \|\mathcal A_{\epsilon,L}-B_2\|_{\mathrm{row},1}
                               \le C\epsilon^2.
\end{equation}
For every fixed nonnegative integer \(m\) there is also the uniform
finite-order expansion
\begin{equation}\label{eq:v37-filter-expansion}
 \mathcal A_{\epsilon,L,e}
 =\sum_{j=0}^{m}\frac{\epsilon^{2j}}{(2j)!}
       \mathbb E D^{2j+1}t_e(0)[X_e^{\otimes 2j},\,\cdot\,]
       +O_m(\epsilon^{2m+2})
\end{equation}
in row norm. This does not assert convergence of the infinite Taylor
series of the smooth completion.
\end{theorem}
\begin{proof}
Realize the entire finite Gaussian as \(Z_L=Tg\), with \(g\) a
standard real Gaussian of dimension equal to its rank; rectangular
\(T\) is allowed. Gaussian integration by parts gives, for each
coordinate \(g_j\),
\[
 \mathbb E[f_{\epsilon,e}(X_e)g_j]
       =\mathbb E\,\partial_{g_j}f_{\epsilon,e}(X_e)
       =\sum_i(\mathcal A_{\epsilon,L})_{e,i}T_{ij}.
\]
This identifies its orthogonal projection onto the span of \(g\).
The mean is zero by color invariance. Subtracting the projection
gives \eqref{eq:v37-split}, its zero first chaos, and zero mean
ambient gradient. Uniform first-derivative bounds and the finite
carrier prove the row norm bound. Simultaneous rotation invariance
of the Gaussian shows that the averaged derivative intertwines the
standard three-dimensional color representation with itself, so
each color block is scalar.

Taylor expand \(D\widetilde t_e(\epsilon X_e)\) at zero through
degree \(2m+1\). All its derivatives at zero are those of \(t_e\).
Every odd term has zero Gaussian expectation. The remainder is
bounded by \(C_m\epsilon^{2m+2}|X_e|^{2m+2}\); the carrier
dimension and its covariance norm are uniformly bounded. This proves
\eqref{eq:v37-filter-expansion}. The case with the first nonzero
correction proves \eqref{eq:v37-filter-small}. This argument uses a
finite Taylor remainder of a globally smooth function, not a
term-by-term expectation of an unproved infinite series.
\end{proof}

In particular the one-step linear correction is the complete finite
Gaussian integral \eqref{eq:v37-filter}, with an exponentially small
local completion ambiguity. The latter follows by applying the same
tail argument to the difference of the first derivatives, which is
supported outside the common inner ball. It is not in general a
scalar orientation multiplier, or an operator determined only by the
already blocked field.

There is a useful normalization check. Let
\(\mathcal S_{\epsilon,L;\mu\nu}\) be the spatial row sum of
\(\mathcal A_{\epsilon,L}\), on a constant fine input of
orientation \(\nu\). At \(L=2\), the earlier mixed-source
calculation gives
\begin{equation}\label{eq:v37-mass}
 \mathcal S_{\epsilon,2}=2I_4+2\epsilon^2\Gamma_2+O(\epsilon^4),
 \qquad
 \Gamma_2=\frac1{9216}
 \begin{pmatrix}
 -5369&0&0&0\\
 0&-4937&144&144\\
 0&0&-4265&336\\
 0&0&0&-3185
 \end{pmatrix}.
\end{equation}
The factor two is necessary: an ambient constant input is twice
\(\iota_2s\). More explicitly the order-\(\epsilon^2\) row
applied to a constant \(b\) is
\(\tfrac12\mathbb E D^3F_2(0)[Z_2,Z_2,2\iota_2b]
=2\Gamma_2b\). For \(L>2\) the analogous primitive correction
uses that scale's incoming covariance; it is not the cumulative
terminal \(\Gamma_L\) of V22. Equation \eqref{eq:v37-mass} is a
calibration of the new exact filter, not a new computation of the
already known matrix.

\subsection{Inverse-fourth-power covariance for the untruncated remainder}

\begin{lemma}[Twice-interpolated Gaussian covariance]
\label{lem:v37-interpolation}
Let \(X,Y\) be finite jointly centered Gaussian vectors, possibly
singular, with cross covariance \(K_{ij}=\mathbb E X_iY_j\).
Let \(u,v\) be smooth scalar functions of at most linear growth
with bounded first and second derivatives. Suppose
\(\mathbb E u(X)=\mathbb E v(Y)=0\) and
\(\mathbb E\partial_i u(X)=\mathbb E\partial_jv(Y)=0\).
Let \(\mathbb E_t\) preserve the two marginal covariances and replace
their cross covariance by \(tK\). Then
\begin{equation}\label{eq:v37-interpolation}
 \mathbb E[u(X)v(Y)]
 =\int_0^1(1-t)\sum_{i,j,k,l}K_{ij}K_{kl}
           \mathbb E_t[\partial_{ik}u(X)\partial_{jl}v(Y)]\,dt.
\end{equation}
\end{lemma}
\begin{proof}
The block covariance with cross term \(tK\) is the convex combination
of the original joint covariance and the independent marginal
covariance, so it is positive semidefinite. First add a positive
multiple of the identity to its diagonal blocks. Differentiation of
the Gaussian density with respect to its covariance, followed by
integration by parts, gives
\[
 \frac{d}{dt}\mathbb E_t[u(X)v(Y)]
       =\sum_{i,j}K_{ij}\mathbb E_t[\partial_i u(X)\partial_jv(Y)].
\]
There is no factor \(1/2\): the two symmetric off-diagonal covariance
entries cancel the usual \(1/2\). Differentiating once more yields
the integrand in \eqref{eq:v37-interpolation}. At \(t=0\), the
zeroth and first derivative are zero by the hypotheses in the
unregularized limit. Taylor's formula with integral remainder gives
the identity. Letting the added diagonal covariance decrease to zero
is justified by Gaussian moment bounds and bounded derivatives.
Thus no inverse of a singular covariance is used. This is the smooth
finite-dimensional covariance-differentiation formula commonly called
Price's theorem; see
\href{https://arxiv.org/abs/1710.03576}{F. Voigtlaender,
\emph{A general version of Price's theorem}} for a broader formulation.
The argument above provides the particular singular extension needed
here directly.
\end{proof}

\begin{theorem}[Covariance gain after removal of the whole first chaos]
\label{thm:v37-covariance}
For any output components at \(e=(r,\mu)\), \(f=(s,\nu)\),
\begin{equation}\label{eq:v37-covariance}
 |\mathbb E R_{\epsilon,L,e}^{\alpha}R_{\epsilon,L,f}^{\gamma}|
       \le C\epsilon^2
                (1+\operatorname{dist}_{L/2}(r,s))^{-4}.
\end{equation}
This estimate includes all higher Gaussian chaoses of the completed
nonlinear primitive.
\end{theorem}
\begin{proof}
As a function of its local Gaussian input, a remainder component is
\[
 u_e(x)=f_{\epsilon,e}(x)-\mathcal A_{\epsilon,L,e}x.
\]
It has zero mean and zero mean ambient gradient. Its Hessian is
\begin{equation}\label{eq:v37-hessian}
             D^2u_e(x)=\epsilon D^2\widetilde t_e(\epsilon x),
             \qquad \sup_x\|D^2u_e(x)\|\le C\epsilon.
\end{equation}
Every cross covariance between the two fixed carriers is bounded
by \(C(1+\operatorname{dist}_{L/2}(r,s))^{-2}\). The change from
fine to coarse distance and the bounded carrier offsets only alter
the constant; for nearby carriers the covariance bound at zero
suffices. Apply Lemma \ref{lem:v37-interpolation}. There are at most
\(D_0^4\) terms; each contains two such cross covariances and two
Hessians of size \(C\epsilon\). This proves
\eqref{eq:v37-covariance}, without an assumption on a largest
nonzero chaos degree.
\end{proof}

An exact covariance identity follows from the global first-chaos
orthogonality:
\begin{equation}\label{eq:v37-covariance-split}
 \operatorname{Cov}(Y_{\epsilon,L})
   =\mathcal A_{\epsilon,L}C_L\mathcal A_{\epsilon,L}^{\mathsf t}
                       +\operatorname{Cov}(R_{\epsilon,L}).
\end{equation}
Color identities are implicit in \(C_L\) here. The remainder is
not independent Gaussian noise. In particular this equality of
covariances is not a closure equation for the full law.

\subsection{Dimension-free moments from the Hessian rather than a chaos cutoff}

\begin{lemma}[A second-order Gaussian moment estimate]
\label{lem:v37-moments}
Let \(g\) be a standard real Gaussian of any finite dimension. If
\(U(g)\) is a smooth scalar function with
\(\mathbb E U=0\), \(\mathbb E\nabla U=0\), and
\(\sup_g\|D^2U(g)\|_{\rm HS}\le H\), then
\begin{equation}\label{eq:v37-moments}
                       \|U\|_q\le CqH,\qquad q\ge2.
\end{equation}
The constant is independent of the dimension. The assertion applies
to the at-most-linear-growth functions used here.
\end{lemma}
\begin{proof}
For a smooth Hilbert-valued function \(h\), an independent standard
Gaussian \(g'\), and \(0\le t\le\pi/2\), the vectors
\(g\cos t+g'\sin t\) and \(-g\sin t+g'\cos t\) are independent
standard Gaussians. Integrating the derivative along this rotation,
then using conditional Jensen and Minkowski, gives
\[
 \|h(g)-\mathbb Eh(g)\|_q
       \le C\sqrt q\,\|\,\|Dh(g)\|_{\rm HS}\,\|_q.
\]
Indeed for any fixed linear map \(A\),
\(\|Ag'\|_{L^q}\le C\sqrt q\|A\|_{\rm HS}\); this follows
by diagonalizing \(AA^{\mathsf t}\) and applying Minkowski in
\(L^{q/2}\) to the sum of squared Gaussian coordinates. The rotation
interval has fixed length. First use this inequality for \(h=U\),
then for the Hilbert-valued \(h=\nabla U\), whose mean is zero.
The two factors \(\sqrt q\) give \eqref{eq:v37-moments}.
Smooth truncation and Gaussian integrability justify the rotation
argument whenever needed; the functions in this section already
have bounded first derivatives and bounded Hessians.
\end{proof}

Let \(L=2pM\), with dyadic \(p\ge1\), and let \(B_p\) be the
actual remaining \emph{linear} block from side \(pM\) to side
\(M\). For a fixed output orientation its nonnegative row weights
\(a_r\) have mass \(p\), maximum at most \(Cp^{-3}\), and support
of at most \(Cp^4\) sites in a box of diameter at most \(Cp\).
They are the tent/box rows of V26, not a normalized probability
average.

\begin{theorem}[A full nonlinear remainder suppressed by later linear averaging]
\label{thm:v37-smoothing}
For each fixed output link and \(q\ge2\),
\begin{equation}\label{eq:v37-smoothing}
 \|(B_pR_{\epsilon,L})_e\|_q
                \le Cq\epsilon\,p^{-1}\sqrt{1+\log p}.
\end{equation}
The constant is independent of \(L,p,M\). There is also a common
\(c>0\) such that
\begin{equation}\label{eq:v37-exponential}
 \mathbb E\exp\left(
       \frac{c p |(B_pR_{\epsilon,L})_e|}
            {\epsilon\sqrt{1+\log p}}\right)\le2.
\end{equation}
Both statements concern the entire smooth nonlinear remainder.
\end{theorem}
\begin{proof}
Work first with one real output component, and set
\(U=\sum_r a_r R_{\epsilon,L,(r,\mu)}^{\alpha}\).
It has zero mean and zero first chaos, so in the representation
\(Z_L=Tg\) it satisfies \(\mathbb E\nabla_gU=0\).
Let \(H_r\) be the ambient fine-coordinate Hessian of the local
remainder, extended by zero off its carrier and with repeated
entries aggregated. Then
\[
 D_g^2U=T^{\mathsf t}\left(\sum_r a_rH_r\right)T.
\]
Writing \(C=TT^{\mathsf t}\), the square of the Hilbert--Schmidt
norm is
\[
 \operatorname{Tr}\left[
       \left(\sum_r a_rH_r\right)C
       \left(\sum_s a_sH_s\right)C\right].
\]
Expansion in carrier coordinates, followed by absolute values,
\eqref{eq:v37-hessian}, and the Green bound gives the pointwise
estimate, uniformly in \(g\),
\begin{equation}\label{eq:v37-hessian-sum}
 \|D_g^2U\|_{\rm HS}^2
 \le C\epsilon^2\sum_{r,s}a_ra_s
                   (1+\operatorname{dist}_{pM}(r,s))^{-4}.
\end{equation}
This estimate does not use a bound on the full covariance operator
norm, which would grow with the volume.

For each \(r\) in the row support, the sum of the inverse-fourth
power kernel over that support is at most \(C(1+\log p)\).
To check this, group sites by periodic maximum-norm distance: a
four-dimensional shell of radius \(j\) has at most
\(C(j+1)^3\) sites, and only radii up to \(Cp\) occur. Small
tori reduce the number of distinct sites, or can be treated with
bounded periodic multiplicity. Hence
\[
 \begin{aligned}
 \sum_{r,s}a_ra_s(1+\operatorname{dist}_{pM}(r,s))^{-4}
 &\le \left(\max_s a_s\right)\left(\sum_r a_r\right)
                      C(1+\log p)\\
 &\le Cp^{-2}(1+\log p).
 \end{aligned}
\]
Combine this with \eqref{eq:v37-hessian-sum} and Lemma
\ref{lem:v37-moments}. Summing the three component estimates gives
\eqref{eq:v37-smoothing} for the Lie-vector norm. Finally expand
the exponential and use the moment bound for integer \(q\), with
\(q^q/q!\le e^q\) times a fixed harmless factor. Choosing \(c\)
sufficiently small proves \eqref{eq:v37-exponential}. The cases
\(p=1\) and \(M=1\) are included; \(1+\log p\) prevents a
spurious zero in the first case.
\end{proof}

In particular the \(p=1\) estimate is a uniform local
\(\|R_{\epsilon,L,e}\|_q\le Cq\epsilon\) bound. It provides
a quantitative all-order replacement for merely writing an
uncontrolled local \(O(\epsilon)\) after the Gaussian first-chaos
projection.

\subsection{A one-step distributional statement on the admitted event}

Define, on the common free Gaussian realization,
\begin{equation}\label{eq:v37-linear-comparison}
 W_{\epsilon,L,p}=B_pY_{\epsilon,L},\qquad
 G_{\epsilon,L,p}=B_p\mathcal A_{\epsilon,L}Z_L.
\end{equation}
The latter is an exactly Gaussian field, with the covariance specified
by its displayed filter. Theorem \ref{thm:v37-smoothing} gives for
each output link
\[
 \|W_{\epsilon,L,p}-G_{\epsilon,L,p}\|_q
                \le Cq\epsilon p^{-1}\sqrt{1+\log p}.
\]
The cross covariance between \(G\) and this error is exactly zero,
not just small, because the error has no first chaos. For any fixed
finite list of output links the same coupling estimate holds with
a constant depending only on that list's size.

On \(\mathcal E_{\epsilon,L}\), \(W\) equals
\(\epsilon^{-1}B_p\mathfrak T_L(\epsilon Z_L)\), one actual
nonlinear step followed by \(\log_2p\) linear steps. Thus, for a
test function \(\varphi\) of one output link with
\(\|\varphi\|_\infty\le1\) and Lipschitz constant at most one,
\begin{equation}\label{eq:v37-law}
 \begin{aligned}
 &\left|\mathbb E\left[
   \varphi\bigl(\epsilon^{-1}B_p\mathfrak T_L(\epsilon Z_L)_e\bigr)
                         \mid\mathcal E_{\epsilon,L}\right]
               -\mathbb E\varphi(G_{\epsilon,L,p,e})\right|\\
 &\hspace{12mm}\le
       C\epsilon p^{-1}\sqrt{1+\log p}+C L^4e^{-c/\epsilon^2}.
 \end{aligned}
\end{equation}
To prove this, compare the unconditioned completed output with \(G\)
by the coupling estimate, and then condition the completed output
on \(\mathcal E_{\epsilon,L}\). For a bounded function that last
change costs at most \(2\mathbb P(\mathcal E_{\epsilon,L}^c)\).
The actual logarithmic block is evaluated only on the event where
it is admitted. This avoids an implicit definition on bad Gaussian
samples.

For calibration, the difference filter
\(B_p(\mathcal A_{\epsilon,L}-B_2)\) has a row supported in a
box of side \(Cp\), total absolute mass at most \(C\epsilon^2p\),
and maximum coefficient at most \(C\epsilon^2p^{-3}\). This
follows from the finite primitive carrier, bounded incidence, and
\eqref{eq:v37-filter-small}. Its \(\ell^{4/3}\) norm is at most
\(C\epsilon^2\); the row dual-Sobolev estimate already used in
V26 therefore gives
\begin{equation}\label{eq:v37-free-calibration}
 \|(G_{\epsilon,L,p}-B_{2p}Z_L)_e\|_q
                                   \le C\sqrt q\,\epsilon^2,
                       \qquad q\ge2.
\end{equation}
Here \(B_pB_2=B_{2p}\) is the actual composition identity, with
the appropriate fine and coarse lattices. This confirms the
normalization and the Gaussian weak-coupling limit of this one-step
experiment. It is not an interacting continuum-limit theorem.

\subsection{Why this estimate cannot simply be iterated as a Gaussian law}

There are two separate missing implications. First,
\(B_p\mathfrak T_L\) in \eqref{eq:v37-law} is not the complete
nested nonlinear block. Subsequent nonlinear steps couple the
first-chaos field to its higher-chaos remainder. Orthogonality of
the two fields does not prevent a first-chaos return from their
products. Already for a scalar standard Gaussian \(g\),
\[
 \mathsf P_1(g^2-1)=0,\qquad
                   \mathsf P_1\{g(g^2-1)\}=2g.
\]
For the actual root this return includes the
\(D^2E^{[2]}[Y,Q]\) term retained in V26. Resetting the nonlinear
remainder to zero after each step would lose precisely such terms
and their accumulated source renormalization. The spatial smoothing
bound proved here applies to later linear averages; it is not a
bound for inserting the remainder into all later nonlinear maps.

Second, all probabilities and projections in this section are under
the free transverse Gaussian. The Wilson interaction, the finite
constraint, the gauge determinant, fine Haar factor and the normal
source-coordinate factor change that law. A tail estimate under the
free reference does not automatically survive its normalized measured
tilt. Neither covariance closeness nor a local completion ambiguity
controls that change of measure.

The next missing object is therefore an all-order, scale-local bound
for the \emph{generated conditional law}, preserving enough mixed
derivative information to estimate these first-chaos returns and
the source-normal density factor. One must retain the non-Gaussian
remainder together with that density, not reinitialize an independent
Gaussian at the next scale. The present result supplies an actual
nonlinear first-step estimate on a high-probability free event, with
constants independent of the number of cells. It is a stronger input
to that problem than a finite jet, but it does not assume the missing
measured stability theorem.

\subsection{Exact checks, preservation, and the remaining goal}

The certificate
\begin{center}\small
\path{scripts/verify_ym_f14_v37_nonlinear_gaussian_block.py}
\end{center}
has 5,627 bytes and SHA--256
\begin{center}\scriptsize\ttfamily
19B9A0AC67C7244C4CB0149B6407B93F8BD65AFAC84A3002DA1C4A3AB50DEA16.
\end{center}
It verifies the projection and twice-interpolated covariance algebra
on a singular rank-two three-coordinate Gaussian. For the explicit
polynomial used in that check, the canonical first-chaos row is
\((3,0,1)\), and the remaining cross covariance under interpolation
is exactly \(73t^2+8t^3+24t^4\). The integral of its second
derivative against \(1-t\) is 105. An independent rectangular
Gaussian representation verifies the whitened Hessian trace identity
without a covariance inverse.

The same script evaluates the literal contracted primitive on all
four constant orientation inputs and confirms the factor
\(2\Gamma_2\) in \eqref{eq:v37-mass}. Twelve exact tests at
\(p=1,2,4\) check \(B_pB_2=B_{2p}\) in all four orientations,
together with the row mass and maximum coefficient. The exact
maximum-norm shell count is \(64j^3+16j\), which verifies the
four-dimensional logarithmic summation convention. These are finite
algebra checks. The smooth completion, Gaussian probability bounds,
all-order remainder estimate, and independence of lattice size follow
from the proofs, not from numerical sampling or the finite checks.

V36 has 834,910 bytes and SHA--256
\begin{center}\scriptsize\ttfamily
6A2006E298E8A98E4B32499044FB6F1E56B7FA7BF52CE9008B168F6643A61FE6.
\end{center}
Its preceding body is preserved, apart from the title version and
this terminal insertion. V37 replaces V36 as the single active main
note; archives, monographs and companion files are preserved. V36
constructed a nonlinear source chart; this section adds a distinct
nonlinear estimate addressing its probability-domain limitation.

\begin{firewall}
The new results concern the complete smooth primitive under the free
Gaussian reference: an exact finite-carrier first-chaos filter, a
high-probability event on which the primitive is literal, an
inverse-fourth-power covariance remainder bound, and dimension-free
all-order moment and exponential estimates after later linear
blocking. They do not establish measured stability through an
unbounded sequence of nonlinear RG steps, an interacting continuum
Yang--Mills state, or a positive physical Hamiltonian gap. The full
goal remains unproved and active.
\end{firewall}
\section{V38: measured local regression, a nonzero Wilson-score return, and transfer to the actual fixed box}
\label{sec:v38}

\subsection{What must change when the law is not Gaussian}

V37 subtracts the complete first Gaussian chaos of a nonlinear
observable. That subtraction equals its averaged derivative only
because the reference law is Gaussian. This section keeps the actual
non-Gaussian local measured law and replaces that identity by its
correct linear regression. The regression coefficient is exponentially
local, although it need not have a finite carrier. A cubic-action
score contributes to its first correction. We compute a nonzero
instance using the literal side-two Wilson action and balanced root.

The existence, convexity, local moments, spatial covariance response,
and wall estimates for one admitted measured shell were already
established in Shell V7 from the f12 local inputs. They are not
reproved or claimed for a new class of generated RG interactions.
The new work is an observable-level consequence for that same law:
an exact orthogonal linear component, a controlled full nonlinear
remainder, and transfer of both back to the original fixed box.
No global likelihood-ratio bound, or volume-independent probability
of the whole box, is assumed. The scheduled companion review remains
V40; the source files are unchanged.

\subsection{The precise measured input and the observable carrier}

Fix a retained field \(c\) in the smaller local chart of Shell V7.
Use its independent fibre coordinates \(z\), minimum \(z_*(c)\),
and \(t=\beta^{-1/2}\). The number of coordinates per spatial
site and all primitive incidences are bounded. Set
\(\eta=(z-z_*)/t\). The already constructed whole-space extension
has normalized law
\begin{equation}\label{eq:v38-law}
 d\mu_t^{\rm ext}(\eta)=Z_t^{-1}e^{-U_t(\eta)}d\eta,
 \qquad
 U_t=\tfrac12\eta^{\mathsf t}H\eta+
          \sum_y\{t^{-2}\widehat r_y(t\eta)-\widehat d_y(t\eta)\}.
\end{equation}
The retained argument \(c\) is suppressed. Here \(\widehat r_y\)
is the cutoff action remainder after degree two, and
\(\widehat d_y\) is the cutoff difference of the complete specified
local log density. Thus the latter retains the prescribed Haar and
inverse-pivot factors in this chart. No additional determinant is
removed by the observable decomposition below.

The inputs from Shell V7, with constants uniform in volume, the
admitted retained field, and \(|t|\le t_0\), are
\begin{equation}\label{eq:v38-stock}
 mI\preceq U_t''\preceq MI,\qquad
 \sup_i\mathbb E_t|\eta_i|^q\le C_q,
 \qquad
 |\operatorname{Cov}_t(F,G)|
 \le C\sum_{i,j}e^{-\nu d(i,j)}
                 \|\partial_iF\|_2\|\partial_jG\|_2.
\end{equation}
The upper Hessian bound follows also directly from the fixed local
cutoff derivatives and bounded incidence in that source. Local noise
derivatives through order four and their first fibre gradients have
uniform fixed-moment bounds. Normalized differentiation gives the
corresponding finite connected-cumulant formulas. These facts concern
\eqref{eq:v38-law}, not the massless all-fine covariance of V37.
The signed family satisfies \(U_{-t}(-\eta)=U_t(\eta)\).

Let \(o_y(v,c)\) be a local smooth scalar observable with
\(o_y(0,c)=0\), fixed carrier and bounded fixed derivatives, completed
smoothly as needed away from its original chart. For example one may
use a local Wilson action or force insertion whose stencil bounds are
already part of Shell V7. Define
\begin{equation}\label{eq:v38-observable}
 F_{t,y}(\eta)=\frac{o_y(t\eta,c)}{t},\qquad
 F_{0,y}(\eta)=B_y\eta,\qquad B_y=D_vo_y(0,c).
\end{equation}
Vector observables of a fixed number of components are handled
componentwise. The root source itself, once fixed as the conditioning
variable, is constant on its fibre and gives only a trivial test;
the nontrivial applications here are varying fine observables and
measured insertions. No new independence assumption on their law is
made.

For a spatial matrix define \(\|K\|_a\) to be the maximum of its
weighted absolute row and column sums, with weight \(e^{a d(i,j)}\).
For rectangular observable/detail matrices the distance joins the
observable carrier center to the detail site. These norms are
submultiplicative when the site multiplicities are bounded.

\subsection{Covariance invertibility and local regression on the extension}

Write
\[
 m_t=\mathbb E_t\eta,\qquad
 \Sigma_t=\operatorname{Cov}_t(\eta,\eta),\qquad
 J_{t;y,i}=\operatorname{Cov}_t(F_{t,y},\eta_i).
\]
All inverses in this subsection act on the independent fibre
coordinates, not on a singular ambient gauge space.

\begin{lemma}[An exponentially local inverse covariance]
\label{lem:v38-inverse}
For a fixed \(a>0\), possibly smaller than \(\nu\),
\begin{equation}\label{eq:v38-inverse}
 M^{-1}I\preceq\Sigma_t\preceq m^{-1}I,
 \qquad \|\Sigma_t\|_a+\|\Sigma_t^{-1}\|_a\le C.
\end{equation}
Its first four noise derivatives, as well as those of its inverse
and of \(J_t\), are bounded in a common smaller weighted norm.
\end{lemma}
\begin{proof}
The upper covariance bound follows from the Poincare inequality
with constant \(m^{-1}\). For the lower bound, integration by parts
on the whole-space law gives
\[
 \mathbb E[(\eta-m_t)(\nabla U_t)^{\mathsf t}]=I,
 \qquad \mathbb E[\nabla U_t(\nabla U_t)^{\mathsf t}]
                         =\mathbb EU_t''\preceq MI.
\]
Cauchy--Schwarz, tested against any vector \(v\), gives
\(|v|^4\le(v^{\mathsf t}\Sigma_tv)M|v|^2\), proving the lower
bound. The covariance response in \eqref{eq:v38-stock} gives
\(|\Sigma_{t;ij}|\le Ce^{-\nu d(i,j)}\).

For completeness, its inverse locality follows without identifying
\(\Sigma_t^{-1}\) with the averaged Hessian. Fix a site \(j\),
put \(W_{ii}=e^{\sigma d(i,j)}\), and conjugate \(\Sigma_t\)
by \(W\). The row and column sums of
\(W\Sigma_tW^{-1}-\Sigma_t\) are at most \(C\sigma\), for
\(0<\sigma\le\nu/2\), since
\(|e^{\sigma u}-1|\le\sigma|u|e^{\sigma|u|}\) and
\(|d(i,j)-d(k,j)|\le d(i,k)\). Choose \(MC\sigma<1/2\).
The Neumann series and \(\|\Sigma_t^{-1}\|_{2\to2}\le M\)
give \(\|(W\Sigma_tW^{-1})^{-1}\|_{2\to2}\le2M\).
Thus \(|(\Sigma_t^{-1})_{ij}|\le2M e^{-\sigma d(i,j)}\).
Summing at a smaller exponential rate proves the weighted bound.

Noise differentiation of a normalized covariance inserts derivatives
of \(U_t\) in connected cumulants. For example,
\[
 \partial_t\operatorname{Cov}_t(A,B)
 =\operatorname{Cov}_t(A',B)+\operatorname{Cov}_t(A,B')
                              -\kappa_t(A,B,U_t').
\]
Through four derivatives there are at most six entries: the two
external observables and at most four local noise insertions.
Their local moments and gradient bounds follow from the stated
stock and the integral Taylor formula for \eqref{eq:v38-observable}.
Shell V6's product-covariance argument, already used in V7, bounds
each connected cumulant by a spatial tree. Summing the internal
noise sites preserves exponential decay between the two external
sites after reducing its rate. There are only finitely many
partitions at these orders. This proves the weighted derivative
bounds for \(\Sigma_t,J_t\). Differentiating the inverse identity,
beginning with \((\Sigma^{-1})'=-\Sigma^{-1}\Sigma'\Sigma^{-1}\),
and using the weighted algebra proves the inverse derivative bounds.
No volume-extensive disconnected moment is estimated in place of
these connected expressions.
\end{proof}

Define the exact measured regression and the averaged derivative by
\begin{equation}\label{eq:v38-regression}
 P_t=J_t\Sigma_t^{-1},\qquad
 A_{t;y,i}=\mathbb E_t\partial_iF_{t,y}.
\end{equation}
They generally differ.

\begin{theorem}[A local orthogonal decomposition without Gaussian replacement]
\label{thm:v38-regression}
On the extension law there is the exact identity
\begin{equation}\label{eq:v38-decomposition}
 F_{t,y}=\mathbb E_tF_{t,y}+P_{t,y}(\eta-m_t)+R_{t,y},
 \qquad
 \mathbb E_tR_{t,y}=0,\quad
 \operatorname{Cov}_t(R_{t,y},\eta_i)=0\quad\hbox{for every }i.
\end{equation}
For some fixed \(a>0\),
\begin{equation}\label{eq:v38-small}
 \|P_t-B\|_a+\|A_t-B\|_a+\|A_t-P_t\|_a\le Ct^2.
\end{equation}
The coefficients \(P_t,A_t\) are even in \(t\), and their
expansions through degree two have \(O(t^4)\) remainders in this
weighted norm. The first term involving \(\eta\) in
\eqref{eq:v38-decomposition} retains its actual non-Gaussian law.
\end{theorem}
\begin{proof}
Equation \eqref{eq:v38-regression} proves the covariance orthogonality
directly. Lemma \ref{lem:v38-inverse} proves the locality of \(P_t\).
The averaged derivative has the fixed observable carrier. Its first
four noise derivatives have the same normalized-cumulant bounds,
now with one local observable and at most four noise insertions.

Under \(t\mapsto-t\), the law sends \(\eta\mapsto-\eta\), and
\(F_{-t,y}(-\eta)=-F_{t,y}(\eta)\). Hence \(\Sigma_t,J_t,P_t\)
are even. The derivative \(D_\eta F_{t,y}=D_vo_y(t\eta)\)
transforms evenly as well, so \(A_t\) is even. At zero the law is
Gaussian with covariance \(H^{-1}\), \(F_{0,y}=B_y\eta\), and
\(P_0=A_0=B\). Taylor's formula with the uniform second and fourth
noise derivative bounds proves \eqref{eq:v38-small} and the stated
remainders. This argument does not identify regression with an
averaged derivative under a non-Gaussian law.
\end{proof}

\subsection{The first correction contains the measured score}

At a fixed retained field write the Gaussian Taylor data as
\[
 U_t=\tfrac12\eta^{\mathsf t}H\eta+tV_1(\eta)+O(t^2),
 \qquad
 F_t=B\eta+tQ(\eta)+t^2T(\eta)+O(t^3),
\]
where \(Q,T\) are the homogeneous quadratic and cubic observable
coefficients. Let \(\gamma\) have covariance \(C=H^{-1}\).
The actual first score is
\begin{equation}\label{eq:v38-score}
 V_1(\eta)=\sum_y\left\{
       \tfrac16D^3 f_y(z_*)[\eta,\eta,\eta]
                              -D\ell_y(z_*)[\eta]\right\}.
\end{equation}
It includes the linear part of the complete specified local log
density. Dropping that part would change the regression coefficient.

\begin{proposition}[Measured regression versus the averaged derivative]
\label{prop:v38-score}
The second coefficients satisfy
\begin{equation}\label{eq:v38-score-coefficient}
 \begin{aligned}
 P_t&=B+t^2\left\{\mathbb E_\gamma DT
                 -\kappa_\gamma(Q,\eta,V_1)C^{-1}\right\}+O(t^4),\\
 A_t&=B+t^2\left\{\mathbb E_\gamma DT
                 -\operatorname{Cov}_\gamma(DQ,V_1)\right\}+O(t^4),\\
 A_t-P_t&=t^2\operatorname{Cov}_\gamma(Q,DV_1)+O(t^4).
 \end{aligned}
\end{equation}
The displayed rows and remainders obey the weighted locality bounds
proved above. Connected covariances include their mean subtractions.
\end{proposition}
\begin{proof}
Separate the linear part exactly:
\[
 J_t=B\Sigma_t+t\operatorname{Cov}_t(Q,\eta)
                         +t^2\operatorname{Cov}_t(T,\eta)+O(t^3).
\]
Gaussian parity gives \(\operatorname{Cov}_\gamma(Q,\eta)=0\).
Its first noise derivative is
\(-\kappa_\gamma(Q,\eta,V_1)\). Gaussian integration by parts gives
\(\operatorname{Cov}_\gamma(T,\eta)=(\mathbb E_\gamma DT)C\).
Multiplication by \(\Sigma_t^{-1}\), and the already proved evenness
and fourth-derivative remainder, proves the first line. Notice that
the perturbation of \(\Sigma_t\) cancels in \(B\Sigma_t\Sigma_t^{-1}\).
For the second line, differentiate the expectation of
\(D_\eta F_t=B+tDQ+t^2DT+\cdots\) under the same normalized law.

Finally \(V_1\) is odd and has Gaussian mean zero. Integration by
parts in
\(\mathbb E_\gamma[(Q-\mathbb E_\gamma Q)\eta V_1]\)
gives the exact row identity
\[
 \kappa_\gamma(Q,\eta,V_1)C^{-1}
       =\operatorname{Cov}_\gamma(DQ,V_1)
                        +\operatorname{Cov}_\gamma(Q,DV_1).
\]
This proves the third line. In particular its last covariance is
not a normalization constant which can be absorbed without changing
the transported observable.
\end{proof}

\subsection{Returning the regression packet to the original measured box}

Let \(\mu_t^{\rm box}\) be the actual normalized law
\(e^{-\beta f+\ell}dz\) on Shell V7's original fixed box
\(\mathcal D\), expressed in \(\eta=(z-z_*)/t\). Use the
original observable there, rather than its arbitrary continuation.
Denote its \(\Sigma,J,P,A\) by a superscript \({\rm box}\).
This is a local-box conclusion; the compact complement is not included.

\begin{lemma}[Local observable and covariance transfer without global conditioning cost]
\label{lem:v38-box-transfer}
For some \(a,b>0\) and all sufficiently small positive \(t\),
\begin{equation}\label{eq:v38-box-transfer}
 \|\Sigma_t^{\rm box}-\Sigma_t^{\rm ext}\|_a
 +\|J_t^{\rm box}-J_t^{\rm ext}\|_a
 +\|A_t^{\rm box}-A_t^{\rm ext}\|_a
                                      \le Ce^{-b/t^2}.
\end{equation}
Consequently the inverse covariances and regressions differ by the
same type of weighted bound. In particular
\(\|P_t^{\rm box}-B\|_a+\|A_t^{\rm box}-P_t^{\rm box}\|_a
\le Ct^2\), and \eqref{eq:v38-score-coefficient} holds on the
original box with the same second coefficients and an \(O(t^4)\)
remainder.
\end{lemma}
\begin{proof}
Apply the exact wall interpolation and same-box potential interpolation
of Shell V7, now with external observable marks. For local \(X,Y\)
independent of wall strength,
\[
 \partial_\lambda\mathbb E_\lambda X
       =-\sum_i\operatorname{Cov}_\lambda(X,\psi_i),\qquad
 \partial_\lambda\operatorname{Cov}_\lambda(X,Y)
       =-\sum_i\kappa_\lambda(X,Y,\psi_i).
\]
Shell V7's connected wall lemma applies to these two- and three-entry
expressions: its proof requires local moment and first-gradient
bounds, not that an external mark literally be a potential derivative.
Use marks \(\eta_i,\eta_j\), or \(F_{t,y},\eta_j\), or a local
entry of \(D_\eta F_{t,y}\). Their required moments are uniformly
bounded in the rescaled coordinates. Their gradients in the original
\(z\) coordinates cost at most a fixed power of \(t^{-1}\).
The wall bound retains both \(e^{-b/t^2}\) and an integrable
\((1+\lambda)^{-21/20}\) factor. Its spatial connected tree joins
the two external centers and the wall site. Sum the wall site,
integrate its strength, and absorb the fixed power of \(t^{-1}\)
into a smaller rarity exponent. The result is exponentially small
with exponential decay between the two external centers.

On the original box the interpolation from the extended potential
to \(\beta f-\ell\) inserts the local defect
\((1-\chi_y)(\beta r_y-d_y)\). The same-box local tail and
covariance estimates of Shell V7 give it exponentially small moments
and first gradients, up to fixed powers of \(\beta\). The same
two- or three-entry connected argument proves the corresponding
observable and covariance comparison. If the observable continuation
differs from the original observable outside its inner chart, their
local difference and its first gradient carry that same rare factor;
the covariance response controls its pairing with a distant mark.
Thus the estimate applies to the original observable, not only to
its completion.

Reducing the spatial exponent once more and summing rows and columns
proves \eqref{eq:v38-box-transfer}. All sums are rooted at an
external mark, so no total-volume factor remains. In particular this
argument never divides by an asserted uniformly positive probability
of \(\mathcal D\).

Lemma \ref{lem:v38-inverse} bounds the extension's inverse in a
weighted matrix algebra. For small \(t\), the perturbation in
\eqref{eq:v38-box-transfer} is small in that same algebra. Its
Neumann series gives the boxed inverse and its exponentially small
difference. The identity \(P=J\Sigma^{-1}\) proves the regression
comparison. This also gives a uniform lower boxed covariance bound,
without using a whole-space score integration-by-parts formula at
the boundary. The extension's expansions transfer because
\(e^{-b/t^2}=O(t^k)\) for every fixed \(k\). No noise derivative
of a moving box boundary was assumed in this argument.
\end{proof}

\subsection{Full nonlinear remainder bounds on the measured box}

In this subsection either the extension or the actual boxed law is
allowed, and \(P,A,m_t,R_t\) always refer to that same law. The
boxed potential after rescaling is uniformly strongly convex on a
convex box. It has the same dimension-independent Poincare and
logarithmic Sobolev estimates as the extension, with constants
depending on the common convexity lower bound. One way to see the
latter is the diffusion entropy identity and the gradient estimate
\(\Gamma(P_su)\le e^{-2ms}P_s\Gamma(u)\); the Neumann boundary
term on a convex box is nonnegative. Approximation by smooth convex
domains handles corners. Integration in \(s\) gives
\(\operatorname{Ent}(h^2)\le2m^{-1}\mathbb E|\nabla h|^2\).

For a centered scalar \(u\), this implies
\begin{equation}\label{eq:v38-Lp}
       \|u\|_q\le\sqrt{(q-1)/m}\,
                           \|\,|\nabla u|\,\|_q,\qquad q\ge2.
\end{equation}
For clarity, differentiate \(\|u\|_q^2\) in \(q\), apply the
logarithmic Sobolev inequality to \(|u|^{q/2}\), and then Holder.
The derivative is at most \(m^{-1}\|\,|\nabla u|\,\|_q^2\).
Integrate from two and use Poincare for the initial value. Smooth
regularization at zeros and truncation give the general case. This
is the usual moment argument of
\href{https://www.intlpress.com/site/pub/files/_fulltext/journals/mrl/1994/0001/0001/MRL-1994-0001-0001-a009.pdf}
{Aida--Stroock, \emph{Moment estimates derived from Poincare and
logarithmic Sobolev inequalities}}; the calculation specifies the
normalization used here.

\begin{theorem}[Measured nonlinear remainder with no finite-chaos hypothesis]
\label{thm:v38-remainder}
For the exact remainder in \eqref{eq:v38-decomposition}, on either
specified law, there are constants independent of volume such that
\begin{equation}\label{eq:v38-remainder}
 \|R_{t,y}\|_q\le Cq|t|,\qquad
 |\operatorname{Cov}(R_{t,y},R_{t,z})|
                        \le Ct^2e^{-a d(y,z)},\qquad q\ge2.
\end{equation}
For every deterministic finite scalar weight vector \(b=(b_y)\),
\begin{equation}\label{eq:v38-sum}
 \left\|\sum_y b_yR_{t,y}\right\|_q
                          \le Cq|t|\,\|b\|_2.
\end{equation}
Thus for a four-dimensional tent/box row with mass \(p\) and
maximum \(Cp^{-3}\),
\begin{equation}\label{eq:v38-smoothing}
 \left\|\sum_y b_yR_{t,y}\right\|_q\le Cq|t|p^{-1},
 \qquad
 \mathbb E\exp\left(\frac{c p|\sum_y b_yR_{t,y}|}{|t|}\right)\le2.
\end{equation}
These are estimates on full nonlinear observables under the measured
law, not under a Gaussian density substituted for it.
\end{theorem}
\begin{proof}
Taylor's integral formula on each local carrier gives
\[
 |F_{t,y}-B_y\eta|\le C|t|\,|\eta_{S_y}|^2,
 \qquad
 D_\eta^2R_{t,y}=tD_v^2o_y(t\eta).
\]
Equation \eqref{eq:v38-Lp} applied to each centered coordinate,
whose mean is uniformly bounded, gives
\(\|\eta_i\|_q\le C\sqrt q\). These coordinate moments and
\(\|P_t-B\|_a\le Ct^2\) prove the first estimate in \eqref{eq:v38-remainder},
including the mean subtraction. Also
\[
 \|\partial_iR_{t,y}\|_2\le C|t|e^{-a d(y,i)}.
\]
Inside the finite carrier use
\(D_vo_y(t\eta)-B_y=O(|t|\,|\eta_{S_y}|)\); outside it only
the exponentially local \(P_t-B\) contributes. Insert these
gradient bounds into the spatial covariance response of the same
law and sum the two exponential kernels. This proves the covariance
estimate after reducing \(a\).

To obtain all moments of a sum, put \(U=\sum_yb_yR_{t,y}\).
Its mean is zero, but its mean gradient need not vanish. The exact
identity is
\[
 \mathbb E\nabla U=b^{\mathsf t}(A_t-P_t),\qquad
                         |\mathbb E\nabla U|\le Ct^2\|b\|_2,
\]
by the row and column bounds. Bounded carrier incidence gives the
pointwise Hilbert--Schmidt estimate
\[
                   \sup_\eta\|D^2U(\eta)\|_{\rm HS}
                                      \le C|t|\|b\|_2=:H_b.
\]
Indeed each matrix entry receives only boundedly many local Hessian
contributions; Cauchy--Schwarz on those contributions and summation
of the entries give \(C t^2\sum_yb_y^2\). The nonlocal linear
regression term has zero Hessian.

Componentwise Poincare gives
\(\mathbb E|\nabla U-\mathbb E\nabla U|^2\le m^{-1}H_b^2\).
Apply \eqref{eq:v38-Lp} to the centered scalar norm of this vector;
its gradient is bounded by \(H_b\). It follows that
\[
 \|\,|\nabla U-\mathbb E\nabla U|\,\|_q
                                        \le C\sqrt q\,H_b.
\]
Another application of \eqref{eq:v38-Lp} to \(U\) proves
\(\|U\|_q\le C(\sqrt q\,t^2+q|t|)\|b\|_2\), hence
\eqref{eq:v38-sum} on the fixed small-noise interval. For the stated
averaging row, \(\|b\|_2^2\le(\max b_y)\sum_yb_y\le Cp^{-2}\).
The exponential moment follows by expanding its series and using
the integer moment estimates, exactly as in V37. The logarithm
present in V37 is absent here because this admitted eliminated-shell
law is gapped and exponentially mixing. This does not imply that
the retained physical theory has a mass gap.
\end{proof}

\subsection{A nonzero score return from the literal Wilson action}

The distinction \(A_t\ne P_t\) is not only formal. We give a
finite calibration on the side-two torus. This particular probe
uses the local zero-harmonic Coulomb law \emph{before} imposing a
nonlinear root constraint. It is not claimed to be that constrained
law. Its purpose is to compute the Wilson contribution to the
observable regression explicitly in the live fine-detail space.

Let \(V_2\) be the mean-zero Coulomb space, \(\Pi_2\) its scalar
projection, and \(C_2\) the scalar transverse covariance, with three
independent colors. Use the local quotient density and the literal
Wilson action in this chart, with an admitted symmetric local
completion for the finite coefficient calculation. At this fixed
regulator the quadratic Hessian on \(V_2\) is positive; a sufficiently
small convex completion gives the local noise expansion used here.
Only its coefficients at zero enter the calibration. The Haar and
logarithmic gauge density have zero linear term at the origin:
their linear derivative would be a color-rotation-invariant linear
functional on the link colors, and no nonzero such functional
exists. Consequently \(V_1=S_3\), the cubic Wilson action. There
is no root-normal determinant in this probe, since a nonlinear root
constraint has not been imposed. This distinction keeps the measured
ledger consistent.

For \(F_2(tZ_2)/t=tQ+t^2T+\cdots\), the linear root coefficient
vanishes on \(V_2\). Define the intrinsic score row, one matching
color at a time, by
\begin{equation}\label{eq:v38-Delta}
 \Delta_{\mu,(x,\nu)}
   =\operatorname{Cov}_\gamma\left(
       Q_\mu^3,(\Pi_2\nabla S_3)_{x,\nu}^3\right).
\end{equation}
The other color blocks are the same diagonal block by rotation
invariance.

\begin{proposition}[An exact nonzero Wilson-score row]
\label{prop:v38-Delta}
Put \(\sigma_j(x)=(-1)^{x_j}\), and choose
\[
 (a_\mu,b_\mu)=(3,4),(3,4),(2,4),(2,3)
                   \quad\hbox{for }\mu=1,2,3,4.
\]
Then
\begin{equation}\label{eq:v38-Delta-value}
 \Delta_{\mu,(x,\nu)}
  =\frac{\delta_{\mu\nu}}{768}
    \{-3\sigma_{a_\mu}(x)-7\sigma_{b_\mu}(x)
                       -2\sigma_{a_\mu}(x)\sigma_{b_\mu}(x)\}.
\end{equation}
These rows already belong to the mean-zero Coulomb space, and
\begin{equation}\label{eq:v38-Delta-norm}
                     \Delta C_2\Delta^{\mathsf t}
                                      =\frac5{12288}I_4.
\end{equation}
In this finite probe the regression and averaged-derivative rows
therefore differ by \(P_t-A_t=-t^2\Delta+O(t^4)\). In particular
their difference, applied to one color of the centered input, has
\(L^2\) norm
\(t^2\sqrt{5/12288}+O(t^4)\) in each output orientation.
\end{proposition}
\begin{proof}
Here is an explicit finite contraction specification for
\eqref{eq:v38-Delta-value}. For a plaquette use its four oriented
Lie inputs \(X_1=A_{x,\mu}\), \(X_2=A_{x+e_\mu,\nu}\),
\(X_3=-A_{x+e_\nu,\mu}\), \(X_4=-A_{x,\nu}\). Its cubic
Wilson coefficient, with \(S=\sum(1-\tfrac12\Tr U_p)\), is
\begin{equation}\label{eq:v38-Wilson-cubic}
               \sum_{1\le r<s<u\le4}X_r\cdot(X_s\times X_u).
\end{equation}
Indeed the linear plaquette logarithm is \(\sum_rX_r\), its
quadratic logarithm is \(\sum_{r<s}X_r\times X_s\), and the
cubic part of \(1-\cos|\log U_p|\) is their scalar product.
For each ordered triple the three cyclic contributions have signs
\(+,-,+\), giving coefficient one in
\eqref{eq:v38-Wilson-cubic}.

Index the 64 scalar links lexicographically by \((x,\mu)\).
Let \(T_\mu\) be the exact antisymmetric quadratic-root bank from
V14, so \(Q_\mu^3=2\sum_{i,j}(T_\mu)_{ij}Z_i^1Z_j^2\).
Expanding \eqref{eq:v38-Wilson-cubic} over the 96 plaquettes gives
384 signed distinct link triples. For each derivative link \(h\),
write
\(\partial_{Z_h^3}S_3=\sum_{k,l}(D_h)_{kl}Z_k^1Z_l^2\).
A triple \(a\,Z_i\cdot(Z_j\times Z_k)\) adds \(+a\) to
\((D_i)_{jk},(D_j)_{ki},(D_k)_{ij}\) and \(-a\) to their
transposes. Gaussian independence of the two colors gives the
ambient row
\begin{equation}\label{eq:v38-contraction}
       d_{\mu,h}=2\sum_{k,l}(C_2T_\mu C_2)_{kl}(D_h)_{kl},
                        \qquad \Delta=d\Pi_2.
\end{equation}
For full specification of the remaining finite input,
\[
 (C_2)_{(x,\mu),(y,\nu)}
 =\frac1{64}\sum_{k\in\{0,1\}^4\setminus\{0\}}
  \frac{(-1)^{k\cdot(x-y)}}{m(k)}
             \left(\delta_{\mu\nu}-\frac{k_\mu k_\nu}{m(k)}\right),
\]
where \(m(k)=\sum_jk_j\) is the Hamming weight. The projection
\(\Pi_2\) is the same sum with prefactor \(1/16\) and without
the outer factor \(1/m(k)\). Thus
\eqref{eq:v38-contraction}, the stated triple rule and the literal
quadratic bank specify only rational finite sums. Their evaluation
gives \eqref{eq:v38-Delta-value}; the certificate below checks all
256 scalar entries and an independent sixth-moment contraction.

Every row in \eqref{eq:v38-Delta-value} has zero mean, has only
orientation \(\mu\), and is independent of \(x_\mu\). Hence it
is already Coulomb. The three Fourier characters in that row have
scalar eigenvalues \(4,4,8\). Its variance is therefore
\[
 \frac{16}{768^2}\left(\frac9{4}+\frac{49}{4}+\frac4{8}\right)
                                      =\frac5{12288}.
\]
Distinct orientation rows are covariance-orthogonal: any shared
character has zero wave-vector components in both of those link
orientations, so the transverse off-diagonal symbol vanishes.
This proves \eqref{eq:v38-Delta-norm}. Proposition
\ref{prop:v38-score}, read in orthonormal coordinates on \(V_2\),
gives the stated row difference and norm. In this probe
\(\operatorname{Cov}_\gamma(DQ,S_3)=0\), also seen from color
parity, but \(\operatorname{Cov}_\gamma(Q,D S_3)\) is the nonzero
row just calculated. Thus the missing score term cannot be repaired
by merely reusing the Gaussian averaged derivative.
\end{proof}

\subsection{Verification and the new boundary of the argument}

The exact certificate is
\begin{center}\small
\path{scripts/verify_ym_f14_v38_measured_regression.py}.
\end{center}
It has 6,416 bytes and SHA--256
\begin{center}\scriptsize\ttfamily
832D1604DED8D235C13B27B15D2CECF9A6DB2644AE169DA853AC4C9B50C0A161.
\end{center}
Eight independent rational quaternion evaluations of the complete
plaquette products check the cubic Wilson polynomial. Twelve direct
sixth-moment columns check the contraction against
\(\mathbb E[Q_\mu^3Z_h^3S_3]\), separately from differentiating
the cubic score. The entire projected row and its covariance matrix
are checked exactly. A further normalized scalar calculation retains
both the partition denominator and mean subtractions; its regression
and averaged-gradient second coefficients are \(-2/35\) and
\(334/35\), with non-Gaussian difference \(48/5\). The scalar
test includes a stabilizing quartic term and does not integrate an
uncut cubic exponential. These checks certify finite algebra, not
the uniform-volume conclusions of the analytic proofs.

The acquired packet now retains the full measured conditional law:
its mean, covariance, exponentially local linear regression, and
orthogonal nonlinear observable remainder. The local wall argument
returns it to the actual fixed box without a global Gaussian
likelihood comparison. This is a substantive extension of V37's
free-reference estimate, but not yet an invariant RG class. In
particular, neither the regression coefficient nor the covariance
specifies the full retained law. Later nonlinear operations and the
next measured density still have to be controlled with their mixed
responses. The shell gap used here is the admitted eliminated-fibre
gap, not a physical Hamiltonian gap.

The remaining upstream task is to prove that successive generated
conditional laws keep the locality, convexity or an adequate
replacement, and observable-response bounds used here, while also
controlling the compact complement. The fixed-box theorem cannot
be applied to a new generated potential merely because its first
few coefficients agree with the current one. Nor does the finite
zero-harmonic probe impose the actual nonlinear root constraint
which it deliberately leaves free.

V37 has 862,290 bytes and SHA--256
\begin{center}\scriptsize\ttfamily
5904EFB35531E3647F08B219C8A1F20C0E650D03936562BC3112F0CDC940A0B9.
\end{center}
Its preceding body is preserved apart from the title version and
this terminal insertion. V38 replaces V37 as the single active main
note; all archives and companion files are preserved.

\begin{firewall}
The new results are a measured, exponentially local regression
packet for full nonlinear local observables, its transfer to the
actual admitted fixed box, dimension-independent remainder bounds
under that law, and an explicit nonzero cubic-Wilson score return.
They do not prove uniform stability of every generated RG law,
global compact entry, an interacting four-dimensional continuum
Yang--Mills state, or its positive physical mass gap. The full goal
remains unproved.
\end{firewall}
\section{V39: the missing quadratic stability test, its Ward obstruction, and a local nonlinear gauge criterion}
\label{sec:v39}

\subsection{Why a local remainder estimate is not yet a kinetic estimate}

V38 gives full nonlinear observable bounds under the admitted local
measured law. To use such results at the next elimination, the
generated action must still have the required quadratic control.
The retained precision is massless in the ultraviolet comparison:
a small absolute Hessian error need not be small relative to its
curl energy. Shell V7 already warns that its Cartesian localization
does not by itself supply the full Ward identity. We make this
remaining test explicit rather than applying that local theorem to
an unverified next-scale action.

There are two new structural results. First, an analytic two-sided
linear Ward kernel factors through two curls except for a completely
antisymmetric linear three-index term. A finite-range self-adjoint
lattice example shows that this exception cannot simply be omitted.
Second, a genuinely local \emph{nonlinear} gauge-invariant interaction
has a curvature-factorized Hessian without that exception. Gauge
averaging converts a controlled compact interaction decomposition
into this form without changing its total action. The resulting
weighted norm gives an explicit sufficient quadratic stability
criterion, including a separate estimate for winding terms.

The old tree--chord curvature-to-defect bound in the closure monograph
concerns compact localization and repair. It is not being counted
again here. The new use of axial fillings is a Hessian bound for
arbitrary local gauge-invariant interaction terms. Likewise, the
all-level Gaussian Schur locality of Shell V4--V5 is inherited; this
section does not claim a new Gaussian fixed point or identify it
with a physical mass gap. The next companion checkpoint is V40.

\subsection{Analytic Ward kernels: an exact double-curl homotopy}

Let \(x\in\mathbb C^4\) lie in a ball about zero. For a vector
\(a\in\mathbb C^4\), define the six-component curl symbol
\[
                (D(x)a)_{ij}=x_i a_j-x_j a_i,\qquad i<j.
\]
All transposes in the analytic identities below are ordinary
transposes. Reality and Hermitian symmetry are imposed separately
when the identity is applied on real momentum.

\begin{theorem}[The complete linear obstruction to analytic double-curl factorization]
\label{thm:v39-homotopy}
Let \(N(x)\) be an analytic \(4\)-by-\(4\) matrix satisfying
\begin{equation}\label{eq:v39-Ward}
                        N(x)x=0,\qquad x^{\mathsf t}N(x)=0.
\end{equation}
Then \(N(0)=0\), and its linear Taylor term is
\[
             L_{ij}(x)=\sum_k T_{ijk}x_k,
\]
where \(T_{ijk}\) is completely antisymmetric. There is an explicit
analytic six-by-six matrix \(W(x)\) such that
\begin{equation}\label{eq:v39-factor}
                  N(x)=L(x)+D(x)^{\mathsf t}W(x)D(x).
\end{equation}
On every smaller ball,
\begin{equation}\label{eq:v39-factor-bound}
 \|W(x)\|\le C\sup_{0\le t\le1}\|D^2N(tx)\|,
 \qquad
 \|N(x)-L(x)\|\le C|x|^2
                         \sup_{0\le t\le1}\|D^2N(tx)\|.
\end{equation}
The construction involves no division by \(|x|^2\) and no singular
transverse projector.
\end{theorem}
\begin{proof}
The lowest-order Ward identity gives \(N(0)x=0\) for every \(x\),
so \(N(0)=0\). At linear order write \(N_{ij}=T_{ijk}x_k\).
The right Ward identity makes \(T\) antisymmetric in its last two
indices, and the left one makes it antisymmetric in its first and
third indices. Together these imply complete antisymmetry.

Write \(S=(N+N^{\mathsf t})/2\), \(A=(N-N^{\mathsf t})/2\).
Both satisfy the two Ward identities. For any matrix \(M\), set
\begin{equation}\label{eq:v39-curvature}
 \mathcal C_{ikjl}[M]
   =\partial_k\partial_l M_{ij}
    -\partial_i\partial_l M_{kj}
    -\partial_k\partial_j M_{il}
    +\partial_i\partial_j M_{kl}.
\end{equation}
This tensor is antisymmetric in \(i,k\) and in \(j,l\). Define,
for \(i<k\), \(j<l\),
\begin{equation}\label{eq:v39-homotopy}
 W_{ik,jl}(x)=\int_0^1
       \left\{t(1-t)\mathcal C_{ikjl}[S](tx)
                  +\frac{1-t^3}{3}\mathcal C_{ikjl}[A](tx)\right\}dt.
\end{equation}
The integrand is analytic and uniformly bounded on smaller balls,
so the same is true of \(W\).

To verify the identity, expand in homogeneous Taylor matrices
\(S_m,A_m\) of degree \(m\). Twice differentiating their Ward
identities and using Euler's identity gives
\[
 \begin{aligned}
 \sum_{k,l}x_kx_l\mathcal C_{ikjl}[S_m]
                                      &=m(m+1)(S_m)_{ij},\\
 \sum_{k,l}x_kx_l\mathcal C_{ikjl}[A_m]
                                      &=(m-1)(m+2)(A_m)_{ij}.
 \end{aligned}
\]
For example the first second-derivative term contributes
\(m(m-1)M_{ij}\); the two middle terms contribute
\(2(m-1)M_{ij}\); and the last contributes
\(M_{ij}+M_{ji}\). This gives both factors. The linear antisymmetric
term has zero second derivative. For \(m\ge2\), the two radial
integrals in \eqref{eq:v39-homotopy} are respectively
\[
 \int_0^1t^{m-2}t(1-t)dt=\frac1{m(m+1)},\qquad
 \int_0^1t^{m-2}\frac{1-t^3}{3}dt=\frac1{(m-1)(m+2)}.
\]
Antisymmetry in each index pair converts contraction with
\(x_kx_l\) into \(D(x)^{\mathsf t}W(x)D(x)\), with pairs
enumerated only once. Thus the factorization holds term by term.
Uniform convergence on smaller balls justifies summation. Finally,
each entry in \eqref{eq:v39-curvature} is a sum of four second
derivatives; the two scalar weights have integrals \(1/6\) and
\(1/4\). These facts and \(\|D(x)\|\le C|x|\) prove the bounds.
\end{proof}

The same formula \eqref{eq:v39-homotopy} applied to an arbitrary
analytic matrix defines a double-curl projection
\(\mathcal P N=D^{\mathsf t}W[N]D\). It is idempotent: its
output has both Ward identities, zero constant and linear terms,
and the theorem reconstructs such an output exactly. This is an
algebraic projection, not an authorization to replace a measured
action by a projected action. Any such replacement would require
its own exact density and counterterm ledger.

For a lattice Hessian symbol \(K(\theta)\), put
\(q_i=e^{i\theta_i}-1\), \(E_{ii}=e^{i\theta_i/2}\), and
\(x_i=2\sin(\theta_i/2)\). The midpoint-frame matrix
\(N(x)=E^{-1}K(\theta)E\), with
\(\theta_i=2\arcsin(x_i/2)\) near zero, satisfies
\eqref{eq:v39-Ward} whenever \(Kq=0\) and \(q^*K=0\) hold
as analytic identities. The change of variables is regular near
zero. If \(K\) is Hermitian on real momentum, the construction
preserves that property. In particular, if the linear term vanishes,
\eqref{eq:v39-factor-bound} gives a low-momentum bound by the
Maxwell curl form. At momenta bounded away from zero, an absolute
symbol bound and the Ward identities give the same relative bound
on transverse polarizations, with a fixed larger constant.

Exact midpoint reflection symmetry \(N(-x)=N(x)\) excludes \(T\).
Positive semidefiniteness of \(N(x)\) on a real neighborhood also
excludes it: testing the quadratic form at \(sx\) and \(-sx\)
and sending \(s\downarrow0\) forces its linear term to vanish.
Neither condition is assumed automatically for a generated ordered-root
packet. In particular the positivity one is trying to propagate
cannot be used circularly as an unproved input.

\subsection{A finite-range Ward kernel with a dangerous linear term}

\begin{proposition}[Linear Ward identities and locality alone do not suffice]
\label{prop:v39-obstruction}
There is a real finite-range, self-adjoint lattice quadratic kernel
with both exact linear Ward identities whose relative Maxwell norm
grows linearly with the side length on its lowest nonzero modes.
It is not claimed to be a nonlinear non-Abelian gauge-invariant
action.
\end{proposition}
\begin{proof}
Use Laurent variables \(z_i=e^{i\theta_i}\), \(q_i=z_i-1\),
\(q_i^\sharp=z_i^{-1}-1\). For indices \(i,j\in\{1,2,3\}\), put
\begin{equation}\label{eq:v39-obstruction}
 K_{ij}(z)=z_i f(z)\sum_{k=1}^3\epsilon_{ijk}(z_k-1),
 \qquad f(z)=\frac18\prod_{k=1}^3(1+z_k^{-1}),
\end{equation}
and set all entries involving orientation four to zero. These are
finite Laurent polynomials with real coefficients. Antisymmetry
gives \(Kq=0\) and \(q^\sharp K=0\), the latter read as a row
identity, because \(q_i^\sharp z_i=-q_i\). Also
\(f(z^{-1})=z_1z_2z_3f(z)\), which proves
\(K(z^{-1})^{\mathsf t}=K(z)\). Thus the real-space kernel is
self-adjoint.

In the midpoint frame it is
\[
 N_{ij}(x)=i\prod_{k=1}^3\cos(\theta_k/2)
                              \sum_{k=1}^3\epsilon_{ijk}x_k.
\]
Its nonzero linear three-index tensor is exactly the obstruction
in Theorem \ref{thm:v39-homotopy}. On the axis
\(\theta=(\vartheta,0,0,0)\), the transverse \((2,3)\) block
has off-diagonal entry \(i\sin\vartheta\) and eigenvalues
\(\pm\sin\vartheta\). The Maxwell symbol there is
\(4\sin^2(\vartheta/2)\) times the transverse identity. Hence
the relative norm is
\begin{equation}\label{eq:v39-growth}
                  \frac12|\cot(\vartheta/2)|.
\end{equation}
At \(\vartheta=2\pi/N\) this is asymptotic to \(N/(2\pi)\).
The example is indefinite and fails midpoint reflection symmetry.
It disproves the inference from only linear Ward identities and
short range to an \(O(|q|^2)\) remainder. It is not evidence for
a parity-odd term in the actual Yang--Mills action.
\end{proof}

A still simpler error occurs if Ward symmetry is not known at all.
The on-site Hessian error \(\delta I\) has range zero and absolute
norm \(\delta\), but its relative norm on the lowest transverse
mode is
\[
                       \frac{\delta}{4\sin^2(\pi/N)}.
\]
With a negative sign it can destroy the infrared lower bound; with
a positive sign it introduces a mass-like term not controlled by
the curl norm. Thus an absolute \(O(\beta^{-2})\) local remainder
does not, by itself, give an \(O(\beta^{-2})\) kinetic remainder
uniformly in \(N\). This example concerns an uncontrolled
symmetry-breaking error, not physical mass generation.

\subsection{Why a local nonlinear gauge-invariant interaction is stronger}

Work first in a non-winding cubical subgraph \(B_R\) of the lattice,
with integer side length \(R\ge1\). Define the linear plaquette curl
by
\[
 (dA)_{x,\mu\nu}=A_{x,\mu}+A_{x+e_\mu,\nu}
                          -A_{x+e_\nu,\mu}-A_{x,\nu}.
\]
All edge and plaquette norms use unnormalized counting measure and
the standard \(\mathfrak{su}(2)\) norm. Gauge invariance in the
next theorem is the full group identity, including transformations
at boundary vertices, not just the two linear Ward identities of
a Hessian.

\begin{theorem}[A local nonlinear gauge criterion for the Hessian]
\label{thm:v39-local}
Let \(V(U)\) be a scalar \(C^2\) gauge-invariant function of links
in \(B_R\), invariant under all vertex \(\SU(2)\) transformations.
Let
\[
 h=\|D^2_A V((e^{A_e})_e)|_{A=0}\|_{2\to2}.
\]
Then its first derivative at the identity is zero, and its Hessian
has an exact finite double-curvature representation. In particular
\begin{equation}\label{eq:v39-local-bound}
                 |D^2V(I)[A,A]|
                         \le96h R^5\sum_{p\subset B_R}|(dA)_p|^2.
\end{equation}
The constant is deliberately nonoptimal. It does not depend on the
ambient torus volume or on positivity of the individual interaction.
\end{theorem}
\begin{proof}
Choose the axial tree with root at one corner of \(B_R\): the path
to a vertex \(x\) first takes all direction-one steps, then all
direction-two steps, and so on. Every configuration near the
identity can be gauged to tree links equal to the identity. Write
\(H_e(U)\) for the resulting chord holonomy. Gauge invariance gives
\(V(U)=v((H_e(U))_{e\ {
m chord}})\). At the identity, the
residual root transformation conjugates every chord simultaneously.
There is no nonzero invariant linear functional on an
\(\mathfrak{su}(2)\) color vector, so \(Dv(I)=0\). This also
proves \(DV(I)=0\).

Let \(Q_eA\) be the first variation of the logarithm of the chord
holonomy. It is the circulation along the tree path to the start
vertex, the chord, and the reversed tree path from its endpoint.
Because \(Dv(I)=0\), the second derivative of the nonlinear chord
map drops out of the chain rule, leaving the exact identity
\begin{equation}\label{eq:v39-chord-Hessian}
                 D^2V(I)[A,A]=(QA)^{\mathsf t}H_v(QA),
                         \qquad H_v=D^2v(I).
\end{equation}
Setting tree links to the identity and varying just the chords is
an isometric coordinate restriction of the original link carrier,
so \(\|H_v\|\le h\).

For an edge \((x,\mu)\), commute its final \(\mu\)-step backwards
through the higher-direction portions of the axial path. Each
commutation supplies one oriented unit plaquette. The resulting
filling of its fundamental cycle has
\(\sum_{\nu>\mu}x_\nu\le3R\) distinct plaquettes, each with
coefficient \(\pm1\). More explicitly, for every \(\nu>\mu\)
and \(0\le s<x_\nu\), use minus the \((\mu,\nu)\) plaquette
whose base has coordinates \(x_j\) for \(j<\nu\), coordinate
\(s\) for \(j=\nu\), and zero for \(j>\nu\). Its boundary sum
is exactly the tree--edge--tree cycle. Thus \(Q=Cd\) for a finite
plaquette-filling matrix \(C\).

Cauchy--Schwarz on each filling gives
\(|Q_eA|^2\le3R\sum_{p\subset B_R}|(dA)_p|^2\). There are at
most \(4R(R+1)^3\le32R^4\) edges. Summing over chords in
\eqref{eq:v39-chord-Hessian} proves
\eqref{eq:v39-local-bound}, and also supplies the exact factorization
\(D^2V(I)=d^{\mathsf t}C^{\mathsf t}H_vCd\). No expansion
of a merely quadratic gauge transformation was substituted for
the nonlinear invariance assumption.
\end{proof}

The theorem excludes the linear obstruction for such an interaction
without requiring a reflection symmetry of its chosen tree or
individual coefficients. For a translation-invariant sum of uniformly
local interactions, its Fourier Hessian is therefore quadratically
soft at zero momentum. This is the additional information missing
from the finite-range example \eqref{eq:v39-obstruction}.

\subsection{Gauge averaging and a checkable interaction norm}

Suppose on a finite periodic lattice the actual scalar compact
action admits a decomposition
\begin{equation}\label{eq:v39-decomposition}
                   W(U)=\sum_X\Phi_X(U),
\end{equation}
convergent in \(C^2\), and \(W\) is gauge invariant. Individual
terms need not be invariant. Associate to each term an owner site
\(z_X\) and integer radius \(R_X\ge1\). When
\(R_X<N/4\), require its support to fit in an unfolded box of side
\(R_X\) rooted at that owner. Terms with \(R_X\ge N/4\) may
wind and will be estimated separately.

For each term define the pure-gauge-orbit Hessian stock
\begin{equation}\label{eq:v39-orbitstock}
 h_X=\sup_g\left\|D_A^2
       \Phi_X\bigl((g_x e^{A_{x,\mu}}g_{x+e_\mu}^{-1})_{x,\mu}\bigr)
                                      \big|_{A=0}\right\|_{2\to2},
 \qquad
 \varepsilon_a=\sup_z\sum_{X:z_X=z}e^{aR_X}h_X.
\end{equation}
The Hessian acts on the original support coordinates. This is a
bound along an entire compact pure-gauge orbit, not a bound only
in one small logarithmic representative.

\begin{theorem}[Gauge-invariant local decomposition gives critical quadratic control]
\label{thm:v39-interaction}
If \(\varepsilon_a<\infty\), there is a gauge-invariant
decomposition of the same \(W\), without enlarging the supports,
whose short, non-winding part obeys
\begin{equation}\label{eq:v39-critical}
 |D^2W_{<}(I)[A,A]|
                  \le C_a\varepsilon_a\|dA\|_2^2.
\end{equation}
Its long part obeys the absolute bound
\begin{equation}\label{eq:v39-winding}
 |D^2W_{\ge}(I)[A,A]|
          \le N^4\varepsilon_a e^{-aN/4}\|A\|_2^2.
\end{equation}
Consequently, on mean-zero Coulomb fields,
\begin{equation}\label{eq:v39-critical-total}
                       |D^2W(I)[A,A]|
                              \le C'_a\varepsilon_a\|dA\|_2^2,
\end{equation}
uniformly in \(N\). Retained harmonic fields are not included in
this last Poincare conversion.
\end{theorem}
\begin{proof}
Average each \(\Phi_X\) over the independent Haar gauge variables
at its incident vertices, writing the result as \(\Psi_X\). The
support does not enlarge, since those vertex transformations act
only on the links on which the term already depends. Haar
invariance makes \(\Psi_X\) gauge invariant. Averaging the total
action over the full gauge group gives
\[
                      \sum_X\Psi_X=W,
\]
because \(W\) itself is invariant. The stated \(C^2\) convergence
justifies interchanging summation, differentiation and Haar
integration. The definition \eqref{eq:v39-orbitstock} gives
\(\|D^2\Psi_X(I)\|\le h_X\). Equivalently, at a pure-gauge
configuration left and right translations act by norm-preserving
adjoint maps in exponential tangent coordinates.

Apply Theorem \ref{thm:v39-local} to each short term. At any fixed
plaquette, the possible owner sites of a radius-\(r\) box containing
it number at most \(C(r+2)^4\). Hence its total curvature
coefficient is bounded by
\[
       C\varepsilon_a\sum_{r\ge1}(r+2)^4r^5e^{-ar}<\infty,
\]
which proves \eqref{eq:v39-critical}. For the long terms, use their
Hessian norms directly. There are \(N^4\) owner sites, and the
sum of \(h_X\) over their long terms is at most
\(N^4\varepsilon_a e^{-aN/4}\). This proves
\eqref{eq:v39-winding}.

On mean-zero Coulomb fields the exact free lattice inequality is
\[
 \|A\|_2^2\le\frac{\|dA\|_2^2}{4\sin^2(\pi/N)}
                                      \le\frac{N^2}{16}\|dA\|_2^2,
                         \qquad N\ge2.
\]
Thus the additional relative coefficient is at most
\(\varepsilon_a N^6e^{-aN/4}/16\), bounded uniformly in \(N\).
This proves \eqref{eq:v39-critical-total}. Gauge averaging has
preserved the total compact action exactly; it has not replaced
the fixed-box law by an assumed gauge-invariant law.
\end{proof}

The winding qualification is essential. For instance a Polyakov
loop potential \(1-\tfrac12\Tr\prod_{x_1=0}^{N-1}U_{x,1}\)
is gauge invariant, but on the constant field \(A_{x,1}=aE_3\)
its value is \(1-\cos(Na)\), with a nonzero quadratic term.
That field has zero plaquette curl. Such a winding interaction
cannot satisfy the short-box conclusion; it belongs in
\eqref{eq:v39-winding} and in the separately retained harmonic
sector. The physical-space proof also avoids promoting a Ward
identity known only on a finite Fourier grid to an analytic
identity at every complex momentum. Finite versus infinite
lattice Poincare identities have their own corrections; see, for
context, \href{https://arxiv.org/abs/hep-lat/0309022}
{Kadoh--Kikukawa--Nakayama, \emph{Solving the local cohomology
problem in U(1) chiral gauge theories within a finite lattice}}.
No non-Abelian continuum conclusion is imported from that work.

\subsection{The resulting stability interface and the missing verification}

Suppose a declared retained reference Hessian \(H_0\) has the
known bound \(\langle A,H_0A\rangle\ge\kappa\|dA\|_2^2\)
on the stated mean-zero gauge quotient. If a perturbing \emph{actual
compact action} has the stock \eqref{eq:v39-orbitstock}, then
Theorem \ref{thm:v39-interaction} gives
\begin{equation}\label{eq:v39-stability}
 \langle A,(H_0+D^2W(I))A\rangle
             \ge(\kappa-C'_a\varepsilon_a)\|dA\|_2^2.
\end{equation}
This is a volume-independent quadratic perturbation criterion. It
is sufficient, not a claim that the currently generated action
already satisfies its hypothesis. It applies to an indefinite
interaction correction as well as a positive one.

In particular, an \(O(\beta^{-2})\) estimate in the orbit-weighted
interaction stock would give an \(O(\beta^{-2})\) relative kinetic
error. Shell V7's existing Cartesian-chart marked norm is not that
stock: it neither bounds each interaction on its whole pure-gauge
orbit nor supplies a gauge-invariant compact interaction decomposition
with the same constants. Haar averaging cannot manufacture those
missing estimates by notation. V38's local wall comparison keeps
the measured law correct on its box, but does not supply this
global orbit extension or the compact complement.

The upstream task is therefore sharper than ``obtain another local
moment bound'': propagate a complete measured interaction decomposition
with the orbit Hessian norm, retain its relevant and marginal
coefficients, and control its long or winding part. The analytic
homotopy gives a direct check of a proposed quadratic Ward argument;
the nonlinear local-gauge theorem gives a sufficient mechanism for
excluding its linear obstruction. Neither theorem controls higher
generated vertices, proves the needed decomposition at every RG
step, or establishes physical clustering.

\subsection{Exact verification and append-only record}

The certificate is
\begin{center}\small
\path{scripts/verify_ym_f14_v39_ward_curvature.py}.
\end{center}
It has 5,582 bytes and SHA--256
\begin{center}\scriptsize\ttfamily
7BAFF9ECB29BBD1FEB1E3360C2C70D1BEC1ABE25CF365E17ADD460495F1BAFC9.
\end{center}
It checks the double-curl homotopy on an exact polynomial Ward
matrix containing symmetric and antisymmetric higher terms plus
a nonzero linear three-form. All sixteen reconstructed entries,
Hermitian symmetry, and the projection's idempotence on a separate
non-Ward input are checked algebraically. The finite-range example
has both Laurent Ward identities and its full self-adjointness
verified, together with the exact axial relative norm.

The axial filling is checked on every edge of boxes of sides one,
two and three: respectively 32, 216 and 768 edges. Their non-tree
chord counts are 17, 136 and 513. Each sparse plaquette boundary
equals the actual tree--edge--tree cycle, and each filling has at
most \(3R\) unit coefficients. These finite checks certify the
algebra and normalization. The analytic and uniform interaction-norm
claims use the proofs above, not numerical extrapolation.

V38 has 892,093 bytes and SHA--256
\begin{center}\scriptsize\ttfamily
723D15FAB08FB41EF27550AD0331E8E6763A101307EEB9CAF7B315E1B9E7EC63.
\end{center}
Its preceding body is preserved apart from the title version and
this terminal insertion. V39 replaces V38 as the single active
main note. Archives and companion files are preserved; V40 is
the next scheduled companion review.

\begin{firewall}
The new results classify the linear obstruction in analytic Ward
Hessians, exhibit an exact finite-range counterexample to an
insufficient stability inference, and prove a nonlinear gauge-local
interaction criterion giving a volume-independent relative curl
bound. The required orbit-weighted norm has not been propagated
for the complete generated Yang--Mills packet. Global compact
entry, an interacting continuum state, and its positive physical
Hamiltonian gap remain unproved.
\end{firewall}
\section{V40: a covariant local shell, uniformly convex activation, and curvature control of its measured remainder}
\label{sec:v40}

\subsection{Companion checkpoint and the precise next object}

The compulsory V40 review finds the same active companions:
SM--Einstein V72, NS V26, and Shell Mixing V16. Their respective
SHA--256 digests are
\begin{center}\scriptsize\ttfamily
F44F9745D43379E1672C5550411B58E97195A4C9278592D221DF5D3F644DC1F5,\\
82C887F8C2C69E4F28FAD7C3DC8DE5B9099CB5A4BF431EBA1FFF3E91935978AD,\\
BF138E63BA826BC030D3A8B87E04FDD3B1EC712D269BB064674E96509B46E710.
\end{center}
SM V72 derives its covariance flow in a fixed-mode, fixed-time
many-copy limit; it does not bound an increasing Yang--Mills cutoff.
NS V26 gives the displayed rank-one Oseen contractions, not a
Yang--Mills gauge-local interaction norm. Those results do not remove
V39's missing hypothesis. Shell V16 retains the complete measured
one-loop convention. Its proposed noncommuting quartic calculation
has already been completed in the present V30 and is not repeated.
The supplied suggestions were reread with that later result in view.

The new target is narrower than an all-scale claim but directly
addresses the identified obstruction. We construct an explicitly
coarse-gauge-invariant \emph{local contribution to the same completed
compact pushforward}. For this radial local contribution of the
primitive parity-tree block, we turn Shell V7's connected remainder
estimate into a decomposition with the pure-gauge-orbit Hessian
stock of V39. This proves a relative curl bound for its exact
two-loop remainder at the identity.

The radial domain below is not Shell V7's Cartesian cube. Its
definition, its measure, and the required extension of the wall
argument are given explicitly. No equality with that cube or with
the full compact partition function is asserted. In particular,
the compact complement is not discarded by redefining the object.

\subsection{A background frame in the actual forty-five-group fibre}

Use the monograph's fixed parity tree in each sixteen-site cell.
There are seventeen internal non-tree groups and twenty-eight
nonpivot boundary groups per cell; the four other boundary groups
are the solved pivots. Let \(i\) index these forty-five independent
groups, and let \(a(i),b(i)\) be their endpoint cells. If a group is
internal, set \(B_i(C)=I\). If it crosses the coarse edge \(e\),
set \(B_i(C)=C_e\). Write
\begin{equation}\label{eq:v40-frame}
                  Y_i=B_i(C)e^{\eta_i},\qquad
                           \eta_i\in\mathfrak{su}(2).
\end{equation}
This is a change of fibre coordinates, not a change of the balanced
root source \(C\). At \(\eta=0\), the fourteen nonroot products
over each coarse edge all equal \(C_e\); its solved pivot is
therefore \(C_e\) as well.

For a coarse gauge transformation \(g\),
\[
 C_e^g=g_{e^-}C_eg_{e^+}^{-1},\qquad
 B_i(C^g)=g_{a(i)}B_i(C)g_{b(i)}^{-1},\qquad
 \eta_i^g=\operatorname{Ad}_{g_{b(i)}}\eta_i.
\]
Thus \eqref{eq:v40-frame} intertwines the actual endpoint action.
Each relative root product transforms by conjugation at the terminal
cell. The completed logarithm is conjugation equivariant, so
uniqueness of the compact pivot transports the solved pivot too.
Its inverse Haar Jacobian is unchanged: the input and output
translations preserve Haar measure. This uses the existing pivot
theorem, including its two duplicate root paths, not a linearized
substitute for it.

Let \(J_G(\eta)\) be the fine exponential Haar density and write
\begin{equation}\label{eq:v40-measure}
 \begin{aligned}
 f(C,\eta)&=S_0(Y(C,\eta),B(Y(C,\eta),C)),\\
 \ell(C,\eta)&=\log j(Y(C,\eta),C)+\sum_i\log J_G(\eta_i).
 \end{aligned}
\end{equation}
Here the symbol \(B(Y,C)\) denotes solved pivots, distinguished by
its arguments from the elementary frame \(B_i(C)\).
Every fine plaquette term of \(f\), every pivot-density term, and
every Haar term in \(\ell\) is a local scalar invariant under the
simultaneous transformation of \(C\) and \(\eta\). All carriers
and incidences are fixed. These statements hold before summing over
the lattice.

Choose a fixed sufficiently small \(r>0\), and put
\begin{equation}\label{eq:v40-radial-law}
 \mathcal D_r=\prod_i\{\eta_i:|\eta_i|<r\},\qquad
 Z_\beta^{\rm rad}(C)=\int_{\mathcal D_r}
                  e^{-\beta f(C,\eta)+\ell(C,\eta)}\,d\eta,
 \qquad W_\beta^{\rm rad}=-\log Z_\beta^{\rm rad}.
\end{equation}
The density is relative to \emph{coarse Haar}. No additional coarse
Haar factor is to be inserted. At fixed \(C\), multiplication by
\(B_i(C)\) preserves fine Haar. In the joint variables \((Y,C)\)
the coordinate change is triangular, so it introduces no extra
source determinant either. The actual inverse-pivot density in
\eqref{eq:v40-measure} is retained once.

The domain is invariant under the blockwise adjoint action. Hence
\(Z_\beta^{\rm rad}(C^g)=Z_\beta^{\rm rad}(C)\), even for large
coarse gauge transformations. This is an exact finite-volume
identity. The product cube \([-r,r]^{3\times45|\Lambda|}\) lacks
that invariance; it is not silently being identified with
\(\mathcal D_r\).

The original completed compact integral splits exactly into the
image of \(\mathcal D_r\) under \eqref{eq:v40-frame} and its
complement. Both contributions are positive and gauge invariant.
Their ratio is not bounded here. Although its description in the
old \(Y\) chart moves with \(C\), the domain in
\eqref{eq:v40-radial-law} is fixed. All source derivatives of the
frame occur in \(f\) and \(\ell\); no moving-boundary score has
been suppressed.

\subsection{Activation which does not destroy the eliminated-space floor}

On a sufficiently small fixed maximum-norm chart \(C=e^c\), the
primitive stencils in \eqref{eq:v40-measure} have the inherited
uniform derivative bounds. At \((c,\eta)=(0,0)\), their fibre
Hessian is the actual parity-tree free Hessian. Its positive floor
and fixed incidence give, after decreasing \(r\) and the retained
radius \(\delta\),
\begin{equation}\label{eq:v40-local-floor}
 hI\preceq f_{\eta\eta}(C,\eta)\preceq MI,
 \qquad \|\ell_{\eta\eta}(C,\eta)\|\le L,
           \quad \eta\in\overline{\mathcal D_r},\quad |c|_\infty<\delta,
\end{equation}
with positive constants independent of volume. The frame is local
and equals the old free frame at zero, so this assertion follows
by bounded-incidence perturbation of that same Hessian. It does not
use the false infrared floor for all fine gauge modes.

Assign a scalar activation \(s_i\in[0,1]\) to each independent
group, let \((D_s\eta)_i=s_i\eta_i\), and fix \(h_0>0\). Define
\begin{equation}\label{eq:v40-activation}
 f_s(C,\eta)=f(C,D_s\eta)+\frac{h_0}{2}
                         \sum_i(1-s_i^2)|\eta_i|^2,
 \qquad \ell_s(C,\eta)=\ell(C,D_s\eta).
\end{equation}
The integrals always use \(d\eta\) on \(\mathcal D_r\).
There is no factor \(\prod_i s_i^3\): this is an interpolation of
potentials, not the singular change of variables \(\eta\mapsto
D_s\eta\). Its endpoint \(s=\boldsymbol1\) is exactly
\eqref{eq:v40-radial-law}. At binary \(s\), inactive groups are
pinned at zero in the physical stencils and have independent
auxiliary Gaussian integrals. These Gaussian factors are kept in
the normalization.

\begin{lemma}[Uniformly convex, covariant activation family]
\label{lem:v40-activation}
For every activation pattern,
\begin{equation}\label{eq:v40-activation-floor}
 \begin{aligned}
 (f_s)_{\eta\eta}
   &=D_s f_{\eta\eta}(C,D_s\eta)D_s+h_0(I-D_s^2)\\
   &\succeq \min(h,h_0)I.
 \end{aligned}
\end{equation}
Its upper and finite-range bounds are uniform. The measured
potential \(\beta f_s-\ell_s\) has a uniform floor \(m\beta I\)
for all sufficiently large \(\beta\). Its classical and measured
minima lie in \(\prod_i\{ |\eta_i|<r/2\}\), after a common
decrease of \(\delta\). The classical minimum and its fixed-order
responses to coarse or activation marks have exponential tree
locality, uniformly over all activation patterns. The integrals
and their stationary coefficients are coarse-gauge invariant.
\end{lemma}
\begin{proof}
The displayed Hessian identity is the ordinary chain rule, since
\(D_s\) is constant in \(\eta\). Its quadratic form is bounded
below by \(h|D_sv|^2+h_0|(I-D_s^2)^{1/2}v|^2\), which proves
the floor. The density Hessian is bounded uniformly and costs only
\(L\). Range, incidence, and all needed local derivatives remain
bounded; an activation derivative inserts a local factor of
\(\eta_i\), not a volume-sized sum.

At \(c=0\), the classical minimum is zero for every \(s\).
The inverse of the Hessian at zero has uniformly summable
exponential rows, by the same finite-range coercive inverse
argument as for the original fibre. Its maximum-norm implicit
function argument therefore works uniformly on the compact
activation cube. Concretely, write the critical equation using
that inverse and the nonlinear remainder of the gradient. The
remainder has maximum-norm Lipschitz constant \(O(r+\delta)\)
by fixed incidence; decreasing those radii gives a contraction.
The inhomogeneous gradient at zero is \(O(\delta)\). This puts
the solution in the inner half-domain. The measured minimum is an
additional \(O(\beta^{-1})\) maximum-norm perturbation. Strong
convexity gives uniqueness.

Repeated differentiation of the critical equation separates one
inverse Hessian applied to a sum of local terms involving lower
responses. Joining its exponentially decaying links, retaining
half their decay for the marked terminal tree, and summing the
remaining links proves the asserted response bounds at each of
the finitely many required orders. The number of marks, rather
than the number of activated groups, determines the constants.
Finally, \(D_s\) commutes with each block adjoint action and the
pinning terms are radial. All changes of variables have determinant
one. Uniqueness transports the minimum, Hessian, and density, and
proves covariance of the stationary coefficients.
\end{proof}

\subsection{The radial wall and the four-mark remainder estimate}

Let \(n=135|\Lambda|\) be the real fibre dimension. For the
activated radial integral write its exact expansion as
\begin{equation}\label{eq:v40-expansion}
 W_{\beta,s}^{\rm rad}(C)
 =\beta G_s(C)+\frac n2\log(\beta/2\pi)
                  +Q_s(C)+\beta^{-1}\mathcal B_s(C)
                  +R_{\beta,s}(C),
 \qquad Q_s=\tfrac12\log\det H_s-\ell_s(C,\eta_{*,s}).
\end{equation}
The coefficient \(\mathcal B_s\) is the complete normalized
Laplace coefficient, including the density and moving minimum.
Equivalently, if the rescaled exponent has first two noise terms
\(tV_1+t^2V_2\), then
\(\mathcal B_s=\mathbb E V_2-\tfrac12\operatorname{Var}(V_1)\).
This is Shell V7's coefficient in the present coarse-Haar convention.
The residual in \eqref{eq:v40-expansion} is defined by the exact
integral, not by an uncut cubic approximation.

\begin{lemma}[Radial-domain and activation extension of the measured bound]
\label{lem:v40-remainder}
There are \(a>0\) and \(K_k<\infty\), independent of volume,
activation pattern, and \(\beta\ge\beta_0\), such that
\begin{equation}\label{eq:v40-marked}
 \sup_{c,s}\sup_{u_1}\sum_{u_2,\ldots,u_k}
 e^{a\tau(u_1,\ldots,u_k)}
       |\partial_{u_1}\cdots\partial_{u_k}R_{\beta,s}(e^c)|
                         \le K_k\beta^{-2},\qquad 1\le k\le4.
\end{equation}
Each mark is either a coarse link/color coordinate or one of the
local activations. Both types are anchored at their actual cells.
In particular two activation and two coarse marks fit within the
four-mark budget; no sixth retained derivative is being inferred
from a fourth-derivative estimate.
\end{lemma}
\begin{proof}
We verify the changes to Shell V7 rather than applying its cube
theorem to a different domain. Lemma \ref{lem:v40-activation}
supplies its uniform local analytic stencils, interior minimum,
coercivity, and marked-center estimates, now with \((c,s)\) as the
retained parameter set. The usual centered stencil cutoff gives
a whole-space strongly convex extension: its nonlinear action
Hessian is \(O(r)\), its density Hessian is bounded, and the
quadratic part has the uniform floor just proved. On the radial
domain, interpolation with the actual potential remains strongly
convex and has the same interior measured minimum. The cutoff can
be chosen to agree on a common neighborhood of that minimum.
All constants are uniform in \(s\).

Only the hard-wall step changes geometrically. For each
three-component group use
\[
 \psi_i(\eta)=\tfrac12(|\eta_i|-r)_+^2,
 \qquad E_{\beta,s}(C,\eta)+\lambda\sum_i\psi_i(\eta),
                                                        \quad \lambda\ge0.
\]
These walls are convex and fixed in \((c,s)\). Outside the ball,
with \(n_i=\eta_i/|\eta_i|\), their Hessians are
\begin{equation}\label{eq:v40-radial-Hessian}
 D^2\psi_i=n_in_i^{\mathsf t}
       +\frac{|\eta_i|-r}{|\eta_i|}(I-n_in_i^{\mathsf t})
                         \succeq n_in_i^{\mathsf t}.
\end{equation}
The gradient of \(\psi_i\) lies in this normal direction. The
covariance resolvent's quadratic form therefore gives its wall
gradient the additional \((m\beta+\lambda)^{-1/2}\) factor.
For clarity, the source is in the range of the multiplication
projection \(1_{|\eta_i|>r}n_in_i^{\mathsf t}\); the vector
resolvent is coercive by \(m\beta I+\lambda\) times this
projection. Its other leg costs \(\beta^{-1/2}\). The spatial
weight conjugation is scalar on each three-component block and
commutes with the wall Hessian. Its perturbation affects only
the original finite-range off-diagonal Hessian. The same reduced
exponential weight thus preserves both estimates. No angular
derivative of the projection is taken.

The wall-layer estimates also have the same powers as in one
component, up to harmless powers of \(\beta\):
\begin{equation}\label{eq:v40-radial-layer}
 \begin{aligned}
 \|\psi_i\|_q&\le C_q\beta^{b_q}e^{-c_q\beta}
                         (\beta+\lambda)^{-1-1/(2q)},\\
 \|\nabla\psi_i\|_q&\le C_q\beta^{b_q}e^{-c_q\beta}
                         (\beta+\lambda)^{-1/2-1/(2q)}.
 \end{aligned}
\end{equation}
Here is the finite-dimensional normal-coordinate reason for the
unchanged powers. The inherited local concentration proof uses
strong convexity and weighted inverse rows, and continues to hold
for the positive block walls \eqref{eq:v40-radial-Hessian}.
It puts each fixed block near its inner minimum with all centered
moments \(O_q(\beta^{-q/2})\), uniformly in \(\lambda\).
The same statement holds when this one block is fixed at zero:
the remaining Hessian is a principal restriction, with the same
floor, local bounds, and interior minimum argument.

For its marginal density, subtract the explicit term
\(\lambda\psi_i\) from the effective potential. Integrating the
other variables preserves the \(m\beta\) strong convexity in this
fixed block, by the convex marginal argument. Its gradient at zero
is \(O(\beta\delta+\sqrt\beta+1)\), by the preceding conditional
moments and the local gradient formula. For small fixed \(\delta\)
and large \(\beta\), its outward derivative at \(|\eta_i|=r\)
is therefore positive. A layer of width \(\beta^{-1}\) just
inside that sphere, the local gradient bound there, and the local
exponential tail bound give surface density at most
\(C\beta^b e^{-c\beta}\). Strong convexity along every outward
ray then bounds the density a distance \(u\ge0\) outside by
this surface bound times \(e^{-(m\beta+\lambda)u^2/2}\).
The radial volume factor is \((r+u)^2\). Integrating \(u^{2q}\)
or \(u^q\), and using \(\beta+\lambda\ge1\), proves
\eqref{eq:v40-radial-layer}. This is a codimension-one layer;
one must not replace its exponent by a three-dimensional tail
exponent. Fixed powers of \(\beta\) can be absorbed by reducing
\(c_q\).

Consequently the connected wall estimate through five entries
has the same integrable strength power \((1+\lambda)^{-21/20}\)
as in Shell V7. In a differentiated product, either the wall
itself is differentiated, when its normal resolvent gain is used,
or it is not, when its stronger undifferentiated layer moment is
used. In a cumulant partition whose changed moment omits the wall,
the remaining moment containing the wall supplies that same rare
factor. The longest-edge cut retains the spatial tree. These are
all cases in the five-entry argument, and none uses a scalar
rather than a three-component wall except for the layer and
resolvent estimates just checked. Integration in \(\lambda\)
therefore returns to the actual product of balls with exponentially
small anchored errors through four marks. The actual/extended
same-domain interpolation has the same rare local defects and
convex-domain covariance proof, so its error is exponentially
small as well. Smooth convex approximations justify the
\(C^{1,1}\) walls; no retained derivative differentiates them.

For the whole-space part, signed noise \(t=\beta^{-1/2}\) and
reflection \((t,\eta)\mapsto(-t,-\eta)\) give the same even
normalized correction \(\mathcal W(t,c,s)\) as Shell V7. Its
fourth-noise derivative with at most four retained marks is a
sum of connected moments with at most eight local entries. The
uniform local moments, spatial covariance responses, and marked
minimum bounds established above give their anchored tree sums
exactly as in that proof. Activation marks cause only additional
local polynomial factors of bounded degree, whose fixed moments
are included in these estimates. Taylor's exact remainder is
\[
 \mathcal W(t,c,s)-t^2\mathcal B_s(c)
       =\frac{t^4}{6}\int_0^1(1-v)^3
                         \partial_t^4\mathcal W(vt,c,s)\,dv.
\]
Differentiate this identity in the marked parameters, sum its
connected bounds, and add the two exponentially small domain
errors. This proves \eqref{eq:v40-marked}. The estimates are for
normalized integrals; an absolute extensive moment is never used.
\end{proof}

\subsection{From connected responses to local invariant interaction germs}

Let \(R_S\) denote \(R_{\beta,s}\) at the binary activation
\(s=1_S\), and define
\(\overline R_S(C)=R_S(C)-R_S(I)\).
At \(S=\varnothing\), the integral is the product of auxiliary
truncated Gaussians times
\(e^{-\beta f(C,0)+\ell(C,0)}\). Its residual is a vacuum
constant, so \(\overline R_\varnothing=0\).

Order the independent groups once and for all. For group \(i\),
let \(P_i\) be its predecessors. If \(B_R(i)\) is the periodic
cell-graph ball of integer radius \(R\), put
\begin{equation}\label{eq:v40-increments}
 \begin{aligned}
 A_{i,R}&=P_i\cap B_R(i),\\
 F_{i,R}(C)&=\overline R_{A_{i,R}\cup\{i\}}(C)
                               -\overline R_{A_{i,R}}(C),\\
 \Phi_{i,0}&=F_{i,0},\qquad
 \Phi_{i,R}=F_{i,R}-F_{i,R-1}\quad(R\ge1).
 \end{aligned}
\end{equation}
Several independent groups at the same cell cause only a fixed
multiplicity. At radius equal to the finite graph diameter, all
predecessors have been included.

\begin{theorem}[Local gauge-invariant remainder germs]
\label{thm:v40-decomposition}
The definitions above give the exact finite-volume identity
\begin{equation}\label{eq:v40-telescope}
          R_{\beta,\boldsymbol1}(C)-R_{\beta,\boldsymbol1}(I)
                              =\sum_i\sum_{R\ge0}\Phi_{i,R}(C),
\end{equation}
where the radius sum terminates. Each term depends only on coarse
links in a ball of radius \(R+s_0\) around \(i\), for a fixed
stencil allowance \(s_0\). It is a gauge-invariant germ on a
neighborhood of the pure-gauge orbit of those links, including
transformations at its boundary vertices. For some \(a_1>0\),
\begin{equation}\label{eq:v40-orbit-norm}
 \sup_g\left\|D_c^2\Phi_{i,R}((g_xe^{c_e}g_y^{-1})_e)
                                      \big|_{c=0}\right\|_{2\to2}
                              \le K\beta^{-2}e^{-a_1R}.
\end{equation}
Constants are uniform in the finite torus. The decomposition is
not asserted to represent the remainder outside its admitted
stationary-branch domain.
\end{theorem}
\begin{proof}
First telescope over \(R\), then over the fixed group ordering.
This proves \eqref{eq:v40-telescope} for arbitrary set functions;
independence of neighboring active groups is not assumed.

If only groups in \(S\) are active, stencils not meeting \(S\)
are independent of every integration variable. Their factors
cancel when the integral is divided by the all-inactive integral.
Only stencils meeting \(S\), and their fixed retained carriers,
remain. The inactive Gaussian normalizers also cancel. The same
locality follows for differences of stationary coefficients,
either directly from their finite active-variable formulas or
by uniqueness of their fixed-volume asymptotic coefficients.
Thus \(F_{i,R}\), and hence \(\Phi_{i,R}\), have the stated
finite support. All their exact integrals are gauge invariant;
uniqueness and covariance of the local minima make their
coefficients invariant on the gauge-saturated branch domain.

The case \(R=0\) follows by integrating one activation derivative
and using two coarse marks in Lemma \ref{lem:v40-remainder}.
For \(R\ge1\), add the groups of
\(A_{i,R}\setminus A_{i,R-1}\) one at a time. The difference
of the increment at \(i\) is exactly the double integral of
\(\partial_{s_i}\partial_{s_j}R_{\beta,s}\) for the newly
added group \(j\). Two coarse derivatives of this expression
use four marks in total. Their tree joins \(i,j\) and the two
coarse links, and therefore has length at least \(d(i,j)\ge R\).

One minor uniformity point matters: the intermediate activation
patterns vary with \(j\), so a supremum cannot simply be moved
inside a previously summed norm. Instead
\eqref{eq:v40-marked} first gives the uniform pointwise bound
\(K\beta^{-2}e^{-a\tau(i,j,e,f)}\). Retain half the decay
for \(R\) and sum the other half over \(j,e,f\). With one
anchor fixed, exponential tree sums of three further terminals
on a four-dimensional lattice are finite, uniformly on tori.
This proves
\[
       \sum_{e,f}\sup_{|c|_\infty<\delta}
                  |\partial_e\partial_f\Phi_{i,R}(e^c)|
                         \le K\beta^{-2}e^{-a_1R},
\]
after decreasing the rate once more if needed. Color multiplicity
is fixed. The sum bounds the matrix operator norm. Finally exact
gauge invariance gives
\(\Phi_{i,R}((g_xe^{c_e}g_y^{-1})_e)=\Phi_{i,R}((e^{c_e})_e)\)
on the corresponding germ. It transports this Hessian bound to
every pure-gauge representative with no loss.
\end{proof}

Only a germ is needed for V39's local Hessian theorem. To be
explicit, fix one non-winding carrier box, use its axial tree,
and gauge nearby links to tree links equal to \(I\). For that
fixed finite box this operation stays in the admitted germ on a
possibly smaller neighborhood of \(I\). Residual simultaneous
conjugation makes the first chord derivative zero. The chain rule
and plaquette fillings in Theorem \ref{thm:v39-local} consequently
apply without modification. They require no uniform neighborhood
size as the box grows, since the Hessian is evaluated at \(I\).
This does not prove a common nonlinear domain for the sum of
all carriers. If desired an individual germ can be extended
smoothly using conjugation-invariant cutoffs of its tree-chord
holonomies, but such extensions are not substitutions in the
physical compact action.

\subsection{The resulting volume-independent quadratic estimate}

\begin{theorem}[Critical quadratic control of the radial local two-loop remainder]
\label{thm:v40-critical}
For the primitive completed root-path block, with the specified
radial contribution and its exact measured remainder, there is
\(K<\infty\) independent of \(\beta\ge\beta_0\) and torus
side \(N\ge2\) such that
\begin{equation}\label{eq:v40-critical}
       |D^2R_{\beta,\boldsymbol1}(I)[A,A]|
                          \le K\beta^{-2}\|dA\|_2^2
\end{equation}
for every mean-zero Coulomb coarse field \(A\). The sum of its
non-winding terms has the same bound for arbitrary \(A\),
without a gauge condition. Harmonic fields are excluded from
the conversion of the winding tail to curl energy.
\end{theorem}
\begin{proof}
A carrier in Theorem \ref{thm:v40-decomposition} fits, when it
does not wrap, in a box of side
\(L_R=2(R+s_0)+1\). For \(L_R<N/4\), apply the germ version
of Theorem \ref{thm:v39-local} with
\(h\le K\beta^{-2}e^{-a_1R}\). A fixed plaquette belongs
to at most \(C(L_R+2)^4\) such boxes at this radius; the number
of group types per cell is fixed. Its total coefficient is at most
\[
      K\beta^{-2}\sum_{R\ge0}
                   (L_R+2)^4L_R^5e^{-a_1R}<\infty.
\]
This proves the non-winding conclusion. It is precisely an
orbit-weighted interaction estimate: for any \(0<b<a_1/2\),
the V39 radius weight \(e^{bL_R}\) is summable against
\eqref{eq:v40-orbit-norm}, with stock \(O(\beta^{-2})\).

For the remaining terms \(L_R\ge N/4\), use their Hessian norms
directly. Summation over owners and radii gives
\[
       |D^2R_{\rm long}(I)[A,A]|
           \le K\beta^{-2}N^4 e^{-a_2N}\|A\|_2^2
\]
for some \(a_2>0\), adjusting \(K\) to cover the finitely
many small side lengths compared with \(s_0\). The mean-zero
Coulomb inequality \(\|A\|_2^2\le N^2\|dA\|_2^2/16\)
then gives a relative factor \(KN^6e^{-a_2N}\), uniformly
bounded. The finite telescope identifies the sum of these
Hessians with the Hessian of the original radial remainder.
No positivity of that remainder, reflection symmetry of the
ordered tree, or analytic continuation of a finite-grid Ward
identity has been assumed.
\end{proof}

Consequently, if the \emph{retained truncated packet}
\(\beta G+Q+\beta^{-1}\mathcal B\) has an independently
verified mean-zero Coulomb Hessian lower bound
\(\kappa_\beta\|dA\|_2^2\), then adding this exact remainder
reduces its lower coefficient by at most \(K\beta^{-2}\).
This last implication does not assert that the truncated packet
already has its required uniform lower bound, or that its
relevant and marginal coefficients can be discarded.

\subsection{Checks, preservation, and what remains upstream}

The exact checker is
\begin{center}\small
\path{scripts/verify_ym_f14_v40_covariant_activation.py}.
\end{center}
It has 6,712 bytes and SHA--256
\begin{center}\scriptsize\ttfamily
3299ABEB44B7EB638ED7564F5AA58F281538D6EAB78D87C7564A3892D54BD356.
\end{center}
On a side-four fine torus it checks all 1,024 endpoint frame
identities, the 240 tree, 64 pivot, and 720 independent-group
counts, and all 1,536 Wilson plaquette conjugacies using rational
unit quaternions. It checks the fourteen relative-source
conjugacies and an exact adjoint isometry. The pivots in this
frame test are arbitrary endpoint-covariant links; the checker
does not claim to solve generic nonlinear compact pivot equations.
Their solution covariance follows from uniqueness as proved above.

An independent stable two-site quartic calculation checks the
activation normalization and connected subtraction. Its precision is
\[
 H_s=\begin{pmatrix}1+2s_0^2&-s_0s_1\\
                       -s_0s_1&1+s_1^2\end{pmatrix}
       =I+D_s\left(\begin{pmatrix}3&-1\\-1&2\end{pmatrix}-I\right)D_s
                              \succeq I.
\]
For the quartic terms
\((1+u)s_0^4\eta_0^4+(1+v)s_1^4\eta_1^4\), the coefficient
of \(\beta^{-2}\) at full activation is
\[
 -\frac{768(1+u)^2+3888(1+v)^2+456(1+u)(1+v)}{625}.
\]
The mixed Hessian is \(-456/625\), not zero. Gaussian eighth
moments are independently checked by all 105 pairings. In
particular, the radius-one correction is not merely the mixed
monomial: it also contains changes in the one-site variances.
A separate arbitrary six-site set function verifies the full
twofold telescope without probabilistic factorization. These
are finite algebra checks, not numerical certification of the
uniform analytic theorem or of a mass gap.

V39 has 914,775 bytes and SHA--256
\begin{center}\scriptsize\ttfamily
BCE9EA53E80FD280ACCD84F9C05CD8ADE953919D95C60752E83DC41C6975E79C.
\end{center}
Its mathematical body is retained, with only the title version
changed and this terminal section appended. V40 replaces V39
as the single active main file. No archive or companion is
altered. The next compulsory companion review is V45.

\begin{firewall}
The primitive radial local shell now has an explicitly covariant
measured law, an activation family preserving its eliminated-space
floor, an exponentially local invariant remainder decomposition
at the pure-gauge orbit, and an \(O(\beta^{-2})\) relative curl
Hessian bound. These are not estimates on the unrestricted
compact complement or on a uniformly controlled family of
nonlinear iterates. The next genuine obstruction is to propagate
the full incoming action and density, including \(Q\) and
\(\mathcal B\), in compatible covariant norms and to control
entry into this branch. A nontrivial interacting continuum state
and its positive physical Hamiltonian mass gap remain unproved.
\end{firewall}
\section{V41: the complete primitive local Hessian and a second-shell convexity gate without a Wilson reset}
\label{sec:v41}

\subsection{What must be added to V40}

V40 bounds the exact measured remainder for a specified covariant
radial contribution. It leaves the lower bound for the retained
packet conditional. This section supplies the missing primitive
packet estimates and then uses the \emph{whole} resulting local
action at one further elimination. In particular, its one-loop
term is not replaced by a fitted coupling, its two-loop coefficient
is not discarded, and its remainder is not set to zero.

The result is still local in field space and finite in the number
of elimination steps. The new quantitative conclusions are:
the classical Hessian is between one and four times the coarse
curl form; both complete quantum coefficients have invariant local
Hessian decompositions; the exact first local action has a
volume-uniform positive curl comparison at sufficiently large
\(\beta\); and this actual generated action gives a uniformly
convex second detail domain. No all-scale norm induction or
compact-complement estimate is supplied by these conclusions.
The next companion checkpoint remains V45.

\subsection{Exact energy normalization of the balanced linear block}

Use the normalization
\(S(U)=\sum_p(1-\tfrac12\Tr U_p)\), with
\([E_a,E_b]=2\epsilon_{abc}E_c\). Thus
\(D^2S(I)[a,a]=\|d_fa\|_2^2\). All edge and plaquette sums
below use counting measure. Let the fine side be \(2N\), the
coarse side \(N\), and let \(L\) be the derivative of the actual
balanced rooted block at \(I\).

For a fine field \(a\), let \(t_{r,\delta}(a)\) be its
circulation along the fixed cell tree from \(2r\) to
\(2r+\delta\), where \(\delta\in\{0,1\}^4\). Put
\(\bar t_r=16^{-1}\sum_\delta t_{r,\delta}\). Differentiating
the sixteen rooted chronological paths, including the two
duplicates, gives the exact formula
\begin{equation}\label{eq:v41-linear-root}
 (La)_{r,\mu}=\frac1{16}\sum_{\delta\in\{0,1\}^4}
       \bigl(a_{2r+\delta,\mu}+a_{2r+\delta+e_\mu,\mu}\bigr)
                         +\bar t_r-\bar t_{r+e_\mu}.
\end{equation}
The last two terms are a coarse gradient. They cannot be omitted
from the root itself, but cancel from its curl and from its sum
over the periodic lattice.

\begin{lemma}[The curl averaging and the boundary-copy lift]
\label{lem:v41-curl}
For \(\mu<\nu\),
\begin{equation}\label{eq:v41-curl-average}
 (d_cLa)_{r,\mu\nu}
   =\frac1{16}\sum_{\delta\in\{0,1\}^4}
                  \sum_{u,v=0}^1
       (d_fa)_{2r+\delta+ue_\mu+ve_\nu,\mu\nu}.
\end{equation}
Consequently
\begin{equation}\label{eq:v41-curl-contraction}
                         \|d_cLa\|_2^2\le\|d_fa\|_2^2.
\end{equation}
Define \(\mathsf R A\) to be zero on internal fine links and
equal to \(A_{r,\mu}\) on every \(\mu\)-boundary link out of
cell \(r\). Then
\begin{equation}\label{eq:v41-lift}
              L\mathsf R=I,\qquad
              \|d_f\mathsf R A\|_2^2=4\|d_cA\|_2^2.
\end{equation}
Also, for each orientation,
\begin{equation}\label{eq:v41-mean}
                 \sum_r(La)_{r,\mu}
                            =\frac18\sum_xa_{x,\mu}.
\end{equation}
In particular, every field in \(\ker L\) has zero fine mean.
\end{lemma}
\begin{proof}
Insert \eqref{eq:v41-linear-root} into the four terms of a coarse
plaquette. The tree terms cancel. Each remaining length-two edge
circulation bounds the two-by-two fine plaquette surface, which
proves \eqref{eq:v41-curl-average} by telescoping its internal
edges. The nonnegative weights in one coarse plaquette row have
total mass four. Their column sum, after summing all coarse
plaquettes of the same orientation, is \(1/4\): for any fine
plaquette there are four choices of \((u,v)\), and for each one
the parity offset and coarse anchor are uniquely determined.
Weighted Cauchy--Schwarz therefore gives
\(\sum_r|\sum_pw_{rp}F_p|^2\le4\sum_{r,p}w_{rp}|F_p|^2
=\sum_p|F_p|^2\). This proves the contraction, separately
for every orientation and color.

Every rooted two-link path of \(\mathsf R A\) crosses precisely
one boundary of its prescribed coarse edge, and all its tree
links vanish. Hence its average is \(A\). A fine
\((\mu,\nu)\)-plaquette has nonzero curl under this lift only
if both corresponding parity bits are one. In that case it is
the coarse plaquette curl. The other two parity bits are free,
so there are four copies of each coarse plaquette. This proves
the energy identity. Finally, summing the averaged two-link
formula counts each fine edge twice with weight \(1/16\);
the tree gradient sums to zero. This proves \eqref{eq:v41-mean}.
\end{proof}

\begin{theorem}[Classical primitive Schur form with explicit comparison constants]
\label{thm:v41-classical}
Let \(G\) be the classical minimum in the actual primitive
radial shell of V40. Its Hessian satisfies, for every coarse
field \(A\),
\begin{equation}\label{eq:v41-classical}
 \|d_cA\|_2^2\le D^2G(I)[A,A]\le4\|d_cA\|_2^2.
\end{equation}
No mean-zero or gauge condition is required in this statement.
\end{theorem}
\begin{proof}
The full fine action has zero first derivative at \(I\). Thus
the second variation of its constrained minimum is the quadratic
minimum
\[
              D^2G(I)[A,A]=\inf_{La=A}\|d_fa\|_2^2.
\]
This is the finite constrained Schur identity, not a replacement
of the nonlinear root away from zero. Fine gauge transformations
equal to the identity at the cell roots remove the fifteen tree
directions without changing either \(La\) or the curl energy.
The remaining forty-five detail coordinates have the inherited
strict Hessian floor, so their minimum is exactly the one
obtained by differentiating the local stationary branch. The
lower bound follows from \eqref{eq:v41-curl-contraction}; the
admissible lift \(\mathsf R A\) gives the upper bound.
\end{proof}

Two related exact facts will be useful for the coefficient
decomposition. The group-valued boundary-copy configuration of
V40 has fine action \(f(C,0)=4S_c(C)\): the same four-copy
plaquette argument now uses matrix products, and every other
fine plaquette equals \(I\). Its density \(\ell(C,0)\) is a
constant independent of \(C\). Indeed all fourteen nonroot
products equal \(C_e\). At fixed such products, the completed
pivot map near \(B=C_e\) is exactly
\[
 C_e e^u\exp(-14u/16)=C_e e^{u/8}.
\]
Its inverse Haar Jacobian at \(u=0\) is \(8^3\), independently
of \(C_e\); each independent fine Haar factor is \(J_G(0)\).
This proves constancy with the actual measure, not by deleting
a source-dependent Jacobian.

\subsection{Four-mark locality of the complete quantum coefficients}

Retain V40's activation \(s\), minimum \(\eta_{*,s}\), and
complete density. Write \(\Sigma_s=H_s^{-1}\). The available
uniform inverse and marked-minimum estimates imply, after
reducing a positive rate at each of the finitely many orders,
\begin{equation}\label{eq:v41-marked-kernels}
 \begin{aligned}
 |\partial_I(H_s)_{ab}|&\le C_{|I|}
          1_{d(a,b)\le r_0}\,e^{-\nu_{|I|}\tau(a,I)},\\
 |\partial_I(\Sigma_s)_{ab}|&\le C_{|I|}
                            e^{-\nu_{|I|}\tau(a,b,I)},
                                                    \qquad |I|\le4.
 \end{aligned}
\end{equation}
Marks in this section are again either coarse coordinates or
activation parameters. Empty marked sets have their usual
unmarked interpretation. Component multiplicities are fixed.

\begin{lemma}[The measured coefficient stocks, uniformly over activation]
\label{lem:v41-coefficients}
For one \(a>0\) and each \(1\le k\le4\), the anchored norms
of \(G_s,Q_s,\mathcal B_s\) in the union of these marked
parameters are bounded by constants independent of volume and
activation. In particular,
\begin{equation}\label{eq:v41-coefficient-norms}
              \|Q_s\|_{a,k}+\|\mathcal B_s\|_{a,k}\le C_k.
\end{equation}
These are bounds on the full measured coefficients of
\eqref{eq:v40-expansion}, not their precision-only parts.
\end{lemma}
\begin{proof}
The first line of \eqref{eq:v41-marked-kernels} follows by the
local chain rule at the marked minimum. Differentiating
\(H_s\Sigma_s=I\) repeatedly expresses every marked inverse
as a finite sum of products of unmarked inverses and marked
Hessians. Their exponential links join both output indices and
all marked terminals. Splitting the rate between the terminal
tree and the internal sums proves the second line. This also
proves the same marked local-tensor bounds for the action and
density derivatives at the minimum.

For the classical term, the first derivative of the minimized
action is its explicit retained derivative, because
\(\partial_\eta f_s=0\). Subsequent differentiation uses
only the local tensors and the marked minimum responses.
Every term is a connected tree rooted at an external mark.
Summing it with a smaller exponential weight proves the bounds
for \(G_s\). Local explicit terms, such as the derivative of
the Gaussian pinning coefficient, fit the same estimate.

For the determinant, one may use the following exact finite
derivative formula. If \(I\) is a nonempty labelled set and
the sum runs over its ordered partitions into nonempty blocks,
then
\begin{equation}\label{eq:v41-logdet-marks}
 \partial_I\log\det H_s
  =\sum_{(I_1,\ldots,I_m)}\frac{(-1)^{m-1}}m
       \Tr\bigl(\Sigma_s H_{s,I_1}\cdots
                                      \Sigma_s H_{s,I_m}\bigr).
\end{equation}
This is the mixed coefficient of
\(\Tr\log(I+H_s^{-1}\Delta H_s)\); only its finite Taylor
polynomial of order \(|I|\) is used. Each trace is a closed
connected walk with every marked terminal attached to one of
its vertices. Since at least one mark is present, this graph
can be rooted there. Insert \eqref{eq:v41-marked-kernels}, keep
part of its decay for the external tree and sum its internal
vertices. There is no free volume factor. The local-density
term \(-\ell_s(\eta_{*,s},C)\) has the same rooted chain-rule
bound. These two terms prove the estimate for \(Q_s\).

For clarity, the complete inherited Wick expression is
\begin{equation}\label{eq:v41-Wick}
 \begin{aligned}
 \mathcal B_s={}&\tfrac18 V_{abcd}\Sigma_{ab}\Sigma_{cd}
                    -\tfrac12(L_2)_{ab}\Sigma_{ab}\\
 &-\tfrac1{12}T_{abc}T_{def}\Sigma_{ad}\Sigma_{be}\Sigma_{cf}\\
 &-\tfrac12(p-\vartheta/2)^{\mathsf t}\Sigma(p-\vartheta/2),
 \qquad \vartheta_i=T_{abi}\Sigma_{ab},
 \end{aligned}
\end{equation}
where \(T=f_{\eta^3}\), \(V=f_{\eta^4}\),
\(p=\ell_\eta\), and \(L_2=\ell_{\eta^2}\), all at the
activated moving minimum. The first two terms are connected
one-vertex contractions; the last two are connected two-vertex
contractions, including the one-cross-line tadpole and density
terms. Formula \eqref{eq:v41-Wick} is Shell V7's measured Wick
formula, not a newly claimed coefficient. Its purpose here is
to identify every graph to which the locality estimate applies.

Distribute at most four marks among its finite number of tensors
and covariance lines. The differentiated kernel bounds above
keep each contraction connected to every external mark. The
same internal-vertex summation proves the anchored bounds for
\(\mathcal B_s\). This reasoning includes derivatives of
\(\Sigma_s\), the minimum, fine Haar, and the pivot density;
none is held fixed as the retained source varies. Only a fixed
number of derivatives of the primitive analytic stencils is
required, within their already available bounds.
\end{proof}

\begin{proposition}[Invariant local decomposition of both coefficients]
\label{prop:v41-invariant-coefficients}
For \(F_s=Q_s\) and for \(F_s=\mathcal B_s\), there are
local invariant germs \(\Phi^F_{i,R}\), of the same carriers
as in V40, such that
\begin{equation}\label{eq:v41-coefficient-telescope}
 F_{\boldsymbol1}(C)-F_{\boldsymbol1}(I)
     =\sum_{i,R}\Phi^F_{i,R}(C),\qquad
 \sup_g\|D^2\Phi^F_{i,R}(g e^c g^{-1})|_{c=0}\|
                                                    \le C e^{-aR}.
\end{equation}
The endpoint notation in the norm means the linkwise transformation
\(g_xe^{c_e}g_y^{-1}\), as in V40. The norm acts on all
coordinates of the term's original support.
\end{proposition}
\begin{proof}
At zero activation, \(H_0=h_0I\) and the density at the
minimum is the constant \(\ell(C,0)\) computed above. Thus
\(Q_0\) is a vacuum constant. All cubic, quartic, and density
derivatives in inactive integration variables vanish, so
\(\mathcal B_0=0\). Therefore the vacuum-subtracted functions
have zero all-inactive base, just as the remainder in V40.

Use precisely the predecessor ordering and radius increments
\eqref{eq:v40-increments}, with \(F\) in place of \(R\).
The finite double telescope is an exact identity. Locality
of each increment follows from the active-variable factorization;
all terms independent of active variables cancel before taking
stationary coefficients. Gauge invariance follows from the same
equivariant minima and normalized finite integrals.

For a radius increment, the double activation interpolation
introduces two marks. A Hessian in the coarse field introduces
two more, so \eqref{eq:v41-coefficient-norms} applies. The
pointwise tree bound precedes the supremum over its varying
activation patterns. Retain part of its decay for the separation
\(R\), then sum the other terminals with the remaining decay.
This is exactly the supremum-ordering argument in the proof
of Theorem \ref{thm:v40-decomposition}, now with size one
instead of \(\beta^{-2}\). Gauge covariance transports the
Hessian to the whole pure-gauge orbit. No global compact
extension of these stationary germs is inferred.
\end{proof}

\subsection{The first exact local action is quadratically stable}

\begin{theorem}[Complete primitive local curl comparison]
\label{thm:v41-complete}
There are finite \(K_Q,K_B,K_R\), independent of volume and
\(\beta\ge\beta_0\), for which
\begin{equation}\label{eq:v41-quantum-critical}
 \begin{aligned}
 |D^2Q(I)[A,A]|&\le K_Q\|d_cA\|_2^2,\\
 |D^2\mathcal B(I)[A,A]|&\le K_B\|d_cA\|_2^2,\\
 |D^2R_\beta(I)[A,A]|&\le K_R\beta^{-2}\|d_cA\|_2^2
 \end{aligned}
\end{equation}
on mean-zero Coulomb fields. Consequently the \emph{exact}
radial local action obeys
\begin{equation}\label{eq:v41-full-comparison}
 \begin{aligned}
 (\beta-K_Q-K_B/\beta-K_R/\beta^2)\|d_cA\|_2^2
       &\le D^2W_\beta^{\rm rad}(I)[A,A]\\
       &\le(4\beta+K_Q+K_B/\beta+K_R/\beta^2)\|d_cA\|_2^2.
 \end{aligned}
\end{equation}
For a volume-independent \(\beta_*\), this gives the simpler
comparison
\begin{equation}\label{eq:v41-simple-comparison}
 \frac\beta2\|d_cA\|_2^2\le D^2W_\beta^{\rm rad}(I)[A,A]
                                      \le5\beta\|d_cA\|_2^2.
\end{equation}
Both comparisons extend to every mean-zero field, without a
Coulomb condition. Moreover, in the original retained chart,
the anchored field derivatives of
\(\beta^{-1}W_\beta^{\rm rad}\) through order four are bounded
uniformly in volume and \(\beta\ge\beta_*\).
\end{theorem}
\begin{proof}
Apply V39's local curvature theorem to the non-winding germs
in \eqref{eq:v41-coefficient-telescope}. The radius sum is
bounded by \(C\sum_R(R+1)^9e^{-aR}\), including the number
of owners covering a fixed plaquette. The winding tail has
absolute Hessian norm at most \(CN^4e^{-a'N}\), and the
mean-zero Coulomb Poincare conversion costs at most \(N^2/16\).
This proves the first two bounds; the third is V40. Sum these
with Theorem \ref{thm:v41-classical} in the exact measured
expansion. All vacuum terms are constants in \(C\).
For example it suffices to take
\(\beta_*\ge\max(1,\beta_0,2(K_Q+K_B+K_R))\) for
\eqref{eq:v41-simple-comparison}.

The exact action is gauge invariant. At \(I\), its first
derivative vanishes by simultaneous color conjugation. Its
Hessian therefore annihilates every infinitesimal gauge gradient,
with both legs, by differentiating that invariance. Decompose a
mean-zero field as \(A=A_\perp+d_c\phi\), where \(A_\perp\)
is mean-zero Coulomb. The Hessian and curl are unchanged, proving
the extension. This projection is used only in the estimate;
it does not alter the root map or the integrated measure.

Finally combine the coarse-only marked estimates for
\(G,Q,\mathcal B\) with V40's remainder estimates and divide
the exact expansion by \(\beta\). The logarithmic normalization
has no field derivatives. This proves the uniform derivative
stocks through order four on a fixed smaller retained chart.
\end{proof}

This comparison is not a physical mass gap. Its reference is the
massless curl form; its smallest nonzero transverse Fourier
eigenvalue is \(4\sin^2(\pi/N)\). The theorem concerns the
Hessian of a localized Euclidean effective action, not the
spectrum of the physical Hamiltonian or a globally log-concave
compact Yang--Mills measure. Harmonic modes were excluded from
the winding-tail Poincare argument, not erased from the law.

\subsection{A genuine second local elimination without changing the incoming action}

Suppose the first coarse side is \(2M\), with \(M\ge2\).
Let \(\mathcal X_2(v,z)\) be the exact next primitive
parity-tree fibre section, with coarse source \(D=e^v\) and
forty-five independent groups per new cell in V40's background
frame. Thus \(\mathcal X_2(0,0)=I\). Keep the entire first
action and define
\begin{equation}\label{eq:v41-second-action}
 F_\beta(v,z)=\frac1\beta\left\{
       W_\beta^{\rm rad}(\mathcal X_2(v,z))-W_\beta^{\rm rad}(I)
                                      \right\}.
\end{equation}
This action is allowed to depend on \(\beta\). Let \(\ell_2\)
be the actual second fine Haar and inverse-pivot logarithmic
density, again with output measured relative to coarse Haar.
For a sufficiently small fixed product of balls \(\mathcal D_2\),
the two-stage localized contribution is exactly
\begin{equation}\label{eq:v41-second-integral}
 Z_\beta^{(2),\rm loc}(D)=e^{-W_\beta^{\rm rad}(I)}
       \int_{\mathcal D_2}e^{-\beta F_\beta(v,z)+\ell_2(v,z)}\,dz.
\end{equation}
Equivalently the integrand is
\(Z_\beta^{\rm rad}(\mathcal X_2(v,z))e^{\ell_2(v,z)}\).
This identity retains every coefficient and the first exact
remainder. It is a positive two-stage contribution obtained
by imposing the two specified local filters in the compact
pushforward, not the unfiltered two-stage partition function.

\begin{theorem}[Second-shell local convexity with the actual incoming packet]
\label{thm:v41-second}
For all sufficiently large \(\beta\), there are retained and
detail radii independent of \(M\) and \(\beta\) on which
\(F_\beta\) has a uniform positive detail Hessian floor.
Its weighted derivative norms through order four are uniform.
After a common further decrease of the retained radius it has
a unique nearby detail minimum in the inner half-domain, and
the measured potential \(\beta F_\beta-\ell_2\) is uniformly
strongly convex there. The inverse detail Hessian has uniform
exponential spatial decay. The minimum's retained responses
through order three have exponential tree bounds.

This is a second-shell convexity and regularity gate, not a
proved second-shell two-loop expansion or an invariant norm
ball for arbitrarily many eliminations.
\end{theorem}
\begin{proof}
Let \(E=D_z\mathcal X_2(0,0)\) in Lie coordinates. This is
the actual forty-five-coordinate detail embedding for the next
balanced root, so \(L_2E=0\). By \eqref{eq:v41-mean},
\(E\zeta\) has zero mean. The inherited geometrical detail
estimate is
\[
                     \|d_cE\zeta\|_2^2\ge h_{\rm fib}\|\zeta\|_2^2
\]
for one \(h_{\rm fib}>0\) independent of volume. This is the
fixed block-map fibre inequality, not an assertion that the
incoming action is Wilson.

The first derivative of \(W_\beta^{\rm rad}\) vanishes at
\(I\). Consequently the second derivative of
\eqref{eq:v41-second-action} at zero has no constraint-acceleration
term, and Theorem \ref{thm:v41-complete} gives
\begin{equation}\label{eq:v41-second-floor}
 D_z^2F_\beta(0,0)[\zeta,\zeta]
   =\beta^{-1}D^2W_\beta^{\rm rad}(I)[E\zeta,E\zeta]
                         \ge\tfrac12h_{\rm fib}\|\zeta\|_2^2.
\end{equation}
We used its extension to all mean-zero fields: \(E\zeta\)
need not itself be Coulomb.

The next fibre section has fixed analytic stencils and bounded
incidence. Choose the second radii so that its image lies in
the first action's fixed retained chart; the uniform local
coordinate bounds make this choice independent of volume.
Composing the first action's weighted derivative
bounds with that section, by the finite chain rule through
order four, preserves a smaller positive exponential weight
and uniform constants. For example, summing the third derivative
with its first index fixed bounds the row sum of a Hessian
increment by
\[
 \max_i\sum_j| (F_\beta)_{z_iz_j}(v,z)
                         -(F_\beta)_{z_iz_j}(0,0)|
                              \le K(|v|_\infty+|z|_\infty).
\]
The same holds for columns and with a smaller spatial weight.
Schur's row/column bound controls the operator norm. Thus on
a common small domain the lower floor is at least
\(h_{\rm fib}/4\). Away from zero the full chain rule includes
\(DW_\beta^{\rm rad}\,D_z^2\mathcal X_2\); this term is
part of the just-estimated increment and is not dropped.
The second density has uniformly bounded Hessian, so for large
\(\beta\) its subtraction leaves, for example, a measured
floor \(\beta h_{\rm fib}/8\).

The detail Hessian is generally exponentially local, not
finite-range. This is enough for its inverse estimate. Conjugate
component \(i\) by \(e^{\sigma d(i,j_0)}\). The weighted row
stock bounds the change in the Hessian by \(C\sigma\) in
operator norm. Choose a fixed \(\sigma>0\) so that this is
less than half the proved floor. The resolvent identity gives
a uniform conjugated inverse norm and hence exponential decay
of its entries, with summable rows.

At \((v,z)=(0,0)\) the detail gradient is zero. Its value at
\(z=0\) and small \(v\) is \(O(|v|_\infty)\), uniformly,
by the mixed Hessian row bound. Use the inverse Hessian at
zero in the maximum-norm critical equation. The nonlinear
gradient remainder has Lipschitz norm bounded by
\(K(|v|_\infty+|z|_\infty)\). On the chosen small detail
ball it is a contraction; decreasing the retained radius
puts its fixed point in the inner half-domain. Strong convexity
gives uniqueness. The measured minimum is obtained with the
additional bounded gradient perturbation \(-\ell_2/\beta\).
Finally, differentiate the critical equation up to three times.
Each response is an inverse Hessian applied to local or
exponentially local derivatives and lower responses. Their
connected exponential sums prove the asserted tree bounds.
The available fourth derivatives are sufficient for these
three response orders.
\end{proof}

The theorem pays for one genuinely nonprimitive incoming action.
It does not establish all of Shell V7's next-step hypotheses:
the new action is not a fixed finite-range Wilson stencil sum,
the required higher derivative and complete loop-filtration
norms have not yet been propagated here, and the constants
have not been shown to stay controlled with the number of
steps. Its \(\beta\)-dependence must be retained in any next
stationary expansion, as in Shell V8's incoming-packet rules.
This is the next substantive target, together with compact
entry, rather than repeating the primitive remainder theorem.

\subsection{Executed finite checks and append-only record}

The checker is
\begin{center}\small
\path{scripts/verify_ym_f14_v41_complete_local_shell.py}.
\end{center}
It has 6,642 bytes and SHA--256
\begin{center}\scriptsize\ttfamily
4B6A7F26CA1E9B0E0F1917EB1995A7A940F010A688F3A34F9A0D29F4608BD7F6.
\end{center}
On coarse sides two and three it builds the literal linear
tree-dressed root rows and checks all 96 and 486 coarse curl
identities against the independently assembled fine plaquette
surfaces. Their exact row masses are four and column masses
\(1/4\). It checks all 64 and 324 right-inverse rows, every
one of the 1,536 and 7,776 fine plaquette rows under the
boundary-copy lift, and the four global mean identities.
These are exact rational identities, not sampled field tests.

A separate noncommuting two-by-two matrix test checks the
trace-log and inverse identities through total marked degree
four, against the scalar determinant and direct multiplication.
For a finite two-variable measured polynomial test it compares
all fifteen marked coefficients of the normalized Gaussian
formula with the complete tadpole-square form. Its mixed
degree-four coefficient is \(-436831/450000\). That number is
a regression value of the declared finite test, not a Yang--Mills
running coefficient. Neither this test nor the geometric
checks numerically certify the uniform analytic estimates;
those use the proofs and their stated inherited inputs above.

V40 has 943,556 bytes and SHA--256
\begin{center}\scriptsize\ttfamily
00C74C503038D5581B2C930349279D3E064CADCA5433088B9120EE3D5FDB9AD8.
\end{center}
Its preceding body is preserved apart from the title version
and this terminal insertion. V41 replaces V40 as the single
active main note; archives and companions remain untouched.

\begin{firewall}
The complete first radial local action, including its actual
one-loop, two-loop, and remainder terms, now has a uniform
quadratic curl comparison. Its unchanged generated action
also gives a common second detail domain with a proved
convexity floor and four-derivative control. No replacement
by a primitive Wilson action was made. An all-scale controlled
incoming packet, domination or repair of the omitted compact
sectors, an interacting continuum state, and a positive
physical Hamiltonian mass gap remain unproved.
\end{firewall}
\section{V42: a uniform two-level measured expansion and an obstruction to flattening all levels into one unweighted fibre}
\label{sec:v42}

\subsection{Keep the joint variables before estimating the generated action}

V41 admits a second local elimination with the actual first
action, but does not yet provide all its higher quasilocal
derivatives. There is a direct way to prove the two-level
expansion without assuming those missing estimates: retain
both layers of integration variables. At two fixed levels
the original microscopic action and the full density then
remain fixed finite-range stencils on a larger cell carrier.

This section proves a volume-uniform measured expansion through
two loops for that exact joint local integral, and identifies
its coefficients with the incoming-packet formula of Shell V8.
The coefficient composition identity itself is inherited; the
new conclusion is its uniform remainder for these explicit
compatible radial domains. We then prove why this fixed-depth
construction cannot be promoted to an all-depth uniform
convexity theorem in the same unweighted coordinates. The
counterexample uses the literal parity-tree detail variables,
not an unrelated massless scalar comparison.

\subsection{One joint density for two actual levels}

Let the final lattice have side \(M\ge2\), the intermediate
lattice side \(2M\), and the original lattice side \(4M\).
Write \(X_1(C,\eta)\) for the first completed fibre section,
and \(X_2(D,\zeta)\) for the second, in V40's background frames.
Both maps retain the original balanced group-valued source.
The joint section and its measured data are
\begin{equation}\label{eq:v42-joint-data}
 \begin{aligned}
 X_{12}(D,\eta,\zeta)&=X_1(X_2(D,\zeta),\eta),\\
 f_{12}(D,\eta,\zeta)&=S_0(X_{12}(D,\eta,\zeta)),\\
 \ell_{12}(D,\eta,\zeta)
       &=\ell_1(X_2(D,\zeta),\eta)+\ell_2(D,\zeta).
 \end{aligned}
\end{equation}
Each \(\ell_j\) contains precisely its stage's independent
fine Haar factors and inverse-pivot Jacobian. Since the first
pushforward is relative to intermediate Haar, that same
measure is the second stage's input Haar; it is not inserted
a second time into \(\ell_1\). The normalized gauge volume
removed by each finite tree is one. Gauge invariance of the
first radial contribution justifies the intermediate tree
retraction, with its first fibre variables transported as well.

Fix product-of-ball radii \(r_1,r_2\) and set
\begin{equation}\label{eq:v42-joint-integral}
 Z^{12}_\beta(D)=\int_{\mathcal D_{r_2}}\int_{\mathcal D_{r_1}}
             e^{-\beta f_{12}(D,\eta,\zeta)
                              +\ell_{12}(D,\eta,\zeta)}\,d\eta\,d\zeta.
\end{equation}
Fubini gives exactly
\[
 Z^{12}_\beta(D)=\int_{\mathcal D_{r_2}}
           Z^{1,\rm rad}_\beta(X_2(D,\zeta))e^{\ell_2(D,\zeta)}\,d\zeta.
\]
Thus no first-stage remainder has been deleted. The two
integration domains are fixed in their displayed coordinates;
the source dependence of both background frames is retained
in \eqref{eq:v42-joint-data}.

We choose a common compatible member of the small-radius
families of V40--V41: first choose \(r_1\) sufficiently small
for the joint Hessian estimate below, then choose \(r_2\)
so that \(X_2(D,\zeta)\) stays in the first admitted retained
chart. Finally decrease the final retained radius. The radii
are fixed independently of \(M\) and \(\beta\). This is a
selection within the previously unspecified sufficiently-small
radius families. Changing an already fixed numerical radius
would change the exact local contribution; no identity between
different such localizations is asserted.

Each final cell has \(16\cdot45+45=765\) independent group
coordinates, or \(2295\) real Lie coordinates. Equivalently,
the original \(4^4\)-site cell has \(1024\) fine links, of
which \(255\) tree gauges and four final source groups are
removed. The fixed finite number of components is large but
does not depend on the total volume. Composition of the two
fixed stencils leaves a fixed finite range on the final cell
lattice, with bounded incidence and derivatives of every
fixed order on a common small chart.

Final coarse gauge transformations lift to cell-constant
transformations at the intermediate level. On both sets of
noise groups they act by the appropriate endpoint adjoints,
which are orthogonal and independent of the integration
variables. Each plaquette and density stencil in
\eqref{eq:v42-joint-data} is therefore invariant, and the
product of radial domains is invariant. This is an exact
gauge identity for the joint local integral.

\subsection{Uniform positivity of the two-level classical fibre}

\begin{lemma}[A controlled joint Schur factorization]
\label{lem:v42-joint-floor}
At \(D=I,\eta=\zeta=0\), write the joint action Hessian as
\[
             \mathbb H=\begin{pmatrix}A&B\\B^{\mathsf t}&D_0\end{pmatrix}.
\]
There are \(h_1,h_2>0\) and \(K<\infty\), independent of
\(M\), such that
\[
 A\succeq h_1I,\qquad
 D_0-B^{\mathsf t}A^{-1}B\succeq h_2I,\qquad
 \|A^{-1}B\|\le K.
\]
In particular
\begin{equation}\label{eq:v42-joint-floor}
             \mathbb H\succeq
        \frac{\min(h_1,h_2)}{2(1+K)^2}I.
\end{equation}
On sufficiently small compatible radial domains, the joint
action and measured potential have uniform positive fibre
Hessian floors, and their minima lie in the inner half-domains
for small final source and large \(\beta\).
\end{lemma}
\begin{proof}
The first block is the original primitive detail Hessian, with
its proved eliminated-space floor. Eliminating its quadratic
variation gives the Hessian of \(G_1\circ X_2\). By V41,
\(D^2G_1(I)\) dominates intermediate curl energy. Restricting
to the second detail embedding gives its geometric fibre
floor \(h_2\). This is a classical Schur complement; it
does not replace \(G_1\) by a Wilson action. Fixed incidence
bounds \(B\) in operator norm, and \(\|A^{-1}\|\le h_1^{-1}\),
which gives \(K\).

For \(w=u+A^{-1}Bv\), the quadratic form is
\[
 u^{\mathsf t}Au+2u^{\mathsf t}Bv+v^{\mathsf t}D_0v
       =w^{\mathsf t}Aw+v^{\mathsf t}
                      (D_0-B^{\mathsf t}A^{-1}B)v.
\]
Also \(|u|^2+|v|^2\le2(1+K)^2(|w|^2+|v|^2)\).
This proves \eqref{eq:v42-joint-floor}. The local third
derivative and bounded incidence make its Hessian perturbation
\(O(r_1+r_2+|\log D|_\infty)\) in row and column norms.
Choose those radii so that this loses less than half the
displayed floor. The density Hessian is uniformly bounded,
so \(\beta f_{12}-\ell_{12}\) has a positive floor of
order \(\beta\) for large \(\beta\).

The unperturbed critical point is zero. The uniform weighted
inverse, maximum-norm implicit equation, and small retained
gradient put the perturbed critical point in the inner
half-domains, exactly as in V40's local-input proof. Different
fixed radii for the two component types only change the
constants through their positive minimum and finite maximum.
Strong convexity gives uniqueness. The same argument with
the bounded density gradient gives the measured minimum.
\end{proof}

The determinant factorization
\(\det\mathbb H=\det A\det(D_0-B^{\mathsf t}A^{-1}B)\)
is only the Gaussian part of the joint measure. Both density
terms in \eqref{eq:v42-joint-data} remain present in every
coefficient below.

\subsection{The two-level remainder is now uniform in volume}

\begin{theorem}[Measured two-loop expansion of the exact joint radial contribution]
\label{thm:v42-expansion}
Let \(n_{12}=2295M^4\). For the compatible fixed domains
above there are \(a>0\), \(\beta_2<\infty\), and finite
constants independent of \(M\) and \(\beta\ge\beta_2\)
such that
\begin{equation}\label{eq:v42-expansion}
 -\log Z^{12}_\beta(D)
   =\beta G_{12}(D)+\frac{n_{12}}2\log(\beta/2\pi)
        +Q_{12}(D)+\beta^{-1}\mathcal B_{12}(D)+R_{12,\beta}(D),
\end{equation}
with
\begin{equation}\label{eq:v42-remainder}
 \|R_{12,\beta}\|_{a,k}\le C_k\beta^{-2},\quad 1\le k\le4,
 \qquad n_{12}^{-1}\sup_D|R_{12,\beta}(D)|\le C_0\beta^{-2}.
\end{equation}
The classical and both complete quantum coefficients have
uniform anchored norms through four final-source derivatives.
All are in the same final coarse-Haar convention.
\end{theorem}
\begin{proof}
The local-input proof of the radial theorem in V40 uses a
fixed number of components per cell, finite-range analytic
stencils with bounded incidence, a uniform positive action
Hessian, a controlled interior minimum, and the true density.
It does not use the number forty-five in any identity specific
to that proof. Here those inputs have all been verified on
the 765-group joint carrier by \eqref{eq:v42-joint-data} and
Lemma \ref{lem:v42-joint-floor}.

More explicitly, activate each joint group by
\(f_{12,s}(D,z)=f_{12}(D,D_sz)
 +(h_0/2)\sum_i(1-s_i^2)|z_i|^2\), and set
\(\ell_{12,s}(D,z)=\ell_{12}(D,D_sz)\), where
\(z=(\eta,\zeta)\). Its Hessian is
\(D_s(f_{12})_{zz}D_s+h_0(I-D_s^2)\); hence the activation
family retains a uniform floor. Coarse and activation marks
are local on the final cell lattice. The inverse and moving
minimum responses have the same exponential tree bounds.
The centered convex extension and signed-noise Taylor formula
then use at most eight connected entries for four source
marks, just as in V40. Each entry has the required fixed
moments and local derivative bounds; increasing the component
number changes constants but not volume powers.

The wall for a group at level \(j\) is
\(\tfrac12(|z_i|-r_j)_+^2\). Its normal Hessian projection,
wall-sensitive resolvent gain, and codimension-one layer
estimate are those proved in V40. With two fixed positive
radii the same integrable strength power \(21/20\) works
uniformly over group types. The hard-wall passage and
actual/extended interpolation thus return to precisely
\eqref{eq:v42-joint-integral}, with exponentially small
anchored errors. Adding these to the fourth-noise Taylor
remainder proves \eqref{eq:v42-remainder}. The zeroth-order
bound is a pressure-per-coordinate estimate, not a positive
volume-uniform probability for the whole local event.

The marked minimum, trace-log, and connected Wick graph
arguments of V41 apply to the same verified joint kernels
and give the coefficient bounds. This proves the theorem
from local data before integrating either level out; it does
not assume higher quasilocal derivatives of the intermediate
logarithm which V41 had not yet estimated.
\end{proof}

For identification with the sequential calculation, put
\[
 f=G_1\circ X_2,\qquad q=Q_1\circ X_2-\ell_2,
 \qquad b=\mathcal B_1\circ X_2.
\]
Let \(\zeta_*(D)\) minimize \(f\), let
\(K=f_{\zeta\zeta}(\zeta_*,D)\), and
\(\Sigma=K^{-1}\). Every tensor in the following formula
is evaluated there:
\begin{equation}\label{eq:v42-sequential}
 \begin{aligned}
 G_{12}&=f(\zeta_*,D),\qquad
 Q_{12}=q(\zeta_*,D)+\tfrac12\log\det K,\\
 \mathcal B_{12}&=b(\zeta_*,D)
  +\mathbb E_\Sigma\left[\tfrac1{24}f_{\zeta^4}[\xi^4]
                              +\tfrac12q_{\zeta^2}[\xi^2]\right]\\
 &\hspace{1em}-\tfrac12\mathbb E_\Sigma
       \left[\tfrac16f_{\zeta^3}[\xi^3]+q_\zeta[\xi]\right]^2.
 \end{aligned}
\end{equation}
These are the inherited measured composition formulas of
Shell V8. Their application here is justified as follows.
The exact domains have the Fubini representation above, and
both minima lie in their common interior. At each fixed
volume, ordinary local Laplace expansion in either order
therefore identifies the same coefficient germs by uniqueness.
The joint theorem supplies the uniform remainder; the
fixed-volume coefficient identity is not being used as its
substitute. The field-independent Gaussian dimensions add:
\(135(2M)^4+135M^4=2295M^4\). In particular, incoming
\(Q_1\) contributes its gradient and Hessian to the second
two-loop term, and incoming \(\mathcal B_1\) contributes its
value at the outer classical minimum. Neither can be reset
to zero.

\subsection{Curvature control of the complete two-level local action}

\begin{proposition}[Two-level critical Hessian bounds]
\label{prop:v42-critical}
For mean-zero final fields, there are constants independent
of \(M\) such that
\[
 |D^2Q_{12}(I)[A,A]|+|D^2\mathcal B_{12}(I)[A,A]|
                                     \le C\|dA\|_2^2,
 \qquad |D^2R_{12,\beta}(I)[A,A]|
                                     \le C\beta^{-2}\|dA\|_2^2.
\]
The classical term satisfies, for arbitrary final fields,
\begin{equation}\label{eq:v42-classical-comparison}
          \|dA\|_2^2\le D^2G_{12}(I)[A,A]\le16\|dA\|_2^2.
\end{equation}
Consequently, for sufficiently large \(\beta\), uniformly
in \(M\),
\begin{equation}\label{eq:v42-complete-comparison}
 \frac\beta2\|dA\|_2^2\le
             D^2[-\log Z^{12}_\beta](I)[A,A]
                                      \le17\beta\|dA\|_2^2
\end{equation}
on mean-zero fields. These are Hessian bounds at the
pure-gauge orbit, not nonlinear bounds throughout the compact
configuration space.
\end{proposition}
\begin{proof}
At zero joint activation the twice-copied boundary field has
action \(16S_c(D)\), and each pivot factor is evaluated with
all relative products equal to \(I\). The density is constant
in \(D\), the inactive Hessian is \(h_0I\), and the two-loop
coefficient is zero. The vacuum-subtracted \(Q,\mathcal B,R\)
thus have zero inactive bases. Use the exact predecessor and
radius telescope of V40 on the joint groups. Two activation
marks and two field marks give exponentially small Hessian
norms of its radius increments. Each increment is a local
nonlinear gauge-invariant germ, because the joint stencils
and radial domains are invariant. V39's tree-filling Hessian
bound handles the non-winding terms. The exponentially small
winding tail is converted on mean-zero Coulomb fields by
Poincare, and gauge Ward symmetry extends the result to all
mean-zero fields, exactly as in V41.

For the classical term, the linear constraint is \(L_2L_1a=A\).
Applying V41's curl contraction twice gives the lower bound
for its quadratic minimum. The twice-applied boundary-copy
right inverse has energy \(4^2\|dA\|_2^2\), giving the upper
bound. Sum these estimates in \eqref{eq:v42-expansion} and
increase the fixed large-\(\beta\) threshold to obtain
\eqref{eq:v42-complete-comparison}.
\end{proof}

This proves a complete two-level local packet, rather than
only a positive second Hessian. It does not yet construct
the continuum law. In particular, the two-level constants
and compatible radii have not been shown uniform as the
number of levels increases.

\subsection{What the same argument gives at any fixed depth}

For a fixed integer \(m\), the joint number of independent
groups per final cell is
\begin{equation}\label{eq:v42-depth-count}
              45\sum_{j=0}^{m-1}16^j=3(16^m-1).
\end{equation}
The exact iterated radial contribution has a joint density
formed by the original fine action and the sum of all
stage densities. At that fixed depth, these are still
finite-range stencils with finite-order bounds independent
of the number of final cells.

\begin{proposition}[Fixed-depth, not depth-uniform, extension]
\label{prop:v42-fixed-depth}
For every fixed \(m\), sufficiently small compatible radii
and a sufficiently large \(\beta_m\) give the same type of
volume-uniform four-mark measured expansion and critical
coefficient bounds. Its classical Hessian lies between
\(\|dA\|_2^2\) and \(4^m\|dA\|_2^2\).
The radii, constants, and threshold in this statement may
depend on \(m\).
\end{proposition}
\begin{proof}
Induct on \(m\). The previous joint Hessian is the first
Schur block. The new complement is the derivative of the
previous classical minimum restricted to the next detail
space. Repeated curl contraction bounds that previous
classical form below by the current curl form; the new
detail floor is therefore strictly positive. At each fixed
depth, the mixed derivative norm is finite by fixed
incidence. Lemma \ref{lem:v42-joint-floor} gives a positive
joint floor, possibly smaller than the previous one.
Choose the compatible radii accordingly and apply the
verified local-input theorem. The classical upper bound
comes from \(m\) boundary-copy lifts. This induction proves
the stated finite-depth assertion, but contains no bound
on its constants as \(m\to\infty\).
\end{proof}

The following test shows that this missing qualification is
essential for the unweighted joint coordinates actually used
above. It is not merely an artifact of a loose Schur estimate.
Existence at each fixed depth also does not select a common
nested cutoff family: decreasing an earlier radius when a
new depth is considered changes the filtered contribution.
Consistency of such choices requires its own comparison;
no projective family of measures follows from this proposition.

\subsection{An actual raw-history field whose Rayleigh quotient tends to zero}

Let \(\mathbb H_m\) be the Hessian of the \emph{classical}
joint action at zero on the \(m\)-level fibre with final
source fixed to \(I\). Its coordinate norm is the unweighted
sum \(\sum_{j=1}^m\|\eta^{(j)}\|_2^2\) of the independent
Lie coordinates in the boundary-copy frames of V40. No
Gaussian harmonic re-centering or level weight is applied.

\begin{lemma}[Exact detail cost of a one-direction profile]
\label{lem:v42-profile-cell}
Consider a scalar-color fine field
\(a_{x,4}=a(x_1)E_3\), with all other orientations zero.
For one cell put \(a_k=a(2r_1+k)\), \(0\le k\le3\), and
\(t=(a_0+a_1-a_2-a_3)/4\). The sum of squares of its
forty-five independent detail coordinates is exactly
\begin{equation}\label{eq:v42-profile-cost}
          11(a_0-a_1)^2+3t^2+4(a_1-a_2-t)^2.
\end{equation}
In particular, on an affine four-point profile with slope
\(s\), this is \(14s^2\). The coarse field is gauge equivalent
to the same one-direction form with new profile
\(a_+(r_1)=a(2r_1)+a(2r_1+1)\).
\end{lemma}
\begin{proof}
The tree circulation at parity \(\delta\) is
\(\delta_4a(2r_1+\delta_1)\). After tree retraction, the
orientation-four boundary link is \(2a(2r_1+\delta_1)\),
and its coarse value is \(a_0+a_1\). The seven nonpivot
boundary links of that orientation give seven copies of
\((a_0-a_1)^2\). The four internal orientation-one links
with \(\delta_1=0,\delta_4=1\) give four further copies.
The coarse orientation-one value is \(t\). Its three
nonpivot boundary links with \(\delta_4=0\) contribute
\(3t^2\); its four boundary links with \(\delta_4=1\)
contribute \(4(a_1-a_2-t)^2\). All other free components
are zero. The four pivot coordinates are not counted as
independent details. This proves \eqref{eq:v42-profile-cost}.
For \(a_k=b+ks\), one has \(t=-s\), so the result is
\(11s^2+3s^2\).

The extra coarse orientation-one component is the gradient
of \((a_0+a_1)/4\). It can be removed by a periodic gauge
transformation without changing orientation four. At the
linear level the independent detail coordinates are
unchanged by adding an input gauge gradient: tree retraction
leaves only the root endpoint gauge action, and subtracting
the boundary-copy coarse frame cancels it. Equivalently,
the nonlinear detail transforms by an endpoint adjoint,
whose derivative at zero detail is the identity on this
variation. Thus the same calculation may be applied at
every subsequent level to the gauge-equivalent profile.
\end{proof}

\begin{theorem}[No depth-independent floor in the unweighted raw history]
\label{thm:v42-history-obstruction}
For every \(m\ge4\), on a fine torus of side \(2^{m+1}\)
with final side two, the actual joint fibre Hessian satisfies
\begin{equation}\label{eq:v42-history-obstruction}
                  \lambda_{\min}(\mathbb H_m)
                                  \le\frac{16}{7(m-3)}.
\end{equation}
Consequently these unweighted joint Hessians cannot possess
a common positive lower bound independent of depth.
\end{theorem}
\begin{proof}
Put \(p=2^m\), \(N=2p\), and on the fine lattice choose
the one-direction profile
\begin{equation}\label{eq:v42-triangle}
 a(k)=\min(k\bmod p,\ p-(k\bmod p))-p/4,
              \qquad a_{x,4}=a(x_1)E_3.
\end{equation}
Its sum over each period is zero, and every fine difference
has absolute value one. Hence
\(\|d_fa\|_2^2=N^4\). After \(j-1\) blocks, ignoring only
the gauge gradient identified in Lemma
\ref{lem:v42-profile-cell}, the input profile is
\[
       a_j(k)=\sum_{u=0}^{q-1}a(qk+u),\qquad q=2^{j-1}.
\]
After all \(m\) blocks its two retained values are both
zero, since each is a full period sum. Its final gradient
component is zero as well. Thus it defines a tangent vector
\(z_m=(\eta^{(1)},\ldots,\eta^{(m)})\) in the actual final
zero-source fibre.

On an affine piece of the original triangle, four consecutive
values of \(a_j\) have slope \(\pm q^2\). For a fixed level
there are four triangle interfaces around the fine torus.
At most two output-cell anchors per interface have a
four-point window crossing that interface. Thus at most
eight values of the output coordinate \(r_1\) are bad.
For \(j\le m-3\), one has \(q\le p/16\), and the number
of output \(r_1\) values is \(N/(2q)=p/q\ge16\). At least
half of them are therefore affine good cells. The other
three cell coordinates are unrestricted.

By \eqref{eq:v42-profile-cost}, every good cell contributes
\(14q^4\) to the detail-coordinate norm. Therefore
\[
 \|\eta^{(j)}\|_2^2
       \ge\frac12\left(\frac{N}{2q}\right)^4 14q^4
       =\frac7{16}N^4,\qquad 1\le j\le m-3.
\]
The other levels have nonnegative costs. The joint tangent
reconstructs a fine field gauge equivalent to \(a\), so its
joint quadratic action is exactly \(N^4\). It follows that
\[
 \frac{z_m^{\mathsf t}\mathbb H_m z_m}{\|z_m\|_2^2}
          \le\frac{N^4}{(7/16)(m-3)N^4}
          =\frac{16}{7(m-3)}.
\]
The variational principle proves
\eqref{eq:v42-history-obstruction}. The triangle's amplitude
grows with \(p\), but this is a tangent calculation: it may
be multiplied by an arbitrarily small scalar to enter any
declared local chart, without changing its Rayleigh quotient.
No large-field configuration has been inserted into a
small-field Hessian theorem.
\end{proof}

The theorem gives an upper bound tending to zero, not a
matching asymptotic law for the smallest eigenvalue. It also
does not forbid a depth-dependent \(\beta\), level weights,
or Gaussian innovation coordinates. The already inherited
uniform one-shell Gaussian estimates of Shell V5 use the
appropriate harmonic organization; they are not contradicted.
What is excluded is using the fixed-depth proof above with
one unchanged positive action floor in this raw all-level
coordinate norm. The affine profile pays a nonzero detail
cost at every early level, while its original curvature
energy is paid only once. This identifies the concrete
conditioning issue that a multiscale change of coordinates
must repair.

\subsection{Exact verification, provenance, and next direction}

The checker is
\begin{center}\small
\path{scripts/verify_ym_f14_v42_joint_shell.py}.
\end{center}
It has 8,179 bytes and SHA--256
\begin{center}\scriptsize\ttfamily
57FA03D44D87242B179788384FD97A28E974345281BCF3B0D67287857DEE92AF.
\end{center}
It composes the actual tree-dressed linear root rows on fine
side eight and final side two. All 96 final curl rows agree
with independently assembled four-by-four fine surfaces.
Their exact row mass is sixteen and column mass \(1/16\).
All 24,576 fine plaquette rows under the twice-copied lift
give the energy factor sixteen. The right-inverse and
global mean identities are checked independently.

A coupled two-variable measured polynomial test compares
direct joint Gaussian contractions with sequential Laplace
coefficients, retaining a nonlinear conditional minimum,
a moving precision, and a nonconstant density. The two
expressions agree exactly as rational functions of a retained
parameter and through its first four derivatives. At zero
that test has incoming \(Q' =11/15\),
\(Q''=299/252\), and complete two-loop coefficient
\(-114673/907200\). Dropping the incoming quantum packet
changes it by \(-18037/226800\). These are declared finite
regression values, not Yang--Mills running coefficients.

An independent enumeration of the forty-five free groups
verifies \eqref{eq:v42-profile-cost} as an exact polynomial
in four arbitrary profile values. The triangular history
is then evaluated without a full high-dimensional lattice
allocation at depths \(4,6,8,10,12\), checking its terminal
zero source and every guaranteed affine-level lower bound.
For example the actual test Rayleigh quotients at depths
four and twelve are \(64/151\) and
\(8388608/77597953\). The all-depth theorem is proved by
the interface count above, not inferred from these samples.

V41 has 969,139 bytes and SHA--256
\begin{center}\scriptsize\ttfamily
53D79839915A68D9A4A1C78BD4D9BFED08AB100B253A127CE415E9E1A6C656C6.
\end{center}
Its body is preserved except for the title version and
this terminal addition. V42 replaces V41 as the only
active main file. Archives and companion notes are preserved;
the next companion checkpoint remains V45.

\begin{firewall}
The exact two-level filtered compact contribution now has
a volume-uniform measured two-loop expansion and critical
Hessian control. At any fixed depth the same type of result
is available with depth-dependent constants. The explicit
triangle proves that uniform raw-history coercivity cannot
be supplied by flattening more levels into the same norm.
The next task is to carry the measured nonlinear packet
in a multiscale coordinate system and norm which avoid
this demonstrated accumulation, while retaining compact
entry as a separate unresolved requirement. No interacting
continuum construction or physical Hamiltonian mass gap
has been proved.
\end{firewall}
\section{V43: exact conditional-centre coordinates, a uniform free history norm, and the surviving Wilson interaction and cutoff boundary}
\label{sec:v43}

\subsection{Why the next step changes coordinates}

V42 proved a measured two-level theorem, but also exhibited an
actual history whose unweighted raw-coordinate Rayleigh quotient
tends to zero with depth. Increasing the number of raw levels in
the same convexity argument cannot repair that obstruction.
The next step is therefore to identify the harmonic innovation
coordinates of the parallel Gaussian programme inside the
\emph{actual nonlinear reconstruction}, including its density
and domains. This is an exact change of integration variables,
not a change of the compact block rule.

The all-level Gaussian Schur bounds and spatial locality below
are inherited from Shell V4--V5 in the current Shell V16 file.
They are not re-proved as a new Gaussian fixed-point theorem.
The archived f6 V48 parameter-dependent Morse construction also
already gives a local quadratic normal form with a nonconstant
Jacobian. We do not duplicate it. The present map uses only
successive conditional translations, has Jacobian exactly one,
and \emph{does not} remove the nonlinear action. Its new roles
here are the identification with V42's measured history, the
quantitative cost of returning to the raw norm, and an explicit
actual Wilson source--detail vertex which survives centering.

All nonlinear statements initially concern a fixed finite depth
and compatible local branches. A pointwise bound uniform in
depth will be proved for the centred Hessian at the trivial
configuration. No common nonlinear neighbourhood is inferred
from that pointwise statement.

\subsection{The triangular map on the actual reconstruction}

Fix a depth \(m\), a finite torus divisible by \(2^m\), and final
coarse data \(C_m\). Write \(C_j\) for the level-\(j\) link
configuration and \(\eta_j\) for the forty-five independent
Lie-algebra coordinates per level-\(j\) cell. As in V40--V42,
the chronological root, solved pivots, parity tree, and
boundary-copy frames are part of the map
\[
                   C_{j-1}=X_j(C_j,\eta_j),\qquad
                            j=m,m-1,\ldots,1.
\]
The coordinate dimensions may differ with \(j\). Lebesgue
measure on all these declared independent real coordinates
is denoted \(d\eta\).

Let \(G_0=S_0\) be the original Wilson action in the
normalization \(D^2S_0(I)[a,a]=\|d_fa\|_2^2\). Define
recursively the \emph{classical} functions
\begin{align}
 f_j(C,\eta)&=G_{j-1}(X_j(C,\eta)),&
 h_j(C)&=\mathop{\rm argmin}_{\eta} f_j(C,\eta),
                                                        \label{eq:v43-fh}\\
 G_j(C)&=f_j(C,h_j(C)).\notag
\end{align}
Each minimum here is the unique small local minimum on the
fixed-depth compatible branch provided by V42, with its
source domain reduced if necessary. It is not the minimum
of a quantum-corrected action. We restrict the history domain
so all displayed evaluations are on those branches. The
fixed-depth local contribution of V42 can be put in this form
by choosing its radii and final source chart accordingly.

\begin{theorem}[Exact triangular conditional centering]
\label{thm:v43-triangular}
For fixed \(C_m\), set
\begin{equation}\label{eq:v43-xi}
                         \xi_j=\eta_j-h_j(C_j).
\end{equation}
This is a smooth bijection between the compatible raw domain
and its image, with a smooth inverse and
\begin{equation}\label{eq:v43-unitjac}
             \det\frac{\partial(\eta_1,\ldots,\eta_m)}
                           {\partial(\xi_1,\ldots,\xi_m)}=1.
\end{equation}
It is gauge covariant as a germ on the corresponding
gauge-saturated domains. Define
\begin{equation}\label{eq:v43-loss}
 \Delta_j(C,\xi)=f_j(C,h_j(C)+\xi)-G_j(C).
\end{equation}
Then the exact action, not merely its Taylor polynomial, is
\begin{equation}\label{eq:v43-telescope}
 S_0(C_0)=G_m(C_m)+
          \sum_{j=1}^m\Delta_j(C_j(\xi_{j+1},\ldots,\xi_m;C_m),
                              \xi_j).
\end{equation}
Every summand has a double zero in its own innovation:
\begin{align}
 \Delta_j(C,0)&=0,& D_\xi\Delta_j(C,0)&=0,\notag\\
 \Delta_j(C,\xi)
  &=\int_0^1(1-t)
       D_\eta^2 f_j(C,h_j(C)+t\xi)[\xi,\xi]\,dt .
                                                        \label{eq:v43-integral-loss}
\end{align}
The integral identity holds wherever that line segment remains
in the branch chart.
\end{theorem}
\begin{proof}
Start with \(\eta_m=\xi_m+h_m(C_m)\), reconstruct
\(C_{m-1}=X_m(C_m,\eta_m)\), and continue downwards.
At step \(j\), the already reconstructed \(C_j\) depends only
on \(\xi_{j+1},\ldots,\xi_m,C_m\); therefore
\(\eta_j=\xi_j+h_j(C_j)\) has derivative the identity in
\(\xi_j\) and zero derivatives in all finer \(\xi_i\), \(i<j\).
This constructs the inverse and proves the block triangular
Jacobian with identity diagonal. No equality of the block
dimensions is required. Keeping \(C_m\) as an additional
unchanged coordinate would also give determinant one.

The reconstruction is covariant, the radial coordinate
domains are preserved by endpoint adjoints, and the classical
action is invariant. Uniqueness of each local minimum
therefore sends \(h_j(C)\) to the same endpoint adjoints as
\(\eta_j\). Their difference transforms in that way too.
This proves covariance on the transported branches, without
asserting that one exponential chart is globally gauge
invariant.

Subtract \(G_j(C_j)\) from
\(G_{j-1}(C_{j-1})=f_j(C_j,\eta_j)\) and sum over \(j\).
The classical values telescope, giving
\eqref{eq:v43-telescope}. The first derivative in \(\xi\)
vanishes by the defining critical equation for \(h_j\).
Taylor's formula with an integral remainder gives
\eqref{eq:v43-integral-loss}.
\end{proof}

Let \(\ell_{\rm joint}(C_m,\eta)\) denote exactly the sum
of the stage logarithmic densities in V42: the independent
group Haar factors and the inverse pivot Haar Jacobians,
relative to final coarse Haar measure. Let
\(\Omega_m(C_m)\) be the image of the raw integration domain.
The exact local partition function is consequently
\begin{equation}\label{eq:v43-measured}
 \begin{aligned}
 Z_m^{\rm loc}(C_m)
   ={}&e^{-\beta G_m(C_m)}
      \int_{\Omega_m(C_m)}
       \exp\left\{-\beta\sum_{j=1}^m
                    \Delta_j(C_j(\xi_{>j};C_m),\xi_j)\right.\\
       &\hspace{42mm}\left.
              +\ell_{\rm joint}(C_m,\eta(\xi;C_m))\right\}
                                                       \,d\xi .
 \end{aligned}
\end{equation}
There is no new log-Jacobian from centering. There is also
no cancellation of the \emph{old} Haar or pivot density:
their values and derivatives must be composed with the
triangular inverse. Since \eqref{eq:v43-measured} changes
variables in the original joint integral before any level
is integrated, it does not reset the incoming \(Q\) or
\(B\) packet. Sequential expansion must still reproduce
the complete measured coefficients of V42.

\begin{proposition}[The centred Hessian at the nested classical branch]
\label{prop:v43-block-hessian}
Let \(C_j^*(C_m)\) be the configurations obtained by setting
all innovations to zero, and put
\[
 H_j(C)=D_\eta^2 f_j(C,h_j(C)).
\]
At fixed \(C_m\), the Hessian in all \(\xi\)'s of
\(S_0(C_0(\xi;C_m))\), evaluated at \(\xi=0\), is
\begin{equation}\label{eq:v43-block-hessian}
                \operatorname{diag}_{1\le j\le m}
                                      H_j(C_j^*(C_m)).
\end{equation}
\end{proposition}
\begin{proof}
For every \(C\), both \(\Delta_j(C,0)\) and
\(D_\xi\Delta_j(C,0)\) vanish identically. Thus their
derivatives with respect to \(C\) vanish as well.
In the second derivative of the \(j\)-th summand of
\eqref{eq:v43-telescope}, every term containing fewer than
two derivatives in its own \(\xi_j\) is zero at the origin.
The remaining term is \(H_j(C_j^*)\). Summing proves the
diagonal assertion, including the moving intermediate
configurations.
\end{proof}

\subsection{Identification with the inherited harmonic packet}

At the trivial configuration let \(K_j=D^2G_j(I)\) on the
level-\(j\) link variables, with gauge zero modes retained.
Let \(E_j\) be the actual independent-fibre inclusion and
\(R_j\) the boundary-copy right inverse for step \(j\),
including the tree embedding when they are viewed as maps
to all links. Thus
\[
       a_{j-1}=R_ja_j+E_j\eta_j,\qquad
                    L_jE_j=0,\quad L_jR_j=I.
\]
The finite-dimensional notation below includes all cells.
It is equivalent under cell Fourier transformation to the
fixed \(49\)- and \(45\)-coordinate matrices of Shell V4.
Define
\begin{align}
 H_j&=E_j^*K_{j-1}E_j,&
 B_j&=E_j^*K_{j-1}R_j,\notag\\
 M_j&=-H_j^{-1}B_j,&
 T_j&=R_j+E_jM_j.
                                                        \label{eq:v43-harmonic}
\end{align}
Here \(M_j\) is a conditional-centre derivative, not a
martingale increment. Shell V4's \(H_{j-1}\) is the present
\(H_j\); the shift is only an indexing convention.

\begin{theorem}[A depth-independent free norm on the actual history]
\label{thm:v43-free-history}
In the canonical curl normalization, for every depth and
admissible finite volume,
\begin{equation}\label{eq:v43-free-floor}
   \kappa I\preceq H_j\preceq M I,\qquad
                  \kappa=\frac{24}{11635},\quad M<\infty,
\end{equation}
where \(M\) is independent of \(j\) and volume. The linearization
of \eqref{eq:v43-xi} is exactly
\(\xi_j=\eta_j-M_ja_j\). It gives the identities
\begin{align}
 E_j^*K_{j-1}T_j&=0,&
 K_j&=T_j^*K_{j-1}T_j,\notag\\
 \langle a_0,K_0a_0\rangle
     &=\langle a_m,K_ma_m\rangle+
                              \sum_{j=1}^m\langle\xi_j,H_j\xi_j\rangle .
                                                        \label{eq:v43-free-telescope}
\end{align}
In particular the nonlinear centred joint action at
\(C_m=I,\xi=0\) has Hessian between
\(\kappa I\) and \(MI\) in the direct sum of the innovation
Euclidean norms, uniformly in depth and volume.

The kernels of \(H_j,H_j^{-1},M_j,T_j\) have uniform
exponentially weighted row and column bounds in the
\emph{current rescaled cell units}. In the untruncated
Gaussian model at fixed final source, the \(\xi_j\) are
independent centred Gaussians with covariance
\((\beta H_j)^{-1}\).
\end{theorem}
\begin{proof}
At \(I\), the action is critical and the classical minimum
is zero. Differentiating the critical equation
\(D_\eta f_j(C,h_j(C))=0\) gives
\[
                         H_jDh_j(I)+B_j=0.
\]
Consequently \(Dh_j(I)=M_j\). The same equation gives
\(E_j^*K_{j-1}T_j=B_j+H_jM_j=0\). Substituting
\(a_{j-1}=T_ja_j+E_j\xi_j\) into its quadratic energy
therefore eliminates the cross term. Minimizing in
\(\xi_j\) proves the Schur identity for \(K_j\), and
summing these energy identities proves
\eqref{eq:v43-free-telescope}. Inductively this is the
same Gaussian Schur tower as the parallel programme,
not a different action substituted at each step.

Shell V4's complete-momentum fibre comparison now gives
the lower bound \(\kappa=24/11635\) for these very \(H_j\).
Its upper comparison, the bounded primitive curl operator,
and \(\|E_j\|^2\le15\) give a finite common \(M\).
Colour tensoring and finite-torus periodization preserve
the comparison. Shell V5 supplies uniformly bounded
holomorphic symbols on a common strip for \(K_{j-1}\),
the cell precision, \(H_j,H_j^{-1}\), and the harmonic
lift. The finite-range \(E_j,R_j\) and
\eqref{eq:v43-harmonic} give the same strip property
for \(M_j\). Fourier contour displacement on a narrower
strip gives exponential coefficient decay; summation
and periodization give the stated weighted bounds.
No inverse of the retained massless kernel has been used.

At final source \(I\), all nested minima are \(I\).
Proposition~\ref{prop:v43-block-hessian} therefore identifies
the nonlinear joint Hessian there with the displayed
Gaussian blocks. Finally the quadratic energy and unit
Jacobian factor the Gaussian density on the whole
innovation space, proving the independence assertion
for that untruncated reference law.
\end{proof}

The last assertion is not independence under
\eqref{eq:v43-measured}: its nonlinear losses depend on
coarser innovations, its density is not constant, and
its domain is generally not a product. Also the weighted
locality statement is spatial locality of each one-step
harmonic operator in its own units. It is not a bound
on arbitrary products of harmonic lifts through an
unbounded number of levels.

\subsection{The exact price of returning to the raw norm}

For fixed zero final source, let \(U_m\) be the linear map
from the raw history \(\eta\) to the innovations \(\xi\).
It is block triangular with identity diagonal. By
\eqref{eq:v43-free-telescope}, V42's raw Hessian
\(\mathbb H_m\) satisfies
\begin{equation}\label{eq:v43-congruence}
 \mathbb H_m=U_m^*\operatorname{diag}(H_1,\ldots,H_m)U_m,
 \qquad
 \det\mathbb H_m=\prod_{j=1}^m\det H_j.
\end{equation}
The determinants here are on the final-source fibre;
there is no determinant of a retained gauge zero mode.

\begin{proposition}[Unit Jacobian does not imply uniform norm equivalence]
\label{prop:v43-norm-price}
On V42's tori of side \(2^{m+1}\), for \(m\ge4\),
\begin{equation}\label{eq:v43-norm-price}
                 \|U_m^{-1}\|\ge
                  \left(\frac{7\kappa(m-3)}{16}\right)^{1/2}.
\end{equation}
Thus centering repairs the free quadratic history norm
without giving a uniformly bounded inverse coordinate
map in the old unweighted Euclidean history norm.
\end{proposition}
\begin{proof}
Use precisely the actual triangular-profile tangent
\(z_m\) of Theorem~\ref{thm:v42-history-obstruction},
not a new abstract matrix example. With \(N=2^{m+1}\),
that proof gives
\[
 z_m^*\mathbb H_mz_m=N^4,\qquad
                    \|z_m\|_2^2\ge(7/16)(m-3)N^4.
\]
The free floor and \eqref{eq:v43-congruence} imply
\(\kappa\|U_mz_m\|_2^2\le N^4\).
Taking the ratio of these two inequalities gives
\eqref{eq:v43-norm-price}.
\end{proof}

This is a statement about a global history norm, not
about the variance of each individual raw coordinate.
It neither contradicts bounded one-level harmonic
operators nor proves divergence of a specified local
moment. It does rule out transferring estimates on a
fixed ball between these two history norms with one
depth-independent bi-Lipschitz constant.

\subsection{Which nonlinear jets remain after centering}

The following formulas record the derivative convention
needed for this coordinate change. They specialize the
already established conditional-minimum calculus; the
implicit function theorem itself is not a new result.
For a single stage write its variables in the order
\((\eta,c)\), take the reference point to be zero, and let
\[
 H=f_{\eta\eta}(0),\quad B=f_{\eta c}(0),\quad
 M=-H^{-1}B,\quad V(c)=(Mc,c),\quad
                  \mathcal T=D^3f(0),\quad\mathcal F=D^4f(0).
\]
All derivatives in these definitions are true derivatives.
Let \(\mathcal T_\eta[v,w]\) denote the detail covector
whose value on \(u\) is \(\mathcal T[(u,0),v,w]\);
use the analogous convention for \(\mathcal F_\eta\).
If the centre has ordinary homogeneous expansion
\(h(c)=Mc+h_{[2]}(c)+h_{[3]}(c)+O(|c|^4)\), then
\begin{align}
 h_{[2]}(c)
   &=-\tfrac12H^{-1}\mathcal T_\eta[V(c),V(c)],\notag\\
 h_{[3]}(c)
   &=-H^{-1}\left\{
        \mathcal T_\eta[(h_{[2]}(c),0),V(c)]
           +\tfrac16\mathcal F_\eta[V(c),V(c),V(c)]\right\}.
                                                        \label{eq:v43-centre-jets}
\end{align}
Indeed expansion of \(f_\eta(h(c),c)=0\) in degrees
one, two and three gives respectively
\(HM+B=0\), \(Hh_{[2]}+\mathcal T_\eta[V,V]/2=0\),
and the cubic equation displayed above.

Writing \(H(c)=f_{\eta\eta}(h(c),c)\), differentiation gives
\begin{equation}\label{eq:v43-Hprime}
                  DH(0)[v][u,w]
                    =\mathcal T[(u,0),(w,0),V(v)].
\end{equation}
Consequently the centred loss, through total degree three,
is
\begin{equation}\label{eq:v43-centred-cubic}
 \Delta(c,\xi)
   =\tfrac12H[\xi,\xi]
      +\tfrac16\mathcal T[(\xi,0),(\xi,0),(\xi,0)]
      +\tfrac12DH(0)[c][\xi,\xi]+O((|c|+|\xi|)^4).
\end{equation}
For a smooth branch this follows by expanding the integral
formula \eqref{eq:v43-integral-loss}. In particular neither
a pure cubic detail vertex nor a source--two-detail vertex
is removed by a conditional translation.

For \(j<k\), the exact cross-level consequence at the nested
zero branch, with final source fixed, is
\begin{equation}\label{eq:v43-mixed-level}
 D_{\xi_k}D_{\xi_j}^2(S_0\circ X_{\rm joint})(0)[v,u,w]
       =DH_j(C_j^*)[D_{\xi_k}C_j(0)v][u,w].
\end{equation}
Here \(X_{\rm joint}\) means the reconstruction expressed
in \(\xi\) coordinates. To prove the formula, differentiate
\eqref{eq:v43-telescope}. A finer summand \(i<j\) still
has its double zero in \(\xi_i=0\); a coarser summand
\(i>j\) is independent of \(\xi_j\). Only summand \(j\)
remains, and its second detail derivative is
\(H_j(C_j)\). At \(I\), the linear map
\(D_{\xi_k}C_j\) consists of the insertion \(E_k\)
followed by the relevant harmonic lifts \(T_{k-1},\ldots,T_{j+1}\).
For \(k=j+1\) this is just \(E_{j+1}\). This also explains
why one-step spatial locality alone does not yet supply
an all-depth nonlinear vertex bound.

\subsection{A literal Wilson mixed vertex in the centred coordinates}

The following finite test uses the supplied nonabelian
root, not a scalar analogue. Work on a fine torus of side
two with one retained coarse cell. In a colour basis with
\([E_a,E_b]=2\epsilon_{abc}E_c\), choose the flat retained
source
\[
                    C_1=\exp(cE_3),\qquad
                         C_2=C_3=C_4=I.
\]
The boundary-copy configuration is flat: its internal
links are \(I\), and its orientation-one boundary links
are \(C_1\). Its Wilson action is zero. Hence its
primitive classical centre is \(h_1(C)=0\) along this
small source line, by nonnegativity and local uniqueness.
The solved-pivot chart is in the fully active logarithmic
root regime near this background.

Let \(e=((0,0,0,1),1)\) and \(f=((0,0,0,1),2)\) be two
of the actual independent internal links. Set
\(\eta_e=aE_1,\eta_f=bE_2\), with all other independent
coordinates zero, and solve the four pivots as usual.

\begin{proposition}[A nonzero actual source--detail--detail coefficient]
\label{prop:v43-actual-vertex}
For this primitive constrained Wilson action,
\begin{equation}\label{eq:v43-actual-vertex}
       \left.\partial_c\partial_a\partial_b
                         f_1(C,\eta)\right|_{c=a=b=0}=1.
\end{equation}
The same coefficient occurs in its centred loss. It is
a classical cubic response, not a running coupling or a
beta-function coefficient.
\end{proposition}
\begin{proof}
Use right-trivialized tangent coordinates at this flat
source. Differentiating the actual root gives the
orientation-one pivot tangent
\begin{equation}\label{eq:v43-pivot-tangent}
 v_{p_1}
    =-\tfrac12\sum_{\text{free internal orientation one}}
                 (I+\operatorname{Ad}_{C_1^{-1}})v_i
      -\sum_{\text{free boundary orientation one}}v_i.
\end{equation}
To see the coefficient directly, each internal link
occurs in two chronological paths, on opposite sides of
the background \(C_1\); in right coordinates their
derivatives are \(v_i\) and
\(\operatorname{Ad}_{C_1^{-1}}v_i\).
The pivot root derivative is \(1/8\), while a single
path has weight \(1/16\). Every nonpivot boundary link
has two identical occurrences and therefore coefficient
minus one. This proves \eqref{eq:v43-pivot-tangent}.
The other orientations have \(C_\mu=I\).

In particular, writing \(p_\mu=(e_\mu,\mu)\),
the two selected variations have
\[
 v_{p_1}=-aE_1+caE_2+O(c^2a),\qquad
                              v_{p_2}=-bE_2.
\]
The \(caE_2\) term is the moving fibre tangent. It cannot
be discarded just because the background action is zero.

For a reproducible incidence evaluation, let \(d\) be the
scalar \(96\)-by-\(64\) curl matrix of this fine torus.
Let \(E_0\) be its \(64\)-by-\(45\) free embedding:
the free rows form the identity and each pivot row is
minus the sum of all free rows of that orientation.
Let \(P\) have value \(1/2\) in the orientation-one
pivot row at each free internal orientation-one column,
and zero elsewhere. Put
\[
 J=\operatorname{ad}_{E_3}
     =\begin{pmatrix}0&-2&0\\2&0&0\\0&0&0\end{pmatrix}.
\]
At a plaquette \((x;1,\nu)\) with \(x_1=1\), the
right-trivialized linear curvature is
\[
 \operatorname{Ad}_{C_1}
          (v_{x,1}+v_{x+e_1,\nu}-v_{x+e_\nu,1})-v_{x,\nu}.
\]
At all other plaquettes it is the usual scalar curl
applied to the colour components. Let \(T\) be the scalar
matrix consisting of those first three signed terms
on those plaquettes, and zero on other rows. The
linear-curvature map and its source derivative at zero
are therefore
\[
                 F_0=dE_0,\qquad F_1=dP+TE_0,
\]
with colour factors \(I\) and \(J\), respectively.
Every background plaquette is \(I\), so the Hessian of
the Wilson action is the Gram matrix of its linear
curvature. The background action gradient vanishes,
and hence no coordinate-acceleration term is missing
from this Hessian. Its source derivative is
\begin{equation}\label{eq:v43-Gram-derivative}
                       (F_0^*F_1-F_1^*F_0)\otimes J.
\end{equation}

Here is the complete scalar overlap needed for the selected
two columns. Write \(u=e-p_1\) and \(v=f-p_2\), now
regarding link labels as unit scalar link vectors. Direct
plaquette incidence gives
\[
        \langle dp_1,dv\rangle=1,\qquad
        \langle Tu,dv\rangle=-1,\qquad
        \langle du,Tv\rangle=-1.
\]
For the first equality, \(p_1\) and \(f\) share no
plaquette, whereas \(p_1\) and \(p_2\) share just the
\((1,2)\) plaquette based at \((1,1,0,0)\), with opposite
signs. For each of the last two equalities, the only
common nonzero row is that same \((1,2)\) plaquette based
at \((1,1,0,0)\). In that row \(Tu=du=1\) and
\(Tv=dv=-1\), proving both overlaps. These statements
also follow immediately by writing its four signed edges.

Since the \(e\)-column of \(P\) is \(p_1/2\) and its
\(f\)-column is zero, \eqref{eq:v43-Gram-derivative}
on colours \(E_1,E_2\) has value
\[
       \langle dp_1,dv\rangle
       +2\{\langle Tu,dv\rangle-\langle du,Tv\rangle\}
                                  =1+2(-1+1)=1.
\]
Thus the pivot part is one and the
plaquette-transport parts cancel. This proves
\eqref{eq:v43-actual-vertex}. Along the declared source
line \(h_1(C)=0\) and \(G_1(C)=0\), so subtracting the
classical minimum and centering do not change it.
\end{proof}

This witness uses a retained flat source, not a specified
coarser-shell detail mode. It certifies that a source
dependence of the centred fibre Hessian really occurs
for the supplied Wilson block. Formula
\eqref{eq:v43-mixed-level} describes its cross-level
insertion when such a source variation is generated by
coarser innovations; a particular nonzero cross-level
projection is not asserted by the one-cell test alone.
The earlier V38 cubic bank already contains additional
actual Wilson interactions. No new claim that the
centred action is Gaussian is compatible with these data.

\subsection{Transport the cutoff, not only the density}

Suppose the compatible raw domain uses the actual radial
group-coordinate bounds \(|\eta_{j,i}|<r_j\), together
with any branch-admissibility restrictions already
declared. Its centred image is exactly
\begin{equation}\label{eq:v43-domain}
 \Omega_m(C_m)=
  \left\{\xi:
    |\xi_{j,i}+h_{j,i}(C_j(\xi_{>j};C_m))|<r_j
                       \text{ for every }j,i\right\},
\end{equation}
with those same restrictions transported too. These
inequalities couple levels and depend on the source.
On a specified set \(\mathcal A\) of histories on which
all branches exist and
\(|h_{j,i}(C_j)|\le r_j/2\), the triangle inequality gives
the precise useful inclusions
\begin{equation}\label{eq:v43-domain-sandwich}
 \mathcal A\cap\prod_{j,i}B(r_j/2)
       \ \subset\ \Omega_m(C_m)\cap\mathcal A
       \ \subset\ \mathcal A\cap\prod_{j,i}B(3r_j/2).
\end{equation}
The set \(\mathcal A\) may itself depend on the coarser
innovations. These are not unconditional product-domain
inclusions.

There is a simple exact boundary calculation showing why
replacing \eqref{eq:v43-domain} by a fresh rectangle is
not an innocuous use of the unit Jacobian.

\begin{proposition}[A curvature lost by discarding a transported wall]
\label{prop:v43-wall}
Take the two-scalar Gaussian density
\[
           \exp\{-\tfrac{\beta}{2}[(\eta-\alpha a)^2+a^2]\}
                  \,\mathbf 1_{\{|\eta|<r,\ |a|<R\}},
\]
where \(\beta,r,R>0\) and \(\alpha\ne0\).
The innovation \(\xi=\eta-\alpha a\) has Jacobian one.
After integrating \(\xi\), the marginal of \(a\), apart
from its Gaussian factor and a constant, contains
\[
 p_\beta(a)=
  \Phi(\sqrt\beta(r-\alpha a))
       -\Phi(\sqrt\beta(-r-\alpha a)),
\]
where \(\Phi\) and \(\phi\) are the standard normal
distribution function and density. Its extra potential
has strictly positive curvature at zero:
\begin{equation}\label{eq:v43-wall-curvature}
 [-\log p_\beta]''(0)=
   \frac{2\alpha^2\beta^{3/2}r\,
                       \phi(\sqrt\beta r)}
        {2\Phi(\sqrt\beta r)-1}>0.
\end{equation}
Replacing the transported domain by \(|\xi|<r\)
sets this extra curvature to zero.
\end{proposition}
\begin{proof}
The exact new bound is \(|\xi+\alpha a|<r\), and
the action is \(\beta(\xi^2+a^2)/2\). Integrating the
normal density on this shifted interval gives
\(p_\beta(a)\). The symmetric interval gives
\(p_\beta'(0)=0\), while direct endpoint differentiation
gives
\[
             p_\beta''(0)=
                    -2\alpha^2\beta^{3/2}r\phi(\sqrt\beta r).
\]
The second derivative of \(-\log p_\beta\) proves
\eqref{eq:v43-wall-curvature}. A fixed centred interval
instead gives the constant factor \(p_\beta(0)\).
\end{proof}

For any fixed positive margin \(\delta\), on
\(|\alpha a|\le r-\delta\), the difference \(1-p_\beta\)
and any fixed finite number of its source derivatives
are bounded by a polynomial in \(\beta\) times
\(e^{-\beta\delta^2/2}\), for \(\beta\ge1\). This follows
from the normal tail integral and derivatives of its
endpoint density. The same holds for \(-\log p_\beta\)
and its derivatives once \(\beta\) is large enough that
\(p_\beta\ge1/2\). The margin cannot be omitted:
if \(\alpha a=r\) lies within the allowed \(a\)-interval,
then \(p_\beta(a)\to1/2\), so
\(-\log p_\beta(a)\to\log2\).
This is an illustrative boundary identity, not a
numerical prediction for the Yang--Mills wall. It proves
the need for a margin and a comparison before discarding
transported boundaries even in an exactly Gaussian action.

\subsection{Verification and the next unresolved estimate}

The exact certificate is
\begin{center}\small
\path{scripts/verify_ym_f14_v43_centered_innovations.py}.
\end{center}
It has 7,344 bytes and SHA--256
\begin{center}\scriptsize\ttfamily
53D21706ED63E8A8565444A488946DB6323550860E32E01673E15C8DE94B90CE.
\end{center}
It reconstructs all \(64\) links, \(96\) plaquettes and
\(45\) independent scalar coordinates of the actual
one-cell flat-source test. The differentiated scalar
antisymmetric bank has \(336\) nonzero upper-triangular
pairs. The selected colour coefficient is exactly one,
with pivot contribution one and transport contribution
zero. A second evaluation differentiates the exact
\(\operatorname{Ad}_{\exp(cE_3)}\) rotation in the literal
plaquette tangent, independently of the matrix
antisymmetrization. The three scalar overlaps used in
the proof are individually checked.

A separate declared three-scalar calibration checks a
nonlinear triangular map with determinant one, centred
Hessian \(\operatorname{diag}(3,2,5)\), raw determinant
\(30\), and surviving mixed cubic equal to one.
The centre-jet formulas are checked against a direct
critical-equation series, with ordinary quadratic and
cubic coefficients \(-14/125\) and \(-5899/253125\).
The transported-wall curvature is verified symbolically.
These finite algebra checks do not establish the
all-depth estimates; the inherited Gaussian bounds,
the exact identities, and the norm comparison above
give the stated analytic conclusions.

V42 has 994,756 bytes and SHA--256
\begin{center}\scriptsize\ttfamily
401A07F4070145F856894B409D042AB0286C1087D1EDC805E5F3D7A26FF20C38.
\end{center}
Its body is preserved except for the title version and
this terminal addition. V43 replaces V42 as the sole
active main file. Archives and companions are unchanged;
the next scheduled SM/NS comparison remains V45.

\begin{firewall}
The free conditioning obstruction in the raw history has
now been resolved by an exact identification with
harmonic innovations: their pointwise zero-background
joint Hessian has a depth-independent positive floor.
This is not a nonlinear all-depth convexity theorem.
The explicit inverse-map lower bound prevents a free
transfer of the old raw-domain estimates, and the actual
Wilson vertex shows that mixed interactions remain.
The next acquisition is a uniform nonlinear
conditional-centre and Hessian-response budget in these
coordinates, together with admissible-history and
transported-wall estimates for the actual measure.
Compact entry, a cofinal interacting continuum limit,
and a positive physical Hamiltonian mass gap remain
unproved.
\end{firewall}
\section{V44: longitudinal cancellation, an exact uniformly controlled coarse Coulomb source, and a nonlinear classical history chart}
\label{sec:v44}

\subsection{The upstream step and the three newly supplied sources}

V43 controls the free innovation norm, but does not provide
a common nonlinear history neighbourhood. Shell V12 already
proves a uniform nonlinear balanced root map on a small
global counting-\(\ell^4\) ball. Its subsequent actual
minimum theorem is restricted to one coarse cell: on
several coarse cells the balanced output need not be
logarithmically Coulomb, and linear transverse projection
is not an exact gauge transformation. The new estimate
below controls precisely that output's longitudinal part.
It permits an exact nonlinear gauge correction with
constants independent of both refinement and coarse volume.
The resulting classical constrained-action theorem uses a
global critical-energy ball, not a thermodynamically typical
fluctuation domain.

During this iteration three additional files were supplied.
Their relevant current theorem statements and proofs,
source audits, and scope conclusions were inspected:
\begin{center}\small
\path{Renewal_SM_Source_Selection_Research_Notes_V17.tex},\\
\path{Renewal_Native_Regenerative_Stationarity_Closure_Notes_V21.tex},\\
\path{Renewal_Exact_Action_Einstein_Bridge_Research_Notes_V16.tex}.
\end{center}
They remain in the supplied Downloads locations and are
not edited or merged into the active companion lineages.
Their contents are mathematical source material, not
instructions changing this programme.

The SM source-selection note, V17, proves a sharp five-total-slot
bound for a determinant-sensitive completely positive port
on its specified native representation, and constructs a
two-input, three-output comparison instrument attaining it.
Its audit also shows that the unchanged full-record source
does not thereby acquire that covariance. The pertinent
lesson for YM is to identify the actual measured operation:
an available representation or a comparison instrument is
not a supplied physical transfer law. No determinant-port
theorem is used as a YM spectral estimate here.

Native V21 proves, for its normalized three-spatial-dimensional
Wilson Hamiltonian and zero initial real Hamilton--Jacobi
phase, a transverse caustic by
\(t=\pi h/(4\sqrt3)\). Its positive phase-space alternative
is established for the quadratic diagnostic, with an
unsubtracted all-mode uncertainty cost of order
\(\eta h^{-4}\). Its proof retains harmonic modes and
explicitly does not construct the nonlinear compact singlet
limit. Thus a global real-phase Hessian on a fixed time
interval is not an appropriate replacement for the missing
YM continuum construction. The elliptic classical minimization
below does not require such a phase or send a physical
Planck parameter to zero.

Bridge V16 proves that its actual prepared scalar source
excites a bare diagonal soft sector, while the retained
flux excludes every nonzero limiting vector in that sector
from finite flux energy. Its full-energy compression
therefore cannot be replaced by the bare diagonal inverse;
nor does that compression determine the Schur complement
after the remaining components are allowed to respond.
Here this reinforces the need to retain the source-normal
motion and its quadratic action. We do so explicitly.
No bridge flux operator or scalar inverse is identified
with a YM fibre Hessian.

These are concrete scope checks, not independent verification
of the entire three new source lineages. Their hashes are
recorded at the end. They support retaining the exact action,
gauge correction and measure distinction rather than changing
the goal to a semiclassical or source-selection problem.

\subsection{Spaces and inherited analytic inputs}

Let the coarse torus have side \(N\ge1\), the fine torus
side \(NL\), and \(L=2^m\). All lattice norms use
\emph{unnormalized counting measure}. To avoid an ambiguity,
in this section \(D\) is the scalar-to-link forward
gradient, \(D^*\) its divergence adjoint, and \(d\) is
the link-to-plaquette curl. Thus \(dD=0\).
Let \(\Pi\) be the orthogonal projection onto
\(\ker D^*\), preserving constant links, and put
\(P_\parallel=I-\Pi\). Subscripts \(f,c\) distinguish
fine and coarse lattices.

The fine Coulomb norm and coarse source norm are
\[
 \|A\|_{\mathcal E}=\|d_fA\|_2+\|A\|_4,\qquad
 \|c\|_{\mathcal Y_N}=\|d_cc\|_2+N|\overline c|,
                              \quad D_c^*c=0.
\]
For a mean-zero site field set
\(\|\varphi\|_{\mathcal G_N}=\|D_c\varphi\|_2\).
The volume-independent lattice Sobolev inequality from
Shell V9 gives
\begin{equation}\label{eq:v44-Sobolev}
 \|\varphi\|_4\le C_S\|\varphi\|_{\mathcal G_N},
 \qquad
 \|w\|_{\mathcal E}\le(1+C_S)\|d_fw\|_2
             \quad(D_f^*w=0,\ \overline w=0).
\end{equation}
These are counting-norm statements; in particular
\(\|\varphi\|_\infty\le\|\varphi\|_4\).
For \(N=1\) the mean-zero site space is zero.

Write \(B_L\) for the composed balanced \emph{linear}
average and \(\mathcal F_{L,N}\) for the actual nonlinear
balanced root logarithm. The latter uses the original
ordered paths and only the explicit coarse gauge balancing
of Shell V10. Shell V12 proves, on a common small complex
\(\ell^4\) ball,
\begin{equation}\label{eq:v44-root-input}
 \mathcal F_{L,N}(A)=B_LA+\mathcal Z_{L,N}(A),\qquad
 \|\mathcal Z_{L,N}(A)\|_2\le C\|A\|_4^2,\qquad
 \|\mathcal F_{L,N}(A)\|_4\le\tfrac54\|A\|_4.
\end{equation}
The remainder is holomorphic with uniform fixed-order
mixed derivatives on smaller balls, independently of
\(L,N\). Exact real orbit equivalence with the literal
nested root is part of that input. It is not equality
of arbitrary representative-valued densities.

Let \(\mathcal Q_L=\Pi_cB_L\) on fine Coulomb fields.
Shell V9 supplies the harmonic lift \(\mathcal H_{L,N}\)
and the detail space
\[
 \mathcal V_{L,N}
   =\{w:D_f^*w=0,\ \mathcal Q_Lw=0\},\qquad
                         \|w\|_V=\|d_fw\|_2.
\]
Every detail has zero mean. Uniformly in \(L,N\),
\begin{align}
 \mathcal Q_L\mathcal H_{L,N}c&=c,&
 \|\mathcal H_{L,N}c\|_{\mathcal E}&\le K_0\|c\|_{\mathcal Y_N},
                                                        \label{eq:v44-split}\\
 \langle d_f\mathcal H_{L,N}c,d_fw\rangle&=0,&
 \|d_f\mathcal H_{L,N}c\|_2^2
                                  &=\langle c,S^{(m)}c\rangle.\notag
\end{align}
Also \(\|\mathcal Q_LA\|_{\mathcal Y_N}\le K_0\|A\|_{\mathcal E}\).
The quadratic form \(S^{(m)}\) is bounded by a common
constant times \(\|c\|_{\mathcal Y_N}^2\); it is not
inverted on retained gauge or harmonic zero modes.

Finally the complete Wilson action has the inherited
holomorphic decomposition
\begin{equation}\label{eq:v44-Wilson-input}
 \mathcal S(A)=\tfrac12\|d_fA\|_{\rm bil}^2+\mathcal N(A),
 \qquad
 | \mathcal N(A)|\le C\|A\|_{\mathcal E}^3,\quad
 \|D\mathcal N(A)\|\le C\|A\|_{\mathcal E}^2,\quad
 \|D^2\mathcal N(A)\|\le C\|A\|_{\mathcal E}.
\end{equation}
The derivative norms are the corresponding energy
multilinear norms. The suffix denotes bilinear
complex extension, not a Hermitian norm in a holomorphic
formula. Positivity statements below concern real fields.

\subsection{The longitudinal average gains the required energy bound}

The whole average \(B_L\) cannot be bounded from critical
energy to counting \(\ell^2\) by ignoring soft momenta.
Its \emph{longitudinal part on fine Coulomb inputs} has
an additional cancellation.

\begin{theorem}[Uniform longitudinal estimate for the actual average]
\label{thm:v44-longitudinal}
There is a constant \(C_\parallel<\infty\), independent
of \(L,N\), such that
\begin{equation}\label{eq:v44-longitudinal}
              \|P_{\parallel,c}B_LA\|_2
                              \le C_\parallel\|d_fA\|_2,
                           \qquad D_f^*A=0.
\end{equation}
One valid constant is
\begin{equation}\label{eq:v44-long-constant}
 C_\parallel^2=\frac{\pi^4}{144}+64\mathcal A_2,\qquad
 \mathcal A_2=
 \sum_{M\in\mathbb Z^4\setminus\{0\}}
       |M|_2^{-2}\prod_{\nu=1}^4(1+|M_\nu|)^{-1}<\infty.
\end{equation}
The finiteness of this alias sum is the already proved
Shell V9 envelope, not a new summability claim.
\end{theorem}
\begin{proof}
For a nonzero coarse momentum \(k\in[-\pi,\pi]^4\),
use centered aliases \(p_M=(k+2\pi M)/L\), and write
\[
 q(p)_\mu=e^{ip_\mu}-1,\quad
 \lambda(p)=|q(p)|^2,\quad
 D_L(t)=\sum_{r=0}^{L-1}e^{irt},\quad
 b_M=L^{-4}\prod_\nu D_L(p_{M,\nu}).
\]
The actual linear symbol is
\[
           W_M=b_M\operatorname{diag}_\mu D_L(p_{M,\mu}).
\]
In counting-unitary Fourier coordinates the alias row
is \(L^{-2}W_M\), not \(W_M\). This is the same \(L^{-4}\)
factor in its covariance Gram as in Shell V4's exact
Gaussian recursion.

Since \(q(k)_\mu=D_L(p_\mu)q(p)_\mu\), there is the
exact identity
\begin{equation}\label{eq:v44-long-cancellation}
 \Pi(p)W_M^*q(k)
   =\overline b_M\Pi(p)
              \operatorname{diag}_\mu
                         \bigl(|D_L(p_\mu)|^2-L^2\bigr)q(p).
\end{equation}
Subtracting \(L^2q(p)\) is legitimate because
\(\Pi(p)q(p)=0\). This is where fine Coulomb data enter;
the estimate is not asserted for arbitrary fine inputs.

At the principal alias \(p=k/L\), the exact moment
\[
 \sum_{r,t=0}^{L-1}(r-t)^2=\frac{L^2(L^2-1)}6
\]
and \(1-\cos x\le x^2/2\) imply
\[
 0\le L^2-|D_L(k_\mu/L)|^2
      \le \frac{L^2|k_\mu|^2}{12}.
\]
As \(|b_0|\le1\), \eqref{eq:v44-long-cancellation}
therefore bounds the principal energy-weighted row Gram by
\begin{equation}\label{eq:v44-principal}
 L^{-4}\frac{\|\Pi(p_0)W_0^*q(k)\|^2}
                        {\lambda(p_0)\lambda(k)}
       \le\frac{|k|_2^4}{144\lambda(k)}
       \le\frac{\pi^4}{144}.
\end{equation}
For the last inequality use
\(\lambda(k)\ge4|k|_2^2/\pi^2\) and \(|k|_2^2\le4\pi^2\).

For a nonprincipal centered alias, Shell V9's elementary
geometric-sum bounds give
\[
 \|W_M\|\le16L\prod_\nu(1+|M_\nu|)^{-1},
 \qquad
                    \lambda(p_M)\ge4|M|_2^2/L^2.
\]
Dropping only the two orthogonal projectors, its
energy-weighted row Gram is at most
\[
 64|M|_2^{-2}\prod_\nu(1+|M_\nu|)^{-2}.
\]
Summation is bounded by \(64\mathcal A_2\).

For completeness, on transverse fine coefficients
\(a_M=\Pi(p_M)a_M\), weighted Cauchy--Schwarz gives
\[
 \left|P_{\parallel,c}L^{-2}\sum_MW_Ma_M\right|^2
 \le\left\{
    L^{-4}\sum_M
      \frac{\|\Pi(p_M)W_M^*q(k)\|^2}
           {\lambda(p_M)\lambda(k)}\right\}
                         \sum_M\lambda(p_M)|a_M|^2.
\]
Insert the preceding bounds and sum over coarse momenta.
At \(k=0\), \(P_{\parallel,c}=0\); thus no inverse is
assigned to the constant fine mode or to a retained zero
mode. All geometric sums use their continuous values at
zero arguments. Parseval proves
\eqref{eq:v44-longitudinal}, with the same argument
componentwise in colour and for complex Coulomb fields.
\end{proof}

\subsection{An exact coarse gauge correction on a mixed norm domain}

For a coarse link logarithm \(y\), introduce
\[
       \|y\|_{\mathcal W_N}
            =\|y\|_4+\|P_{\parallel,c}y\|_2.
\]
This is a genuine norm; it is stronger than the bare
counting-\(\ell^4\) norm. Let
\(\Delta_c=D_c^*D_c\) on mean-zero site fields.
The operator \(\Delta_c^{-1}D_c^*\), as a map from link
\(\ell^2\) to \(\mathcal G_N\), has norm one. This
follows by the orthogonal-gradient projection identity.
No uniform \(\ell^2\)-to-\(\ell^2\) bound for
\(\Delta_c^{-1}\) is asserted.

For small \(y,\varphi\), define the literal gauge word
\begin{equation}\label{eq:v44-gauge-word}
 \mathfrak C(y,\varphi)_{x,\mu}
       =\log\bigl(e^{-\varphi_x}e^{y_{x,\mu}}
                                  e^{\varphi_{x+e_\mu}}\bigr)
       =y_{x,\mu}+(D_c\varphi)_{x,\mu}
                                  +\mathcal R(y,\varphi)_{x,\mu}.
\end{equation}
Only the common identity logarithm is used.

\begin{lemma}[The gauge remainder lies in the contracting output norm]
\label{lem:v44-gauge-remainder}
On fixed sufficiently small balls, independently of \(N\),
\begin{align}
 \|\mathcal R(y,\varphi)\|_2
      &\le C(\|y\|_4+\|\varphi\|_{\mathcal G_N})
                                      \|\varphi\|_{\mathcal G_N},
                                                        \label{eq:v44-gauge-R}\\
 \|D_\varphi\mathcal R(y,\varphi)\|_{\mathcal G_N\to\ell^2}
      &\le C(\|y\|_4+\|\varphi\|_{\mathcal G_N}).\notag
\end{align}
The maps are holomorphic on the corresponding complex
balls, and \(\mathcal R(y,0)=0\).
\end{lemma}
\begin{proof}
On one link put \(u=\varphi_x\) and
\(v=\varphi_{x+e_\mu}-\varphi_x\). Factor its group word as
\[
       (e^{-u}e^ye^u)(e^{-u}e^{u+v}).
\]
The logarithm of the first factor is \(y+O(|u||y|)\);
that of the second is \(v+O(|u||v|)\).
The latter error vanishes when either \(u=0\) or \(v=0\).
The local logarithm of their product has the additional
error \(O(|y||v|)\), with uniformly bounded analytic
coefficients on a fixed finite-dimensional polydisc.
Consequently
\[
 |\mathcal R(y,\varphi)_e|
             \le C\{|u||y_e|+|u||v|+|y_e||v|\}.
\]
This formulation retains a derivative in the pure-gauge
term; replacing it by a raw \(|\varphi|^2\) estimate
would lose that information.

Each site belongs to a fixed number of oriented links.
Apply Holder with two counting-\(\ell^4\) factors,
\eqref{eq:v44-Sobolev}, and
\(\|D_c\varphi\|_4\le\|D_c\varphi\|_2\).
These prove the first estimate, with constants independent
of the number of sites. Differentiating the same local
analytic factors, or their convergent power series,
gives the second estimate by the identical Holder bounds.
Higher factors are bounded in maximum norm by the small
balls. Sobolev also keeps \(\varphi\) pointwise in the
primitive polydisc. The argument applies to the complex
coefficient norms, proving holomorphy.
\end{proof}

\begin{theorem}[Uniform local Coulomb retraction on the mixed domain]
\label{thm:v44-gauge}
There are \(r,C>0\), independent of \(N\), such that for
\(\|y\|_{\mathcal W_N}<r\) there is a unique sufficiently
small mean-zero solution \(\varphi_N(y)\) of
\[
                  D_c^*\mathfrak C(y,\varphi_N(y))=0.
\]
It is holomorphic on a common complex mixed-norm ball.
Writing
\(\Gamma_N(y)=\mathfrak C(y,\varphi_N(y))\), one has
\begin{align}
 \|\varphi_N(y)\|_{\mathcal G_N}
       &\le2\|P_{\parallel,c}y\|_2,\notag\\
 \Gamma_N(y)&=\Pi_cy+\mathcal K_N(y),&
 \|\mathcal K_N(y)\|_2&\le C\|y\|_{\mathcal W_N}^2.
                                                        \label{eq:v44-gauge-bounds}
\end{align}
Both maps have uniform fixed-order mixed derivatives
on smaller balls. On real data \(\Gamma_N\) is an
actual coarse gauge transformation; it is not simply
\(\Pi_c y\). Constant global colour conjugations are
retained and commute with this construction.
\end{theorem}
\begin{proof}
The exact fixed-point equation on mean-zero site fields is
\begin{equation}\label{eq:v44-gauge-FP}
 \varphi=-\Delta_c^{-1}D_c^*y
                   -\Delta_c^{-1}D_c^*\mathcal R(y,\varphi).
\end{equation}
Its linear term has \(\mathcal G_N\) norm exactly
\(a=\|P_{\parallel,c}y\|_2\). By
Lemma~\ref{lem:v44-gauge-remainder}, for \(r\) small the
right side maps the \(2a\)-ball into itself and has
Lipschitz constant at most \(1/2\) there. More generally
it is a contraction on one fixed small ball, so it
also gives uniqueness among all sufficiently small
solutions. When \(a=0\), the solution is exactly zero.
The same estimates on complex balls and uniform Picard
iteration give a holomorphic map; Cauchy estimates give
the derivative assertions.

Apply \(D_c\) to \eqref{eq:v44-gauge-FP}:
\(D_c\varphi=-P_{\parallel,c}(y+\mathcal R)\).
Substitute this identity in \eqref{eq:v44-gauge-word}.
It follows that
\[
                  \mathcal K_N(y)
                         =\Pi_c\mathcal R(y,\varphi_N(y)).
\]
The counting-\(\ell^2\) norm of \(\Pi_c\) is one,
so \eqref{eq:v44-gauge-bounds} follows. Real site
exponentials are genuine \(\mathrm{SU}(2)\) gauges.
The equations, mean-zero condition, and norms are
equivariant under constant conjugation; uniqueness
proves the final statement.
\end{proof}

This theorem does not assert a global slice on all
compact configurations, or on a bare \(\ell^4\) ball
without the longitudinal condition. Its mixed domain
is exactly what Theorem~\ref{thm:v44-longitudinal}
will supply for the actual root output.

\subsection{The gauge correction has a nonzero quadratic jet}

Let \(\varphi_{[1]}=-\Delta_c^{-1}D_c^*y\), and for
an edge \(e=(x,x+e_\mu)\) put
\begin{equation}\label{eq:v44-gauge-Q2}
 \mathcal R_{[2]}(y,\varphi)_e
   =-\tfrac12[\varphi_x+\varphi_{x+e_\mu},y_e]
                      -\tfrac12[\varphi_x,\varphi_{x+e_\mu}].
\end{equation}
Direct multiplication of the three exponential series,
or BCH through degree two, gives
\begin{align}
 \varphi_{[2]}
     &=-\Delta_c^{-1}D_c^*
                       \mathcal R_{[2]}(y,\varphi_{[1]}),\notag\\
 \Gamma_N(y)
     &=\Pi_cy+
       \Pi_c\mathcal R_{[2]}(y,\varphi_{[1]})
                                +O(\|y\|_{\mathcal W_N}^3).
                                                        \label{eq:v44-gauge-jets}
\end{align}
The remainders here are bounded in the output norms
of Theorem~\ref{thm:v44-gauge}. These are ordinary
homogeneous coefficients, not second derivatives.

For an exact finite check take \(N=2\) and the coarse
field \(y\) with only
\[
 y_{(0,0,0,0),1}=E_1,\qquad
 y_{(1,0,0,0),2}=E_2
\]
nonzero, interpreting \(ty\) as a small one-parameter
input. In the convention \([E_a,E_b]=2\epsilon_{abc}E_c\),
the quadratic coefficient of \(\Gamma_2(ty)\) on its
first listed edge in colour \(E_3\) is \(343/3072\).
Its full squared counting norm is \(79139/2654208\).
The certificate below obtains these exact rational
values by the mean-zero finite Laplacian inverse and
also checks all \(64\) literal quaternion gauge words.
In particular, a linear projection would omit a
genuinely nonzero correction.

The signs can be checked without those numerical values.
For a pure gauge
\(e^{y_e(t)}=e^{-t\chi_x}e^{t\chi_{x+e_\mu}}\),
with mean-zero \(\chi\), the first two input coefficients are
\[
        y_{[1]}=D_c\chi,\qquad
        y_{[2],e}=-\tfrac12[\chi_x,\chi_{x+e_\mu}].
\]
Since \(\varphi_{[1]}=-\chi\), substitution in
\eqref{eq:v44-gauge-Q2} gives
\(\mathcal R_{[2]}(y_{[1]},-\chi)_e
       =[\chi_x,\chi_{x+e_\mu}]/2\).
Thus the complete second projected coefficient
\(\Pi_c(y_{[2]}+\mathcal R_{[2]})\) is zero, as
required by the exact gauge word. Deleting the
quadratic term would fail this cancellation.

\subsection{The actual multicell source now has a common critical domain}

\begin{theorem}[Uniform nonlinear source on the exact coarse orbit]
\label{thm:v44-source}
For fine Coulomb fields in one fixed small critical
ball, define
\begin{equation}\label{eq:v44-source}
                 \mathcal O_{L,N}(A)
                    =\Gamma_N(\mathcal F_{L,N}(A)).
\end{equation}
This map is holomorphic, takes values in the coarse
Coulomb space, and, independently of \(L,N\), satisfies
\begin{equation}\label{eq:v44-source-remainder}
 \mathcal O_{L,N}(A)=\mathcal Q_LA+\mathcal Z^\circ_{L,N}(A),
 \qquad
 \|\mathcal Z^\circ_{L,N}(A)\|_{\mathcal Y_N}
                                  \le C\|A\|_{\mathcal E}^2.
\end{equation}
All fixed-order mixed remainder derivatives on smaller
balls have uniform energy-to-source bounds; in particular
\[
 \|D\mathcal Z^\circ_{L,N}(A)\|_{\mathcal E\to\mathcal Y_N}
                  \le C\|A\|_{\mathcal E},\qquad
                     D\mathcal O_{L,N}(0)=\mathcal Q_L .
 \]
For real inputs \(e^{\mathcal O_{L,N}(A)}\) lies on
the actual literal nested-root coarse gauge orbit.
\end{theorem}
\begin{proof}
Let \(y=\mathcal F_{L,N}(A)\). By
\eqref{eq:v44-root-input} and
Theorem~\ref{thm:v44-longitudinal},
\begin{equation}\label{eq:v44-output-mixed}
 \|y\|_{\mathcal W_N}
    \le\tfrac54\|A\|_4+C_\parallel\|d_fA\|_2+C\|A\|_4^2.
\end{equation}
Thus a common small critical ball is mapped into the
mixed domain of the exact gauge theorem. In counting
\(\ell^2\),
\[
 \mathcal O_{L,N}(A)-\mathcal Q_LA
       =\Pi_c\mathcal Z_{L,N}(A)+\mathcal K_N(y),
 \qquad
 \|\Pi_c\mathcal Z_{L,N}(A)+\mathcal K_N(y)\|_2
                                           \le C\|A\|_{\mathcal E}^2.
\]
For a coarse Coulomb remainder \(z\),
\[
 \|z\|_{\mathcal Y_N}
     =\|d_cz\|_2+N|\overline z|
       \le(4+N^{-1})\|z\|_2\le5\|z\|_2.
\]
Here the mean estimate uses \(N^4\) sites and the
curl norm is at most four. This proves
\eqref{eq:v44-source-remainder}.
All constituent maps are holomorphic on the common
complex domain. Cauchy estimates applied to the
quadratically vanishing remainder give the displayed
first-derivative and all fixed-order bounds, without
a volume-dependent norm equivalence.
Finally the balancing in \(\mathcal F\) and the
correction \(\Gamma\) are both actual coarse gauge
transformations on real inputs. Their composition
proves the orbit statement.
\end{proof}

The improvement over Shell V12 is not another proof
of its root-map estimate. It is the longitudinal
cancellation and exact controlled gauge correction
which make that estimate an onto source chart in
the Gaussian coarse Coulomb coordinates for every
number of cells.

\subsection{The full nonlinear constrained action, including source-normal motion}

For \(w\in\mathcal V_{L,N}\), seek the normal variable
\(q=q_{L,N}(c,w)\in\ker D_c^*\) from
\begin{equation}\label{eq:v44-normal-equation}
 \mathcal O_{L,N}(\mathcal H_{L,N}q+w)=c,\qquad
 q=c-\mathcal Z^\circ_{L,N}(\mathcal H_{L,N}q+w).
\end{equation}
The notation \(q\) here is a source field, not the
Fourier difference symbol of the earlier proof.

\begin{lemma}[Uniform exact source chart]
\label{lem:v44-chart}
On common small complex balls in
\(\mathcal Y_N\oplus\mathcal V_{L,N}\),
\eqref{eq:v44-normal-equation} has a unique small
solution, holomorphic with uniform fixed-order
derivative bounds, and
\begin{equation}\label{eq:v44-normal-bound}
 q_{L,N}(c,w)=c+O((\|c\|_{\mathcal Y_N}+\|w\|_V)^2).
\end{equation}
The field
\(A_{L,N}(c,w)=\mathcal H_{L,N}q_{L,N}(c,w)+w\)
is an exact real root-orbit constraint chart.
\end{lemma}
\begin{proof}
The critical splitting \eqref{eq:v44-split} and
\eqref{eq:v44-Sobolev} bound its input energy by
\(C(\|q\|_{\mathcal Y_N}+\|w\|_V)\).
Theorem~\ref{thm:v44-source} then makes the second
equation in \eqref{eq:v44-normal-equation} a
contraction in \(q\) on common sufficiently small
balls, with Lipschitz constant at most \(1/2\).
Its correction is quadratic, which gives
\eqref{eq:v44-normal-bound}. Complex Picard
iteration and Cauchy estimates prove the analytic
assertions. The first equation and the exact real
orbit identification prove the final assertion.
\end{proof}

Let \(F_{L,N}(c,w)=\mathcal S(A_{L,N}(c,w))\).
Harmonic orthogonality gives the exact identity
\begin{align}
 F_{L,N}(c,w)
   &=\tfrac12\|w\|_V^2+
                \tfrac12\langle c,S^{(m)}c\rangle
                                +\mathcal P_{L,N}(c,w),\notag\\
 \mathcal P_{L,N}(c,w)
   &=\tfrac12\{\langle q,S^{(m)}q\rangle
                            -\langle c,S^{(m)}c\rangle\}
                          +\mathcal N(\mathcal H_{L,N}q+w).
                                                        \label{eq:v44-full-chart-action}
\end{align}
All pairings have their bilinear extensions for
complex fields. The first term in the second line
\emph{must be kept}. Unlike the one-cell constant
source case, a multicell source generally has
nonzero quadratic curl energy, and its normal
coordinate \(q(c,w)\) moves with the detail.

\begin{theorem}[Uniform local nonlinear classical minimum at all depths and volumes]
\label{thm:v44-minimum}
There are \(\delta,R,C>0\), independent of \(L,N\),
such that for \(\|c\|_{\mathcal Y_N}<\delta\) the
exact chart action has a unique stationary point
\(w^*_{L,N}(c)\) in \(\|w\|_V<R\). It is the strict
minimum on that real detail ball, and
\begin{align}
 \|w^*_{L,N}(c)\|_V
       &\le C\|c\|_{\mathcal Y_N}^2,&
 \|A^*_{L,N}(c)\|_{\mathcal E}
       &\le C\|c\|_{\mathcal Y_N},\notag\\
 \tfrac12\|u\|_V^2
       &\le D_w^2F_{L,N}(c,w)[u,u]
                                  \le\tfrac32\|u\|_V^2 .
                                                        \label{eq:v44-nonlinear-floor}
\end{align}
The Hessian estimate holds throughout smaller common
real source and detail balls, not only at the
stationary point. The branch and its classical value
\(\mathcal G_{L,N}(c)=\mathcal S(A^*_{L,N}(c))\)
are holomorphic on smaller common complex balls.
Every fixed source derivative is uniformly bounded
in the stated energy and source norms. In particular,
\begin{equation}\label{eq:v44-classical-value}
 \mathcal G_{L,N}(c)
       =\tfrac12\langle c,S^{(m)}c\rangle
                                      +O(\|c\|_{\mathcal Y_N}^3).
\end{equation}
\end{theorem}
\begin{proof}
Put \(r_1=\|c\|_{\mathcal Y_N}+\|w\|_V\).
The normal bound and the uniform quadratic-form
bound for \(S^{(m)}\) give
\[
 |\langle q,S^{(m)}q\rangle-\langle c,S^{(m)}c\rangle|
              \le C\|q-c\|_{\mathcal Y_N}
                          (\|q\|_{\mathcal Y_N}+\|c\|_{\mathcal Y_N})
              \le Cr_1^3.
\]
The complete Wilson bound gives the same cubic
bound for the other term in \(\mathcal P\).
Thus \(\mathcal P\) is a uniformly bounded
holomorphic cubic remainder on a common ball.
Its convergent homogeneous expansion, or Cauchy
estimates after scaling the variables, yields
\begin{equation}\label{eq:v44-Pbounds}
 |\mathcal P|\le Cr_1^3,\qquad
 \|D_w\mathcal P\|_{V^*}\le Cr_1^2,\qquad
 \|D_w^2\mathcal P\|_{V\to V^*}\le Cr_1.
\end{equation}
This argument differentiates the moving \(q\) in
both terms of \eqref{eq:v44-full-chart-action}.
It does not omit its second derivatives.

The energy Riesz map \(\mathcal J:V\to V^*\)
and its inverse have norm one. Stationarity is
exactly
\[
                       w=-\mathcal J^{-1}
                                      D_w\mathcal P_{L,N}(c,w).
\]
Choose \(R\) and then \(\delta\) using only the
common constants in \eqref{eq:v44-Pbounds}.
The right side preserves the \(R\)-ball and
contracts by at most \(1/2\). It also preserves
a ball of radius \(C'\|c\|_{\mathcal Y_N}^2\)
for sufficiently large fixed \(C'\), proving the
sharper stationary bound. The second derivative
is \(\mathcal J+D_w^2\mathcal P\); reducing the
balls so its perturbation norm is at most \(1/2\)
proves \eqref{eq:v44-nonlinear-floor}. Convexity
on every line segment proves the minimum statement.
Complex contraction gives the holomorphic branch,
and common-ball Cauchy estimates give all fixed
source derivatives. Substitution proves
\eqref{eq:v44-classical-value}.
\end{proof}

One elementary calibration makes the normal-action
term visible. For the declared scalar constraint
\[
             O(q,w)=q+\alpha qw+\gamma w^2=c
\]
and action \(\tfrac12w^2+\tfrac{\sigma}{2}q^2\),
the exact normal solve is
\(q=(c-\gamma w^2)/(1+\alpha w)\). At \(w=0\)
the detail force is \(-\sigma\alpha c^2\) and
the detail Hessian is
\[
                         1-2\sigma\gamma c
                                   +3\sigma\alpha^2c^2 .
\]
Freezing \(q=c\) would give zero force and Hessian
one. This is a declared finite check of
\eqref{eq:v44-full-chart-action}, not a replacement
of the actual root or a Yang--Mills coefficient.

\subsection{Local gauge transversality and the relation to V43 histories}

The previous statements do not rely on a global
compact Coulomb slice. Their local gauge interpretation
can nevertheless be checked on the whole small
critical neighbourhood, rather than only at \(A=0\).

\begin{lemma}[Uniform infinitesimal gauge transversality in energy]
\label{lem:v44-FP}
On any finite four-torus, let \(A\) be a sufficiently
small link field in counting \(\ell^4\). The derivative
at \(\psi=0\) of
\[
             \psi\longmapsto
                D^*\log(e^{-\psi_x}e^{A_{x,\mu}}
                                         e^{\psi_{x+e_\mu}})
\]
is invertible from mean-zero site \(\mathcal G\)
to its energy dual, with inverse norm at most two
after comparison with the energy Riesz map. The
smallness threshold is independent of the torus.
\end{lemma}
\begin{proof}
The link derivative is \(D\psi+\mathcal L_A\psi\).
Analytic differentiation of the local gauge word
gives
\[
 |(\mathcal L_A\psi)_e|
                   \le C|A_e|(|\psi_x|+|\psi_{x+e_\mu}|).
\]
For mean-zero site fields \(\zeta,\psi\), Holder
and the lattice Sobolev bound give
\[
 |\langle D\zeta,\mathcal L_A\psi\rangle|
       \le C\|A\|_4\|D\zeta\|_2\|D\psi\|_2.
\]
Consequently the derivative, compared with
\(\Delta=D^*D:\mathcal G\to\mathcal G^*\), is
an operator perturbation of norm at most
\(C\|A\|_4\). A Neumann series proves the
claim. This uses no counting-\(\ell^2\) lower
bound on the Laplacian.
\end{proof}

For each finite regulator, the ordinary local
implicit function theorem now identifies the
logarithmic Coulomb manifold as a local gauge
slice modulo constant colour rotations near
the fields used above. No global orbit uniqueness
or uniform chart for every gauge group element
is inferred. The source retraction is equivariant
under those constant rotations. Thus a nearby
variation in the literal fixed-root fibre can
first be gauge-represented on the fine Coulomb
slice; its coarse Coulomb source differs from
the prescribed one at most by a residual
constant conjugation, which can be compensated
by a fine constant conjugation. Exact root
covariance and gauge invariance of \(\mathcal S\)
preserve its action. This supplies the local
root-orbit minimum interpretation of
Theorem~\ref{thm:v44-minimum}. It is not a
uniqueness theorem among all compact wells.

At zero final source, let \(D_m=\operatorname{diag}(H_1,\ldots,H_m)\)
be V43's innovation precision. Reconstruct a fine
linear field from those innovations through the
actual raw history, and then apply \(\Pi_f\);
denote this linear map by \(\mathcal I_m\).
It maps into \(\mathcal V_{L,N}\), since the
literal linear root differs from \(B_L\) by a
coarse gradient, and fine gauge changes under
\(B_L\) also become coarse gradients. The
energy identity of V43 gives
\begin{equation}\label{eq:v44-history-isometry}
       \|\mathcal I_m\xi\|_V^2=\langle\xi,D_m\xi\rangle,
             \qquad \kappa I\preceq D_m\preceq MI .
\end{equation}
This proves injectivity. Both real spaces have
dimension \(9N^4(L^4-1)\): subtract the coarse
Coulomb dimension \(9N^4+3\) from the fine
dimension \(9(NL)^4+3\), or count V43's independent
history groups. Hence the map is onto.

\begin{proposition}[A nonlinear classical chart in the harmonic history norm]
\label{prop:v44-history-chart}
The exact chart
\[
                 (c,\xi)\longmapsto
                          A_{L,N}(c,\mathcal I_m\xi)
\]
has a common domain in the source and unweighted
innovation norms, for every \(m,N\). On smaller
real balls its classical detail Hessian satisfies
\begin{equation}\label{eq:v44-history-floor}
 \frac{\kappa}{2}\|u\|_2^2
   \le D_\xi^2F_{L,N}(c,\mathcal I_m\xi)[u,u]
                              \le\frac{3M}{2}\|u\|_2^2 .
\end{equation}
Its stationary innovation has norm
\(O(\|c\|_{\mathcal Y_N}^2)\), uniformly.
\end{proposition}
\begin{proof}
The upper bound in \eqref{eq:v44-history-isometry}
puts a fixed innovation ball in the uniform \(V\)
ball of the chart. Its inverse has norm at most
\(\kappa^{-1/2}\) from \(V\) to innovations.
Apply Theorem~\ref{thm:v44-minimum} and the linear
chain rule, then use the two bounds in
\eqref{eq:v44-history-isometry}. The stationary
bound follows from the inverse norm.
\end{proof}

This is not an assertion that V43's nonlinear
successive conditional translations themselves
obey \eqref{eq:v44-history-floor}. The new chart
uses their \emph{linear harmonic} coordinates on
the fine Coulomb detail, and then solves the
full source-normal equation. It agrees with
that innovation description at the tangent
level but is a different nonlinear parametrization.
In particular, V43's unit-Jacobian assertion
cannot be transferred to it.

\subsection{The measure and probability boundary is still essential}

All the new coercivity statements concern the
classical action on a small global critical
ball. Haar densities, the original pivot
Jacobians, the new source-normal Jacobian,
and the local gauge-slice density would have
to be transported together before using this
chart in a measured integral. A gauge-orbit
equality of actions alone is not such a
density computation. Likewise V43's exact
cutoff domain must be pulled back, not
replaced by a new product of innovation balls.

The rarity of this global domain is already
quantified in V23, including the discussion
after Theorem~\ref{thm:v23-smallball}: under
the free reference,
\(\beta\|w\|_V^2\) has a chi-square law with
the number of detail degrees of freedom.
The present multicell dimension is
\(9N^4(L^4-1)\). Thus a common radius for the
classical minimum is not a probability bound
covering typical ultraviolet fluctuations.
That obstruction is retained, not claimed
anew or removed by \eqref{eq:v44-history-floor}.

The new sources emphasize the same separation:
Native V21's quadratic uncertainty floor is
not removed by centering means; SM V17's
structural algebra does not manufacture a
covariant physical operation; Bridge V16's
positive compression does not replace its
full constrained response. None authorizes
identifying the present elliptic detail
floor with a physical Hamiltonian mass gap.

\subsection{Certificates, source fingerprints, and the next acquisition}

The verification script is
\begin{center}\small
\path{scripts/verify_ym_f14_v44_coulomb_source.py}.
\end{center}
It has 6,960 bytes and SHA--256
\begin{center}\scriptsize\ttfamily
BAE5A3B6244AF06BD8247E4AE26D043828670789A959E29ADC4B3FDC1F8F6204.
\end{center}
Exact rational arithmetic constructs the complete
coarse \(N=2\) gradient, mean-zero Green matrix,
and Coulomb projector. It checks both divergence
equations in \eqref{eq:v44-gauge-jets}, all
\(64\) quaternion gauge-word expansions, the
nonzero quadratic correction and its norm,
and pure-gauge cancellation. It also verifies
the geometric-sum second moment and the
moving-source scalar action calibration.

Separate floating-point regressions use all
aliases at \(L=2,4,8,16\), with generic,
near-zero, and near-face four-dimensional
momenta. They compare the direct longitudinal
row with \eqref{eq:v44-long-cancellation}
and test the displayed principal and
nonprincipal envelopes at those samples.
They are not interval certificates over all
momenta. The uniform theorem is proved by
the explicit analytic envelope and Parseval,
not extrapolated from the samples.

The newly supplied Downloads files have the
following unchanged fingerprints, in the
order listed at the start of V44:
\begin{center}\scriptsize
\begin{tabular}{ll}
SM source selection V17 & 742,509 bytes\\
\multicolumn{2}{l}{\ttfamily A3AB81EED2BAC59CE4E7B3F97EAFA36B9B5F7090589D60FAAAC57C6AF05CE6D5}\\[2pt]
Native stationarity V21 & 802,207 bytes\\
\multicolumn{2}{l}{\ttfamily C2D452CA7A031E2AF8B319073A79FE820C37007BDFC5442B8158916539E1A580}\\[2pt]
Exact-action bridge V16 & 584,434 bytes\\
\multicolumn{2}{l}{\ttfamily 53C50B2CFBA4A05B8749CF165EF1ADEEC8B7B1133FA9FB6AD6D45C8D0F5EB49F}
\end{tabular}
\end{center}
The proof inputs taken from Shell V9--V12 were
re-read in the current Shell V16 file. The
related V34 gauge correction and V23 global-ball
obstruction in this lineage were checked to
avoid presenting them as new results. The
new mathematical estimates are the longitudinal
bound, its exact uniform mixed-domain Coulomb
correction, and the resulting multicell
nonlinear actual-source chart and minimum.

V43 has 1,024,036 bytes and SHA--256
\begin{center}\scriptsize\ttfamily
09710413F63905AF94FED1B4CFFAD6FD59B400B4630EC645A6984280144913BE.
\end{center}
Its body is preserved except for the title
version and this terminal addition. V44
replaces V43 as the sole active main file.
Archives, the regular companions, and the
new supplied files are unchanged. V45 is
the next scheduled full companion checkpoint.

\begin{firewall}
A depth- and volume-independent nonlinear
classical chart now exists for the actual
small root orbit on a global critical-energy
domain. Its action has a uniform detail
Hessian bound in harmonic history coordinates,
with the moving source action retained.
This closes a specific deterministic
multicell gauge/source obstruction, not the
quantum programme. The next requirement is
to transport the complete measured density
and control statistically relevant local
fluctuation domains; the small global ball
does not supply those facts. Compact entry,
an interacting cofinal quantum continuum
construction, and a positive physical mass
gap remain unproved.
\end{firewall}
\section{V45: the exact multicell quotient measure, its controlled normal determinant, and an identity-step test excluding a false mass term}
\label{sec:v45}

\subsection{Companion checkpoint and the next measured question}

V44 supplies a common nonlinear classical source chart, not
the corresponding probability law. The next calculation must
transport the original measure through that chart. The fine
logarithmic Coulomb density is already derived in Shell V16
and recorded in Section~\ref{sec:v19-live}; it is not an
unspecified input and is not re-derived as a new result here.
The new issue on several coarse cells is that the
\emph{coarse reference measure} has a Coulomb gauge
determinant too. A one-cell reference formula cannot be
extended by inserting only a product-Haar logarithm.

This is the scheduled V45 companion checkpoint. The current
active regular companions remain SM--Einstein V72, NS V26,
and Shell V16, with the same byte counts and SHA--256 values
recorded at the preceding checkpoints. Their latest relevant
sections were re-read. SM V72 transports empirical covariance
moments in a many-copy scaling with a fixed physical mode
set; it does not give a spatial-refinement-uniform quantum
YM measure. NS V26 gives pointwise rank-one kernel-derivative
contractions and comparisons for its Oseen problem; no
such coefficient is substituted for a gauge-fibre bound.
Shell V16 provides the actual logarithmic quotient density,
its primed gauge convention, and the complete one-cell
flat calculation used below as a nonduplication boundary.

The three Downloads sources added during V44 were also
checked for changes. They remain SM source selection V17,
Native stationarity V21, and Exact-action bridge V16.
Their implications and limitations are the ones recorded
in V44: a structural algebra is not an actual covariant
instrument, a quadratic phase-space comparison is not a
nonlinear quantum continuum theorem, and a compression
bound is not the inverse of the full coupled operator.
No new companion result closes the remaining YM measure
or physical mass-gap problem. The next regular checkpoint
is V50.

\subsection{Both fine and coarse quotient reference densities}

Retain V44's notation: the coarse side is \(N\), the
fine side is \(NL\), \(L=2^m\), \(D\) is the vertex
gradient, and \(\mathcal C_M=\ker D_M^*\) is the
linear Coulomb link space on a torus of side \(M\).
Its real dimension, including constant links, is
\[
                         s_M=9M^4+3.
\]
For \(M=1\), the mean-zero vertex gauge space is zero.
At general \(M\), a prime on a vertex determinant
removes exactly the three constant site gauge
parameters, not constant link holonomies.

The inherited local product-Haar quotient density on
\(\mathcal C_M\), relative to its fixed Euclidean
volume, is
\begin{align}
 \rho_M(a)&=J_{{\rm H},M}(a)\,\mathfrak d_M(a),\notag\\
 J_{{\rm H},M}(a)
      &=\prod_{\text{all links }e}
                         \left(\frac{\sin|a_e|}{|a_e|}\right)^2,\notag\\
 \mathfrak d_M(a)
      &=\det{}'\!\left(\Delta_M^{-1/2}\mathcal M_a
                                      \Delta_M^{-1/2}\right),
 \qquad \mathcal M_a=D_M^*R_a,\quad
                         \Delta_M=D_M^*D_M,               \label{eq:v45-rho}\\
 (R_a\omega)_{xy}
      &=\sigma(\operatorname{ad}_{a_{xy}})(\omega_y-\omega_x)
                         +\tfrac12[a_{xy},\omega_x+\omega_y],
 \qquad \sigma(z)=\frac z2\coth\frac z2 .\notag
\end{align}
The values at zero are removable, and \(\rho_M(0)=1\)
in this normalization. Real positivity holds on the
small critical branches in use. The uniform energy
gauge bound is inherited from Shell V16, consistently
with Lemma~\ref{lem:v44-FP}; it is not a lower bound
on the unscaled smallest Laplacian eigenvalue.
Any field-independent Haar-to-slice normalization
is absorbed once into the chosen constant multiple
of Euclidean slice volume. Equivalently those
constants can be kept in the linear prefactor below;
they are not discarded from a normalized identity.

Formula \eqref{eq:v45-rho} is used at \emph{both}
\(M=NL\) and \(M=N\). In particular
\[
                         \rho_1(c)=J_{{\rm H},1}(c),
\]
because there is no primed one-site gauge determinant.
This is the precise reason that the earlier one-cell
formula contains only the coarse product-Haar factor.

All measure statements below are local on a regular
gauge-saturated branch, with compatible localizations.
Constant colour conjugations are retained. A test may
be averaged over that compact group; no division by
a singular stabilizer volume is performed. We assume
the branch and any cutoff are invariant under these
constant rotations. This is satisfied by compatible
radial localizations of the actual history. The
claims do not classify global Coulomb copies or
other compact wells.

\subsection{The exact finite disintegration in V44's chart}

Write \(\mathcal H=\mathcal H_{L,N}\),
\(\mathcal O=\mathcal O_{L,N}\), and
\(\mathcal V=\mathcal V_{L,N}\). The linear
isomorphism
\[
                  (q,w)\longmapsto\mathcal Hq+w
                   \quad\hbox{from }\mathcal C_N\oplus
                         \mathcal V\hbox{ to }\mathcal C_{NL}
\]
has a positive field-independent volume Jacobian
\(j_{L,N}^{\rm lin}\). We use fixed source volume
\(dq\) and fixed energy-orthonormal detail volume
\(dw_V\); the conversion from these volumes to
Euclidean fine-slice volume is included in that
constant, not asserted to be one.

On the exact V44 source chart define
\begin{equation}\label{eq:v45-normal}
 A(c,w)=\mathcal Hq_{L,N}(c,w)+w,\qquad
 \mathcal N_{L,N}(c,w)=D\mathcal O(A(c,w))\,\mathcal H .
\end{equation}
The source-normal matrix acts on \(\mathcal C_N\).
It is not the gauge operator \(\mathcal M_a\).
The implicit normal equation gives
\[
                  D_cq_{L,N}(c,w)=\mathcal N_{L,N}(c,w)^{-1}.
\]
Its real determinant is positive near the trivial
configuration, on the uniform domain specified below.

Let \(\chi_{L,N}\) denote an actual compatible
localization or the pullback of a declared history
cutoff, supported in the chosen branch. Dependence
on all reconstructed intermediate levels is allowed.
In particular it is not reset to a product cutoff
in \(w\).
Integrals written over \(\mathcal V\) below mean
integration on the chart domain with the localized
integrand extended by zero outside it.

\begin{theorem}[Exact local pushforward relative to coarse quotient Haar]
\label{thm:v45-measure}
The local pushed Wilson density relative to
\(\rho_N(c)\,dc\) is
\begin{equation}\label{eq:v45-measure}
 \begin{aligned}
 Z^{\rm loc}_{L,N}(c)
   =\frac{j_{L,N}^{\rm lin}}{\rho_N(c)}
       \int_{\mathcal V}
        &e^{-\beta\mathcal S(A(c,w))}
         \chi_{L,N}(c,w)\\[-2pt]
        &{}\times
        \frac{\rho_{NL}(A(c,w))}
                        {\det\mathcal N_{L,N}(c,w)}\,dw_V .
 \end{aligned}
\end{equation}
This is an exact finite-regulator identity on the
declared branch. Equivalently, it is the local
coarse-orbit density of the original rooted
pushforward, lifted invariantly to the coarse
Coulomb slice. It preserves the original action,
the complete Haar law, and the cutoff.
\end{theorem}
\begin{proof}
Begin with the inherited quotient formula
\(\rho_{NL}(A)\,dA_{\mathcal C_{NL}}\) for the
fine product-Haar law, multiplying by its exact
Wilson weight and localization. The linear
splitting contributes \(j_{L,N}^{\rm lin}\).
In \((q,w)\) coordinates the map
\[
             (q,w)\longmapsto(\mathcal O(\mathcal Hq+w),w)
\]
has a block triangular differential, with diagonal
blocks \(\mathcal N_{L,N}\) and the identity.
Its inverse therefore contributes
\((\det\mathcal N_{L,N})^{-1}\). Cross derivatives
of \(q(c,w)\) do not give a second determinant.
Integration in \(w\) gives the density relative
to \(dc\). Dividing by \(\rho_N(c)\) gives
\eqref{eq:v45-measure}.

For the orbit interpretation, the balanced root
and V44's exact Coulomb correction have the same
coarse gauge orbit as the literal rooted map.
Hence all gauge-invariant coarse tests have
identical values before integration. The fine
law and localization, and both lifted coarse
laws, are invariant under constant conjugation.
Averaging any coarse-slice test over those
rotations reduces it to the same orbit test.
This identifies the invariant local pushforwards.
This argument identifies measures; it does not
erase a representative-coordinate Jacobian.
\end{proof}

Thus the full logarithmic density in the
\((c,w)\) chart, apart from the cutoff and the
field-independent linear constant, is
\begin{equation}\label{eq:v45-complete-ell}
 \ell^{\rm quot}_{L,N}(c,w)
    =\log\rho_{NL}(A(c,w))
       -\log\det\mathcal N_{L,N}(c,w)-\log\rho_N(c).
\end{equation}
The stage pivot factors of the old tree
description must not be multiplied into this
formula again. They and the present normal
Jacobian are parts of two descriptions of
the same pushforward. Equality is at the
level of the transported measure, not a
term-by-term equality of their individual
determinants. An additional incoming density,
if independently part of an iterated law,
must still be composed with \(A(c,w)\);
it is not reset by \eqref{eq:v45-measure}.

\subsection{A uniform source-normal determinant budget}

The complete fine density is not uniformly
controlled by classical invertibility alone.
There is, however, a uniform and quantified
bound for the source-normal part just introduced.
Use the fixed Hilbert norm
\[
 \|c\|_{\mathcal Y_N,2}
       =(\|d_cc\|_2^2+N^2|\overline c|^2)^{1/2},
 \qquad
 \|c\|_{\mathcal Y_N,2}\le\|c\|_{\mathcal Y_N}
                         \le\sqrt2\|c\|_{\mathcal Y_N,2}.
\]
Changing to an orthonormal basis for this norm
conjugates the source-normal matrix by a
fixed isomorphism. Its determinant is unchanged.
Set
\[
      r(c,w)=\|c\|_{\mathcal Y_N}+\|w\|_V,\qquad
      b_{L,N}(c,w)=\log\det\mathcal N_{L,N}(c,w).
\]

\begin{theorem}[Uniform invertibility and normalized marked log determinant]
\label{thm:v45-normal}
There are common \(r_0,C,C_k>0\), independent
of \(L,N\), such that on \(r(c,w)<r_0\),
\begin{equation}\label{eq:v45-normal-op}
 \|\mathcal N_{L,N}(c,w)-I\|_{\mathcal Y_N,2\to\mathcal Y_N,2}
                     \le Cr(c,w)<\tfrac12,\qquad
                         \|\mathcal N_{L,N}^{-1}\|\le2 .
\end{equation}
The determinant is positive for real arguments,
and has a holomorphic logarithm vanishing at
zero on the common complex ball. For real data,
with \(s_N=9N^4+3\),
\begin{equation}\label{eq:v45-normal-det}
       (1-Cr)^{s_N}\le\det\mathcal N_{L,N}(c,w)
                                          \le(1+Cr)^{s_N}.
\end{equation}
For any fixed number \(k\ge2\) of source or
detail directions \(v_1,\ldots,v_k\), measured
in the joint norm \(\mathcal Y_N\oplus V\),
\begin{align}
 N^{-4}|b_{L,N}(c,w)|&\le Cr(c,w)^2,\notag\\
 N^{-4}|Db_{L,N}(c,w)[v]|&\le Cr(c,w)\|v\|,\notag\\
 N^{-4}|D^kb_{L,N}(c,w)[v_1,\ldots,v_k]|
                  &\le C_k\prod_{i=1}^k\|v_i\|.
                                                        \label{eq:v45-normal-marks}
\end{align}
These are global critical-norm and pressure
bounds. They are not anchored spatial
interaction estimates.
\end{theorem}
\begin{proof}
Theorem~\ref{thm:v44-source} gives
\(D\mathcal O(0)\mathcal H=I\) and uniform
derivatives of its nonlinear remainder.
Lemma~\ref{lem:v44-chart} gives uniform
derivatives of \(A(c,w)\), with
\(\|A(c,w)\|_{\mathcal E}\le Cr(c,w)\).
The chain rule and the uniform lift bound
therefore prove \eqref{eq:v45-normal-op},
including all fixed-order derivatives of
the source-normal matrix. The equivalence
of the two source norms costs only \(\sqrt2\).
A Neumann series gives the inverse estimate
and the trace-log series defines a common
holomorphic logarithm.

In a \(\mathcal Y_{N,2}\)-orthonormal basis,
all singular values of \(\mathcal N\) lie
between \(1-Cr\) and \(1+Cr\). The determinant
has positive sign by continuity from one,
so their product gives
\eqref{eq:v45-normal-det}.
For \(k\) marked derivatives the finite
logarithmic derivative is a sum over
ordered nonempty partitions
\((I_1,\ldots,I_j)\) of the \(k\) labels:
\[
 D_1\cdots D_k\log\det\mathcal N
   =\sum_{j=1}^k\frac{(-1)^{j-1}}j
       \sum_{(I_1,\ldots,I_j)}
       \operatorname{Tr}\bigl(
          \mathcal N^{-1}D_{I_1}\mathcal N
                   \cdots\mathcal N^{-1}D_{I_j}\mathcal N\bigr).
\]
The formula follows by differentiating
\(\operatorname{Tr}\log\mathcal N\), or
the convergent local trace-log series.
Each trace has absolute value at most
\(s_N\) times the product of its operator
norms. Since \(s_N/N^4\le12\), this proves
the last bound in \eqref{eq:v45-normal-marks}.

It remains to justify the improved orders
in the first two bounds. Every source map,
normal solve, and harmonic lift commutes
with constant colour conjugation. Consequently
\(\mathcal N\) is conjugated on its source
space by such a rotation, and \(b\) is an
invariant scalar function of \((c,w)\).
There is no invariant linear functional
on their Lie-colour components: averaging
one adjoint \(\mathfrak{su}(2)\) component
over global rotations gives zero.
Thus \(Db(0)=0\), as well as \(b(0)=0\).
Integrating the uniform second-derivative
bound along the radial segment proves
the first two bounds.
\end{proof}

In particular, the lower bound is exponential
in the \emph{coarse} source dimension, not a
dimension-free lower bound for the determinant
as a number. This distinction is necessary.
Uniform normal inverse and marked bounds do
not imply uniform control of every fine
trace in \(\log\rho_{NL}\), or of the
complete fluctuation determinant.

\subsection{The finite measured one-loop formula in the multicell chart}

Suppose the localization is one on a
neighbourhood of V44's classical minimum
\(w^*(c)\), and is supported inside its
strictly convex local chart. Let
\[
 A^*(c)=A(c,w^*(c)),\qquad
 H_V(c)=D_w^2[\mathcal S(A(c,w))]_{w=w^*(c)}
\]
in energy-orthonormal detail coordinates.
At any fixed finite regulator, the ordinary
local Laplace expansion of
\eqref{eq:v45-measure} gives, up to
source-independent constants,
\begin{equation}\label{eq:v45-one-loop}
 \begin{aligned}
 Q^{\rm quot}_{L,N}(c)={}&
       \tfrac12\log\det H_V(c)
       +\log\det\mathcal N_{L,N}(c,w^*(c))\\
       &-\log J_{{\rm H},NL}(A^*(c))
        -\log\mathfrak d_{NL}(A^*(c))\\
       &+\log J_{{\rm H},N}(c)+\log\mathfrak d_N(c).
 \end{aligned}
\end{equation}
For example, after subtracting the zero-source
value of this expression,
\[
 -\log Z^{\rm loc}_{L,N}(c)
       =\beta\mathcal G_{L,N}(c)
          +\frac{\dim V}{2}\log\frac{\beta}{2\pi}
                   +Q^{\rm quot}_{L,N}(c)+O_{L,N}(\beta^{-1})
\]
with a corresponding source-independent
normalization understood.
To verify the coefficient, expand the
action quadratically at the unique minimum,
rescale the detail by \(\beta^{-1/2}\),
and evaluate the smooth amplitude
\(\rho_{NL}/\det\mathcal N\) there.
The Gaussian integral contributes
\((\det H_V)^{-1/2}\), while the outside
reference factor is \(1/\rho_N(c)\).
Taking the negative logarithm gives
exactly \eqref{eq:v45-one-loop}.

This paragraph asserts a fixed-regulator
coefficient identity, not a uniform remainder
as \(L,N,\beta\) change together. At \(N=1\),
\(\mathfrak d_N=1\) and the identity reduces,
after the fixed coordinate normalizations,
to the inherited Shell V16 assembly.
Its new coarse determinant is essential
when \(N>1\). No new one-cell flat sum,
fifth modified determinant theorem, or
two-loop composition formula is claimed.
Those calculations remain inherited inputs.

\subsection{An exact identity-block test and its nonzero coarse gauge factor}

The zero-step case \(L=1\) is a useful
unit test of the measure, before any
asymptotic expansion. Then \(\mathcal V=0\),
\(\mathcal H=I\), and, on the coarse
Coulomb slice,
\[
        \mathcal O_{1,N}(A)=\Gamma_N(A)=A,\qquad
        \mathcal N_{1,N}=I,\qquad \rho_{NL}=\rho_N .
 \]
Thus the exact identity-step effective
density is \(e^{-\beta\mathcal S(c)}\)
relative to the same quotient Haar reference.
There is no induced potential.

If one were to use the fine Coulomb density
but divide only by \(J_{{\rm H},N}(c)\)
instead of \(\rho_N(c)\), the result would
instead be
\[
               e^{-\beta\mathcal S(c)}\mathfrak d_N(c).
\]
This is a density relative to a different
reference, not the correct physical
effective potential. Its negative logarithm
would contain the spurious term
\(-\log\mathfrak d_N(c)\). The following
actual finite computation makes the error
visible on a flat field.

\begin{proposition}[The coarse side-two gauge determinant]
\label{prop:v45-coarse-gauge}
On a coarse torus of side \(N=2\), set
\(c_1=xE_3\), \(c_2=c_3=c_4=0\) at every
site, with \([E_a,E_b]=2\epsilon_{abc}E_c\).
For sufficiently small real or complex \(x\),
\begin{equation}\label{eq:v45-coarse-gauge}
 \mathfrak d_2(c)
     =\prod_{r=0}^3
          \left(\frac{r+x\cot x}{r+1}\right)^{2\binom3r},
 \qquad
 \log\mathfrak d_2(c)
       =-\frac52x^2-\frac{175}{432}x^4+O(x^6).
\end{equation}
The full coarse quotient density has
\begin{equation}\label{eq:v45-coarse-rho}
             \log\rho_2(c)
                =-\frac{47}{6}x^2
                         -\frac{1259}{2160}x^4+O(x^6).
\end{equation}
The Wilson action is zero on this family.
Correct identity blocking consequently gives
zero effective potential, while the incomplete
reference would give
\(\frac52x^2+\frac{175}{432}x^4+O(x^6)\).
\end{proposition}
\begin{proof}
Fourier momenta on the side-two torus have
components zero or \(\pi\). The homogeneous
odd gauge term in \(D^*R_c\) is proportional
to \(\sin k_1\), so it vanishes exactly.
The neutral colour has eigenvalue
\(\lambda(k)=4\sum_\mu\mathbf1_{\{k_\mu=\pi\}}\).
For the two charged colours,
\(\operatorname{ad}_{xE_3}\) has eigenvalues
\(\pm2ix\), and
\(\sigma(\pm2ix)=x\cot x\).
If \(k_1=0\), the determinant ratio is one.
If \(k_1=\pi\) and exactly \(r\) other
components are \(\pi\), each charged ratio
is \((r+x\cot x)/(r+1)\), with multiplicity
\(\binom3r\). The primed constant mode is
absent and never appears in this last case.
This proves the product formula.

Use \(x\cot x=1-x^2/3-x^4/45+O(x^6)\) and
the exact sums
\[
 \sum_{r=0}^3\frac{\binom3r}{r+1}=\frac{15}4,\qquad
 \sum_{r=0}^3\frac{\binom3r}{(r+1)^2}=\frac{103}{48}.
\]
The quadratic logarithmic coefficient is
\(-\frac23\frac{15}4=-5/2\).
The fourth coefficient is
\[
                 -\frac2{45}\frac{15}4
                   -\frac19\frac{103}{48}=-\frac{175}{432}.
\]
There are sixteen orientation-one links,
so their product-Haar logarithm is
\[
        32\log(\sin x/x)
                  =-\frac{16}3x^2-\frac8{45}x^4+O(x^6).
\]
Adding gives \eqref{eq:v45-coarse-rho}.
All other orientations are identity links,
so every plaquette is identity and the
action is zero. The identity-step
calculation preceding the proposition
now proves the final assertion.
\end{proof}

For \(0<|x|<\pi/2\), \(0<x\cot x<1\),
so the erroneous potential is genuinely
positive, not an artefact of truncating
its Taylor series. Nevertheless it has
been manufactured by inconsistent reference
densities. It is not a mass, a confinement
estimate, or a valid quantum correction.
Keeping coarse harmonic holonomies and
the correct quotient reference is essential
even in this finite normalization test.

\subsection{What is now controlled, and what is not}

The source-normal determinant in the
exact multicell law is nonzero on a
common critical domain and has a uniform
marked pressure budget. The correct
coarse quotient reference is now explicit,
and the identity-step test detects its
omission. These are new measured inputs
to V44's classical construction.

They do not bound the complete fine
density or its renormalized local
fluctuation integral. Already
Proposition~\ref{prop:v20-Haar} exhibits
a growing actual fine-Haar mixed
coefficient on the one-cell branch;
Shell V16 also shows a low gauge trace
which survives despite convergence
of its higher modified determinant.
Those factors must be combined with
the full constrained fluctuation
determinant, not bounded separately
and then declared uniformly finite.
The one-cell flat cancellation is
inherited and does not by itself
settle inhomogeneous multicell sources.

Nor does a uniform classical Hessian
turn a global critical ball into a
typical fluctuation event. V23's
chi-square obstruction remains in
force. The exact localization
\(\chi_{L,N}\) in
\eqref{eq:v45-measure} is retained
for this reason. The next substantive
estimate must control the combined
measured response on statistically
relevant domains, with all source
derivatives and cutoff contributions.
No auxiliary determinant floor is
substituted for physical-time
transfer contraction.

\subsection{Verification and preservation}

The exact finite certificate is
\begin{center}\small
\path{scripts/verify_ym_f14_v45_quotient_measure.py}.
\end{center}
It has 4,571 bytes and SHA--256
\begin{center}\scriptsize\ttfamily
224D26F7620263B28B557015BC4E1CEAC0B76467BC7CF2FFBA26DF5538F8968A.
\end{center}
It constructs the \(64\)-link,
\(16\)-vertex coarse gradient and
mean-zero Green matrix. The odd
homogeneous gauge term is checked
to vanish as an exact matrix.
The second and fourth trace-log
coefficients are computed from
the full colour matrix and
independently from the Fourier
product. The full coarse-Haar
sum, correct identity cancellation,
and false-reference potential are
then checked rationally.
A separate two-source example
verifies the triangular normal
Jacobian, its inverse, and invariance
of the determinant under a fixed
change of source norm. None of
these finite tests certifies an
infinite-dimensional determinant
or a physical spectral gap.

The regular companion checkpoint
fingerprints remain:
\begin{center}\scriptsize
\begin{tabular}{ll}
SM--Einstein V72 & 607,372 bytes\\
\multicolumn{2}{l}{\ttfamily F44F9745D43379E1672C5550411B58E97195A4C9278592D221DF5D3F644DC1F5}\\[2pt]
NS V26 & 278,283 bytes\\
\multicolumn{2}{l}{\ttfamily 82C887F8C2C69E4F28FAD7C3DC8DE5B9099CB5A4BF431EBA1FFF3E91935978AD}\\[2pt]
Shell V16 & 623,261 bytes\\
\multicolumn{2}{l}{\ttfamily BF138E63BA826BC030D3A8B87E04FDD3B1EC712D269BB064674E96509B46E710}
\end{tabular}
\end{center}
The three newly supplied files retain
the fingerprints recorded in V44.
All these sources and all archives
are preserved.

V44 has 1,060,717 bytes and SHA--256
\begin{center}\scriptsize\ttfamily
1997961852E2E2F060999B4E978DF703F2B3F05CC6C0C04F6FED7F1BF7FA3B31.
\end{center}
Its body is preserved except for the
title version and this terminal
addition. V45 replaces V44 as the
only active main note.

\begin{firewall}
The exact multicell measure is
now written in the new classical
chart with both quotient reference
densities and its source-normal
Jacobian. The latter has uniform
global marked pressure bounds,
and an actual identity-block
countercheck rules out a spurious
flat-holonomy potential from
omitting the coarse gauge factor.
The combined ultraviolet density,
typical-domain fluctuation control,
compact entry, interacting
continuum construction, and
positive physical Hamiltonian
mass gap remain unproved.
\end{firewall}
\section{V46: the full measured law on the critical ball, a quantitative Gaussian regime, and its interacting boundary layer}
\label{sec:v46}

\subsection{Context and the distinction from the previous Gaussian diagnostic}

V44 supplies a common classical chart and V45 its exact
multicell measure. Neither statement says when the integral
is concentrated near the classical minimum. V23 already
proved a free-Gaussian obstruction to a global small-field
domain. The new question is about the \emph{normalized,
nonlinear, fully measured integral itself}, not another
test of the free field.

We answer that question on a declared global energy ball.
First, the fine quotient density has a quadratic budget
of order \((NL)^2\), sharper than counting every fine
gauge dimension. Together with the source-normal factor
this gives an explicit comparison of the actual local
law to its Gaussian reference in a sufficiently cold
joint limit. Second, if the inverse temperature grows
more slowly than the detail dimension, the same exact
law has a universal exponential boundary layer. A
source-dependent change of radius then changes the
quadratic coarse stiffness by a term of order that
dimension. The complete Haar and gauge factors are
retained in all three conclusions.

This does not repeat V30's already evaluated measured
quartic logarithm, Shell V16's flat determinant
cancellation, or V43's one-variable moving-cutoff
example. It gives a dimension-explicit probability
theorem and a boundary response for V44's actual new
multicell chart. It also specifies where that chart
cannot replace the local shell domains of V40--V42.
The V45 companion review remains current; the next
scheduled checkpoint is V50. The three supplied
related manuscripts are used with the scope
restrictions recorded in V44, not as additional
instructions or as proofs of a YM continuum limit.

\subsection{An improved budget for the full fine quotient density}

On a torus of side \(M\ge2\), let \(\Delta_M^{-1}\)
mean the inverse on mean-zero vertex functions and
zero on constants. Its scalar diagonal satisfies
\begin{equation}\label{eq:v46-green}
 g_M=M^{-4}\sum_{k\ne0}
       \left(4\sum_{\mu=1}^4\sin^2(\pi k_\mu/M)\right)^{-1}
                                                    \le C.
\end{equation}
Indeed, for centered representatives
\(\lambda_M(k)\ge16|k|^2/M^2\), and counting
four-dimensional annuli gives
\(\sum_{0<|k|\le M}|k|^{-2}\le CM^2\).
Changing the constant in the radius covers the
centered cube. This argument works for every
integer \(M\), not only dyadic \(M\).

Use counting norms on links. For the endpoint-sum
operator \(S\omega_{xy}=\omega_x+\omega_y\), put
\[
 U=D_M\Delta_M^{-1/2},\qquad
 V=S\Delta_M^{-1/2}.
\]
The operator \(U\) is an isometry from the primed
vertex space into link space. All these operators
act diagonally in colour. In the notation of
\eqref{eq:v45-rho}, the normalized gauge operator is
\begin{equation}\label{eq:v46-gauge-split}
 \begin{aligned}
 I+G_M(a)&=\Delta_M^{-1/2}\mathcal M_a\Delta_M^{-1/2},\\
 G_M(a)&=G_1(a)+G_{\mathrm e}(a),\\
 G_1(a)&=\tfrac12U^*\operatorname{ad}_a V,\qquad
 G_{\mathrm e}(a)=U^*E_aU,\\
 (E_a)_e&=\sigma(\operatorname{ad}_{a_e})-I.
 \end{aligned}
\end{equation}
The subscript \(\mathrm e\) denotes the even
function of \(a\), not a link index.

\begin{lemma}[Fine density controlled by the squared link norm]
\label{lem:v46-density}
On a common sufficiently small complex counting-\(\ell^4\)
ball, the logarithm continued from \(\rho_M(0)=1\)
satisfies
\begin{equation}\label{eq:v46-density}
 |\log\rho_M(a)|
      \le C\bigl(\|a\|_2^2+\|a\|_4^4\bigr)
      \le CM^2\|a\|_4^2 .
\end{equation}
It is real on real fields in the local branch.
The constants are independent of \(M\).
The same conclusion for \(M=1\) follows directly
from its four Haar factors.
\end{lemma}
\begin{proof}
The inherited critical-norm gauge estimate gives
\(\|G_M(a)\|\le C(\|a\|_4+\|a\|_4^2)<1/2\).
The pointwise analytic expansion of \(\sigma\)
gives
\[
 \|G_{\mathrm e}\|_1\le C\sum_e|a_e|^2,\qquad
 \|G_{\mathrm e}\|_2^2\le C\sum_e|a_e|^4.
\]
Here subscripts on operators are Schatten norms.
Both inequalities follow by factoring through
the block-diagonal multiplication \(E_a\), with
the contractions \(U,U^*\) on its sides.

For the odd part, the trace vanishes exactly:
\(\operatorname{Tr}G_1=0\). Each spatial diagonal
contraction multiplies the colour matrix
\(\operatorname{ad}_{a_e}\), whose trace is zero.
This does not assume that \(a\) is homogeneous.
Moreover,
\[
 \begin{aligned}
 \|G_1\|_2^2
 &\le\tfrac14\|\operatorname{ad}_a V\|_2^2\\
 &=\tfrac14\sum_e
       \|\operatorname{ad}_{a_e}\|_{\mathrm{HS}}^2
                    (S\Delta_M^{-1}S^*)_{ee}
 \le C\sum_e|a_e|^2 .
 \end{aligned}
\]
The last inequality uses
\((S\Delta_M^{-1}S^*)_{ee}\le4g_M\), by
positivity and translation invariance of the Green
matrix, and \eqref{eq:v46-green}.
These estimates apply to complex coefficients too,
with Hermitian norms used only for the bounds.

Expand the holomorphic trace logarithm. For \(j\ge2\),
\[
 |\operatorname{Tr}G_M^j|
                 \le\|G_M\|^{j-2}\|G_M\|_2^2.
\]
The vanished odd first trace, the trace-norm bound
on \(G_{\mathrm e}\), and the geometric series prove
\[
 |\log\mathfrak d_M(a)|
                    \le C(\|a\|_2^2+\|a\|_4^4).
\]
The analytic logarithm of each fine Haar factor
is bounded by \(C|a_e|^2\) on the same ball.
Finally, there are \(4M^4\) links, so
\(\|a\|_2^2\le2M^2\|a\|_4^2\).
Smallness of \(\|a\|_4\) absorbs its fourth power.
This proves \eqref{eq:v46-density}.
\end{proof}

The lemma is a growth bound, not the cancellation of
the fine density against the fluctuation determinant.
In particular \(M^2\) still grows with the cutoff.
The surviving low ghost trace in Shell V16 is
consistent with, and can saturate the order of,
this bound.

Return to the exact V44--V45 chart, with
\begin{equation}\label{eq:v46-dimensions}
 M=LN,\qquad
 d=\dim\mathcal V=9N^4(L^4-1),\qquad
 b=M^2+N^4,\qquad L\ge2.
\end{equation}
The letter \(d\) in this section, without a
subscript or argument, is a dimension, not the curl.
Write the complete logarithmic amplitude as
\[
 \ell(c,w)=\log\rho_M(A(c,w))
       -\log\det\mathcal N_{L,N}(c,w)-\log\rho_N(c).
\]
It is precisely \eqref{eq:v45-complete-ell};
no old inverse-pivot factors are multiplied into it.
V44 bounds \(\|A(c,w)\|_4\) by
\(C(\|c\|_{\mathcal Y_N}+\|w\|_V)\).
The coarse Sobolev estimate similarly bounds
\(\|c\|_4\) by \(C\|c\|_{\mathcal Y_N}\).
Lemma~\ref{lem:v46-density}, Theorem~\ref{thm:v45-normal},
and Cauchy estimates on a smaller common complex ball
therefore give
\begin{equation}\label{eq:v46-amplitude-marks}
 \begin{aligned}
 |\ell(c,w)|&\le Cb(\|c\|_{\mathcal Y_N}+\|w\|_V)^2,\\
 \|D\ell(c,w)\|&\le Cb(\|c\|_{\mathcal Y_N}+\|w\|_V),\\
 \|D^k\ell(c,w)\|&\le C_kb,\qquad k\ge2.
 \end{aligned}
\end{equation}
In particular \(\ell(0,0)=D\ell(0,0)=0\).
For the first derivative this also follows by
integrating the second bound from the origin;
absence of an invariant linear colour functional
ensures that no linear term has been hidden.
Operator norms use the joint critical source and
energy-detail norms. The factor \(N^2\) from the
coarse density has been absorbed into \(N^4\).

A useful exact count is
\begin{equation}\label{eq:v46-ratio}
 \frac bd=
 \frac{L^2}{9N^2(L^4-1)}+\frac1{9(L^4-1)}
                       \longrightarrow0
                \quad(L\to\infty),
\end{equation}
uniformly over \(N\ge1\). This improvement over
a bound of order \(d\) is what permits the sharp
interacting boundary-layer limit below.

\subsection{A specified exact local probability law}

Set
\[
 F(w)=\mathcal S(A(0,w)),\qquad \ell_0(w)=\ell(0,w).
 \]
All gradients and Hessians in \(w\) are taken in
the fixed energy-orthonormal coordinates of
\(\mathcal V\). By Theorem~\ref{thm:v44-minimum},
on a common ball of radius \(r_0>0\),
\begin{equation}\label{eq:v46-action}
 \begin{gathered}
 F(0)=0,\quad DF(0)=0,\quad D^2F(0)=I,\qquad
 \tfrac12 I\preceq D^2F(w)\preceq\tfrac32 I,\\
 \|\nabla F(w)-w\|\le C\|w\|_V^2.
 \end{gathered}
\end{equation}
The last estimate follows from its uniform
third derivative, not just from convexity.
For \(0<r<r_0\) and \(\beta\ge0\), define
\begin{equation}\label{eq:v46-law}
 \begin{aligned}
 I_r&=\int_{\|w\|_V<r}e^{-\beta F(w)+\ell_0(w)}\,dw_V,\\
 d\mu_{\beta,L,N,r}(w)
     &=I_r^{-1}\mathbf1_{\{\|w\|_V<r\}}
                     e^{-\beta F(w)+\ell_0(w)}\,dw_V .
 \end{aligned}
\end{equation}
This is V45's actual density with the \emph{declared}
hard energy-ball localization. The field-independent
linear volume constant cancels on normalization.
The ball is invariant under constant colour rotations
and defines a compatible local orbit cutoff.
It is not asserted to be the pullback of V40's
product of radial group balls, V43's history cutoff,
or the full completed compact integral.

The hard boundary is part of this probability law.
We do not extend \(F\) or the source chart outside
the ball, and do not erase its surface terms.

\subsection{A dimension-explicit Gaussian comparison for the complete law}

Let \(\gamma_{\beta,d}=N(0,\beta^{-1}I_d)\)
and let \(\gamma_{\beta,d,r}\) be its conditioning
on \(\|w\|_V<r\). Relative entropy is
\(D(\mu\Vert\nu)=\int\log(d\mu/d\nu)\,d\mu\);
total variation means \(\sup_E|\mu(E)-\nu(E)|\).

\begin{lemma}[Gaussian entropy inequality with the boundary retained]
\label{lem:v46-lsi}
For a positive smooth density \(f\) of integral one
against \(\gamma_{\beta,d,r}\),
\begin{equation}\label{eq:v46-lsi}
 \int f\log f\,d\gamma_{\beta,d,r}
       \le\frac1{2\beta}\int\frac{|\nabla f|^2}{f}
                                      \,d\gamma_{\beta,d,r}.
\end{equation}
No Neumann condition is imposed on \(f\).
\end{lemma}
\begin{proof}
Use the reflecting semigroup \(P_t\) for
\(\Delta-\beta w\cdot\nabla\) on the ball.
Its invariant probability is the conditioned
Gaussian. The convex-boundary gradient estimate is
\[
 |\nabla P_tf|\le e^{-\beta t}P_t|\nabla f|.
\]
This standard estimate, including initial functions
without a Neumann condition, is the case
\(K=-\beta,\ \sigma=0\) of Theorem 1.1 in
\href{https://ems.press/journals/dm/articles/8965235}
{F.-Y. Wang, \emph{Semigroup properties for the second
fundamental form}, Documenta Mathematica 15 (2010), 527--543}.
The boundary is convex and the drift derivative is
\(-\beta I\), so its hypotheses hold.

Weighted Cauchy--Schwarz and invariance give
\[
 \int\frac{|\nabla P_tf|^2}{P_tf}\,d\gamma_{\beta,d,r}
 \le e^{-2\beta t}
            \int\frac{|\nabla f|^2}{f}\,d\gamma_{\beta,d,r}.
\]
For positive times the Neumann condition of \(P_tf\)
eliminates the boundary term in entropy dissipation.
Integrating that identity from zero to infinity
proves \eqref{eq:v46-lsi}, since \(P_tf\) tends to
one. Smooth positive approximation justifies
the initial endpoint.
\end{proof}

\begin{theorem}[Quantitative comparison without discarding the amplitude]
\label{thm:v46-gaussian}
There are common constants \(C,K\) such that, for
\(\beta\ge Kb\),
\begin{equation}\label{eq:v46-entropy}
 D(\mu_{\beta,L,N,r}\Vert\gamma_{\beta,d,r})
       \le C\left(\frac{d(d+2)}{\beta}
                         +\frac{b^2d}{\beta^2}\right).
\end{equation}
Consequently
\begin{equation}\label{eq:v46-tv}
 \begin{aligned}
 \|\mu_{\beta,L,N,r}-\gamma_{\beta,d}\|_{\rm TV}
 \le{}& C\left(\frac{d(d+2)}{\beta}
                          +\frac{b^2d}{\beta^2}\right)^{1/2}\\
 &+\min\left\{1,\,
       \exp\left(-\frac{\beta r^2}{4}
                              +\frac d2\log2\right)\right\}.
 \end{aligned}
\end{equation}
For fixed \(r>0\), the sufficient joint condition
\(\beta/d^2\to\infty\) makes this total variation
distance tend to zero. No assertion of optimality
of this sufficient condition is made.
\end{theorem}
\begin{proof}
Put \(U_\beta=\beta F-\ell_0\). Equations
\eqref{eq:v46-amplitude-marks} and
\eqref{eq:v46-action} give, for \(K\) large enough,
\[
 D^2U_\beta\succeq(\beta/4)I,\qquad
 \nabla U_\beta(0)=0,\qquad
 w\cdot\nabla U_\beta(w)\ge(\beta/4)|w|^2.
\]
Integrate the divergence of
\(|w|^{2k-2}w e^{-U_\beta(w)}\) on the ball.
Its outward flux is nonnegative, not zero.
It follows that
\[
 \frac\beta4\,\mathbb E_\mu |w|^{2k}
       \le(d+2k-2)\mathbb E_\mu |w|^{2k-2}.
\]
In particular,
\begin{equation}\label{eq:v46-moments}
 \mathbb E_\mu |w|^2\le\frac{4d}{\beta},\qquad
 \mathbb E_\mu |w|^4\le\frac{16d(d+2)}{\beta^2}.
\end{equation}
These are estimates under the full nonlinear law.

For \(f=d\mu/d\gamma_{\beta,d,r}\) the normalization
constant disappears upon differentiation:
\[
 \nabla\log f=-\beta(\nabla F-w)+\nabla\ell_0,\qquad
 |\nabla\log f|\le C\beta|w|^2+Cb|w|.
\]
Use \((u+v)^2\le2u^2+2v^2\), then
\eqref{eq:v46-moments}, to bound its squared
expectation by \(C\{d(d+2)+b^2d/\beta\}\).
Lemma~\ref{lem:v46-lsi} proves
\eqref{eq:v46-entropy}.

Pinsker's inequality gives the square-root term
in \eqref{eq:v46-tv}. It follows, for example,
by the log-sum inequality on the set where the
first density exceeds the second, and the bound
of binary relative entropy below by twice the
squared difference of its probabilities.
The distance between a probability and its
conditioning on a set is the probability of
the complement. For \(\xi\sim N(0,I_d)\),
\[
 \mathbb P\{|\xi|^2\ge\beta r^2\}
       \le e^{-\beta r^2/4}\mathbb E e^{|\xi|^2/4}
       =e^{-\beta r^2/4}2^{d/2}.
\]
The triangle inequality proves the claim.
Finally \(b\le Cd\) by \eqref{eq:v46-ratio};
the stated joint condition makes both terms vanish.
\end{proof}

This is a comparison of normalized probabilities,
not a source-marked expansion of their partition
functions. It applies at zero coarse source, with
the specified cutoff, and proves nothing about
the mass of its compact complement. Its constants
are uniform in depth and volume, but its
temperature requirement depends on the full
detail dimension.

\subsection{The exact interacting boundary layer at slower temperature scaling}

Here \(\beta\ge0\) is arbitrary; in particular
we do not assume the convexity of \(U_\beta\)
used in Theorem~\ref{thm:v46-gaussian}.
Choose \(C_0\) so that
\(|\nabla\ell_0(w)|\le C_0b|w|\), and define
\begin{equation}\label{eq:v46-exponents}
 p_-=d-(\tfrac32\beta+C_0b)r^2,\qquad
 p_+=d+C_0br^2 .
\end{equation}
When needed below assume \(p_->0\).

\begin{theorem}[Two-sided radial bounds and a universal boundary law]
\label{thm:v46-boundary}
For \(0<s<1\) and \(p_->0\), the exact law
\eqref{eq:v46-law} satisfies
\begin{equation}\label{eq:v46-radial}
 s^{p_+}\le
       \mu_{\beta,L,N,r}\{|w|\le sr\}
                                      \le s^{p_-}.
\end{equation}
Its logarithmic surface pressure
\(\Xi_r=r\partial_r\log I_r\) obeys
\begin{equation}\label{eq:v46-pressure}
 p_-\le\Xi_r\le p_+.
\end{equation}
Fix \(r\in(0,r_0)\). Along any sequence
\[
 L\to\infty,\qquad N\ge1,\qquad \beta/d\to0,
\]
one has \(\Xi_r/d\to1\) and, in distribution
under the exact nonlinear measured laws,
\begin{equation}\label{eq:v46-exponential}
                -d\log(|w|/r)\ \longrightarrow\
                         \operatorname{Exp}(1).
\end{equation}
The value at \(w=0\) is immaterial, since that
point has probability zero.
\end{theorem}
\begin{proof}
For a unit energy vector \(\theta\), write
\(v(t,\theta)=e^{-\beta F(t\theta)+\ell_0(t\theta)}\).
Convexity of \(F\), its stationary origin, and
\eqref{eq:v46-action} imply
\[
 0\le\partial_t F(t\theta)\le\tfrac32t.
\]
Hence
\[
 -(\tfrac32\beta+C_0b)t
       \le\partial_t\log v(t,\theta)\le C_0bt .
\]
Let \(a(t)=t^{d-1}\int_{\mathbb S^{d-1}}v(t,\theta)\,d\theta\)
be the unnormalized radial density. Differentiation
under this finite smooth angular integral preserves
the same bounds. For \(0<t\le r\),
\[
 \frac{p_--1}{t}\le\partial_t\log a(t)
                                      \le\frac{p_+-1}{t}.
\]
Integrating between \(st\) and \(t\) gives
\[
 s^{p_+-1}a(t)\le a(st)\le s^{p_--1}a(t).
\]
Multiply by \(s\), integrate \(t\) from zero
to \(r\), and divide by \(I_r=\int_0^r a(t)\,dt\).
This proves \eqref{eq:v46-radial}.

The divergence theorem, now applied to
\(w e^{-\beta F+\ell_0}\), gives the exact identity
\begin{equation}\label{eq:v46-virial}
 \Xi_r=d-\beta\mathbb E_\mu[w\cdot\nabla F(w)]
                         +\mathbb E_\mu[w\cdot\nabla\ell_0(w)] .
\end{equation}
The same pointwise bounds prove
\eqref{eq:v46-pressure}. This step explicitly
retains the nonzero boundary flux.

Equation \eqref{eq:v46-ratio} and \(\beta/d\to0\)
imply \(p_\pm/d\to1\), uniformly in the allowed
coarse sides. For each \(u\ge0\), apply
\eqref{eq:v46-radial} to \(s=e^{-u/d}\).
The survival probability of the variable in
\eqref{eq:v46-exponential} lies between
\(e^{-up_+/d}\) and \(e^{-up_-/d}\), both tending
to \(e^{-u}\). This proves the distributional limit
and the pressure assertion.
\end{proof}

For example, the upper bound also gives
\[
 \mathbb E_\mu\frac{r-|w|}{r}
        =\int_0^1\mu\{|w|/r\le s\}\,ds
                       \le\frac1{p_-+1}.
\]
Thus this is an interacting boundary layer of
relative thickness of order \(1/d\), not
concentration at the classical minimum.
It does not describe the angular law.
All density factors enter its proof through the
actual amplitude estimate, not an assumed
Gaussian replacement.

At fixed \(N\), the illustrative scaling
\(\beta=O(\log L)\) satisfies the hypothesis.
This observation does not assert that such a
coupling trajectory has already been constructed.
It shows that V30's logarithmic coefficient cannot
justify applying a global-ball saddle expansion
in that scaling. The much colder sufficient
Gaussian regime of Theorem~\ref{thm:v46-gaussian}
does not resolve it.

\subsection{A gauge-invariant cutoff change can alter the coarse stiffness}

For small coarse source \(c\) and a variable radius
\(R\), use the same chart and write
\[
 Z(c,R)=j_{L,N}^{\rm lin}
     \int_{|w|<R}e^{-\beta\mathcal S(A(c,w))+\ell(c,w)}\,dw_V,
 \qquad W(c,R)=-\log Z(c,R).
\]
This includes the coarse quotient reference through
\(\ell\). Let \(h(c)\) be a smooth invariant
nonnegative function with \(h(0)=Dh(0)=0\), and set
\[
                  r_\tau(c)=r e^{-\tau h(c)},\qquad \tau\ge0.
\]
This is a new, explicitly source-dependent
localization, not an assertion that the original
cutoff has that form.

\begin{proposition}[Exact cutoff contribution to the second source derivative]
\label{prop:v46-cutoff}
For any two coarse directions \(u,v\),
\begin{equation}\label{eq:v46-cutoff}
 D_c^2\{W(c,r_\tau(c))-W(c,r)\}_{c=0}[u,v]
                          =\tau\Xi_r D^2h(0)[u,v].
\end{equation}
Taking \(h(c)=\mathcal S_N(e^c)\), the actual
coarse Wilson action, gives
\[
 D^2h(0)[u,v]=\langle d_cu,d_cv\rangle.
\]
Under the cofinal hypotheses of
Theorem~\ref{thm:v46-boundary}, the coefficient
of this quadratic curl form in
\eqref{eq:v46-cutoff}, divided by \(d\), tends
to \(\tau\).
\end{proposition}
\begin{proof}
For fixed finite regulator, polar coordinates
make \(Z(c,R)\) smooth for \(R>0\) inside the
chart. At zero, \(Dr_\tau=0\) and
\(D^2r_\tau=-\tau rD^2h(0)\). The second chain
rule for the difference in
\eqref{eq:v46-cutoff} therefore leaves only
the radius derivative:
\[
 (-\partial_R\log Z(0,r))(-\tau rD^2h(0))
                          =\tau\Xi_r D^2h(0).
\]
The constants in \(Z(0,r)=j_{L,N}^{\rm lin}I_r\)
do not affect its radius derivative.
The quadratic Wilson expansion proves the
displayed curl formula. Finally use
\(\Xi_r/d\to1\).
\end{proof}

Both compared densities have the same value at
zero source; normalizing their effective potentials
by that value does not remove this difference.
For \(N\ge2\), choosing a coarse transverse direction
with nonzero curl makes the quadratic effect nonzero.
For \(N=1\) the coarse Wilson action starts at fourth
order, and this particular choice of \(h\) has no
quadratic effect; no one-cell quadratic term is claimed.
The example is gauge invariant and vanishes on
linear pure gauges, so Ward identities alone
cannot reject it. It changes a permitted
curvature stiffness, not a physical mass term.
Its order \(d\) is much larger than the previously
computed logarithmic one-loop coefficient.
It belongs to the specified change of cutoff,
not to a renormalization of the unchanged full
compact model.

\subsection{Finite certificate and the revised next step}

The exact checker
\begin{center}\small
\path{scripts/verify_ym_f14_v46_radial_measure.py}
\end{center}
has 5,885 bytes and SHA--256
\begin{center}\scriptsize\ttfamily
8BE6C78AB091153035ECEC82BF7A631922876DCFEDEF098EF52AD9384CF5ED0B.
\end{center}
It constructs a transverse, mean-zero,
inhomogeneous colour field on the side-two torus.
Specifically it sums oriented plaquette boundaries:
colour \(E_1\) on \((0;1,2)\) with weight one,
colour \(E_2\) on \((0;3,4)\) with weight \(2/3\),
and colour \(E_3\) on \(((1,0,1,0);1,3)\)
with weight \(1/2\).
An oriented boundary has coefficients \(+1,+1,-1,-1\)
on its four directed links.
For this field \(a\), the exact values are
\[
 \sum_e|a_e|^2=\frac{61}{9},\qquad
 g_2=\frac{103}{768},\qquad
 \operatorname{Tr}G_1(a)=0,\qquad
 \operatorname{Tr}G_1(a)^2=\frac{2989}{5184}.
\]
The matrix \(G_1\) is nonzero; no special
homogeneous cancellation was used.
Its squared Hilbert--Schmidt upper bound from
the endpoint formula is \(1769/432\).
Writing \(ta\) for the field, the checker finds
\[
 [t^2]\log\mathfrak d_2(ta)=-\frac{13969}{10368},
 \qquad
 [t^2]\log\rho_2(ta)=-\frac{37393}{10368}.
\]
These last coefficients concern the fine
quotient density alone, not its normal-source
Jacobian or its fluctuation determinant.
Using the similar matrix
\(\Delta_2^{-1}\mathcal M\) avoids square roots
in the exact trace computations.

Twenty further symbolic identities verify the
radial divergence sign and Gaussian moving-radius
response in even dimensions \(2,4,6,8\), with
moment powers \(k=1,2,3\). A uniform-ball check
gives exactly the stiffness \(\tau d\).
The script also checks the detail dimensions
and the budget ratio. These are finite algebraic
certificates; the uniform estimates, entropy
comparison, and exponential boundary law are
proved above, not inferred from sampled dimensions.

The upstream consequence is now more precise.
The global classical chart supports a quantitative
semiclassical probability comparison in a specified
cold regime. It does not supply a statistically
interior saddle in the slower regime just proved.
Rewriting its Gaussian determinant or adjusting
its global radius cannot bridge that difference.
The next useful direction is the actual local
shell law: retain generated interactions and
incoming densities while controlling local
domain changes and compact complements in a
norm compatible with iteration. V40--V42 provide
fixed-depth local inputs; V43's history interactions
and transported cutoffs still prevent assuming
an all-depth product law.

The predecessor V45 has 1,083,853 bytes and
SHA--256
\begin{center}\scriptsize\ttfamily
B5FBCC4A364A3DA19E134CF12321D52E7AD559038095780818795411F49A0CBE.
\end{center}
Its complete body is preserved except for the
title version and this terminal insertion.
V46 replaces it as the single active main note.
No companion, archive, supplied manuscript,
or closure monograph has been edited.

\begin{firewall}
The probability comparison and interacting
boundary law concern the stated local cutoff
of the exact finite YM pushforward. They do
not control its compact complement, give
cutoff-uniform marked effective interactions,
construct a nontrivial interacting continuum
measure, establish reflection-positive physical
transfer, or prove a positive Hamiltonian
mass gap. The full Yang--Mills objective
remains unproved and active.
\end{firewall}
\section{V47: local radius deformations, exponentially small boundary pressure, and the exact curvature cost of cutoff motion}
\label{sec:v47}

\subsection{The upstream choice after the global-ball calculation}

V46 proved that a global energy-ball cutoff can produce
a stiffness correction of order the entire detail dimension.
Its next-step recommendation was to return to the actual
local shell law. This section makes that distinction
quantitative: a bounded local change of the shell radii
has exponentially small boundary pressure, with the
complete measure kept. Local curvature-dependent radii
then have an exact, exponentially small quadratic
effect on the coarse action.

This is a cutoff-shape result inside the admitted local
domain, not the missing theorem about the full compact
complement. The latter is not being renamed. The audit
also found that f12 V13 already proves global section
coercivity and a volume-dependent matched complement
bound, while f12 V28 controls fixed compact blocks
with protected exterior data. Neither is repeated.
V40's radial-wall estimates and V42's two-level local
inputs are inherited; the new steps are the exact
radius-coordinate ownership, its marked comparison,
and the induced curvature response.

We work first with the actual forty-five-group
primitive shell of V40. Thus the action includes
the solved balanced-root pivots and its amplitude
includes their inverse Haar Jacobians. No Wilson
reset or independent-noise replacement is made.
The companion checkpoint remains scheduled for V50.

\subsection{Radius changes as an exact fixed-domain change of variables}

Let \(\Lambda\) be the coarse torus and
\(\mathcal I\) its \(45|\Lambda|\) independent group
indices. Write \(J=|\mathcal I|\), and retain
\(f(C,\zeta)\), \(\ell(C,\zeta)\), and the
endpoint-covariant frame \(Y_i=B_i(C)e^{\zeta_i}\)
of \eqref{eq:v40-frame}--\eqref{eq:v40-measure}.
All densities in this section are relative to
coarse group Haar, as in V40, not a new coarse
Coulomb reference.

Fix \(\theta_0>0\) and a small \(r>0\), so that
all radii \(re^{t_i}\), \(|t_i|\le\theta_0\),
remain within the common local chart. Decrease the
coarse chart and increase \(\beta_0\), if necessary,
so that every admitted classical and measured
minimum lies inside the ball of radius
\(re^{-\theta_0}/4\) in each physical coordinate.
The activation-uniform minimum argument in
Lemma~\ref{lem:v40-activation} permits this common
choice.

Set \((D_t\eta)_i=e^{t_i}\eta_i\). For independent
radius parameters \(t=(t_i)\), the integral is
\begin{equation}\label{eq:v47-radius-integral}
 \begin{aligned}
 Z_\beta(C,t)
 &=\int_{\prod_i\{|\zeta_i|<re^{t_i}\}}
                         e^{-\beta f(C,\zeta)+\ell(C,\zeta)}\,d\zeta\\
 &=\int_{\mathcal D_r}
       \exp\left\{-\beta f(C,D_t\eta)+\ell(C,D_t\eta)
                                           +3\sum_i t_i\right\}\,d\eta,
 \qquad W_\beta(C,t)=-\log Z_\beta(C,t).
 \end{aligned}
\end{equation}
The factor \(3\) is the real dimension of
\(\mathfrak{su}(2)\). The fine exponential Haar
factors are already inside \(\ell(C,D_t\eta)\);
the additional term \(3\sum t_i\) is the ordinary
linear dilation Jacobian and is also required.

If \(t=t(C)\), the joint change
\((C,\eta)\mapsto(C,D_{t(C)}\eta)\) is block
triangular. Its off-diagonal source derivatives
introduce no further determinant. Consequently
\eqref{eq:v47-radius-integral} still holds for
source-dependent radii, including the displayed
Jacobian.

For activation \(s\in[0,1]^{\mathcal I}\), use
V40's physical potentials
\[
 f_s(C,\zeta)=f(C,D_s\zeta)
       +\frac{h_0}{2}\sum_i(1-s_i^2)|\zeta_i|^2,\qquad
 \ell_s(C,\zeta)=\ell(C,D_s\zeta).
\]
Their correctly rescaled versions on \(\mathcal D_r\) are
\begin{equation}\label{eq:v47-activated}
 \begin{aligned}
 f_{s,t}(C,\eta)
     &=f(C,D_sD_t\eta)
            +\frac{h_0}{2}\sum_i(1-s_i^2)e^{2t_i}|\eta_i|^2,\\
 \ell_{s,t}(C,\eta)&=\ell(C,D_sD_t\eta)+3\sum_i t_i .
 \end{aligned}
\end{equation}
In particular the inactive Gaussian pin must also
be dilated. These are changes of coordinates
for the activated integrals, not new interpolations
whose measure has been normalized differently.

\subsection{The marked wall estimate survives the dilation}

Use a mark at a coarse link, an activation, or a
radius parameter, anchored at its actual coarse
cell. For \(k\ge1\), the anchored tree norm is
the supremum over the first mark, followed by
the sum over the other \(k-1\) marks weighted
by \(e^{aT}\), where \(T\) is their shortest
connecting tree length. As in V40, fixed
component multiplicities are included.

\begin{lemma}[Radius-uniform local wall comparison]
\label{lem:v47-wall}
There is a whole-space comparison integral
\(Z_{\beta,s}^{\rm ext}(C)\), chosen in the
physical \(\zeta\) coordinates, independent of
\(t\), such that
\[
 \Delta_{\beta,s}(C,t)
       =-\log Z_{\beta,s}(C,t)+\log Z_{\beta,s}^{\rm ext}(C)
\]
has, for some common \(a,\kappa>0\), \(B<\infty\),
and \(K_k<\infty\),
\begin{equation}\label{eq:v47-wall}
 J^{-1}\sup_{C,s,t}|\Delta_{\beta,s}(C,t)|
                          \le K_0(1+\beta)^B e^{-\kappa\beta},
 \qquad
 \|\Delta_{\beta,s}\|_{a,k}
                          \le K_k(1+\beta)^B e^{-\kappa\beta},
 \quad 1\le k\le4.
\end{equation}
Here \(|t|_\infty\le\theta_0\) and
\(\beta\ge\beta_0\). Constants are independent
of the number of coarse cells and of the
activation pattern.
\end{lemma}
\begin{proof}
We apply the exponential \emph{domain-comparison}
part of Lemma~\ref{lem:v40-remainder}, not its
subsequent \(O(\beta^{-2})\) Taylor remainder.
The following verifications justify the extra
radius marks and the common comparison integral.

First, the transformed fibre Hessian is
\[
 D_\eta^2f_{s,t}
 =D_t\{D_s f_{\zeta\zeta}(C,D_sD_t\eta)D_s
                                      +h_0(I-D_s^2)\}D_t .
\]
It therefore retains the action floor
\(e^{-2\theta_0}\min(h,h_0)I\). Its upper,
range, and row-sum bounds are uniform. The
density Hessian remains a bounded perturbation.
Each \(t_i\) derivative inserts only the local
operator \(\zeta_i\cdot\partial_{\zeta_i}\)
and, for a first amplitude derivative, the
constant \(3\). Fixed repeated derivatives
have the same local supports and bounded
polynomial degrees. The interior minimum
condition is uniform by the radius choice.
The marked inverse-Hessian and minimum
estimates of V40 thus apply to these marks.

Second, choose the centered convex extension
used in that proof for \(f_s,\ell_s\) in the
physical \(\zeta\) variables, uniformly in
\(s\). For each \(t\), use its exact pullback
under \(D_t\), including \(3\sum t_i\).
Do not construct an unrelated extension after
each dilation. The transformed whole-space
integral then equals \(Z_{\beta,s}^{\rm ext}(C)\)
by the invertible linear change of variables.
All its radius derivatives vanish exactly.
The extension and actual potentials agree in
a common neighborhood of the moving minimum.
Their discrepancy has the same rare local
supports as in V40, uniformly in \(t\).

Third, in the \(\eta\) coordinates the walls
remain the fixed functions
\(\psi_i(\eta)=\tfrac12(|\eta_i|-r)_+^2\).
None of the new parameter derivatives acts
on a wall or creates a distribution supported
on a moving sphere. Its normal Hessian
projection and the extra wall-sensitive
resolvent factor are unchanged. Spatial
weight conjugations commute with \(D_t\)
and with this normal projection; the
off-diagonal finite-range perturbation
has only changed by fixed bounded factors.
The inward margin used for the marginal
surface-density estimate is uniform.

Consequently the layer estimates
\eqref{eq:v40-radial-layer} and their
five-entry connected estimate have the
same integrable strength bound
\((1+\lambda)^{-21/20}\), up to fixed
polynomials in \(\beta\). At most four
parameter marks and one wall entry are
needed here. The local insertions from
radius derivatives fit the same finite
moment bounds. Integrate in wall strength
and perform the actual/extended same-domain
interpolation exactly as in V40. The
pointwise longest-edge bounds, with some
decay retained for the marked tree, give
the uniform anchored estimates in
\eqref{eq:v47-wall}. Summing the unmarked
wall and extension errors over their \(J\)
owners gives the pressure bound.

This proof uses the already established
local wall comparison after checking its
changed inputs. It does not assume that
the compact exterior has small probability,
or that an extensive unmarked error is
bounded independently of volume.
\end{proof}

\begin{theorem}[Exact local surface pressure]
\label{thm:v47-pressure}
For each group and every admitted \(C,s,t\),
put
\[
 \Pi_i(C,s,t)=\partial_{t_i}\log Z_{\beta,s}(C,t).
\]
It is positive and has the exact physical
boundary representation
\begin{equation}\label{eq:v47-pressure}
 \Pi_i=
 \frac{re^{t_i}}{Z_{\beta,s}(C,t)}
 \int_{\prod_{j\ne i}\{|\zeta_j|<re^{t_j}\}}
 \int_{|\zeta_i|=re^{t_i}}
             e^{-\beta f_s(C,\zeta)+\ell_s(C,\zeta)}
                                    \,d\sigma_i\,d\zeta_{\ne i}.
\end{equation}
Uniformly in volume,
\begin{equation}\label{eq:v47-pressure-bound}
                  0<\Pi_i\le\epsilon_\beta,\qquad
        \epsilon_\beta=K(1+\beta)^B e^{-\kappa\beta}.
\end{equation}
Up to three further parameter marks, with \(i\)
as one anchored terminal, have the corresponding
exponentially small tree bounds.
\end{theorem}
\begin{proof}
Differentiate the single physical ball radius
in the first line of \eqref{eq:v47-radius-integral}.
Its positive integrand is smooth to the
boundary and its sphere has positive area;
this proves the formula and strict positivity.
By Lemma~\ref{lem:v47-wall},
\(\partial_{t_i}[-\log Z_{\beta,s}^{\rm ext}]=0\),
so \(\Pi_i=-\partial_{t_i}\Delta_{\beta,s}\).
The first marked estimate proves
\eqref{eq:v47-pressure-bound}. The remaining
claims use its two-, three-, and four-mark
estimates with the radius mark retained.
\end{proof}

For example, changing independent constant
radii from \(t\) to \(t'\) along their straight
segment costs at most
\begin{equation}\label{eq:v47-pressure-change}
 |W_{\beta,s}(C,t')-W_{\beta,s}(C,t)|
                         \le\epsilon_\beta\sum_i|t_i'-t_i|.
\end{equation}
This is small per local group, not necessarily
small as a total partition-function ratio.
No probability that all local coordinates
simultaneously stay in their balls is inferred
uniformly over arbitrarily large volumes.

\subsection{Local source-dependent cutoffs and their interaction norm}

Let \(t_i(C)\) be bounded gauge-invariant
scalar functions, each carried by a fixed
coarse-link neighborhood of group \(i\).
Assume \(|t_i|\le\theta_0\), uniform local
derivatives through order four, and bounded
carrier overlap. These properties are
understood on the small chart and its
pure-gauge orbit germs. Define the exact
cutoff change
\[
 \mathcal T_s(C)=
           W_{\beta,s}(C,t(C))-W_{\beta,s}(C,0).
\]
The full-space parts in
Lemma~\ref{lem:v47-wall} cancel, even when
the radii depend on \(C\). Thus the chain
rule, finite carrier overlap, and
\eqref{eq:v47-wall} give
\begin{equation}\label{eq:v47-shape-marks}
 \|\mathcal T_s\|_{a',k}\le C_k\epsilon_\beta,
                      \quad 1\le k\le4,\qquad
 J^{-1}\sup|\mathcal T_s|\le C_0\epsilon_\beta
\end{equation}
for a possibly smaller \(a'>0\).
Marks here are coarse or activation marks.
Every term in the chain rule has at most
four marks on \(\Delta\); no fifth derivative
has been extracted from the fourth-order
input.

\begin{proposition}[Gauge-local quadratic stability under radius changes]
\label{prop:v47-shape}
The source-normalized difference
\(\mathcal T_{\boldsymbol1}(C)-\mathcal T_{\boldsymbol1}(I)\)
has a decomposition into local gauge-invariant
germs with exponentially decaying carriers and
Hessian stock \(O(\epsilon_\beta)\).
For every mean-zero Coulomb coarse direction
\(A\) on a torus of side \(N\ge2\),
\begin{equation}\label{eq:v47-shape-curl}
 |D^2\mathcal T_{\boldsymbol1}(I)[A,A]|
                      \le C\epsilon_\beta\|d_cA\|_2^2.
\end{equation}
The non-winding part of the decomposition obeys
the same estimate without the Coulomb or
mean-zero conditions.
\end{proposition}
\begin{proof}
At zero activation the physical integral
factorizes into inactive Gaussian balls
times a common source factor. Hence
\[
 \mathcal T_{\boldsymbol0}(C)
     =-\sum_i\log
       \frac{I_{\beta h_0,3}(re^{t_i(C)})}
            {I_{\beta h_0,3}(r)},\qquad
 I_{a,3}(R)=\int_{|z|<R}e^{-a|z|^2/2}\,dz .
\]
It is a sum of local invariant terms. Its
marked estimates follow either directly
or from \eqref{eq:v47-shape-marks}.
This source-dependent inactive normalization
must not be replaced by a constant.

Subtract this all-inactive expression and
the value at \(I\), and apply V40's
ordered-activation and increasing-radius
telescope to the resulting set function.
At zero activation that function is zero.
At binary activation, inactive Gaussian
factors cancel from each increment. Only
physical stencils meeting the active set,
and the corresponding local radius
carriers, remain. This gives the same
finite-support property as
\eqref{eq:v40-increments}, with a larger
fixed carrier allowance.

For an increment at group \(i\) changed
by a group \(j\) at distance \(R\), two
activation marks and two source marks
in \eqref{eq:v47-shape-marks} give a
pointwise bound \(C\epsilon_\beta e^{-a'T}\).
Keep half the decay for \(R\), and sum
the remaining decay over the other
terminals, as in the proof of
Theorem~\ref{thm:v40-decomposition}.
The resulting carrier Hessian norm is
\(C\epsilon_\beta e^{-a''R}\).
Gauge invariance follows from the
endpoint frame, invariant radii, and
radial pins in the exact integrals.
No property of the auxiliary extension
is needed for this invariance.

Apply the local orbit Hessian theorem
of V39 to every non-winding carrier.
Its polynomial filling and incidence
costs are summable against the exponential
radius decay. The remaining winding
tail has bound
\(C\epsilon_\beta N^4e^{-cN}\|A\|_2^2\).
The mean-zero Coulomb Poincare inequality
absorbs it into a uniform multiple of
\(\epsilon_\beta\|d_cA\|_2^2\).
The finite telescope proves
\eqref{eq:v47-shape-curl}.
\end{proof}

Thus these local shape changes decay faster
than every power of \(\beta^{-1}\), in this
quadratic interaction norm. This conclusion is about
admitted local cutoffs, not about changing
the compact block map or its incoming law.
It also does not control marked norms of
arbitrary order or an unbounded number
of successive shells.

\subsection{An exact curvature-dependent family}

There is a more explicit conclusion for radii
which agree at the trivial source. Let
\[
               t_i(C)=-\alpha h_i(C),\qquad \alpha\ge0,
\]
where each \(h_i\ge0\) is local and gauge
invariant, \(h_i(I)=Dh_i(I)=0\).
For fixed \(\alpha\), decrease the source
chart if needed to keep \(|t_i|\le\theta_0\).
Use \(s=\boldsymbol1\) and define
\(\Pi_i^0=\Pi_i(I,\boldsymbol1,0)\).

\begin{theorem}[The exact quadratic shape response]
\label{thm:v47-curvature}
The change in the source Hessian is exactly
\begin{equation}\label{eq:v47-curvature}
 D^2\mathcal T_{\boldsymbol1}(I)[u,v]
                =\alpha\sum_i\Pi_i^0 D^2h_i(I)[u,v].
\end{equation}
In particular, assign to every one of the
forty-five group indices in cell \(x\)
\begin{equation}\label{eq:v47-curvature-owner}
        h_i(C)=\frac1{45}\sum_{\mu<\nu}s_{x,\mu\nu}(C),
               \qquad \sum_i h_i(C)=\mathcal S_N(C),
\end{equation}
where \(s_{x,\mu\nu}\) are the actual coarse
Wilson plaquette terms. Then for every coarse
direction \(A\), without a mean-zero restriction,
\begin{equation}\label{eq:v47-curvature-bound}
 0\le D^2\mathcal T_{\boldsymbol1}(I)[A,A]
                     \le\alpha\epsilon_\beta\|d_cA\|_2^2.
\end{equation}
For the translation-invariant periodic construction,
the exact coefficient is
\begin{equation}\label{eq:v47-average-pressure}
 D^2\mathcal T_{\boldsymbol1}(I)[u,v]
       =\alpha\overline\Pi_\beta
                              \langle d_cu,d_cv\rangle,\qquad
 \overline\Pi_\beta=\frac1{45}\sum_{\text{group types }a}\Pi_a^0,
       \quad 0<\overline\Pi_\beta\le\epsilon_\beta .
\end{equation}
\end{theorem}
\begin{proof}
At \(I\), \(t(I)=0\), \(Dt(I)=0\), and
\(D^2t_i(I)=-\alpha D^2h_i(I)\).
In the second derivative of
\(W_\beta(C,t(C))-W_\beta(C,0)\), all pure
source derivatives cancel, and every mixed
term containing \(Dt\) vanishes. Only
\[
 \sum_i\partial_{t_i}W_\beta(I,0)D^2t_i(I)
\]
remains. Since \(\partial_{t_i}W=-\Pi_i^0\),
this is \eqref{eq:v47-curvature}.

The plaquette Hessians at \(I\) are
nonnegative and sum to the quadratic curl
form. Equation \eqref{eq:v47-curvature-owner}
is an exact partition of their ownership.
Use \(0<\Pi_i^0\le\epsilon_\beta\) to prove
\eqref{eq:v47-curvature-bound}.
Coarse translations preserve the measure,
the parity-tree convention within each cell,
and the equal base radii. They make
\(\Pi_i^0\) depend only on the group type,
not on \(x\). Summing the forty-five
identical curvature shares in each cell
then proves \eqref{eq:v47-average-pressure}.
\end{proof}

For \(N\ge2\) and a direction of nonzero curl,
the last coefficient is strictly positive
when \(\alpha>0\), but exponentially small
as \(\beta\to\infty\). At \(N=1\), this
particular Wilson curvature family begins
at fourth order; no quadratic one-cell
effect is claimed. The absence of winding
loss in \eqref{eq:v47-curvature-bound} comes
from its exact local plaquette formula,
not from extending the general harmonic
conclusion of Proposition~\ref{prop:v47-shape}.

Compare this with V46's global radius
\(r e^{-\alpha\mathcal S_N(C)}\): that
change had coefficient \(\alpha\Xi_r\),
with \(\Xi_r\) of order the full detail
dimension in the stated cofinal regime.
Here curvature is apportioned to fixed
local radius owners, and the complete
law gives \(\overline\Pi_\beta\) exponentially
small. A global ball is not interchangeable
with this local construction.

\subsection{The dilation Jacobian prevents a false marginal correction}

If the term \(3\sum_i t_i\) is incorrectly
omitted from the fixed-domain integral,
its value is exactly
\[
        Z_\beta^{\rm wrong}(C,t)
               =e^{-3\sum_i t_i}Z_\beta(C,t),\qquad
        W_\beta^{\rm wrong}=W_\beta+3\sum_i t_i .
\]
For the curvature ownership
\eqref{eq:v47-curvature-owner}, this makes
the error in the second source derivative
\begin{equation}\label{eq:v47-false-stiffness}
 D^2\{W_\beta^{\rm wrong}(C,t(C))
                      -W_\beta(C,t(C))\}_{I}[u,v]
                         =-3\alpha\langle d_cu,d_cv\rangle.
\end{equation}
It survives as \(\beta\to\infty\), unlike
the genuine boundary correction. At the
trivial source the two normalizers agree,
so subtracting their zero-source values
does not repair the error.

This is an error of coordinate measure,
not a new physical running-coupling term.
It preserves the linear gauge Ward
identity and can therefore masquerade
as an allowed curvature correction if
the Jacobian is not checked. No mass
gap can be inferred from its sign.

The same ownership check is visible in
the whole-space stationary coefficient.
At its moving minimum the detail Hessian
changes from \(H\) to \(D_tHD_t\);
half its logarithmic determinant increases
by \(3\sum t_i\). The logarithmic amplitude
increases by exactly the same amount.
They cancel. More generally, the exact
whole-space change of variables fixes
every normalized Laplace coefficient,
not merely this determinant.

\subsection{Verification and the remaining nonlocal obligation}

The certificate
\begin{center}\small
\path{scripts/verify_ym_f14_v47_local_radius.py}
\end{center}
has 6,159 bytes and SHA--256
\begin{center}\scriptsize\ttfamily
1AE131960FE4289C98D19DA7785C00E09A95FA34582DCAB7814B15DF129F1AD3.
\end{center}
It first checks the actual radial \(\SU(2)\)
Haar dilation
\[
 e^{3t}\left(\frac{\sin(e^t u)}{e^t u}\right)^2u^2\,du
                         =e^t\sin^2(e^t u)\,du,
\]
and the boundary-divergence identity for
the test weight
\(e^{-2\beta(1-\cos(e^tu))}\). That one-group
weight is an algebra test, not the interacting
forty-five-group marginal. The script also
checks the sign and size of
\eqref{eq:v47-false-stiffness}.

An independent coupled Gaussian check has
three vector variables \(x,y,z\in\mathbb R^3\)
and scalar precision
\[
 H_0=\begin{pmatrix}3&-1&0\\-1&3&-1\\0&-1&2\end{pmatrix}.
\]
Its leading minors are \(3,8,13\). With
cubic insertion \(x\cdot(y\times z)\),
quartic insertion \(|x|^2|y|^2\), and
quadratic logarithmic amplitude \(|x|^2/5\),
exact Wick pairing gives respectively
\[
 \mathbb E\{x\cdot(y\times z)\}^2=\frac6{13},\qquad
 \mathbb E(|x|^2|y|^2)=\frac{294}{169},\qquad
 \mathbb E(|x|^2/5)=\frac3{13}.
\]
The normalized two-loop coefficient is
\(294/169-3/13-3/13=216/169\).
The checker verifies its invariance, and
that of its three separate contractions,
under independent positive dilations of
the three vectors with the transformed
covariance and measure. This finite
coefficient calculation does not
exponentiate an uncut cubic polynomial.

Finally, on the side-two coarse torus,
the \(96\)-plaquette, \(64\)-link curl
matrix verifies the forty-five-fold
curvature ownership exactly. Arbitrary
rational pressure weights, constant by
group type across cells, check
\eqref{eq:v47-average-pressure}. Those
weights are not numerical estimates
of the interacting YM pressures;
their positivity and uniform exponential
bound are established analytically above.

The same dilation and wall argument
applies to the joint two-level radial
integral of V42, using its verified
765-group local carrier and its two
fixed radii. It retains
\(\ell_1(X_2(D,\zeta),\eta)+\ell_2(D,\zeta)\)
from \eqref{eq:v42-joint-data}, not a
replacement one-level density. Constants
then belong to that fixed-depth problem.
V42's raw-history obstruction and V43's
surviving cross-level interactions forbid
claiming them uniform at arbitrary depth.

We have therefore removed one local
ambiguity: the permitted radius changes
have a controlled exponentially small
quadratic effect in the actual measured
shell. The next unresolved comparison
must enlarge beyond these admitted local
domains while keeping generated laws and
their spatial responses. Repeating local
radius changes, global coercivity, or
formal one-loop subtractions will not
supply that compact-complement and
all-depth control.

V46 has 1,107,534 bytes and SHA--256
\begin{center}\scriptsize\ttfamily
ED03DD081894F1C900F10103D3195CC5926B899CB0337C40F92B75FF0908029A.
\end{center}
Its body is preserved except for the
title version and this terminal addition.
V47 is the sole active main note.
All companions, archives, supplied
manuscripts and monographs are preserved.

\begin{firewall}
Local cutoff-shape stability is not
cutoff removal for the compact model.
No all-depth generated-law estimate,
nontrivial interacting continuum
construction, reflection-positive
physical transfer, or positive
Hamiltonian mass gap is proved here.
The full Yang--Mills objective remains
unproved and active.
\end{firewall}
\section{V48: an explicit strong-mixing region for the full compact completed block, with its generated interactions retained}
\label{sec:v48}

\subsection{Context: a compact endpoint, not a new local cutoff estimate}

V46--V47 distinguished global and local fluctuation
cutoffs. We now leave those restricted integrals
and work with the \emph{entire} compact completed
pushforward. The new result is an explicit
strong-coupling region for that actual law:
with the existing Hilbert--Schmidt geodesic
cutoff \(\rho=1/64\), every
\[
                         0\le\beta\le10^{-8}
\]
has a volume-uniform weighted conditional-influence
bound strictly below one. A coarse plaquette
probe also has variance greater than \(1/5\).
All solved pivots, inverse Haar factors, and
generated interactions are retained.

This does not put the weak-coupling continuum
trajectory into this region. Small \(\beta\)
here is a strong-coupling endpoint, whereas
the ultraviolet construction studied earlier
requires a different regime and a controlled
flow between them. The earlier monograph
already gives abstract terminal criteria and
a strong-electric slab example. This section
verifies a compact criterion for the specific
smooth \(\SU(2)\) root block used in these
notes, without a Haar reset of its output.
It does not assert reflection positivity of
that output or identify a physical clock.
The next companion checkpoint remains V50.

\subsection{Separate the seven boundary inputs before estimating the measure}

Work on a coarse torus of side \(N\ge2\) and
fine side \(2N\). Let \(X_x\in\SU(2)^{17}\)
be the internal non-tree variables in cell
\(x\). For each coarse edge \(e=(x,\mu)\),
let \(Y_e\in\SU(2)^7\) be its seven nonpivot
boundary groups, and let \(B_e\) be its pivot.
All Haar measures below are normalized.

The inherited sixteen chronological two-link
paths have two root products \(B_e\). The
other fourteen occur in seven pairs:
\begin{equation}\label{eq:v48-paired-products}
 P_{e,a,\ell}=L_{e,a,\ell}(X)\,Y_{e,a}\,
                                  R_{e,a,\ell}(X),
       \qquad 1\le a\le7,\quad \ell=1,2.
\end{equation}
The fixed factors involve only the internal
groups in the two endpoint cells. Indeed a
two-link path crosses exactly one boundary
link. Its starting parity in direction
\(\mu\) decides whether the internal step
comes before or after that link. The seven
nonzero transverse parities each occur
twice; the zero transverse parity gives
the two root paths. Tree dressings are
identity on this slice. In particular no
\(P_{e,a,\ell}\) contains a pivot.

For fixed \(X,Y\), solve the actual completed
equation
\[
 C_e=B_e\exp\left\{\frac1{16}
       \sum_{a=1}^7\sum_{\ell=1}^2
                    g_\rho(B_e^{-1}P_{e,a,\ell})\right\},
 \qquad
 g_\rho(U)=\chi(d(I,U)/\rho)\log U .
\]
The cutoff is the inherited smooth
nonincreasing one, supported in \(d(I,U)<\rho\).
Denote the inverse Haar Jacobian by
\(j_e(X,Y_e,C_e)>0\). It has no inputs from
other boundary families. Define
\begin{equation}\label{eq:v48-kernel}
 \begin{aligned}
 k_e(X,C_e)&=\int_{\SU(2)^7}j_e(X,Y_e,C_e)\,dY_e,\\
 d\nu_e(Y_e\mid X,C_e)&=
              k_e(X,C_e)^{-1}j_e(X,Y_e,C_e)\,dY_e .
 \end{aligned}
\end{equation}
For every fixed \(X\),
\(\int k_e(X,C_e)\,dC_e=1\), by changing
from \(C_e\) back to \(B_e\) at fixed \(Y_e\).
The \(\nu_e\)'s are therefore genuine
conditional probability measures.

We do \emph{not} assume \(k_e=1\).
After integrating \(Y\) but retaining \(X,C\),
the exact joint density at zero Wilson
coupling is \(\prod_e k_e(X,C_e)\).
Estimating these factors is different from
bounding the pointwise product of all
inverse-pivot Jacobians.

\subsection{Rare source caps control the averaged Jacobian at the working radius}

Let \(R=15\rho/8\). For fixed \(X,C_e\), call
pair \(a\) exposed if at least one of its
two products has distance at most \(R\)
from \(C_e\). Let \(M_e\) count exposed pairs.
Each product in \eqref{eq:v48-paired-products}
is Haar distributed when its own \(Y_{e,a}\)
is Haar, and the seven exposure events are
independent conditional on \(X,C_e\).
No independence between the two products
of one pair is asserted.

On the round \(\SU(2)\) sphere of radius
\(\sqrt2\), normalized Haar volume gives
\begin{equation}\label{eq:v48-cap}
 {\rm Haar}(B_R)
   =\frac1\pi\left(\theta-\frac12\sin2\theta\right)
   \le\frac{R^3}{3\pi\sqrt2},\qquad \theta=R/\sqrt2 .
\end{equation}
The bound follows by integrating
\((2/\pi)\sin^2 t\le(2/\pi)t^2\).
At \(\rho=1/64\), the union bound for a
pair, using \(\pi>3\) and \(\sqrt2>7/5\),
is
\begin{equation}\label{eq:v48-p}
 \mathbb P\{\text{pair }a\text{ exposed}\}
 \le2\,{\rm Haar}(B_R)
 <\frac{10R^3}{63}
 =\frac{1875}{469762048}
 <p,\qquad p=\frac1{250000}.
\end{equation}

\begin{lemma}[A Jacobian bound resolved by exposed-pair count]
\label{lem:v48-jacobian}
On \(M_e=m\ge1\), the inverse Haar Jacobian
satisfies
\begin{equation}\label{eq:v48-jacobian}
 j_e\le J_m,\qquad
 J_m=
 \left(1-\frac m8-\frac{m^2\rho^2}{256}\right)^{-3}
 \left(1-\frac{m^2\rho^2}{768}\right)^{-2}.
\end{equation}
On \(M_e=0\), \(B_e=C_e\) and \(j_e=1\).
At \(\rho=1/64\), \(J_1<2\) and
\(\max_{1\le m\le7}J_m=J_7<513\).
\end{lemma}
\begin{proof}
The pivot equation always gives
\(d(B_e,C_e)\le7\rho/8\).
If a source is active at the solved pivot,
then \(d(B_e,P)<\rho\), hence
\(d(C_e,P)<15\rho/8\). All active
sources belong to exposed pairs. On
\(M_e=m\), there are at most \(2m\)
such sources, so the same equation
improves the displacement to
\[
                         d(B_e,C_e)\le r_m:=m\rho/8.
\]
If \(m=0\), every cutoff term and its
derivative vanish, giving the stated
identity pivot and Jacobian.

For \(m\ge1\), use the inherited pivot
potential
\[
 \mathcal V_C(B)=\tfrac12d(B,C)^2
        -\frac1{16}\sum_{a,\ell}\Phi_\rho(d(B,P_{a,\ell})),
 \qquad \Phi_\rho'(r)=r\chi(r/\rho).
\]
At the solved pivot, its Hessian is
bounded below by
\[
 \left(1-\frac{r_m^2}{4}-\frac{2m}{16}\right)I
       =:\kappa_m I>0.
\]
The squared-distance transverse Hessian
is at least \(1-r_m^2/4\) on this small
ball. Each active \(\Phi_\rho\) has Hessian
at most \(I\), because
\(\chi+s\chi'\le1\) and its spherical
transverse eigenvalue is at most one.
Inactive terms have zero Hessian.

Differentiate the critical equation
\(\nabla_B\mathcal V_C=0\) with respect
to \(C\), keeping the products fixed.
The source derivative is the differential
of \(\log_B C\). Consequently the
absolute determinant of \(dB/dC\) is
the determinant of that differential
divided by \(\det\operatorname{Hess}_B\mathcal V_C\).
The former is
\[
       \left(\frac{r/\sqrt2}{\sin(r/\sqrt2)}\right)^2,
                       \qquad r=d(B,C).
\]
This is the inverse exponential-volume
Jacobian, so both determinants use the
same Riemannian Haar volumes. Using
\(\sin z\ge z-z^3/6\) gives the upper
bound \((1-r_m^2/12)^{-2}\).
The Hessian determinant is at least
\(\kappa_m^3\), proving
\eqref{eq:v48-jacobian}. The final
inequalities are rational substitutions
at \(\rho=1/64\), checked below.
\end{proof}

\begin{theorem}[Uniformly small deviation of the full averaged boundary kernel]
\label{thm:v48-kernel}
For every compact \(X,C_e\), at the
unchanged working radius \(\rho=1/64\),
\begin{equation}\label{eq:v48-kernel-bound}
       |k_e(X,C_e)-1|\le\delta_*<\frac1{35000},
 \qquad
 \delta_*=
       7p+513\{(1+p)^7-1-7p\}.
\end{equation}
This is uniform over volume and over all
retained or internal exterior values.
\end{theorem}
\begin{proof}
On \(M_e=1\), positivity and \(J_1<2\)
give \(|j_e-1|\le1\). On \(M_e\ge2\),
the bound \(J_7<513\) gives the safe
estimate \(|j_e-1|\le513\).
For each \(m\), independence between
the seven pair variables and
\eqref{eq:v48-p} imply
\(\mathbb P(M_e=m)\le\binom7m p^m\).
Integrating \(|j_e-1|\), and using its
zero value at \(M_e=0\), proves the
displayed expression for \(\delta_*\).
Direct rational arithmetic gives its
strict upper bound. No derivative of
the cutoff profile is bounded from below,
and no typicality under a low-temperature
Wilson law is used.
\end{proof}

The last qualification matters: the
Haar integration here is an exact
step in constructing a new conditional
reference. It is not the assertion
that the same source caps remain rare
under an arbitrary generated law.

\subsection{Reintroduce the complete Wilson interaction by an exact conditional expansion}

Write \(S_0=\sum_p s_p\) for the actual
fine Wilson action. The conservative
bound \(0\le s_p\le4\) is sufficient;
it covers the Wilson normalization in
use without a rescaling of \(\beta\).
With \(B=B(X,Y,C)\), the exact compact
joint density after integrating \(Y\) is
\begin{equation}\label{eq:v48-exact-joint}
 \prod_e k_e(X,C_e)\ \mathcal Z_\beta(X,C),
 \qquad
 \mathcal Z_\beta(X,C)
       =\mathbb E_{\otimes_e\nu_e(\cdot\mid X,C_e)}
                                  e^{-\beta S_0(X,Y,B)} .
\end{equation}
Its \(X\)-integral is exactly the
completed coarse density. The normalizer
of the whole model is still included
when this density is made a probability.

At fixed \(X,C\), set
\[
          f_p(Y)=e^{-\beta s_p(X,Y,B)}-1,\qquad
                        |f_p|\le u:=4\beta .
\]
Let \(E(p)\) be the boundary families
whose \(Y_e\)'s enter that plaquette
after solving its pivots. A plaquette
has at most four boundary links, so
\(|E(p)|\le4\). Join two distinct
plaquettes when their \(E\)-sets
intersect. This graph has maximum
degree at most
\begin{equation}\label{eq:v48-degree}
                              D=192 .
\end{equation}
Indeed a boundary family has eight
fine boundary links, each incident
to six plaquettes, hence occurs in
at most \(48\) plaquettes. Four
families give fewer than \(4\cdot48\)
possible neighbors.

A polymer \(\gamma\) is a nonempty
connected set in this graph, with
activity
\begin{equation}\label{eq:v48-activity}
        z_\gamma(X,C)
             =\mathbb E_{\otimes\nu_e}\prod_{p\in\gamma}f_p,
                  \qquad |z_\gamma|\le u^{|\gamma|}.
\end{equation}
Polymers are incompatible if they
overlap or are joined by a graph edge.
Every polymer is incompatible with
itself. Connected components of any
set of plaquettes have disjoint
hidden \(Y_e\) inputs, so conditional
independence gives the exact identity
\begin{equation}\label{eq:v48-polymer-identity}
                  \mathcal Z_\beta(X,C)
                   =\sum_{\text{compatible families }\mathcal F}
                                \prod_{\gamma\in\mathcal F}z_\gamma(X,C).
\end{equation}
The empty family contributes one.
Plaquettes with \(E(p)=\varnothing\)
are isolated singleton polymers, not
deleted terms. The internal variables
may be shared by distinct polymers;
they are fixed parameters in this
conditional expansion, not independent
random variables being factorized.

\subsection{A convergent complete interaction with an explicit weighted norm}

Place a retained variable \(X_x\) at
the coarse cell \(x\), and \(C_e\)
at its edge midpoint. Use their
periodic \(\ell^1\) distance, or the
corresponding infinite-lattice distance.
A plaquette's retained carrier \(V(p)\)
contains its endpoint cells and its
coarse edge variables. It has at most
eight variables and diameter at most
two. Any \(X_x\) belongs to at most
\(16\cdot24=384\) such carriers: a fine
cell has sixteen vertices, and a
vertex is incident to twenty-four
plaquettes. A \(C_e\) belongs to at
most \(48\) carriers.

Fix
\begin{equation}\label{eq:v48-weights}
 \mu=\log2,\qquad \sigma=2\mu+1,\qquad
                    v=e^{1+\sigma}u=4e^2u .
\end{equation}
At \(\beta\le10^{-8}\), \(e^2<8\)
gives \(v\le128\beta\).

\begin{lemma}[Absolute generated-interaction bound]
\label{lem:v48-cluster}
The logarithm of
\eqref{eq:v48-polymer-identity} has an
absolutely convergent connected-cluster
expansion. Write its negative as
\(\sum_{\mathcal A}\Psi_{\mathcal A}(X,C)\),
with repeated polymer occurrences retained
in the cluster index. If
\(m(\mathcal A)\) is the sum of their
polymer sizes, then
\begin{equation}\label{eq:v48-pinned}
 \sup_p\sum_{\mathcal A:\,p\text{ occurs}}
           e^{\sigma m(\mathcal A)}
                   \|\Psi_{\mathcal A}\|_\infty\le2v.
\end{equation}
Consequently, for the retained carrier
\(V(\mathcal A)\),
\begin{equation}\label{eq:v48-interaction}
 \sup_a\sum_{\mathcal A:\,a\in V(\mathcal A)}
 |V(\mathcal A)|e^{\mu\operatorname{diam}V(\mathcal A)}
                     \|\Psi_{\mathcal A}\|_\infty
                             \le6144v\le786432\beta .
\end{equation}
All suprema include every compact \(X,C\).
\end{lemma}
\begin{proof}
A connected set of \(n\) graph vertices
containing a prescribed vertex has at
most \((4D)^{n-1}\) possibilities: choose
a canonical rooted plane spanning tree,
bound its shapes by \(4^{n-1}\), and
assign each child one of at most \(D\)
neighbors. Thus
\[
 \sum_{\gamma\ni p}u^{|\gamma|}
                       e^{(1+\sigma)|\gamma|}
                    \le\frac{v}{1-4Dv}\le2v .
\]
A polymer incompatible with \(\gamma_0\)
contains a vertex in the closed graph
neighborhood of \(\gamma_0\), of size
at most \((D+1)|\gamma_0|\). Therefore
\[
 \sum_{\gamma\not\sim\gamma_0}
     u^{|\gamma|}e^{\sigma|\gamma|}e^{|\gamma|}
               \le2(D+1)v|\gamma_0|\le|\gamma_0|.
\]
Both numerical inequalities hold at
\(v\le128\cdot10^{-8}\).
This verifies the Kotecky--Preiss
condition with size weight \(|\gamma|\)
for the extra weighted activities.
The standard pinned absolute-cluster
bound is used in precisely the form
(2.14)--(2.15) of
\href{https://arxiv.org/abs/math-ph/0605041}
{R. Fernandez and A. Procacci,
\emph{Cluster expansion for abstract
polymer models. New bounds from an old
approach} (2007)}.
Summing that bound over root polymers
containing \(p\) proves
\eqref{eq:v48-pinned}. The supremum
norm is legitimate because
\eqref{eq:v48-activity} is uniform
in the retained parameters.

For a connected polymer cluster,
its retained carriers connect through
shared coarse boundary families or
shared plaquettes. Hence
\[
 |V(\mathcal A)|\le8m(\mathcal A),\qquad
 \operatorname{diam}V(\mathcal A)\le2m(\mathcal A).
\]
Since \(m\le e^m\), the corresponding
weighted carrier sum is at most
eight times the pinned sum, for
each plaquette carrier containing
the specified retained variable.
There are at most \(384\) such
plaquettes. Equation
\eqref{eq:v48-pinned} yields
\(8\cdot384\cdot2v=6144v\), proving
\eqref{eq:v48-interaction}.
\end{proof}

Thus the exact retained joint potential is
\begin{equation}\label{eq:v48-complete-potential}
                \mathcal U_\beta(X,C)
                  =-\sum_e\log k_e(X,C_e)
                           +\sum_{\mathcal A}\Psi_{\mathcal A}(X,C),
\end{equation}
up to a single overall normalization
constant. This is an equality, not a
quadratic or finite-order approximation.
The parameter-dependent conditional
reference in \eqref{eq:v48-kernel}
remains inside every activity.

\subsection{Uniform conditional contraction for the complete compact law}

For a compact product specification,
let \(A_{ab}\) be an upper bound on
the total variation change of its
conditional law at variable \(a\)
when only exterior variable \(b\)
is changed. Set \(A_{aa}=0\).
The weighted row norm is
\[
                   \|A\|_\mu=
                    \sup_a\sum_{b\ne a}e^{\mu d(a,b)}A_{ab}.
\]
Here \(X_x\), though it has seventeen
group components, is one product-valued
site. Total variation bounds do not
cost its internal dimension.

We use the elementary tilt bound:
if a conditional density is changed
by \(e^h\) and
\(\sup h-\inf h\le R\), its total
variation changes by at most
\(\tanh(R/4)\le R/4\).
To verify the sharp first bound,
write the normalized density ratio
as \(r\in[m,M]\), with mean one and
\(M/m\le e^R\). Its positive deviation
has expectation at most
\((M-1)(1-m)/(M-m)\). Maximizing this
quantity gives \(m=e^{-R/2}\),
\(M=e^{R/2}\), and the stated bound.

\begin{theorem}[A full compact strong-mixing region at the specified completion]
\label{thm:v48-terminal}
For \(\rho=1/64\), \(N\ge2\), and
\(0\le\beta\le10^{-8}\), the exact
joint law \eqref{eq:v48-exact-joint}
has an influence matrix with
\begin{equation}\label{eq:v48-alpha}
 \|A\|_{\log2}
     \le\frac{32}{34999}+786432\beta
                            <\frac1{100}.
\end{equation}
The actual completed coarse law,
after integrating all \(X\)'s as well,
has a weighted single-edge conditional
influence norm less than \(1/99\).
The bounds hold for every compact
conditioned exterior, not only for
good fields or a protected collar.
\end{theorem}
\begin{proof}
Put \(\delta=1/35000\) and
\(\ell_*=\log((1+\delta)/(1-\delta))\).
The range of each \(-\log k_e\)
has width at most \(\ell_*\).
Its mixed rectangular oscillation
for a pair of its three variables
is at most \(2\ell_*\), contributing
at most \(\ell_*/2\) to influence.
A cell \(X_x\) occurs in eight edge
factors, with two other variables
per factor, at distance at most
one. Counting multiplicities also
covers the parallel periodic edges
when \(N=2\). Their weighted row
contribution is at most
\[
       8e^\mu\ell_*=16\ell_*
            \le\frac{32\delta}{1-\delta}
            =\frac{32}{34999}.
\]
A coarse edge variable occurs in
only one such factor and has a
smaller bound.

For any additional interaction
\(\Psi\), its mixed rectangular
oscillation is at most \(4\|\Psi\|_\infty\).
The tilt bound therefore charges
at most \(\|\Psi\|_\infty\) for
each other variable in its carrier.
Summing and using
\eqref{eq:v48-interaction} proves
\eqref{eq:v48-alpha}. The final
strict numerical inequality is
rational at \(\beta=10^{-8}\).

It remains to check the coarse
marginal rather than assume
contraction survives marginalization
with the same matrix. Fix a target
coarse edge \(e\), condition all
other coarse edges, and leave
\[
                       F_e=\{X_x:\ x\in\Lambda\}\cup\{C_e\}
\]
free. Changing one fixed coarse
edge \(f\ne e\) produces forcing
at a free site \(a\) bounded by
\(A_{af}\). The usual coupled
conditional-update comparison
gives discrepancy vector
\((I-A_{F_eF_e})^{-1}A_{F_ef}\).
For clarity, the comparison follows
by maximally coupling each single-site
conditional update: at stationarity
its discrepancy probabilities obey
\(d\le A_{F_eF_e}d+A_{F_ef}\).
Finite compact positive densities
permit the stationary coupling,
and the matrix inverse is its
convergent nonnegative series.
This is the classical comparison
mechanism also treated in
\href{https://arxiv.org/abs/1308.4117}
{P. Rebeschini and R. van Handel,
\emph{Comparison theorems for Gibbs measures}}.

Retain its row at \(C_e\), sum over
\(f\ne e\), and insert the distance
weights. The triangle inequality
for distance and restriction to a
subset of paths bound the result
by
\[
                 \sum_{n\ge1}\|A\|_\mu^n
                    \le\frac{\|A\|_\mu}{1-\|A\|_\mu}
                                 <\frac1{99}.
\]
This is a bound on the actual coarse
conditionals, since the free-site
law used above integrates precisely
the hidden internal groups of
\eqref{eq:v48-exact-joint}.
\end{proof}

The cluster interaction is absolutely
summable with an exponential weight.
It therefore defines an infinite-lattice
joint specification. Its finite-region
boundary comparisons use the same
strict contraction bound; influences
from a boundary receding to infinity
vanish exponentially. Compactness gives
existence, and the comparison gives
uniqueness. Periodic finite-volume
laws converge locally to that state.
The uniform conditional comparison
with the \(X\)'s left free similarly
defines a unique quasilocal coarse
specification and its coarse marginal.

In particular bounded coarse cylinder
observables \(F,G\) satisfy a covariance
bound of the form
\[
 |\operatorname{Cov}(F,G)|
       \le K|S_F||S_G|\|F\|_\infty\|G\|_\infty
                      e^{-(\log2)d(S_F,S_G)}
\]
with a universal \(K\), uniformly on
the finite tori and in the limit.
One obtains this by applying the same
comparison to a bounded local tilt
by \(G\), differentiating at zero,
and summing the weighted resolvent
from the support of \(F\) to that of
\(G\). This is spatial mixing in
coarse-cell units, not a declaration
of a continuum Hamiltonian gap.

\subsection{A noncollapsed gauge-invariant compact probe}

\begin{proposition}[Uniform coarse plaquette variance]
\label{prop:v48-probe}
In the preceding strong-mixing region,
for every coarse plaquette \(p\),
\begin{equation}\label{eq:v48-variance}
        \operatorname{Var}\left(\tfrac12
                      \operatorname{Re}\operatorname{Tr}C_p\right)
                                                   >\frac15 .
\end{equation}
The same estimate holds in the unique
infinite-volume coarse state.
\end{proposition}
\begin{proof}
At a fixed coarse edge \(e\), the
joint single-site conditional has
log-density oscillation at most
\[
 \eta_*=\ell_*+192v
       \le\frac2{34999}
                    +192\cdot128\cdot10^{-8}
                                           <\frac1{1000}.
\]
The first term is its sole boundary
kernel factor. For the generated
terms, at most \(48\) plaquette
carriers contain \(C_e\), and
\eqref{eq:v48-pinned} gives total
supremum norm at most \(96v\);
oscillation costs a factor two.
The normalized conditional density
relative to Haar is consequently
between \(e^{-\eta_*}\) and
\(e^{\eta_*}\), for every fixed
internal field.

The coarse conditional with the
internal fields integrated is a
mixture of these single-edge
conditionals, with the normalized
law of \(X\) given the other coarse
edges as mixing law. The same
pointwise density bounds therefore
hold for it.

Condition all coarse edges except
one edge of \(p\). For \(N\ge2\)
that edge occurs exactly once in
the plaquette. Under Haar variation
of this edge, the normalized trace
has mean zero and variance \(1/4\):
it is a coordinate of a uniform
unit quaternion, up to Haar-preserving
multiplication or inversion.
If a density is bounded below by
\(e^{-\eta_*}\), minimization over
the centering constant gives its
variance at least \(e^{-\eta_*}/4\).
The conditional-variance identity
gives the same lower bound for the
unconditional variance. Finally
\(e^{-\eta_*}\ge1-\eta_*>999/1000\),
which proves \eqref{eq:v48-variance}.
The local limit preserves these
bounded-observable moments.
\end{proof}

This establishes noncollapse of a
specific compact gauge-invariant
probe in the endpoint region. It
does not identify a renormalized
continuum field or a nonzero
continuum coupling.

\subsection{Certificate, preservation, and the remaining bridge}

The certificate
\begin{center}\small
\path{scripts/verify_ym_f14_v48_compact_terminal.py}
\end{center}
has 8,506 bytes and SHA--256
\begin{center}\scriptsize\ttfamily
93510C56E956463DFBE18ACF923259B31E83662128817921B283774BB4C0FA52.
\end{center}
It verifies rationally the cap bound,
all seven \(J_m\) upper bounds,
\(\delta_*<1/35000\), the
Kotecky--Preiss numerical margins,
the influence bound, and the
plaquette-variance margin. Decimal
displays of \(J_m\) are diagnostics;
none of these inequalities is
decided by floating-point arithmetic.

On coarse sides two and three it
enumerates the literal chronological
paths and checks two root products
and seven pairs per edge. It verifies
the seventeen internal groups per
cell and the dependency supports
of all \(1536\) and \(7776\) fine
plaquettes. The observed maxima
are \(36\) plaquettes per boundary
family, \(216\) per retained
carrier variable, and graph degrees
\(123\) and \(131\), respectively.
The proofs use the larger analytic
bounds \(48,384,192\), valid at every
volume, rather than extrapolating
those observed maxima.

A separate exact conditional example
with three independent hidden
Bernoulli blocks and four factors
checks the polymer identity. One
factor has no hidden input and is
retained. Direct integration, the
expanded product, and the compatible
polymer sum all give \(46291/56250\).
This finite test checks the
conditional factorization logic;
the infinite cluster theorem and
uniform spatial estimates are
proved with their explicit hypotheses
above.

The new endpoint is stronger than
a criterion waiting for an unspecified
complete interaction: the compact
interaction in
\eqref{eq:v48-complete-potential}
is actually constructed and lies
in the verified region. Conversely,
it is only a one-step completed
block of a Wilson input in the
stated strong-coupling interval.
An arbitrary effective input from
many weak-coupling shells has not
been shown to satisfy these bounds.
Its generated interactions cannot
be replaced by \(\beta S_0\) in
\eqref{eq:v48-exact-joint}.
The remaining task is that
renormalization and physical-transfer
bridge, not another verification
of this endpoint.

V47 has 1,130,582 bytes and SHA--256
\begin{center}\scriptsize\ttfamily
99D04DBAA97BF00EC2D61DD7D5E68E12FD615CDAC62586C6686FB7335717988F.
\end{center}
Its entire body is preserved except
for the title version and this
terminal insertion. V48 replaces
it as the sole active main note;
the archives, companions, supplied
manuscripts, and monographs remain
unchanged.

\begin{firewall}
The explicit compact strong-mixing
region and noncollapsed plaquette
probe do not construct the
four-dimensional ultraviolet
continuum theory, show that its
actual renormalization trajectory
enters this region, or prove the
required reflection-positive
physical-time contraction.
The full Yang--Mills mass-gap
objective remains unproved and active.
\end{firewall}
\section{V49: leaf cancellation, a small complete coarse interaction, and a second full compact block without resetting the action}
\label{sec:v49}

\subsection{Context and the precise advance beyond V48}

V48 constructs a one-step compact endpoint, but its
joint-law influence bound is not an estimate in a
norm suitable for inserting the generated action
into another block. This section supplies such an
insertion, for two actual completed blocks. The
fine, intermediate, and final sides are \(4N,2N,N\),
respectively, with \(N\ge2\). The completion remains
the same one with \(\rho=1/64\).

The first new ingredient is an exact cancellation,
not a smaller pointwise estimate on the solved
pivot. After the internal groups are integrated,
every zero-coupling connected edge contribution
with a leaf vanishes. At intermediate side
\(2N\ge4\), surviving contributions therefore have
at least four edges. This produces a much smaller
complete interaction than the joint-law estimate
of V48. A weighted cluster argument controls the
additional Wilson terms and then accepts the
entire resulting interaction at the next block.

For the conservative interval
\[
                         0\le\beta\le10^{-15},
\]
the final compact law has weighted single-edge
conditional influence less than \(3/200\) at
weight \(\mu=\log2\), and its normalized
plaquette trace has variance greater than
\(1/5\). Neither step uses a fluctuation cutoff,
an interaction truncation, or a replacement
of the intermediate action by a Wilson action.
These are fixed-lattice strong-coupling results,
not entry of the ultraviolet trajectory into
this region.

The three supplied Downloads notes retain the
fingerprints recorded in V44--V47. Their
source-selection, stationarity, and exact-action
results do not identify the missing YM flow
or physical transfer. No additional conclusion
is imported from them. The scheduled full
SM/NS/Shell companion checkpoint is V50.

\subsection{The averaged kernel has an exact endpoint cancellation}

Use \(X_x,Y_e,C_e,k_e,\nu_e\) from V48 and set
\[
                  h_e(X_x,X_y,C_e)=k_e(X_x,X_y,C_e)-1,
                 \qquad e=(x,y),\qquad \delta=\frac1{35000}.
\]
Here \(X_x\) is the product of the seventeen
internal groups. The inherited estimates and
normalization give
\begin{equation}\label{eq:v49-kernel-stock}
 \|h_e\|_\infty<\delta,\qquad
                       \int h_e(X_x,X_y,C_e)\,dC_e=0 .
\end{equation}
The product Haar measure of \(X_x\) is invariant
under simultaneous conjugation of all its
components. Residual gauge covariance gives
\begin{equation}\label{eq:v49-kernel-covariance}
 h_e(g_xX_xg_x^{-1},g_yX_yg_y^{-1},g_xC_eg_y^{-1})
                           =h_e(X_x,X_y,C_e).
\end{equation}
To see this directly, apply the same fine
gauge transformation at every vertex of one
cell. Its tree links remain identity, internal
groups conjugate, and every boundary group
transforms at its two endpoints. The completed
pivot equation is covariant and Haar volume
is invariant. Integrating the seven nonpivot
groups therefore gives
\eqref{eq:v49-kernel-covariance}.

\begin{lemma}[Leaf cancellation for the actual averaged kernel]
\label{lem:v49-leaf}
For fixed \(X_y,C_e\),
\[
                    \int h_e(X_x,X_y,C_e)\,dX_x=0,
\]
and the analogous statement holds at the other
endpoint. Consequently, if a nonempty edge
set \(\gamma\) has a degree-one vertex, then
\begin{equation}\label{eq:v49-leaf-zero}
 H_\gamma(C):=
       \int\prod_{e\in\gamma}h_e(X,C_e)
                          \prod_{x\in V(\gamma)}dX_x=0 .
\end{equation}
\end{lemma}
\begin{proof}
After integrating \(X_x\), covariance under
\(g_x\) makes the result invariant under
arbitrary left multiplication of \(C_e\).
It is thus constant in \(C_e\). Its Haar
integral is zero by
\eqref{eq:v49-kernel-stock}, so that constant
is zero. At the other endpoint use right
multiplication. If a vertex is a leaf in
\(\gamma\), its internal variable occurs
in exactly one factor. Integrating it first
proves \eqref{eq:v49-leaf-zero}.
\end{proof}

This assertion does not remove a bridge
joining two cyclic subgraphs: neither
endpoint of such a bridge need be a leaf.
A nonzero connected contribution must itself
have minimum degree at least two; a graph
with a cycle and an attached leaf still
vanishes. A survivor need not be a simple
cycle. Nor does this assert \(k_e=1\).

\subsection{A complete zero-coupling coarse potential with a fourth-order smallness bound}

At zero Wilson coupling the exact coarse
density relative to product Haar is
\[
                         Z_0(C)=\int\prod_e(1+h_e(X,C_e))\,dX .
\]
Expand this finite product and decompose
each selected edge set into connected
components, where edges are adjacent when
they share a vertex. Components with
disjoint vertex sets have independent
internal Haar inputs. Thus \(Z_0\) is
exactly the hard-core polymer partition
function with activities \(H_\gamma\),
where polymers are nonempty connected
edge sets and compatibility is vertex
disjointness. Self-incompatibility is
included. By the lemma, polymers with
leaves have zero activity, and
\[
                              \|H_\gamma\|_\infty\le\delta^{|\gamma|}.
\]
Each activity is a gauge-invariant function
of the coarse edges in \(\gamma\).

For a potential presented as a summable
family of gauge-invariant interactions
\(\Phi_\alpha(C_{S_\alpha})\) on nonempty
connected edge carriers, define
\begin{equation}\label{eq:v49-input-norm}
 \|\Phi\|_{(3)}
       :=\sup_e\sum_{\alpha:\,e\in S_\alpha}
                   e^{3|S_\alpha|}\|\Phi_\alpha\|_\infty .
\end{equation}
Indices with repeated polymer occurrences
remain distinct. Constants can be retained
separately in the overall normalization;
no field-dependent term is discarded.

\begin{proposition}[The actual compact Haar image lies near Haar in an interaction norm]
\label{prop:v49-zero-potential}
On a torus of side at least four, the
exact potential \(-\log Z_0\) has a connected, gauge-invariant
interaction expansion satisfying
\begin{equation}\label{eq:v49-cycle-budget}
 \|\Phi^{(0)}\|_{(3)}
 \le B_4:=\frac{q(56q)^3}{1-56q}
       =\frac{14641}{12468750000},
                  \qquad q=\frac{55}{35000}.
\end{equation}
The same bound holds for its infinite-lattice
local interaction limit. In infinite volume
this means a Gibbs specification, not the
logarithm of a global density relative to
an infinite Haar product.
\end{proposition}
\begin{proof}
The coarse edge-adjacency graph has degree
at most fourteen: each endpoint meets
seven other edges. The number of connected
\(n\)-edge sets containing a specified edge
is at most \(56^{n-1}\), by the rooted
plane spanning-tree count used in V48.
A connected graph with minimum degree
two contains a cycle. The coarse graph
on the stated tori has girth at least
four, so a surviving polymer has
\(n\ge4\). This last step would not
hold on the side-two torus, with its
parallel periodic edges.

Put an extra weight \(e^{3|\gamma|}\)
on the activities and use the
Kotecky--Preiss size \(|\gamma|\).
Since \(e^4<55\), the sum of inflated
activities over polymers containing a
specified edge is at most
\[
             \sum_{n\ge4}56^{n-1}(\delta e^4)^n
                              \le B_4 .
\]
A polymer incompatible with \(\gamma_0\)
contains an edge in the closed edge
neighborhood of \(\gamma_0\), whose
cardinality is at most \(15|\gamma_0|\).
The condition \(15B_4<1\) verifies the
cluster criterion. The pinned absolute
cluster bound therefore gives
\[
 \sup_e\sum_{\mathcal A:\,e\in S_{\mathcal A}}
 e^{3m(\mathcal A)}\|\Phi_{\mathcal A}^{(0)}\|_\infty
                                                     \le B_4,
\]
where \(m\) counts edge occurrences with
multiplicity. Since
\(|S_{\mathcal A}|\le m(\mathcal A)\),
this proves \eqref{eq:v49-cycle-budget}.
Vertex-connected clusters have connected
edge carriers. Gauge invariance is
preserved by products, internal Haar
integration, and the logarithmic cluster
coefficients.
\end{proof}

The cluster input here and below is the
standard pinned Kotecky--Preiss bound,
in the form (2.14)--(2.15) of
\href{https://arxiv.org/abs/math-ph/0605041}
{Fernandez--Procacci, \emph{Cluster expansion
for abstract polymer models}}.
The remaining combinatorial reductions
and all smallness conditions used here
are given explicitly.

\subsection{A conditional cell-polymer lemma that keeps the intermediate law}

The convenient normalized reference at a
fixed output \(C\) is
\begin{equation}\label{eq:v49-reference}
             dP_C(X,Y)=\prod_xdX_x
                       \prod_e d\nu_e(Y_e\mid X_x,X_y,C_e).
\end{equation}
This is a probability measure. The
boundary families are independent
\emph{conditional on \(X\)}, not in
general after \(X\) is integrated.

An atom \(i\) is a bounded factor \(f_i\)
with a nonempty cell carrier \(R_i\).
It may depend on the internal groups
in \(R_i\), on boundary families with
both endpoints there, and on the
corresponding coarse sources. Include
the source dependence of each \(\nu_e\)
when specifying its coarse-edge carrier
\(T_i\). Let \(r_i=|R_i|\) and
\(\|f_i\|_\infty\le w_i\). Atoms on
disjoint cell carriers are independent
under \(P_C\): first integrate their
distinct boundary variables, and then
their distinct internal Haar variables.
This is the required independence.
Sharing an internal cell prevents that
factorization even if the boundary
families themselves are different.

\begin{lemma}[Two combinatorial sums with tunable cell weights]
\label{lem:v49-cell-cluster}
Assign each atom a nonnegative charge
\(t_i\). Suppose that for some \(a,b>0\),
\begin{equation}\label{eq:v49-q-general}
 Q:=\sup_x\sum_{i:\,x\in R_i}
       w_i\exp\{t_i+(a+b)r_i\}\le\min(a,b).
\end{equation}
Then
\[
                         Z(C)=\mathbb E_{P_C}\prod_i(1+f_i)
\]
has an absolutely convergent logarithmic
interaction expansion, up to its constant
part. For its complete cluster index
\(\mathcal A\), count every atom occurrence,
including multiplicities from the
logarithm. The negative logarithm terms
\(\Psi_{\mathcal A}\) satisfy
\begin{equation}\label{eq:v49-cell-pinned}
 \sup_x\sum_{\mathcal A:\,x\in R(\mathcal A)}
     \exp\!\left\{\sum_{i\ {\rm occurring}}t_i\right\}
                         \|\Psi_{\mathcal A}\|_\infty\le Q.
\end{equation}
The partition identity is for finite
volumes and uniformly summable atom
families there. Its interaction bounds
also pass to infinite-volume local
limits, without asserting a global
infinite-volume density.
\end{lemma}
\begin{proof}
Join two atoms when their cell carriers
intersect. Expand the product and group
each selected set into connected
components. Disjoint components have
disjoint cell carriers and hence
factorize under \(P_C\). The exact
polymer activity for a connected atom
set \(\gamma\) is
\[
                 z_\gamma=\mathbb E_{P_C}\prod_{i\in\gamma}f_i,
                    \qquad |z_\gamma|\le\prod_{i\in\gamma}w_i.
\]
Polymer incompatibility is intersection
of their cell carriers, including
self-incompatibility.

Here is a useful bound on the first
connected-set sum. Put
\(v_i=w_i e^{t_i+br_i}\).
Connected atom sets rooted at \(i\)
are bounded by rooted spanning trees.
Allowing repeated labels only enlarges
the positive tree sum. Its finite-depth
majorants obey
\[
            T_i^{(0)}=v_i,\qquad
            T_i^{(n+1)}
              =v_i\exp\!\left\{\sum_{j:\,R_j\cap R_i\ne\varnothing}
                                      T_j^{(n)}\right\}.
\]
The exponential supplies the factorial
for unordered children; equivalently
one sums labelled rooted trees with
the \(1/(n-1)!\) normalization.
Every connected set supplies at least
one spanning tree, so this is an
upper bound, not a factorization
assumption on its activity.
Inductively
\[
 T_i^{(n)}\le v_i e^{ar_i},
\]
because the exponent on the right is
at most \(r_iQ\le ar_i\).
Consequently the sum of connected
polymer majorants with a cell \(x\)
in their carrier, and with weight
\(\exp\{\sum(t_i+br_i)\}\), is at most
\[
                  \sum_{i:\,x\in R_i}v_i e^{ar_i}\le Q.
\]
Rooting at every possible atom
containing \(x\) only overcounts.

For the second sum inflate
\(z_\gamma\) by
\(\exp\{\sum_{i\in\gamma}t_i\}\)
and use polymer size
\(b\sum_{i\in\gamma}r_i\).
The preceding bound makes the sum
over polymers incompatible with
\(\gamma_0\) at most
\[
             |R(\gamma_0)|Q
                 \le b\sum_{i\in\gamma_0}r_i.
\]
This is exactly the Kotecky--Preiss
condition. Its pinned bound, summed
over root polymers containing \(x\),
proves \eqref{eq:v49-cell-pinned}.
Uniform absolute majorants allow
truncation of a countable family and
passage to its limit. This proof uses
no positivity of individual activities.
\end{proof}

When the \(R_i\)'s are connected and
\(T_i\) spans \(R_i\), with the allowed
exception \(R_i=\{x\},T_i=\varnothing\),
each nonempty output carrier
\(\bigcup T_i\) is connected. Empty
source carriers contribute constants
or connect other atoms through their
single cell; those atoms are retained.
For the actual blocked Wilson
plaquettes these geometric conditions
hold: project the fine plaquette
walk onto cells and erase its internal
steps. Its projected walk has either
one cell and no edge, two cells and
one edge, or four cells and four edges.
Inverse pivots and conditional
\(\nu_e\)'s use exactly those endpoint
cells and boundary families.

\subsection{The first complete Wilson output satisfies a reusable small-interaction bound}

In \eqref{eq:v49-reference}, take two
types of atoms:
\[
       h_e,\qquad f_p=e^{-\beta s_p(X,Y,B(X,Y,C))}-1 .
\]
The exact first-block density is
\begin{equation}\label{eq:v49-first-exact}
               Z_\beta(C)
                 =\mathbb E_{P_C}
                          \prod_e(1+h_e)\prod_p(1+f_p).
\end{equation}
Indeed \(\prod_e(1+h_e)\,dP_C
=dX\prod_ej_e\,dY_e\).
Thus both the original Wilson action
and every inverse Haar factor occur
once, with the global normalizer
retained upon forming the probability.

\begin{theorem}[A small full generated action, not a Wilson projection]
\label{thm:v49-first-action}
On an output torus of side at least
four, for \(0\le\beta\le10^{-15}\),
the exact action
\(-\log Z_\beta(C)\) has a connected
gauge-invariant interaction family
\(\Phi_\beta\) with
\begin{equation}\label{eq:v49-first-small}
 \|\Phi_\beta\|_{(3)}
 \le B_4+\frac{\beta}{R}Q_R
 \le\frac{685051993}{194824218750000}
                         <\frac1{250000},
 \qquad R=10^{-11},
\end{equation}
where
\begin{equation}\label{eq:v49-reference-budget}
                   Q_R=\frac{256}{35000}+1610612736R
                              =\frac{8005024}{341796875}<\frac1{10}.
\end{equation}
\end{theorem}
\begin{proof}
Use Lemma \ref{lem:v49-cell-cluster}
with \(a=b=1/10\) and
\(t_i=3|T_i|\). An edge atom has
\(r_i=2,|T_i|=1,w_i=\delta\);
eight such atoms meet a cell.
A plaquette atom has
\(r_i\le4,|T_i|\le4\).
At most \(384\) fine plaquettes
have a specified cell in their
projected carrier, by the sixteen
vertices times twenty-four incident
plaquettes bound.

At the reference value \(R\), use
\(w_p=8R\). This also majorizes
\(e^{4R}-1\); for example
\(e^x-1\le x/(1-x)\le2x\)
for \(0\le x\le1/2\).
The atom budget is at most
\[
          8\delta e^{17/5}
                  +384\,(8R)e^{64/5}
             <256\delta+384\,(8R)2^{19}=Q_R.
\]
The stated bounds on exponentials
are checked by rational Taylor
majorants below. Since \(Q_R<1/10\),
the cell-polymer lemma applies
uniformly in every compact source.
For a cluster with nonempty coarse
carrier \(S\),
\[
                    |S|\le\sum_{i\ {\rm occurring}}|T_i|.
\]
Anchoring at an endpoint cell of a
specified coarse edge therefore
converts \eqref{eq:v49-cell-pinned}
into the norm \eqref{eq:v49-input-norm}.

To improve the pure-kernel part,
mark every occurrence of a plaquette
atom. The unmarked logarithmic
clusters are exactly the expansion
of \(Z_0\); they are bounded by
Proposition \ref{prop:v49-zero-potential},
not by the coarser \(Q_R\).
Every remaining cluster has at least
one marked occurrence. For real
\(0\le\beta\le R\),
\[
             |f_p|\le4\beta
                       \le(\beta/R)(8R).
\]
In the absolute cluster majorant
each marked occurrence supplies this
factor. Since \(\beta/R\le1\), a
cluster with one or more marks
is bounded by \(\beta/R\) times
its reference majorant. The total
marked contribution to the norm
is consequently at most
\((\beta/R)Q_R\). This argument keeps
mixed and repeated occurrences; it
is not first-order perturbation theory.

All source-dependent terms are
gauge invariant. This follows from
residual gauge covariance of \(P_C\),
the \(h_e\)'s, and the pulled-back
Wilson plaquette factors, followed
by integration over the relevant
internal and boundary groups.
The carrier argument following
Lemma \ref{lem:v49-cell-cluster}
gives connected supports. Finally,
substitution of \(\beta=10^{-15}\)
gives the rational bound in
\eqref{eq:v49-first-small}.
\end{proof}

Thus the effect of the first block
is described by its complete
interaction, including the entire
zero-coupling part. No assertion
that the completion fixes product
Haar exactly is needed or made.

\subsection{A second completed block accepts this entire generated action}

We first give a reusable one-step
acceptance theorem. Its input is
an arbitrary gauge-invariant law
on the fine torus of side \(2N\),
presented as
\begin{equation}\label{eq:v49-general-input}
       d\mu(U)=Z^{-1}
          \exp\!\left\{-\sum_\alpha\Phi_\alpha(U_{S_\alpha})\right\}dU,
              \qquad \|\Phi\|_{(3)}\le B:=\frac1{250000},
\end{equation}
with nonempty connected edge
carriers and the full family
summed. Constants are included
in \(Z\). No Wilson form is assumed.

\begin{theorem}[Compact blocking of a complete small generated interaction]
\label{thm:v49-acceptance}
For \(N\ge2\), the exact completed
pushforward of
\eqref{eq:v49-general-input}, at
\(\rho=1/64\), has a summable
coarse interaction \(\widehat\Phi\)
satisfying
\begin{equation}\label{eq:v49-final-interaction}
 \sup_e\sum_{\mathcal A:\,e\in S_{\mathcal A}}
 |S_{\mathcal A}|e^{(\log2)\operatorname{diam}S_{\mathcal A}}
                       \|\widehat\Phi_{\mathcal A}\|_\infty
 \le3Q_2<\frac3{200},
 \qquad
           Q_2=\frac{136}{35000}+256B=\frac{537}{109375}<\frac1{200}.
\end{equation}
Distances are periodic \(\ell^1\)
distances between edge midpoints.
The bound controls the actual
single-edge conditional influence
matrix in that same weighted row
norm. Every single-edge conditional
density relative to Haar lies
between \(e^{-2Q_2}\) and \(e^{2Q_2}\).
\end{theorem}
\begin{proof}
Gauge-fix the cell trees at this
new level and use a new copy of
the exact reference
\eqref{eq:v49-reference}. The
factorization of the input action
is now
\[
                  \prod_\alpha(1+F_\alpha),
             \qquad F_\alpha=e^{-\Phi_\alpha(U(X,Y,C))}-1 .
\]
Together with the new \(h_e\)
factors this is exactly the
transformed compact density.
The earlier \(h\)'s and all Wilson
corrections have already become
part of \(\Phi\); they are not
removed or counted twice.

Let \(n_\alpha=|S_\alpha|\).
Each \(\|\Phi_\alpha\|_\infty\le B\),
so
\[
                         \|F_\alpha\|_\infty
                                      \le2\|\Phi_\alpha\|_\infty .
\]
Project the connected fine carrier
onto the new cells. Its cell set
\(R_\alpha\) is connected and has
\(r_\alpha\le2n_\alpha\). Its
boundary edges span \(R_\alpha\)
unless it is contained in one
cell, in which case it has no
coarse-edge source. Tree dressings,
solved pivots, and the conditional
reference use only these cells
and boundary families.

Fix
\[
 \mu=\log2,\qquad \lambda=\tfrac12,\qquad
 a=b=\tfrac1{10},\qquad
 c=\mu+\lambda+a+b=\log2+\tfrac7{10}.
\]
In the cell-polymer lemma choose
\(t_i=(\mu+\lambda)r_i\).
For new edge atoms the budget
is at most
\[
                      8\delta e^{2c}<136\delta ,
\]
since \(e^{2c}=4e^{7/5}<17\).
For an input atom,
\(e^{cr_\alpha}\le e^{2cn_\alpha}
\le e^{3n_\alpha}\), because
\(2c<3\).
If a specified cell belongs to
\(R_\alpha\), some fine edge in
\(S_\alpha\) has an endpoint in
that cell. There are at most
\(16\cdot8=128\) such fine edges.
Hence their total budget is at most
\[
 \begin{split}
 \sum_{\alpha:\,x\in R_\alpha}
                   2\|\Phi_\alpha\|_\infty e^{cr_\alpha}
 &\le2\sum_{\substack{\text{fine edges }f\\
                         f\text{ touches cell }x}}
       \ \sum_{\alpha:\,f\in S_\alpha}
                   e^{3n_\alpha}\|\Phi_\alpha\|_\infty\\
 &\le256B .
 \end{split}
\]
This proves the atom budget
\(Q\le Q_2<1/200<\min(a,b)\),
including every generated input
interaction.

For an output logarithmic cluster
write \(r=\sum r_i\), with
occurrences counted. Its connected
cell union has at most \(r\)
vertices. Its coarse-edge carrier
has at most \(4r\) edges and
diameter at most \(r\). The first
bound uses degree eight, and the
second follows by a spanning-tree
path through the cell union,
including the two half-edge ends.
Both statements allow the parallel
periodic edges at \(N=2\).
Therefore
\[
 |S|e^{\mu\operatorname{diam}S}
      \le4r e^{\mu r}
      \le\frac8e e^{(\mu+1/2)r}
                                      <3e^{(\mu+1/2)r}.
\]
Anchor at an endpoint cell of the
specified edge and apply
\eqref{eq:v49-cell-pinned}. This
gives \eqref{eq:v49-final-interaction}.

For a real interaction \(\Psi\),
the mixed rectangular oscillation
in two edge variables is at most
\(4\|\Psi\|_\infty\). The tilt
inequality proved in V48 bounds
the corresponding total-variation
influence by \(\|\Psi\|_\infty\).
Summing over the other edges in
each carrier proves the asserted
weighted row bound. At one edge,
the unweighted sum of incident
interaction norms is at most
\(Q_2\), directly from
\eqref{eq:v49-cell-pinned}.
Its conditional log-density
oscillation is at most \(2Q_2\).
Normalization relative to Haar
gives the displayed two-sided
conditional-density estimate.
\end{proof}

\begin{theorem}[Two actual full compact blocks, with a noncollapsed final probe]
\label{thm:v49-two-blocks}
Start with the Wilson law on the
four-dimensional torus of side
\(4N\), \(N\ge2\), at
\(0\le\beta\le10^{-15}\). Apply the
same completed block twice, each
time with \(\rho=1/64\), retaining
the complete pushforward law.
The final law has weighted
conditional-influence norm less
than \(3/200\) at \(\mu=\log2\).
Its normalized trace on every
final plaquette \(p\) obeys
\begin{equation}\label{eq:v49-final-variance}
             \Var\!\left(\tfrac12\operatorname{Re}\Tr C_p\right)
                                      \ge\frac{e^{-2Q_2}}4>\frac15 .
\end{equation}
The finite-volume laws have a
unique infinite-volume local
limit with a quasilocal
specification and exponential
clustering of bounded cylinder
observables in final-cell units.
\end{theorem}
\begin{proof}
The intermediate side is \(2N\ge4\),
so the girth-four bound in
Theorem \ref{thm:v49-first-action}
applies. Its complete potential
satisfies \eqref{eq:v49-general-input}.
Theorem \ref{thm:v49-acceptance}
then applies to the second step,
with no change of input action.
Associativity of pushforward, or
Fubini's theorem for the positive
compact integrals, identifies this
law with the literal composition
of the two completed maps.

For the variance, condition all
final edges except one edge of
\(p\). That edge occurs once
because \(N\ge2\). Under its
Haar variation the normalized
trace has variance \(1/4\).
The conditional density lower
bound \(e^{-2Q_2}\), minimization
over the centering constant,
and the conditional-variance
identity give
\eqref{eq:v49-final-variance}.
Since \(2Q_2<1/100\),
\(e^{-2Q_2}>99/100\), which is
more than the required margin.

The absolute cluster sums give
uniformly summable interactions;
local terms converge as the
periodic boundary recedes.
The weighted conditional
influence is strictly less
than one. The same coupled
conditional-update comparison
and weighted resolvent used in
V48 give uniqueness, exponential
boundary insensitivity, and
exponential clustering. Compactness
supplies existence. All these
bounds are for the full compact
specification, including arbitrary
conditioned coarse exteriors.
\end{proof}

\subsection{What has, and has not, become iterable}

The first action estimate now
controls an entire generated
interaction in a stated norm,
and the second block demonstrably
accepts it. This removes an
actual one-to-two-step gap left
by V48. The general acceptance
theorem also applies to any
other gauge-invariant input
that is proved to satisfy
\eqref{eq:v49-general-input}.

However, the conclusion
\eqref{eq:v49-final-interaction}
is an exponentially weighted
spatial-influence bound. It
is not the input norm
\eqref{eq:v49-input-norm} with
the same numerical constant.
It is therefore not legitimate
to feed that conclusion into
the acceptance theorem indefinitely.
For example, an attempt to
produce the output weight
\(e^{3|S|}\) directly from the
general atom lemma charges
\(3|T_\alpha|+(a+b)r_\alpha\),
which the current simple bound
only controls by
\((3+2/5)n_\alpha\), not
\(3n_\alpha\). The support-size
weight loses a positive amount
under this estimate. Resolving
that loss needs additional
geometric or gauge cancellation,
not an assertion that ordinary
mixing is preserved under
arbitrary conditioning.

Even a stable strong-coupling
basin under all further blocks
would not prove that the
weak-coupling ultraviolet
trajectory enters it. The
nonlinear flow and physical
reflection/clock comparison
remain independent obligations.
The present variance bound
is for a compact lattice probe,
not a proof of continuum
interaction or field normalization.

\subsection{Exact checks and single-version preservation}

The certificate
\begin{center}\small
\path{scripts/verify_ym_f14_v49_two_compact_blocks.py}
\end{center}
has 10,296 bytes and SHA--256
\begin{center}\scriptsize\ttfamily
7CFBF3589E8F18336E2CA93E4574BB67E9A2D9E826DD053BF21C51E67F3E3B01.
\end{center}
It checks the displayed budgets
with rational arithmetic.
Exponential inequalities use
rational Taylor sums and a
geometric upper bound for the
remaining tail. Decimal values
are only diagnostics.

For coarse sides two and four,
it enumerates all \(1536\) and
\(24576\) fine plaquettes,
respectively. Their projected
cell/edge cardinalities are
exactly \((1,0),(2,1),(4,4)\);
all nonempty edge carriers
span their projected cells.
It finds \(216\) plaquettes
and \(96\) fine edges incident
to a cell. The proof uses the
larger analytic bounds \(384\)
and \(128\), rather than
extrapolating these enumerations.

A separate nonabelian finite-group
test uses the standard character
\(\chi\) of \(S_3\) and
\[
 h(x,y,c)=\frac1{100}
     \left\{\chi(xcy^{-1}c^{-1})
                        -\frac{\chi(x)\chi(y)}2\right\}.
\]
It verifies endpoint covariance,
zero conditional Haar integral,
and the leaf cancellation exactly.
Every nonempty proper subset of
a square has zero activity, while
the complete square gives
\(3/1600000000\) at identity
holonomy and \(-1/1600000000\)
at a reflection. An attached
leaf integrates to zero.
This test illustrates why
leaf cancellation does not
mean the complete kernel is
Haar. It is not a replacement
for the SU(2) kernel proved
above.

Finally, an exact conditional
model with three internal bits,
two boundary bits, and five
atoms verifies the cell-based
polymer identity. The boundary
bits are independent only
conditionally; the common
internal bit creates a nonzero
correlation. Direct integration,
product expansion, and the
twenty-eight-polymer partition
sum all give
\(523150687/537600000\).
An atom with empty source
carrier is retained. These
finite certificates check the
algebra and constants; the
uniform cluster and comparison
theorems require the analytic
proofs above.

V48 has 1,155,992 bytes and SHA--256
\begin{center}\scriptsize\ttfamily
2BACDFC4F5A33C326D3EC8841878A446FA4AFDCB171089DBD30EBA6FE9C2C0C7.
\end{center}
Its body is preserved except
for the title version and
this terminal insertion.
V49 replaces V48 as the sole
active main note. Archives,
companions, supplied source
files, and monographs are
unchanged.

\begin{firewall}
The new result concerns two
full compact completed blocks
of a strong-coupling Wilson
input, not the four-dimensional
continuum limit. No all-depth
invariant interaction ball,
entry from an ultraviolet
trajectory, reflection-positive
physical transfer, or positive
physical Hamiltonian mass gap
has been established here.
The full Yang--Mills objective
remains unproved and active.
\end{firewall}
\section{V50: a support-halving root reference, an all-depth reference contraction, and quantitative comparison with the actual completion}
\label{sec:v50}

\subsection{Companion checkpoint and the upstream question}

V49 completed two actual compact blocking steps,
but its output estimate did not return to its
input interaction norm. We investigate that
failure upstream. There are two different maps
in this section, and their conclusions must
not be interchanged:
\[
 \mathcal D:\hbox{ exact two-link root decimation},
 \qquad
 \mathcal B_\rho:\hbox{ the actual smooth completed root map},
                    \quad \rho=1/64 .
\]
The new all-depth contraction theorem is for
\(\mathcal D\). We also prove estimates for
\(\mathcal B_\rho\), including its normalized
conditional fibre and actual plaquette responses.
This is a comparison construction, not a
replacement of the programme's completed map
or a claim that its all-depth estimate is closed.

The compulsory V50 companion review was performed
on the current files, not on remembered version
numbers. The latest relevant files remain
SM V72, NS V26, and Shell V16, with SHA--256:
\begin{center}\scriptsize\ttfamily
\begin{tabular}{ll}
SM V72 & F44F9745D43379E1672C5550411B58E97195A4C9278592D221DF5D3F644DC1F5\\
NS V26 & 82C887F8C2C69E4F28FAD7C3DC8DE5B9099CB5A4BF431EBA1FFF3E91935978AD\\
Shell V16 & BF138E63BA826BC030D3A8B87E04FDD3B1EC712D269BB064674E96509B46E710
\end{tabular}
\end{center}
SM V72 gives empirical covariance transport
in a declared many-copy scaling with fixed
physical mode sets and bounded times, not a
growing-cutoff YM estimate. NS V26 supplies
explicit rank-one Oseen contractions; its
references to future V27 grid certification
are not treated as an executed estimate here.
Shell V16 retains the complete measured
flat-root one-loop coefficient and its
summable refinement error. Its proposed
noncommuting quartic calculation has already
been carried out in the present lineage,
notably V30, so it is not repeated.

The three supplied Downloads notes also retain
their recorded hashes:
\begin{center}\scriptsize\ttfamily
\begin{tabular}{ll}
Source Selection V17 & A3AB81EED2BAC59CE4E7B3F97EAFA36B9B5F7090589D60FAAAC57C6AF05CE6D5\\
Stationarity V21 & C2D452CA7A031E2AF8B319073A79FE820C37007BDFC5442B8158916539E1A580\\
Exact Action V16 & 53C50B2CFBA4A05B8749CF165EF1ADEEC8B7B1133FA9FB6AD6D45C8D0F5EB49F
\end{tabular}
\end{center}
Their mathematical scope remains as assessed
in V44--V49. None supplies a missing invariant
YM interaction ball, a continuum trajectory,
or a physical-time comparison. All these
sources were read as mathematical data,
not as new instructions. The next scheduled
companion checkpoint is V55.

\subsection{The root reference in ungauge-fixed Haar coordinates}

On a fine torus of side \(2N\), \(N\ge4\),
define
\begin{equation}\label{eq:v50-root-map}
 (\mathcal DU)_{x,\mu}
       =U_{2x,\mu}U_{2x+\widehat\mu,\mu}.
\end{equation}
The distinguished two-link root paths are
edge-disjoint. Their intermediate vertices
are distinct from the retained even vertices.
Under product Haar, the outputs are therefore
independent Haar groups. Conditionally on
their values \(C_{x,\mu}\), one can write
the two links of each path as
\[
                          V_{x,\mu},\quad
                                      V_{x,\mu}^{-1}C_{x,\mu},
\]
with independent Haar \(V\)'s and with every
other fine link independent Haar.
Denote this conditional expectation by
\(\mathbb E_0[\cdot\mid C]\).

For a fine edge set \(S\), let \(K(S)\)
be the coarse edges whose \emph{both}
root-path links belong to \(S\).
If \(F\) depends on fine links in \(S\),
then
\begin{equation}\label{eq:v50-full-pair-support}
                     \mathbb E_0[F\mid C]
                         \text{ depends only on }C_{K(S)} .
\end{equation}
Indeed, if only the first link of a root
pair is used, it is Haar independently
of \(C\); if only the second is used,
the same follows by the Haar change of
variable \(V^{-1}C\). The independent
conditional pairs allow this argument
one pair at a time, even for a
nonfactorized \(F\).

If \(F\) is fine-gauge invariant, its
conditional expectation is coarse-gauge
invariant. This follows either from the
conditional representation or from
equivariance of \(\mathcal D\).
A gauge-invariant function supported
on a forest of coarse edges is constant:
a gauge transformation at a leaf
removes dependence on its incident
edge, and one successively removes
all leaves. Consequently a nonconstant
conditional expectation needs a cycle
in \(K(S)\). Since \(N\ge4\), this
requires at least four coarse edges
and hence at least eight fine edges
in \(S\). No assertion about the
girth-two side-two torus is made.

\subsection{A simultaneous-rounding lemma removes the support-size loss}

For \(\varepsilon\in\{0,1\}^4\), define
the periodic vertex map
\begin{equation}\label{eq:v50-rounding}
          R_\varepsilon z
               =\left\lfloor\frac{z+\varepsilon}{2}\right\rfloor
                                                     \pmod N .
\end{equation}
Map each fine edge to its coarse edge
if its endpoints have different
images, and delete it if they coincide.
All retained even vertices map to
their specified coarse vertices.

\begin{lemma}[Connected support halving with retained root paths]
\label{lem:v50-halving}
For every finite connected fine edge
set \(S\), there exists an offset
\(\varepsilon(S)\) for which the
coarse image \(T(S)\) is connected,
possibly empty, and
\begin{equation}\label{eq:v50-halving}
                   K(S)\subseteq T(S),\qquad
                             |T(S)|\le\frac{|S|}{2}.
\end{equation}
For any prescribed coarse edge \(e\),
there are at most \(128\) pairs
\((\varepsilon,f)\) with \(f\) a
fine edge mapped to \(e\) by
\(R_\varepsilon\).
\end{lemma}
\begin{proof}
Each fine edge is noncollapsed for
exactly eight of the sixteen offsets:
only the offset bit in its own
orientation decides this. Counting
edge occurrences, the average image
size is therefore at most \(|S|/2\).
Choose an offset minimizing the
number of distinct image edges,
with any fixed tie-breaking order.
The image of a connected graph
under edge collapse is connected.
Every complete root pair maps to
its coarse edge for every offset,
because its endpoints are \(2x\)
and \(2x+2\widehat\mu\). This
proves the inclusion as well.

For fixed \(\varepsilon\) and
\(e=(x,\mu)\), a fine edge with
that noncollapsed image has its
longitudinal base coordinate fixed:
\(z_\mu=2x_\mu+1-\varepsilon_\mu\)
modulo \(2N\). Each of the other
three coordinates has two choices.
There are thus eight preimages
per offset, or \(16\cdot8=128\)
with offsets counted. This upper
bound does not require that the
preimages be distinct across
different offsets.
\end{proof}

The carrier \(T(S)\) can contain
edges on which the conditional
expectation is actually constant.
They are only connecting carriers;
their inclusion changes no function.
The same offset is used for the
entire connected set, not chosen
independently for its individual
edges.

\subsection{A complete nonlinear contraction theorem for the root reference}

Use the interaction norm of V49:
\[
 \|\Phi\|_{(3)}
       =\sup_f\sum_{\alpha:\,f\in S_\alpha}
                         e^{3|S_\alpha|}\|\Phi_\alpha\|_\infty .
\]
Assume that the input interactions
are real, continuous, individually
gauge invariant, and carried by
nonempty connected fine edge sets.
Constants are kept in the overall
normalization. A countable family
is interpreted by its uniform
absolute sums on each finite
torus.

\begin{theorem}[An invariant interaction ball for exact root decimation]
\label{thm:v50-root-contraction}
If
\[
                          B:=\|\Phi\|_{(3)}\le\frac1{160},
\]
then the exact pushforward by
\(\mathcal D\) has a complete
connected, gauge-invariant
interaction presentation
\(\mathcal R_0\Phi\), modulo
constants, such that
\begin{equation}\label{eq:v50-contraction}
                 \|\mathcal R_0\Phi\|_{(3)}
                         \le4096e^{-44/5}B<\frac58 B .
\end{equation}
This is a bound on the full
logarithm, not just its linear
or quadratic part.
\end{theorem}
\begin{proof}
Introduce atoms
\[
                         f_\alpha=e^{-\Phi_\alpha}-1,
       \qquad n_\alpha=|S_\alpha|,\qquad
       \|f_\alpha\|_\infty\le2\|\Phi_\alpha\|_\infty .
\]
The last bound follows from
\(\|\Phi_\alpha\|_\infty\le B\)
and \(e^B<2\). The exact coarse
density relative to Haar is
proportional to
\[
                   \mathbb E_0\!\left[
                              \prod_\alpha(1+f_\alpha)\mid C\right].
\]
For the overlap carrier of an
atom use its \emph{fine vertex}
set \(R_\alpha=V(S_\alpha)\),
with \(r_\alpha\le2n_\alpha\).
Two disjoint such carriers cannot
use the two separate links of
one conditional root pair:
those links share their middle
vertex. Hence atoms on disjoint
vertex carriers are independent
under the conditional Haar
reference. The two-sum lemma
of V49 applies to these carriers
just as it did to cell carriers.

Set
\[
             a=b=\frac1{10},\qquad
                           t_\alpha=\frac{13}{5}n_\alpha .
\]
At a fine vertex \(v\), use its
eight incident edges to bound
the atom sum:
\[
 \begin{split}
 \sum_{\alpha:\,v\in R_\alpha}
 2\|\Phi_\alpha\|_\infty
               e^{(13/5)n_\alpha+(1/5)r_\alpha}
 &\le
 2\sum_{f:\,v\in f}\ 
    \sum_{\alpha:\,f\in S_\alpha}
                 e^{3n_\alpha}\|\Phi_\alpha\|_\infty\\
 &\le16B\le\frac1{10}.
 \end{split}
\]
Thus V49's connected-set sum and
pinned logarithmic-cluster sum
give, for their full terms
\(\Psi_{\mathcal A}\),
\begin{equation}\label{eq:v50-fine-pinned}
 \sup_v\sum_{\mathcal A:\,v\in R(\mathcal A)}
         e^{(13/5)m(\mathcal A)}
                         \|\Psi_{\mathcal A}\|_\infty\le16B .
\end{equation}
Here \(m\) is the sum of all input
edge-carrier sizes, counted with
the atom occurrences in the
logarithmic cluster. Repeated
occurrences are not deleted.
This application uses the same
standard pinned Kotecky--Preiss
bound as V49, in
\href{https://arxiv.org/abs/math-ph/0605041}
{Fernandez--Procacci, equations
(2.14)--(2.15)}.

The union \(S(\mathcal A)\) of
the fine edge carriers is
connected, because each polymer
and each incompatible cluster
connects through fine vertices.
Each polymer activity depends
only on complete root pairs in
its own fine carrier. Therefore
\(\Psi_{\mathcal A}\) depends
only on coarse edges in
\(K(S(\mathcal A))\), and is
gauge invariant. Subtract its
coarse Haar average, adding
that scalar to the normalization.
Call the result
\(\widetilde\Psi_{\mathcal A}\).
Then
\[
             \|\widetilde\Psi_{\mathcal A}\|_\infty
                         \le2\|\Psi_{\mathcal A}\|_\infty .
\]
If the centered term is nonzero,
its coarse support cannot be a
forest. The preceding root-pair
argument implies
\[
                         |S(\mathcal A)|\ge8,\qquad
                                      m(\mathcal A)\ge8.
\]
Give this term the connected
carrier \(T(S(\mathcal A))\)
from Lemma \ref{lem:v50-halving}.
It has size at most \(m/2\).
Consequently
\[
 e^{3|T|}\|\widetilde\Psi_{\mathcal A}\|_\infty
 \le2e^{(3/2)m}\|\Psi_{\mathcal A}\|_\infty
 \le2e^{-44/5}
             e^{(13/5)m}\|\Psi_{\mathcal A}\|_\infty .
\]

For a specified coarse edge
\(e\in T\), some fine edge of
\(S(\mathcal A)\) is a preimage
of \(e\) under the chosen
offset. There are at most
128 possible offset/preimage
pairs. Anchor each at the
base vertex of that fine edge
and use \eqref{eq:v50-fine-pinned}.
Overcounting only enlarges the
absolute sum. This gives
\[
                 2\cdot128\cdot16B\,e^{-44/5},
\]
which is \eqref{eq:v50-contraction}.
The numerical strict inequality
follows from
\(e^{44/5}>32768/5\), verified
by a rational Taylor lower sum.
All cluster terms, including
their constant parts, have
been accounted for.
\end{proof}

The proof of the two-sum lemma
is already given in V49; it
is not a new unverified
combinatorial hypothesis here.
Its only probabilistic input
in the present application is
the explicit conditional
root-pair factorization above.

\begin{proposition}[All-depth consequence for the reference map only]
\label{prop:v50-root-depth}
On a fine torus of side \(2^mN\),
\(N\ge4\), an input with
\(\|\Phi\|_{(3)}\le1/160\) has,
after \(m\) exact root-decimation
steps, an interaction presentation
with
\begin{equation}\label{eq:v50-root-depth}
                  B_m\le(5/8)^m\|\Phi\|_{(3)} .
\end{equation}
The same statement holds for
the associated infinite-lattice
local specifications. The
complete generated interaction
is used at every step.
\end{proposition}
\begin{proof}
The output class in
Theorem \ref{thm:v50-root-contraction}
has exactly the same connected
carrier norm, continuity, and
gauge invariance as its input.
Its norm is smaller. Induction
therefore applies at every
intermediate side, all of which
are at least four. Uniform
cluster sums and local finite
volume limits give the analogous
infinite-lattice interactions;
no global density relative to
an infinite Haar product is
asserted.
\end{proof}

For example, in the standing
normalization
\[
                 S_0(U)=\sum_p(1-\tfrac12\operatorname{Re}\Tr U_p),
\]
put the constant part in the
normalizer. The Wilson interaction
then has norm at most
\(6e^{12}\beta\), since an edge
meets six plaquettes. The
explicit range \(0\le\beta\le10^{-9}\)
lies inside the reference ball:
\[
                     6e^{12}\beta
                              <6\cdot2^{18}10^{-9}<1/160.
\]
This range is not being claimed
for repeated completed blocks.

Because a connected edge
carrier has diameter at most
its cardinality,
\[
             |S|e^{(\log2)\operatorname{diam}S}\le e^{3|S|}.
\]
The influence argument of V48
therefore bounds the reference
output's weighted Dobrushin
norm by \(B_m\). Its single-edge
conditional log-density
oscillation is at most \(2B_m\).
For any fixed set of \(k\)
coarse edges, sequential
coupling to Haar gives total
variation at most \(kB_m\);
one first uses the single-site
tilt bound and then mixtures
for partially conditioned laws.
Thus these reference outputs
approach product Haar locally
as \(m\) grows. This is an
infrared statement for a
strong-coupling lattice region,
not an ultraviolet continuum
construction or a physical
Hamiltonian mass-gap theorem.

\subsection{Return to the actual completed map: a normalized fibre comparison}

Return now to \(\mathcal B_\rho\)
at the unchanged \(\rho=1/64\).
V48's exposed-pair proof gives
the stronger integrated estimate
used there, not just its
consequence for \(k_e\):
\begin{equation}\label{eq:v50-j-lone}
                  \int|j_e(X,Y_e,C_e)-1|\,dY_e<\delta,
                 \qquad \delta=\frac1{35000}.
\end{equation}
Consequently, with
\[
          \epsilon=\frac{\delta}{1-\delta}=\frac1{34999},
       \qquad p=\frac1{250000},\qquad
                                  \tau=\epsilon+7p,
\]
the conditional reference
\(\nu_e=k_e^{-1}j_e\,dY_e\)
has total variation at most
\(\epsilon\) from Haar.
Indeed
\(\int|j_e-k_e|\le2\delta\).
Also, under Haar \(Y_e\),
the probability of any exposed
pair is at most \(7p\), and
outside that event the solved
pivot is exactly \(B_e=C_e\).

Let \(F_\rho(X,Y)\) be a bounded
observable in the actual solved
fibre and let \(F_0(X,Y)\) be
the same fine observable with
all relevant pivots set to
\(C_e\). Suppose the observable
uses \(q\) boundary families
and \(r\) internal cells,
including the endpoint cells
needed by those families.
At fixed \(X,C\), couple the
relevant \(\nu_e\)'s to their
Haar references. A coupling
failure or an exposed Haar
family is the only possible
discrepancy in the resulting
fine configurations. Therefore
\begin{equation}\label{eq:v50-local-reference}
 \left|\mathbb E_{\otimes\nu_e}F_\rho
                          -\mathbb E_{\otimes dY_e}F_0\right|
                  \le2q\tau\,\|F\|_\infty .
\end{equation}
This estimate uses no derivative
of \(F\) and no bound on a
derivative of the cutoff profile.
It is important not to mistake
this normalized conditional
reference for the complete
conditional fibre law.

The complete zero-coupling
conditional fibre given \(C\)
has internal marginal
\begin{equation}\label{eq:v50-internal-marginal}
                 d\lambda_C(X)
                  =Z_0(C)^{-1}\prod_e k_e(X,C_e)\,dX,
\end{equation}
and conditional boundary law
\(\otimes_e\nu_e\). This retains
the entire product of averaged
Jacobians.

\begin{proposition}[Comparison for the complete normalized zero-coupling fibre]
\label{prop:v50-full-fibre}
With the preceding supports,
\begin{equation}\label{eq:v50-full-fibre}
 \left|\mathbb E_{\rho,C}F-\mathbb E_{0,C}F\right|
 \le
       \left(2q\tau+\frac{8r\epsilon}{1-8\epsilon}\right)
                                                    \|F\|_\infty .
\end{equation}
Here \(\mathbb E_{\rho,C}\) is
the actual normalized completed
fibre at zero Wilson coupling,
and \(\mathbb E_{0,C}\) is the
pure-root conditional Haar law,
expressed in the same cell-tree
coordinates. The bound is
uniform in volume and in every
compact coarse source.
\end{proposition}
\begin{proof}
Put
\(\ell=\log((1+\delta)/(1-\delta))\le2\epsilon\).
In the internal marginal
\eqref{eq:v50-internal-marginal},
a cell appears in eight kernel
factors. Its conditional
log-density oscillation relative
to product internal Haar is at
most \(8\ell\). The single-cell
total-variation discrepancy is
therefore at most
\(\tanh(2\ell)\le2\ell\le4\epsilon\).
Two cells interact only through
their shared edge factors.
The mixed rectangular oscillation
per factor is at most \(2\ell\),
so their unweighted influence
row sum is at most
\(8(\ell/2)\le8\epsilon<1\).
Parallel periodic edges are
counted with multiplicity.

Coupled conditional updates
against product Haar, followed
by the nonnegative resolvent,
bound the discrepancy at each
internal cell by
\[
                   \frac{4\epsilon}{1-8\epsilon}
                                             =\frac4{34991}.
\]
The conditional expectation
\(\mathbb E_{\otimes\nu_e}F_\rho\)
is a function of the specified
\(r\) cells only, with norm at
most \(\|F\|_\infty\). Its
expectation changes by at most
twice that norm times the sum
of these cell discrepancies.
This gives the second term of
\eqref{eq:v50-full-fibre}.
Integrating
\eqref{eq:v50-local-reference}
under internal Haar gives
the first term. Cell-tree gauge
fixing has no additional Haar
Jacobian for gauge-invariant
fine observables, so the root
reference is exactly the
ungauge-fixed conditional law
used earlier.
\end{proof}

In particular this compares
the linear response to a
bounded local perturbation
of the input action. It is
not a global total-variation
comparison of arbitrarily
large volumes: the number
of cells and families used
by the observable is retained.

\subsection{An actual nonpivot character gain and two normalized plaquette responses}

\begin{lemma}[Uniform nontrivial representation gain in a boundary family]
\label{lem:v50-character}
Let \(\pi\) be a nontrivial
irreducible unitary representation
of \(\SU(2)\), and let \(Y_{e,a}\)
be any of the seven nonpivot
groups in family \(e\). For
all \(X,C_e\),
\begin{equation}\label{eq:v50-character}
             \left\|\int\pi(Y_{e,a})\,d\nu_e\right\|_{\rm op}
                                                        \le\epsilon.
\end{equation}
The same bound holds for an
inverse or for fixed unitary
left and right factors.
\end{lemma}
\begin{proof}
The Haar average of \(\pi\)
is the orthogonal projection
onto invariant vectors, which
is zero for this representation.
Thus the integral in question is
\[
                  k_e^{-1}\int(j_e-1)\pi(Y_{e,a})\,dY_e .
\]
Unitary operator norm one,
\eqref{eq:v50-j-lone}, and
\(k_e\ge1-\delta\) give the
claim. Inversion and fixed
unitary factors preserve the
same argument. No dimension
factor is needed.
\end{proof}

Consider a fine plaquette \(p\)
whose four vertices project to
four distinct coarse cells.
Its four links are boundary
links in four distinct
families. At least two are
nonpivot groups. To check the
last assertion, its two varying
coordinates start at odd parity.
If either of the two remaining
transverse parities is nonzero,
all four boundary links are
nonpivot. If both are zero,
exactly two are nonpivot and
two are distinguished pivots.
This accounts for all parity
types; no typical-field
assumption is used.

Let
\[
                          \chi_p(U)=\tfrac12\operatorname{Re}\Tr U_p .
\]
At fixed \(X,C\), the four
boundary families are
independent under their
normalized \(\nu_e\)'s.
Expectations of the ordered
matrix product therefore
factor, without commuting
its factors. Each nonpivot
factor has norm at most
\(\epsilon\), whereas a
pivot factor has norm at
most one. It follows that
\begin{equation}\label{eq:v50-plaquette-mean}
                       |\mathbb E_{\rho,C}\chi_p|\le\epsilon^2 .
\end{equation}
The estimate remains valid
after integration against
the actual internal marginal
\(\lambda_C\). It is a bound
for the complete conditional
fibre, not just the auxiliary
reference.

\begin{proposition}[First and second normalized action responses for an actual plaquette insertion]
\label{prop:v50-plaquette-response}
Add the single input action
\(t(1-\chi_p)\), with the
plaquette \(p\) just described,
to the zero-coupling Haar law.
Let \(w_t(C)\) be the negative
logarithm of the \emph{normalized}
coarse density produced by
the actual completed map.
Then
\begin{align}
 \left|\partial_tw_t(C)\big|_{t=0}\right|
       &\le\epsilon^2=\frac1{1224930001},\label{eq:v50-score-one}\\
 \left|\partial_t^2w_t(C)\big|_{t=0}\right|
       &\le8\tau+\epsilon^4<\frac1{2000}.
                                                   \label{eq:v50-score-two}
\end{align}
Also
\begin{equation}\label{eq:v50-fine-variance}
                        \Var_{\rho,C}(\chi_p)>\frac15 .
\end{equation}
The bounds hold for every
compact \(C\).
\end{proposition}
\begin{proof}
With \(Z_t(C)\) the full
unnormalized fibre integral
and \(Z_{\rm fine}(t)\) its
coarse Haar integral, write
\[
                       w_t(C)=-\log Z_t(C)+\log Z_{\rm fine}(t).
\]
The original unconditional
fine Haar plaquette has
\(\mathbb E\chi_p=0\) and
\(\Var(\chi_p)=1/4\), since
its holonomy is Haar and
the normalized trace is a
unit-quaternion coordinate.
Differentiation of these
positive compact integrals
therefore gives exactly
\[
              w'_0(C)=-\mathbb E_{\rho,C}\chi_p,\qquad
              w''_0(C)=\frac14-\Var_{\rho,C}(\chi_p).
\]
The first bound follows from
\eqref{eq:v50-plaquette-mean}.

In the root reference at
fixed \(X,C\), the holonomy
of this four-family plaquette
is Haar: it contains an
independent Haar nonpivot
group occurring once. Thus
\(\mathbb E_{0,Y}(\chi_p^2)=1/4\)
for every fixed \(X,C\).
Apply \eqref{eq:v50-local-reference}
with \(q=4\) and \(F=\chi_p^2\),
whose norm is at most one.
This gives
\[
               \left|\mathbb E_{\rho,Y}(\chi_p^2)-\frac14\right|
                                                         \le8\tau .
\]
Because the reference value
is constant in \(X\), averaging
against \(\lambda_C\) requires
no additional internal-marginal
comparison. Subtracting the
squared mean, bounded by
\(\epsilon^4\), proves
\eqref{eq:v50-score-two}.
The same estimates give
\[
                   \Var_{\rho,C}(\chi_p)
                              \ge\frac14-8\tau-\epsilon^4>\frac15.
\]
The last inequalities are
checked rationally below.
\end{proof}

This proposition concerns one
specified fine plaquette
insertion about zero input
coupling. It is not a bound
in the full interaction norm
on simultaneous arbitrarily
many insertions. In particular,
the small first conditional
moment is compatible with
an order-one conditional
second moment.

\subsection{The remaining obstruction is a specific complete-interaction estimate}

The support-size loss in
V49 is not inevitable for
root decimation. Theorem
\ref{thm:v50-root-contraction}
closes an invariant ball
for that exact nonlinear
reference map and all of
its generated interactions.
It uses conditional Haar
elimination, the girth
constraint, and a single
rounding of each entire
connected fine carrier.

For the actual completed map,
the support statement
\eqref{eq:v50-full-pair-support}
is no longer exact: an
unretained source can enter
through a pivot or through
the conditional density
\(\nu_e\). We now have
derivative-free local
comparison bounds and an
operator-norm gain for each
eligible nontrivial group
representation. They provide
quantitative inputs for
estimating the difference
\(\mathcal R_\rho-\mathcal R_0\),
but do not yet give that
difference a closed bound
in \(\|\cdot\|_{(3)}\).

An occurrence cannot earn
an independent character
gain merely because it is
written twice. For a single
hidden group,
\[
                      \pi(Y)\pi(Y)^{-1}=I
\]
even though \(\int\pi(Y)\,dY=0\).
Two factors using the same
family can therefore lose
their separate Haar
cancellations. The actual
plaquette calculation makes
this issue quantitative:
its mean is at most
\(\epsilon^2\), while its
conditional variance exceeds
\(1/5\). The normalization
subtracts the Haar vacuum
variance in \(w''_0\); it
does not authorize dropping
all repeated-source terms.
Nonlinear overlapping
insertions must be grouped
and centered before claiming
a multiplicative small factor.

The next complete estimate
must exploit these gains
without charging an
unbounded support-size loss
to every completed-map
correction. Neither a
pointwise first-response
bound nor a small norm for
the root reference proves
that estimate. Beyond it,
entry of the ultraviolet
trajectory and calibrated
physical reflected transfer
remain separate requirements.

\subsection{Certificate and preservation}

The certificate
\begin{center}\small
\path{scripts/verify_ym_f14_v50_root_support_and_completion.py}
\end{center}
has 8,850 bytes and SHA--256
\begin{center}\scriptsize\ttfamily
2B8966298935E5AD35401ED5B8936800FB08EEE8C0635D2119A8F4D6B12F7B3C.
\end{center}
It verifies all numerical
margins with rational
arithmetic and Taylor
upper or lower sums for
the exponential. The
displayed decimal contraction
bound is diagnostic only.

On fine side eight it
checks all \(16384\) fine
edges against all sixteen
roundings. Each edge is
noncollapsed eight times;
one specified coarse edge
has \(128\) preimages counted
with offsets and \(54\)
distinct preimage edges.
Every complete root pair
has the same coarse image
under every rounding.
All \(4095\) nonempty edge
subsets of an embedded
unit cube satisfy the
support-halving inequality;
an eight-edge root square
saturates it. The general
proof is the edgewise
averaging argument, not
extrapolation from these
finite examples.

The ninety-six fine
plaquette parity/orientation
types are also enumerated.
Among the twenty-four
four-cell types, six have
two nonpivots and eighteen
have four. A finite \(S_3\)
test checks the conditional
root-pair algebra, the
eight-edge root-loop
identity, the zero Haar
average of a primitive
crossing plaquette, and
the failure of independent
gains for \(Y Y^{-1}\).
This is a structural
regression, not a replacement
of \(\SU(2)\) or a
certificate of a continuum
mass gap.

V49 has 1,184,549 bytes and SHA--256
\begin{center}\scriptsize\ttfamily
9E95E99B56C132DF7CA00935070794CFD5AB63A1BE1872129A80689D6F1E964D.
\end{center}
Its complete body is retained,
apart from the title version
and this terminal addition.
V50 is the single active
main note. All companion
files, archives, supplied
notes, and monographs remain
unchanged.

\begin{firewall}
The all-depth invariant ball
proved here is for exact
root decimation, not for
the actual smooth completed
map. The latter gains
quantitative conditional
comparison and plaquette
response estimates, not a
closed all-depth nonlinear
flow. No nontrivial
four-dimensional continuum
construction or positive
physical Hamiltonian mass
gap has been proved.
The full Yang--Mills goal
remains unproved and active.
\end{firewall}
\section{V51: grouped family defects and an all-depth compact interaction ball for a smaller fixed completion radius}
\label{sec:v51}

\subsection{Context: change the completion parameter explicitly, not the Wilson measure}

V50 identified the obstruction to assigning
an independent small character factor to
every insertion: several insertions can
use the same hidden boundary family.
We now group all uses of that family
into one signed integration operator.
Its norm is small independently of
the number of insertions to which it
is applied. A local footprint allows
that operator to be included in the
same connected-cluster expansion as
the full generated input interaction.

The resulting sufficient constants
are very conservative. They do not
close the invariant ball at the
earlier completion radius \(1/64\).
For this new theorem choose instead
\begin{equation}\label{eq:v51-radius}
                              \rho_\star=2^{-100}.
\end{equation}
This is the same smooth completion
family, with the same cutoff profile
rescaled to a smaller \emph{fixed}
radius at every step. It is neither
the discontinuous choice \(\rho=0\)
nor a limit replacing the completion
by root decimation.

The Wilson input measure is unchanged.
The global blocking observable has
changed. The previous results at
\(\rho=1/64\), including their claim
boundaries, remain as written.
The new result is an all-depth
invariant strong-coupling interaction
ball for the actual completed map
\(\mathcal B_{\rho_\star}\), with
every generated interaction retained.
It is not a proof that the earlier
radius has this property or that a
weak-coupling ultraviolet trajectory
enters this ball.

The main source is the verified V50
state. No newer main file was present.
The companion checkpoint was completed
in V50; the next scheduled review
remains V55. No source document is
being used as a set of instructions.

\subsection{Uniform small-radius bounds for a complete boundary family}

Fix the same nonincreasing smooth
cutoff profile \(\chi\), equal to
one in the inner half of its
support. Allow \(0<\rho\le1/64\).
The root equation is
\[
 C_e=B_e\exp\left\{\frac1{16}
              \sum_{a=1}^7\sum_{\ell=1}^2
                         g_\rho(B_e^{-1}P_{e,a,\ell})\right\}.
\]
V48's inverse-pivot argument applies
throughout this range. For completeness,
every solution satisfies
\(d(B_e,C_e)\le7\rho/8\).
On the geodesic ball of radius \(\rho\)
around \(C_e\), the pivot potential
has Hessian at least
\[
                             (1/8-\rho^2/4)I>0.
\]
Its outward derivative on that
boundary is positive, since the
source gradient has norm at most
\(7\rho/8\). It has an interior
minimum, its critical point is
unique, and differentiation gives
the smooth global inverse.
Shrinking \(\rho\) does not
weaken any of these estimates.
As before, no lower bound on
\(\chi'\) is needed.

For fixed internal groups and
coarse source, expose a boundary
pair if one of its two products
lies within \(15\rho/8\) of the
coarse source. Normalized Haar
and the cap calculation in V48
give a pair probability less than
\[
           \frac{10}{63}(15\rho/8)^3
                  =\frac{1875}{1792}\rho^3< p_\rho,
                     \qquad p_\rho:=2\rho^3 .
\]
The seven pair events are
independent under their seven
Haar inputs. The inverse Haar
Jacobian bounds \(J_1<2\) and
\(J_m<513\), \(1\le m\le7\),
remain valid because the
V48 expressions increase with
\(\rho\). Therefore
\begin{equation}\label{eq:v51-jacobian-small}
 \int|j_e-1|\,dY_e
 \le\delta_\rho:=
 7p_\rho+513\{(1+p_\rho)^7-1-7p_\rho\}
                          <8p_\rho .
\end{equation}
Indeed \(p_\rho\le1/131072\);
the polynomial \(\delta_\rho/p_\rho\)
has nonnegative coefficients,
and its value at that endpoint
is less than eight. The endpoint
inequality is checked rationally.

\subsection{Use original internal Haar links to keep the fine support visible}

Temporarily undo the cell-tree
gauge fixing. Keep all thirty-two
internal fine links of each
cell as independent Haar
variables \(I_x\). There are
eight boundary links per
coarse edge \(e=(x,\mu)\).
Seven form \(Y_e\), and one
is the root pivot link.
Let \(A_e=U_{2x,\mu}\) be the
first link of the distinguished
root path. If its root product
is \(B_e\), the pivot link is
\[
                                  A_e^{-1}B_e .
\]
The internal tree variables
are retained here, rather
than incorrectly counted as
new non-gauge degrees of freedom.
Integrating their normalized
gauge orbits gives exactly
the earlier seventeen-group
cell-tree description for
gauge-invariant integrands.

The dressed nonroot products
depend on the internal links
of the two endpoint cells and
on the seven boundary inputs
in that one family. They do
not involve any other pivot.
At fixed \(I\), the Haar
left/right dressings of a
boundary input preserve the
probability and Jacobian
estimates above.

For a bounded integrand \(F\),
possibly containing many
input insertions, define
the unnormalized family
operators
\begin{align}
 (\mathsf T^\rho_eF)(I,C_e)
    &=\int j_e(I,Y_e,C_e)
            F(I,Y_e,B_e(I,Y_e,C_e))\,dY_e,
                                      \label{eq:v51-T-rho}\\
 (\mathsf T^0_eF)(I,C_e)
    &=\int F(I,Y_e,C_e)\,dY_e,
                    \qquad
 \mathsf D_e:=\mathsf T^\rho_e-\mathsf T^0_e .
                                      \label{eq:v51-D}
\end{align}
Other family coordinates are
left as parameters. The
argument \(B_e\) in \(F\)
means that the actual fine
pivot link is \(A_e^{-1}B_e\).
Thus \(\mathsf T^0\) is
precisely the conditional
root Haar operator of V50,
not a reset of an input
interaction.

\begin{lemma}[One signed defect for an entire family]
\label{lem:v51-defect}
Uniformly over internal
fields and coarse sources,
\begin{equation}\label{eq:v51-defect-norm}
                 \|\mathsf D_e\|_{\infty\to\infty}
                    \le d_\rho:=\delta_\rho+14p_\rho
                                               <44\rho^3 .
\end{equation}
Different family operators
commute at fixed \(I\), and
the product of \(k\) distinct
defects has norm at most
\(d_\rho^k\).
There is one such factor
per selected family, not
per input occurrence.
\end{lemma}
\begin{proof}
In the difference of the
two integrals add and
subtract \(F(I,Y_e,B_e)\)
with Haar weight. The part
with \(j_e-1\) is at most
\(\delta_\rho\|F\|_\infty\).
In the other part the
integrands coincide unless
some pair is exposed:
outside that event \(B_e=C_e\).
The union bound gives
probability at most \(7p_\rho\),
and the difference is at
most \(2\|F\|_\infty\).
This proves
\(\delta_\rho+14p_\rho<22p_\rho\).

At fixed internal fields,
each operator integrates
only its own seven \(Y_e\)'s
and substitutes only its
own pivot. Its kernel
depends on no other family's
integration variables or
pivot. The operators are
therefore tensor-product
operators on these family
coordinates. They commute,
and their norms multiply.
This remains true when
the integrand is a product
of arbitrarily many factors
sharing those coordinates.
\end{proof}

In particular
\(\mathsf D_e1=k_e-1\).
No separate \(k_e\) or
\(\nu_e\) is to be multiplied
into this formulation:
the full unnormalized
Jacobian already occurs
in \(\mathsf T^\rho_e\).
The defect is signed;
it is not assigned a
probability interpretation.

\subsection{The finite footprint of a defect}

Let \(\mathcal A_e\) be the
fine edge set consisting of
all internal edges in the
two endpoint cells of \(e\),
together with all eight
boundary links of family \(e\).
It is connected and has
\begin{equation}\label{eq:v51-footprint}
                  |\mathcal A_e|=72,\qquad
                         |V(\mathcal A_e)|=32.
\end{equation}
The two internal four-cubes
have thirty-two edges each,
and the boundary family
adds eight. Every fine
vertex belongs to at most
eight such footprints,
one for each incident
coarse edge of its cell.
All coefficients of
\(\mathsf D_e\), all of
its substitutions, and
all tree dressings needed
for them use this footprint.
It includes both links
of the distinguished root
pair for \(e\).

These large footprints
are deliberate conservative
carriers. The numerical
radius in \eqref{eq:v51-radius}
pays their cost. No claim
of an optimal radius is
made.

\subsection{An exact mixed expansion: input atoms and family operators}

Let the full input action
be \(\sum_\alpha\Phi_\alpha\),
with the real continuous
gauge-invariant interactions
and connected carriers of
V50. Put
\[
 f_\alpha=e^{-\Phi_\alpha}-1,\qquad
 n_\alpha=|S_\alpha|,\qquad
                    B=\|\Phi\|_{(3)}.
\]
The exact unnormalized
completed output density is
\begin{equation}\label{eq:v51-exact-expansion}
 Z_\rho(C)=
  \int dI\,
       \prod_e(\mathsf T^0_e+\mathsf D_e)
                           \prod_\alpha(1+f_\alpha).
\end{equation}
The products of family
operators act on the
entire product of input
atoms before internal
integration. Expanding
both products selects a
set of input atoms and
a set of distinct family
defects. Unselected
families use \(\mathsf T^0\).

Assign an input atom the
fine vertex carrier
\(V(S_\alpha)\), and a
defect atom the carrier
\(V(\mathcal A_e)\).
Join atoms when these
carriers intersect.
For a connected selected
set \(\gamma\), let
\(z_\gamma(C)\) be its
joint integral, with
root operators on all
other relevant families.

\begin{lemma}[Mixed component factorization and its source support]
\label{lem:v51-mixed}
The connected components
of any selected atom
set factorize exactly.
Thus \eqref{eq:v51-exact-expansion}
is a hard-core polymer
partition function on
connected mixed atom
sets, with incompatibility
given by intersection of
fine vertex carriers,
including self-incompatibility.
Its activity bound is
\begin{equation}\label{eq:v51-mixed-bound}
 \|z_\gamma\|_\infty
       \le d_\rho^{\,|\gamma_D|}
                    \prod_{\alpha\in\gamma_I}\|f_\alpha\|_\infty .
\end{equation}
If
\[
 S(\gamma)=
       \bigcup_{\alpha\in\gamma_I}S_\alpha
           \ \cup\!
       \bigcup_{e\in\gamma_D}\mathcal A_e ,
\]
then \(z_\gamma\) is
gauge invariant and
depends only on the
coarse edges in \(K(S(\gamma))\),
the complete-root-pair
carrier defined in V50.
\end{lemma}
\begin{proof}
For fixed \(I\), the defect
operators have the tensor
bound of
Lemma \ref{lem:v51-defect}.
All other family integrations
have norm one. Internal
Haar integration also has
norm one. This proves
\eqref{eq:v51-mixed-bound}
even when many input
factors use the same
family.

Disjoint fine vertex
carriers cannot share
an internal Haar link,
a boundary variable,
or the two separate
links of a root pair:
those two links share
their intermediate fine
vertex. A defect's full
coefficient and substitution
dependence has been included
in \(\mathcal A_e\).
Consequently the integrals
of disconnected components
factor under the independent
Haar variables and conditional
root pairs. This proves
the exact polymer identity,
not just its majorant.

For the source assertion,
a selected defect can
introduce \(C_e\) directly,
but its footprint contains
the complete root pair
for \(e\). For any other
coarse edge \(f\), the
only possible source
dependence comes through
its root substitution.
If only one link of that
pair occurs in the union
of the input carriers
and defect footprints,
the same Haar change
of variable as in V50
eliminates \(C_f\).
This includes the case
in which the first link
of a pair enters a
defect coefficient and
the second enters an
input atom: in that
case both have correctly
been charged to the
union carrier. No smaller
support is assumed.

Finally expand each
selected \(\mathsf D_e\)
as the difference of
the completed and root
operators. Each summand
is the unnormalized
coarse fibre density
for a map choosing the
completed output on
some edges and the root
output on the others,
with a gauge-invariant
fine integrand. Every
such mixed output map
is gauge covariant,
and Haar measure is
gauge invariant. Each
summand, hence the
activity, is coarse
gauge invariant.
\end{proof}

The lemma also justifies
the use of V49's two-sum
argument for these
\emph{operator} atoms.
That argument needs the
exact component partition
identity, the product
activity majorant, and
the stated overlap
carriers. It does not
require the defects
themselves to be positive
scalar functions.

\subsection{An invariant norm for the complete smooth pushforward}

\begin{theorem}[Small-radius compact invariant ball]
\label{thm:v51-basin}
Fix \(\rho=\rho_\star=2^{-100}\)
and a fine torus of side
\(2N\), \(N\ge4\). If
\[
                         \|\Phi\|_{(3)}=B\le\frac1{320},
\]
then the actual completed
pushforward has a complete
connected, continuous,
gauge-invariant interaction
presentation \(\mathcal R_{\rho_\star}\Phi\)
with
\begin{equation}\label{eq:v51-affine}
 \|\mathcal R_{\rho_\star}\Phi\|_{(3)}
          \le\frac58B+\frac1{32000}<\frac1{320}.
\end{equation}
More precisely the additive
term can be replaced by
\((5/16)d_{\rho_\star}e^{968/5}\).
This is an invariant-ball
estimate for the full
nonlinear law. It is not
a Lipschitz contraction
theorem between two input
laws.
\end{theorem}
\begin{proof}
Use \(a=b=1/10\) in
V49's two-sum argument.
For an input atom choose
the charge
\[
                          t_\alpha=(13/5)n_\alpha .
\]
Its norm is at most
\(2\|\Phi_\alpha\|_\infty\),
and its vertex carrier
has size at most
\(2n_\alpha\). As in V50,
the total inflated input
weight at a fine vertex
is at most \(16B\).

For a defect choose
edge-size charge \(72\)
and vertex size \(32\).
Its inflated weight is
\[
       d_\rho\exp\{(13/5)72+(1/5)32\}
                                      =d_\rho e^{968/5}.
\]
At most eight defects
meet a fine vertex.
The combined atom
budget is therefore
\begin{equation}\label{eq:v51-Q}
                     Q\le16B+8d_\rho e^{968/5}.
\end{equation}
At the chosen radius,
the explicit estimates
are
\[
 \begin{split}
 d_{\rho_\star}e^{968/5}
       &<44\,2^{-300}2^{280}
                      =\frac{11}{262144}<\frac1{10000},\\
 Q&<\frac1{20}+\frac8{10000}<\frac1{10}.
 \end{split}
\]
The exponential bound
\(e^{968/5}<2^{280}\)
is certified by a
rational Taylor upper
sum and geometric tail.
Both connected-set and
logarithmic-cluster
criteria consequently
hold uniformly.

For a logarithmic cluster
\(\mathcal C\), let
\[
 m(\mathcal C)=
       \sum_{\text{input occurrences }\alpha}n_\alpha
                +72\,\#\{\text{defect occurrences}\}.
\]
The occurrences include
all repetitions arising
in the logarithm. The
pinned bound gives
\begin{equation}\label{eq:v51-pinned}
 \sup_v\sum_{\mathcal C:\,v\in R(\mathcal C)}
       e^{(13/5)m(\mathcal C)}
                        \|\Psi_{\mathcal C}\|_\infty\le Q.
\end{equation}
The union of the input
edge carriers and defect
footprints in this
cluster is connected.
Its source support is
contained in its complete
root pairs, by
Lemma \ref{lem:v51-mixed}.
Subtract the coarse
Haar mean of each
logarithmic term, keeping
that scalar in the
normalizer. Gauge
invariance makes a
centered term zero if
its complete-root-pair
support is a forest.
Every nonzero term
thus has \(m\ge8\).

Apply V50's simultaneous
rounding to the entire
connected union, including
the defect footprints.
It gives a connected
coarse carrier \(T\)
of size at most \(m/2\)
containing the source
support. Centering costs
at most two in norm, so
\[
        e^{3|T|}\|\widetilde\Psi_{\mathcal C}\|_\infty
        \le2e^{-44/5}e^{(13/5)m}
                                      \|\Psi_{\mathcal C}\|_\infty .
\]
For an output edge in
\(T\), choose a fine
preimage under that
rounding. There are
at most 128 choices
with offsets counted.
Anchoring each at its
fine base vertex and
using \eqref{eq:v51-pinned}
therefore gives
\[
 \begin{split}
 \|\mathcal R_\rho\Phi\|_{(3)}
       &\le256e^{-44/5}
                         (16B+8d_\rho e^{968/5})\\
       &<\frac58B+\frac5{16}d_\rho e^{968/5}\\
       &<\frac58B+\frac1{32000}.
 \end{split}
\]
The last quantity is
strictly below \(1/320\)
when \(B\le1/320\).
All carrier sums are
absolute. Continuity
follows from the smooth
finite-dimensional
family kernels and
uniform convergence;
gauge invariance follows
termwise from the
preceding lemma.
\end{proof}

The standard logarithmic
cluster input is the
pinned Kotecky--Preiss
estimate used and proved
applicable in V49--V50;
its reference form is
\href{https://arxiv.org/abs/math-ph/0605041}
{Fernandez--Procacci,
equations (2.14)--(2.15)}.
The new analytic work
here is the grouped
signed-family operator,
its local support, and
the resulting mixed
activity bound.

\subsection{All-depth completed compact laws, without a reset}

\begin{theorem}[Arbitrary finite depth in the new completed strong-coupling basin]
\label{thm:v51-depth}
Let the finest torus have
side \(2^mN\), \(N\ge4\),
and let its full input
interaction satisfy
\(B_0\le1/320\).
Apply \(\mathcal B_{\rho_\star}\)
at each of the \(m\)
steps, passing the
complete normalized
output law to the next.
Then the resulting
interaction norms obey
\begin{equation}\label{eq:v51-depth}
 B_j\le(5/8)^jB_0+
                \frac{1-(5/8)^j}{12000},
                       \qquad 0\le j\le m.
\end{equation}
In particular they stay
in the same invariant
ball at every depth,
uniformly over the
finest volume and \(m\).
\end{theorem}
\begin{proof}
The output class in
Theorem \ref{thm:v51-basin}
is the same class as
its input, including
the connected edge
carrier norm. All
intermediate torus
sides are at least
four. Iteration is
therefore legitimate.
The scalar recurrence
\(B_{j+1}\le(5/8)B_j+1/32000\)
gives \eqref{eq:v51-depth}
by summing its geometric
series. The underlying
measures are the exact
composed pushforwards,
by Fubini and the
positivity of the
unexpanded compact
integrals. No generated
input is projected
onto a Wilson form.
\end{proof}

In the normalization
\[
                        S_0=\sum_p(1-\tfrac12\operatorname{Re}\Tr U_p),
\]
the Wilson input has
\(B_0\le6e^{12}\beta\)
after its constant is
put in the normalizer.
The explicit interval
\[
                              0\le\beta\le10^{-9}
\]
lies inside \(B_0<1/320\),
since \(6e^{12}10^{-9}
<6\cdot2^{18}10^{-9}<1/320\).
This supplies a nonempty
interval for the actual
completed, not merely
root-decimated, law
at the new radius.

The norm bound controls
weighted single-edge
conditional influence,
as in V50, because
\[
                   |S|e^{(\log2)\operatorname{diam}S}\le e^{3|S|}.
\]
Thus every positive-depth
output has influence
norm less than \(1/320\).
Its single-edge
conditional density
relative to Haar lies
between \(e^{-2B_j}\)
and \(e^{2B_j}\).
For every output
plaquette \(p\),
\begin{equation}\label{eq:v51-probe}
       \Var\!\left(\tfrac12\operatorname{Re}\Tr C_p\right)
                  \ge\frac{e^{-2B_j}}4>\frac15 .
\end{equation}
The proof is the same
conditional Haar
variance argument as
V48--V50, now with a
depth-uniform bound.

For each fixed depth,
the uniformly summable
interactions have an
infinite-volume local
specification. Compactness
gives existence, and the
strict weighted influence
bound gives uniqueness,
quasilocality, and
exponential clustering
of bounded cylinder
observables, with the
displayed constants
uniform over depth.
These statements concern
spatial separation in
the output lattice's
units. They do not
identify a physical
time semigroup.

The affine estimate does
not say that the outputs
converge to Haar or even
that the entire sequence
has a limit as depth
goes to infinity. Its
additive term retains
the completion's own
possible non-Haar bias.
A Lipschitz estimate
between two input laws,
or a separate scale
convergence argument,
would be needed for
those stronger claims.

\subsection{What the radius change preserves, and what it does not prove}

At every fixed finite
block, the two smooth
completion maps with
radii \(\rho_\star\)
and \(1/64\) agree
on a sufficiently
small neighborhood of
the identity configuration:
all source-to-root
distances there are
less than \(\rho_\star/2\),
so every cutoff is one
for both maps. They
therefore have identical
fixed-order jets at
the identity. This
does not give a
depth-uniform chart
radius or identify
their global finite-field
renormalization
trajectories.

No change was made
to the fine Wilson
probability law. The
choice of completion
is a choice of global
blocking coordinates
and coarse observables.
Nevertheless, the new
radius must be carried
explicitly into any
later comparison with
the local ultraviolet
construction or physical
observables. The enormous
footprint constant
\(e^{968/5}\) explains
the conservative radius;
failure of this sufficient
bound at \(1/64\) is
not a counterexample
to stability at that
earlier radius.

The current new result
closes an actual
all-depth compact
terminal basin for a
specified positive
completion radius.
It does not put the
weak-coupling continuum
trajectory into that
basin, establish an
interacting continuum
limit, or compare the
spatial contraction
with a reflection-positive
physical-time transfer.
Those remain the
requirements relevant
to the full YM mass gap.

\subsection{Exact checks and single-version preservation}

The certificate
\begin{center}\small
\path{scripts/verify_ym_f14_v51_family_defect_basin.py}
\end{center}
has 9,836 bytes and SHA--256
\begin{center}\scriptsize\ttfamily
13CE93D93586CDD7F7CFB7D25EB088E49764B35F16A63D9C3667CA7A684F0135.
\end{center}
It reuses the retained
V50 rational exponential
bound routine. All
numerical decisions
are exact rational
inequalities, including
the endpoint bound
\(\delta_\rho<8p_\rho\),
the 800-term exponential
upper sum and its
geometric tail, and
the invariant-ball
margin.

On coarse side four,
it enumerates all
1024 actual family
footprints. Each has
72 fine edges and
32 fine vertices,
contains only its own
complete root pair,
and is connected.
Every fine vertex
belongs to eight
footprints. Their
rounding geometry
is checked against
V50's support-halving
argument; the proof
uses the general
incidence counts,
not an extrapolation
from one volume.

A finite conditional
kernel regression has
three internal bits,
two boundary bits,
two pivot bits, four
input atoms, and two
family defect atoms.
The resulting sixty
connected polymers
reproduce the direct
kernel integral for
all four coarse sources.
Their coarse average
equals the original
partition function
\(4158837/4480000\),
so the normalization
has not been replaced.
Every activity satisfies
the product operator
majorant and every
disconnected selected
set factorizes over
its actual carriers.

Two input atoms share
the same family in
this regression. Their
one-defect activity is
\[
                         -113/16000000.
\]
Halving the defect
strength halves this
activity, rather than
dividing it by four.
The check therefore
detects the invalid
per-occurrence gain
that motivated this
section. The finite
kernel is only a
structural regression;
the SU(2) analytic
operator and cluster
estimates are proved
above.

V50 has 1,212,518 bytes and SHA--256
\begin{center}\scriptsize\ttfamily
2FDC7DD76440CC1807E1DD6775DB7D1761A3A4FCA967A6A410553480DB81B57B.
\end{center}
Its entire body is
preserved except for
the title version and
this terminal insertion.
V51 replaces V50 as
the single active
main note. Archives,
companions, supplied
notes, and monographs
remain unchanged.

\begin{firewall}
An all-depth compact
strong-coupling basin
has been proved here
for the smaller fixed
completion radius
\(\rho_\star=2^{-100}\).
The same theorem has
not been proved at
\(1/64\). No ultraviolet
entry, nontrivial
interacting continuum
theory, calibrated
physical transfer, or
positive physical
Hamiltonian mass gap
has been established.
The full Yang--Mills
goal remains unproved
and active.
\end{firewall}
\section{V52: a slice-local completed source, exact reflected-history reconstruction, and a calibrated strong-coupling lattice transfer}
\label{sec:v52}

\subsection{Direction, scope, and non-duplication}

This continuation, completed on 15 September 2026, addresses the
physical-time bottleneck exposed by V51. The smaller-radius
four-dimensional compact basin is retained, not replaced or silently
declared reflection positive. We introduce an additional, explicitly
different source map: apply the completed rooted block to spatial
links on each time slice and leave time unchanged. Its full history
law is the exact pushforward of the fine Wilson law. This gives an
exact reflected-transfer identification, whereas a one-step Markov
reset of coarse snapshots generally does not.

The recently supplied SM source-selection V17, regenerative
stationarity V21, and exact-action Einstein bridge V16 have been
reviewed in this lineage. Their useful constraint here is to retain
the actual selected law, physical clock, and compression map.
None supplies the nonlinear Yang--Mills physical-singlet estimate
missing from this programme. In particular, neither a quadratic
regenerative frequency nor an inverse of a compressed matrix is
inserted as the Yang--Mills Hamiltonian. Shell Mixing V16 and the
suggestions remain context; the one-loop calculation already
completed in V30 is not repeated. The V50 companion checkpoint
remains current; the next scheduled SM/NS/Shell review is V55.

The closure monograph already proves the abstract implication from
a complete mixing estimate, reflection positivity, a specified
physical clock, and a dense local algebra to a gap. It also proves
the fixed-one-step-contraction continuum-collapse warning.
Those abstract statements are not new results here. The additions
are an actual time-local completed source, its exact history
isometry, and a positive physical reflected norm for that source.
For orientation we give a self-contained pure-gauge transfer
construction and an explicit conservative small-coupling constant.
The existence of a lattice strong-coupling gap is classical, not
a claimed solution of the continuum problem.

\subsection{The actual fine Wilson transfer and its clock}

Fix a finite spatial torus
\(\Lambda_M=(\mathbb Z/M\mathbb Z)^3\), \(M\ge4\), and let
\(\mathcal Q_M=\SU(2)^{E_s}\) be its oriented spatial-link
configuration space, with one independent group per positive link.
Write \(dU\) for product normalized Haar measure. The local spatial
gauge action is
\[
 (g\cdot U)_{(x,\mu)}=g_xU_{(x,\mu)}g_{x+\hat\mu}^{-1}.
\]
Let \(\mathsf P_G\) be Haar averaging over all these vertex gauges,
and set
\[
 \mathcal H_G=\mathsf P_G L^2(\mathcal Q_M,dU).
\]
No division by a gauge-orbit volume or stabilizer is made.
Our normalization of the four-dimensional action is
\[
 S_0(U)=\sum_p(1-\chi(U_p)),\qquad
             \chi(W)=\tfrac12\operatorname{Re}\Tr W .
\]
Let \(S_s\) be the sum over spatial plaquettes on one slice, and,
for \(\beta>0\), define
\[
 k_\beta(W)=e^{-\beta(1-\chi(W))},\qquad
 (\mathsf M_\beta f)(U)=e^{-\beta S_s(U)/2}f(U),
\]
\[
 (\mathsf C_\beta f)(U)=
       \int\prod_{e\in E_s}k_\beta(U_eV_e^{-1})f(V)\,dV.
\]

\begin{proposition}[Fine positive transfer, including the ground law]
\label{prop:v52-fine-transfer}
The operator
\begin{equation}\label{eq:v52-T}
 \mathsf T_\beta=
       \mathsf M_\beta\mathsf C_\beta\mathsf P_G\mathsf M_\beta
\end{equation}
is compact and self-adjoint, is positive and injective on
\(\mathcal H_G\), and has a simple largest eigenvalue
\(\lambda_0>0\) with a continuous strictly positive
gauge-invariant eigenfunction \(h\), normalized by
\(\int h^2\,dU=1\). Its integral kernel is
\begin{equation}\label{eq:v52-K}
 \begin{split}
 K_\beta(U,V)
  ={}&e^{-\beta(S_s(U)+S_s(V))/2}\\
   &{}\times\int\prod_xdg_x
          \prod_{e\in E_s}k_\beta\bigl(U_e(g\cdot V)_e^{-1}\bigr).
 \end{split}
\end{equation}
The probability measure and transition
\begin{equation}\label{eq:v52-ground}
 \pi_\beta(dU)=h(U)^2\,dU,\qquad
 (\mathsf Pf)(U)=
   \frac{1}{\lambda_0h(U)}
           \int K_\beta(U,V)h(V)f(V)\,dV
\end{equation}
give the stationary reversible spatial-slice process of the
infinite-time Wilson law. On the physical space
\[
             \mathcal H_{\rm phys}=L^2(\pi_\beta)^G
\]
the transfer \(\mathsf P\) is positive and injective, the vacuum
is the simple eigenvector \(1\), and, for a specified physical
time spacing \(a>0\),
\begin{equation}\label{eq:v52-H}
 H_a=-a^{-1}\log\mathsf P,\qquad
          e^{-naH_a}=\mathsf P^n
\end{equation}
is a well-defined nonnegative self-adjoint Hamiltonian.
\end{proposition}

\begin{proof}
Let \(\chi_k\) denote the irreducible \(\SU(2)\) character with
highest weight \(k\) and dimension \(k+1\), so
\(\chi_1(W)=\Tr W=2\chi(W)\). The tensor-product identity
\[
       \chi_1^n=\sum_{k=0}^n m_{n,k}\chi_k,\qquad
       m_{n,k}\ge0,\quad m_{k,k}=1,\quad
       \sum_km_{n,k}(k+1)=2^n
\]
follows inductively from
\(\chi_1\chi_k=\chi_{k+1}+\chi_{k-1}\), with the last term
omitted at \(k=0\). Thus
\[
 \begin{split}
 k_\beta
    &=e^{-\beta}\sum_{n\ge0}
                     \frac{(\beta/2)^n}{n!}\chi_1^n
      =\sum_{k\ge0}a_k(\beta)\chi_k,\\
 a_k(\beta)
    &=e^{-\beta}\sum_{n\ge0}
                    \frac{(\beta/2)^n}{n!}m_{n,k}>0 .
 \end{split}
\]
Absolute uniform convergence follows from
\(|\chi_k|\le k+1\) and the dimension identity. Character
orthogonality gives convolution eigenvalue \(a_k/(k+1)>0\)
on every matrix-coefficient space. Taking products over the
finitely many links proves positivity and injectivity of
\(\mathsf C_\beta\). Simultaneous left and right gauge
transformations preserve its kernel, so it commutes with
\(\mathsf P_G\). The positive, bounded, invertible
multiplication \(\mathsf M_\beta\) is gauge invariant. Hence
\[
 \langle f,\mathsf T_\beta f\rangle
       =\langle\mathsf M_\beta f,
                    \mathsf C_\beta\mathsf M_\beta f\rangle>0
          \quad(0\ne f\in\mathcal H_G).
\]
On the unreduced space there is a kernel from
\(1-\mathsf P_G\); injectivity is asserted only on
\(\mathcal H_G\).

Formula \eqref{eq:v52-K} follows by performing the gauge average.
It is a continuous strictly positive symmetric kernel on a
compact product space, so its operator is compact. A normalized
maximizer of the top Rayleigh quotient can be taken real and
nonnegative: replacing a real function by its absolute value
does not decrease the quotient. Applying the strictly positive
kernel makes the eigenfunction continuous and strictly positive.
If there were another top eigenfunction orthogonal to it, a
real such eigenfunction would have positive and negative sets
of positive measure. Replacing it by its absolute value would
strictly increase its quotient, a contradiction. This also
proves simplicity. The range of \(\mathsf T_\beta\) is gauge
invariant, so \(h\) is gauge invariant.

For \(e=(x,\mu)\), the argument in the temporal kernel is
\[
              U_e g_{x+\hat\mu}V_e^{-1}g_x^{-1},
\]
exactly the holonomy around the corresponding temporal
plaquette of one slab. The \(g_x\)'s are its temporal links.
Multiplying slab kernels combines the two spatial half-actions
into the full spatial action on each intervening slice.
Temporal periodic closure is the trace of a power of the
transfer. Sending the temporal length to infinity gives its
simple top-eigenvector state; equivalently the successive
kernel factors telescope to the transition in
\eqref{eq:v52-ground}. One may take this limit first in finite
temporal cylinders with positive boundary weights; compactness
and the simple isolated top eigenvalue give the same local limit.
This verifies the actual Wilson law, including its scalar
normalization, rather than only a formal transfer symbol.

The eigenfunction equation gives \(\mathsf P1=1\), and symmetry
of \(K_\beta\) gives detailed balance for \(\pi_\beta\).
Multiplication by \(h\) is a unitary
\[
 L^2(\pi_\beta)^G\longrightarrow\mathcal H_G,\qquad
             f\longmapsto hf,
\]
and conjugates \(\mathsf P\) to
\(\mathsf T_\beta/\lambda_0\). Its spectrum is in \([0,1]\)
and its zero eigenspace is trivial on the physical space.
Spectral calculus therefore defines the densely defined
logarithm in \eqref{eq:v52-H}; zero may be a spectral
accumulation point, but has no eigenprojection there.
\end{proof}

The standard transfer-positivity background is
\href{https://bib-pubdb1.desy.de/record/396349/files/7611148.pdf}
{M.~L\"uscher, Construction of a Selfadjoint, Strictly Positive
Transfer Matrix for Euclidean Lattice Gauge Theories,
DESY 76/54; Commun.\ Math.\ Phys.\ 54 (1977), 283--292}.
The proof above spells out the pure-gauge specialization and
the selected ground law needed below. No uniform-in-volume
upper or lower bound on \(h\) has been assumed.
The exclusion of \(\beta=0\) in the Hamiltonian assertion matters:
at zero coupling the physical transfer has rank one, so its
logarithm does not define a densely defined finite-energy
Hamiltonian on the full nonconstant time-zero space.

\subsection{A spatial companion completion, not a new time step}

Let \(M=2^mN\), \(N\ge4\). On each spatial slice use the
three-dimensional rooted block with cells of side two, root
at \(2x\), and a fixed translated cell tree. There are eight
chronological two-link paths for each coarse spatial edge:
two duplicate root paths \(B_e\) and three pairs
\(P_{a,\ell}\), \(1\le a\le3\), \(1\le\ell\le2\).
Use the same smooth cutoff logarithm \(g_\rho\) as in V51 and
define the completed output by
\begin{equation}\label{eq:v52-spatial-map}
 C_e=B_e\exp\left\{\frac18
             \sum_{a=1}^3\sum_{\ell=1}^2
                           g_\rho(B_e^{-1}P_{a,\ell})\right\},
            \qquad \rho=\rho_\star=2^{-100}.
\end{equation}
Denote this map by \(\mathcal B_{\rm sp}\), and its \(m\)-fold
composition by \(\mathcal B_{\rm sp}^{\,m}\). The pivot
construction is the inherited compact completion, now with
root weight \(1/4\) instead of \(1/8\); the positive lower
Hessian bound is consequently stronger. It is smooth,
globally defined, and gauge covariant at the retained roots.

For the actual fine stationary process set
\begin{equation}\label{eq:v52-observe}
                    C_t=\mathcal B_{\rm sp}^{\,m}(U_t),
                         \qquad t\in\mathbb Z .
\end{equation}
There is no use of neighboring time slices in this formula.
Consequently it commutes exactly with site reflection
\(t\mapsto-t\) and with every integer time translation.
The observed spatial spacing is \(2^ma\), but the time spacing
is still \(a\). No factor \(2^m\) is inserted into
\eqref{eq:v52-H}. This is an anisotropic observation family on
the original theory, not a claimed four-dimensional coarse
Wilson action or the four-dimensional pushforward of V51.

The following auxiliary Haar statement will be used only to
prove a positive norm for an actual physical source.

\begin{lemma}[Haar noncollapse for the spatial companion]
\label{lem:v52-spatial-haar}
For product Haar fine spatial links on a torus of side
\(2^mN\), \(N\ge4\), the exact \(m\)-fold pushforward under
\(\mathcal B_{\rm sp}\) has a continuous gauge-invariant
connected interaction with
\[
 \|\Phi_m^{\rm Haar}\|_{(3)}
 :=\sup_e\sum_{\alpha:e\in S_\alpha}
          e^{3|S_\alpha|}\|\Phi_{m,\alpha}^{\rm Haar}\|_\infty
 \le \frac{1-(5/8)^m}{12000}.
\]
In particular, for the normalized trace \(F_m\) of any
coarse elementary spatial plaquette, viewed as a function of
the fine links,
\begin{equation}\label{eq:v52-haar-var}
                   \Var_{\rm Haar}(F_m)>\frac15 .
\end{equation}
\end{lemma}

\begin{proof}
We verify the dimension-dependent constants needed to apply
the grouped-family proof of V51, rather than assuming that its
four-dimensional counts are unchanged. Retain the original
independent spatial Haar links. A cell has twelve internal
links, seven of which form the tree; an outgoing boundary
family has four links, one pivot and three nonpivots.
Conditionally on all internal links and \(C_e\), the three
nonpivot groups are independent Haar groups. Each of its
two chronological products is a left and right translate of
the corresponding nonpivot group.

The global displacement is at most \(3\rho/4\). An active
source is within \(7\rho/4<15\rho/8\) of \(C_e\).
The Haar-cap estimate of V48 therefore bounds the exposure
probability of a given pair by
\(p_\rho=2\rho^3\). Exposure events of different pairs are
independent under this conditional Haar reference. If exactly
\(r\) pairs are exposed, the displacement improves to
\(r\rho/4\), and the same Hessian and exponential-volume
calculation as in Lemma \ref{lem:v48-jacobian} gives
\begin{equation}\label{eq:v52-spatial-J}
 J_r^{(3)}
  =\left(1-\frac r4-\frac{r^2\rho^2}{64}\right)^{-3}
   \left(1-\frac{r^2\rho^2}{192}\right)^{-2},
                   \qquad 1\le r\le3 .
\end{equation}
The first denominator is the Hessian lower bound:
the active fraction is \(2r/8\), and the displacement
contribution is \((r\rho/4)^2/4\). The second factor is the
inverse exponential-volume bound at that displacement.
On no exposure the pivot is exactly \(C_e\) and its inverse
Haar Jacobian \(j_e\) is one. For \(\rho\le1/64\),
\[
                 J_1^{(3)}<3,\qquad J_3^{(3)}<65 .
\]
In particular the four-dimensional estimate \(J_1<2\)
must not be copied into this calculation.
Since \(p_\rho\le1/131072\), integration over the three
nonpivot groups gives
\[
 \delta_\rho^{(3)}:=
       \sup_{I,C_e}\int |j_e-1|\,dY_e
 \le6p_\rho+65(3p_\rho^2+p_\rho^3)<7p_\rho .
\]
For the difference \(\mathsf D_e\) of the unnormalized
completed-family operator and its root-family operator,
split \(j_eF(B_{\rm solved})-F(C_e)\) into the Jacobian
term and the substitution term. The latter vanishes off
the union of the three exposure events and has norm at
most \(2\|F\|_\infty\). Hence
\begin{equation}\label{eq:v52-spatial-defect}
 \|\mathsf D_e\|_{\infty\to\infty}
    \le\delta_\rho^{(3)}+6p_\rho
        <26\rho^3\le d_\rho,\qquad d_\rho:=44\rho^3 .
\end{equation}
There is one operator per family acting on the entire
input product. No small factor is allocated separately to
repeated occurrences of the same nonpivot group.

A defect footprint consists of both endpoint cells'
internal links and their intervening boundary family:
it has \(28=2\cdot12+4\) edges and \(16\) vertices.
It is connected, contains its complete root pair, and
contains no other complete boundary root pair. At most
six such footprints meet a fine vertex. With charge
\((13/5)\) per footprint edge and \(1/5\) per footprint
vertex, its inflated weight is bounded by
\[
                     d_\rho e^{(13/5)28+(1/5)16}
                          =d_\rho e^{76}
                          \le d_\rho e^{968/5}.
\]

Expand each input factor as \(1+(e^{-\Phi_\alpha}-1)\),
and each boundary operator as its root operator plus
\(\mathsf D_e\). Connected components are defined by
overlap of fine vertex carriers. Disjoint components
share neither independent Haar links nor halves of the
same conditional root pair. Thus their integrals
factor exactly, and each connected activity is bounded
by the product of its input norms and family-defect
norms. For an unselected family, a lone half of its
root pair integrates out its coarse source. For a
selected defect, the whole root pair is included in
its footprint. Consequently the activity depends only
on complete root pairs in its union carrier.
Mixed maps that choose the completed or root output
edge by edge are gauge covariant, proving termwise
gauge invariance. These are precisely the component
and support assertions required by
Lemma \ref{lem:v51-mixed}; the argument works here
without conditional independence being imposed after
the internal variables have been integrated out.

For an input interaction of norm \(B\le1/320\), its
inflated atom sum at a vertex is at most \(12B\le16B\),
because spatial fine degree is six and
\(\|e^{-\Phi}-1\|_\infty\le2\|\Phi\|_\infty\).
V49's two-sum lemma with \(a=b=1/10\) thus applies
with the conservative budget
\[
             Q\le16B+8d_\rho e^{968/5}<1/10
                   \quad\text{at }\rho=\rho_\star.
\]
For a logarithmic cluster let \(\ell\) be the sum of
input carrier sizes plus \(28\) per defect occurrence.
The pinned cluster sum has weight \(e^{(13/5)\ell}\).
After subtracting a cluster's coarse Haar mean, its
nonconstant gauge-invariant support must contain a
cycle. Coarse side at least four implies at least four
coarse edges, hence at least eight fine edges and
\(\ell\ge8\).

For completeness, the rounding maps here are
\[
 R_\epsilon z=\left\lfloor\frac{z+\epsilon}{2}\right\rfloor
             \pmod N,\qquad \epsilon\in\{0,1\}^3 .
\]
Each fine edge survives for four of the eight offsets.
The connected union therefore has, for some single
offset, a connected coarse image of size at most half
its fine edge size, containing all its complete root
pairs. For a fixed output edge there are at most four
fine preimages per offset, hence \(32\le128\) choices.
Subtracting the mean costs at most two. The same
spare exponential and pinned sum as in V51 give
\[
 \begin{split}
 B'&\le256e^{-44/5}
                  (16B+8d_{\rho_\star}e^{968/5})\\
   &<\frac58B+\frac1{32000}.
 \end{split}
\]
All sums are absolute; scalar means remain in the
normalizer. This proves the required three-dimensional
invariant ball with the conservative V51 constants.
Starting with the actual Haar input \(B_0=0\) and
iterating the exact pushforward proves the displayed
bound for \(B_m\).

Condition on every coarse link except one plaquette
edge. Its interaction has oscillation at most \(2B_m\),
so its conditional density is at least \(e^{-2B_m}\)
relative to Haar. The normalized trace of its plaquette
has Haar mean zero and variance \(1/4\), whatever the
fixed other three links. Using
\(\Var X=\inf_c\mathbb E|X-c|^2\) and then total
variance gives
\[
             \Var_{\rm Haar}(F_m)
                    \ge\frac14 e^{-2B_m}>\frac15.
\]
Here the variance on the left is for the pullback
to the original fine Haar law, equivalently for the
exact completed pushforward, not for a reset Haar
law on the coarse links.
\end{proof}

\subsection{The exact physical space of complete observed histories}

Let \(\omega_\beta\) denote the stationary law in
Proposition \ref{prop:v52-fine-transfer}.
Let \(\mathcal A_{m,+}\) be the linear span of bounded
continuous functions of finitely many
\(C_0,\ldots,C_r\) that are gauge invariant independently
on each slice. Products of gauge-invariant single-slice
cylinder observables suffice. Let \(\theta\) replace
each time \(t\) by \(-t\), and define the antilinear
reflection \(\Theta F=\overline{F\circ\theta}\).
Time translation is
\[
   (\tau_nF)(C_0,C_1,\ldots)=F(C_n,C_{n+1},\ldots).
\]

\begin{theorem}[History isometry and the unchanged physical transfer]
\label{thm:v52-history}
For every \(\beta>0\), finite spatial volume, and finite
spatial blocking depth \(m\), define
\begin{equation}\label{eq:v52-j}
   j_mF(U_0)=
       \mathbb E_{\omega_\beta}
          [F(C_0,\ldots,C_r)\mid U_0].
\end{equation}
Then
\begin{equation}\label{eq:v52-OS}
 \omega_\beta((\Theta F)G)
             =\langle j_mF,j_mG\rangle_{L^2(\pi_\beta)}
             \quad(F,G\in\mathcal A_{m,+}),
\end{equation}
and
\begin{equation}\label{eq:v52-shift}
                      j_m(\tau_nF)=\mathsf P^n j_mF .
\end{equation}
Thus the reflected quotient and completion of the
observed history algebra is isometric to
\[
 \mathcal H_m=
       \overline{\{j_mF:F\in\mathcal A_{m,+}\}}
                        ^{\,L^2(\pi_\beta)}
             \subseteq\mathcal H_{\rm phys}.
\]
This is a reducing subspace for \(\mathsf P\), contains
the vacuum \(1\), and its reconstructed transfer and
Hamiltonian are exactly
\[
                \mathsf P_m=\mathsf P|_{\mathcal H_m},
                  \qquad H_{a,m}=H_a|_{\mathcal H_m}.
\]
In particular the time unit is the original \(a\).
For a single-slice observable \(F(C_0)\),
\begin{equation}\label{eq:v52-os-variance}
 \|[F-\omega_\beta(F)]\|_{\rm OS}^2
             =\Var_{\pi_\beta}
                    \bigl(F(\mathcal B_{\rm sp}^{\,m}U_0)\bigr).
\end{equation}
\end{theorem}

\begin{proof}
The fine two-sided stationary Markov process has past
and future conditionally independent given \(U_0\).
By reversibility, its conditional past has the same
transition law as its conditional future. Both
statements remain valid when the measured functions
are pulled back through a map of their own time slice.
Consequently conditioning the product
\((\Theta F)G\) on \(U_0\) gives
\(\overline{j_mF(U_0)}j_mG(U_0)\), proving
\eqref{eq:v52-OS}. Gauge covariance of the spatial
map and gauge invariance of the kernels make \(j_mF\)
gauge invariant. Applying the Markov property first
from time zero to time \(n\), and then to the
remaining history, gives \eqref{eq:v52-shift}.

Equation \eqref{eq:v52-OS} identifies the null space
with \(\ker j_m\), so \(j_m\) extends to an isometry
onto exactly the closed space \(\mathcal H_m\).
Equation \eqref{eq:v52-shift} makes this space
invariant under \(\mathsf P\). A closed invariant
subspace of a bounded self-adjoint operator is
reducing: if \(v\) is orthogonal to the subspace,
then \(\langle\mathsf Pv,w\rangle
=\langle v,\mathsf Pw\rangle=0\) for every \(w\)
in the subspace. Spectral calculus therefore
restricts the logarithm as asserted, with domain
\(D(H_a)\cap\mathcal H_m\). The constant history
maps to \(1\). Finally for a time-zero function
the conditional expectation in \eqref{eq:v52-j}
is the function itself; subtracting its expectation
gives \eqref{eq:v52-os-variance}.
\end{proof}

The use of site reflection and inclusion of time zero
are explicit. If one uses only strictly positive
integer times, the resulting closed image is unchanged
at these finite cutoffs: it contains
\(\mathsf P\mathcal H_m\), which is dense in
\(\mathcal H_m\) because \(\mathsf P|_{\mathcal H_m}\)
is self-adjoint and injective.

There is no assertion that
\(\mathcal H_m=\mathcal H_{\rm phys}\) for \(m>0\).
That would require a separate completeness statement
for the observed sources. Nor is
\(\mathcal H_m\) identified with the space of
functions of \(C_0\) alone. The complete histories
are essential to the reducing-subspace conclusion.

\begin{proposition}[A positive reversible example where snapshots lose memory]
\label{prop:v52-memory}
A strictly positive reversible four-state transition
can have a two-state deterministic observation whose
history reflected space has dimension four, and whose
two-step correlations are not those obtained by
iterating its one-step snapshot conditional kernel.
\end{proposition}

\begin{proof}
In the uniform four-state probability space use the
orthonormal functions
\[
 {\bf1}=(1,1,1,1),\quad g=(1,1,-1,-1),\quad
 h_1=(1,-1,1,-1),\quad k=(1,-1,-1,1).
\]
Here \(g h_1=k\) pointwise. Put
\[
 u=w=\tfrac14,\qquad v=z=\tfrac18,\qquad
 P_{ij}=\tfrac14\{
   1+u g_i g_j+v(g_i(h_1)_j+(h_1)_i g_j)
                      +w(h_1)_i(h_1)_j+z k_i k_j\}.
\]
Every entry is strictly positive; rows sum to one
and the matrix is symmetric. In the displayed
orthonormal basis its nonconstant block is
\[
                    \begin{pmatrix}u&v&0\\v&w&0\\0&0&z\end{pmatrix}.
\]
Thus its eigenvalues are \(1,3/8,1/8,1/8\).
Observe only \(Y_t=g(X_t)\in\{-1,1\}\).
For histories \(1,Y_0,Y_1,Y_0Y_1\), their conditional
images at \(X_0\) are respectively
\[
                   1,\quad g,\quad ug+v h_1,\quad u1+v k .
\]
Their reflected Gram matrix is
\[
 \begin{pmatrix}
 1&0&0&1/4\\
 0&1&1/4&0\\
 0&1/4&5/64&0\\
 1/4&0&0&5/64
 \end{pmatrix},
                 \qquad \det=1/4096>0.
\]
The history reflected space is therefore all four
dimensions, despite there being only two snapshot
states. The one-step conditional kernel of \(Y\)
has nonconstant eigenvalue \(u=1/4\). Its iterated
two-step correlation would be \(u^2=1/16\), whereas
the actual correlation is
\[
          \mathbb E[Y_0Y_2]
                 =\langle g,P^2g\rangle=u^2+v^2=5/64.
\]
The difference \(1/64\) proves the claimed memory.
\end{proof}

This example is an exact finite regression for the
history distinction. It is not a model from which
an \(\SU(2)\) estimate is extrapolated.

\subsection{An explicit gap for the fine physical lattice transfer}

\begin{theorem}[Spatial-volume-uniform strong-coupling lattice gap]
\label{thm:v52-lattice-gap}
For the fine isotropic \(\SU(2)\) Wilson normalization
above, if
\[
                            0<\beta\le1/72,
\]
then, on every finite spatial torus of side \(M\ge4\),
\begin{equation}\label{eq:v52-fine-gap}
 \|\mathsf P|_{1^\perp}\|_{\mathcal H_{\rm phys}}
                         \le\frac12,\qquad
 H_a|_{1^\perp}\ge\frac{\log2}{a}.
\end{equation}
The same bound holds for the restriction to
\(\mathcal H_m\ominus\mathbb C1\), for every finite
spatial blocking depth admitted above.
\end{theorem}

\begin{proof}
Work in the unfixed four-dimensional fine-link Gibbs
specification on the spatial torus and infinite
integer time. A link belongs to six plaquettes;
each such plaquette involves three other links.
Use link-midpoint \(\ell^1\) distance, with periodic
spatial coordinates and unbounded time. Other
links in a common plaquette have distance at most
one.

After dropping its scalar, each plaquette potential
is \(-\beta\chi(U_p)\) with norm at most \(\beta\).
Changing one neighboring link changes the
single-link log density by a function with
oscillation at most \(4\beta\) per common
plaquette. The elementary tilted-probability
bound used in V48 gives total-variation
influence at most \(\tanh\beta\le\beta\)
per such plaquette. Summing if more than one
plaquette contains the pair is valid. For the
single-link influence matrix \(A\), at
\(\mu=\log2\),
\begin{equation}\label{eq:v52-dobrushin}
          \sup_e\sum_f A_{ef}e^{\mu d(e,f)}
                     \le6\cdot3\cdot2\,\beta
                     =36\beta\le\frac12.
\end{equation}

The coupling-resolvent proof used in V48 applies
directly to this fine specification. In detail,
propagating a boundary disagreement gives the
series of influence paths
\(\sum_{r\ge0}A^r=(I-A)^{-1}\).
The triangle inequality for midpoint distance
and \eqref{eq:v52-dobrushin} bound its weighted
row sums by \(2\), uniformly in spatial volume
and temporal cylinder length. Boundary
influences tend to zero as their distance
increases, giving uniqueness of the
infinite-time Gibbs law. The same comparison
applied to a bounded local tilt yields
covariance decay between any two bounded
local functions, with a finite prefactor
depending on their supports and sup norms,
and decay \(e^{-\mu d}\). Approximation by
finite cylinders proves the assertion for
the infinite-time law. The Perron law of
Proposition \ref{prop:v52-fine-transfer}
is a Gibbs law for this specification on
gauge-invariant slice observables, hence is
this unique marginal law.

For a bounded continuous gauge-invariant function
\(f\) of the spatial slice, write
\(f_0=f-\pi_\beta f\). Its fine link support is
finite for each \(M\), even if it is the whole
spatial slice. All links of its translate by
\(2n\) time units have time separation \(2n\).
The preceding bound therefore gives
\begin{equation}\label{eq:v52-correlation}
 0\le\langle f_0,\mathsf P^{2n}f_0\rangle_{\pi_\beta}
       =\operatorname{Cov}_{\omega_\beta}
                 (\overline{f(U_0)},f(U_{2n}))
       \le K_f e^{-2\mu n},
\end{equation}
where \(K_f<\infty\) is independent of \(n\).
No volume-uniform bound on \(K_f\) is required
for the following spectral conclusion.

Let \(\nu_f\) be the positive spectral measure
of \(\mathsf P\) at \(f_0\). If it assigned
positive mass to \([r,1]\) for some
\(r>e^{-\mu}=1/2\), its \(2n\)-th moment
would be at least
\(r^{2n}\nu_f([r,1])\), contradicting
\eqref{eq:v52-correlation} as \(n\to\infty\).
Thus \(\nu_f((1/2,1])=0\).
Bounded continuous gauge-invariant functions
are dense in \(L^2(\pi_\beta)^G\), by
continuous approximation followed by gauge
averaging. The spectral projection onto
\((1/2,1]\) consequently vanishes on
\(1^\perp\). Positivity of \(\mathsf P\)
proves its asserted norm bound, and
spectral calculus gives the Hamiltonian
bound. Finally \(\mathcal H_m\) is a
reducing subspace containing \(1\), by
Theorem \ref{thm:v52-history}, so the
same inequalities restrict to its
nonvacuum part.
\end{proof}

This proof identifies a physical lattice transfer,
not an auxiliary Gibbs sampler. It uses the complete
fine Wilson specification, not the output interaction
of a guessed coarse Markov chain. The estimate is
uniform in spatial volume but is stated for the
finite-volume Hamiltonians actually constructed
above. No infinite-spatial-volume injectivity or
continuum Hamiltonian is inferred solely by taking
a limit of these finite-volume logarithms.

\subsection{A nonzero physical source at each fixed spatial depth}

\begin{lemma}[A local density comparison for the actual Perron law]
\label{lem:v52-density}
For any \(\beta>0\) and any finite set \(S\) of \(k\)
spatial fine links on one time slice, the marginal
of the actual law \(\pi_\beta\) satisfies
\begin{equation}\label{eq:v52-density}
                   (\pi_\beta)_S\ge e^{-12\beta k}\,dU_S
\end{equation}
as measures. The constant is independent of spatial
volume and of the other links.
\end{lemma}

\begin{proof}
In a finite four-dimensional Wilson cylinder,
condition on all links except one link \(e\).
The six plaquettes containing \(e\) each have
action contribution in \([0,2]\). The
single-link log-density oscillation is
therefore at most \(12\beta\), giving a
conditional density at least \(e^{-12\beta}\)
relative to normalized Haar. Conditioning
on fewer of the other links takes mixtures
of these conditional probabilities and
preserves this inequality. Enumerate \(S\)
and apply the inequality successively,
conditioning each link on its predecessors.
The product lower bound is
\(e^{-12\beta k}\).
It passes to the infinite-time Perron
marginal by integration against nonnegative
continuous functions. This argument is
valid without a mixing assumption and
without any bound on the pointwise
ground eigenfunction \(h\).
\end{proof}

\begin{theorem}[Fixed-depth noncollapse in the actual physical reflected norm]
\label{thm:v52-physical-source}
Let \(F_m(C_0)\) be the normalized trace of one
elementary spatial plaquette after \(m\) completed
spatial blocks, with fine side \(M=2^mN\),
\(N\ge4\). For every \(\beta>0\),
\begin{equation}\label{eq:v52-source-lower}
 \|[F_m-\omega_\beta(F_m)]\|_{\rm OS}^2
                >\frac15\exp\{-144\beta\,8^m\}>0.
\end{equation}
The bound is uniform in \(N\) for fixed \(m,\beta\).
For \(0<\beta\le1/72\) this is a nonzero
physical vector in a reducing history space whose
Hamiltonian has the gap in
\eqref{eq:v52-fine-gap}.
\end{theorem}

\begin{proof}
Put \(L=2^m\). A one-step completed spatial
edge depends only on internal links of its
two endpoint cells and the boundary family
between them. Inductively, the pullback
of a blocked plaquette is supported on
links with both endpoints in the four
spatial cells of side \(L\) at its four
vertices, including the faces joining
these cells. Indeed an internal finer
edge has both endpoint subcells in its
parent cell; an edge across a parent
face has one endpoint subcell in each
of the two adjacent parent cells.
Iterating cannot introduce a third
parent cell for that edge. Taking
the four plaquette edges proves the
stated four-cell containment.

There are \(4L^3\) fine vertices in
that union. With three positive
spatial links per vertex, a carrier
\(S_m\) for the pullback therefore
satisfies
\[
                        |S_m|\le12L^3=12\,8^m.
\]
Use Lemma \ref{lem:v52-density} on
this carrier. For every constant \(c\),
\[
 \mathbb E_{\pi_\beta}|F_m-c|^2
        \ge e^{-12\beta|S_m|}
                         \mathbb E_{\rm Haar}|F_m-c|^2 .
\]
Taking the infimum over \(c\), then
applying Lemma \ref{lem:v52-spatial-haar},
gives
\[
 \Var_{\pi_\beta}(F_m)
        \ge e^{-12\beta|S_m|}\Var_{\rm Haar}(F_m)
        >\tfrac15e^{-144\beta8^m}.
\]
Theorem \ref{thm:v52-history} identifies
this actual time-zero variance with
the physical reflected norm. Haar
was used only as a comparison measure
in the displayed inequality; the
selected state is still \(\pi_\beta\).
\end{proof}

\subsection{What has and has not crossed the physical interface}

There are three distinct conclusions, with different
parameter ranges. For every \(\beta>0\), the
slice-local history reconstruction is exactly a
restriction of the actual fine physical transfer.
For every fixed \(m\), the blocked spatial
plaquette has a strictly positive physical
reflected norm, uniformly in spatial volume.
For \(0<\beta\le1/72\), that restricted
physical lattice Hamiltonian has gap at
least \((\log2)/a\). These statements do
not require the fine law to be reset to
Haar at any step.

The original V51 map blocks time as well as
space. This section proves no reflection
positivity or transfer intertwining for
that four-dimensional completed law.
It instead supplies a time-local companion
source on the same fine Wilson theory.
A relation strong enough to transfer
continuum conclusions between the two
source constructions remains to be proved.

The lower bound \eqref{eq:v52-source-lower}
decays with \(m\). It is not a uniform
lower frame bound along increasing
depth, and it does not establish
cofinal physical source noncollapse.
Even the possibly larger history
space \(\mathcal H_m\) has not been
shown to exhaust the full fine
physical space.

Finally, the fine time spacing in this
construction is \(a\), not the coarse
spatial spacing \(2^ma\). Holding a
dimensionless contraction bounded
strictly below one as \(a\to0\) does
not create a finite nontrivial
continuum Hamiltonian sector. The
closure monograph's already proved
clock-collapse criterion applies:
fixed-positive-time centered
correlations then collapse, unless
the proposed limiting nonvacuum
entrance fails its noncollapse or
continuity requirements. The
large quantity \((\log2)/a\) must
not be advertised as a continuum
mass-gap solution.

The upstream target is therefore a
nontrivial continuum trajectory with
physical-scale reflected source
control, not another proof that
an uncalibrated compact heat bath
contracts. The slice-local
construction removes a reflection
ambiguity for one concrete source
family; it does not construct the
ultraviolet trajectory, prove the
continuum axioms, or provide
cutoff-uniform physical
noncollapse along that trajectory.

\subsection{Exact checks and provenance}

The executable certificate is
\begin{center}\small
\path{scripts/verify_ym_f14_v52_slice_time_bridge.py}.
\end{center}
It uses exact rational arithmetic
and no external numerical package.
It checks the character-tensor
identities through power eighteen,
the three-dimensional Jacobian
bounds, exposure and defect
budgets, footprint charge \(76\),
and the fine influence value
\(36(1/72)=1/2\). The character
positivity proof for all powers
is the induction above, not
an extrapolation of these finite
checks.

The geometric recursion gives
pullback-carrier sizes \(4,64,640,5632\)
at depths \(0,1,2,3\), respectively,
all within \(12\,8^m\), and verifies
containment in the four parent
cells at each tested depth.
The all-depth containment proof
is given above. In the four-state
example, the certificate sums
the reflected path probabilities
directly, both at zero and at
one delayed time step, and
checks equality with the
conditional-image Gram formulas.
It verifies determinant
\(1/4096\), the positive transfer
eigenvalues, and the snapshot
memory discrepancy \(1/64\).

V51 has 1,236,214 bytes and SHA--256
\begin{center}\scriptsize\ttfamily
3AF95484722F8D364BFA400E0D3408FA2400186DB88D2962639DFB8C2EE53310.
\end{center}
Its complete body is preserved,
apart from the title version
and this terminal insertion.
V52 replaces V51 as the one
active main research note.
The archived lineages,
companion programmes, supplied
notes, and monographs remain
unchanged.

\begin{firewall}
A calibrated physical
\emph{lattice} transfer and
an explicit strong-coupling
lattice gap have now been
identified for the slice-local
companion history construction.
This advances the physical
interface beyond V51's compact
spatial estimates. No full
four-dimensional interacting
continuum Yang--Mills theory,
nor its positive continuum
Hamiltonian mass gap, has
been proved. The full
Yang--Mills goal remains
unproved and active.
\end{firewall}
\section{V53: the actual Wilson time-zero interaction, all-depth physical noncollapse, and a concrete continuum-time obstruction}
\label{sec:v53}

\subsection{Why the next step is the actual static marginal}

V52 identifies the spatial companion's complete history space
with a reducing space of the actual physical lattice transfer.
Its positive reflected norm bound, however, deteriorates as
\(\exp(-144\beta8^m)\) with spatial blocking depth \(m\).
The present continuation, dated 15 September 2026, removes that
depth dependence in an explicit small-\(\beta\) region. The new
input is a complete connected interaction for the actual
Wilson time-zero marginal, obtained by integrating all other
times. It is not a three-dimensional Wilson reset and is
not a substitute ground state.

The resulting direction change is important. A uniform
nonzero physical source and V52's dimensionless contraction,
taken together, give a discontinuous physical-time limit
as the isotropic fine spacing tends to zero. Thus the
strong-coupling basin does not merely lack a proof of
continuum noncollapse: for this explicit noncollapsed
source family it fails the required time continuity.
The programme must leave that fine-lattice regime to
construct a nontrivial strongly continuous continuum
sector. Improving the same terminal basin further
cannot by itself supply the ultraviolet construction.

The finite temporal Wilson-to-Hamiltonian construction
is already retained in the closure monograph's account
of f1 V10--V11. It is not repeated here. Likewise the
abstract clock-collapse proposition is already proved
there. The new work is the actual Wilson marginal
estimate and its application to the concrete completed
source, closing the possible escape through vanishing
time-zero variance in this small-coupling regime.
The next companion-programme checkpoint remains V55.

\subsection{A conditional bulk integration with all source terms retained}

Use the action and physical normalization of V52:
\[
 S_0(U)=\sum_p(1-\chi(U_p)),\qquad
                 \chi(W)=\tfrac12\operatorname{Re}\Tr W.
\]
For the moment work on
\[
  (\mathbb Z/M\mathbb Z)^3\times(\mathbb Z/T\mathbb Z),
                         \qquad M,T\ge4.
\]
Retain exactly the spatial links at time zero,
denoted \(U^{(0)}\), and integrate every other fine
link, including all temporal links, with product
Haar reference. Write \(E_0\) for the retained edge
set. After canceling the field-independent Wilson
action constant, the unnormalized retained density is
\begin{equation}\label{eq:v53-Z}
 Z_{M,T}(U^{(0)})
       =\int\prod_p(1+f_p(U))\,dU_{E\setminus E_0},
       \qquad f_p=e^{\beta\chi(U_p)}-1 .
\end{equation}
The actual marginal is \(Z_{M,T}\,dU^{(0)}\)
divided by its Haar integral. In particular, no
source-dependent part of the bulk integral in
\eqref{eq:v53-Z} has been replaced by a constant.

For a plaquette atom \(p\), let \(R_p\) be its
four fine vertices and \(S_p\) its four fine edges.
Join two atoms if their vertex carriers intersect.
For a connected set \(\gamma\) of distinct atoms define
\[
 z_\gamma(U^{(0)})
       =\int\prod_{p\in\gamma}f_p(U)\,
                                      dU_{E\setminus E_0}.
\]
Only the hidden links on its carrier need be
integrated in this formula; unused Haar factors
have mass one.

\begin{lemma}[Exact source-dependent polymer representation]
\label{lem:v53-polymers}
The expression \eqref{eq:v53-Z} is exactly the
hard-core polymer partition function with activities
\(z_\gamma\), where compatibility means disjoint
fine vertex carriers, including self-incompatibility.
For \(0<\beta\le1\),
\[
                \|z_\gamma\|_\infty\le(2\beta)^{|\gamma|}.
\]
Each activity is a continuous gauge-invariant
function of the retained edges in
\(\bigcup_{p\in\gamma}S_p\cap E_0\).
Its source support need not itself be connected.
\end{lemma}

\begin{proof}
Expand the finite product in \eqref{eq:v53-Z}.
The connected components of a selected plaquette
set have disjoint fine vertex sets and hence
disjoint hidden Haar variables. Their integrals
factor, even though the retained sources are
fixed. Summing these components is exactly the
asserted polymer partition identity.
The bound follows from
\[
              \|f_p\|_\infty\le e^\beta-1\le2\beta .
\]
Gauge transformations at time-zero vertices act
on the retained spatial links in the usual
three-dimensional manner. Apply the same fine
gauge transformation to every incident hidden
link as a Haar change of integration variable.
Each plaquette factor is gauge invariant and
the hidden measure is unchanged. This proves
gauge invariance of the activity in its retained
arguments. Continuity follows from compact
integration. The same argument applies to
each complete logarithmic cluster term.
\end{proof}

We use the already proved two-sum estimate of
Lemma \ref{lem:v49-cell-cluster}, now with a
fine vertex as its indexing site. The independence
and exact component identity needed for that
estimate have just been checked for this different
conditional law. Its standard pinned-cluster
input is the Koteck\'y--Preiss bound, recorded
in equations (2.14)--(2.15) of
\href{https://arxiv.org/html/math-ph/0605041v2}
{Fern\'andez--Procacci, Cluster expansion for abstract
polymer models}. No positivity of individual
activities is required.

Each fine vertex belongs to exactly twenty-four
plaquettes on the stated tori. Put \(a=b=1/10\),
assign charge \(13\) to each plaquette occurrence,
and define
\begin{equation}\label{eq:v53-Q}
                   Q_\beta=48\beta e^{69/5}.
\end{equation}
Indeed the inflated atom sum is bounded by
\[
 24(2\beta)\exp\{13+(a+b)4\}
                     =48\beta e^{69/5}.
\]
For
\begin{equation}\label{eq:v53-beta}
                          0<\beta\le10^{-12},
\end{equation}
the explicit estimates are
\begin{equation}\label{eq:v53-budgets}
 e^{69/5}<10^6,\qquad
 Q_\beta<\frac{48}{10^6}<\frac1{10},
       \qquad 10Q_\beta<\frac1{2000}.
\end{equation}
The exponential inequality is certified by
a rational Taylor sum with a geometric tail.

Let \(\mathcal C\) denote a complete logarithmic
cluster, and let \(k(\mathcal C)\) count every
plaquette occurrence, including repetitions
coming from the logarithm. Write
\(\Psi_\mathcal C(U^{(0)})\) for its contribution
to the negative logarithm of \eqref{eq:v53-Z}.
The two-sum estimate gives
\begin{equation}\label{eq:v53-pinned}
 \sup_{v}\sum_{\mathcal C:\,v\in R(\mathcal C)}
       e^{13k(\mathcal C)}
                        \|\Psi_\mathcal C\|_\infty
                  \le Q_\beta.
\end{equation}
This is a complete cluster sum, not a truncation
at any fixed power of \(\beta\). All the
fine-volume bounds so far are independent of
both \(M\) and \(T\).

\subsection{Projecting the full bulk carriers onto the retained slice}

Let \(\mathfrak p(x,t)=x\) project fine vertices
onto the spatial torus, collapse temporal
edges to vertices, and send spatial edges
at every time to their spatial copies.
For a cluster, let \(\widehat S(\mathcal C)\)
be the edge image of its entire fine edge
union under this projection.
It may include edges that the function does
not actually use; this connected enlargement
is deliberate.

\begin{lemma}[Projection and a summable temporal anchor]
\label{lem:v53-projection}
For every nonconstant logarithmic cluster term,
\(\widehat S(\mathcal C)\) is a nonempty
connected spatial edge carrier containing its
retained source support, and
\[
                        |\widehat S(\mathcal C)|
                                      \le4k(\mathcal C).
\]
If a spatial edge \(e\) belongs to this carrier,
some fine lift \(e_t\) of \(e\) lies in the
cluster, and
\[
                   d_T(t,0)\le k(\mathcal C),
\]
where \(d_T\) is circular distance on the
temporal torus. After subtracting each
cluster's retained Haar mean, the resulting
complete static interaction has, for every
integer \(K\ge0\), the occurrence-tail estimate
\begin{equation}\label{eq:v53-tail}
 \sup_e\sum_{\substack{\mathcal C:\,
                 e\in\widehat S(\mathcal C)\\k(\mathcal C)\ge K}}
     e^{3|\widehat S(\mathcal C)|}
       \|\widetilde\Psi_\mathcal C\|_\infty
               \le10Q_\beta e^{-K/2}.
\end{equation}
\end{lemma}

\begin{proof}
The fine union of a logarithmic cluster is
vertex connected. Each constituent plaquette
contains spatial edges, so its spatial
projection is nonempty. Projection preserves
connectedness, even when temporal paths
collapse to vertices. Its edge count is at
most the number of fine edge occurrences,
which is at most \(4k\).

A nonconstant term must contain a retained
spatial edge at time zero: if the fine union
contains none, its entire activity and
cluster contribution are independent of
\(U^{(0)}\). Follow a path in the connected
plaquette-overlap union from that edge to
the chosen lift \(e_t\). Each plaquette spans
at most one temporal unit. A path using
no more than the number of distinct
plaquettes has temporal displacement at
most that number, which is no greater
than \(k\). This proves the distance bound,
including on the temporal torus.

Put
\[
 \widetilde\Psi_\mathcal C
    =\Psi_\mathcal C-
                      \int\Psi_\mathcal C(U^{(0)})\,dU^{(0)}.
\]
This subtracts only a scalar, kept in the
normalization, and costs at most a factor
two in sup norm. For \(k\ge K\) and a
possible anchor time \(t\), the projection
and distance estimates give
\[
 \begin{split}
 e^{3|\widehat S(\mathcal C)|}
                \|\widetilde\Psi_\mathcal C\|_\infty
   &\le2e^{12k}\|\Psi_\mathcal C\|_\infty\\
   &\le2e^{-K/2}e^{-d_T(t,0)/2}
                     e^{13k}\|\Psi_\mathcal C\|_\infty .
 \end{split}
\]
For a fixed projected edge \(e\), sum over
all its fine lifts \(e_t\), anchor at the
base vertex of each lift, and use
\eqref{eq:v53-pinned}. This overcounts
clusters, which is harmless for the
absolute bound. The temporal anchor sum
is uniformly finite:
\[
 \sum_{t\in\mathbb Z/T\mathbb Z}e^{-d_T(t,0)/2}
       \le1+2\sum_{j\ge1}e^{-j/2}
       =\frac{e^{1/2}+1}{e^{1/2}-1}<5 .
\]
The last inequality follows already from
\(e^{1/2}>3/2\). This proves
\eqref{eq:v53-tail}.
\end{proof}

The temporal anchor is necessary. Merely projecting
every fine carrier to time zero and summing the
fine pinned bound over all possible lift times
would give an unbounded temporal-volume factor.
The nonconstant source's contact with time zero,
together with one unit of unused exponential
charge, supplies both the temporal summation
and the occurrence tail.

\subsection{Identification with the actual Perron vacuum}

\begin{theorem}[A small complete interaction for the Wilson time-zero law]
\label{thm:v53-static}
For every \(M\ge4\) and \eqref{eq:v53-beta}, the
infinite-time Wilson marginal \(\pi_{\beta,M}\)
on its spatial slice has a complete real,
continuous, gauge-invariant interaction
\(\Phi^{\rm stat}_{\beta,M}\), on connected
spatial edge carriers, with
\begin{equation}\label{eq:v53-static-norm}
 B_{\rm stat}:=\|\Phi^{\rm stat}_{\beta,M}\|_{(3)}
       \le10Q_\beta<\frac1{2000}.
\end{equation}
All constants are independent of \(M\). It is
the actual Perron law of V52:
\begin{equation}\label{eq:v53-vacuum}
 \begin{split}
 \pi_{\beta,M}(dU)
   &=\frac{\exp[-\sum_\alpha
                         \Phi^{\rm stat}_{\beta,M,\alpha}(U)]}
           {Z^{\rm stat}_{\beta,M}}\,dU
        =h_{\beta,M}(U)^2\,dU,\\
 h_{\beta,M}(U)
   &=(Z^{\rm stat}_{\beta,M})^{-1/2}
          \exp[-\tfrac12\sum_\alpha
                         \Phi^{\rm stat}_{\beta,M,\alpha}(U)] .
 \end{split}
\end{equation}
The finite-temporal-period interaction can be
chosen consistently with this expansion; for
\(T\ge8\), after matching common cluster indices,
\begin{equation}\label{eq:v53-temporal-convergence}
 \|\Phi^{\rm stat}_{\beta,M,T}
          -\Phi^{\rm stat}_{\beta,M}\|_{(3)}
      \le20Q_\beta
                        e^{-\lfloor T/4\rfloor/2}.
\end{equation}
\end{theorem}

\begin{proof}
At finite \(T\), use the centered cluster
terms from the preceding lemma as the
interaction. Their sum differs from
\(-\log Z_{M,T}\) by a constant only.
The \(K=0\) case of \eqref{eq:v53-tail}
gives its norm bound. On a fixed finite
spatial torus the sum of the sup norms
of all nonconstant terms is at most
\(|E_0|10Q_\beta\), so the potential
itself is absolutely and uniformly
convergent.

Consider a nonconstant cluster with fewer
than \(K=\lfloor T/4\rfloor\) occurrences.
It touches time zero and every one of
its vertices has circular time distance
less than \(K\). Its connected support
has a unique unwrapped lift touching
time zero on the infinite temporal
cylinder; there is no temporal winding.
Its hidden Haar integral, cluster
coefficient, projected carrier, and
retained Haar mean agree with the
infinite-cylinder expressions.
Conversely such a small infinite-cylinder
cluster gives the corresponding finite
one. Thus all the smaller nonconstant
terms match exactly.

Define the infinite-cylinder expansion
by these finite cluster integrals and
their usual logarithmic coefficients,
not by a formal global infinite-time
partition function. Its local pinned
bounds follow from the finite ones
by increasing the allowed temporal
range. Its static norm tail obeys
\eqref{eq:v53-tail} with ordinary
integer time distance. The two
unmatched tails, each starting at
\(K\), give
\eqref{eq:v53-temporal-convergence}.
This also proves the bound and
uniform convergence claimed for the
infinite-time static potential.

At fixed \(M\), the global sup-norm
difference of the potentials is
bounded by \(|E_0|\) times their
interaction-norm difference. The
normalized finite-time densities
therefore converge to the first
density in \eqref{eq:v53-vacuum}.
Independently, Proposition
\ref{prop:v52-fine-transfer} identifies
the limit of these same periodic
Wilson marginals as
\(h_{\beta,M}^2\,dU\). Equality of
the two limits proves the claimed
identification. The positive
normalized square root is unique,
giving the formula for \(h\).
\end{proof}

One useful consequence is a genuinely local
ground-state comparison. If \(U\) and \(V\)
differ on one spatial fine link only, then
\begin{equation}\label{eq:v53-ground-ratio}
                  e^{-B_{\rm stat}}
                  \le\frac{h_{\beta,M}(V)}{h_{\beta,M}(U)}
                  \le e^{B_{\rm stat}} .
\end{equation}
Indeed the sum of the potential differences
has magnitude at most \(2B_{\rm stat}\);
the logarithm of \(h\) contains half that
potential. This does not give a global
volume-independent ratio of \(h\) to a
constant: changing all spatial links
can accumulate an extensive exponent.

The effective interaction in
\eqref{eq:v53-vacuum} includes every
bulk-generated source term. Its leading
spatial plaquette action is not declared
to be its complete action. The following
small regression makes the distinction
particularly explicit.

\begin{proposition}[Bulk integration changes even a one-plaquette marginal]
\label{prop:v53-cube}
In a finite \(\mathbb Z_2\) gauge cube,
retain the four bottom-face links and
integrate the other eight links with
Haar reference. Assign face weights
\(1+t_i\chi_i\), with \(i=0\) the
bottom face and \(i=1,\ldots,5\) the
other faces. If \(C\) is the retained
bottom holonomy and
\(b=t_1t_2t_3t_4t_5\), then the exact
retained density before normalization is
\begin{equation}\label{eq:v53-cube}
                         (1+t_0C)(1+bC).
\end{equation}
After Haar normalization, its character
coefficient is
\[
                       \frac{t_0+b}{1+t_0b},
\]
which is not \(t_0\) when \(b\ne0\)
and \(|t_0|<1\).
\end{proposition}

\begin{proof}
Haar integration on a hidden
\(\mathbb Z_2\) edge kills a selected
face product unless that edge has
even total incidence. On the cube,
the only surviving selected face
sets are the empty set, the bottom
face, the other five faces, and all
six faces. The five-face holonomy
product is \(C\), since the product
over all six faces is one. Thus
the integral is
\[
                     1+t_0C+bC+t_0b,
\]
which proves \eqref{eq:v53-cube}.
The Haar mean is \(1+t_0b\).
\end{proof}

This is a structural regression, not an
\(\SU(2)\) coefficient calculation. In
particular, its fifth-order coefficient
is not imported into the Wilson theory.
Its purpose is to detect a forbidden
replacement of the integrated bulk
by a source-independent scalar.

\subsection{Uniform physical noncollapse through every spatial block}

Let \(\mathcal B_{\rm sp}\) be exactly
the three-dimensional completed map
of \eqref{eq:v52-spatial-map}, at the
unchanged radius \(\rho_\star=2^{-100}\).
Its output time spacing is still the
original \(a\). Let
\[
         \mu_{\beta,M,j}
              =(\mathcal B_{\rm sp}^{\,j})_*
                                     \pi_{\beta,M}
\]
be the actual spatial marginal after
\(j\) completed blocks, with the original
Perron law as its input.

\begin{theorem}[All-depth noncollapse in the actual physical history space]
\label{thm:v53-depth}
Assume \eqref{eq:v53-beta}, let \(M=2^mN\),
and require \(N\ge4\). Every law
\(\mu_{\beta,M,j}\), \(0\le j\le m\),
has a complete connected continuous
gauge-invariant interaction with
\begin{equation}\label{eq:v53-depth-norm}
 B_j\le
       (5/8)^j\,10Q_\beta
                  +\frac{1-(5/8)^j}{12000}
                          <\frac1{2000}.
\end{equation}
For the normalized trace \(F_j\) of
any elementary spatial plaquette at
level \(j\), its centered history
vector satisfies
\begin{equation}\label{eq:v53-uniform-var}
 \begin{split}
 \|[F_j-\omega_\beta(F_j)]\|_{\rm OS}^2
      &=\Var_{\mu_{\beta,M,j}}(F_j)\\
      &\ge \frac14e^{-2B_j}
           >c_0,\qquad c_0:=\frac{999}{4000}>\frac15 .
 \end{split}
\end{equation}
The constants are uniform in \(M,N,m,j\).
The history transfer remains the restriction
of the original physical transfer from V52.
\end{theorem}

\begin{proof}
The proof of Lemma
\ref{lem:v52-spatial-haar} established,
before specializing to Haar input,
the complete three-dimensional
interaction estimate
\[
                         B'\le\frac58B+\frac1{32000}
                                  \quad(B\le1/320).
\]
Its input class is a connected,
continuous, gauge-invariant
interaction on the full compact
spatial links. Theorem
\ref{thm:v53-static} verifies this
class for the actual starting law,
with \(B_0\le10Q_\beta<1/2000\).
The entire generated interaction,
including the scalar normalization,
is used at each succeeding block.
Solving the affine recurrence gives
\eqref{eq:v53-depth-norm}, since
its fixed bound is \(1/12000\)
and both endpoint bounds are
less than \(1/2000\). All final
and intermediate spatial sides
are at least four.

Condition on all level-\(j\)
links except one plaquette edge.
The conditional log-density
oscillation is at most \(2B_j\),
so the conditional density is
at least \(e^{-2B_j}\) relative
to Haar on that link. For fixed
other plaquette edges,
the normalized trace has Haar
variance \(1/4\). Minimization
over its centering constant
and total variance yield
\(\Var(F_j)\ge e^{-2B_j}/4\).
Since \(2B_j<1/1000\) and
\(e^{-x}\ge1-x\), this is
strictly greater than
\((1-1/1000)/4=c_0\).

Finally, Theorem
\ref{thm:v52-history} applies
to these very time-slice
observations of the same
stationary fine process.
It identifies the variance
with the physical reflected
norm and identifies the
history transfer with
the reducing restriction
of the original transfer.
No new conditional Markov
kernel on coarse snapshots
has been introduced.
\end{proof}

This improves the depth-dependent
V52 bound in the stated smaller
coupling interval. It is not a
claim that the sequence of
blocked laws converges as
depth increases. The affine
interaction bound is not a
Lipschitz comparison between
two laws, and no new
projective consistency or
full-history completeness
claim is made.

\subsection{The physical-time obstruction is now witnessed by an actual source}

For clarity consider a fixed coarse
spatial torus of side \(N\ge4\) and
a fixed physical coarse spacing
\(\ell>0\). Choose depths \(m_r\to\infty\),
fine spatial sides \(M_r=2^{m_r}N\),
and isotropic fine spacings
\[
                         a_r=\ell\,2^{-m_r}.
\]
Allow any coupling sequence
\(0<\beta_r\le10^{-12}\).
Let \(F_r\) be the same coarse
plaquette trace after \(m_r\)
spatial companion blocks, and
let \(v_r\) denote its centered
physical history vector.

\begin{theorem}[Nonzero static norms but no strongly continuous continuum-time limit]
\label{thm:v53-time-obstruction}
For this actual sequence of fine
Wilson laws and completed sources,
\begin{equation}\label{eq:v53-time-bounds}
 \begin{split}
 c_0<\|v_r\|^2&\le1,\\
 0\le\langle v_r,e^{-tH_{a_r}}v_r\rangle
                &\le 2^{-t/a_r}\|v_r\|^2
                         \longrightarrow0
                      \qquad(t>0).
 \end{split}
\end{equation}
Consequently no subsequential limit
of these centered kernels, retaining
their time-zero norms, can be
represented by a nonzero vector
and a strongly continuous
nonnegative-Hamiltonian semigroup.
Already across one fine time step,
the stationary observed process obeys
\begin{equation}\label{eq:v53-increment}
 \mathbb E_{\omega_{\beta_r}}
      |F_r(C_1)-F_r(C_0)|^2
                                \ge\|v_r\|^2>c_0 .
\end{equation}
The physical duration of that
step is \(a_r\to0\).
\end{theorem}

\begin{proof}
Theorem \ref{thm:v53-depth} gives
the lower norm bound. The upper
bound follows from \(|F_r|\le1\).
V52's fine physical gap applies,
because \(10^{-12}<1/72\), and
gives
\[
                  H_{a_r}|_{1^\perp}
                              \ge(\log2)/a_r .
\]
The source lies in its centered
reducing history space. Spectral
calculus proves the second
inequality in
\eqref{eq:v53-time-bounds}.
The same limit holds using
integer-time correlations
with \(n_r=\lfloor t/a_r\rfloor\).

Every subsequence has a further
subsequence along which
\(\|v_r\|^2\) tends to a value
\(V\in[c_0,1]\). Its limiting
kernel has value \(V>0\) at
zero and zero at every
positive physical time.
If that kernel were
\(\langle v,e^{-tH}v\rangle\)
with \(e^{-tH}\) strongly
continuous and \(\|v\|^2=V\),
continuity at zero would give
\[
             \lim_{t\downarrow0}
                    \langle v,e^{-tH}v\rangle
                            =V,
\]
contradicting its positive-time
values.

For the one-step statement,
stationarity and the actual
fine Markov representation
give
\[
 \mathbb E|F_r(C_1)-F_r(C_0)|^2
     =2\|v_r\|^2
           -2\langle v_r,\mathsf P_rv_r\rangle .
\]
The centered transfer bound
\(\mathsf P_r\preceq1/2\)
from V52 proves
\eqref{eq:v53-increment}.
This identity uses the fine
process and does not require
the observed snapshots to
be Markov.
\end{proof}

The semigroup contradiction is an
application of the previously proved
physical-clock criterion, not a new
abstract replacement for it. Its
new input is the explicit lower
bound for the actual cofinally
blocked source. The alternative
that this plaquette's time-zero
norm merely vanishes is now
excluded in this regime.

Scalar renormalization cannot
repair this conclusion while
keeping finite nonzero centered
norms. More generally, for any
centered vectors \(w_r\) in the
same physical spaces, if
\(\|w_r\|^2\to V\in(0,\infty)\),
the same gap gives
\[
          |\langle w_r,e^{-tH_{a_r}}w_r\rangle|
                            \le2^{-t/a_r}\|w_r\|^2
                            \longrightarrow0
                       \quad(t>0).
\]
This includes vectors from
complete observation histories;
enlarging the history space
does not lower the spectrum
of a reducing restriction.
Assigning a different clock
would change the limiting
question, not prove the
required limit for the
original isotropic Wilson
time coordinate.

\subsection{What this changes in the programme}

The actual static Wilson law is
now controlled uniformly in
spatial volume inside an explicit
small-coupling region. Spatial
completion of this law has
uniform nonzero physical probes
at every finite depth. This
is stronger than using compact
Haar variance as a comparison
at each fixed depth, and it
identifies a local interaction
for the actual Perron vacuum.

The same region is nevertheless
not a nontrivial continuum-time
entrance for these isotropically
refined sources. That failure
is concrete and not cured by
another terminal contraction
estimate. An eventual continuum
construction must supply a
different fine-coupling trajectory
and a physical-rate comparison
that does not force this
instantaneous decorrelation.
The estimates here do not
prove that trajectory, its
asymptotics, its ultraviolet
renormalization, or a positive
finite physical mass scale.

This also does not rule out a
strong-coupling effective law
at a \emph{fixed physical coarse
scale} reached from a different
fine ultraviolet regime. Such
an application would need the
actual selected coarse history
law and a correctly matched
physical-time transfer. The
fine-regime obstruction proved
above must not be extended to
that materially different
claim without its missing
comparison.

Nothing in this section proves
reflection positivity of the
original four-dimensional V51
completion. All physical
statements use the slice-local
companion from V52. Nor does
one nonzero plaquette establish
density of the observed algebra
in a full continuum physical
space. Those distinctions,
the continuum axioms, and
the nonperturbative physical
mass gap remain open.

\subsection{Verification and predecessor preservation}

The exact certificate is
\begin{center}\small
\path{scripts/verify_ym_f14_v53_static_marginal.py}.
\end{center}
It imports only the rational
exponential enclosure routine
from the V50 certificate and
otherwise uses the Python
standard library. It checks
\(e^{69/5}<10^6\), the pinned
budget, temporal anchor sum,
static norm threshold, affine
recurrence through depth one
hundred, and the rational
variance floor \(999/4000\).
The induction in the proof,
not those hundred tests,
establishes the all-depth
statement.

On a \(4^3\times8\) cylinder
it verifies all twenty-four
plaquette incidences at every
fine vertex and checks
projection connectivity,
edge counts, and temporal
distance for all anchored
vertex-connected plaquette
sets of sizes one, two,
and three. Their counts
are \(1,72,5724\).
The geometric proof above
covers all sizes and
logarithmic repetitions.

For the finite gauge cube
the certificate sums all
\(256\) hidden configurations
for each of the sixteen
retained configurations.
It checks all sixty
connected face polymers,
the exact component
partition identity, and
the four surviving face
subsets. With
\(t_i=1/(7+i)\), the
normalizing scalar is
\(665281/665280\) and
the actual normalized
character coefficient
is \(95047/665281\),
not \(1/7\). This
regression is explicitly
\(\mathbb Z_2\), while
the \(\SU(2)\) analytic
bounds are proved in
the preceding sections.

V52 has 1,272,731 bytes and SHA--256
\begin{center}\scriptsize\ttfamily
574CBC1E10D839B402B01090DE3F1CCF2ED61F00D72B5BDAD329C2D1C6A81BA7.
\end{center}
Its entire body is retained
except for the title version
and this terminal insertion.
V53 replaces V52 as the
single active main note.
Archives, supplied notes,
companion programmes, and
monographs remain unchanged.

\begin{firewall}
Uniform all-depth physical
plaquette noncollapse has
been established for the
slice-local companion in
the stated strong-coupling
region. The same actual
source makes the failure
of a strongly continuous
isotropic continuum-time
limit explicit there.
This is progress in
locating and excluding
a circular route, not
a construction of the
full continuum theory.
The four-dimensional
interacting Yang--Mills
mass-gap goal remains
unproved and active.
\end{firewall}
\section{V54: exact commuting zero-mode normalization and logarithmic toron scales at fixed Wilson volume}
\label{sec:v54}

\subsection{The upstream question after the strong-coupling obstruction}

V53 excludes the current small-\(\beta\) fine-lattice
regime as a nontrivial strongly continuous
continuum-time entrance. We therefore return
upstream to weak coupling, meaning
\(\beta\to\infty\) in the normalization
\(S_0=\sum_p(1-\tfrac12\operatorname{Re}\Tr U_p)\).
The bare Wilson chessboard and averaged
conditional-score estimates already identified
in V31 are not repeated. Neither is the
one-loop coefficient computed in V30.

The missing question treated here is different:
what normalization and amplitude scale does the
commuting constant-field sector actually have?
The zero-mean Gaussian charts exclude precisely
these harmonic variables. Integrating them as
a uniformly normalized twelve-dimensional
quartic reference is not valid in four
dimensions. We first prove the exact
critical logarithm for the compact
four-link Wilson integral, and then
derive its counterpart for the actual
Wilson law on every fixed finite
four-torus. The second result does
not assume that the constant-link
restriction is the marginal of the
full theory.

The critical dimension of the bosonic
\(\SU(2)\) matrix integral is classical;
see
\href{https://arxiv.org/abs/hep-th/0101071}
{P.~Austing and J.~F.~Wheater,
The Convergence of Yang--Mills Integrals,
hep-th/0101071}.
The result below gives the constants
for this compact Wilson normalization,
the corresponding logarithmic radius
law, and the fixed-volume reduction
from the full lattice integral.
The finite-volume constants are not
asserted to be uniform as the torus
grows. That distinction is essential
to the still-open continuum goal.
This section is dated 15 September
2026; the next companion checkpoint
is V55.

\subsection{The exact compact commutator action}

Write an \(\SU(2)\) element as a unit quaternion,
\[
 U=sI+i\,v\cdot\sigma,\qquad s^2+|v|^2=1,
\]
where \(\sigma\) are the Pauli matrices.
The normalized character is \(\chi(U)=s\).
For four such elements let \(A\) be the
real \(4\)-by-\(3\) matrix with rows \(v_\mu\),
and put
\begin{equation}\label{eq:v54-Q}
 Q(A)=\sum_{\mu<\nu}|v_\mu\times v_\nu|^2
     =\frac12\{(\operatorname{tr}A^TA)^2
                         -\operatorname{tr}((A^TA)^2)\}.
\end{equation}
Here the lower-case trace is a real matrix
trace, not the fundamental Wilson trace.

\begin{lemma}[Exact quaternion identity and Haar density]
\label{lem:v54-quaternion}
For \(U=sI+i\,v\cdot\sigma\) and
\(V=tI+i\,w\cdot\sigma\),
\begin{equation}\label{eq:v54-commutator}
             1-\chi(UVU^{-1}V^{-1})=2|v\times w|^2 .
\end{equation}
The pushforward of normalized \(\SU(2)\)
Haar measure under \(U\mapsto v\) is
\begin{equation}\label{eq:v54-Haar}
             \frac{1}{\pi^2}
               \frac{\mathbf1_{\{|v|<1\}}}{\sqrt{1-|v|^2}}\,d^3v.
\end{equation}
On one chosen scalar-sign sheet near
\(U=I\) or \(U=-I\), the density is
half of \eqref{eq:v54-Haar}.
\end{lemma}

\begin{proof}
Quaternion multiplication gives
\[
 (s,v)(t,w)=(st-v\cdot w,\;sw+tv+v\times w).
\]
Multiplying by the inverses \((s,-v)\)
and \((t,-w)\), and taking the scalar
part, gives \(1-2(|v|^2|w|^2-(v\cdot w)^2)\).
This proves \eqref{eq:v54-commutator};
summing the elementary cross-product
identity gives \eqref{eq:v54-Q}.

Normalized Haar is surface measure
on \(S^3\), divided by \(2\pi^2\).
Projection of either scalar-sign
hemisphere to \(v\) has surface
Jacobian \(1/\sqrt{1-|v|^2}\).
Adding the two hemispheres proves
\eqref{eq:v54-Haar}. The equator
has zero Haar measure.
\end{proof}

In particular the compact four-link
partition function
\begin{equation}\label{eq:v54-compact-Z}
 Z_{\rm c}(\beta)=
     \int_{\SU(2)^4}
       \exp\!\left[-\beta\sum_{\mu<\nu}
                (1-\chi(U_\mu U_\nu U_\mu^{-1}U_\nu^{-1}))\right]
                    \prod_{\mu=1}^4dU_\mu
\end{equation}
has action exactly \(2Q(A)\) in
quaternion-vector coordinates.
There are no omitted higher-order
terms in that action identity.
The density and the compact domain
are still present.

\subsection{A normal calculation on the rank-one sphere}

Put \(r=|A|_{\rm F}\) and write
\(A=r\omega\), \(\omega\in S^{11}\).
Let \(q=Q|_{S^{11}}\), and define
\[
                        \mathcal J(t)
                         =\int_{S^{11}}e^{-tq(\omega)}\,d\sigma(\omega).
\]
All spheres and Euclidean matrix
spaces in this subsection use
their ordinary Riemannian volumes,
without probabilistic normalization.

\begin{lemma}[Angular asymptotic with its coefficient]
\label{lem:v54-angular}
As \(t\to\infty\),
\begin{equation}\label{eq:v54-angular}
                \mathcal J(t)=4\pi^6t^{-3}+O(t^{-7/2}),
           \qquad
                \mathcal J(t)\le C(1+t)^{-3}.
\end{equation}
\end{lemma}

\begin{proof}
The nonnegative expression \(Q(A)\)
vanishes exactly when \(A\) has rank
at most one. On the unit sphere its
zero set is the compact smooth manifold
\[
 \mathcal K
    =\{uv^T:u\in S^3,\ v\in S^2\}
    \cong(S^3\times S^2)/\{(u,v)\sim(-u,-v)\}.
\]
It has dimension five. The differential
\(\delta u\,v^T+u\,\delta v^T\)
has squared norm
\(|\delta u|^2+|\delta v|^2\), because
the two summands are orthogonal.
Hence
\[
               \operatorname{vol}(\mathcal K)
                =\tfrac12\operatorname{vol}(S^3)
                              \operatorname{vol}(S^2)=4\pi^3 .
\]

At \(uv^T\), the normal space within
\(S^{11}\) consists of matrices \(B\)
such that \(u^TB=0\) and \(Bv=0\).
Its dimension is \(3\cdot2=6\).
For small \(B\), use the normal
coordinate
\[
            \omega=\frac{uv^T+B}{\sqrt{1+|B|^2}} .
\]
The largest singular line of a
nearby rank-one matrix is simple.
It recovers \(u,v\), modulo their
simultaneous sign, and recovers
\(B\) from the orthogonal remaining
part. Thus these coordinates give
a smooth normal-bundle neighborhood
of \(\mathcal K\). Their volume
Jacobian is \(1+O(|B|)\), uniformly
over the compact base.

By orthogonal invariance, take \(u,v\)
to be the first coordinate vectors.
The remaining \(B\) is a \(3\)-by-\(2\)
matrix in the opposite block. Direct
use of \eqref{eq:v54-Q} gives
\begin{equation}\label{eq:v54-normal}
 Q(uv^T+B)=|B|^2+\det(B^TB),\qquad
 q(\omega)=
       \frac{|B|^2+\det(B^TB)}{(1+|B|^2)^2}.
\end{equation}
The normal Hessian of \(q\) is
therefore \(2I_6\), and
\(q(\omega)=|B|^2+O(|B|^4)\).
Shrink the normal neighborhood
so \(c|B|^2\le q\le C|B|^2\).
Off that neighborhood \(q\) has
a strictly positive minimum.

In each normal fibre substitute
\(B=t^{-1/2}Y\). The Jacobian error
is bounded by \(Ct^{-1/2}|Y|\);
the exponent error is bounded by
\(Ct^{-1}|Y|^4\) with a common
integrable Gaussian majorant
after restricting to a smaller
normal neighborhood. Its complement
is exponentially small. Integration
therefore gives
\[
 \mathcal J(t)
  =t^{-3}\operatorname{vol}(\mathcal K)
             \left\{\int_{\mathbb R^6}e^{-|Y|^2}\,dY
                               +O(t^{-1/2})\right\}.
\]
The Gaussian integral is \(\pi^3\),
proving the first assertion.
The same normal lower bound and
the compact off-neighborhood
estimate give the second.
\end{proof}

\begin{proposition}[The critical quartic radial logarithm]
\label{prop:v54-radial}
For every fixed \(R>0\),
\begin{equation}\label{eq:v54-radial}
 \int_{|A|<R}e^{-bQ(A)}\,d^{12}A
       =\pi^6 b^{-3}\log b+O_R(b^{-3})
                      \qquad(b\to\infty).
\end{equation}
More precisely, if \(bR_b^4\to\infty\),
\begin{equation}\label{eq:v54-moving-radius}
 \int_{|A|<R_b}e^{-bQ(A)}\,d^{12}A
      =\pi^6b^{-3}\log(bR_b^4)+O(b^{-3}),
\end{equation}
with an error uniform once \(bR_b^4\ge1\).
A bounded Lipschitz amplitude \(w(A)\)
on a fixed small neighborhood changes
the leading coefficient to
\(\pi^6w(0)\). A smooth cutoff equal
to one near zero has the same
leading coefficient.
\end{proposition}

\begin{proof}
Polar coordinates and \(t=br^4\)
give the exact identity
\[
 \int_{|A|<R}e^{-bQ(A)}\,dA
       =\frac{b^{-3}}4
                     \int_0^{bR^4}t^2\mathcal J(t)\,dt .
\]
The integrand at large \(t\) is
\(4\pi^6/t+O(t^{-3/2})\).
Its error is integrable at infinity,
and the lower part of the integral
is finite. This proves both radial
formulas.

For an amplitude, subtract \(w(0)\).
Its error is at most \(Cr\). By
\eqref{eq:v54-angular}, the absolute
error integral is bounded by
\[
 C\int_0^R r^{12}(1+br^4)^{-3}\,dr
                                      =O_R(b^{-3}).
\]
Indeed the part below \(b^{-1/4}\)
is \(O(b^{-13/4})\), and the
remaining integrand is at most
\(Cb^{-3}\). Changing a cutoff
only on an annulus separated
from zero also costs \(O(b^{-3})\).
\end{proof}

The uncut twelve-dimensional quartic
integral is consequently divergent:
at large radial cutoff its
divergence is logarithmic.
The normal Gaussian has six
coordinates while the rank-one
cone has radial volume proportional
to \(r^5dr\); normal integration
contributes \(r^{-6}\), leaving
\(dr/r\). This is the specific
normalization problem, not a
failure of the compact integral
to exist.

\subsection{The complete compact four-link law and its logarithmic scale}

\begin{theorem}[Compact partition asymptotic and radius law]
\label{thm:v54-compact}
For \eqref{eq:v54-compact-Z},
\begin{equation}\label{eq:v54-compact-asymptotic}
           Z_{\rm c}(\beta)
              =\frac{\beta^{-3}\log\beta}{8\pi^2}
                              +O(\beta^{-3}).
\end{equation}
Let \(\nu_\beta\) be its normalized
probability law and let
\[
                        R(U)=
                  \left(\sum_{\mu=1}^4|v_\mu|^2\right)^{1/2}.
\]
Then for \(0<s<1/4\),
\begin{equation}\label{eq:v54-compact-radius}
           \nu_\beta\{R\le\beta^{-s}\}
                                  \longrightarrow1-4s .
\end{equation}
Equivalently,
\[
                      -\frac{\log R}{\log\beta}
                       \ \Longrightarrow\ 
                         \operatorname{Uniform}[0,1/4].
\]
The unscaled radius tends to zero
in probability, but for every
fixed \(K<\infty\),
\begin{equation}\label{eq:v54-quartic-escape}
            \nu_\beta\{\beta^{1/4}R\le K\}
                                  =O_K(1/\log\beta)\to0 .
\end{equation}
\end{theorem}

\begin{proof}
The zero-action set is the set of
commuting four-tuples. Away from
the sixteen central tuples
\((\pm I,\pm I,\pm I,\pm I)\),
it is a smooth six-dimensional
manifold: write the tuple as
\[
 U_\mu=\cos\theta_\mu\,I+
                  i\sin\theta_\mu\,n\cdot\sigma ,
\]
with four circle angles and
\(n\in S^2\), modulo the discrete
simultaneous sign identification.
At least one \(\sin\theta_\mu\)
is nonzero, so the two axis
directions remain genuine.

The action has a positive quadratic
form in its six normal directions.
One direct verification uses left
Lie variations \(X_\mu\). At a
commuting tuple the differential
of the \((\mu,\nu)\) commutator is
\[
       (I-\operatorname{Ad}U_\nu)X_\mu
                +(\operatorname{Ad}U_\mu-I)X_\nu .
\]
The four Cartan variations are
in its kernel. In the transverse
complex color line put
\(w_\mu=e^{2i\theta_\mu}-1\).
The remaining differential is
\(w_\mu z_\nu-w_\nu z_\mu\).
When \(w\ne0\), its kernel is
the one-dimensional complex span
of \(w\), corresponding to the
two axis directions. Thus its
real rank is six. Since the
Wilson action is a sum of
nonnegative squared deviations
to second order, its normal
Hessian is positive.

On any compact part of this
regular zero set separated from
the central tuples, smooth normal
coordinates therefore bound the
action below by \(c|y|^2\) in
six normal variables. Integrating
there gives \(O(\beta^{-3})\).
Outside small normal neighborhoods
of the zero set, compactness gives
a positive action minimum and an
exponentially small contribution.
These statements use smooth group
coordinates; the apparent
equatorial singularity in
\eqref{eq:v54-Haar} introduces
no additional zero-mode stratum.

Near each central tuple use its
four quaternion-vector charts.
Their product Haar density at
zero is \((2\pi^2)^{-4}\),
and the action is exactly
\(2Q(A)\). Proposition
\ref{prop:v54-radial} gives the
logarithmic coefficient per
central tuple:
\[
       (2\pi^2)^{-4}\,
                     \pi^6(2\beta)^{-3}\log\beta
           =\frac{\beta^{-3}\log\beta}{128\pi^2}.
\]
All sixteen tuples have the same
coefficient. Adding them and the
regular \(O(\beta^{-3})\) part
proves \eqref{eq:v54-compact-asymptotic}.
This keeps the Haar normalization
and both scalar-sign sheets.

For \(R\le\beta^{-s}\), all four
links lie in those central
charts for sufficiently large
\(\beta\). Apply
\eqref{eq:v54-moving-radius}
with \(b=2\beta\) and
\(R_b=\beta^{-s}\).
The logarithm is
\((1-4s)\log\beta+\log2\);
the amplitude error remains
\(O(\beta^{-3})\).
Dividing by the full partition
function proves
\eqref{eq:v54-compact-radius}.

The region \(R\ge\epsilon>0\)
has only a regular zero set,
and hence probability
\(O_\epsilon(1/\log\beta)\).
For \(R\le K\beta^{-1/4}\),
the exact radial integral has
bounded upper endpoint
\(2K^4\) after substitution,
so it is \(O_K(\beta^{-3})\).
This proves
\eqref{eq:v54-quartic-escape}
and the endpoint assertions
needed for the stated weak
convergence of the logarithmic
radius.
\end{proof}

The full normalized law is invariant
under multiplying any one \(U_\mu\)
by \(-I\). The central labels thus
have equal limiting weights \(1/16\).
This is not gauge fixing those
labels to the identity.

\subsection{Reduction from the actual finite four-dimensional Wilson law}

Now fix \(M\ge4\), put \(V=M^4\),
and let \(Z_M(\beta)\) be the
Wilson partition function on the
periodic four-torus of side \(M\),
with \(4V\) independent oriented
links and normalized product Haar:
\begin{equation}\label{eq:v54-full-Z}
                  Z_M(\beta)=\int e^{-\beta S_0(U)}\,dU.
\end{equation}
This is a finite Euclidean torus,
not the infinite-time Perron law
of V52--V53.

Divide out only the \emph{based}
gauge group, whose transformation
at one chosen root is the identity.
A global spanning-tree Haar change
sets its \(V-1\) tree links to
identity and integrates their
gauge variables with mass one.
The resulting compact smooth
space has dimension
\begin{equation}\label{eq:v54-dimensions}
 D_M=3(4V-(V-1))=9V+3.
\end{equation}
The residual simultaneous
conjugation at the root is
not divided out. In particular,
no singular orbit-volume
denominator is introduced at
central holonomy.

\begin{lemma}[Flat strata and their normal ranks]
\label{lem:v54-flat-strata}
In the based-gauge quotient, the
zero-action set is parametrized
by commuting four-tuples of
based fundamental-cycle holonomies.
It has sixteen central points.
Away from them it is a smooth
six-dimensional manifold with
normal Hessian rank \(9V-3\).
At each central point its
Hessian kernel has dimension
twelve and its positive
nonzero-mode rank is
\begin{equation}\label{eq:v54-d}
                              d_M=9(V-1).
\end{equation}
\end{lemma}

\begin{proof}
Zero Wilson action means every
plaquette holonomy is identity.
Parallel transport along
contractible lattice paths is
then path independent. The
four fundamental holonomies
commute, and determine a
flat configuration up to
based gauge. Conversely
any commuting tuple defines
such a flat configuration.
Its regular and central
strata are exactly those
described in the proof of
Theorem \ref{thm:v54-compact}.

At a flat connection the
quadratic action is a positive
constant times the squared
covariant lattice curl.
Its kernel is the space of
closed adjoint one-cochains.
For completeness, diagonalize
the commuting holonomies.
On a charged complex color
line the twisted Fourier
curl has symbol
\(p_\mu z_\nu-p_\nu z_\mu\).
Whenever \(p\ne0\), closed
one-cochains are its
one-dimensional span and
are gradients. If the
charged holonomy character
is nontrivial in at least
one direction, every such
Fourier symbol is nonzero.
The neutral Cartan line
has exactly four constant
harmonic one-cochains.

At a regular flat connection
the full gauge-gradient
rank is \(3V-1\), because
its stabilizer has dimension
one. Adding the four neutral
harmonics gives full Hessian
kernel dimension \(3V+3\).
Subtracting the based-gauge
dimension \(3(V-1)\) leaves
six. At central holonomy
the adjoint twists are all
trivial: the gradient rank
is \(3V-3\) and there are
twelve harmonic one-cochains.
Subtracting based gauge
leaves twelve. Subtraction
from \eqref{eq:v54-dimensions}
gives the asserted normal
ranks. The sum-of-squares
Hessian is strictly positive
on each normal complement.
\end{proof}

\begin{lemma}[A measured central chart with the exact constant-field action]
\label{lem:v54-central-chart}
At the central point with trivial
fundamental holonomies, there is
a local Coulomb chart consisting
of four constant Lie vectors
\(a_\mu\in\mathbb R^3\) and
\(d_M\) mean-zero transverse
coordinates \(\xi\), with
representative links
\begin{equation}\label{eq:v54-central-representative}
             U_{x,\mu}
                    =\exp\{i(a_\mu+\xi_{x,\mu})\cdot\sigma\},
         \qquad d^*\xi=0,\qquad
                  \sum_x\xi_{x,\mu}=0 .
\end{equation}
The quotient Haar measure in
these coordinates is
\(J_M(a,\xi)\,d^{12}a\,d^{d_M}\xi\),
with \(J_M\) smooth and strictly
positive at the origin.
For \(a\) sufficiently small
and \(M\) fixed,
\[
 \partial_\xi S_0(a,0)=0,\qquad
 \partial_\xi^2S_0(a,\xi)\succeq c_M I
\]
on a common small neighborhood,
with \(c_M>0\). At \(\xi=0\)
the action is exactly
\begin{equation}\label{eq:v54-central-action}
                 S_0(a,0)=2VQ(v(a)),\qquad
       v_\mu(a)=\frac{\sin|a_\mu|}{|a_\mu|}a_\mu .
\end{equation}
Integrating the nonzero modes
on this chart, with a smooth
cutoff equal to one near its
origin, gives
\begin{equation}\label{eq:v54-amplitude}
 \beta^{-d_M/2}
       \int e^{-2\beta VQ(v)}
               \{A_M(v)+O_M(\beta^{-1/2})\}\,d^{12}v ,
\end{equation}
locally near \(v=0\), where
\(A_M\) is smooth,
\(A_M(0)>0\), and the error
is uniform on a fixed smaller
source neighborhood. The same
amplitude at zero applies to
all sixteen central points.
\end{lemma}

\begin{proof}
At the identity, the derivative
of a based gauge transformation
in logarithmic link coordinates
is the lattice gradient.
The equation that their
divergence vanish has derivative
the scalar lattice Laplacian
on gauge parameters with the
root value fixed. It is an
isomorphism onto the mean-zero
site functions. The finite
implicit function theorem
therefore gives a local
section \(d^*\log U=0\).
Its linear space splits into
the twelve constant coordinates
and the mean-zero transverse
space of dimension \(d_M\).
Together with the based-gauge
coordinates it is a nonsingular
change of finite-dimensional
variables. Disintegrating
normalized Haar gives the
smooth positive density \(J_M\).
We do not quotient its
residual global conjugation.

At a spatially and temporally
constant field, the derivative
of \(S_0\) in logarithmic
coordinates is constant in
the lattice base point for
each direction, by translation
invariance. Pairing it with
a mean-zero \(\xi\) is zero.
This proves the exact
stationarity in \(\xi\),
even for noncommuting \(a_\mu\).
At \(a=\xi=0\) the transverse
nonzero-mode Hessian is
positive by Lemma
\ref{lem:v54-flat-strata}.
Finite-dimensional continuity
gives the stated common floor
after shrinking the chart,
with \(M\) held fixed.

Each constant plaquette is
the commutator of the same
two links, and each direction
pair occurs at \(V\) base
points. Lemma
\ref{lem:v54-quaternion}
therefore proves
\eqref{eq:v54-central-action}
exactly. In particular the
nonzero-mode minimizing
center is \(\xi=0\), and
its minimized source action
has not been replaced by
only a Taylor polynomial.

For the integration statement,
factor out \(e^{-\beta S_0(a,0)}\)
and set \(\xi=\beta^{-1/2}y\).
Uniform Taylor expansion of
the action gives
\[
 \beta(S_0(a,\beta^{-1/2}y)-S_0(a,0))
        =\tfrac12 y^TH_M(a)y
                       +O_M(\beta^{-1/2}|y|^3).
\]
On \(|y|\le\beta^{1/10}\)
the remainder can be absorbed
in a fixed Gaussian majorant;
integration of its absolute
error is \(O_M(\beta^{-1/2})\).
Outside that range the
normal lower bound gives
an exponentially small
Gaussian tail. Expanding
the smooth measure density
has the same integrated
error. Thus the leading
amplitude before changing
source coordinates is
\[
           (2\pi)^{d_M/2}
                     \frac{J_M(a,0)}{\sqrt{\det H_M(a)}} .
\]
The local diffeomorphism
\(a\mapsto v(a)\) has
derivative the identity at
zero and a smooth positive
Jacobian nearby. Including
its inverse determinant
gives \(A_M\) in
\eqref{eq:v54-amplitude}.
Cutoffs away from \(\xi=0\)
cost exponentially small
terms on the smaller
source neighborhood.

Finally, each central
holonomy choice has a
flat representative made
of central link signs,
with their product around
every plaquette equal to
one. Multiplication by
such a fixed central
cocycle is an exact
Haar-preserving symmetry
of the Wilson action.
It transports this chart
and its normalization to
the other fifteen points.
\end{proof}

The constants \(c_M\), chart radius,
amplitude norms, and Gaussian
error constants in this lemma
are all finite-volume constants.
In particular the nonzero-mode
curl floor decreases when the
torus grows. The lemma must
not be read as a uniform
statement on the ultraviolet
sequence of V44.

\begin{theorem}[Actual fixed-volume Wilson partition asymptotic]
\label{thm:v54-full-partition}
For every fixed \(M\ge4\), with
\(V=M^4\), there is a constant
\(C_M>0\) such that
\begin{equation}\label{eq:v54-full-asymptotic}
 Z_M(\beta)=
     C_M\,\beta^{-(9V-3)/2}\log\beta
                       +O_M(\beta^{-(9V-3)/2}) .
\end{equation}
In the source-coordinate
normalization of Lemma
\ref{lem:v54-central-chart},
\begin{equation}\label{eq:v54-CM}
                       C_M=\frac{2\pi^6}{V^3}A_M(0).
\end{equation}
The normalized law concentrates
near the sixteen central
flat configurations in the
based-gauge quotient, with
equal limiting weights.
\end{theorem}

\begin{proof}
Cover the central points
by disjoint small coordinate
neighborhoods as in Lemma
\ref{lem:v54-central-chart}.
For one such neighborhood,
Proposition \ref{prop:v54-radial}
applied to
\eqref{eq:v54-amplitude}
gives the contribution
\[
       \beta^{-d_M/2}
         \left\{
           \frac{\pi^6A_M(0)}{8V^3}\,
                       \beta^{-3}\log\beta
                        +O_M(\beta^{-3})
         \right\}.
\]
A source-dependent smooth
amplitude changes only the
nonlogarithmic error, as
proved in that proposition.
The uniform
\(O_M(\beta^{-1/2})\)
amplitude remainder times
the quartic integral is
\[
                 O_M(\beta^{-7/2}\log\beta)
                              =O_M(\beta^{-3}).
\]

After those central
neighborhoods are removed,
the remaining zero set
is compact and regular.
Its normal dimension is
\(9V-3\), by Lemma
\ref{lem:v54-flat-strata}.
A finite normal-coordinate
cover bounds its contribution
by \(O_M(\beta^{-(9V-3)/2})\).
Away from the full zero
set the action is bounded
strictly below by a
positive constant, giving
an exponentially small
contribution. This includes
the chart cutoff transition
regions; no unrelated
global density has been
substituted.

Since
\[
                 \frac{d_M}{2}+3=\frac{9V-3}{2},
\]
the sixteen equal central
logarithmic contributions
dominate the regular part.
Their sum gives
\eqref{eq:v54-full-asymptotic}
and \eqref{eq:v54-CM}.
The same comparison with
any fixed central
neighborhood proves
concentration there.
The exact central-cocycle
symmetry proves equality
of their limiting weights.
\end{proof}

This theorem concerns the full
finite Wilson law, including
the nonzero-mode measure and
determinant. The four-link
calculation by itself would
not have justified that
identification.

\subsection{A gauge-invariant radius for the actual retained toron}

Fix a root and the four straight
fundamental lattice cycles based
there. Write \(W_\mu(U)\) for
their holonomies and define
\begin{equation}\label{eq:v54-holonomy-radius}
       R_H(U)=
          \left\{\sum_{\mu=1}^4
                       (1-\chi(W_\mu(U))^2)\right\}^{1/2}.
\end{equation}
This function is invariant under
gauge transformations and under
independent center changes of
the cycle holonomies. On a
flat configuration it vanishes
exactly at central holonomy.

\begin{theorem}[Logarithmic radius law for the full finite Wilson measure]
\label{thm:v54-full-radius}
Let \(\mathbb P_{\beta,M}\) be the
normalized law from
\eqref{eq:v54-full-Z}, with
\(M\ge4\) fixed. Then
\begin{equation}\label{eq:v54-full-radius}
        -\frac{\log R_H}{\log\beta}
              \ \Longrightarrow\
                        \operatorname{Uniform}[0,1/4]
                  \qquad(\beta\to\infty).
\end{equation}
In particular \(R_H\to0\) in
probability, but for every
fixed \(K<\infty\),
\[
            \mathbb P_{\beta,M}
                \{\beta^{1/4}R_H\le K\}\longrightarrow0 .
\]
The twelve harmonic variables
do not have a tight
\(\beta^{-1/4}\)-rescaled
probability reference.
\end{theorem}

\begin{proof}
In a sufficiently small central
chart, the holonomy at
\(\xi=0\) is, up to its
central sign,
\(\exp(iM a_\mu\cdot\sigma)\).
For fixed \(M\) its radius
is comparable to \(|a|\)
on a small enough source
ball. Smoothness of a
finite link product gives
constants \(c_M,C_M'>0\)
such that
\[
 c_M|a|-C_M'|\xi|
             \le R_H(U(a,\xi))
             \le C_M'(|a|+|\xi|).
\]
The same comparison holds
with \(|v(a)|\) in place
of \(|a|\).

The normal quadratic lower
bound from Lemma
\ref{lem:v54-central-chart}
implies that, inside these
charts, the portion with
\(|\xi|>\beta^{-3/8}\)
is bounded by a fixed
Gaussian prefactor times
\(\exp(-c_M\beta^{1/4})\).
After normalization by
\eqref{eq:v54-full-asymptotic}
this portion tends to
zero. The complement
of the central charts
already has probability
\(O_M(1/\log\beta)\).

Fix \(0<s<1/4\).
On \(|\xi|\le\beta^{-3/8}\),
the event \(R_H\le\beta^{-s}\)
is sandwiched between
source balls of radii
\(c\beta^{-s}\) and
\(C\beta^{-s}\), for
fixed positive constants
depending on \(M\).
For either radius,
the exact radial formula
and the measured amplitude
in \eqref{eq:v54-amplitude}
give central mass
\[
  C_M(1-4s)\,
       \beta^{-(9V-3)/2}\log\beta
                 +O_M(\beta^{-(9V-3)/2}).
\]
A constant multiple in
the radius changes the
logarithm only by a
bounded additive term.
Dividing by the partition
function proves
\[
        \mathbb P_{\beta,M}
                    \{R_H\le\beta^{-s}\}
                            \longrightarrow1-4s.
\]

For the endpoints, fixed
positive radius has only
regular-flat mass outside
the central neighborhoods,
so \(R_H\to0\). At radius
\(K\beta^{-1/4}\) the
source integral's rescaled
radial upper endpoint is
bounded, giving a mass
without the logarithmic
factor; its normalized
probability tends to
zero. The same normal
Gaussian tail and outside-
chart estimates apply.
These facts establish
\eqref{eq:v54-full-radius}.
\end{proof}

Thus a typical logarithmic exponent
is distributed over an interval,
not concentrated at a single
power of \(\beta\). For example,
the limiting probability that
\(R_H\le\beta^{-1/8}\) is
one half. An uncut quartic
matrix integral cannot serve
as the normalized replacement
for this compact harmonic
sector. Nonzero-mode Gaussian
integration is compatible
with the theorem; extending
that Gaussian assertion to
the harmonic variables is
the invalid step.

\subsection{Consequences for the measured programme and its limits}

At fixed volume the scalar free
energy contains
\[
 -\log Z_M(\beta)
      =\frac{9V-3}{2}\log\beta
                   -\log\log\beta-\log C_M
                         +O_M(1/\log\beta).
\]
The \(-\log\log\beta\) term comes
from integrating the retained
commuting zero-mode sector.
It is not another local
one-loop \(F^2\) coefficient,
and it must not be inserted
into the beta function as
one. Conversely it cannot
be dropped if the purpose
is to normalize this entire
finite compact law.

This result addresses a part
left out of the zero-mean
classical and Gaussian charts:
the true harmonic integration
and its singular flat stratum.
It also shows why declaring
a single quartic amplitude
scale for that sector would
give the wrong entrance
probability. A measured
weak-coupling continuation
must keep the sector law,
or prove its replacement
with a quantitative error.

There are three indispensable
limits on the conclusion.
First, all statements involving
the full \(M^4\) torus hold
with \(M\) fixed before
\(\beta\to\infty\). Neither
the error constants nor
the chart size have been
controlled along
\(\beta\asymp\log M\).
Second, the holonomy radius
is a wrapping finite-volume
observable, not a dense
algebra of local
infinite-volume probes.
Third, these are static
Euclidean integrals. No
physical transfer gap,
infinite-time Perron
asymptotic, or continuum
Hamiltonian is derived
from this radius law.

The next upstream quantitative
question is the uniform
coupling of this retained
sector to the growing
nonzero-mode system. The
positive finite-volume
normal determinant by
itself does not answer
it. The present theorem
supplies a normalization
and scaling test that
such a comparison must
pass; it does not replace
the comparison or settle
the Yang--Mills goal.

\subsection{Exact verification and preservation}

The certificate is
\begin{center}\small
\path{scripts/verify_ym_f14_v54_commuting_zero_modes.py}.
\end{center}
It uses rational arithmetic
and the Python standard
library. It verifies the
quaternion commutator
identity on \(729\)
rational unit-quaternion
pairs, the exact rank-one
normal formula on \(729\)
rational normal matrices,
and the rank of the
four-dimensional wedge
symbol at all \(81\)
vectors with entries
in \(\{-1,0,1\}\).
The symbolic identities
and all-parameter rank
proofs are given above;
these finite tests are
regressions, not their
substitutes.

An independent rational
normalization check
integrates the Vandermonde
polynomial on the ordered
three-eigenvalue simplex.
The ordered integral is
\(1/480\), and the full
one is \(1/80\), agreeing
with the angular coefficient
\(4\pi^6\). The proof of
that coefficient in this
note uses the explicit
normal geometry and does
not require a singular-value
integration formula.

The certificate checks the
full based-gauge dimension,
central nonzero-mode
dimension, and regular-flat
normal dimension at sides
four, eight, and sixteen.
It also checks the
coefficient ledger:
one central four-link
corner contributes
\(1/(128\pi^2)\), all
sixteen contribute
\(1/(8\pi^2)\), and
the radial power in
dimension four is
exactly \(-1\).
No numerical sampling
is used to assert an
asymptotic probability
or a continuum theorem.

V53 has 1,300,388 bytes and SHA--256
\begin{center}\scriptsize\ttfamily
C4C691569187C489C7DC2D2CB26855DA748A768F2003BA0172B0FC37E35B6575.
\end{center}
Its full body is preserved
except for the title
version and this terminal
insertion. V54 replaces
V53 as the single active
main note. Supplied
notes, archived lineages,
companion programmes,
and monographs remain
unchanged.

\begin{firewall}
The compact harmonic
normalization and its
logarithmic scale law
have been proved here,
including their reduction
from the actual Wilson
law at each fixed finite
four-torus. No cofinal
ultraviolet estimate,
interacting continuum
construction, or positive
continuum Hamiltonian
mass gap has been
established. The full
Yang--Mills goal remains
unproved and active.
\end{firewall}
\section{V55: the complete four-holonomy mean amplitude, a uniform determinant limit, and a mixed commuting response}
\label{sec:v55}

\subsection{Companion checkpoint and the direction actually changed}

The V55 checkpoint inspected the current SM--Einstein V72,
NS Stability V26, and Shell Mixing V16. Their source
fingerprints are unchanged from V50:
\begin{center}\scriptsize
\begin{tabular}{ll}
SM--Einstein V72 & 607,372 bytes\\
\multicolumn{2}{l}{\ttfamily F44F9745D43379E1672C5550411B58E97195A4C9278592D221DF5D3F644DC1F5}\\[2pt]
NS Stability V26 & 278,283 bytes\\
\multicolumn{2}{l}{\ttfamily 82C887F8C2C69E4F28FAD7C3DC8DE5B9099CB5A4BF431EBA1FFF3E91935978AD}\\[2pt]
Shell Mixing V16 & 623,261 bytes\\
\multicolumn{2}{l}{\ttfamily BF138E63BA826BC030D3A8B87E04FDD3B1EC712D269BB064674E96509B46E710}
\end{tabular}
\end{center}

SM V72 distinguishes transport of a distribution
from closure at a point-concentrated state.
That distinction is relevant to V54's
logarithmically distributed toron scale,
but the auxiliary many-copy SM limit
supplies no Yang--Mills cutoff estimate.
NS V26 computes particular rank-one
tensor norms; those do not control
the full noncommuting YM fluctuation
operator and are not imported.

Shell V15--V16 already proves the
literal Wilson and gauge determinant
tail estimates and the complete
one-holonomy root coefficient.
We do not repeat those results.
Its \emph{intermediate mean-fixed}
mode cancellation is instead
adapted to all four commuting
holonomies and identified with
V54's actual Gaussian amplitude.
We retain the change from
logarithmic to quaternion
source volume. The resulting
whole determinant, including
its low-order terms, has a
uniform four-variable limit
and a mixed quartic coefficient
not determined by the axial
formula alone.

This section, dated 15 September
2026, supplies one volume-uniform
part of V54's normalization.
It is still a nonzero-mode
Gaussian amplitude on the
commuting locus, not a uniform
remainder for the full
interacting integral. The
next companion checkpoint
is V60.

\subsection{The actual mean-fixed amplitude and its source volume}

Keep the based-gauge Coulomb slice
and all twelve constant source
coordinates of V54. Let
\(\mathcal A_M^{\log}(a)\) denote
the leading nonzero-mode
Gaussian amplitude before the
change \(a\mapsto v(a)\).
Thus, with the fixed Euclidean
normal volume used there,
\[
 \mathcal A_M^{\log}(a)
       =(2\pi)^{d_M/2}
                \frac{J_M(a,0)}{\sqrt{\det H_M(a)}} .
\]
The amplitude denoted \(A_M(v)\)
in \eqref{eq:v54-amplitude} uses
the twelve-dimensional
quaternion-vector source
volume \(d^{12}v\).
These are two expressions of
the same measured integral,
not two choices of its law.

Consider the commuting branch
\begin{equation}\label{eq:v55-background}
            a_\mu=\frac{\theta_\mu}{M}E_3,
                 \qquad \mu=1,\ldots,4 ,
\end{equation}
where \(E_3=i\sigma_3\) in
the quaternion convention
of V54. The four physical
cycle angles are
\(\theta=(\theta_1,\ldots,\theta_4)\).
All plaquettes on this branch
are identity. The mean-fixed
normal Hessian and the gauge
Jacobian are positive for a
sufficiently small common real
angle neighborhood; below we
may work on
\[
             \mathcal D_{\rho_0}
                 =\{z\in\mathbb C^4:|z_\mu|<\rho_0\},
                       \qquad \rho_0=1/64 .
\]
The formulas specify the
analytic continuations of
the amplitude germs at zero.
At each fixed real background
they remain the local
Laplace coefficients of
the actual finite integral.
No common nonlinear
fluctuation radius or
common \(\beta\)-remainder
is inferred from this
analytic continuation.

Use centered integer representatives
\(\mathcal K_M\) for \(\mathbb Z/M\mathbb Z\),
and put
\begin{equation}\label{eq:v55-dispersion}
 \Lambda_M(n,z)=
       4M^2\sum_{\mu=1}^4
               \sin^2\!\left(\frac{\pi n_\mu+z_\mu}{M}\right),
       \qquad n\in\mathcal K_M^4 .
\end{equation}
For complex \(z\), the square
is the holomorphic square,
not an absolute square.
Define
\begin{equation}\label{eq:v55-FM}
 F_M(z)=\sum_{\substack{n\in\mathcal K_M^4\\ n\ne0}}
                    \log\frac{\Lambda_M(n,z)}{\Lambda_M(n,0)}
\end{equation}
on the branch through zero.
Exactly the Fourier index
\(n=0\) is removed for all
sources. Its charged
coordinates are retained
harmonic variables, not
Gaussian variables that
can be reinstated when
their eigenvalues become
nonzero.

\begin{theorem}[Complete four-holonomy amplitude]
\label{thm:v55-amplitude}
On the stated branch,
\begin{equation}\label{eq:v55-mean}
 -\log\frac{\mathcal A_M^{\log}(z/M)}
                    {\mathcal A_M^{\log}(0)}
      =2F_M(z)
          -2\sum_{\mu=1}^4
                   \log\frac{\sin(z_\mu/M)}{z_\mu/M}.
\end{equation}
For V54's fixed quaternion
source volume the corresponding
formula is
\begin{equation}\label{eq:v55-quaternion}
 \widehat A_M(z):=
       \frac{A_M(v(z/M))}{A_M(0)}
         =e^{-2F_M(z)}
                       \prod_{\mu=1}^4\cos(z_\mu/M)^{-1}.
\end{equation}
These formulas contain the full
fine Haar and primed gauge
factors. No root constraint,
six-dimensional root return,
or coarse-Haar subtraction
belongs in this mean-fixed
amplitude.
\end{theorem}

\begin{proof}
For one charged complex color
component and a nonzero Fourier
mode \(k=2\pi n/M\), write
\[
 q_\mu=e^{ik_\mu}-1,\quad
 q_{+,\mu}=e^{i(k_\mu+2z_\mu/M)}-1,\quad
 \lambda_0=|q|^2,\quad \lambda_+=|q_+|^2
\]
first for real \(z\), and then
extend the resulting identities
holomorphically. The
logarithmic-to-right tangent
matrix is
\[
 E=\operatorname{diag}_\mu
       \left(e^{iz_\mu/M}
                   \frac{\sin(z_\mu/M)}{z_\mu/M}\right).
\]
The flat right-coordinate
Wilson Hessian is
\(H_+=\lambda_+I-q_+q_+^*\).
The logarithmic gauge scalar
in the ordinary Coulomb
plane \(q^\perp\) is
\[
 \begin{split}
 m(k)&=q^*E^{-1}q_+\\
 &=\sum_\mu\left\{
     |q_\mu|^2\frac{z_\mu}{M}\cot\frac{z_\mu}{M}
                    +\frac{2z_\mu}{M}\sin k_\mu\right\}.
 \end{split}
\]
The quotient density is the
one in \eqref{eq:v45-rho};
its determinant excludes
constant site gauges only.

If \(C_0\) is an orthonormal
basis of \(q^\perp\), the
codimension-one volume
identity gives
\[
 \det(C_0^*E^*H_+EC_0)
          =|\det E|^2\lambda_+^2
                                  \frac{|m(k)|^2}{\lambda_0}.
\]
To verify it, complete
\(C_0\) by \(q/|q|\), and
complete the target
transverse basis by
\(q_+/|q_+|\). The
transverse volume of
\(EC_0\) is
\[
       |\det E|\,
           \frac{|q^*E^{-1}q_+|}{|q||q_+|}.
\]
The three nonzero
eigenvalues of \(H_+\)
are \(\lambda_+\).
Multiplying these factors
proves the determinant
identity. This is the
same identity as Shell
V16's one-axis calculation,
now with all four diagonal
entries of \(E\) present.

Combining half the real
Hessian logarithm, the
fine Haar logarithm,
and the gauge determinant
cancels their tangent and
gauge factors mode by
mode. The two real
charged colors give
\(2\sum_{n\ne0}
\log(\Lambda_M(n,z)/\Lambda_M(n,0))\);
the opposite charge is
equivalent under \(n\mapsto-n\).
Neutral colors are
source independent.
There are \(M^4\) fine
Haar factors for each
direction but only
\(M^4-1\) integrated
Fourier modes. The
remaining constant-source
exponential volume is
the last term in
\eqref{eq:v55-mean}.
This verifies the actual
mean-fixed measure,
including its low-order
terms, before taking a
limit.

The real positive branches extend
through the displayed complex
polydisc. The free lower bound is
\[
                         \Lambda_M(n,0)\ge16|n|^2.
\]
A direct sine expansion gives
\[
 |\Lambda_M(n,z)-\Lambda_M(n,0)|
                    \le32\pi\rho_0|n|+32\rho_0^2
                         <8|n|^2\quad(n\ne0).
\]
Similarly the expression
for \(m\) is a perturbation
of \(\lambda_0\) with
relative size at most
\(C\rho_0+C\rho_0^2\);
directly bounding
\(\sum|\sin k_\mu|
\le2\sqrt{\lambda_0}\)
and \(\lambda_0\ge16/M^2\)
keeps it nonzero at
this radius. The
real restricted Hessian
is positive because
the gauge kernel is
transverse to \(EC_0\).
These facts choose the
branches in the formulas.

Finally, for one
three-component logarithm,
\[
 \det D_a v(a)
       =\cos|a|
                    \left(\frac{\sin|a|}{|a|}\right)^2 .
\]
The radial derivative
is \(\cos|a|\), and
the two tangent
derivatives are
\(\sin|a|/|a|\).
The amplitude for
\(d^{12}v\) is therefore
the amplitude for
\(d^{12}a\) divided by
the product of these
four determinants.
Applying this to
\eqref{eq:v55-mean}
cancels the sine factors
and proves
\eqref{eq:v55-quaternion}.
\end{proof}

Formula \eqref{eq:v55-quaternion}
evaluates the twelve-dimensional
amplitude at commuting arguments.
It does not declare the
four-dimensional commuting
locus to carry the full
retained measure. Nor is
it the density obtained
by pulling that full
measure back to physical
angle coordinates. All
field-independent source
volume factors remain in
V54's \(A_M(0)\) and
partition coefficient
\(C_M\); only an explicitly
normalized ratio is used
here.

\subsection{A complete heat representation with the harmonic mode removed}

Put
\[
 K_M(t,z)=\sum_{n\in\mathcal K_M^4}e^{-t\Lambda_M(n,z)},\qquad
 D_M(t,z)=K_M(t,z)-K_M(t,0)
                         -e^{-t\Lambda_M(0,z)}+1 .
\]
The subtraction in \(D_M\)
removes precisely the
source-dependent \(n=0\)
term, and restores the
free \(n=0\) term. Thus
\[
 D_M(t,z)=
      \sum_{n\ne0}
              \{e^{-t\Lambda_M(n,z)}
                            -e^{-t\Lambda_M(n,0)}\}.
\]
Frullani's formula gives
the exact finite identity
\begin{equation}\label{eq:v55-heat-F}
                     F_M(z)=-\int_0^\infty D_M(t,z)\,\frac{dt}{t}.
\end{equation}
Near zero this integral
must be estimated after
the complete mode sum.
Separate absolute bounds
on its low Born terms
would miss their
cancellation.

Let \(p_l(x)\) be the
probability that the
nearest-neighbor symmetric
walk on \(\mathbb Z^4\),
with eight equally likely
steps, is at \(x\) after
\(l\) steps. Expanding the
lattice heat semigroup gives
the exact winding identity
\begin{equation}\label{eq:v55-winding}
 K_M(t,z)-K_M(t,0)
   =M^4\sum_{l\ge0}
       e^{-8M^2t}\frac{(8M^2t)^l}{l!}
           \sum_{w\in\mathbb Z^4\setminus\{0\}}
                         p_l(Mw)(e^{2iw\cdot z}-1).
\end{equation}
A closed lattice path has
total lifted displacement
\(Mw\), so its twist is
exactly \(e^{2iw\cdot z}\).
There is no contribution
with \(l<M\).

\begin{lemma}[Uniform control of the small heat-time integral]
\label{lem:v55-small-time}
For \(z\) in the closed
polydisc of radius \(\rho_0\)
and all \(M\ge4\), there
are constants independent
of \(M\) such that, for
all sufficiently small
\(\delta>0\),
\begin{equation}\label{eq:v55-small-time}
 \begin{split}
 \int_0^\delta
     |K_M(t,z)-K_M(t,0)|\,\frac{dt}{t}
       &\le C\{\delta^2+e^{-c/\sqrt\delta}\},\\
 \int_0^\delta |D_M(t,z)|\,\frac{dt}{t}
       &\le C\delta .
 \end{split}
\end{equation}
\end{lemma}

\begin{proof}
We first record an elementary
walk estimate:
\begin{equation}\label{eq:v55-walk-bound}
                   p_l(x)\le C l^{-2}e^{-c|x|^2/l},
                               \qquad l\ge2 .
\end{equation}
The Fourier symbol of one
step is \(\frac14\sum_\mu\cos k_\mu\).
Its absolute value is
bounded strictly below one
away from neighborhoods
of \(0\) and
\((\pi,\pi,\pi,\pi)\);
near those points its
deficit is quadratic.
Integrating the \(l\)-th
absolute power of that
symbol gives
\(\sup_xp_l(x)\le C l^{-2}\).
Each coordinate has
bounded mean-zero
increments, whose
exponential generating
function is bounded
by \(e^{Cs^2l}\).
Optimizing in \(s\)
and using a union bound
gives
\(\mathbb P(|X_l|\ge R)
\le C e^{-cR^2/l}\).
Split a walk into two
halves. In the convolution
for \(p_l(x)\), either
the first endpoint or
the remaining displacement
has magnitude at least
\(|x|/2\). Bound the
other half by its
supremum and use the
tail estimate for that
half. This proves
\eqref{eq:v55-walk-bound}.

Integrating the Poisson
coefficient in
\eqref{eq:v55-winding}
before summing gives,
for \(l\ge1\),
\[
 \int_0^\delta
    e^{-8M^2t}\frac{(8M^2t)^l}{l!}\,\frac{dt}{t}
     =\frac1l\,
        \mathbb P\{\operatorname{Pois}(8M^2\delta)\ge l\}.
\]
Split the sum at
\(l=M^2\sqrt\delta\).
For the lower portion,
\eqref{eq:v55-walk-bound}
and
\(|e^{2iw\cdot z}-1|
\le2e^{2\rho_0|w|_1}\)
give a bound by
\[
                  CM^4
                   \sum_{M\le l\le M^2\sqrt\delta}
                         l^{-3}e^{-cM^2/l}.
\]
The winding Gaussian
absorbs the fixed
complex exponential
uniformly when
\(l/M^2\le\sqrt\delta\).
If this range is
nonempty, \(M\ge\delta^{-1/2}\).
Putting \(s=l/M^2\)
and comparing the
increasing small-\(s\)
function \(s^{-3}e^{-c/s}\)
with its integral bounds
the display by
\(Ce^{-c'/\sqrt\delta}\).

For the upper portion,
the Poisson Chernoff
bound is
\[
 \mathbb P\{\operatorname{Pois}(8M^2\delta)\ge l\}
                     \le(8e\sqrt\delta)^l .
\]
A winding endpoint
satisfies
\(|w|_1\le l/M\).
Since the sum of the
endpoint probabilities
is at most one, this
portion is bounded by
\[
 CM^4\sum_{l\ge M}\frac{q^l}{l},
           \qquad
 q=8e\,e^{\rho_0/2}\sqrt\delta .
\]
Take \(\delta\) small
enough that \(q<1/2\).
The expression is at
most \(CM^3q^M\).
Uniformly over \(M\ge4\)
this is \(O(q^4)=O(\delta^2)\).
This proves the first
bound.

The remaining zero-mode
term has
\(|\Lambda_M(0,z)|\le C|z|^2\).
Consequently
\(|e^{-t\Lambda_M(0,z)}-1|\le Ct\)
for \(0<t\le\delta\),
uniformly in the
same polydisc and
in \(M\). Its
integral contributes
at most \(C\delta\),
proving the second
bound.
\end{proof}

The winding argument controls
the complete ultraviolet part.
It does not require an
unevaluated cancellation of
separately divergent second
or fourth traces.

\subsection{The uniform analytic limit and a quantitative rate}

For the unit continuum torus,
let
\[
 \begin{split}
 \Lambda_\infty(n,z)
       &=\sum_{\mu=1}^4(2\pi n_\mu+2z_\mu)^2,\\
 K_\infty(t,z)
       &=\sum_{n\in\mathbb Z^4}e^{-t\Lambda_\infty(n,z)},\\
 D_\infty(t,z)
       &=K_\infty(t,z)-K_\infty(t,0)
                                  -e^{-4t z\cdot z}+1 .
 \end{split}
\]
Define the grouped integral
\begin{equation}\label{eq:v55-Finfty}
                      F_\infty(z)=
                         -\int_0^\infty D_\infty(t,z)\,\frac{dt}{t}.
\end{equation}

\begin{theorem}[Complete determinant convergence through zero holonomy]
\label{thm:v55-limit}
The function \(F_\infty\) is
holomorphic on \(\mathcal D_{\rho_0}\),
with \(F_\infty(0)=0\), and
on every smaller closed
polydisc
\begin{equation}\label{eq:v55-rate}
                 \sup_z |F_M(z)-F_\infty(z)|
                                   \le C M^{-1/2}.
\end{equation}
Every fixed source derivative
has the same rate on a
further smaller polydisc.
In particular the normalized
actual quaternion-source
amplitude satisfies
\begin{equation}\label{eq:v55-amplitude-limit}
              \widehat A_M(z)\longrightarrow e^{-2F_\infty(z)}
\end{equation}
locally uniformly with
all fixed derivatives,
at rate \(O(M^{-1/2})\).
On each smaller real
angle neighborhood there
are \(0<c<C<\infty\),
independent of \(M\),
such that
\[
                             c\le\widehat A_M(\theta)\le C .
\]
\end{theorem}

\begin{proof}
Poisson summation for the
Gaussian gives
\begin{equation}\label{eq:v55-Poisson}
 K_\infty(t,z)
     =(4\pi t)^{-2}
       \sum_{w\in\mathbb Z^4}
                e^{-|w|^2/(4t)}e^{2iw\cdot z}.
\end{equation}
The \(w=0\) term cancels
in \(K_\infty(t,z)-K_\infty(t,0)\).
The remaining terms
are exponentially small
as \(t\downarrow0\),
uniformly on the
complex polydisc.
The removed harmonic
term is \(O(t)\).
Thus the continuum
analogue of the
second estimate in
\eqref{eq:v55-small-time}
holds.

For \(n\ne0\), the
dispersion bound used
above gives
\[
 \operatorname{Re}\Lambda_M(n,z)\ge8|n|^2,\qquad
 \operatorname{Re}\Lambda_\infty(n,z)\ge8|n|^2
\]
on this polydisc.
It follows that the
nonzero-mode expressions
for both \(D\)'s
are bounded by
\(Ce^{-ct}\) for
\(t\ge1\).
These estimates already
prove convergence and
holomorphy of
\eqref{eq:v55-Finfty}
and uniform boundedness
of all the grouped
integrals.

To obtain a rate, Taylor's
formula for the sine on
its fixed complex
neighborhood gives
\[
 |\Lambda_M(n,z)-\Lambda_\infty(n,z)|
             \le C M^{-2}(|n|+1)^4
                  \quad(n\in\mathcal K_M^4).
\]
Both endpoints and
their linear interpolation
retain the displayed
real-part lower bound.
Hence their exponential
difference is at most
\[
          C M^{-2}t(|n|+1)^4 e^{-ct|n|^2}
                   \qquad(n\ne0).
\]
Summing in four dimensions
gives \(CM^{-2}t^{-3}\)
for \(0<t\le1\).
The omitted continuum
modes outside
\(\mathcal K_M^4\)
contribute at most
\(Ct^{-2}e^{-ctM^2}\).
For \(t\ge1\) the
integrated difference
is \(O(M^{-2})\).

Split the heat integral
at \(\delta=M^{-1/2}\)
and at one. For all
sufficiently large \(M\),
the small-time bounds
give \(C\delta\).
The middle range gives
\[
       CM^{-2}\int_\delta^1t^{-4}\,dt
                         \le CM^{-2}\delta^{-3}
                         =CM^{-1/2}.
\]
The missing-mode term
is exponentially small
because \(\delta M^2=M^{3/2}\).
The finitely many smaller
\(M\)'s are absorbed in
the constant. This proves
\eqref{eq:v55-rate}
uniformly on the complex
polydisc. Cauchy estimates
give the differentiated
rate on smaller discs.

In \eqref{eq:v55-quaternion},
the extra logarithm
\(-\sum_\mu\log\cos(z_\mu/M)\)
is \(O(M^{-2})\) with
all fixed source derivatives
on smaller discs.
Exponentiating the
uniformly bounded
logarithms proves
\eqref{eq:v55-amplitude-limit}
and the positive real
upper and lower bounds.
\end{proof}

Limits of twisted discrete
Laplacian determinants are
classical; a primary reference
is
\href{https://arxiv.org/abs/1607.01566}
{F.~Friedli, The bundle Laplacian on discrete tori}.
The proof here retains the
specified harmonic subtraction
through zero holonomy, supplies
a common complex neighborhood
and a conservative rate, and
identifies the complete Wilson
mean amplitude with its actual
source-volume factor. It is
not an application of a
fixed nonzero-twist formula
without checking its
zero-holonomy endpoint.

\subsection{A four-variable identity and the mixed commuting quartic}

There is a useful additional
constraint on the limiting
function. Write
\(\Delta_z=\sum_{\mu=1}^4\partial_{z_\mu}^2\),
understood first in real
angle coordinates.

\begin{proposition}[Biharmonic identity near the retained zero mode]
\label{prop:v55-biharmonic}
On the real neighborhood of zero,
and by analytic continuation
on its complex neighborhood,
\begin{equation}\label{eq:v55-biharmonic}
                           \Delta_z^2F_\infty(z)=\frac{16}{\pi^2}.
\end{equation}
\end{proposition}

\begin{proof}
For real \(\theta\notin\pi\mathbb Z^4\),
define the periodic regularized
full determinant logarithm
\[
 \mathcal D(\theta)=
      -\int_0^\infty
          \left\{K_\infty(t,\theta)
             -\mathbf1_{\{t\le1\}}(4\pi t)^{-2}\right\}\frac{dt}{t}.
\]
The small-time subtraction
removes the \(w=0\)
Poisson term, and the
large-time integral
converges off zero
holonomy. This periodic
function is locally
integrable through its
logarithmic singularities.
For \(w\ne0\), its
Fourier coefficient
on the period-\(\pi\)
torus is
\[
 -\frac1{16\pi^2}
       \int_0^\infty t^{-3}e^{-|w|^2/(4t)}\,dt
                             =-\frac1{\pi^2|w|^4}.
\]
The exchange for this
coefficient is justified
by its integrable
positive Gaussian
majorant; the constant
coefficient is irrelevant
to the following derivatives.
Thus, as periodic
distributions,
\[
 \Delta_\theta^2\mathcal D
        =\frac{16}{\pi^2}
             -16\pi^2\sum_{m\in\mathbb Z^4}\delta_{\pi m}.
\]
Here
\(\sum_w e^{2iw\cdot\theta}
=\pi^4\sum_m\delta_{\pi m}\).

Near zero, remove the
one eigenvalue
\(\Lambda_\infty(0,\theta)=4|\theta|^2\).
Frullani's formula
shows directly from
the heat definitions
that
\[
 F_\infty(\theta)
     =\mathcal D(\theta)
                 -\log(4|\theta|^2)-C_0
\]
off zero, with \(C_0\)
chosen so the smooth
extension has value
zero at the origin.
In four dimensions,
\(\Delta\log|x|^2=4|x|^{-2}\)
off zero and
\(\Delta |x|^{-2}=-4\pi^2\delta_0\).
The latter coefficient
follows by the normal
flux through a small
three-sphere of area
\(2\pi^2r^3\).
Consequently
\[
                  \Delta^2\log(4|x|^2)=-16\pi^2\delta_0.
\]
The retained-mode
subtraction cancels
the central delta
but not the constant.
This gives
\eqref{eq:v55-biharmonic}
as distributions near
zero, hence as analytic
functions by Theorem
\ref{thm:v55-limit}.
\end{proof}

Hypercubic symmetries and independent
angle sign changes imply an expansion
\begin{equation}\label{eq:v55-quartic}
 F_\infty(z)=
       c_2\sum_\mu z_\mu^2
         +a_4\sum_\mu z_\mu^4
         +b_4\sum_{\mu<\nu}z_\mu^2z_\nu^2
         +O(|z|^6).
\end{equation}
The axial coefficient is already
accessible from Shell V16's
mean-fixed one-axis formula.
For clarity put
\[
 S_j=\sum_{m\in\mathbb Z^3\setminus\{0\}}
                     (\cosh(2\pi|m|)-1)^{-j}.
\]
These are exponentially
convergent positive sums.
The inherited one-axis
Chebyshev product gives
\[
 F_\infty(x,0,0,0)
      =2\log\frac{\sin x}{x}
       +\sum_{m\ne0}
             \log\left(1+\frac{1-\cos2x}{\cosh(2\pi|m|)-1}\right).
\]
In particular
\begin{equation}\label{eq:v55-axial}
             c_2=-\frac13+2S_1,\qquad
             a_4=-\frac1{90}-\frac23S_1-2S_2 .
\end{equation}
These are coefficients of \(F_\infty\),
not of Shell V16's root-constrained
action. The latter contains additional
source-return and reference terms.

\begin{theorem}[Mixed commuting quartic fixed by the complete determinant]
\label{thm:v55-mixed}
The coefficient in
\eqref{eq:v55-quartic} is
\begin{equation}\label{eq:v55-b4}
             b_4=\frac1{3\pi^2}-2a_4
                 =\frac1{3\pi^2}+\frac1{45}
                              +\frac43S_1+4S_2>0.
\end{equation}
The coefficients of the complete
finite functions \(F_M\) converge
to \eqref{eq:v55-axial} and
\eqref{eq:v55-b4} at the
rate in Theorem
\ref{thm:v55-limit}.
\end{theorem}

\begin{proof}
In four variables,
\[
 \Delta_z^2\sum_\mu z_\mu^4=96,\qquad
 \Delta_z^2\sum_{\mu<\nu}z_\mu^2z_\nu^2=48.
\]
Applying Proposition
\ref{prop:v55-biharmonic}
at zero therefore gives
\[
                       96a_4+48b_4=16/\pi^2.
\]
Solve for \(b_4\) and use
\eqref{eq:v55-axial}.
The differentiated uniform
convergence proves the
last assertion.
\end{proof}

All four backgrounds in this
calculation commute and have
zero curvature. The positive
mixed coefficient in
\eqref{eq:v55-b4} is a
finite-volume holonomy
response, not the
noncommuting mixed
\(F^2\) calculation
already completed in V30.
It is neither a local
gluon mass nor a new
beta-function coefficient.
The nonzero constant in
\eqref{eq:v55-biharmonic}
also explains why
discarding every low
polynomial term from
a convergent modified
determinant would lose
information here.

\subsection{What is now uniform and what remains upstream}

V54's \(M\)-dependent Gaussian
amplitude is now normalized
and uniformly controlled
on the whole small commuting
four-holonomy family at
the physical scaling
\(a_\mu=\theta_\mu E_3/M\).
The limit includes the
low-order terms, admits
all fixed source
derivatives, and has
an explicit mixed
commuting quartic
coefficient. The
quaternion source-volume
factor was kept at
every finite \(M\).

This does not control
the same amplitude in
transverse noncommuting
source directions.
In particular a
logarithm multiplying
the commutator quartic
can vanish on every
commuting background.
Nor does it bound
the full finite-\(\beta\)
normal integral relative
to this amplitude
uniformly in \(M\).
V54's source-normal
Gaussian remainder
still has fixed-volume
constants.

Thus no joint limit
with \(\beta\) of
order \(\log M\),
no full retained
toron probability law
in that limit, and
no physical mass
gap follows here.
The next required
comparison must retain
the transverse source
response and the
interacting nonzero
modes together. A
flat amplitude or
an auxiliary determinant
floor cannot replace
that estimate.

\subsection{Executed checks and append record}

The exact certificate is
\begin{center}\small
\path{scripts/verify_ym_f14_v55_flat_toron_amplitude.py}.
\end{center}
It expands the full
nonzero mode sums
independently on sides
two and four. At
side two it obtains
\[
 [z_1^2]F_2=\frac{13}{48},\qquad
 [z_1^4]F_2=-\frac{323}{1536},\qquad
 [z_1^2z_2^2]F_2=\frac7{256}.
\]
For the actual logarithmic
mean amplitude, its
negative logarithm has
coefficients
\[
                              5/8,\qquad -4837/11520.
\]
After changing to the
quaternion source volume,
they become \(5/12\)
and \(-109/256\).
The difference is
the measured source
Jacobian, not a
choice of renormalization.
Side two is used only
as this finite regression;
the uniform theorems
are stated for \(M\ge4\).

The source-volume cancellation
is checked through sixth
order. The finite
mode quartic coefficients
also check the
one-axis formula from
the companion without
using it to generate
the mode-sum side.
The finite biharmonic
values need not equal
the continuum constant:
for example they are
\(-151/8\) at side
two and
\(-6186618259/1728720000\)
at side four.
The limiting identity
was proved from the
complete heat kernel,
not inferred from
those values.

The winding regression
compares an independently
iterated finite seam-twisted
walk with the lifted
\(\mathbb Z^4\) walk.
For all sixteen central
twist choices and
lengths zero through
ten, it checks \(176\)
exact closed-walk
identities. Twists are
invisible below winding
length four and all
fifteen nontrivial
choices are detected
at that first possible
length. These algebraic
checks do not replace
the uniform small-time
estimate or its
complex-domain proof.

V54 has 1,331,920 bytes and SHA--256
\begin{center}\scriptsize\ttfamily
A00211A2CE755C1BAEEC9FFEAAAB9EA7CCB2EC08CAFC9D7C5C981E6D4010786F.
\end{center}
Its entire body is
preserved except for
the title version and
this terminal insertion.
V55 replaces V54 as
the single active
main note. The
checkpoint companion
files, archived
lineages, supplied
notes, and monographs
remain unchanged.

\begin{firewall}
A complete mean-fixed
Gaussian amplitude now
has a volume-uniform
four-holonomy limit
with its actual
source normalization.
This is not a
uniform interacting
weak-coupling integral,
a continuum theory,
or a physical
Hamiltonian gap.
The full Yang--Mills
mass-gap objective
remains unproved
and active.
\end{firewall}
\section{V56: scale-local transverse Gaussian regulators and their actual Wilson action error}
\label{sec:v56}

\subsection{The upstream move and the non-duplication audit}

This section was added on 15 September 2026. V55 proves convergence
of the complete Gaussian amplitude on a four-angle commuting
family. It does not estimate the interacting normal integral
uniformly in the cutoff. V23 already proves that a fixed global
counting-\(\ell^4\) ball has negligible free probability on the
logarithmic weak-coupling scale; V46 also treats the actual
interacting law restricted to its global energy ball. Repeating
either calculation would not remove that bottleneck.

There is also an important positive prior result. Archive f7 V63,
in \emph{Correlated Gaussian multi-block tail}, proves an abstract
bad-set estimate from a bounded covariance, and its subsequent
local marginal density criterion explains how such estimates
could pass to an interacting law. The closure monograph already
contains a positive conditional Gaussian shell decomposition.
Neither abstract statement is claimed as new here.

The new input is a scale-dependent, derivative-sensitive regulator
for the \emph{actual zero-background transverse Maxwell covariance}.
We verify both its operator bound and its local trace bound as
the number of links in a block grows. The local trace, rather
than the dimension of that block, is the crucial improvement.
We then prove a deterministic Wilson plaquette estimate for the
sum of the coarser fields on the resulting local domain.
This supplies a concrete free reference and an actual nonlinear
action bound, not the missing interacting probability estimate.

The heat decomposition below is an auxiliary identity, not the
conditional shell of the literal balanced root map. It has no
asserted finite spatial range. For context, positive \emph{finite-range}
decompositions have a distinct construction in
\href{https://arxiv.org/abs/1206.2212}
{R. Bauerschmidt, \emph{A simple method for finite range decomposition
of quadratic forms and Gaussian fields}}. No finite-range conclusion
from that paper is imported into this heat construction.
The V55 companion checkpoint remains current; the next required
SM/NS/shell review is V60.

\subsection{An exact positive decomposition with the gauge modes removed correctly}

Let \(M=2^J\ge4\), with counting measure on the four-dimensional
periodic lattice. Use the forward gradient \(D\), the link curl
\(d\), and the componentwise scalar Laplacian
\(\Delta=D^*D\). On one-forms the exact identity is
\[
                  d^*d+DD^*=\Delta I_4.
\]
There are three real colour components. Write
\[
 {\cal H}_M=\{u:D^*u=0,\ \overline u_\mu=0\text{ for every }\mu\},
 \qquad \dim{\cal H}_M=9(M^4-1).
\]
The removed twelve harmonic coordinates are retained source
variables in the full programme, not assigned a Gaussian mass.
In this section only, their background value is zero.
The normalized Gaussian on \({\cal H}_M\) has action
\(\frac12\|du\|_2^2\) and covariance \(C_M\).

At momentum \(k=2\pi n/M\ne0\), put
\[
 q_\mu=e^{ik_\mu}-1,\qquad
 \lambda=|q|^2,\qquad
 \Pi(k)=I_4-\frac{q q^*}{\lambda}.
\]
Thus \(C_M(k)=\Pi(k)/\lambda\), independently in each colour.
All operators defined below vanish at \(k=0\).
Let \(h(s)=(1+s)e^{-s}\). Set \(L_j=2^j\), \(0\le j<J\),
and define covariance multipliers
\begin{equation}\label{eq:v56-shells}
 \begin{aligned}
 c_{\mathrm u}(\lambda)&=\frac{1-h(\lambda)}{\lambda},
       &L_{\mathrm u}&=1,\\
 c_j(\lambda)&=\frac{h(L_j^2\lambda)-h(4L_j^2\lambda)}{\lambda},
       &0&\le j<J,\\
 c_J(\lambda)&=\frac{h(M^2\lambda)}{\lambda},
       &L_J&=M.
 \end{aligned}
\end{equation}
The matrix covariance of level \(j\) is \(c_j(\lambda)\Pi(k)\);
the same convention applies to the extra ultraviolet level
\(\mathrm u\). Denote these operators by \(C_j\).

\begin{lemma}[Exact free measure, including the ultraviolet and infrared ends]
\label{lem:v56-decomposition}
The operators \(C_j\) are nonnegative and
\[
                     C_M=C_{\mathrm u}+\sum_{j=0}^{J}C_j.
\]
Consequently, independent real Gaussians \(Z_j\) with these
covariances satisfy
\[
 A=\beta^{-1/2}\left(Z_{\mathrm u}+\sum_{j=0}^{J}Z_j\right)
       \ \sim\ N(0,\beta^{-1}C_M).
\]
Every summand is transverse and mean zero.
\end{lemma}
\begin{proof}
On the full one-form space, let \(K=d^*d\). It is nonnegative
and commutes with \(\Delta\). The operator identity
\[
             C_M=\int_0^\infty t e^{-t\Delta}K\,dt
\]
holds on each nonzero Fourier mode, since \(K(k)=\lambda\Pi(k)\)
and \(\int_0^\infty t\lambda e^{-t\lambda}dt=1/\lambda\).
Both sides vanish on gradients and constants. Split this integral
into \([0,1]\), the intervals \([L_j^2,4L_j^2]\), and
\([M^2,\infty)\). Integration gives \eqref{eq:v56-shells}.
Each integral is positive; equivalently \(h'(s)=-se^{-s}\le0\).
The endpoint terms telescope. Covariances of independent Gaussians
add, proving the law, including its normalization.
\end{proof}

\subsection{Uniform derivative bounds, not a rank bound}

For an integer \(r\ge0\), \(\nabla^r u\) denotes all ordered
forward differences of order \(r\), in all directions, link
components and colours. In particular
\[
             \|\nabla^r u\|_2^2=\langle u,\Delta^r u\rangle.
\]
For a field at length \(L\), define its global quadratic regulator
\begin{equation}\label{eq:v56-global-regulator}
 {\cal Q}_{L}(u)
       =\sum_{r=0}^{4}L^{2r-2}\|\nabla^r u\|_2^2.
\end{equation}
This regulator is used for estimates, not inserted as a mass
or as a new term in the Wilson action.

\begin{lemma}[The scale-uniform operator and diagonal estimates]
\label{lem:v56-operator}
There are constants \(a_r<\infty\), independent of \(M,j\), such that
\begin{equation}\label{eq:v56-diagonal}
 \mathbb E|\nabla^r Z_j(x)|^2\le a_r L_j^{-2r-2},
                  \qquad 0\le r\le4.
\end{equation}
The same statement holds at level \(\mathrm u\), with \(L_{\mathrm u}=1\).
Moreover, if \(T_Lu=(L^{r-1}\nabla^r u)_{r=0}^4\), then
\begin{equation}\label{eq:v56-op}
            \|T_{L_j}C_j^{1/2}\|^2\le B_0,\qquad B_0=2^{15}.
\end{equation}
\end{lemma}
\begin{proof}
The square of the operator norm is bounded by the supremum of
\[
       P_L(\lambda)c_j(\lambda),\qquad
       P_L(\lambda)=L^{-2}\sum_{r=0}^{4}(L^2\lambda)^r.
\]
For an ordinary dyadic level, write \(s=L^2\lambda\).
The product is
\[
        \left(\sum_{r=0}^{4}s^r\right)\frac{h(s)-h(4s)}s.
\]
If \(0<s\le1\), use
\(h(s)-h(4s)=\int_s^{4s}te^{-t}dt\le15s^2/2\);
the product is at most \(75/2\). If \(s\ge1\), it is at most
\(10s^4e^{-s}<50\). The latter bound follows by maximizing
\(s^4e^{-s}\) at \(s=4\) and, for example, bounding \(e^4\)
below by its Taylor sum through degree eight.
At the infrared end the same large-\(s\) bound applies because
the centered-frequency estimate gives \(M^2\lambda\ge16\).
At level \(\mathrm u\), \(0<\lambda\le16\) and
\[
 c_{\mathrm u}(\lambda)\le\min(\lambda/2,\lambda^{-1}).
\]
For \(\lambda\le1\) this gives \(5/2\); for \(\lambda\ge1\)
it gives \(5\lambda^3\le20480<2^{15}\).
The transverse projection has norm one, so these bounds prove
\eqref{eq:v56-op}, including all colour components.

For the diagonal estimate, first note the Fourier bound, for
every fixed integer \(m\ge0\) and \(1\le L\le M\),
\begin{equation}\label{eq:v56-heat-moment}
 M^{-4}\sum_{n\ne0}\lambda(n)^m e^{-L^2\lambda(n)}
                              \le b_m L^{-4-2m}.
\end{equation}
Indeed, centered indices satisfy
\[
             16|n|^2/M^2\le\lambda(n)\le4\pi^2|n|^2/M^2.
\]
Comparison with four-dimensional annuli, after rescaling by
\(M/L\ge1\), proves the displayed sum with a finite constant \(b_m\).
One can bound each annulus by its number of integer points
times a polynomial times \(e^{-c|n|^2L^2/M^2}\); the resulting
geometric-Gaussian series converges uniformly.

At an ordinary level,
\[
 c_j(\lambda)=\lambda\int_{L^2}^{4L^2}te^{-t\lambda}dt
                         \le \frac{15}{2}L^4\lambda e^{-L^2\lambda}.
\]
Apply \eqref{eq:v56-heat-moment} with \(m=r+1\).
Translation invariance turns the trace per site into the
diagonal; the three colours and at most four transverse
components only change the constant.
At the infrared end set \(s=M^2\lambda\); the sum over
centered nonzero \(n\) of
\(s^{r-1}(1+s)e^{-s}\) is uniformly finite, including \(r=0\)
since \(s\ge16\). Its prefactor is \(M^{-2r-2}\).
At the ultraviolet end the finite spectral range and
\(c_{\mathrm u}\le\lambda/2\) give a uniform diagonal bound.
\end{proof}

The low-frequency behavior \(c_j(\lambda)=O(L_j^4\lambda)\)
is essential here. The \emph{sum of all coarser covariances}
does not obey \eqref{eq:v56-op} with \(L=L_j\):
its \(L_j^{-2}C\) eigenvalue grows like \((M/L_j)^2\).
Thus one must not replace a single level by the unsplit
massless tail in the following product estimate.

\subsection{Local regulators with a bounded overlap cost}

At each ordinary level partition the torus into aligned
dyadic vertex cubes \(B\) of side \(L_j\). At level
\(\mathrm u\) use the unit cubes again, now for its independent
field; level \(J\) has one cube. Let \(B^+\) be the vertices
within periodic coordinate distance \(10L_j\) of \(B\).
Forward differences are anchored at the indicated vertices.
Set
\begin{equation}\label{eq:v56-local-regulator}
 Q_{j,B}(u)=
       \sum_{r=0}^{4}L_j^{2r-2}
                  \sum_{x\in B^+}|\nabla^r u(x)|^2.
\end{equation}
These enlargements have size at most \((23L_j)^4\), and at
one scale each vertex belongs to at most \(K_*=23^4\) of them.
The deliberately loose bound also covers periodic wraparound
and the infrared end.

\begin{theorem}[Actual derivative-regulator bad-set bound]
\label{thm:v56-badsets}
There are finite constants \(A_*,B_*>0\), independent of
\(M\), the level, and the chosen blocks, with
\(B_*=K_*B_0\), such that, for any collection \({\cal X}_j\)
of distinct blocks at one level,
\begin{equation}\label{eq:v56-trace-op}
 \begin{split}
 \left\|\operatorname{Cov}
       \bigl((L_j^{r-1}\nabla^r Z_j|_{B^+})_{B\in{\cal X}_j,r\le4}
          \bigr)\right\|&\le B_*,\\
 \operatorname{Tr}\operatorname{Cov}
       \bigl((L_j^{r-1}\nabla^r Z_j|_{B^+})_{B\in{\cal X}_j,r\le4}
          \bigr)&\le A_*|{\cal X}_j|.
 \end{split}
\end{equation}
Put \(\eta=(4B_*)^{-1}\), \(a_*=2\eta A_*\), and
\(p_R=\exp(a_*-\eta R^2)\). When \(p_R<1\), any finite set
\({\cal X}\) of distinct level-block pairs satisfies
\begin{equation}\label{eq:v56-bad-product}
 \mathbb P\{Q_{j,B}(\beta^{-1/2}Z_j)>R^2/\beta
                  \text{ for every }(j,B)\in{\cal X}\}
                                      \le p_R^{|{\cal X}|}.
\end{equation}
\end{theorem}
\begin{proof}
Duplicating coordinates in the enlarged blocks costs at most
their overlap \(K_*\) in the squared operator norm. Thus
\eqref{eq:v56-op} proves the first inequality.
For each block, \eqref{eq:v56-diagonal} bounds its trace by
\[
 \sum_{r=0}^{4}
        L_j^{2r-2}(23L_j)^4 a_rL_j^{-2r-2}
                             \le23^4\sum_{r=0}^{4}a_r.
\]
This proves the second inequality with this value as \(A_*\).
In particular it does not cost the number \(O(L_j^4)\)
of coordinates in the block.

For completeness, apply the determinant identity used in f7 V63
to this particular covariance \(K_{\cal X}\). For
\(0\le2\eta K_{\cal X}\le I/2\),
\[
 \log\mathbb E\exp\!\left(\eta\sum_{B\in{\cal X}_j}
                                      Q_{j,B}(Z_j)\right)
 =-\tfrac12\log\det(I-2\eta K_{\cal X})
 \le2\eta\operatorname{Tr}K_{\cal X}
 \le a_*|{\cal X}_j|.
\]
Zero covariance eigenvalues cause no problem. Exponential
Markov inequality gives the desired bound for one scale.
The independent level fields from Lemma~\ref{lem:v56-decomposition}
then multiply those bounds for arbitrary \({\cal X}\).
Spatial block independence has not been assumed.
\end{proof}

This is a joint-event bound, not a claim of stochastic domination
under every increasing event or of a uniform conditional bound
after arbitrary other block values have been fixed.

\begin{proposition}[Multiscale bad clusters and the logarithmic schedule]
\label{prop:v56-clusters}
Join nearest-neighbor blocks at each scale, and join each block
to its dyadic parent. Join each level-\(\mathrm u\) block to the
same level-zero block. This graph has maximum degree at most \(25\).
If \({\cal C}(v)\) is the connected bad cluster of a fixed node,
empty when that node is good, then
\begin{equation}\label{eq:v56-clusters}
          \mathbb P\{|{\cal C}(v)|\ge k\}
                  \le p_R(625p_R)^{k-1},\qquad k\ge1.
\end{equation}
In particular the bound is uniform in the number of scales and
the torus size whenever \(625p_R<1\).
\end{proposition}
\begin{proof}
A node has at most eight horizontal neighbors, sixteen children,
and one parent. The extra ultraviolet nodes do not increase
that bound. A connected set of size \(k\) containing \(v\)
has a spanning-tree traversal of length \(2k-2\) starting at
\(v\). Fixing an ordering to select one traversal for each
set, its visited vertex set recovers that set. Thus there
are at most \(25^{2k-2}\) such sets.
If the cluster has size at least \(k\), it contains one of them.
Use \eqref{eq:v56-bad-product} and a union bound. This is an
application of the archived bad-polymer counting argument to
the newly verified regulator, not a new abstract polymer lemma.
\end{proof}

There are exactly \((31M^4-1)/15<3M^4\) level-block nodes.
Consequently the event that all nodes are good has probability
at least \(1-3M^4p_R\). This is sometimes a useful finite-volume
bound, but is not the required thermodynamic good event.
For the local construction use instead
\begin{equation}\label{eq:v56-schedule}
 R^2=\kappa\log\beta,\qquad
 \epsilon=\frac{R}{\sqrt\beta},\qquad
 p_R=e^{a_*}\beta^{-\eta\kappa},\qquad \beta>1.
\end{equation}
If \(\beta(M)\ge c\log M\) for a fixed \(c>0\), and
\(\eta\kappa>1\), a fixed block's entire chain of coarser
ancestors is good with probability tending to one: its
failure probability is at most \((J+2)p_R\).
This conclusion only uses the sufficient logarithmic schedule;
it does not identify the coefficient of an interacting running
coupling. A global absence of bad blocks is neither assumed
nor deduced with this choice of \(R\).

\subsection{A local Sobolev estimate and summing the coarser fields}

\begin{lemma}[The regulator controls the quantities the plaquette uses]
\label{lem:v56-Sobolev}
Uniformly for \(1\le L\le M\), let \(B\) be a dyadic \(L\)-cube
and use \eqref{eq:v56-local-regulator}. On the vertices within
distance \(2L\) of \(B\),
\begin{equation}\label{eq:v56-sup}
       L\|u\|_\infty+L^2\|\nabla u\|_\infty
                                     \le C Q_B(u)^{1/2}.
\end{equation}
The norms on the left include the finite vector components.
\end{lemma}
\begin{proof}
The global Fourier Cauchy--Schwarz inequality gives
\[
 \|f\|_\infty
 \le C L^{-2}
       \left\{\sum_{r=0}^{3}L^{2r}\|\nabla^r f\|_2^2\right\}^{1/2}.
\]
Indeed, with a unitary Fourier transform the squared prefactor
is bounded by
\[
 M^{-4}\sum_n(1+L^2\lambda(n))^{-3}\le C L^{-4}.
\]
The constant mode contributes \(M^{-4}\le L^{-4}\); the other
terms follow by four-dimensional annular summation. Expanding
\((1+L^2\lambda)^3\) changes only fixed binomial constants.

Choose a sampled smooth cutoff equal to one within distance
\(2L\) of \(B\), supported within distance \(4L\), with its
differences of order \(s\le4\) bounded by \(C_sL^{-s}\).
For large \(L\) relative to \(M\) one may take the constant
cutoff one; then \(B^+\) is the whole torus. In the other
cases a product of one-dimensional cutoffs on the periodic
coordinate intervals has the asserted bounds.
The discrete product rule and its shifted factors through
order four only use vertices inside \(B^+\).
Applying the global estimate to the cutoff times \(u\)
proves the first term of \eqref{eq:v56-sup}. Apply it
separately to the cutoff times each component of \(\nabla u\);
the terms through order three now use derivatives of \(u\)
through order four. Multiplying by \(L^2\) gives exactly
the weights in \(Q_B\), proving the second term.
\end{proof}

Let \(a_j=\beta^{-1/2}Z_j\). For an ordinary level \(k\) and
one of its blocks \(B\), put
\[
              A_{\ge k}=\sum_{j=k}^{J}a_j.
\]
Suppose every ancestor \(B_j\) of \(B\) satisfies
\(Q_{j,B_j}(a_j)\le\epsilon^2\). The containing cubes are
nested, and their length-\(L_j\) neighborhoods contain
the vertices of every plaquette anchored in \(B\).
Lemma~\ref{lem:v56-Sobolev} and two geometric series give
\begin{equation}\label{eq:v56-tail-local}
 \begin{aligned}
 \|A_{\ge k}\|_{\infty,B+\text{one plaquette}}
       &\le C\epsilon\sum_{j=k}^{J}L_j^{-1}
                                      \le C'\epsilon/L_k,\\
 \|\nabla A_{\ge k}\|_{\infty,B+\text{one plaquette}}
       &\le C\epsilon\sum_{j=k}^{J}L_j^{-2}
                                      \le C'\epsilon/L_k^2.
 \end{aligned}
\end{equation}
For the finest full field include its level-\(\mathrm u\)
block as well; it adds one more term at length one.
There is no factor \(J\) in these deterministic bounds.
In particular the local counting-\(\ell^4\) norm of the field
and counting-\(\ell^2\) norm of its curl on \(B\) are both
at most \(C\epsilon\), although their global counterparts
need not be small.

\subsection{The nonlinear Wilson action bound on this domain}

Choose an orthonormal Lie-algebra identification in which
\(\frac12\operatorname{Tr}\exp A=\cos|A|\) for \(A\in\mathfrak{su}(2)\).
For each positively oriented plaquette \(p\), let
\(U_p(A)\) be the literal ordered product of its four
exponentiated links, with the two inverse orientations.
Let \(S_B(A)\) sum \(1-\frac12\operatorname{Tr}U_p(A)\)
over all six plaquettes anchored at each vertex of \(B\).
The linear oriented sum of their link logarithms is \((dA)_p\).

\begin{theorem}[Scale-uniform error for the actual plaquette action]
\label{thm:v56-Wilson}
There are constants \(C,\epsilon_0>0\), independent of \(M,L\),
such that every real link field satisfying
\[
       \|A\|_\infty\le K\epsilon/L,\qquad
       \|\nabla A\|_\infty\le K\epsilon/L^2
\]
on the plaquettes anchored in an \(L\)-cube, with
\(0<\epsilon\le\epsilon_0\), obeys
\begin{equation}\label{eq:v56-Wilson-error}
 \left|S_B(A)-\frac12\sum_{p\text{ anchored in }B}|(dA)_p|^2\right|
                                       \le C\epsilon^3.
\end{equation}
Here \(K\) is any fixed constant; \(C,\epsilon_0\) may depend
on \(K\), but not on the scale. In particular on the ancestor-good
event above, with the schedule \eqref{eq:v56-schedule},
\begin{equation}\label{eq:v56-beta-error}
 \beta\left|S_B(A_{\ge k})
          -\tfrac12\|dA_{\ge k}\|_{2,B}^2\right|
       \le C\frac{R^3}{\sqrt\beta}
       =C\frac{(\kappa\log\beta)^{3/2}}{\sqrt\beta}
                                                   \longrightarrow0.
\end{equation}
\end{theorem}
\begin{proof}
For four small oriented logarithms \(A_1,\ldots,A_4\),
analyticity of the finite product and the local logarithm gives
\[
 \log(e^{A_1}e^{A_2}e^{A_3}e^{A_4})
                 =\sum_i A_i+{\cal R},\qquad
 |{\cal R}|\le C\left(\sum_i|A_i|\right)^2.
\]
One may obtain this bound directly by subtracting the value and
first derivative of this analytic map at the origin and
integrating its bounded second derivative on the line segment.
It is uniform for the fixed four-factor product. The derivative
assumption implies \(|(dA)_p|\le C\epsilon/L^2\);
the link bound implies \(|{\cal R}_p|\le C\epsilon^2/L^2\).
Both facts are used: bounding the linear sum merely by four
link sizes would lose two powers of the scale in the action.

If \(F_p=\log U_p=(dA)_p+{\cal R}_p\), the exact SU(2) trace
and the Taylor remainder of cosine imply
\[
 \left|1-\tfrac12\operatorname{Tr}U_p-\tfrac12|(dA)_p|^2\right|
 \le C\left(|(dA)_p||{\cal R}_p|+|{\cal R}_p|^2+|F_p|^4\right)
 \le C\epsilon^3 L^{-4}.
\]
The last step uses \(L\ge1\), \(\epsilon\le\epsilon_0\le1\).
There are \(6L^4\) anchored plaquettes. Summing proves
\eqref{eq:v56-Wilson-error}; then insert
\(\epsilon=R/\sqrt\beta\) to obtain \eqref{eq:v56-beta-error}.
This proof bounds the literal compact Wilson plaquette, not an
action in which its cubic or quartic terms were discarded.
\end{proof}

\begin{proposition}[Four local marks keep the scale-uniform bound]
\label{prop:v56-marks}
Let \({\cal N}_B(A)=S_B(A)-\frac12\|dA\|_{\mathrm{bil},B}^2\)
denote the holomorphic extension, with a bilinear square.
On the finite set of link variables needed by \(B\) and one
additional layer of forward differences use the norm
\[
        \|A\|_{X_L}=L\|A\|_\infty+L^2\|\nabla A\|_\infty.
\]
The suprema include all those variables and differences; enlarging
that set by a fixed number of lattice steps changes only constants.
On a common smaller ball, for arbitrary marks \(h_1,\ldots,h_r\),
\[
 |D^r{\cal N}_B(A)[h_1,\ldots,h_r]|
 \le
 \begin{cases}
 C_r\|A\|_{X_L}^{\,3-r}\prod_i\|h_i\|_{X_L},&0\le r\le3,\\
 C_4\prod_i\|h_i\|_{X_L},&r=4 .
 \end{cases}
\]
At \(A=0\) the derivatives of orders below three vanish.
For \(\|X\|_{X_L}\le R\), \(R/\sqrt\beta\) small, the marked
function \(\beta{\cal N}_B(\beta^{-1/2}X)\) therefore has bounds
\(C_r\beta^{-1/2}R^{3-r}\) through order three, and
\(C_4\beta^{-1}\) at order four, for unit \(X_L\)-norm marks.
\end{proposition}
\begin{proof}
The local product logarithm and the trace are holomorphic on a
common complex ball. The proof of
Theorem~\ref{thm:v56-Wilson} remains a bound for their absolute
complex values when Hermitian norms are used only to estimate
the bilinear expressions. Thus
\(|{\cal N}_B(A)|\le C\|A\|_{X_L}^3\), uniformly in \(L,M\),
on that ball. Multivariable Cauchy estimates on a ball with
radius comparable to \(\|A\|_{X_L}\) give the first three
marked bounds; continuity and the vanished lower Taylor
coefficients cover \(A=0\). For the fourth derivative use
a fixed-radius ball inside the analytic domain instead.
Finally the chain rule contributes \(\beta^{1-r/2}\)
and proves the stated bounds in the normalized variable.
\end{proof}

For a fixed finest block at \(\beta(M)\ge c\log M\), choose
\(\eta\kappa>1\). With probability tending to one under the
free Gaussian law, all its ancestors satisfy
\eqref{eq:v56-beta-error} at once. Thus a local multiscale
domain and a vanishing local nonlinear error coexist on a
logarithmic coupling schedule. This does not contradict
V23's global small-ball obstruction: the field is allowed
to occupy arbitrarily many locally controlled blocks.

\subsection{What still fails to transfer, and the next estimate}

There are three precise limits to this progress.
First, \(A_{\ge k}\) is a heat-smoothed field on the fine
lattice. Theorem~\ref{thm:v56-Wilson} does not assert that its
action is the action obtained by integrating out the finer
fields. The mixed vertices in that integration are still present.
Second, the probabilities in
\eqref{eq:v56-bad-product} are free probabilities.
The exact Haar, quotient, source-normal and interacting density
must change them when the full measured law is used.
Third, all these regulators concern the mean-zero normal field;
they do not normalize the retained toron sector of V54--V55.

On a finite local Coulomb chart where the quotient density
\(\rho_M(A)\) is positive and the chart is one-to-one, the
reference change would contain the full weight
\begin{equation}\label{eq:v56-needed-weight}
 W(A)=\mathbf1_{\cal D}(A)\rho_M(A)
       \exp\!\left[-\beta\left\{S(A)-\tfrac12\|dA\|_2^2\right\}\right].
\end{equation}
Reweighting the independent heat fields by \(W(\sum_j a_j)\)
is an exact identity for that \emph{restricted} chart integral.
It is not a global chart construction or an estimate of its
normalizer. Local goodness of the heat components has not
proved positivity of a global Faddeev--Popov chart on their sum.

The localized density/pressure criterion needed to pass such
a reference bound to the full interacting law is already
identified in f7 V63 and is not reintroduced as a proved
hypothesis. The useful next target is to verify it for actual
mixed Wilson vertices and local quotient coordinates, retaining
the compact exterior, instead of paying a norm for
\eqref{eq:v56-needed-weight} over the entire volume.
Summing \eqref{eq:v56-beta-error} over all blocks gives an
extensive bound, so that substitution alone would not prove
the localized criterion. Neither the positive Gaussian
decomposition nor the bad-cluster bound supplies reflection
positivity for a new interacting coarse law or a calibrated
Hamiltonian gap.

\subsection{Verification and preservation}

The companion certificate
\begin{center}\small
\path{scripts/verify_ym_f14_v56_scale_local_regulator.py}
\end{center}
checks the nonzero Fourier curl/projection identities with exact
rational complex arithmetic at sides two and four, the symbolic
heat telescope including both ends, the polynomial operator-norm
majorants, the bounded-overlap and node-count arithmetic, and
the second and third coefficients of nonconstant four-link
SU(2) plaquettes by independent formal quaternion multiplication.
Those finite checks do not replace the all-volume estimates
or the analytic plaquette remainder proof.
The certificate passed: 270 nonzero Fourier modes, 288 symbolic
heat telescopes, 27 one-dimensional overlap checks (whose Cartesian
products give the four-dimensional check), and 81 nonconstant
plaquette jets. Eighty of the latter have a nonzero cubic term.
The two explicitly constructed level graphs have 529 and 8,465
nodes and maximum degrees 21 and 25, respectively.
The 9,222-byte certificate has SHA--256
\begin{center}\scriptsize\ttfamily
E82521CBFDCEFFFF5BD0A0194D1B50D7545ADBF4F5213B94C5441B27EA1A6941.
\end{center}

The predecessor V55 has 1,358,401 bytes and SHA--256
\begin{center}\scriptsize\ttfamily
E8FFC91EC007EF47C3C03FE9B0EEA01B3EFF37F9DC8B54F9CC01E22EFA3CD035.
\end{center}
Its body is preserved except for the title version and this
terminal insertion. V56 replaces V55 as the single active main
note. The protected companion files, supplied notes, archives
and monographs are unchanged.

\begin{firewall}
The new result is a verified local derivative regulator for
the actual transverse Gaussian reference, with volume-uniform
joint bad-set bounds and a scale-uniform nonlinear Wilson
action error on ancestor-good blocks. It is not an interacting
bad-set estimate, a construction of continuum Yang--Mills,
or a physical mass-gap proof. The full objective remains
unproved and active.
\end{firewall}
\section{V57: a quantitative finite-amplitude limit for the complete renormalized mean-fixed Gaussian coefficient}
\label{sec:v57}

\subsection{What is inherited and what is actually new}

This section was added on 15 September 2026. V56 supplies local
Gaussian domains and Wilson action bounds but does not supply
the normalized interacting probability estimate. The fixed-block
convex extensions in V38--V42 and the archived compact conditional
estimates already cover several simpler versions of that question.
We do not reintroduce one of those versions as its solution.

An upstream audit also finds that V21 already computes the
noncommuting mean-fixed quartic logarithm, its angular integral,
and its finite part. In particular its coefficient is
\(-2E_4/(3\pi^2)\), not an unevaluated instance of V30's
source-transport calculation. The coefficient, the exact
24-point angular rule, and the finite-part formula are reused
below; they are not new results of V57.

The new result sums \emph{all} higher source degrees on a common
twelve-dimensional complex neighborhood. It proves the
quantitative limit of the full renormalized mean-fixed amplitude,
with error \(O(M^{-2}(1+\log M))\) and every fixed source derivative.
In particular it extends V55 beyond the commuting locus and
improves that section's conservative convergence rate. This is
still a Gaussian integration coefficient of the actual fixed-mean
law, not its interacting normal integral at jointly varying
\(\beta,M\).

The notation \(E_4(c)=2\sum_{\mu<\nu}|c_\mu\times c_\nu|^2\)
is the existing homogeneous classical action. For complex \(c\)
all its squares and dot products are extended bilinearly;
Hermitian norms will be used only for estimates. Let
\(|c|^2=\sum_\mu|c_\mu|^2\) on real tuples, with the same
polynomial extension where it occurs in analytic formulas.
The V55 companion checkpoint remains current; the next is V60.

\subsection{The exact mean-fixed object, not the root-constrained one}

Use the mean-fixed coefficient \(\mathcal M_M(c)\) of
\eqref{eq:v21-mean}, normalized to vanish at zero, on a torus
of side \(M\ge4\). Its fine homogeneous background is \(a=c/M\).
All twelve constant link coordinates remain sources. The
mean-zero transverse Gaussian has dimension \(9(M^4-1)\).
V54 identifies this coefficient with its actual local
normal integration amplitude:
\begin{equation}\label{eq:v57-amplitude-identity}
 \exp[-\mathcal M_M(c)]
       =\frac{\mathcal A_M^{\log}(c/M)}{\mathcal A_M^{\log}(0)}.
\end{equation}
Here the amplitude uses twelve-dimensional logarithmic source
volume. There is no nonlinear root normal determinant, relative
root return, or coarse-Haar subtraction in this formula.

For clarity the finite symbol underlying V21 can be used at
nonzero source, not only as a fourth-order formal expression.
Let \(k\) be a nonzero point of the Brillouin torus, put
\[
 q_\mu=e^{ik_\mu}-1,\quad
 \lambda=|q|^2,\quad
 \Pi_k=I-\frac{q q^*}{\lambda},\quad P_k=\Pi_k\otimes I_3,
\]
and let \(H(k;a)\) be the literal Wilson Hessian used in V20--V21.
It is obtained by differentiating twice the four exponentials
of each plaquette, then applying their Fourier phases. Thus
it includes the curvature part of the Hessian at a noncommuting
constant field. Define
\begin{equation}\label{eq:v57-mode}
 \begin{aligned}
 X(k,a)&=\lambda^{-1}P_kH(k;a)P_k-P_k,\\
 Y(k,a)&=\lambda^{-1}\left\{
        \sum_\mu|q_\mu|^2\sigma(\operatorname{ad}_{a_\mu})
            -i\sum_\mu\sin k_\mu\operatorname{ad}_{a_\mu}\right\}-I_3,\\
 g(k,a)&=\tfrac12\Tr\log(I+X(k,a))-\Tr\log(I+Y(k,a)).
 \end{aligned}
\end{equation}
The first trace may be taken on all twelve orientation-colour
coordinates: \(X\) is zero on the free longitudinal subspace.
Each logarithm is continued from zero source. The exact expression is
\begin{equation}\label{eq:v57-exact}
 \mathcal M_M(c)=
       \sum_{n\in\mathcal K_M^4\setminus\{0\}}g(2\pi n/M,c/M)
        -2M^4\sum_\mu
                  \log\operatorname{sinc}(|c_\mu|/M).
\end{equation}
The last functions mean their analytic power series in
\(c_\mu\cdot c_\mu\). This is the complete fine-Haar and
logarithmic-gauge convention of V21, not a replacement measure.

\subsection{The two parity identities that improve the cutoff estimate}

\begin{lemma}[Reality, reflection, and simultaneous scaling]
\label{lem:v57-parity}
For real \(k\ne0\) and sufficiently small real \(a\),
\begin{equation}\label{eq:v57-parity}
             g(-k,a)=g(k,a),\qquad g(-k,-a)=g(k,a).
\end{equation}
Consequently \(g(k,a)\) is even in total source degree.
The identities also hold on the corresponding analytic extensions.
\end{lemma}
\begin{proof}
The real-space Wilson Hessian has real coefficients and is
symmetric, so \(H(-k;a)=\overline{H(k;a)}\). The same relation
holds for the free projections and the real logarithmic gauge
operator. On a positive real branch their determinants are
positive real numbers. Complex conjugation therefore leaves
their real logarithms unchanged, proving the first identity.

For the second identity reflect every lattice coordinate.
A positive link logarithm transforms as
\[
          a_\mu(x)\longmapsto-a_\mu(-x-\hat\mu).
\]
A mode \(e^{ikx}b_\mu\) consequently becomes a mode at \(-k\)
with amplitude \(R_\mu(k)b_\mu\), where
\(R_\mu(k)=-e^{-ik_\mu}\). The Wilson action and the mean
constraint are preserved. Moreover
\(R(k)q(k)=q(-k)\), so the free Coulomb subspace transforms
into the correct free Coulomb subspace at \(-k\).
It follows that
\[
       H(-k;-a)=R(k)H(k;a)R(k)^*,\qquad
       P_{-k}=R(k)P_kR(k)^* .
\]
The scalar gauge symbol in \eqref{eq:v57-mode} is unchanged
under \((k,a)\mapsto(-k,-a)\), since \(\sigma\) is even.
Determinants give the second identity.
Combining the two gives source evenness. Identity of analytic
germs extends each equality to complex sources and to the
one-variable momentum continuations used below.
\end{proof}

\begin{lemma}[A uniformly analytic simultaneous momentum-source germ]
\label{lem:v57-joint}
There are \(r_0,v_0>0\) such that, uniformly for \(n\in S^3\),
\[
                    G_n(r,v)=g(rn,rv)
\]
extends holomorphically to \(|r|<r_0,\ |v|<v_0\), including
\(r=0\). It is even separately in \(r\) and in total \(v\)-degree.
Writing \(T_{\le4}^v\) for its source Taylor polynomial through
degree four, the remainder
\(U_n=(I-T_{\le4}^v)G_n\) satisfies on smaller fixed domains
\begin{equation}\label{eq:v57-joint-remainder}
 |U_n(r,v)|\le C|v|^6,\qquad
 |U_n(r,v)-U_n(0,v)|\le C|r|^2|v|^6.
\end{equation}
\end{lemma}
\begin{proof}
Here all negative-frequency factors such as \(e^{-ik_\mu}\)
are continued analytically; a complex conjugate of the variable
\(r\) is never inserted. We have
\[
 q(rn)=irn+O(r^2),\qquad \lambda(rn)=r^2+O(r^4).
\]
Thus the free projection has a removable analytic extension
at \(r=0\), bounded uniformly in the real unit vector \(n\).

The finite plaquette derivatives are analytic in \(k,a\).
The source-free Hessian is \(\lambda P_k\). Its source-linear
coefficient vanishes at zero momentum and is \(O(|k||a|)\).
Every source coefficient of degree at least two is bounded by
its convergent local analytic series. These facts are also
the direct literal-Hessian inputs of Lemma~\ref{lem:v21-symbol}.
After \((k,a)=(rn,rv)\), the full Hessian is divisible by \(r^2\);
the free-normalized perturbation is bounded by
\(C(|v|+|v|^2)\) on a common complex domain.
The displayed gauge symbol has the same removable division,
and its perturbation obeys the same bound.
Decrease \(v_0,r_0\) so these norms are below \(1/2\).
The convergent matrix logarithms prove the uniform analytic
extension, including the point \(r=0\).

The simultaneous reflection in Lemma~\ref{lem:v57-parity}
gives \(G_n(-r,v)=G_n(r,v)\). Its source evenness gives
\(G_n(r,-v)=G_n(r,v)\). Uniform multivariable Cauchy bounds
therefore start \(U_n\) at degree six in \(v\), and start
\(U_n(r,v)-U_n(0,v)\) at degree two in \(r\).
Summing the two convergent power series on smaller domains
proves \eqref{eq:v57-joint-remainder}.
\end{proof}

These estimates also give a common ball \(|c|<c_0\) on which
\eqref{eq:v57-exact} is analytic for every \(M\ge4\).
Indeed \(|a|/|k|=|c|/(2\pi|n|)\) for centered nonzero modes.
For small momenta use the preceding lemma; away from zero
use compactness of the Brillouin region and ordinary matrix
continuity. On real sources the free-relative Hessian and
gauge operator remain positive on this ball. Their unscaled
lowest eigenvalues are not asserted to be bounded below
independently of \(M\).

\subsection{All source degrees beyond four have a common limit and rate}

For a continuum momentum \(p=2\pi n\ne0\), the limit
\[
                g_{\rm cont}(p,c)
                   =G_{p/|p|}(0,c/|p|)
\]
is a specified finite-dimensional determinant. Explicitly, let
\[
 T_\mu=\operatorname{ad}_{c_\mu},\quad T_p=\sum_\mu p_\mu T_\mu,
 \quad
 B_{\mu\nu}=-\delta_{\mu\nu}\sum_\rho T_\rho^2
                         +2T_\nu T_\mu-T_\mu T_\nu .
\]
With \(P_p=(I-pp^T/|p|^2)\otimes I_3\), put
\begin{equation}\label{eq:v57-cont-symbol}
 g_{\rm cont}(p,c)=
 \frac12\Tr\log\!\left[
 I+|p|^{-2}P_p\{-2iI_4\otimes T_p+B\}P_p\right]
       -\Tr\log\!\left[I_3-i|p|^{-2}T_p\right].
\end{equation}
This is the ordinary, mean-fixed Coulomb symbol of V21.
In particular its ghost is not the background-covariant
Laplacian. The distinction between conventional and
background-field gauge fixing is discussed in
\href{https://columbia.edu/cu/neurotheory/Larry/AbbottNP81a.pdf}
{L. F. Abbott, \emph{The Background Field Method Beyond One Loop},
Sections 2--3}; no coefficient is imported from that comparison.

\begin{theorem}[A summable finite-amplitude remainder, not only finite jets]
\label{thm:v57-high}
On a common smaller complex source ball the series
\begin{equation}\label{eq:v57-high-limit}
 {\cal R}_\infty(c)=
   \sum_{n\in\mathbb Z^4\setminus\{0\}}
       (I-T_{\le4}^c)g_{\rm cont}(2\pi n,c)
\end{equation}
converges normally and defines a holomorphic function.
If
\[
 {\cal R}_M(c)=
      \mathcal M_M(c)-\mathcal M_{M,[2]}(c)-\mathcal M_{M,[4]}(c),
\]
then
\begin{equation}\label{eq:v57-high-rate}
       |{\cal R}_M(c)-{\cal R}_\infty(c)|
            \le C|c|^6\,\frac{1+\log M}{M^2}.
\end{equation}
All fixed derivatives converge on smaller balls with that
cutoff rate, with derivative-dependent constants.
\end{theorem}
\begin{proof}
The bound at \(r=0\) in \eqref{eq:v57-joint-remainder}
majorizes the continuum summand by \(C|c|^6|n|^{-6}\).
This is summable in four dimensions. The same majorant
holds for each finite-grid remainder: near \(k=0\) use
Lemma~\ref{lem:v57-joint}; away from zero ordinary uniform
source Cauchy estimates give \(C|c/M|^6\le C'|c|^6|n|^{-6}\).
Uniform convergence on complex balls gives holomorphy.

For \(|p|<r_0M/2\), set \(r=|p|/M\), \(v=c/|p|\).
Taylor subtraction in \(c\) is precisely subtraction in \(v\).
The second inequality of \eqref{eq:v57-joint-remainder} gives
\[
 \left|(I-T_{\le4}^c)
       \{g(p/M,c/M)-g_{\rm cont}(p,c)\}\right|
                  \le \frac{C|c|^6}{M^2|p|^4}.
\]
The sum over these nonzero momenta is bounded by
\(C|c|^6M^{-2}(1+\log M)\).
For the remaining finite-grid momenta, and for the continuum
momenta outside that region or outside the grid, use the
\(|n|^{-6}\) majorant; its tail is \(O(M^{-2})\).

The remaining fine-Haar term after degree four is at most
\[
       C M^4\sum_\mu |c_\mu/M|^6
                         \le C|c|^6M^{-2}.
\]
Its series is even and analytic uniformly on the selected ball.
This proves \eqref{eq:v57-high-rate}. Finally apply the
multivariable Cauchy formula on a slightly larger common
source ball to obtain each fixed derivative estimate.
\end{proof}

The conclusion would not follow from merely knowing convergence
of an operator logarithm with its first four \emph{Born powers}
removed. Those powers contain arbitrarily high source degrees
when the operator is nonlinear in the source. The simultaneous
scaling argument above controls the actual source remainder.

\subsection{Quantitative versions of the inherited low-degree limits}

V21 already gives the exact quadratic coefficient
\(\mathcal M_{M,[2]}(c)=a_M|c|^2\) and the quartic finite part
\[
 \mathcal M_{M,[4]}(c)
       +\frac{2}{3\pi^2}E_4(c)\log M
       \longrightarrow
 {\cal P}_{4}^{\rm mean}(c)
       :=\Lambda_\infty^{\rm ren}(c)
                       +\frac1{90}\sum_\mu(c_\mu\cdot c_\mu)^2.
\]
The last Haar polynomial stays here: unlike the root-source
formula, this mean-coordinate expression has no coarse-Haar
factor to cancel it.

\begin{lemma}[Low-degree rate with the actual lattice finite part]
\label{lem:v57-low}
In any fixed norm on the indicated polynomial spaces,
\begin{equation}\label{eq:v57-low-rate}
 \begin{aligned}
 |a_M-a_\infty|&\le CM^{-2},\\
 \left\|\mathcal M_{M,[4]}
       +\frac{2}{3\pi^2}E_4\log M-{\cal P}_{4}^{\rm mean}\right\|
                                  &\le C M^{-2}(1+\log M).
 \end{aligned}
\end{equation}
The limiting polynomial is exactly V21's finite-part formula,
not a newly chosen subtraction scheme.
\end{lemma}
\begin{proof}
For the quadratic term use \eqref{eq:v21-quadratic}.
On \(|m|\le cM\), with a sufficiently small fixed \(c>0\),
its lattice dispersion differs from \(2\pi|m|\) by at most
\(C|m|^3/M^2\), by Taylor expansion of sine and inverse
hyperbolic sine. Both dispersions are bounded below by
\(c'|m|\). The derivative of \((\cosh x-1)^{-1}\) on
these intervals is bounded by \(Ce^{-c''|m|}\).
Thus the summed difference is \(O(M^{-2})\). The remaining
dispersion and continuum tails are exponentially small in \(M\).
The explicit \(M^{-2}\) term in \(a_M\) has the same rate.

For the quartic term let \(f(k;c)\) be exactly
\eqref{eq:v21-f}, and \(h(k;c)=|k|^{-4}F(k/|k|;c)\) its
already computed singular part. Lemma~\ref{lem:v57-parity}
makes \(f\) even in \(k\). Its entries are finite Taylor
coefficients of analytic plaquette stencils with powers of
\(\lambda(k)\) in the denominator. Their polar Taylor
expansions are uniform in direction. The order \(-3\)
term in Lemma~\ref{lem:v21-symbol} is therefore zero by
evenness. For \(j=0,1,2\),
\begin{equation}\label{eq:v57-integrable-bounds}
            |D_k^j(f-h)|\le C_j |k|^{-2-j}.
\end{equation}
To justify the derivative statement, expand the finite analytic
numerators and \(\lambda(k)/|k|^2=1+O(|k|^2)\) to the
required two additional orders on the compact unit sphere.
Odd homogeneous terms vanish; differentiating the analytic
remainders gives the stated bounds.

Use the same radial cutoff \(\eta\) and remainder
\(R=f-\eta h\) as in Lemma~\ref{lem:v21-sum}.
On the central grid cell, the omitted integral of \(R\)
is \(O(M^{-2})\) by \eqref{eq:v57-integrable-bounds}.
For the other cells, midpoint Taylor expansion has no
first-order integral. Its error on a cell of side
\(s=2\pi/M\) is at most
\(Cs^6\sup|D^2R|\). Summing \eqref{eq:v57-integrable-bounds}
outside the finitely many cells nearest zero gives
\[
       Cs^6\sum_{0<|n|\le CM}|sn|^{-4}
                                    \le CM^{-2}(1+\log M).
\]
The finitely many nearest cells contribute \(O(M^{-2})\)
directly. Periodic cells at the Brillouin boundary use the
periodic continuation of \(R\), so no boundary quadrature
term is introduced.

The singular-part comparison in Lemma~\ref{lem:v21-sum}
also improves by one power: the centered unit-cell moment
of its first derivative is zero, and \(D^2h=O(|x|^{-6})\).
Consequently
\[
  \left|h(n)-\int_{n+[-1/2,1/2]^4}h(x)\,dx\right|
                                      \le C|n|^{-6}.
\]
The tail of that cell-difference series is \(O(M^{-2})\).
For the cutoff-weight error, supported where \(|n|\asymp M\),
the constant-times-linear term again integrates to zero.
The remaining terms contain either \(D^2\eta_M\,h\) or
\(D\eta_M\,Dh\), where \(\eta_M(x)=\eta(2\pi|x|/M)\).
Their summed bound is \(O(M^{-2})\).
The corresponding radial integral gives exactly the same
\(\log M\) and finite constant as in V21, for all sufficiently
large \(M\). This proves the second rate in
\eqref{eq:v57-low-rate}; finitely many smaller \(M\) are
absorbed in the constant.
\end{proof}

\subsection{The full renormalized amplitude and its source normalization}

\begin{theorem}[The twelve-source finite-amplitude limit]
\label{thm:v57-full}
Define
\begin{equation}\label{eq:v57-renormalized}
 \begin{aligned}
 \widehat{\mathcal M}_M(c)&=
       \mathcal M_M(c)+\frac{2}{3\pi^2}E_4(c)\log M,\\
 \mathcal M_\infty^{\rm ren}(c)&=
       a_\infty|c|^2+{\cal P}_{4}^{\rm mean}(c)+{\cal R}_\infty(c).
 \end{aligned}
\end{equation}
These are holomorphic on one common small complex ball.
For every multi-index \(\alpha\), on a smaller common ball,
\begin{equation}\label{eq:v57-full-rate}
 \sup_c\left|\partial_c^\alpha
       \{\widehat{\mathcal M}_M(c)-\mathcal M_\infty^{\rm ren}(c)\}
         \right|\le C_\alpha\frac{1+\log M}{M^2}.
\end{equation}
Equivalently the completely normalized Gaussian amplitude obeys
\begin{equation}\label{eq:v57-ren-amplitude}
 M^{-2E_4(c)/(3\pi^2)}
       \frac{\mathcal A_M^{\log}(c/M)}{\mathcal A_M^{\log}(0)}
                 \longrightarrow e^{-\mathcal M_\infty^{\rm ren}(c)}
\end{equation}
with that rate and all fixed derivatives.
For real sources in a smaller ball the renormalized ratios
are bounded above and below by positive constants independent
of \(M\). The same limit and rate hold for V54's quaternion
source-volume amplitude.
\end{theorem}
\begin{proof}
Combine Lemma~\ref{lem:v57-low} with
Theorem~\ref{thm:v57-high} in the exact source decomposition.
This proves \eqref{eq:v57-full-rate} without derivatives.
Cauchy estimates on a larger
common ball give every fixed derivative. Bounded holomorphic
exponentiation and \eqref{eq:v57-amplitude-identity} give
\eqref{eq:v57-ren-amplitude}. Positivity on the real branch
and a common bound for its logarithm give the two-sided bounds.

For quaternion source coordinates
\(v_\mu(a)=\operatorname{sinc}(|a_\mu|)a_\mu\), the exact
source Jacobian from V54--V55 is
\[
       J_{\rm src}(a)=
          \prod_\mu\cos|a_\mu|\,\operatorname{sinc}^2|a_\mu|.
\]
The quaternion negative-log amplitude differs from
\(\mathcal M_M(c)\) by \(\log J_{\rm src}(c/M)\).
This analytic function and each fixed source derivative are
\(O(M^{-2})\) on the common ball. Its value at zero is zero.
It does not change the stated limit or rate.
\end{proof}

On the commuting family \(c_\mu=z_\mu E_3\) the counterterm
vanishes. Uniqueness of the limit and V55 then give
\[
           \mathcal M_\infty^{\rm ren}(z_\mu E_3)=2F_\infty(z).
\]
Thus the four-angle amplitude of V55 is recovered with the
stronger \(O(M^{-2}(1+\log M))\) rate. For a fixed small real
noncommuting tuple, in contrast, the \emph{unrenormalized}
Gaussian ratio has the growth
\[
 \frac{\mathcal A_M^{\log}(c/M)}{\mathcal A_M^{\log}(0)}
   =M^{2E_4(c)/(3\pi^2)}
       e^{-\mathcal M_\infty^{\rm ren}(c)}
            \{1+O_c(M^{-2}(1+\log M))\}.
\]
Uniformity on the flat locus alone could not have detected
this growth away from it.

\subsection{What this says about a one-loop source weight}

For precision, define the \emph{one-loop source expression}
\[
 {\cal W}^{(1)}_{\beta,M}(c)=
  \exp[-\beta S_0(c/M,0)]\,
             \frac{\mathcal A_M^{\log}(c/M)}{\mathcal A_M^{\log}(0)}.
\]
This is a defined expression from the classical action and
the Gaussian normal coefficient, not a notation for the full
normal integral. The exact constant-link action in V54 gives
\[
             S_0(c/M,0)=E_4(c)+O(M^{-2}|c|^6).
\]
Combining that identity with Theorem~\ref{thm:v57-full} yields
\begin{equation}\label{eq:v57-one-loop-weight}
 \log{\cal W}^{(1)}_{\beta,M}(c)
   =-\left(\beta-\frac{2}{3\pi^2}\log M\right)E_4(c)
       -\mathcal M_\infty^{\rm ren}(c)
        +O\!\left(\frac{\beta+1+\log M}{M^2}\right)
\end{equation}
on smaller source balls, with fixed derivative versions.
In particular this error tends to zero if \(\beta=o(M^2)\).
The identity is nonperturbative in the small source \(c\);
it is still only a one-loop statement in the normal fluctuations.

The subtraction in \eqref{eq:v57-one-loop-weight} is the
mean-coordinate coefficient already identified in V21.
It is not, by itself, a physical running coupling.
V30 includes additional measured source transport, and
V19 records how changing that source also changes the
next loop order and the observable. No source replacement
has been made here. Likewise the logarithmic toron
normalization of V54 cannot be replaced by a normalized
unbounded quartic matrix integral: the commuting directions
still have \(E_4=0\).

The next missing estimate is now sharper. It must control
the actual fully measured interacting normal integral relative
to the renormalized amplitude on a jointly varying cutoff
and coupling sequence. V56's local domains may be useful for
that estimate, but the present determinant theorem does not
transfer their free bad-set probabilities to the interacting
law, cover the compact exterior, or construct its continuum
time evolution.

\subsection{Executed checks and preservation}

The certificate
\begin{center}\small
\path{scripts/verify_ym_f14_v57_analytic_amplitude.py}
\end{center}
reuses the existing V20 literal plaquette differentiator.
It checks the reflected Fourier-offset identity and real
Hessian symmetry for 364 stencil entries containing 874
source coefficients; 70 midpoint monomials; and independent
sixth/eighth source and second mesh-order jets on flat-mode
regressions. These checks support the new symmetry and
rate arguments. They do not repeat V21's angular calculation
as a new result, or prove a continuum statement by sampling.
The 5,817-byte certificate has SHA--256
\begin{center}\scriptsize\ttfamily
2339C3CBDA6BCF4AA7BD2AB1461D0423356C8EBBB66E329A4BED38B653A475B1.
\end{center}
The imported V20 checker is unchanged, with SHA--256
\begin{center}\scriptsize\ttfamily
D709195859864718EE92648D6883D8ED50CF14108561B98050D8965809DECB77.
\end{center}

The predecessor V56 has 1,384,894 bytes and SHA--256
\begin{center}\scriptsize\ttfamily
349F96354BBC696758E609BA285460BFCA59C7FC9A4F7C462625FCA526CFE061.
\end{center}
Its complete body is retained apart from the title version
and this terminal insertion. V57 replaces V56 as the single
active main note. Companion files, supplied manuscripts,
archived notes and monographs remain unchanged.

\begin{firewall}
The complete mean-fixed Gaussian amplitude now has an
identified, quantitative finite-amplitude limit after its
already known quartic subtraction. This extends beyond
the commuting locus and beyond finitely many Taylor jets.
It does not prove a joint nonperturbative normal-integral
estimate, a nontrivial continuum Yang--Mills theory, or a
positive physical Hamiltonian mass gap. The full objective
remains unproved and active.
\end{firewall}
\section{V58: the complete two-loop vacuum density and an extensive cubic fluctuation}
\label{sec:v58}

\subsection{Context, acquisition, and the order of limits}

This section was added on 15 September 2026. V57 controls the
mean-fixed Gaussian amplitude, including its finite source
dependence. The next question is what happens when the actual
normal integral is expanded beyond that amplitude. A bound
which treats every fluctuation as a globally small vector is
poorly adapted to the growing number of fine links. We therefore
first identify the connected vacuum correction per lattice site.

The new acquisition is a thermodynamic-limit theorem for the
complete first non-Gaussian coefficient at zero retained mean.
It includes the Wilson quartic vertex, the square of the Wilson
cubic vertex, the fine Haar density, and both quadratic
Faddeev--Popov contributions. The cubic covariance is absolutely
summable, with a strictly positive limiting variance density.
This is not a new derivation of V38's cubic plaquette polynomial
or V41's general measured Wick formula: both are reused below.
What is new is the all-volume summability and the identified
limit of their complete vacuum specialization.

There are two distinct limits. First, at each fixed lattice
side \(M\), take the coefficient in the low-temperature
expansion of a specified local integral. Then let \(M\) grow
in that coefficient. We do not interchange these limits in
the interacting integral. In particular, a bounded coefficient
per site does not imply a uniform remainder at
\(\beta\) of order \(\log M\).

The source-selection, native-stationarity and exact-action
manuscripts supplied by the user remain subject to V44's
scope audit. They help distinguish an exact measured law
from a comparison law, but do not provide the missing
Yang--Mills fluctuation estimate. No text inside them is
treated as an instruction. The last regular SM/NS/Shell
review was V55; the next is V60.

\subsection{The fixed-mean local integral and its coefficient}

Write \(V=M^4\), initially \(M\ge4\), and use unit lattice
spacing and counting inner products. Let
\[
 \mathcal H_M=\{a:D^*a=0,\ \sum_xa_\mu(x)=0
                         \text{ for each }\mu\},
 \qquad \dim_{\mathbb R}\mathcal H_M=9(V-1).
\]
There are three real colour components. The twelve constant
link components are fixed to zero, not integrated. Constant
colour rotations are retained; no additional global-colour
quotient or harmonic mass is inserted.

For momentum \(k\in(2\pi/M)\mathbb Z^4/(2\pi\mathbb Z)^4\), set
\begin{equation}\label{eq:v58-free}
 q_\mu(k)=e^{ik_\mu}-1,\quad
 \lambda(k)=\sum_\mu|q_\mu(k)|^2,\quad
 C(k)=\frac1{\lambda(k)}
       \left(I-\frac{q(k)q(k)^*}{\lambda(k)}\right),\quad k\ne0,
 \qquad C(0)=0.
\end{equation}
The free real Gaussian \(Z\) on \(\mathcal H_M\) has
covariance \(C_M\otimes I_3\), where
\[
 C_{M,\mu\nu}(x)=\frac1V\sum_k e^{ik\cdot x}C_{\mu\nu}(k).
\]
Thus its precision is the restriction of the free Wilson
Hessian, not a massive comparison operator.

Choose a colour basis with
\([u,v]=2u\times v\). A simultaneous opposite colour convention
changes the sign of the cubic vertex but none of the
coefficients below. Write the literal action and full
quotient logarithm as
\begin{equation}\label{eq:v58-jets}
\begin{aligned}
 S(tZ)&=t^2S_2(Z)+t^3S_3(Z)+t^4S_4(Z)+O(t^5),
       &S_2(Z)&=\tfrac12\|dZ\|_2^2,\\
 \ell(tZ)&=\log\rho_M(tZ)=t^2\ell_2(Z)+O(t^3).
\end{aligned}
\end{equation}
Here \(\rho_M\) is precisely the fine Coulomb quotient
density in \eqref{eq:v45-rho}, normalized by \(\rho_M(0)=1\).
Splitting its linear Coulomb coordinates into their mean
and mean-zero parts has only a field-independent Jacobian.
Consequently no root-return determinant, inverse-pivot
factor or coarse Haar density belongs in this particular
fixed-mean integral.

For each fixed \(M\), take a nonnegative smooth cutoff
\(\chi_M\), equal to one near zero and supported in a
sufficiently small chart where zero is the unique action
minimum. Define
\[
 I_M(\beta)=\int_{\mathcal H_M}
       \chi_M(a)e^{-\beta S(a)}\rho_M(a)\,da,\qquad
 I_M^{\rm G}(\beta)=\int_{\mathcal H_M}
                         e^{-\beta S_2(a)}\,da.
\]
All assertions about this localization are at fixed \(M\).
It is not a localization of uniform probability under
the full compact measure.

\begin{proposition}[Inherited measured formula at the vacuum]
\label{prop:v58-coefficient}
For every such fixed \(M\),
\begin{equation}\label{eq:v58-B}
 -\log\frac{I_M(\beta)}{I_M^{\rm G}(\beta)}
       =\frac{B_M}{\beta}+O_M(\beta^{-2}),\qquad
 B_M=\mathbb E S_4-\tfrac12\mathbb E S_3^2-\mathbb E\ell_2.
\end{equation}
The coefficient does not depend on the chosen cutoff.
\end{proposition}
\begin{proof}
After \(a=\beta^{-1/2}Z\), expand the normalized integrand
through order \(\beta^{-1}\). It is
\[
 1-\beta^{-1/2}S_3+
 \beta^{-1}\{\tfrac12S_3^2-S_4+\ell_2\}.
\]
There is no linear colour-invariant density term.
Odd homogeneous Gaussian moments vanish, in particular
\(\mathbb E S_3=0\). Taking the negative logarithm gives
\eqref{eq:v58-B}, the zero-source specialization of
\eqref{eq:v41-Wick}. For completeness, at fixed \(M\)
the action is bounded below by a positive quadratic form
on a smaller chart. Split the integral into that chart
and its complement in the cutoff support. The latter is
exponentially small in \(\beta\). Taylor expansion through
degree six on the former, with Gaussian domination and
vanishing odd terms, gives the stated remainder.
Cutoffs agreeing near zero have exponentially small
differences relative to the Gaussian normalization.
These constants and neighborhoods may depend on \(M\).
\end{proof}

\subsection{A uniform transverse Green estimate}

Let \(d_M(x,0)\) be the nearest periodic Euclidean distance.
An ordered lattice derivative means any product of forward
or backward unit differences; bounded fixed shifts of the
argument are harmless.

\begin{lemma}[Massless transverse kernels]
\label{lem:v58-Green}
For each fixed derivative order \(r\ge0\), there is \(A_r\)
independent of \(M\) such that
\begin{equation}\label{eq:v58-decay}
 |\nabla^r C_{M,\mu\nu}(x)|
       \le A_r(1+d_M(x,0))^{-2-r}.
\end{equation}
At every fixed offset these kernels converge to the
corresponding derivatives of
\[
 C_{\infty,\mu\nu}(x)=
 \int_{[-\pi,\pi]^4}e^{ik\cdot x}C_{\mu\nu}(k)
                                      \frac{d^4k}{(2\pi)^4}.
\]
\end{lemma}
\begin{proof}
Away from zero the symbol is smooth and periodic. Near zero
it obeys
\[
 |\partial^\alpha C(k)|\le A_\alpha|k|^{-2-|\alpha|}.
\]
A product of \(r\) difference multipliers improves this
to symbol order \(r-2\). Use a smooth dyadic partition
near zero. On an annulus of radius \(L^{-1}\), rescaling
and integration by parts give for its Fourier coefficient
\[
 |K_{L,r}(x)|\le A_{N,r}L^{-2-r}
                                   (1+|x|/L)^{-N}.
\]
The part away from zero satisfies the same bound with
\(L\) comparable to one.

For the finite torus first multiply the symbol by a smooth
cutoff \(\eta(Mk)\) which vanishes in a ball about zero and
is one at every nonzero grid point. Its transition can be
placed strictly below the shortest nonzero momentum.
Only scales \(1\le L\le cM\) then occur; on the terminal
annulus the cutoff has exactly the derivative bounds used
above. Sampling a smooth periodic symbol periodizes its
Fourier coefficients. For \(N>4\),
\[
 \sum_{z\in\mathbb Z^4}
 (1+|x+Mz|/L)^{-N}
       \le A_N(1+d_M(x,0)/L)^{-N},\qquad L\le cM,
\]
with constants allowed to depend on the fixed \(c\).
Take also \(N>2+r\). Summing the dyadic estimates, first
over \(L\le1+d_M(x,0)\) and then over the larger scales,
proves \eqref{eq:v58-decay}.

For fixed-offset convergence remove a ball of radius
\(\delta\) in momentum. The remaining Riemann sums converge.
The deleted integrals and normalized sums are bounded by
\(A_r\delta^{2+r}\), using annulus counting and the omitted
zero mode. Let \(\delta\) tend to zero. This argument
controls the full transverse matrix, including its
\(q q^*/\lambda^2\) term.
\end{proof}

\subsection{The actual cubic vertex has summable covariance}

For the plaquette \((x;\mu,\nu)\), \(\mu<\nu\), let its
four oriented logarithm vectors be
\[
 v_1=Z_\mu(x),\quad v_2=Z_\nu(x+e_\mu),\quad
 v_3=-Z_\mu(x+e_\nu),\quad v_4=-Z_\nu(x).
\]
V38's literal cubic polynomial can equivalently be written
\begin{equation}\label{eq:v58-cubic}
 s_{3,p}=\sum_{i<j<k}\det(v_i,v_j,v_k)
       =l_p\cdot B_{2,p},\qquad
 l_p=\sum_i v_i=(dZ)_p,\qquad
 B_{2,p}=\sum_{i<j}v_i\times v_j.
\end{equation}
Indeed, expansion of \(l_p\cdot B_{2,p}\) gives one copy
of each ordered triple determinant. This identity is
used to retain a discrete curl at each cubic vertex.
Set \(s_{3,x}=\sum_{\mu<\nu}s_{3,(x;\mu,\nu)}\);
then \(S_3=\sum_xs_{3,x}\).

\begin{theorem}[Extensive, nonzero cubic variance]
\label{thm:v58-cubic}
Uniformly in \(M\),
\begin{equation}\label{eq:v58-cubic-decay}
 |\mathbb E(s_{3,0}s_{3,x})|
       \le A(1+d_M(x,0))^{-8}.
\end{equation}
In particular,
\begin{equation}\label{eq:v58-sigma}
 \sigma_{3,M}^2:=V^{-1}\mathbb E S_3^2
   =\sum_{x\in\mathbb T_M^4}\mathbb E(s_{3,0}s_{3,x})
       \longrightarrow
 \sigma_3^2:=\sum_{x\in\mathbb Z^4}
                   \mathbb E_\infty(s_{3,0}s_{3,x}),
 \qquad 0<\sigma_3^2<\infty .
\end{equation}
The infinite sum is absolutely convergent. Its absolute
spatial moments of every order \(s<4\) are finite.
\end{theorem}
\begin{proof}
Every cubic term carries a colour epsilon tensor.
Contracting any two legs of one vertex gives a colour
Kronecker delta and hence zero. Thus each cubic is already
Wick ordered: only the three-cross-line pairings contribute
to its covariance with another cubic.

Use the \(l_p\cdot B_{2,p}\) expression. If the curl legs
are paired to one another, that covariance has two lattice
differences and decays as \((1+d)^{-4}\). The other two
covariances each decay as \((1+d)^{-2}\). If the two curl
legs belong to distinct cross-lines, those two lines
each decay as \((1+d)^{-3}\), and the remaining line as
\((1+d)^{-2}\). In both cases the product is bounded by
\(A(1+d)^{-8}\), by Lemma~\ref{lem:v58-Green}.
There are only finitely many plaquette and Wick terms.
This proves \eqref{eq:v58-cubic-decay}.

Translation invariance gives the finite equality in
\eqref{eq:v58-sigma}. Fixed-offset convergence follows
by Wick's rule and Lemma~\ref{lem:v58-Green}; the uniform
tail is \(O(R^{-4})\), since a four-dimensional annulus
has \(O(R^4)\) sites. This proves convergence and the
moment assertion. Strict positivity requires a further
argument, since the individual covariances need not
be nonnegative.

Use the unitary Fourier transform and order the three
colours once and for all. There is a finite trigonometric
polynomial tensor \(K\) such that
\begin{equation}\label{eq:v58-vertex}
 S_3=V^{-1/2}\!
 \sum_{p+q+r=0}\sum_{\mu,\nu,\rho}
 K_{\mu\nu\rho}(p,q,r)
     \widehat Z^1_\mu(p)\widehat Z^2_\nu(q)
                                      \widehat Z^3_\rho(r).
\end{equation}
Its explicit definition is as follows. For each of the
six oriented plaquette paths, select three of its four
positions and assign colours \(1,2,3\) in every order.
Multiply the permutation sign, the three link signs,
and the phases at the three link origins. Sum them in
the corresponding three orientation slots. This gives
144 primitive terms. There is no extra \(3!\) in the
variance, because the colours in \eqref{eq:v58-vertex}
have already been ordered.

If \(\mathcal C=C(p)\otimes C(q)\otimes C(r)\), put
\[
 [K,\mathcal C K]_{\rm cov}
       :=\sum_{I,J}K_I\overline{K_J}\mathcal C_{IJ}
       =\langle\overline K,\mathcal C\,\overline K\rangle
       \ge0 .
\]
This notation fixes the conjugations of link Fourier
phases. Wick contraction gives
\begin{equation}\label{eq:v58-sigma-integral}
 \sigma_3^2=
 \int\!\!\int
 [K(p,q,-p-q),
   \{C(p)\otimes C(q)\otimes C(-p-q)\}K(p,q,-p-q)]_{\rm cov}
                         \frac{d^4p\,d^4q}{(2\pi)^8}.
\end{equation}
To justify the formula, the finite counterpart is
\(V^{-2}\sum_{p,q}\) of the displayed integrand.
It is bounded by
\(A f(p)f(q)f(-p-q)\), where \(f=\lambda^{-1}\)
off zero and \(f(0)=0\) on a grid.
Both the integral and normalized grid \(L^{3/2}\) norms
of \(f\) are uniformly finite. Those norms restricted
to \(|k|<\delta\) are \(O(\delta^{2/3})\).
Young's inequality and then Hölder's inequality give
\[
 \int f(p)f(q)f(-p-q)\,dp\,dq
       \le \|f\|_{3/2}^3
\]
with normalized torus measure, and the same estimate
holds on every finite group. Applying this bound with
one restricted factor removes all three singular sets
uniformly. On their complement ordinary Riemann sums
converge. This proves \eqref{eq:v58-sigma-integral}.

For positivity choose
\[
 p=(\pi/2,0,0,0),\qquad q=(0,\pi/2,0,0),\qquad
 r=(-\pi/2,-\pi/2,0,0).
\]
Give colours \(1,2,3\) respectively the spatial
polarizations \(e_2,e_3,e_3\). Each is transverse
to its own \(q(k)\). Only the \((2,3)\) plaquette
can contain the three polarizations. With colour
amplitudes \(u,v,w\), its oriented vectors are
\[
 uE_1,\quad ivE_2-iwE_3,\quad -uE_1,\quad
                                      -vE_2-wE_3 .
\]
The two nonzero triple determinants in
\eqref{eq:v58-cubic} each contribute \(-2iuvw\).
Consequently \(K_{233}(p,q,r)=-4i\).
The restriction of the tensor to the three transverse
spaces is nonzero. All three nonzero-momentum
covariances are positive definite on those spaces.
Continuity gives a set of positive measure where the
integrand in \eqref{eq:v58-sigma-integral} is strictly
positive. This proves \(\sigma_3^2>0\).
\end{proof}

\subsection{The Wilson quartic term is local and has a limit}

Let \(c_{ij}=\mathbb E(v_i^a v_j^a)\) for a fixed colour
\(a\), with no colour sum, on one oriented plaquette.
This includes the signs in \(v_3,v_4\).
Define the polynomial
\begin{equation}\label{eq:v58-quartic-polynomial}
\begin{aligned}
 W(c)={}&\frac58\sum_i c_{ii}^2\\
 &+\sum_{i<j}\left\{
   \frac94c_{ii}c_{jj}+\frac32c_{ij}^2
       +\frac52(c_{ii}+c_{jj})c_{ij}
       +\sum_{k\ne i,j}
           \left(\frac92c_{kk}c_{ij}+3c_{ki}c_{kj}\right)
                         \right\}\\
 &+9c_{12}c_{34}-3c_{13}c_{24}+9c_{14}c_{23}.
\end{aligned}
\end{equation}

\begin{lemma}[Literal quartic expectation]
\label{lem:v58-quartic}
Writing \(s_{4,M}=\mathbb E S_{4,p}\) on any one
positively oriented plaquette, one has
\[
 s_{4,M}=-W(c_M),\qquad
 V^{-1}\mathbb E S_4=6s_{4,M},\qquad
 s_{4,M}\longrightarrow s_{4,\infty}=-W(c_\infty).
\]
These quantities are uniformly bounded.
\end{lemma}
\begin{proof}
In the colour convention of \eqref{eq:v58-cubic},
the four quaternion exponentials have scalar parts
\(1-t^2|v_i|^2/2+t^4|v_i|^4/24\) and vector parts
\(tv_i-t^3v_i|v_i|^2/6\), through the required orders.
The scalar fourth coefficient of their product has
four sorts of terms:
\[
 \sum_i\frac{|v_i|^4}{24},\quad
 \sum_{i<j}\frac{|v_i|^2|v_j|^2}{4},\quad
 \sum_{i<j}(v_i\cdot v_j)
   \left\{\frac{|v_i|^2+|v_j|^2}{6}
                         +\sum_{k\ne i,j}\frac{|v_k|^2}{2}\right\},
\]
and
\[
 (v_1\cdot v_2)(v_3\cdot v_4)
 -(v_1\cdot v_3)(v_2\cdot v_4)
 +(v_1\cdot v_4)(v_2\cdot v_3).
\]
For colour-isotropic Gaussians, Wick's rule gives
\[
\begin{aligned}
 \mathbb E|v_i|^4&=15c_{ii}^2,\\
 \mathbb E(|v_i|^2|v_j|^2)&=9c_{ii}c_{jj}+6c_{ij}^2,\\
 \mathbb E\{|v_k|^2(v_i\cdot v_j)\}
                                  &=9c_{kk}c_{ij}+6c_{ki}c_{kj},\\
 \mathbb E\{(v_i\cdot v_j)(v_k\cdot v_l)\}
                         &=9c_{ij}c_{kl}+3c_{ik}c_{jl}+3c_{il}c_{jk}.
\end{aligned}
\]
Substitution gives \(W\). The Wilson action is one minus
the scalar part, explaining its negative sign.
Hypercubic symmetry and translation invariance give
the factor six. The covariance entries involve only
fixed offsets, so their boundedness and convergence
follow from Lemma~\ref{lem:v58-Green}.
\end{proof}

\subsection{Both gauge-density contributions and the Haar term}

Retain the primed vertex space, and distinguish its
endpoint-sum operator \(S_{\rm end}\) from the Wilson action.
As in \eqref{eq:v46-gauge-split}, put
\[
 U=D\Delta^{-1/2},\qquad W=S_{\rm end}\Delta^{-1/2},\qquad
 G_1(Z)=\tfrac12U^*\operatorname{ad}_Z W,\qquad
 G_2(Z)=\tfrac1{12}U^*(\operatorname{ad}_Z)^2U.
\]
The expansion
\(\sigma(A)=I+A^2/12+O(A^4)\), together with all fine
Haar factors, gives exactly
\begin{equation}\label{eq:v58-density}
 \ell_2(Z)=-\tfrac13\sum_e|Z_e|^2
                 +\operatorname{Tr}G_2(Z)
                 -\tfrac12\operatorname{Tr}G_1(Z)^2.
\end{equation}
Here \(\operatorname{Tr}G_1=0\), but this does not erase
the square of \(G_1\).

\begin{lemma}[Complete quadratic density per site]
\label{lem:v58-density}
Let
\[
 g_M=\frac1V\sum_{k\ne0}\frac1{\lambda(k)},\qquad
 j_M=\frac1V\mathbb E\operatorname{Tr}G_1(Z)^2.
\]
Then
\begin{equation}\label{eq:v58-density-average}
 -V^{-1}\mathbb E\ell_2
       =\left(\frac92-\frac3{2V}\right)g_M+\frac12j_M,
 \qquad 0\le j_M\le24g_M^2.
\end{equation}
Moreover \(g_M\to g_\infty=\int\lambda(k)^{-1}d^4k/(2\pi)^4\)
and \(j_M\to j_\infty<\infty\). An explicit formula is
\begin{equation}\label{eq:v58-j-integral}
 j_\infty=24\int\!\!\int
 \frac{w(p,q)^T C(p-q)\overline{w(p,q)}}
      {\lambda(p)\lambda(q)}
                             \frac{d^4p\,d^4q}{(2\pi)^8},
 \qquad
 w_\mu(p,q)=\tfrac12\overline{q_\mu(p)}(1+e^{iq_\mu}).
\end{equation}
In the last exponential \(q_\mu\) is the \(\mu\)-component
of the momentum \(q\); \(q_\mu(p)=e^{ip_\mu}-1\) is the
previously defined difference symbol.
\end{lemma}
\begin{proof}
First \(G_1(Z)\) is self-adjoint on a real Coulomb field.
Indeed,
\[
 G_1-G_1^*=
 \tfrac12\Delta^{-1/2}
   \{D^*\operatorname{ad}_Z S_{\rm end}
                +S_{\rm end}^*\operatorname{ad}_Z D\}
                                                    \Delta^{-1/2}.
\]
For each edge the off-diagonal vertex blocks in braces
cancel. The remaining diagonal at a vertex is
\(2\operatorname{ad}_{D^*Z}\), hence zero. It follows
that \(j_M\) is a squared Hilbert--Schmidt expectation.

The free covariance has total trace \(9Vg_M\).
The scalar gradient projection \(UU^*\) has diagonal
\((V-1)/(4V)\) on every link: its trace is \(V-1\),
and translations and coordinate permutations make
all \(4V\) diagonal entries equal.
Since \(\operatorname{tr}_{\rm colour}
(\operatorname{ad}_z)^2=-8|z|^2\),
\[
 V^{-1}\mathbb E\operatorname{Tr}G_2
                  =-\tfrac32(1-V^{-1})g_M.
\]
The Haar contribution in \eqref{eq:v58-density} is
\(-3g_M\) per site. Combining these identities proves
the equality in \eqref{eq:v58-density-average}.

In unitary Fourier coordinates the odd gauge block is
\[
 G_1(p,q)=
 \frac1{\sqrt V\sqrt{\lambda(p)\lambda(q)}}
         \sum_\mu w_\mu(p,q)
                         \operatorname{ad}_{\widehat Z_\mu(p-q)}.
\]
Its two vertex momenta are nonzero, and \(C(0)=0\)
omits the retained harmonic field. Taking its squared
Hilbert--Schmidt norm and contracting colours gives
\begin{equation}\label{eq:v58-j-sum}
 j_M=\frac{24}{V^2}\sum_{p,q\ne0}
       \frac{w(p,q)^T C(p-q)\overline{w(p,q)}}
                              {\lambda(p)\lambda(q)}.
\end{equation}
The colour factor is \(3\cdot8=24\).
As \(|w(p,q)|^2\le\lambda(p)\) and
\(\|C(k)\|\le\lambda(k)^{-1}\), each summand without
the factor \(24/V^2\) is at most
\[
 \frac1{\lambda(q)\lambda(p-q)}.
\]
With the zero convention, summing over all \(p,q\)
and changing variables to \(k=p-q\) bounds it by
\(g_M^2\). This proves the inequality.

To pass to the limit, remove \(|q|<\delta\) and
\(|p-q|<\delta\). The normalized \(\lambda^{-1}\)
sums on either deleted set are \(O(\delta^2)\),
uniformly in \(M\), and their partner sums are bounded.
On the complement the summand is bounded. Its possible
directional discontinuity at \(p=0\) is on a set of
measure zero, and ordinary Riemann sums converge there
after a further arbitrarily small neighborhood is
removed. Let that neighborhood shrink, then let
\(\delta\) shrink. This proves \eqref{eq:v58-j-integral}.
The same one-momentum argument proves convergence
of \(g_M\).
\end{proof}

\subsection{The vacuum-density limit and the next genuine bottleneck}

\begin{theorem}[Identified complete two-loop vacuum density]
\label{thm:v58-vacuum}
For the coefficient \(B_M\) of the actual fixed-mean
local integral in Proposition~\ref{prop:v58-coefficient},
\begin{equation}\label{eq:v58-density-limit}
\begin{aligned}
 b_M:=\frac{B_M}{V}
   &=6s_{4,M}-\tfrac12\sigma_{3,M}^2
           +\left(\tfrac92-\tfrac3{2V}\right)g_M+\tfrac12j_M,\\
 b_M&\longrightarrow
 b_\infty:=6s_{4,\infty}-\tfrac12\sigma_3^2
                                +\tfrac92g_\infty+\tfrac12j_\infty .
\end{aligned}
\end{equation}
Every term in the last line is finite and identified
by an absolutely convergent sum, an integrable
Brillouin-zone integral, or the explicit finite
polynomial \eqref{eq:v58-quartic-polynomial}.
In particular \(\sup_{M\ge4}|B_M|/M^4<\infty\).
No sign is asserted for the total \(b_\infty\).
\end{theorem}
\begin{proof}
Substitute Theorem~\ref{thm:v58-cubic} and
Lemmas~\ref{lem:v58-quartic}--\ref{lem:v58-density}
into \eqref{eq:v58-B}. Each convergence and bound
has been proved independently before taking this sum.
\end{proof}

There is also a useful quantitative warning about the
scale of the first score. Under the free law,
\begin{equation}\label{eq:v58-score-variance}
 \operatorname{Var}\{\beta^{-1/2}S_3(Z)\}
       =\frac{M^4}{\beta}\sigma_{3,M}^2.
\end{equation}
Since \(\sigma_3^2>0\), this variance diverges along
every sequence \(\beta=o(M^4)\), including logarithmic
coupling sequences. This is a statement about the
Gaussian first score, not a theorem of divergence
in total variation for the actual interacting law.
It makes precise why control of a finite local
nonlinearity does not make its unnormalized global
sum a small random variable.

The scalar vacuum coefficient can cancel in ratios
of normalized observables. To use that cancellation
one still needs estimates of the connected remainder
and of insertions near the observable; subtracting
\(Vb_M/\beta\) on paper is not such an estimate.
The next upstream target is therefore a measured,
normalized connected expansion on V56's local
fluctuation domains, together with control of their
compact complement. The fixed-\(M\) cutoff in
\eqref{eq:v58-B} has not been shown suitable for
that joint limit. Source-dependent higher-loop
counterterms, ultraviolet entry into an interacting
continuum law, and a physical time-evolution limit
remain additional requirements.

The target itself is still the interacting
four-dimensional existence and mass-gap problem,
as distinguished in the
\href{https://www.claymath.org/millennium/yang-mills-the-maths-gap/}
{Clay Mathematics Institute problem description}.
A vacuum-pressure coefficient, even with all-volume
control, does not establish it.

\subsection{Exact checks, preservation, and scope}

The certificate
\begin{center}\small
\path{scripts/verify_ym_f14_v58_vacuum_density.py}
\end{center}
uses exact rational arithmetic, including rational
real and imaginary parts. It independently expands
the quaternion plaquette through degree four
(315 quartic monomials), checks
\eqref{eq:v58-quartic-polynomial} on twelve rational
Gram covariances, and verifies the nonzero transverse
vertex against that literal expansion.
It constructs the 144 primitive colour-ordered cubic
terms and computes the finite Fourier contractions.
An independent spatial Hilbert--Schmidt contraction
checks the odd gauge-square term.

On the side-two torus, used only as a finite algebra
regression, the exact results are
\[
 g_2=\frac{103}{768},\qquad
 s_{4,2}=-\frac{186793}{884736},\qquad
 \sigma_{3,2}^2=\frac{2711}{18432},\qquad
 j_2=\frac{667}{12288},\qquad
 b_2=-\frac{11833}{16384}.
\]
No finite lattice is used as evidence for a limit
without the analytic arguments above.
The 9,667-byte certificate has SHA--256
\begin{center}\scriptsize\ttfamily
7857456A20EDE9A17D66A645A9B68513D73C95E624882C7619FB1753309C29FE.
\end{center}
It imports only the unchanged rational-complex
helpers from V56's certificate, whose SHA--256 is
\begin{center}\scriptsize\ttfamily
E82521CBFDCEFFFF5BD0A0194D1B50D7545ADBF4F5213B94C5441B27EA1A6941.
\end{center}

The predecessor V57 has 1,407,702 bytes and SHA--256
\begin{center}\scriptsize\ttfamily
A6C985CDF4FD57651E4A4DDE623AE3021B29AF7D14966C5826B7BB91CAC40994.
\end{center}
Its body is preserved apart from the title version
and this terminal insertion. V58 replaces V57 as
the sole active main note. Supplied manuscripts,
companion files, monographs and archives are unchanged.

\begin{firewall}
The complete zero-source non-Gaussian vacuum coefficient
now has an identified finite limit per lattice site.
The actual cubic Wilson fluctuation has summable
covariance and a strictly positive extensive variance.
These are coefficient and free-reference theorems,
not a uniform interacting remainder theorem.
They do not construct a nontrivial continuum
Yang--Mills theory or prove a positive physical
Hamiltonian mass gap. The full objective remains
unproved and active.
\end{firewall}
\section{V59: normalized local plaquette coefficients and the complete massless quadratic insertion}
\label{sec:v59}

\subsection{Context and the new estimates}

This section was added on 15 September 2026. The previous turn
made concrete progress: V58 identified a finite thermodynamic
limit of the complete vacuum coefficient per site. That
coefficient is nevertheless extensive before normalization.
The present task is to control what remains after it cancels
in a local observable.

We prove three related results for the actual mean-fixed
local Coulomb integral. First, its complete quadratic
measure term is a bounded operator with a summable kernel.
Second, the first two correction coefficients of normalized
observables of finitely many nonlinear plaquette logarithms
have finite, explicitly specified infinite-volume limits.
Third, the complete quadratic connected insertion has a
kernel decaying as distance to the power \(-6\). V21's
already evaluated constant-source identity then implies
that its zero-momentum limit vanishes; its midpoint symbol
is bounded by \(C|k|^2(1+|\log|k||)\).

These are new all-volume estimates on the massless free
reference, not a repetition of the conditional connected
identities in archived eighth-lineage V59, or of V40--V42's
uniformly convex local boxes. We reuse the normalization
identity rather than calling it a new theorem. Likewise,
V21's constant-source calculation is not recomputed as
new work: it supplies an exact sum rule for the connected
kernel constructed here.

The distinctions from V58 remain essential. The local
interacting expansion is initially taken at fixed \(M\).
We subsequently pass to infinite volume in its
\emph{coefficients}, not in a claimed uniformly small
remainder. No full compact-measure comparison or
physical continuum time evolution is obtained.
V55 remains the last regular companion checkpoint;
V60 will review SM, NS and Shell again.

\subsection{The full quadratic measure is a summable operator}

Use V58's notation, including \(V=M^4\), the real
mean-zero Coulomb space \(\mathcal H_M\), covariance
\(C_M\otimes I_3\), and the omitted constant modes.
Let \(G_M=\Delta_M^{-1}\) on primed scalar vertices,
extended by zero on constants. On scalar links define
\[
 A_M=DG_MD^*,\qquad
 B_M^{\rm end}=S_{\rm end}G_MS_{\rm end}^*.
\]
The symbol \(B_M^{\rm end}\) is not the vacuum
coefficient \(B_M\) of V58. Write \(\circ\) for
entrywise matrix multiplication.

\begin{theorem}[Complete quadratic-density kernel]
\label{thm:v59-measure}
On real Coulomb fields, including constant fields,
\begin{equation}\label{eq:v59-Q}
 \ell_2(a)=-\tfrac12\langle a,\mathcal Q_Ma\rangle,\qquad
 \mathcal Q_M=
 \left\{\left(1-\frac1{3V}\right)I
                      +2A_M\circ B_M^{\rm end}\right\}\otimes I_3.
\end{equation}
For links \(e=(x,\mu)\), \(f=(y,\nu)\), its scalar kernel
obeys
\begin{equation}\label{eq:v59-Q-decay}
 |(\mathcal Q_M)_{ef}|\le C(1+d_M(x,y))^{-6},\qquad
 \left(1-\frac1{3V}\right)I\preceq\mathcal Q_M\preceq CI .
\end{equation}
There is a bounded translation-invariant infinite-lattice
operator \(\mathcal Q_\infty\), obtained by the same
kernel formula with local identity coefficient one.
Its entries are fixed-offset limits of those of
\(\mathcal Q_M\), and its kernel is absolutely summable.
\end{theorem}
\begin{proof}
V58 gives
\[
 \ell_2(a)=-\tfrac13\sum_e|a_e|^2
       +\tfrac1{12}\operatorname{Tr}
                      U^*(\operatorname{ad}_a)^2U
       -\tfrac12\operatorname{Tr}G_1(a)^2,
 \quad U=D\Delta^{-1/2}.
\]
The diagonal of \(UU^*=A_M\) is
\((V-1)/(4V)\), and
\(\operatorname{tr}_{\rm colour}(\operatorname{ad}_z)^2
=-8|z|^2\). The first two terms are therefore
\[
 -\left(\frac12-\frac1{6V}\right)\sum_e|a_e|^2.
\]
By the Coulomb self-adjointness identity proved in V58,
\(\operatorname{Tr}G_1(a)^2=\|G_1(a)\|_{\rm HS}^2\).
Expand \(G_1=\tfrac12U^*\operatorname{ad}_a W\),
where \(W=S_{\rm end}\Delta^{-1/2}\). The spatial
Hilbert--Schmidt contraction at two links is
\((UU^*)_{ef}(WW^*)_{ef}\), while
\[
 \operatorname{tr}_{\rm colour}
       (\operatorname{ad}_{a_e})^*\operatorname{ad}_{a_f}
                           =8a_e\cdot a_f .
\]
Consequently
\[
 \operatorname{Tr}G_1(a)^2
    =2\sum_{e,f}(A_M)_{ef}(B_M^{\rm end})_{ef}\,a_e\cdot a_f,
\]
which proves \eqref{eq:v59-Q}.

The scalar version of Lemma~\ref{lem:v58-Green}, with
symbol \(\lambda^{-1}\), gives
\[
 |(A_M)_{ef}|\le C(1+d_M(x,y))^{-4},\qquad
 |(B_M^{\rm end})_{ef}|\le C(1+d_M(x,y))^{-2}.
\]
The first estimate uses two differences; endpoint sums
are only fixed shifts in the second. Their product
proves the kernel estimate. Its row and column sums
are bounded, giving the operator upper bound.

Both \(A_M\) and \(B_M^{\rm end}\) are positive
semidefinite. If \(A_M=\sum_i u_i u_i^*\) and
\(B_M^{\rm end}=\sum_j v_jv_j^*\), their entrywise
product is
\(\sum_{i,j}(u_i\circ v_j)(u_i\circ v_j)^*\)
in these real coordinates. It is positive semidefinite,
which proves the lower bound. Fixed-offset convergence
and the uniform summable tails construct
\(\mathcal Q_\infty\) and give its asserted properties.
\end{proof}

The positive local term in \(\mathcal Q_M\) is
\emph{not} a Yang--Mills mass. It is only one part
of a perturbative insertion, in gauge-fixed
coordinates. The Wilson contractions at the
same order must be added before making any
statement about its long-distance behavior.

\subsection{Local curvature observables and the exact normalized coefficients}

Fix a finite list \(\mathcal P\) of oriented plaquettes,
embedded in all sufficiently large tori without changing
its offsets. Let \(LZ=((dZ)_p)_{p\in\mathcal P}\), including
the three colour components. For each component \(\alpha\)
put
\[
 h_{\alpha,M}=C_M L_\alpha^*,\qquad
 h_{\alpha,\infty}=C_\infty L_\alpha^*.
\]
The colour identity is implicit. If the patch lies in
a fixed neighborhood of the origin, V58 gives
\begin{equation}\label{eq:v59-external}
 |h_{\alpha,M}(x)|\le C(1+d_M(x,0))^{-3}.
\end{equation}
A local undifferentiated link insertion instead has
bound \(C(1+d_M(x,0))^{-2}\). This distinction will
be used, not suppressed.

Let \(F\) be a \(C^8\) function of the finitely many
components of \(LZ\), with every derivative through
order eight bounded by \(K(1+|y|)^d\), for some
fixed \(K,d\). Bounded smooth functions and polynomial
energy observables are both allowed. Gaussian
expectations of their derivative bounds are uniform
in \(M\), since the covariance on a fixed patch is
uniformly bounded.

To retain the observable's own nonlinearity, define
\[
 Y_{p,t}(Z)=t^{-1}\log\!\left(
 e^{tv_1}e^{tv_2}e^{tv_3}e^{tv_4}\right)
       =(LZ)_p+t\,b_{2,p}(Z)+t^2b_{3,p}(Z)+O(t^3),
\]
on the principal local logarithm branch. The \(v_i\)
are the oriented vectors in \eqref{eq:v58-cubic}.
These logarithms are identified with their colour
vectors in the same basis. The quadratic and cubic
polynomials can be specified without a matrix logarithm:
start \(A=B=C=0\), append each \(v=v_i\), and update
\begin{equation}\label{eq:v59-BCH}
\begin{aligned}
 A_{\rm new}&=A+v,\\
 B_{\rm new}&=B+A\times v,\\
 C_{\rm new}&=C+B\times v+
       \tfrac13\{A\times(A\times v)+v\times(v\times A)\}.
\end{aligned}
\end{equation}
The right sides use the old \(A,B,C\).
After four updates, \(A=LZ\), \(B=b_2\), \(C=b_3\).
This is the degree-three BCH formula with
\([a,b]=2a\times b\); it also follows directly
from the quaternion product.

Define the local functions
\begin{equation}\label{eq:v59-observable-jets}
\begin{aligned}
 R_0(Z)&=F(LZ),\\
 R_1(Z)&=DF(LZ)[b_2(Z)],\\
 R_2(Z)&=DF(LZ)[b_3(Z)]
                  +\tfrac12D^2F(LZ)[b_2(Z),b_2(Z)].
\end{aligned}
\end{equation}
Use a translation-, hypercubic- and colour-invariant
cutoff of the fixed-\(M\) kind allowed in V58;
a sufficiently small radial cutoff in the linear
Coulomb coordinates is one choice. Its support is
also chosen inside all the plaquette logarithm
branches in this patch. Let \(\mu_{M,t}\) be the
probability law of \(Z=a/t\) in that exact local
integral, with \(t=\beta^{-1/2}>0\).

\begin{theorem}[Normalized nonlinear plaquette coefficients]
\label{thm:v59-local}
At each fixed \(M\),
\begin{equation}\label{eq:v59-local-expansion}
 \mathbb E_{\mu_{M,t}}F(Y_t)
       =A_{0,M}+tA_{1,M}+t^2A_{2,M}+O_M(t^3),
\end{equation}
where all expectations on the right below are free
Gaussian expectations on \(\mathcal H_M\):
\begin{equation}\label{eq:v59-local-coefficients}
\begin{aligned}
 A_{0,M}&=\mathbb E R_0,\\
 A_{1,M}&=\mathbb E R_1-\mathbb E(R_0S_3),\\
 A_{2,M}&=\mathbb E R_2-\mathbb E(R_1S_3)
       -\operatorname{Cov}(R_0,S_4-\ell_2)
                    +\tfrac12\operatorname{Cov}(R_0,S_3^2).
\end{aligned}
\end{equation}
For each fixed patch and \(F\), these three coefficients
are uniformly bounded in \(M\) and converge to finite
infinite-lattice limits. The limiting coefficients
are given by the absolutely convergent connected
contractions specified in the proof.

If \(F(-y)=F(y)\), then \(A_{1,M}=0\) and the remainder
in \eqref{eq:v59-local-expansion} improves to \(O_M(t^4)\).
Neither remainder bound is asserted to be uniform in \(M\).
\end{theorem}
\begin{proof}
The normalized fixed-\(M\) expansion of V58 gives
the weight
\[
 1-tS_3+t^2\{\tfrac12S_3^2-S_4+\ell_2\}+O(t^3)
\]
under the Gaussian integral. Multiply by
\(R_0+tR_1+t^2R_2\) and divide by its expectation
with observable one. Since \(\mathbb E S_3=0\),
the disconnected order-\(t^2\) term cancels exactly
and gives \eqref{eq:v59-local-coefficients}.
Fixed-dimensional localization and Gaussian domination
justify the expansion as in V58; polynomial growth
of \(F\) does not affect that argument.
If \(F\) is even, \(R_0\) is even, \(R_1\) odd,
and \(R_2\) even. More generally changing
\((t,Z)\) to \((-t,-Z)\) leaves the rescaled
weight unchanged and sends \(Y_t\) to \(-Y_t\).
All odd asymptotic coefficients vanish. Expanding
one further order gives the improved remainder.

It remains to prove the new all-volume assertions.
All Wick orderings below use the actual covariance
\(C_M\), not an independent-site covariance.
For Gaussian coordinates and a smooth local function,
integration by parts gives
\begin{equation}\label{eq:v59-Wick-IBP}
 \mathbb E\{F(LZ):Z_{i_1}\cdots Z_{i_n}:\}
 =\sum_{\alpha_1,\ldots,\alpha_n}
       \left(\prod_{j=1}^n(h_{\alpha_j,M})_{i_j}\right)
          \mathbb E(\partial_{\alpha_1}\cdots
                               \partial_{\alpha_n}F)(LZ).
\end{equation}
One obtains this identity by successive Gaussian
integration by parts; the Wick subtractions cancel
all contractions among the displayed \(Z\)'s.
Approximation by bounded smooth functions justifies
it for the stated polynomial-growth class.

Every local cubic Wilson vertex is already Wick
ordered, as proved in V58. Thus
\(\mathbb E(R_0S_3)\) has three external kernels
of type \eqref{eq:v59-external}. Its absolute
spatial sum is bounded by
\(\sum_x C(1+d_M(x,0))^{-9}\).
The functions \(R_1,R_2\) themselves involve only
the fixed patch, and have uniform Gaussian moments.

For the quartic term, local Wick expansion has
degrees four, two, and zero. Its internal contractions
are uniformly bounded fixed-offset covariances.
The zero-degree term cancels in the covariance.
The two remaining external sums have respectively
four and two factors of \eqref{eq:v59-external},
and hence decay as powers \(-12\) and \(-6\).
They are uniformly summable.

For the density term, \eqref{eq:v59-Q} and
\eqref{eq:v59-Wick-IBP} give the explicit formula
\begin{equation}\label{eq:v59-density-insertion}
 \operatorname{Cov}(R_0,\ell_2)=
 -\tfrac12\sum_{\alpha,\gamma}
  \mathbb E(\partial_\alpha\partial_\gamma F)(LZ)
         \langle h_{\alpha,M},\mathcal Q_Mh_{\gamma,M}\rangle .
\end{equation}
Its double kernel sum is absolutely convergent,
uniformly in \(M\): \(h_\alpha\) is uniformly
square summable, and the absolute kernel of
\(\mathcal Q_M\) is uniformly summable.

The only nontrivial two-vertex sum is
\(\operatorname{Cov}(R_0,S_3^2)\). If \(r\) lines
are contracted between two cubic Wick vertices,
the product formula has coefficients
\[
 r!\binom3r^2=1,\ 9,\ 18,\ 6,\qquad r=0,1,2,3 .
\]
There remain \(6-2r\) Wick legs. For \(r=3\)
there are none, and this is exactly the vacuum
term removed by the covariance. For \(r=0\)
there are three external legs at each vertex;
for \(r=1\), two at each vertex; for \(r=2\),
one at each vertex. The corresponding absolute
majorants, up to fixed shifts and constants, are
\begin{equation}\label{eq:v59-majorants}
\begin{array}{c|c}
 r&\text{majorant at vertices }x,y\\ \hline
 0&(1+d_M(x,0))^{-9}(1+d_M(y,0))^{-9}\\
 1&(1+d_M(x,0))^{-6}(1+d_M(y,0))^{-6}
                                  (1+d_M(x,y))^{-2}\\
 2&(1+d_M(x,0))^{-3}(1+d_M(y,0))^{-3}
                                  (1+d_M(x,y))^{-4}.
\end{array}
\end{equation}
The first two double sums are bounded directly.
For the third set
\(u_M(x)=(1+d_M(x,0))^{-3}\) and
\(v_M(x)=(1+d_M(x,0))^{-4}\).
In counting-measure norms both have uniformly
bounded \(\ell^{3/2}\) norms. Young and Hölder give
\begin{equation}\label{eq:v59-Young}
 \sum_{x,y}u_M(x)u_M(y)v_M(x-y)
       \le \|u_M\|_{3/2}^2\|v_M\|_{3/2}\le C .
\end{equation}
These inequalities hold on every finite torus
with the same constants as on the infinite lattice.
Replacing either \(u_M\) by its part outside
distance \(R\) bounds the corresponding tail by
\(CR^{-1/3}\). Indeed
\(\|u_M\mathbf1_{\{d_M>R\}}\|_{3/2}\le CR^{-1/3}\).
This gives a uniform vanishing tail in both
vertex positions, without assuming exponential
decay or a massive covariance.

Finally consider \(\mathbb E(R_1S_3)\).
Apply three Gaussian integrations by parts at
the cubic vertex. The local polynomial \(b_2\)
in \(R_1=DF(LZ)[b_2]\) has degree two, so at least
one of these derivatives must hit the factor
depending on \(LZ\). There is therefore at least
one external curvature kernel of decay \(-3\);
the other two can be undifferentiated local-link
kernels of decay \(-2\). The worst bound is
\(C(1+d_M(x,0))^{-7}\). Remaining local polynomial
factors have uniform Gaussian moments. This
establishes summability while retaining the
nonlinear observable, rather than replacing
it by \(F(LZ)\).

All finite-patch covariances converge by V58.
Their Gaussian expectations of the continuous
derivative functions and polynomial factors
converge by uniform Gaussian moment bounds.
For the spatial sums, first restrict all vertices
to a fixed radius. Their individual contractions
then converge. The tail estimates just proved,
including \eqref{eq:v59-Young}, remove the
restriction uniformly. This constructs the
claimed absolutely convergent limits and
completes the proof.
\end{proof}

For example,
\[
 F(y)=\exp\!\left(-\sum_{p\in\mathcal P}s_p|y_p|^2\right),
 \qquad s_p\ge0,
\]
gives bounded local tests of the actual plaquette
logarithms. They are invariant under conjugation
at each plaquette base. The theorem concerns
their expectation in the specified mean-fixed
local law, not in an already constructed
infinite-volume interacting law.

\subsection{An explicit connected generating coefficient}

Let \(\mathcal T_M,\mathcal U_M\) be the real symmetric
polarizations of the literal homogeneous polynomials:
\[
 S_3(a)=\mathcal T_M(a,a,a),\qquad
 S_4(a)=\mathcal U_M(a,a,a,a).
\]
For this subsection only, take the source observable
to be the \emph{linear} curvature characteristic
function and define
\[
 \Psi_M(J,t)=\log\mathbb E_{\mu_{M,t}}e^{i\langle J,LZ\rangle},
 \qquad h=C_ML^*J,\qquad \nu=\langle J,LC_ML^*J\rangle .
\]
The logarithm is continued from its nonzero Gaussian
value at \(t=0\). The distinction from the nonlinear
observable of Theorem~\ref{thm:v59-local} is explicit.

\begin{proposition}[No volume term and no sixth-degree source term]
\label{prop:v59-characteristic}
At fixed \(M\),
\begin{equation}\label{eq:v59-characteristic}
\begin{aligned}
 \Psi_M(J,t)
   &=-\tfrac12\nu+it\,\mathcal T_M(h,h,h)
                           +t^2\{\mathfrak R_M(h)+\mathfrak V_M(h)\}
                           +O_M(t^3),\\
 \mathfrak R_M(h)
   &=6\mathbb E\mathcal U_M(Z,Z,h,h)
       -\tfrac92\mathbb E\mathcal T_M(Z,Z,h)^2
                              +\tfrac12\langle h,\mathcal Q_Mh\rangle,\\
 \mathfrak V_M(h)
   &=-\mathcal U_M(h,h,h,h)
                              +\tfrac92\mathbb E\mathcal T_M(Z,h,h)^2 .
\end{aligned}
\end{equation}
The displayed source polynomials have degrees
two, three and four only. All their coefficients
have finite infinite-volume limits for the fixed
local curvature source. Bounds are uniform on
compact sets of \(J\), as are the bounds on each
fixed derivative of these coefficient polynomials.
\end{proposition}
\begin{proof}
The finite-dimensional Gaussian shift identity is
\[
 \mathbb E\{e^{i\langle J,LZ\rangle}P(Z)\}
       =e^{-\nu/2}\mathbb E P(Z+ih)
\]
for every polynomial \(P\). It follows either by
differentiating the Gaussian exponential moment
formula or by completing its square. Use it in
the normalized expansion of V58 and then take
the logarithm.

All single contractions of \(\mathcal T_M\) vanish
by the colour identity. Consequently
\[
\begin{aligned}
 \mathcal T_M(Z+ih,Z+ih,Z+ih)
  ={}&\mathcal T_M(Z,Z,Z)
       +3i\mathcal T_M(Z,Z,h)\\
    &-3\mathcal T_M(Z,h,h)-i\mathcal T_M(h,h,h).
\end{aligned}
\]
The four terms have Gaussian chaos degrees
three, two, one and zero. Distinct degrees
are orthogonal for Gaussian expectation.
In the second cumulant the last, constant term
cancels. Subtracting the unshifted variance
therefore leaves
\[
 -9\mathbb E\mathcal T_M(Z,Z,h)^2
                       +9\mathbb E\mathcal T_M(Z,h,h)^2 .
\]
This cancellation removes the putative
sixth-degree source term. Expansion of the
shifted quartic expectation gives
\[
 \mathbb E S_4(Z+ih)-\mathbb E S_4(Z)
     =-6\mathbb E\mathcal U_M(Z,Z,h,h)
                           +\mathcal U_M(h,h,h,h),
\]
and the density shift is
\(\mathbb E\ell_2(Z+ih)-\mathbb E\ell_2(Z)
=\tfrac12\langle h,\mathcal Q_Mh\rangle\).
Their signs in the normalized logarithm give
\eqref{eq:v59-characteristic}.

The estimates in Theorem~\ref{thm:v59-local}
apply directly to these contractions.
For instance \(\mathbb E\mathcal T_M(Z,Z,h)^2\)
has the \(r=2\) majorant in
\eqref{eq:v59-majorants}. The term with one
Gaussian and two \(h\)'s at each vertex is
bounded by two summable \(|h|^2\) factors and
a bounded covariance. The remaining terms
are local sums of at least two \(h\)'s.
Polynomial dependence on the finite vector
\(J\) gives the last assertion.
\end{proof}

\subsection{The complete quadratic insertion and its zero-momentum cancellation}

The expression \(\mathfrak R_M(h)\) in
\eqref{eq:v59-characteristic} also makes sense
for an arbitrary real link field \(h\), using
the operator extension in \eqref{eq:v59-Q}.
This extension agrees with the actual density
on Coulomb fields; no assertion about a
non-Coulomb gauge density is needed.
It defines a real symmetric, translation-invariant
quadratic kernel \(\mathcal R_M\) by
\[
 \mathfrak R_M(h)=\langle h,\mathcal R_Mh\rangle .
\]
Global colour invariance makes this kernel
scalar in colour.

\begin{theorem}[Summable self-energy kernel with no constant mass term]
\label{thm:v59-selfenergy}
The complete kernel satisfies
\begin{equation}\label{eq:v59-R-kernel}
 |(\mathcal R_M)_{\mu\nu}(x)|
                  \le C(1+d_M(x,0))^{-6}.
\end{equation}
Its fixed-offset limit \(\mathcal R_\infty\) is
absolutely summable. In Fourier coordinates,
\begin{equation}\label{eq:v59-R-zero}
 \widehat{\mathcal R}_M(0)=\frac{a_M}{M^2}I_4,\qquad
 \widehat{\mathcal R}_\infty(0)=0,
\end{equation}
where \(a_M\) is the already identified bounded
coefficient in \eqref{eq:v21-quadratic}.

Let \(E(k)=\operatorname{diag}(e^{ik_\mu/2})\)
and use midpoint link coordinates:
\[
 \widetilde{\mathcal R}_\infty(k)
       =E(k)^*\widehat{\mathcal R}_\infty(k)E(k).
\]
This matrix is even in \(k\), and for
\(0<|k|\le1/2\),
\begin{equation}\label{eq:v59-R-symbol}
 \|\widetilde{\mathcal R}_\infty(k)\|
                    \le C|k|^2(1+|\log|k||).
\end{equation}
No sign or nonperturbative spectral gap is
asserted for this kernel.
\end{theorem}
\begin{proof}
The quartic tadpole
\(6\mathbb E\mathcal U_M(Z,Z,h,h)\)
has a finite-range kernel with uniformly
bounded coefficients, since both external
legs lie in one plaquette. The density part
has \eqref{eq:v59-Q-decay}.

For the cubic-square part use the polarization
of \(s_{3,p}=(dZ)_p\cdot B_{2,p}(Z)\).
Each term of \(\mathcal T_M(Z,Z,h)\) contains
one \(h\), two Gaussian fields, and one
discrete curl applied either to \(h\) or to
one of those Gaussian fields. A single
internal contraction vanishes by colour,
so its square has exactly two cross
covariances between the two vertices.

If both curls act on Gaussian fields, the
two-line kernel has two derivatives in total:
it is bounded by \(C(1+d)^{-6}\).
If one curl acts on \(h\), sum by parts at
that external vertex. The derivative moves
to the product of cross covariances.
There are again two derivatives in total.
If both curls act on the two external
\(h\)'s, sum by parts at both vertices;
two differences act on a product of two
undifferentiated covariances. The discrete
product rule includes only bounded shifts,
and in every case Lemma~\ref{lem:v58-Green}
gives the same decay \(-6\).
These are finite sums of local terms, with
no boundary contribution on the torus.
This proves \eqref{eq:v59-R-kernel}.
The argument also gives fixed-offset
convergence. Its summable bound has tail
\(O(R^{-2})\), uniformly in \(M\), so the
infinite kernel and its zero Fourier value
are well defined.

To identify that value, take \(h_\mu(x)=c_\mu\)
constant, temporarily without requiring
zero mean. On the mean-zero normal space
the literal Hessian at the background \(s h\)
has the expansion
\[
 H(sh)=H(0)+6s\,\mathcal T_M(h,\cdot,\cdot)
                  +12s^2\,\mathcal U_M(h,h,\cdot,\cdot)
                  +O(s^3).
\]
The coefficient of \(s^2\) in
\(\tfrac12\log\det H(sh)-\ell(sh)\),
relative to its value at zero, is precisely
\(\mathfrak R_M(h)\). To see the second
term, the trace-log square is
\(-9\operatorname{Tr}(C_M T_h C_M T_h)\),
where \(T_h=\mathcal T_M(h,\cdot,\cdot)\);
Wick's rule gives
\(\mathbb E\mathcal T_M(Z,Z,h)^2
=2\operatorname{Tr}(C_M T_h C_M T_h)\).
The mean of that quadratic polynomial is
zero by colour. The other two terms are
the first trace and the density contribution.

This is exactly the quadratic part of
V21's complete mean-fixed amplitude,
with physical source \(Mc\) in its
notation, not a root source. Therefore
\[
 \mathfrak R_M(h)=a_M M^2\sum_\mu|c_\mu|^2 .
\]
For a translation-invariant quadratic
kernel the left side is
\(V\sum_{\mu,\nu}c_\mu\cdot
       \widehat{\mathcal R}_{M,\mu\nu}(0)c_\nu\).
Polarization proves the first identity
in \eqref{eq:v59-R-zero}. The boundedness
of \(a_M\), already proved in V21, and
the uniform summable kernel tails imply
the second. This uses the known
constant-source cancellation as a sum
rule; it does not estimate the Haar,
ghost and Wilson pieces separately.

Finally full spatial inversion acts on
link fields by
\[
 h_\mu(x)\longmapsto-h_\mu(-x-e_\mu).
\]
It preserves the literal action, the
Gaussian covariance and \(\mathcal Q_M\).
Thus the infinite kernel obeys
\[
 \mathcal R_{\infty,\mu\nu}(x)
   =\mathcal R_{\infty,\mu\nu}(-x-e_\mu+e_\nu).
\]
The displacement between link midpoints
is \(\delta=x+(e_\mu-e_\nu)/2\); inversion
sends \(\delta\) to \(-\delta\).
Consequently the midpoint symbol is even
and its exponential Fourier factors can
be replaced by cosines. Subtract its
zero value from \eqref{eq:v59-R-zero}.
For \(K=|k|^{-1}\),
\[
\begin{aligned}
 |\widetilde{\mathcal R}_{\infty,\mu\nu}(k)|
 &\le C\sum_x(1+|x|)^{-6}
                      \min\{2,|k|^2(1+|x|)^2\}\\
 &\le C|k|^2\sum_{|x|\le K}(1+|x|)^{-4}
                         +C\sum_{|x|>K}(1+|x|)^{-6}\\
 &\le C|k|^2(1+\log K).
\end{aligned}
\]
There are only four orientation indices.
This proves \eqref{eq:v59-R-symbol}.
\end{proof}

The result identifies a perturbative cancellation
that would be lost by exponentiating the positive
measure term alone. It does not forbid
nonperturbative mass generation, and it is
not a statement about the spectrum of
gauge-invariant physical states.

\subsection{A check on a genuine one-plaquette energy observable}

For one plaquette \(p\) define
\[
 \mathcal W_p(a)=1-\tfrac12\operatorname{tr}U_p(a),
 \qquad
 \mathcal L_p(a)=\tfrac12|\log U_p(a)|^2
\]
on the local branch. Both are conjugation-invariant.

\begin{proposition}[Normalized energy coefficients]
\label{prop:v59-energy}
For the symmetric fixed-\(M\) local law,
\begin{equation}\label{eq:v59-energy}
\begin{aligned}
 \mathbb E_\mu\{\beta\mathcal W_p(a)\}
    &=\frac34(1-V^{-1})-\frac{b_M}{6\beta}
                                                 +O_M(\beta^{-2}),\\
 \mathbb E_\mu\{\beta\mathcal L_p(a)\}
    &=\frac34(1-V^{-1})
        +\frac1\beta\left\{-\frac{b_M}{6}
                    +\frac58\left(\frac{1-V^{-1}}2\right)^2\right\}
                                                 +O_M(\beta^{-2}).
\end{aligned}
\end{equation}
Hence these coefficients converge respectively to
\(3/4\), \(-b_\infty/6\), and
\(-b_\infty/6+5/32\), with the roles indicated
by the two lines. This is a coefficientwise
statement, not a joint-limit asymptotic.
\end{proposition}
\begin{proof}
The cutoff is independent of \(\beta\), so
\(-\partial_\beta\log I_M(\beta)=\mathbb E_\mu S\).
At fixed \(M\), the differentiable Laplace expansion gives
\[
 \mathbb E_\mu(\beta S)=\tfrac92(V-1)-B_M/\beta
                                             +O_M(\beta^{-2}).
\]
Divide by \(6V\), using the symmetries of
the chosen cutoff. Equivalently, this
coefficient follows directly from
\eqref{eq:v59-local-coefficients} and the
Gaussian identity
\(\operatorname{Cov}(S_2,P_n)=\tfrac n2\mathbb E P_n\)
for homogeneous \(P_n\).

The free plaquette curvature has covariance
\(\kappa_M I_3\), where
\(\kappa_M=(1-V^{-1})/2\). Indeed the scalar
trace of \(dC_Md^*\) is \(3(V-1)\);
symmetry divides this equally among \(6V\)
plaquettes. Thus
\(\mathbb E|(dZ)_p|^4=15\kappa_M^2\).
On the principal branch,
\[
 \beta\mathcal W_p(a)
   =\tfrac12|Y_{p,t}|^2-\tfrac{t^2}{24}|Y_{p,t}|^4+O(t^4|Y_{p,t}|^6).
\]
The fixed-\(M\) Gaussian leading expectation
of the quartic term is \(15\kappa_M^2/24\).
This yields the second line of
\eqref{eq:v59-energy}. V58 supplies convergence
of \(b_M\).
\end{proof}

\subsection{What has and has not been transferred}

The volume-dependent vacuum factor is now
cancelled inside actual normalized local
coefficients. Their massless contractions
are absolutely summable, including the
nonlinear observable and the complete
quadratic measure. The connected quadratic
kernel has an identified zero-momentum
cancellation and a logarithmic second-order
momentum bound.

None of these conclusions controls the
remainder of \eqref{eq:v59-local-expansion}
uniformly along a refining lattice and
running coupling. Establishing that
estimate requires a controlled normalized
expansion on local typical-field domains,
not merely more coefficients. The compact
exterior, compatibility of the domains
under integration, and physical source
and time matching must also be retained.
The present work does not replace these
requirements by an assumed massive
Gaussian covariance.

\subsection{Executed exact checks and preservation}

The certificate
\begin{center}\small
\path{scripts/verify_ym_f14_v59_local_insertions.py}
\end{center}
performs twenty rational quaternion checks
of \eqref{eq:v59-BCH}. On the actual side-two
torus it constructs the full 64-link
density kernel and compares it with four
independent primed gauge-matrix trace
calculations on exact Coulomb fields.
It also verifies the colour-ordered
Gaussian shift cancellation and the
Wick-product coefficients \(1,9,18,6\).
Those finite algebra tests support,
but do not replace, the summability proofs.

For a constant unit background in one
orientation and colour at \(M=2\), the
three quadratic contributions are
\[
 6\mathbb E\mathcal U_2(Z,Z,h,h)=-\frac{16}{3},
 \qquad
 9\mathbb E\mathcal T_2(Z,Z,h)^2=0,
 \qquad
 \tfrac12\langle h,\mathcal Q_2h\rangle=\frac{47}{6}.
\]
Their sum is \(5/2\), so the zero-momentum
symbol is \(5/32=a_2/4\), with \(a_2=5/8\).
These terms are calculated from the
literal quartic plaquette and cubic
polynomials and the density kernel,
independently of substituting V21's
answer. The exact local-energy
coefficients in \eqref{eq:v59-energy}
at \(M=2\) are
\[
 \frac{45}{64},\qquad
 \frac{11833}{98304},\qquad
 \frac{25333}{98304}.
\]

The 12,217-byte certificate has SHA--256
\begin{center}\scriptsize\ttfamily
0AAF2695D84E06197AF1D54E30811D7783EE51B6320A6268E8A63B60B0388C73.
\end{center}
It imports the unchanged V56 rational-complex
helpers and V58 finite plaquette routines;
their fingerprints are recorded in V58.
The predecessor V58 has 1,432,791 bytes
and SHA--256
\begin{center}\scriptsize\ttfamily
AF39D741A42C1A84EC5C076FDC067479BC4B4F97F71FA5E8A0B1B210168E0E30.
\end{center}
Its complete body is retained apart from
the title version and this terminal
addition. V59 replaces V58 as the single
active main note. The supplied manuscripts,
companion files, monographs and archived
lineages remain unchanged.

\begin{firewall}
We have proved convergence of complete
normalized local plaquette coefficients
and an infrared bound for the complete
quadratic connected insertion. We have
not proved a uniform nonperturbative
remainder, constructed an interacting
continuum Yang--Mills theory, or
established a positive physical
Hamiltonian mass gap. The original
objective remains unproved and active.
\end{firewall}
\section{V60: a probabilistic gauge floor, its soft-mode logarithm, and the full density beyond a Taylor coefficient}
\label{sec:v60}

\subsection{Checkpoint and why the argument moves upstream}

This section was added on 15 September 2026. V59 was a
progress turn: it proved normalized local coefficient
limits and identified the complete quadratic insertion.
Those statements still expand the interacting law at
fixed volume. Before trying to make its remainder uniform,
one must know where the actual logarithmic gauge density
is regular on fields typical of the reference law.
A global critical-norm ball does not answer that
probability question.

The scheduled companion review was carried out before
selecting this direction. The current regular files are
unchanged: SM--Einstein V72, NS V26, and Shell V16.
Their terminal mathematical sections were re-read.
SM V72 proves an empirical covariance limit in a
declared many-copy, fixed-mode model; it does not
give a spatial-refinement estimate for this gauge
operator. NS V26's rank-one Oseen-kernel contractions
do not transfer as a numerical gauge-fibre bound.
Shell V16 supplies the exact logarithmic gauge
density, but expressly leaves a joint
nonperturbative coupling/refinement remainder
and Gribov coverage open. No companion result
removes the present task.

The new results below concern the actual free
mean-zero Coulomb field. Its odd Faddeev--Popov
matrix has logarithmic matrix variance, with an
identified soft-mode coefficient. A matrix
Gaussian estimate then gives a high-probability
positive branch for the full, untruncated gauge
operator when \(\beta\) is sufficiently large
compared with \((\log M)^2\). On the same event,
the full logarithmic quotient density is
approximated per site with an explicit remainder.
This is a probabilistic result about an
untruncated function, not just another coefficient.

The sufficient coupling window is not silently
substituted for the desired interacting continuum
trajectory. A lower bound also shows a genuine
logarithmic soft-mode cost. The gap between that
necessary scale and the sufficient square-logarithmic
scale, and transfer from the free law to the
interacting law, are left explicit.

\subsection{The exact gauge operator and its Gaussian variance}

Retain the unit-spacing torus, \(V=M^4\), and
V58's free real Gaussian \(Z\) on \(\mathcal H_M\).
Let \(d_{\rm g}=3(V-1)\) be the real primed
vertex dimension. For real Coulomb fields set
\begin{equation}\label{eq:v60-gauge}
\begin{aligned}
 \mathcal F_M(a)
   &=\Delta^{-1/2}\mathcal M_a\Delta^{-1/2}
       =I+G_1(a)+G_{\rm e}(a),\\
 G_1(a)&=\tfrac12U^*\operatorname{ad}_a W,\qquad
 U=D\Delta^{-1/2},\quad W=S_{\rm end}\Delta^{-1/2},\\
 G_{\rm e}(a)&=U^*\{\sigma(\operatorname{ad}_a)-I\}U,\qquad
 \sigma(A)=\tfrac A2\coth(\tfrac A2).
\end{aligned}
\end{equation}
All vertex inverses are primed. As proved in V58,
\(G_1(a)\) is self-adjoint when \(D^*a=0\);
the even term is self-adjoint as well.

Use momenta \(p,q\) in the finite Brillouin zone,
the difference symbol \(q_\mu(p)=e^{ip_\mu}-1\),
and covariance \(C(k)\) from \eqref{eq:v58-free}.
Define
\begin{equation}\label{eq:v60-phi}
\begin{aligned}
 w_\mu(p,q)&=\tfrac12\overline{q_\mu(p)}(1+e^{iq_\mu}),\\
 \Phi_M(p,q)&=
   \frac{w(p,q)^T C(p-q)\overline{w(p,q)}}
                       {\lambda(p)\lambda(q)},\qquad p,q\ne0,\\
 \nu_M(p)&=\frac8V\sum_{q\ne0}\Phi_M(p,q).
\end{aligned}
\end{equation}
Set the summand to zero if \(p=q\), because
the constant link field is omitted. In the
exponential above \(q_\mu\) is the component
of the momentum \(q\), not a difference symbol.

\begin{theorem}[Exact matrix variance and its logarithmic upper bound]
\label{thm:v60-variance}
The Fourier blocks of the matrix variance are
\begin{equation}\label{eq:v60-variance}
 \mathbb E\,G_1(Z)^2\big|_p=\nu_M(p)I_3,\qquad
 \sigma_M^2:=\|\mathbb E\,G_1(Z)^2\|
       \le\frac8V\sum_{q\ne0}\lambda(q)^{-2}
       \le C(1+\log M).
\end{equation}
In particular, with \(\mathcal L_M=1+\log M\),
\begin{equation}\label{eq:v60-matrix-tail}
 \mathbb P\{\|G_1(Z)\|\ge s\}
       \le 2d_{\rm g}
             \exp\!\left(-\frac{s^2}{2C\mathcal L_M}\right),
 \qquad s\ge0 .
\end{equation}
\end{theorem}
\begin{proof}
The Fourier block of \(G_1\) is
\[
 G_1(p,q)=
 \frac1{\sqrt V\sqrt{\lambda(p)\lambda(q)}}
       \sum_\mu w_\mu(p,q)
                        \operatorname{ad}_{\widehat Z_\mu(p-q)}.
\]
Translation invariance makes \(\mathbb E G_1^2\)
diagonal in momentum. Colour invariance makes
each block scalar. For a unit colour vector,
\[
 \sum_{a=1}^3
       (\operatorname{ad}_{E_a})^*\operatorname{ad}_{E_a}
                                                        =8I_3.
\]
Together with the covariance of \(\widehat Z\),
this gives \(\nu_M\). The factor is eight for
one input colour, not the factor 24 used
for a full colour trace in V58.

The numerator in \eqref{eq:v60-phi} is nonnegative,
\(|w(p,q)|^2\le\lambda(p)\), and
\(\|C(k)\|\le\lambda(k)^{-1}\). Hence
\[
 0\le\Phi_M(p,q)
              \le\frac1{\lambda(q)\lambda(p-q)}.
\]
Cauchy--Schwarz for the momentum sum, with
\(\lambda^{-1}\) extended by zero at zero,
bounds its convolution by
\(V^{-1}\sum_{q\ne0}\lambda(q)^{-2}\).
For centered lattice representatives,
\(\lambda(2\pi n/M)\ge16|n|^2/M^2\).
Four-dimensional annulus counting gives
\(\sum_{0<|n|\le CM}|n|^{-4}\le C(1+\log M)\),
proving \eqref{eq:v60-variance}.

Expand \(Z\) in any real orthonormal Gaussian
coordinate system on its covariance range:
\(G_1(Z)=\sum_j\xi_j A_j\), with independent
standard real Gaussians \(\xi_j\).
Each fixed \(A_j\) is self-adjoint, since
the corresponding field is Coulomb, and
\(\sum_j A_j^2=\mathbb E G_1(Z)^2\).
The norm inequality in Theorem 4.1,
equation (4.4), of
\href{https://tropp.caltech.edu/papers/Tro11-User-Friendly-preprint.pdf}
{J.~A. Tropp, \emph{User-friendly tail bounds for sums
of random matrices}, Foundations of Computational
Mathematics 12 (2012), 389--434}
therefore applies in dimension \(d_{\rm g}\).
It gives \eqref{eq:v60-matrix-tail} with the
variance bound just proved. No independence of
matrix entries or of neighboring links is assumed.
\end{proof}

As a consistency identity with V58,
\[
 \frac1V\mathbb E\operatorname{Tr}G_1(Z)^2
                    =\frac3V\sum_{p\ne0}\nu_M(p)=j_M.
\]
The bounded average \(j_M\) does not imply
a bounded maximum \(\sigma_M^2\).

\subsection{A soft mode has a genuine logarithmic cost}

Let \(p_M=(2\pi/M,0,0,0)\), and let \(f_M\)
be the unit complex Fourier vertex vector at
that momentum in any fixed unit colour.
Complexification does not change the norm
of the real self-adjoint operator.

\begin{theorem}[Soft variance and concentration of a test-column norm]
\label{thm:v60-soft}
As \(M\to\infty\),
\begin{equation}\label{eq:v60-soft}
 \nu_M(p_M)=\frac3{4\pi^2}\log M+O(1),\qquad
 \operatorname{Var}\{\|G_1(Z)f_M\|^2\}\le C.
\end{equation}
Consequently
\begin{equation}\label{eq:v60-soft-probability}
 \frac{\|G_1(Z)f_M\|^2}{\log M}
                    \longrightarrow\frac3{4\pi^2}
                    \quad\hbox{in }L^2\hbox{ and in probability}.
\end{equation}
Thus \(\|\beta^{-1/2}G_1(Z)\|\to0\) in probability
can hold only if \(\beta/\log M\to\infty\).
If \(\beta=o(\log M)\), that norm diverges in
probability.
\end{theorem}
\begin{proof}
Only the first component of \(w(p_M,q)\)
is nonzero, and
\[
 \frac{|w_1(p_M,q)|^2}{\lambda(p_M)}
                                      =\cos^2(q_1/2).
\]
The exact expression is therefore
\begin{equation}\label{eq:v60-soft-exact}
 \nu_M(p_M)=\frac8V\sum_{q\ne0}
     \frac{\cos^2(q_1/2)}{\lambda(q)\lambda(p_M-q)}
       \left(1-\frac{|q_1(p_M-q)|^2}{\lambda(p_M-q)}\right),
\end{equation}
with the term \(q=p_M\) omitted.
Fix a sufficiently small \(\delta>0\).
On \(3|p_M|\le|q|\le\delta\), Taylor expansion
of the sine squares and their first derivatives
gives, for the summand without \(8/V\),
\[
 \frac{1-q_1^2/|q|^2}{|q|^4}
          +O\!\left(|p_M||q|^{-5}+|q|^{-2}\right).
\]
Both error sums after multiplication by \(V^{-1}\)
are bounded independently of \(M\).
There are only a bounded number of modes with
\(|q|<3|p_M|\); their normalized contribution
is bounded by the previous convolution estimate
term by term. The part \(|q|>\delta\) is also
bounded, since \(|p_M|\to0\).

Write \(q=2\pi n/M\) in the remaining sum.
On each coordinate-permutation-invariant
integer ball,
\[
 \sum_n\frac{1-n_1^2/|n|^2}{|n|^4}
                      =\frac34\sum_n|n|^{-4}.
\]
Unit-cell comparison with the radial integral
gives
\(\sum_{1\le|n|\le R}|n|^{-4}=2\pi^2\log R+O(1)\).
The comparison error is summable because
the gradient is \(O(|n|^{-5})\).
Combining the factors \(8\), \(3/4\),
\((2\pi)^{-4}\), and \(2\pi^2\) proves the
first assertion of \eqref{eq:v60-soft}.

For the variance assertion put \(X=G_1(Z)f_M\).
It is a centered, possibly non-circular,
complex Gaussian vector. Its Hermitian
covariance \(H=\mathbb E XX^*\) is block
diagonal in the output momentum \(q\).
Its pseudo-covariance \(P=\mathbb E XX^T\)
has a block only when \(q+q'=2p_M\);
this is the reality pairing of
\(\widehat Z(q-p_M)\).
Self-adjointness of \(G_1\) gives
\(\Phi_M(p,q)=\Phi_M(q,p)\).
Hence the trace \(s_q\) of the \(q\)-block
of \(H\) is bounded by
\[
 s_q\le\frac{C}{V\lambda(q)\lambda(q-p_M)}.
\]
For \(q=2\pi n/M\) this is at most
\(C(1+|n|)^{-4}\) outside a fixed finite
set, and bounded on that set; zero
denominators correspond to omitted modes.
Centered periodic distances to \(n\)
and \(n-e_1\) are comparable outside
that finite set. Thus \(\sum_q s_q^2\le C\).

Complex Gaussian Wick contraction gives
\[
 \operatorname{Var}\|X\|^2=\operatorname{Tr}H^2+\|P\|_{\rm HS}^2.
\]
The first term is at most \(\sum_q s_q^2\).
For the paired pseudo-covariance blocks,
Cauchy--Schwarz bounds their squared
Hilbert--Schmidt norms by \(s_qs_{2p_M-q}\).
The pairing is an involution, so their
sum is at most \(\sum_qs_q^2\).
This proves the uniform variance bound.
The expectation is \(\nu_M(p_M)\), giving
\eqref{eq:v60-soft-probability}.

Finally
\(\|\beta^{-1/2}G_1(Z)\|^2
\ge\beta^{-1}\|G_1(Z)f_M\|^2\).
The convergence just proved yields the
last two statements, using a bounded
subsequence of \(\beta/\log M\) if that
ratio does not tend to infinity.
\end{proof}

This theorem does not say that the full
gauge operator is singular when
\(\beta\) is a sufficiently large constant
times \(\log M\). It says that its odd
perturbation cannot then tend to zero
in operator norm. The square-logarithmic
sufficient estimate below is not claimed
to be necessary.

\subsection{An exact small-operator domain, without a global critical-norm ball}

For a real Coulomb field \(a=tZ\), \(t>0\),
put
\[
 r=t\max_e|Z_e|,\qquad
 \alpha=t\|G_1(Z)\|,\qquad
 H=\mathcal F_M(tZ)-I.
\]
There are fixed \(r_*>0\) and \(c_*>0\)
such that the following statement holds
whenever \(r\le r_*\) and
\(\alpha+c_*r^2\le1/2\).

\begin{proposition}[Full positive gauge branch and density remainder]
\label{prop:v60-deterministic}
On this domain,
\begin{equation}\label{eq:v60-floor}
 \tfrac12 I\preceq\mathcal F_M(stZ)\preceq\tfrac32 I,
                  \qquad 0\le s\le1 .
\end{equation}
The complete local quotient-density expression
\[
 \rho_M(a)=
       \prod_e\left(\frac{\sin|a_e|}{|a_e|}\right)^2
                                      \det\mathcal F_M(a)
\]
is positive along that segment. Its real logarithm
continued from zero satisfies the nonperturbative bound
\begin{equation}\label{eq:v60-density-remainder}
 \left|\log\rho_M(tZ)-t^2\ell_2(Z)\right|
          \le C(\alpha+r^2)t^2\sum_e|Z_e|^2 .
\end{equation}
All constants are independent of \(M\).
\end{proposition}
\begin{proof}
On a real colour vector of length \(b\),
\(\operatorname{ad}_a\) has eigenvalues
\(0,\pm2ib\). The corresponding values
of \(\sigma\) are \(1,b\cot b,b\cot b\).
For \(b\le r_*<\pi/2\),
\[
 -Cb^2I\preceq\sigma(\operatorname{ad}_a)-I\preceq0,
 \qquad
 \sigma(\operatorname{ad}_a)
      =I+\tfrac1{12}(\operatorname{ad}_a)^2+O(b^4).
\]
Since \(U\) is an isometry,
\(\|G_{\rm e}(tZ)\|\le Cr^2\).
Together with the odd part this gives
\(\|H\|\le\alpha+Cr^2\le1/2\);
the same bound holds all along the segment.
It proves \eqref{eq:v60-floor}. All Haar
factors are positive as well.

Write \(E=G_{\rm e}(tZ)\) and
\(G_2(Z)=\tfrac1{12}U^*(\operatorname{ad}_Z)^2U\).
The block-diagonal pointwise expansion,
factored through the contraction \(U\), gives
\[
\begin{aligned}
 \|E\|_1&\le Ct^2\sum_e|Z_e|^2,\\
 |\operatorname{Tr}(E-t^2G_2(Z))|
                  &\le Ct^4\sum_e|Z_e|^4
                   \le Cr^2t^2\sum_e|Z_e|^2,\\
 \|E\|_{\rm HS}^2&\le Cr^2t^2\sum_e|Z_e|^2 .
\end{aligned}
\]
Moreover, the exact Hilbert--Schmidt
bound in V46, or the summable operator
of V59, gives
\[
 t^2\|G_1(Z)\|_{\rm HS}^2
                              \le Ct^2\sum_e|Z_e|^2.
\]
The linear trace vanishes:
\(\operatorname{Tr}G_1(Z)=0\).
Since \(H=tG_1(Z)+E\), its squared trace
differs from \(t^2\operatorname{Tr}G_1(Z)^2\)
by at most
\[
 2\alpha\|E\|_1+\|E\|_{\rm HS}^2
                \le C(\alpha+r^2)t^2\sum_e|Z_e|^2.
\]
For self-adjoint \(H\) with \(\|H\|\le1/2\),
the convergent trace-log series gives
\[
 \left|\operatorname{Tr}\log(I+H)
              -\operatorname{Tr}H+\tfrac12\operatorname{Tr}H^2\right|
                         \le C\|H\|\operatorname{Tr}H^2.
\]
The preceding Hilbert--Schmidt estimates
bound the right side by the same expression.
Finally
\[
 2\log(\sin b/b)=-b^2/3+O(b^4)
\]
has a uniform remainder for \(b\le r_*\).
Its sum is at most \(Cr^2t^2\sum_e|Z_e|^2\).
Adding these estimates, and retaining the
complete \(\ell_2\) of V58, proves
\eqref{eq:v60-density-remainder}.
\end{proof}

The density in this proposition is the
actual local Faddeev--Popov pullback
expression, not a replacement Jacobian.
Positivity means that its gauge derivative
is locally invertible along the segment.
It does \emph{not} prove that the whole
set maps injectively to gauge orbits,
exclude other Gribov copies, or justify
counting this set once in a global
compact quotient integral.

\subsection{A high-probability event and a uniform full-density estimate}

Put
\[
 \mathcal L_M=1+\log M,\qquad
 \mathcal H_{M,u}=\mathcal L_M+u,\qquad
 d_M=-V^{-1}\mathbb E\ell_2(Z)
        =\left(\tfrac92-\tfrac3{2V}\right)g_M+\tfrac12j_M .
\]
The last equality is V58's complete density
coefficient, not just a Haar expectation.

\begin{lemma}[Four simultaneous free-field controls]
\label{lem:v60-event}
For \(M\ge4\), \(u\ge3\), there is an event
\(\mathcal E_{M,u}\), independent of \(\beta\),
with probability at least \(1-5e^{-u}\), on which
\begin{equation}\label{eq:v60-event}
\begin{aligned}
 \max_e|Z_e|^2&\le C\mathcal H_{M,u},\\
 \|G_1(Z)\|^2&\le C\mathcal L_M\mathcal H_{M,u},\\
 V^{-1}\|Z\|_2^2&\le C(1+u/M^2),\\
 V^{-1}|\ell_2(Z)-\mathbb E\ell_2(Z)|
                         &\le C\{\sqrt{\mathcal L_Mu}+u\}/M^2 .
\end{aligned}
\end{equation}
\end{lemma}
\begin{proof}
Every real link component has variance
bounded by a constant, by the free Green
diagonal estimate. A union bound over
the \(12V\) real components and the
scalar Gaussian tail gives the first
inequality with failure probability
at most \(e^{-u}\).
Equation \eqref{eq:v60-matrix-tail},
with threshold squared
\(2C\mathcal L_M\{\log(2d_{\rm g})+u\}\),
gives the second with the same failure
probability.

For a centered real Gaussian with covariance
\(\mathsf C=C_M\otimes I_3\),
\[
 \operatorname{Tr}\mathsf C\le CV,\quad
 \|\mathsf C\|\le CM^2,\quad
 \operatorname{Tr}\mathsf C^2\le CV\mathcal L_M .
\]
These follow from its nonzero eigenvalues
\(\lambda(k)^{-1}\), of multiplicity nine.
The Gaussian determinant moment formula
at parameter \(1/(4\|\mathsf C\|)\) yields
\[
 \mathbb P\{\|Z\|^2>2\operatorname{Tr}\mathsf C
                              +4\|\mathsf C\|u\}\le e^{-u}.
\]
This proves the third bound.

By V59, \(-2\ell_2(Z)=\langle Z,\mathcal Q_MZ\rangle\),
with \(0\preceq\mathcal Q_M\preceq CI\).
The eigenvalues of
\(K=\mathsf C^{1/2}\mathcal Q_M\mathsf C^{1/2}\)
obey
\[
 \|K\|\le CM^2,\qquad
 \operatorname{Tr}K^2\le C V\mathcal L_M.
\]
For completeness, if
\(Y=\sum_i\lambda_i(\xi_i^2-1)\),
\(\lambda_i\ge0\), the determinant moment
formula implies
\[
 \log\mathbb E e^{sY}
       \le \frac{s^2\sum_i\lambda_i^2}
                         {1-2s\max_i\lambda_i},
       \qquad 0<s<(2\max_i\lambda_i)^{-1},
\]
and
\(\log\mathbb E e^{-sY}\le s^2\sum_i\lambda_i^2\).
Chernoff optimization gives a two-sided
deviation at most
\(2\sqrt{u\sum_i\lambda_i^2}+2u\max_i\lambda_i\)
outside an event of probability \(2e^{-u}\).
Apply this to \(K\) and divide by \(2V\).
This gives the fourth inequality.
The four failure probabilities sum to
at most \(5e^{-u}\).
\end{proof}

\begin{theorem}[Full gauge density on typical free fields]
\label{thm:v60-probability}
There are universal \(K_0,C>0\) such that, if
\begin{equation}\label{eq:v60-coupling}
 \beta\ge K_0\mathcal L_M\mathcal H_{M,u},
                         \qquad M\ge4,\quad u\ge3,
\end{equation}
then on \(\mathcal E_{M,u}\) the full
operator \(\mathcal F_M(s\beta^{-1/2}Z)\)
has the floor \eqref{eq:v60-floor}
for every \(0\le s\le1\). Its complete
real logarithmic density satisfies
\begin{equation}\label{eq:v60-probability}
\begin{aligned}
 \left|\frac1V\log\rho_M(\beta^{-1/2}Z)
                                      +\frac{d_M}{\beta}\right|
 \le\frac C\beta\Bigg[
  &(1+u/M^2)
      \left\{\sqrt{\frac{\mathcal L_M\mathcal H_{M,u}}{\beta}}
                          +\frac{\mathcal H_{M,u}}{\beta}\right\}\\
  &+\frac{\sqrt{\mathcal L_Mu}+u}{M^2}\Bigg].
\end{aligned}
\end{equation}
The probability of this simultaneous conclusion
is at least \(1-5e^{-u}\) under the actual free
Coulomb Gaussian.
\end{theorem}
\begin{proof}
On \(\mathcal E_{M,u}\),
\[
 r^2\le C\mathcal H_{M,u}/\beta,\qquad
 \alpha\le C
        \sqrt{\mathcal L_M\mathcal H_{M,u}/\beta}.
\]
Choose \(K_0\) large enough to ensure
\(r\le r_*\) and \(\alpha+c_*r^2\le1/2\).
Proposition~\ref{prop:v60-deterministic}
then proves the full-operator floor.
Divide \eqref{eq:v60-density-remainder}
by \(V\), insert the third inequality
of \eqref{eq:v60-event}, and then use
its fourth inequality to replace
\(\ell_2(Z)/V\) by \(-d_M\).
This is exactly \eqref{eq:v60-probability}.
No truncation of \(\sigma\), the determinant,
or the Haar factors is used on the left.
\end{proof}

For example, take \(u=p(1+\log M)\) with
fixed \(p\ge3\). The exceptional probability
is \(O(M^{-p})\), and a sufficiently large
constant times \((1+\log M)^2\) suffices
for positivity. If
\(\beta/(1+\log M)^2\to\infty\), then on
events of probability tending to one,
\[
 \frac{\beta}{V}\log\rho_M(\beta^{-1/2}Z)
       \longrightarrow -d_\infty,\qquad
 d_\infty=\tfrac92g_\infty+\tfrac12j_\infty>0 .
\]
This follows from \eqref{eq:v60-probability}
and V58's convergence of \(g_M,j_M\).
The logarithm is only being evaluated
on the positive events just proved;
no global logarithm through determinant
zeros is presumed.

\subsection{The remaining gap is now a specific operator and measure problem}

The result improves the gauge-domain
criterion from a deterministic global
critical-norm smallness condition to a
high-probability statement using the
actual covariance. It also bounds the
untruncated local density beyond its
quadratic coefficient.

There are nevertheless two separate
unpaid steps. The soft-mode theorem
shows a necessary logarithmic scale
for a vanishing odd operator, while
the general matrix tail gives a
square-logarithmic sufficient scale.
It has not been shown that this extra
logarithm is intrinsic; improving it
requires using more of this matrix's
spatial and momentum structure.
Nor does the lower bound prove a
failure of positivity at a large
fixed logarithmic coupling.

More importantly, the event probability
is Gaussian. Multiplying by the full
Wilson and density weight and then
normalizing can amplify rare sets.
The extensive density asymptotic above
is not a bound on that normalized
change of measure. Gauge-slice
injectivity, other compact sectors,
and compatibility of local domains
under integration also remain open.
The present square-logarithmic window
is a proved intermediate regime,
not a claimed replacement for an
interacting continuum trajectory.

\subsection{Checkpoint fingerprints and exact verification}

The reviewed regular companion files retain
the following fingerprints:
\begin{center}\scriptsize
\begin{tabular}{ll}
SM--Einstein V72 & 607,372 bytes\\
\multicolumn{2}{l}{\ttfamily F44F9745D43379E1672C5550411B58E97195A4C9278592D221DF5D3F644DC1F5}\\[2pt]
NS V26 & 278,283 bytes\\
\multicolumn{2}{l}{\ttfamily 82C887F8C2C69E4F28FAD7C3DC8DE5B9099CB5A4BF431EBA1FFF3E91935978AD}\\[2pt]
Shell V16 & 623,261 bytes\\
\multicolumn{2}{l}{\ttfamily BF138E63BA826BC030D3A8B87E04FDD3B1EC712D269BB064674E96509B46E710}
\end{tabular}
\end{center}
The three supplied Downloads manuscripts
and suggestions file are preserved;
their unchanged scope audit is recorded
in V44. Attached mathematical text has
not been treated as a user instruction.
The next regular checkpoint is V65.

The certificate
\begin{center}\small
\path{scripts/verify_ym_f14_v60_gaussian_gauge.py}
\end{center}
uses exact rational complex arithmetic.
It checks the colour factor eight
versus the full trace factor 24,
all fifteen nonzero variance rows at
side two, and five selected complete
rows at side four. It verifies 1,500
pairwise Coulomb symmetry identities,
the convolution upper bound, and the
independent single-soft-direction
formula. The exact soft variances are
\[
 \nu_2(p_2)=\frac{55}{1536},\qquad
 \nu_4(p_4)=\frac{13813}{172032}.
\]
The side-two trace average is
\(3V^{-1}\sum_p\nu_2(p)=667/12288=j_2\),
agreeing with V58. The checker also
verifies the even Taylor coefficients
of \(r\cot r\) and the full logarithmic
Haar factor through degree eight,
and the normalization \(3/(4\pi^2)\)
of the analytic soft logarithm.
No finite row sample is used to
prove its asymptotic or a mass gap.

The 5,294-byte certificate has SHA--256
\begin{center}\scriptsize\ttfamily
A97317B44C5AEDD0B756417A6569DD59BB9ED344A4788C4E92ADE4AACAC75458.
\end{center}
It reuses unchanged V56 and V58
arithmetic helpers. The predecessor
V59 has 1,462,439 bytes and SHA--256
\begin{center}\scriptsize\ttfamily
816985194D834A826631106838B1A271FEBB8DEBA612A76AEF1A35D9E7422795.
\end{center}
Its complete body is retained apart
from the title version and this
terminal insertion. V60 replaces V59
as the sole active main note.
Companions, supplied manuscripts,
monographs and archives are unchanged.

\begin{firewall}
A high-probability positive local
gauge branch and an untruncated
density estimate are proved in a
specified coupling window under
the free Coulomb law. The soft
matrix variance has an identified
logarithmic coefficient. These
results do not establish a globally
injective gauge chart, an interacting
change-of-measure bound, a nontrivial
continuum Yang--Mills theory, or a
positive physical Hamiltonian mass
gap. The full objective remains
unproved and active.
\end{firewall}
\section{V61: the exact discrete div--curl structure removes the extra gauge logarithm}
\label{sec:v61}

\subsection{The upstream improvement and its precise scope}

V60 left a concrete discrepancy: a soft gauge mode forces a
logarithmic variance, but the general matrix-Gaussian estimate
paid a second logarithm to control the operator norm.
This section resolves that discrepancy for the actual
mean-zero free Coulomb field. The additional logarithm is
not needed. The reason is the divergence-free structure of
the link field, not a smaller matrix dimension or a
replacement covariance.

The new deterministic estimate expresses the odd
Faddeev--Popov form through the bounded mean oscillation
of a canonical two-form potential. An exact shifted
discrete div--curl identity reduces that estimate to
uniform Riesz-transform commutator bounds. The two-form
potential is logarithmically correlated in four dimensions,
but its normalized oscillation on each cube has a
uniform covariance-operator bound. A Gaussian quadratic
tail and a union bound then cost only one logarithm
for the \emph{squared} gauge norm.

The existing exact density, its Hilbert--Schmidt estimates,
the soft-mode lower bound, and the quadratic concentration
from V58--V60 are used without recalculation. The generic
shell paraproduct estimates in archived f8 V28--V30 are
not new results here; the acquisition is the exact
Faddeev--Popov reduction with its lattice shifts and
the resulting free-probability estimate. The source
and companion audit of V60 remains in force, with
the next regular checkpoint at V65.

All results below retain the existing \(\SU(2)\)
normalization \(\operatorname{ad}_a v=2a\times v\),
the omitted harmonic link modes, and the free
Gaussian law of \eqref{eq:v58-free}. They do not
claim a bound under the normalized interacting
Wilson law.

\subsection{A canonical two-form potential on the finite torus}

Let \(X_M=(\mathbb Z/M\mathbb Z)^4\), \(V=M^4\), and use
counting-measure inner products. Write
\[
 T_\mu f(x)=f(x+e_\mu),\qquad
 D_\mu=T_\mu-I,\qquad
 D_\mu^*=-T_\mu^{-1}D_\mu,\qquad
 \Delta=\sum_{\mu=1}^4D_\mu^*D_\mu.
\]
The same scalar \(\Delta\) acts componentwise on
links and two-forms. Inverses and fractional inverses
annihilate the zero Fourier mode. The exterior
derivative from links to the six plaquette orientations is
\[
 (da)_{\mu\nu}=D_\mu a_\nu-D_\nu a_\mu,\qquad \mu<\nu.
\]
Its adjoint is with respect to the sum over those six
orientations, not a doubled sum over ordered pairs.

\begin{lemma}[Exact potential and its free covariance]
\label{lem:v61-potential}
If \(a\) is a mean-zero Coulomb link field, set
\begin{equation}\label{eq:v61-potential}
 B=d\Delta^{-1}a.
\end{equation}
Then \(d^*B=a\). For \(a=Z\), the covariance
of the eighteen real components of \(B\) has nonzero
Fourier symbol
\begin{equation}\label{eq:v61-B-covariance}
 \mathsf K_B(k)=
       \lambda(k)^{-2}\Pi_2(k)\otimes I_3,\qquad
 \Pi_2(k)=\lambda(k)^{-1}d(k)d(k)^*,
\end{equation}
where \(\Pi_2(k)\) is an orthogonal rank-three
projection on the six two-form orientations.
In particular,
\[
 0\preceq\mathsf K_B(k)\preceq\lambda(k)^{-2}I_{18},
 \qquad \operatorname{Tr}\mathsf K_B(k)=9\lambda(k)^{-2}.
\]
The zero Fourier covariance is zero.
\end{lemma}
\begin{proof}
The commuting differences give the link Hodge identity
\[
 d^*d+DD^*=\Delta I_4.
\]
For example, in the \(\nu\) component the first term
is \(\Delta a_\nu-D_\nu\sum_\mu D_\mu^*a_\mu\).
Thus \(D^*a=0\) and the absence of the constant
link modes imply \(d^*d\Delta^{-1}a=a\).

At nonzero momentum put
\(q_\mu=e^{ik_\mu}-1\), \(\lambda=\sum_\mu|q_\mu|^2\),
and \(\Pi=I_4-qq^*/\lambda\). Then
\[
 d^*d=\lambda\Pi,\qquad d\Pi=d,\qquad
 (dd^*)^2=\lambda dd^*.
\]
The rank of \(d\) is three, so \(dd^*/\lambda\)
is the asserted projection. The link covariance
is \(\Pi/\lambda\) per colour. Applying
\(d/\lambda\) to each side gives
\[
 (d/\lambda)(\Pi/\lambda)(d^*/\lambda)
                         =dd^*/\lambda^3=\Pi_2/\lambda^2.
\]
This proves all claims.
\end{proof}

The harmonic exclusion is essential: a nonzero constant
link field cannot be recovered as \(d^*B\) on this torus.
Adding a constant two-form to \(B\), however, changes
neither \(d^*B\) nor the seminorm used next.

Let \(\mathcal C_M\) be the periodic axis-parallel cubes
\[
 Q=x_0+\{0,\ldots,R-1\}^4,\qquad x_0\in X_M,\quad 1\le R\le M.
\]
There are at most \(M^5\) such cubes. For a vector-valued
function \(B\) define
\begin{equation}\label{eq:v61-BMO}
 \|B\|_{\mathrm{BMO}_2}
 =\sup_{Q\in\mathcal C_M}
       \left\{|Q|^{-1}\sum_{x\in Q}|B(x)-B_Q|^2\right\}^{1/2},
 \qquad B_Q=|Q|^{-1}\sum_{x\in Q}B(x).
\end{equation}
The norm inside the sum includes all eighteen components.
The whole torus is among the cubes, while singleton
oscillations vanish.

\subsection{Uniform commutators on counting-measure tori}

Define scalar discrete Riesz transforms
\[
 R_\mu=D_\mu\Delta^{-1/2},\qquad
 R_\mu^*=-T_\mu^{-1}R_\mu.
\]
Translations, all \(R_\mu\), and their adjoints commute.
Their \(L^2\) norms are at most one.

\begin{lemma}[The commutator input, uniformly in the side length]
\label{lem:v61-commutator}
For a scalar real function \(b\), multiplication \(M_b\),
and any
\[
 A\in\{R_\mu^*,\,R_\mu^*T_\nu^{-1}:1\le\mu,\nu\le4\},
\]
there is a constant independent of \(M\ge4\) such that
\begin{equation}\label{eq:v61-commutator}
 \|[A,M_b]\|_{2\to2}\le C\|b\|_{\mathrm{BMO}_2}.
\end{equation}
The scalar version of \eqref{eq:v61-BMO} is used here.
\end{lemma}
\begin{proof}
Equip \(X_M\) with the periodic maximum metric \(d_M\)
and counting measure. Its doubling constant is bounded
independently of \(M\). For distinct points at distance
\(r\), the cardinality of a metric ball of radius \(r\)
is comparable to \(r^4\), up to universal constants
and the saturation at \(V\).

We first verify uniform Calder\'on--Zygmund bounds.
If \(K_{A,M}\) denotes the convolution kernel, then
\begin{equation}\label{eq:v61-CZ-kernel}
 \begin{aligned}
 |K_{A,M}(x)|&\le C(1+d_M(x,0))^{-4},\\
 |D_\eta K_{A,M}(x)|&\le C(1+d_M(x,0))^{-5}.
 \end{aligned}
\end{equation}
Here and below changing a fixed one-step shift
only changes \(C\).
For detail, the symbol of \(R_\mu\) is
\(q_\mu/\sqrt\lambda\), smooth away from the origin
on the Brillouin torus, with derivatives of order
\(j\) bounded by \(C_j|k|^{-j}\) near the origin.
Cut it off in a ball of radius less than the
smallest nonzero lattice momentum, and set its
value at zero to zero. This leaves all sampled
Fourier values unchanged. Decompose the cutoff
symbol into smooth annuli of radius comparable
to \(\ell^{-1}\), \(1\le\ell\lesssim M\).
Integration by parts on each annulus bounds its
Fourier coefficients, for any fixed \(N\), by
\[
 C_N\ell^{-4}(1+|x|/\ell)^{-N}.
\]
A discrete difference supplies the additional
symbol factor \(q_\eta\), so its coefficients
have the same bound with \(\ell^{-5}\).
The lowest annulus has the same scaled derivative
bounds, including the cutoff derivatives.

Sampling these smooth periodic symbols gives the
periodization of their Fourier coefficients.
For \(N>6\), summing the nearest image over
dyadic \(\ell\) gives \((1+d_M)^{-4}\), or
\((1+d_M)^{-5}\) for a difference. The other
images contribute respectively at most
\[
 C\sum_{\ell\lesssim M}\ell^{N-4}M^{-N}
       \le CM^{-4},\qquad
 C\sum_{\ell\lesssim M}\ell^{N-5}M^{-N}
       \le CM^{-5}.
\]
These are bounded by the claimed right sides.
Multiplication of the symbol by the fixed
translation phases and taking adjoints preserve
the bounds. This proves \eqref{eq:v61-CZ-kernel}.
Telescoping along a shortest path of length
\(h\ll r\) gives kernel smoothness
\(Ch/r^5\); the size estimate handles the
remaining comparable distances. Both kernel
variables obey the same estimate. A diagonal
kernel term commutes with \(M_b\).

We now apply the unweighted \(p=2\) case of
Anderson--Dami\'an, Corollary 3.8, equation (3.7),
in \emph{Calder\'on--Zygmund operators and commutators
in spaces of homogeneous type: weighted inequalities},
\href{https://arxiv.org/pdf/1401.2061}{arXiv:1401.2061}.
Its homogeneous-type definitions allow finite
counting spaces; no nonatomic or infinite-measure
hypothesis is imposed. The commutator constant
depends on the doubling, kernel and \(L^2\)
bounds, all uniform here.

To compare the BMO conventions, each metric ball,
and each bounded-shape dyadic cube used in that
theorem, lies in a member \(Q\) of \(\mathcal C_M\)
with comparable cardinality. At macroscopic
scales one may use \(Q=X_M\). If \(E\subset Q\)
and \(|Q|\le C|E|\), then
\[
 |E|^{-1}\sum_E|b-b_E|
 \le 2|E|^{-1}\sum_E|b-b_Q|
 \le 2\sqrt{|Q|/|E|}
          \left\{|Q|^{-1}\sum_Q|b-b_Q|^2\right\}^{1/2}.
\]
Thus the BMO seminorm in the cited commutator
bound is controlled by \eqref{eq:v61-BMO}.
Singleton bottom cubes contribute zero, and
the top cube is the whole torus.
The theorem therefore yields
\eqref{eq:v61-commutator} with no volume factor.
\end{proof}

\subsection{The exact gauge bilinear form is a shifted div--curl form}

\begin{theorem}[Deterministic odd-gauge control by a two-form oscillation]
\label{thm:v61-divcurl}
For every real mean-zero Coulomb field \(a\), with
\(B=d\Delta^{-1}a\),
\begin{equation}\label{eq:v61-gauge-BMO}
 \|G_1(a)\|_{2\to2}\le C\|B\|_{\mathrm{BMO}_2}.
\end{equation}
Here \(G_1\) is the actual odd Faddeev--Popov operator
of \eqref{eq:v60-gauge}.
\end{theorem}
\begin{proof}
It suffices to bound real bilinear forms; real
and complex operator norms agree. For mean-zero
vertex colour fields \(f,g\), put
\(\phi=\Delta^{-1/2}f\), \(\psi=\Delta^{-1/2}g\).
The endpoint sum is \(S_\mu=I+T_\mu\).
The factor \(1/2\) in \(G_1\) cancels the factor
two in \(\operatorname{ad}_a\), giving exactly
\[
 \langle f,G_1(a)g\rangle
    =\sum_{x,\mu}a_\mu(x)\cdot
                    \{S_\mu\psi(x)\times D_\mu\phi(x)\}.
\]
Define the link current
\[
 J_\mu(x)=\psi(x)\times\phi(x+e_\mu)
                  -\psi(x+e_\mu)\times\phi(x).
\]
A direct expansion gives
\[
 S_\mu\psi\times D_\mu\phi
                         =J_\mu+D_\mu(\psi\times\phi).
\]
The last term pairs to zero with \(a\), since
\(D^*a=0\). Consequently
\begin{equation}\label{eq:v61-current}
 \langle f,G_1(a)g\rangle
       =\langle a,J\rangle=\langle B,dJ\rangle.
\end{equation}

The required discrete curl identity, with every
translation retained, is
\begin{equation}\label{eq:v61-discrete-curl}
 \begin{aligned}
 (dJ)_{\mu\nu}
   ={}&D_\mu\psi\times T_\mu D_\nu\phi
           +T_\nu D_\mu\psi\times D_\nu\phi\\
     &-D_\nu\psi\times T_\nu D_\mu\phi
           -T_\mu D_\nu\psi\times D_\mu\phi .
 \end{aligned}
\end{equation}
For verification, use
\(D_\mu(uv)=D_\mu u\,T_\mu v+uD_\mu v\)
in \(D_\mu J_\nu-D_\nu J_\mu\); the mixed
second differences cancel because the translations
commute. Expanding the remaining first differences
gives the four displayed terms. No continuum
Leibniz rule is substituted.

Fix one scalar component \(b\) of \(B_{\mu\nu}\)
and one scalar colour pairing in
\eqref{eq:v61-discrete-curl}. Its bilinear form
is \(\langle g,K_b f\rangle\), where
\[
 \begin{aligned}
 K_b={}&R_\mu^*M_bT_\mu R_\nu
          +R_\mu^*T_\nu^{-1}M_bR_\nu\\
       &-R_\nu^*M_bT_\nu R_\mu
          -R_\nu^*T_\mu^{-1}M_bR_\mu .
 \end{aligned}
\]
The operator obtained with a constant multiplier
is zero:
\[
 R_\mu^*T_\mu R_\nu
   +R_\mu^*T_\nu^{-1}R_\nu
   -R_\nu^*T_\nu R_\mu
   -R_\nu^*T_\mu^{-1}R_\mu=0.
\]
Indeed, replacing \(R_\eta^*\) by
\(-T_\eta^{-1}R_\eta\) cancels the first
and third terms, and also the second and fourth.
Moving each multiplier to the far left therefore
gives the exact commutator identity
\begin{equation}\label{eq:v61-commutator-identity}
 \begin{aligned}
 K_b={}&[R_\mu^*,M_b]T_\mu R_\nu
            +[R_\mu^*T_\nu^{-1},M_b]R_\nu\\
       &-[R_\nu^*,M_b]T_\nu R_\mu
            -[R_\nu^*T_\mu^{-1},M_b]R_\mu .
 \end{aligned}
\end{equation}
By Lemma~\ref{lem:v61-commutator}, translations
are isometries and each rightmost Riesz transform
is a contraction, so
\(\|K_b\|\le C\|b\|_{\mathrm{BMO}_2}\).
Expand the cross products with the three-colour
alternating tensor and sum over the six plaquette
orientations. There are only a fixed number
of terms, each scalar BMO seminorm is bounded
by the vector seminorm, and Cauchy--Schwarz
over colours gives
\[
 |\langle f,G_1(a)g\rangle|
      \le C\|B\|_{\mathrm{BMO}_2}\|f\|_2\|g\|_2.
\]
This proves \eqref{eq:v61-gauge-BMO}.
\end{proof}

\subsection{The free two-form has a one-logarithm squared oscillation}

\begin{lemma}[Uniform covariance of a normalized cube oscillation]
\label{lem:v61-cube}
For \(B=d\Delta^{-1}Z\), a cube \(Q\) of side \(R\),
and \(m=|Q|=R^4\), let \(P_Q\) subtract the
mean of each component on \(Q\), and set
\[
 Y_Q=m^{-1/2}P_Q(B|_Q),\qquad
 K_Q=\operatorname{Cov}(Y_Q).
\]
Then, uniformly in \(M\ge4\), \(Q\), and \(1\le R\le M\),
\begin{equation}\label{eq:v61-cube-covariance}
 \|K_Q\|\le C,\qquad
 \operatorname{Tr}K_Q\le C\log(2R).
\end{equation}
\end{lemma}
\begin{proof}
The map \(A_Q=m^{-1/2}P_Q\,\operatorname{res}_Q\)
has norm at most \(m^{-1/2}\).
Split \(\mathsf K_B\) into the positive Fourier
parts \(0<|k|<R^{-1}\) and \(|k|\ge R^{-1}\),
using shortest representatives in the Brillouin
cube. Throughout that cube
\(\lambda(k)\ge c|k|^2\).
For the high part, \(\mathsf K_B(k)\preceq CR^4I\),
and hence
\[
 \|A_Q\mathsf K_{B,\mathrm{hi}}A_Q^*\|
                 \le CR^4/m=C.
\]

For a plane wave on \(Q\), the exact oscillation factor is
\[
 \theta_Q(k)
 =m^{-1}\sum_{x\in Q}
       \left|e^{ikx}-m^{-1}\sum_{y\in Q}e^{iky}\right|^2
 =1-\left|m^{-1}\sum_{x\in Q}e^{ikx}\right|^2.
\]
Pairwise differences of the phases, with the lifted
coordinates \(0,\ldots,R-1\), give
\[
 0\le\theta_Q(k)\le C\min\{1,R^2|k|^2\}.
\]
This also holds for a cube crossing a periodic seam.
With unitary Fourier normalization and
\eqref{eq:v61-B-covariance}, the trace is exactly
\begin{equation}\label{eq:v61-cube-trace}
 \operatorname{Tr}K_Q
       =\frac9V\sum_{k\ne0}\lambda(k)^{-2}\theta_Q(k).
\end{equation}
Consequently the low-frequency trace is at most
\[
 \frac{CR^2}{V}
       \sum_{0<|k|<R^{-1}}|k|^{-2}\le C.
\]
The final inequality follows by dyadic annuli:
an annulus of radius \(s\ge c/M\) has at most
\(CVs^4\) sampled momenta, and its contribution
is at most \(CR^2s^2\). Summing below \(R^{-1}\)
is bounded. If no low momenta exist the assertion
is immediate. A positive matrix has norm at most
its trace, giving the low operator bound.

The high trace is at most
\[
 \frac C V\sum_{|k|\ge R^{-1}}|k|^{-4}
                                  \le C\log(2R),
\]
by the same annular count, with a bounded contribution
per dyadic annulus. This proves both bounds.
For \(R=1\), \(P_Q=0\), so the actual covariance
vanishes; for \(R=M\) the proof includes the whole
torus without an infinite-volume extension.
\end{proof}

\begin{theorem}[The improved free gauge tail]
\label{thm:v61-gauge-tail}
There is a constant \(C\) such that for \(M\ge4\)
and \(u\ge0\),
\begin{equation}\label{eq:v61-BMO-tail}
 \mathbb P\left\{
       \|d\Delta^{-1}Z\|_{\mathrm{BMO}_2}^2
                              >C(1+\log M+u)\right\}\le e^{-u},
\end{equation}
and
\begin{equation}\label{eq:v61-gauge-tail}
 \mathbb P\left\{
        \|G_1(Z)\|>C\sqrt{1+\log M+u}\right\}\le e^{-u}.
\end{equation}
In particular,
\(\mathbb E\|G_1(Z)\|^2\le C(1+\log M)\).
\end{theorem}
\begin{proof}
If a finite real centered Gaussian vector \(Y\)
has covariance \(K\), diagonalization and its
Gaussian determinant moment formula give
\[
 \mathbb P\{\|Y\|^2>
                   2\operatorname{Tr}K+4\|K\|v\}\le e^{-v},
                   \qquad v\ge0.
\]
For example, for \(\|K\|>0\) use the exponential
parameter \(1/(4\|K\|)\) and
\(-\log(1-x)\le2x\) for \(0\le x\le1/2\).
The zero-covariance case is immediate.
Apply this to each \(Y_Q\) in
Lemma~\ref{lem:v61-cube}, with \(v=u+5\log M\).
Since there are at most \(M^5\) cubes and
\(\operatorname{Tr}K_Q\le C(1+\log M)\),
a union bound proves \eqref{eq:v61-BMO-tail}.
No independence between cubes is needed.
Theorem~\ref{thm:v61-divcurl} proves
\eqref{eq:v61-gauge-tail}, and integration
of that tail proves the expectation bound.
\end{proof}

\begin{proposition}[The exact threshold for a vanishing odd operator]
\label{prop:v61-sharp-scale}
Along any sequence \(M\to\infty\) and \(\beta_M>0\),
\begin{equation}\label{eq:v61-iff}
 \beta_M^{-1/2}\|G_1(Z_M)\|
                 \longrightarrow0\ \hbox{in probability}
 \quad\Longleftrightarrow\quad
 \frac{\beta_M}{\log M}\longrightarrow\infty.
\end{equation}
For large \(M\),
\[
 c\log M\le\mathbb E\|G_1(Z_M)\|^2\le C\log M.
\]
\end{proposition}
\begin{proof}
Sufficiency follows from the expectation bound
and Markov's inequality. For necessity use the
unit soft vector of Theorem~\ref{thm:v60-soft}:
\[
 \frac{\|G_1(Z_M)f_{\mathrm{soft},M}\|^2}{\log M}
              \longrightarrow\frac{3}{4\pi^2}
                  \quad\hbox{in }L^2.
\]
The operator norm dominates this vector norm.
If \(\beta_M/\log M\) does not tend to infinity,
choose a subsequence on which it is bounded
above by a finite positive constant. On that
subsequence the scaled operator norm is bounded
away from zero with probability tending to one.
This contradicts convergence in probability
to zero. The same soft-vector expectation
gives the lower expectation estimate.
\end{proof}

The coefficient \(3/(4\pi^2)\) belongs to the
specified soft-vector variance; no exact
limiting coefficient for the full operator
norm is asserted. Also, vanishing operator
norm and a fixed positive gauge floor are
different requirements. The latter can
hold at a sufficiently large \emph{fixed}
multiple of \(\log M\), as proved next.

\subsection{The full local gauge density in the logarithmic window}

Retain V60's notation
\[
 \mathcal L_M=1+\log M,\qquad
 \mathcal H_{M,u}=\mathcal L_M+u,\qquad
 d_M=-V^{-1}\mathbb E\ell_2(Z)
        =\left(\tfrac92-\tfrac3{2V}\right)g_M+\tfrac12j_M.
\]
The exact density \(\rho_M\) includes the
Faddeev--Popov determinant and all Haar factors.

\begin{theorem}[Full positivity and density with one logarithm]
\label{thm:v61-density}
There are constants \(K_1,C>0\) such that,
for \(M\ge4\), \(u\ge3\), and
\begin{equation}\label{eq:v61-beta}
                    \beta\ge K_1\mathcal H_{M,u},
\end{equation}
there is an event \(\mathcal E'_{M,u}\),
independent of \(\beta\), with
\(\mathbb P(\mathcal E'_{M,u})\ge1-5e^{-u}\),
on which
\begin{equation}\label{eq:v61-floor}
 \tfrac12 I\preceq
      \mathcal F_M(s\beta^{-1/2}Z)
                  \preceq\tfrac32 I,\qquad 0\le s\le1,
\end{equation}
and
\begin{equation}\label{eq:v61-density}
 \begin{aligned}
 \left|\frac1V\log\rho_M(\beta^{-1/2}Z)
                                     +\frac{d_M}{\beta}\right|
 \le\frac C\beta\Bigg[
   &(1+u/M^2)
       \left\{\sqrt{\frac{\mathcal H_{M,u}}{\beta}}
                          +\frac{\mathcal H_{M,u}}{\beta}\right\}\\
   &+\frac{\sqrt{\mathcal L_Mu}+u}{M^2}\Bigg].
 \end{aligned}
\end{equation}
The operator and the logarithmic density on the
left are untruncated.
\end{theorem}
\begin{proof}
Use the same first, third and fourth controls
as in Lemma~\ref{lem:v60-event}, but replace
its matrix-tail event by
\eqref{eq:v61-gauge-tail}. Their intersection
satisfies
\[
 \begin{aligned}
 \max_e|Z_e|^2&\le C\mathcal H_{M,u},&
 \|G_1(Z)\|^2&\le C\mathcal H_{M,u},\\
 V^{-1}\|Z\|_2^2&\le C(1+u/M^2),&
 V^{-1}|\ell_2(Z)-\mathbb E\ell_2(Z)|
       &\le C\{\sqrt{\mathcal L_Mu}+u\}/M^2 .
 \end{aligned}
\]
The failure probabilities are respectively
\(e^{-u},e^{-u},e^{-u},2e^{-u}\), so their
sum is at most \(5e^{-u}\).

Put \(t=\beta^{-1/2}\),
\(r=t\max_e|Z_e|\), \(\alpha=t\|G_1(Z)\|\).
On this event
\[
 r^2\le C\mathcal H_{M,u}/\beta,\qquad
 \alpha\le C\sqrt{\mathcal H_{M,u}/\beta}.
\]
Choosing \(K_1\) sufficiently large ensures
\(r\le r_*\) and \(\alpha+c_*r^2\le1/2\),
with the constants in
Proposition~\ref{prop:v60-deterministic}.
That proposition applies to the full even
function of \(\operatorname{ad}_{tZ}\)
and proves \eqref{eq:v61-floor} along the
entire segment.

The same proposition gives the deterministic
untruncated-density remainder
\[
 |\log\rho_M(tZ)-t^2\ell_2(Z)|
                      \le C(\alpha+r^2)t^2\|Z\|_2^2.
\]
Divide by \(V\), insert the bounds on
\(\alpha,r,\|Z\|_2^2/V\), and then use
the displayed concentration of \(\ell_2/V\)
to replace it by \(-d_M\). This is
\eqref{eq:v61-density}.
\end{proof}

For fixed \(p\ge3\), take \(u=p\mathcal L_M\).
Then
\[
 \beta\ge K_p(1+\log M)
 \quad\Longrightarrow\quad
 \mathbb P\{\text{\eqref{eq:v61-floor} holds}\}
                              \ge1-O(M^{-p}).
\]
This is a logarithmic, not square-logarithmic,
sufficient coupling condition under the free law.
If \(\beta/\log M\to\infty\), the error in
\eqref{eq:v61-density}, multiplied by \(\beta\),
tends to zero for this choice of \(u\). Hence
\[
 \frac{\beta}{V}\log\rho_M(\beta^{-1/2}Z)
          \longrightarrow-d_\infty,\qquad
 d_\infty=\tfrac92g_\infty+\tfrac12j_\infty>0,
\]
on positive events of probability tending to one.
As in V60, the logarithm is asserted only on
those events, not across possible determinant
zeros elsewhere.

At \(\beta=b\log M\) with a fixed sufficiently
large \(b\), the theorem instead gives a
controlled relative density error of order
\(b^{-1/2}+b^{-1}\), with constants depending
on the selected probability exponent. It
does not make that error vanish as \(M\)
increases at fixed \(b\).

\subsection{What this removes and what must be paid next}

The gauge-operator part of V60's two-logarithm
barrier is now removed. Both the upper bound
and the soft test put the squared odd norm
on the logarithmic scale. The improvement
uses exactly the Coulomb covariance and
the lattice gauge derivative already in
the programme; no independence approximation,
extra mass, or changed gauge operator has
been introduced.

The statement remains a free-field domain
estimate. A normalized change of measure
can amplify the exceptional event. In
particular, an extensive pointwise
log-density estimate is not an estimate
on the normalizing Wilson integral,
and cannot by itself be exponentiated
to obtain a uniform interacting probability.
The extensive cubic fluctuation of V58
is still present.

Likewise, positivity of the gauge derivative
proves local invertibility, not global
injectivity of a slice or absence of other
compact sectors. The omitted harmonic
variables cannot be reintroduced into
\eqref{eq:v61-potential} by declaration.
They require the separate compact
normalization and compatibility analysis.
Finally the sufficient constant \(K_p\)
has not been identified with a prescribed
physical running-coupling coefficient.

The next useful upstream target is therefore
an interacting, normalized and domain-compatible
control of these local gauge and field regulators,
including the actual Wilson weight and sector
multiplicity. Repeating the free coefficient
expansion would not supply that missing
estimate. Neither the present norm bound
nor the earlier auxiliary contractions
are a physical Hamiltonian mass-gap proof.

\subsection{Exact verification and append preservation}

The exact-arithmetic certificate
\begin{center}\small
\path{scripts/verify_ym_f14_v61_divcurl_bmo.py}
\end{center}
checks all 96 pointwise shifted curl identities
at side two and all 1,536 at side four.
It constructs \(a=d^*B\), verifies \(D^*a=0\),
and checks the equality of the original gauge
bilinear form, the current pairing, and the
two-form pairing. Their test values are
\[
 -234/77\quad(M=2),\qquad 15154/385\quad(M=4).
\]
There are twelve exact commutator-word checks
at each side length, including constant
multiplier cancellation. These finite word
checks replace \(R_\mu\) by \(D_\mu\):
they certify the shift algebra and cancellation,
while the analytic proof retains the common
commuting \(\Delta^{-1/2}\) factors and proves
the actual operator estimate.

The certificate also verifies the canonical
potential recovery, rank-three two-form
projection and covariance at all fifteen
nonzero momenta for \(M=2\), and all 255
for \(M=4\). The exact full-torus mean-square
oscillations of the eighteen-component
potential are
\[
 \mathbb E\|Y_{X_2}\|^2=865/4096,\qquad
 \mathbb E\|Y_{X_4}\|^2=20316969/80281600.
\]
Intermediate cube traces are checked from
\eqref{eq:v61-cube-trace}; singleton traces
are zero. These finite identities are
regressions, not numerical evidence used
in place of the uniform harmonic-analysis
or probability proofs.

The 8,134-byte certificate has SHA--256
\begin{center}\scriptsize\ttfamily
D98E981154BC60FDC37C5C7F253D3FD498AD2D3FA878C85F0CD40FE84859153B.
\end{center}
It uses unchanged V56 and V58 arithmetic helpers.
The predecessor V60 has 1,485,471 bytes and
SHA--256
\begin{center}\scriptsize\ttfamily
5CE9CEA377F4CEDAC910CA120B3C4F4AC5FC118F554A684EBFC453DC7676A5E5.
\end{center}
Its full body is retained apart from the title
version and this terminal insertion. V61
replaces V60 as the sole active main note.
Supplied manuscripts, companion notes,
monographs and archives are unchanged.

\begin{firewall}
This section proves a logarithmic-scale
free Coulomb gauge bound, the sharp criterion
for its odd operator to vanish, and improved
full local density control on positive
free-probability events. It does not prove
interacting domain stability, a globally
single-counted quotient construction,
a nontrivial four-dimensional continuum
Yang--Mills theory, or its positive physical
Hamiltonian mass gap. The full objective
remains unproved and active.
\end{firewall}
\section{V62: exact geometric log-concavity and normalized gauge-domain concentration}
\label{sec:v62}

\subsection{Which part of the change of measure is now resolved}

V61 proves the sharp logarithmic scale for the free
odd gauge operator, but a free-probability estimate
does not automatically survive normalization after
reweighting. We now treat the complete geometric
factor \(\rho_M\), including its actual determinant
and every Haar factor, by a different method.
The even gauge matrix is matrix-concave throughout
the principal logarithm ball. Consequently its
positive domain is convex, and the negative
logarithm of the complete geometric density is
strongly convex there.

This yields concentration under a \emph{normalized,
non-Gaussian} reference law containing the full
geometric weight. Its gauge floor is an output,
not a floor imposed in the definition of that law.
The nonlinear Wilson remainder is not absorbed
by a declaration of convexity: its exact remaining
Radon--Nikodym factor is displayed at the end.
Thus one actual component of the measure-transfer
problem is resolved, not the full interacting
problem or the mass gap.

The local convex extensions in V38--V42 and
archived f12 V41 concern different, restricted
conditional potentials. They do not prove the
global-in-dimension geometric statement below.
We reuse V45--V46's complete quotient density,
V59's quadratic operator, and V61's div--curl
estimate; none is recalculated. This is not
the companion checkpoint, which remains V65.

\subsection{The full even gauge matrix is matrix-concave}

Continue with the \(\SU(2)\) convention
\(\operatorname{ad}_a v=2a\times v\). For a single
link, put \(r=|a|\). The exact even matrix in
\eqref{eq:v60-gauge} is
\[
 \sigma(\operatorname{ad}_a)
 =r\cot r\,I_3+(1-r\cot r)\frac{aa^T}{r^2},
\]
with its continuous value \(I_3\) at \(a=0\).
We use the exact product and partial-fraction identities
\begin{equation}\label{eq:v62-cot}
 \frac{\sin r}{r}=\prod_{n\ge1}\left(1-\frac{r^2}{n^2\pi^2}\right),
 \qquad
 r\cot r=1-2\sum_{n\ge1}\frac{r^2}{n^2\pi^2-r^2}.
\end{equation}
These are equations 4.22.1 and 4.22.3 of the
\href{https://dlmf.nist.gov/4.22}{NIST Digital Library
of Mathematical Functions}. They and their
derivatives converge uniformly on compact
subintervals of \(|r|<\pi\).

\begin{lemma}[A positive Hessian behind each pole]
\label{lem:v62-perspective}
For \(c>0\), fixed \(v\in\mathbb R^3\), and \(|a|<c\),
the function
\[
                 q_{c,v}(a)=\frac{|a\times v|^2}{c^2-|a|^2}
\]
is convex. Consequently
\(a\mapsto\sigma(\operatorname{ad}_a)\)
is matrix-concave on \(\{|a|<\pi\}\).
The negative logarithmic Haar factor
\[
                   H(a)=-2\log(\sin|a|/|a|)
\]
satisfies \(\Hess H(a)\succeq(2/3)I_3\) there.
\end{lemma}
\begin{proof}
Write \(w=a\times v\), \(d=c^2-|a|^2>0\).
The variational identity
\[
 q_{c,v}(a)=\sup_{z\in\mathbb R^3}
       \{2z\cdot(a\times v)-|z|^2(c^2-|a|^2)\}
\]
expresses \(q_{c,v}\) as a supremum of convex
quadratics in \(a\). Explicitly, if \(b=h\times v\),
direct differentiation gives
\begin{equation}\label{eq:v62-perspective-Hessian}
 D^2q_{c,v}(a)[h,h]
 =\frac2d\left|b+\frac{2(a\cdot h)}d w\right|^2
                       +\frac{2|h|^2|w|^2}{d^2}\ge0.
\end{equation}
By \eqref{eq:v62-cot},
\[
 v^T\sigma(\operatorname{ad}_a)v
       =|v|^2-2\sum_{n\ge1}q_{n\pi,v}(a).
\]
Each summand is convex. Locally uniform
convergence preserves the resulting concavity
for every \(v\), which is matrix concavity.

The product in \eqref{eq:v62-cot} gives
\[
 H(a)=-2\sum_{n\ge1}\log(1-|a|^2/(n^2\pi^2)).
\]
For the summand with \(c=n\pi\), its directional
Hessian is
\[
 \frac{4|h|^2}{c^2-|a|^2}
       +\frac{8(a\cdot h)^2}{(c^2-|a|^2)^2}
                         \ge \frac{4|h|^2}{c^2}.
\]
Summing and using \(\sum_{n\ge1}(n\pi)^{-2}=1/6\)
proves the Haar lower bound.
\end{proof}

This argument applies to the full meromorphic
function before its first pole, not just to
its quadratic Taylor term.

\subsection{A convex positive domain and the complete geometric potential}

Let \(\mathcal H_M\) be the real mean-zero Coulomb
link space of V58, with dimension \(9(V-1)\).
Fix a radius \(0<r_0<\pi\), independent of \(M\)
and \(\beta\), and define
\begin{equation}\label{eq:v62-domain}
 \mathcal D_{M,r_0}
   =\{a\in\mathcal H_M:\ \max_e|a_e|<r_0,\
                                      \mathcal F_M(a)\succ0\}.
\end{equation}
Only strict positivity, not a quantitative
spectral floor, is imposed here. The primed
vertex gauge operator is that of
\eqref{eq:v60-gauge}.

\begin{theorem}[Exact log-concavity of the geometric density]
\label{thm:v62-convexity}
The set \(\mathcal D_{M,r_0}\) is convex, contains
a neighbourhood of zero in \(\mathcal H_M\),
and is invariant under simultaneous colour
rotations. On it the actual normalized-at-zero
geometric factor
\begin{equation}\label{eq:v62-rho}
 \rho_M(a)=\det\mathcal F_M(a)
                   \prod_e\left(\frac{\sin|a_e|}{|a_e|}\right)^2
\end{equation}
has potential \(W_M=-\log\rho_M\) satisfying
\begin{equation}\label{eq:v62-Hessian}
 W_M(0)=0,\qquad DW_M(0)=0,\qquad
       D^2W_M(a)[h,h]\ge\tfrac23\|h\|_2^2.
\end{equation}
In particular
\begin{equation}\label{eq:v62-domination}
                 0<\rho_M(a)\le e^{-\|a\|_2^2/3}\le1.
\end{equation}
Extending \(W_M\) by \(+\infty\) outside this
domain preserves convexity.
\end{theorem}
\begin{proof}
The map \(G_1(a)\) is linear and self-adjoint
on the Coulomb space. The other term in
\(\mathcal F_M(a)=I+G_1(a)+
U^*\{\sigma(\operatorname{ad}_a)-I\}U\)
is matrix-concave by Lemma~\ref{lem:v62-perspective}
and compression by the fixed isometry \(U\).
Thus \(\mathcal F_M\) is matrix-concave.
If \(a,b\) have positive operators, then for
\(0\le s\le1\),
\[
 \mathcal F_M((1-s)a+sb)
       \succeq(1-s)\mathcal F_M(a)+s\mathcal F_M(b)\succ0.
\]
Intersecting this convex positivity set with
the link balls and the linear space
\(\mathcal H_M\) proves convexity.
Since \(\mathcal F_M(0)=I\), continuity gives
a neighbourhood of zero. All factors are
equivariant under a common adjoint colour rotation.

For a line \(a+sh\) inside the domain, write
\[
 F=\mathcal F_M(a),\qquad F'=D\mathcal F_M(a)[h],\qquad
 F''=D^2\mathcal F_M(a)[h,h]\preceq0.
\]
The exact identity
\begin{equation}\label{eq:v62-logdet-Hessian}
 \frac{d^2}{ds^2}\big[-\log\det\mathcal F_M(a+sh)\big]_{s=0}
       =\operatorname{Tr}(F^{-1}F'F^{-1}F')
                            -\operatorname{Tr}(F^{-1}F'')\ge0
\end{equation}
proves convexity of the negative logarithmic
determinant. The first trace is the squared
Hilbert--Schmidt norm of \(F^{-1/2}F'F^{-1/2}\);
the second has the asserted sign.
Adding the Haar Hessians proves the lower
bound in \eqref{eq:v62-Hessian}.

The derivative at zero of the even term and
of the Haar potential vanishes, while
\(\operatorname{Tr}G_1(h)=0\) by V60.
Therefore \(DW_M(0)=0\). Integrating the
Hessian bound twice along the segment from
zero to \(a\), which remains in the domain,
gives \(W_M(a)\ge\|a\|_2^2/3\).
The extension by infinity is convex because
the domain is convex; its values on its
boundary are immaterial for the integrals below.
\end{proof}

Convexity of this coordinate domain is not
injectivity of its map to gauge orbits.
No assertion about the number of other
Gribov representatives follows from this theorem.

\subsection{A normalized non-Gaussian law with the full geometric weight}

Let \(\gamma_M\) be the free Gaussian on
\(\mathcal H_M\), with covariance
\(\mathsf C_M=C_M\otimes I_3\). Its inverse
on \(\mathcal H_M\) is \(\Delta\).
Put \(t=\beta^{-1/2}\) and define
\begin{equation}\label{eq:v62-reference}
 \begin{aligned}
 \mathcal Z^{\mathrm g}_{M,\beta}
   &=\mathbb E_{\gamma_M}
        [\mathbf1_{\mathcal D_{M,r_0}}(tZ)\rho_M(tZ)],\\
 d\Gamma_{M,\beta}(Z)
   &=\frac{\mathbf1_{\mathcal D_{M,r_0}}(tZ)\rho_M(tZ)}
                      {\mathcal Z^{\mathrm g}_{M,\beta}}\,d\gamma_M(Z).
 \end{aligned}
\end{equation}
The superscript \(\mathrm g\) means geometric.
The denominator is positive because the
domain contains a neighbourhood of zero,
and finite by \eqref{eq:v62-domination}.
This definition includes the entire positive
domain inside the fixed principal link balls,
not V61's smaller high-probability event.

\begin{theorem}[Normalized sub-Gaussian domination]
\label{thm:v62-Laplace}
Set
\begin{equation}\label{eq:v62-screened}
 m_\beta^2=\frac{2}{3\beta},\qquad
 \mathsf C_{M,\beta}
       =(\mathsf C_M^{-1}+m_\beta^2 I_{\mathcal H_M})^{-1}.
\end{equation}
For every real \(J\in\mathcal H_M\),
\begin{equation}\label{eq:v62-Laplace}
 \mathbb E_{\Gamma_{M,\beta}}Z=0,\qquad
 \mathbb E_{\Gamma_{M,\beta}}e^{\langle J,Z\rangle}
             \le \exp\{\tfrac12\langle J,\mathsf C_{M,\beta}J\rangle\}.
\end{equation}
For every linear map \(A:\mathcal H_M\to\mathbb R^n\),
write \(K=A\mathsf C_{M,\beta}A^*\). Then
\begin{equation}\label{eq:v62-quadratic-mgf}
 \mathbb E_{\Gamma_{M,\beta}}e^{s\|AZ\|^2}
        \le\det(I-2sK)^{-1/2},
       \qquad 0\le s<(2\|K\|)^{-1},
\end{equation}
and for \(u\ge0\),
\begin{equation}\label{eq:v62-quadratic-tail}
 \Gamma_{M,\beta}\{\|AZ\|^2>
                         2\operatorname{Tr}K+4\|K\|u\}\le e^{-u}.
\end{equation}
The zero-\(K\) cases have their continuous interpretation.
\end{theorem}
\begin{proof}
Simultaneous colour rotations preserve the
domain, Gaussian, determinant and Haar factors.
The mean of every link colour vector is thus
fixed by all rotations in \(\mathrm{SO}(3)\),
and is zero. This does not require the generally
unavailable symmetry \(\rho_M(a)=\rho_M(-a)\).

Let \(B_\beta=\mathsf C_{M,\beta}^{-1}\). On
the scaled domain define
\[
 \widetilde W(Z)=W_M(tZ)-\tfrac12m_\beta^2\|Z\|^2,
\]
and extend it by infinity. By
\eqref{eq:v62-Hessian}, \(\widetilde W\)
is convex. Its value and gradient at zero
vanish, so it is nonnegative. The density
in \eqref{eq:v62-reference} is proportional to
\[
          \exp\{-\tfrac12\langle Z,B_\beta Z\rangle-\widetilde W(Z)\}.
\]
For \(b\in\mathcal H_M\), consider its Gaussian
convolution
\[
 \Phi(b)=\int_{\mathcal H_M}
       \exp\{-\tfrac12\langle Z-b,B_\beta(Z-b)\rangle
                                      -\widetilde W(Z)\}\,dZ.
\]
The integrand is jointly log-concave in
\((Z,b)\). Marginalization preserves log-concavity
by Pr\'ekopa's Theorem 6 in
\emph{On logarithmic concave measures and functions},
Acta Sci. Math. 34 (1973), 335--343,
\href{https://rutcor.rutgers.edu/Prekopa/pdf/SCIENT2.pdf}
{author-hosted full text}. Its finite-dimensional
statement permits convex supports through an
extended-valued convex potential, as used here.
Thus \(\log\Phi\) is concave.

Differentiation of the Gaussian convolution gives
\(\nabla\log\Phi(0)=B_\beta\mathbb E_\Gamma Z=0\).
All derivatives are justified by Gaussian
domination; the support is also bounded at
each fixed \(M,\beta\). Concavity implies
\(\Phi(b)\le\Phi(0)\).
Completing the square now gives the exact identity
\[
 \mathbb E_\Gamma e^{\langle J,Z\rangle}
 =e^{\langle J,B_\beta^{-1}J\rangle/2}
           \frac{\Phi(B_\beta^{-1}J)}{\Phi(0)}
 \le e^{\langle J,B_\beta^{-1}J\rangle/2}.
\]
This proves \eqref{eq:v62-Laplace} with the
normalizing denominator already included.

For an independent standard Gaussian
\(\eta\in\mathbb R^n\), Tonelli's theorem and
\eqref{eq:v62-Laplace} give
\[
 \mathbb E_\Gamma e^{s\|AZ\|^2}
 =\mathbb E_\eta\mathbb E_\Gamma
          e^{\sqrt{2s}\langle A^*\eta,Z\rangle}
 \le\mathbb E_\eta e^{s\langle\eta,K\eta\rangle}
 =\det(I-2sK)^{-1/2}.
\]
The same elementary Chernoff estimate as in
V61, at \(s=1/(4\|K\|)\), proves
\eqref{eq:v62-quadratic-tail}.
No independence of the interacting coordinates
of \(\Gamma\) is assumed.
\end{proof}

In particular \(\operatorname{Cov}_\Gamma Z
\preceq\mathsf C_{M,\beta}\preceq\mathsf C_M\),
by differentiating the Laplace inequality
at zero. This is a covariance upper bound,
not an assertion that \(\Gamma\) is Gaussian.

\subsection{Cube oscillations and an interior gauge floor after normalization}

The screened comparison covariance in
\eqref{eq:v62-screened} has nonzero link symbol
\[
                  \frac{\Pi(k)}{\lambda(k)+m_\beta^2}\otimes I_3.
\]
Consequently for \(B=d\Delta^{-1}Z\) its
comparison two-form covariance has symbol
\begin{equation}\label{eq:v62-twoform}
              \frac{\Pi_2(k)}
                 {\lambda(k)(\lambda(k)+m_\beta^2)}\otimes I_3.
\end{equation}
This is the comparison matrix supplied by
the Laplace bound, not an equality for
the actual non-Gaussian covariance.

\begin{proposition}[Scale-resolved normalized oscillation bounds]
\label{prop:v62-cubes}
For a periodic cube \(Q\) of side \(R\), let
\(A_QZ=|Q|^{-1/2}P_Q((d\Delta^{-1}Z)|_Q)\)
as in V61, and \(K_{Q,\beta}=A_Q\mathsf C_{M,\beta}A_Q^*\).
For \(\beta\ge1\),
\begin{equation}\label{eq:v62-cube}
 \|K_{Q,\beta}\|\le\frac{C}{1+R^2/\beta},\qquad
 \operatorname{Tr}K_{Q,\beta}
                  \le C\log(2+\min\{R,\sqrt\beta\}).
\end{equation}
The unrefined V61 bounds
\(\|K_{Q,\beta}\|\le C\),
\(\operatorname{Tr}K_{Q,\beta}\le C\log(2R)\)
hold for every \(\beta>0\).
Simultaneously for all periodic cubes,
outside an event of \(\Gamma_{M,\beta}\)-probability
at most \(e^{-u}\), one has
\begin{equation}\label{eq:v62-cube-tail}
 \frac1{|Q|}\sum_Q|B-B_Q|^2
 \le C\left\{\log(2+\min\{R,\sqrt\beta\})
                     +\frac{u+5\log M}{1+R^2/\beta}\right\}
\end{equation}
when \(\beta\ge1\).
\end{proposition}
\begin{proof}
Use the frequency split and plane-wave
oscillation factor of Lemma~\ref{lem:v61-cube},
now with the symbol \eqref{eq:v62-twoform}.
On \(|k|\ge R^{-1}\), its norm is at most
\(CR^4/(1+m_\beta^2R^2)\).
The normalized restriction has squared
norm at most \(R^{-4}\), giving the
corresponding high-frequency operator bound.
The low-frequency trace is at most
\[
 \frac{CR^2}{V}\sum_{0<|k|<R^{-1}}
                          \frac1{|k|^2+m_\beta^2}
                       \le\frac{C}{1+m_\beta^2R^2}.
\]
To see the last estimate directly, the
annulus of radius \(s\le R^{-1}\)
contributes at most \(CR^2s^4/(s^2+m_\beta^2)\);
the sum is bounded by its largest-scale
term up to a universal factor.
The low operator norm is bounded by this trace.

For the high trace, the annular contributions
are bounded by \(Cs^2/(s^2+m_\beta^2)\)
for \(R^{-1}\lesssim s\lesssim1\).
Above \(m_\beta\) there is a bounded
contribution per annulus; below it the
sum is geometric. Adding the low part
gives \(C\log(2+\min\{R,m_\beta^{-1}\})\),
equivalent to \eqref{eq:v62-cube}.
Since \(\mathsf C_{M,\beta}\preceq\mathsf C_M\),
the unrefined estimates follow directly
from V61.
Apply \eqref{eq:v62-quadratic-tail} with
parameter \(u+5\log M\) and take the union
over at most \(M^5\) cubes to obtain
\eqref{eq:v62-cube-tail}.
\end{proof}

\begin{theorem}[The normalized geometric law lies inside a safe gauge domain]
\label{thm:v62-floor}
For \(M\ge4\), \(u\ge3\), put
\(\mathcal H_{M,u}=1+\log M+u\).
Under \(\Gamma_{M,\beta}\), with probability
at least \(1-3e^{-u}\),
\begin{equation}\label{eq:v62-event}
 \begin{aligned}
 \max_e|Z_e|^2&\le C\mathcal H_{M,u},\\
 \|G_1(Z)\|^2&\le C\mathcal H_{M,u},\\
 V^{-1}\|Z\|_2^2&\le C(1+u/M^2).
 \end{aligned}
\end{equation}
If \(\beta\ge K(r_0)\mathcal H_{M,u}\), this event
also satisfies
\begin{equation}\label{eq:v62-interior}
 \max_e|\beta^{-1/2}Z_e|\le r_0/2,\qquad
 \tfrac12 I\preceq\mathcal F_M(s\beta^{-1/2}Z)
                      \preceq\tfrac32 I\quad(0\le s\le1),
\end{equation}
and
\begin{equation}\label{eq:v62-density-remainder}
 \begin{aligned}
 \frac1V\left|\log\rho_M(\beta^{-1/2}Z)-\beta^{-1}\ell_2(Z)\right|
 \le\frac C\beta(1+u/M^2)
      \left\{\sqrt{\frac{\mathcal H_{M,u}}{\beta}}
                           +\frac{\mathcal H_{M,u}}{\beta}\right\}.
 \end{aligned}
\end{equation}
\end{theorem}
\begin{proof}
Each real link coordinate has a sub-Gaussian
Laplace bound with variance proxy at most
the bounded free diagonal. The two-sided
scalar tail and a union over \(12V\)
coordinates prove the first inequality
with failure probability at most \(e^{-u}\).
The unrefined cube bound in
Proposition~\ref{prop:v62-cubes}, together
with \eqref{eq:v62-quadratic-tail}, proves
the V61 BMO tail under \(\Gamma\).
The deterministic div--curl theorem
then proves the second inequality, with
the same failure probability.
For the third apply
\eqref{eq:v62-quadratic-tail} to \(A=I\),
using \(\operatorname{Tr}\mathsf C_{M,\beta}\le CV\)
and \(\|\mathsf C_{M,\beta}\|\le CM^2\).
The union has failure probability at most
\(3e^{-u}\).

On that event, \(r=t\max_e|Z_e|\) and
\(\alpha=t\|G_1(Z)\|\) obey the same bounds
as in V61. Increase \(K(r_0)\) until
\(r\le\min\{r_0/2,r_*\}\) and
\(\alpha+c_*r^2\le1/2\).
Proposition~\ref{prop:v60-deterministic}
then proves both \eqref{eq:v62-interior}
and \eqref{eq:v62-density-remainder}.
Its proof is pointwise and makes no
Gaussian assumption.
\end{proof}

Taking \(u=p(1+\log M)\) proves that a sufficiently
large fixed multiple of \(\log M\) suffices
for an interior floor with probability
\(1-O(M^{-p})\) under the \emph{normalized}
geometric law. Positivity alone was imposed
in \eqref{eq:v62-domain}; the distance from
the link boundary and the quantitative floor
in \eqref{eq:v62-interior} were not.

The free expectation \(-V^{-1}\mathbb E_{\gamma_M}\ell_2=d_M\)
has not been substituted for an expectation under
\(\Gamma\). Accordingly \eqref{eq:v62-density-remainder}
retains \(\ell_2(Z)\) itself. Free Gaussian
concentration around \(d_M\) cannot be copied
unchanged into this theorem.

\subsection{The geometric normalizer has the correct extensive order}

\begin{proposition}[Untruncated geometric pressure bounds]
\label{prop:v62-pressure}
For fixed \(r_0\in(0,\pi)\), there are
\(K,c,C>0\) such that for \(M\ge4\) and
\(\beta\ge K(1+\log M)\),
\begin{equation}\label{eq:v62-pressure}
             c\,\frac V\beta
                 \le-\log\mathcal Z^{\mathrm g}_{M,\beta}
                 \le C\,\frac V\beta.
\end{equation}
No asymptotic coefficient for this pressure
is asserted here.
\end{proposition}
\begin{proof}
Equation \eqref{eq:v62-domination} gives
\[
 \mathcal Z^{\mathrm g}_{M,\beta}
 \le\mathbb E_{\gamma_M}e^{-\|Z\|^2/(3\beta)}
 =\det_{\mathcal H_M}(I+m_\beta^2\mathsf C_M)^{-1/2}.
\]
There are nine nonzero free eigenvalues
\(\lambda(k)^{-1}\) per momentum, and
\(\lambda(k)\le16\). Therefore
\[
 -\log\mathcal Z^{\mathrm g}_{M,\beta}
 \ge\frac92\sum_{k\ne0}\log(1+m_\beta^2/\lambda(k))
                                      \ge cV/\beta
\]
for \(\beta\ge1\).

For the other direction use free probability
only to obtain a normalizer lower bound.
The first two free controls in V61 imply
that for sufficiently small fixed
\(\varepsilon>0\), and
\(\beta\ge K(1+\log M)\), there is an event
\(E\) with
\[
 \gamma_M(E)\ge1-\delta,\qquad
 \delta\le2e^{-c_0\beta},
\]
on which \(r=t\max_e|Z_e|\) and
\(\alpha=t\|G_1(Z)\|\) are small enough
for \(tZ\in\mathcal D_{M,r_0}\) and
the deterministic V60 density estimate.
For specificity choose the tail parameter
\(u=\varepsilon\beta-(1+\log M)\);
increase \(K\) until \(u\ge\varepsilon\beta/2\).
Then \(r^2\le C\varepsilon\),
\(\alpha\le C\sqrt\varepsilon\), and
the two failures have the displayed bound.

Since \(|\ell_2(Z)|\le C\|Z\|^2\) by V59,
the deterministic remainder implies on \(E\)
\[
                    0\le W_M(tZ)\le C\beta^{-1}\|Z\|^2.
\]
Conditional Jensen, using
\(\mathbb E_{\gamma_M}\|Z\|^2\le CV\), gives
\[
 \mathcal Z^{\mathrm g}_{M,\beta}
 \ge \gamma_M(E)
       \exp\{-\mathbb E_{\gamma_M}[W_M(tZ)\mid E]\}
 \ge(1-\delta)\exp\{-CV/(\beta(1-\delta))\}.
\]
Choose \(K\) so that \(\delta\le1/2\).
Since \(e^{-c_0\beta}\le C/\beta\),
taking negative logarithms proves the
upper bound in \eqref{eq:v62-pressure}.
\end{proof}

This illustrates why a naive probability-ratio
transfer would still be inadequate: the geometric
normalizer itself is of order \(\exp(-cV/\beta)\).
The normalized concentration theorem did not
divide by this small number. It used the
actual convexity of the full density instead.

\subsection{The exact Wilson factor that remains}

Let \(S\) be the actual fine Wilson action with
the normalization \(S_2(Z)=\frac12\langle Z,\Delta Z\rangle\).
On the same positive chart domain define the
normalized local Wilson pullback \(\mu^+_{M,\beta}\)
by the density proportional to
\(\mathbf1_{\mathcal D}(tZ)\rho_M(tZ)e^{-\beta S(tZ)}\,dZ\).
There is the exact identity
\begin{equation}\label{eq:v62-Wilson-change}
 \frac{d\mu^+_{M,\beta}}{d\Gamma_{M,\beta}}(Z)
 =\frac{e^{-\mathcal R_{M,\beta}(Z)}}
          {\mathbb E_{\Gamma_{M,\beta}}e^{-\mathcal R_{M,\beta}}},
 \qquad
 \mathcal R_{M,\beta}(Z)=\beta S(\beta^{-1/2}Z)-S_2(Z).
\end{equation}
This follows by cancelling the common geometric
weight and Gaussian action. At each fixed
finite \(M,\beta\) all these integrals exist;
the scaled link domain is bounded. It is not
yet a global, single-counted compact quotient
integral.

The Haar and Faddeev--Popov factors are no
longer unpaid parts of a \emph{reference}
probability bound: they are fully included
in \(\Gamma\), with its own normalization.
What is still missing is uniform control
of the normalized factor
\eqref{eq:v62-Wilson-change}, its compatible
local domains under integration, and the
other compact sectors and possible multiplicities.
The fixed-volume expansion
\(\mathcal R=\beta^{-1/2}S_3+\beta^{-1}S_4+\cdots\)
does not provide that control. V58's extensive
cubic variance remains a warning against
using it as a small global random error.

The parameter \(m_\beta^2=2/(3\beta)\) in
\eqref{eq:v62-screened} is a comparison
curvature of this geometric reference law.
It is not a physical gluon mass or a
Hamiltonian spectral gap. In particular,
V59's complete quadratic insertion includes
the Wilson contributions and has the
zero-momentum cancellation
\eqref{eq:v59-R-zero}. Dropping the remaining
factor and interpreting geometric screening
as the full interacting mass would contradict
the already identified normalization problem.

The next substantive estimate should therefore
address \eqref{eq:v62-Wilson-change} relative
to this exact normalized reference, rather
than reopening a free gauge-floor argument
or replacing the full density by its Taylor
coefficient. The full continuum construction
and physical mass gap remain the unchanged goal.

\subsection{Verification and preservation}

The certificate
\begin{center}\small
\path{scripts/verify_ym_f14_v62_geometric_measure.py}
\end{center}
executes exact rational checks. It verifies
120 directional perspective Hessians,
their variational optimizers, and the
corresponding Haar Hessian lower bounds.
It independently checks the negative-logdet
Hessian formula against a literal determinant
Taylor jet in a three-dimensional rational
finite-pole matrix model. This model tests
the algebra; it is explicitly not presented
as the full lattice Faddeev--Popov matrix
or as a replacement for the infinite
cotangent sum.

At all fifteen nonzero momenta of side two,
and all 255 of side four, it checks the
screened link inverse and the two-form
covariance identity \eqref{eq:v62-twoform}.
For the test value \(\beta=7\), its
two-form comparison traces per volume are
\[
 \frac{8320374657}{40164999680}\quad(M=2),\qquad
 \frac{1711225527547023}{6931193664778240}\quad(M=4).
\]
No sampled covariance is mistaken for the
actual covariance of \(\Gamma\).
The normalized probability and all-volume
conclusions are supplied by the analytic
proofs, not by these finite regressions.

The 7,097-byte certificate has SHA--256
\begin{center}\scriptsize\ttfamily
816E79C06133F9CA3CE4BE2D43D2A717C1897612A0BF39CAB0E8755A94415D66.
\end{center}
It reuses unchanged V56 and V58 arithmetic helpers.
The predecessor V61 has 1,511,002 bytes and
SHA--256
\begin{center}\scriptsize\ttfamily
42E608F813F83B0465BACED7E36A604E52C8621AF33051BD6E2306DF5AA5B58E.
\end{center}
Its complete body is preserved apart from
the title version and this terminal insertion.
V62 replaces V61 as the sole active main
note. Supplied sources, companion notes,
monographs and archives are unchanged.

\begin{firewall}
The complete geometric density is proved
log-concave on its positive principal
Coulomb domain. Its normalized reference
law has explicit sub-Gaussian and
quadratic concentration bounds, a
logarithmic-window interior gauge floor,
and an extensive geometric-pressure bound.
The nonlinear normalized Wilson factor,
global quotient counting, continuum
construction, and physical Yang--Mills
mass gap are not proved. The full goal
remains unachieved and active.
\end{firewall}
\section{V63: a positive gauge operator does not make the full Wilson density convex}
\label{sec:v63}

\subsection{The route tested after the normalized geometric estimate}

V62 incorporates the complete geometric density
into a normalized, log-concave reference law.
A tempting next step would be to prove that
the nonlinear Wilson factor preserves global
convexity, or is a small relative Hessian
perturbation throughout the same positive
gauge domain. This section tests that step
against the actual finite-lattice action.
It is false on that domain.

We construct a mean-zero, non-Abelian alternating
background with a uniformly positive
Faddeev--Popov operator. The exact Wilson
Hessian compressed to a specified Coulomb
polarization is a two-dimensional discrete
Laplacian minus an explicitly negative
zeroth-order term. The complete geometric
Hessian has a volume-uniform upper bound
on these backgrounds, so it cannot remove
the negative directions in the weak-coupling
sequence constructed below.

The result is stronger than a fixed large-field
counterexample: the link amplitude tends to
zero and the normalized gauge operator tends
to the identity, while the number of negative
directions grows. It is not a physical
instability claim. The backgrounds need not
be stationary points, and failure of convexity
does not imply failure of concentration.
An explicit reference-probability bound
at the end keeps that distinction quantitative.

The commuting checkerboard calculations of
archived f11 V94--V99 concern a different
single-colour constrained minimizing lift
and its effective determinant coefficient.
Here two noncommuting link colours are
used, and the object is the raw fine
Wilson Hessian inside V62's actual
positive Coulomb domain. V54's quaternion
commutator identity, V59's summable
quadratic-density operator, and V62's
convexity and normalized Laplace bounds
are reused, not reproved. The next
regular companion checkpoint remains V65.

\subsection{A mean-zero background with an explicit gauge floor}

Use orientations \(1,2,3,4\), an even torus
side \(M\ge4\), and the unit colour basis
\(E_1,E_2,E_3\), with
\([E_i,E_j]=2\epsilon_{ijk}E_k\).
Let
\begin{equation}\label{eq:v63-background}
 \begin{aligned}
 \varepsilon(x)&=(-1)^{x_3},\\
 a^{(\alpha)}_1(x)&=\alpha\varepsilon(x)E_1,\qquad
 a^{(\alpha)}_2(x)=\alpha\varepsilon(x)E_2,\\
 a^{(\alpha)}_3(x)&=a^{(\alpha)}_4(x)=0,
                         \qquad 0<\alpha\le\tfrac14 .
 \end{aligned}
\end{equation}
Every nonzero link component is independent
of its own coordinate, so its contribution
to \(D^*a\) vanishes separately. The
alternating sign has zero mean because
\(M\) is even. Thus this is an element
of the same mean-zero Coulomb space
\(\mathcal H_M\), with no restored
harmonic links.

\begin{proposition}[Uniform positivity for the actual gauge matrix]
\label{prop:v63-gauge}
For the background \eqref{eq:v63-background},
\begin{equation}\label{eq:v63-gauge}
 \|G_1(a^{(\alpha)})\|\le\alpha,\qquad
 (\alpha\cot\alpha-\alpha)I
       \preceq\mathcal F_M(a^{(\alpha)})
       \preceq(1+\alpha)I .
\end{equation}
In particular
\begin{equation}\label{eq:v63-gauge-floor}
 \tfrac{23}{32}I\preceq\mathcal F_M(a^{(\alpha)})
                  \preceq\tfrac54 I,\qquad
 \|\mathcal F_M(a^{(\alpha)})-I\|
                                    \le\alpha+\tfrac12\alpha^2.
\end{equation}
If \(\alpha<r_0\), the background lies
strictly inside \(\mathcal D_{M,r_0}\)
of V62.
\end{proposition}
\begin{proof}
Multiplication by \(\varepsilon\) shifts
Fourier momentum by \(\pi e_3\); write
\(k'=k+\pi e_3\) modulo the Brillouin torus.
For \(\mu=1,2\), it commutes with
\(D_\mu\) and the endpoint sum, and
\[
 D_\mu^*(I+T_\mu)
                          \ \longleftrightarrow\ -2i\sin k_\mu.
\]
It follows directly from
\eqref{eq:v60-gauge} that the odd operator
maps momentum \(k\) to \(k'\) by the
colour matrix
\begin{equation}\label{eq:v63-gauge-block}
 \mathcal B_\alpha(k)
 =\frac{-2i\alpha}{\sqrt{\lambda(k)\lambda(k')}}
          [\,\sin k_1 E_1+\sin k_2 E_2\,]_\times ,
\end{equation}
where \([v]_\times w=v\times w\).
The matrix is Hermitian and the same
in the reverse block. If a primed
denominator is zero, the numerator
vanishes and that constant-mode
coupling is omitted.

Put \(h(k)=\sum_{\mu\ne3}4\sin^2(k_\mu/2)\)
and \(b(k)=4\sin^2(k_3/2)\).
Then
\[
 \lambda(k)\lambda(k')
    =(h+b)(h+4-b)=h^2+4h+b(4-b)\ge4h,
\]
while \(\sin^2k_1+\sin^2k_2\le h\).
The norm of \([v]_\times\) is \(|v|\),
so \(\|\mathcal B_\alpha(k)\|\le\alpha\).
The two-momentum blocks are orthogonal,
proving the odd norm bound.

The even link matrices are independent
of the alternating sign. Their eigenvalues
are \(1\) and \(\alpha\cot\alpha\) on
the two nonzero orientations, and \(1\)
on the others. Since \(U^*U=I\),
\[
       \alpha\cot\alpha\,I
           \preceq U^*\sigma(\operatorname{ad}_{a^{(\alpha)}})U
           \preceq I.
\]
Adding the odd term proves
\eqref{eq:v63-gauge}. For \(0\le\alpha\le1/4\),
\[
 1\ge\alpha\cot\alpha\ge\cos\alpha
                                      \ge1-\alpha^2/2.
\]
Hence the lower endpoint is at least
\(1-\alpha-\alpha^2/2\ge23/32\),
and the deviation from the identity
is at most \(\alpha+\alpha^2/2\).
The strict link-radius condition and
the positive floor give the final claim.
\end{proof}

This is a direct Fourier computation of
the actual gauge operator, not an appeal
to typical free fields. It applies
uniformly in the torus side.

\subsection{An exact nonlinear Wilson path and its transverse compression}

Let \(f=f(x_3,x_4)\) be a real function
on the two-dimensional coordinate torus
with zero mean. Extend it independently
of \(x_1,x_2\), and define a link direction
\begin{equation}\label{eq:v63-polarization}
 H(f)_1(x)=f(x)E_1,\qquad
 H(f)_2(x)=-f(x)E_2,\qquad
 H(f)_3(x)=H(f)_4(x)=0.
\end{equation}
This direction is also mean-zero and
Coulomb. In particular, no constant
harmonic variation is being used.
Its squared norm is \(2\sum_{x\in X_M}f(x)^2\).

\begin{theorem}[Exact Wilson action and Hessian compression]
\label{thm:v63-Wilson}
For the Wilson action \(S\) of V58 and
every real \(s\),
\begin{equation}\label{eq:v63-exact-action}
 \begin{aligned}
 S(a^{(\alpha)}+sH(f))
 =\sum_{x\in X_M}\Big[
   &2\{\sin^2\alpha-\sin^2(sf(x))\}^2\\
   &+2\{1-\cos(2\alpha)\cos(sD_3f(x))\}\\
   &+2\{1-\cos(sD_4f(x))\}\Big].
 \end{aligned}
\end{equation}
Consequently
\begin{equation}\label{eq:v63-Wilson-Hessian}
 \begin{aligned}
 D^2S(a^{(\alpha)})[H(f),H(f)]
   =2\sum_{x\in X_M}\big[
       \cos(2\alpha)|D_3f(x)|^2+|D_4f(x)|^2
                                      -4\sin^2\alpha\,|f(x)|^2\big].
 \end{aligned}
\end{equation}
Under the norm induced by \(H\), the
Hessian compression to this polarization
is therefore
\begin{equation}\label{eq:v63-compression}
             L_\alpha=\cos(2\alpha)\Delta_3+\Delta_4-4\sin^2\alpha\,I
\end{equation}
on the mean-zero two-dimensional torus.
Its Fourier eigenvalues are
\begin{equation}\label{eq:v63-eigenvalues}
 \ell_\alpha(n_3,n_4)
 =4\cos(2\alpha)\sin^2(\pi n_3/M)
       +4\sin^2(\pi n_4/M)-4\sin^2\alpha,
                    \qquad(n_3,n_4)\ne(0,0).
\end{equation}
\end{theorem}
\begin{proof}
For the \(12\) plaquette, both relevant
fields are constant along coordinates
one and two. Its group word is the
commutator of
\[
 \exp\{(\alpha\varepsilon+sf)E_1\},\qquad
 \exp\{(\alpha\varepsilon-sf)E_2\}.
\]
The exact quaternion identity
\eqref{eq:v54-commutator} gives action
\[
 2\sin^2(\alpha\varepsilon+sf)
       \sin^2(\alpha\varepsilon-sf)
                  =2(\sin^2\alpha-\sin^2(sf))^2.
\]
On the \(13\) and \(23\) plaquettes
the fields are in a single colour
within each plaquette. Their angles
are respectively
\(2\alpha\varepsilon-sD_3f\) and
\(2\alpha\varepsilon+sD_3f\).
Their actions sum to
\(2-2\cos(2\alpha)\cos(sD_3f)\).
The \(14\) and \(24\) angles are
\(-sD_4f\) and \(sD_4f\), and
their actions sum to \(2-2\cos(sD_4f)\).
The \(34\) plaquette is the identity.
These are all six orientations, proving
the exact path formula.

Differentiate twice at \(s=0\).
The three contributions are respectively
\[
 -8\sin^2\alpha f^2,\qquad
 2\cos(2\alpha)|D_3f|^2,\qquad
 2|D_4f|^2,
\]
giving
\eqref{eq:v63-Wilson-Hessian}.
Discrete summation by parts and
\(\|H(f)\|^2=2\sum_x f^2\) identify
the compressed operator. Fourier
diagonalization of its two scalar
Laplacians gives \eqref{eq:v63-eigenvalues}.
\end{proof}

For example, take
\(f(x)=\cos(2\pi x_4/M)\). The Rayleigh
quotient is exactly
\begin{equation}\label{eq:v63-soft-negative}
 \frac{D^2S(a^{(\alpha)})[H(f),H(f)]}{\|H(f)\|^2}
                =4\sin^2(\pi/M)-4\sin^2\alpha .
\end{equation}
It is negative whenever \(\pi/M<\alpha\le1/4\),
even though \eqref{eq:v63-gauge-floor} holds.
The number of negative directions in the
specified compression is exactly the
number of nonzero momentum pairs for
which \(\ell_\alpha<0\). This gives a
lower bound on the negative index of
the full Wilson Hessian. It is not
a claim that the polarization is an
invariant subspace of the full Hessian.

There is a concrete arithmetic witness:
let
\[
 M=32,\qquad
 \alpha=2\arctan(1/16),\qquad
 \sin\alpha=32/257,\quad\cos\alpha=255/257.
\]
Then \(\alpha<1/8\), the background belongs
to \(\mathcal D_{32,1/4}\), and
\[
 \mathcal F_{32}(a^{(\alpha)})\succeq\tfrac{111}{128}I.
\]
Using \(\pi^2<10\) in \eqref{eq:v63-soft-negative}
gives the strictly negative rational upper bound
\begin{equation}\label{eq:v63-rational-negative}
 4\sin^2(\pi/32)-4\sin^2\alpha
       <\frac{40}{32^2}-\frac{4096}{66049}
                   =-\frac{194043}{8454272}.
\end{equation}
This comparison concerns the Wilson action
itself. The full geometric density is
included in the next step.

\subsection{The complete geometric Hessian cannot grow with the volume}

\begin{lemma}[Uniform upper and lower geometric curvature on a fixed floor]
\label{lem:v63-geometric-Hessian}
Fix \(r_1<\pi\) and \(\kappa>0\). Whenever
\(a\in\mathcal H_M\), \(\max_e|a_e|\le r_1\),
and \(\mathcal F_M(a)\succeq\kappa I\),
the full geometric potential of V62 satisfies
\begin{equation}\label{eq:v63-geometric-Hessian}
 \tfrac23\|h\|^2
       \le D^2W_M(a)[h,h]\le C(r_1,\kappa)\|h\|^2,
                             \qquad h\in\mathcal H_M.
\end{equation}
The upper constant is independent of \(M\).
\end{lemma}
\begin{proof}
The lower bound is V62. For the upper
bound retain \(F=\mathcal F_M(a)\).
The regularity of the full single-link
matrix \(\sigma(\operatorname{ad}_a)\)
on \(|a|\le r_1<\pi\) gives
\[
 \|D\sigma_a[h]\|_{\mathrm{HS}}\le C_{r_1}|h|,
 \qquad
 \|D^2\sigma_a[h,h]\|_1\le C_{r_1}|h|^2.
\]
Matrix concavity also gives
\(D^2\sigma_a[h,h]\preceq0\).
For the block-diagonal first derivative,
compression by \(U\) implies
\[
 \|U^*\operatorname{diag}(D\sigma_{a_e}[h_e])U\|_{\mathrm{HS}}^2
       \le C_{r_1}\sum_e|h_e|^2.
\]
For clarity, bound its Hilbert--Schmidt
norm by that of
\(\operatorname{diag}(D\sigma_{a_e}[h_e])U\).
The squared norm is a trace against
\(UU^*\); the scalar link diagonal of
this projection is \((V-1)/(4V)\le1\).
Thus the bound has no hidden dimension
or volume factor. The same diagonal
trace calculation gives
\[
            -\operatorname{Tr}F''\le C_{r_1}\|h\|^2.
\]

The first derivative of \(F\) also
contains \(G_1(h)\). By V59,
\[
 \|G_1(h)\|_{\mathrm{HS}}^2
  =\left\langle h,
       \left\{\mathcal Q_M-\left(1-\frac1{3V}\right)I\right\}h
    \right\rangle\le C\|h\|^2.
\]
Combining this with the preceding
compression bound gives
\(\|F'\|_{\mathrm{HS}}^2\le C_{r_1}\|h\|^2\).
In the exact formula
\eqref{eq:v62-logdet-Hessian},
\[
 \operatorname{Tr}(F^{-1}F'F^{-1}F')
                    \le\kappa^{-2}\|F'\|_{\mathrm{HS}}^2,\qquad
 -\operatorname{Tr}(F^{-1}F'')
                    \le\kappa^{-1}(-\operatorname{Tr}F'').
\]
Both are uniformly bounded by
\(C_{r_1}(\kappa^{-2}+\kappa^{-1})\|h\|^2\).
The single-link Haar Hessians are
uniformly bounded above on the same
closed ball. Adding them proves
\eqref{eq:v63-geometric-Hessian}.
\end{proof}

No Taylor truncation of the determinant
or Haar density has been used. This
upper bound is the reason a factor of
\(\beta\) multiplying the Wilson action
cannot be compensated by the geometric
curvature in the following sequence.

\subsection{Nonconvexity of the full normalized local Wilson potential}

Normalization adds a constant to the
negative log density and does not
change its Hessian. In the rescaled
coordinates \(Z=\sqrt\beta\,a\), the
potential of \(\mu^+_{M,\beta}\) from
\eqref{eq:v62-Wilson-change} is, up
to that constant,
\begin{equation}\label{eq:v63-full-potential}
       \mathcal U_{M,\beta}(Z)
          =\beta S(\beta^{-1/2}Z)+W_M(\beta^{-1/2}Z).
\end{equation}
It is smooth in the interior of the
scaled domain. Its Hessian there is exactly
\begin{equation}\label{eq:v63-full-Hessian}
 D^2\mathcal U_{M,\beta}(Z)[h,h]
       =D^2S(a)[h,h]+\beta^{-1}D^2W_M(a)[h,h].
\end{equation}

\begin{theorem}[Vanishing link amplitudes, a nearly identity gauge matrix,
and growing negative index]
\label{thm:v63-nonconvex}
Fix \(r_0\in(0,\pi)\). There is a fixed
\(R>0\), independent of \(M,\beta\), with
the following property. For every
sequence of even \(M\to\infty\) and
\[
                \beta_M\longrightarrow\infty,\qquad
                \beta_M=o(M^2),
\]
set
\[
 \alpha_M=R/\sqrt{\beta_M},\qquad
 a_M=a^{(\alpha_M)},\qquad
 Z_M=\sqrt{\beta_M}\,a_M .
\]
For all sufficiently large \(M\),
\(a_M\in\mathcal D_{M,r_0}\),
\(\max_e|a_{M,e}|\to0\), and
\(\|\mathcal F_M(a_M)-I\|\to0\).
Nevertheless the full Hessian
\(D^2\mathcal U_{M,\beta_M}(Z_M)\)
has at least \(cR^2M^2/\beta_M\)
strictly negative directions. On
the exhibited subspace its quadratic
form is at most
\begin{equation}\label{eq:v63-negative-band}
                     -\frac{R^2}{2\beta_M}\|h\|^2.
\end{equation}
In particular the full local Wilson
density is not log-concave throughout
V62's positive domain in this regime.
\end{theorem}
\begin{proof}
Use Lemma~\ref{lem:v63-geometric-Hessian}
with \(r_1=1/4\), \(\kappa=1/2\), and
write its upper constant as \(C_{\mathrm g}\).
Choose \(R^2\ge\max\{1,C_{\mathrm g}\}\).
Eventually \(\alpha_M\le1/4\) and
\(\alpha_M<r_0/2\). Proposition~\ref{prop:v63-gauge}
then gives membership, the uniform floor,
and convergence of the gauge matrix
to the identity.

Consider real mean-zero functions
\(f(x_3,x_4)\) whose Fourier momenta obey
\[
             \lambda_3(n_3)+\lambda_4(n_4)
                                      \le R^2/(2\beta_M).
\]
Since \(\cos(2\alpha_M)\le1\) and
\(\sin^2\alpha_M\ge\alpha_M^2/2\)
for these small amplitudes, the exact
compression satisfies
\[
 D^2S(a_M)[H(f),H(f)]
       \le-\frac{3R^2}{2\beta_M}\|H(f)\|^2.
\]
Equations \eqref{eq:v63-full-Hessian}
and \eqref{eq:v63-geometric-Hessian}
add at most \(C_{\mathrm g}/\beta_M\)
times the squared norm. The chosen
\(R\) proves \eqref{eq:v63-negative-band}
on the entire subspace, not merely
on each basis vector separately.

To count its dimension, take the
integer square
\[
 |n_3|,|n_4|\le N_M,\qquad
 N_M=\left\lfloor\frac{RM}{4\pi\sqrt{\beta_M}}\right\rfloor,
 \qquad(n_3,n_4)\ne(0,0).
\]
The bound \(\sin u\le|u|\) gives
\(\lambda_3+\lambda_4\le8\pi^2N_M^2/M^2
\le R^2/(2\beta_M)\).
Since \(\beta_M\to\infty\), this square
eventually has no Nyquist identifications;
since \(\beta_M=o(M^2)\), \(N_M\to\infty\).
The conjugation-symmetric Fourier set
has real dimension \((2N_M+1)^2-1\),
which is at least \(cR^2M^2/\beta_M\).
The embedding \(f\mapsto H(f)\) is
injective. The min--max principle for
a real symmetric finite matrix gives
the same lower bound on the negative
index of the full Hessian.
\end{proof}

For the single soft choice
\(f(x)=\cos(2\pi x_4/M)\),
\(\langle H(f),\Delta H(f)\rangle
=4\sin^2(\pi/M)\|H(f)\|^2\).
The theorem consequently gives
\begin{equation}\label{eq:v63-relative-failure}
 \frac{D^2\mathcal U_{M,\beta_M}(Z_M)[H(f),H(f)]}
           {\langle H(f),\Delta H(f)\rangle}
             \le-\frac{R^2M^2}{8\pi^2\beta_M}
                         \longrightarrow-\infty.
\end{equation}
Thus there is not even a volume-uniform
absolute relative-Hessian bound in the
free Cameron--Martin metric on this
whole admissible domain. In particular,
a small relative \(C^2\) perturbation
argument cannot transfer V62 globally.
This conclusion includes
\(\beta_M=b(1+\log M)\) for any fixed
positive \(b\), including sufficiently
large constants in V62's concentration
window.

The compression also identifies the
scale of the exhibited obstruction:
momenta of order at most
\(\beta_M^{-1/2}\), or wavelengths
of order at least \(\sqrt{\beta_M}\).
Modes for which \(\ell_{\alpha_M}>0\)
have positive Wilson curvature within
this particular polarization, and
the geometric term is positive there.
This is only a statement about the
specified compression, not a general
ultraviolet convexity theorem.

\subsection{The coherent witness is rare under the proved reference law}

The preceding theorem is a pointwise
obstruction to convexity, not a
probabilistic obstruction to concentration.
For the same field \(Z_M\), both
nonzero link components have magnitude
\(R\), independent of \(\beta_M\), and
\[
       \|Z_M\|^2=2VR^2,\qquad
       \Delta Z_M=4Z_M,\qquad
       \langle Z_M,\Delta Z_M\rangle=8VR^2.
\]
Define its aligned half-space
\begin{equation}\label{eq:v63-halfspace}
 \mathcal A_M=
   \left\{Z:\ \langle\Delta Z_M,Z\rangle
                          \ge\tfrac12\langle Z_M,\Delta Z_M\rangle\right\}.
\end{equation}
It contains \(Z_M\), and also the
Euclidean ball of radius
\(\|Z_M\|/2\) centered at \(Z_M\).

\begin{proposition}[An exponentially rare reference neighbourhood]
\label{prop:v63-rarity}
For the exact normalized geometric
law \(\Gamma_{M,\beta}\) of V62,
\begin{equation}\label{eq:v63-rarity}
                         \Gamma_{M,\beta}(\mathcal A_M)\le e^{-VR^2}.
\end{equation}
The bound holds for every \(\beta>0\);
the comparison does not assert it
under the full local Wilson law.
\end{proposition}
\begin{proof}
Apply \eqref{eq:v62-Laplace} to
the source \(J=\Delta Z_M\).
Its variance proxy is at most
\[
 \langle\Delta Z_M,\mathsf C_M\Delta Z_M\rangle
                     =\langle Z_M,\Delta Z_M\rangle=8VR^2,
\]
because the mean-zero Coulomb covariance
is \(\Delta^{-1}\), and the screened
comparison is no larger. The half-space
threshold is \(4VR^2\). Optimizing the
one-sided sub-Gaussian Chernoff bound
gives
\[
 \exp\left\{-\frac{(4VR^2)^2}{2(8VR^2)}\right\}
                                      =e^{-VR^2}.
\]
For the ball inclusion, use
\(\Delta Z_M=4Z_M\) and
Cauchy--Schwarz.
\end{proof}

This does not estimate the union of
all possible nonconvex regions, nor
their probability after the normalized
Wilson reweighting. It does show
concretely why a coherent counterexample
to a global Hessian bound is not
evidence of typical instability.

\subsection{The direction of the programme after this obstruction}

The positive gauge floor, small link
radius, and exact log-concavity of the
geometric factor all coexist with
negative curvature of the complete
local Wilson potential. Therefore
none can be used as an unstated
global convexity hypothesis in the
Pr\'ekopa step of V62.
The negative directions here lie in
the mean-zero Coulomb space and are
not removed by the omitted harmonic
sector or by another primed determinant.

What remains viable is a probability-based
localization or a normalized integration
of the problematic modes, with the
actual nonlinear factor
\eqref{eq:v62-Wilson-change} retained.
The explicit \(\sqrt\beta\) wavelength
scale tells such an argument where
this particular Hessian mechanism
becomes relevant. A cutoff or convex
extension that excludes the mechanism
would need its own quantitative
Wilson-weighted error and compatibility
estimate; it cannot silently replace
the original measure.

The normalized marginal may behave
differently from its pointwise
integrand Hessian. No sign for that
marginal has been inferred here.
Likewise the zero-momentum cancellation
in V59 and the geometric comparison
curvature in V62 remain distinct from
a physical mass. The present result
eliminates a false global-convexity
route and identifies a concrete
scale-dependent interaction to control.
It neither proves nor disproves the
full continuum Yang--Mills mass gap.

\subsection{Exact verification and append preservation}

The certificate
\begin{center}\small
\path{scripts/verify_ym_f14_v63_wilson_nonconvexity.py}
\end{center}
uses exact rational quaternion and
Fourier arithmetic. At the rational
trigonometric point
\(\sin\alpha=32/257\), \(\cos\alpha=255/257\),
it constructs the actual link and
plaquette Taylor jets and verifies
all 96 plaquettes at side two and
all 1,536 at side four, separately
by orientation. The first-order
terms of the paired \(13\) and \(23\)
plaquettes cancel exactly. The total
second derivatives for its two
mean-zero test functions are
\[
 \frac{129951712}{3236401}\quad(M=2),\qquad
 \frac{1680049664}{3236401}\quad(M=4).
\]
These high-frequency finite test
functions need not themselves have
negative curvature; they verify the
exact quadratic form. The separately
certified side-32 soft inequality
\eqref{eq:v63-rational-negative}
is the explicit negative witness.

The checker verifies the Hermitian
colour blocks and gauge norm bound
for all primed nonzero-denominator
momenta: fourteen at side two and
254 at side four. The maximal
squared odd norms divided by
\(\alpha^2\) in those tests are
\(0\) and \(1/3\), respectively.
It checks all four and sixteen
two-dimensional compression modes,
while explicitly excluding the
constant mode from the actual
Coulomb witness. It also checks
the background precision
\(\langle Z_M,\Delta Z_M\rangle/V=8R^2\).

The 7,164-byte certificate has SHA--256
\begin{center}\scriptsize\ttfamily
002C5DAD78DEBF6656529509C9F215982E4B41876FF4F7224EE7FBCA2387DF2C.
\end{center}
It reuses unchanged V54 quaternion
and V56/V58 Fourier helpers.
The predecessor V62 has 1,535,523 bytes
and SHA--256
\begin{center}\scriptsize\ttfamily
ACB8859C8349205DC667EA059B8E1672AD0080AD4738345E2F9786DA65655663.
\end{center}
Its full body is retained apart
from the title version and this
terminal insertion. V63 replaces
V62 as the sole active main note.
Companion notes, supplied sources,
monographs and archives are unchanged.

\begin{firewall}
An exact mean-zero lattice family
disproves global convexity of the
full local Wilson potential on the
positive gauge domain, even for
vanishing link amplitudes and a
gauge matrix approaching the identity.
Its negative-index scale and a
reference-rarity bound are proved.
This is not a failure theorem for
concentration, continuum existence
or a physical mass gap. The normalized
Wilson transfer, global quotient
construction and full mass-gap goal
remain unproved and active.
\end{firewall}
\section{V64: local Wilson coercivity, exact high-frequency conditionals, and a boundary-retaining low marginal}
\label{sec:v64}

\subsection{A positive step after the global-convexity obstruction}

V63 rules out convexity of the full Wilson
potential throughout the positive gauge
domain. It does not rule out coercivity
in sufficiently high-frequency directions.
We now prove such coercivity for arbitrary
small link backgrounds, not just for the
special polarization used in that example.
The possible negative contribution is
controlled by a finite-range potential
quadratic in the background.

This yields an exact normalized
high-frequency conditional concentration
theorem for the \emph{full} local Wilson
pullback, including its nonlinear action,
Haar factors and gauge determinant.
The conditional mean is retained rather
than set to zero. Integrating the high
variables also gives a semiconvexity
bound for the actual low marginal, with
the moving domain indicator included.
These results do not provide a positive
curvature floor for the remaining low
variables, a cutoff-coverage estimate
under the full Wilson law, or a
physical mass gap.

The G{\aa}rding-type statements in
archived f6 V66 and V80 are conditional
transfer-form criteria. They are not
the fine Wilson Hessian estimate proved
here. Likewise V38--V42 concern root
detail boxes and measured local shells,
not this spectral disintegration of
the complete Coulomb-coordinate law.
We reuse the actual cubic vertex from
V58 and the exact convex gauge domain
from V62. The next companion checkpoint
is V65.

\subsection{The differentiated cubic retains a difference of the variation}

In this section \(M\ge4\) need not be even.
The even-side family of V63 is used
only when checking sharpness of the
quadratic loss.

Write \(D_\nu h_\mu(x)=h_\mu(x+e_\nu)-h_\mu(x)\)
and
\[
 \|\nabla h\|_2^2=\sum_{x,\mu,\nu}|D_\nu h_\mu(x)|^2.
\]
For a plaquette \(p=(x;\mu,\nu)\), let
\(v_i(h)\), \(1\le i\le4\), be its four
oriented link vectors as in V58, and put
\[
 q_p(h)=\sum_{i<j}v_i(h)\times v_j(h).
\]
Thus the exact homogeneous cubic is
\(S_3(h)=\langle dh,q(h)\rangle\).
The derivative \(Dq(h)[a]\) is a bilinear
finite-stencil expression in \(h,a\).

\begin{lemma}[A differentiated cubic identity with no derivative of the background]
\label{lem:v64-cubic}
For real link fields \(a,h\),
\begin{equation}\label{eq:v64-cubic-identity}
 D^2S_3(a)[h,h]
   =2\{\langle a,d^*q(h)\rangle
                         +\langle dh,Dq(h)[a]\rangle\}.
\end{equation}
In particular,
\begin{equation}\label{eq:v64-cubic-global}
 |D^2S_3(a)[h,h]|
          \le C\|a\|_\infty\|h\|_2\|\nabla h\|_2 .
\end{equation}
No derivative norm of \(a\) is required.
\end{lemma}
\begin{proof}
Homogeneity and polarization give
\(D^2S_3(a)[h,h]=2DS_3(h)[a]\).
Differentiating
\(S_3(h)=\langle dh,q(h)\rangle\)
and using the exact adjoint relation
\(\langle da,q\rangle=\langle a,d^*q\rangle\)
proves \eqref{eq:v64-cubic-identity}.

Each term of \(d^*q(h)\) is a
coordinate difference of a quadratic
product of shifted components of \(h\).
The discrete product rule
\[
        D_\nu(fg)=(D_\nu f)T_\nu g+fD_\nu g
\]
therefore writes it as a fixed finite
sum of products \(h\,D_\nu h_\mu\),
with their actual lattice shifts.
Pairing with \(a\) and applying
Cauchy--Schwarz and bounded stencil
overlap gives
\[
 |\langle a,d^*q(h)\rangle|
               \le C\|a\|_\infty\|h\|_2\|\nabla h\|_2.
\]
Also
\(\|Dq(h)[a]\|_2\le C\|a\|_\infty\|h\|_2\)
and \(\|dh\|_2\le\|\nabla h\|_2\).
These prove the second bound.
The shifts are finite translations,
so the constants are independent
of the torus side, including wrapped
stencils.
\end{proof}

For a constant link field \(h\),
\(dh=0\) and \(q(h)\) is constant
on each plaquette orientation, so
\(d^*q(h)=0\). This explains the
derivative of \(h\) in the bound.
It is an algebraic observation,
not a restoration of harmonic
variables to the physical fluctuation
space.

\subsection{An exact local G{\aa}rding estimate for the Wilson Hessian}

For a link \(e=(x,\mu)\), define the
nonnegative finite-range potential
\begin{equation}\label{eq:v64-local-potential}
 \mathcal V_a(e)=
     \sum_{\nu=1}^4
       \sum_{\substack{y\in X_M\\d_\infty(x,y)\le4}}|a_\nu(y)|^2 .
\end{equation}
The radius four is a fixed safe
enlargement of the plaquette stencils.
Its cardinality is uniformly bounded
by \(4\cdot9^4\), also on small
periodic tori. Set \(r_1=1/4\).

\begin{theorem}[Full, untruncated Wilson coercivity with a local quadratic loss]
\label{thm:v64-Garding}
There is a constant \(C\), independent
of \(M\), such that for
\(\|a\|_\infty\le r_1\), \(h\in\mathcal H_M\),
and \(0<\epsilon\le1/2\),
\begin{equation}\label{eq:v64-Garding}
 \left|D^2S(a)[h,h]-\langle h,\Delta h\rangle\right|
 \le \epsilon\langle h,\Delta h\rangle
            +\frac C\epsilon\sum_e\mathcal V_a(e)|h_e|^2 .
\end{equation}
In particular a universal \(C_0\ge1\)
can be fixed so that, whenever
\(\|a\|_\infty\le r\le r_1\),
\begin{equation}\label{eq:v64-global-Garding}
 \tfrac12\langle h,\Delta h\rangle-C_0r^2\|h\|^2
 \le D^2S(a)[h,h]
 \le\tfrac32\langle h,\Delta h\rangle+C_0r^2\|h\|^2 .
\end{equation}
The action on the left is the actual
Wilson action, not its quartic
polynomial approximation.
\end{theorem}
\begin{proof}
For a single plaquette, Taylor's
formula applied to its Hessian gives
\[
 D^2S_p(a)[h,h]
   =D^2S_{2,p}[h,h]+D^2S_{3,p}(a)[h,h]+R_p(a;h).
\]
More explicitly, the remainder is
the integral of
\((1-s)D^4S_p(sa)[a,a,h,h]\) for
\(0\le s\le1\), with \(a,h\)
restricted to the four plaquette
links. The full product of four
matrix exponentials is smooth on
this fixed finite-dimensional
compact argument set. Its fourth
derivatives have a uniform bound.
Consequently
\[
 |R_p(a;h)|
       \le C\left(\sum_{e\in p}|a_e|^2\right)
                    \left(\sum_{e\in p}|h_e|^2\right),
\]
and bounded overlap gives
\begin{equation}\label{eq:v64-remainder}
              \left|\sum_pR_p(a;h)\right|
                         \le C\sum_e\mathcal V_a(e)|h_e|^2 .
\end{equation}
This is an estimate of the exact
remainder, without changing the action.

We localize the proof of
Lemma~\ref{lem:v64-cubic}.
After its summation by parts and
the discrete product rule, every
absolute term is bounded by a
constant times
\[
              |a_\rho(y)|\,|h_\mu(z)|\,|D_\eta h_\nu(w)|,
\]
where the indices range over a
fixed finite set and the locations
belong to a common bounded stencil.
In particular \(d_\infty(y,z)\le4\).
Young's inequality bounds each term
by
\[
 \delta|D_\eta h_\nu(w)|^2
                    +C\delta^{-1}|a_\rho(y)|^2|h_\mu(z)|^2.
\]
Each derivative occurs only a
uniformly bounded number of times.
Choose \(\delta\) to make their
total at most \(\epsilon\|\nabla h\|_2^2\).
The remaining terms are at most
\(C\epsilon^{-1}\sum_e\mathcal V_a(e)|h_e|^2\).
Thus
\[
 |D^2S_3(a)[h,h]|
 \le\epsilon\|\nabla h\|_2^2
             +C\epsilon^{-1}\sum_e\mathcal V_a(e)|h_e|^2.
\]

On the Coulomb space, the exact
Hodge identity \(d^*d+DD^*=\Delta I\)
gives
\(\|dh\|_2^2=\|\nabla h\|_2^2=\langle h,\Delta h\rangle\).
Since \(D^2S_2[h,h]=\|dh\|_2^2\),
combining the cubic and remainder
estimates proves \eqref{eq:v64-Garding}.
Finally \(\mathcal V_a(e)\le4\cdot9^4r^2\).
Take \(\epsilon=1/2\), enlarge the
constant to \(C_0\), and obtain
\eqref{eq:v64-global-Garding}.
\end{proof}

The quadratic order \(r^2\) of the
possible negative term cannot in
general be replaced by \(o(r^2)\).
Indeed, in V63 take \(\alpha=r\to0\)
with \(Mr\to\infty\), and use its
soft variation. The Wilson Rayleigh
quotient divided by \(r^2\) tends
to \(-4\), whereas the free
Laplacian quotient divided by \(r^2\)
tends to zero. This uses the
already proved example rather
than a new large-field heuristic.

\subsection{The exact high-frequency conditional of the full Wilson law}

Fix \(0<r\le r_1\) and retain the
actual convex positive domain
\[
 \mathcal D_{M,r}
   =\{a\in\mathcal H_M:\ \max_e|a_e|<r,\
                                      \mathcal F_M(a)\succ0\}.
\]
Let \(\mu^+_{M,\beta,r}\) be the
full local Wilson pullback on this
domain. In \(Z=\sqrt\beta\,a\)
coordinates its density is proportional
to \(\mathbf1_{\sqrt\beta\mathcal D_{M,r}}
e^{-\mathcal U(Z)}\), where
\[
 \mathcal U(Z)=\beta S(\beta^{-1/2}Z)+W_M(\beta^{-1/2}Z),
       \qquad m_\beta^2=2/(3\beta).
\]
V62 and \eqref{eq:v64-global-Garding}
give the full quadratic-form bound
\begin{equation}\label{eq:v64-full-lower}
             D^2\mathcal U(Z)\succeq
                    \tfrac12\Delta+(m_\beta^2-C_0r^2)I
\end{equation}
at every interior point. A quantitative
gauge floor is not assumed here;
the geometric Hessian lower bound
holds throughout the positive domain.

Choose a spectral threshold
\begin{equation}\label{eq:v64-threshold}
                     4C_0r^2\le\tau\le16,
\end{equation}
and split the real Coulomb space
orthogonally as
\[
 \mathcal H_M=\mathcal H_<\oplus\mathcal H_\ge,\qquad
 \mathcal H_<=\mathbf1_{\{\Delta<\tau\}}\mathcal H_M,\qquad
 \mathcal H_\ge=\mathbf1_{\{\Delta\ge\tau\}}\mathcal H_M.
\]
Constants remain omitted. The theorem
is useful when this high-frequency
space is nontrivial; \(r\) can be
chosen small enough that \(\tau<16\).
Write \(Z=y+\eta\) in this decomposition.
For every \(y\) in the projection of
the open domain, its high fibre is
\[
 \mathcal D_y^\ge
   =\{\eta\in\mathcal H_\ge:\ y+\eta\in\sqrt\beta\mathcal D_{M,r}\}.
\]
This is a bounded nonempty open convex
set in its relative high space.

\begin{theorem}[Full nonlinear conditional concentration]
\label{thm:v64-conditional}
Define the actual normalized conditional
\begin{equation}\label{eq:v64-conditional}
 d\nu_y(\eta)=
   \frac{\mathbf1_{\mathcal D_y^\ge}(\eta)e^{-\mathcal U(y+\eta)}}
        {\int_{\mathcal D_y^\ge}e^{-\mathcal U(y+\zeta)}\,d\zeta}\,d\eta,
 \qquad \bar\eta_y=\int\eta\,d\nu_y .
\end{equation}
Let
\[
       A_\ge=\tfrac14\Delta|_{\mathcal H_\ge}+m_\beta^2I,\qquad
       K_\ge=A_\ge^{-1}.
\]
For every real high source \(J\),
\begin{equation}\label{eq:v64-conditional-Laplace}
       \int e^{\langle J,\eta-\bar\eta_y\rangle}\,d\nu_y
                              \le e^{\langle J,K_\ge J\rangle/2}.
\end{equation}
The constants and comparison operator
do not depend on the retained value
\(y\) or on the conditional normalizer.
For a linear map \(B\) on the high
space, set \(K_B=BK_\ge B^*\). Then
\begin{equation}\label{eq:v64-conditional-tail}
 \nu_y\left\{\|B(\eta-\bar\eta_y)\|^2>
                        2\operatorname{Tr}K_B+4\|K_B\|u\right\}\le e^{-u},
                  \qquad u\ge0 .
\end{equation}
All nonlinear Wilson and geometric
factors in \eqref{eq:v64-conditional}
are retained exactly.
\end{theorem}
\begin{proof}
For \(h\in\mathcal H_\ge\),
\(\langle h,\Delta h\rangle\ge\tau\|h\|^2\).
Using \eqref{eq:v64-threshold} in
\eqref{eq:v64-full-lower} gives
\[
       D^2_{\eta\eta}\mathcal U(y+\eta)[h,h]
          \ge\tfrac14\langle h,\Delta h\rangle
                                +m_\beta^2\|h\|^2
          =\langle h,A_\ge h\rangle.
\]
Thus the conditional potential,
minus its \(A_\ge\)-quadratic part,
is convex on the actual fibre,
and remains convex when extended
by infinity.

Translate by the true conditional
mean: \(x=\eta-\bar\eta_y\).
The potential
\(\mathcal U(y+\bar\eta_y+x)-\frac12\langle x,A_\ge x\rangle\),
with its translated support, is
convex. Apply the same Gaussian
convolution argument as in
Theorem~\ref{thm:v62-Laplace}.
Marginal log-concavity of that
convolution and its zero logarithmic
gradient at the origin follow from
\(\int x\,d\nu_y=0\).
Completing the square proves
\eqref{eq:v64-conditional-Laplace}.
The convex remainder need not
have its pointwise minimum at zero;
only this probabilistic centering
is used. The bounded support
justifies all differentiations.

The independent Gaussian-source
integration from V62 yields
\[
 \int e^{s\|Bx\|^2}\,d\nu_y
                  \le\det(I-2sK_B)^{-1/2}
 \quad\text{for }0\le s<(2\|K_B\|)^{-1}.
\]
The same determinant Chernoff bound
then proves \eqref{eq:v64-conditional-tail}.
\end{proof}

In particular the conditional covariance
is bounded above by \(K_\ge\), whose
norm is at most \((\tau/4+m_\beta^2)^{-1}\).
The conditional mean generally depends
on \(y\); fixing \(y\) need not preserve
global colour symmetry. Replacing
\(\bar\eta_y\) by zero would be an
unproved additional source estimate.

These are conditionals of the full
\(\mu^+_{M,\beta,r}\), not of the
geometric reference \(\Gamma\).
If one starts with a larger chosen
radius \(r_0>r\), applying this result
amounts to conditioning further on
the smaller link domain. The theorem
does not assert that this extra
conditioning has small Wilson probability
cost.

\subsection{An exact low marginal with its moving boundary included}

Define
\[
 Z_\ge(y)=\int_{\mathcal H_\ge}
       \mathbf1_{\sqrt\beta\mathcal D_{M,r}}(y+\eta)
                                  e^{-\mathcal U(y+\eta)}\,d\eta,
 \qquad V_{\mathrm{eff}}(y)=-\log Z_\ge(y),
\]
with \(V_{\mathrm{eff}}=+\infty\) where
the fibre has zero mass. Addition of
the full partition constant would
normalize the low density and does
not affect the following statement.
Put
\begin{equation}\label{eq:v64-low-floor}
             L_<=\tfrac12\Delta|_{\mathcal H_<}
                               +(m_\beta^2-C_0r^2)I.
\end{equation}
This operator is allowed to have
negative eigenvalues.

\begin{theorem}[Boundary-retaining semiconvexity of the complete marginal]
\label{thm:v64-marginal}
Under \eqref{eq:v64-threshold}, the
extended-valued function
\begin{equation}\label{eq:v64-marginal}
                    V_{\mathrm{eff}}(y)-\tfrac12\langle y,L_<y\rangle
\end{equation}
is convex. Thus the actual low
marginal retains the stated Hessian
lower bound in the distributional
sense on the open projection of
\(\sqrt\beta\mathcal D_{M,r}\).
No
smoothness through a moving fibre
boundary is assumed.
\end{theorem}
\begin{proof}
On the full joint space define
\[
 \widetilde{\mathcal U}(y+\eta)
       =\mathcal U(y+\eta)-\tfrac12\langle y,L_<y\rangle
\]
inside the scaled domain and extend
it by infinity outside.
By \eqref{eq:v64-full-lower},
its full joint Hessian is bounded
below by zero on the low space
and by
\(\frac12\Delta|_{\mathcal H_\ge}
 +(m_\beta^2-C_0r^2)I\succeq A_\ge\)
on the high space, as a block
quadratic-form lower bound.
Consequently this joint function
is convex. The extension is convex
because the actual domain is convex.

The function \(e^{-\widetilde{\mathcal U}(y+\eta)}\)
is therefore jointly log-concave,
including its zero values outside
the domain. Pr\'ekopa's marginal
theorem, Theorem 6 in
\href{https://rutcor.rutgers.edu/Prekopa/pdf/SCIENT2.pdf}
{On logarithmic concave measures and functions},
gives log-concavity of
\[
 \int e^{-\widetilde{\mathcal U}(y+\eta)}\,d\eta
                       =e^{\langle y,L_<y\rangle/2}Z_\ge(y).
\]
Taking its negative logarithm proves
\eqref{eq:v64-marginal}.
The integration is over the fixed
ambient high space, with the full
joint domain indicator inside the
integrand. No derivative of a moving
boundary has been dropped.
\end{proof}

For example, its explicit Jensen
form, for finite-potential \(y_0,y_1\)
and \(0<s<1\), is
\[
 \begin{aligned}
 V_{\mathrm{eff}}((1-s)y_0+sy_1)
 \le{}&(1-s)V_{\mathrm{eff}}(y_0)+sV_{\mathrm{eff}}(y_1)\\
     &-\tfrac12s(1-s)\langle y_0-y_1,L_<(y_0-y_1)\rangle .
 \end{aligned}
\]
The sign of \(L_<\) has not been
changed to make this a convexity
claim for \(V_{\mathrm{eff}}\) itself.
The same argument may be applied
to a sequence of high-coordinate
integrations of this fixed joint
law; Fubini retains the same final
marginal and domain indicator.
It is not an all-scale locality
or renormalization induction.

\subsection{What the quadratic loss says about the remaining low sector}

As another consequence of
\eqref{eq:v64-full-lower}, the number
of strictly negative eigenvalues
of the full finite Hessian is at most
\begin{equation}\label{eq:v64-index}
 9\,\#\{k\ne0:\ \lambda(k)<2C_0r^2-2m_\beta^2\}
                              \le C(1+Mr)^4.
\end{equation}
The first bound is the min--max
principle applied to the diagonal
comparison operator. The second
uses \(\lambda(k)\ge c|k|^2\)
for shortest momentum representatives
and elementary lattice-point counting.
If the displayed threshold is
nonpositive, the first count is zero.
This bounds an index, not the
Fourier support of every negative
eigenvector; low--high mixing can
still occur.

For a link radius \(r=R/\sqrt\beta\),
the sufficient high cutoff is
\(\tau\ge4C_0R^2/\beta\).
Its wavelength scale matches the
\(\sqrt\beta\) mechanism of V63 up
to fixed constants when \(R\) is
fixed. The remaining low quadratic
floor is
\[
          \tfrac12\Delta|_{\mathcal H_<}
                         +\frac{2/3-C_0R^2}{\beta}I,
\]
which can be negative. No low
gap has been manufactured.

A fixed bound \(|Z_e|<R\) at
every link is not asserted to be
typical uniformly in the volume.
Using V62's reference concentration
for a simultaneous link bound
instead gives \(R^2\) of order
\(1+\log M+u\); this supplies only
reference-law coverage, not coverage
under the fully reweighted Wilson law.
The local form
\eqref{eq:v64-Garding}, before taking
the global maximum, is retained
precisely so that a later argument
can address spatially nonuniform
fields without automatically paying
that maximum everywhere.

The substantive new transfer is
the normalized high conditional
of the actual nonlinear law and
the exact low marginal lower bound.
What remains unpaid is control
of that low marginal, its conditional
means and interactions, the Wilson
probability cost of localization,
and compatibility with compact
sector counting and physical
reflection/transfer reconstruction.
The sharp Fourier split is also
nonlocal in space; no local Markov
or finite-range RG structure has
been inferred from it.

\subsection{Exact verification and append preservation}

The certificate
\begin{center}\small
\path{scripts/verify_ym_f14_v64_wilson_coercivity.py}
\end{center}
constructs exact rational mean-zero
Coulomb fields as adjoint curls.
It checks the differentiated cubic
identity against the literal triple
determinants on all 96 plaquettes
at side two and all 1,536 at side
four. Its total cubic second
derivatives are
\[
             -289578/6859,\qquad 1005720/6859.
\]
It independently verifies the
adjoint summation by parts,
the constant-field cubic cancellation,
and equality of curl and full
gradient energies on its Coulomb
variations. The two energies are
\[
                 484108/361,\qquad 6806496/361.
\]
No numerical regression is used
to infer a G{\aa}rding constant.

A separate rational quadratic
sign test verifies the high
conditional covariance order and
the low Schur-complement lower
bound. Its low marginal still
has negative curvature
\(-421/3963\) in one direction,
confirming that the bound is
semiconvex rather than a hidden
positive-floor assertion.
This quadratic test is not
substituted for the Wilson action
or for its moving integration
domain; those are handled in
the analytic theorems.

The 6,450-byte certificate has SHA--256
\begin{center}\scriptsize\ttfamily
EF65F3707D639B3ABD3D0D8557161CA601EC478EFED005BEF9597044CB2E0C97.
\end{center}
It uses unchanged V54 and V62
arithmetic helpers and their
unchanged V56/V58 dependencies.
The predecessor V63 has 1,558,299 bytes
and SHA--256
\begin{center}\scriptsize\ttfamily
864D6FD0329E4A29917FA60CCC2408F053A7E8FC7B366576ED7E60D8B1A4C9AE.
\end{center}
Its entire body is retained apart
from the title version and this
terminal insertion. V64 replaces
V63 as the sole active main note.
Supplied sources, companions,
monographs and archives are unchanged.

\begin{firewall}
The full Wilson Hessian has a
proved local quadratic-loss bound.
The exact normalized high-frequency
conditionals are strongly log-concave,
and the complete low marginal
has a boundary-retaining semiconvexity
bound. These are statements on
the specified positive coordinate
domain, not a proof that it covers
the compact interacting measure
with controlled multiplicity or
probability. The remaining low
sector, continuum construction
and physical Yang--Mills mass
gap are unproved. The full
goal remains active.
\end{firewall}
\section{V65: homogeneous-source centering and full Wilson cutoff stability}
\label{sec:v65}

\subsection{Companion checkpoint and the particular gap now addressed}

V64 proves concentration about the \emph{actual conditional mean}
of the high modes, but does not locate that mean. It also explicitly
leaves unpaid the probability cost of imposing a smaller link
domain. This section closes those two issues on a specified family
of fibres: homogeneous retained means, with every other retained
low mode fixed to zero. The conditional includes the full nonlinear
Wilson action, Faddeev--Popov determinant and fine Haar factors.
It is not a new Gaussian or finite-order replacement.

At this V65 checkpoint the latest active companion files remain
SM--Einstein V72, NS Stability V26, and Yang--Mills Shell Mixing
V16. They are unchanged from the preceding review. Their relevant
scope distinctions are as follows.
\begin{itemize}
\item SM V72 derives its covariance transport from an auxiliary
many-copy quantum limit, at fixed spatial cutoff and fixed times.
Its point-concentrated initial data cannot be replaced by a general
distribution while retaining a closed equation for its mean.
Here the corresponding lesson is to establish conditional centering
before using a concentration estimate. The SM result is not a
finite-species continuum theorem.
\item NS V26 evaluates rank-one first and second Oseen derivative
norms. Those identities supply no bound for the present Wilson
conditional law. No NS numerical constant is imported into a
Yang--Mills inequality.
\item Shell V16 retains the complete measured flat-holonomy
one-loop coefficient, including the source-normal return and
fine/coarse factors appropriate to its root constraint. That
coefficient is not a joint nonlinear coupling--volume remainder.
We keep the complete weight of the \emph{present} mean-fixed
integral, and compare its source-normalized values. We do not
identify this integral with the different root-constrained one.
\end{itemize}
The supplied source-selection V17, native-stationarity V21 and
exact-action Einstein bridge V16 also remain unchanged. Their
previously recorded scope audits are not repeated. In particular,
stationarity or an auxiliary contraction is not relabelled as a
physical Yang--Mills gap. The next companion checkpoint is V70.

The new acquisition is an exact probability estimate at finite
coupling and a bound for normalized partition ratios. The
homogeneous-background stationarity already used in
\eqref{eq:v21-mean} is not being recomputed. What must now be
proved is the mean of a \emph{non-Gaussian conditional distribution},
followed by a controlled change of its integration domain.

\subsection{An affine mean-fixed domain with a quantitative base floor}

Work on the side-\(M\) four-torus, \(M\ge4\), with all conventions
of V58--V64. For \(c=(c_1,\ldots,c_4)\in(\mathbb R^3)^4\), set
\begin{equation}\label{eq:v65-source}
 a_c(x,\mu)=c_\mu/M,\qquad
 |c|^2=\sum_{\mu=1}^4|c_\mu|^2,\qquad c_*=\tfrac14.
\end{equation}
The colour vectors \(c_\mu\) need not commute. This is a fixed
external harmonic mean, not a Gaussian harmonic integration
variable. The fluctuations will still belong to the mean-zero
Coulomb space \(\mathcal H_M\).

\begin{lemma}[A uniform floor for small homogeneous means]
\label{lem:v65-base}
For \(|c|\le c_*\), the primed operator in
\eqref{eq:v60-gauge} satisfies
\begin{equation}\label{eq:v65-base}
 \|G_1(a_c)\|\le |c|/2,\qquad
 \|\mathcal F_M(a_c)-I\|
       \le |c|/2+|c|^2/(2M^2)\le65/512.
\end{equation}
In particular \(\mathcal F_M(a_c)\succeq(447/512)I\).
\end{lemma}
\begin{proof}
Put \([v]_\times w=v\times w\). Since
\(\operatorname{ad}_{a}=2[a]_\times\), the literal Fourier
block of the odd term at nonzero momentum is
\begin{equation}\label{eq:v65-constant-block}
 G_1(a_c)(k)
   =-\frac{2i}{M\lambda(k)}
               \left[\sum_{\mu=1}^4\sin(k_\mu)c_\mu\right]_\times .
\end{equation}
Indeed
\((e^{-ik_\mu}-1)(1+e^{ik_\mu})=-2i\sin k_\mu\).
The complex operator norm of \([v]_\times\) is \(|v|\), and
Cauchy--Schwarz gives
\[
 \left|\sum_\mu\sin(k_\mu)c_\mu\right|
       \le |c|\left(\sum_\mu\sin^2 k_\mu\right)^{1/2}
       \le |c|\sqrt{\lambda(k)}.
\]
The smallest nonzero Laplacian eigenvalue obeys
\[
 \lambda_{\min}=4\sin^2(\pi/M)\ge16/M^2.
\]
Taking the supremum of the block norms proves the odd bound.

For \(r=|a_e|\le1/4\), the eigenvalues of
\(\sigma(\operatorname{ad}_{a_e})\) are \(1,r\cot r,r\cot r\).
The inequalities
\(\tan r\ge r\) and \(r\cot r\ge\cos r\ge1-r^2/2\) show that
\begin{equation}\label{eq:v65-even}
 0\preceq I-\sigma(\operatorname{ad}_{a_e})
                                      \preceq (r^2/2)I.
\end{equation}
Compression by the fixed isometry \(U\) in
\eqref{eq:v60-gauge} bounds the even term by
\(\max_e|a_e|^2/2\). At \(a_c\) this is at most
\(|c|^2/(2M^2)\). Insert \(|c|\le1/4\), \(M\ge4\),
to obtain the stated constants.
\end{proof}

Choose fixed radii and a spectral threshold satisfying
\begin{equation}\label{eq:v65-parameters}
 0<r_-<r_+\le\tfrac14,\qquad
       4C_0r_+^2\le\tau\le1,\qquad
       c_*/M\le r_-/4,
\end{equation}
where \(C_0\) is the universal constant in
\eqref{eq:v64-global-Garding}. These choices are possible by first
choosing \(r_+\) sufficiently small. The same threshold is used for
both radii and for every source. Let
\[
 \mathcal H_\ge=\mathbf1_{\{\Delta\ge\tau\}}\mathcal H_M,
            \qquad t=\beta^{-1/2}.
\]
This high space is nonzero for \(M\ge4\) and \(\tau\le1\).
All nonzero modes of \(\mathcal H_M\) below \(\tau\) are fixed to
zero, not integrated out.

Define two subsets of this fixed Euclidean high space:
\begin{equation}\label{eq:v65-domains}
\begin{aligned}
 D_+(c)&=\{\eta\in\mathcal H_\ge:
       \max_e|a_c+t\eta_e|<r_+,\ 
                  \mathcal F_M(a_c+t\eta)\succ0\},\\
 D_-(c)&=\{\eta\in\mathcal H_\ge:
       \max_e|a_c+t\eta_e|<r_-,\
                  \mathcal F_M(a_c+t\eta)\succ\tfrac12 I\}.
\end{aligned}
\end{equation}
Here \(a_c+t\eta_e\) means \(a_c(e)+t\eta_e\).
The notation emphasizes that the two domains have the same
retained mean and the same fluctuation variables.

\begin{lemma}[Affine convexity and the actual high conditional]
\label{lem:v65-affine}
The sets \(D_-(c)\subset D_+(c)\) are bounded nonempty open convex
sets in \(\mathcal H_\ge\). With
\begin{equation}\label{eq:v65-law}
\begin{aligned}
 w_c(\eta)&=\exp[-\beta S(a_c+t\eta)]\,\rho_M(a_c+t\eta),\\
 Z_\pm(c)&=\int_{D_\pm(c)}w_c(\eta)\,d\eta,\qquad
 d\nu_{\pm,c}(\eta)
       =Z_\pm(c)^{-1}\mathbf1_{D_\pm(c)}(\eta)w_c(\eta)\,d\eta ,
\end{aligned}
\end{equation}
both normalizers are finite and strictly positive. On either
domain the Hessian of \(-\log w_c\), restricted to the high space,
is bounded below by
\begin{equation}\label{eq:v65-high-floor}
 A_\ge=\tfrac14\Delta|_{\mathcal H_\ge}
                        +\frac{2}{3\beta}I.
\end{equation}
\end{lemma}
\begin{proof}
The proofs of matrix concavity and the Haar Hessian bound in V62
use \(D^*a=0\), but do not require the spatial average of \(a\)
to vanish. They therefore apply on the affine Coulomb slice
\(a_c+\mathcal H_\ge\). In particular,
\(\mathcal F_M\) is matrix-concave there. Its strict superlevel
sets above \(0\) and above \(I/2\) are convex. Intersecting with
the convex link balls proves convexity. The base floor in
Lemma~\ref{lem:v65-base} and
\(\max_e|a_c(e)|\le r_-/4\) place \(\eta=0\) strictly inside
both sets. Boundedness follows from the link bounds.

The Wilson G{\aa}rding proof in V64 applies to arbitrary small
backgrounds \(a\); only the variation \(h\) is required to be
mean-zero and Coulomb. Combining that estimate with the affine
geometric Hessian bound gives
\[
 D_\eta^2[-\log w_c](\eta)[h,h]
 \ge \tfrac12\langle h,\Delta h\rangle
           +\left(\frac{2}{3\beta}-C_0r_+^2\right)\|h\|^2.
\]
For \(h\in\mathcal H_\ge\), the threshold condition makes this
at least \(\langle h,A_\ge h\rangle\). This uses the complete
Wilson action, not its Hessian at \(a_c\) alone.

On the bounded link closure the Haar factors and finite-dimensional
determinant are bounded above. The Wilson factor is also bounded
there. Thus the integrals are finite. The strictly positive
continuous integrand near \(\eta=0\) gives strict positivity.
\end{proof}

The factor \(\rho_M\) here is exactly
\eqref{eq:v62-rho}, with all primed vertex inverses unchanged.
The linear harmonic source constraint has a source-independent
coordinate Jacobian, which cancels from the comparisons below.
There is no nonlinear root constraint in
\eqref{eq:v65-law}; adding its return determinant or a coarse Haar
factor would change the integral. Conversely, this formula does
not compute the root-constrained shell coefficient. It is the
full local coordinate weight on the specified affine high fibre,
not an assertion of global gauge-orbit coverage or multiplicity.

\subsection{Translation symmetry determines the non-Gaussian center}

\begin{theorem}[Exact centering without colour symmetry of the source]
\label{thm:v65-centering}
For every source and parameter choice above,
\begin{equation}\label{eq:v65-centering}
             \mathbb E_{\nu_{+,c}}\eta
                   =\mathbb E_{\nu_{-,c}}\eta=0.
\end{equation}
Consequently, for either sign and every real \(J\in\mathcal H_\ge\),
\begin{equation}\label{eq:v65-Laplace}
 \mathbb E_{\nu_{\pm,c}}e^{\langle J,\eta\rangle}
       \le \exp\left(\tfrac12\langle J,K_\ge J\rangle\right),
               \qquad K_\ge=A_\ge^{-1}.
\end{equation}
Extended by zero on the other mean-zero modes,
\begin{equation}\label{eq:v65-covariance-order}
                       0\preceq K_\ge\preceq4\mathsf C_M.
\end{equation}
The operator on the right is a comparison covariance, not an
identification of the actual conditional covariance.
\end{theorem}
\begin{proof}
Every torus translation acts orthogonally on \(\mathcal H_\ge\).
It commutes with the spectral projection and leaves \(a_c\)
unchanged. It permutes plaquettes and links, preserving the Wilson
action and all fine Haar factors. It conjugates the primed gauge
operator by the corresponding vertex translation, and therefore
preserves its determinant and both gauge inequalities in
\eqref{eq:v65-domains}. Thus the domain, Lebesgue measure and
complete density of each \(\nu_{\pm,c}\) are translation invariant.

The expectation exists by bounded support and is fixed by every
translation. Hence it is spatially constant in each orientation
and colour component. But it also belongs to \(\mathcal H_\ge\),
which has no spatially constant mode. It must be zero. This
argument does not require the \(c_\mu\) to be parallel, or a colour
rotation that fixes all of them.

The Hessian bound in Lemma~\ref{lem:v65-affine} on a convex support
allows exactly the centered Laplace argument of
Theorem~\ref{thm:v64-conditional}. Its true mean is now known to
be zero, so that theorem gives \eqref{eq:v65-Laplace}. Finally,
on a high eigenmode the comparison eigenvalue is
\[
       \left(\lambda/4+2/(3\beta)\right)^{-1}\le4/\lambda.
\]
On omitted modes it is zero. This proves the operator order.
\end{proof}

No integration by parts across a moving boundary was used to
establish \eqref{eq:v65-centering}. No identity for the expectation
of a nonlinear force was closed by replacing it with the force
of the expectation. The exact symmetry acts on the probability
law itself.

\subsection{Uniform tails under the full normalized Wilson conditional}

Write \(L_M=1+\log M\).
\begin{proposition}[Absolute link and gauge tails]
\label{prop:v65-tails}
There is a universal \(C\) such that, uniformly over
\(|c|\le c_*\), either sign, and \(u\ge0\),
\begin{equation}\label{eq:v65-tails}
\begin{aligned}
 \nu_{\pm,c}\{\max_e|\eta_e|^2>C(L_M+u)\}&\le e^{-u},\\
 \nu_{\pm,c}\{\|G_1(\eta)\|^2>C(L_M+u)\}&\le e^{-u}.
\end{aligned}
\end{equation}
These are estimates for the normalized interacting conditional,
not for a Gaussian reference before reweighting.
\end{proposition}
\begin{proof}
In dimension four the diagonal entries of \(\mathsf C_M\) are
bounded uniformly in \(M\), as established in V58. The scalar
Laplace estimate \eqref{eq:v65-Laplace} and
\eqref{eq:v65-covariance-order} give a uniform sub-Gaussian tail
for each real link-colour coordinate. A union bound over the
\(12M^4\) coordinates, absorbing this count into \(C L_M\),
proves the first inequality.

For completeness, the second estimate does not require independent
Gaussian Fourier coefficients. If \(B\) is any real linear map
on the high space, independent auxiliary Gaussian-source
integration in \eqref{eq:v65-Laplace} gives
\begin{equation}\label{eq:v65-quadratic}
 \mathbb E_{\nu_{\pm,c}}e^{s\|B\eta\|^2}
       \le \det(I-2s B K_\ge B^*)^{-1/2},
       \qquad 0\le s<(2\|B K_\ge B^*\|)^{-1}.
\end{equation}
When the norm is zero the expression has its evident limiting
meaning. In detail, for a standard real auxiliary Gaussian \(g\),
\(e^{s\|B\eta\|^2}=\mathbb E_g
e^{\sqrt{2s}\langle B^*g,\eta\rangle}\);
Tonelli and \eqref{eq:v65-Laplace} prove the formula.

For a side-\(R\) periodic cube \(Q\), use
\[
 B_Q=|Q|^{-1/2}P_Q\operatorname{res}_Q d\Delta^{-1}
                                      \big|_{\mathcal H_\ge}.
\]
The order \eqref{eq:v65-covariance-order} and
Lemma~\ref{lem:v61-cube} imply
\[
 \|B_QK_\ge B_Q^*\|\le C,\qquad
       \operatorname{Tr}(B_QK_\ge B_Q^*)\le C\log(2R).
\]
The determinant bound \eqref{eq:v65-quadratic} yields the same
quadratic tail as \eqref{eq:v64-conditional-tail}, with these
comparison matrices. A union bound over at most \(M^5\) cubes
therefore gives
\[
 \nu_{\pm,c}\{\|d\Delta^{-1}\eta\|_{\mathrm{BMO}_2}^2
                                  >C(L_M+u)\}\le e^{-u}.
\]
Apply the deterministic div--curl bound
\eqref{eq:v61-gauge-BMO}, valid for each mean-zero Coulomb
configuration \(\eta\). This proves the gauge tail.
\end{proof}

\begin{theorem}[The outer Wilson conditional reaches the inner safe domain]
\label{thm:v65-coverage}
There is a constant \(K(r_-)\), independent of \(M,\beta,c,u\),
such that under \eqref{eq:v65-parameters}, \(u\ge3\), and
\begin{equation}\label{eq:v65-beta}
                        \beta\ge K(r_-)(L_M+u),
\end{equation}
one has, simultaneously for all individual sources \(|c|\le c_*\),
\begin{equation}\label{eq:v65-coverage}
                    \nu_{+,c}(D_-(c))\ge1-2e^{-u}.
\end{equation}
More precisely, outside an event of \(\nu_{+,c}\)-probability
at most \(2e^{-u}\), the entire fixed-mean segment
\(a_s=a_c+s\beta^{-1/2}\eta\), \(0\le s\le1\), satisfies
\begin{equation}\label{eq:v65-segment}
 \max_e|a_s(e)|\le3r_-/4,\qquad
        \tfrac{23}{32}I\preceq\mathcal F_M(a_s)
                                      \preceq\tfrac{41}{32}I.
\end{equation}
Uniformity in \(c\) means a uniform bound for this family of
conditional laws; no common random coupling of all sources is
asserted.
\end{theorem}
\begin{proof}
Use the intersection of the two good events in
\eqref{eq:v65-tails} for \(\nu_{+,c}\). Its complement has
probability at most \(2e^{-u}\). Enlarge \(K(r_-)\) to ensure
on this intersection that
\[
 t\max_e|\eta_e|\le r_-/2,\qquad t\|G_1(\eta)\|\le1/8.
\]
For \(0\le s\le1\), the background condition in
\eqref{eq:v65-parameters} gives the first bound in
\eqref{eq:v65-segment}. Linearity and
Lemma~\ref{lem:v65-base} give
\[
 \|G_1(a_s)\|
       \le\|G_1(a_c)\|+t\|G_1(\eta)\|\le1/4 .
\]
By \eqref{eq:v65-even} and the isometry of \(U\),
\(\|G_{\rm e}(a_s)\|\le r_-^2/2\le1/32\).
All these operators are self-adjoint on the affine Coulomb slice.
Thus \(\|\mathcal F_M(a_s)-I\|\le9/32\), proving the stated
spectral interval. At \(s=1\) both strict inner-domain conditions
hold, with a margin. This proves \eqref{eq:v65-coverage}.
\end{proof}

For any prescribed \(p>0\), taking
\(u=pL_M+3\) gives, whenever \(\beta\ge K_p(r_-)L_M\),
\begin{equation}\label{eq:v65-polynomial-window}
 \sup_{|c|\le c_*}\nu_{+,c}(D_-(c)^c)
                               \le2e^{-p-3}M^{-p}.
\end{equation}
This is a joint coupling--volume estimate, not an asymptotic
expansion with \(M\) fixed first. Its constant is not identified
with a physical running-coupling coefficient. The high threshold
and outer small-field domain remain part of its hypotheses.

\subsection{Exact conditional normalization and source-amplitude stability}

\begin{theorem}[Nested cutoff ratios with no partition-denominator loss]
\label{thm:v65-ratios}
Under Theorem~\ref{thm:v65-coverage}, set \(\delta=2e^{-u}<1\).
For every \(|c|\le c_*\),
\begin{equation}\label{eq:v65-ratio}
 1-\delta\le \frac{Z_-(c)}{Z_+(c)}\le1,\qquad
 \|\nu_{-,c}-\nu_{+,c}\|_{\mathrm{TV}}
              =1-\frac{Z_-(c)}{Z_+(c)}\le\delta ,
\end{equation}
where total variation means \(\sup_A|\mu(A)-\nu(A)|\).
Define the source-normalized high-integral amplitudes
\begin{equation}\label{eq:v65-amplitude}
              \mathcal Q_\pm(c)=-\log Z_\pm(c)+\log Z_\pm(0).
\end{equation}
Then
\begin{equation}\label{eq:v65-amplitude-error}
 \sup_{|c|\le c_*}|\mathcal Q_-(c)-\mathcal Q_+(c)|
                      \le-\log(1-\delta)\le4e^{-u}.
\end{equation}
Equivalently,
\begin{equation}\label{eq:v65-normalized-ratio}
 1-\delta\le
       \frac{Z_-(c)/Z_-(0)}{Z_+(c)/Z_+(0)}
                                      \le(1-\delta)^{-1}.
\end{equation}
\end{theorem}
\begin{proof}
The domains are nested and the unnormalized weight \(w_c\) is
identical on their intersection. Therefore the exact identities
\[
 \frac{Z_-(c)}{Z_+(c)}=\nu_{+,c}(D_-(c)),\qquad
          \nu_{-,c}=\nu_{+,c}(\,\cdot\,\mid D_-(c))
\]
hold. Conditioning a probability measure on an event of
probability \(1-\delta_c\) changes its total variation by exactly
\(\delta_c\): the mass removed outside the event is \(\delta_c\),
and the mass added inside it is also \(\delta_c\). This proves
\eqref{eq:v65-ratio}, including equality before the bound.

Let
\[
 d(c)=-\log[Z_-(c)/Z_+(c)].
\]
The ratio bound places every \(d(c)\) in the same interval
\([0,-\log(1-\delta)]\). Since
\(\mathcal Q_-(c)-\mathcal Q_+(c)=d(c)-d(0)\), the difference
has absolute value at most the length of that interval.
For \(u\ge3\), \(\delta<1/2\) and
\(-\log(1-\delta)\le2\delta=4e^{-u}\).
Dividing the ratios at \(c\) and at zero proves
\eqref{eq:v65-normalized-ratio}.
\end{proof}

The source action \(\beta S(a_c)\) and the separate geometric
normalizers may be large. They have not been bounded below by a
volume-independent constant, nor discarded. The common-weight
conditional identity gives a relative error directly under the
normalized full Wilson law, so no division of a Gaussian bad-event
bound by an exponentially small unknown partition function occurs.

\begin{proposition}[Fixed local polynomial observables]
\label{prop:v65-local}
Let \(P\) be a fixed polynomial in a fixed finite number of real
link-colour components of \(\eta\), with coefficients independent
of \(M,\beta,c\). Its locations may vary with \(M\). Under the same
conditions there is \(C_P\), independent of those parameters,
such that
\begin{equation}\label{eq:v65-local}
 |\mathbb E_{\nu_{-,c}}P-\mathbb E_{\nu_{+,c}}P|
                                                \le C_Pe^{-u/2}.
\end{equation}
\end{proposition}
\begin{proof}
The scalar sub-Gaussian estimates used above bound all fixed
coordinate moments uniformly. H{\"o}lder's inequality bounds each
monomial moment of \(P^2\), so
\(\|P\|_{L^2(\nu_{+,c})}\le C_P\). Put \(E=D_-(c)\) and
\(\delta_c=\nu_{+,c}(E^c)\). The exact conditioning formula gives
\[
 \mathbb E_{\nu_{-,c}}P-\mathbb E_{\nu_{+,c}}P
   =\frac{\delta_c\mathbb E_{\nu_{+,c}}P
                -\mathbb E_{\nu_{+,c}}(P\mathbf1_{E^c})}
          {1-\delta_c}.
\]
Cauchy--Schwarz bounds its absolute value by
\((\delta_c+\sqrt{\delta_c})\|P\|_2/(1-\delta_c)\).
Use \(0\le\delta_c\le2e^{-u}<1/2\).
\end{proof}

\subsection{What has and has not moved upstream}

The following distinctions are essential for further use.
\begin{itemize}
\item V64 controlled fluctuations about an unspecified center
on a general high fibre. Translation invariance now supplies
that center \emph{exactly} for homogeneous means with zero
nonconstant retained field. It does not supply the center at
an arbitrary nonhomogeneous low field.
\item The excluded inner-radius and gauge-floor region now has
small probability under an actual full Wilson conditional.
Only the transition from \(D_+(c)\) to \(D_-(c)\) is controlled.
The outer positive small-field condition itself has not been
removed from the compact interacting measure.
\item The source-normalized estimate is uniform in the values
of \(c\), including noncommuting choices. Uniform small error
of function values does not imply small derivatives of that
error. In particular, no quadratic or quartic source-jet
estimate, no noncommuting shell coefficient, and no
renormalized coupling follows by differentiating
\eqref{eq:v65-amplitude-error}.
\item The domains move with \(c\). Their contribution was retained
through exact set conditioning, not dropped from a differentiated
integral. Smooth cutoff or explicit boundary estimates would
still be needed for a source-derivative argument.
\item A fixed sharp high-frequency projection is spatially
nonlocal. Neither a local all-scale RG map nor control of the
low marginal has been deduced.
\end{itemize}
The next upstream issue is therefore the response of the high
conditional mean to nonhomogeneous retained fields, together with
the law of those retained fields. Merely repeating the free
Gaussian gauge floor would not address it. An eventual physical
mass-gap argument still needs the continuum state, positivity
and reconstruction, and estimates in physical units.

\subsection{Exact verification and preservation record}

The certificate
\begin{center}\small
\path{scripts/verify_ym_f14_v65_homogeneous_cutoff.py}
\end{center}
uses an explicitly noncommuting rational four-colour source of
squared norm
\[
                   22296487451910/949136745811969<1/16.
\]
It checks the literal homogeneous odd-gauge Fourier blocks at
all 15 nonzero momenta for side two and all 255 for side four,
including Hermiticity, their squared cross-product matrices and
the norm bound of Lemma~\ref{lem:v65-base}. The maximum squared
norms in these two finite regressions are respectively
\[
                              0,\qquad 205/279752.
\]
It also verifies the high comparison inverse, its free covariance
order, and exclusion of translation-fixed modes, with rational
test parameters \(\beta=7,\tau=3\). This latter threshold is an
algebra test, not a claim that the theorem's
\eqref{eq:v65-parameters} or coupling window holds at that test
point. There are 15 and 247 high momenta, and the comparison
traces per volume are
\[
       25623297/5657080,\qquad
       446257500741297/99813791059840.
\]

A separate twelve-state cyclic-orbit model, at three rational
source values, checks exact centering and the nested-normalization
identities. Its maximum excluded probability is \(7/57\).
This is solely a finite structural regression, not a Wilson
probability calculation. The nonlinear uniform estimate is the
analytic theorem proved above; the finite checks do not certify
its universal constants.

The 5,244-byte certificate has SHA--256
\begin{center}\scriptsize\ttfamily
B6912B26CE94225A2C8AB013064C0F94FAE324691B98C44F3B658206292FE7DB.
\end{center}
It uses unchanged V56/V58 arithmetic helpers. The predecessor V64
has 1,578,528 bytes and SHA--256
\begin{center}\scriptsize\ttfamily
FB7A3C0F729BC82418D1E27A4BA991F5154B9C11213DD57A1A02D3F59A983120.
\end{center}
Its complete body is retained apart from the title version and
this terminal insertion. V65 replaces V64 as the sole active
main note. Supplied sources, companions, monographs and archives
are unchanged.

\begin{firewall}
The complete nonlinear Wilson high conditional at a small
homogeneous retained mean has an exact center, a joint
coupling--volume inner-domain probability bound, and controlled
source-normalized cutoff dependence. These do not cover general
retained low fields or the full compact measure. No continuum
Yang--Mills construction or physical mass gap is proved.
The full goal remains active.
\end{firewall}
\section{V66: nonhomogeneous low inputs and controlled conditional means}
\label{sec:v66}

\subsection{From a zero center to a quantified nonlinear response}

V65 supplies an exact zero center when the retained field is
homogeneous. A general nonconstant retained field breaks that
symmetry, and V64's concentration about an unspecified center
does not by itself locate the field. We now derive response
estimates for such retained inputs. They concern the complete
finite-coupling Wilson density and have no hidden dimension factor.

There is a second issue: a coordinate cutoff generally moves
when the retained field changes. We first differentiate on an
explicit fixed high carrier on which the complete density remains
valid. We then compare that law to a specified moving cutoff by
exact conditioning, including a dimension-independent estimate
for the change in its mean. The fixed carrier is not silently
identified with the original moving one.

The main new cancellation is representation-theoretic. A
translation-equivariant linear map cannot send a low Fourier
character to a disjoint high character. Consequently the high
conditional mean has zero first derivative at zero nonconstant
input, even with a noncommuting homogeneous mean retained.
The response is therefore at least quadratic on the fixed carrier.
This is an exact normalized nonlinear statement, not the
conditional Gaussian calculation of V15 or another cubic
coefficient calculation.

The probability of a suitable ball of retained inputs under the
low marginal is still not known. The results below do not assert
that this ball contains typical low fields of the interacting
theory. That is a remaining upstream requirement, not a change
of the full Yang--Mills goal. The next companion checkpoint
remains V70.

\subsection{Uniform derivatives of the complete coordinate potential}

Let \(0<R\le1/4\), and let \(a\) be a real Coulomb field,
possibly with a fixed harmonic mean, satisfying
\[
             \max_e|a_e|\le R,\qquad \mathcal F_M(a)\succeq\kappa I .
\]
All variations in this subsection are in \(\mathcal H_M\);
in particular they have zero harmonic mean. Write
\[
 \mathcal U_{c,\beta}(z)
       =\beta S(a_c+\beta^{-1/2}z)+W_M(a_c+\beta^{-1/2}z).
\]
The point under consideration is \(a=a_c+\beta^{-1/2}z\).

\begin{lemma}[Mixed and third derivative bounds without a volume factor]
\label{lem:v66-derivatives}
There are constants \(C_{\mathrm S}\), \(C_{\mathrm F}(R)\),
\(C_{\mathrm W,2}(R,\kappa)\), and \(C_{\mathrm W,3}(R,\kappa)\),
independent of \(M,\beta,c\), with the following properties.
First,
\begin{equation}\label{eq:v66-geometric-derivatives}
\begin{aligned}
 \|D\mathcal F_M(a)[h]\|&\le C_{\mathrm F}(R)\|h\|_2,\\
 |D^jW_M(a)[h_1,\ldots,h_j]|
       &\le C_{\mathrm W,j}(R,\kappa)\prod_{i=1}^j\|h_i\|_2,
                                      \qquad j=2,3.
\end{aligned}
\end{equation}
The first inequality does not require the lower bound
\(\mathcal F_M(a)\succeq\kappa I\).
For the spectral split at \(\tau\), \(v\in\mathcal H_<\) and
\(h\in\mathcal H_\ge\), one has
\begin{equation}\label{eq:v66-mixed}
 |D^2\mathcal U_{c,\beta}(z)[v,h]|
               \le B_\beta\|v\|_2\|h\|_2,\qquad
 B_\beta=C_{\mathrm S}R+\beta^{-1}C_{\mathrm W,2}(R,\kappa).
\end{equation}
For any three mean-zero Coulomb variations,
\begin{equation}\label{eq:v66-third}
\begin{aligned}
 |D^3\mathcal U_{c,\beta}(z)[h_1,h_2,h_3]|
     &\le T_\beta\prod_{i=1}^3\|h_i\|_2,\\
 T_\beta&=C_{\mathrm S}\beta^{-1/2}
                    +C_{\mathrm W,3}(R,\kappa)\beta^{-3/2}.
\end{aligned}
\end{equation}
\end{lemma}
\begin{proof}
The Hilbert--Schmidt bound for \(G_1(h)\) in V59 gives
\(\|G_1(h)\|_{\mathrm{HS}}\le C\|h\|_2\).
The even part of \(\mathcal F_M\) is a compression by the
fixed isometry \(U\) of a block-diagonal single-link function.
Its first derivative has Hilbert--Schmidt norm at most
\(C_R\|h\|_2\), since all single-link derivatives are uniformly
bounded on \(|a_e|\le R<\pi\). Thus
\[
          \|D\mathcal F_M(a)[h]\|_{\mathrm{HS}}\le C_R\|h\|_2.
\]
This proves the first assertion, even without a gauge floor.
For the higher derivatives, trace-norm contraction under
compression and Cauchy--Schwarz on the link sum give
\[
 \|D^2\mathcal F_M(a)[h,v]\|_1\le C_R\|h\|_2\|v\|_2,\qquad
 \|D^3\mathcal F_M(a)[h,v,w]\|_1
                                \le C_R\|h\|_2\|v\|_2\|w\|_2 .
\]
For the latter inequality one may use
\(\sum_e|h_e||v_e||w_e|
\le\|h\|_2\|v\|_2\|w\|_\infty\le\|h\|_2\|v\|_2\|w\|_2\).
These estimates do not use the number of links.

In the second and third derivatives of \(-\log\det F\),
where \(F=\mathcal F_M(a)\), every inverse has norm at most
\(\kappa^{-1}\). The third derivative consists of two traces
with three first derivatives, three with one first and one
second derivative, and one with the third derivative of \(F\).
Bound two first derivatives in Hilbert--Schmidt norm and the
remaining factors in operator norm; use trace norm for a
second or third derivative of \(F\). The preceding estimates
bound every term by the product of the three variation norms.
The same link-sum argument applies to the Haar potential.
The second derivative is also precisely the bound of
Lemma~\ref{lem:v63-geometric-Hessian}, extended to the affine
slice as in V65. This proves
\eqref{eq:v66-geometric-derivatives}.

Set \(\epsilon=R\) in the G{\aa}rding estimate
\eqref{eq:v64-Garding}. Since \(\Delta\preceq16I\) and
\(\mathcal V_a\le C R^2\), it gives, on \(\mathcal H_M\),
\[
                 \|D^2S(a)-\Delta\|_{2\to2}\le C_{\mathrm S}R .
\]
The low and high spectral spaces are orthogonal for \(\Delta\),
so its mixed block vanishes. The full rescaling gives
\(D^2\mathcal U=D^2S+\beta^{-1}D^2W_M\), proving
\eqref{eq:v66-mixed}.

Every third derivative of a single Wilson plaquette is bounded
on this fixed compact link-argument set. Summing its trilinear
bound over plaquettes costs only a bounded overlap constant:
discrete H{\"o}lder and \(\|\cdot\|_3\le\|\cdot\|_2\) bound
each product of three shifted fields by their three
\(\ell^2\) norms. Hence
\(|D^3S(a)[h_1,h_2,h_3]|\le C_{\mathrm S}\prod_i\|h_i\|_2\).
Finally \(\beta(\beta^{-1/2})^3=\beta^{-1/2}\), while the
geometric factor scales as \(\beta^{-3/2}\). This proves
\eqref{eq:v66-third}.
\end{proof}

We also fix a universal constant \(C_{\mathrm G}\) such that
\begin{equation}\label{eq:v66-G-linear}
                    \|G_1(h)\|\le C_{\mathrm G}\|h\|_2
                              \quad(h\in\mathcal H_M),
\end{equation}
as follows from the same Hilbert--Schmidt estimate. None of
these derivative bounds is a physical screening-mass assertion.

\subsection{A common carrier with a protected source neighbourhood}

Fix \(r>0\) and \(\tau\) such that
\begin{equation}\label{eq:v66-parameters}
 R=4r\le\tfrac14,\qquad
              4C_0R^2\le\tau\le1,\qquad c_*/M\le r/8,
                          \qquad |c|\le c_*=\tfrac14.
\end{equation}
Use the mean-zero low and high spaces of V64 at this threshold.
Take \(\beta\ge1\), \(t=\beta^{-1/2}\), and put
\begin{equation}\label{eq:v66-guard}
 \gamma=\min\{r/8,\ (64C_{\mathrm G})^{-1},
                              (16C_{\mathrm F}(R))^{-1}\}.
\end{equation}
For nonconstant retained input \(y\in\mathcal H_<\), require
\begin{equation}\label{eq:v66-source-ball}
                              t\|y\|_2\le\gamma .
\end{equation}
The input \(y\) is in fluctuation coordinates; the physical
link-coordinate shift is \(ty\).

Define the fixed high carrier
\begin{equation}\label{eq:v66-carrier}
 E_c=\{\eta\in\mathcal H_\ge:
       \max_e|a_c(e)+t\eta_e|<2r,\ 
                   \mathcal F_M(a_c+t\eta)\succ\tfrac12 I\}.
\end{equation}
It depends on \(c,M,\beta\), but not on \(y\). Let
\begin{equation}\label{eq:v66-law}
\begin{aligned}
 U_y(\eta)&=\mathcal U_{c,\beta}(y+\eta),\\
 \mathcal Z_c(y)&=\int_{E_c}e^{-U_y(\eta)}\,d\eta,\qquad
 d\mu_{c,y}(\eta)
       =\mathcal Z_c(y)^{-1}\mathbf1_{E_c}(\eta)e^{-U_y(\eta)}\,d\eta,\\
 m_c(y)&=\mathbb E_{\mu_{c,y}}\eta .
\end{aligned}
\end{equation}

\begin{lemma}[Validity of the full density on the common carrier]
\label{lem:v66-carrier}
The carrier \(E_c\) is bounded, open, nonempty, convex and
translation invariant. For every \(\eta\in E_c\) and every
input in \eqref{eq:v66-source-ball},
\begin{equation}\label{eq:v66-carrier-floor}
 \max_e|a_c(e)+t(y+\eta)_e|<R,\qquad
              \mathcal F_M(a_c+t(y+\eta))\succ\tfrac7{16}I.
\end{equation}
All \(\mathcal Z_c(y)\) are finite and positive. The full
high Hessian satisfies
\begin{equation}\label{eq:v66-precision}
 D^2_{\eta\eta}U_y\succeq A_\ge
       =\tfrac14\Delta|_{\mathcal H_\ge}+\tfrac{2}{3\beta}I
       \succeq\alpha_\beta I,\qquad
                    \alpha_\beta=\tau/4+2/(3\beta).
\end{equation}
The constants \(B_\beta,T_\beta\) in
Lemma~\ref{lem:v66-derivatives} may be used throughout this
family with \(\kappa=7/16\).
\end{lemma}
\begin{proof}
Matrix concavity of \(\mathcal F_M\) proves convexity exactly
as in V65. The base estimate
\eqref{eq:v65-base}, together with
\(\max_e|a_c(e)|\le r/8\), places zero strictly inside the
carrier. Boundedness and translation invariance follow from
its definition and the Fourier invariance of the high space.

Along the segment \(a_c+t\eta+sty\), \(0\le s\le1\), the
link radius is at most \(2r+\gamma<R\). Integrating the first
inequality of \eqref{eq:v66-geometric-derivatives} gives
\[
 \|\mathcal F_M(a_c+t\eta+ty)-\mathcal F_M(a_c+t\eta)\|
                \le C_{\mathrm F}(R)t\|y\|_2\le1/16.
\]
This proves the asserted gauge floor. Finite positive
normalization follows on the bounded carrier, with its
strictly positive density near zero. The affine full Wilson
G{\aa}rding lower bound and the threshold in
\eqref{eq:v66-parameters} give \eqref{eq:v66-precision}.
\end{proof}

This carrier uses the complete \(\rho_M\), not a polynomial
in its logarithm. It is a specified coordinate restriction,
not a claim that all configurations of a larger positive
branch belong to it.

\subsection{A covariance identity with no differentiated boundary}

We will use the following finite-dimensional consequence of
strong log-concavity:
\begin{equation}\label{eq:v66-Poincare}
 \operatorname{Var}_\mu(f)\le\alpha^{-1}
                         \mathbb E_\mu\|\nabla f\|^2
 \quad\hbox{if }D^2U\succeq\alpha I
              \hbox{ on a convex support}.
\end{equation}
The unconstrained variance inequality is the classical
Brascamp--Lieb inequality; a primary formulation is
Carlen--Cordero-Erausquin--Lieb,
\emph{Asymmetric Covariance Estimates of Brascamp--Lieb Type},
equation (1.3),
\href{https://arxiv.org/pdf/1106.0709}{arXiv:1106.0709}.
For clarity, the convex-support extension here follows by
extending \(U-\alpha|\eta|^2/2\) as a lower-semicontinuous
convex function with value \(+\infty\) off the support, taking
its Moreau regularizations and then smooth convex
approximations. Their addition to \(\alpha|\eta|^2/2\)
retains the same Hessian lower bound. A common affine lower
bound supplies Gaussian domination, and the resulting
measures converge to the restricted measure. Passing the
variance inequality to the limit gives
\eqref{eq:v66-Poincare}; smooth functions on the bounded
carrier suffice here. Thus a hard convex boundary does not
require discarding a surface term from a source derivative.

\begin{theorem}[Exact first and second normalized responses]
\label{thm:v66-response}
On the interior of the input ball, \(m_c\) is smooth. For
\(v,w\in\mathcal H_<\) put
\[
 \phi_v=D_yU_y[v],\qquad
          \psi_{vw}=D_y^2U_y[v,w],\qquad
          \widetilde\phi_v=\phi_v-\mathbb E_{\mu_{c,y}}\phi_v.
\]
Vector covariances are understood componentwise. Then
\begin{equation}\label{eq:v66-response}
\begin{aligned}
 Dm_c(y)[v]
       &=-\operatorname{Cov}_{\mu_{c,y}}(\eta,\phi_v),\\
 D^2m_c(y)[v,w]
       &=-\operatorname{Cov}_{\mu_{c,y}}(\eta,\psi_{vw})
         +\mathbb E_{\mu_{c,y}}
                   [(\eta-m_c(y))\widetilde\phi_v\widetilde\phi_w].
\end{aligned}
\end{equation}
Uniformly in the input ball,
\begin{equation}\label{eq:v66-response-bounds}
\begin{aligned}
 \|Dm_c(y)[v]\|_2
       &\le (B_\beta/\alpha_\beta)\|v\|_2,\\
 \|D^2m_c(y)[v,w]\|_2
       &\le Q_\beta\|v\|_2\|w\|_2,\\
 Q_\beta&=T_\beta/\alpha_\beta
                         +2B_\beta^2/\alpha_\beta^{3/2}.
\end{aligned}
\end{equation}
There is no factor \(\sqrt{\dim\mathcal H_\ge}\) in these bounds.
\end{theorem}
\begin{proof}
The integration carrier is fixed in \(y\). Its closure is
bounded, and the source neighbourhood has a positive gauge
margin by \eqref{eq:v66-carrier-floor}. For each finite \(M,\beta\)
all needed derivatives of the integrand are bounded on a
slightly smaller source neighbourhood and this closure.
Differentiation under the integral and quotient rule therefore
give both identities in \eqref{eq:v66-response}. The
normalizer derivatives produce exactly the centered covariance
and third joint cumulant displayed there.

For a unit high vector \(J\), set \(X_J=\langle J,\eta\rangle\).
The mixed and third derivative bounds give
\[
 \|\nabla_\eta\phi_v\|\le B_\beta\|v\|_2,\qquad
 \|\nabla_\eta\psi_{vw}\|
                              \le T_\beta\|v\|_2\|w\|_2.
\]
Poincar\'e and Cauchy--Schwarz imply
\[
 |\operatorname{Cov}(X_J,\phi_v)|
                           \le B_\beta\|v\|_2/\alpha_\beta,
 \qquad
 |\operatorname{Cov}(X_J,\psi_{vw})|
                      \le T_\beta\|v\|_2\|w\|_2/\alpha_\beta.
\]
For the last term of the second identity, use
\[
 \mathbb E[(X_J-\mathbb E X_J)\widetilde\phi_v\widetilde\phi_w]
       =\operatorname{Cov}(X_J,\widetilde\phi_v\widetilde\phi_w).
\]
Another application of \eqref{eq:v66-Poincare}, now to the
product, bounds its absolute value by
\[
 \frac1{\alpha_\beta}
 \left(\mathbb E\left\|
     \widetilde\phi_v\nabla\phi_w+
                         \widetilde\phi_w\nabla\phi_v\right\|^2\right)^{1/2}.
\]
Use the triangle inequality in \(L^2\) and the estimate
\[
 \|\widetilde\phi_v\|_{L^2}
             \le B_\beta\|v\|_2/\sqrt{\alpha_\beta}.
\]
The resulting bound is
\[
                2B_\beta^2\|v\|_2\|w\|_2/\alpha_\beta^{3/2}.
\]
Take the supremum over unit \(J\) to obtain the norm estimates.
Bounding each coordinate separately and then summing would
incorrectly lose the dimension-independent conclusion.
\end{proof}

No source derivative of the indicator of a moving domain is
missing from \eqref{eq:v66-response}: its indicator is precisely
that of the fixed \(E_c\). The moving domain will be treated
below through a different, exact identity.

\subsection{Spectral selection removes the linear displacement}

\begin{theorem}[No linear low-to-high response, and higher selection]
\label{thm:v66-selection}
One has
\begin{equation}\label{eq:v66-zero-linear}
                            m_c(0)=0,\qquad Dm_c(0)=0.
\end{equation}
Consequently, with \(q=\|y\|_2\),
\begin{equation}\label{eq:v66-mean-bound}
\begin{aligned}
 \|m_c(y)\|_2&\le b_\beta(q),\\
 b_\beta(q)&=\min\left\{
       \frac{B_\beta}{\alpha_\beta}q,\
                              \frac{Q_\beta}{2}q^2\right\}.
\end{aligned}
\end{equation}
After complexification, consider the \(j\)-th derivative at zero.
Inputs of characters \(k_1,\ldots,k_j\) can produce only the
character \(k_1+\cdots+k_j\), modulo the reciprocal torus.
If that character is not high, the output is zero.
For an integer \(p\ge1\), it follows in particular that all
derivatives of orders \(1\le j\le p\) vanish on inputs from
\begin{equation}\label{eq:v66-deep-low}
                    \mathbf1_{\{\Delta<\tau/p^2\}}\mathcal H_M .
\end{equation}
This last statement is an exact selection rule; no
volume-uniform bound on a derivative of order \(p+1\) is
asserted here.
\end{theorem}
\begin{proof}
At \(y=0\), the complete law on \(E_c\) is translation
invariant. Its mean belongs to the high space and is fixed
by all translations, so it vanishes, as in V65.
For a torus translation \(T_x\), translation covariance of
the full density and invariance of \(E_c\) give
\[
                    m_c(T_xy)=T_xm_c(y).
\]
Differentiate at zero. The linear map \(Dm_c(0)\) intertwines
the translation representations on the low and high spaces.
On complex Fourier components this preserves each character.
The two spaces have disjoint characters, since the threshold
is a function of momentum alone. Thus this map is zero.
Colour and orientation multiplicities do not change the
disjointness. A noncommuting \(c\) still has zero spatial
momentum and leaves the argument intact.

Integrating the first derivative bound from zero proves the
linear estimate in \eqref{eq:v66-mean-bound}. Integrating the
second derivative bound twice, using both identities in
\eqref{eq:v66-zero-linear}, gives
\[
 m_c(y)=\int_0^1(1-s)D^2m_c(sy)[y,y]\,ds,
\]
and hence the quadratic estimate.

Differentiating equivariance \(j\) times shows that the output
character is the product of the input characters. This is
the asserted momentum sum rule. With
\(q_\mu(k)=e^{ik_\mu}-1\), one has
\(q_\mu(k+\ell)=q_\mu(k)+e^{ik_\mu}q_\mu(\ell)\), so
\[
 \sqrt{\lambda(k_1+\cdots+k_j)}
                       \le\sum_{i=1}^j\sqrt{\lambda(k_i)}.
\]
For inputs in \eqref{eq:v66-deep-low} and \(j\le p\), this
is strictly less than \(\sqrt{\tau}\).
The output is therefore not high and must vanish.
Multilinearity extends the conclusion from characters to
arbitrary inputs in that subspace.
\end{proof}

Two ordinary low characters can sum to a high one. Thus
\eqref{eq:v66-zero-linear} cannot be extended to all quadratic
responses by the same symmetry. The finite actual Wilson
cubic check at the end of this section illustrates that
such a coupling is allowed. It is not an evaluation of the
integrated conditional mean.

\subsection{Concentration centered at the now controlled mean}

\begin{proposition}[Uniform centered high tails for nonconstant inputs]
\label{prop:v66-tails}
Let \(\xi=\eta-m_c(y)\), \(L_M=1+\log M\).
Uniformly in all sources in \eqref{eq:v66-source-ball},
\begin{equation}\label{eq:v66-Laplace}
 \mathbb E_{\mu_{c,y}}e^{\langle J,\xi\rangle}
       \le e^{\langle J,K_\ge J\rangle/2},\qquad
 K_\ge=A_\ge^{-1}\preceq4\mathsf C_M .
\end{equation}
There is a universal \(C\) such that
\begin{equation}\label{eq:v66-tails}
\begin{aligned}
 \mu_{c,y}\{\max_e|\xi_e|^2>C(L_M+u)\}&\le e^{-u},\\
 \mu_{c,y}\{\|G_1(\xi)\|^2>C(L_M+u)\}&\le e^{-u},
                                                   \qquad u\ge0 .
\end{aligned}
\end{equation}
Also
\[
 \max_e|m_c(y)_e|\le b_\beta(\|y\|_2),\qquad
          \|G_1(m_c(y))\|\le C_{\mathrm G}b_\beta(\|y\|_2).
\]
\end{proposition}
\begin{proof}
The precision bound \eqref{eq:v66-precision} on a convex
carrier gives the centered Laplace estimate by the
Gaussian-convolution argument of V64. Its Fourier comparison
is the same as in V65. The scalar link bounds and
Gaussian-source quadratic-map argument in
Proposition~\ref{prop:v65-tails} therefore apply to \(\xi\).
They use the deterministic div--curl bound, not independence
of the interacting Fourier coordinates. The two mean bounds
follow from \eqref{eq:v66-mean-bound} and
\eqref{eq:v66-G-linear}.
\end{proof}

\subsection{An actual moving cutoff and its mean comparison}

The moving domain to which we now compare is explicitly
\begin{equation}\label{eq:v66-moving}
 D_c(y)=\{\eta\in\mathcal H_\ge:
       \max_e|a_c(e)+t(y+\eta)_e|<r,\
        \mathcal F_M(a_c+t(y+\eta))\succ\tfrac34I\}.
\end{equation}
This has a quantitative gauge floor in its definition.
It is not the whole positive branch with only
\(\mathcal F_M\succ0\) imposed.
Use exactly the same unnormalized density as above:
\begin{equation}\label{eq:v66-moving-law}
\begin{aligned}
 Z_c(y)&=\int_{D_c(y)}e^{-U_y(\eta)}\,d\eta,\\
 d\nu_{c,y}(\eta)&=Z_c(y)^{-1}
                  \mathbf1_{D_c(y)}(\eta)e^{-U_y(\eta)}\,d\eta,\qquad
 n_c(y)=\mathbb E_{\nu_{c,y}}\eta .
\end{aligned}
\end{equation}

\begin{theorem}[Moving-cutoff coverage and dimension-independent mean error]
\label{thm:v66-moving}
There is \(K_0(r)\) such that the following conditions suffice:
\begin{equation}\label{eq:v66-window}
 u\ge3,\qquad
 \beta\ge K_0(r)(L_M+u),\qquad
      t\|y\|_2\le\gamma,\qquad t\,b_\beta(\|y\|_2)\le\gamma .
\end{equation}
Then \(D_c(y)\subset E_c\) and, with \(\delta=2e^{-u}\),
\begin{equation}\label{eq:v66-moving-ratio}
 1-\delta\le Z_c(y)/\mathcal Z_c(y)\le1,\qquad
                \|\nu_{c,y}-\mu_{c,y}\|_{\mathrm{TV}}\le\delta .
\end{equation}
Moreover,
\begin{equation}\label{eq:v66-moving-mean}
 \|n_c(y)-m_c(y)\|_2
           \le\sqrt{\frac{\delta}{\alpha_\beta(1-\delta)}},
 \qquad
 \|n_c(y)\|_2
           \le b_\beta(\|y\|_2)
                      +\sqrt{\frac{\delta}{\alpha_\beta(1-\delta)}}.
\end{equation}
For fixed \(r,\tau\), a simpler sufficient window is
\begin{equation}\label{eq:v66-simple-window}
            \beta\ge K(r,\tau)\,[L_M+u+\|y\|_2^2].
\end{equation}
All statements are uniform for the allowed \(c,y,M,\beta\).
\end{theorem}
\begin{proof}
If \(\eta\in D_c(y)\), subtracting \(ty\) makes the link
radius at most \(r+\gamma<2r\). The first derivative bound
for \(\mathcal F_M\), along this segment of radius less
than \(R\), changes the gauge operator by at most \(1/16\).
Thus the base operator at \(a_c+t\eta\) is strictly above
\((3/4-1/16)I\), and in particular above \(I/2\). This proves
the actual inclusion \(D_c(y)\subset E_c\).

Use the intersection of the two good events in
\eqref{eq:v66-tails}. Its complement has probability at most
\(\delta\). Increase \(K_0(r)\) so that on this event
\[
 t\max_e|\xi_e|\le r/8,\qquad t\|G_1(\xi)\|\le1/64 .
\]
Write
\[
 a=a_c+t(y+m_c(y)+\xi).
\]
Each of its four displayed terms contributes at most
\(r/8\) to the link radius, by
\eqref{eq:v66-parameters}, \eqref{eq:v66-guard},
\eqref{eq:v66-window} and the preceding tail event.
Hence \(\max_e|a_e|\le r/2<r\).
For the odd operator, \eqref{eq:v65-base} and the guard give
\[
 \|G_1(a)\|
       \le 1/8+1/64+1/64+1/64=11/64 .
\]
The even term has norm at most \(R^2/2\le1/32\), by
\eqref{eq:v65-even}. Thus
\(\|\mathcal F_M(a)-I\|\le13/64\), giving the strict floor
\(\mathcal F_M(a)\succeq(51/64)I\succ(3/4)I\).
The good event is therefore contained in \(D_c(y)\), and
\(\mu_{c,y}(D_c(y))\ge1-\delta\).

The common-weight identity now reads exactly
\[
       Z_c(y)/\mathcal Z_c(y)=\mu_{c,y}(D_c(y)),\qquad
       \nu_{c,y}=\mu_{c,y}(\,\cdot\,\mid D_c(y)).
\]
It proves \eqref{eq:v66-moving-ratio}. For the stronger
mean comparison, put
\(\delta_y=\mu_{c,y}(D_c(y)^c)\le\delta\).
For every unit high vector \(J\),
\[
 \langle J,n_c(y)-m_c(y)\rangle
    =-\frac{\mathbb E_{\mu_{c,y}}
            [\langle J,\xi\rangle\mathbf1_{D_c(y)^c}]}
             {1-\delta_y}.
\]
The linear variable has mean zero and variance at most
\(\alpha_\beta^{-1}\); the indicator has variance
\(\delta_y(1-\delta_y)\). Cauchy--Schwarz therefore bounds
the absolute value by
\(\sqrt{\delta_y/(\alpha_\beta(1-\delta_y))}\).
Taking the supremum over unit \(J\), rather than summing
coordinate errors, proves \eqref{eq:v66-moving-mean}.

Finally, for \(\beta\ge1\), \(B_\beta\) is bounded by a
constant depending on \(r\), and
\(\alpha_\beta\ge\tau/4\). Thus
\(b_\beta(q)\le C(r,\tau)q\).
Choosing \(K(r,\tau)\) large enough in
\eqref{eq:v66-simple-window} implies all three remaining
conditions in \eqref{eq:v66-window}.
\end{proof}

At \(y=0\), both means are exactly zero by translation
invariance; the error estimate is not needed there.
For nonzero \(y\), \eqref{eq:v66-moving-mean} is a bound
for the actual moving-cutoff mean. It does not assert
that differentiating the moving cutoff produces no
boundary contribution.

\begin{proposition}[Source-normalized amplitudes on the moving cutoff]
\label{prop:v66-amplitude}
For the same allowed input family, define
\[
 \mathcal A_c(y)=-\log\mathcal Z_c(y)+\log\mathcal Z_c(0),
             \qquad
 A_c(y)=-\log Z_c(y)+\log Z_c(0).
\]
Then
\begin{equation}\label{eq:v66-amplitude}
               |A_c(y)-\mathcal A_c(y)|
                        \le-\log(1-2e^{-u})\le4e^{-u}.
\end{equation}
\end{proposition}
\begin{proof}
Both \(y\) and zero satisfy the hypotheses with the same
constants. Define
\[
                d_c(y)=-\log[Z_c(y)/\mathcal Z_c(y)].
\]
All these values lie in
\([0,-\log(1-2e^{-u})]\).
The displayed difference is \(d_c(y)-d_c(0)\), whose
absolute value is at most the length of this interval.
\end{proof}

Taking \(u=p(1+\log M)+3\) makes the amplitude error
\(O(M^{-p})\) and the mean-comparison error
\(O_\tau(M^{-p/2})\), under
\(\beta\ge K_p(r,\tau)(1+\log M+\|y\|_2^2)\).
These estimates are joint in coupling and volume.
They are not obtained by dividing a reference bad-event
probability by an unknown full partition function.

\subsection{Scope of the new low-field neighbourhood}

The center is no longer left as an unspecified function
on this family: \eqref{eq:v66-mean-bound} bounds it
nonlinearly, and \eqref{eq:v66-moving-mean} transfers the
bound to an explicitly moving physical link cutoff.
The main unresolved implications are nevertheless substantial.
\begin{itemize}
\item The input control is in the full \(\ell^2\) norm
of the retained fluctuation field. At a logarithmic
coupling window it allows a ball of squared norm of
order \(\log M\), with constants as specified above.
It does not control typical extensive low fields.
No estimate for the interacting low marginal placing
it in this ball has been proved.
\item The source-dependent cutoff in
\eqref{eq:v66-moving} imposes a gauge floor \(3/4\).
The common carrier imposes a floor \(1/2\) at the
unshifted mean. Concentration under either restricted
law does not determine the mass of its carrier under
the larger law with only \(\mathcal F_M\succ0\), still
less under the full compact gauge measure.
\item The exact zero first derivative belongs to the
fixed-carrier mean. The moving-carrier result is a
finite-value estimate with an explicit small error.
Neither this error nor the amplitude error can be
differentiated without an additional boundary argument.
\item The parameters \(r,\tau\) are fixed and the
projection is a sharp Fourier projection. The constants
are not proved uniform as \(\tau\downarrow0\).
No spatially local multiscale map has been constructed.
\end{itemize}
The remaining upstream work is to control a spatially
appropriate norm or distribution of the retained field,
and to pay the larger-domain comparison. A physical
continuum construction and Hamiltonian mass gap still
require those inputs and the reconstruction interfaces.
An auxiliary conditional Poincar\'e constant is not
that mass gap.

\subsection{Exact checks and append preservation}

The certificate
\begin{center}\small
\path{scripts/verify_ym_f14_v66_low_source_response.py}
\end{center}
has three independent roles.
First, a four-state rational weighted model verifies
the normalized first and second derivative identities,
including the covariance and centered third cumulant.
Its two derivatives are
\[
                         -446/289,\qquad 2782/4913.
\]
A separate cyclic low/high example has zero first
derivative and second derivative \(-7\), showing that
first-order selection is not general vanishing of the
nonlinear mean. These weighted models are not Wilson
probability measures.

Second, at side four and the algebraic test threshold
\(\tau=3\), there are 8 low and 247 high nonzero momenta.
The certificate checks all 1,976 relative-character
averages for the disjoint spaces. Among the 64 ordered
low pairs, 56 can sum to a high character, with 28
distinct high outputs. This is a finite check of the
selection mechanism, not a numerical inference of its
general proof. The test threshold is not claimed to
satisfy the theorem's \(\tau\le1\) window.

Third, it uses actual side-four Wilson plaquettes and
the cubic identity already established in V64. With
\[
 y_1(x)=\cos(\pi x_2/2)E_2,\qquad
 y_2(x)=\cos(\pi x_1/2)E_1,\qquad y_3=y_4=0,
\]
and \(h=d^*B\), where the only nonzero two-form is
\[
 B_{12}(x)=\cos(\pi(x_1+x_2+\epsilon)/2)E_3,
                                \qquad \epsilon=0,1,
\]
the respective exact cubic polarizations are
\[
                    D^3S(0)[y,y,h]=0,\qquad -1024 .
\]
The low field has Laplacian eigenvalue \(2\), and the
high field eigenvalue \(4\). The check evaluates all
1,536 plaquettes for each polarization and agrees with
an independent odd cubic difference. This is an actual
allowed Wilson vertex, not a claim that its nonzero
value survives every contribution to the normalized
integrated mean.

The 6,858-byte certificate has SHA--256
\begin{center}\scriptsize\ttfamily
3734F44BB0C1AA8590AE329FEFDEAC4A4B6671DF7B0CBFA34A6FD61BD69D769C.
\end{center}
It imports unchanged V56, V58 and V64 helpers and their
unchanged dependencies. The predecessor V65 has
1,602,671 bytes and SHA--256
\begin{center}\scriptsize\ttfamily
B8BE4FD7B9EF42611A4A99FA54B2ED6513F6435BF3034E6B3373F5B599C86440.
\end{center}
Its entire body is retained apart from the title
version and this terminal insertion. V66 replaces V65
as the sole active main note. Source documents,
companions, monographs and archives remain unchanged.

\begin{firewall}
The full high conditional now has a quantitative
nonhomogeneous-input response on a controlled low-field
neighbourhood, with exact linear selection and an
explicit moving-cutoff mean comparison. This neither
controls the entire low marginal nor removes the
positive coordinate-sector restriction. The continuum
Yang--Mills construction and physical mass gap remain
unproved. The full goal remains active.
\end{firewall}
\section{V67: extensive retained inputs and interacting bad-link estimates}
\label{sec:v67}

\subsection{Replacing the global source ball without dropping the gauge term}

V66 makes the high conditional mean quantitative, but its
sufficient source guard uses the total \(\ell^2\) norm.
At fixed low-frequency threshold, that guard is much smaller
than a typical large-volume free low field. Repeating that
global restriction would not address the retained-field
bottleneck.

Here we replace it by a bound on the maximum link amplitude
and on the \emph{actual odd gauge operator}. The new guard
allows an extensive total field norm. On this enlarged
family the full high conditional still has uniform
coercivity and response bounds. We prove relative-entropy
control and a joint bad-link estimate under that normalized
interacting conditional.

There is an explicit exceptional set: links where the
conditional mean is large. Its cardinality is bounded, but
its geometry under the incoming low law is not yet
controlled. Outside this set, joint bad events and
connected bad clusters have exponential bounds. This is
not the assertion that all links are good. It is also not
the free heat-shell statement of V56: the present
conditional contains the exact nonlinear Wilson and
geometric weights. The fixed positive carrier and the
sharp high projection remain limitations. The next
companion checkpoint remains V70.

\subsection{A nonextensive source guard for the full gauge operator}

Retain \eqref{eq:v66-parameters}, so
\[
 R=4r\le\tfrac14,\qquad 4C_0R^2\le\tau\le1,\qquad
                  |c|\le\tfrac14,\quad c_*/M\le r/8 .
\]
The fixed carrier \(E_c\) is exactly
\eqref{eq:v66-carrier}. Put \(t=\beta^{-1/2}\), \(\beta\ge1\).

\begin{lemma}[A maximum-norm bound for the even gauge increment]
\label{lem:v67-gauge-increment}
There is a universal \(C_\sigma\) such that, if the
segment from \(a\) to \(a+b\) is Coulomb and has maximum
link radius at most \(R\le1/4\), then
\begin{equation}\label{eq:v67-gauge-increment}
 \|\mathcal F_M(a+b)-\mathcal F_M(a)\|
       \le \|G_1(b)\|+C_\sigma R\|b\|_\infty .
\end{equation}
No gauge floor is needed for this difference estimate.
\end{lemma}
\begin{proof}
The map \(a_e\mapsto\sigma(\operatorname{ad}_{a_e})\)
is a smooth even matrix function on the fixed ball
\(|a_e|\le1/4\), and its first derivative vanishes
at zero. The bounded second derivative on this ball
therefore gives
\[
 \|D[\sigma(\operatorname{ad}_{a_e})][b_e]\|
                                   \le C_\sigma|a_e||b_e|.
\]
Integrate along each link segment, take the maximum
over links, and compress by the isometry \(U\) in
\eqref{eq:v60-gauge}. The odd increment is exactly
\(G_1(b)\) by linearity. This proves the result
without a sum over the volume.
\end{proof}

Define the retained-source set
\begin{equation}\label{eq:v67-guard}
\begin{split}
 \mathcal Y_{M,\beta}
   =\{y\in\mathcal H_<:\;&t\|y\|_\infty\le r/8,\\
       &t(\|G_1(y)\|+C_\sigma R\|y\|_\infty)\le1/16\}.
\end{split}
\end{equation}
It is a convex, symmetric, translation-invariant
neighbourhood of zero in the low space. The operator
norm in this definition is global; we are not
claiming that maximum link size alone controls the
gauge operator. The term \emph{nonextensive} refers to
the absence of a total \(\ell^2\) norm in the guard.

\begin{theorem}[The complete high law extends to the enlarged guard]
\label{thm:v67-enlarged}
For every \(y\in\mathcal Y_{M,\beta}\), define
\(\mu_{c,y}\), \(\mathcal Z_c(y)\), and \(m_c(y)\) by
the literal integral \eqref{eq:v66-law} on \(E_c\).
The integrand is well-defined and positive, and
\begin{equation}\label{eq:v67-enlarged-floor}
 \max_e|a_c(e)+t(y+\eta)_e|<R,\qquad
       \mathcal F_M(a_c+t(y+\eta))\succ\tfrac7{16}I
                         \quad(\eta\in E_c).
\end{equation}
With the same \(A_\ge,K_\ge,\alpha_\beta,B_\beta,T_\beta\)
as V66, using geometric floor \(7/16\), one has
\begin{equation}\label{eq:v67-enlarged-precision}
 D^2_{\eta\eta}U_y\succeq A_\ge\succeq\alpha_\beta I,
       \qquad K_\ge=A_\ge^{-1}\preceq4\mathsf C_M.
\end{equation}
The first and second response identities and bounds
\eqref{eq:v66-response}--\eqref{eq:v66-response-bounds}
hold on this enlarged source set. In particular,
\begin{equation}\label{eq:v67-mean}
 m_c(0)=0,\qquad Dm_c(0)=0,\qquad
          \|m_c(y)\|_2\le b_\beta(\|y\|_2),
\end{equation}
where \(b_\beta\) is defined in
\eqref{eq:v66-mean-bound}. For any \(y,z\) in the guard,
\begin{equation}\label{eq:v67-mean-Lipschitz}
 \|m_c(y)-m_c(z)\|_2
                      \le(B_\beta/\alpha_\beta)\|y-z\|_2 .
\end{equation}
The centered Laplace bound \eqref{eq:v66-Laplace}
also holds. These assertions do not impose
\(t\|y\|_2\le\gamma\).
\end{theorem}
\begin{proof}
For \(\eta\in E_c\), the segment
\(a_c+t\eta+sty\), \(0\le s\le1\), has radius less
than \(2r+r/8<R\). Lemma~\ref{lem:v67-gauge-increment}
and \eqref{eq:v67-guard} change the gauge operator
by at most \(1/16\). Its base floor is \(1/2\), so
\eqref{eq:v67-enlarged-floor} follows.
The bounded positive integral and the high precision
bound follow by exactly the affine G{\aa}rding argument
of V66.

All mixed and third derivatives used in the response
proof depend on the pointwise link radius and this
gauge floor, not on the total norm of the retained
input. They therefore have the same constants.
The indicator of \(E_c\) is still independent of
the input. The normalized covariance proof of V66
applies on the interior of the guard and by continuity
at its boundary. The mean and derivative at zero
vanish by the same disjoint-character argument.
The segment between any two allowed inputs remains
inside the guard by convexity; integrating the
first derivative bound gives
\eqref{eq:v67-mean-Lipschitz}.
Finally strong log-concavity on the fixed convex
carrier gives the centered Laplace comparison.
\end{proof}

Only the source guard and the range of validity of
these high-law estimates have been extended.
V66's whole moving-cutoff coverage theorem is
\emph{not} automatically extended: it additionally
requires the mean to be small everywhere. We replace
that global conclusion below by local bad-set bounds.

\subsection{The guard permits the extensive size of free low fields}

This subsection is a reference-scale audit, not
identification of the interacting low marginal.
Let \(\gamma_<\) be the centered Gaussian on
\(\mathcal H_<\) with covariance
\[
 \mathsf C_<=
       \mathbf1_{\{0<\Delta<\tau\}}\mathsf C_M,
                  \qquad Y\sim\gamma_< .
\]

\begin{proposition}[Reference probability and extensive low energy]
\label{prop:v67-reference}
There is \(K_{\mathrm g}(r)\) such that
\begin{equation}\label{eq:v67-reference-guard}
 \beta\ge K_{\mathrm g}(r)(1+\log M+u)
       \quad\Longrightarrow\quad
       \gamma_<(\mathcal Y_{M,\beta})\ge1-2e^{-u},
                                       \qquad u\ge0 .
\end{equation}
For fixed \(0<\tau\le1\), as \(M\to\infty\),
\begin{equation}\label{eq:v67-reference-energy}
\begin{aligned}
 \frac1V\mathbb E\|Y\|_2^2
       &=\frac9V\sum_{0<\lambda(k)<\tau}\lambda(k)^{-1}
                                                   \longrightarrow e_\tau>0,\\
 e_\tau&=\frac9{(2\pi)^4}
       \int_{\{k\in[-\pi,\pi]^4:\lambda(k)<\tau\}}
                                                \frac{dk}{\lambda(k)},\\
 \operatorname{Var}\!\left(\frac{\|Y\|_2^2}{V}\right)
       &=\frac{18}{V^2}\sum_{0<\lambda(k)<\tau}\lambda(k)^{-2}
                                        \le C(1+\log M)/V .
\end{aligned}
\end{equation}
In particular \(\|Y\|_2^2/V\to e_\tau\) in \(L^2\), and
\begin{equation}\label{eq:v67-low-energy-coefficient}
                   e_\tau=\frac9{16\pi^2}\tau+O(\tau^2)
                                           \quad(\tau\downarrow0).
\end{equation}
For fixed \(\gamma>0\) and \(\beta=o(V)\),
\begin{equation}\label{eq:v67-old-guard-cost}
                \gamma_<\{\|Y\|_2^2\le\gamma^2\beta\}\longrightarrow0 .
\end{equation}
\end{proposition}
\begin{proof}
The covariance \(\mathsf C_<\) is bounded above by
\(\mathsf C_M\). Its link diagonal is uniformly
bounded, giving
\[
 \gamma_<\{\|Y\|_\infty^2>C(1+\log M+u)\}\le e^{-u}
\]
by the scalar Gaussian tail and a union bound.
For the gauge term, the cube covariance matrices of
\(d\Delta^{-1}Y\) are bounded above by the corresponding
matrices in Lemma~\ref{lem:v61-cube}. The same
quadratic Gaussian and deterministic div--curl
argument therefore gives
\[
 \gamma_<\{\|G_1(Y)\|^2>C(1+\log M+u)\}\le e^{-u}.
\]
Increase \(K_{\mathrm g}(r)\) so their common good
event satisfies both parts of \eqref{eq:v67-guard}.

Each nonzero momentum contributes nine real
transverse-colour degrees of freedom when conjugate
momenta are counted in the usual real Fourier
space. Gaussian trace identities give the first
and third expressions in
\eqref{eq:v67-reference-energy}.
Near zero \(\lambda(k)\) is comparable to \(|k|^2\).
Dyadic counting gives a uniformly integrable
\(|k|^{-2}\) singularity for the first Riemann sum,
and at most \(C V(1+\log M)\) for the
\(\lambda^{-2}\) sum. Away from zero ordinary
Riemann convergence applies. For \(0<\tau\le1\)
the level set \(\lambda=\tau\) is smooth and has
zero four-dimensional measure, so the sharp
spectral boundary does not alter that argument.

For the small-\(\tau\) coefficient use
\(\lambda(k)=|k|^2+O(|k|^4)\). Replacing the domain
by \(|k|<\sqrt\tau\) and the integrand by
\(|k|^{-2}\) changes the integral by \(O(\tau^2)\).
The leading integral is \(\pi^2\tau\), which proves
\eqref{eq:v67-low-energy-coefficient}.
The variance bound proves \(L^2\) convergence.
Since the threshold \(\gamma^2\beta/V\) tends to
zero while \(e_\tau>0\), it also proves
\eqref{eq:v67-old-guard-cost}.
\end{proof}

Thus the new guard can have probability tending
to one on a logarithmic coupling schedule even
though the previous global source ball has
probability tending to zero. This comparison is
only about \(\gamma_<\). Replacing the actual low
marginal by this reference is not an authorized
mathematical step.

\subsection{Entropy cost of the exact high conditional}

Write \(D(\mu\Vert\nu)\) for relative entropy,
with natural logarithms.

\begin{theorem}[Dimension-independent conditional entropy response]
\label{thm:v67-entropy}
For any \(y,z\in\mathcal Y_{M,\beta}\),
\begin{equation}\label{eq:v67-entropy}
 \max\{D(\mu_{c,y}\Vert\mu_{c,z}),
                  D(\mu_{c,z}\Vert\mu_{c,y})\}
                  \le\frac{B_\beta^2}{2\alpha_\beta}\|y-z\|_2^2 .
\end{equation}
For any probability law \(\pi\) on the guard, let
\[
 \mathbb P_\pi(dy,d\eta)=\pi(dy)\mu_{c,y}(d\eta),
       \qquad \overline\mu_c=\int\mu_{c,y}\,\pi(dy).
\]
Then
\begin{equation}\label{eq:v67-information}
\begin{aligned}
 I_{\mathbb P_\pi}(Y;\eta)
       +D(\overline\mu_c\Vert\mu_{c,0})
       &=\int D(\mu_{c,y}\Vert\mu_{c,0})\,\pi(dy)\\
       &\le\frac{B_\beta^2}{2\alpha_\beta}
                                      \mathbb E_\pi\|Y\|_2^2 .
\end{aligned}
\end{equation}
\end{theorem}
\begin{proof}
Set \(V_{yz}(\eta)=U_y(\eta)-U_z(\eta)\).
The mixed derivative bound, integrated along
the segment of inputs from \(z\) to \(y\), gives
\[
                 \|\nabla_\eta V_{yz}\|
                                      \le B_\beta\|y-z\|_2.
\]
For \(0\le s\le1\) define
\[
 \Phi(s)=\log\int_{E_c}
                    \exp[-U_z(\eta)-sV_{yz}(\eta)]\,d\eta.
\]
The interpolated potential is a convex combination
of \(U_z\) and \(U_y\), so its high Hessian is
still at least \(\alpha_\beta I\).
The Poincar\'e inequality \eqref{eq:v66-Poincare}
therefore yields
\[
       \Phi''(s)=\operatorname{Var}_s(V_{yz})
                      \le(B_\beta^2/\alpha_\beta)\|y-z\|_2^2.
\]
The normalized entropy identities are
\[
 \begin{aligned}
 D(\mu_{c,y}\Vert\mu_{c,z})
         &=\Phi'(1)-\Phi(1)+\Phi(0)
                                  =\int_0^1s\Phi''(s)\,ds,\\
 D(\mu_{c,z}\Vert\mu_{c,y})
         &=\Phi(1)-\Phi(0)-\Phi'(0)
                                  =\int_0^1(1-s)\Phi''(s)\,ds .
 \end{aligned}
\]
This proves the factor \(1/2\) in
\eqref{eq:v67-entropy}. The common carrier and
its positive densities justify these derivatives;
source-only constants cancel in relative entropy.

The identity in \eqref{eq:v67-information} is the
chain rule obtained by inserting
\(\overline\mu_c\) between \(\mu_{c,y}\) and
\(\mu_{c,0}\) inside the logarithm and integrating.
Apply \eqref{eq:v67-entropy} at \(z=0\).
\end{proof}

The estimate is extensive when the input norm is
extensive; it is not a small total-variation
comparison in a large volume. Dividing by \(V\)
gives an entropy-density bound in terms of
\(\mathbb E_\pi\|Y\|_2^2/V\). No such incoming
moment bound for the full Yang--Mills low marginal
has been assumed to be proved.

\subsection{A sparse mean budget and an interacting joint bad-set bound}

Let \(\mathcal L_M\) be the set of \(4V\) positively
oriented links. For a fixed allowed input \(y\), put
\[
 m=m_c(y),\qquad \xi=\eta-m,\qquad
                          a=a_c+t(y+\eta).
\]
Define the deterministic mean-exceptional set and
the random physical bad-link set by
\begin{equation}\label{eq:v67-bad-definitions}
 \mathcal B_y=\{e:t|m_e|>r/4\},\qquad
                      \mathcal A_y(\eta)=\{e:|a_e|\ge r\}.
\end{equation}
The first set is a function of the full conditional
mean, not an independent random polymer.

Choose a universal \(C_*\) such that for every
finite set \(S\subset\mathcal L_M\), restriction
to its three-colour link coordinates satisfies
\begin{equation}\label{eq:v67-restriction}
 K_S=\operatorname{res}_S K_\ge\operatorname{res}_S^*,
       \qquad \operatorname{Tr}K_S\le C_*|S|,
       \qquad \|K_S\|\le\alpha_\beta^{-1}.
\end{equation}
Such a constant exists: \(K_\ge\preceq4\mathsf C_M\)
has a uniformly bounded link diagonal, and
\(K_\ge\preceq\alpha_\beta^{-1}I\).

\begin{theorem}[Joint bad-link suppression outside the explicit mean set]
\label{thm:v67-badsets}
Uniformly over the enlarged guard,
\begin{equation}\label{eq:v67-mean-budget}
 |\mathcal B_y|
     \le\frac{16\,b_\beta(\|y\|_2)^2}{\beta r^2}
     \le\frac{16B_\beta^2}{\alpha_\beta^2\beta r^2}\|y\|_2^2 .
\end{equation}
If \(\beta r^2\ge16C_*\), set
\begin{equation}\label{eq:v67-activity}
             a_\beta=\alpha_\beta\beta r^2/32,
                         \qquad p_\beta=e^{-a_\beta}.
\end{equation}
Then every link set \(S\) obeys
\begin{equation}\label{eq:v67-bad-product}
 \mu_{c,y}\{S\subset\mathcal A_y\}
                            \le p_\beta^{\,|S\setminus\mathcal B_y|}.
\end{equation}
The probability is that of the complete normalized
nonlinear Wilson conditional on \(E_c\).
\end{theorem}
\begin{proof}
Every exceptional link contributes more than
\(\beta r^2/16\) to \(\|m\|_2^2\). Sum over
these links and use \eqref{eq:v67-mean}.

The deterministic retained background obeys
\(\max_e|a_c(e)+ty_e|\le r/4\).
At a nonexceptional link, \(t|m_e|\le r/4\).
Thus a bad link outside \(\mathcal B_y\) must have
\[
                             |\xi_e|\ge\sqrt\beta\,r/2 .
\]
For a set \(S\) disjoint from \(\mathcal B_y\), the
event that every link is bad implies
\[
             \|\operatorname{res}_S\xi\|_2^2
                                      \ge |S|\beta r^2/4 .
\]
The centered Laplace comparison, followed by an
independent auxiliary Gaussian integration as in
\eqref{eq:v65-quadratic}, gives
\[
 \mathbb E_{\mu_{c,y}}
           e^{s\|\operatorname{res}_S\xi\|^2}
                       \le\det(I-2sK_S)^{-1/2}.
\]
This comparison requires neither independent
interacting coordinates nor equality of their
covariance with \(K_S\). At \(s=\alpha_\beta/4\),
\(2s\|K_S\|\le1/2\). The scalar inequality
\(-\log(1-x)\le2x\) gives
\[
 \log\det(I-2sK_S)^{-1/2}
                    \le2s\operatorname{Tr}K_S
                    \le(C_*\alpha_\beta/2)|S|.
\]
Exponential Markov therefore bounds the event by
\[
 \exp\!\left[
       -\alpha_\beta\left(\frac{\beta r^2}{16}
                                    -\frac{C_*}{2}\right)|S|
        \right]
                      \le e^{-a_\beta|S|}.
\]
The last step uses \(\beta r^2\ge16C_*\).
For a general set \(S\), discard its exceptional
links from the event to obtain
\eqref{eq:v67-bad-product}.
\end{proof}

The auxiliary Gaussian argument for quadratic
forms under a linear-source Laplace bound is also
recorded in Hsu--Kakade--Zhang,
\emph{A tail inequality for quadratic forms of
subgaussian random vectors}, Remark 2.3,
\href{https://www.cs.columbia.edu/~djhsu/papers/quadratic-ecp.pdf}
{Electronic Communications in Probability 17 (2012), paper 52}.
Here the comparison matrix, normalization and
mean-exceptional set have all been specified
for the Wilson conditional itself.

\subsection{Counts and connected clusters, with the exceptional set retained}

Throughout the remaining probability estimates assume
\(\beta r^2\ge16C_*\), as in Theorem~\ref{thm:v67-badsets}.
Let \(N_y=|\mathcal L_M\setminus\mathcal B_y|\) and
\(N_{\rm bad}=|\mathcal A_y\setminus\mathcal B_y|\).

\begin{proposition}[A binomial generating bound without independence]
\label{prop:v67-counts}
For \(\theta\ge0\) and \(1\le n\le N_y\),
\begin{equation}\label{eq:v67-counts}
\begin{aligned}
 \mathbb E_{\mu_{c,y}}e^{\theta N_{\rm bad}}
          &\le[1+(e^\theta-1)p_\beta]^{N_y},\\
 \mu_{c,y}\{N_{\rm bad}\ge n\}
          &\le \binom{N_y}{n}p_\beta^n
                       \le(eN_yp_\beta/n)^n .
\end{aligned}
\end{equation}
In particular,
\begin{equation}\label{eq:v67-density}
 \mathbb E_{\mu_{c,y}}\frac{|\mathcal A_y|}{4V}
       \le\frac{4B_\beta^2}{\alpha_\beta^2\beta r^2}
                                     \frac{\|y\|_2^2}{V}+p_\beta .
\end{equation}
\end{proposition}
\begin{proof}
For the indicators \(I_e=\mathbf1_{\{e\in\mathcal A_y\}}\)
on the nonexceptional links, expand the finite product
\[
 e^{\theta N_{\rm bad}}
        =\prod_{e\notin\mathcal B_y}[1+(e^\theta-1)I_e].
\]
All coefficients are nonnegative. Apply
\eqref{eq:v67-bad-product} to each joint indicator,
and sum by cardinality. A union bound over
the \(n\)-element subsets proves the count tail.
The expectation is at most \(N_yp_\beta\);
add the cardinality bound for \(\mathcal B_y\)
and divide by \(4V\) to obtain
\eqref{eq:v67-density}.
\end{proof}

Give \(\mathcal L_M\) the adjacency in which distinct
links are neighbours when they share a plaquette.
A link belongs to six plaquettes, each with three
other links, so its degree is at most eighteen.
For \(e\notin\mathcal B_y\), let
\(\mathcal C_y(e)\) be its connected component in
the graph of bad links after the vertices
\(\mathcal B_y\) have been removed. If \(e\) is
not bad, this component is empty.

\begin{proposition}[Conditional bad clusters away from the mean set]
\label{prop:v67-clusters}
For \(n\ge1\),
\begin{equation}\label{eq:v67-clusters}
 \mu_{c,y}\{|\mathcal C_y(e)|\ge n\}
             \le18^{\,2(n-1)}e^{-a_\beta n}.
\end{equation}
In particular this is a geometric cluster tail
if \(a_\beta>2\log18\).
\end{proposition}
\begin{proof}
A rooted connected set of \(n\) links can be
encoded by the depth-first walk of a deterministically
chosen spanning tree, of length \(2(n-1)\).
The visited set recovers the original set, so
this is an injective encoding into at most
\(18^{2(n-1)}\) walks. Any component of size at
least \(n\) contains a rooted connected subset
of size exactly \(n\). Its links are all outside
\(\mathcal B_y\), and
\eqref{eq:v67-bad-product} bounds the joint event
by \(e^{-a_\beta n}\). Take the union bound.
\end{proof}

This proposition does not bound a component
allowed to travel through \(\mathcal B_y\).
Its geometry cannot be erased by calling its
cardinality small.

\subsection{What an incoming low law must still provide}

For any \(\pi\) supported on the guard, set
\[
                         \rho_\pi=V^{-1}\mathbb E_\pi\|Y\|_2^2 .
\]
The actual conditional results integrate to
\begin{equation}\label{eq:v67-integrated}
\begin{aligned}
 V^{-1}I_{\mathbb P_\pi}(Y;\eta)
                   &\le (B_\beta^2/(2\alpha_\beta))\rho_\pi,\\
 \mathbb E_{\mathbb P_\pi}\frac{|\mathcal A_Y|}{4V}
                   &\le\frac{4B_\beta^2}{\alpha_\beta^2\beta r^2}
                                                      \rho_\pi+p_\beta .
\end{aligned}
\end{equation}
These hold for any such incoming law, not only
a Gaussian one. But a bounded \(\rho_\pi\) controls
an expected fraction, not all spatial joint events.
More precisely, \eqref{eq:v67-bad-product} gives
\begin{equation}\label{eq:v67-overlap}
 \mathbb P_\pi\{S\subset\mathcal A_Y\}
       \le e^{-a_\beta|S|}
                    \mathbb E_\pi e^{a_\beta|S\cap\mathcal B_Y|}.
\end{equation}
Thus the missing local input can be stated
without an extensive global denominator: it is
an estimate for the exponential overlap of
the \emph{actual conditional-mean set}
\(\mathcal B_Y\) with arbitrary test sets \(S\).
For example, a proved bound
\[
             \mathbb E_\pi e^{a_\beta|S\cap\mathcal B_Y|}
                                          \le e^{b|S|},\qquad b<a_\beta,
\]
would yield an unconditional joint bad-set
exponent \(a_\beta-b\). No such inequality
for the actual Yang--Mills incoming law is
claimed here. The cardinality estimate
\eqref{eq:v67-mean-budget} alone cannot prove it:
a deterministic exceptional set may contain
a long connected region while occupying a
small fraction of the torus.

There are also two separate unpaid measure
questions. The actual incoming low marginal
has not been shown to lie in
\(\mathcal Y_{M,\beta}\) with high probability.
And the high carrier \(E_c\) remains restricted
by a quantitative positive gauge floor.
The estimates do not give its weight inside
the unrestricted positive branch or the full
compact measure. The free calculation above
is not a substitute for either question.

This moves the earlier obstruction from a
global \(\ell^2\) source ball to an explicit
spatial conditional-mean problem. A next useful
estimate must relate that set locally to the
retained field or control its overlap under
the true low marginal. No all-scale local
RG map, reflection-positive continuum state,
or physical Hamiltonian gap follows yet.

\subsection{Exact finite checks and preservation}

The certificate
\begin{center}\small
\path{scripts/verify_ym_f14_v67_conditional_bad_links.py}
\end{center}
uses actual side-four Fourier covariances at
the algebraic test values \(\beta=7,\tau=3\).
The low real dimension is \(72\); its free
mean norm square per volume and the variance
of that quantity are
\[
                              9/64,\qquad 9/16384.
\]
For three selected links it verifies exact
positivity of the scalar colour block of the restricted high comparison,
of \(I-\alpha_\beta K_S\), and of the
Gaussian-source determinant matrix, by all
principal minors. The three-colour operator is its tensor
product with \(I_3\). The comparison floor is
\(71/84\), and its three-colour selected trace is
\[
                         1338772502223891/399255164239360.
\]
The threshold \(3\) is only a finite algebra
regression and is not asserted to satisfy
the theorem's \(\tau\le1\) window.

On the actual side-two and side-four
plaquette graphs the maximum link degrees
are \(15\) and \(18\). The rooted connected
set counts at sizes one, two and three are
respectively
\[
                         (1,15,231),\qquad (1,18,375).
\]
Every link is incident to six plaquettes.
These checks certify the local combinatorics
used by the analytic counting argument.

A dependent finite-set model with fifteen
outcomes checks all 256 joint subset bounds
with two fixed exceptional links. Its
nonexceptional count generating function at
two is \(4\), while the binomial majorant
is \(4096/729\). It is a combinatorial
regression, not a Wilson distribution.
A separate rational quadratic model checks
the entropy factor \(1/2\); its linear-tilt
relative entropy is \(189/58\).
Neither model proves a uniform Wilson
probability bound by numerical experiment;
those bounds were proved above.

The 7,690-byte certificate has SHA--256
\begin{center}\scriptsize\ttfamily
46DB2D86B208FCD96F33FF908C1BCFCE41E79DD4A4CCA7BB10B288C8852113A7.
\end{center}
It uses unchanged V56, V58 and V62 helpers.
The predecessor V66 has 1,632,189 bytes
and SHA--256
\begin{center}\scriptsize\ttfamily
5F54EFA116CE67B7AD6E292C8CD6F6BF3BD982C784E1A652194D52FF6B03479F.
\end{center}
Its body is retained apart from the title
version and this terminal insertion. V67
replaces V66 as the sole active main note.
Sources, companions, monographs and archives
are unchanged.

\begin{firewall}
Extensive retained inputs are now allowed
under a nonextensive amplitude-and-gauge
guard. The complete high conditional has
entropy and joint bad-link estimates, with
the mean-exceptional set retained explicitly.
Its spatial distribution under the true
low marginal, the larger-domain comparison,
and the physical continuum Yang--Mills
mass gap remain unproved. The full goal
remains active.
\end{firewall}
\section{V68: a localized massive-curl extraction of the exact Wilson law}
\label{sec:v68}

\subsection{Going upstream from the mean-exceptional set}

V67 proves an interacting bad-link bound, but leaves the
spatial distribution of the conditional mean unresolved.
The sharp spectral projection used there is itself
nonlocal. This section replaces that projection by a
smooth Gaussian extraction with an explicitly localized
transverse factor. The change is an exact representation
of the same restricted coordinate law, not a replacement
of the Wilson action by a massive free action.

V56 already supplied positive smooth heat covariances
and free bad-block estimates. We do not repeat that
decomposition. The new object is a rational massive-curl
factor, together with an exact nonlinear conditional
coercivity theorem in its independent two-form noise
coordinates. The retained marginal keeps the full
normalizer and moving domain. Its lower-form sign
will show explicitly why the extraction parameter
is not a Yang--Mills mass gap.

In this section the retained and extracted Gaussian
fields have overlapping Fourier support. They are
\emph{not} the disjoint low and high spaces of
V64--V67. Accordingly, the zero linear response
proved there by disjoint characters is not imported.
The next companion checkpoint remains V70.

\subsection{An exact positive split with a local transverse numerator}

Let \(\mathcal H_M\) remain the real mean-zero
Coulomb link space, with free covariance
\(\mathsf C_M=\Delta^{-1}\) on that space, and
dimension \(9(V-1)\). On the full link space
its symbol is \(\Pi(k)/\lambda(k)\), with the
zero mode set to zero. Let
\[
              \mathcal X_M=\mathbb R^{6V}\otimes\mathbb R^3
\]
be the space of all real two-form coordinates,
with its usual Euclidean inner product. Its
dimension is \(18V\). All occurrences of
\(\Delta\) on forms below mean the componentwise
lattice Laplacian. For \(0<m\le1\), define
\begin{equation}\label{eq:v68-extraction}
 \mathsf T_m=d^*(\Delta+m^2)^{-1}:\mathcal X_M\longrightarrow\mathcal H_M,
              \qquad \mathsf C_m^{\mathrm h}=\mathsf T_m\mathsf T_m^* .
\end{equation}
The massive inverse here is unprimed. Its
constant-mode part is annihilated by \(d^*\).

\begin{lemma}[Positive exact Gaussian extraction]
\label{lem:v68-split}
The map \(\mathsf T_m\) is onto \(\mathcal H_M\), and
\begin{equation}\label{eq:v68-split}
\begin{aligned}
 \mathsf C_m^{\mathrm h}(k)
     &=\frac{\lambda(k)}{(\lambda(k)+m^2)^2}\Pi(k)\otimes I_3,\\
 \mathsf C_m^\ell(k)
     &:=\mathsf C_M(k)-\mathsf C_m^{\mathrm h}(k)\\
     &=\frac{m^2(2\lambda(k)+m^2)}
             {\lambda(k)(\lambda(k)+m^2)^2}\Pi(k)\otimes I_3
                                            \qquad(k\ne0).
\end{aligned}
\end{equation}
Both covariances vanish at zero and are positive
definite on \(\mathcal H_M\).
For independent normalized Gaussians
\[
 Y\sim N(0,\mathsf C_m^\ell),\qquad
              \Xi\sim N(0,I_{\mathcal X_M}),
\]
the sum \(Z=Y+\mathsf T_m\Xi\) has exactly the
original law \(N(0,\mathsf C_M)\). Moreover,
\begin{equation}\label{eq:v68-norms}
             \|\mathsf T_m\|\le(2m)^{-1},\qquad
                       \|\Delta^{1/2}\mathsf T_m\|\le1.
\end{equation}
\end{lemma}
\begin{proof}
The Hodge identity gives \(d^*d=\Delta\Pi\)
on links, and \(D^*d^*=0\). Thus the range
of \(d^*\) is mean-zero and Coulomb. For
\(h\in\mathcal H_M\), the two-form
\[
                     (\Delta+m^2)\Delta^{-1}dh
\]
maps to \(h\) under \(\mathsf T_m\), proving
surjectivity. The same identity gives the
first symbol in \eqref{eq:v68-split}.
Subtract it from \(1/\lambda\) on each
transverse mode to obtain the displayed
strictly positive remainder.
Covariances of independent Gaussians add,
including their normalized laws.
The singular values of \(\mathsf T_m\)
on its active Fourier components are
\(\sqrt\lambda/(\lambda+m^2)\). Their squares
are at most \(1/(4m^2)\), and multiplication
by \(\lambda\) gives a number at most one.
\end{proof}

The extracted covariance has the alternative
full-space symbol
\begin{equation}\label{eq:v68-local-numerator}
           \mathsf C_m^{\mathrm h}(k)
              =\frac{\lambda(k)I-q(k)q(k)^*}
                     {(\lambda(k)+m^2)^2}\otimes I_3 .
\end{equation}
The transverse numerator is a local difference
operator. In particular there is no residual
\(1/\lambda\) inside a Coulomb projector in
this formula. A naive choice
\((\Delta+m^2)^{-1}\Pi\) would not have this
cancellation at zero momentum.

\subsection{Uniform exponential locality of the extraction factor}

Use the periodic nearest-neighbour graph distance
between base sites of form coordinates. For sets
of such coordinates \(A,B\), denote this distance
by \(\operatorname{dist}_M(A,B)\).

\begin{theorem}[Localized extraction in operator norm]
\label{thm:v68-locality}
Uniformly in \(M\ge4\), \(0<m\le1\), and coordinate sets,
\begin{equation}\label{eq:v68-locality}
\begin{aligned}
 \|\mathbf1_A\mathsf T_m\mathbf1_B\|_{2\to2}
       &\le\frac4m\,e^{-m\operatorname{dist}_M(A,B)/4},\\
 \|\mathbf1_A\mathsf C_m^{\mathrm h}\mathbf1_B\|_{2\to2}
       &\le\frac{16}{m^2}\,e^{-m\operatorname{dist}_M(A,B)/4}.
\end{aligned}
\end{equation}
These are statements about the extraction operator
and its Gaussian covariance, not yet about the
off-diagonal covariance of the interacting law.
\end{theorem}
\begin{proof}
Empty coordinate sets make the conclusions immediate.
Let \(f\) be a real function on the torus with
\(|f(x+e_\mu)-f(x)|\le1\), and let \(W\) multiply
every component based at \(x\) by \(e^{\omega f(x)}\).
The same weight is used on every form degree.
Put \(\omega=m/4\) and
\(s_{x\mu}=\omega(f(x)-f(x+e_\mu))\).
Summing over positively oriented nearest-neighbour
edges, the exact real-part identity is
\begin{equation}\label{eq:v68-weighted-energy}
\begin{aligned}
 \operatorname{Re}\langle u,W\Delta W^{-1}u\rangle
   &=\|\nabla u\|_2^2\\
   &\quad-2\sum_{x,\mu}(\cosh s_{x\mu}-1)
                    \operatorname{Re}\langle u(x),u(x+e_\mu)\rangle .
\end{aligned}
\end{equation}
Using \(2|\langle u(x),u(y)\rangle|
\le|u(x)|^2+|u(y)|^2\), the correction is bounded
below by \(-8(\cosh\omega-1)\|u\|_2^2\).
For \(\omega\le1\), \(\cosh\omega-1\le\omega^2\).
Consequently
\begin{equation}\label{eq:v68-accretivity}
 \operatorname{Re}\langle u,W(\Delta+m^2)W^{-1}u\rangle
             \ge\|\nabla u\|_2^2+(m^2/2)\|u\|_2^2 .
\end{equation}
If \(u=[W(\Delta+m^2)W^{-1}]^{-1}g\), this gives
\[
                  \|u\|_2\le2m^{-2}\|g\|_2,\qquad
                  \|\nabla u\|_2\le\sqrt2\,m^{-1}\|g\|_2 .
\]

Conjugation changes \(d^*\) only in its shifted
terms. Each output link has three such two-form
terms, and each input two-form contributes to
two output orientations. Their coefficients
change by at most \(e^\omega-1\le2\omega\).
Thus
\[
              \|Wd^*W^{-1}-d^*\|
                                  \le2\sqrt6\,\omega .
\]
The Fourier Hodge inequality
\(\|d^*u\|_2\le\|\nabla u\|_2\) now gives
\[
 \|W\mathsf T_m W^{-1}g\|_2
       \le\left(\frac{\sqrt2}{m}
                    +\frac{4\sqrt6\,\omega}{m^2}\right)\|g\|_2
       <\frac4m\|g\|_2 .
\]
The same proof holds for \(-f\). Taking an
adjoint therefore also bounds
\(\|W\mathsf T_m^*W^{-1}\|\) by \(4/m\).
Multiplication yields the bound \(16/m^2\)
for \(W\mathsf C_m^{\mathrm h}W^{-1}\).

Finally take \(f(x)\) to be the periodic
distance from the base sites of \(B\).
It is zero on \(B\), at least
\(\operatorname{dist}_M(A,B)\) on \(A\), and
has the required nearest-neighbour Lipschitz
constant. Removing the weights proves
\eqref{eq:v68-locality}.
\end{proof}

This is a finite weighted-conjugation argument.
The general transformation approach to exponential
localization has classical background in
Combes--Thomas,
\href{https://link.springer.com/article/10.1007/BF01646473}
{Communications in Mathematical Physics 34 (1973), 251--270}.
The torus constants used above are proved directly;
none is imported from a continuum resolvent estimate.

\subsection{The lifted measure is exactly the same local Wilson measure}

Fix
\begin{equation}\label{eq:v68-parameters}
 0<R\le\tfrac14,\quad M\ge4,\quad |c|\le c_*=\tfrac14,
       \quad c_*/M\le R/4,\quad
             \beta\ge1,\quad C_0R^2\le m^2\le1 .
\end{equation}
Here \(a_c(x,\mu)=c_\mu/M\), as before.
No commutativity of the \(c_\mu\) is assumed.
Define the total-field domain in \(\mathcal H_M\)
\begin{equation}\label{eq:v68-domain}
 \mathcal D_c=\{z:\max_e|a_c(e)+tz_e|<R,\
                   \mathcal F_M(a_c+tz)\succ0\}.
\end{equation}
Only positivity is imposed, not a quantitative
gauge floor. By V62 and its affine extension,
this is an open convex bounded set. The base
estimate of V65 places \(z=0\) strictly inside it.
Let
\begin{equation}\label{eq:v68-interaction}
 \mathcal R_c(z)=\beta S(a_c+tz)+W_M(a_c+tz)
                                        -\tfrac12\langle z,\Delta z\rangle
                  \quad(z\in\mathcal D_c).
\end{equation}
Denote by \(\gamma_m^\ell\) the normalized retained
Gaussian and by \(\gamma_{\mathcal X}\) the standard
two-form Gaussian.

\begin{theorem}[Exact normalized lift, including its retained weight]
\label{thm:v68-lift}
The probability law
\begin{equation}\label{eq:v68-joint}
\begin{split}
 d\mathbb P_{c,m}(y,\xi)
       =\mathfrak Z_c^{-1}
          \mathbf1_{\mathcal D_c}(y+\mathsf T_m\xi)
          e^{-\mathcal R_c(y+\mathsf T_m\xi)}
             d\gamma_m^\ell(y)\,d\gamma_{\mathcal X}(\xi)
\end{split}
\end{equation}
pushes forward under \(z=y+\mathsf T_m\xi\) to
\begin{equation}\label{eq:v68-target}
 d\mu_c(z)=
    \mathfrak Z_c^{-1}\mathbf1_{\mathcal D_c}(z)
                           e^{-\mathcal R_c(z)}d\gamma_M(z).
\end{equation}
In Lebesgue coordinates this is precisely the
normalized weight
\(\mathbf1_{\mathcal D_c}(z)\rho_M(a_c+tz)
e^{-\beta S(a_c+tz)}\,dz\).
Its Gaussian-reference normalizer is
\[
      \mathfrak Z_c=
        \mathbb E_{\gamma_M}
              [\mathbf1_{\mathcal D_c}e^{-\mathcal R_c}],
\]
independent of the extraction parameter \(m\).
The actual retained marginal is
\begin{equation}\label{eq:v68-retained}
\begin{aligned}
 d\pi_{c,m}(y)&=\mathfrak Z_c^{-1}Z_{c,m}(y)\,d\gamma_m^\ell(y),\\
 Z_{c,m}(y)&=\int
       \mathbf1_{\mathcal D_c}(y+\mathsf T_m\xi)
             e^{-\mathcal R_c(y+\mathsf T_m\xi)}d\gamma_{\mathcal X}(\xi).
\end{aligned}
\end{equation}
For every \(y\in\mathcal H_M\), \(Z_{c,m}(y)\) is
finite and strictly positive.
\end{theorem}
\begin{proof}
Lemma~\ref{lem:v68-split} identifies the law of the
sum under the product Gaussian. Multiplying any
bounded test function of that sum by the same
nonnegative function
\(\mathbf1_{\mathcal D_c}e^{-\mathcal R_c}\) proves
the pushforward identity and equality of normalizers.
The Gaussian density cancels the subtracted
quadratic term in \eqref{eq:v68-interaction},
giving the full Lebesgue-coordinate weight.

The total-field domain is bounded. The original
finite-dimensional Wilson, Haar and determinant
factors make the nonnegative reweighting function
bounded above there; at the gauge boundary the
determinant may vanish, not diverge.
It is strictly positive near zero. Thus its
Gaussian integral is finite and positive.
For every fixed \(y\), surjectivity of
\(\mathsf T_m\) gives a noise vector with
\(y+\mathsf T_m\xi=0\). A neighbourhood maps
inside \(\mathcal D_c\), so the conditional
normalizer is positive. Its finiteness follows
from the same upper bound on the reweighting
function. Integrating first in \(\xi\) proves
\eqref{eq:v68-retained}.
\end{proof}

The new noise variables are auxiliary Gaussian
coordinates, not extra physical fields. Indeed
\[
 \dim\ker\mathsf T_m=18V-9(V-1)=9V+9.
\]
Decomposing \(\mathcal X_M\) orthogonally into
this kernel and its complement, the interaction
and domain depend only on the complement.
The kernel component remains an independent
standard Gaussian in every conditional and in
the joint law. Its normalized integral is one.
There is no omitted source-dependent determinant
from these redundant coordinates.

\subsection{Uniform nonlinear coercivity for every retained field}

For arbitrary \(y\in\mathcal H_M\) let
\[
 \mathcal E_y=\{\xi\in\mathcal X_M:
                            y+\mathsf T_m\xi\in\mathcal D_c\}.
\]
It is a nonempty open convex cylinder. On it
define the full conditional potential
\begin{equation}\label{eq:v68-conditional-potential}
             V_y(\xi)=\tfrac12\|\xi\|_2^2
                                +\mathcal R_c(y+\mathsf T_m\xi).
\end{equation}
The actual conditional \(\nu_{c,m,y}\) has
Lebesgue density proportional to
\(\mathbf1_{\mathcal E_y}e^{-V_y}\).

\begin{theorem}[The complete massive-curl conditional has a uniform floor]
\label{thm:v68-conditional}
Put \(g_\beta=2/(3\beta)\). For every retained
field \(y\), without a source-size guard,
\begin{equation}\label{eq:v68-noise-floor}
 D_\xi^2V_y\succeq
          A_m:=\tfrac12I_{\mathcal X_M}
                        +g_\beta\mathsf T_m^*\mathsf T_m
                  \succeq\tfrac12I_{\mathcal X_M}.
\end{equation}
If \(\bar\xi_y=\mathbb E_{\nu_{c,m,y}}\Xi\), then
\begin{equation}\label{eq:v68-noise-Laplace}
 \mathbb E_{\nu_{c,m,y}}
        e^{\langle J,\Xi-\bar\xi_y\rangle}
                          \le e^{\langle J,A_m^{-1}J\rangle/2}.
\end{equation}
The physical extracted fluctuation
\(\eta=\mathsf T_m\Xi\), centered at
\(\bar\eta_y=\mathsf T_m\bar\xi_y\), has comparison
\begin{equation}\label{eq:v68-extracted-comparison}
\begin{aligned}
 K_m&=\mathsf T_m A_m^{-1}\mathsf T_m^*
       =\left\{\tfrac12(\mathsf C_m^{\mathrm h})^{-1}
                                      +g_\beta I\right\}^{-1}
                                  \quad\hbox{on }\mathcal H_M,\\
 K_m&\preceq2\mathsf C_m^{\mathrm h}\preceq2\mathsf C_M,\\
 K_m(k)&=\frac{\lambda(k)}
              {(\lambda(k)+m^2)^2/2+g_\beta\lambda(k)}
                                                  \Pi(k)\otimes I_3 .
\end{aligned}
\end{equation}
Also
\begin{equation}\label{eq:v68-noise-Poincare}
 \operatorname{Var}_{\nu_{c,m,y}}F
                          \le2\mathbb E_{\nu_{c,m,y}}\|\nabla_\xi F\|^2 .
\end{equation}
\end{theorem}
\begin{proof}
On \(\mathcal D_c\), the affine full Wilson
G{\aa}rding estimate and the geometric Hessian
bound give
\[
 D^2\mathcal R_c(z)
                   \succeq-\tfrac12\Delta
                                      +(g_\beta-C_0R^2)I .
\]
Set \(k_R=C_0R^2\le m^2\). Thus
\[
 D_\xi^2V_y
   \succeq I-\mathsf T_m^*(\tfrac12\Delta+k_R I)\mathsf T_m
                                     +g_\beta\mathsf T_m^*\mathsf T_m .
\]
On an active two-form mode with Laplacian
eigenvalue \(\lambda\), the subtracted operator
has eigenvalue at most
\[
 \frac{\lambda(\lambda/2+m^2)}{(\lambda+m^2)^2}
            =\frac12-\frac{m^4}{2(\lambda+m^2)^2}\le\frac12.
\]
On the kernel of \(\mathsf T_m\) it is zero.
This proves \eqref{eq:v68-noise-floor} for
the complete nonlinear potential, not merely
at its minimizer or at zero.

Subtract the \(A_m\) quadratic form and extend
the remainder by infinity outside the actual
convex cylinder. Strong log-concavity gives
the centered Laplace comparison by the
Gaussian-convolution argument of V64.
One can factor off the independent Gaussian
kernel first; the remaining active fibre is
bounded at each fixed \(M,m,\beta,y\).
Thus all centering and differentiations are
justified as in the earlier bounded-fibre proof.
The convex-support Poincar\'e inequality of
V66 gives \eqref{eq:v68-noise-Poincare}.

Apply the Laplace bound with
\(J=\mathsf T_m^*h\). Diagonalizing
\(\mathsf T_m^*\mathsf T_m\) on its active
range proves all formulas in
\eqref{eq:v68-extracted-comparison}.
\end{proof}

At \(y=0\) the physical extracted mean is exactly
zero: translation invariance makes it constant,
while it lies in \(\mathcal H_M\), whose constant
mode is absent. For arbitrary \(y\), the mean
is retained in the estimates. In particular,
uniform centered concentration is not a proof
of a uniformly small conditional mean.

The localized comparison covariance must not
be confused with entrywise covariance control
of the interacting field. Matrix order alone
does not give that implication. For example,
\[
 0\preceq
 \begin{pmatrix}1/2&1/4\\1/4&1/2\end{pmatrix}
 \preceq I_2
\]
has a nonzero off-diagonal entry while its
upper comparison \(I_2\) has none.
An actual spatial response or covariance
estimate needs further work.

\subsection{The retained marginal keeps the soft-sector obstruction}

The exact marginal in \eqref{eq:v68-retained}
has, up to a constant, Lebesgue potential
\[
 \mathcal V_{c,m}(y)=
       \tfrac12\langle y,(\mathsf C_m^\ell)^{-1}y\rangle
                                      -\log Z_{c,m}(y).
\]
This includes its complete high normalizer
and the moving domain.
For a nonzero Laplacian eigenvalue write
\begin{equation}\label{eq:v68-schur-data}
 h_m(\lambda)=\frac{\lambda}{(\lambda+m^2)^2},\qquad
 \ell_m(\lambda)=\frac1\lambda-h_m(\lambda),\qquad
 b_R(\lambda)=-\lambda/2+g_\beta-k_R .
\end{equation}

\begin{theorem}[A boundary-retaining lower form for the true retained law]
\label{thm:v68-marginal}
On \(\mathcal H_M\), define the Fourier multiplier
\begin{equation}\label{eq:v68-retained-floor}
\begin{aligned}
 L_m(\lambda)
     &=\ell_m(\lambda)^{-1}
                    +\frac{b_R(\lambda)}{1+h_m(\lambda)b_R(\lambda)}\\
     &=\frac{\lambda/2+g_\beta-k_R}
        {\lambda\,\ell_m(\lambda)
                       [1+h_m(\lambda)b_R(\lambda)]}.
\end{aligned}
\end{equation}
Then
\begin{equation}\label{eq:v68-marginal-convex}
            \mathcal V_{c,m}(y)-\tfrac12\langle y,L_m y\rangle
                                          \quad\hbox{is convex}.
\end{equation}
In particular its distributional Hessian is
bounded below by \(L_m\). The denominator in
\eqref{eq:v68-retained-floor} is positive.
The sign of this proved lower multiplier is
exactly the sign of
\begin{equation}\label{eq:v68-soft-sign}
                            \lambda/2+g_\beta-k_R ,
\end{equation}
independent of the extraction parameter \(m\).
\end{theorem}
\begin{proof}
Let \(B_R=-\Delta/2+(g_\beta-k_R)I\)
on the total-field space. The full joint
potential in Lebesgue \((y,\xi)\) coordinates
has Hessian bounded below by the constant
block matrix
\[
 \begin{pmatrix}
   (\mathsf C_m^\ell)^{-1}+B_R & B_R\mathsf T_m\\
   \mathsf T_m^*B_R & I+\mathsf T_m^*B_R\mathsf T_m
 \end{pmatrix}.
\]
Its lower-right block is at least \(I/2\)
by the same calculation as in
Theorem~\ref{thm:v68-conditional}.
Take its Schur complement. Since all total-field
multipliers involved are functions of \(\Delta\),
the resulting scalar multiplier is
\[
 \ell_m^{-1}+b_R-\frac{h_m b_R^2}{1+h_m b_R}
                         =\ell_m^{-1}+\frac{b_R}{1+h_m b_R}=L_m .
\]
Subtracting \(\operatorname{diag}(L_m,0)\)
therefore makes this block matrix positive
semidefinite. The exact joint potential
minus \(\langle y,L_my\rangle/2\) is convex
on the joint domain.

That domain is the preimage of the convex
\(\mathcal D_c\) under \(y+\mathsf T_m\xi\).
Include its indicator by extending the
potential by infinity. Pr\'ekopa marginal
log-concavity, in the same form used in V64,
then gives \eqref{eq:v68-marginal-convex}
by integration over the fixed whole space
\(\mathcal X_M\). This does not differentiate
or discard a moving boundary.

Finally \(\ell_m>0\) and
\(1+h_mb_R\ge1/2\). Using
\(\ell_m+h_m=1/\lambda\) gives the second
formula in \eqref{eq:v68-retained-floor},
and hence the sign assertion.
\end{proof}

When \(\lambda<2(k_R-g_\beta)\), this lower
form is negative. That is a limitation of
the lower bound, not a proof that the actual
retained Hessian has negative curvature.
In either case a positive physical infrared
floor has not been manufactured by the
massive extraction.

For an independent check of the interpretation,
turn off the interaction and the domain
restriction in the Gaussian representation.
The noise conditional is then exactly
standard Gaussian, with Poincar\'e constant
one for every \(M\). Yet the recombined field
has covariance
\(\mathsf C_m^\ell+\mathsf C_m^{\mathrm h}
=\mathsf C_M\), with inverse soft eigenvalue
\(4\sin^2(\pi/M)\to0\).
The parameter \(m\) disappears from that
recombined free law. A uniform auxiliary
noise floor is therefore not a physical
mass-gap statement.

\subsection{The new interface and the remaining inputs}

The acquisition is an exponentially localized
transverse extraction map, an exact normalized
lift of the full restricted Wilson law,
and a uniform nonlinear conditional floor
for \emph{every} retained value \(y\).
There is no hard Fourier projection and no
retained-source size guard in this conditional
coercivity theorem. Its domain is the actual
positive total-field domain rather than the
fixed quantitative-floor carrier of V67.

The unresolved requirements are not removed.
\begin{itemize}
\item The retained marginal is
\(Z_{c,m}(y)\,d\gamma_m^\ell(y)/\mathfrak Z_c\),
not the retained Gaussian alone. Its normalizer,
soft modes and induced interactions require
control. The preceding theorem gives only
the stated semiconvex lower form.
\item The fibre moves with \(y\). A source-response
formula must retain that boundary or replace it
with a proved comparison. V66's fixed-carrier
derivative formulas do not apply unchanged.
\item The two Gaussian factors have overlapping
Fourier supports. Translation symmetry no longer
forbids a linear low-to-high response by disjoint
characters. That cancellation is not assumed.
\item Localization of \(\mathsf T_m\) is not
locality of the full conditional inverse Hessian.
In particular the Faddeev--Popov determinant
still contributes its true interactions.
\item The coercivity condition
\(C_0R^2\le m^2\) must be paid at each scale.
This note does not prove uniform control as
\(m\downarrow0\) at a fixed link radius, nor
that the full interacting measure occupies a
shrinking-radius domain with sufficient weight.
\item The initial positive coordinate domain
is still a restriction. Compact-sector coverage,
gauge-orbit multiplicity, continuum construction,
reflection positivity and physical transfer
calibration remain outside these theorems.
\end{itemize}
The localized factor now offers an appropriate
starting point for a spatial mean or boundary
response estimate. Merely renaming \(m\) as
the desired Yang--Mills mass would bypass the
retained soft sector rather than solve it.

\subsection{Exact regression checks and preservation}

The certificate
\begin{center}\small
\path{scripts/verify_ym_f14_v68_massive_curl_split.py}
\end{center}
uses \(m^2=1/5\) and \(g_\beta=2/21\).
It checks all 15 nonzero Fourier blocks at
side two and all 255 at side four, including
transversality, the exact covariance sum,
the two-form projection, the noise lower form,
and its screened comparison inverse.
The total, active and independent-kernel
noise dimensions are
\[
                   (288,135,153),\qquad (4608,2295,2313).
\]
At the test loss \(k_R=m^2\), the smallest
active lower eigenvalue checked is
\(3281/6561>1/2\).
These finite parameters verify the scalar
coercivity algebra; no value of the universal
Wilson constant \(C_0\) is inferred from them.

On the actual torus the script separately
checks \eqref{eq:v68-weighted-energy} with
rational weights
\((17/16)^{\operatorname{dist}_M(x,0)}\)
and complex rational fields. The two real
weighted Laplacian energies are
\[
                  47265483/806344,\qquad 110511459/100793,
\]
while the unweighted gradient energies are
\[
                              31596/539,\qquad 590896/539.
\]
The exact degree-eight loss bound is \(1/68\).
The proof of the uniform exponential estimate
is analytic, not a fit to these finite energies.

A scalar quadratic lift, distinct from the
Wilson density, checks normalization with
extracted and retained variances \(100/441\)
and \(41/1764\), whose sum is \(1/4\).
With its chosen negative quadratic interaction,
the projected variance and mean are \(3/11\)
and \(6/55\); the normalizing determinant
ratio is \(11/12\) and the source exponent
is \(6/275\). This checks that the Gaussian
lift preserves the complete normalized law,
not only its quadratic covariance before
reweighting.

Fifteen further rational checks verify the
Schur lower-form identity and independence
of its sign from the extraction mass.
They include a negative soft lower bound;
that test is not labelled as actual negative
curvature of a Wilson marginal.

The 8,681-byte certificate has SHA--256
\begin{center}\scriptsize\ttfamily
994EA9A903FD79538268C77026973897CE03ECBB5D92B68A55E2F04BAD25EA6C.
\end{center}
It uses unchanged V56, V58 and V62 helpers.
The predecessor V67 has 1,657,097 bytes
and SHA--256
\begin{center}\scriptsize\ttfamily
4845D4D313A0F2DB68F0D58F5D6EE9DA6FDB239ABCE0C0581DB3FC27C4F0A3AA.
\end{center}
Its entire body is retained apart from the
title version and this terminal insertion.
V68 replaces V67 as the sole active main
note. Sources, companions, monographs and
archives are unchanged.

\begin{firewall}
The same restricted nonlinear Wilson measure
now has a localized massive-curl Gaussian
representation with a uniform conditional
noise floor. The complete retained marginal
still carries the soft-sector problem.
No full compact-to-continuum Yang--Mills
construction or physical mass gap is proved.
The full goal remains active.
\end{firewall}
\section{V69: fixed-domain response and the exact retained Gaussian channel}

\subsection{Context and the new question}

V68 replaced the disjoint spectral split by a massive-curl
extraction whose physical range is all of \(\mathcal H_M\).
It proved a uniform conditional floor for every retained value,
but left a moving-domain issue for source derivatives and a
nontrivial retained marginal. This note uses that surjectivity
twice. First, changing from the extracted field to the total field
makes the conditional domain fixed. Second, Gaussian regression
before the full Wilson reweighting gives an exact sampling
representation of the retained marginal after reweighting.

The resulting source formulas are for the complete restricted
Wilson law, including its Haar and Faddeev--Popov factors.
They do not require discarding a boundary contribution.
They also expose a limitation: the localized retained channel
is not, in general, a positive spatial average. We prove actual
retained-marginal amplitude estimates at an enlarged radius,
and keep that radius distinct from the smaller guard required
by the earlier contraction programme.

All parameter assumptions are those in
\eqref{eq:v68-parameters}. The homogeneous field \(c\) is fixed
throughout the differentiations below. These are retained-field
source derivatives, not derivatives with respect to the
holonomy \(c\). No assertion about \(c\)-dependent boundaries
or source jets is inferred from them.

\subsection{A fixed-domain exponential family for every retained input}

For this section abbreviate the positive operators on
\(\mathcal H_M\) by
\begin{equation}\label{eq:v69-covariances}
 C=\mathsf C_M=\Delta^{-1},\qquad
 H=\mathsf C_m^{\mathrm h},\qquad
 L=\mathsf C_m^\ell,\qquad H+L=C,\qquad s=m^2.
\end{equation}
Inverses in this notation act only on \(\mathcal H_M\);
they do not invert a constant or a longitudinal zero mode.
Put \(D=\mathcal D_c\), \(d_M=\dim\mathcal H_M\), and
\begin{equation}\label{eq:v69-pinned-potential}
\begin{aligned}
 U_c(z)&=\beta S(a_c+tz)+W_M(a_c+tz),\\
 Q_m&=H^{-1}-\Delta=2sI+s^2\Delta^{-1},\\
 \Phi_{c,m}(z)&=U_c(z)+\tfrac12\langle z,Q_mz\rangle
                                                    \quad(z\in D),\\
 N_{c,m}(b)&=\int_D e^{-\Phi_{c,m}(z)+\langle b,z\rangle}\,dz,
       \qquad \psi_{c,m}(b)=\log N_{c,m}(b).
\end{aligned}
\end{equation}
Lebesgue measure here is the Euclidean measure on
\(\mathcal H_M\). The pin \(Q_m\) is strictly positive,
but contains an inverse Laplacian; it is not a local
physical mass term.

\begin{theorem}[Exact fixed-domain conditional and its source derivatives]
\label{thm:v69-fixed-domain}
The conditional law of the total field \(Z=Y+\eta\)
given \(Y=y\), under the actual joint law of V68, is
\begin{equation}\label{eq:v69-theta}
 d\theta_y(z)=N_{c,m}(H^{-1}y)^{-1}
       \mathbf1_D(z)e^{-\Phi_{c,m}(z)+\langle H^{-1}y,z\rangle}\,dz .
\end{equation}
Its complete high normalizer satisfies the identity
\begin{equation}\label{eq:v69-normalizer}
 Z_{c,m}(y)=\kappa_H e^{-\langle y,H^{-1}y\rangle/2}
                                  N_{c,m}(H^{-1}y),
 \qquad
 \kappa_H=(2\pi)^{-d_M/2}(\det_{\mathcal H_M}H)^{-1/2}.
\end{equation}
Both this normalizer and its logarithm are real analytic
on all of \(\mathcal H_M\). Define
\[
 \bar z_y=\mathbb E_{\theta_y}Z,\qquad
 \Gamma_y=\operatorname{Cov}_{\theta_y}(Z),\qquad
                  \bar\eta_y=\bar z_y-y .
\]
Then, with Euclidean gradients and derivatives on \(\mathcal H_M\),
\begin{equation}\label{eq:v69-source-response}
\begin{aligned}
 \nabla_y\log Z_{c,m}(y)&=H^{-1}\bar\eta_y,\\
 D_y^2\log Z_{c,m}(y)&=H^{-1}\Gamma_yH^{-1}-H^{-1},\\
 D_y\bar z_y&=\Gamma_yH^{-1},\\
 D_y\bar\eta_y&=\Gamma_yH^{-1}-I.
\end{aligned}
\end{equation}
These are identities for the normalized conditional with
its actual domain restriction.
\end{theorem}
\begin{proof}
The pushforward of the normalized two-form Gaussian under
\(\mathsf T_m\) is \(N(0,H)\), with full support on
\(\mathcal H_M\). Consequently the defining high normalizer
in \eqref{eq:v68-retained} is
\[
 \kappa_H\int_{\mathcal H_M}
       \mathbf1_D(y+\eta)
       \exp\{-\mathcal R_c(y+\eta)
                     -\tfrac12\langle\eta,H^{-1}\eta\rangle\}\,d\eta .
\]
Change variables by \(z=y+\eta\). This is a translation
of the same full-dimensional space with Jacobian one.
Using \(\mathcal R_c=U_c-\langle z,\Delta z\rangle/2\),
the exponent becomes
\[
 -\Phi_{c,m}(z)+\langle H^{-1}y,z\rangle
                              -\tfrac12\langle y,H^{-1}y\rangle .
\]
This proves both \eqref{eq:v69-theta} and
\eqref{eq:v69-normalizer}, including the Gaussian determinant.

The set \(D\) is bounded, open and nonempty. The positive
weight \(e^{-\Phi_{c,m}}\) is integrable there: the full
Wilson, Haar and determinant factors are bounded on the
bounded coordinate closure, and the determinant may vanish
at its boundary rather than diverge. For any compact set
of complex sources, the absolute values of all derivatives
of \(e^{\langle b,z\rangle}\) are uniformly bounded on \(D\)
by an integrable constant times \(e^{-\Phi_{c,m}}\).
Thus \(N_{c,m}\) is entire as a finite-dimensional Laplace
transform. It is strictly positive at every real source,
so its real logarithm is real analytic locally at every
such source. No zero-free claim for complex sources is made.

Differentiating this fixed integral gives
\(\nabla_b\psi=\bar z\) and \(D_b^2\psi=\Gamma\).
The chain rule with \(b=H^{-1}y\), including the explicit
quadratic prefactor in \eqref{eq:v69-normalizer}, proves
all four identities. No derivative of an indicator and
no integration by parts at a moving boundary was used.
\end{proof}

This change of variables was unavailable for the earlier
strict low/high spectral decomposition: its extracted
range did not equal the whole total-field space.
It is the overlap in V68, not an assumption that the
boundary is negligible, which supplies the new identity.

\subsection{Smooth marginal curvature and a precise mean-response norm}

Write
\begin{equation}\label{eq:v69-pinned-floor}
 k_R=C_0R^2,\quad g_\beta=2/(3\beta),\quad
 B_R=-\tfrac12\Delta+(g_\beta-k_R)I,\quad
 P_m=H^{-1}+B_R.
\end{equation}
The scalar estimate used in V68 gives
\begin{equation}\label{eq:v69-pinned-comparison}
 P_m\succeq\tfrac12H^{-1}+g_\beta I\succ0.
\end{equation}

\begin{theorem}[Pointwise retained curvature and global mean control]
\label{thm:v69-curvature}
For every \(y\in\mathcal H_M\),
\begin{equation}\label{eq:v69-covariance}
                   0\prec\Gamma_y\preceq P_m^{-1}\preceq2H.
\end{equation}
The actual retained Lebesgue potential
\[
       \mathcal V_{c,m}(y)
             =\tfrac12\langle y,L^{-1}y\rangle-\log Z_{c,m}(y)
\]
is real analytic and satisfies
\begin{equation}\label{eq:v69-retained-derivatives}
\begin{aligned}
 \nabla\mathcal V_{c,m}(y)
       &=(L^{-1}+H^{-1})y-H^{-1}\bar z_y,\\
 D^2\mathcal V_{c,m}(y)
       &=L^{-1}+H^{-1}-H^{-1}\Gamma_yH^{-1},\\
 L_m&\preceq D^2\mathcal V_{c,m}(y)\preceq L^{-1}+H^{-1}.
\end{aligned}
\end{equation}
Here \(L_m\) is exactly the lower multiplier in
\eqref{eq:v68-retained-floor}. The lower bound is now a
pointwise smooth Hessian inequality, with the same sign
limitation as in V68.

In the norm \(\|v\|_{H^{-1}}=\|H^{-1/2}v\|_2\), the
physical extracted mean map is globally nonexpansive:
\begin{equation}\label{eq:v69-mean-Lipschitz}
 \|\bar\eta_{y_1}-\bar\eta_{y_2}\|_{H^{-1}}
                                  \le\|y_1-y_2\|_{H^{-1}}.
\end{equation}
In particular \(\bar\eta_0=0\) and
\(\|\bar\eta_y\|_{H^{-1}}\le\|y\|_{H^{-1}}\).
For fixed finite \(M,m,\beta,c\), no global Lipschitz
constant strictly smaller than one is possible in
\eqref{eq:v69-mean-Lipschitz}.
\end{theorem}
\begin{proof}
The full Wilson and geometric lower forms give
\[
 D^2\Phi_{c,m}
   =D^2U_c+H^{-1}-\Delta
   \succeq H^{-1}-\tfrac12\Delta+(g_\beta-k_R)I=P_m .
\]
On a nonzero mode,
\[
 P_m-\tfrac12H^{-1}-g_\beta I
   =s-k_R+\frac{s^2}{2\lambda}\ge0 .
\]
Linear tilting does not change this Hessian.
The strong log-concavity covariance inequality on a convex
support, in the form proved and used in V66--V68, yields
\(\Gamma_y\preceq P_m^{-1}\). Equivalently one can whiten
by \(P_m^{1/2}\) and apply the convex-support Poincar\'e
inequality to every linear observable. Extending the
convex remainder by infinity retains the boundary.
Since the density is positive on the full-dimensional
open set \(D\), the variance of every nonzero linear
observable is strictly positive. This proves
\eqref{eq:v69-covariance}.

The derivative formulas follow from
\eqref{eq:v69-source-response}. Inserting the covariance
upper bound gives the lower operator
\[
 L^{-1}+H^{-1}-H^{-1}P_m^{-1}H^{-1}.
\]
All operators in this last expression commute. At
eigenvalue \(\lambda>0\), with \(h,\ell,b_R\) as in
\eqref{eq:v68-schur-data}, it equals
\[
 \ell^{-1}+h^{-1}-\frac{h^{-2}}{h^{-1}+b_R}
                    =\ell^{-1}+\frac{b_R}{1+hb_R}=L_m(\lambda).
\]
Positivity of \(\Gamma_y\) gives the upper bound.

For the mean derivative, conjugation into the stated
metric gives the symmetric operator
\begin{equation}\label{eq:v69-conjugated-response}
 H^{-1/2}(D_y\bar\eta_y)H^{1/2}
                     =H^{-1/2}\Gamma_yH^{-1/2}-I.
\end{equation}
Its spectrum lies in \([-1,1]\) by
\eqref{eq:v69-covariance}. Integrating this operator norm
bound along the segment between any two retained inputs
proves \eqref{eq:v69-mean-Lipschitz}. Translation invariance
at \(y=0\) makes \(\bar z_0\) constant, and membership in
\(\mathcal H_M\) makes it zero.

Finally fix \(v\ne0\). Boundedness of \(D\) implies
\(\sup_y\|\bar z_y\|_{H^{-1}}<\infty\) at the fixed
finite parameters. Therefore
\[
 \lim_{u\to\infty}
       \frac{\|\bar\eta_{uv}-\bar\eta_0\|_{H^{-1}}}
            {u\|v\|_{H^{-1}}}
   =\lim_{u\to\infty}
       \frac{\|\bar z_{uv}-uv\|_{H^{-1}}}{u\|v\|_{H^{-1}}}=1.
\]
A strictly smaller global constant would contradict this
limit. This obstruction comes from the bounded total-field
domain and is not an assertion of negative Wilson curvature.
\end{proof}

The metric in this theorem is important.
Its precision is \(\Delta+2sI+s^2\Delta^{-1}\).
It is neither a local maximum norm nor uniformly equivalent
to the Euclidean norm as the volume grows.
Even at the matrix level, covariance order does not remove
this distinction. For example, take
\[
 H=\begin{pmatrix}1&0\\0&1/100\end{pmatrix},\qquad
 \Gamma=\begin{pmatrix}1&2/25\\2/25&1/100\end{pmatrix}.
\]
Then \(0\prec\Gamma\prec2H\), whereas
\[
 \Gamma H^{-1}-I=\begin{pmatrix}0&8\\2/25&0\end{pmatrix},\qquad
 H^{-1/2}(\Gamma H^{-1}-I)H^{1/2}
                  =\begin{pmatrix}0&4/5\\4/5&0\end{pmatrix}.
\]
The Euclidean operator norm is eight and the conjugated
norm is \(4/5\). This is a comparison-matrix example,
not an identification of a particular Wilson covariance.
The theorem supplies no entrywise spatial response estimate.

\subsection{Gaussian regression survives the complete Wilson reweighting}

Define on \(\mathcal H_M\)
\begin{equation}\label{eq:v69-channel-operators}
\begin{aligned}
 \mathsf A_m&=LC^{-1}=I-H\Delta
        =s(2\Delta+s)(\Delta+s)^{-2}\\
        &=2s(\Delta+s)^{-1}-s^2(\Delta+s)^{-2},\\
 \Sigma_m&=L-LC^{-1}L=LH C^{-1}=H\mathsf A_m .
\end{aligned}
\end{equation}
We use the scalar expression on the right to extend
\(\mathsf A_m\) to all link fields. In particular its
constant-mode multiplier is one. The extension of
\(\Sigma_m\) by \(H\mathsf A_m\) still vanishes on
constant and longitudinal modes. For \(k\ne0\),
\begin{equation}\label{eq:v69-channel-symbols}
\begin{aligned}
 a_m(\lambda)&=\frac{s(2\lambda+s)}{(\lambda+s)^2},\\
 \Sigma_m(k)&=\frac{s\lambda(2\lambda+s)}{(\lambda+s)^4}
                                               \Pi(k)\otimes I_3.
\end{aligned}
\end{equation}
Consequently, on \(\mathcal H_M\),
\begin{equation}\label{eq:v69-regression-identities}
\begin{gathered}
 0\prec\Sigma_m\preceq H\preceq C,\qquad
                       \|\Sigma_m\|\le(4m^2)^{-1},\\
 \Sigma_m^{-1}=L^{-1}+H^{-1},\qquad
 \Sigma_m^{-1}\mathsf A_m=H^{-1},\qquad
 \mathsf A_m^*\Sigma_m^{-1}\mathsf A_m=Q_m .
\end{gathered}
\end{equation}

\begin{theorem}[The true retained marginal is an exact affine Gaussian channel]
\label{thm:v69-channel}
Sample \(Z\) from the actual restricted Wilson law \(\mu_c\)
in \eqref{eq:v68-target}, and independently sample
\(\varepsilon\sim N(0,\Sigma_m)\) on \(\mathcal H_M\).
Then
\begin{equation}\label{eq:v69-channel}
             Y=\mathsf A_m Z+\varepsilon,\qquad
             \eta=(I-\mathsf A_m)Z-\varepsilon
\end{equation}
has precisely the physical joint law \((Y,\eta)\)
obtained from V68. In particular the complete retained
marginal, not its reference Gaussian, has density
\begin{equation}\label{eq:v69-channel-density}
     \pi_{c,m}(y)=\int_D \gamma_{\Sigma_m}(y-\mathsf A_m z)\,d\mu_c(z),
\end{equation}
where \(\gamma_{\Sigma_m}\) denotes the normalized
Lebesgue Gaussian density on \(\mathcal H_M\).
\end{theorem}
\begin{proof}
Start before reweighting with independent
\(Y\sim N(0,L)\) and \(\eta\sim N(0,H)\), and put
\(Z=Y+\eta\). Then
\(\operatorname{Cov}(Z)=C\) and
\(\operatorname{Cov}(Y,Z)=L\).
The centered Gaussian
\(\varepsilon=Y-LC^{-1}Z\) has covariance
\(L-LC^{-1}L=\Sigma_m\) and has zero cross covariance
with \(Z\). Joint Gaussianity makes \(\varepsilon\)
independent of \(Z\). These statements are identities
of normalized probability laws.

Now multiply the joint law by the same nonnegative
function \(\mathbf1_D(Z)e^{-\mathcal R_c(Z)}\) as in V68,
and divide by its integral \(\mathfrak Z_c\).
It depends only on \(Z\). Hence the independent
\(\varepsilon\) factor is unchanged and the \(Z\)
factor becomes exactly \(\mu_c\).
Solving the definitions for \(Y,\eta\) proves
\eqref{eq:v69-channel}. Integrating \(Z\) proves
\eqref{eq:v69-channel-density}.

For completeness, \(0<a_m(\lambda)<1\) for \(\lambda>0\).
Thus \(\Sigma_m=H\mathsf A_m\preceq H\), and the bound
on its norm follows from \(h_m(\lambda)\le1/(4m^2)\).
The three inverse identities in
\eqref{eq:v69-regression-identities} follow by
substituting \(C=H+L\), with commuting positive operators.
They also give an independent check of the pinned potential:
expanding the Gaussian exponent in
\eqref{eq:v69-channel-density} produces
\(Q_m\) in the \(z\)-quadratic term and \(H^{-1}y\)
in its source term, exactly as in
\eqref{eq:v69-pinned-potential}.
\end{proof}

The Gaussian regression identity has standard background
in Rasmussen--Williams,
\href{https://gaussianprocess.org/gpml/chapters/RWA.pdf}
{Gaussian Processes for Machine Learning, Appendix A.2,
equation (A.6)}. The argument above proves the particular
operators and the reweighting step explicitly.
It does not assert that the interacting total field
or its \(Z\)-given-\(Y\) conditional is Gaussian.
The unknown nonlinear Wilson law remains the channel input.

\subsection{Locality of the channel and failure of positive averaging}

\begin{proposition}[Localized channel and a uniform maximum-norm bound]
\label{prop:v69-locality}
With the coordinate distance of V68, for coordinate sets \(E,F\),
\begin{equation}\label{eq:v69-locality}
\begin{aligned}
 \|\mathbf1_E\mathsf A_m\mathbf1_F\|_{2\to2}
      &\le8e^{-m\operatorname{dist}_M(E,F)/4},\\
 \|\mathbf1_E\Sigma_m\mathbf1_F\|_{2\to2}
      &\le\frac{128}{m^2}e^{-m\operatorname{dist}_M(E,F)/4}.
\end{aligned}
\end{equation}
The scalar extension of \(\mathsf A_m\) preserves constants and
\begin{equation}\label{eq:v69-infinity}
 \|\mathsf A_m v\|_\infty\le\chi\|v\|_\infty,
                   \qquad \chi=1+2e^{-2},
\end{equation}
where the maximum norm takes the Euclidean colour norm
at each link. This extension need not preserve positivity
of scalar fields.
\end{proposition}
\begin{proof}
For the weights \(W=e^{mf/4}\) used in V68, its
accretivity estimate gives
\(\|W(\Delta+s)^{-1}W^{-1}\|\le2/s\).
The resolvent expression in
\eqref{eq:v69-channel-operators} therefore gives
\[
       \|W\mathsf A_mW^{-1}\|\le4+4=8 .
\]
Also \(\|WHW^{-1}\|\le16/m^2\) by V68.
Multiplication bounds the weighted \(\Sigma_m\)
by \(128/m^2\). Taking \(f\) to be distance from
the base sites of \(F\) proves
\eqref{eq:v69-locality}, including for the full
extensions just specified.

Let \(P_u=e^{-u\Delta}\) on the scalar torus.
It is positive and preserves constants: since
\(\Delta=8I-\sum_{\mu=1}^4(T_\mu+T_\mu^{-1})\),
this follows directly from the nonnegative exponential
series for the commuting shifts. Hence \(P_u\)
contracts the stated maximum norm on vector-valued links.
Laplace transforms of the two resolvents give
\begin{equation}\label{eq:v69-signed-heat}
       \mathsf A_m=\int_0^\infty s(2-su)e^{-su}P_u\,du .
\end{equation}
The signed scalar weight has integral one and
absolute integral
\[
 \int_0^\infty |2-v|e^{-v}\,dv=1+2e^{-2}.
\]
This proves preservation of constants and
\eqref{eq:v69-infinity}.

To prove that positivity cannot be assumed, take
the actual torus of side four, \(s=1\), and
\(x=(2,2,2,2)\). In one coordinate its heat kernel
at the antipode is \((1-e^{-2u})^2/4\).
The four-dimensional product kernel and
\eqref{eq:v69-signed-heat} give the channel entry
\begin{equation}\label{eq:v69-negative-kernel}
\begin{aligned}
 a_{1,4}(x)
   &=\frac1{256}\sum_{j=0}^{8}(-1)^j\binom8j
       \left\{\frac2{1+2j}-\frac1{(1+2j)^2}\right\}\\
   &=-\frac{7902592}{83770862175}<0 .
\end{aligned}
\end{equation}
This is a finite-torus counterexample, with the
same scalar Laplacian as the construction.
\end{proof}

Thus the spatial kernel is a localized signed kernel,
not in general a Markov averaging kernel. Positivity of
its Fourier multiplier is a different statement.
One cannot apply matrix concavity of \(\mathcal F_M\)
as if \(\mathsf A_m a\) were a convex combination
of translated connections. Nor does the proved
\(\chi>1\) bound furnish a shrinking-radius map.
The counterexample concerns this full scalar extension;
it does not assert a particular negative direction
for the interacting Wilson measure.

\subsection{Actual retained tails, with the radius enlargement made explicit}

Let \(N_{\mathrm{link}}=4M^4\). Choose a universal
constant \(C_{\mathrm d}\) so that for every link set \(S\),
\begin{equation}\label{eq:v69-trace}
          \Tr(\mathbf1_S C\mathbf1_S)\le C_{\mathrm d}|S|.
\end{equation}
The colour components are included in this trace.
Such a constant is the uniform diagonal Green bound
already proved in V58. By \eqref{eq:v69-regression-identities}
the same bound holds for \(\Sigma_m\).

\begin{theorem}[Full retained-marginal amplitude bounds]
\label{thm:v69-retained-tails}
There is a universal \(C_1\) such that for every \(u\ge0\),
\begin{equation}\label{eq:v69-retained-max}
 \pi_{c,m}\!\left\{
      \max_e|a_c(e)+tY_e|>
          \chi R+C_1\sqrt{\frac{1+\log M+u}{\beta}}
                    \right\}\le e^{-u}.
\end{equation}
More generally, if \(\rho>0\) and
\(\beta\rho^2\ge4C_{\mathrm d}\), then every link set \(S\)
satisfies
\begin{equation}\label{eq:v69-retained-joint}
 \pi_{c,m}\!\left\{
       |a_c(e)+tY_e|\ge\chi R+\rho
                               \text{ for every }e\in S\right\}
           \le\exp\{-\tfrac12m^2\beta\rho^2|S|\}.
\end{equation}
These estimates are for the actual marginal, including
its high normalizer and the input Wilson restriction.
\end{theorem}
\begin{proof}
Use the independent sampling representation of
Theorem~\ref{thm:v69-channel}. Since the channel preserves
the homogeneous \(a_c\),
\begin{equation}\label{eq:v69-affine-field}
       a_c+tY=\mathsf A_m(a_c+tZ)+t\varepsilon .
\end{equation}
The input law is supported on \(D\), so its input
maximum is less than \(R\). The deterministic channel
part in \eqref{eq:v69-affine-field} is therefore bounded
by \(\chi R\).

Each real component of \(\varepsilon\) has variance
at most \(C_{\mathrm d}\). The scalar Gaussian tail
and a union bound over \(12M^4\) real components give
\[
 \mathbb P\{\max_e|\varepsilon_e|>b\}
           \le24M^4\exp\{-b^2/(6C_{\mathrm d})\}.
\]
Taking \(b^2=6C_{\mathrm d}(\log24+4\log M+u)\)
and enlarging a universal constant proves
\eqref{eq:v69-retained-max}.

For the joint bound, the event on the left of
\eqref{eq:v69-retained-joint} implies
\(\|\varepsilon_S\|_2^2\ge\beta\rho^2|S|\)
by the reverse triangle inequality, simultaneously
for every realization of the input. Write
\(\Sigma_S=\mathbf1_S\Sigma_m\mathbf1_S\) as a covariance
on its \(3|S|\) coordinates. It may be singular.
Its trace is at most \(C_{\mathrm d}|S|\) and
its norm is at most \(1/(4m^2)\).
The Gaussian quadratic moment formula, also valid with
zero covariance eigenvalues, and
\(-\log(1-v)\le v/(1-v)\) yield
\[
\begin{split}
 \log\mathbb E e^{m^2\|\varepsilon_S\|_2^2}
    &=-\tfrac12\log\det(I-2m^2\Sigma_S)\\
    &\le2m^2\Tr\Sigma_S
     \le2m^2C_{\mathrm d}|S|.
\end{split}
\]
Markov's inequality gives exponent
\(-m^2(\beta\rho^2-2C_{\mathrm d})|S|\).
The stipulated lower bound on \(\beta\rho^2\)
proves \eqref{eq:v69-retained-joint}.
For \(S=\varnothing\) both sides equal one.
No independence of different noise links is assumed.
\end{proof}

This removes the retained-reference-measure substitution
for these particular amplitude events. It does not
establish the small incoming guard of V67: the threshold
is larger than the input radius, and the conclusion
depends on an input law already restricted to \(D\).
Taking \(R\) fixed while sending \(m\) to zero would
also violate the conditional-coercivity hypothesis
\(C_0R^2\le m^2\).

\begin{proposition}[The gauge-linear tail is transferred, not created]
\label{prop:v69-gauge-transfer}
For a universal \(C_2\), every \(K\ge0\) and \(u\ge0\),
\begin{equation}\label{eq:v69-gauge-transfer}
\begin{split}
 &\pi_{c,m}\!\left\{
    \|G_1(a_c+tY)\|>
          \chi K+C_2\sqrt{(1+\log M+u)/\beta}\right\}\\
 &\hspace{12mm}\le
       \mu_c\{\|G_1(a_c+tZ)\|>K\}+e^{-u}.
\end{split}
\end{equation}
All \(G_1\) norms are operator norms on the normalized
gauge space used in V60--V61.
\end{proposition}
\begin{proof}
The channel is scalar convolution with kernel whose
absolute sum is at most \(\chi\). The linear map \(G_1\)
is translation covariant: translating a connection
unitarily conjugates its gauge operator. Consequently
\[
 \|G_1(\mathsf A_m v)\|
                     \le\chi\|G_1(v)\|.
\]
Linearity and \eqref{eq:v69-affine-field} now imply
\[
 \|G_1(a_c+tY)\|
          \le\chi\|G_1(a_c+tZ)\|+t\|G_1(\varepsilon)\|.
\]
The centered Gaussian \(\varepsilon\) has covariance
\(\Sigma_m\preceq C\). The div--curl/BMO proof of
\eqref{eq:v61-gauge-tail} uses Gaussian quadratic
covariance bounds; they remain valid under this
covariance domination, as detailed in V62.
It therefore gives
\[
 \mathbb P\{\|G_1(\varepsilon)\|>
                     C_2\sqrt{1+\log M+u}\}\le e^{-u}.
\]
Splitting according to the input event proves the claim.
\end{proof}

The defining positivity \(\mathcal F_M(a_c+tZ)\succ0\)
does not provide the small value of
\(\|G_1(a_c+tZ)\|\) required by this last bound.
The input probability in \eqref{eq:v69-gauge-transfer}
has not been estimated away.

\subsection{What this changes upstream, and what it does not}

The moving retained-source boundary is now accounted for
exactly, not by an additional boundary hypothesis.
The actual retained marginal is also an explicit
localized Gaussian channel of the same restricted
interacting total field. Its normalization is not
replaced by a free normalization. These statements
are new interfaces; they do not replace the unresolved
infrared estimate by a different name.

Several tempting closures are now ruled out within
this route. The conditional extracted-mean map cannot
be a global strict contraction on all inputs in its
natural metric because the total-field domain is
bounded. The channel is not generally positive spatial
averaging. Its true retained amplitude bound is at
\(\chi R+\rho\), not at a smaller radius. Exponential
locality of the channel and its independent noise is
not exponential decay of the correlated input \(\mu_c\).
Finally the exact pointwise marginal Hessian has only
the same proved lower form \(L_m\) as in V68, with its
uncontrolled soft sign.

The next substantive input must therefore control the
interacting total-field distribution or its actual
spatial response, including the gauge-linear event,
at a scale-compatible radius. A bounded-support
identity or the free reference covariance alone does
not supply that input. Compact-sector coverage,
continuum construction, reflection positivity and
physical transfer calibration are still separate
requirements. The next scheduled SM/NS/shell companion
checkpoint is V70; this V69 does not alter those files.

\subsection{Exact certificates and preservation}

The new certificate is
\begin{center}\small
\path{scripts/verify_ym_f14_v69_fixed_domain_channel.py}.
\end{center}
It checks both \(m^2=1/5\) and \(m^2=1\) on every
nonzero Fourier block of the actual tori of sides
two and four: 30 and 510 exact block tests respectively.
These verify the covariance sum, regression
cross covariance, retained covariance, inverse
identities, pin precision and noise comparisons.
The finite parameters test algebra, not the value
of the unspecified universal Wilson constant \(C_0\).

An independent exact Fourier sum agrees with the
heat-kernel computation of the negative entry
\(-7902592/83770862175\) in
\eqref{eq:v69-negative-kernel}.
A separate scalar quadratic model, explicitly not
the Wilson density, checks the source prefactor and
the response formulas. With \(\lambda=4\), \(s=1/5\),
interaction curvature \(-1/3\), linear source \(2/5\)
and retained input \(y=2/7\), it gives
\[
 H=100/441,\quad L=41/1764,\quad Q_m=41/100,\quad
                         \Gamma_y=300/1223.
\]
The total mean and its extracted counterpart are
\(498/1223\) and \(1040/8561\); their derivatives are
\(1323/1223\) and \(100/1223\).
The exact retained variance and mean computed both
ways are \(50143/2139291\) and \(82/8085\).
The matrix response example independently checks
the Euclidean norm eight and the metric norm \(4/5\).
These finite regressions supplement, not substitute
for, the analytic proofs of uniform estimates.

The 5,735-byte certificate has SHA--256
\begin{center}\scriptsize\ttfamily
0FCB2D15A84EC38341A5F9932418CA678011D677A6BF01FF77A9FC7514D27AFB.
\end{center}
It uses the unchanged V56 and V58 helpers.
The predecessor V68 has 1,682,407 bytes and SHA--256
\begin{center}\scriptsize\ttfamily
F4E829A8986B5D24A47A1A9E93E0FB0114D7BD6DBF4294B4698AF2A07ADFEC21.
\end{center}
Its complete body is preserved apart from the title
version and this terminal insertion. V69 replaces V68
as the single active main note; archives, attachments,
companion notes and monographs are unchanged.

\begin{firewall}
This note proves exact retained-source identities,
a sharp global nonexpansive bound in a specified
Gaussian precision norm, a localized Gaussian-channel
representation of the full restricted retained law,
and actual retained amplitude tails at an enlarged
radius. It does not prove a shrinking-radius
renormalization step, a uniform physical infrared
floor, an unrestricted continuum Yang--Mills theory,
or its Hamiltonian mass gap. The full goal remains active.
\end{firewall}
\section{V70: soft-mode observation and an exact variance--entropy budget}

\subsection{Companion checkpoint and direction}

The scheduled fifth-version checkpoint was performed before
this continuation. The current active companions are unchanged
from V65:
\begin{center}\small
\path{Renewal_SM_Einstein_Continuum_Research_Notes_V72.tex}\\
\path{Renewal_NS_Stability_Research_Notes_V26.tex}\\
\path{Renewal_Yang_Mills_Shell_Mixing_Research_Notes_V16.tex}.
\end{center}
Their SHA--256 digests, in the same order, are
\begin{center}\scriptsize\ttfamily
F44F9745D43379E1672C5550411B58E97195A4C9278592D221DF5D3F644DC1F5\\
82C887F8C2C69E4F28FAD7C3DC8DE5B9099CB5A4BF431EBA1FFF3E91935978AD\\
BF138E63BA826BC030D3A8B87E04FDD3B1EC712D269BB064674E96509B46E710.
\end{center}
The latest mathematical sections and their scope statements
were inspected again, rather than inferred from their filenames.

SM V72 derives covariance transport in a specified many-copy
model at fixed mode cutoff and fixed times. Closure of means
requires point-concentrated limiting data; a distribution of
sources must otherwise remain a distribution. NS V26 gives
rank-one formulas for Oseen-kernel derivatives. Its later
certification references do not furnish a Yang--Mills
probability estimate. Shell V16 evaluates the fully measured
flat-root one-loop coefficient, including the actual source
normal and tangent-return factors. Its coefficient limit is
not a joint nonperturbative volume/coupling estimate. None
supplies the missing interacting infrared bound here.

The direction taken is consequently a distribution-level
one. V69's exact observation interpretation is extended
to a pin acting only on the soft modes needed by the
proved Wilson lower form. We retain the full random
source law, derive the exact variance and entropy it
carries, and state the covariance estimate that would
control this cost. A finite-volume estimate is proved,
but its volume dependence is displayed and compared
with a stronger elementary bounded-perturbation estimate.
Neither is represented as a thermodynamic mass gap.
The next companion checkpoint is V75.

\subsection{A minimal positive spectral pin for the known Wilson lower form}

Keep the finite-dimensional space \(\mathcal H_M\),
the fixed homogeneous \(c\), and the convex domain
\(D=\mathcal D_c\) of V68. Set
\[
 \mu=\mu_c,\qquad
 U(z)=\beta S(a_c+\beta^{-1/2}z)+W_M(a_c+\beta^{-1/2}z),
 \quad k=C_0R^2,\quad g=2/(3\beta).
\]
Thus \(d\mu\) is proportional to
\(\mathbf1_D e^{-U}\,dz\), with every original factor
retained, and
\begin{equation}\label{eq:v70-original-floor}
                   D^2U\succeq\tfrac12\Delta+(g-k)I .
\end{equation}
Use the nonemptiness and small-coordinate assumptions
of V68, with \(k\le1\). No extraction mass is needed
for the construction in this note.

Define positive spectral operators
\begin{equation}\label{eq:v70-soft-pin}
 \mathsf J=\tfrac14\Delta+gI,\qquad
 \mathsf K=(kI-\tfrac14\Delta)_+,\qquad
 E=\operatorname{ran}\mathsf K,\qquad
 \mathsf B=\mathsf K^{1/2}:\mathcal H_M\longrightarrow E.
\end{equation}
The notation \(E\) here is a Euclidean observation
space, not the earlier two-form noise space.
If \(\mathsf K=0\), take \(E=\{0\}\) and interpret
all its Gaussian factors as the unit probability law.

\begin{lemma}[Soft-only convexification]
\label{lem:v70-pin}
For every \(b\in E\), the potential
\begin{equation}\label{eq:v70-endpoint-potential}
             U(z)+\tfrac12\|\mathsf Bz\|^2-\langle b,\mathsf Bz\rangle
                                                        \quad(z\in D)
\end{equation}
has Hessian at least \(\mathsf J\succ0\).
Among nonnegative multipliers \(q(\Delta)\) which
raise the particular lower form in
\eqref{eq:v70-original-floor} to \(\mathsf J\),
\(\mathsf K\) is the smallest at each eigenvalue.
This is minimality for that lower-form comparison,
not a claim about the optimal pin for the actual Hessian.
Its real rank \(r=\dim E\) obeys
\begin{equation}\label{eq:v70-rank}
 r=9\,\#\{n\ne0:\lambda(n)<4k\}
             \le\min\{9(M^4-1),\,9(1+M\sqrt{k})^4\}.
\end{equation}
\end{lemma}
\begin{proof}
At an eigenvalue \(\lambda>0\), adding
\(\kappa(\lambda)=(k-\lambda/4)_+\) to the
known lower form gives
\[
       \lambda/2+g-k+\kappa(\lambda)\ge\lambda/4+g.
\]
Conversely, any nonnegative \(q(\lambda)\) which
makes this inequality true must satisfy
\(q(\lambda)\ge(k-\lambda/4)_+\).
Linear sources have zero Hessian.

There are nine real transverse-colour degrees of freedom
per nonzero Fourier frequency in the real dimension count.
For centered integer representatives \(n_\mu\),
\[
 \lambda(n)=4\sum_{\mu=1}^4\sin^2(\pi n_\mu/M)
                                      \ge16|n|^2/M^2.
\]
Thus a contributing \(n\) has
\(|n|<M\sqrt{k}/2\). Counting the enclosing integer
cube proves \eqref{eq:v70-rank}. Conjugate frequencies
are counted with their real degrees of freedom, not
as additional independent complex fields.
\end{proof}

The rank can still grow proportionally to the volume
at fixed \(R\). The sharp spectral cutoff also is
not a spatially local blocking map. The present
construction is an exact observation and convexification
device, not a replacement for V68's localized extraction.

\subsection{The exact nonlinear posterior and its true source law}

For \(0\le\tau\le1\), \(b\in E\), define
\begin{equation}\label{eq:v70-posterior}
\begin{aligned}
 \mathcal N_\tau(b)
   &=\mathbb E_\mu
       \exp\{\langle b,\mathsf BZ\rangle-\tfrac{\tau}{2}\|\mathsf BZ\|^2\},\\
 d\theta_{\tau,b}(z)
   &=\mathcal N_\tau(b)^{-1}
       e^{\langle b,\mathsf Bz\rangle-\tau\|\mathsf Bz\|^2/2}\,d\mu(z).
\end{aligned}
\end{equation}
The integral is strictly positive and finite.
Boundedness of \(D\) justifies all real source
derivatives, as in V69. In particular
\begin{equation}\label{eq:v70-posterior-derivative}
 \nabla_b\log\mathcal N_\tau(b)=\theta_{\tau,b}(\mathsf BZ),
 \qquad
 D_b^2\log\mathcal N_\tau(b)
                 =\operatorname{Cov}_{\theta_{\tau,b}}(\mathsf BZ).
\end{equation}

\begin{theorem}[A normalized soft observation of the actual Wilson field]
\label{thm:v70-observation}
Sample \(Z\sim\mu\), independently of a standard
Brownian motion \(W_\tau\) on \(E\), and set
\begin{equation}\label{eq:v70-observation}
                    X=\mathsf BZ,\qquad O_\tau=\tau X+W_\tau.
\end{equation}
The conditional law of \(Z\) given the observation
path up to \(\tau\) is \(\theta_{\tau,O_\tau}\).
For \(\tau>0\), the actual observation density is
\begin{equation}\label{eq:v70-source-law}
 p_\tau(b)=(2\pi\tau)^{-r/2}
                    e^{-\|b\|^2/(2\tau)}\mathcal N_\tau(b).
\end{equation}
It integrates to one. Consequently
\begin{equation}\label{eq:v70-mixture}
                       \mu=\mathbb E\,\theta_{\tau,O_\tau}
\end{equation}
as an identity of probability measures.
At \(\tau=1\) every conditional in this mixture
has the uniform lower form \(\mathsf J\).
\end{theorem}
\begin{proof}
Given \(Z=z\), the observation at time \(\tau\)
is Gaussian with mean \(\tau\mathsf Bz\) and
covariance \(\tau I_E\). Expanding its normalized
density gives the factor
\[
 (2\pi\tau)^{-r/2}e^{-\|b\|^2/(2\tau)}
        e^{\langle b,\mathsf Bz\rangle-\tau\|\mathsf Bz\|^2/2}.
\]
Bayes' rule proves the conditional law at the
endpoint and \eqref{eq:v70-source-law}.
The Gaussian exponential moment gives
\(\int p_\tau=1\) directly, without replacing
\(\mathcal N_\tau\) by a constant.

For observations at finitely many time increments,
the likelihood ratio for a fixed \(X=x\) against
zero drift is
\[
       \exp\{\langle x,O_\tau\rangle-\tfrac{\tau}{2}\|x\|^2\}.
\]
The sum of increment terms telescopes to the endpoint.
Passing along increasing finite grids, by continuity
and conditional-expectation convergence, gives the
same posterior for the whole observed path.
Boundedness of \(X\) ensures integrability of these
likelihoods at every finite time.
The mixture identity is the law of total expectation.
The endpoint lower form follows from
Lemma~\ref{lem:v70-pin}.
\end{proof}

This observation construction does not assume that
the original measure can already be sampled efficiently.
It is an identity defining a joint probability law.
Its source marginal is \(p_\tau\), not a free Gaussian.
The fixed total-field domain \(D\) is present in
every posterior integral.

It is also compatible with V69 rather than a change
of the original theory. In that note
\(Q_m=H^{-1}-\Delta\), and the actual source
\(H^{-1}Y\), conditional on \(Z\), has mean \(Q_mZ\)
and covariance \(Q_m\). The identities
\(H^{-1}\mathsf A_m=Q_m\) and
\(H^{-1}\Sigma_mH^{-1}=Q_m\) prove this directly.
Whitening by \(Q_m^{-1/2}\) gives the time-one
observation \(Q_m^{1/2}Z+G\).
The new choice \(\mathsf K\) observes only the soft
subspace needed to raise the proved lower form.

\subsection{Exact mean, covariance, variance and entropy evolution}

Let \(\mathcal F_\tau\) be the observation filtration,
and abbreviate \(\theta_\tau=\theta_{\tau,O_\tau}\).
Write
\[
 a_\tau=\theta_\tau Z,\qquad
 \Gamma_\tau=\operatorname{Cov}_{\theta_\tau}(Z),\qquad
 c_\tau(f)=\operatorname{Cov}_{\theta_\tau}(f,X).
\]
The last expression is a vector in \(E\).
Throughout, expectations outside \(\theta_\tau\)
refer to the true joint observation law.

\begin{theorem}[The exact conditional-variance and entropy costs]
\label{thm:v70-budgets}
The process
\[
                 I_\tau=O_\tau-\int_0^\tau\theta_uX\,du
\]
is a standard Brownian motion in the observation
filtration. For every bounded real \(f\),
\begin{equation}\label{eq:v70-filter}
          d(\theta_\tau f)=\langle c_\tau(f),dI_\tau\rangle,
       \qquad da_\tau=\Gamma_\tau\mathsf B^*dI_\tau.
\end{equation}
In particular
\begin{equation}\label{eq:v70-covariance-drift}
          \frac{d}{d\tau}\mathbb E\Gamma_\tau
                             =-\mathbb E(\Gamma_\tau\mathsf K\Gamma_\tau).
\end{equation}
For \(0\le T\le1\),
\begin{equation}\label{eq:v70-variance-budget}
 \Var_\mu f
   =\mathbb E\Var_{\theta_T}f
                         +\int_0^T\mathbb E\|c_\tau(f)\|^2\,d\tau.
\end{equation}
If \(f\) is bounded above and bounded strictly
away from zero, with
\(\operatorname{Ent}_\nu f=\nu(f\log f)-\nu f\log(\nu f)\), then
\begin{equation}\label{eq:v70-entropy-budget}
 \operatorname{Ent}_\mu f
   =\mathbb E\operatorname{Ent}_{\theta_T}f
      +\frac12\int_0^T
          \mathbb E\frac{\|c_\tau(f)\|^2}{\theta_\tau f}\,d\tau .
\end{equation}
No prior or posterior is replaced by a Gaussian
in these identities.
\end{theorem}
\begin{proof}
The observation drift in its own filtration is
\(\theta_\tau X\). Indeed, conditional expectation
of \(X\) at that time is \(\theta_\tau X\), and
future increments of the original Brownian motion
have zero conditional mean. Thus \(I_\tau\) is
a continuous martingale with quadratic variation
\(\tau I_E\); the Brownian martingale characterization
makes it a Brownian motion in \(\mathcal F_\tau\).

Apply It\^o's rule to the normalized fixed-domain
ratio in \eqref{eq:v70-posterior}, using
\(dO_\tau=\theta_\tau X\,d\tau+dI_\tau\).
The derivative in \(b\) is \(c_\tau(f)\).
The explicit \(-\tau\|X\|^2/2\) derivative,
the quadratic variation, and the normalization
drift cancel. This gives \eqref{eq:v70-filter}.
Equivalently the normalized density has differential
\[
              d\theta_\tau(z)
                =\theta_\tau(z)\langle\mathsf Bz-\theta_\tau X,dI_\tau\rangle .
\]
All integrands are bounded at the fixed finite
parameters, so the displayed martingales and
quadratic variations are integrable.

Apply this formula to \(Z\) and \(ZZ^*\).
The latter conditional expectation is a martingale.
Subtract the product \(a_\tau a_\tau^*\).
Its It\^o drift is
\(\Gamma_\tau\mathsf B^*\mathsf B\Gamma_\tau\),
which proves \eqref{eq:v70-covariance-drift}.

For \(h_\tau=\theta_\tau f\), the It\^o drift of
\(h_\tau^2\) is \(\|c_\tau(f)\|^2\).
Since \(\mathbb E\theta_\tau f^2=\mu f^2\),
integration proves \eqref{eq:v70-variance-budget}.
For positive \(f\), the drift of \(h_\tau\log h_\tau\)
is \(\|c_\tau(f)\|^2/(2h_\tau)\).
The expectation of \(\theta_\tau(f\log f)\) is constant.
This proves \eqref{eq:v70-entropy-budget}.
The assumed lower bound on \(f\) justifies the
It\^o derivatives and integrability without stopping.
\end{proof}

These are specialization and verification of a standard
stochastic-observation method, not a claim of a new
general localization theory. For its general
variance/entropy framework see Chen--Eldan,
\href{https://arxiv.org/pdf/2203.04163}
{Localization Schemes, Sections 3.1--3.3}.
The Wilson pin, actual domain, finite-rank observation
and explicit comparison constants used here are
specified and proved in this note.

\subsection{The precise covariance budget that transfers back to the original law}

Suppose a measurable function \(\alpha:[0,1]\to[0,\infty)\)
has the following bound for all sources:
\begin{equation}\label{eq:v70-budget-hypothesis}
 \operatorname{Cov}_{\theta_{\tau,b}}(X)
              \preceq\alpha(\tau)I_E
   \quad\text{for every }b\in E,\quad
   \text{for almost every }\tau\in[0,1].
\end{equation}
Set
\begin{equation}\label{eq:v70-integrated-budget}
                      \mathcal A=\int_0^1\alpha(\tau)\,d\tau .
\end{equation}
The condition is on the complete interacting
posteriors, including all linear tilts in the
observed soft variables. The endpoint convexity
theorem alone does not give this bound on
the earlier interval.

\begin{theorem}[Variance and entropy retained through the pinning interval]
\label{thm:v70-conservation}
If \(\mathcal A<\infty\), then
\begin{equation}\label{eq:v70-conservation}
\begin{aligned}
 \mathbb E\Var_{\theta_1}f&\ge e^{-\mathcal A}\Var_\mu f,\\
 \mathbb E\operatorname{Ent}_{\theta_1}f
                          &\ge e^{-\mathcal A}\operatorname{Ent}_\mu f.
\end{aligned}
\end{equation}
The first holds for bounded real \(f\); the second
first holds for bounded strictly positive \(f\)
and extends to nonnegative bounded \(f\) by approximation.
\end{theorem}
\begin{proof}
For any unit vector \(v\in E\), Cauchy--Schwarz gives
\[
 |\operatorname{Cov}_{\theta}(f,\langle v,X\rangle)|^2
          \le\Var_\theta f\,
                    \Var_\theta\langle v,X\rangle .
\]
Taking the supremum over \(v\) and using
\eqref{eq:v70-budget-hypothesis} bounds
\(\|c_\tau(f)\|^2\) by
\(\alpha(\tau)\Var_{\theta_\tau}f\).
Thus \(r(\tau)=\mathbb E\Var_{\theta_\tau}f\)
is absolutely continuous with
\(r'(\tau)\ge-\alpha(\tau)r(\tau)\), by the
exact variance identity. Gronwall's inequality
proves the first conclusion.

For the entropy bound, fix \(\tau,b\).
Every linear tilt in \(X\) changes \(b\) to
another allowed source. Integrating the covariance
bound along that source line yields
\[
 \log\theta_{\tau,b}
       e^{\langle v,X-\theta_{\tau,b}X\rangle}
                          \le\tfrac12\alpha(\tau)\|v\|^2
                                                   \quad(v\in E).
\]
Let \(\widehat\theta=f\theta/(\theta f)\).
The entropy variational inequality, obtained by
Jensen's inequality under \(\widehat\theta\), gives
\[
 \langle v,\widehat\theta X-\theta X\rangle
       -\tfrac12\alpha(\tau)\|v\|^2
                  \le\frac{\operatorname{Ent}_\theta f}{\theta f}.
\]
Optimizing in \(v\) gives
\begin{equation}\label{eq:v70-entropic-covariance}
       \frac{\|\operatorname{Cov}_\theta(f,X)\|^2}{\theta f}
                            \le2\alpha(\tau)\operatorname{Ent}_\theta f.
\end{equation}
When \(\alpha(\tau)=0\), \(X\) is constant under
the posterior and the same statement is immediate.
The exact entropy identity now gives
\[
 \frac{d}{d\tau}\mathbb E\operatorname{Ent}_{\theta_\tau}f
                   \ge-\alpha(\tau)\mathbb E\operatorname{Ent}_{\theta_\tau}f.
\]
Another Gronwall argument proves the second
conclusion. Replacing \(f\) by \(f+\varepsilon\)
and sending \(\varepsilon\downarrow0\) gives the
bounded nonnegative extension.
\end{proof}

The all-source qualification is substantive for
the entropy step. A covariance estimate at \(b=0\)
alone would not control the tilted Laplace transform
used in \eqref{eq:v70-entropic-covariance}.

\subsection{Endpoint inequalities on the actual convex domain}

\begin{lemma}[The endpoint functional inequalities retain the boundary]
\label{lem:v70-endpoint-LSI}
For every \(b\), smooth real \(f\), and smooth
positive \(h\) with the indicated finite integrals,
\begin{equation}\label{eq:v70-endpoint-LSI}
\begin{aligned}
 \Var_{\theta_{1,b}}f
       &\le\theta_{1,b}
                    \langle\nabla f,\mathsf J^{-1}\nabla f\rangle,\\
 \operatorname{Ent}_{\theta_{1,b}}h
       &\le\tfrac12\theta_{1,b}
           \frac{\langle\nabla h,\mathsf J^{-1}\nabla h\rangle}{h}.
\end{aligned}
\end{equation}
Gradients are on \(\mathcal H_M\).
\end{lemma}
\begin{proof}
Whiten by \(x=\mathsf J^{1/2}z\). The transformed
potential minus \(\|x\|^2/2\), extended by
infinity outside the transformed convex support,
has a proper lower-semicontinuous convex version.
The closure changes no density on the interior;
the boundary of the full-dimensional convex set
has zero Lebesgue measure.

Here is one way to retain this support in the
usual strong-convexity argument. Approximate the
convex remainder by its Moreau envelopes and
then smooth by convolution. Its affine minorant
provides a common Gaussian-integrable upper
bound for the resulting densities. The smooth
potentials on the whole space have Hessian
at least the identity. For their reversible
diffusion, the gradient/curvature identity gives
Fisher-information decay at rate two. Integrating
the entropy dissipation yields
\(\operatorname{Ent}h\le\frac12\int|\nabla h|^2/h\).
The same argument, or its expansion at a constant,
gives the Poincar\'e inequality with constant one.
Convergence of the normalized densities and the
common integrable bound pass these inequalities
to bounded smooth tests on the convex support.
Truncation and Sobolev approximation give the
stated finite-integral extensions.
Undoing the whitening proves
\eqref{eq:v70-endpoint-LSI}.
\end{proof}

The convex-support version of this classical
Bakry--\'Emery criterion is also stated in
Schlichting,
\href{https://arxiv.org/pdf/1812.06464}
{Poincar\'e and Log--Sobolev Inequalities for Mixtures,
Theorem A1}.
The preceding argument explains the treatment
of the possibly vanishing determinant boundary
in the present application.

\begin{theorem}[Functional inequalities for the original restricted Wilson law]
\label{thm:v70-global-functional}
Under \eqref{eq:v70-budget-hypothesis},
\begin{equation}\label{eq:v70-global-functional}
\begin{aligned}
 \Var_\mu f
    &\le e^{\mathcal A}\,
            \mu\langle\nabla f,\mathsf J^{-1}\nabla f\rangle,\\
 \operatorname{Ent}_\mu h
    &\le\tfrac12e^{\mathcal A}\,
             \mu\frac{\langle\nabla h,\mathsf J^{-1}\nabla h\rangle}{h}.
\end{aligned}
\end{equation}
These estimates concern the actual \(\mu\), not
only its convex conditionals.
\end{theorem}
\begin{proof}
Use \eqref{eq:v70-conservation}, apply the endpoint
inequalities inside the posterior expectation,
and then use \(\mathbb E\theta_1=\mu\).
The gradient integrands are functions of \(z\)
only and \(\mathsf J\) is independent of the
random source. This proves both inequalities,
first for bounded smooth tests and then by the
same approximations as above.
\end{proof}

The terminal source marginal itself still displays
the original soft-sector limitation. Its negative
logarithmic density is
\(\|b\|^2/2-\log\mathcal N_1(b)\), up to a constant,
so \eqref{eq:v70-posterior-derivative} and the endpoint
covariance inequality give
\begin{equation}\label{eq:v70-source-curvature}
 D_b^2[-\log p_1(b)]
   =I_E-\operatorname{Cov}_{\theta_{1,b}}(X)
                   \succeq I_E-\mathsf B\mathsf J^{-1}\mathsf B^* .
\end{equation}
On an observed mode, its lower multiplier is
\[
 1-\frac{k-\lambda/4}{\lambda/4+g}
                         =\frac{\lambda/2+g-k}{\lambda/4+g}.
\]
The sign of this bound is precisely the old soft
sign. A negative lower bound is not a proof of
negative curvature of the actual source law.
It does show why simply assigning the free Gaussian
curvature to \(p_1\) would remove an unpaid term.

\subsection{An exact auxiliary sampler, with its missing rate quantified}

Define a Markov kernel \(P\) on \(D\) as follows:
from \(z\), draw \(b=\mathsf Bz+G\), where
\(G\sim N(0,I_E)\), and then draw
\(z'\sim\theta_{1,b}\).
This is a formal exact kernel even if an exact
implementation of the second draw is unavailable.

\begin{theorem}[A rate for the actual two-stage auxiliary kernel]
\label{thm:v70-sampler}
The kernel is \(\mu\)-reversible and positive
as an operator on \(L^2(\mu)\). It satisfies
\begin{equation}\label{eq:v70-sampler-Dirichlet}
 \langle f,(I-P)f\rangle_\mu
                         =\mathbb E\Var_{\theta_1}f.
\end{equation}
Consequently
\begin{equation}\label{eq:v70-sampler-rate}
 \operatorname{gap}(P)\ge e^{-\mathcal A},\qquad
 \operatorname{Ent}_\mu(P h)
              \le(1-e^{-\mathcal A})\operatorname{Ent}_\mu h .
\end{equation}
The second assertion holds for nonnegative
bounded \(h\). This is an auxiliary update rate,
not a rate of physical Euclidean-time transfer.
\end{theorem}
\begin{proof}
Let \(T:L^2(\mu)\to L^2(p_1)\) be
\((Tf)(b)=\theta_{1,b}f\). Its adjoint is
conditional expectation over \(b\) given \(z\).
The specified two-stage kernel is \(P=T^*T\).
Thus it is self-adjoint and positive, preserves
constants, and has invariant law \(\mu\).
Moreover
\[
 \langle f,(I-P)f\rangle
                  =\mu f^2-\mathbb E(\theta_1 f)^2
                  =\mathbb E\Var_{\theta_1}f.
\]
The first conservation inequality proves the
spectral-gap bound, with its displayed convention.

Conditional Jensen gives
\(\operatorname{Ent}_\mu(T^*q)\le\operatorname{Ent}_{p_1}q\).
The exact entropy chain rule gives
\[
 \operatorname{Ent}_{p_1}(Th)
       =\operatorname{Ent}_\mu h
                             -\mathbb E\operatorname{Ent}_{\theta_1}h .
\]
Take \(q=Th\) and use the second conservation
inequality. This proves the entropy contraction.
\end{proof}

\subsection{What is proved without a new infrared estimate}

The finite support gives an explicit admissible
covariance budget. If \(z,z'\in D\), then
\[
 \|z-z'\|_2^2
       \le4\beta R^2(4M^4)=16\beta R^2M^4,
 \qquad
 \|\mathsf Bz-\mathsf Bz'\|^2\le16k\beta R^2M^4 .
\]
Every scalar projection of \(X\) therefore has
range at most \(4\sqrt{k\beta R^2M^4}\).
The elementary variance bound by one quarter
of the squared range holds for every probability
measure on that support, including every posterior
and every tilt. Hence
\begin{equation}\label{eq:v70-diameter-budget}
       \alpha(\tau)=4k\beta R^2M^4,\qquad
                         \mathcal A_{\mathrm{diam}}=4k\beta R^2M^4
\end{equation}
is proved without any further hypothesis.
The auxiliary gap lower bound is thus positive
at every fixed volume, but this estimate deteriorates
exponentially in volume at fixed \(R,\beta\).
The inequality does not prove that the actual gap
has that deterioration.

For the differential functional inequalities,
there is a better elementary benchmark. Put
\begin{equation}\label{eq:v70-oscillation}
 \psi(z)=\tfrac12\langle z,\mathsf Kz\rangle,\qquad
       \mathcal B=\operatorname{osc}_D\psi
                                     \le2k\beta R^2M^4 .
\end{equation}
Indeed \(0\in D\), so the infimum of \(\psi\) is
zero. Since \(z\) is mean-zero in every orientation,
the cross term with the homogeneous \(a_c\) vanishes:
\[
 \beta^{-1}\|z\|_2^2
   =\sum_e|a_c(e)+\beta^{-1/2}z_e|^2-\sum_e|a_c(e)|^2
                                             \le4R^2M^4.
\]
Now use \(\mathsf K\preceq kI\) to obtain
\eqref{eq:v70-oscillation}.

\begin{proposition}[A stronger elementary finite-volume baseline]
\label{prop:v70-baseline}
The inequalities in \eqref{eq:v70-global-functional}
remain true with \(e^{\mathcal A}\) replaced by
\begin{equation}\label{eq:v70-minimum-cost}
              e^{\min\{\mathcal A,\mathcal B\}}.
\end{equation}
In particular \(e^{2k\beta R^2M^4}\) is an
unconditional proved coefficient for both of
those differential inequalities.
\end{proposition}
\begin{proof}
Let \(\nu=\theta_{1,0}\).
Then \(d\mu/d\nu\) is a constant times \(e^\psi\),
whose supremum-to-infimum ratio is \(e^{\mathcal B}\).
If \(a\le d\mu/d\nu\le b\), the variational formula
for variance gives
\(\Var_\mu f\le b\Var_\nu f\).
For entropy, use its nonnegative integrand
\[
 \operatorname{Ent}_\mu h
    =\inf_{q>0}\mu\{h\log(h/q)-h+q\}
                     \le b\operatorname{Ent}_\nu h .
\]
Apply the endpoint functional inequality for \(\nu\),
and compare its nonnegative gradient integrand
back to \(\mu\) with factor \(1/a\).
The ratio \(b/a\) can be taken arbitrarily
close to \(e^{\mathcal B}\). Taking the better
of this estimate and Theorem~\ref{thm:v70-global-functional}
proves the proposition.
\end{proof}

Thus the crude diameter budget is not presented
as an improvement on bounded perturbation.
The additional content of the observation route
is its exact source-variance/entropy accounting,
its auxiliary two-stage rate, and its opportunity
to replace the volume-sized loss by a genuinely
interacting all-tilt covariance estimate.
Such an improvement has not yet been proved.
Reducing the rank of the pin does not, by itself,
bound the norm of the covariance of its collective
coordinates.

\subsection{A sharp Gaussian test and the unchanged physical obstruction}

A scalar Gaussian model makes the budget visible
without any inequality slack. This is a regression
model, not a Wilson density. Let \(Z\sim N(0,v)\)
and \(X=\sqrt{\kappa}\,Z\), with \(v,\kappa>0\).
The posterior variance is
\[
 v_\tau=\frac{v}{1+\tau\kappa v},\qquad
 \alpha(\tau)=\frac{\kappa v}{1+\tau\kappa v},\qquad
       \mathcal A=\log(1+\kappa v).
\]
For the linear test \(f(z)=z\),
\[
         \mathbb E\Var_{\theta_1}f
                       =v_1=e^{-\mathcal A}v.
\]
The conservation factor is exactly attained.
For the positive test
\(h(z)=\exp(jz-j^2v/2)\), Gaussian integration gives
\[
 \operatorname{Ent}_\mu h=\tfrac12j^2v,\qquad
 \mathbb E\operatorname{Ent}_{\theta_1}h=\tfrac12j^2v_1.
\]
So the entropy factor is also exactly attained.
These Gaussian identities follow directly by
integration; bounded-support approximation is
not being used to conceal a change of input law.

The two-stage sampler in this model is Gaussian
autoregression with linear coefficient
\(\varrho=\kappa v/(1+\kappa v)\) and stationary
variance \(v\). Its Hermite modes have eigenvalues
\(\varrho^n\). This follows by applying the kernel
to the Gaussian exponential generating function;
the Hermite polynomials form a complete orthogonal
basis. Therefore
\[
                   \operatorname{gap}(P)
                           =1-\varrho=(1+\kappa v)^{-1}
                           =e^{-\mathcal A}.
\]
As \(v\to\infty\), the pinned conditional precision
\(v^{-1}+\kappa\) stays bounded below by \(\kappa\),
but the actual auxiliary gap tends to zero.
The diverging covariance budget accounts for the
loss; the conditional lower floor alone does not.

For Yang--Mills the corresponding unproved estimate
is \eqref{eq:v70-budget-hypothesis} with a suitably
controlled integral across the soft pinning interval
and across the intended cutoff and volume regime.
Neither free Fourier covariance nor the endpoint
convexity inequality supplies it. The posterior
law at \(\tau<1\) still contains the full nonlinear
Wilson action and the actual positive gauge domain.

Even a volume-uniform \(\mathcal A\) would not
alone prove the physical mass gap. The differential
Dirichlet form here uses \(\mathsf J^{-1}\), a
preconditioned, nonlocal auxiliary metric.
Its Euclidean conversion has lower rate only
\[
           e^{-\mathcal A}\bigl(\lambda_{\min}(\Delta)/4+g\bigr),
\]
which may vanish in the volume/coupling limit.
The kernel \(P\) is likewise an auxiliary resampling
operator, not reflection-positive physical transfer.
Coverage of the unrestricted compact theory,
continuum noncollapse, and calibrated physical-time
control are still required.

\subsection{Exact checks and preservation}

The certificate
\begin{center}\small
\path{scripts/verify_ym_f14_v70_soft_observation.py}
\end{center}
uses exact rational and complex-rational arithmetic.
For each of \(k=1/5,3/5,1\), it checks all 15
nonzero Fourier blocks at side two and all 255
at side four. It verifies the positive pin,
its exact active set, transversality, the raised
lower form and the corresponding scalar source
curvature identity. At side four the pin ranks
are \(0,72,72\), with traces \(0,36/5,36\).
These are finite algebra checks, not a fitted
value of the universal Wilson constant \(C_0\).

An independent four-atom calculation checks the
normalized It\^o likelihood cancellation,
the field covariance drift, and the variance
and entropy drifts for a nonlinear positive
test. Its test mean is \(49/5\), its observed
covariance vector is \((223/100,114/25)\),
and the two drifts are
\[
              -51533/2000,\qquad -51533/39200.
\]
This finite law tests the general filtering
algebra, not a discretized Yang--Mills integral.
The scalar Gaussian check uses \(v=3/2\),
\(\kappa=1/4\), \(j=2/3\), and yields endpoint
variance \(12/11\), budget exponential \(11/8\),
exact auxiliary gap \(8/11\), and entropy values
\(1/3\) and \(8/33\).
The continuum-time martingale arguments and the
uniform-in-source functional inequalities are
proved analytically above, not inferred from
these finite tests.

The 5,964-byte certificate has SHA--256
\begin{center}\scriptsize\ttfamily
1B8794086A7930DCE3AFA3BDA108A63F2BFCDF5E8484166899986FB9344674FD.
\end{center}
Its V56 and V58 helpers are unchanged.
The predecessor V69 has 1,709,602 bytes and SHA--256
\begin{center}\scriptsize\ttfamily
A01110F2C14FE20A8F2C26665B22C4D172BC7B24317211211FACD60B81836251.
\end{center}
Its body is preserved except for the title version
and this terminal insertion. V70 is the single
active main note. Companion files, attached
sources, monographs and archives are unchanged.

\begin{firewall}
The complete restricted Wilson law now has a
soft-only observation decomposition with uniformly
convex endpoint conditionals, exact variance and
entropy identities, and proved finite-volume
functional inequalities and auxiliary update bounds.
The unconditional constants supplied here are
volume-dependent; a uniform interacting covariance
budget is not proved. No full continuum Yang--Mills
construction or physical Hamiltonian mass gap follows.
The original goal remains active.
\end{firewall}
\section{V71: the unpinning equation and a nonlinear Wilson cumulant sign}

\subsection{Context: testing the proposed covariance closure}

V70 isolated an integrated, all-source covariance bound
which would transfer its convex endpoint estimates to
the original restricted Wilson law. This note tests two
possible ways of proving that bound. The first is to
propagate the endpoint covariance bound backwards through
the Gaussian pinning equation. We prove that the direct
matrix comparison closes through the whole interval
only in the regime where the earlier Wilson lower form
was already positive. Thus that argument alone does not
resolve the infrared problem.

The second test uses a genuinely nonlinear coefficient
of the full Wilson measure. V59 already derived the
connected cubic/quartic source polynomial, so it is not
recomputed as a new result here. We adapt that identity
to a partially pinned covariance and justify the
fixed-volume expansion on the actual positive gauge
domain. For a Cartan-polarized soft mode, the cubic
exchange vanishes and the fourth cumulant is strictly
positive at sufficiently weak coupling. This gives
an explicit sign in the covariance evolution and rules
out a proposed nonpositive-fourth-cumulant shortcut.
It is not a proof that all covariance-control methods
fail, nor a mass-gap calculation.

\subsection{The exact backwards heat equation in the source}

Retain the notation \(\mathcal N_\tau(b)\),
\(\theta_{\tau,b}\), \(X=\mathsf BZ\), \(E\),
\(\mathsf K\), and \(\mathsf J\) of V70. In this
subsection \(\mu\), including its volume, coupling
and domain, is fixed. Put
\[
 \psi_\tau(b)=\log\mathcal N_\tau(b),\qquad
 Q_\tau(b)=D_b^2\psi_\tau(b)
                   =\operatorname{Cov}_{\theta_{\tau,b}}(X).
\]
Derivatives and the Laplacian below are on \(E\),
not on the physical lattice.

\begin{proposition}[An exact source equation with its nonlinear covariance term]
\label{prop:v71-source-PDE}
For \(0\le\tau\le1\),
\begin{equation}\label{eq:v71-source-PDE}
\begin{aligned}
 \partial_\tau\mathcal N_\tau
          &=-\tfrac12\Delta_b\mathcal N_\tau,\\
 \partial_\tau\psi_\tau
          &=-\tfrac12\Delta_b\psi_\tau-\tfrac12|\nabla_b\psi_\tau|^2,\\
 \partial_\tau Q_\tau
          &=-\tfrac12\Delta_b Q_\tau
             -(\nabla_b\psi_\tau\cdot\nabla_b)Q_\tau-Q_\tau^2.
\end{aligned}
\end{equation}
There is also the exact positive integral identity
\begin{equation}\label{eq:v71-source-convolution}
 \mathcal N_\tau(b)
      =\mathbb E_{G\sim N(0,I_E)}
                  \mathcal N_1(b+\sqrt{1-\tau}\,G).
\end{equation}
Every derivative in these formulas retains the
full Wilson weight and its fixed total-field domain.
\end{proposition}
\begin{proof}
The observed vector \(X\) is bounded at the fixed
finite parameters. Differentiating its defining
Laplace integral is therefore legitimate.
The \(\tau\) derivative inserts \(-|X|^2/2\);
the source Laplacian inserts \(|X|^2\).
This gives the first equation. Division by
\(\mathcal N_\tau>0\) gives the second.
Taking two source derivatives gives the third;
in components the Hessian of
\(|\nabla\psi|^2/2\) is
\((D^2\psi)^2+\nabla\psi\cdot\nabla(D^2\psi)\).

For the convolution identity insert the defining
integral for \(\mathcal N_1\). The Gaussian moment
of \(\exp(\sqrt{1-\tau}\langle G,X\rangle)\)
is \(\exp((1-\tau)|X|^2/2)\), which changes
the pin coefficient from one to \(\tau\).
Tonelli's theorem applies to these positive
integrands. Boundedness of \(X\) also proves
finiteness for every real \(b\).
\end{proof}

These deterministic identities are consistent with
the stochastic covariance identity of V70, but
they are not that identity with the random source
held fixed. In particular the diffusion and drift
terms in \eqref{eq:v71-source-PDE} cannot be dropped.
The all-tilt covariance requirement used for entropy
in V70 is the same kind of requirement discussed
in Chen--Eldan,
\href{https://arxiv.org/pdf/2203.04163}
{Localization Schemes, Lemma 40}.
It is not a consequence of a covariance bound
at a single source.

\subsection{What the endpoint comparison actually propagates}

\begin{theorem}[The direct backwards comparison and its exact threshold]
\label{thm:v71-backwards}
Suppose \(Q_1(b)\preceq M_*\) for every \(b\),
where \(M_*\) is a constant nonnegative matrix.
For every \(\tau\) with
\((1-\tau)M_*\prec I_E\),
\begin{equation}\label{eq:v71-backwards}
            Q_\tau(b)
                  \preceq M_*[I_E-(1-\tau)M_*]^{-1}
                                                    \quad(b\in E).
\end{equation}
For the proved endpoint matrix
\begin{equation}\label{eq:v71-endpoint-M}
                   M_*=\mathsf B\mathsf J^{-1}\mathsf B^*,
\end{equation}
write \(a=\|M_*\|\). If \(a<1\), V70's covariance
budget can be bounded by
\begin{equation}\label{eq:v71-budget}
       \alpha(\tau)=\frac{a}{1-a+a\tau},
                 \qquad \mathcal A\le-\log(1-a).
\end{equation}
But \(a<1\) holds exactly when the already known
lower form
\begin{equation}\label{eq:v71-old-positive}
                         \tfrac12\Delta+(g-k)I
\end{equation}
is positive on \(\mathcal H_M\).
Thus this sufficient closure adds no new
positivity regime beyond that lower form.
\end{theorem}
\begin{proof}
Let \(r=1-\tau>0\), and use
\eqref{eq:v71-source-convolution} in Lebesgue
coordinates on \(E\). For fixed \(b\), introduce
the normalized law
\[
 d\zeta_b(y)\ \propto\
          \exp\{-|y-b|^2/(2r)+\psi_1(y)\}\,dy .
\]
It is integrable: the bounded input implies at
most linear growth of \(\psi_1(y)\).
Its potential has Hessian at least
\(r^{-1}I_E-M_*\succ0\).
The covariance inequality for a strongly
log-concave measure gives
\(\operatorname{Cov}_{\zeta_b}(y)
\preceq(r^{-1}I_E-M_*)^{-1}\).
Differentiating the Gaussian kernel in the
same integral gives the exact formula
\[
 Q_\tau(b)=-r^{-1}I_E+r^{-2}\operatorname{Cov}_{\zeta_b}(y).
\]
Substitution and elementary matrix algebra prove
\eqref{eq:v71-backwards}. The case \(r=0\) is
the assumed endpoint bound.

The matrix in \eqref{eq:v71-endpoint-M}
is an endpoint upper bound by V70.
On an observed Laplacian mode its eigenvalue is
\[
                   \frac{k-\lambda/4}{\lambda/4+g},
                                             \qquad \lambda<4k.
\]
This is less than one precisely when
\(\lambda/2+g-k>0\). On the unobserved modes,
\(\lambda\ge4k\) already gives
\(\lambda/2+g-k\ge k+g>0\).
This proves the equivalence, including the
trivial empty observation space.
Taking norms in \eqref{eq:v71-backwards} and
integrating the displayed scalar upper bound
proves \eqref{eq:v71-budget}.
\end{proof}

If \(a\ge1\), the pole in this comparison is
not a singularity of the true posterior family.
The bounded-support Laplace integrals remain
finite and smooth for the entire interval.
It is a failure of this upper comparison.
Conversely, obtaining a better actual endpoint
matrix or exploiting more nonlinear information
could improve the argument; the theorem does
not exclude those possibilities.

\subsection{Removing one real soft coordinate while retaining every other pin}

To test a nonlinear term without conflating it
with a trace over all observation directions,
take \(c=0\). Let
\(\lambda_1=4\sin^2(\pi/M)\), choose
\begin{equation}\label{eq:v71-fixed-volume-assumptions}
 M\ge4,\quad 0<R\le\tfrac14,\quad
              k=C_0R^2<\lambda_1/2,
\end{equation}
and choose a real unit vector \(v\in\mathcal H_M\)
with
\[
 \Delta v=\lambda v,\qquad
                \lambda<4k,\qquad
                \kappa=k-\lambda/4>0.
\]
Every colour vector \(v_e\) is assumed to be a
scalar multiple of the same unit Lie-algebra
vector \(E_3\). Denote the orthogonal rank-one
projector by \(P_v=v\otimes v\), and put
\[
          \mathsf K_\perp=\mathsf K-\kappa P_v\succeq0,
             \qquad X_v(z)=\sqrt{\kappa}\langle v,z\rangle .
\]
This subsection keeps \(M,R\) fixed as
\(\beta\to\infty\). It does not impose a
uniform-in-volume weak-coupling remainder.

Write the connection-coordinate domain and its
rescaled version as
\begin{equation}\label{eq:v71-domain}
\begin{aligned}
 \Omega_R&=\{a\in\mathcal H_M:
                \max_e|a_e|<R,\ \mathcal F_M(a)\succ0\},\\
 D_\beta&=\sqrt{\beta}\,\Omega_R .
\end{aligned}
\end{equation}
They are nonempty open convex sets containing
zero. No worst-case bound for a nonzero \(c\)
is needed in this \(c=0\) example.
The normalized partially pinned posterior is
\begin{equation}\label{eq:v71-one-coordinate}
\begin{split}
 d\nu_{\beta,\tau,j}(z)\ \propto\
 &\mathbf1_{D_\beta}(z)\rho_M(\beta^{-1/2}z)
                   e^{-\beta S(\beta^{-1/2}z)}\\
 &\cdot\exp\{-\tfrac12\langle z,\mathsf K_\perp z\rangle
                  -\tfrac{\tau}{2}X_v(z)^2+jX_v(z)\}\,dz ,
                \qquad 0\le\tau\le1 .
\end{split}
\end{equation}
At \(\tau=1\) this is exactly V70's full pinned
posterior with source \(b=jv\). At \(\tau=0\)
only this one real coordinate is unpinned;
the other pins remain.

\begin{lemma}[An exact scalar variance equation]
\label{lem:v71-scalar-equation}
At \(j=0\), the law of \(X_v\) under
\(\nu_{\beta,\tau,0}\) is symmetric. Put
\[
 \mathsf v_{\beta,\tau}=\nu_{\beta,\tau,0}(X_v^2),\qquad
 \kappa_{4,\beta,\tau}
      =\nu_{\beta,\tau,0}(X_v^4)-3\mathsf v_{\beta,\tau}^2.
\]
Then
\begin{equation}\label{eq:v71-scalar-equation}
          \partial_\tau\mathsf v_{\beta,\tau}
                   =-\mathsf v_{\beta,\tau}^2
                                     -\tfrac12\kappa_{4,\beta,\tau}.
\end{equation}
\end{lemma}
\begin{proof}
A global colour rotation taking \(E_3\) to
\(-E_3\) preserves \(S,\rho_M,\Omega_R\), the
colour-scalar \(\mathsf K\), and the projector
\(P_v\). It changes \(X_v\) to \(-X_v\).
This proves symmetry and zero mean.

Differentiate the normalized integral in
\eqref{eq:v71-one-coordinate} at \(j=0\).
The inserted quantity is \(-X_v^2/2\), so
\[
 \partial_\tau\nu_{\beta,\tau,0}(X_v^2)
       =-\tfrac12\{\nu_{\beta,\tau,0}(X_v^4)
                         -\nu_{\beta,\tau,0}(X_v^2)^2\}.
\]
Substitution of the fourth cumulant gives
\eqref{eq:v71-scalar-equation}. Bounded support
at each fixed \(\beta\) justifies the derivative.
It is also the scalar source equation in
\eqref{eq:v71-source-PDE}, differentiated twice,
with its odd source derivatives zero at \(j=0\).
\end{proof}

A positive fourth cumulant makes the variance
increase faster than the Gaussian Riccati square
when the pin is removed. A negative fourth cumulant
has the opposite effect. This scalar assertion
does not identify the sign of a trace over mixed
fourth cumulants in the full matrix source equation.

\subsection{A justified pinned weak-coupling expansion on the actual domain}

Use the already established homogeneous Wilson
polynomials \(S_3,S_4\) and the complete geometric
quadratic \(\ell_2\):
\[
 S(a)=\tfrac12\langle a,\Delta a\rangle+S_3(a)+S_4(a)+O_M(\|a\|^5),
 \qquad
 \log\rho_M(a)=\ell_2(a)+O_M(\|a\|^3)
                                                       \quad(a\to0).
\]
These are the literal action and density expansions
of V58--V59, not a replacement measure.
The free covariance appropriate to
\eqref{eq:v71-one-coordinate} is
\begin{equation}\label{eq:v71-pinned-free}
 C_\tau=\{\Delta+\mathsf K_\perp+\tau\kappa P_v\}^{-1},
       \qquad h_\tau=C_\tau(\sqrt{\kappa}\,v)
                         =\frac{\sqrt{\kappa}}{\lambda+\tau\kappa}\,v.
\end{equation}

\begin{lemma}[Fixed-volume expansion with source derivatives and the true cutoff]
\label{lem:v71-Laplace}
Under \eqref{eq:v71-fixed-volume-assumptions},
the normalized logarithmic moment generating function
of \(X_v\) under \(\nu_{\beta,\tau,0}\) has its
Gaussian perturbation expansion through order
\(\epsilon^2\), where \(\epsilon=\beta^{-1/2}\),
with error \(O_{M,R,j_0}(\epsilon^3)\) in four
source derivatives on every fixed real interval
\(|j|\le j_0\). The error is uniform for
\(\tau\in[0,1]\).
The coefficients are computed with \(G\sim N(0,C_\tau)\)
and the normalized weight
\begin{equation}\label{eq:v71-expanded-weight}
       1-\epsilon S_3(G)
           +\epsilon^2\{\tfrac12S_3(G)^2-S_4(G)+\ell_2(G)\}.
\end{equation}
In particular, the domain boundary and the complete
Haar/gauge density have been included in this assertion.
\end{lemma}
\begin{proof}
The classical Wilson lower form on \(\Omega_R\)
is at least \((\lambda_1/2-k)I\).
The segment from zero to every \(a\in\Omega_R\)
lies in this convex domain. Since \(S(0)=0\)
and \(DS(0)=0\),
\begin{equation}\label{eq:v71-classical-tail}
              S(a)\ge\tfrac12(\lambda_1/2-k)\|a\|^2 .
\end{equation}
Thus zero is the unique classical minimum on
this domain. The extra pin is nonnegative.
The finite-dimensional density \(\rho_M(a)\)
is bounded above on its bounded coordinate
closure, including where its determinant vanishes.
For fixed real \(|j|\le j_0\), the unnormalized
integrands are consequently bounded by a
constant times
\[
       \exp\{-\tfrac12(\lambda_1/2-k)\|z\|^2
                                     +j_0\sqrt{\kappa}\|z\|\}.
\]
This bound is integrable, uniformly in \(\beta\)
and \(\tau\), and remains so after multiplying
by any fixed polynomial.

Choose a fixed Euclidean ball
\(\{|a|<\delta\}\subset\Omega_R\) on which
the action and density are analytic.
The portion of \(D_\beta\) outside
\(\{|z|<\delta/\epsilon\}\) has a Gaussian
tail \(O_{M,R,j_0}(e^{-c/\epsilon^2})\);
the same is true after four source derivatives.
The analogous polynomial-weighted Gaussian
tails in the reference integrals have this
bound as well.

On the inner ball, Taylor expansion in
\(\epsilon\) is justified uniformly after
integration. More explicitly, the apparent
singularity of \(S(\epsilon z)/\epsilon^2\)
is removed by the identity
\[
 \frac{S(u z)}{u^2}
    =\int_0^1(1-r)\,D^2S(ru z)[z,z]\,dr .
\]
Its first three \(u\) derivatives are bounded
by fixed constants times \(\|z\|^{3}\),
\(\|z\|^{4}\), and \(\|z\|^{5}\), respectively,
when \(0\le u\le\epsilon\) and
\(\|\epsilon z\|<\delta\). The density derivatives
have analogous polynomial bounds. Inequality
\eqref{eq:v71-classical-tail} supplies the
Gaussian bound for the differentiated exponential.
Taylor's integral remainder is therefore
bounded by \(\epsilon^3\) times an integrable
polynomial--Gaussian function, also after
four source derivatives.

The expansion is exactly
\eqref{eq:v71-expanded-weight}. Its Gaussian
quadratic includes the fixed partial pin,
so the covariance is \(C_\tau\).
The leading normalizers are continuous and
strictly positive on the compact set of
\(\tau,j\) under consideration. Dividing the
integrals and taking their real logarithms
therefore preserves the stated differentiable
error. This proves the normalized assertion.
\end{proof}

The small-radius condition in this lemma is
used to justify a global finite-volume Laplace
expansion on \(D_\beta\), not as a new
thermodynamic assumption. Its constants need
not remain controlled when \(M\) grows.
No unproved suppression of other compact
sectors is used because those sectors were
not included in the specified local law.

\subsection{The actual Wilson fourth cumulant is positive in a soft Cartan direction}

\begin{theorem}[A nonlinear source sign with all measure factors retained]
\label{thm:v71-positive-fourth}
Under the preceding assumptions,
\begin{equation}\label{eq:v71-positive-fourth}
 \kappa_{4,\beta,\tau}
   =\frac{1}{\beta}
       \frac{\kappa^2}{(\lambda+\tau\kappa)^4}
                        \sum_p |(dv)_p|^4
                        +O_{M,R}(\beta^{-3/2}),
                           \qquad 0\le\tau\le1 .
\end{equation}
The remainder is uniform in \(\tau\).
The leading coefficient is strictly positive.
Consequently there is a finite
\(\beta_0(M,R,v)\) such that
\(\kappa_{4,\beta,\tau}>0\) for every
\(\beta\ge\beta_0\) and every \(\tau\in[0,1]\).
This is a statement about the full nonabelian
integral in \eqref{eq:v71-one-coordinate};
only the tested source direction is Cartan-polarized.
\end{theorem}
\begin{proof}
We use the connected coefficient of V59,
not its thermodynamic local-curvature limit.
Here is the precise covariance substitution
needed for this application. For a source
\(\ell\), let \(h=C_\tau\ell\), and define
the linear functional \(g_h\) by
\[
                   D^2S_3(z)[h,h]=\langle g_h,z\rangle .
\]
The Gaussian shift \(G\mapsto G+jh\), followed
by the normalized logarithm of
\eqref{eq:v71-expanded-weight}, gives the
fourth-cumulant coefficient
\begin{equation}\label{eq:v71-connected-coefficient}
       -24S_4(h)+3\langle g_h,C_\tau g_h\rangle .
\end{equation}
For the fourth-degree term, V59's variance
formula reduces to
\[
 [j^4]\left\{\tfrac12\operatorname{Var}S_3(G+jh)\right\}
                       =\tfrac18\langle g_h,C_\tau g_h\rangle .
\]
Indeed the \(j^2\) term of the cubic is
\(\frac12\langle g_h,G\rangle\); the cubic
\(j^3\) term is constant and disappears
from the variance. The quartic action
contributes \(-S_4(h)\), and the quadratic
\(\ell_2(G+jh)\) contributes no \(j^4\)
term. Multiplication by \(4!=24\) gives
\eqref{eq:v71-connected-coefficient}.
Lemma~\ref{lem:v71-Laplace} justifies the
same operation on the actual integrals.

For the present \(h=h_\tau\), all its link
values have the same Lie-algebra direction.
The literal Wilson cubic is a sum of
oriented scalar triple products. Its first
derivative at \(h\) in any direction has
two collinear legs and is zero. Homogeneity
then gives \(g_h=0\). This is a cancellation
in the full cubic exchange, not a deletion
of nonabelian fluctuations.

Along the deterministic connection \(u h\),
all four link exponentials in a plaquette
commute. Their product is the exponential
of the scalar curl. Thus the exact
plaquette action is \(1-\cos(u(dh)_p)\)
in the existing normalization, and
\[
                        S_4(h)=-\tfrac1{24}\sum_p|(dh)_p|^4 .
\]
Substitute \(h_\tau\) from
\eqref{eq:v71-pinned-free} in
\eqref{eq:v71-connected-coefficient}.
This proves \eqref{eq:v71-positive-fourth}.
Since \(v\) is a nonzero mean-zero Coulomb
eigenvector with \(\lambda>0\),
\(\|dv\|^2=\lambda\|v\|^2=\lambda\);
at least one curl entry is nonzero.
The coefficient therefore has a positive
minimum on the compact pin interval.
The uniform remainder proves the eventual
strict sign.
\end{proof}

The geometric quadratic does affect the
two-point function at this order. Its
absence from the connected fourth coefficient
does not authorize deleting it from the
normalized measure, the covariance, or a
source quadratic. The cancellation here is
specific to the degree of the connected
observable being computed.

\subsection{An explicit side-four coefficient and the direction of unpinning}

For a fully specified example, enlarge the valid
universal G{\aa}rding constant if necessary so that
\(C_0\ge12\), which preserves its earlier inequalities.
Take
\begin{equation}\label{eq:v71-example-parameters}
 M=4,\qquad R=\sqrt{3/(4C_0)},\qquad
          k=3/4,\quad \lambda_1=2,\quad \kappa=1/4.
\end{equation}
Then \(R\le1/4\) and \(k<\lambda_1/2\).
This choice specifies a permitted radius in
terms of a valid constant; it does not assign
a numerical value to the optimal Wilson constant.

Let the unnormalized real mode be
\[
 v^{(0)}_{(x,1)}=\cos(\pi x_2/2)E_3,\qquad
                       v^{(0)}_{(x,\mu)}=0\quad(\mu\ne1),
       \qquad v=v^{(0)}/\sqrt{128}.
\]
It is mean-zero and Coulomb, since its sole
nonzero orientation is independent of \(x_1\).
Direct differences give
\[
 \Delta v=2v,\qquad \|v^{(0)}\|^2=128,\qquad
 \sum_p|(dv^{(0)})_p|^4=256,\qquad
                         \sum_p|(dv)_p|^4=1/64.
\]
Only the \((1,2)\) plaquettes contribute.
The soft pin has real rank 72 and trace 18;
the scalar experiment removes only one of
these real directions.

Theorem~\ref{thm:v71-positive-fourth} now becomes
\begin{equation}\label{eq:v71-example-cumulant}
       \kappa_{4,\beta,\tau}
             =\frac{1}{4\beta(8+\tau)^4}
                                  +O_R(\beta^{-3/2}).
\end{equation}
The leading Gaussian variance of \(X_v\)
is \(1/(8+\tau)\).
Combining the exact scalar equation with
\eqref{eq:v71-example-cumulant} gives
\begin{equation}\label{eq:v71-example-defect}
 \partial_\tau\mathsf v_{\beta,\tau}
                       +\mathsf v_{\beta,\tau}^{\,2}
        =-\frac{1}{8\beta(8+\tau)^4}
                                  +O_R(\beta^{-3/2}).
\end{equation}
In particular the fourth-cumulant and
Riccati-defect coefficients at the full pin
are \(1/26244\) and \(-1/52488\).
The square on the left is the \emph{actual}
variance squared; its order-\(\beta^{-1}\)
correction has not been omitted.

Reversing the pin clock changes the sign
of the derivative. Thus in this actual
Wilson example the leading nonlinear
correction increases the growth of the
variance under unpinning. It does not
provide a negative correction with which
to close a Gaussian upper Riccati bound.
The example is inside a finite-volume
classically convex regime; it therefore
does not show an instability or a failure
of every covariance bound. It shows that
the favorable fourth-cumulant sign is
not guaranteed even there.

\subsection{Consequences for the next infrared step}

The two tests distinguish a circular estimate
from a model-specific nonlinear fact.
Backwards propagation using only the known
endpoint covariance matrix reaches the whole
interval precisely in an already positive
lower-form regime. The new fully weighted
Wilson example then rules out assuming
nonpositive directional fourth cumulants
as a general repair of that argument.
It does not claim that the whole mixed
fourth-cumulant tensor has one sign, or that
a positive directional coefficient forces
the operator covariance budget to diverge.

A further useful estimate must retain
nonlinear mode coupling, source dependence
or spatial structure beyond these lower
forms and sign assumptions. In particular
the exact matrix evolution contains a
source Laplacian and a drift, and the
unpinning of one direction is not the
simultaneous unpinning of all soft modes.
The finite-\(M\) expansion above is not
uniform in the cofinal volume/coupling
regime and supplies no physical correlation
length by itself.

This changes the proposed route rather than
declaring the covariance problem solved.
The unresolved requirements remain an
interacting infrared estimate with its true
normalization, control of the unrestricted
compact sectors, a nontrivial continuum
limit, and a physically calibrated gap.
The next scheduled companion checkpoint
remains V75.

\subsection{Exact certificate and preservation}

The certificate
\begin{center}\small
\path{scripts/verify_ym_f14_v71_soft_cumulant.py}
\end{center}
reuses V58's literal quaternion plaquette
polynomial. On all 96 plaquettes of side
two and all 1,536 plaquettes of side four,
it verifies the Cartan quartic and the
vanishing of every component of the literal
cubic gradient at that Cartan input.
It also checks the Coulomb condition,
mode normalization, curl energies and
fourth-power sums. The unnormalized
quartic action in both tested modes is
\(-32/3\); the normalized fourth-power
sums are respectively \(1\) and \(1/64\).
Side two is an algebra regression, not
an additional parameter case of the
\(M\ge4\) theorem.

A separate scalar polynomial Wick test
computes the connected fourth cumulant
from normalized raw moments, including
the nonzero perturbed mean. It agrees
with the quartic and cubic-exchange formula
at \(346112/6615\), and its arbitrary
quadratic density contribution cancels.
This scalar regression is not substituted
for the Wilson integral.
Exact side-four source checks give
the rank 72, pin trace 18, endpoint
variance \(1/9\), and the coefficients
in \eqref{eq:v71-example-cumulant} and
\eqref{eq:v71-example-defect}.
Two finite symmetric laws, with opposite
fourth-cumulant signs, independently
check the scalar heat-equation identity.

No finite-\(\beta\) quadrature is asserted
to certify a sign for the Yang--Mills law.
That eventual sign follows from the
proved fixed-volume differentiable
Laplace expansion and its strictly
positive leading coefficient.
The 6,546-byte certificate has SHA--256
\begin{center}\scriptsize\ttfamily
77100ED5A10C8558A25C95C0158C024DFCA7678F128ABF53E4F0E64F84FFAA0B.
\end{center}
The predecessor V70 has 1,740,141 bytes
and SHA--256
\begin{center}\scriptsize\ttfamily
E889DB062130146EA2BAAAF9959178EA5673CD7B8DDD9317623991DE032D185F.
\end{center}
Its entire body is retained apart from
the title version and this terminal
insertion. V71 replaces V70 as the
single active main note; all companion
notes, attached sources, monographs,
archives and earlier certificates are
unchanged.

\begin{firewall}
This note proves the exact source evolution,
the limitation of a direct backwards
covariance comparison, and a positive
nonlinear fourth cumulant of a specified
fully weighted Wilson posterior at fixed
volume and sufficiently large coupling
parameter \(\beta\). It does not prove
a uniform covariance budget, a continuum
Yang--Mills construction, or a physical
Hamiltonian mass gap. The full goal
remains active.
\end{firewall}
\section{V72: complete soft-shell cumulants and the covariance return}

\subsection{Context and the new calculation}

V71 proved a positive fourth cumulant in one
Cartan-polarized direction. It explicitly did not
determine the mixed fourth-cumulant trace in the
simultaneous soft-mode evolution. This note computes
that trace for the complete side-four soft space:
all 72 real modes, including both other colours
and every allowed polarization and phase.
The cubic exchange is retained with the actual
Coulomb covariance and its fixed harmonic modes.

The separate contributions have different signs.
The sum of the quartic contact contributions is
negative, but the positive cubic exchange is
larger. Their complete mixed trace is positive.
Integrating the resulting first-order covariance
equation then gives a normalized inverse-covariance
return from the full pin to the original restricted
law. Its unknown endpoint correction cancels.
This is a finite-volume nonlinear coefficient,
not a continuum mass or a uniform covariance bound.

The calculation uses the existing V58 literal
vertices, V59 connected coefficient, and V71
fixed-domain Laplace justification. It does not
repeat the flat-root coefficient calculation of
the shell companion. In particular this remains
the specified mean-fixed local law, not a root
constraint with a new normal-return determinant.

\subsection{All soft modes are now unpinned together}

Fix exactly the side-four parameters from
\eqref{eq:v71-example-parameters}:
\[
 M=4,\quad c=0,\quad C_0\ge12,\quad
 R=\sqrt{3/(4C_0)},\quad k=3/4,\quad
             \mathsf K=(kI-\Delta/4)_+.
\]
The smallest positive Laplacian eigenvalue is
two, and the next is four. Hence
\begin{equation}\label{eq:v72-soft-space}
       E=\{\Delta=2\}\cap\mathcal H_4,\qquad
       \dim_{\mathbb R}E=72,\qquad
              \mathsf K=\tfrac14 P_E .
\end{equation}
Here \(P_E\) is the orthogonal projector on this
whole real eigenspace. Use
\[
        \mathsf B=\mathsf K^{1/2},\qquad X=\mathsf BZ .
\]
For \(\tau\in[0,1]\), let
\begin{equation}\label{eq:v72-full-posterior}
 d\theta_{\beta,\tau,b}(z)\ \propto\
       e^{\langle b,\mathsf Bz\rangle
                         -\tau\langle z,\mathsf Kz\rangle/2}\,d\mu_\beta(z),
                      \qquad b\in E,
\end{equation}
where \(\mu_\beta\) is the complete restricted
Wilson law on \(D_\beta\) in V71.
At \(\tau=0,b=0\) this is \(\mu_\beta\);
at \(\tau=1\) it is V70's full pinned posterior.
In contrast with V71's scalar experiment,
every soft pin varies with \(\tau\).

Write \(f_{\mathrm c}(n)=\cos(\pi n/2)\) and
\(f_{\mathrm s}(n)=\sin(\pi n/2)\).
An orthonormal real basis of \(E\) is
\begin{equation}\label{eq:v72-basis}
 u_{\nu\mu}^{a,\sigma}(x,\alpha)
     =\frac{\delta_{\alpha\mu}}{\sqrt{128}}\,
                       f_\sigma(x_\nu)E_a,\qquad
 \nu\ne\mu,\quad a\in\{1,2,3\},\quad
                             \sigma\in\{\mathrm c,\mathrm s\}.
\end{equation}
The indices \(\nu,\mu\) run from one to four.
There are \(4\cdot3\cdot3\cdot2=72\) vectors.
Their norm and orthogonality follow by summing
the side-four sine and cosine products. Each
has mean zero, is Coulomb because its field
is independent of its own link coordinate,
and has Laplacian eigenvalue two.
There are eight nonzero momenta of eigenvalue
two, each with nine transverse-colour real
degrees of freedom in the Fourier dimension
count. Thus this is the full space, not
a sample of its polarizations.

Fix the same tested mode as in V71,
\[
                         v=u_{2,1}^{3,\mathrm c}.
\]
For \(u\in E\), let \(X_u=\langle u,X\rangle\).
All expectations and cumulants below are at
\(b=0\) under \(\theta_{\beta,\tau,0}\).
Translation invariance implies \(\mathbb E Z=0\):
the mean is constant and belongs to the
mean-zero space. Consequently \(\mathbb E X=0\).
Define
\begin{equation}\label{eq:v72-trace-definition}
\begin{aligned}
 Q_{\beta,\tau}&=\operatorname{Cov}(X)\quad\hbox{on }E,\\
 \mathcal T_{\beta,\tau}(v)
    &=\sum_{u\ {\rm in}\ \eqref{eq:v72-basis}}
                   \operatorname{cum}_4(X_v,X_v,X_u,X_u).
\end{aligned}
\end{equation}
For centered random variables \(A,B\), the convention is
\[
       \operatorname{cum}_4(A,A,B,B)
          =\mathbb E A^2B^2-\mathbb E A^2\,\mathbb E B^2
                                      -2(\mathbb E AB)^2 .
\]
The trace in \eqref{eq:v72-trace-definition}
does not depend on the orthonormal basis of \(E\).

\subsection{The full mixed vertex, with both parts of the connected coefficient}

The free covariance for simultaneous pinning is
\begin{equation}\label{eq:v72-free-covariance}
 C_\tau=(\Delta+\tau\mathsf K)^{-1}
                     \quad\hbox{on }\mathcal H_4,\qquad
 C_\tau\mathsf B^*u=d_\tau u,\qquad
                         d_\tau=\frac{2}{8+\tau}.
\end{equation}
As usual, extend \(C_\tau\) by zero on longitudinal
and constant link modes when applying it to
an unconstrained polynomial gradient.

For real link fields \(f,h\), define the full
linear representative \(g_{f,h}\) by
\begin{equation}\label{eq:v72-mixed-gradient}
                    D^2S_3(z)[f,h]=\langle g_{f,h},z\rangle .
\end{equation}
This identity can be evaluated on the full
link space before taking the Coulomb projection.
Every basis vector in \eqref{eq:v72-basis}
has one fixed colour. Therefore
\[
                           g_{v,v}=g_{u,u}=0
\]
by the two-collinear-leg cancellation of V71.
This does not imply \(g_{v,u}=0\) when their
colours are different.

\begin{lemma}[Mixed connected coefficient for each real soft mode]
\label{lem:v72-mixed-coefficient}
For every \(u\) in \eqref{eq:v72-basis},
\begin{equation}\label{eq:v72-mixed-coefficient}
 \operatorname{cum}_4(X_v,X_v,X_u,X_u)
     =\beta^{-1}d_\tau^4\{A(v,u)+B(v,u)\}
                                +O_R(\beta^{-3/2}),
\end{equation}
uniformly in \(\tau\in[0,1]\), where
\begin{equation}\label{eq:v72-contact-exchange}
\begin{aligned}
 A(v,u)&=-4[s^2t^2]S_4(sv+tu),\\
 B(v,u)&=2\langle g_{v,u},C_Mg_{v,u}\rangle,\qquad
                                  C_M=\Delta^{-1}\text{ on }\mathcal H_4.
\end{aligned}
\end{equation}
The second expression is nonnegative.
Its internal covariance is independent of
the pin time in this particular soft space.
\end{lemma}
\begin{proof}
Apply V71's fixed-volume Laplace argument to
the nonnegative pin \(\tau\mathsf K\) and to a
source in the finite-dimensional space \(E\).
The same classical Gaussian tail bound,
compact source sets and polynomial derivative
bounds prove the expansion in four source
derivatives uniformly in \(\tau\).
The actual positive domain and complete
\(\rho_M\) are unchanged.

In V59's connected fourth-degree source
polynomial, with covariance \(C_\tau\), the
quartic expression for a shift \(H\) is
\[
                -S_4(H)+\tfrac18\langle g_{H,H},C_\tau g_{H,H}\rangle .
\]
For \(H=sv+tu\), the two pure diagonal
gradients vanish, so
\(g_{H,H}=2st\,g_{v,u}\).
Taking \(\partial_s^2\partial_t^2\), and
then inserting the common external factor
\(d_\tau^4\), gives the contact coefficient
\(-4[s^2t^2]S_4\) and the exchange
\(2\langle g_{v,u},C_\tau g_{v,u}\rangle\).
The quadratic Haar/gauge density contributes
no connected fourth source degree at this
order. It remains part of the measure and
of the two-point correction.

The frequencies of \(v\) are \(\pm e_2\)
in integer side-four Fourier indices.
Those of \(u\) are \(\pm e_\nu\).
The mixed cubic gradient has only sums of
one frequency from each field. Every such
sum is either zero or has Laplacian
eigenvalue four: for different axes it
has two components equal to \(\pm1\);
for the same axis its nonzero component
is two. On the nonzero modes of eigenvalue
four, \(\mathsf K=0\) and \(C_\tau=C_M\).
On zero modes both covariances are zero
because the harmonic field is fixed.
This proves the asserted pin-time
independence and nonnegativity.
\end{proof}

In particular, there is an internal nonzero
mode in the exchange even though both
external source legs belong to the lowest
soft space. Keeping only the literal
quartic action is not the complete
connected four-point calculation.

\subsection{Exact evaluation over the 72 real modes}

\begin{proposition}[The complete contact and exchange sums]
\label{prop:v72-exact-sums}
For the normalized basis in \eqref{eq:v72-basis},
\begin{equation}\label{eq:v72-exact-sums}
 \sum_u A(v,u)=-\tfrac14,\qquad
 \sum_u B(v,u)=\tfrac78,\qquad
                   \sum_u\{A(v,u)+B(v,u)\}=\tfrac58 .
\end{equation}
The decomposition by colour is
\[
\begin{array}{c|r|r|r|r}
 \text{colour sector}&\text{real modes}&\sum A&\sum B&\sum(A+B)\\ \hline
 a=3&24&1/16&0&1/16\\
 a=1,2&48&-5/16&7/8&9/16\\ \hline
 \text{all}&72&-1/4&7/8&5/8 .
\end{array}
\]
\end{proposition}
\begin{proof}
For the same colour, the deterministic
plaquette fields commute. Therefore
\[
             A(v,u)=\sum_p |(dv)_p|^2|(du)_p|^2,\qquad B(v,u)=0.
\]
Only the four vectors with
\((\nu,\mu)=(2,1)\) or \((1,2)\), and either
real phase, contribute to this sum. Each
gives \(A=1/64\). This proves the first row.

For either one of the other two colours,
the exhaustive finite table is as follows.
Multiplicities in this table are for one
fixed colour; they sum to 24.
\[
\begin{array}{c|c|r|r|r}
 (\nu,\mu)&\sigma&\text{multiplicity}& A(v,u)&B(v,u)\\ \hline
 (2,1)&\mathrm c&1&1/192&0\\
 (2,1)&\mathrm s&1&1/64&0\\
 (2,3),(2,4)&\mathrm c&2&-1/16&0\\
 (2,3),(2,4)&\mathrm s&2&0&0\\
 (3,4),(4,3)&\mathrm c,\mathrm s&4&-1/32&0\\
 (1,2)&\mathrm c,\mathrm s&2&1/64&1/32\\
 (3,1),(4,1)&\mathrm c,\mathrm s&4&1/96&1/32\\
 (1,3),(1,4),(3,2),(4,2)&\mathrm c,\mathrm s&8&0&1/32 .
\end{array}
\]
Here is the exact finite evaluation rule for
the table, together with a direct exchange
check. It also specifies the arithmetic
used in the accompanying certificate.
For the four oriented plaquette vectors
\(w_1,\ldots,w_4\), put
\(a_i=|w_i|^2\) and \(d_{ij}=w_i\cdot w_j\).
The literal quaternion expansion used in
V58 gives
\begin{equation}\label{eq:v72-local-quartic}
\begin{aligned}
 -s_4(w_1,w_2,w_3,w_4)
  ={}&\tfrac1{24}\sum_i a_i^2
          +\tfrac14\sum_{i<j}a_i a_j
          +\tfrac16\sum_{i<j}(a_i+a_j)d_{ij}\\
     &+\tfrac12\sum_{i<j}\sum_{\ell\notin\{i,j\}}a_\ell d_{ij}
          +d_{12}d_{34}-d_{13}d_{24}+d_{14}d_{23}.
\end{aligned}
\end{equation}
The cubic is
\(\sum_{i<j<\ell}w_i\cdot(w_j\times w_\ell)\).
Insert \(w_i=s\,v_i^{(0)}+t\,u_i^{(0)}\),
where the superscript denotes the raw sine
or cosine field, without \(1/\sqrt{128}\).
Extract the indicated mixed coefficient
and sum over all \(4^4\cdot6\) positively
oriented plaquettes. Divide the fourth
degree answer by \(128^2\). The sine and
cosine entries on this torus are integers
\(0,1,-1\), so this is an exact rational
evaluation of \eqref{eq:v72-local-quartic}.
It gives the contact column above.

For exchange, differentiate the displayed
cubic in a test field to form the raw
\(g_{v^{(0)},u^{(0)}}\), including its
oriented shifts. Use Fourier series
\[
 g(x)=\sum_n\widehat g(n)e^{2\pi i n\cdot x/4},\qquad
 \langle g,C_Mg\rangle
      =256\sum_{n\ne0}
                  \widehat g(n)^*\,\frac{\Pi(n)}{\lambda(n)}\,\widehat g(n).
\]
The spatial gradient can equivalently
be reconstructed using V58's complete
colour-ordered cubic vertex. Its incoming
frequencies are the two raw cosine or
sine pairs, with coefficients \(1/2\)
or \(\pm i/2\).
Every internal nonzero frequency has
\(\lambda=4\), as just proved.
This yields the exchange column, equal
to \(1/32\) in fourteen spatial/phase
cases per other colour and zero in the
remaining cases.

For example take \(u=u_{1,2}^{1,\mathrm c}\).
The raw mixed gradient has colour \(E_2\);
two of its Fourier rows are
\[
\begin{aligned}
 \widehat g(1,1,0,0)
      &=\tfrac12(-1+i,\ 1-i,\ 0,\ 0)E_2,\\
 \widehat g(1,3,0,0)
      &=\tfrac12(1-i,\ -1-i,\ 0,\ 0)E_2 .
\end{aligned}
\]
The other two rows are their conjugates at
the opposite momenta. Each row is transverse
and has squared norm one. Their raw
covariance contraction is therefore 256,
and the normalized exchange is
\(2\cdot256/128^2=1/32\).
For the same pair, the raw mixed quartic
coefficient is \(-64\), giving
\(-4(-64)/128^2=1/64\).
The same finite rules for each of the
listed disjoint cases give the full table.
The certificate checks every local
gradient against its Fourier reconstruction,
not only these two example rows.

Summing the table gives \(-5/32\) for
contact and \(7/16\) for exchange for
one other colour. The two other colours
give identical contributions. Adding
the same-colour row proves
\eqref{eq:v72-exact-sums}.
\end{proof}

There is one important zero-mode check.
For \(u=u_{2,1}^{1,\mathrm s}\), the raw
mixed gradient is purely harmonic with
Fourier row
\[
                   \widehat g(0)=(0,2,0,0)E_2 .
\]
Its exchange in the present law is zero,
not an inverse zero eigenvalue and not
an omitted positive term. The harmonic
connection is fixed, not integrated.
The squared norms of these normalized
harmonic gradients sum to \(1/8\) over
the two other colours, but they contribute
nothing to \(C_M\). Allowing fluctuating
harmonics would be a different integral
requiring its own treatment.

\subsection{The complete nonlinear covariance drift}

\begin{theorem}[The mixed trace and the simultaneous-pinning defect]
\label{thm:v72-trace}
For the actual posterior
\eqref{eq:v72-full-posterior}, uniformly
for \(0\le\tau\le1\) at the fixed side-four
parameters,
\begin{equation}\label{eq:v72-trace}
       \mathcal T_{\beta,\tau}(v)
                 =\frac{10}{\beta(8+\tau)^4}
                                      +O_R(\beta^{-3/2}).
\end{equation}
Consequently
\begin{equation}\label{eq:v72-drift}
 \partial_\tau\langle v,Q_{\beta,\tau}v\rangle
                  +\langle v,Q_{\beta,\tau}^{\,2}v\rangle
            =-\frac{5}{\beta(8+\tau)^4}
                                      +O_R(\beta^{-3/2}).
\end{equation}
For all sufficiently large \(\beta\), the
complete trace is strictly positive on
the whole pin interval.
\end{theorem}
\begin{proof}
Sum \eqref{eq:v72-mixed-coefficient} over
the finite basis. Equations
\eqref{eq:v72-exact-sums} and
\eqref{eq:v72-free-covariance} give
\[
         d_\tau^4\,\tfrac58
                    =\frac{16}{(8+\tau)^4}\,\tfrac58
                    =\frac{10}{(8+\tau)^4}.
\]
The uniform remainder survives the finite sum.
The leading coefficient has a positive
minimum on \([0,1]\), proving the eventual sign.

At \(b=0\), the source gradient of
\(\psi_{\beta,\tau}\) vanishes by the
translation argument above. The exact
matrix source equation in V71 thus gives
\[
 \partial_\tau\langle v,Q_{\beta,\tau}v\rangle
   =-\langle v,Q_{\beta,\tau}^{\,2}v\rangle
                            -\tfrac12\Delta_b(D_v^2\psi_{\beta,\tau})(0).
\]
The last source Laplacian is precisely
the cumulant trace in
\eqref{eq:v72-trace-definition}. This
proves \eqref{eq:v72-drift}.
\end{proof}

The contact part alone would have given
\(-4/(\beta(8+\tau)^4)\) for the trace,
with the opposite sign. The exchange
contributes \(14/(\beta(8+\tau)^4)\).
Thus its inclusion is essential to the
sign of the complete answer.
At \(\tau=1\), the full trace coefficient
\(10/6561\) is forty times the directional
coefficient \(1/26244\) of V71. This
comparison concerns the leading coefficients
of the same full-pinned posterior; the
unpinning paths away from that endpoint
are different.

No assertion that \(Q_{\beta,\tau}\)
is a scalar matrix is needed in
\eqref{eq:v72-drift}. In particular its
squared-matrix term has not been replaced
by the square of one diagonal entry.

\subsection{A normalized inverse-covariance return with no unknown endpoint coefficient}

Put
\[
                 q_{\beta,\tau}=\langle v,Q_{\beta,\tau}v\rangle .
\]
This is the variance of the observed
soft coordinate \(X_v\), not the variance
of the unscaled field coordinate.

\begin{theorem}[The complete first nonlinear inverse-covariance difference]
\label{thm:v72-return}
At the fixed side-four parameters, uniformly
for \(0\le\tau\le1\),
\begin{equation}\label{eq:v72-return}
 q_{\beta,\tau}^{-1}-q_{\beta,1}^{-1}
      =\tau-1-\frac5\beta
                  \left(\frac1{8+\tau}-\frac19\right)
                                  +O_R(\beta^{-3/2}).
\end{equation}
In particular the original restricted law
and the full pinned law satisfy
\begin{equation}\label{eq:v72-return-endpoints}
             q_{\beta,0}^{-1}
                   =q_{\beta,1}^{-1}-1-\frac5{72\beta}
                                         +O_R(\beta^{-3/2}).
\end{equation}
The endpoint two-point correction, including
its Haar and gauge contributions, cancels
from this difference rather than being set
to zero.
\end{theorem}
\begin{proof}
The fixed-volume expansion gives, on the
finite observation space,
\begin{equation}\label{eq:v72-covariance-expansion}
 Q_{\beta,\tau}
        =q_0(\tau)I_E+\beta^{-1}D_\tau+O_R(\beta^{-3/2}),
                       \qquad q_0(\tau)=\frac1{8+\tau},
\end{equation}
uniformly in \(\tau\). There is no
\(\beta^{-1/2}\) covariance term: at this
order a centered quadratic observable
is paired with the odd cubic Gaussian
polynomial. The complete density starts
at quadratic order in \(\beta^{-1/2}\).
The coefficient \(D_\tau\) includes all
order-\(\beta^{-1}\) action, exchange and
geometric terms. Its entries are continuous
functions of \(\tau\), by the same Gaussian
integrability as in the expansion proof.

Let \(p_\tau=\langle v,D_\tau v\rangle\).
Matrix multiplication in
\eqref{eq:v72-covariance-expansion} gives
\[
 \langle v,Q_{\beta,\tau}^{\,2}v\rangle
        =q_0(\tau)^2+2\beta^{-1}q_0(\tau)p_\tau
                                      +O_R(\beta^{-3/2}).
\]
Off-diagonal entries of \(D_\tau\) have
not been assumed zero; they enter a
square only at the next order.
Insert these uniform expansions into
the integrated form of
\eqref{eq:v72-drift}. Identifying the
order-\(\beta^{-1}\) coefficient proves
that \(p_\tau\) obeys
\begin{equation}\label{eq:v72-coefficient-ODE}
         p_\tau'=-\frac{2p_\tau}{8+\tau}
                              -\frac5{(8+\tau)^4}.
\end{equation}
This argument does not differentiate an
uncontrolled asymptotic remainder.

Multiplying by \((8+\tau)^2\) and integrating
from the full-pin endpoint gives
\[
 (8+\tau)^2p_\tau
                    =81p_1+\frac5{8+\tau}-\frac59 .
\]
Finally \(q_0(\tau)\) is bounded away
from zero on the interval. The uniform
inverse expansion is therefore
\[
 q_{\beta,\tau}^{-1}
          =8+\tau-\beta^{-1}(8+\tau)^2p_\tau
                                          +O_R(\beta^{-3/2}).
\]
Subtract the endpoint expansion. The
quantity \(81p_1\) cancels and yields
\eqref{eq:v72-return}. At \(\tau=0\),
\(5(1/8-1/9)=5/72\), proving
\eqref{eq:v72-return-endpoints}.
\end{proof}

For comparison in the original field
coordinate, \(X_v=\frac12\langle v,Z\rangle\).
If \(V_{\beta,\tau}=\Var_{\theta_{\beta,\tau,0}}
\langle v,Z\rangle=4q_{\beta,\tau}\), then
\[
 V_{\beta,0}^{-1}
       =V_{\beta,1}^{-1}-\tfrac14-\frac5{288\beta}
                                      +O_R(\beta^{-3/2}).
\]
This conversion fixes the normalization;
neither of these inverse variances has
been identified with a physical Hamiltonian
energy.

\subsection{What this resolves and where the programme still needs a new estimate}

The mixed terms left open in V71 are now
computed for the full finite soft space.
They do not reverse the positive directional
sign: the nonabelian exchange instead
increases the complete traced coefficient.
The integrated return is stronger than a
sign test, because it gives the first
nonlinear change of the actual normalized
two-point precision without requiring an
independent evaluation of its endpoint
counterterm.

The result is still at fixed side four
and at the small radius for which the
classical tail bound in V71 applies.
It does not extend that tail bound to
larger radii or to an unrestricted
compact measure. It does not establish
a remainder uniform as the torus grows,
and it says nothing by itself about
the all-source supremum in V70's entropy
budget. Positivity of this trace at zero
source is neither divergence nor boundedness
of that all-source supremum.

The immediate remaining issue is scale
control of the complete nonlinear
return, including its spatial and
source dependence. The finite calculation
shows concretely that neither the
quartic contact term alone nor the
one-direction cumulant can replace
that return. Compact-sector coverage,
continuum noncollapse and physical
transfer calibration remain required
for the full mass-gap objective.
The next companion checkpoint is V75.

\subsection{Exact finite verification and preservation}

The new certificate is
\begin{center}\small
\path{scripts/verify_ym_f14_v72_mixed_soft_cumulants.py}.
\end{center}
It verifies that the 72 real basis vectors
are orthogonal, have raw squared norm
128, and exhaust the eigenvalue-two
soft space. For each basis vector it
evaluates every one of the 1,536 literal
plaquettes. The local cubic gradient
is then independently reconstructed
from V58's Fourier vertex at every
site, orientation and colour: 221,184
component equalities are checked.
The covariance contraction retains
the Coulomb projection and excludes
the fixed harmonic sector.

The quartic polynomial is evaluated
with integer coefficients after
multiplication by 24. Only repeated
literal input patterns are cached;
there are 235 such patterns in this
calculation. No plaquette or source
direction is replaced by a sample.
The script asserts every class of
the displayed table and both colour
sums, then checks the total trace
coefficient \(10/6561\) and covariance
defect coefficient \(-5/6561\) at the
endpoint.

The integrated-coefficient identity
is checked independently at four
rational pin times and three arbitrary
endpoint corrections, \(-2/7,0,5/11\).
In each case the inverse-covariance
difference coefficient is \(-5/72\).
These are algebra tests of cancellation,
not estimates of the actual endpoint
correction \(p_1\).

The 11,071-byte certificate has SHA--256
\begin{center}\scriptsize\ttfamily
D68499067061F03D1D89B3F276F71C24F47129B6387F0D0B6FDC58181B09FA9C.
\end{center}
Its V56 and V58 helpers are unchanged.
The predecessor V71 has 1,765,145 bytes
and SHA--256
\begin{center}\scriptsize\ttfamily
D4B97040E8A0A8DDF0DB9B5A0AD39E7C227CF655AC47508A9DF14F1AFE0643D1.
\end{center}
Its body is preserved except for the
title version and this terminal insertion.
V72 is the single active main note;
earlier mathematical sections, archives,
companion notes, attached sources and
monographs are unchanged.

\begin{firewall}
This note evaluates the complete mixed
soft-shell fourth-cumulant trace and
derives a finite-volume normalized
inverse-covariance return, with both
contact and cubic exchange retained.
The analytic remainder is fixed-volume.
No uniform interacting infrared bound,
continuum Yang--Mills construction, or
physical Hamiltonian mass gap is proved.
The original goal remains active.
\end{firewall}
\section{V73: the volume limit of the soft trace and the cost of its domain}

\subsection{Context: separating a uniform coefficient from a uniform expansion}

V72 determined the complete lowest-shell cumulant
trace on a side-four torus. This note computes
the corresponding coefficient for every \(M\ge4\).
The calculation is symbolic, not an extrapolation
from a list of small tori. Summing both real phases
before simplifying also handles the side-four
frequency coincidence.

The normalized leading coefficient remains
strictly positive and bounded as \(M\) grows.
That is not the missing uniform remainder.
Indeed, maintaining the classical convexity
used in V71 makes this particular coordinate
domain shrink like \(M^{-1}\). We prove an
independent Gaussian coverage obstruction:
the reference Gaussian assigns asymptotically
zero probability to this domain if
\(\beta=o(M^2\log M)\), and asymptotically full
probability if \(\beta/(M^2\log M)\to\infty\).
The latter assertion is still only Gaussian
coverage, not an interacting expansion theorem.

There is also an important nonduplication
point. In this shrinking-domain regime,
V71 already supplies a uniform all-source
covariance budget. We record its specialization
explicitly, rather than present the new
coefficient as solving that already-controlled
case. The remaining difficulty is to leave
the shrinking classical chart while retaining
normalized control of the actual law.
The full target continues to include the
continuum existence and positive physical
gap requirements in the
\href{https://www.claymath.org/millennium/yang-mills-the-maths-gap/}
{Clay Yang--Mills problem}; an auxiliary
covariance calculation is not that target.

\subsection{A lowest-shell family at every volume}

Set
\begin{equation}\label{eq:v73-parameters}
 \begin{gathered}
 V=M^4,\qquad \lambda=\lambda_M=4\sin^2(\pi/M),\qquad
 C_0\ge12,\\
 R_M^2=\frac{3\lambda}{8C_0},\qquad
 k_M=C_0R_M^2=\frac{3\lambda}{8},\qquad
 \kappa_M=\frac{\lambda}{8}.
 \end{gathered}
\end{equation}
Then \(R_M\le1/4\) and
\(\lambda/2-k_M=\lambda/8>0\).
Let \(\Omega_{R_M}\) be V71's mean-zero
Coulomb domain with the link-radius bound
and the positive actual gauge determinant,
and let \(D_{\beta,M}=\sqrt\beta\,\Omega_{R_M}\).
The measure \(\mu_{\beta,M}\) is the same
complete restricted Wilson law, including
the Haar and gauge density and with the
harmonic connection fixed to zero.

The next positive eigenvalue after \(\lambda\)
is \(2\lambda\). To check this, a momentum
with two nonzero coordinates costs at least
\(2\lambda\); a nonzero one-coordinate
momentum other than the two minimizers
costs at least \(4\lambda-\lambda^2\ge2\lambda\).
Equality in the latter comparison occurs
at \(M=4\). Hence
\begin{equation}\label{eq:v73-soft-pin}
 E_M=\{\Delta=\lambda\}\cap\mathcal H_M,\qquad
 \dim_{\mathbb R}E_M=72,\qquad
 (k_MI-\Delta/4)_+=\kappa_M P_{E_M}.
\end{equation}
The real orthonormal basis is
\[
 u_{\nu\mu}^{a,\sigma}(x,\alpha)
   =\sqrt{\frac2V}\,\delta_{\alpha\mu}
       f_\sigma(2\pi x_\nu/M)E_a,\qquad
 \nu\ne\mu,\quad a=1,2,3,\quad
 \sigma=\mathrm c,\mathrm s ,
\]
where \(f_{\mathrm c}=\cos\) and
\(f_{\mathrm s}=\sin\).
As in V72 fix \(v=u_{2,1}^{3,\mathrm c}\).
Put \(\mathsf K=\kappa_M P_{E_M}\),
\(\mathsf B=\sqrt{\kappa_M}P_{E_M}\), and \(X=\mathsf BZ\).
The simultaneous-pinning posterior is
\begin{equation}\label{eq:v73-posterior}
 d\theta_{\beta,M,\tau,b}(z)\ \propto\
 e^{\langle b,\mathsf Bz\rangle-\tau\langle z,\mathsf Kz\rangle/2}
                                  d\mu_{\beta,M}(z),
 \qquad b\in E_M,\quad 0\le\tau\le1.
\end{equation}
All zero-source means vanish by translation
invariance and the mean-zero constraint.

The appropriate free covariance and source
shift are
\begin{equation}\label{eq:v73-free}
 C_{M,\tau}=(\Delta+\tau\mathsf K)^{-1}\text{ on }\mathcal H_M,
 \qquad C_{M,\tau}\mathsf B^*u=d_{M,\tau}u,\qquad
 d_{M,\tau}=\frac{\sqrt{8/\lambda}}{8+\tau}.
\end{equation}
On the fixed harmonic and longitudinal
spaces this covariance is extended by zero.

\subsection{Exact phase-paired contact and exchange polynomials}

Use \(g_{f,h}\), \(A(v,u)\), and \(B(v,u)\)
from \eqref{eq:v72-mixed-gradient} and
\eqref{eq:v72-contact-exchange}. For a
spatial pair and a colour define
\[
 A_{\nu\mu}^{a}
    =V\sum_{\sigma=\mathrm c,\mathrm s}A(v,u_{\nu\mu}^{a,\sigma}),
 \qquad
 B_{\nu\mu}^{a}
    =V\sum_{\sigma=\mathrm c,\mathrm s}B(v,u_{\nu\mu}^{a,\sigma}).
\]
These definitions include both real
phases, unlike V72's phase-resolved table.

\begin{proposition}[The full lowest-shell coefficient at arbitrary side length]
\label{prop:v73-polynomials}
For \(M\ge4\), the same-colour exchange
vanishes. The only nonzero same-colour
contacts are
\[
                  A_{2,1}^{3}=A_{1,2}^{3}=2\lambda^2.
\]
For either other colour the exhaustive
phase-paired table is
\[
\begin{array}{c|r|r|r}
 (\nu,\mu)&\text{number of spatial pairs}
                         &A_{\nu\mu}^{a}&B_{\nu\mu}^{a}\\ \hline
 (1,2)&1&-16+16\lambda-2\lambda^2&4(4-\lambda)^2\\
 (2,1)&1&8\lambda/3&0\\
 (2,3),(2,4)&2&-16&0\\
 (3,1),(4,1)&2&8\lambda/3&16\\
 (1,3),(1,4),(3,2),(4,2)&4&-16+8\lambda&8(4-\lambda)\\
 (3,4),(4,3)&2&-16&0 .
\end{array}
\]
Consequently
\begin{equation}\label{eq:v73-sums}
 \begin{aligned}
 \sum_{u\in E_M\ {\rm basis}} A(v,u)
      &=\frac{-288+112\lambda}{V},\\
 \sum_{u\in E_M\ {\rm basis}} B(v,u)
      &=\frac{448-128\lambda+8\lambda^2}{V},\\
 \sum_{u\in E_M\ {\rm basis}}\{A(v,u)+B(v,u)\}
      &=\frac{8(\lambda^2-2\lambda+20)}{V}.
 \end{aligned}
\end{equation}
The covariance in every nonzero exchange
line is independent of the pin time.
\end{proposition}
\begin{proof}
The mixed gradient frequencies are
\(\pm e_2\pm e_\nu\). For distinct axes
they have eigenvalue \(2\lambda\). For the
same axis they are zero or have eigenvalue
\(4\lambda-\lambda^2\). Both nonzero
eigenvalues are outside \(E_M\), so the
pin does not alter these covariances.
The zero modes are fixed and have
covariance zero.

Here are finite algebraic rules for all
entries of the table. They give a
direct verification for arbitrary \(M\).
Temporarily omit the factor \(\sqrt{2/V}\)
from each external field. For offsets
\(r,s\), the phase identity is
\[
 \sum_{\sigma=\mathrm c,\mathrm s}
 f_\sigma(\theta(x_\nu+r))f_\sigma(\theta(x_\nu+s))
                                  =\cos(\theta(r-s)),
 \qquad \theta=2\pi/M.
\]
The spatial average of the remaining two
fixed cosine legs is
\(\tfrac12\cos(\theta(r'-s'))\).
This holds also at \(M=4\).
Thus phase summation removes the possible
four-wave alias before any simplification.

Expand the literal quartic
\eqref{eq:v72-local-quartic}.
For each monomial choose which two of
its four labelled positions carry \(v\).
Discard choices with incompatible link
orientations or colours, retain all
oriented-edge signs, and multiply its
coefficient by
\[
 -8\cos(\theta(r'-s'))\cos(\theta(r-s)).
\]
The factor \(-8\) is
\(-4\cdot(2/V)^2\cdot V^2/2\):
the contact derivative, four external
normalizations, the lattice sum and
the extra \(V\) in the table, followed
by the fixed-cosine average.
Sum over the six plaquette orientations.
This is exactly \(A_{\nu\mu}^{a}\).
All offset differences are integers,
so these sums are finite Laurent
polynomials in \(z=e^{i\theta}\).
Use
\(\lambda=2-z-z^{-1}\) to obtain the
displayed contact entries. For example,
the two same-colour entries are
\(2(2-z-z^{-1})^2=2\lambda^2\).

For exchange, let \(W_{\nu\mu}^{a}(p,q)\)
be the four-component gradient obtained
from the complete V58 cubic vertex, with
the fixed colour-\(3\), link-\(1\) leg at
\(p=\pm e_2\), the colour-\(a\), link-\(\mu\)
leg at \(q=\pm e_\nu\), and the remaining
colour leg at \(-p-q\).
The plane-wave coefficient of the fixed
cosine is \(1/2\). The two phases of the
other leg obey
\[
 \sum_{\sigma}
      a_\sigma(\varepsilon)\,
             \overline{a_\sigma(\varepsilon')}
                           =\tfrac12\delta_{\varepsilon,\varepsilon'}.
\]
Consequently every signed pair \((p,q)\)
has weight \(1/8\) in the phase-summed
raw gradient norm, and there are no
cross terms. Even when \(+2e_2=-2e_2\)
on the side-four torus, the two other-leg
signs differ and their cross term is zero.

With \(r=p+q\) and
\(D(r)_j=z^{r_j}-1\), the resulting rule is
\begin{equation}\label{eq:v73-exchange-rule}
 B_{\nu\mu}^{a}
    =\sum_{\substack{p=\pm e_2,\ q=\pm e_\nu\\r\ne0}}
       \left\{\frac{|W_{\nu\mu}^{a}(p,q)|^2}{\lambda(r)}
           -\frac{|D(r)^*W_{\nu\mu}^{a}(p,q)|^2}{\lambda(r)^2}\right\}.
\end{equation}
The normalization is
\(2\cdot(2/V)^2\cdot V^2\cdot(1/8)=1\).
The literal cubic formula
\(\sum_{i<j<\ell}w_i\cdot(w_j\times w_\ell)\)
computes every \(W\) in this expression.
In these particular equal-shell pairs
\[
                          D(r)^*W_{\nu\mu}^{a}(p,q)=0.
\]
This identity is verified before simplifying
the denominator; it is not an assumption
that a full-space gradient is always transverse.

For a representative nonzero case,
\((\nu,\mu,a)=(1,2,1)\) and \(p=e_2,q=e_1\)
give
\[
 W=(t,-t,0,0)E_2,\qquad
                  t=(z+1)(z-z^{-1}).
\]
It is transverse since the first two
components of \(D(r)\) coincide.
Its squared norm is
\(2\lambda(4-\lambda)^2\).
Division by \(2\lambda\) and summation
over the four signed pairs gives
\(4(4-\lambda)^2\), the first exchange
entry. For \((\nu,\mu,a)=(3,1,1)\),
one such row is
\[
                       (0,-2(z-1),2(z-1),0)E_2.
\]
It gives four after covariance contraction,
and sixteen after the four sign choices.
Applying the same literal rule to the
remaining disjoint pairs gives the
exchange column in the table.
All these are identities of Laurent
polynomials after multiplying by the
explicit nonzero denominator
\(\lambda(r)^2\).

For one other colour the contact sum is
\(-144+56\lambda-2\lambda^2\), and the
exchange sum is \(224-64\lambda+4\lambda^2\).
There are two such colours.
Adding the same-colour contact \(4\lambda^2\)
proves \eqref{eq:v73-sums}.
\end{proof}

The purely harmonic output is also checked.
For \((\nu,\mu,a)=(2,1,1)\), with
\(p=e_2,q=-e_2\), the raw vertex is
\[
                         (0,-2(z-z^{-1}),0,0)E_2.
\]
The two signed harmonic vertex norms
sum to \(8\lambda(4-\lambda)\) for
one other colour. Including the phase
weights and normalized external fields,
the sum of squared harmonic gradients
over both other colours is
\[
                           \frac{8\lambda(4-\lambda)}{V}.
\]
It equals \(1/8\) at \(M=4\), as in V72,
but has exchange covariance zero.
No singular inverse at \(r=0\) has been used.

\subsection{A bounded nonzero leading coefficient in four dimensions}

Write
\[
 Q_{\beta,M,\tau}=\operatorname{Cov}_{\theta_{\beta,M,\tau,0}}(X),
 \qquad q_{\beta,M,\tau}=\langle v,Q_{\beta,M,\tau}v\rangle ,
\]
and use V72's complete traced fourth
cumulant \(\mathcal T_{\beta,M,\tau}(v)\).
Define
\begin{equation}\label{eq:v73-gamma}
        \Gamma_M=
             \frac{256(\lambda_M^2-2\lambda_M+20)}
                         {M^4\lambda_M^2}.
\end{equation}

\begin{theorem}[The general-volume trace and the precision difference]
\label{thm:v73-return}
For each fixed \(M\ge4\), uniformly in
\(\tau\in[0,1]\) as \(\beta\to\infty\),
\begin{equation}\label{eq:v73-trace}
 \begin{aligned}
 \mathcal T_{\beta,M,\tau}(v)
       &=\frac{2\Gamma_M}{\beta(8+\tau)^4}
                                         +O_{M,C_0}(\beta^{-3/2}),\\
 \partial_\tau q_{\beta,M,\tau}
                  +\langle v,Q_{\beta,M,\tau}^2v\rangle
       &=-\frac{\Gamma_M}{\beta(8+\tau)^4}
                                         +O_{M,C_0}(\beta^{-3/2}).
 \end{aligned}
\end{equation}
Moreover
\begin{equation}\label{eq:v73-return}
 q_{\beta,M,\tau}^{-1}-q_{\beta,M,1}^{-1}
    =\tau-1-\frac{\Gamma_M}{\beta}
                \left(\frac1{8+\tau}-\frac19\right)
                                         +O_{M,C_0}(\beta^{-3/2}).
\end{equation}
The leading coefficient satisfies
\begin{equation}\label{eq:v73-limit}
 \frac{304}{\pi^4}\le\Gamma_M\le5,\qquad
 \Gamma_4=5,\qquad
                \lim_{M\to\infty}\Gamma_M=\frac{320}{\pi^4}.
\end{equation}
No uniform-in-\(M\) remainder is asserted.
\end{theorem}
\begin{proof}
V71's fixed-volume Laplace justification
applies because \(\lambda/2-k_M>0\).
Use the complete V59 connected coefficient
and the pure-colour cancellations exactly
as in V72. Equation \eqref{eq:v73-free}
and Proposition~\ref{prop:v73-polynomials}
give
\[
 d_{M,\tau}^4
   \sum_u\{A(v,u)+B(v,u)\}
   =\frac{64}{\lambda^2(8+\tau)^4}
           \frac{8(\lambda^2-2\lambda+20)}{V}
   =\frac{2\Gamma_M}{(8+\tau)^4}.
\]
This proves the trace expansion.
The exact source equation
\eqref{eq:v71-source-PDE}, with its source
gradient zero at \(b=0\), proves the
second line of \eqref{eq:v73-trace}.

The leading covariance on \(E_M\) is
\((8+\tau)^{-1}I_{E_M}\).
Its order-\(\beta^{-1/2}\) term vanishes
by Gaussian parity. Denote the tested
order-\(\beta^{-1}\) entry by \(p_{M,\tau}\).
As in V72, insert the uniform fixed-\(M\)
expansions into the integrated covariance
equation. This gives
\[
 p_{M,\tau}'=-\frac{2p_{M,\tau}}{8+\tau}
                            -\frac{\Gamma_M}{(8+\tau)^4},
 \qquad
 (8+\tau)^2p_{M,\tau}
      =81p_{M,1}+\frac{\Gamma_M}{8+\tau}-\frac{\Gamma_M}{9}.
\]
Expansion of the reciprocal and subtraction
at \(\tau=1\) proves \eqref{eq:v73-return}.
The complete endpoint coefficient cancels;
Haar and gauge terms have not been discarded.
No scalar-matrix identity for the exact
covariance is needed.

For the bounds, \(0<\lambda_M\le2\) implies
\[
                 19\le\lambda_M^2-2\lambda_M+20\le20.
\]
Also
\[
  M^4\lambda_M^2=16\{M\sin(\pi/M)\}^4
                        \in[1024,16\pi^4].
\]
The lower endpoint follows because
\(\sin x/x\) decreases on \(0<x\le\pi/4\):
its derivative has numerator
\(x\cos x-\sin x<0\), by differentiating
that numerator from zero.
Thus \(M\sin(\pi/M)\) increases from
its value at \(M=4\), and is bounded
above by \(\pi\).
These inequalities prove both bounds
in \eqref{eq:v73-limit}.
Finally \(M^2\lambda_M\to4\pi^2\)
and \(\lambda_M\to0\), giving the limit.
\end{proof}

This is a statement about the iterated
limit of a coefficient. For instance,
\[
 \lim_{M\to\infty}\lim_{\beta\to\infty}
       \beta\,\mathcal T_{\beta,M,\tau}(v)
                       =\frac{640}{\pi^4(8+\tau)^4}.
\]
It is not a statement about an arbitrary
joint sequence \((\beta,M)\).

For comparison, for \(M\ge5\) the same
single-direction Cartan coefficient has
\[
 \sum_p|(dv)_p|^4=\frac{3\lambda_M^2}{2V}.
\]
Indeed, the fourth-power cosine average
is \(3/8\) when \(M\) does not divide four.
V71's formula therefore gives the
single-direction leading cumulant
\[
                          \frac{96}{\beta V(8+\tau)^4}.
\]
It tends to zero after multiplication by
\(\beta\), whereas the complete mixed
trace does not. The same-colour part of
the full trace is
\(256/(\beta V(8+\tau)^4)\) at leading
order and also vanishes in that limit.
The surviving traced coefficient is
genuinely from the other-colour pairs,
with their contact and exchange both retained.
The exceptional directional coefficient
at \(M=4\) remains exactly the one in V71.

If \(\mathsf V_{\beta,M,\tau}
=\Var_{\theta_{\beta,M,\tau,0}}\langle v,Z\rangle\),
then \(\mathsf V=q/\kappa_M\). Thus
\begin{equation}\label{eq:v73-field-return}
 \mathsf V_{\beta,M,\tau}^{-1}
       -\mathsf V_{\beta,M,1}^{-1}
   =\frac{\lambda_M}{8}(\tau-1)
      -\frac{\lambda_M\Gamma_M}{8\beta}
           \left(\frac1{8+\tau}-\frac19\right)
                                  +O_{M,C_0}(\beta^{-3/2}).
\end{equation}
The displayed nonlinear correction to
this difference is proportional to
\(\lambda_M\), not a positive constant
mass term. This does not determine the
unknown absolute endpoint precision,
and no inverse variance here has been
identified with a physical energy.

\subsection{The all-source bound in this regime was already available}

\begin{record}[Specialization of V71, not a new covariance-closure theorem]
\label{rec:v73-old-budget}
For the complete posteriors in
\eqref{eq:v73-posterior}, all \(\beta\ge1\),
all \(M\ge4\), and every \(b\in E_M\),
\begin{equation}\label{eq:v73-old-budget}
      \operatorname{Cov}_{\theta_{\beta,M,\tau,b}}(X)
                          \preceq\frac1{1+\tau}I_{E_M}.
\end{equation}
In particular V70's all-source budget
can be bounded by \(\log2\) in this family.
\end{record}
\begin{proof}
The endpoint matrix in
\eqref{eq:v71-endpoint-M} is scalar on
the observed space, with eigenvalue
\[
              a_M=\frac{\lambda_M/8}{\lambda_M/4+2/(3\beta)}
                                                        \le\frac12.
\]
Apply Theorem~\ref{thm:v71-backwards}.
The scalar function
\(a/(1-a+a\tau)\) increases with \(a<1\);
at \(a=1/2\) it is \(1/(1+\tau)\).
Integrating proves the stated budget.
\end{proof}

This uniform bound and the positive nonlinear
coefficient are compatible: the bound is
not the Gaussian covariance equality.
It would be circular to keep seeking this
same all-source bound as though V73 had
not yet paid it in the shrinking-radius
regime. It does not cover a fixed positive
radius as \(M\to\infty\).
The resulting functional inequality still
has an unscreened Laplacian scale in
its spatial metric, as explained in V70.
It is not a physical mass-gap estimate.

\subsection{An upstream obstruction: the Gaussian probability of the actual domain}

Let \(\gamma_{M,\tau}=N(0,C_{M,\tau})\)
on \(\mathcal H_M\).
This is the Gaussian reference in the
proved coefficient expansion, not the
interacting posterior itself.
Put
\[
 N_M=\lfloor M/2\rfloor^2M^2,\qquad
 c_*=\frac4{65},\qquad
 p(x)=\mathbb P\{|N(0,1)|\le x\}.
\]

\begin{theorem}[Coverage scale of the shrinking classical domain]
\label{thm:v73-coverage}
Uniformly in \(0\le\tau\le1\),
\begin{equation}\label{eq:v73-box-upper}
 \gamma_{M,\tau}(D_{\beta,M})
      \le p\!\left(2\sqrt{\frac{\beta R_M^2}{c_*}}\right)^{N_M}.
\end{equation}
Consequently, for sequences \(\beta_M\ge1\),
\begin{align}
 \frac{\beta_M}{M^2\log M}\longrightarrow0
 &\quad\Longrightarrow\quad
     \sup_{\tau\in[0,1]}
             \gamma_{M,\tau}(D_{\beta_M,M})\longrightarrow0,
                                      \label{eq:v73-coverage-zero}\\
 \frac{\beta_M}{M^2\log M}\longrightarrow\infty
 &\quad\Longrightarrow\quad
     \inf_{\tau\in[0,1]}
             \gamma_{M,\tau}(D_{\beta_M,M})\longrightarrow1.
                                      \label{eq:v73-coverage-one}
\end{align}
No critical constant at order \(M^2\log M\)
is asserted.
\end{theorem}
\begin{proof}
Choose unit oriented plaquette cycles in
colour \(E_3\) in the \(1,2\) plane,
with bases
\[
 x_1,x_2\in\{0,2,\ldots,2\lfloor M/2\rfloor-2\},
 \qquad x_3,x_4\in\{0,\ldots,M-1\}.
\]
For each base \(x\), the four signed link
entries are \(1/2,1/2,-1/2,-1/2\) along
\[
       (x,1),\quad(x+e_1,2),\quad(x+e_2,1),\quad(x,2).
\]
Denote these real fields by \(w_j\).
Each has zero divergence, zero mean,
squared norm one and link \(\ell^1\)
norm two. Their link supports are
disjoint, so they form \(N_M\)
orthonormal vectors in \(\mathcal H_M\).
This construction works for odd as well
as even side lengths, without crossing
the unused final row in the first
two coordinates.

On \(\mathcal H_M\),
\[
     \Delta+\tau\mathsf K\preceq(16+1/4)I=\frac{65}{4}I,
       \qquad C_{M,\tau}\succeq c_* I.
\]
For \(G\sim\gamma_{M,\tau}\) the vector
\(Y_j=\langle w_j,G\rangle\) therefore has
covariance at least \(c_*I_{N_M}\).
It has a representation in distribution
\[
                      Y=\sqrt{c_*}\,\xi+\eta ,
\]
where the coordinates of \(\xi\) are
independent standard real Gaussians
and \(\eta\) is an independent centered
Gaussian with the remaining positive
semidefinite covariance.

Membership in \(D_{\beta,M}\) implies
\(\max_e|G_e|<T=\sqrt\beta R_M\), and hence
\(|Y_j|<2T\) for every \(j\).
For a real Gaussian, the probability of
an interval of fixed length is largest
when it is centered at its mean. This
follows directly by differentiating the
integral of its symmetric decreasing
density over the shifted interval.
Conditional on \(\eta\), the \(\xi_j\)
are independent. Thus
\[
 \mathbb P\{\max_j|Y_j|\le2T\mid\eta\}
                 \le p(2T/\sqrt{c_*})^{N_M}.
\]
Averaging proves \eqref{eq:v73-box-upper}.
The gauge-positivity condition can only
decrease the probability and was not
needed for this upper bound.

For \(x\ge0\), integrating the Gaussian
density just on
\([x,x+(1+x)^{-1}]\) gives the elementary
tail lower bound
\[
                   1-p(x)\ge c\,\frac{e^{-x^2/2}}{1+x}
\]
for an absolute \(c>0\). The additional
exponent across this interval is bounded
by an absolute constant.
Since \(p(x)^N\le\exp[-N(1-p(x))]\),
the right side of \eqref{eq:v73-box-upper}
tends to zero whenever
\(\beta R_M^2=o(\log N_M)\).
Here \(N_M\) is comparable to \(M^4\)
and \(R_M^2\) is comparable to \(M^{-2}\),
with constants depending only on \(C_0\).
This proves \eqref{eq:v73-coverage-zero}.

For the converse, \(C_{M,\tau}\preceq C_M\).
The uniformly bounded diagonal covariance
from V58 and a union bound over all
link-colour coordinates give
\[
 \mathbb P_{\gamma_{M,\tau}}
       \{\max_e|G_e|^2>C(1+\log M+u)\}\le e^{-u}.
\]
The covariance comparison also applies
to every cube oscillation covariance
in Lemma~\ref{lem:v61-cube}.
The same Gaussian determinant bound
and union bound in
Theorem~\ref{thm:v61-gauge-tail} therefore give
\[
 \mathbb P_{\gamma_{M,\tau}}
       \{\|G_1(G)\|^2>C(1+\log M+u)\}\le e^{-u},
\]
uniformly in \(\tau\).
This is a covariance-order application
of the existing Gaussian argument;
no interacting tail estimate is assumed.

Take \(u=2\log M\). If
\(\beta_M/(M^2\log M)\to\infty\), the
first bound implies
\(\max_e|\beta_M^{-1/2}G_e|<R_M\)
with probability tending to one.
The second gives
\(\|\beta_M^{-1/2}G_1(G)\|\to0\)
on the same events.
By \eqref{eq:v65-even}, the even gauge
correction has norm at most
\(\max_e|\beta_M^{-1/2}G_e|^2/2\),
which tends to zero as well.
Thus the actual normalized gauge matrix
\(\mathcal F_M(\beta_M^{-1/2}G)\)
is positive, in fact above \(I/2\),
with probability tending to one.
Both defining conditions of
\(D_{\beta_M,M}\) hold, proving
\eqref{eq:v73-coverage-one}.
\end{proof}

For example, every sequence
\(\beta_M=O(\log M)\) lies on the
zero-coverage side of this theorem.
This does not require identifying such
a sequence with a physical renormalization
trajectory. It is a statement about
the exact Gaussian references and
the specified shrinking domains.

\subsection{What changes in the research direction}

The complete leading soft coefficient is
now controlled at every volume, and its
four-dimensional limit is nonzero.
There is no leading-coefficient infrared
divergence in this 72-mode experiment.
There is equally no proof that the
fixed-\(M\) remainder is harmless in a
joint limit. In the low-coverage regime,
one cannot justify that claim by saying
the Gaussian reference is mostly inside
the classical chart.

The upper-coverage regime is not a
uniform interacting perturbation theorem
either. Small Gaussian probability of
the excluded set does not by itself
control reweighting by the full Wilson
interaction, normalized source derivatives,
or an extensive action remainder.
The theorem supplies a necessary scale
warning for this particular coverage
argument and a sufficient Gaussian
coverage condition, not a new global
construction.

Accordingly, further enumeration within
this shrinking lowest-shell family is
not the missing step. The programme must
return to the fixed-radius retained-field
problem, for example the actual marginal
and Schur complement isolated in
V68--V69, where soft directions are not
removed by making the domain smaller.
The missing estimate must control that
normalized interacting marginal across
scales and must ultimately be connected
to compact-sector coverage and a physical
transfer operator. Neither a bounded
source-space budget on a shrinking chart
nor the coefficient in
\eqref{eq:v73-field-return} supplies that
connection. The next scheduled companion
checkpoint remains V75.

\subsection{Exact verification and preservation}

The new certificate is
\begin{center}\small
\path{scripts/verify_ym_f14_v73_shell_scaling.py}.
\end{center}
It uses rational Laurent polynomials,
with conjugation \(z\mapsto z^{-1}\).
All contact terms are extracted from
the unchanged V58 literal quaternion
quartic. Each of the 96 cubic vertex
rows for the two other colours is
also reconstructed independently from
the literal plaquette cubic before
checking the complete Fourier vertex.
The transverse contractions and harmonic
outputs are checked before division.
All divisions are exact polynomial
identities after conversion to
\(\lambda=2-z-z^{-1}\).

The script asserts every phase-paired
table entry, both colour sums and the
total polynomial. At \(\lambda=2\)
it compares all 36 spatial-colour
phase pairs against V72's independent
side-four local calculation, including
the frequency coincidence.
It checks the integrated precision
identity at three rational values of
\(\lambda\), four pin times and three
arbitrary endpoint corrections.
Those rational values test algebra,
not hypothetical torus eigenvalues.
It also verifies disjoint support,
zero divergence, zero mean and unit
norm for the compact cycles at
sides \(4,5,6,9\); the general
construction is proved above.

The 11,823-byte certificate has SHA--256
\begin{center}\scriptsize\ttfamily
B09EF5E1B6B28EDE220246EE3449C41BE8DAD14C8B3CA0795365EB822A08E958.
\end{center}
Its V58 and V72 inputs and their helpers
are unchanged. The predecessor V72
contains 1,788,065 bytes and has SHA--256
\begin{center}\scriptsize\ttfamily
FDFC950AC3771201EECDB751C144EF8FD54DCA6AB28E08A900EF222122F7CF93.
\end{center}
Its body is preserved except for the
title version and this terminal insertion.
V73 is the single active main note.
Archives, companions, attached sources
and monographs are unchanged.

\begin{firewall}
This note proves an arbitrary-volume
lowest-shell coefficient identity,
its iterated weak-coupling/volume limit,
and a Gaussian coverage scale for the
actual shrinking coordinate domain.
It does not prove a uniform interacting
remainder or coverage of the full compact
Yang--Mills measure. The previously known
all-source bound in this shrinking regime
is recorded as such, not claimed anew.
No continuum Yang--Mills construction or
positive physical Hamiltonian gap is proved.
The original goal remains active.
\end{firewall}
\section{V74: a positive retained channel and its Coulomb locality tradeoff}

\subsection{Context: returning to a fixed-radius domain}

V73 shows why the shrinking classical chart
is not a substitute for a fixed-radius
interacting estimate. We therefore return
to the exact retained-law construction
of V68--V69. The signed kernel in V69
lost a factor \(\chi>1\) in the radius
and could not be inserted into matrix
concavity of the gauge operator.

This note changes the extraction covariance,
not the input Wilson law. Its deterministic
retained map is a positive resolvent
average. It preserves the actual radius
and positive gauge domain exactly, before
the independent noise is added. The
conditional potential has a local scalar
pin and remains uniformly convex in its
Gaussian metric at fixed radius.
These properties lead to an actual-marginal
gauge-floor transfer estimate that requires
an input gauge floor, rather than the
stronger small odd-gauge norm in V69.

There is a cost, which is proved rather
than suppressed: the extracted link
covariance and channel-noise covariance
are not uniformly exponentially local.
For a nontrivial scalar positive average
this obstruction is structural in the
mean-fixed Coulomb representation. Taking
the curl does restore exponential locality
of the independent noise covariance.
The true marginal retains its unresolved
soft curvature. No new physical mass
is inferred from the auxiliary scalar pin.

\subsection{A positive split with a local pin}

Retain the fixed homogeneous \(a_c=c/M\),
the total-field domain \(D=\mathcal D_c\),
and the full potential \(U_c\) from
\eqref{eq:v68-domain} and
\eqref{eq:v69-pinned-potential}.
Use the parameter assumptions of V68,
now with
\begin{equation}\label{eq:v74-parameters}
       0<R<\tfrac14,\qquad
       k_R=C_0R^2,\qquad 2k_R\le s\le1,\qquad
       m=\sqrt s,\qquad g_\beta=\frac2{3\beta}.
\end{equation}
Thus \(R\) can be fixed independently
of \(M\). The original affine-domain
nonemptiness assumptions are unchanged.
All inverses in the following covariance
formulas act on \(\mathcal H_M\).

Define a new pair, denoted with hats to
distinguish it from the massive-curl split:
\begin{equation}\label{eq:v74-split}
 \begin{aligned}
 \widehat H_s&=(\Delta+s)^{-1},&
 \widehat L_s&=\Delta^{-1}-\widehat H_s
                          =s\Delta^{-1}(\Delta+s)^{-1},\\
 \widehat A_s&=\widehat L_s\Delta
                          =s(\Delta+s)^{-1},&
 \widehat\Sigma_s&=\widehat H_s\widehat A_s
                          =s(\Delta+s)^{-2}.
 \end{aligned}
\end{equation}
For the covariances, the full-space
extensions carry the Coulomb projector
and vanish on the fixed harmonic sector:
\[
 \widehat H_s(k)=\frac{\Pi(k)}{\lambda(k)+s}\otimes I_3,\qquad
 \widehat\Sigma_s(k)=\frac{s\,\Pi(k)}{(\lambda(k)+s)^2}\otimes I_3
                                                   \quad(k\ne0).
\]
In contrast the scalar convolution
\(\widehat A_s=s(\Delta+s)^{-1}\) is
extended to every link field, with
constant-mode multiplier one. This
distinction is essential for averaging
and for the locality claim below.

\begin{proposition}[Exact normalized channel for the unchanged Wilson input]
\label{prop:v74-channel}
Let \(Z\sim\mu_c\) be the complete
restricted Wilson law from V68 and
independently let
\(\varepsilon\sim N(0,\widehat\Sigma_s)\)
on \(\mathcal H_M\). Then
\begin{equation}\label{eq:v74-channel}
       Y=\widehat A_s Z+\varepsilon,\qquad
       \eta=(I-\widehat A_s)Z-\varepsilon
\end{equation}
is the normalized lift obtained by
splitting the free covariance as
\(\widehat L_s+\widehat H_s=\Delta^{-1}\)
and reweighting by the original
\(\mathbf1_D e^{-\mathcal R_c}\).
Its retained marginal is exactly
\[
 \widehat\pi_{c,s}(y)
       =\int_D\gamma_{\widehat\Sigma_s}
                    (y-\widehat A_s z)\,d\mu_c(z).
\]
In particular the normalizer of the
input law is unchanged.

Given \(Y=y\), the total field has
fixed-domain conditional
\begin{equation}\label{eq:v74-conditional}
 d\widehat\theta_y(z)\ \propto\
       \mathbf1_D(z)\exp\left\{-U_c(z)-\frac s2\|z\|^2
                           +\langle(\Delta+s)y,z\rangle\right\}\,dz .
\end{equation}
For every \(y\), its potential Hessian
has lower form
\begin{equation}\label{eq:v74-floor}
       \widehat P_s
           =\tfrac12\Delta+(s+g_\beta-k_R)I
           \succeq\tfrac12(\Delta+s)+g_\beta I .
\end{equation}
Thus its centered Laplace covariance
comparison is \(\widehat P_s^{-1}\),
and in particular
\(\widehat\Gamma_y=\operatorname{Cov}_{\widehat\theta_y}Z
\preceq2\widehat H_s\).
\end{proposition}
\begin{proof}
All four operators in \eqref{eq:v74-split}
are positive on \(\mathcal H_M\), and
\[
 \widehat\Sigma_s^{-1}
       =\widehat L_s^{-1}+\widehat H_s^{-1},\qquad
 \widehat\Sigma_s^{-1}\widehat A_s=\widehat H_s^{-1},\qquad
 \widehat A_s^*\widehat\Sigma_s^{-1}\widehat A_s=sI .
\]
Apply the normalized regression and
reweighting argument of
Theorem~\ref{thm:v69-channel} to this
positive covariance pair. The independent
residual Gaussian is unaffected by a
reweighting depending only on \(Z\).
This proves the joint law, marginal
and equality of input normalizers.

Changing from \(\eta\) to \(z=y+\eta\)
on the full space \(\mathcal H_M\)
has Jacobian one and fixes the domain.
The pin is now exactly
\(\widehat H_s^{-1}-\Delta=sI\),
which gives \eqref{eq:v74-conditional}.
The full lower form from V64 and V62 is
\[
          D^2 U_c\succeq\tfrac12\Delta+(g_\beta-k_R)I.
\]
Adding \(sI\), and using \(s\ge2k_R\),
proves \eqref{eq:v74-floor}. Linear
tilts do not change this form.
The convex-support argument used in
V68--V69 gives the centered Laplace
comparison and covariance order.
No source-size bound or discarded
boundary term is involved.
\end{proof}

The input law is still the one with
strict gauge positivity, not a law
newly restricted to a quantitative floor.
The output \(Y\) is unbounded because
of its independent Gaussian noise.
Exact preservation below concerns the
deterministic average, and the later
probability estimate accounts for that noise.

\subsection{The deterministic average preserves the gauge domain}

Let \(\mathcal U_x\) denote translation
on the primed vertex gauge space and
\(\mathcal T_x\) translation on connections.
The operators satisfy
\(\mathcal F_M(\mathcal T_xa)
=\mathcal U_x\mathcal F_M(a)\mathcal U_x^*\).

\begin{theorem}[Positive averaging, gauge floors and geometric gain]
\label{thm:v74-domain}
The scalar kernel of \(\widehat A_s\)
has positive entries \(p_{s,M}(x)\)
with sum one. For every real affine
Coulomb connection \(a\) with
\(\max_e|a_e|<R\),
\begin{equation}\label{eq:v74-Jensen}
 \begin{aligned}
  \max_e|(\widehat A_s a)_e|&<R,\\
  \mathcal F_M(\widehat A_s a)
       &\succeq\sum_x p_{s,M}(x)\,
                 \mathcal U_x\mathcal F_M(a)\mathcal U_x^* .
 \end{aligned}
\end{equation}
Consequently any positive gauge floor
of \(a\) is preserved by this deterministic
map, including mere strict positivity.
It preserves \(a_c\) and the affine
Coulomb slice.

For the complete geometric potential,
\begin{equation}\label{eq:v74-geometric-gain}
 W_M(\widehat A_s a)
       \le W_M(a)-\frac13
                    \{\|a\|_2^2-\|\widehat A_s a\|_2^2\}.
\end{equation}
For the full interacting potential on \(D\),
put \(B_R=\tfrac12\Delta+(g_\beta-k_R)I\).
Then
\begin{equation}\label{eq:v74-energy-Jensen}
 U_c(\widehat A_s z)
       \le U_c(z)-\tfrac12
                    \langle z,B_R(I-\widehat A_s^2)z\rangle .
\end{equation}
The last quadratic form need not be positive.
\end{theorem}
\begin{proof}
The scalar heat kernel is positive and
preserves constants, as proved in V69.
Its resolvent integral gives
\[
 \widehat A_s=\int_0^\infty s e^{-su}e^{-u\Delta}\,du
                   =\sum_x p_{s,M}(x)\mathcal T_x .
\]
The integral of the scalar weight is one.
On a connected finite torus every heat
kernel entry is positive for \(u>0\),
by the nearest-neighbour path expansion.
Thus each \(p_{s,M}(x)>0\).
The radius bound follows by convexity
of the Euclidean colour ball. The
finite set of translated input link
norms has a maximum strictly below \(R\).
Commutation with differences and
preservation of constants prove the
affine-slice assertion.

V62 and its affine extension prove matrix
concavity of \(\mathcal F_M\) on this
link ball. Apply it to the finite convex
combination of translated connections,
then use translation covariance. This
proves the second line of
\eqref{eq:v74-Jensen}. If
\(\mathcal F_M(a)\succeq\delta I\), every
translated operator has that same floor;
their average does too. For strict
positivity choose the positive minimum
eigenvalue at the fixed finite input.

The geometric Hessian is at least
\((2/3)I\). Its strong-convexity Jensen
inequality, together with translation
invariance of \(W_M\), therefore gives
\[
 W_M(\widehat A_s a)
     \le W_M(a)-\frac13\sum_x p_{s,M}(x)
                         \|\mathcal T_xa-\widehat A_s a\|_2^2 .
\]
The variance identity for a convex
combination is
\[
 \sum_xp_{s,M}(x)\|\mathcal T_xa-\widehat A_s a\|_2^2
                          =\|a\|_2^2-\|\widehat A_s a\|_2^2 .
\]
All intermediate connections lie in the
positive domain just proved, so the
strong-convexity argument applies to the
actual geometric potential. This proves
\eqref{eq:v74-geometric-gain}.

Finally \(U_c(z)-\langle z,B_Rz\rangle/2\)
is convex on \(D\). Translations of
\(z\) remain there, \(U_c\) is translation
invariant because \(a_c\) is homogeneous,
and \(B_R\) commutes with translations
and with \(\widehat A_s\).
Ordinary Jensen applied to this convex
function yields \eqref{eq:v74-energy-Jensen}.
This argument allows an indefinite \(B_R\);
it does not replace it by a positive form.
\end{proof}

For clarity the mode coefficient in the
last energy difference is
\[
     (\lambda/2+g_\beta-k_R)
          \left(1-\frac{s^2}{(\lambda+s)^2}\right),\qquad
     1-\frac{s^2}{(\lambda+s)^2}
                      =\frac{\lambda(\lambda+2s)}{(\lambda+s)^2}
                      \le\frac{2\lambda}{s}.
\]
Thus the possibly unfavorable part has
the explicit estimate
\[
          U_c(\widehat A_s z)\le U_c(z)
                              +\frac{k_R}{s}\langle z,\Delta z\rangle.
\]
The geometric gain does not prove a
gain for the full Wilson density.
Nor is a deterministic density comparison
an identity of pushforward normalizers.
The exact normalized marginal is still
the one in Proposition~\ref{prop:v74-channel}.

\subsection{Noise estimates for the actual retained law}

The following local smoothness bound will
allow gauge positivity itself, not a
small odd-gauge operator, to be transferred.

\begin{lemma}[An explicit even-gauge perturbation bound]
\label{lem:v74-even}
If \(\max_e|b_e|\le R\),
\(\max_e|e_e|\le r\), and \(R+r\le1/4\),
then
\begin{equation}\label{eq:v74-even}
 \|G_{\rm e}(b+e)-G_{\rm e}(b)\|
                                  \le3(R+r)r .
\end{equation}
For Coulomb \(b,e\) the full normalized
gauge operators consequently satisfy
\[
 \mathcal F_M(b+e)
    \succeq\mathcal F_M(b)
                      -\{\|G_1(e)\|+3(R+r)r\}I .
\]
\end{lemma}
\begin{proof}
The partial-fraction representation
\eqref{eq:v62-cot}, also recorded in
\href{https://dlmf.nist.gov/4.22#E3}{NIST DLMF 4.22.3},
gives on one colour block
\[
 \sigma(\operatorname{ad}_a)
       =I-2\sum_{n\ge1}
                   \frac{|a|^2I-aa^T}{n^2\pi^2-|a|^2}.
\]
For \(|a|\le q\le1/4\), differentiation
in direction \(h\) gives the norm bound
\[
 \|D\sigma(\operatorname{ad}_a)[h]\|
 \le |h|\left\{
      8q\sum_{n\ge1}\frac1{n^2\pi^2-q^2}
       +4q^3\sum_{n\ge1}\frac1{(n^2\pi^2-q^2)^2}\right\}.
\]
The notation on the left denotes the
derivative of \(a\mapsto\sigma(\operatorname{ad}_a)\).
Use \(n^2\pi^2-q^2\ge n^2\pi^2/2\),
\(\sum(n\pi)^{-2}=1/6\), and
\(\sum(n\pi)^{-4}\le1/(6\pi^2)\le1/54\).
The two sums in braces are at most
\(8q/3+8q^3/27\le(145/54)q<3q\).
The locally uniformly differentiable
series justifies these operations.
Integrate along the segment from
\(b_e\) to \(b_e+e_e\) with \(q=R+r\).
Compression by the fixed isometry \(U\)
in \eqref{eq:v60-gauge} gives
\eqref{eq:v74-even} without a volume
factor. Add the linear odd term and
use self-adjointness on the Coulomb
slice to obtain the lower form.
\end{proof}

\begin{proposition}[Uniform noise tails with the actual channel covariance]
\label{prop:v74-noise}
There is an absolute \(C\) such that,
for \(0<s\le1\), \(M\ge4\), \(u\ge0\),
and \(\varepsilon\sim N(0,\widehat\Sigma_s)\),
\begin{equation}\label{eq:v74-noise-tails}
 \begin{aligned}
 \mathbb P\{\max_e|\varepsilon_e|^2
           >Cs\log(e/s)(1+\log M+u)\}&\le e^{-u},\\
 \mathbb P\{\|G_1(\varepsilon)\|^2
           >C(1+\log M+u)\}&\le e^{-u}.
 \end{aligned}
\end{equation}
For a set \(S\) of links,
\(\operatorname{Tr}(\mathbf1_S\widehat\Sigma_s\mathbf1_S)
\le D_s|S|\), where \(D_s=Cs\log(e/s)\)
may be chosen with a larger absolute \(C\),
and \(\|\widehat\Sigma_s\|\le1/s\).
\end{proposition}
\begin{proof}
The diagonal trace is bounded by a constant
times
\[
                \frac1{M^4}\sum_{k\ne0}\frac{s}{(\lambda(k)+s)^2}.
\]
In four dimensions the number of nonzero
modes with \(\lambda(k)\le t\) is at most
\(CM^4t^2\); if \(t\) is below the first
positive eigenvalue the count is zero.
This follows by counting integer points
in the shortest-representative Fourier
ball, using \(\lambda(k)\ge c|k|^2\).
The part \(\lambda\le s\) contributes
at most \(Cs\). Every dyadic band
\(2^js<\lambda\le2^{j+1}s\) also contributes
at most \(Cs\), and only
\(C(1+\log(1/s))\) such bands intersect
the spectrum \([0,16]\).
This proves the diagonal and trace bounds.
A scalar Gaussian tail and a union
bound over \(12M^4\) real coordinates
give the first line of
\eqref{eq:v74-noise-tails}.

Spectrally,
\(\widehat\Sigma_s\preceq C_M\) and
\(s/(\lambda+s)^2\le1/s\).
Covariance domination in each cube
oscillation, followed by the determinant
tail and the deterministic div--curl
bound of V61, proves the second line
exactly as in V69. No independence of
the physical noise links is assumed.
\end{proof}

\begin{theorem}[A unit-factor radius and gauge-floor transfer for the actual marginal]
\label{thm:v74-transfer}
Let \(\widehat\pi_{c,s}\) be the actual
retained marginal, and put
\[
 L_{M,u}=1+\log M+u,\qquad
 r=C\sqrt{\frac{s\log(e/s)L_{M,u}}{\beta}},\qquad
 e_{\rm g}=C\sqrt{\frac{L_{M,u}}{\beta}}+3(R+r)r .
\]
Choose \(C\) large enough for
Proposition~\ref{prop:v74-noise}.
Suppose \(R+r\le1/4\) and
\(\delta>e_{\rm g}\).
Define the good output event by
\[
 \mathcal G=
 \{y:\max_e|a_c(e)+ty_e|\le R+r,\ 
       \mathcal F_M(a_c+ty)\succeq(\delta-e_{\rm g})I\}.
\]
The gauge expression here is evaluated
only inside the displayed link ball.
Then
\begin{equation}\label{eq:v74-floor-transfer}
 \widehat\pi_{c,s}(\mathcal G^c)
     \le\mu_c\{\lambda_{\min}\mathcal F_M(a_c+tZ)<\delta\}
                                                       +2e^{-u}.
\end{equation}
In particular the deterministic radius
factor is one, not the \(\chi>1\) of V69.

There is also the actual joint bad-link
bound, for any \(r_0>0\) with
\(\beta r_0^2\ge4D_s\),
\begin{equation}\label{eq:v74-joint-bad}
 \widehat\pi_{c,s}
  \{|a_c(e)+tY_e|>R+r_0\text{ for every }e\in S\}
             \le\exp[-s\beta r_0^2|S|/8].
\end{equation}
\end{theorem}
\begin{proof}
By the exact channel and preservation
of the homogeneous part,
\[
                a_c+tY=\widehat A_s(a_c+tZ)+t\varepsilon .
\]
If the input gauge floor is at least
\(\delta\), Theorem~\ref{thm:v74-domain}
puts the deterministic term in the same
radius-\(R\) domain with that same floor.
Outside a noise event of probability
at most \(2e^{-u}\), Proposition~\ref{prop:v74-noise}
gives
\[
           \max_e|t\varepsilon_e|\le r,\qquad
           \|G_1(t\varepsilon)\|\le C\sqrt{L_{M,u}/\beta}.
\]
Lemma~\ref{lem:v74-even} then puts the
output in \(\mathcal G\).
Splitting according to the input-floor
event proves \eqref{eq:v74-floor-transfer}.
The input probability has not been
replaced by a Gaussian probability.

For the joint event, the deterministic
radius bound implies
\(\|\varepsilon_S\|_2^2>\beta r_0^2|S|\).
Use the Gaussian quadratic moment with
parameter \(s/4\). Since
\(\|\widehat\Sigma_{s,S}\|\le1/s\),
\[
 \log\mathbb E e^{s\|\varepsilon_S\|_2^2/4}
     =-\tfrac12\log\det(I-s\widehat\Sigma_{s,S}/2)
     \le\tfrac s2\operatorname{Tr}\widehat\Sigma_{s,S}
     \le\tfrac s2D_s|S|.
\]
Markov's inequality gives exponent
\(-\tfrac s4(\beta r_0^2-2D_s)|S|\),
which is bounded above by the exponent
in \eqref{eq:v74-joint-bad}.
Zero eigenvalues of the restricted
covariance cause no change in this
determinant identity. The empty set
case gives equality one.
\end{proof}

This estimate transfers a quantitative
input gauge floor but does not create one.
The actual input only has strict
positivity; the probability on the
right of \eqref{eq:v74-floor-transfer}
remains to be controlled. In particular
one cannot put \(\delta=0\) and infer
that the added Gaussian noise stays
inside the positive domain.
The output radius is \(R+r\), not a
contracted radius. Nevertheless the
factor-\(\chi\) loss and the need for
a small input odd-gauge norm in V69
have both been removed from this transfer.

\subsection{Why positive averaging cannot retain the old link locality}

\begin{theorem}[A scalar averaging--Coulomb locality obstruction]
\label{thm:v74-no-go}
On \(\mathbb Z^d\), \(d\ge2\), let \(p\)
be a symmetric probability kernel with
\(\sum_x|x|^2p(x)<\infty\), and let
\(a(k)=\sum_xp(x)e^{ik\cdot x}\).
Write \(\lambda(k)=\sum_j|e^{ik_j}-1|^2\)
and \(\Pi(k)\) for the associated
transverse projector at \(k\ne0\).
If the covariance symbol
\begin{equation}\label{eq:v74-no-go-symbol}
                 H(k)=\frac{1-a(k)}{\lambda(k)}\Pi(k)
\end{equation}
has a continuous matrix extension to
\(k=0\), then \(p=\delta_0\).
In particular a nontrivial positive
scalar average with finite second
moment cannot give an absolutely
summable extracted link kernel in
this covariance-splitting form.
\end{theorem}
\begin{proof}
Transversality gives \(q(k)^*H(k)=0\).
Along \(k=tu\), divide by \(t\) and let
\(t\to0\). Since \(q(tu)/t\to iu\),
continuity would imply
\(u^*H(0)=0\) for every real \(u\).
Thus \(H(0)=0\).
Taking traces in \eqref{eq:v74-no-go-symbol}
then gives
\[
                         \frac{1-a(tu)}{\lambda(tu)}\longrightarrow0.
\]
Symmetry and the finite second moment
permit dominated differentiation of
the characteristic function at zero:
\[
 1-a(tu)=\frac{t^2}{2}\sum_xp(x)(u\cdot x)^2+o(t^2),
 \qquad \lambda(tu)=t^2|u|^2+o(t^2).
\]
It follows that every directional
second moment vanishes. Sum over a
coordinate basis to get
\(\sum_x|x|^2p(x)=0\).
Nonnegativity of \(p\) forces all its
mass to be at zero.
Finally an absolutely summable matrix
kernel has a continuous Fourier symbol,
by uniform convergence of its Fourier
series. This proves the last assertion.
\end{proof}

This theorem is specific to a scalar
convolution average and the stated
mean-fixed Coulomb covariance split.
It does not rule out a change of
variables, a map between different
lattices, or a different observable
space. In particular it is not a
restriction on all real-space
renormalization schemes.

For the present channel the failure
is explicit. At fixed \(s>0\), along
the first momentum axis
\((\widehat H_s)_{22}(k)\to1/s\),
whereas along the second axis it is
identically zero. The same two limits
hold for \((\widehat\Sigma_s)_{22}\).
No choice of a zero-mode value repairs
that directional discontinuity.
Thus neither infinite-volume kernel
is absolutely summable.

This also rules out a uniform-in-\(M\)
exponential entrywise bound for these
finite-torus covariances at fixed \(s\).
Indeed, their finite Fourier sums
converge at each fixed separation to
the infinite-volume Fourier integral:
the symbols are bounded and continuous
except at one point, hence Riemann
integrable. A uniform exponential
bound would pass to that limit and
make its kernel absolutely summable,
contradicting the discontinuity.
This is not a claim that
\(\widehat A_s\) itself is nonlocal;
its scalar resolvent kernel is
exponentially local and positive.

\subsection{The independent curl noise is exponentially local}

\begin{proposition}[Locality after taking the exterior derivative]
\label{prop:v74-curl-locality}
Let \(d\varepsilon\) be the linear
lattice curl of the channel noise.
On all real two-form coordinates,
\begin{equation}\label{eq:v74-curl-covariance}
 \operatorname{Cov}(d\varepsilon)
       =s\,d(\Delta+s)^{-2}d^*
       =s\,\mathsf T_m^*\mathsf T_m,
 \qquad \mathsf T_m=d^*(\Delta+m^2)^{-1}.
\end{equation}
For coordinate sets \(E,F\), the
following bounds are uniform in
\(M\ge4\) and \(0<m\le1\):
\begin{equation}\label{eq:v74-locality}
 \begin{aligned}
 \|\mathbf1_E\widehat A_s\mathbf1_F\|
          &\le2e^{-m\operatorname{dist}_M(E,F)/4},\\
 \|\mathbf1_E\operatorname{Cov}(d\varepsilon)\mathbf1_F\|
          &\le16e^{-m\operatorname{dist}_M(E,F)/4}.
 \end{aligned}
\end{equation}
\end{proposition}
\begin{proof}
The identities \(d\Pi=d\), \(\Pi d^*=d^*\),
and commutation with the componentwise
resolvent remove the projector from
\(d\widehat\Sigma_s d^*\).
The zero mode is also annihilated by
the derivatives. This proves
\eqref{eq:v74-curl-covariance}.

Use the weights \(W=e^{mf/4}\) of V68.
Its weighted resolvent bound is
\(\|W(\Delta+s)^{-1}W^{-1}\|\le2/s\),
which gives the first bound with
weighted norm two.
The proof of Theorem~\ref{thm:v68-locality}
gives weighted norms at most \(4/m\)
for both \(\mathsf T_m\) and its
adjoint, using the same base-site
weight on the two form degrees.
Multiplication and \(s=m^2\) give
weighted norm at most sixteen for
the curl covariance. Choose \(f\)
as distance from \(F\) and remove
the weights, as in V68.
\end{proof}

The curl here is a linear observable
of the independent channel noise.
It is not the full nonlinear plaquette
field, and its locality is not decay
of connected correlations under the
interacting input. The exact relation
\[
                     dY=\widehat A_s\,dZ+d\varepsilon
\]
still contains the unknown input
curvature correlations.
This identifies a possible observable
space in which to continue, but does
not assume those correlations are
already controlled.

\subsection{The retained curvature and the remaining task}

The general fixed-domain derivative
formula in V69 applies with hats.
The actual retained potential obeys
\[
 D^2\widehat{\mathcal V}_{c,s}(y)
     =\widehat L_s^{-1}+\widehat H_s^{-1}
              -\widehat H_s^{-1}\widehat\Gamma_y\widehat H_s^{-1}.
\]
Substituting only the proved conditional
covariance bound
\(\widehat\Gamma_y\preceq\widehat P_s^{-1}\)
gives the explicit lower multiplier
\begin{equation}\label{eq:v74-retained-floor}
 \widehat{\mathcal L}_s(\lambda)
   =\frac{(\lambda+s)^2(\lambda/2+g_\beta-k_R)}
                 {s(\lambda/2+s+g_\beta-k_R)}.
\end{equation}
Indeed,
\[
 \widehat L_s^{-1}+\widehat H_s^{-1}
       =\frac{(\lambda+s)^2}{s}
\]
on a transverse mode, and subtraction
of \((\lambda+s)^2/
(\lambda/2+s+g_\beta-k_R)\) gives
\eqref{eq:v74-retained-floor}.
The denominator is positive, so its
sign is exactly the old soft-sector
sign. The positive average and local
pin have not proved new positivity of
the actual marginal.

As a normalization check, with a free
unrestricted Gaussian input the
conditional covariance is massive,
\((\Delta+s)^{-1}\), while the retained
covariance is
\[
                         \widehat L_s(\lambda)
                                    =\frac{s}{\lambda(\lambda+s)}
                                    \sim\lambda^{-1}\quad(\lambda\to0).
\]
This comparison model is not the
restricted Wilson law. It demonstrates
directly why the conditional pin alone
cannot be reported as a physical gap.

The new progress is a domain-preserving
positive extraction interface for the
same fixed-radius input, a quantitative
actual-marginal floor transfer, and an
explicit locality tradeoff with a
curl-noise alternative. The next
unpaid inputs are the distribution of
the actual input gauge floor and the
normalized interacting curvature
response across scales. A successful
continuation must estimate them rather
than insert them as assumptions or
replace the input by its Gaussian
reference. Compact-sector coverage,
continuum construction and physical
transfer calibration remain required.
The next companion checkpoint is V75.

\subsection{Exact checks and preservation}

The new certificate is
\begin{center}\small
\path{scripts/verify_ym_f14_v74_positive_channel.py}.
\end{center}
For both \(s=1/5\) and \(s=1\), on
each actual Fourier block of sides
two and four, it checks the complete
covariance sum, regression identities,
local pin, conditional lower form and
factored retained Schur form. There
are 540 nonzero block checks.
It verifies projector cancellation in
the six-orientation curl covariance
as an exact rational-complex matrix
identity.

Independently, all 544 scalar kernel
entries in these four parameter cases
are obtained by full Fourier summation.
They are positive, sum to one in each
case, and satisfy the local resolvent
equation at every site, including
periodic neighbour coincidences.
For example at \(M=4,s=1\) the
smallest kernel entry is
\(128/109395>0\), contrasting with the
signed kernel of V69.

A rational three-pole matrix model
checks the Jensen gap by exact
\(LDL^T\) factorization, together with
the averaging variance identity.
That model is not substituted for the
full cotangent or the full lattice
gauge operator. A separate quadratic
interaction checks the normalized
output covariance and mean. The
infinite-volume obstruction, domain
preservation and probability bounds
have their analytic proofs above;
they are not inferred from the finite
kernel tests.

The 6,466-byte certificate has SHA--256
\begin{center}\scriptsize\ttfamily
5F0D8BB1D71FBD267179D6CA118C87C755CF8A9CD3245CCE79C4F45191F90584.
\end{center}
Its V56 and V58 inputs are unchanged.
The predecessor V73 has 1,814,492 bytes
and SHA--256
\begin{center}\scriptsize\ttfamily
6DEB30306292CD49D8A7F160DB02A7028DA9D0FCF90ACA1E2432BD2A789E5609.
\end{center}
Its body is preserved except for the
title version and this terminal insertion.
V74 is the single active main note.
The archives, companion programmes,
attached sources and monographs
remain unchanged.

\begin{firewall}
The input Wilson measure has not been
replaced. The new exact auxiliary
split preserves its normalizer and
has a positive deterministic retained
average, but neither its noise nor
its conditional scalar pin is a new
physical field or a Yang--Mills mass.
The actual marginal has only the
displayed lower-form sign and
input-dependent gauge-floor transfer.
No uniform interacting infrared
closure, continuum construction or
physical Hamiltonian mass gap is proved.
The original objective remains active.
\end{firewall}
\section{V75: an actual determinant-barrier estimate and its dimension cost}

\subsection{The scheduled companion review and the remaining question}

The V75 checkpoint reopens the latest
active companion files in the folder.
Their versions and contents are unchanged
from the preceding checkpoint:
\begin{itemize}
 \item SM--Einstein V72 derives its
 covariance flow in a declared many-copy
 limit, at fixed mode cutoff and fixed
 times. Its concentrated initial-data
 result does not provide a growing-cutoff
 Yang--Mills gauge-floor estimate.
 \item NS Stability V26 computes
 contractions and rank-one norms of
 Oseen-kernel derivatives. Those
 objects are not the present gauge
 operator, and no NS constant is
 imported into a Yang--Mills estimate.
 \item Shell Mixing V16 assembles the
 complete measured flat-root one-loop
 coefficient, with its root return,
 and proves its coefficient limit.
 It explicitly does not prove a joint
 interacting \(\beta,L\) remainder or
 a thermodynamic small-field measure.
\end{itemize}
The companion calculations remain useful
for maintaining the distinctions between
an auxiliary limit and an actual law,
between a tested direction and a full
operator, and between a coefficient and
its nonperturbative remainder.
They supply no replacement for the
input probability in
\eqref{eq:v74-floor-transfer}.
The already-completed mixed source
calculations are not repeated here.

V74 preserved a positive input gauge
floor under deterministic averaging and
quantified its loss under the added
noise. It left the actual probability
of an input close to the gauge boundary
unestimated. This note proves such an
estimate directly under the full
restricted Wilson law. The determinant
is retained, not replaced by a Gaussian
reference or a fixed spectral floor.

The resulting bound is quadratic in
the small gauge eigenvalue, but its
constant includes the dimension and
a controlled adverse radial Wilson
contribution. It therefore does not
close the fixed-radius continuum
problem. A separate consequence,
the expected fraction of small gauge
eigenvalues, has a volume-independent
bound at fixed \(\beta,R\). That
distinction will matter for the next
step.

\subsection{The complete fixed-radius law in connection coordinates}

Use the affine setup and fixed radius
of V74. It is convenient to work with
\[
 h=tZ,\qquad a=a_c+h,\qquad
 \mathcal A_c=\{h\in\mathcal H_M:
            \max_e|a_c(e)+h_e|<R,\ 
            \mathcal F_M(a_c+h)\succ0\}.
\]
This is a change of scale in the same
law. Its constant Lebesgue Jacobian
cancels in normalized probabilities.
Write
\begin{equation}\label{eq:v75-density}
 \begin{gathered}
 F(h)=\mathcal F_M(a_c+h),\qquad f(h)=\lambda_{\min}F(h),\\
 \mathcal H(a)=-2\sum_e\log\left(\frac{\sin|a_e|}{|a_e|}\right),\\
 d\mu_c(h)=\mathcal Z_c^{-1}\mathbf1_{\mathcal A_c}(h)
       \det F(h)\,e^{-\beta S(a_c+h)-\mathcal H(a_c+h)}\,dh .
 \end{gathered}
\end{equation}
Here \(\mathcal Z_c\) is the corresponding
Lebesgue normalizer; it is not a newly
chosen normalization convention.
The physical coordinate and primed
gauge dimensions are
\[
                    n=9(V-1),\qquad d_{\rm g}=3(V-1),\qquad V=M^4.
\]
Let \(\kappa_0>0\) be a known base floor,
\begin{equation}\label{eq:v75-base}
                         F(0)\succeq\kappa_0 I.
\end{equation}
For \(c=0\) take \(\kappa_0=1\).
Under V65's homogeneous assumptions one
may take \(\kappa_0=447/512\).
No such floor is assumed for random \(h\).

The following explicit constants will
be used:
\begin{equation}\label{eq:v75-constants}
 \mathcal J=4\beta C_0R^4V,\qquad
                     K=n+d_{\rm g}+\mathcal J.
\end{equation}
If \(c=0\), every statement below also
holds with the improved value
\(\mathcal J=(4/3)\beta C_0R^4V\).
In particular \(K\ge2\) throughout the
present lattice application.

\begin{lemma}[Radial inequalities with all density factors retained]
\label{lem:v75-radial}
On \(\mathcal A_c\),
\begin{equation}\label{eq:v75-radial}
 \begin{aligned}
 DF(h)[h]&\preceq F(h)-\kappa_0 I,\\
 D\mathcal H(a_c+h)[h]&\ge0,\\
 \beta DS(a_c+h)[h]&\ge-\mathcal J .
 \end{aligned}
\end{equation}
All these bounds concern the actual
functions in \eqref{eq:v75-density}.
\end{lemma}
\begin{proof}
Matrix concavity and the tangent
inequality at \(h\) give
\[
                         F(0)\preceq F(h)-DF(h)[h].
\]
Insert \eqref{eq:v75-base} to obtain
the first assertion.

At the homogeneous \(a_c\), the
gradients of both the scalar Haar
potential and the Wilson action
are homogeneous link fields.
Their pairings with any mean-zero
variation \(h\) therefore vanish.
Convexity of \(\mathcal H\), in fact
its Hessian lower bound \(2I/3\),
then gives
\[
                 D\mathcal H(a_c+h)[h]\ge\frac23\|h\|^2\ge0.
\]
The segment \(a_c+uh\), \(0\le u\le1\),
stays inside the radius-\(R\) link
ball. V64's actual Wilson lower
form gives
\[
 DS(a_c+h)[h]
     \ge\tfrac12\langle h,\Delta h\rangle-C_0R^2\|h\|^2.
\]
The homogeneous part is orthogonal
to \(h\), so
\[
       \|h\|^2=\|a_c+h\|^2-\|a_c\|^2\le4R^2V .
\]
Dropping the nonnegative kinetic
term proves the third bound.

For \(c=0\), the radius along the
segment is \(uR\), rather than \(R\).
Integrating the sharper lower form
there gives
\[
 DS(h)[h]\ge\tfrac12\langle h,\Delta h\rangle
                              -\frac{C_0R^2}{3}\|h\|^2 .
\]
This proves the improved constant.
No assertion that the Wilson action
is radially nondecreasing at fixed
radius has been made.
\end{proof}

\subsection{Actual inverse moments of the gauge operator}

\begin{theorem}[A determinant-barrier moment inequality for the full law]
\label{thm:v75-moments}
Under \eqref{eq:v75-density},
\begin{equation}\label{eq:v75-trace-moment}
                     \kappa_0\,\mathbb E_{\mu_c}\Tr F^{-1}\le K.
\end{equation}
For \(0<p\le1\),
\begin{equation}\label{eq:v75-low-moments}
                 \mathbb E_{\mu_c}(\kappa_0/f)^p\le K^p,
\end{equation}
and for \(1\le p<2\),
\begin{equation}\label{eq:v75-moments}
                 \mathbb E_{\mu_c}(\kappa_0/f)^p
                                      \le\frac{K^p}{2-p}.
\end{equation}
These are normalized interacting
estimates, with no assumed gap from
the gauge boundary.
\end{theorem}
\begin{proof}
For \(0<\epsilon<\kappa_0\), restrict
temporarily to
\[
 \mathcal A_{c,\epsilon}
     =\{h:\max_e|a_c+h_e|<R,\ F(h)\succ\epsilon I\}.
\]
This convex set contains zero.
Along each ray \(h=r\theta\) it is
an interval \(0<r<r_{\epsilon,\theta}\).
All factors are bounded and regular
away from the spectral boundary.
Radial integration of the derivative
of \(r^n\) times the density has
a nonnegative endpoint contribution.
Equivalently the outward radial
boundary flux is nonnegative.
Thus, for any potential \(U\) whose
positive density is being integrated
on this set,
\begin{equation}\label{eq:v75-IBP}
                 \int DU(h)[h]e^{-U(h)}\,dh
                                     \le n\int e^{-U(h)}\,dh .
\end{equation}
This is an inequality, not an
integration-by-parts identity with
the radius boundary omitted.
The contribution at \(r=0\) vanishes.
One may perform this calculation
ray by ray, so no smoothness of the
intersection boundary is required.

First take
\(U=\beta S+\mathcal H-\log\det F\).
The first inequality in
\eqref{eq:v75-radial} implies
\[
 -\Tr(F^{-1}DF[h])
                       \ge\kappa_0\Tr F^{-1}-d_{\rm g}.
\]
The other two radial bounds show
\[
             DU[h]\ge\kappa_0\Tr F^{-1}-d_{\rm g}-\mathcal J.
\]
Insert this into \eqref{eq:v75-IBP}.
Let \(\epsilon\downarrow0\) by monotone
convergence of the nonnegative
integrals on the left. The right
side is bounded by \(K\mathcal Z_c\).
Division by the unchanged normalizer
proves \eqref{eq:v75-trace-moment}.

To prove higher moments, let
\(0\le q<1\) and multiply the actual
density by \(f^{-q}\), still first
on \(\mathcal A_{c,\epsilon}\).
Its potential is \(U_q=U+q\log f\).
The minimum eigenvalue is locally
Lipschitz, differentiable almost
everywhere along rays, and bounded
below by \(\epsilon\) here.
At a differentiability point choose
a minimum-eigenvalue vector realizing
its radial derivative, and complete
it to an orthonormal eigenbasis of
\(F\). In a degenerate minimum space
one may diagonalize the compressed
radial derivative first.
The coefficient of that vector in
the logarithmic determinant derivative
is changed from \(-1\) to \(-(1-q)\);
all remaining coefficients stay at
\(-1\).

For every vector in the chosen basis,
\eqref{eq:v75-radial} bounds the
radial diagonal entry of \(DF\)
by its eigenvalue minus \(\kappa_0\).
Consequently
\[
 -\Tr(F^{-1}DF[h])+q\,\frac{Df[h]}f
        \ge\frac{(1-q)\kappa_0}{f}-(d_{\rm g}-q).
\]
The same Haar and Wilson bounds yield,
from \eqref{eq:v75-IBP},
\begin{equation}\label{eq:v75-recursion}
 (1-q)\kappa_0
       \int_{\mathcal A_{c,\epsilon}}f^{-q-1}\,d\mu_c
       \le (K-q)
       \int_{\mathcal A_{c,\epsilon}}f^{-q}\,d\mu_c .
\end{equation}
The notation \(d\mu_c\) keeps the
original normalization on both sides;
renormalizing the truncated law is
not necessary.

The first-moment result and concavity
of \(x^q\) show
\(\mathbb E(\kappa_0/f)^q\le K^q\).
This supplies the finite right-hand
integral needed when
\(\epsilon\downarrow0\) in
\eqref{eq:v75-recursion}. Monotone
convergence then proves
\[
       \mathbb E(\kappa_0/f)^{q+1}
             \le\frac{K-q}{1-q}K^q
             \le\frac{K^{q+1}}{1-q}.
\]
This is \eqref{eq:v75-moments} with
\(p=q+1\). The same concavity
argument proves \eqref{eq:v75-low-moments}.
The possible nondifferentiability of
an ordered eigenvalue causes no
additional term: the truncated
radial functions are absolutely
continuous and the preceding
inequalities hold almost everywhere.
\end{proof}

The physical dimension \(n\), the
primed determinant dimension
\(d_{\rm g}\), and the radial Wilson
cost \(\mathcal J\) have distinct
origins in this proof. None is
removed by normalizing the density.

\subsection{A direct quadratic probability bound, without a logarithmic loss}

Optimizing Markov's inequality in
\eqref{eq:v75-moments} would introduce
an unnecessary logarithmic loss near
the endpoint \(p=2\). The radial
geometry gives a direct probability
estimate instead.

\begin{theorem}[Small gauge eigenvalues under the actual Wilson law]
\label{thm:v75-tail}
For every \(\delta>0\),
\begin{equation}\label{eq:v75-tail}
 \mu_c\{f<\delta\}
       \le \min\left\{1,\ 
                   2e^2\left(\frac{K\delta}{\kappa_0}\right)^2\right\}.
\end{equation}
The event and normalization are those
of the complete measure in
\eqref{eq:v75-density}.
\end{theorem}
\begin{proof}
Fix a ray \(h=r\theta\) of the
convex domain, with endpoint \(r_*\).
The function \(f(r)=f(r\theta)\) is
positive and concave, with
\(f(0)\ge\kappa_0\).
Write its unnormalized radial density
as \(f(r)b(r)\), where
\[
 b(r)=r^{n-1}\frac{\det F(r\theta)}{f(r)}
             e^{-\beta S(a_c+r\theta)-\mathcal H(a_c+r\theta)}.
\]
The same spectral derivative
calculation used above, now omitting
one minimum-eigenvalue factor, gives
almost everywhere
\begin{equation}\label{eq:v75-ray-growth}
                         r(\log b)'(r)\le K-2.
\end{equation}
Indeed the Jacobian contributes
\(n-1\), the remaining determinant
eigenvalues contribute at most
\(d_{\rm g}-1\), the adverse Wilson
term at most \(\mathcal J\), and
the Haar term is nonpositive.
This estimate does not require a
positive lower bound for the other
eigenvalues independent of the ray.

It suffices first to take
\(\delta\le\kappa_0/(2K)\).
If the ray has no point with
\(f<\delta\), its bad integral is
zero. Otherwise there is a first
crossing \(r_\delta\) with
\(f(r_\delta)=\delta\), after which
the concave function stays below
\(\delta\). Its supporting slope
at that crossing is at most
\(-(\kappa_0-\delta)/r_\delta\).
Thus for \(r\ge r_\delta\),
\[
 f(r)\le\delta-\frac{\kappa_0-\delta}{r_\delta}(r-r_\delta),
 \qquad
        r_*-r_\delta\le\frac{\delta r_\delta}{\kappa_0-\delta}.
\]
These statements also hold if the
ray ends first at the link-radius
boundary, with a strictly positive
limiting gauge eigenvalue.

Let \(x=\delta/\kappa_0\).
Integrating the last linear upper
bound, and applying
\eqref{eq:v75-ray-growth}, gives
\begin{equation}\label{eq:v75-bad-ray}
 \int_{r_\delta}^{r_*}f(r)b(r)\,dr
     \le
       \frac{\delta^2r_\delta}{2(\kappa_0-\delta)}
                (1-x)^{-(K-2)}b(r_\delta).
\end{equation}
For comparison use the earlier ray
segment
\[
       r=t r_\delta,\qquad
         1-\frac1K\le t\le1-\frac1{2K}.
\]
Concavity gives \(f(tr_\delta)\ge
(1-t)\kappa_0+t\delta\ge\kappa_0/(2K)\).
Integrating \eqref{eq:v75-ray-growth}
backwards gives
\[
 b(tr_\delta)\ge t^{K-2}b(r_\delta)
                                   \ge e^{-1}b(r_\delta).
\]
For the last inequality use
\(\log(1-1/K)\ge-1/(K-1)\) and \(K\ge2\).
This earlier interval has length
\(r_\delta/(2K)\), so the full ray
integral is at least
\begin{equation}\label{eq:v75-good-ray}
               \frac{\kappa_0r_\delta}{4eK^2}b(r_\delta).
\end{equation}
The ratio of \eqref{eq:v75-bad-ray}
to \eqref{eq:v75-good-ray} is at most
\[
                  2e(Kx)^2(1-x)^{-(K-1)}
                                  \le2e^2(Kx)^2 .
\]
Here \(x\le1/(2K)\) implies
\[
 -(K-1)\log(1-x)
        \le\frac{(K-1)x}{1-x}
        \le\frac{K-1}{2K-1}<1.
\]
This estimate is uniform in the ray.
Integrating the inequality between
bad and total ray integrals over
the unit sphere proves the bound
for this range of \(\delta\), with
the actual global normalizer.

If \(\delta>\kappa_0/(2K)\), the
displayed quadratic upper bound
exceeds \(e^2/2>1\). The trivial
probability bound supplies that
range. Taking the minimum with one
proves \eqref{eq:v75-tail}.
\end{proof}

\subsection{A volume-independent estimate for spectral density, not a global floor}

\begin{proposition}[Expected fraction of small primed gauge eigenvalues]
\label{prop:v75-density}
Let \(N_\delta(F)\) count eigenvalues
of \(F\) below \(\delta\), with
multiplicity. Then
\begin{equation}\label{eq:v75-spectral-density}
 \mathbb E_{\mu_c}\frac{N_\delta(F)}{d_{\rm g}}
       \le \frac{\delta}{\kappa_0}
                       \left(4+2\beta C_0R^4\right)
                                                     \qquad(M\ge4).
\end{equation}
In the zero homogeneous sector one
may replace \(2\beta C_0R^4\) by
\((2/3)\beta C_0R^4\).
\end{proposition}
\begin{proof}
Every eigenvalue below \(\delta\)
contributes more than \(1/\delta\)
to \(\Tr F^{-1}\), so
\(N_\delta(F)\le\delta\Tr F^{-1}\).
Apply \eqref{eq:v75-trace-moment}.
Since \(n=3d_{\rm g}\),
\[
 \frac K{d_{\rm g}}
       =4+\frac{4\beta C_0R^4V}{3(V-1)}
       \le4+2\beta C_0R^4.
\]
Use the improved radial cost when
\(c=0\) to get the last assertion.
\end{proof}

At fixed \(\beta,R\), this bound is
uniform in volume. It does not say
that the probability of even one
small eigenvalue is small. Nor does
it identify a spatial set of bad
links or blocks: a low-dimensional
spectral subspace can be spatially
delocalized. Its use in a local
renormalization argument requires
additional information on the
associated spectral projectors.

\subsection{Why the determinant alone does not supply a uniform minimum}

A one-ray comparison makes the
dimension dependence concrete.
For an integer \(n\ge1\), let
\[
                 d\nu_n(r)=n(n+1)r^{n-1}(1-r)\,dr,\qquad 0<r<1.
\]
The factor \(1-r\) is a single
linearly vanishing determinant
factor. Its exact properties are
\begin{equation}\label{eq:v75-beta-model}
 \begin{aligned}
 \nu_n\{1-r<\delta\}
       &=1-(1-\delta)^n(1+n\delta),\qquad 0\le\delta\le1,\\
 \mathbb E_{\nu_n}(1-r)^{-1}&=n+1,\\
 \mathbb E_{\nu_n}(1-r)^{-p}
       &=\frac{(n+1)!}{\prod_{j=0}^{n-1}(2-p+j)}
                                                \quad(0\le p<2).
 \end{aligned}
\end{equation}
The formulas follow by integrating
the elementary beta density; the
last integral can also be evaluated
by \(n-1\) integrations by parts.
At \(p=2\) the integral diverges
logarithmically at \(r=1\).
For a fixed \(0<\delta<1\), its
bad probability tends to one as
\(n\to\infty\), despite the vanishing
determinant factor. For small
\(\delta\), it is
\[
               \frac{n(n+1)}2\delta^2+O_n(\delta^3).
\]
Thus neither a dimension-free
minimum-eigenvalue bound nor a
second inverse moment follows
from a simple determinant barrier
by itself.

This is an explicitly separate
radial comparison law, not a claim
that the Yang--Mills measure has
this distribution. It explains
why the dimension costs and the
endpoint \(p<2\) in the proved
estimates should not be silently
removed.

\subsection{Testing the new estimate against V74's actual requirement}

Equation \eqref{eq:v75-tail} now
provides a proved input for the
right-hand side of
\eqref{eq:v74-floor-transfer}.
Substitution gives, under V74's
same output-event assumptions,
\begin{equation}\label{eq:v75-combined}
 \widehat\pi_{c,s}(\mathcal G^c)
      \le \min\left\{1,\ 
                 2e^2(K\delta/\kappa_0)^2\right\}+2e^{-u}.
\end{equation}
This is an unconditional estimate
for this restricted Wilson input,
not a closure with an unproved
input-floor hypothesis. Its value
must nevertheless be tested.

For the proved noise allowance in
V74, \(\delta\) must exceed
\[
                    e_{\rm g}\ge C\sqrt{(1+\log M+u)/\beta}.
\]
At fixed \(R>0\), the present
worst-case constant has
\(K\ge4\beta C_0R^4V\).
Consequently every allowed
\(\delta\) obeys
\[
 \frac{K\delta}{\kappa_0}
       \ge \frac{4CC_0R^4}{\kappa_0}\,
                    V\sqrt{\beta(1+\log M+u)}.
\]
The zero homogeneous improvement
only changes the factor four to
\(4/3\). This lower bound diverges
as \(M\to\infty\), for any sequence
\(\beta\ge1\) at fixed \(R\).
Therefore the right side of
\eqref{eq:v75-combined} becomes
trivial in this use. This is a
failure of the available estimates
to close, not a proof that the
actual output gauge floor fails.

The new nonperturbative information
is still substantive: the full law
has controlled inverse gauge moments,
a quadratic bound on the boundary
probability with explicit costs, and
a bound uniform in volume for the
expected fraction of small gauge
eigenvalues at fixed parameters.
The next step must improve the
adverse radial Wilson estimate or
replace a global minimum by controlled
spatial spectral information. It
must not merely optimize the moment
exponent or hide \(K\) in a constant.
Neither the companion coefficient
limits nor the normalized
determinant factor currently remove
that requirement.

\subsection{Exact checks and checkpoint preservation}

The new certificate is
\begin{center}\small
\path{scripts/verify_ym_f14_v75_gauge_barrier.py}.
\end{center}
Forty rational finite-pole matrix
tests check the radial concavity
remainder and the inverse-trace
inequality, with a noncommuting
affine symmetric term retained.
Sixteen diagonal spectral tests
check the minimum-eigenvalue
subtraction for \(0\le q<1\).
These finite-pole models do not
replace the full cotangent in the
analytic proof.

For radial dimensions \(1,2,3,9,27\),
the script verifies normalization,
the exact inverse first moment,
the weighted radial recursion,
and the \(p=3/2\) bound using
rational arithmetic after squaring.
At four rational thresholds per
dimension, the closed tail formula
in \eqref{eq:v75-beta-model} agrees
with an independently expanded
polynomial integral. These are
comparison-law identities, not
sampling of the interacting
Yang--Mills measure.

The 4,707-byte certificate has SHA--256
\begin{center}\scriptsize\ttfamily
811306FC2A6E60FAE95F3F2CC432346AAFF842221A36890253AEDC42AAC3D66D.
\end{center}
The input helper and its dependencies
are unchanged. The predecessor V74
has 1,841,333 bytes and SHA--256
\begin{center}\scriptsize\ttfamily
FB1EE330D8B8416B7E88B72C5BB33265924E4A791A32750D3DB8C03132C5EB95.
\end{center}
Its body is preserved except for the
title version and this terminal
insertion. V75 is the single active
main note.

The companion checkpoint confirms
SM V72, 607,372 bytes, with SHA--256
\begin{center}\scriptsize\ttfamily
F44F9745D43379E1672C5550411B58E97195A4C9278592D221DF5D3F644DC1F5;
\end{center}
NS V26, 278,283 bytes, with SHA--256
\begin{center}\scriptsize\ttfamily
82C887F8C2C69E4F28FAD7C3DC8DE5B9099CB5A4BF431EBA1FFF3E91935978AD;
\end{center}
and Shell Mixing V16, 623,261 bytes,
with SHA--256
\begin{center}\scriptsize\ttfamily
BF138E63BA826BC030D3A8B87E04FDD3B1EC712D269BB064674E96509B46E710.
\end{center}
They are read-only inputs and have
not been altered. The other attached
sources, archives and monographs
remain unchanged. The next scheduled
companion checkpoint is V80.

\begin{firewall}
The probability estimates in this
note are for the complete normalized
Wilson density on the stated
positive fixed-radius coordinate
domain. They do not cover the other
compact sectors. Their explicit
dimension and radial-action costs
do not prove the uniform minimum
needed for V74's noisy return.
The spectral-density bound is not
a spatial localization theorem.
No continuum construction or
physical Yang--Mills mass gap is
proved. The original objective
remains active.
\end{firewall}
\section{V76: the infrared gauge subspace and an exact determinant reduction}

\subsection{Context: where can a small gauge eigenvalue occur?}

V75 bounds the inverse gauge trace and
the expected number of small eigenvalues,
but warns that a spectral count is not
a collection of localized bad links.
This note addresses that missing
geometric information before attempting
another probability estimate.

The actual normalized gauge operator
has a deterministic smoothing identity
in its odd part, because the connection
is Coulomb. We use it to show that
every vector in a small-eigenvalue
spectral subspace has Laplacian energy
of order \(R^2\). This bounds the
dimension of that subspace and limits
its concentration on a small set of
sites. The conclusion is an infrared
restriction, not exponential localization.

A fixed split of the primed gauge
space then separates a uniformly
positive high-momentum block from
an exact low-momentum Schur complement.
The full determinant and positivity
domain factor through this complement.
Its matrix concavity survives the
elimination, so V75's barrier argument
applies to the reduced determinant.
The physical-coordinate dimension
still enters that argument and is
not replaced by the smaller gauge
dimension. No physical mass gap is
inferred from this reduction.

\subsection{The Coulomb identity supplies a weighted operator bound}

Let \(\mathcal V_M\) be the real primed
vertex colour space, of dimension
\(d_{\rm g}=3(V-1)\).
The Laplacian in this section acts
on that space. Fix any real Coulomb
connection \(a\) with
\(\max_e|a_e|\le R\), where \(0<R\le1/4\), not
necessarily with a positive gauge
operator. Set
\[
                  F=\mathcal F_M(a)=I+G_1(a)+G_{\rm e}(a).
\]
All spectral projections below are
on the primed vertex space. They
are different from the physical
link projections in V68--V74.

\begin{lemma}[A deterministic weighted gauge bound]
\label{lem:v76-weighted}
With the exact operators of
\eqref{eq:v60-gauge},
\begin{equation}\label{eq:v76-weighted}
 \begin{aligned}
 \|G_1(a)\Delta^{1/2}\|
        =\|\Delta^{1/2}G_1(a)\|&\le4R,\\
 \|\Delta^{1/2}G_{\rm e}(a)\|&\le2R^2,\\
 \|\Delta^{1/2}(F-I)\|&\le b_R:=4R+2R^2.
 \end{aligned}
\end{equation}
\end{lemma}
\begin{proof}
The endpoint-sum operator obeys
\[
 S_{\rm end}^*S_{\rm end}=16I-\Delta
\]
on scalar vertex coordinates, and
the same identity holds in each
colour. On a Fourier mode it is
the equality
\(\sum_\mu|1+e^{ik_\mu}|^2=16-\lambda(k)\).
Thus \(\|S_{\rm end}\|\le4\).
Since \(U=D\Delta^{-1/2}\) is an
isometry and
\(\|\operatorname{ad}_a\|\le2R\),
\[
 G_1(a)\Delta^{1/2}
                 =\tfrac12 U^*\operatorname{ad}_a S_{\rm end}
\]
has norm at most \(4R\).
The omitted constant-mode projection
in this display acts as the identity
on the primed input.
The Coulomb condition makes \(G_1(a)\)
self-adjoint. Taking the adjoint
therefore gives the left-weighted
bound as well. This step would not
be justified for a non-Coulomb input.

By \eqref{eq:v65-even},
\(\|\sigma(\operatorname{ad}_a)-I\|\le R^2/2\).
Also \(\Delta^{1/2}U^*=D^*\) and
\(\|D^*\|\le4\).
Consequently
\[
 \Delta^{1/2}G_{\rm e}(a)
                   =D^*\{\sigma(\operatorname{ad}_a)-I\}U
\]
has norm at most \(2R^2\).
The last inequality follows by
addition. No derivative of the
connection and no probabilistic
gauge-floor assumption were used.
\end{proof}

\subsection{Small gauge eigenspaces have bounded kinetic energy}

For \(0\le\delta\le1/4\), let
\[
                         P_\delta(a)=\mathbf1_{(-\infty,\delta]}(F).
\]
On the actual positive domain this
is just the small positive spectral
subspace; the definition also covers
negative eigenvalues outside it.

\begin{theorem}[A deterministic infrared restriction]
\label{thm:v76-infrared}
For every input specified above,
\begin{equation}\label{eq:v76-kinetic}
 \|\Delta^{1/2}P_\delta(a)\|
          \le\frac{b_R}{1-\delta}\le6R.
\end{equation}
In particular every \(u\) in this
spectral subspace satisfies
\[
                       \langle u,\Delta u\rangle\le36R^2\|u\|^2,
\]
and
\begin{equation}\label{eq:v76-rank}
 \operatorname{rank}P_\delta(a)
       \le3\,\#\{k\ne0:\lambda(k)\le36R^2\}
       \le3(1+3MR)^4 .
\end{equation}
For any free cutoff \(\Lambda>0\),
\begin{equation}\label{eq:v76-high-tail}
 \|\mathbf1_{\{\Delta\ge\Lambda\}}P_\delta(a)\|
                                                 \le\frac{6R}{\sqrt\Lambda}.
\end{equation}
\end{theorem}
\begin{proof}
On the range of \(P_\delta\), the
operator \(I-F\) is invertible and
its inverse has norm at most
\((1-\delta)^{-1}\).
The exact identity is
\[
 P_\delta=-(F-I)P_\delta
                  \{(I-F)|_{\operatorname{Ran}P_\delta}\}^{-1}.
\]
Here the restricted inverse is
extended by zero on
\(\operatorname{Ran}P_\delta^\perp\),
so both sides act on the full
primed vertex space.
Apply \(\Delta^{1/2}\) on the left
and use Lemma~\ref{lem:v76-weighted}.
Since \(R\le1/4\) and \(\delta\le1/4\),
\((4R+2R^2)/(1-\delta)\le6R\).
This proves the operator estimate
and the Rayleigh-quotient bound for
every linear combination in the
subspace, not merely individual
eigenvectors.

The min--max principle for the
positive primed Laplacian now bounds
the dimension by its eigenvalue
count below \(36R^2\).
With shortest integer Fourier
representatives,
\[
                       \lambda(2\pi n/M)\ge16|n|^2/M^2.
\]
Thus each coordinate of a counted
integer vector has absolute value
at most \(3MR/2\); the elementary
box count is at most \((1+3MR)^4\).
The factor three is the colour
multiplicity. The zero mode is
still excluded. Finally the
spectral inequality
\(\mathbf1_{\{\Delta\ge\Lambda\}}
\preceq\Delta/\Lambda\) proves
\eqref{eq:v76-high-tail}.
\end{proof}

The bound does not say that
\(P_\delta\) is a free Fourier
projection: \(F\) and \(\Delta\)
need not commute. It supplies an
energy and a high-frequency tail
for the actual spectral projection.

\begin{proposition}[A spatial nonconcentration consequence]
\label{prop:v76-nonconcentration}
For a set \(B\) of vertex sites
and every \(\Lambda>0\),
\begin{equation}\label{eq:v76-local-norm}
 \|\mathbf1_BP_\delta(a)\|
       \le \frac9{16}\Lambda |B|^{1/2}
                              +\frac{6R}{\sqrt\Lambda}.
\end{equation}
In particular, if a unit vector
\(u\in\operatorname{Ran}P_\delta(a)\)
satisfies \(\|\mathbf1_Bu\|\ge3/4\),
then
\begin{equation}\label{eq:v76-support}
                            |B|\ge\frac1{324^2R^4}.
\end{equation}
\end{proposition}
\begin{proof}
Let \(N_\Lambda\) count scalar
nonzero Fourier modes with
\(\lambda<\Lambda\).
If this count is nonzero, then
\(M\sqrt\Lambda\ge4\), and the
same integer-box count gives
\[
 N_\Lambda\le(1+M\sqrt\Lambda/2)^4
                           \le\frac{81}{256}V\Lambda^2.
\]
If it is zero this inequality is
automatic.
For the free low projection
\(P_{<\Lambda}^{\,0}\), Fourier
Cauchy--Schwarz, applied to the
three-component colour vector,
gives
\[
 \|\mathbf1_B P_{<\Lambda}^{\,0}\|
                  \le\sqrt{|B|N_\Lambda/V}
                  \le\frac9{16}\Lambda |B|^{1/2}.
\]
There is no extra colour factor
in this operator norm: the same
scalar projection acts separately
in each colour.
Decompose \(P_\delta\) by the free
low and high projections and use
\eqref{eq:v76-high-tail}.
This proves \eqref{eq:v76-local-norm}.
Take \(\Lambda=144R^2\).
The right side is
\(81R^2|B|^{1/2}+1/2\).
The stipulated lower bound \(3/4\)
then gives \eqref{eq:v76-support}.
\end{proof}

The numerical constant is not
claimed sharp. Its role is to
make a genuine spatial statement:
these vectors cannot have most of
their norm on sets arbitrarily
small compared with \(R^{-4}\).
It is a nonconcentration bound,
not decay away from a localization
centre or a construction of
independent bad blocks.

\subsection{A uniformly positive high gauge block}

Fix the background-independent
gauge projections
\begin{equation}\label{eq:v76-gauge-split}
       P=\mathbf1_{\{0<\Delta<64R^2\}},\qquad Q=I-P,
       \qquad r_{\rm g}=\operatorname{rank}P.
\end{equation}
They act only on the primed gauge
space and commute with translations.
Their rank obeys
\[
                   r_{\rm g}
                       =3\,\#\{0<\lambda<64R^2\}
                       \le3(1+4MR)^4.
\]
Write the actual matrix in this
orthogonal decomposition as
\[
                     F(a)=
                       \begin{pmatrix}A(a)&B(a)^*\\B(a)&D(a)\end{pmatrix}.
\]
In this section \(D(a)=QF(a)Q\)
is a gauge block, not a
total-field coordinate domain.
Define
\[
 \eta_R=\tfrac12+\tfrac12R^2,\qquad
 c_R=1-\eta_R=\tfrac12-\tfrac12R^2,\qquad
 \rho_R=\eta_R/c_R.
\]
The constants satisfy
\(c_R\ge15/32\) and \(\rho_R\le17/15\).

\begin{theorem}[Exact positivity and determinant reduction]
\label{thm:v76-Schur}
Throughout the full affine Coulomb
link ball \(\max_e|a_e|<R\), including
points where \(F\) is not positive,
\begin{equation}\label{eq:v76-high-block}
                 D(a)\succeq c_R I,\qquad
                 \|D(a)-I\|\le\eta_R,\qquad
                 \|B(a)\|\le\eta_R .
\end{equation}
The Schur complement
\begin{equation}\label{eq:v76-Schur}
                    \mathsf S(a)=A(a)-B(a)^*D(a)^{-1}B(a)
\end{equation}
is therefore well-defined and satisfies
\begin{equation}\label{eq:v76-determinant}
                    F(a)\succ0\ \Longleftrightarrow\ \mathsf S(a)\succ0,
             \qquad \det F(a)=\det D(a)\det\mathsf S(a).
\end{equation}
Moreover \(a\mapsto\mathsf S(a)\) is
matrix-concave on that link ball.
If \(F(a_c)\succeq\kappa_0I\),
then \(\mathsf S(a_c)\succeq\kappa_0I\)
on the low gauge space.
\end{theorem}
\begin{proof}
The right-weighted odd bound gives
\[
                  \|G_1(a)Q\|\le\frac{4R}{\sqrt{64R^2}}=\frac12.
\]
Self-adjointness gives the same
bound for \(QG_1(a)\).
Since \(\|G_{\rm e}(a)\|\le R^2/2\),
\[
                         \|Q(F(a)-I)\|\le\eta_R.
\]
Compressing, or restricting to
the other input block, proves
\eqref{eq:v76-high-block}.
The usual triangular multiplication,
which is an exact finite matrix
identity, gives
\[
 F=
 \begin{pmatrix}I&B^*D^{-1}\\0&I\end{pmatrix}
 \begin{pmatrix}\mathsf S&0\\0&D\end{pmatrix}
 \begin{pmatrix}I&0\\D^{-1}B&I\end{pmatrix}.
\]
This proves the positivity equivalence
and determinant formula.

For every low vector \(x\),
\[
             \langle x,\mathsf S(a)x\rangle
                 =\inf_y\langle(x,y),F(a)(x,y)\rangle,
 \qquad y_{\rm min}=-D(a)^{-1}B(a)x.
\]
The infimum is finite because of
the uniform high-block floor.
For each fixed \(x,y\), matrix
concavity of \(F\) makes the
expression concave in \(a\).
Taking the infimum preserves the
concavity inequality: the infimum
of a convex combination is at
least the convex combination of
the two infima. This proves
matrix concavity of \(\mathsf S\).
At the homogeneous base, the
same variational formula and
\(F(a_c)\succeq\kappa_0I\) give
\[
  \langle x,\mathsf S(a_c)x\rangle
       \ge\inf_y\kappa_0(\|x\|^2+\|y\|^2)
       =\kappa_0\|x\|^2.
\]
\end{proof}

An explicit derivative formula
records the sign of the elimination
term. Put
\[
          J_a=\begin{pmatrix}I\\-D(a)^{-1}B(a)\end{pmatrix}.
\]
For a fixed connection variation,
denote derivatives along its line
by dots. Differentiation of the
minimizing quadratic form gives
\begin{equation}\label{eq:v76-Schur-Hessian}
 \begin{aligned}
 \dot{\mathsf S}&=J_a^*\dot FJ_a,\\
 \ddot{\mathsf S}
       &=J_a^*\ddot FJ_a
          -2(Q\dot FJ_a)^*D(a)^{-1}(Q\dot FJ_a)\preceq0 .
 \end{aligned}
\end{equation}
Indeed the derivative of the
minimizing high component is
\(-D^{-1}Q\dot FJ_a\), while
\(QFJ_a=0\). Substitution gives
the two formulas and their sign.
The second term is retained,
not dropped as a harmless
normalizer correction.

If \(r_{\rm g}=0\), the residual
matrix is empty, its determinant
is one, and the entire gauge
matrix already has the uniform
floor \(c_R\) throughout the
link ball. Statements about its
residual minimum eigenvalue
below are only used when
\(r_{\rm g}>0\).

\subsection{The full density keeps the exact high determinant}

On the affine link ball define
\[
           W_{\rm h}(a)=\mathcal H(a)-\log\det D(a),
\]
where \(\mathcal H\) is the actual
Haar potential from V75.
The high block is uniformly
positive there, so this is a
smooth function even beyond
the positive domain of the
full gauge matrix.

\begin{proposition}[An exact reduced barrier for the actual law]
\label{prop:v76-reduced-law}
The function \(W_{\rm h}\) is convex,
has Hessian at least \(2I/3\) on
the mean-zero affine variations,
and is translation invariant.
At the homogeneous \(a_c\),
its first derivative pairs to
zero with every mean-zero
variation. The actual law of
V75 is exactly
\begin{equation}\label{eq:v76-reduced-law}
 d\mu_c(h)=\mathcal Z_c^{-1}
     \mathbf1_{\{\max_e|a_c+h_e|<R,\ \mathsf S(a_c+h)\succ0\}}
     \det\mathsf S(a_c+h)\,
             e^{-\beta S(a_c+h)-W_{\rm h}(a_c+h)}\,dh .
\end{equation}
There is no change to its normalizer.
For \(r_{\rm g}>0\), let
\[
 f_{\rm red}(h)=\lambda_{\min}\mathsf S(a_c+h),\qquad
                   K_{\rm red}=n+r_{\rm g}+\mathcal J,
\]
with the same radial Wilson cost
\(\mathcal J\) as in V75.
Then
\begin{equation}\label{eq:v76-reduced-moments}
 \begin{aligned}
 \kappa_0\mathbb E_{\mu_c}\Tr\mathsf S^{-1}&\le K_{\rm red},\\
 \mathbb E_{\mu_c}(\kappa_0/f_{\rm red})^p
             &\le\frac{K_{\rm red}^p}{2-p}
                                             \quad(1\le p<2),\\
 \mu_c\{f_{\rm red}<\delta\}
             &\le\min\{1,\,2e^2(K_{\rm red}\delta/\kappa_0)^2\}.
 \end{aligned}
\end{equation}
\end{proposition}
\begin{proof}
The compression \(D(a)\) is
matrix-concave and uniformly
positive on this whole ball.
The logarithmic determinant
Hessian identity used in V62
shows that \(-\log\det D(a)\)
is convex. Adding the Haar
potential gives the asserted
Hessian floor.
The fixed projection \(Q\)
commutes with translations;
its determinant and the Haar
potential are translation
invariant. Their gradient at
a homogeneous input is
therefore homogeneous, so
its pairing with a mean-zero
variation is zero.

Substitute \eqref{eq:v76-determinant}
into \eqref{eq:v75-density}.
Both the positivity event and
the density coincide pointwise,
which proves
\eqref{eq:v76-reduced-law}.
No integration over a new
physical variable and no
field-dependent change of
gauge basis has been made.

For the moment and probability
claims, all hypotheses used
in the proofs of V75 now
hold with \(F\) replaced by
\(\mathsf S\): matrix concavity,
the base floor \(\kappa_0\),
a convex positive domain,
and a remaining geometric
potential whose radial
derivative is nonnegative.
The adverse radial Wilson
cost is unchanged.
The determinant size is now
\(r_{\rm g}\), but radial
integration is still in the
physical \(n\)-dimensional
space. Applying those proved
arguments gives precisely
\(K_{\rm red}=n+r_{\rm g}+\mathcal J\)
and \eqref{eq:v76-reduced-moments}.
\end{proof}

Thus only the determinant-size
term has been reduced. In
particular \(n=9(V-1)\) remains
extensive, and the factor
\(\mathcal J\) has not improved.
This representation cannot be
reported as an \(r_{\rm g}\)-
dimensional physical probability
law or as a removal of the
dimension obstruction found
in V75.

For completeness, on the
positive domain the exact
inverse is
\begin{equation}\label{eq:v76-inverse}
                F^{-1}=J_a\mathsf S^{-1}J_a^*
                              +\begin{pmatrix}0&0\\0&D^{-1}\end{pmatrix}.
\end{equation}
Since \(\|D^{-1}B\|\le\rho_R\),
the triangular factorization gives
\begin{equation}\label{eq:v76-minimum-comparison}
 \frac{\min\{f_{\rm red},c_R\}}{(1+\rho_R)^2}
               \le\lambda_{\min}F\le f_{\rm red},
 \qquad (1+\rho_R)^2\le(32/15)^2.
\end{equation}
The lower bound follows by
bounding the inverse triangular
factor by \(1+\rho_R\). For
the upper bound, evaluate the
full Rayleigh quotient on
the minimizing lift \(J_ax\)
of a unit lowest vector of
\(\mathsf S\); its norm is
at least one.
These comparisons retain the
coupling between the gauge
blocks. They do not turn the
reduced determinant into a
uniform bound on the full
minimum at growing volume.

\subsection{Local spectral weight under the interacting law}

Return to the positive law
\(\mu_c\) and homogeneous \(a_c\).
For any set \(B\) of vertex
sites, let
\[
          \mathcal W_{B,\delta}(a)
               =\Tr\bigl(\mathbf1_BP_\delta(a)\mathbf1_B\bigr).
\]
The projection is extended
by zero on the constant
vertex colour modes when
writing its diagonal in
site coordinates.

\begin{proposition}[An actual local spectral-weight bound]
\label{prop:v76-local-weight}
For \(0<\delta\le1/4\),
\begin{equation}\label{eq:v76-local-weight}
 \mathbb E_{\mu_c}\mathcal W_{B,\delta}
    \le |B|\min\left\{
        3,\ \frac{3(1+3MR)^4}{V},\
                       \frac{\delta K}{\kappa_0 V}\right\},
\end{equation}
where \(K\) is the proved V75
constant, with its improved
zero-homogeneous-sector choice
when applicable.
\end{proposition}
\begin{proof}
The law, the gauge matrix
and its spectral projection
are translation covariant.
The expected colour trace
of the diagonal of
\(P_\delta\) is consequently
the same at every site.
Summing over all \(V\) sites
gives the exact identity
\[
 \mathbb E\mathcal W_{B,\delta}
                     =\frac{|B|}{V}\,
                           \mathbb E\operatorname{rank}P_\delta.
\]
The rank is at most \(3V\),
is bounded deterministically
by \eqref{eq:v76-rank}, and,
on the positive domain,
is at most \(\delta\Tr F^{-1}\).
The latter statement also
holds for eigenvalues equal
to \(\delta\), which are
included in \(P_\delta\).
Apply \eqref{eq:v75-trace-moment}
and take the minimum.
\end{proof}

This is an estimate for local
spectral weight of the actual
interacting input, not a
reference-measure substitution.
It still gives no product
bound for separated sets \(B\),
no independence of bad blocks,
and no decay estimate for the
Schur-complement entries.
The free gauge projection
used in the reduction is a
sharp Fourier projection;
its spatial locality has
not been assumed.

\subsection{What changes in the programme}

The gauge-boundary obstruction
is now confined, in a precise
operator sense, to modes with
kinetic scale \(O(R^2)\).
Their number is deterministically
bounded, their spatial
concentration is restricted,
and their expected local
spectral weight is controlled.
The high gauge determinant
has a uniform floor throughout
the entire fixed-radius link
ball and has been separated
without deleting its density
factor.

The remaining barrier is the
actual low-gauge Schur
complement coupled to the
full physical field. Its
dimension is of order \(R^4V\)
when \(MR\) is large, not a
fixed finite number.
The next estimate needs
spatial or scale-dependent
control of this residual
operator under the true law,
or an improved radial action
estimate. The reduction does
not supply that control by
itself. In particular the
physical-coordinate entropy
cost \(n\) in V75 cannot
be removed just by counting
fewer gauge eigenvalues.

This gauge-space factorization
is not the root coarea or
source-return determinant
in the shell companion, and
it is not the physical
link-field split of V68.
No extra copy of either
determinant is introduced.
Compact-sector coverage,
continuum construction and
physical transfer calibration
remain separate requirements.
The next companion checkpoint
remains V80.

\subsection{Exact checks and preservation}

The new certificate is
\begin{center}\small
\path{scripts/verify_ym_f14_v76_gauge_soft_space.py}.
\end{center}
It constructs coexact real
link fields on actual tori
of sides two and four and
builds the full unnormalized
odd gauge operator
\(D^*\operatorname{ad}_aS_{\rm end}/2\).
All 204 and 11,468 nonzero
matrix entries, respectively,
satisfy the exact adjoint
identity. The fields are
checked to be Coulomb and
mean zero; the operator
annihilates all three
constant colour vectors.
The endpoint-sum identity
is checked on all 272
Fourier modes of the two
tori, including the zero
mode for that scalar identity.

The free low-gauge ranks and
the displayed block constants
are checked at four rational
radii. These are not samples
of the nonlinear gauge
spectrum. A declared
rational concave block-matrix
family, evaluated at three rational
parameter values, independently checks
the determinant factorization,
the complete inverse formula,
and the Schur Hessian:
a direct differentiated
inverse agrees with
\eqref{eq:v76-Schur-Hessian}.
The general kinetic,
nonconcentration and
probability statements have
the analytic proofs above.

The 7,151-byte certificate has SHA--256
\begin{center}\scriptsize\ttfamily
7E3421A390E0885774B4AFF4E455938C3373E73E9DBE76F54FFDE33FFD94532D.
\end{center}
Its V58 and V62 inputs and
their helpers are unchanged.
The predecessor V75 contains
1,862,418 bytes and has SHA--256
\begin{center}\scriptsize\ttfamily
6DD94F3B1731975D651AFC8EE6750DC4465E9E71EA54BE1BE26B00F129EE5B19.
\end{center}
Its body is preserved except
for the title version and
this terminal insertion.
V76 is the single active
main note; companion notes,
archives, attached sources
and monographs are unchanged.

\begin{firewall}
The new results bound the kinetic
energy of small gauge modes.
They give an exact determinant
reduction with a positive high
block, and a local spectral-weight
estimate under the interacting law.
None is a physical Hamiltonian
gap. The residual Schur
complement remains correlated
with all physical coordinates,
and the extensive radial cost
has not disappeared.
No uniform interacting infrared
closure, continuum Yang--Mills
construction or physical mass
gap is proved. The original
objective remains active.
\end{firewall}
\section{V77: a spatially local gauge elimination with its full determinant}

\subsection{Context and the change of split}

V76 confines small gauge modes to low
kinetic energies, but its sharp Fourier
projection does not supply spatial
locality of the residual operator.
Here we change the gauge-space split,
not the interacting law. We work with
the unnormalized, nearest-neighbour
gauge matrix and eliminate fields
whose mean vanishes on each spatial
block. A block Poincare inequality
gives a high-block floor throughout
the small-link ball. A weighted
inverse estimate then proves
exponential locality in block distance,
uniformly in the total volume.

The resulting coarse gauge matrix
is an exact Schur complement. Its
entries are exponentially local,
and the eliminated logarithmic
determinant has exponentially small
mixed derivatives for separated
insertions. Neither assertion
allows that determinant to be
discarded. An exact nonzero
separated derivative below makes
this distinction concrete.

One new cost also has to be paid.
A fixed block partition is not
invariant under one-site translations.
Consequently that symmetry argument
does not establish a zero derivative
of the high determinant at a
nonzero homogeneous background
in every mean-zero direction.
We retain an explicit bound for
that term. At zero background
the derivative is zero by the
colour trace identity, not by
an unavailable translation symmetry.

This is not the classical physical
link-field reduction in V40--V44
or the shell companion. The variables
being eliminated here are gauge
parameters inside the Faddeev--Popov
determinant. The physical field is
not integrated out by this operation.
The full goal remains the
nontrivial continuum theory with
a positive physical Hamiltonian gap,
as distinguished in the
\href{https://www.claymath.org/millennium/yang-mills-the-maths-gap/}
{official Yang--Mills problem statement}.

\subsection{A local Garding estimate for the actual gauge matrix}

Use a four-dimensional torus of
side \(M=N\ell\), where
\[
     N\ge1,\qquad \ell\ge2,\qquad
       0<R\le1/4,\qquad R\ell\le1/3.
\]
The block side \(\ell\) divides
\(M\). In particular this choice
requires \(R\le1/6\) and an
admissible divisor of \(M\);
it is not asserted for every
radius and torus allowed in V76.
For fixed \(R,\ell\), it applies
to all volumes with \(M\) a
multiple of \(\ell\).
All vertex and link norms
use counting measure. Let \(a\)
be any real Coulomb connection
with \(\max_e|a_e|\le R\);
positivity of the full gauge
matrix is not assumed.
On the full real vertex colour
space set
\begin{equation}\label{eq:v77-local-M}
 \mathcal M(a)=D^*\sigma(\operatorname{ad}_a)D
                  +\tfrac12D^*\operatorname{ad}_aS_{\rm end}.
\end{equation}
On the primed space,
\(\mathcal F_M(a)=\Delta^{-1/2}
\mathcal M(a)\Delta^{-1/2}\).
These are exactly the V60 operators.
The three constant vertex colour
vectors are in the kernel of
\(\mathcal M(a)\), and the Coulomb
condition makes the matrix self-adjoint.

\begin{lemma}[Local coefficients and a lower form]
\label{lem:v77-Garding}
Throughout this ball,
\begin{equation}\label{eq:v77-Garding}
          \mathcal M(a)\succeq\tfrac12\Delta-9R^2 I,
                   \qquad \|\mathcal M(a)\|\le20 .
\end{equation}
Each off-diagonal nearest-neighbour
link contribution has norm at
most \(1+R\). The sum of the
off-diagonal contribution norms
in any row or column is at most
\(8(1+R)\le10\).
\end{lemma}
\begin{proof}
The even-block bound from
\eqref{eq:v65-even} gives
\[
     (1-R^2/2)I\preceq
        \sigma(\operatorname{ad}_{a_e})\preceq I.
\]
The odd quadratic form has
absolute value at most
\(4R\|Du\|\|u\|\), by
\(\|\operatorname{ad}_a\|\le2R\)
and \(\|S_{\rm end}\|\le4\).
For nonnegative \(x,y\),
\[
        4Rxy\le\frac{15}{32}x^2
                          +\frac{128}{15}R^2y^2.
\]
Since \(1-R^2/2\ge31/32\),
this proves
\[
 \langle u,\mathcal M(a)u\rangle
       \ge\tfrac12\|Du\|^2
                  -\tfrac{128}{15}R^2\|u\|^2,
\]
which is stronger than the
stated lower bound.
Also
\(\|D^*\sigma D\|\le16\)
and
\(\|D^*\operatorname{ad}_aS_{\rm end}/2\|
\le16R\le4\), proving the norm bound.

For a directed link \(e=(x,y)\),
write \(\sigma_e=\sigma(\operatorname{ad}_{a_e})\).
Its two off-diagonal contributions are
\[
       \mathcal M_{xy}^{(e)}
                =-\sigma_e-\tfrac12\operatorname{ad}_{a_e},
       \qquad
       \mathcal M_{yx}^{(e)}
                =-\sigma_e+\tfrac12\operatorname{ad}_{a_e}.
\]
They are adjoints and have norm
at most \(1+R\).
The odd diagonal sum is
\(\operatorname{ad}_{D^*a}/2=0\).
There are eight incident link
contributions at every vertex.
These assertions count periodic
coincident links with their
multiplicity, so side-two tori
cause no exception.
Finally \(D\) annihilates constants,
and the odd term on a constant
colour vector is
\(\operatorname{ad}_{D^*a}\) applied
to that vector, hence zero.
\end{proof}

\subsection{The block constraint and its uniform high floor}

Partition the torus into the
\(N^4\) disjoint coordinate cubes
\[
       B_b=\ell b+\{0,\ldots,\ell-1\}^4,
                    \qquad b\in(\mathbb Z/N\mathbb Z)^4.
\]
Let \(P_B\) be orthogonal averaging
on block \(B\), in each colour,
and set
\[
       P=\sum_B P_B,\qquad Q=I-P.
\]
Thus \(Q\) consists of vertex
fields with zero average on
every block. If \(P_{\rm c}\)
is the global constant projection,
the primed space is
\[
       \mathcal V_M
         =\operatorname{Ran}(P-P_{\rm c})
                          \oplus\operatorname{Ran}Q.
\]
Its low and high dimensions are
\[
       r_{\rm b}=3(N^4-1),\qquad
       d_{\rm h}=3(M^4-N^4).
\]
For example, taking
\(\ell=\lfloor1/(3R)\rfloor\ge2\)
and compatible volumes gives
\(r_{\rm b}\le3(6R)^4M^4\).
This is an extensive but
low-density gauge space when
\(R\) is small, not a fixed
number of physical variables.

\begin{lemma}[A high block positive beyond the gauge-positive domain]
\label{lem:v77-high}
Every \(u\in\operatorname{Ran}Q\) obeys
\begin{equation}\label{eq:v77-Poincare}
                       \|u\|^2\le\frac{\ell^2}{4}\|Du\|^2 .
\end{equation}
The restricted operator
\begin{equation}\label{eq:v77-A}
        A(a)=Q\mathcal M(a)Q\big|_{\operatorname{Ran}Q}
\end{equation}
therefore satisfies
\begin{equation}\label{eq:v77-high}
          A(a)\succeq\tfrac14 Q\Delta Q
                         \succeq\ell^{-2}I,
             \qquad \|A(a)^{-1}\|\le\ell^2 .
\end{equation}
\end{lemma}
\begin{proof}
The path Laplacian on
\(\{0,\ldots,\ell-1\}\), with
only its internal edges, has
eigenvalues
\[
           4\sin^2\frac{\pi j}{2\ell},
                        \qquad j=0,\ldots,\ell-1.
\]
For example the eigenvectors
\(\cos(\pi j(x+1/2)/\ell)\)
verify this by substitution at
both endpoints and interior sites;
their orthogonality follows from
the finite cosine sums.
The first positive eigenvalue
is at least \(4/\ell^2\), using
\(\sin t\ge2t/\pi\) on
\([0,\pi/2]\).
The four-dimensional internal
block Laplacian is the tensor
sum of four such operators.
Its kernel is precisely the
block constants and its first
positive eigenvalue has the
same lower bound.
Sum the resulting inequalities
over the disjoint blocks and
then include the nonnegative
remaining edge energies.
This proves \eqref{eq:v77-Poincare}.

By Lemma~\ref{lem:v77-Garding},
on this subspace the lower
form coefficient of \(\Delta\) is
\[
          \tfrac12-\tfrac94 R^2\ell^2
                                \ge\tfrac14.
\]
Another application of
\eqref{eq:v77-Poincare} gives
the floor and inverse bound.
No assumption on a random
minimum eigenvalue of
\(\mathcal F_M(a)\) has entered.
\end{proof}

\subsection{Exponential locality of the high inverse}

Use the nearest-neighbour graph
of blocks and its periodic graph
distance \(d_{\rm b}\).
For a set \(\mathcal X\) of blocks,
\(\chi_{\mathcal X}\) denotes
multiplication by its vertex
indicator. It commutes with \(P\)
and \(Q\). An inverse on
\(\operatorname{Ran}Q\) is
extended by zero on
\(\operatorname{Ran}P\).

\begin{theorem}[A volume-independent high inverse estimate]
\label{thm:v77-local-inverse}
For nonempty sets of blocks
\(\mathcal X,\mathcal Y\),
\begin{equation}\label{eq:v77-local-inverse}
 \|\chi_{\mathcal X}A(a)^{-1}\chi_{\mathcal Y}\|
       \le2\ell^2
              \exp\!\left(-\frac{d_{\rm b}(\mathcal X,\mathcal Y)}
                                      {5\ell}\right).
\end{equation}
The estimate holds throughout
the full Coulomb link ball,
including where the remaining
low gauge matrix is not positive.
\end{theorem}
\begin{proof}
Let \(f\) be a real, blockwise
constant function whose difference
across neighbouring blocks has
absolute value at most one.
Put \(W=e^{\theta f}\), with
\(\theta=1/(5\ell)\).
It preserves \(\operatorname{Ran}Q\)
and commutes with \(Q\).
Although it need not preserve the
full primed space, no such
preservation is needed for this
restricted inverse.

For the Hermitian part of the
weighted full matrix, each
off-diagonal entry is multiplied
by \(\cosh(\theta(f_x-f_y))\).
The diagonal is unchanged.
The row and column estimates
of Lemma~\ref{lem:v77-Garding}
and the operator-valued Schur
bound imply
\[
 \|\operatorname{Re}(W\mathcal M W^{-1})
                           -\mathcal M\|
       \le10(\cosh\theta-1)
       \le10\theta^2 .
\]
Here \(0<\theta\le1\), and
\(\cosh\theta-1\le\theta^2\).
After compression to \(Q\),
\[
 \operatorname{Re}(WAW^{-1})
       \succeq(\ell^{-2}-10\theta^2)I
       \succeq(2\ell^2)^{-1}I .
\]
A finite matrix with Hermitian
part at least \(\gamma I\) has
inverse norm at most \(\gamma^{-1}\):
Cauchy--Schwarz applied to its
quadratic form gives
\(\|Tu\|\ge\gamma\|u\|\).
Consequently
\[
                      \|WA^{-1}W^{-1}\|\le2\ell^2.
\]
Take \(f(B)=d_{\rm b}(B,\mathcal Y)\).
The weight is one on \(\mathcal Y\)
and at least
\(\exp(\theta d_{\rm b}(\mathcal X,\mathcal Y))\)
on \(\mathcal X\). This gives
\eqref{eq:v77-local-inverse}.
\end{proof}

The displayed rate is
\(1/(5\ell)\) per block step.
It corresponds to a length of
order \(\ell^2\) in fine-lattice
units; a length of order
\(\ell\) has not been proved.
For fixed admissible \(R,\ell\)
it is nevertheless uniform in
\(M\). This is a deterministic
matrix inverse, not a covariance
of the physical field.

\subsection{A local residual matrix and an exact determinant ratio}

Define on the full block-constant
space
\begin{equation}\label{eq:v77-Schur}
 \mathcal S(a)
       =P\mathcal M(a)P
           -P\mathcal M(a)Q A(a)^{-1}Q\mathcal M(a)P .
\end{equation}
Let \(A_0=A(0)\) and
\(\mathcal S_0=\mathcal S(0)
\big|_{\operatorname{Ran}(P-P_{\rm c})}\).
The latter is positive definite.
The normalized residual matrix is
\begin{equation}\label{eq:v77-normalized}
       F_{\rm b}(a)=\mathcal S_0^{-1/2}
          \mathcal S(a)\big|_{\operatorname{Ran}(P-P_{\rm c})}
                                      \mathcal S_0^{-1/2}.
\end{equation}
An empty residual determinant
when \(N=1\) is understood as one.

\begin{theorem}[Local Schur matrix and exact gauge density]
\label{thm:v77-Schur}
The matrix \(\mathcal S(a)\)
is self-adjoint, is matrix-concave
as a function of the Coulomb
connection, and annihilates
the three global constant
colour vectors.
Writing
\(d=d_{\rm b}(\mathcal X,\mathcal Y)\),
its blocks satisfy
\begin{equation}\label{eq:v77-Schur-local}
 \|P_{\mathcal X}\mathcal S(a)P_{\mathcal Y}\|
   \le20\,\mathbf1_{\{d\le1\}}
          +800\ell^2
                 \exp\!\left(-\frac{(d-2)_+}{5\ell}\right),
 \qquad P_{\mathcal X}=\chi_{\mathcal X}P .
\end{equation}
Moreover
\begin{equation}\label{eq:v77-det-ratio}
 \begin{aligned}
 \mathcal F_M(a)\succ0
          &\ \Longleftrightarrow\ F_{\rm b}(a)\succ0,\\
 \det\mathcal F_M(a)
          &=\frac{\det A(a)}{\det A_0}\det F_{\rm b}(a).
 \end{aligned}
\end{equation}
If \(\mathcal F_M(a_c)\succeq\kappa_0I\),
then \(F_{\rm b}(a_c)\succeq\kappa_0I\).
\end{theorem}
\begin{proof}
The high floor makes the
usual minimizing formula valid
without positivity of the full
matrix:
\[
 \langle x,\mathcal S(a)x\rangle
       =\inf_{y\in\operatorname{Ran}Q}
              \langle x+y,\mathcal M(a)(x+y)\rangle,
                          \qquad x\in\operatorname{Ran}P.
\]
The minimizer is
\(-A^{-1}Q\mathcal M Px\).
The same concave-infimum argument
proved in V76 applies:
\(\mathcal M(a)\) is matrix-concave
because its even part is so
and its odd part is linear.
This proves concavity of
\(\mathcal S\), and fixed
congruence proves it for \(F_{\rm b}\).
The constant kernel of
\(\mathcal M\) proves the
asserted constant kernel of
\(\mathcal S\).

The first term in
\eqref{eq:v77-Schur} vanishes
between block sets at distance
greater than one and has norm
at most 20 otherwise.
The maps \(Q\mathcal M P_{\mathcal Y}\)
and \(P_{\mathcal X}\mathcal M Q\)
have norm at most 20 and are
supported in the one-block
neighbourhoods of their respective
sets. Their distance is at
least \((d-2)_+\).
Apply Theorem~\ref{thm:v77-local-inverse}
between them. This proves
\eqref{eq:v77-Schur-local}.

Triangular factorization on the
primed decomposition gives
\[
 \det{}'\mathcal M(a)
      =\det A(a)\,
         \det\!\left(\mathcal S(a)
                        \big|_{\operatorname{Ran}(P-P_{\rm c})}\right).
\]
At \(a=0\), the same identity
gives
\(\det{}'\Delta=\det A_0\det\mathcal S_0\).
The free Schur form is strictly
positive on the primed low space:
zero energy in its minimizing
formula would make \(x+y\)
globally constant, hence \(x=0\)
on that space.
Since
\(\det\mathcal F_M=\det{}'\mathcal M/\det{}'\Delta\),
division proves the exact ratio.
Positivity follows from the
same triangular factorization
and the high floor.

Finally
\(\mathcal M(a_c)\succeq\kappa_0\Delta\)
on the full space, including
its constant kernel.
Taking infima in the two
quadratic forms gives
\(\mathcal S(a_c)\succeq\kappa_0\mathcal S(0)\).
The fixed congruence proves
the normalized base floor.
\end{proof}

Locality in \eqref{eq:v77-Schur-local}
is asserted for the unnormalized
\(\mathcal S\) in block coordinates.
It is not automatically inherited
by \(F_{\rm b}\): the factors
\(\mathcal S_0^{-1/2}\) are not
assumed local. Likewise the
constant kernel is a gauge
parameter identity, not a
physical particle mass statement.

\subsection{The eliminated determinant has controlled, nonzero interactions}

Define its dimensionless potential
\begin{equation}\label{eq:v77-high-potential}
                 T_{\rm h}(a)
                     =-\log\frac{\det A(a)}{\det A_0}.
\end{equation}
This is smooth and convex on
the full affine Coulomb link
ball. Let \(v,w\) be Coulomb
link variations. Denote by
\(\mathcal B_v,\mathcal B_w\)
the blocks containing endpoints
of links in their supports.
Set
\[
 L_R=16+48R,\qquad
 d_v=3(\ell^4-1)|\mathcal B_v|,\qquad
 d_w=3(\ell^4-1)|\mathcal B_w|.
\]

\begin{proposition}[Exponential decay of separated mixed derivatives]
\label{prop:v77-high-Hessian}
If \(v,w\) have disjoint link
supports, then
\begin{equation}\label{eq:v77-high-Hessian}
 \begin{split}
 |D^2T_{\rm h}(a)[v,w]|
       \le{}&4\ell^4\min(d_v,d_w)L_R^2
                        \|v\|_\infty\|w\|_\infty\\
       &{}\times
         \exp\!\left(-\frac{2d_{\rm b}(\mathcal B_v,\mathcal B_w)}
                                     {5\ell}\right).
 \end{split}
\end{equation}
All constants are independent
of total volume for fixed
insertion supports and block size.
\end{proposition}
\begin{proof}
The linkwise derivative estimate
proved in Lemma~\ref{lem:v74-even}
gives
\(\|D\sigma(\operatorname{ad}_a)[v]\|
\le3R\|v\|_\infty\).
Using \(\|D\|,\|S_{\rm end}\|\le4\),
\[
       \|D\mathcal M(a)[v]\|
                         \le(48R+16)\|v\|_\infty.
\]
After multiplication by \(Q\)
the derivative
\(A_v=QD\mathcal M(a)[v]Q\)
is supported on
\(\operatorname{Ran}(Q\chi_{\mathcal B_v})\),
of dimension \(d_v\).
The analogous statements hold
for \(w\).
The linkwise dependence of
the coefficients implies
\(D^2\mathcal M(a)[v,w]=0\)
for disjoint link supports.
Thus the exact determinant
derivative is
\[
                 D^2T_{\rm h}(a)[v,w]
                        =\Tr(A^{-1}A_vA^{-1}A_w).
\]
Insert the two support
projections. The resulting trace
contains two high inverse
operators between the specified
sets. Its absolute value is
at most their norm product,
\(\|A_v\|\|A_w\|\), and the
smaller support dimension.
Apply \eqref{eq:v77-local-inverse}
to both inverse factors.
This proves the displayed bound.
Convexity of \(T_{\rm h}\) follows
separately from matrix concavity
and positivity of \(A\);
individual mixed derivatives
need not have a fixed sign.
\end{proof}

\begin{record}[An exact separated-block derivative is not zero]
\label{rec:v77-nonzero}
On the side-four torus with
\(\ell=2\), let \(v\) be the
unit \(E_3\)-coloured plaquette
current based at \(0\), in
directions 1 and 2. Let \(w\)
be its translate by
\((2,2,2,2)\). They are
mean-zero Coulomb link variations,
each contained in one block,
and their block distance is four.
At the zero connection,
\begin{equation}\label{eq:v77-nonzero}
                       D^2T_{\rm h}(0)[v,w]=\frac1{50176}>0 .
\end{equation}
\end{record}
\begin{proof}
The derivative of \(\sigma\)
at zero is zero. Write
\[
 J=\begin{pmatrix}0&-1&0\\1&0&0\\0&0&0\end{pmatrix}.
\]
The odd derivative is
\(K_v\otimes J\), with scalar
skew matrix having edge entries
\((K_v)_{xy}=-v_e\),
\((K_v)_{yx}=v_e\).
The diagonal is zero by the
Coulomb condition.
Let \(G\) be the scalar
extension by zero of
\((Q\Delta Q)^{-1}\).
The mixed derivative is
\[
                  -2\Tr(GK_vGK_w),
\]
because \(\Tr J^2=-2\).

Here is a finite explicit
formula for every required entry
of \(G\). Write \(x=2b+s\),
\(y=2c+t\), with all four
vectors in \(\{0,1\}^4\).
The block Fourier and internal
Walsh bases give
\[
 G(x,y)=\frac1{256}
  \sum_{\substack{k,\rho\in\{0,1\}^4\\\rho\ne0}}
  \frac{(-1)^{k\cdot(b-c)+\rho\cdot(s-t)}}
       {\sum_{\mu=1}^4
          \{2-(1+(-1)^{k_\mu})(-1)^{\rho_\mu}\}}.
\]
All denominators in this sum
are positive. The exclusion
\(\rho=0\) is exactly the
block-zero-mean condition, not
merely removal of the global
constant mode.

Order the four current sites
as \(0,e_2,e_1,e_1+e_2\),
and use the translated order
for the other current.
The cross Green matrix and
the two identical scalar
current matrices are
\[
 G_{vw}=-\frac1{430080}
 \begin{pmatrix}
 533&323&323&53\\
 323&533&53&323\\
 323&53&533&323\\
 53&323&323&533
 \end{pmatrix},
 \qquad
 K=\begin{pmatrix}
 0&1&-1&0\\
 -1&0&0&1\\
 1&0&0&-1\\
 0&-1&1&0
 \end{pmatrix}.
\]
Substitution gives
\(-2\Tr(G_{vw}KG_{vw}^{T}K)=1/50176\).
This is rational finite
matrix multiplication and is
also checked independently
by the certificate.
\end{proof}

This example excludes exact
factorization of the high
determinant into independent
single-block terms. It does
not assert a lower bound on
long-distance interactions at
arbitrary volume. The uniform
result is the upper bound
\eqref{eq:v77-high-Hessian}.

\subsection{The same interacting law and the partition derivative cost}

Return to the actual physical
coordinate law of V75,
with \(a=a_c+h\).
The fixed partition introduces
no new physical integration
variable. Put
\[
                       W_{\rm b}(a)=\mathcal H(a)+T_{\rm h}(a).
\]
Equations \eqref{eq:v75-density}
and \eqref{eq:v77-det-ratio}
give exactly
\begin{equation}\label{eq:v77-law}
 \begin{split}
 d\mu_c(h)={}&\mathcal Z_c^{-1}
  \mathbf1_{\{\max_e|a_c(e)+h_e|<R,\ F_{\rm b}(a_c+h)\succ0\}}\\
 &{}\times\det F_{\rm b}(a_c+h)
          e^{-\beta S(a_c+h)-W_{\rm b}(a_c+h)}\,dh .
 \end{split}
\end{equation}
The normalizer \(\mathcal Z_c\)
is unchanged. The free determinant
in \(T_{\rm h}\), as well as
the normalization by
\(\mathcal S_0\), is essential
to this exact statement.

\begin{proposition}[A reduced barrier with an explicit partition cost]
\label{prop:v77-barrier}
The potential \(W_{\rm b}\)
has Hessian at least \(2I/3\)
on physical Coulomb variations.
For homogeneous \(a_c\) in
the link ball, every admissible
mean-zero \(h\) satisfies
\[
 DW_{\rm b}(a_c+h)[h]\ge-\mathcal J_{\rm part},
 \qquad
 \mathcal J_{\rm part}=
 \begin{cases}
 0,&a_c=0,\\
 2R\ell^2d_{\rm h}L_R,&\text{in general}.
 \end{cases}
\]
Assume \(N>1\), put
\[
 f_{\rm b}(h)=\lambda_{\min}F_{\rm b}(a_c+h),\qquad
 K_{\rm b}=n+r_{\rm b}+\mathcal J+\mathcal J_{\rm part},
\]
and retain V75's base floor
\(\kappa_0\) and Wilson radial
cost \(\mathcal J\), including
its improved choice when
\(a_c=0\). Then
\begin{equation}\label{eq:v77-barrier}
 \begin{aligned}
 \kappa_0\mathbb E_{\mu_c}\Tr F_{\rm b}^{-1}
                                &\le K_{\rm b},\\
 \mathbb E_{\mu_c}(\kappa_0/f_{\rm b})^p
                                &\le\frac{K_{\rm b}^p}{2-p},
                                      &&1\le p<2,\\
 \mu_c\{f_{\rm b}<\delta\}
                       &\le\min\{1,2e^2(K_{\rm b}\delta/\kappa_0)^2\}.
 \end{aligned}
\end{equation}
\end{proposition}
\begin{proof}
The high determinant is convex,
and the Haar potential has
the inherited Hessian floor.
Hence
\[
 DW_{\rm b}(a_c+h)[h]
                  \ge DW_{\rm b}(a_c)[h]+\tfrac23\|h\|^2.
\]
The Haar derivative at homogeneous
\(a_c\) pairs to zero with
mean-zero \(h\).
The remaining derivative is
\[
                  DT_{\rm h}(a_c)[h]
                         =-\Tr(A(a_c)^{-1}D A(a_c)[h]).
\]
The inverse bound \(\ell^2\),
the derivative bound
\(L_R\|h\|_\infty\), and
\(\|h\|_\infty\le2R\)
give its absolute value at
most \(2R\ell^2d_{\rm h}L_R\).
This proves the general case.

At zero background,
\(A_0\) is scalar in colour.
The even derivative is zero,
and every odd colour matrix
\(\operatorname{ad}_{h_e}\)
has trace zero. Thus
\(DT_{\rm h}(0)[h]=0\) exactly,
proving the improved case.
A fixed partition commutes
only with block translations;
we have not used it as if it
commuted with every translation.

For the remaining claims,
\(F_{\rm b}\) is concave,
has the stipulated base floor,
and is positive precisely
on the original convex domain.
In the radial arguments of
V75, the adverse Wilson
derivative is bounded by
\(\mathcal J\); the additional
geometric adverse derivative
is now bounded by
\(\mathcal J_{\rm part}\).
Their sum replaces
\(\mathcal J\) in every
truncated integration and
one-dimensional ray estimate.
The determinant dimension is
\(r_{\rm b}\); the radial
physical dimension remains
\(n=9(M^4-1)\).
The proofs of V75 therefore
give exactly
\eqref{eq:v77-barrier},
including the boundary flux
and the direct quadratic
probability bound.
\end{proof}

The new geometric representation
does not remove the extensive
physical dimension. Its general
partition cost is deliberately
crude and is not claimed to
improve V75's numerical bound.
At zero background the cost
vanishes, but \(n\) and the
Wilson radial cost remain.
When \(N=1\), there is no
residual gauge determinant;
the high floor already holds
for the whole primed matrix,
and the residual-minimum
statements are unnecessary.

\subsection{What is now local, and what still prevents closure}

The bottleneck has changed in
one definite respect. V76 had
an exact but spatially uncontrolled
Fourier residual. We now have
an exact real-space residual
whose unnormalized matrix entries
decay exponentially, together
with a corresponding locality
bound for separated derivatives
of the high determinant.
Both statements hold on every
admissible small-link configuration,
not only under a Gaussian
reference law or after conditioning
on an unproved gauge floor.

They do not prove that the
residual determinant is weakly
coupled. Derivatives of
\(-\log\det F_{\rm b}\) involve
its inverse, whose small spectrum
is precisely the unresolved
part. The normalization
\(\mathcal S_0^{-1/2}\) must
also be kept distinct from
local unnormalized coefficients.
No product bound for separated
bad events and no exponential
decay of physical correlations
has been obtained.

The next useful estimate must
control that residual low
operator under \(\mu_c\), or
use the local unnormalized form
to obtain a multiscale estimate
without assuming its conclusion
as a uniform floor. Repeating
the radial argument with only
a smaller determinant rank
does not accomplish this.
Compact-sector coverage,
uniform interacting continuum
limits and physical transfer
calibration remain necessary.
The next companion checkpoint
is still V80.

\subsection{Exact verification and predecessor preservation}

The new certificate is
\begin{center}\small
\path{scripts/verify_ym_f14_v77_local_gauge_schur.py}.
\end{center}
It verifies the path Poincare
forms by exact rational
\(LDL^T\) elimination for
\(\ell=2,\ldots,12\), together
with the radius and weighted
coercivity constants.
These finite checks supplement
the all-\(\ell\) proof above.

It then builds all 16 free
block Fourier matrices for
the actual side-four torus,
with two-site cells. In a
nonorthonormal high basis it
keeps the Gram determinant
explicitly, checking the
determinant ratio, high inverse
and full inverse identity
at every nonzero coarse momentum.
The zero coarse momentum has
the required constant kernel.
To avoid relying only on the
special side-four phase cancellation,
three actual side-eight coarse
momenta with phases
\((1,0,0,0)\), \((1,1,0,0)\),
and \((1,2,3,0)\), in units
of \(\pi/2\), are checked
by realifying their complex
matrices. Their nonzero
high--low couplings give scalar
Schur entries
\(2/3,\ 16/11,\ 84/23\),
respectively.

For the separated actual
gauge derivatives at zero,
the high Green entries are
computed independently by
block inverse and by the
explicit Walsh sum.
The antipodal value in
\eqref{eq:v77-nonzero} is
reproduced exactly. Two
other declared current pairs
give \(1089/1254400\) and
\(-1273/716800\). Mixed
derivatives can therefore
have either sign, consistently
with convexity of the complete
high potential.
No finite-pole replacement
of \(r\cot r\) is used:
the finite matrix checks are
at the actual zero connection
and its exact first derivatives.

The certificate contains 9,045 bytes
and has SHA--256
\begin{center}\scriptsize\ttfamily
27957C85E37D9B7781135807CB1EC6DD9207A57696B6DDD5B2F501601D7D7674.
\end{center}
Its inherited helpers are unchanged.
The predecessor V76 contains
1,884,064 bytes and has SHA--256
\begin{center}\scriptsize\ttfamily
72BD31BE8AFC6851C07EBFC365F887F10EF155AE5B4C9A4B225F0A8A23CAA898.
\end{center}
That body is preserved except
for the title version and
this terminal insertion.
V77 replaces V76 as the
single active main note;
the complete prior content
remains here. Companion files,
archives and supplied sources
are not modified.

\begin{firewall}
The new estimates concern a
local gauge-parameter elimination.
They preserve the full
interacting density and prove
volume-independent spatial
decay for the eliminated
inverse and specified derivatives.
They do not prove decay of
physical correlations or a
physical Hamiltonian gap.
The low residual spectrum,
continuum construction and
mass-gap objective remain open.
\end{firewall}
\section{V78: a translation-averaged determinant barrier without the partition cost}

\subsection{Context: repair the symmetry before further coarse graining}

V77 proves a local gauge-space
elimination, but its fixed block
partition produces an additional
radial cost at a nonzero homogeneous
background. Before attempting to
iterate that estimate, this note
removes the artificial partition
cost without changing the measure.
We average the exact logarithmic
determinant decompositions over
all translates of the partition.
This restores one-site translation
invariance. It is neither an
average of approximate densities
nor a replacement of the input
law by independent coarse copies.

The residual density is a
geometric mean of determinants
of several different Schur
matrices for the same connection.
Its effective determinant rank
in the radial argument is the
rank of one coarse system,
not that rank multiplied by
the number of partitions.
To prove a quadratic probability
bound, we use the geometric
mean of their minimum eigenvalues.
An energy-normalized comparison
then connects this scalar
barrier to the minimum of
the original gauge matrix
with constants independent
of the torus size.

The physical-coordinate dimension
and Wilson radial cost remain.
Thus this repairs one genuine
loss in V77, but does not
close the infrared probability
estimate or the continuum
mass-gap problem. In particular
it supplies no decrease of the
small-field radius under a
renormalization step.

\subsection{All translated partitions and the unchanged density}

Retain V77's assumptions
\[
 M=N\ell,\qquad N>1,\qquad \ell\ge2,\qquad
       0<R\le1/4,\qquad R\ell\le1/3.
\]
Let
\[
       \mathcal T=\{0,\ldots,\ell-1\}^4,\qquad
       m_{\rm p}=|\mathcal T|=\ell^4,\qquad
       r=3(N^4-1).
\]
For \(s\in\mathcal T\), translate
the entire V77 partition by
\(s\), modulo \(M\), and attach
a subscript \(s\) to its
operators \(P_s,Q_s,A_s,\mathcal S_s\).
The normalized residual matrix
from \eqref{eq:v77-normalized}
will be denoted
\[
 B_s(a)=\mathcal S_{s,0}^{-1/2}
   \mathcal S_s(a)\big|_{\operatorname{Ran}(P_s-P_{\rm c})}
                                      \mathcal S_{s,0}^{-1/2}.
\]
Here \(\mathcal S_{s,0}\) is
the free residual, not the
residual at \(a_c\).
All these low matrices have
the same dimension \(r\).
They need not commute and
act on different low spaces.
No identification of their
eigenvectors is made.

Define
\begin{equation}\label{eq:v78-average}
 \begin{aligned}
 \overline T(a)&=-\frac1{m_{\rm p}}
       \sum_{s\in\mathcal T}\log\frac{\det A_s(a)}{\det A_s(0)},\\
 \overline W(a)&=\mathcal H(a)+\overline T(a),\\
 \mathcal D(a)&=\prod_{s\in\mathcal T}\det B_s(a)^{1/m_{\rm p}}
                           \quad\hbox{when }\mathcal F_M(a)\succ0.
 \end{aligned}
\end{equation}
The high potential is smooth
on the whole Coulomb link
ball, not just the full
gauge-positive domain.

\begin{theorem}[An exact symmetric decomposition]
\label{thm:v78-symmetric}
The function \(\overline W\)
is invariant under every
one-site translation and has
Hessian at least \(2I/3\)
on the physical Coulomb
variations. At every homogeneous
background \(a_c\), its derivative
pairs to zero with every
mean-zero physical variation.
On the positive domain,
\begin{equation}\label{eq:v78-det}
                    \det\mathcal F_M(a)
                              =e^{-\overline T(a)}\mathcal D(a).
\end{equation}
Consequently the actual law
of V75 is exactly
\begin{equation}\label{eq:v78-law}
 d\mu_c(h)=\mathcal Z_c^{-1}
  \mathbf1_{\mathcal A_c}(h)\,
  \mathcal D(a_c+h)
        e^{-\beta S(a_c+h)-\overline W(a_c+h)}\,dh ,
\end{equation}
with the same domain and
normalizer as before.
\end{theorem}
\begin{proof}
Each translated high block
has the V77 uniform floor.
Its negative logarithmic
determinant ratio is convex.
Averaging and adding the
Haar potential proves the
Hessian assertion.

A one-site translation permutes
the partitions in \(\mathcal T\),
possibly relabelling whole blocks.
The high operators and their
free reference operators are
unitarily conjugated under
that relabelling. Their
determinant ratios are unchanged.
The average, and hence
\(\overline W\), is therefore
translation invariant.
At homogeneous \(a_c\), average
any variation over all torus
translations. A mean-zero link
variation has zero such average,
separately in each orientation
and colour. Translation
invariance of the derivative
then gives
\(D\overline W(a_c)[h]=0\).
Convexity in particular yields
\begin{equation}\label{eq:v78-radial-geometric}
             D\overline W(a_c+h)[h]\ge\tfrac23\|h\|^2\ge0.
\end{equation}

For every \(s\), V77 gives
the same exact identity
\[
 \log\det\mathcal F_M(a)
     =\log\frac{\det A_s(a)}{\det A_s(0)}
                                      +\log\det B_s(a).
\]
Average these identities and
exponentiate to obtain
\eqref{eq:v78-det}.
Substitution into the original
density proves \eqref{eq:v78-law}.
The full positivity event is
equivalent to \(B_s(a)\succ0\)
for each \(s\), since all
high blocks are positive.
Thus no extra intersection
event or altered normalization
has been inserted.
\end{proof}

The index \(s\) is an algebraic
averaging index, not a new
physical field. The fractional
exponents sum to one. In
particular no power of the
original density other than
one appears in
\eqref{eq:v78-law}.

\subsection{The averaged high potential remains spatially controlled}

Let \(v,w\) be Coulomb link
variations with disjoint link
supports. Let \(V_v,V_w\)
be the sets of their endpoint
vertices, let
\(n_v=|V_v|\), \(n_w=|V_w|\),
and let \(d_{\rm f}(V_v,V_w)\)
be their fine-lattice graph
distance. Retain
\(L_R=16+48R\).

\begin{proposition}[No multiplicity loss in the local derivative bound]
\label{prop:v78-local}
For nonzero such variations,
\begin{equation}\label{eq:v78-local}
 \begin{split}
 |D^2\overline T(a)[v,w]|
   \le{}&12\ell^4(\ell^4-1)\min(n_v,n_w)L_R^2
                                      \|v\|_\infty\|w\|_\infty\\
    &{}\times\exp\!\left(
           -\frac{2(d_{\rm f}(V_v,V_w)-4\ell)_+}{5\ell^2}
                     \right).
 \end{split}
\end{equation}
There is no factor \(m_{\rm p}\).
\end{proposition}
\begin{proof}
For each translated partition,
the number of blocks touched
by \(V_v\) is at most \(n_v\),
and similarly for \(w\).
The V77 support dimension
is therefore at most
\(3(\ell^4-1)n_v\).
If \(d_s\) is the minimum
block distance between their
touched blocks, choose endpoint
vertices in a pair of blocks
attaining \(d_s\). The periodic
coordinate metric gives
\[
             d_{\rm f}(V_v,V_w)
                             \le\ell d_s+4(\ell-1).
\]
Thus
\(d_s\ge(d_{\rm f}(V_v,V_w)-4\ell)_+/\ell\).
Substitute both bounds into
\eqref{eq:v77-high-Hessian}.
The resulting estimate is
the same for every \(s\).
The absolute value of their
average is bounded by that
common estimate, proving
\eqref{eq:v78-local}.
\end{proof}

This is locality of a high
determinant derivative, not
independence of separated
physical fluctuations. The
averaged residual factor
\(\mathcal D\) is still present.

\subsection{A common scalar barrier for the fractional determinants}

Let the known homogeneous
base floor be
\(\mathcal F_M(a_c)\succeq\kappa_0I\).
By V77, \(B_s(a_c)\succeq\kappa_0I\)
for every \(s\).
On the actual positive domain set
\begin{equation}\label{eq:v78-g}
       f_s(h)=\lambda_{\min}B_s(a_c+h),\qquad
       g(h)=\prod_{s\in\mathcal T}f_s(h)^{1/m_{\rm p}}.
\end{equation}
The geometric mean is of
positive real scalars; it is
not a determinant or a matrix
geometric mean.

\begin{lemma}[Concavity and the radial derivative identities]
\label{lem:v78-g}
The function \(g\) is positive
and concave on the actual
domain, with \(g(0)\ge\kappa_0\).
At almost every point of a
radial line,
\begin{equation}\label{eq:v78-radial}
 \begin{aligned}
 DB_s(a_c+h)[h]&\preceq B_s(a_c+h)-\kappa_0I,\\
 D\log g(h)[h]&=\frac1{m_{\rm p}}
                     \sum_s\frac{Df_s(h)[h]}{f_s(h)},\\
 D\log\!\left(\frac{\mathcal D(a_c+h)}{g(h)}\right)[h]
                                                    &\le r-1 .
 \end{aligned}
\end{equation}
For \(0\le q<1\), also
\begin{equation}\label{eq:v78-weighted}
 \begin{split}
 &-\frac1{m_{\rm p}}\sum_s
       \Tr\!\left(B_s^{-1}DB_s[h]\right)
                  +qD\log g[h]\\
 &\hspace{12mm}\ge
      (1-q)\kappa_0\frac1{m_{\rm p}}\sum_s\frac1{f_s}-(r-q)
       \ge\frac{(1-q)\kappa_0}{g}-(r-q).
 \end{split}
\end{equation}
\end{lemma}
\begin{proof}
Every \(B_s\) is matrix-concave,
so its minimum eigenvalue is
concave. The weighted geometric
mean \(G(x)=\prod_s x_s^{w_s}\),
where \(w_s=1/m_{\rm p}\),
is increasing and concave on
the positive orthant. Directly,
\[
 \frac{D^2G(x)[z,z]}{G(x)}
     =\left(\sum_s w_s z_s/x_s\right)^2
                            -\sum_s w_s(z_s/x_s)^2\le0.
\]
Its composition with the
positive concave \(f_s\)
is therefore concave.
This proves the first assertion.

The tangent inequality for
matrix concavity gives
\[
 B_s(a_c)\preceq B_s(a_c+h)-DB_s(a_c+h)[h],
\]
which proves the first line
of \eqref{eq:v78-radial}.
The second line is the scalar
product rule, valid almost
everywhere on rays.
At such a differentiability
point choose a minimum
eigenvector realizing
\(Df_s[h]\), and complete
it to an orthonormal
eigenbasis of \(B_s\).
The first inequality bounds
the diagonal derivative in
each eigenvector of eigenvalue
\(\lambda_i\) by
\(\lambda_i-\kappa_0\).
Omitting one minimum factor
therefore gives
\[
 D(\log\det B_s-\log f_s)[h]
       \le r-1-\kappa_0
                        \left(\Tr B_s^{-1}-\frac1{f_s}\right)
       \le r-1.
\]
Average this inequality to
obtain the final line.

For \eqref{eq:v78-weighted},
the coefficient of that one
minimum-eigenvalue derivative
is \(-(1-q)\), and every
other coefficient is \(-1\).
All are nonpositive when
\(q<1\). The same diagonal
inequalities imply, for each \(s\),
\[
 -\Tr(B_s^{-1}DB_s[h])
          +q\,Df_s[h]/f_s
       \ge (1-q)\kappa_0/f_s-(r-q).
\]
Average, and apply the
arithmetic--geometric mean
inequality to the positive
numbers \(1/f_s\).
This proves the result.
At eigenvalue crossings the
radial minimum functions are
locally Lipschitz and
differentiable almost everywhere;
the inequalities extend by
their absolutely continuous
radial representatives.
\end{proof}

In particular the permitted
range is \(q<1\), not
\(q<1/m_{\rm p}\).
The same scalar logarithmic
weight removes a fraction
\(q\) of the chosen minimum
factor in \emph{every} partition.
This is the reason that
the number of partitions
does not worsen the
inverse-moment endpoint.

\subsection{The actual probability bound with the reduced rank}

Let \(\mathcal J\) be precisely
the V75 Wilson radial cost:
\[
 \mathcal J=4\beta C_0R^4M^4,
 \quad\hbox{or }\quad
 \mathcal J=\tfrac43\beta C_0R^4M^4
                                      \ \hbox{if }a_c=0.
\]
Define
\begin{equation}\label{eq:v78-K}
                    K_{\rm av}=n+r+\mathcal J,\qquad n=9(M^4-1).
\end{equation}
There is no \(\mathcal J_{\rm part}\).

\begin{theorem}[A symmetric reduced determinant barrier]
\label{thm:v78-probability}
Under the actual law \(\mu_c\),
\begin{equation}\label{eq:v78-moments}
 \begin{aligned}
 \kappa_0\mathbb E_{\mu_c}
       \frac1{m_{\rm p}}\sum_s\Tr B_s^{-1}
                                       &\le K_{\rm av},\\
 \mathbb E_{\mu_c}(\kappa_0/g)^p
                                       &\le K_{\rm av}^{p},
                                                 &&0<p\le1,\\
 \mathbb E_{\mu_c}(\kappa_0/g)^p
                                       &\le\frac{K_{\rm av}^{p}}{2-p},
                                                 &&1\le p<2,\\
 \mu_c\{g<\delta\}
                &\le\min\{1,\,2e^2(K_{\rm av}\delta/\kappa_0)^2\},
                                                 &&\delta>0 .
 \end{aligned}
\end{equation}
Furthermore, for every individual
translated partition,
\begin{equation}\label{eq:v78-each-trace}
                    \kappa_0\mathbb E_{\mu_c}\Tr B_s^{-1}
                                             \le K_{\rm av}.
\end{equation}
\end{theorem}
\begin{proof}
For the unweighted trace bound,
radial differentiation of the
negative logarithm of
\eqref{eq:v78-law}, together
with \eqref{eq:v78-radial-geometric},
gives
\[
 DU(h)[h]\ge
      \kappa_0\frac1{m_{\rm p}}\sum_s\Tr B_s^{-1}
                                              -r-\mathcal J.
\]
Perform the V75 radial
integration on the convex
truncation \(g>\varepsilon\)
inside the link ball, with
\(0<\varepsilon<\kappa_0\).
Concavity of \(g\) makes
this truncation convex and
it contains the origin.
The outward radial boundary
flux is nonnegative. Thus
the integrated radial derivative
is at most \(n\) times the
truncated mass. Let
\(\varepsilon\) decrease to
zero. Monotone convergence
for the nonnegative trace
integrand proves the first
bound with \(K_{\rm av}\).
The geometric--arithmetic
mean inequality gives
\(1/g\le m_{\rm p}^{-1}\sum_s1/f_s
\le m_{\rm p}^{-1}\sum_s\Tr B_s^{-1}\).
Jensen then proves the
range \(0<p\le1\).

For the density weighted by
\(g^{-q}\), \(0\le q<1\),
the radial logarithmic
potential is \(U+q\log g\).
By \eqref{eq:v78-weighted},
the same truncated integration
gives, after exhaustion,
\begin{equation}\label{eq:v78-recursion}
 (1-q)\kappa_0
       \mathbb E_{\mu_c}g^{-q-1}
                  \le(K_{\rm av}-q)\mathbb E_{\mu_c}g^{-q}.
\end{equation}
The right side is finite
by the already proved
low moments. Jensen bounds
it by
\((K_{\rm av}-q)(K_{\rm av}/\kappa_0)^q\).
With \(p=1+q\) this proves
the stated inverse moments.
The nonnegative radial flux
is retained in both integrations.

For the quadratic tail,
fix a ray \(h=t\theta\).
Write its exact density
as \(g(t)b(t)\), where
\[
 b(t)=t^{n-1}
         \frac{\mathcal D(a_c+t\theta)}{g(t\theta)}
            e^{-\beta S(a_c+t\theta)-\overline W(a_c+t\theta)}.
\]
Lemma~\ref{lem:v78-g} and
the radial action estimates give
\[
                         t(\log b)'(t)\le K_{\rm av}-2 .
\]
Thus the complete hypotheses
of the V75 one-ray argument
hold with \(f,K\) replaced
by \(g,K_{\rm av}\).
Explicitly, at a first crossing
\(t_\delta\) with \(g=\delta\),
for \(x=\delta/\kappa_0\le1/(2K_{\rm av})\)
the bad integral is at most
\[
 \frac{\delta^2t_\delta}{2(\kappa_0-\delta)}
              (1-x)^{-(K_{\rm av}-2)}b(t_\delta),
\]
whereas the earlier interval
\[
 \left[(1-1/K_{\rm av})t_\delta,\,
                  (1-1/(2K_{\rm av}))t_\delta\right]
\]
contributes at least
\(\kappa_0t_\delta b(t_\delta)/(4eK_{\rm av}^2)\).
Their ratio is at most
\(2e^2(K_{\rm av}x)^2\).
Integration over rays uses
the original normalizer;
larger \(\delta\) use the
trivial bound one.
This proves the fourth line.

Finally \(\mu_c\) is invariant
under every torus translation,
because \(a_c\) is homogeneous.
Translation of the physical
field maps \(\Tr B_s^{-1}\)
to the trace for another
translated partition, with
unitarily conjugated free
normalizations. All their
expectations are therefore
equal. The bound for the
average is the bound for
each one, proving
\eqref{eq:v78-each-trace}.
No independence assertion is used.
\end{proof}

The last assertion improves
the individual trace conclusion
in V77 for homogeneous
backgrounds: the partition cost
is gone even for a fixed
partition's expected trace.
It does not by itself replace
the individual minimum by \(g\)
in every higher-moment statement.
The following comparison makes
that distinction quantitative.

\subsection{Energy normalization compares all residual minima to the full one}

For a fixed partition let
\(\langle u,v\rangle_\Delta=\langle Du,Dv\rangle\)
on the primed vertex space.
Its high subspace is
\(\operatorname{Ran}Q_s\).
The complementary low subspace
is its \(\Delta\)-orthogonal
complement, consisting of free
harmonic lifts of block constants:
\[
           E_s x=x-A_s(0)^{-1}Q_s\Delta x,
                      \qquad x\in\operatorname{Ran}(P_s-P_{\rm c}).
\]
Here the high inverse is
extended by zero as in V77.
The lift satisfies
\[
 Q_s\Delta E_sx=0,\qquad
             E_s^*\Delta E_s=\mathcal S_{s,0}.
\]
This orthogonal decomposition
in the free energy inner
product is not the Euclidean
block-constant decomposition.

\begin{lemma}[Uniform comparison near the common gauge boundary]
\label{lem:v78-comparison}
Put
\[
 \eta=2R\ell+R^2/2,\qquad c=1-\eta,
           \qquad \eta\le3/4,\quad c\ge1/4 .
\]
Let \(f(a)=\lambda_{\min}\mathcal F_M(a)>0\).
For every translated partition,
\begin{equation}\label{eq:v78-comparison}
        f(a)\le f_s(a-a_c),\qquad
        f(a)\ge c^2\min\{f_s(a-a_c),c\}
                 \ge\frac1{16}\min\{f_s(a-a_c),1/4\}.
\end{equation}
Consequently, if \(f(a)<1/64\),
\begin{equation}\label{eq:v78-common}
                       f(a)\le g(a-a_c)\le16f(a).
\end{equation}
The constants are independent
of \(M\) and of the partition.
\end{lemma}
\begin{proof}
Write the odd part of
\(\mathcal M\) as
\(O=D^*\operatorname{ad}_aS_{\rm end}/2\).
The Coulomb weighted estimate
\eqref{eq:v76-weighted} is
equivalently
\[
                       \|O\Delta^{-1/2}\|\le4R.
\]
For any primed \(v\) and
any block-zero-mean \(u\),
the even bound and the
block Poincare inequality give
\[
 \begin{split}
 |\langle u,(\mathcal M-\Delta)v\rangle|
  &\le (R^2/2)\|u\|_\Delta\|v\|_\Delta
                                      +4R\|u\|\|v\|_\Delta\\
  &\le (R^2/2+2R\ell)\|u\|_\Delta\|v\|_\Delta.
 \end{split}
\]
Use energy-orthonormal coordinates
on the high space and on
the free harmonic lifts.
The matrix of the full form
in these coordinates is
unitarily equivalent to
\(\mathcal F_M(a)\).
Its block representation is
\[
           \begin{pmatrix} L_s&C_s^*\\C_s&D_s\end{pmatrix},
       \qquad
       \|D_s-I\|\le\eta,\quad\|C_s\|\le\eta,\quad D_s\succeq cI.
\]
Here \(D_s=A_s(0)^{-1/2}A_s(a)A_s(0)^{-1/2}\)
in the high coordinates.
Eliminating its high variable
produces exactly \(B_s(a)\):
the free harmonic lift differs
from the block-constant lift
only by an allowed high vector,
so both minimizations have
the same fixed block average.
The low energy normalization
is precisely \(\mathcal S_{s,0}\).

The triangular factorization
with Schur complement \(B_s\)
has inverse triangular norm
at most
\[
                 1+\|D_s^{-1}C_s\|
                              \le1+\eta/c=1/c.
\]
Its lower eigenvalue bound
is therefore
\(f\ge c^2\min(f_s,c)\).
The reverse inequality \(f\le f_s\)
follows either by testing on
the minimizing high lift of
a unit low eigenvector, whose
full energy-coordinate norm
is at least one, or from
the variational comparison
\(\mathcal M\ge f\Delta\).
This proves \eqref{eq:v78-comparison}.

If \(f<1/64\), the last
inequality excludes \(f_s\ge1/4\)
and yields \(f_s\le16f\)
for every \(s\). Taking
geometric means, together
with \(f_s\ge f\), proves
\eqref{eq:v78-common}.
\end{proof}

\begin{proposition}[A full gauge-floor consequence, with its constants visible]
\label{prop:v78-full-tail}
For \(0<\delta<1/64\),
\begin{equation}\label{eq:v78-full-tail}
 \mu_c\{\lambda_{\min}\mathcal F_M(a_c+h)<\delta\}
      \le\min\{1,\,2e^2(16K_{\rm av}\delta/\kappa_0)^2\}.
\end{equation}
For \(1\le p<2\), one also has
\[
 \mathbb E_{\mu_c} f^{-p}
        \le64^p+
          \frac{16^pK_{\rm av}^{p}}{(2-p)\kappa_0^p}.
\]
The same bound holds with
\(f\) replaced by any \(f_s\).
\end{proposition}
\begin{proof}
On \(f<\delta<1/64\),
\eqref{eq:v78-common} gives
\(g<16\delta\).
Apply Theorem~\ref{thm:v78-probability}.
For the moment bound split
at \(f=1/64\).
Above it the inverse is
bounded by \(64\); below
it \(f^{-1}\le16g^{-1}\).
The moment bound for \(g\)
proves the displayed estimate.
Since \(f_s\ge f\), it also
bounds the individual residual
inverse moments.
\end{proof}

This comparison does not
assert that the new full-floor
constant is numerically better
than V75. The factor 16 can
more than offset the smaller
determinant-rank term.
The useful new conclusions
are restoration of the exact
symmetry, elimination of
\(\mathcal J_{\rm part}\), and
a uniform identification of
the residual small-spectrum
regime with the original one.

\subsection{A finite actual check of the averaging operation}

For the antipodal current pair
of Record~\ref{rec:v77-nonzero},
consider every one of the
16 translated two-site
partitions on the side-four
torus. Their exact high
mixed derivatives at zero
are not all equal. Their
minimum, maximum and average
are respectively
\begin{equation}\label{eq:v78-finite-average}
          -\frac3{200704},\qquad
           \frac1{28224},\qquad
           \frac{23}{3612672}.
\end{equation}
The original aligned partition
has the V77 value \(1/50176\).
Thus logarithmic averaging
restores symmetry but does
not remove the interaction.

These values follow from a
finite formula requiring no
approximation to the gauge
matrix. Let \(G_0(x,y)\)
be the exact high Green
matrix in V77's Walsh sum.
For translated partition \(s\),
\[
        G_s(x,y)=G_0(x-s,y-s),\qquad
        D^2T_s(0)[v,w]=-2\Tr(G_sK_vG_sK_w).
\]
Average the last expression
over \(s\in\{0,1\}^4\).
Translation of both currents
permutes these 16 values.
The certificate evaluates
these rational sums, and
verifies this permutation
for three nontrivial translations.
Two other declared current
pairs from the V77 certificate
have averages
\(63323/90316800\) and
\(-1487/2867200\).
No expectation under the
interacting law is inferred
from these zero-connection
derivative checks.

\subsection{What this changes, and the scale at which it still stops}

The exact block reduction can
now be used without a preferred
partition in the geometric
potential and without its
artificial homogeneous-background
radial penalty. The probability
argument for its fractional
residual determinant is complete,
including inverse moments
up to \(p<2\) and a direct
quadratic tail.

It still has
\(K_{\rm av}=9(M^4-1)+r+\mathcal J\).
The physical entropy term
has not become a coarse
dimension. Further deterministic
use of V77 on larger blocks
also still requires the
admissibility inequality
\(R\ell\le1/3\); translation
averaging does not change it.
There is no proved flow of
the actual retained field
into smaller effective radii,
and no permission to treat
the several residual matrices
as independent samples.

The next missing input must
address the residual spectrum
under the interacting law
or the physical-field radial
cost itself. Iterating exact
determinant identities alone
cannot be reported as that
input. The continuum construction,
compact-sector coverage and
physical Hamiltonian mass
gap remain unproved.
The next scheduled companion
inspection remains V80.

\subsection{Exact checks and preservation}

The new certificate is
\begin{center}\small
\path{scripts/verify_ym_f14_v78_translation_averaged_barrier.py}.
\end{center}
It evaluates the actual
side-four translated-partition
derivatives above. Separately,
declared rational concave
diagonal matrix families
check 16 weighted radial
inequalities, the geometric-mean
concavity formula, and the
remaining-determinant derivative.
Their weights are
\(1/2,1/3,1/6\), and the
tested \(q\) values include
\(3/4\) and \(9/10\), exceeding
every small individual weight
while remaining below one.
Weighted arithmetic--geometric
mean is checked after raising
to the common denominator,
so all arithmetic is rational.

Three other declared positive
matrix examples, each viewed
in three rotated high--low
decompositions, verify the
common determinant and minimum
comparisons. These are not
lattice gauge samples.
The constants for the
energy-normalized high bound
are checked for
\(\ell=2,\ldots,24\).
The all-volume statements
and actual-law probability
bounds have the analytic
proofs above, not a sampling
certificate.

The certificate has 7,926 bytes and SHA--256
\begin{center}\scriptsize\ttfamily
184BBDBF43B92CF5664082E9571619751680109DCCA5D68A20D053DF471613E4.
\end{center}
Its inherited helpers are unchanged.
The predecessor V77 has
1,910,974 bytes and SHA--256
\begin{center}\scriptsize\ttfamily
EA48D551D373A4801D175898F81857B799D71E78130CCB1858562B629808D79E.
\end{center}
Its complete body is preserved
apart from the title version
and this terminal insertion.
V78 is the single active
main research note; V77's
contents remain in it.
No companion, archived note,
monograph or attached source
is modified.

\begin{firewall}
The symmetric decomposition
is exact and its probability
bound uses the actual
interacting measure.
It removes a partition cost,
not the physical-volume cost.
The averaged gauge barrier
is not a physical Hamiltonian
or a continuum construction.
The full mass-gap objective
remains active and unproved.
\end{firewall}
\section{V79: gauge stability of the exact flat sector and a physical-amplitude consequence}

\subsection{Context: go upstream to the zero-action configurations}

V78 removes the partition penalty,
but its Wilson radial bound still
uses a worst-case link-amplitude
cost. Before replacing that cost
by a small-fluctuation expectation,
one must understand which
zero-action configurations can
actually occur in the closure
of the gauge-positive chart.
This note addresses that
nonperturbative prerequisite.

For \(SU(2)\), the logarithmic
gauge-fixing functional is a
sum of squared geodesic distances
on the unit three-sphere.
Four ambient tangent variations
give an exact trace identity.
It proves that every nontrivial
periodic pure-gauge Coulomb copy
has a normalized gauge eigenvalue
at most \(-1/3\).
Such copies are therefore
uniformly outside the closed
gauge-positive region.
For flat connections with
nontrivial torus holonomies,
a periodic adapted frame gives
a different conclusion:
gauge stability bounds the
total squared link amplitude
by a holonomy-dependent
quantity of order \(M^2\),
not \(M^4\).

The ambient-variation mechanism
has a classical continuum
analogue in harmonic-map
instability; see
\href{https://www.jstage.jst.go.jp/article/tmj1949/38/2/38_2_259/_article/-char/en}
{Y.~Ohnita, \emph{Stability of harmonic maps and standard minimal immersions},
Tohoku Math.\ J.\ 38 (1986), 259--267}.
We do not import a continuum
theorem as a lattice estimate.
The exact logarithmic lattice
functional and all constants
used below are computed here.
The need to match the gauge
functional to the logarithmic
link definition is also
discussed in Section II.B of
\href{https://arxiv.org/pdf/1010.5120}
{Ilgenfritz et al., arXiv:1010.5120}.
Their simulation results are
not used as proof inputs.

The final consequence is for
the actual zero-homogeneous-sector
measure at fixed volume as
\(\beta\) tends to infinity.
It has no uniform rate in
volume. No continuum limit,
physical correlation decay,
or mass gap follows from
that ordered-limit statement.

\subsection{The logarithmic gauge Hessian is the inherited matrix}

Identify \(SU(2)\) with the
unit quaternions \(S^3\subset
\mathbb R^4\), with the colour
basis of the earlier notes.
For a directed link \(e=(x,y)\),
write
\[
             U_e=e^{a_e},\qquad r_e=|a_e|<\pi.
\]
All logarithms in this section
are principal. The restriction
\(r_e<\pi\) avoids the cut
at \(-I\). Define the
gauge-orbit functional
\begin{equation}\label{eq:v79-gauge-energy}
       \Phi_U(g)=\frac12\sum_{e=(x,y)}
                    |\log(g_x^{-1}U_e g_y)|^2 .
\end{equation}
This is a gauge-fixing
functional, not the Wilson action.

\begin{lemma}[Exact first and second variations]
\label{lem:v79-gauge-Hessian}
For a real periodic vertex
colour field \(\omega\), put
\(g_x(t)=e^{t\omega_x}\).
At \(t=0\),
\begin{equation}\label{eq:v79-log-variation}
 \dot a_e
       =\sigma(\operatorname{ad}_{a_e})(\omega_y-\omega_x)
                       +\tfrac12\operatorname{ad}_{a_e}(\omega_x+\omega_y).
\end{equation}
Consequently
\begin{equation}\label{eq:v79-Hessian}
 \Phi_U'(0)=\langle D^*a,\omega\rangle,\qquad
 \Phi_U''(0)=\langle\omega,\mathcal M(a)\omega\rangle,
\end{equation}
where \(\mathcal M(a)\) is exactly
\eqref{eq:v77-local-M}.
At a Coulomb connection this
is its self-adjoint gauge Hessian.
\end{lemma}
\begin{proof}
For \(A=\operatorname{ad}_{a_e}\),
the left-trivialized velocity
of \(e^{-t\omega_x}e^{a_e}e^{t\omega_y}\)
at zero is
\(\omega_y-e^{-A}\omega_x\).
The left-trivialized exponential
differential is \((1-e^{-A})/A\).
Its inverse gives
\[
 \dot a_e=\frac{A}{1-e^{-A}}\omega_y
                         -\frac{A}{e^A-1}\omega_x.
\]
Using
\(A/(1-e^{-A})=\sigma(A)+A/2\)
and
\(A/(e^A-1)=\sigma(A)-A/2\)
proves \eqref{eq:v79-log-variation}.
The formulas extend through
zero eigenvalues by their
analytic limits.

Since \(\sigma(A)a_e=a_e\)
and \(A a_e=0\),
\(\langle a_e,\dot a_e\rangle
=\langle a_e,\omega_y-\omega_x\rangle\).
The same first-variation
formula holds at every nearby
point of the one-parameter
gauge curve, with the same
generators \(\omega_x\).
Differentiate it once more
and insert
\eqref{eq:v79-log-variation}.
This proves both assertions.
\end{proof}

In particular positivity of
the primed \(\mathcal F_M\)
is positivity of this gauge
Hessian after its three
constant colour directions
are removed. It is not
stability of the physical
Hamiltonian or convexity of
the Wilson action.

\subsection{Four sphere variations and an exact edge identity}

Let \(p,z\in S^3\), let
\(k=e^{vE_3}\), and write
\[
        c=p\cdot(kz)=\cos r,\qquad
        d=p\cdot z,\qquad \gamma=\cos v.
\]
The chosen axis \(E_3\)
may be replaced by any
unit imaginary quaternion.
Let \(e_\alpha\),
\(\alpha=0,1,2,3\), be
the standard orthonormal
basis of \(\mathbb R^4\).
Vary \(p\) and \(z\)
along their sphere geodesics
with initial tangent vectors
\[
       V_\alpha(p)=e_\alpha-p_\alpha p,\qquad
       V_\alpha(z)=e_\alpha-z_\alpha z.
\]
For each \(\alpha\), let
\(c_\alpha(t)=p_\alpha(t)\cdot k z_\alpha(t)\);
here \(p_\alpha(t),z_\alpha(t)\)
denote the varied sphere
points, not their components.

\begin{lemma}[The discrete sphere trace formula]
\label{lem:v79-edge}
For \(0\le r<\pi\),
\begin{equation}\label{eq:v79-edge}
 \begin{aligned}
 \sum_{\alpha=0}^3
       \left.\frac{d^2}{dt^2}
                \frac12\arccos(c_\alpha(t))^2\right|_{t=0}
          &=2(1-d)+4\frac r{\sin r}(c-\gamma),\\
 \sum_{\alpha=0}^3
       |\operatorname{Im}(z^{-1}e_\alpha)
               -\operatorname{Im}(p^{-1}e_\alpha)|^2
          &=6(1-d).
 \end{aligned}
\end{equation}
The convention at \(r=0\)
is \(r/\sin r=1\).
For \(k=I\), the first
line reduces to
\begin{equation}\label{eq:v79-untwisted}
                        2(1-\cos r)\left(1-\frac{2r}{\sin r}\right).
\end{equation}
\end{lemma}
\begin{proof}
Left multiplication by \(k\)
is an orthogonal matrix \(K\)
with
\[
 K+K^T=2\gamma I,\qquad
 K^2-2\gamma K+I=0,\qquad \Tr K=4\gamma.
\]
The sphere geodesic accelerations
are \(-|V_\alpha(p)|^2p\)
and \(-|V_\alpha(z)|^2z\).
Writing \(P_p=I-pp^T\)
and \(P_z=I-zz^T\), direct
differentiation gives
\[
 \sum_\alpha(c_\alpha')^2
                              =2(1-c^2)(1-d),\qquad
 \sum_\alpha c_\alpha''
                              =4\gamma+2cd-6c.
\]
For clarity, the first
derivative vector in ambient
coordinates is
\[
                  Kz-cp+K^Tp-cz.
\]
Each of its two indicated
tangent summands has squared
norm \(1-c^2\), and their
inner product is
\(-(1-c^2)d\), by the
quadratic identity for \(K\).
For the second sum,
\(\Tr(P_pKP_z)=2\gamma+cd\)
and
\(\sum_\alpha|V_\alpha(p)|^2
=\sum_\alpha|V_\alpha(z)|^2=3\).
These observations prove
the two displayed sums.

For \(f(c)=\arccos(c)^2/2\),
\[
        f'(c)=-\frac r{\sin r},\qquad
        f''(c)=\frac{1-r\cot r}{1-c^2}.
\]
Substitution gives the first
line of \eqref{eq:v79-edge}.
Continuity proves it at
\(r=0\).

For the second line, the
three imaginary components
of quaternion multiplication
give
\[
 \sum_\alpha|\operatorname{Im}(p^{-1}e_\alpha)|^2=3,\qquad
 \sum_\alpha
  \operatorname{Im}(p^{-1}e_\alpha)\cdot
  \operatorname{Im}(z^{-1}e_\alpha)=3(p\cdot z).
\]
Expand the square. When
\(k=I\), \(d=c\) and
\(\gamma=1\), giving
\eqref{eq:v79-untwisted}.
\end{proof}

\subsection{Periodic vacuum copies have a uniformly negative gauge direction}

A periodic pure-gauge
connection is one of the form
\[
                      U_{xy}=q_x^{-1}q_y
\]
for a periodic \(q_x\in SU(2)\).
Equivalently, it is flat
with all four fundamental
torus holonomies equal to
the identity. The next
theorem only requires link
angles below the logarithmic
cut, not the radius-\(1/4\)
restriction.

\begin{theorem}[A quantitative exclusion of nontrivial vacuum copies]
\label{thm:v79-vacuum}
Suppose \(U\) is a periodic
pure-gauge connection with
\(r_e<\pi\) and
\(D^*a=0\).
If at least one link is
not the identity, then
\begin{equation}\label{eq:v79-negative}
                          \lambda_{\min}\mathcal F_M(a)\le-\frac13.
\end{equation}
Thus a pure-gauge Coulomb
connection with nonnegative
primed gauge matrix must
have \(U_e=I\), and hence
\(a_e=0\), on every link.
\end{theorem}
\begin{proof}
For each ambient basis vector
choose the periodic gauge
generator
\[
                  \omega^\alpha_x
                      =\operatorname{Im}(q_x^{-1}e_\alpha).
\]
The curve \(q_xe^{t\omega^\alpha_x}\)
is exactly the sphere
geodesic with initial
tangent \(V_\alpha(q_x)\).
Lemmas~\ref{lem:v79-gauge-Hessian}
and \ref{lem:v79-edge} give
\[
 \sum_\alpha
       \langle\omega^\alpha,\mathcal M(a)\omega^\alpha\rangle
       =2\sum_e(1-\cos r_e)
                             \left(1-\frac{2r_e}{\sin r_e}\right)
       \le-2\sum_e(1-\cos r_e),
\]
and
\[
                        \sum_\alpha\|D\omega^\alpha\|^2
                                    =6\sum_e(1-\cos r_e).
\]
The latter sum is positive
by the nontrivial-link
assumption. At a Coulomb
connection the three
constant colour fields lie
in the kernel of \(\mathcal M\).
Subtract the mean of each
\(\omega^\alpha\); neither
quadratic form nor gradient
energy changes.
The minimum generalized
Rayleigh quotient relative
to \(\Delta\), which is
\(\lambda_{\min}\mathcal F_M\),
is at most the ratio of
these two sums. This ratio
is at most \(-1/3\).
\end{proof}

The statement concerns the
vacuum orbit in the smooth
logarithmic chart. It is
not a proof that arbitrary
nonflat gauge orbits have
unique Coulomb representatives.

\begin{proposition}[A controlled exclusion neighbourhood]
\label{prop:v79-tube}
Suppose \(a\) is a nontrivial
pure-gauge Coulomb connection
and \(b\) is another Coulomb
connection, with
\(\max_e|a_e|,\max_e|b_e|\le R\le1/4\).
Set
\[
  \rho=\max_e|b_e-a_e|,\qquad
  L_M=\frac4{\sqrt{\lambda_1(\Delta)}}+3R.
\]
Then
\[
                  \lambda_{\min}\mathcal F_M(b)
                              \le-\frac13+L_M\rho.
\]
In particular \(L_M\rho\le1/12\)
implies a minimum at most
\(-1/4\).
\end{proposition}
\begin{proof}
The difference \(b-a\)
is Coulomb. The V76
weighted odd bound gives
\(\|G_1(b-a)\|\le4\rho/\sqrt{\lambda_1}\).
Along the straight segment
between \(a\) and \(b\),
the link radius is at
most \(R\).
The linkwise derivative
estimate in Lemma~\ref{lem:v74-even}
therefore gives
\(\|G_{\rm e}(b)-G_{\rm e}(a)\|\le3R\rho\).
Use the minimum-eigenvalue
perturbation inequality and
\eqref{eq:v79-negative}.
\end{proof}

This neighbourhood has
fine-link radius of order
\(M^{-1}\). No volume-uniform
radius in the unscaled
link supremum norm is claimed.

\subsection{Nontrivial holonomies require a periodic adapted frame}

Now assume all plaquettes
are the identity, but do
not assume trivial torus
holonomies. On the universal
cover, flatness gives
\[
 U_{xy}=q_x^{-1}q_y,\qquad
       q_{x+Me_\mu}=H_\mu q_x.
\]
One constructs \(q_x\)
by parallel transport from
a fixed vertex. Elementary
plaquette relations and
backtracking make this
transport independent of
the path on the covering
cubical complex.
Periodicity then gives
the displayed left twists;
transport around a fundamental
rectangle shows that all
\(H_\mu\) commute.

Commuting \(SU(2)\) matrices
can be placed in one
maximal torus. After a
constant change of frame write
\[
         H_\mu=e^{\theta_\mu E_3},\qquad
             -\pi\le\theta_\mu\le\pi,\qquad
         q_x^0=\exp\left(E_3\sum_\mu\theta_\mu x_\mu/M\right).
\]
Then \(p_x=(q_x^0)^{-1}q_x\)
is periodic and
\begin{equation}\label{eq:v79-flat-frame}
                U_{x,\mu}=p_x^{-1}k_\mu p_{x+e_\mu},
                     \qquad k_\mu=e^{v_\mu E_3},\quad v_\mu=\theta_\mu/M.
\end{equation}
For central holonomies one
may choose either principal
angle at the endpoint.
Only its squared bound
will be used.

\begin{theorem}[Gauge-stable flat fields have subextensive total amplitude]
\label{thm:v79-flat-amplitude}
Suppose \(M\ge2\), the
connection is flat and
Coulomb, its link angles
satisfy \(r_e\le R\le1/4\),
and \(\mathcal F_M(a)\succeq0\).
Then
\begin{equation}\label{eq:v79-flat-amplitude}
                  \sum_e|a_e|^2
                      \le24M^2\sum_{\mu=1}^4\theta_\mu^2
                      \le96\pi^2M^2.
\end{equation}
If the stronger relative
floor \(\mathcal F_M(a)\succeq I/3\)
holds, the first constant
24 can be replaced by 3.
\end{theorem}
\begin{proof}
Use the periodic generators
\(\omega^\alpha_x=\operatorname{Im}(p_x^{-1}e_\alpha)\).
For the link \((x,\mu)\),
let \(c_e=\cos r_e\),
\(d_e=p_x\cdot p_{x+e_\mu}\),
and \(\gamma_\mu=\cos v_\mu\).
The sum of the four gauge
second variations is
\begin{equation}\label{eq:v79-twisted-trace}
       \sum_e\left\{2(1-d_e)
                 +4\frac{r_e}{\sin r_e}(c_e-\gamma_\mu)\right\}.
\end{equation}
Gauge semidefiniteness makes
this sum nonnegative.

An elementary two-plane
Cauchy--Schwarz inequality gives
\[
       d_e\ge\gamma_\mu c_e
                                   -|\sin v_\mu|\sin r_e.
\]
For example resolve \(p_x\)
along the orthonormal pair
\(k_\mu p_{x+e_\mu}\) and
\(E_3k_\mu p_{x+e_\mu}\);
the first component is
\(c_e\), so the absolute
value of the second is
at most \(\sin r_e\).
Since \(0\le c_e\le1\),
\[
        1-d_e\le(1-c_e)+(1-\gamma_\mu)
                                  +|\sin v_\mu|\sin r_e .
\]
Write \(X_e=r_e/\sin r_e\).
For \(0\le r_e\le1/4\),
\[
       1\le X_e\le2,\qquad
       1-\cos r_e\ge r_e^2/3.
\]
Also \(1-\cos v_\mu\le v_\mu^2/2\)
and \(|\sin v_\mu|\le|v_\mu|\).
The link contribution in
\eqref{eq:v79-twisted-trace}
is consequently at most
\[
 \begin{aligned}
 &-(4X_e-2)(1-c_e)
                 +(2+4X_e)(1-\gamma_\mu)
                         +2|\sin v_\mu|\sin r_e\\
 &\hspace{8mm}\le-\tfrac23 r_e^2+5v_\mu^2+2|v_\mu|r_e
                 \le-\tfrac13r_e^2+8v_\mu^2.
 \end{aligned}
\]
For the last step use
\(2|v|r\le r^2/3+3v^2\).
Summing and using stability
proves
\(\sum_e r_e^2\le24\sum_e v_\mu^2
=24M^2\sum_\mu\theta_\mu^2\).

If \(\mathcal F_M\succeq I/3\),
the sum in
\eqref{eq:v79-twisted-trace}
is at least one third of
the summed gradient energies,
namely \(2\sum_e(1-d_e)\).
Cancel these terms to get
\[
                 \sum_e X_e(1-c_e)
                            \le\sum_e X_e(1-\gamma_\mu).
\]
The left side is at least
\(\sum_e r_e^2/3\); the
right side is at most
\(\sum_e v_\mu^2\).
This proves the improved constant.
\end{proof}

The order \(M^2\) cannot
be replaced by a bounded
quantity for all holonomy
sectors: small commuting
homogeneous flat connections
\(a_\mu=c_\mu/M\) already
have squared amplitude
\(M^2\sum_\mu|c_\mu|^2\)
and the V65 positive
gauge floor.
These examples are not in
the zero-homogeneous sector
unless \(c=0\).
The theorem does not
classify every stable flat
representative in that
fixed-mean sector, and does
not silently set its
holonomies to the identity.

\subsection{A consequence for the actual zero-sector measure at fixed volume}

Take \(a_c=0\) and the
actual measure \(\mu_{0,\beta}\)
of V75, with a fixed
\(M,R\), \(0<R\le1/4\).
Write
\[
       Q(a)=\sum_e|a_e|^2,\qquad B_M=96\pi^2M^2,\qquad n=9(M^4-1).
\]
The closed set
\[
 \mathcal K_R=\{a\in\mathcal H_M:
                     \max_e|a_e|\le R,\ \mathcal F_M(a)\succeq0\}
\]
is compact. It contains
the closure of the actual
positive domain. The Wilson
action is nonnegative, and
vanishes precisely when
every plaquette is the
identity.

\begin{proposition}[Fixed-volume weak-coupling suppression of excessive amplitude]
\label{prop:v79-concentration}
For every \(\varepsilon>0\),
either the set
\(\mathcal K_R\cap\{Q\ge B_M+\varepsilon\}\)
is empty, or there exist
\(\delta_{M,R,\varepsilon}>0\)
and \(C_{M,R,\varepsilon}<\infty\)
such that, for \(\beta\ge1\),
\begin{equation}\label{eq:v79-concentration}
       \mu_{0,\beta}\{Q\ge B_M+\varepsilon\}
          \le C_{M,R,\varepsilon}\,
                      \beta^{n/2}e^{-\beta\delta_{M,R,\varepsilon}}.
\end{equation}
In particular
\begin{equation}\label{eq:v79-moment}
       \limsup_{\beta\to\infty}\mathbb E_{\mu_{0,\beta}}Q
                            \le B_M,\qquad
       \limsup_{\beta\to\infty}
              \frac1{4M^4}\mathbb E_{\mu_{0,\beta}}Q
                            \le\frac{24\pi^2}{M^2}.
\end{equation}
All limits in this proposition
hold with \(M,R\) fixed.
\end{proposition}
\begin{proof}
If the indicated compact
set is nonempty, its minimum
Wilson action
\[
 \delta_{M,R,\varepsilon}
       =\min_{\mathcal K_R\cap\{Q\ge B_M+\varepsilon\}} S(a)
\]
is positive. Otherwise a
minimizer with zero action
would be flat, Coulomb
and gauge semidefinite, in
contradiction with
Theorem~\ref{thm:v79-flat-amplitude}.
Using semidefiniteness in
that theorem is important:
it covers the gauge boundary
as well as the open domain.

The prefactor
\(\det\mathcal F_M(a)e^{-\mathcal H(a)}\)
is continuous and bounded
on this compact set, and
the ambient coordinate volume
is finite. Thus the
unnormalized bad integral
is at most a fixed constant
times \(e^{-\beta\delta_{M,R,\varepsilon}}\).
For the denominator, a
sufficiently small fixed
Euclidean ball about zero
lies in the actual positive
domain; there the gauge
determinant and Haar factor
have positive lower bounds.
Smoothness of the Wilson
action, \(S(0)=0\), and
\(DS(0)=0\) give
\(S(a)\le C_M\|a\|^2\)
on that ball.
Integrate over its
\(\beta^{-1/2}\)-scaled copy.
This proves
\(\mathcal Z_{0,\beta}\ge c_{M,R}\beta^{-n/2}\)
for \(\beta\ge1\).
Taking the ratio gives
\eqref{eq:v79-concentration}
with the original normalizer.

Since \(0\le Q\le4R^2M^4\),
the same estimate yields
\[
 \mathbb E Q\le B_M+\varepsilon
      +4R^2M^4 C_{M,R,\varepsilon}
                    \beta^{n/2}e^{-\beta\delta_{M,R,\varepsilon}}.
\]
First send \(\beta\) to
infinity at fixed \(M,R\),
then send \(\varepsilon\) to
zero. This proves both
statements in
\eqref{eq:v79-moment}.
If the compact set is
empty the corresponding
bad probability is zero,
and the same conclusion
is immediate.
\end{proof}

This is not conditioning
the measure on exact flatness
or on identity holonomies,
both of which would change
the question. It estimates
the original finite-volume
measure using its complete
zero-action set.
The constants and action
separation in
\eqref{eq:v79-concentration}
have not been bounded
uniformly in \(M\).
In particular the ordered
limit implied by
\eqref{eq:v79-moment} is
not a joint continuum or
thermodynamic construction.

\subsection{Returning to the radial Wilson cost without replacing the law}

The preceding result concerns
a physical-coordinate moment,
not another determinant
factorization. Here is the
precise point where it can
enter the barrier calculation.
At zero homogeneous background,
use either
\[
 \mathcal T(a)=\Tr\mathcal F_M(a)^{-1},\quad d_{\rm eff}=3(M^4-1),
\]
or, when V78's block
assumptions hold,
\[
 \mathcal T(a)=\frac1{\ell^4}\sum_s\Tr B_s(a)^{-1},
                         \quad d_{\rm eff}=3(N^4-1).
\]

\begin{proposition}[A virial bound retaining the actual amplitude moment]
\label{prop:v79-virial}
For either choice, at
\(a_c=0\) one has
\begin{equation}\label{eq:v79-virial}
 \mathbb E\mathcal T
       +\frac{\beta}{2}\mathbb E\langle a,\Delta a\rangle
       +\frac23\mathbb E Q
   \le n+d_{\rm eff}
                  +\frac{\beta C_0R^2}{3}\mathbb E Q .
\end{equation}
All expectations are under
\(\mu_{0,\beta}\).
\end{proposition}
\begin{proof}
The sharper zero-background
radial estimate in V75 is
\[
 DS(a)[a]\ge\tfrac12\langle a,\Delta a\rangle
                                  -\tfrac{C_0R^2}{3}Q(a).
\]
It follows by integrating
V64's Hessian lower form
along \(ua\), whose radius
is \(uR\).
The determinant radial
inequality has the exact
base floor one and gives
\(\mathcal T-d_{\rm eff}\).
The remaining geometric
potential has radial
derivative at least
\(2Q/3\): this is the
Haar potential for the
first choice and
\(\overline W\) for the
second.
Apply the same convex
truncated radial integration
as in V75 or V78.
The outward boundary flux
is nonnegative, so the
integrated logarithmic
derivative is at most \(n\).
Retaining the kinetic and
Haar terms instead of
dropping them proves
\eqref{eq:v79-virial}.
\end{proof}

At fixed \(M,R\), one may
insert the explicit moment
bound in the proof of
Proposition~\ref{prop:v79-concentration}
into the right side of
\eqref{eq:v79-virial}.
This gives a rigorously
smaller zero-temperature
amplitude scale than the
worst-case \(4R^2M^4\)
when \(M\) is large.
It does not give the
uniform fluctuation estimate
\(\mathbb E Q=O(M^4/\beta)\),
nor a volume-independent
gauge minimum, nor a
physical mass.

\subsection{The new bottleneck and the next companion inspection}

The exact flat sector is
now more constrained than
the previous radial estimate
recorded. Periodic vacuum
copies are quantitatively
unstable for the inherited
logarithmic gauge Hessian.
Stable flat fields in the
other holonomy sectors have
subextensive squared amplitude.
This gives an actual
finite-volume probability
consequence, with its
normalizer included.

The missing estimate is
quantitative control of
near-flat configurations as
the volume changes.
The compactness separation
\(\delta_{M,R,\varepsilon}\)
does not supply that control;
it could become too small
for a useful joint limit.
Nor has the holonomy-dependent
bound been converted into
a classification of all
fixed-mean stable flat fields.
The next step must address
one of these quantitative
questions without substituting
fixed-volume weak-coupling
concentration for the
interacting continuum theory.
The required SM, NS and
shell companion inspection
is due at the next version,
V80.

\subsection{Exact checks and preservation}

The certificate is
\begin{center}\small
\path{scripts/verify_ym_f14_v79_flat_stability.py}.
\end{center}
It checks 128 actual
rational \(SU(2)\) edge
configurations, both with
and without a frame twist.
For each edge, it compares
the sphere second variation
with the exact logarithmic
gauge-matrix quadratic form.
Both are linear in
\(X=r/\sin r\).
The two coefficients are
rational and are compared
separately, so no numerical
evaluation or approximation
of \(r\cot r\) is involved.
The gradient-energy denominator
and the twist
Cauchy--Schwarz inequality
are checked independently.

An actual periodic side-four
example has
\(q_x=E_3^{x_1}\),
\(U_{x,1}=E_3\), and all
other links equal to \(I\).
Its logarithm is Coulomb,
all plaquettes and fundamental
holonomies are exactly \(I\),
and the summed Hessian
coefficients are \(512,-1024\)
with free-energy denominator
\(1536\).
Since \(X=\pi/2\), the
normalized test ratio is
\((1-\pi)/3\).
This example is below the
logarithmic cut but outside
the radius-\(1/4\) chart;
it checks the general
pure-gauge theorem, not
the small-radius amplitude
bound. Another 100 rational
checks verify the scalar
inequalities used in that
amplitude bound.
The general instability,
holonomy and probability
statements have the analytic
proofs above.

The certificate contains 5,900 bytes and has SHA--256
\begin{center}\scriptsize\ttfamily
8BF29703CBB23A5DDCB47CC3177DCE0D85996915586052036D16F981F23294A2.
\end{center}
Its V54 quaternion helper
is unchanged.
The predecessor V78 contains
1,935,556 bytes and has SHA--256
\begin{center}\scriptsize\ttfamily
1CD31F9AB8AF7CB9EFA569893BC42BAF09FAD1536C31057D903C7537CAEA5FE8.
\end{center}
Its body is preserved except
for the title version and
this terminal insertion.
V79 is the single active
main note; all V78 contents
remain in it. Companion
notes, archives, monographs
and attached sources are
not modified.

\begin{firewall}
The new negative eigenvalue
is a gauge-instability statement,
not a negative physical energy.
The stable-flat amplitude bound
and fixed-volume concentration
do not prove a joint continuum
limit or a physical mass gap.
The original full objective
remains active and unproved.
\end{firewall}
\section{V80: orbit comparison, mesoscopic tree gauges, and an actual joint weak-coupling amplitude estimate}
\label{sec:v80}

\subsection{Companion checkpoint and the quantitative question left by V79}

This section was added on 15 September 2026. The scheduled
five-version inspection was performed before choosing the new
argument. The current SM--Einstein V72, NS V26 and Shell V16
files are unchanged. Their relevant terminal sections were
re-read. The three additional Downloads sources remain the
ones audited in V44; no change in those sources is used here.
In particular, the source-selection representation, the
regenerative quadratic phase-space diagnostic, and the
Einstein-bridge compression are not YM transfer operators.

SM V72's many-copy covariance transport keeps its momentum
cutoff and time interval fixed. NS V26 gives particular
rank-one contractions of the three-dimensional Oseen kernel;
its numerical comparisons do not provide a four-dimensional
YM estimate. Shell V16 retains the complete measured
one-loop density and proves relative gauge control on a
small global counting-\(\ell^4\) ball. It does not prove
that the actual fixed-radius interacting law is concentrated
in that ball, or supply a joint nonperturbative remainder.
These scope checks support continuing with the actual law
and its normalizer. They supply none of the estimates below.
The next companion checkpoint is V85.

V79 controlled the exact flat sector and obtained a
fixed-volume concentration result by compactness. Its
positive action separation was not quantified as the
volume changed. We now address that particular gap.
First, the sphere calculation extends to comparison with
\emph{every} periodic gauge frame, not only a flat one.
Second, explicit tree frames on translated smaller boxes
convert this comparison into an action-density bound.
Finally, a direct lower estimate of the original partition
function gives an actual simultaneous volume/coupling limit
for the mean squared bare-link amplitude.

V39 already used axial fillings for a local Hessian at the
identity. We reuse that geometry, not that theorem's
conclusion: the present filling is a nonlinear group-word
identity for an arbitrary connection. The new ingredients
are gauge stability, translation-averaged tile counts, a
free choice of box size, and the complete interacting
normalizer. No tree gauge is used to replace the measure.

\subsection{A stable small-link representative can be compared with any gauge frame}

Work on the unit-spacing four-torus of side \(M\ge2\), with
\(V=M^4\). Write \(U_{xy}=e^{a_{xy}}\),
\(r_e=|a_e|\), and
\[
                 Q(a)=\sum_e r_e^2.
\]
The original logarithm is Coulomb, \(D^*a=0\), and
\(r_e\le R\le1/4\). Harmonic modes need not vanish in the
next two deterministic results. All gauge matrices and
inverse Laplacians on vertices are primed as in V79.

For an arbitrary periodic site frame \(p_x\in\SU(2)\), set
\begin{equation}\label{eq:v80-comparator}
 k_e=p_xU_{xy}p_y^{-1},\qquad
 v_e=\operatorname{dist}_{\SU(2)}(I,k_e)\in[0,\pi],\qquad
 Q_p(U)=\sum_e v_e^2.
\end{equation}
The distance is the round distance on unit quaternions,
so it agrees with the inherited logarithm norm below
the cut. A logarithm of \(k_e\) at the cut is not needed.

\begin{theorem}[A quantitative comparison with the whole gauge orbit]
\label{thm:v80-orbit}
If \(\mathcal F_M(a)\succeq0\), then for every periodic
frame \(p\),
\begin{equation}\label{eq:v80-orbit}
                         Q(a)\le24Q_p(U).
\end{equation}
If \(\mathcal F_M(a)\succeq I/3\), the constant 24 can
be replaced by 3. Consequently these inequalities also
hold with \(Q_p\) replaced by its minimum over the compact
space of periodic site frames. No smallness, flatness,
Coulomb condition or gauge stability is required of the
comparator \(k\).
\end{theorem}
\begin{proof}
Use the four periodic infinitesimal gauge fields
\(\omega_x^\alpha=\operatorname{Im}(p_x^{-1}e_\alpha)\),
where \(e_0,\ldots,e_3\) is the real quaternion basis.
For a fixed edge put
\[
 c=p_x\cdot k_ep_y=\cos r_e,\qquad d=p_x\cdot p_y,
 \qquad\gamma=\cos v_e,\qquad X=r_e/\sin r_e.
\]
The continuous value of \(X\) at zero is one.
The V79 edge calculation applies to left multiplication
by this \(k_e\): its real matrix satisfies
\(K+K^{\mathsf t}=2\gamma I\) and
\(\operatorname{Tr}K=4\gamma\). It therefore gives
\begin{equation}\label{eq:v80-edge}
 \sum_{\alpha=0}^3\delta^2_{\omega^\alpha}
       \left(\tfrac12r_e^2\right)
       =2(1-d)+4X(c-\gamma),\qquad
 \sum_{\alpha=0}^3|D\omega_e^\alpha|^2=6(1-d).
\end{equation}
More explicitly, for the tangent fields
\(e_\alpha-p_\alpha p\) on each endpoint sphere, the
two derivative sums are
\(\sum_\alpha(c'_\alpha)^2=2(1-c^2)(1-d)\) and
\(\sum_\alpha c''_\alpha=4\gamma+2cd-6c\).
Inserting the derivatives of
\(\tfrac12(\arccos c)^2\) gives \eqref{eq:v80-edge}.
This calculation is edgewise, so the rotation axis of
\(k_e\) may vary arbitrarily from one edge to another.
The central values \(k_e=\pm I\) follow by continuity.

Resolving along the two-plane generated by \(k_ep_y\)
and its quaternion rotation gives
\[
 d\ge\gamma c-|\sin v_e|\sin r_e,
 \qquad
 1-d\le(1-c)+(1-\gamma)+|\sin v_e|\sin r_e.
\]
As in V79, \(1\le X\le2\),
\(1-c\ge r_e^2/3\), and
\(1-\gamma\le v_e^2/2\). Thus
\[
 \begin{aligned}
 2(1-d)+4X(c-\gamma)
 &\le-\tfrac23r_e^2+5v_e^2+2v_er_e\\
 &\le-\tfrac13r_e^2+8v_e^2.
 \end{aligned}
\]
The summed left side is the sum of four quadratic forms
of \(\mathcal M_a\). Gauge semidefiniteness makes it
nonnegative. Subtracting constant components of the
\(\omega^\alpha\)'s changes neither those forms nor their
free gradient energies. Summation proves
\eqref{eq:v80-orbit}.

For the stronger floor, the same sum is at least
\(\frac13\sum_\alpha\|D\omega^\alpha\|^2
=2\sum_e(1-d_e)\). Cancellation in
\eqref{eq:v80-edge} leaves
\(\sum_eX_e(1-c_e)\le\sum_eX_e(1-\gamma_e)\).
The two sides bound \(Q/3\) from below and \(Q_p\)
from above, respectively. The compactness assertion
uses continuity of squared group distance, not a
continuous choice of comparator logarithms.
\end{proof}

This is not a uniqueness theorem for the gauge slice.
It bounds the amplitude of any stable representative
already in the specified chart by the orbit minimum.
In particular, no assumption that a local gauge minimum
is an absolute one has been inserted.

\subsection{Translated mesoscopic boxes and a nonlinear curvature bound}

Let \(1\le\ell\le M\) be an integer and
\(N_\ell=\lceil M/\ell\rceil\). Partition each periodic
coordinate into consecutive intervals of length \(\ell\),
except for a possibly shorter final interval. Their
products partition the torus into rectangular boxes.
Translate the whole partition by \(o\in\mathbb Z_M^4\).
Inside each box, use the rooted axial path taking
directions \(1,2,3,4\) in order, and let \(T_x\) be
its ordered link product. Define the site frame by
\(p_x=T_x\) independently in each box. This is a
well-defined periodic frame; no agreement of different
box frames on crossing links is imposed or required.

Write \(U_p\) for an elementary plaquette and put
\[
 \theta_p=\operatorname{dist}(I,U_p),\qquad
 S(a)=\sum_p(1-\cos\theta_p).
\]
There are six positively oriented plaquettes per site.

\begin{lemma}[An averaged tree-frame bound with explicit boundary cost]
\label{lem:v80-tree}
For an arbitrary \(\SU(2)\) link configuration, the
frames just defined obey
\begin{equation}\label{eq:v80-tree}
 \frac1V\sum_o Q_{p^{(o)}}(U)
 \le4\pi^2N_\ell M^3+
          \frac{3(\ell-1)^3}{2\ell}\sum_p\theta_p^2.
\end{equation}
This statement does not require a logarithmic chart,
a gauge condition, or small plaquette angles.
\end{lemma}
\begin{proof}
There are exactly \(N_\ell M^3\) crossing links in
each direction. Bound each of their squared comparator
angles by \(\pi^2\).

Consider an internal \(\mu\)-edge at local box coordinate
\(x\), where the box has side lengths \(L_j\le\ell\).
Its comparator is the tree--edge--reverse-tree product
\(T_xU_{x,\mu}T_{x+e_\mu}^{-1}\).
Commute the final \(\mu\)-step backwards through the
higher-direction portions of the first tree path.
Each interchange of a \(\nu\)-step followed by a
\(\mu\)-step, \(\nu>\mu\), with the opposite ordering
gives a conjugate of an inverse \((\mu,\nu)\) plaquette.
Indeed, if \(A\) is the product along the preceding
path, the ratio of the two full path products is
\(AU_p^{-1}A^{-1}\). The common suffix cancels exactly.
Multiplying these ratios telescopes to the original
tree--edge--reverse-tree product. This is a group-word
identity, not an abelian Stokes approximation.

The filling has
\(L_e=\sum_{\nu>\mu}x_\nu\le3(\ell-1)\) tiles.
For \(\nu>\mu\), its \((\mu,\nu)\) tiles have base
coordinates \(x_j\) for \(j<\nu\), coordinate
\(s\in\{0,\ldots,x_\nu-1\}\) for \(j=\nu\), and
zero for \(j>\nu\), relative to the box root.
Triangle inequality and bi-invariance of group distance
therefore imply
\[
 v_e^2\le L_e\sum_{p\text{ in its filling}}\theta_p^2
       \le3(\ell-1)\sum_{p\text{ in its filling}}\theta_p^2.
\]
An empty filling has \(v_e=0\).

Fix an orientation \(\mu<\nu\). Summed over all
internal edges of one box \(B\), the number of its
\((\mu,\nu)\) tiles, counted with multiplicity, is
\begin{equation}\label{eq:v80-tiles}
 (L_\mu-1)\frac{L_\nu(L_\nu-1)}2
       \prod_{j\ne\mu,\nu}L_j
 =\frac{|B|}{2}\frac{(L_\mu-1)(L_\nu-1)}{L_\mu}
 \le |B|\frac{(\ell-1)^2}{2\ell}.
\end{equation}
Only \(\mu\)-edges supply tiles with this orientation.
The expression follows by summing \(x_\nu\), allowing
\(L_\mu-1\) internal choices for \(x_\mu\), and
allowing every value of the other coordinates.

The boxes have total volume \(V\). Averaging over all
translations makes the tile multiplicity the same for
every base point at the fixed orientation. There are
\(V\) such plaquettes, so their average multiplicity
is at most \((\ell-1)^2/(2\ell)\). Multiply by
\(3(\ell-1)\), sum orientations, and add the crossing
cost. This proves \eqref{eq:v80-tree}, including the
case of a shorter final interval. For \(\ell=1\)
all links are crossing links and the tile term is zero.
\end{proof}

\begin{theorem}[Quantitative near-flatness in a gauge-stable chart]
\label{thm:v80-action}
Under the hypotheses of Theorem~\ref{thm:v80-orbit}
with \(\mathcal F_M(a)\succeq0\), for every integer
\(1\le\ell\le M\),
\begin{equation}\label{eq:v80-action}
 \frac{Q(a)}V
 \le A\frac{\lceil M/\ell\rceil}{M}
             +B_\ell\frac{S(a)}V,
 \qquad
 A=96\pi^2,\quad
 B_\ell=18\pi^2\frac{(\ell-1)^3}{\ell}.
\end{equation}
In particular
\begin{equation}\label{eq:v80-action-simple}
 \frac{Q(a)}V
 \le96\pi^2\left(\frac1\ell+\frac1M\right)
                         +18\pi^2\ell^2\frac{S(a)}V.
\end{equation}
Both constants can be divided by eight if
\(\mathcal F_M(a)\succeq I/3\).
\end{theorem}
\begin{proof}
Apply \eqref{eq:v80-orbit} to every tree frame and
average. For \(0\le\theta\le\pi\),
\(\sin(\theta/2)\ge\theta/\pi\) by concavity of sine,
and hence
\(\theta^2\le(\pi^2/2)(1-\cos\theta)\).
Inserting this in \eqref{eq:v80-tree} gives
\eqref{eq:v80-action}. Finally
\(\lceil M/\ell\rceil/M\le1/\ell+1/M\) and
\((\ell-1)^3/\ell\le\ell^2\).
The stronger gauge floor changes the comparison factor
from 24 to 3 and no other step.
\end{proof}

For \(\ell=M\), this gives the fully explicit bound
\[
 Q(a)\le96\pi^2M^3+
                18\pi^2\frac{(M-1)^3}{M}S(a).
\]
Its flat-field term is weaker than V79's sharper
\(96\pi^2M^2\) holonomy estimate. The present advantage
is its control of nonzero action and its adjustable
box size. Neither statement supersedes the other.
The frame is used only for a pointwise comparison;
there is no new gauge Jacobian, orbit selection, or
change of integration variables in the actual law.

\subsection{An isotropic Gaussian used only to lower-bound the original normalizer}

Return now to the zero-homogeneous sector
\(a\in\mathcal H_M\), of real dimension \(n=9(V-1)\).
Let \(Z_{\rm iso}\) be a standard Gaussian in this
Euclidean space. Its ambient covariance is the orthogonal
projection onto mean-zero Coulomb links: its Fourier
symbol is \(\Pi(k)\), with \(\Pi(0)=0\).
This is \emph{not} V60's free Gaussian, whose covariance
is \(\Pi(k)/\lambda(k)\). It is not asserted to be the
law of the interacting field.

\begin{lemma}[A bounded isotropic gauge variance]
\label{lem:v80-isotropic}
Let \(d_{\rm g}=3(V-1)\) and let \(G_1\) be the
linear odd gauge matrix in \eqref{eq:v60-gauge}. Then
\begin{equation}\label{eq:v80-isotropic}
 \left\|\mathbb E G_1(Z_{\rm iso})^2\right\|\le10,
 \qquad
 \mathbb P\{\|G_1(Z_{\rm iso})\|\ge t\}
          \le2d_{\rm g}e^{-t^2/20}.
\end{equation}
Also
\begin{equation}\label{eq:v80-link-tail}
 \mathbb P\{\max_e|(Z_{\rm iso})_e|\ge t\}
                                  \le24V e^{-t^2/6}.
\end{equation}
\end{lemma}
\begin{proof}
The exact V60 Fourier calculation, with \(C\) replaced
by \(\Pi\), gives the scalar colour block
\[
 \nu_{\rm iso}(p)=\frac8V\sum_{q\ne0}
   \frac{w(p,q)^T\Pi(p-q)\overline{w(p,q)}}
             {\lambda(p)\lambda(q)},\qquad p\ne0.
\]
Here \(w\) is exactly \eqref{eq:v60-phi}; the factor
eight is the adjoint colour Casimir for one input
colour. Because \(0\preceq\Pi\preceq I\) and
\(|w|^2\le\lambda(p)\),
\[
                  \nu_{\rm iso}(p)
                     \le\frac8V\sum_{q\ne0}\lambda(q)^{-1}.
\]
For centered lattice representatives \(q=2\pi m/M\),
\(\lambda(q)\ge16|m|^2/M^2\). Put \(L=\lfloor M/2\rfloor\).
The shell \(\|m\|_\infty=j\) has at most
\((2j+1)^4-(2j-1)^4=64j^3+16j\le80j^3\) points.
Boundary conventions for an even \(M\) only remove
points from these full integer shells. Consequently
\[
 \begin{aligned}
 \frac1V\sum_{q\ne0}\lambda(q)^{-1}
 &\le\frac1{16M^2}\sum_{0<\|m\|_\infty\le L}|m|^{-2}\\
 &\le\frac{40L(L+1)}{16M^2}\le\frac54.
 \end{aligned}
\]
This proves the variance bound.

Expand \(Z_{\rm iso}\) in any real orthonormal basis
of \(\mathcal H_M\). The corresponding coefficient
matrices of \(G_1\) are self-adjoint by the Coulomb
condition. The norm tail in Theorem 4.1, equation (4.4), of
\href{https://tropp.caltech.edu/papers/Tro11-User-Friendly-preprint.pdf}
{J.~A. Tropp, \emph{User-friendly tail bounds for sums of
random matrices}, Foundations of Computational Mathematics
12 (2012), 389--434}, gives \eqref{eq:v80-isotropic}.
The hypotheses are the real Gaussian expansion and the
variance bound just established, not independent entries
of the gauge matrix.

Every real colour coordinate of a link has variance at
most one, since the ambient covariance is a projection.
If a three-component link has norm at least \(t\), one
coordinate has absolute value at least \(t/\sqrt3\).
The scalar Gaussian tail and a union over \(12V\)
coordinates give \eqref{eq:v80-link-tail}, without any
independence assumption between links.
\end{proof}

For clarity denote the unchanged positive domain and
geometric density by
\[
 \Omega_R=\{a\in\mathcal H_M:\max_e|a_e|<R,
                                \ \mathcal F_M(a)\succ0\},
 \qquad w(a)=\det\mathcal F_M(a)e^{-\mathcal H(a)}.
\]
The actual law, with Euclidean Lebesgue measure on
\(\mathcal H_M\), is still
\begin{equation}\label{eq:v80-law}
 d\mu_{0,\beta}(a)=Z_\beta^{-1}\mathbf1_{\Omega_R}(a)
                 w(a)e^{-\beta S(a)}\,da,
 \qquad
 Z_\beta=\int_{\Omega_R}w(a)e^{-\beta S(a)}\,da.
\end{equation}
This is \eqref{eq:v75-density} at \(a_c=0\), with a
shorter notation for the same normalizer. Let \(Z_0\)
denote its value at \(\beta=0\), keeping the same domain.

\begin{proposition}[Explicit upper and lower normalizers in the same measure]
\label{prop:v80-normalizer}
For \(0<R\le1/4\), \(M\ge2\), and \(\beta\ge1\), set
\begin{equation}\label{eq:v80-threshold}
 b=12\beta+1,\qquad
 b_* =\max\{3R^{-2}\log(96V),\ 160\log(8d_{\rm g})\}.
\end{equation}
If \(b\ge b_*\), then
\begin{equation}\label{eq:v80-normalizer}
 Z_0\le(\pi e R^2)^{n/2},\qquad
 Z_\beta\ge2^{-d_{\rm g}-1}(\pi/b)^{n/2}.
\end{equation}
\end{proposition}
\begin{proof}
On the whole chart, \(G_{\rm e}(a)\preceq0\) by
\eqref{eq:v65-even}. Also \(\operatorname{Tr}G_1(a)=0\):
the spatial coefficients are scalar in colour and
each adjoint colour matrix has zero trace.
On \(\Omega_R\) the inequality \(\log x\le x-1\)
therefore yields
\[
 \log\det\mathcal F_M(a)
 \le\operatorname{Tr}(\mathcal F_M(a)-I)\le0.
\]
The Haar factor is at most one, so \(0<w\le1\).
The domain is contained in the Euclidean ball of radius
\(\rho=2R\sqrt V\). For any \(t>0\), Gaussian
integration bounds the volume of that ball by
\(e^{t\rho^2}(\pi/t)^{n/2}\).
Taking \(t=n/(2\rho^2)\) and using \(8V/n\le1\)
for \(M\ge2\) proves the first bound in
\eqref{eq:v80-normalizer}.

For the lower bound first note the global estimate
\begin{equation}\label{eq:v80-action-upper}
                            S(a)\le12Q(a).
\end{equation}
The distance of a plaquette from the identity is at
most the sum of its four link angles. Thus
\(1-\cos\theta_p\le\theta_p^2/2
\le2\sum_{e\in\partial p}r_e^2\).
Each link occurs in six plaquettes, counted with
multiplicity, which proves \eqref{eq:v80-action-upper}.
For \(r\le1/4\),
\[
 -2\log(\sin r/r)
 \le-2\log(1-r^2/6)
 \le\frac{r^2}{3(1-r^2/6)}\le r^2.
\]
It follows that \(\beta S(a)+\mathcal H(a)\le bQ(a)\)
on the link ball.

Now put \(Y=Z_{\rm iso}/\sqrt{2b}\) and consider
\[
          \mathcal G=\{\max_e|Y_e|<R,
                                  \ \|G_1(Y)\|\le1/4\}.
\]
Lemma~\ref{lem:v80-isotropic} gives
\[
 \mathbb P(\mathcal G^c)
     \le24V e^{-bR^2/3}+2d_{\rm g}e^{-b/160}\le\tfrac12
\]
under \eqref{eq:v80-threshold}. On this event the even
estimate \(-R^2I/2\preceq G_{\rm e}\preceq0\) implies
\[
     \mathcal F_M(Y)\succeq(1-1/4-1/32)I
                         =\tfrac{23}{32}I\succeq\tfrac12 I.
\]
In particular \(\mathcal G\subset\Omega_R\) and
\(\det\mathcal F_M(Y)\ge2^{-d_{\rm g}}\).
The Gaussian density of \(Y\) is
\((b/\pi)^{n/2}e^{-bQ}\). Integrating the original
positive integrand just over this event proves
\[
 Z_\beta\ge2^{-d_{\rm g}}\int_{\mathcal G}e^{-bQ(a)}\,da
          \ge2^{-d_{\rm g}-1}(\pi/b)^{n/2}.
\]
The comparison Gaussian has only evaluated a subset
integral. It has not replaced \(\mu_{0,\beta}\).
\end{proof}

\subsection{An energy bound and amplitude tails under the actual law}

Define the explicit entropy allowance
\begin{equation}\label{eq:v80-allowance}
 \Lambda_{M,R}(\beta)=(d_{\rm g}+1)\log2
       +\frac n2\log\bigl(1+eR^2(12\beta+1)\bigr),
 \qquad \tau_{M,R}(\beta)=\frac{\Lambda_{M,R}(\beta)}{\beta V}.
\end{equation}
The added one inside the logarithm merely enlarges the
upper bound; it is not a new density factor.

\begin{theorem}[Actual action and amplitude estimates with a free box scale]
\label{thm:v80-probability}
Assume the hypotheses and explicit threshold of
Proposition~\ref{prop:v80-normalizer}. Expectations and
probabilities below are under \(\mu_{0,\beta}\).
Then, for every \(u\ge0\),
\begin{equation}\label{eq:v80-energy}
 \frac1V\mathbb E S\le\tau_{M,R}(\beta),\qquad
 \mathbb P\{S/V\ge u\}
            \le\min\{1,e^{\Lambda_{M,R}(\beta)-\beta Vu}\}.
\end{equation}
For every integer \(1\le\ell\le M\),
\begin{equation}\label{eq:v80-moment}
 \frac1V\mathbb E Q
  \le A\frac{\lceil M/\ell\rceil}{M}
                      +B_\ell\tau_{M,R}(\beta).
\end{equation}
For \(2\le\ell\le M\) and \(t\ge0\), one also has
\begin{equation}\label{eq:v80-amplitude-tail}
 \mathbb P\left\{\frac QV\ge
       A\frac{\lceil M/\ell\rceil}{M}
       +B_\ell\frac{\Lambda_{M,R}(\beta)+t}{\beta V}\right\}
                                                   \le e^{-t}.
\end{equation}
Here \(A,B_\ell\) are exactly those of
\eqref{eq:v80-action}; no spectral floor for a typical
input is assumed.
\end{theorem}
\begin{proof}
The normalizer estimates imply
\(\log(Z_0/Z_\beta)\le\Lambda_{M,R}(\beta)\).
Let \(\nu\) be the probability with density
\(Z_0^{-1}\mathbf1_{\Omega_R}w\). Both probabilities
have the same support and the action is bounded at
fixed volume. Their relative entropy is finite and is
\[
 0\le D(\mu_{0,\beta}\Vert\nu)
       =\log(Z_0/Z_\beta)-\beta\mathbb E S.
\]
This proves the mean bound in \eqref{eq:v80-energy}.
For the tail, the numerator restricted to \(S\ge Vu\)
is at most \(e^{-\beta Vu}Z_0\); divide by the same
\(Z_\beta\). The pointwise estimate
\eqref{eq:v80-action}, followed by integration, proves
\eqref{eq:v80-moment}. Since \(B_\ell>0\) for
\(\ell\ge2\), its indicated amplitude event implies
\(S/V\ge(\Lambda+t)/(\beta V)\). Applying the action
tail proves \eqref{eq:v80-amplitude-tail}.
\end{proof}

\begin{proposition}[An explicit simultaneous weak-coupling conclusion]
\label{prop:v80-joint}
If \(0<\tau=\tau_{M,R}(\beta)\le1\), the same law obeys
\begin{equation}\label{eq:v80-optimized}
             \frac1V\mathbb E Q
                  \le168\pi^2\tau^{1/3}+\frac{192\pi^2}{M}.
\end{equation}
Consequently, fix \(0<R\le1/4\) and let \(M\to\infty\)
through integers. For any sequence \(\beta_M\ge1\)
satisfying the explicit threshold
\eqref{eq:v80-threshold} eventually,
\begin{equation}\label{eq:v80-joint}
 \frac1{M^4}\mathbb E_{\mu_{0,\beta_M}}Q\longrightarrow0,
 \qquad
 \frac{Q}{M^4}\longrightarrow0
       \quad\hbox{in probability under }\mu_{0,\beta_M}.
\end{equation}
This is a statement about the actual fixed-radius
restricted law along a joint sequence, not an iterated
fixed-volume limit.
\end{proposition}
\begin{proof}
Choose the deterministic integer
\(\ell=\min\{M,\lceil\tau^{-1/3}\rceil\}\).
If the minimum is the second entry, then
\(1/\ell\le\tau^{1/3}\); if it is \(M\), then
\(1/\ell=1/M\). In either case
\[
 \frac1\ell\le\tau^{1/3}+\frac1M,
 \qquad
 \ell^2\tau\le(\tau^{-1/3}+1)^2\tau\le4\tau^{1/3}.
\]
Use \eqref{eq:v80-action-simple} and the mean action
bound to obtain \eqref{eq:v80-optimized}.

For fixed \(R\),
\[
 \tau_{M,R}(\beta)
 \le\frac{3\log2+\frac92\log(1+eR^2(12\beta+1))}{\beta}.
\]
Here \((d_{\rm g}+1)/V\le3\) and \(n/V\le9\).
The threshold forces \(\beta_M\to\infty\) as
\(M\to\infty\), so this bound tends to zero. Equation
\eqref{eq:v80-optimized} proves convergence of the mean;
Markov's inequality proves the probability assertion.
For example \(\beta_M=M^\varepsilon\), with any fixed
\(\varepsilon>0\), eventually satisfies the threshold.
The sufficient threshold itself only has logarithmic
volume dependence, with the explicit constants in
\eqref{eq:v80-threshold}.
\end{proof}

\subsection{What has moved upstream, and what remains unpaid}

The compactness constant of V79 is no longer needed for
this particular mean-amplitude question. Gauge stability
now gives a pointwise comparison at every action value,
and the original normalizer supplies an energy bound.
The box optimization removes the \(M^2\) loss that
would result from using just one tree on the whole
torus. In the proved regime the squared amplitude per
site is controlled by
\(O((\log\beta/\beta)^{1/3}+M^{-1})\), with constants
explicit above. This is not the sharper fluctuation
scale \(\mathbb E Q=O(V/\beta)\).

There are three important distinctions. First, a small
\emph{average} amplitude is not a bound on the maximum
link, the complete critical counting norm, or the
smallest gauge eigenvalue. For example the elementary
consequence
\[
 \frac1V\mathbb E\#\{e:|a_e|>\varepsilon\}
                   \le\frac{\mathbb E Q}{\varepsilon^2V}
\]
controls a fraction of exceptional links, not their
spatial clusters or their effect on a soft gauge mode.
The nonlocal Coulomb constraint also prevents treating
the deterministic tree boxes as independent random
blocks. No such factorization has been claimed.

Second, inserting \eqref{eq:v80-moment} into V79's
virial inequality still leaves its physical coordinate
dimension and the factor \(\beta R^2\mathbb E Q\).
Thus the present estimate does not yield a
volume-independent barrier inverse moment or close
V74's input-floor requirement. The next useful task
is control of the spatial organization and gauge
effect of the remaining fluctuations, or a stronger
fluctuation estimate; another purely algebraic
factorization would not pay those costs.

Third, \(a_e\) is an unrescaled bare link logarithm.
No physical lattice spacing, field normalization, or
renormalized coupling trajectory has been constructed
by choosing \(\beta_M\) in \eqref{eq:v80-joint}.
The restricted positive chart has not been proved to
represent the full compact gauge theory with the
required continuum axioms. Small bare links can occur
without deciding whether a rescaled continuum field
is nontrivial. Neither that issue nor physical
Hamiltonian spectral support is resolved by this
amplitude estimate. The full YM objective is unchanged.

\subsection{Exact checks and preservation}

The certificate is
\begin{center}\small
\path{scripts/verify_ym_f14_v80_orbit_action.py}.
\end{center}
It checks 96 rational quaternion edge configurations
with arbitrary comparator axes, including eight
central \(-I\) comparators. The original edge is
chosen close to the identity. The exact constant and
\(X=r/\sin r\) coefficients of the sphere and gauge
Hessians agree separately, using V79's unchanged
checker. No transcendental floating value is used
to certify a sign.

For 26 actual noncommuting tree cycles it verifies
all 44 successive plaquette conjugations, the
telescoped group product, and the explicit tile list.
Seven torus/box-size pairs include sides two, three
and four and incomplete final boxes. Their crossing
counts, orientation multiplicities and literal full
translation averages agree with
\eqref{eq:v80-tiles}. These are group and counting
checks, not evaluations of a commuting surrogate.

The isotropic covariance is obtained from the complete
Coulomb projector and checked to be a projection.
The exact gauge variance is evaluated at every
nonzero side-two momentum and five selected side-four
momenta; 1,500 pairwise Coulomb symmetry checks pass.
The one-axis soft variances are \(11/32\) and
\(243/512\), respectively. These are isotropic
variances, not the free variances of V60. The general
bound ten, the Gaussian probability inequality,
normalizer estimates and joint limiting statements
have the analytic proofs above.

The certificate contains 9,922 bytes and has SHA--256
\begin{center}\scriptsize\ttfamily
C0FE2DF95902A71804A3237FAD881E8F1111B7C5157B67363DCC1C7D7B7FA00E.
\end{center}
The predecessor V79 contains 1,959,308 bytes and has SHA--256
\begin{center}\scriptsize\ttfamily
3EB24156548DF879AB2616EC983F04F8DF35B2267C8B3DDB6A1E1E38D76BA3B0.
\end{center}
Its body is preserved except for the title version
and this terminal insertion. V80 is the single
active main note and contains all V79 contents.
The helper certificates, companions, attached sources,
monographs and archives are not modified.

\begin{firewall}
The new result is a quantitative, nonperturbative
amplitude estimate for the stated restricted
finite-volume Wilson law, including a proved joint
weak-coupling sequence. It is not a construction of
the full interacting four-dimensional continuum
Yang--Mills theory and is not a physical mass-gap
proof. Those parts of the original goal remain open.
\end{firewall}
\section{V81: localized gauge stability, a capacity-scale boundary cost, and actual bulk spectral concentration}
\label{sec:v81}

\subsection{The new step and the estimates that are already available}

V80 was a progress turn. Its complete body and certificate
were checked before this continuation. It gives a pointwise
orbit-to-action comparison and a joint weak-coupling estimate
under the actual restricted Wilson law. The comparison uses
sharp box boundaries, paying for every crossing link. That
produces a boundary cost proportional to the boundary area.
The next question is whether the gauge-stability test can
be localized with a smaller cost, without replacing the law
by independent box measures.

It can. Multiplying the four quaternion test fields by a
site cutoff gives an exact identity whose error is the
cutoff's squared gradient. Translated sine cutoffs then
replace the inverse-box-length cost of V80 by an
inverse-square-length cost. The same identity gives an
individual-box estimate in terms of action in an enlarged
box. This is local energetic information, not a cluster
probability estimate.

For the spectral consequence we reuse, rather than
reprove as new work, V59's complete quadratic-density
identity. An explicit constant in its operator bound,
together with the exact even gauge term, controls the
Hilbert--Schmidt distance of the full relative gauge
matrix from the identity. The new amplitude estimate
then yields concentration of its bulk spectrum under
the actual interacting measure.

Use of gauge positivity to bound lattice fields has
precedents; see Appendix D of
\href{https://arxiv.org/pdf/1209.1974}{D.~Zwanziger,
\emph{Some exact properties of the gluon propagator},
Physical Review D 87 (2013), 085039}.
That appendix uses the trace gauge functional and the
anti-Hermitian part of a link. Our gauge functional is
the squared principal logarithm, with the Hessian fixed
in V79. We derive its cutoff identity below and do not
import the other functional's bounds or identify the
two finite-lattice Faddeev--Popov matrices.

The last companion checkpoint is V80 and the next is
V85. No companion, attached manuscript or archive is
modified during this iteration. The full interacting
continuum mass-gap objective remains unchanged.

\subsection{An exact cutoff identity for the four sphere test fields}

Retain the four-torus of side \(M\ge2\), \(V=M^4\),
the Coulomb logarithm \(a\), and \(r_e=|a_e|\le R\le1/4\).
Set \(F=\mathcal F_M(a)\) and \(Q(a)=\sum_e r_e^2\).
The deterministic statements initially allow harmonic
modes. Let \(p_x\in\SU(2)\) be any periodic frame and
\(k_e=p_xU_ep_y^{-1}\) its comparator, with
\(v_e=\operatorname{dist}(I,k_e)\).
At an edge \(e=(x,y)\), put
\[
 c_e=\cos r_e,\quad d_e=p_x\cdot p_y,\quad
 \gamma_e=\cos v_e,\quad X_e=r_e/\sin r_e,
 \quad T_e=1+2r_e\cot r_e.
\]
The values at \(r_e=0\) are obtained continuously.
Let \(\eta\) be a real site function and define
\(\omega_x^\alpha=\operatorname{Im}(p_x^{-1}e_\alpha)\)
and \(\widetilde\omega_x^\alpha=\eta_x\omega_x^\alpha\),
for the four quaternion basis vectors \(e_\alpha\).

\begin{lemma}[Exact localization and a weighted orbit comparison]
\label{lem:v81-cutoff}
The sum of the edge gauge quadratic forms is
\begin{equation}\label{eq:v81-cutoff}
 \sum_{\alpha=0}^3 h_e(\widetilde\omega_x^\alpha,
                                  \widetilde\omega_y^\alpha)
 =\eta_x\eta_y\{2(1-d_e)+4X_e(c_e-\gamma_e)\}
                         +T_e(\eta_y-\eta_x)^2.
\end{equation}
Here \(\sum_e h_e(\omega_x,\omega_y)
=\langle\omega,\mathcal M_a\omega\rangle\), with exactly
the logarithmic gauge Hessian of V79. Moreover
\begin{equation}\label{eq:v81-cutoff-free}
 \sum_{\alpha=0}^3|D\widetilde\omega_e^\alpha|^2
 =6\eta_x\eta_y(1-d_e)+3(\eta_y-\eta_x)^2.
\end{equation}
If \(\eta\ge0\) and \(F\succeq0\), then
\begin{equation}\label{eq:v81-weighted}
 \sum_e\eta_x\eta_y r_e^2
 \le24\sum_e\eta_x\eta_y v_e^2+9\|D\eta\|_2^2.
\end{equation}
For \(F\succeq I/3\), the constants 24 and 9 can
be replaced by 3 and \(3/2\), respectively.
\end{lemma}
\begin{proof}
Write \(\sigma_e=\sigma(\operatorname{ad}_{a_e})\)
for the exact even edge matrix. On real colour vectors,
the edge form from V79 is
\[
 h_e(u,v)=(v-u)\cdot\sigma_e(v-u)
                             +(v-u)\cdot\bigl(a_e\times(u+v)\bigr).
\]
Its two diagonal blocks are both \(\sigma_e\); the
odd term contributes only to the off-diagonal block.
Completeness of the quaternion basis gives
\(\sum_\alpha\omega_x^\alpha(\omega_x^\alpha)^{\mathsf t}=I_3\)
at either endpoint. Since \(\sigma_e\) has eigenvalues
\(1,r_e\cot r_e,r_e\cot r_e\), the summed diagonal
contribution is \(T_e\) at each endpoint.

Denote the uncut summed edge form by \(H_e\).
Multiplication by the endpoint cutoffs changes it to
\[
 T_e(\eta_x^2+\eta_y^2)+\eta_x\eta_y(H_e-2T_e)
                   =\eta_x\eta_yH_e+T_e(\eta_y-\eta_x)^2.
\]
Insert the exact V80 identity
\(H_e=2(1-d_e)+4X_e(c_e-\gamma_e)\).
This proves \eqref{eq:v81-cutoff}; applying the same
calculation to the free form, whose diagonal trace is
three, proves \eqref{eq:v81-cutoff-free}.

For \(r_e\le1/4\), \(T_e\le3\).
V80 proved the edge upper bound
\(H_e\le-r_e^2/3+8v_e^2\).
Multiply it by the nonnegative \(\eta_x\eta_y\), sum,
and use gauge semidefiniteness on the four cutoff
fields. Constant components of those fields can again
be subtracted because \(\mathcal M_a\) annihilates
constant vertices. This gives
\eqref{eq:v81-weighted}.

For the stronger floor subtract one third of
\eqref{eq:v81-cutoff-free} from
\eqref{eq:v81-cutoff}. Positivity implies
\[
 4\sum_e\eta_x\eta_y X_e(1-c_e)
 \le4\sum_e\eta_x\eta_y X_e(1-\gamma_e)
                           +\sum_e(T_e-1)(D\eta_e)^2.
\]
The first sum on the left is at least
\(\frac43\sum_e\eta_x\eta_y r_e^2\), the first sum
on the right is at most \(4\sum_e\eta_x\eta_y v_e^2\),
and \(T_e-1\le2\). This proves the stated constants.
\end{proof}

No derivatives of \(p\) have been estimated in the
cutoff error. Completeness of all four test fields is
what reduces that error to the scalar \(T_e|D\eta_e|^2\).
Dropping some test fields would not justify this formula.

\subsection{Sine cutoffs and the weighted nonlinear filling count}

Take an integer \(3\le\ell\le M\). Define a one-dimensional
function on the periodic coordinate by
\[
 s_\ell(j)=
 \begin{cases}
 \sin(\pi j/\ell),&0\le j\le\ell,\\
 0,&\text{outside this interval},
 \end{cases}
\]
where the two zero endpoints are identified when
\(\ell=M\). Thus it is positive precisely at
\(1,\ldots,\ell-1\). Put
\(\eta^{(o)}_x=\prod_{j=1}^4s_\ell(x_j-o_j)\),
using the indicated translated interval in each
coordinate. Define
\begin{equation}\label{eq:v81-cutoff-norms}
 c_\ell=(\ell/2)^4\cos(\pi/\ell)>0,\qquad
 E_\ell=8(\ell/2)^4\{1-\cos(\pi/\ell)\}.
\end{equation}
The exact scalar identities are
\begin{equation}\label{eq:v81-cutoff-sums}
 \sum_x\eta_x\eta_{x+e_\mu}=c_\ell,\qquad
 \|D\eta\|_2^2=E_\ell,
 \qquad
 \sum_{j=0}^\ell j\,s_\ell(j)^2
                      =\frac\ell2\sum_{j=0}^\ell s_\ell(j)^2.
\end{equation}
Indeed \(\sum s_\ell^2=\ell/2\),
\(\sum s_\ell(j)s_\ell(j+1)=(\ell/2)\cos(\pi/\ell)\),
by the sine recurrence and the finite geometric sum.
Reflection \(j\mapsto\ell-j\) gives the last identity.
Tensor multiplication proves the four-dimensional
formulas, also for \(\ell=M\).

For each origin \(o\), choose the V80 axial tree frame
inside the lifted box \(o+\{0,\ldots,\ell-1\}^4\),
and extend it arbitrarily to the other sites. Every
edge with \(\eta_x^{(o)}\eta_y^{(o)}>0\) is internal
to this box. There are therefore no weighted crossing
links requiring a \(\pi^2\) estimate.

\begin{lemma}[Weighted translated filling count]
\label{lem:v81-fillings}
For these frames and an arbitrary link configuration,
\begin{equation}\label{eq:v81-fillings}
 \sum_o\sum_e\eta_x^{(o)}\eta_y^{(o)}v_{e,o}^2
 \le\frac32\ell(\ell-1)c_\ell\sum_p\theta_p^2.
\end{equation}
The sum over origins is unnormalized. For each fixed
plaquette, its coefficient in the right-hand
majorization of the left side is at most
\(\frac32\ell(\ell-1)c_\ell\).
\end{lemma}
\begin{proof}
Use the exact nonabelian plaquette-word filling of
Lemma~\ref{lem:v80-tree}. Its length for an internal
\(\mu\)-edge at local coordinate \(x\) is
\(L_e=\sum_{\nu>\mu}x_\nu\le3(\ell-1)\), so
\[
 \eta_x\eta_{x+e_\mu}v_e^2
 \le3(\ell-1)\eta_x\eta_{x+e_\mu}
                         \sum_{p\text{ in its filling}}\theta_p^2.
\]
For a fixed orientation \(\mu<\nu\), that edge has
\(x_\nu\) tiles of this orientation. Its total weighted
tile count, summed over edges of one box, is exactly
\[
 \begin{aligned}
 \sum_x\eta_x\eta_{x+e_\mu}x_\nu
 &=\left(\sum_j s_\ell(j)s_\ell(j+1)\right)
   \left(\sum_j s_\ell(j)^2\right)^2
   \left(\sum_j j\,s_\ell(j)^2\right)\\
 &=\frac\ell2 c_\ell.
 \end{aligned}
\]
Weights on the endpoints omitted from an internal
edge sum are zero. Translations spread this weighted
count equally among all plaquette base points: for
a fixed plaquette the sum over all origins equals
the total weighted count at one origin. Multiply
by \(3(\ell-1)\). This proves both assertions.
Keeping the cutoff weights in this count is essential;
replacing them by one would lose an unnecessary
factor from the cutoff normalization.
\end{proof}

\subsection{A global action bound with a capacity-scale cost}

Define the two explicit constants
\begin{equation}\label{eq:v81-constants}
 C_\ell=72\{\sec(\pi/\ell)-1\},\qquad
 D_\ell=18\pi^2\ell(\ell-1).
\end{equation}

\begin{theorem}[Localized stability improves the box-scale comparison]
\label{thm:v81-global}
For a Coulomb field with \(\max_e|a_e|\le R\le1/4\)
and \(F\succeq0\), for every \(3\le\ell\le M\),
\begin{equation}\label{eq:v81-global}
 \frac{Q(a)}V\le C_\ell+D_\ell\frac{S(a)}V
           \le\frac{72\pi^2}{\ell^2}
                            +18\pi^2\ell^2\frac{S(a)}V.
\end{equation}
Under \(F\succeq I/3\), the first term \(C_\ell\)
can be divided by six and the coefficient \(D_\ell\)
by eight.
\end{theorem}
\begin{proof}
Apply \eqref{eq:v81-weighted} to each translated cutoff
and its tree frame. The left side summed over origins
is exactly \(c_\ell Q(a)\), because the edge overlap
is the same in each direction. The error is
\(9VE_\ell\). By Lemma~\ref{lem:v81-fillings}, the
comparator term is at most
\(36\ell(\ell-1)c_\ell\sum_p\theta_p^2\).
Use \(\theta_p^2\le(\pi^2/2)(1-\cos\theta_p)\)
and divide by \(Vc_\ell\). Since
\(9E_\ell/c_\ell=72\{\sec(\pi/\ell)-1\}\), this
proves the first inequality.

For \(\ell\ge3\), \(\cos(\pi/\ell)\ge1/2\) and
\(1-\cos(\pi/\ell)\le\pi^2/(2\ell^2)\), hence
\(\sec(\pi/\ell)-1\le\pi^2/\ell^2\).
Also \(\ell(\ell-1)\le\ell^2\).
The improved-floor constants follow from
Lemma~\ref{lem:v81-cutoff}.
\end{proof}

For a flat field the choice \(\ell=M\) gives
\[
                   Q(a)\le72M^4\{\sec(\pi/M)-1\}
                                =36\pi^2M^2+O(1).
\]
This improves the holonomy-independent large-volume
constant from V79, but does not retain its more
informative dependence on the four individual
holonomy angles. The present new point is control
at nonzero action and a squared inverse-length
boundary cost. No assertion is made for \(M=2\)
using this cutoff, whose edge overlap would vanish.

\subsection{A bound for an individual box, without probabilistic decoupling}

Let \(\mathcal A\) be any set of positively oriented
links. At fixed \(\ell\), let \(\mathcal O_{\mathcal A}\)
be the set of origins whose cutoff edge weight is
nonzero on at least one link of \(\mathcal A\).
Let \(\mathcal P_{\mathcal A}\) be the union of
plaquettes in the tree fillings of all edges with
positive cutoff weight at these origins. These are
geometric sets; they do not depend on the field.
For a plaquette set \(\mathcal P\) write
\(S_{\mathcal P}=\sum_{p\in\mathcal P}(1-\cos\theta_p)\).

\begin{proposition}[A local action estimate with explicit cutoff cost]
\label{prop:v81-local}
Under the same stability hypotheses,
\begin{equation}\label{eq:v81-local}
 \sum_{e\in\mathcal A}r_e^2
       \le D_\ell S_{\mathcal P_{\mathcal A}}
                            +C_\ell|\mathcal O_{\mathcal A}|.
\end{equation}
In particular suppose \(M\ge4\ell\), and let \(B\)
be a translated cube of \(\ell\) vertices in each
direction. Let \(\mathcal A_B\) be its internal links,
and let \(B^+\) be the cube enlarged by \(\ell\)
vertices on each side, of side \(3\ell\). Then
\begin{equation}\label{eq:v81-box}
 \sum_{e\in\mathcal A_B}r_e^2
 \le18\pi^2\ell(\ell-1)S_{B^+}+1152\pi^2\ell^2,
\end{equation}
where \(S_{B^+}\) sums plaquettes contained in \(B^+\).
\end{proposition}
\begin{proof}
Sum the localized inequality only over
\(o\in\mathcal O_{\mathcal A}\). For each target
edge, every nonzero translated weight has been
included; its total is \(c_\ell\). Discard the
other nonnegative terms on the left. Each
plaquette coefficient on the right is bounded
by the full-origin coefficient from
Lemma~\ref{lem:v81-fillings}, since all weights
are nonnegative. Only \(\mathcal P_{\mathcal A}\)
can occur. The total cutoff error is
\(9|\mathcal O_{\mathcal A}|E_\ell\).
The same division and angle-to-action estimate
as above prove \eqref{eq:v81-local}.

For the specified cube, an origin in
\(\mathcal O_{\mathcal A_B}\) lies in a cube with
at most \(2\ell\) coordinates per direction.
Thus \(|\mathcal O_{\mathcal A_B}|\le(2\ell)^4\).
Every corresponding tree box and its fillings
lie in \(B^+\). The condition \(M\ge4\ell\)
allows this geometry to be lifted without
ambiguity across the torus seam. Finally
\(C_\ell\le72\pi^2/\ell^2\), which proves
\eqref{eq:v81-box}.
\end{proof}

In four dimensions the final error has order
\(\ell^2\), the scaling of the Dirichlet energy
of a unit-height cutoff on a box of side \(\ell\).
That is the meaning of ``capacity-scale'' here;
no optimal-capacity theorem is invoked.
The error does not vanish when the local action
vanishes. Nor does this deterministic estimate
make distant boxes independent under the
Coulomb-constrained interacting measure.

\subsection{The improved mean and tail under the unchanged Wilson law}

Now take the actual zero-homogeneous law
\(\mu_{0,\beta}\) of \eqref{eq:v80-law}.
Keep exactly V80's threshold, normalizer and allowance:
\[
 \begin{gathered}
 0<R\le1/4,\quad \beta\ge1,\quad
 12\beta+1\ge\max\{3R^{-2}\log(96V),160\log(8d_{\rm g})\},\\
 d_{\rm g}=3(V-1),\quad n=9(V-1),\\
 \Lambda=(d_{\rm g}+1)\log2+
       \tfrac n2\log(1+eR^2(12\beta+1)),\qquad
 \tau=\Lambda/(\beta V).
 \end{gathered}
\]
The determinant and Haar factor are unchanged.
In particular V80 already proves
\(\mathbb E S/V\le\tau\) and
\(\mathbb P\{S/V\ge u\}\le\min\{1,e^{\Lambda-\beta Vu}\}\).

\begin{theorem}[An actual square-root amplitude rate]
\label{thm:v81-probability}
For \(M\ge3\), every \(3\le\ell\le M\), and
\(t\ge0\), one has
\begin{equation}\label{eq:v81-mean}
       \frac1V\mathbb E Q\le C_\ell+D_\ell\tau,
\end{equation}
and
\begin{equation}\label{eq:v81-tail}
 \mathbb P\left\{Q/V\ge C_\ell+
                     D_\ell\frac{\Lambda+t}{\beta V}\right\}
                                                \le e^{-t}.
\end{equation}
If \(0<\tau\le1/81\), then
\begin{equation}\label{eq:v81-optimized}
             \frac1V\mathbb E Q
                       \le144\pi^2\sqrt\tau+\frac{72\pi^2}{M^2}.
\end{equation}
For fixed \(R\), along every joint sequence allowed
by V80's threshold, this is
\(O(\sqrt{\log\beta/\beta}+M^{-2})\).
\end{theorem}
\begin{proof}
Integrate \eqref{eq:v81-global} and apply V80's
actual mean action bound. Since \(D_\ell>0\),
the event in \eqref{eq:v81-tail} implies
\(S/V\ge(\Lambda+t)/(\beta V)\), proving that tail.

For the optimized estimate choose the deterministic
scale \(\ell=\min\{M,\lceil\tau^{-1/4}\rceil\}\).
The hypotheses make it at least three. One has
\[
      \ell^{-2}\le\sqrt\tau+M^{-2},\qquad
      \ell^2\tau\le(\tau^{-1/4}+1)^2\tau\le4\sqrt\tau.
\]
Insert these in the last inequality of
\eqref{eq:v81-global} and use \(\mathbb E S/V\le\tau\).
The constants are \(72+4\cdot18=144\) and 72.
V80 already showed that at fixed \(R\) its threshold
forces \(\beta\to\infty\) and
\(\tau=O(\log\beta/\beta)\), proving the final statement.
\end{proof}

This improves the exponent and finite-size error of
\eqref{eq:v80-optimized}. The actual-law conclusion,
rather than an estimate under a reference Gaussian,
is inherited through the same proved normalizer
bound. No new partition function is introduced.

\subsection{The full relative gauge matrix is controlled in mean square}

The following explicit version of an inherited
quadratic bound is useful for a different reason:
it converts amplitude into a spectral count.
The Hilbert--Schmidt norm and trace are on the
primed real gauge space.

\begin{lemma}[From amplitude to the complete relative gauge matrix]
\label{lem:v81-HS}
For every real Coulomb field with
\(\max_e|a_e|\le R\le1/4\), including harmonic modes,
\begin{equation}\label{eq:v81-HS}
 \operatorname{Tr}(F-I)^2
 \le(\sqrt{10}+R/\sqrt2)^2Q(a)\le12Q(a).
\end{equation}
This deterministic statement does not require
positivity of \(F\).
\end{lemma}
\begin{proof}
Use the already proved identity in
Theorem~\ref{thm:v59-measure}. With scalar-link
matrices
\(A=D\Delta^{-1}D^*\) and
\(B=S_{\rm end}\Delta^{-1}S_{\rm end}^*\), it states
\[
 \operatorname{Tr}G_1(a)^2
                =2\langle a,((A\circ B)\otimes I_3)a\rangle.
\]
Here \(A\) is an orthogonal projection and \(B\succeq0\).
Positivity of the entrywise product gives
\(0\preceq A\circ B\preceq I\circ B=\operatorname{diag}B\).
For completeness, positivity follows by writing
two positive matrices as sums of rank-one matrices:
their entrywise product is the sum of the rank-one
matrices formed by pairwise entrywise products of
the vectors. Apply this to \((I-A)\circ B\).

Writing \(G=\Delta^{-1}\) on primed scalar vertices,
\[
 B_{ee}=2G(0)+2G(e_\mu)\le4G(0)\le5.
\]
The first inequality is Cauchy--Schwarz for the
positive translation-invariant matrix \(G\);
the last is V80's explicit
\(V^{-1}\sum_{k\ne0}\lambda(k)^{-1}\le5/4\).
Consequently \(\|G_1(a)\|_{\rm HS}^2\le10Q(a)\).
This is an explicit constant in V59's existing
bounded quadratic-density operator, not a new
quadratic-density formula.

For the even term put
\(B_{\rm e}=I-\sigma(\operatorname{ad}_a)\) on colour
links, so \(G_{\rm e}=-U^*B_{\rm e}U\) and
\(U=D\Delta^{-1/2}\) is an isometry.
At an edge \(B_{\rm e}\) has eigenvalues
\(0,1-r_e\cot r_e,1-r_e\cot r_e\).
The inherited even bound gives
\(0\le1-r_e\cot r_e\le r_e^2/2\). Compression by
an isometry cannot increase the Hilbert--Schmidt norm:
\[
 \|G_{\rm e}\|_{\rm HS}^2
 \le\operatorname{Tr}B_{\rm e}^2
 \le\tfrac12\sum_e r_e^4\le\tfrac{R^2}{2}Q(a).
\]
For example the first inequality follows from
\(UU^*\preceq I\) and
\(\operatorname{Tr}(U^*B_{\rm e}U)^2
\le\operatorname{Tr}(U^*B_{\rm e}^2U)
\le\operatorname{Tr}B_{\rm e}^2\).
Triangle inequality for the Hilbert--Schmidt norm
proves \eqref{eq:v81-HS}; the last constant follows
from \(10+\sqrt{20}R+R^2/2<12\) at \(R\le1/4\).
\end{proof}

\begin{theorem}[Actual bulk spectral concentration, with its precise limitation]
\label{thm:v81-spectrum}
Under the actual-law hypotheses above, let
\(\lambda_1(a),\ldots,\lambda_{d_{\rm g}}(a)>0\)
be the eigenvalues of \(F\), and define
\[
 \nu_a=\frac1{d_{\rm g}}\sum_{j=1}^{d_{\rm g}}\delta_{\lambda_j(a)},
 \qquad m_{M,\beta}=V^{-1}\mathbb E Q.
\]
Then
\begin{equation}\label{eq:v81-spectral-square}
 \mathbb E\int(\lambda-1)^2\,d\nu_a(\lambda)
                      \le\frac{12V}{d_{\rm g}}m_{M,\beta}.
\end{equation}
For \(0<\delta<1\), put
\(P_\delta=\mathbf1_{(0,\delta]}(F)\). One has
\begin{equation}\label{eq:v81-count}
 \frac1V\mathbb E\operatorname{rank}P_\delta
                       \le\frac{12}{(1-\delta)^2}m_{M,\beta}.
\end{equation}
For every deterministic vertex set \(B\), extending
\(P_\delta\) by zero on constant vertices,
\begin{equation}\label{eq:v81-local-spectral}
 \mathbb E\operatorname{Tr}(\mathbf1_B P_\delta\mathbf1_B)
             \le\frac{12|B|}{(1-\delta)^2}m_{M,\beta}.
\end{equation}
Along the fixed-\(R\) joint sequences of
Theorem~\ref{thm:v81-probability}, the empirical spectrum
therefore concentrates in mean square at one; in
particular the expected fraction of eigenvalues below
any fixed \(\delta<1\) tends to zero.
\end{theorem}
\begin{proof}
Equation \eqref{eq:v81-spectral-square} is the
spectral theorem applied to \eqref{eq:v81-HS},
followed by expectation. Each eigenvalue in
\((0,\delta]\) contributes at least
\((1-\delta)^2\) to \(\operatorname{Tr}(F-I)^2\),
which proves \eqref{eq:v81-count}.

The actual law and the extended operator are
covariant under torus translations. Thus the
expected colour trace of the diagonal of
\(P_\delta\) is the same at every site, and
\[
 \mathbb E\operatorname{Tr}(\mathbf1_B P_\delta\mathbf1_B)
                  =\frac{|B|}{V}\mathbb E\operatorname{rank}P_\delta.
\]
This proves \eqref{eq:v81-local-spectral}. The new
amplitude estimate makes \(m_{M,\beta}\to0\), and
\(V/d_{\rm g}\) stays bounded. This proves all
limiting assertions without assuming a floor for
the smallest random eigenvalue.
\end{proof}

For \(\delta\le1/2\), the spectral-count coefficient
in \eqref{eq:v81-count} is at most 48. Unlike V76's
worst-case radius count, the bound tends to zero
per site under the proved actual-law sequence.
V76's kinetic and spatial nonconcentration estimates
still describe the exceptional eigenspace; they are
not replaced by this averaged count.

\subsection{The remaining exceptional modes and the physical gap}

The empirical conclusion does not bound
\(\lambda_{\min}(F)\), \(\|F^{-1}\|\), or the
contribution of exceptional modes to
\(-\log\det F\). A vanishing fraction of eigenvalues
can still approach zero arbitrarily fast.
Similarly \eqref{eq:v81-local-spectral} is an
\emph{expected local trace}, not a probability
that a local region contains no dangerous modes,
and not a bound on their long-range correlations.

The local box estimate controls amplitude by nearby
curvature plus a capacity-scale error, but it does
not turn V80's global entropy allowance into a
volume-independent tail for action in one box.
Applying the same global numerator bound to such
a box would still carry the full \(\Lambda\).
The missing spatial estimate must be proved in
the actual constrained law, with its determinant
coupling, rather than inferred from deterministic
tree gauges or from independent reference boxes.

The improved amplitude estimate can be inserted
in V79's virial bound, but its physical dimension
and the adverse term \(\beta R^2\mathbb E Q\)
remain. Thus this iteration does not close the
input-floor condition needed in V74 or justify
a uniform inverse-moment estimate.

Finally, the spectrum here is that of the
\emph{relative gauge Hessian}. Already at \(a=0\)
it is identically one, whereas the unnormalized
gauge Hessian is \(\Delta\), with smallest nonzero
eigenvalue \(4\sin^2(\pi/M)\to0\).
Neither operator is the physical Hamiltonian.
Consequently concentration of \(\nu_a\) at one
is not a physical mass or a substitute for
reflection-positive continuum reconstruction.
The next upstream problem is quantitative control
of the exceptional spatial/spectral part under
the full law, not another interpretation of a
bulk spectral limit as a gap.

\subsection{Exact checks and preservation}

The certificate is
\begin{center}\small
\path{scripts/verify_ym_f14_v81_localized_stability.py}.
\end{center}
It checks 256 localized rational quaternion edge
configurations. For unequal cutoffs, including
one vanishing endpoint, the two exact coefficients
of \(1\) and \(X=r/\sin r\) agree separately with
\eqref{eq:v81-cutoff}; the free identity is checked
independently. The unchanged V54 quaternion
operations and V79 uncut edge checker are reused.

Sine cutoffs of lengths three, four and six are
evaluated exactly in \(\mathbb Q(\sqrt2)\) or
\(\mathbb Q(\sqrt3)\), not with floating trigonometry.
Six torus/cutoff pairs check tensor norms, gradient
costs and weighted tile moments, both when the
cutoff reaches the torus period and when it does
not. Three pairs also enumerate the complete
translation average of the weighted tile bank.
For example \(\ell=3\) gives
\(c_3=81/32\), \(E_3=81/4\), and weighted
orientation count \(243/64\); \(\ell=4\) gives
\(c_4=8\sqrt2\) and \(E_4=128-64\sqrt2\).

For the actual side-two scalar Green matrix, the
certificate constructs the 64-link matrices
\(A=DG D^*\) and \(B=S_{\rm end}G S_{\rm end}^*\).
It checks \(A^2=A\), and an exact rational
semidefinite elimination gives rank 63 for
\(\operatorname{diag}B-A\circ B\).
Here \(B_{ee}=29/96\le4G(0)=103/192\).
This finite check supplements the general
entrywise-product proof. The nonlinear even
bound, all-volume probability statements and
spectral limits are proved analytically above.

The certificate contains 9,546 bytes and has SHA--256
\begin{center}\scriptsize\ttfamily
3DB27A156FB36F9A76E113CD9DC7EA05A6D4D9F0DAA2F7932AD4986A011AC63F.
\end{center}
The predecessor V80 contains 1,986,383 bytes and has SHA--256
\begin{center}\scriptsize\ttfamily
52965DAD8B37E4407CD3EC526A83121E8A5EE46620BA2BBE7473737652A450CE.
\end{center}
Its body is preserved except for the title version
and this terminal insertion. V81 is the single
active main note and retains every V80 section.
All companion files, attached sources, helper
certificates, monographs and archives are unchanged.

\begin{firewall}
The new results are an exact localized gauge
comparison, a sharper actual-law amplitude rate,
and concentration of the bulk relative gauge
spectrum. Exceptional modes, the full interacting
continuum construction, and a positive physical
Hamiltonian mass gap remain unresolved.
The original goal remains active and unproved.
\end{firewall}
\section{V82: exceptional gauge logarithms and a normalized density comparison}
\label{sec:v82}

\subsection{The exceptional-mode question left by V81}

V81 was a progress turn: it supplied a localized gauge
comparison, an improved actual-law amplitude estimate,
and concentration of the bulk relative gauge spectrum.
Its complete file and certificate were verified before
this continuation. A bulk squared-distance estimate
does not itself bound the logarithms of a few very small
eigenvalues. Those logarithms enter the actual gauge
determinant, so their contribution must be paid before
making even a per-site normalization comparison.

We now combine two previously separate inputs. V81
bounds the number of exceptional eigenvalues by the
amplitude. V75 bounds an inverse trace under the
determinant-weighted law. Jensen's inequality on the
exceptional spectral measure converts these inputs into
a logarithmic bound, without replacing all eigenvalues
by the smallest one.

To control a normalization ratio, the same estimates
must hold between its two endpoints. We therefore
introduce an explicitly auxiliary power of the complete
geometric density. It changes the comparison measure,
not the definition of the physical endpoint. All members
keep the identical positive gauge domain. The radial
inverse-trace bound has a factor equal to that power;
its resulting endpoint singularity is only logarithmic
in the estimate we need and is integrable.

The outcome is a proved comparison of normalizers and
forward relative entropies per lattice site in a stated
joint weak-coupling window. It is not removal of the
gauge-domain restriction, equivalence of local observables,
or a continuum reconstruction. The last companion
checkpoint remains V80; the next is V85. No attached
source, companion or archive is changed here.

\subsection{One fixed domain and an auxiliary density power}

Work in the real mean-zero Coulomb link space
\(\mathcal H_M\), with \(M\ge3\), \(V=M^4\),
\(n=9(V-1)\), and primed gauge dimension
\(d_{\rm g}=3(V-1)\). Keep \(0<R\le1/4\) fixed and set
\[
 F(a)=\mathcal F_M(a),\qquad Q(a)=\sum_e|a_e|^2,
 \qquad
 \Omega_R=\{a\in\mathcal H_M:\max_e|a_e|<R,\ F(a)\succ0\}.
\]
Define the full geometric potential by
\begin{equation}\label{eq:v82-geometric}
 \mathcal L(a)=\mathcal H(a)-\log\det F(a),\qquad
 w(a)=e^{-\mathcal L(a)}.
\end{equation}
V80 proves \(0<w\le1\) on this domain, so
\(\mathcal L\ge0\). For \(0\le t\le1\), define
\begin{equation}\label{eq:v82-family}
 \begin{aligned}
 \mathcal Z_t(\beta)
   &=\int_{\Omega_R}e^{-\beta S(a)-t\mathcal L(a)}\,da,\\
 d\nu_t^\beta(a)
   &=\mathcal Z_t(\beta)^{-1}\mathbf1_{\Omega_R}(a)
                      e^{-\beta S(a)-t\mathcal L(a)}\,da.
 \end{aligned}
\end{equation}
Here \(t\) is an interpolation parameter, not physical
time or the bare Wilson coupling. The actual law is
exactly \(\nu_1^\beta=\mu_{0,\beta}\) from V80.
The comparison endpoint \(\nu_0^\beta\) omits the
determinant and Haar density but \emph{still retains}
\(\Omega_R\), the mean-zero constraint, the Coulomb
constraint, and the same Wilson action. It is not
the full compact Wilson measure.

Use the already established threshold
\begin{equation}\label{eq:v82-threshold}
 \beta\ge1,\qquad b=12\beta+1\ge
       \max\{3R^{-2}\log(96V),\ 160\log(8d_{\rm g})\},
\end{equation}
and the V80 allowance
\[
 \Lambda=(d_{\rm g}+1)\log2+
        \frac n2\log(1+eR^2(12\beta+1)),\qquad
 \tau=\frac{\Lambda}{\beta V}.
\]
For \(3\le\ell\le M\), let \(C_\ell,D_\ell\) be
exactly \eqref{eq:v81-constants}, and define
\begin{equation}\label{eq:v82-budgets}
 \begin{gathered}
 m_* =\min\left\{4R^2,\ \min_{3\le\ell\le M}
                                     (C_\ell+D_\ell\tau)\right\},\\
 \alpha_* =\min\{d_{\rm g}/V,\ 48m_*\},\\
 \mathcal J=\tfrac43\beta C_0R^4V,\qquad
 K=n+d_{\rm g}+\mathcal J,\qquad k_*=K/V.
 \end{gathered}
\end{equation}
The constant \(C_0\) is the same Wilson Hessian constant
as in V64--V75. These are upper bounds, not fitted
parameters. In particular \(0<\alpha_*\le d_{\rm g}/V\le k_*\).

\subsection{Uniform amplitude control along the complete interpolation}

\begin{lemma}[The V80 normalizer estimate holds for every density power]
\label{lem:v82-uniform}
Under \eqref{eq:v82-threshold}, for every \(t\in[0,1]\),
\begin{equation}\label{eq:v82-normalizers}
 \mathcal Z_t(0)\le(\pi eR^2)^{n/2},\qquad
 \mathcal Z_t(\beta)\ge2^{-d_{\rm g}-1}(\pi/b)^{n/2}.
\end{equation}
Consequently, writing \(\mathbb E_t\) for expectation
under \(\nu_t^\beta\),
\begin{equation}\label{eq:v82-uniform}
 \frac1V\mathbb E_t S\le\tau,\qquad
 \frac1V\mathbb E_t Q\le m_*,\qquad
 \frac1V\mathbb E_t\operatorname{Tr}(F-I)^2\le12m_*.
\end{equation}
These include the comparison endpoint \(t=0\).
\end{lemma}
\begin{proof}
Since \(0<w^t\le1\), the same Euclidean-ball volume
bound as in V80 proves the upper bound in
\eqref{eq:v82-normalizers}.
Use the very same isotropic comparison Gaussian
\(Y=Z_{\rm iso}/\sqrt{2b}\) and good event from V80.
It has probability at least \(1/2\), lies inside
\(\Omega_R\), and satisfies \(F(Y)\succeq I/2\).
On that event
\[
    (\det F)^t\ge2^{-td_{\rm g}}\ge2^{-d_{\rm g}},
    \qquad \beta S+t\mathcal H\le bQ.
\]
The latter inequality uses \(S\le12Q\),
\(\mathcal H\le Q\), and \(0\le t\le1\).
Integrating on the event proves the lower bound.
The Gaussian is only a way to lower-bound this
integral; it is not any of the laws \(\nu_t^\beta\).

Apply nonnegativity of relative entropy to
\(\nu_t^\beta\) and the same-power law at \(\beta=0\).
It gives
\[
 \beta\mathbb E_t S\le
          \log\{\mathcal Z_t(0)/\mathcal Z_t(\beta)\}\le\Lambda.
\]
The pointwise V81 inequality
\(Q/V\le C_\ell+D_\ell S/V\) holds throughout the
common domain. Take expectations, minimize over
\(\ell\), and include \(Q/V\le4R^2\).
The last assertion of \eqref{eq:v82-uniform} is
the unchanged full-matrix bound
\(\operatorname{Tr}(F-I)^2\le12Q\) from V81.
\end{proof}

\begin{lemma}[The inverse-trace price of a fractional determinant]
\label{lem:v82-inverse}
For \(0<t\le1\), the same law obeys
\begin{equation}\label{eq:v82-inverse}
 t\,\mathbb E_t\operatorname{Tr}F^{-1}
                 \le n+td_{\rm g}+\mathcal J\le K.
\end{equation}
The factor \(t\) cannot be dropped in this argument.
No inverse-trace assertion at \(t=0\) is made.
\end{lemma}
\begin{proof}
This is the radial integration of V75, performed
for the actual density at the specified power.
Matrix concavity, \(F(0)=I\), and the zero-background
radial estimates give
\[
 DF(a)[a]\preceq F(a)-I,\quad
 D\mathcal H(a)[a]\ge0,\quad
 \beta DS(a)[a]\ge-\mathcal J.
\]
Thus the potential
\(U_t=\beta S+t\mathcal H-t\log\det F\) satisfies
\[
 DU_t(a)[a]\ge t\operatorname{Tr}F(a)^{-1}
                                     -td_{\rm g}-\mathcal J.
\]
First integrate on
\(\Omega_{R,\epsilon}=\{a:\max|a_e|<R,\ F(a)\succ\epsilon I\}\),
with \(0<\epsilon<1\). This is a convex domain
containing zero. Along each ray, integration of
the derivative of \(r^n e^{-U_t(r\theta)}\)
has a nonnegative endpoint contribution. Hence
\[
          \int_{\Omega_{R,\epsilon}}DU_t[a]e^{-U_t}\,da
                \le n\int_{\Omega_{R,\epsilon}}e^{-U_t}\,da.
\]
The radius-boundary term is retained with its
nonnegative sign; it is not assumed to vanish.
Insert the preceding lower bound. Let
\(\epsilon\downarrow0\), using monotone convergence
for the nonnegative inverse-trace integral, and
divide by the original \(\mathcal Z_t(\beta)\).
This proves \eqref{eq:v82-inverse} even when
\(t<1\), where no integer determinant multiplicity
is available.
\end{proof}

\subsection{Counting the exceptional modes before taking their logarithms}

For a positive matrix \(F\), write
\[
 N_<(a)=\#\{j:\lambda_j(F)<1/2\},\qquad
 L_<(a)=\sum_{\lambda_j(F)<1/2}-\log\lambda_j(F).
\]
Define the actual expected exceptional density
\(\alpha_t=V^{-1}\mathbb E_tN_<\).
V81's spectral count and \eqref{eq:v82-uniform}
give \(0\le\alpha_t\le\alpha_*\).

\begin{proposition}[The exceptional logarithmic cost is controlled]
\label{prop:v82-logarithms}
For \(0<t\le1\),
\begin{equation}\label{eq:v82-small-log}
 \frac1V\mathbb E_t L_<
 \le\alpha_t\log\left(\frac{k_*}{t\alpha_t}\right)
 \le\alpha_*\log\left(\frac{e k_*}{t\alpha_*}\right).
\end{equation}
The first right side is interpreted as zero when
\(\alpha_t=0\). Moreover
\begin{equation}\label{eq:v82-geometric-mean}
 0\le\frac1V\mathbb E_t\mathcal L
 \le14m_*+\alpha_*\log\left(\frac{e k_*}{t\alpha_*}\right).
\end{equation}
In particular the singularity in this upper bound
as \(t\downarrow0\) is integrable.
\end{proposition}
\begin{proof}
Consider the finite measure on pairs consisting
of a field and an exceptional eigenvalue: its
mass is \(\alpha_t\), since the field measure
is divided by \(V\). Its integral of
\(\lambda^{-1}\) is at most \(k_*/t\) by
Lemma~\ref{lem:v82-inverse}. If \(\alpha_t>0\),
normalize this spectral measure to a probability
and apply concavity of the logarithm. This gives
the first inequality in \eqref{eq:v82-small-log}.
If its mass is zero the logarithmic integral
is zero as well.

The function
\(x\mapsto x\log(e k_*/(tx))\) is increasing
for \(0<x\le\alpha_*\), because its derivative
is \(\log(k_*/(tx))\ge0\).
Here \(\alpha_*\le k_*\) and \(t\le1\).
Enlarging the first logarithm by one and using
\(\alpha_t\le\alpha_*\) proves the second inequality.
No eigenvalue has been replaced by
\(\lambda_{\min}(F)\).

For the geometric potential use
\(h(\lambda)=\lambda-1-\log\lambda\ge0\).
For \(\lambda\ge1/2\),
\begin{equation}\label{eq:v82-scalar}
              h(\lambda)\le(\lambda-1)^2.
\end{equation}
Indeed the difference has derivative
\((2\lambda-1)(\lambda-1)/\lambda\), decreases
up to one, increases thereafter, and vanishes
at one. For \(0<\lambda<1/2\),
\(h(\lambda)\le-\log\lambda\).
Thus
\[
 \mathcal L
 =\mathcal H+(d_{\rm g}-\operatorname{Tr}F)
                         +\sum_jh(\lambda_j)
 \le\mathcal H+(d_{\rm g}-\operatorname{Tr}F)
                         +\operatorname{Tr}(F-I)^2+L_<.
\]
V80 gives \(\mathcal H\le Q\).
Also \(\operatorname{Tr}G_1=0\), and with the
even link matrix \(B_{\rm e}=I-\sigma(\operatorname{ad}_a)\)
from V81,
\[
 0\le d_{\rm g}-\operatorname{Tr}F
     =\operatorname{Tr}(U^*B_{\rm e}U)
     \le\operatorname{Tr}B_{\rm e}
     \le\sum_e r_e^2=Q.
\]
The last step uses its two nonzero eigenvalues
\(1-r_e\cot r_e\le r_e^2/2\).
Together with \(\operatorname{Tr}(F-I)^2\le12Q\)
this proves the pointwise bound
\begin{equation}\label{eq:v82-pointwise}
                         0\le\mathcal L\le14Q+L_<.
\end{equation}
Take expectations and use \eqref{eq:v82-small-log}.
Finally \(\int_0^1\log(1/t)\,dt=1\), proving
the integrability assertion.
\end{proof}

The eigenvalue count and inverse trace are averaged
under the \emph{same} \(\nu_t^\beta\). Substituting
a count from a free Gaussian or an inverse moment
from a different power would not justify this
Jensen estimate. The factor \(1/t\) in
\eqref{eq:v82-inverse} is explicitly included.

\subsection{The normalization ratio, including the endpoint at zero power}

Set
\begin{equation}\label{eq:v82-cost}
 \mathcal B_* =14m_*+
        \alpha_*\left\{\log\left(\frac{e k_*}{\alpha_*}\right)+1\right\}.
\end{equation}
The auxiliary normalizer identity used next is
ordinary thermodynamic integration; for the general
power-density construction see Section 3, equation (5), of
\href{https://www2.stat.duke.edu/~scs/Courses/Stat376/Papers/NormConstants/FrielPettitPath05.pdf}
{N.~Friel and A.~N. Pettitt, \emph{Marginal likelihood
estimation via power posteriors}}.
The application below is analytic, not Monte Carlo
sampling. In particular the possibly singular
zero-power endpoint is justified by the estimates
just proved, rather than assumed regular.

\begin{theorem}[A normalized comparison on the same positive chart]
\label{thm:v82-normalized}
Under \eqref{eq:v82-threshold},
\begin{equation}\label{eq:v82-free-energy}
 0\le\frac1V\log\frac{\mathcal Z_0(\beta)}{\mathcal Z_1(\beta)}
       =\frac1V\int_0^1\mathbb E_t\mathcal L\,dt
       \le\mathcal B_*.
\end{equation}
For the actual law \(\nu_1^\beta=\mu_{0,\beta}\),
\begin{equation}\label{eq:v82-relative-entropy}
 \begin{aligned}
 D(\nu_1^\beta\Vert\nu_0^\beta)
    &=\log\frac{\mathcal Z_0(\beta)}{\mathcal Z_1(\beta)}
                                        -\mathbb E_1\mathcal L,\\
 0&\le V^{-1}D(\nu_1^\beta\Vert\nu_0^\beta)\le\mathcal B_*.
 \end{aligned}
\end{equation}
The direction of this relative entropy is part of
the statement; no reverse-entropy estimate is asserted.
\end{theorem}
\begin{proof}
For every \(t>0\), differentiation under the
integral is justified locally in \(t\):
\(x e^{-t x/2}\) is bounded for \(x\ge0\), the
domain has finite volume, and \(e^{-\beta S}\le1\).
Thus
\[
              \frac d{dt}\log\mathcal Z_t(\beta)
                                      =-\mathbb E_t\mathcal L.
\]
Integrate first from \(\epsilon\) to one, with
\(\epsilon>0\). Since \(0<w^t\le1\) and
\(w^t\to1\) inside \(\Omega_R\), dominated
convergence gives
\(\mathcal Z_\epsilon(\beta)\to\mathcal Z_0(\beta)>0\).
The nonnegative right side has the integrable
majorant in \eqref{eq:v82-geometric-mean}.
Passing to the endpoint yields the identity in
\eqref{eq:v82-free-energy}. Its upper bound is
the integral of that majorant; its lower bound
also follows directly from \(w\le1\).

The exact density ratio is
\[
 \frac{d\nu_1^\beta}{d\nu_0^\beta}
      =e^{-\mathcal L}\frac{\mathcal Z_0(\beta)}{\mathcal Z_1(\beta)}.
\]
Its logarithm is integrable under \(\nu_1^\beta\)
by \eqref{eq:v82-geometric-mean} at \(t=1\).
Taking its expectation proves the identity in
\eqref{eq:v82-relative-entropy}. Relative entropy
is nonnegative, and \(\mathcal L\ge0\), so the
normalizer upper bound implies the entropy upper
bound. No unknown normalization factor is dropped.
\end{proof}

Small \(\mathbb E_1\mathcal L\) alone would not
have proved \eqref{eq:v82-free-energy}. The
intermediate laws and the integrable bound as
\(t\downarrow0\) are essential to controlling
the normalization ratio.

\subsection{Logarithmic spectral convergence under the actual law}

\begin{proposition}[The full logarithmic spectrum is integrable in mean]
\label{prop:v82-log-spectrum}
At the actual endpoint,
\begin{equation}\label{eq:v82-log-spectrum}
 \frac1V\mathbb E_1\operatorname{Tr}|\log F|
 \le14m_*+\alpha_*\log\left(\frac{e k_*}{\alpha_*}\right)
                                                      +12\sqrt{m_*}.
\end{equation}
Here \(\operatorname{Tr}|\log F|=\sum_j|\log\lambda_j(F)|\).
This includes, rather than deletes, arbitrarily
small positive eigenvalues.
\end{proposition}
\begin{proof}
The scalar inequality \(\log\lambda\le\lambda-1\)
gives
\[
 \operatorname{Tr}|\log F|
 =-\log\det F+2\operatorname{Tr}(\log F)_+
 \le\mathcal L+2\operatorname{Tr}(F-I)_+.
\]
Cauchy--Schwarz first over eigenvalues and then
over the field law yields
\[
 \frac2V\mathbb E_1\operatorname{Tr}(F-I)_+
 \le2\sqrt{\frac{d_{\rm g}}V
                  \frac1V\mathbb E_1\operatorname{Tr}(F-I)^2}
 \le12\sqrt{m_*},
\]
since \(d_{\rm g}/V\le3\). Use
\eqref{eq:v82-geometric-mean} at \(t=1\).
\end{proof}

\begin{theorem}[A joint window with vanishing geometric and entropy costs per site]
\label{thm:v82-joint}
Fix \(0<R\le1/4\). Let \(M\to\infty\) and choose
\(\beta_M\) satisfying \eqref{eq:v82-threshold}
eventually and
\begin{equation}\label{eq:v82-window}
                         \frac{\log\beta_M}{M^2}\longrightarrow0.
\end{equation}
Then
\begin{equation}\label{eq:v82-joint}
 \begin{gathered}
 \mathcal B_*\longrightarrow0,\qquad
 V^{-1}\mathbb E_{\mu_{0,\beta_M}}\mathcal L\longrightarrow0,
 \qquad V^{-1}\mathbb E_{\mu_{0,\beta_M}}
                               \operatorname{Tr}|\log F|\longrightarrow0,\\
 V^{-1}\log\frac{\mathcal Z_0(\beta_M)}{\mathcal Z_1(\beta_M)}
       \longrightarrow0,\qquad
 V^{-1}D(\mu_{0,\beta_M}\Vert\nu_0^{\beta_M})\longrightarrow0.
 \end{gathered}
\end{equation}
Every power law \(\beta_M=M^\varepsilon\),
\(\varepsilon>0\), satisfies these conditions
eventually. The additional upper-window condition
\eqref{eq:v82-window} is a sufficient hypothesis
of this proof, not a necessity theorem.
\end{theorem}
\begin{proof}
V80's threshold forces \(\beta_M\to\infty\).
Combining Lemma~\ref{lem:v82-uniform} with
V81's optimized estimate gives
\[
 m_*\le C_R\sqrt{\frac{\log\beta_M}{\beta_M}}
                                      +\frac{72\pi^2}{M^2}
\]
for all sufficiently large \(M\). Thus \(m_*\to0\),
\(\alpha_*\to0\), and
\[
 m_*\log\beta_M
 \le C_R\frac{(\log\beta_M)^{3/2}}{\sqrt{\beta_M}}
                         +72\pi^2\frac{\log\beta_M}{M^2}
                                                    \longrightarrow0.
\]
At fixed \(R,C_0\), \(k_*=O(1+\beta_M)\).
Since \(\alpha_*\le48m_*\), the product
\(\alpha_*\log k_*\) tends to zero.
Also \(\alpha_*\log(1/\alpha_*)\to0\).
These facts prove \(\mathcal B_*\to0\).
All other conclusions follow from
\eqref{eq:v82-geometric-mean},
\eqref{eq:v82-free-energy},
\eqref{eq:v82-relative-entropy}, and
\eqref{eq:v82-log-spectrum}.
\end{proof}

The mean-square empirical spectral convergence
of V81 has therefore been upgraded to a
logarithmic statement in this window. A
vanishing fraction of very small modes no
longer leaves its total expected logarithmic
cost unestimated. The inverse trace itself
is still only bounded by \(Vk_*\) at the
actual endpoint, and \(k_*\) may grow with
\(\beta\). No uniform minimum or inverse
operator norm follows.

\subsection{What this comparison does not permit}

All quantities tending to zero in
\eqref{eq:v82-joint} have been divided by the
number of lattice sites. A vanishing
per-site relative entropy is not a vanishing
total relative entropy, and it does not by
itself transfer local correlations or a
spectral gap. The presently proved upper
bound on the total relative entropy is
\(V\mathcal B_*\), which need not tend to
zero. No total-variation convergence of the
endpoint measures has been deduced.

The endpoint \(\nu_0^\beta\) still has a
global gauge-domain restriction and the
Coulomb constraint. Removing its determinant
and Haar density only in the sense of this
normalization/entropy comparison does not
turn it into independent local factors or
the full compact Yang--Mills measure.
In particular the result is not permission
to use the reference endpoint's unproved
correlation bounds for the actual endpoint.

Nor is the dimensionless per-site free-energy
cost a renormalized physical vacuum-energy
density. No relation between \(M\), physical
lattice spacing, field normalization and an
interacting continuum trajectory has been
constructed by \eqref{eq:v82-window}.
The logarithmic exceptional-mode cost is
now controlled in a concrete regime, but
their spatial correlations and their inverse
spectral effect remain the next unresolved
part of the programme. The physical
Hamiltonian mass-gap objective is unchanged.

\subsection{Exact checks and preservation}

The certificate is
\begin{center}\small
\path{scripts/verify_ym_f14_v82_logarithmic_cost.py}.
\end{center}
It checks ten scalar spectral inequalities
and twelve nonconstant spectral Jensen
comparisons using rigorous rational intervals
for logarithms. Power-of-two range reduction
puts each logarithm in \([1,2]\); there the
\(\operatorname{arctanh}\) series has ratio at
most \(1/9\), and an explicit rational tail
bound encloses the omitted terms. These are
certified interval checks, not floating
sign tests.

Two separate diagnostics check the powers
and the orientation of the entropy identity.
They are not claimed to be YM field models.
For the radial density
\(r^{n-1}(1-r)^t\,dr\), \(0<r<1\), the
exact beta integral gives
\[
 Z_t=\frac{(n-1)!}{\prod_{j=1}^n(t+j)},\qquad
 t\,\mathbb E_t(1-r)^{-1}=n+t,\qquad
 \mathbb E_t[-\log(1-r)]=\sum_{j=1}^n\frac1{t+j}.
\]
Twenty rational power/dimension checks verify
the inverse-trace factor and differentiate
the finite product. Integration from zero
to one telescopes exactly to \(\log(n+1)\),
the logarithm of the endpoint normalizer
ratio. This diagnostic would detect the
incorrect omission of the factor \(t\).

For a finite positive-weight bank with base
weights \((1,2,3)\) and geometric weights
\((1,1/4,1/16)\), the normalizers at powers
zero, one half, and one are exactly
\(6,11/4,27/16\).
The forward relative entropy is
\((107/27)\log2-2\log3\).
Prime-log coefficient comparison verifies
its equality to normalizer cost minus
actual mean geometric loss, including the
normalizer telescope. The actual YM
uniform-amplitude, radial and normalization
proofs are the analytic arguments above,
using the unchanged V75/V80/V81 inputs.

The certificate contains 7,559 bytes and has SHA--256
\begin{center}\scriptsize\ttfamily
3F6E406615EE0ECF1C81CB507BB367104ADC2FB17A7800BC51CB00FF6C375FE2.
\end{center}
The predecessor V81 contains 2,012,157 bytes and has SHA--256
\begin{center}\scriptsize\ttfamily
80A6805252729F558823C553A33E25B07C940E8B7DCCB3490ADCB1531774B1CE.
\end{center}
Its body is preserved except for the title
version and this terminal insertion. V82
is the single active main note and retains
all V81 sections. Previous certificates,
companions, supplied sources, monographs
and archives are not modified.

\begin{firewall}
We have controlled an exceptional logarithmic
gauge cost and an exact normalization/entropy
comparison per site on one fixed positive
chart. This does not construct the full
interacting continuum Yang--Mills theory
or establish a positive physical mass gap.
The full objective remains active and unproved.
\end{firewall}
\section{V83: full fractional gauge inverses and actual spatial kernel bounds}
\label{sec:v83}

\subsection{Moving from a logarithmic cost to an inverse kernel}

V82 was a progress turn. It controlled exceptional
logarithmic eigenvalue costs and a normalization
comparison per site, with all intermediate measures
specified. The remaining issue is not another
interpretation of that entropy bound: it is the
inverse spectral and spatial effect of the small
modes. We now obtain a first, explicitly fractional,
inverse estimate under the actual measure.

The argument has three parts. A min--max comparison
combines the free Laplacian eigenvalue count with
the actual relative gauge count. It does not commute
the two operators. This gives fractional inverse
moments of the full unnormalized gauge Hessian.
Its finite-range coefficients then give a weighted
heat estimate, which, combined with the inverse
moment, yields a spatial bound for a fractional
Green kernel. Finally a spectral truncation proves
convergence of that fractional kernel to its free
lattice counterpart in a specified weak-coupling
sequence.

The new estimates are under the unmodified
\(\mu_{0,\beta}\), not the zero-power reference
of V82 and not a free Gaussian. They concern a
gauge operator with all nonconstant modes retained,
not only V77's high block. Their limitation is
equally explicit: the proved uniform inverse
power is below one. No ordinary inverse bound,
exponential clustering, or physical Hamiltonian
gap is inferred. The next companion checkpoint
remains V85; no companion or supplied source is
changed in this iteration.

\subsection{Operators, probability law, and the available budgets}

Use \(M\ge3\), \(V=M^4\), fixed \(0<R\le1/4\),
and the actual zero-homogeneous law on
\(\Omega_R\) from V80--V82. Assume V82's threshold
\eqref{eq:v82-threshold}. Write \(m_*,\alpha_*,k_*\)
for the explicit budgets in
\eqref{eq:v82-budgets}. Only that threshold is
needed below, not V82's additional logarithmic
window. At the actual endpoint the proved inputs
are
\begin{equation}\label{eq:v83-inputs}
 \begin{gathered}
 V^{-1}\mathbb E Q\le m_*,\qquad
 V^{-1}\mathbb E\operatorname{Tr}F^{-1}\le k_*,\qquad
 V^{-1}\mathbb E N_F(1/2)\le\alpha_*,\\
 \operatorname{Tr}(F-I)^2\le12Q,\qquad
 N_F(s)=\#\{\lambda_j(F)\le s\}.
 \end{gathered}
\end{equation}
Here and throughout this section the expectation
is under \(\mu_{0,\beta}\). All traces in these
inputs are on the primed gauge space of dimension
\(d_{\rm g}=3(V-1)\).

On the full vertex colour space let
\[
 L_a=\mathcal M(a)
       =D^*\sigma(\operatorname{ad}_a)D
                          +\tfrac12D^*\operatorname{ad}_aS_{\rm end}.
\]
Let \(P_0\) be the orthogonal projection onto the
three constant colour fields and \(P'=I-P_0\).
The Coulomb condition makes \(L_a\) self-adjoint
and gives \(L_aP_0=0\). On \(\Omega_R\) it is
positive on \(\operatorname{Ran}P'\), where
\begin{equation}\label{eq:v83-product}
                      L_a=\Delta^{1/2}F(a)\Delta^{1/2}.
\end{equation}
For \(p>0\), define \(\mathcal G_p(a)\) as
\(L_a^{-p}\) on that primed space, extended by
zero on constants. Its \((x,y)\) kernel block
is a real \(3\)-by-\(3\) matrix. In particular,
we do not define it by the generally incorrect
formula \(\Delta^{-p/2}F^{-p}\Delta^{-p/2}\).

\subsection{A noncommuting eigenvalue-count comparison}

For a positive operator on the primed space,
let \(N_A(s)\) count its eigenvalues at most
\(s\), with multiplicity.

\begin{lemma}[Free count and product min--max bound]
\label{lem:v83-count}
For every \(r>0\),
\begin{equation}\label{eq:v83-free-count}
                            N_\Delta(r)\le Vr^2.
\end{equation}
For every \(r,\epsilon>0\) and every input in
\(\Omega_R\),
\begin{equation}\label{eq:v83-product-count}
 N_{L_a}(\epsilon)
                  \le N_\Delta(r)+N_F(\epsilon/r).
\end{equation}
\end{lemma}
\begin{proof}
For centered integer momenta,
\(\lambda(2\pi n/M)\ge16|n|^2/M^2\).
If the nonzero free count is nonempty, then
\(M\sqrt r\ge4\). The same box count as in
V76 consequently gives
\[
 N_\Delta(r)\le3(1+M\sqrt r/2)^4
                  \le\frac{243}{256}Vr^2\le Vr^2.
\]
When the count is empty the assertion is
immediate. The estimate is valid also for
large \(r\); overcounting the finite momentum
box causes no problem.

For the second assertion, take the intersection
of the free spectral subspace \(\Delta>r\)
with the inverse image under \(\Delta^{1/2}\)
of the subspace \(F>\epsilon/r\).
Since \(\Delta^{1/2}\) is invertible on the
primed space, this intersection has codimension
at most \(N_\Delta(r)+N_F(\epsilon/r)\).
Every nonzero vector \(u\) in it satisfies
\[
 \langle u,L_a u\rangle
 =\langle\Delta^{1/2}u,F\Delta^{1/2}u\rangle
 >\frac\epsilon r\langle u,\Delta u\rangle
 >\epsilon\|u\|^2.
\]
The min--max principle proves
\eqref{eq:v83-product-count}. The subspace is
an intersection with a pullback, not a
simultaneous eigenspace of \(F\) and \(\Delta\).
\end{proof}

\begin{theorem}[An actual low-energy count for the unnormalized gauge Hessian]
\label{thm:v83-low-count}
For every \(\epsilon>0\),
\begin{equation}\label{eq:v83-low-count}
 \frac1V\mathbb E N_{L_a}(\epsilon)
 \le\min\{3,\ 4\epsilon^2+\alpha_*,\
                                      2(\epsilon k_*)^{2/3}\}.
\end{equation}
\end{theorem}
\begin{proof}
The dimension gives the first bound. Choose
\(r=2\epsilon\) in
\eqref{eq:v83-product-count} and use the
exceptional count at \(1/2\) to obtain the
second. For every \(s>0\), positivity gives
\(N_F(s)\le s\operatorname{Tr}F^{-1}\).
The general comparison therefore yields
\[
           V^{-1}\mathbb E N_{L_a}(\epsilon)
                                  \le r^2+\epsilon k_*/r.
\]
Minimizing over \(r>0\) gives
\(3\,2^{-2/3}(\epsilon k_*)^{2/3}
\le2(\epsilon k_*)^{2/3}\), proving the result.
\end{proof}

The first scale-sensitive bound has a small
exceptional density; the second has a genuine
power of \(\epsilon\) but a coupling-dependent
coefficient. Both are needed in the next
integration. Neither alone supplies the
uniform statement that follows.

\subsection{Fractional inverse moments under the actual law}

For \(0<q<2/3\), define
\begin{equation}\label{eq:v83-moment-budget}
 \mathfrak C_q
 =3+\frac{4q}{2-q}
   +\frac{2^{3q/2}}{1-3q/2}
                  \alpha_*^{\,1-3q/2}k_*^q.
\end{equation}

\begin{theorem}[Full fractional inverse-trace bound]
\label{thm:v83-moment}
For every \(0<q<2/3\),
\begin{equation}\label{eq:v83-moment}
                   \frac1V\mathbb E\operatorname{Tr}\mathcal G_q(a)
                                               \le\mathfrak C_q.
\end{equation}
In particular, along the explicit sequence
\(\beta_M=M\), for fixed \(R\) and all sufficiently
large \(M\), \(\mathfrak C_{1/4}\) is bounded
independently of \(M\). Its exceptional term
tends to zero with upper order
\begin{equation}\label{eq:v83-quarter-rate}
              O_R\bigl(M^{-1/16}(\log M)^{5/16}+M^{-1}\bigr).
\end{equation}
\end{theorem}
\begin{proof}
Layer-cake integration of the positive eigenvalues
gives, by Tonelli,
\[
 \frac1V\mathbb E\operatorname{Tr}\mathcal G_q
    =q\int_0^\infty\epsilon^{-q-1}
                \frac1V\mathbb E N_{L_a}(\epsilon)\,d\epsilon.
\]
The part above one is at most three. On
\((0,1)\), use
\(\min\{x+\alpha,y\}\le x+\min\{\alpha,y\}\)
with \(x=4\epsilon^2\). The first term
integrates to \(4q/(2-q)\). Enlarge the
remaining integration to \((0,\infty)\) and
use the exact scalar identity
\[
 q\int_0^\infty\epsilon^{-q-1}
          \min\{\alpha,b\epsilon^{2/3}\}\,d\epsilon
     =\frac{\alpha^{1-3q/2}b^{3q/2}}{1-3q/2},
                         \qquad 0<q<2/3.
\]
It follows by splitting at
\(\epsilon=(\alpha/b)^{3/2}\).
Take \(\alpha=\alpha_*\) and \(b=2k_*^{2/3}\).
This proves \eqref{eq:v83-moment} without
assuming integrability beforehand.

For \(\beta=M\), V82's threshold holds
eventually at every fixed \(R>0\).
The explicit V81 estimate gives
\[
 \alpha_*\le C_R\sqrt{\frac{\log M}{M}}+C M^{-2},
                                     \qquad k_*=O_R(M).
\]
At \(q=1/4\), the exceptional factor is
\(\alpha_*^{5/8}k_*^{1/4}\).
For \(0<s<1\), \((x+y)^s\le x^s+y^s\).
Applying this with \(s=5/8\) gives exactly
\eqref{eq:v83-quarter-rate}. The remaining
terms in \(\mathfrak C_{1/4}\) are constants.
\end{proof}

More generally, if \(\beta_M=M^\varepsilon\)
with \(\varepsilon>0\), the exceptional term
tends to zero for
\[
                  0<q<\min\{2/7,\ 2/(\varepsilon+3)\}.
\]
Indeed its two powers are controlled by
\(\beta^{(7q-2)/4}(\log\beta)^{(2-3q)/4}\)
and \(\beta^q/M^{2-3q}\).
This is a sufficient fractional range of the
proved estimates, not a sharpness claim.
The ordinary inverse \(q=1\) is outside it.

\subsection{A weighted heat estimate using the actual local coefficients}

We now use locality of \(L_a\), not of the
relative matrix \(F\). The edge coefficients
in V77 give
\[
 \|L_a\|\le20,\qquad
 \sup_x\sum_{y\ne x}\|(L_a)_{xy}\|\le10,
\]
and the off-diagonal blocks connect only
nearest neighbours. These coefficient bounds
use only \(R\le1/4\). They do not need
V77's later block-divisibility or high-floor
assumptions. On the actual domain \(L_a\succeq0\)
on the full space, with kernel exactly the
constant colour fields.

\begin{lemma}[An off-diagonal heat bound with constants retained]
\label{lem:v83-heat}
Let \(x\ne y\), and let \(d=d_M(x,y)\ge1\)
be their nearest-neighbour torus distance.
For every \(t>0\),
\begin{equation}\label{eq:v83-heat}
 \|(e^{-tL_a}P')_{xy}\|
           \le e^{-d/2}+e^{-d^2/(40t)}+V^{-1}.
\end{equation}
For \(0<q<2/3\), one also has the actual
averaged bound
\begin{equation}\label{eq:v83-heat-long}
 \mathbb E\|(e^{-tL_a}P')_{xy}\|
           \le\left(\frac q e\right)^q\mathfrak C_q t^{-q}.
\end{equation}
Kernel norms are colour-matrix operator norms.
No positivity-preserving or Markov property
of the colour heat kernel is assumed.
\end{lemma}
\begin{proof}
Let \(\phi\) be a real site function with
\(|\phi_x-\phi_y|\le\theta\le1\) on neighbouring
sites, and let \(W=e^\phi\), acting identically
in colour. The Hermitian part of the conjugated
matrix has off-diagonal difference from \(L_a\)
equal to
\[
       \{\cosh(\phi_x-\phi_y)-1\}(L_a)_{xy}.
\]
Its diagonal difference is zero.
For \(|u|\le1\), \(\cosh u-1\le u^2\);
for instance the power series bounds it by
\((6/11)u^2\). The row and column bounds
therefore give
\[
            \operatorname{Re}(WL_aW^{-1})\succeq-10\theta^2 I.
\]
Differentiating the squared norm of a solution
to the conjugated heat equation yields
\[
                 \|W e^{-tL_a}W^{-1}\|\le e^{10\theta^2t}.
\]
Choose
\(\phi_z=\theta\min\{d_M(z,y),d\}\).
Its value is zero at \(y\) and \(\theta d\)
at \(x\), so
\[
                   \|(e^{-tL_a})_{xy}\|
                               \le e^{-\theta d+10\theta^2t}.
\]
Taking \(\theta=1\) when \(t<d/20\), and
\(\theta=d/(20t)\) otherwise, gives respectively
\(e^{-d/2}\) and \(e^{-d^2/(40t)}\).
Finally \(e^{-tL_a}P'=e^{-tL_a}-P_0\), and
\((P_0)_{xy}=I_3/V\). This proves
\eqref{eq:v83-heat}; the constant-mode correction
has not been discarded.

For the second estimate write
\(K_t=e^{-tL_a}P'\succeq0\) and factor it as
\(K_t=K_{t/2}K_{t/2}\).
Block Cauchy--Schwarz gives
\[
                  \|(K_t)_{xy}\|
                       \le\sqrt{\operatorname{tr}(K_t)_{xx}
                                        \operatorname{tr}(K_t)_{yy}}.
\]
The law and operator are translation covariant.
Cauchy--Schwarz over the law therefore bounds
the expectation by
\(V^{-1}\mathbb E\operatorname{Tr}K_t\).
For \(\lambda>0\),
\(e^{-t\lambda}\le(q/e)^q t^{-q}\lambda^{-q}\).
Use \eqref{eq:v83-moment} to finish the proof.
\end{proof}

\subsection{An actual spatial estimate for the full fractional Green kernel}

\begin{theorem}[Fractional inverse kernels decay in space]
\label{thm:v83-spatial}
Let \(0<p<q<2/3\). For \(d=d_M(x,y)\ge1\)
and every \(T>0\),
\begin{equation}\label{eq:v83-spatial-T}
 \begin{aligned}
 \mathbb E\|\mathcal G_p(a;x,y)\|
 \le\frac1{\Gamma(p)}\bigg[
  &\frac{T^p}{p}\{e^{-d/2}+e^{-d^2/(40T)}+V^{-1}\}\\
  &+\frac{(q/e)^q\mathfrak C_q}{q-p}T^{p-q}\bigg].
 \end{aligned}
\end{equation}
Consequently there are constants depending
only on \(p,q\) such that
\begin{equation}\label{eq:v83-spatial}
 \mathbb E\|\mathcal G_p(a;x,y)\|
 \le C_{p,q}(1+\mathfrak C_q)
           \left(\frac{\log(2+d)}{d^2}\right)^{q-p}
       +\frac{C_p}{V}
           \left(\frac{d^2}{\log(2+d)}\right)^p.
\end{equation}
In particular, at fixed \(R\), for \(\beta_M=M\)
and all sufficiently large \(M\),
\begin{equation}\label{eq:v83-eighth}
 \mathbb E\|\mathcal G_{1/8}(a;x,y)\|
       \le C_R\left(\frac{\log(2+d)}{d^2}\right)^{1/8},
                                      \qquad 1\le d\le2M.
\end{equation}
This is an algebraic estimate for a fractional
gauge kernel, not exponential physical clustering.
\end{theorem}
\begin{proof}
For each input in the positive domain, spectral
calculus on the primed space gives
\[
              \mathcal G_p(a)
               =\frac1{\Gamma(p)}\int_0^\infty
                             t^{p-1}e^{-tL_a}P'\,dt.
\]
Use \eqref{eq:v83-heat} up to \(T\) and
\eqref{eq:v83-heat-long} after \(T\).
The first exponential in the short-time bound
is constant in \(t\), and the second is at
most its value at \(T\).
Since \(p<q\), the long-time integral is
finite and equals its displayed upper bound.
These estimates also justify expectation and
kernel integration by absolute integrability.
This proves \eqref{eq:v83-spatial-T}.

Take \(T=d^2/(80\log(2+d))\).
Then \(e^{-d^2/(40T)}=(2+d)^{-2}\), and
the long-time term has the first form in
\eqref{eq:v83-spatial}. The other two
short-time exponential terms are bounded
by the same form after changing a constant:
\(d^{2q}e^{-d/2}/\log(2+d)^q\) is bounded,
and so is
\(d^{2q}(2+d)^{-2}/\log(2+d)^q\), because
\(q<1\). The remaining term is exactly
the finite-volume term in
\eqref{eq:v83-spatial}.

For \(q=1/4\), \(p=1/8\), use the uniform
moment bound along \(\beta=M\).
The last term can be absorbed in the first:
their scale ratio, apart from constants, is
\(d^{2q}/(V\log(2+d)^q)\), bounded for
\(d\le2M\), \(V=M^4\), and \(q<2\).
This proves \eqref{eq:v83-eighth}, with the
constant-mode correction accounted for before
absorption.
\end{proof}

The heat parameter in this proof belongs to
the auxiliary gauge operator. It is not a
calibrated Euclidean physical time. V77's
high-block inverse was exponentially local
because it had a deterministic floor. The
present theorem retains the entire primed
space but obtains only the stated fractional,
averaged algebraic bound.

\subsection{Spectral truncation gives a full fractional-kernel comparison}

Let \(\mathcal G_p^0\) denote \(\Delta^{-p}\)
on primed vertices, extended by zero on
constants. The free count already gives
\begin{equation}\label{eq:v83-free-moment}
 \frac1V\operatorname{Tr}\mathcal G_q^0
                    \le\mathfrak C_q^0:=3+\frac{q}{2-q},
                                      \qquad 0<q<2.
\end{equation}
This follows by the same layer-cake integral,
using \(N_\Delta(s)/V\le\min\{3,s^2\}\).

\begin{theorem}[Actual fractional-kernel comparison with the free torus]
\label{thm:v83-comparison}
For \(0<p<q<2/3\), every \(0<\epsilon\le1\),
and every pair of sites \(x,y\),
\begin{equation}\label{eq:v83-comparison}
 \begin{aligned}
 \mathbb E\|\mathcal G_p(a;x,y)-\mathcal G_p^0(x,y)\|
 \le{}&32\sqrt3\,p\,\epsilon^{-p-1}\sqrt{m_*}\\
      &+(\mathfrak C_q+\mathfrak C_q^0)\epsilon^{q-p}.
 \end{aligned}
\end{equation}
Since \(0<m_*\le4R^2<1\), one may choose
\(\epsilon=m_*^{1/(2(q+1))}\), giving
\begin{equation}\label{eq:v83-comparison-rate}
 \mathbb E\|\mathcal G_p(a;x,y)-\mathcal G_p^0(x,y)\|
 \le(32\sqrt3\,p+\mathfrak C_q+\mathfrak C_q^0)
                         m_*^{(q-p)/(2(q+1))}.
\end{equation}
In particular the right side tends to zero
uniformly in \(x,y\) along \(\beta=M\), for
\(q=1/4\), \(p=1/8\); the amplitude power
then is \(1/20\).
\end{theorem}
\begin{proof}
From \eqref{eq:v83-product} and \(\|\Delta\|\le16\),
the full-space extension obeys
\begin{equation}\label{eq:v83-M-HS}
 \|L_a-\Delta\|_{\rm HS}^2
             \le16^2\|F-I\|_{\rm HS}^2\le3072Q(a).
\end{equation}
Both operators vanish on the same constant
subspace, so there is no additional zero-mode
term in this difference.

Put \(f_\epsilon(s)=\max\{s,\epsilon\}^{-p}\)
for \(s\ge0\). Its scalar Lipschitz constant
is \(p\epsilon^{-p-1}\).
For finite self-adjoint matrices \(A,B\)
and a scalar Lipschitz function \(f\), one has
\[
                  \|f(A)-f(B)\|_{\rm HS}
                              \le\operatorname{Lip}(f)\|A-B\|_{\rm HS}.
\]
To see this, use orthonormal eigenbases
\(u_i,v_j\) of \(A,B\). The squared left
side is
\(\sum_{i,j}|f(\lambda_i)-f(\mu_j)|^2
                  |\langle u_i,v_j\rangle|^2\);
the squared input norm is the same expression
with \(f(\lambda_i)-f(\mu_j)\) replaced by
\(\lambda_i-\mu_j\). This is a
Hilbert--Schmidt inequality, not a claim that
every scalar Lipschitz function is operator
Lipschitz in operator norm.

Apply it to \(L_a,\Delta\) and use
\eqref{eq:v83-M-HS}. For
\(T_\epsilon=f_\epsilon(L_a)-f_\epsilon(\Delta)\),
translation covariance and Cauchy--Schwarz give
\[
 \begin{aligned}
 \mathbb E\|(T_\epsilon)_{xy}\|
 &\le\left(\mathbb E\operatorname{Tr}
                    (\mathbf1_xT_\epsilon^2\mathbf1_x)\right)^{1/2}\\
 &=\left(V^{-1}\mathbb E\operatorname{Tr}T_\epsilon^2\right)^{1/2}
 \le32\sqrt3\,p\,\epsilon^{-p-1}\sqrt{m_*}.
 \end{aligned}
\]
The constant eigenvalue of both truncated
functions is \(\epsilon^{-p}\), and cancels
in \(T_\epsilon\).

The remaining operator
\[
 R_\epsilon(L_a)=\mathcal G_p(a)-f_\epsilon(L_a)
                                               +\epsilon^{-p}P_0
\]
is positive, is zero on constants, and obeys
\(R_\epsilon(L_a)\preceq\epsilon^{q-p}\mathcal G_q(a)\).
This is checked eigenvalue by eigenvalue:
only \(0<\lambda<\epsilon\) contributes and
\(\lambda^{-p}\le\epsilon^{q-p}\lambda^{-q}\).
For a positive block matrix, the norm of a
kernel block is at most the geometric mean
of the two diagonal colour traces.
Translation invariance and
\eqref{eq:v83-moment} therefore bound its
expected kernel norm by
\(\epsilon^{q-p}\mathfrak C_q\).
The free remainder is bounded in the same
way by \(\epsilon^{q-p}\mathfrak C_q^0\).
The constant terms in the two remainders
cancel when taking their difference.
Together these estimates prove
\eqref{eq:v83-comparison}.
The stated choice of \(\epsilon\) makes
both powers equal, giving
\eqref{eq:v83-comparison-rate}.
\end{proof}

For fixed integer separation, the free
fractional kernel also has its usual
infinite-unit-lattice Fourier limit when
\(p<2\). The singular part of its Fourier
sum is uniformly integrable by
\eqref{eq:v83-free-count}; away from zero
ordinary Riemann sums apply. Thus the
actual fractional kernel has that same
unrescaled limit in the stated sequence.
This is not a physically rescaled interacting
continuum limit, and no field-renormalization
factor has been silently set equal to one.

\subsection{What is now controlled, and why the ordinary inverse remains open}

We now have actual, full-space fractional
inverse moments and a spatial estimate for
their kernels. The estimates retain both
noncommutation with \(\Delta\) and the
constant-mode subtraction. They do not
follow merely from V82's small per-site
relative entropy; they use the actual
inverse-trace and small-mode counts directly.

The exponent restriction is substantive.
The low-energy count in
\eqref{eq:v83-low-count} only supplies an
\(\epsilon^{2/3}\) tail at very small
energy, with a coefficient growing in
\(\beta\). Combining it with the shrinking
exceptional density pays a uniform moment
below \(2/7\) in the broad polynomial regime,
not at power one. The heat integration also
requires a strictly higher inverse moment
than the kernel power being estimated.
It cannot be continued to the ordinary
Green operator by declaring the fractional
exponents harmless.

The remaining problem is stronger control
of the very low-energy sector, with its
spatial correlations, at the actual input.
An improved relative count or a genuinely
local inverse estimate could change this
conclusion; a rephrasing of the present
fractional bound cannot. No volume-uniform
minimum of \(F\), ordinary inverse bound,
physical transfer comparison, or positive
Hamiltonian mass gap has been proved.

\subsection{Exact checks and preservation}

The certificate is
\begin{center}\small
\path{scripts/verify_ym_f14_v83_fractional_gauge_green.py}.
\end{center}
It checks 196 instances of the product
eigenvalue-count inequality for strictly
positive noncommuting rational matrices.
Counts use exact symmetric congruence
elimination, including two-by-two pivots
when diagonal pivots vanish; no floating
eigenvalue threshold is used. These are
tests of the general min--max statement,
not surrogate YM configurations.

The free counts are checked on the actual
four-dimensional tori of sides two, three,
four and six, with all nonzero momenta and
the three colour copies. There are 96
threshold checks in total. Actual scalar
Laplacians of sides two, three and four
also verify the weighted Hermitian-part
identity, using weight \((9/8)^{d(x,0)}\).
The exact \(\cosh\) defect on a unit
distance step is \(1/144\), and the maximum
error row sum is \(1/18\). Constant row
sums remain exactly zero before weighting.
The general signed-colour coefficient and
semigroup bounds have the analytic proofs
above.

A separate noncommuting matrix check uses
\(f(s)=\max\{s,1/4\}^{-1/2}\), whose
spectral values are rational in the chosen
example. The Hilbert--Schmidt difference
square is exactly \(265/72\), agreeing
with the independent eigenbasis-overlap
formula. Scalar checks verify the
layer-cake coefficient, the coupling
exponents in \eqref{eq:v83-quarter-rate},
and the comparison exponent \(1/20\).
These exact finite checks supplement,
rather than replace, the all-volume and
actual-law proofs.

The certificate contains 8,953 bytes and has SHA--256
\begin{center}\scriptsize\ttfamily
52FC4249C075233619E8F0AEF147688A8D8BDBF833944F603E70A46E62A1F274.
\end{center}
The predecessor V82 contains 2,033,700 bytes and has SHA--256
\begin{center}\scriptsize\ttfamily
236587F5CD577DC613C673A7D2640753C08166AAA1B78DFE1089803F4B2DC66E.
\end{center}
Its body is preserved except for the
title version and this terminal insertion.
V83 is the single active main note and
retains all V82 contents. Earlier
certificates, companions, supplied sources,
monographs and archives are unchanged.

\begin{firewall}
The new spatial estimate is for a
fractional inverse of the auxiliary
gauge Hessian under the stated restricted
Wilson law. It is not a physical
mass-gap estimate. The interacting
continuum construction and positive
Hamiltonian gap remain unproved;
the original goal remains active.
\end{firewall}
\section{V84: a volume-uniform normalization cost and an amplitude-paid inverse budget}
\label{sec:v84}

\subsection{The upstream step and what is genuinely new}

V83 was a progress turn: it proved full fractional
gauge inverse moments and spatial bounds, retaining
the actual restricted law and the constant modes.
Its ordinary-inverse bottleneck remains real.
Simply increasing the exponent in its last integral
does not resolve it. Here we go upstream to two
costs that entered that integral.

First, the lower normalizer in V80--V83 used an
isotropic Gaussian event of probability at least
one half. Its sufficient condition forced
\(\beta\) to dominate a specified multiple of
\(\log V\). A normalizer per site does not require
such a high-probability event. We construct a
different comparison Gaussian as the codifferential
of independent two-form coordinates. Its useful
event has an explicitly paid exponential-in-volume
probability, and its pushforward has full rank
on the original integration space. The resulting
lower bound works for every \(\beta\ge1\), with
no comparison between \(\beta\) and \(M\).

Second, we insert the resulting amplitude estimate
in the already-proved V79 radial identity.
This replaces the worst-case cost proportional
to \(\beta R^4\) per site by a cost proportional
to \(\beta R^2\mathbb E Q/V\). The combination
enlarges the uniform fractional-moment range
from the sufficient \(q<2/7\) of V83 to
\(q<2/5\), subject to an explicit upper-growth
condition on \(\beta_M\).
In particular logarithmic sequences with any
positive coefficient are now covered.

Neither Gaussian is asserted to be the actual
field law. No two-form variables are added to
that law. The new Gaussian only evaluates a
subset integral with the original Lebesgue
normalization. V61's shifted curl identity,
V79's virial inequality, and V81--V83's subsequent
arguments are reused, not claimed as new.
The scheduled companion checkpoint is V85.

\subsection{A bounded two-form produces a controlled actual gauge matrix}

Use the forward differences, exterior derivative
\(d\), and counting inner products of V61.
The field space \(\mathcal H_M\) is the real
mean-zero Coulomb link space with three colours,
of dimension \(n=9(V-1)\), \(V=M^4\).
The primed vertex gauge dimension is
\(d_{\rm g}=3(V-1)\). Throughout, \(M\ge3\)
and \(0<R\le1/4\) are as in V83.

For an arbitrary real two-form \(B\), not
necessarily the canonical potential, put
\[
 a=d^*B,\qquad
 \rho(B)=\max_{x,\mu<\nu}|B_{\mu\nu}(x)|.
\]
The norm in \(\rho\) is the three-colour
Euclidean norm. There are \(18V\) real
two-form coordinates.

\begin{lemma}[An explicit bounded-potential estimate]
\label{lem:v84-potential}
One has \(a\in\mathcal H_M\), and
\begin{equation}\label{eq:v84-potential}
 \max_e|a_e|\le6\rho(B),\qquad
                    \|G_1(a)\|\le6\rho(B).
\end{equation}
In particular, if \(\rho(B)<R/6\), then
\begin{equation}\label{eq:v84-subset-floor}
 a\in\Omega_R,\qquad
 F(a)\succeq\kappa_R I,\qquad
 \kappa_R:=1-R-\frac{R^2}{2}\ge\frac{23}{32}.
\end{equation}
\end{lemma}
\begin{proof}
The identities \(dD=0\) and telescoping
periodic sums imply \(D^*d^*B=0\) and
zero mean in every link direction and colour.
Each component \((d^*B)_\mu(x)\) is a sum of
three backward differences of two-form
components. Each such difference has norm
at most \(2\rho(B)\). This proves the first
inequality in \eqref{eq:v84-potential}.

For primed vertex fields \(f,g\), set
\(\phi=\Delta^{-1/2}f\),
\(\psi=\Delta^{-1/2}g\).
The proof of \eqref{eq:v61-current} uses
only \(a=d^*B\), not the canonical choice of
the potential. Thus
\[
 \langle f,G_1(a)g\rangle=\langle B,dJ\rangle,
 \qquad
 J_\mu=\psi\times T_\mu\phi-T_\mu\psi\times\phi.
\]
Use the four exact shifted terms in
\eqref{eq:v61-discrete-curl}. Put
\[
 X_\mu=\|D_\mu\psi\|_2,\qquad
 Y_\mu=\|D_\mu\phi\|_2.
\]
Translations preserve these norms, and
the norm of a cross product is bounded
by the product of the two norms.
Consequently
\[
 \begin{aligned}
 |\langle B,dJ\rangle|
 &\le2\rho(B)\sum_{\mu\ne\nu}X_\mu Y_\nu\\
 &\le6\rho(B)
       \left(\sum_\mu X_\mu^2\right)^{1/2}
       \left(\sum_\mu Y_\mu^2\right)^{1/2}\\
 &=6\rho(B)\|g\|_2\|f\|_2.
 \end{aligned}
\]
The middle inequality uses the norm three
of the four-by-four matrix with zeros on
the diagonal and ones off the diagonal.
This proves the second inequality.
It needs neither a BMO commutator theorem
nor an estimate for the canonical potential.

If \(\rho(B)<R/6\), the link radius is
strictly below \(R\). The unchanged even
bound \(-R^2I/2\preceq G_{\rm e}(a)\preceq0\)
and \(F=I+G_1+G_{\rm e}\) give
\eqref{eq:v84-subset-floor}. This is a
deterministic statement about a subset of
the original positive chart.
\end{proof}

\subsection{The full-rank Gaussian pushforward and its determinant}

Let \(\Delta_{\rm sc}\) be the scalar vertex
Laplacian, and write
\[
                D_M=\det{}'\Delta_{\rm sc}.
\]
It is useful to keep its finite-volume
normalization explicit.

\begin{lemma}[Determinant and covariance of the coexact comparison]
\label{lem:v84-gaussian}
Let \(b>0\), and give all \(18V\) real
coordinates of \(B\) independent centered
Gaussian laws of variance \(1/(32b)\).
Then \(A=d^*B\) has a nondegenerate Gaussian
law on \(\mathcal H_M\), with density
\begin{equation}\label{eq:v84-gaussian-density}
 \left(\frac{16b}{\pi}\right)^{n/2}
 D_M^{-9/2}
       \exp\{-16b\langle a,\Delta^{-1}a\rangle\}\,da .
\end{equation}
Moreover
\begin{equation}\label{eq:v84-determinant}
           D_M\ge1,\qquad
 \det_{\mathcal H_M}\Delta=D_M^9,\qquad
              16\langle a,\Delta^{-1}a\rangle\ge Q(a).
\end{equation}
\end{lemma}
\begin{proof}
The link Hodge identity gives
\(d^*d=\Delta\Pi\), where \(\Pi\) is the
orthogonal projection onto mean-zero
Coulomb links when restricted away from
zero momentum. In full space both sides
vanish on constant links and on gradients.
Thus
\[
              \operatorname{Cov}(A)
                  =\frac1{32b}d^*d
                  =\frac1{32b}\Delta\Pi.
\]
On \(\mathcal H_M\) this is
\(\Delta/(32b)\), which is strictly
positive. The pushforward has full rank
on that space despite the kernel of
\(d^*\) in two-form coordinates.

At each nonzero momentum there are three
transverse link directions and three
colours. Hence the determinant on
\(\mathcal H_M\) is the ninth power of
the scalar primed determinant. This also
counts correctly on real space: conjugate
Fourier pairs have the same eigenvalues
and the same total real multiplicity.
The finite-dimensional Gaussian density
formula now proves \eqref{eq:v84-gaussian-density}.

The scalar Laplacian has integer entries,
is positive semidefinite, and has exactly
one zero eigenvalue because the torus
graph is connected. The coefficient of
\(z\) in its integer polynomial
\(\det(zI+\Delta_{\rm sc})\) is the
strictly positive number \(D_M\).
It is therefore a positive integer,
so \(D_M\ge1\). No thermodynamic determinant
limit is used.
Finally \(\|\Delta\|\le16\), so
\(16\Delta^{-1}\succeq I\) on
\(\mathcal H_M\), proving the last assertion.
\end{proof}

The covariance in this lemma is
\(\Delta\Pi/(32b)\), not the free covariance
\(\Pi/\lambda\), and not V80's isotropic
covariance. Its density is written in the
original Euclidean link coordinates.
We will not integrate the physical
measure over the redundant two-form
coordinates, so no extra Jacobian or
two-form gauge quotient is being omitted.

\subsection{A lower normalizer with no volume-dependent coupling threshold}

Keep exactly V82's family on the common
domain \(\Omega_R\):
\[
 \mathcal Z_t(\beta)
    =\int_{\Omega_R}e^{-\beta S(a)}
                   \{\det F(a)e^{-\mathcal H(a)}\}^{t}\,da,
 \qquad 0\le t\le1.
\]
Its normalized law is \(\nu_t^\beta\);
\(\nu_1^\beta=\mu_{0,\beta}\).
The value \(t=0\) still retains the
positive gauge-domain restriction.
Set
\begin{equation}\label{eq:v84-probability}
 \begin{gathered}
 b=12\beta+1,\qquad
 h_{R,b}=R\sqrt{\frac{8b}{27}},\\
 p_{R,b}=\mathbb P\{|N(0,1)|<h_{R,b}\}
      =\sqrt{\frac2\pi}\int_0^{h_{R,b}}e^{-s^2/2}\,ds .
 \end{gathered}
\end{equation}

\begin{theorem}[Same-measure normalizers at every coupling]
\label{thm:v84-normalizer}
For \(M\ge3\), \(0<R\le1/4\), \(\beta\ge1\),
and every \(0\le t\le1\),
\begin{equation}\label{eq:v84-normalizer}
 \begin{aligned}
 \mathcal Z_t(0)&\le(\pi eR^2)^{n/2},\\
 \mathcal Z_t(\beta)&\ge
   \kappa_R^{d_{\rm g}}
   \left(\frac{\pi}{16b}\right)^{n/2}
                D_M^{9/2}p_{R,b}^{18V}\\
 &\ge
   \kappa_R^{d_{\rm g}}
   \left(\frac{\pi}{16b}\right)^{n/2}
                          p_{R,b}^{18V}.
 \end{aligned}
\end{equation}
There is no assumption that \(b\) exceed
any multiple of \(\log V\).
\end{theorem}
\begin{proof}
V80 proved \(0<w=\det F\,e^{-\mathcal H}\le1\)
on \(\Omega_R\). Thus \(w^t\le1\), also
at \(t=0\), and its Euclidean ball-volume
argument gives the stated upper bound.

Give \(B\) the independent Gaussian law
of Lemma~\ref{lem:v84-gaussian}. Consider
the coordinate event
\[
 \mathcal C_R=
 \left\{|B_{\mu\nu}^{c}(x)|<
              \frac{R}{6\sqrt3}
            \text{ for every }x,\mu<\nu,c\right\}.
\]
The coordinates, not the resulting links,
are independent. Their standardized
threshold is exactly \(h_{R,b}\), so
\[
                   \mathbb P(\mathcal C_R)=p_{R,b}^{18V}.
\]
On this event \(\rho(B)<R/6\).
Consequently \(A=d^*B\) lies in the
measurable link set
\[
 \mathcal G_R=\{a\in\mathcal H_M:
                   \max_e|a_e|<R,\ F(a)\succeq\kappa_R I\}
                    \subset\Omega_R.
\]
It follows that
\(\mathbb P\{A\in\mathcal G_R\}\ge p_{R,b}^{18V}\).
This is an inequality between probabilities
of events under the auxiliary Gaussian;
injectivity of \(d^*\) is not required.

On \(\mathcal G_R\),
\((\det F)^t\ge\kappa_R^{td_{\rm g}}
                 \ge\kappa_R^{d_{\rm g}}\).
The proved bounds \(S\le12Q\) and
\(\mathcal H\le Q\) imply
\(\beta S+t\mathcal H\le bQ\).
Use also the last inequality in
\eqref{eq:v84-determinant}. Integrating
the original positive integrand gives
\[
 \begin{aligned}
 \mathcal Z_t(\beta)
 &\ge\kappa_R^{d_{\rm g}}\int_{\mathcal G_R}e^{-bQ(a)}\,da\\
 &\ge\kappa_R^{d_{\rm g}}
       \int_{\mathcal G_R}
         e^{-16b\langle a,\Delta^{-1}a\rangle}\,da\\
 &=\kappa_R^{d_{\rm g}}
       \left(\frac{\pi}{16b}\right)^{n/2}
       D_M^{9/2}\mathbb P\{A\in\mathcal G_R\}.
 \end{aligned}
\]
The probability lower bound and \(D_M\ge1\)
prove \eqref{eq:v84-normalizer}.
Every integral in this display is with
respect to the original link-space
Lebesgue measure.
\end{proof}

For explicit estimates without a Gaussian
distribution function one may use
\begin{equation}\label{eq:v84-scalar-lower}
                 p_{R,b}\ge\frac13\min\{h_{R,b},1\}.
\end{equation}
Indeed on \(0\le s\le\min\{h_{R,b},1\}\)
the integrand is at least \(e^{-1/2}\).
The elementary bounds \(\pi<4\), \(e<3\)
give \(\sqrt{2/(\pi e)}>1/3\).
In particular \(p_{R,b}\ge p_{R,13}>0\)
for every \(\beta\ge1\) at fixed \(R\).

The event probability need not be bounded
away from zero as \(V\) grows.
Its cost per site is the explicit
\(18\log(1/p_{R,b})\). This is precisely
why no high-probability Gaussian
small-field assertion is needed here.

\subsection{Action and amplitude estimates for arbitrary joint weak coupling}

Define the new, explicitly larger if necessary,
entropy allowance and its amplitude consequences:
\begin{equation}\label{eq:v84-budgets}
 \begin{gathered}
 \widehat\Lambda
   =d_{\rm g}\log(\kappa_R^{-1})
      +\frac n2\log(1+16eR^2b)
      +18V\log(p_{R,b}^{-1}),\qquad
 \widehat\tau=\frac{\widehat\Lambda}{\beta V},\\
 \widehat m
   =\min\left\{4R^2,\
       \min_{3\le\ell\le M}
                    (C_\ell+D_\ell\widehat\tau)\right\},\qquad
 \widehat\alpha=\min\{d_{\rm g}/V,48\widehat m\}.
 \end{gathered}
\end{equation}
Here \(C_\ell=72(\sec(\pi/\ell)-1)\) and
\(D_\ell=18\pi^2\ell(\ell-1)\) are exactly
the V81 constants. The notation with hats
distinguishes these budgets from the old
ones; old threshold-dependent statements
are not silently being redefined.

\begin{theorem}[Uniform-in-power actual amplitude control]
\label{thm:v84-amplitude}
For every \(0\le t\le1\), under the unchanged
\(\nu_t^\beta\),
\begin{equation}\label{eq:v84-amplitude}
 \frac1V\mathbb E_t S\le\widehat\tau,\qquad
 \frac1V\mathbb E_t Q\le\widehat m,\qquad
 \mathbb P_t\{S/V\ge u\}
       \le\min\{1,e^{\widehat\Lambda-\beta Vu}\}
                         \quad(u\ge0).
\end{equation}
For fixed \(R\), as \(\beta\to\infty\),
uniformly in \(M\ge3\) and \(t\),
\begin{equation}\label{eq:v84-joint-amplitude}
 \widehat\tau\le C_R\frac{\log(2+\beta)}{\beta},
 \qquad
 \widehat m\le
       C_R\sqrt{\frac{\log(2+\beta)}{\beta}}
                       +\frac{72\pi^2}{M^2}.
\end{equation}
The second inequality holds for all sufficiently
large \(\beta\), with a threshold depending
only on \(R\).
Thus \(\mathbb E Q/V\to0\), and \(Q/V\to0\)
in probability, for every joint sequence
\(M\to\infty\), \(\beta_M\to\infty\).
\end{theorem}
\begin{proof}
The normalizers imply
\[
        \log\{\mathcal Z_t(0)/\mathcal Z_t(\beta)\}
                                  \le\widehat\Lambda.
\]
Nonnegativity of relative entropy between
the same-\(t\) laws at couplings \(\beta\)
and zero gives the mean action bound,
as in V80. Bounding a tail numerator by
\(e^{-\beta Vu}\mathcal Z_t(0)\) gives
the tail. The deterministic V81 estimate
\(Q/V\le C_\ell+D_\ell S/V\) holds
throughout the common positive domain.
Taking expectations and minimizing yields
\eqref{eq:v84-amplitude}.

Since \(d_{\rm g}/V\le3\), \(n/V\le9\),
and \(p_{R,b}\ge p_{R,13}>0\), the
definition of \(\widehat\Lambda\) gives
the first bound in \eqref{eq:v84-joint-amplitude}.
It is independent of volume.
For sufficiently large \(\beta\),
\(\widehat\tau\le1/81\). The deterministic
integer-scale optimization used in
\eqref{eq:v81-optimized}, now with
\(\widehat\tau\), gives
\[
          \widehat m\le
                144\pi^2\sqrt{\widehat\tau}
                            +72\pi^2/M^2.
\]
That optimization only uses the explicit
\(C_\ell,D_\ell\), not V80's probability
threshold. This proves the stated estimate.
The final probability assertion is Markov's
inequality at the actual endpoint \(t=1\).
\end{proof}

This improves the range of a proved actual-law
estimate. It does not prove that the full
unrestricted compact Wilson measure occupies
this chart, or that the chart is a globally
injective gauge parametrization.

\subsection{Paying the radial inverse cost with the actual amplitude}

Let \(C_0\) retain its exact V64 meaning,
and define
\begin{equation}\label{eq:v84-inverse-budget}
 \widehat k
     =\frac{n+d_{\rm g}}V
               +\frac{C_0R^2}{3}\beta\widehat m
     \le12+\frac{C_0R^2}{3}\beta\widehat m .
\end{equation}

\begin{proposition}[An amplitude-paid inverse trace]
\label{prop:v84-inverse}
For \(0<t\le1\),
\begin{equation}\label{eq:v84-virial}
 t\,\mathbb E_t\operatorname{Tr}F^{-1}
  +\frac\beta2\mathbb E_t\langle a,\Delta a\rangle
  +\frac{2t}{3}\mathbb E_t Q
       \le n+td_{\rm g}
                  +\frac{\beta C_0R^2}{3}\mathbb E_t Q .
\end{equation}
In particular
\begin{equation}\label{eq:v84-inverse}
        \frac tV\mathbb E_t\operatorname{Tr}F^{-1}
                         \le\widehat k.
\end{equation}
There is no inverse-trace assertion at \(t=0\).
\end{proposition}
\begin{proof}
At \(t=1\), \eqref{eq:v84-virial} is
exactly V79's full-matrix choice in
\eqref{eq:v79-virial}, not a new identity.
To check its use at all positive powers,
the potential is
\(U_t=\beta S+t\mathcal H-t\log\det F\).
The already-proved zero-background radial
inequalities give
\[
 \begin{aligned}
 DU_t(a)[a]\ge{}&
 t\operatorname{Tr}F^{-1}-td_{\rm g}
  +\frac\beta2\langle a,\Delta a\rangle\\
 &+\frac{2t}{3}Q-\frac{\beta C_0R^2}{3}Q .
 \end{aligned}
\]
Integrate first on the convex sets
\(\Omega_{R,\epsilon}\) used in V75 and V82.
Their outward radial boundary flux is
nonnegative, so \(\int DU_t[a]e^{-U_t}
\le n\int e^{-U_t}\) there.
The bounded amplitude and kinetic terms
pass to the limit by dominated convergence.
The inverse-trace term is nonnegative and
passes by monotone convergence.
Divide by the unchanged normalizer.
This proves \eqref{eq:v84-virial}, with
the radius boundary retained.
Drop the two nonnegative terms and insert
\(\mathbb E_tQ/V\le\widehat m\) to prove
\eqref{eq:v84-inverse}.
\end{proof}

The improvement is in an expectation under
the actual specified power, not the
pointwise assertion \(Q\le V\widehat m\),
which has not been made.
Also \(\widehat m\le4R^2\) shows that
\(\widehat k\) is no larger than the old
worst-case inverse budget \(k_*\) at the
same parameters. Its sharper asymptotic
cost is what the next step uses.

\subsection{A larger uniform fractional range, including logarithmic trajectories}

Retain V83's full unnormalized operator
\(L_a=\Delta^{1/2}F(a)\Delta^{1/2}\) on
primed vertices, with zero extension on
constants, and its fractional inverse
\(\mathcal G_q(a)\).
For \(0<q<2/3\), put
\begin{equation}\label{eq:v84-moment-budget}
 \widehat{\mathfrak C}_q
  =3+\frac{4q}{2-q}
     +\frac{2^{3q/2}}{1-3q/2}
               \widehat\alpha^{\,1-3q/2}\widehat k^q .
\end{equation}

\begin{theorem}[The refined actual fractional bound]
\label{thm:v84-fractional}
At the actual endpoint \(t=1\), for every
\(\beta\ge1\), \(M\ge3\), and \(0<q<2/3\),
\begin{equation}\label{eq:v84-moment}
        \frac1V\mathbb E_{\mu_{0,\beta}}
                    \operatorname{Tr}\mathcal G_q(a)
                         \le\widehat{\mathfrak C}_q .
\end{equation}
Fix \(R\) and \(0<q<2/5\). If
\begin{equation}\label{eq:v84-window}
 M\longrightarrow\infty,\qquad
 \beta_M\longrightarrow\infty,\qquad
                 \frac{\beta_M^q}{M^{\,2-q}}\longrightarrow0,
\end{equation}
then the exceptional term in
\(\widehat{\mathfrak C}_q\) tends to zero.
In particular the fractional moment is
uniformly bounded along that sequence.
Every \(\beta_M=c\log M\), \(c>0\),
eventually satisfies these hypotheses.
\end{theorem}
\begin{proof}
The actual inputs for V83's min--max proof
are now
\[
 \frac1V\mathbb E N_F(1/2)\le\widehat\alpha,\qquad
 \frac1V\mathbb E\operatorname{Tr}F^{-1}\le\widehat k .
\]
The first follows from the unchanged
\(\operatorname{Tr}(F-I)^2\le12Q\) and
the amplitude theorem; the second is
\eqref{eq:v84-inverse} at \(t=1\).
Thus the exact same noncommuting count gives
\[
 \frac1V\mathbb E N_{L_a}(\epsilon)
   \le\min\{3,\ 4\epsilon^2+\widehat\alpha,\
                         2(\epsilon\widehat k)^{2/3}\}.
\]
V83's layer-cake integration proves
\eqref{eq:v84-moment}. Its former coupling
threshold entered only through those inputs,
and both have now been verified for every
\(\beta\ge1\).

For the improvement in the uniform range,
write \(m=\widehat m\) in this proof.
For \(0<q<2/3\), the inequalities
\(\widehat\alpha\le48m\) and
\(\widehat k\le12+C_R\beta m\), and
subadditivity of the \(q\)-th power, give
\begin{equation}\label{eq:v84-rare-factor}
 \widehat\alpha^{1-3q/2}\widehat k^q
   \le C_{R,q}
       \{m^{1-3q/2}+\beta^q m^{\,1-q/2}\}.
\end{equation}
The first term tends to zero under any
joint weak-coupling sequence.
Since \(0<1-q/2<1\), insert
\eqref{eq:v84-joint-amplitude} in the second
term, again using subadditivity:
\begin{equation}\label{eq:v84-rare-rates}
 \beta^q m^{\,1-q/2}
 \le C_{R,q}\left\{
    \beta^{(5q-2)/4}
       \log(2+\beta)^{(2-q)/4}
             +\frac{\beta^q}{M^{\,2-q}}\right\}.
\end{equation}
The coupling power is negative exactly
when \(q<2/5\). The volume term vanishes
by \eqref{eq:v84-window}. This proves
the theorem. At \(q=2/5\) the displayed
logarithmic term is not controlled by
a decaying power, so that endpoint is
not asserted.
\end{proof}

For example, \(q=1/3\) is permitted when
\(\beta_M=o(M^5)\) and \(\beta_M\to\infty\).
The two terms in \eqref{eq:v84-rare-rates}
then are
\[
      \beta^{-1/12}\log(2+\beta)^{5/12},
                      \qquad \beta^{1/3}M^{-5/3}.
\]
This is strictly stronger than merely
putting a new constant in the old
\(q<2/7\) estimate.
It still does not reach inverse power one.

\subsection{Consequences without repeating the kernel and entropy proofs}

All spatial arguments in V83 used only
locality of \(L_a\), positivity on the
primed space, translation covariance of
the actual law, and the moment just
proved. None used a high-probability
Gaussian event once its moment input
was supplied. Therefore
\eqref{eq:v83-spatial-T} and
\eqref{eq:v83-spatial} hold with
\(\mathfrak C_q\) replaced by
\(\widehat{\mathfrak C}_q\), for all
the parameters of \eqref{eq:v84-moment}.
The constant-mode correction is unchanged.

Likewise V83's Hilbert--Schmidt
functional-calculus comparison uses only
\(\mathbb E Q/V\le m\) and the fractional
moment. Thus, for \(0<p<q<2/3\),
\begin{equation}\label{eq:v84-kernel-comparison}
 \mathbb E\|\mathcal G_p(a;x,y)-\mathcal G_p^0(x,y)\|
 \le(32\sqrt3\,p+\widehat{\mathfrak C}_q+\mathfrak C_q^0)
                  \widehat m^{(q-p)/(2(q+1))}.
\end{equation}
This is a comparison of actual kernel
expectations in the sense of the
displayed expected norm, not just a
comparison of spectral distributions.

Taking \(q=1/3\), \(p=1/6\), fixed \(R\),
and any joint sequence
\(\beta_M\to\infty\), \(\beta_M=o(M^5)\),
gives, for all sufficiently large \(M\),
\begin{equation}\label{eq:v84-sixth}
 \begin{aligned}
 \mathbb E\|\mathcal G_{1/6}(a;x,y)\|
  &\le C_R
        \left(\frac{\log(2+d_M(x,y))}{d_M(x,y)^2}\right)^{1/6},
                           &&x\ne y,\\
 \mathbb E\|\mathcal G_{1/6}(a;x,y)-\mathcal G_{1/6}^0(x,y)\|
  &\le C_R\widehat m^{1/16},
                           &&x,y\in X_M.
 \end{aligned}
\end{equation}
In the first bound the finite-volume
term is absorbed exactly as in V83:
its ratio to the displayed distance
factor is bounded by a constant times
\(d_M^{2/3}/(V\log(2+d_M)^{1/3})\),
and \(d_M\le2M\).
For the second bound,
\((q-p)/(2(q+1))=1/16\).
Both conclusions therefore apply to
every logarithmic sequence
\(\beta_M=c\log M\), \(c>0\).

For completeness the power-family
inputs for V82's logarithmic comparison
have also been verified here.
One may replace its budgets by
\[
 \widehat{\mathcal B}
 =14\widehat m+\widehat\alpha
       \left\{\log\left(\frac{e\widehat k}
                              {\widehat\alpha}\right)+1\right\}.
\]
The unchanged integration over \(0<t\le1\)
then gives
\begin{equation}\label{eq:v84-entropy}
 \begin{gathered}
 0\le\frac1V\log\frac{\mathcal Z_0(\beta)}
                          {\mathcal Z_1(\beta)}
                         \le\widehat{\mathcal B},\\
 0\le\frac1V D(\mu_{0,\beta}\Vert\nu_0^\beta)
                         \le\widehat{\mathcal B}.
 \end{gathered}
\end{equation}
Here \(\mathcal Z_0(\beta)\) is the
power-zero normalizer at coupling
\(\beta\), not \(\mathcal Z_1(0)\).
The proof still retains the factor
\(t\) in the inverse estimate and
uses the integrability of \(\log(1/t)\);
no inverse estimate at \(t=0\) has
been introduced.
At fixed \(R\), these per-site bounds
tend to zero whenever
\(\beta_M\to\infty\), \(M\to\infty\),
and \(\log(2+\beta_M)/M^2\to0\).
Indeed \(\widehat k\le C_R(1+\beta)\),
\(\widehat\alpha\le48\widehat m\), and
\(\widehat m\log(2+\beta)\to0\) by
\eqref{eq:v84-joint-amplitude}.
These remain per-site conclusions,
not total-variation convergence.

\subsection{What the logarithmic window does and does not settle}

The new construction removes a sufficient
normalizer threshold, not a physical
infrared cutoff. It does not contradict
V61's sharp free-Gaussian obstruction
to a vanishing odd gauge norm when
\(\beta/\log M\) fails to diverge.
The Gaussian here has a different
covariance, and its good event is only
used to lower-bound a subset integral.
No high-probability claim for that
event under the free or actual law
has been deduced.

Allowing \(\beta=c\log M\) is relevant
to the scale dependence being studied
in the shell programme, but it does
not identify \(c\) with a renormalized
physical coefficient or identify
\(M^{-1}\) with a physical lattice
spacing. Those identifications require
the actual interacting continuum and
observable comparisons, not a choice
of notation. In particular the free
unit-lattice limit of a fractional
gauge kernel is not a free or an
interacting physical continuum limit.

The improved radial cost still leaves
the very-low-energy counting power
\(\epsilon^{2/3}\) in place.
The new budget pays a uniform moment
below \(2/5\), not a moment above one.
A next improvement must address that
spectral tail or a genuinely local
inverse mechanism, and must ultimately
connect gauge-invariant physical
correlations to the Hamiltonian.
Neither an ordinary gauge inverse
bound nor exponential physical
clustering follows from the present
fractional result.

\subsection{Exact checks and preservation}

The certificate is
\begin{center}\small
\path{scripts/verify_ym_f14_v84_curl_normalizer.py}.
\end{center}
It uses six rational three-colour
two-form fields on actual tori of
sides two, three, and four.
For their codifferentials it checks
the Coulomb and mean-zero constraints,
the six-\(\rho\) link estimate, the
shifted four-term curl identity,
and equality of the actual odd
gauge form with the current and
two-form pairings. It also checks
the corresponding six-\(\rho\)
bilinear bound after squaring.
There are 12,708 scalar curl
components checked across the six
fields.

On sides two and three, independently
assembled integer difference matrices
verify \(dD=0\) and the full link
Hodge identity. Fraction-free
determinants of scalar Laplacian
cofactors agree with the primed
Fourier eigenvalue products after
multiplication by the number of
sites. The side-two primed product
is exactly \(22265110462464\).
These are finite checks on the
normalization convention; the
positive-integer determinant proof
above applies to every side length.

Scalar rational checks verify the
Gaussian threshold square
\(8bR^2/27\), the floor
\(\kappa_R\ge23/32\), the elementary
probability lower constant, the
coupling powers at \(q=1/4,1/3,3/8\),
and the comparison power \(1/16\).
No Monte Carlo normalizer or
floating-point sign test is used.
The Gaussian subset probability,
all-volume estimates, and interacting
expectations have the analytic
proofs above; finite checks are not
a substitute for them.

The certificate contains 8,735 bytes and has SHA--256
\begin{center}\scriptsize\ttfamily
D627BB10A1DE5E21C93676E4E3CA698A6996A6CBB3B0F0E0874E75CF27B07718.
\end{center}
The predecessor V83 contains 2,056,572 bytes and has SHA--256
\begin{center}\scriptsize\ttfamily
3F5A0EFFC5ECE4FAA4DF80C4B03560968762F695A1D98233A93B8979D1A6F150.
\end{center}
Its body is preserved except for
the title version and this terminal
insertion. V84 is the single active
main note and retains all V83
contents. Earlier certificates,
companions, supplied sources,
monographs and archives are unchanged.

\begin{firewall}
We have removed an artificial
volume-dependent normalization
threshold and improved a proved
fractional inverse budget under
the same restricted Wilson law.
The full interacting continuum
Yang--Mills construction and its
positive physical Hamiltonian
mass gap remain unproved.
The original goal remains active.
\end{firewall}
\section{V85: local coexact moves and transport of the complete restricted law}
\label{sec:v85}

\subsection{The compulsory companion checkpoint and the next target}

V84 was a progress turn: the actual normalization
estimate no longer requires a volume-dependent
lower threshold on the coupling. Its amplitude
estimate also reduces the radial inverse-trace
cost. The ordinary unrestricted vertex inverse
and the physical mass gap remain uncontrolled.
This checkpoint first reopens the companion
notes and then uses the new budget in a local
problem, rather than changing the exponent
of the same global inverse integral.

The latest SM--Einstein file remains V72
(607,372 bytes), the NS Stability file remains
V26 (278,283 bytes), and Shell Mixing remains
V16 (623,261 bytes). Their fingerprints agree
with the V80 checkpoint. The reopened conclusions
have the following precise scopes:
\begin{itemize}
 \item SM V72 proves fixed-time many-copy
 covariance transport with the momentum cutoff
 fixed. It does not give a cutoff-uniform
 physical finite-species limit or the present
 local gauge inverse.
 \item NS V26 derives rank-one contractions
 of Oseen-kernel derivatives and reports
 numerical comparisons. These are not the
 Yang--Mills operator or a proof of its
 gauge-boundary control. No numerical NS
 constant is imported here.
 \item Shell V16 identifies the complete
 measured one-loop coefficient on a flat
 holonomy family, retaining the normal and
 tangent root factors. Its stated remaining
 noncommuting calculation and its lack of a
 joint nonperturbative remainder are still
 operative. A coefficient limit does not
 replace an estimate for the actual law.
\end{itemize}
The additional parallel YM V5 file was also
reopened. Its conclusion distinguishes compact
flux/holonomy labels from extensive integer
charge, and rules out its proposed uniform
global-charge tail under the stated
tensorization hypothesis. That reinforces
the local-observable direction here; it does
not supply a gauge resolvent bound. This
54,174-byte source has SHA--256
\begin{center}\scriptsize\ttfamily
306F1BB84CFA8A8D47EA569EC17E9545F9867C8D32B548EC9D9C444B4F3C0512.
\end{center}
Source decisions and stopping rules are treated
as source material, not as instructions to
alter this task. All companions remain unchanged.
The next regular checkpoint is V90.

The present acquisition is a local stability
and measure-transport theorem. A compactly
supported coexact change of a Coulomb field
has a gauge increment factored through the
gradients on a finite edge set. This allows
its admissibility and determinant ratio to
be tested by a finite compression of an
ordinary gradient inverse. V84 pays the
mean trace of that compression. We then
include the Wilson and Haar changes and
the possible loss of chart support, obtaining
a bound in total variation for the complete
normalized law under a deterministic local
translation.

\subsection{A local ordinary inverse with an actual trace budget}

Use \(M\ge3\), \(V=M^4\), \(0<R\le1/4\),
\(\beta\ge1\), and exactly V84's
\(\mu=\mu_{0,\beta}\) on \(\Omega_R\).
Write \(\widehat m,\widehat\tau,\widehat k\)
for its proved budgets. On primed vertex
colour fields,
\[
 L_a=\Delta^{1/2}F(a)\Delta^{1/2}\succ0,
 \qquad U=D\Delta^{-1/2},\qquad U^*U=I.
\]
For a finite set \(E\) of positively oriented
edges, let \(D_E=\mathbf1_ED\), with all
three colour coordinates retained, and set
\begin{equation}\label{eq:v85-local-inverse}
 K_E(a)=D_E L_a^{-1}D_E^*
       =\mathbf1_EUF(a)^{-1}U^*\mathbf1_E.
\end{equation}
This is a positive matrix of size \(3|E|\).
It uses inverse power one, not a fractional
inverse; its sources are discrete divergences.
It is a spatial compression, not an inverse
with Dirichlet boundary conditions on \(E\).

\begin{lemma}[The localized gradient trace]
\label{lem:v85-local-trace}
For every deterministic edge set \(E\),
\begin{equation}\label{eq:v85-local-trace}
              \mathbb E_\mu\operatorname{Tr}K_E
                                      \le |E|\widehat k.
\end{equation}
In particular, for a deterministic colour
edge vector \(b\) supported on \(E\), with
\(j=D_E^*b\),
\begin{equation}\label{eq:v85-divergence-source}
          \mathbb E_\mu\langle j,L_a^{-1}j\rangle
                       \le |E|\widehat k\,\|b\|^2.
\end{equation}
The estimates are uniform in volume at
fixed \(R,\beta\).
\end{lemma}
\begin{proof}
On all edges the positive operator
\(UF^{-1}U^*\) has trace
\(\operatorname{Tr}F^{-1}\), since \(U\)
is an isometry on primed vertices.
Translation covariance makes the expected
diagonal colour trace of an edge depend
only on its orientation. The sum of the
four orientation traces is therefore
\[
         V^{-1}\mathbb E_\mu\operatorname{Tr}F^{-1}
                                  \le\widehat k.
\]
Each trace is nonnegative, so each is
bounded by \(\widehat k\). Sum over \(E\)
to obtain \eqref{eq:v85-local-trace}.
Positivity gives
\(\langle b,K_Eb\rangle\le
\|b\|^2\operatorname{Tr}K_E\), proving
the source estimate. Finally
\(\widehat k\le12+(4/3)\beta C_0R^4\)
at fixed parameters, independently of \(M\).
\end{proof}

This statement does not bound \(L_a^{-1}\)
on arbitrary vertex sources. Nor is
\(K_E(a)\) a function of the field on
\(E\) alone: it is computed using the
full-volume inverse. Finite matrix size
must not be confused with local dependence
of its entries on the environment.

\subsection{The exact local increment for a coexact move}

Let \(\mathcal P\) be a deterministic set
of oriented plaquettes and let \(B\) be
a real colour two-form supported there.
Put
\[
       h=d^*B,\qquad
       \rho=\max_{x,\mu<\nu}|B_{\mu\nu}(x)|,\qquad
       E=\bigcup_{p\in\mathcal P}\partial p.
\]
As in V84, \(h\in\mathcal H_M\),
\(\operatorname{supp}h\subset E\), and
\(\max_e|h_e|\le6\rho\).
Set
\begin{equation}\label{eq:v85-size}
             r=6\rho,\qquad \eta=6(1+3R)\rho .
\end{equation}
Only \(0<\rho\le R/12\) is needed below.
The zero move is immediate.

\begin{lemma}[Exact finite-edge factorization]
\label{lem:v85-factor}
Suppose \(a\in\Omega_R\) and
\(\max_{e\in E}|a_e|<R-r\).
Then for every \(0\le s\le1\) there is
a self-adjoint matrix \(H_s(a,B)\) on
the edge-colour space of \(E\) such that
\begin{equation}\label{eq:v85-factor}
 L_{a+sh}-L_a=D_E^*H_s(a,B)D_E,\qquad
                             \|H_s(a,B)\|\le s\eta.
\end{equation}
The factorization is valid before
positivity of \(L_{a+sh}\) is established.
\end{lemma}
\begin{proof}
The radius assumption keeps the whole
segment within the link ball of radius
\(R\): outside \(E\) no link changes.
For arbitrary edge-colour vectors \(u,v\)
define \(H_B\), extended by zero off \(E\),
through
\begin{equation}\label{eq:v85-HB}
 \begin{aligned}
 \langle u,H_Bv\rangle
 =\sum_{x,\mu<\nu} B_{\mu\nu}(x)\cdot
 \big\{&
 v_\mu(x)\times u_\nu(x+e_\mu)
 +v_\mu(x+e_\nu)\times u_\nu(x)\\
 &-v_\nu(x)\times u_\mu(x+e_\nu)
 -v_\nu(x+e_\mu)\times u_\mu(x)\big\}.
 \end{aligned}
\end{equation}
Every edge occurring here lies on a
plaquette in \(\mathcal P\).
The first and fourth terms form a
transpose pair, as do the second and
third, because the matrix \(v\mapsto B\times v\)
is skew-adjoint. Thus \(H_B\) is self-adjoint.
For each edge there are at most two
incident plaquettes for each of the
three other coordinate directions.
Each contributes an off-diagonal block
of norm at most \(\rho\).
The block row and column sums are at
most \(6\rho\), whence
\[
                         \|H_B\|\le6\rho.
\]
Alternatively this follows from the
restricted-edge Cauchy--Schwarz proof
of \eqref{eq:v84-potential}.

Inserting \(u=D\phi\), \(v=D\psi\) in
\eqref{eq:v85-HB} gives exactly V61's
four-term curl identity. Its factor
\(\operatorname{ad}_h v=2h\times v\)
and the half in the odd gauge operator
cancel as before. Hence the odd
unnormalized increment is \(D_E^*H_BD_E\).

The remaining increment is the edge-diagonal
matrix
\[
       J_s(e)=
       \sigma(\operatorname{ad}_{a_e+sh_e})
                        -\sigma(\operatorname{ad}_{a_e}),
              \qquad e\in E.
\]
The pointwise derivative estimate in the
proof of Lemma~\ref{lem:v74-even} is
\(\|D\sigma(\operatorname{ad}_v)[w]\|
\le3R|w|\) for \(|v|\le R\le1/4\).
That proof used the convergent cotangent
partial fractions, also recorded in
\href{https://dlmf.nist.gov/4.22#E3}
{NIST DLMF 4.22.3}.
Here the entire link segment stays in
that ball, so
\[
                 \|J_s\|\le3Rs\max_e|h_e|
                                      \le18Rs\rho .
\]
No extra assumption \(R+r\le1/4\)
is needed: \(R\) itself is the radius
of this guarded segment.
Take \(H_s=sH_B+J_s\). The local formula
for \(L_a\) and the preceding bounds
prove \eqref{eq:v85-factor}.
\end{proof}

\subsection{A local stability test and an exact determinant ratio}

For \(0<\theta<1\) define an event using
the original input \(a\):
\begin{equation}\label{eq:v85-good}
 \mathcal A_{E,\rho,\theta}
 =\left\{a\in\Omega_R:
       \max_{e\in E}|a_e|<R-r,\quad
                    \eta\operatorname{Tr}K_E(a)\le\theta\right\}.
\end{equation}
It is the same event for every two-form
supported on \(\mathcal P\) with norm
at most \(\rho\). This uniformity concerns
deterministic geometric admissibility;
the measure-translation theorem below
will require a deterministic move.

\begin{theorem}[Local admissibility without a global gauge floor]
\label{thm:v85-stability}
On \(\mathcal A_{E,\rho,\theta}\), for
every \(0\le s\le1\),
\begin{equation}\label{eq:v85-stability}
 \begin{gathered}
 a+sh\in\Omega_R,\qquad
 (1-s\theta)L_a\preceq L_{a+sh}\preceq(1+s\theta)L_a,\\
 (1+s\theta)^{-1}K_E(a)
       \preceq K_E(a+sh)
       \preceq(1-s\theta)^{-1}K_E(a).
 \end{gathered}
\end{equation}
At \(s=1\), write \(H=H_1(a,B)\),
\(K=K_E(a)\). Then
\begin{equation}\label{eq:v85-determinant}
 \frac{\det F(a+h)}{\det F(a)}
       =\det_E(I+KH)
       =\det_E(I+K^{1/2}HK^{1/2})>0,
\end{equation}
and
\begin{equation}\label{eq:v85-logdet}
 \left|\log\frac{\det F(a+h)}{\det F(a)}\right|
                   \le\frac{\eta}{1-\theta}\operatorname{Tr}K_E(a).
\end{equation}
The exact local inverse update is
\begin{equation}\label{eq:v85-resolvent}
 K_E(a+h)=K-KH(I+KH)^{-1}K.
\end{equation}
Neither \(H\) nor \(K\) is assumed invertible.
\end{theorem}
\begin{proof}
On primed vertices put
\[
 T_s=L_a^{-1/2}D_E^*H_sD_E L_a^{-1/2}.
\]
By \eqref{eq:v85-factor},
\[
  \|T_s\|\le s\eta\|D_E L_a^{-1/2}\|^2
           =s\eta\|K_E(a)\|
           \le s\eta\operatorname{Tr}K_E(a)\le s\theta .
\]
It follows that \(I+T_s\) is positive
and lies between \(1-s\theta\) and
\(1+s\theta\). Conjugation by \(L_a^{1/2}\)
proves the operator comparison and
strict positivity. The radius guard
has already kept the segment in its
original link domain.
Inverting the comparison and compressing
with \(D_E\) proves the assertion about
\(K_E(a+sh)\).

The determinants of the two \(L\)'s
have the same \(\det{}'\Delta\) factor
as the corresponding determinants of
\(F\). Sylvester's finite-dimensional
determinant identity, applied to
\(T_1\), gives \eqref{eq:v85-determinant}.
It also shows that \(I+KH\) is invertible,
even if \(K\) is singular.
The nonzero eigenvalues can equivalently
be read in the symmetric local matrix
\(K^{1/2}HK^{1/2}\); noncommutation is
not being suppressed.

For \(|u|\le\theta<1\),
\(|\log(1+u)|\le|u|/(1-\theta)\).
In addition
\[
 \|T_1\|_1
   \le\|H\|\,\|D_E L_a^{-1/2}\|_{\rm HS}^2
                       \le\eta\operatorname{Tr}K_E(a).
\]
Apply the scalar inequality to all
eigenvalues of \(T_1\) to obtain
\eqref{eq:v85-logdet}.
Finally multiply out the finite-rank
resolvent identity
\[
 L_{a+h}^{-1}
 =L_a^{-1}
  -L_a^{-1}D_E^*H(I+KH)^{-1}D_E L_a^{-1}.
\]
This identity uses no inverse of \(H\).
Compression by \(D_E\) yields
\eqref{eq:v85-resolvent}.
All vertex inverses here are primed;
the common constant modes have not
been inverted.
\end{proof}

\begin{proposition}[The actual probability of the local guard]
\label{prop:v85-guard}
The unchanged input law satisfies
\begin{equation}\label{eq:v85-guard}
 \mu(\mathcal A_{E,\rho,\theta}^{\,c})
 \le\min\left\{1,\
      |E|\left(\frac{\widehat m}{(R-r)^2}
                         +\frac{\eta\widehat k}{\theta}\right)\right\}.
\end{equation}
\end{proposition}
\begin{proof}
Translation covariance and nonnegativity
give \(\mathbb E|a_e|^2\le\widehat m\)
for each oriented edge, by summing
the four orientations at one site
and using \(\mathbb E Q/V\le\widehat m\).
A union bound and Markov's inequality
therefore bound the radius failure by
\(|E|\widehat m/(R-r)^2\).
Equation \eqref{eq:v85-local-trace} and
Markov's inequality bound the inverse
failure by \(\eta|E|\widehat k/\theta\).
No independence of edges or spectral
variables is used.
\end{proof}

Unlike the previous global minimum test,
this bound charges the size of the
modified patch, not the number of
all lattice sites. It does not assert
that a local inverse is uniformly
bounded pointwise, nor that a move
with a fixed thermal-scale amplitude
already has small failure probability.

\subsection{The actual Wilson and Haar increments are also local}

Let \(\mathcal P(E)\) denote all
plaquettes touching an edge of \(E\);
then \(|\mathcal P(E)|\le6|E|\).
Write \(S_{\mathcal P(E)}(a)\) for
the sum of the original nonnegative
Wilson plaquette actions on that set.

\begin{lemma}[Action-sensitive cost of the same deterministic move]
\label{lem:v85-density-increments}
For \(h=d^*B\) as above,
\begin{equation}\label{eq:v85-action}
 |S(a+h)-S(a)|
 \le24\sqrt3\,\rho
                 \sqrt{|E|S_{\mathcal P(E)}(a)}
                          +432\rho^2|E|.
\end{equation}
This action bound needs no positivity
of the output gauge matrix.
Under the actual input law,
\begin{equation}\label{eq:v85-action-mean}
 \mathbb E_\mu|S(a+h)-S(a)|
       \le |E|\{72\sqrt2\,\rho\sqrt{\widehat\tau}
                                               +432\rho^2\}.
\end{equation}
On the radius guard in
\eqref{eq:v85-good},
\begin{equation}\label{eq:v85-Haar}
              |\mathcal H(a+h)-\mathcal H(a)|
                                     \le12R\rho|E|.
\end{equation}
\end{lemma}
\begin{proof}
Represent \(\SU(2)\) by unit quaternions
with the existing normalization
\(s_p=1-\operatorname{Re}U_p
       =|U_p-1|^2/2\).
The exponential map is one-Lipschitz
from the Lie-algebra norm to the
bi-invariant geodesic distance:
the differential is an average of
isometric adjoint actions, so the
image of the straight segment has
length at most the length of that
segment. Multiplication and inversion
preserve the group distance.
Thus for
\[
 t_p=\sum_{e\in\partial p\cap E}|h_e|
\]
the quaternion chord satisfies
\(|U_p(a+h)-U_p(a)|\le t_p\).
Expanding its squared distance from
the identity gives
\[
                 |s_p(a+h)-s_p(a)|
                  \le\sqrt{2s_p(a)}\,t_p+\tfrac12t_p^2.
\]
This retains the input action in
the linear term; a fixed Lipschitz
bound would lose the weak-coupling
gain needed below.

Only \(\mathcal P(E)\) contributes.
Each plaquette has four links and
each link belongs to six plaquettes,
so
\[
 \sum_{p\in\mathcal P(E)}t_p^2
       \le24\sum_{e\in E}|h_e|^2
       \le864\rho^2|E|.
\]
Cauchy--Schwarz in the linear term
proves \eqref{eq:v85-action}.
The expected action of a plaquette
of any fixed orientation is at most
\(\mathbb E S/V\le\widehat\tau\).
Consequently
\(\mathbb E S_{\mathcal P(E)}
\le6|E|\widehat\tau\).
Jensen's inequality for the square
root gives \eqref{eq:v85-action-mean}.

For the one-link Haar potential
\(H_0(r)=-2\log(\sin r/r)\),
the same convergent partial fractions
give
\[
 H_0'(r)=4r\sum_{n\ge1}(n^2\pi^2-r^2)^{-1}
                          \le\tfrac43r\le2R
                       \quad(0\le r\le R).
\]
The continuous value at zero is
understood. Integrate this derivative
on the guarded segment, and use
\(\sum_{e\in E}|h_e|\le6\rho|E|\).
This proves \eqref{eq:v85-Haar}.
\end{proof}

\subsection{Transport of the complete normalized law in total variation}

Let \(T_h(a)=a+h\) on the original
Euclidean space \(\mathcal H_M\).
In the following theorem \(B\), and
hence \(h\), is deterministic.
Translation preserves this space
because \(h=d^*B\) is mean-zero and
Coulomb, and its Lebesgue Jacobian
is exactly one.
Use the convention
\(\|\nu-\mu\|_{\rm TV}=\sup_A|\nu(A)-\mu(A)|
                      =\tfrac12\int|d\nu-d\mu|\).

\begin{theorem}[A local translation estimate for the actual measure]
\label{thm:v85-TV}
For every deterministic \(B\) supported
on \(\mathcal P\), \(0<\rho\le R/12\),
and \(0<\theta<1\),
\begin{equation}\label{eq:v85-TV}
 \begin{aligned}
 \|(T_h)_*\mu-\mu\|_{\rm TV}
 \le\min\Bigg\{1,\ |E|\Bigg[
 &\frac{\widehat m}{(R-6\rho)^2}
   +\eta\widehat k
                   \left(\frac1\theta+\frac1{1-\theta}\right)\\
 &+72\sqrt2\,\beta\rho\sqrt{\widehat\tau}
       +432\beta\rho^2+12R\rho\Bigg]\Bigg\}.
 \end{aligned}
\end{equation}
The determinant, Wilson action, Haar
potential, chart restriction and
normalization are all those of
\(\mu_{0,\beta}\).
\end{theorem}
\begin{proof}
Extend the actual normalized Lebesgue
density by zero outside \(\Omega_R\):
\[
 f(a)=Z_\beta^{-1}\mathbf1_{\Omega_R}(a)
                       \det F(a)e^{-\beta S(a)-\mathcal H(a)}.
\]
For \(f(a)>0\), define the ratio
\(q_h(a)=f(a+h)/f(a)\), with value
zero if the output lies outside
the support. Translation invariance
of Lebesgue measure gives the exact
overlap identity
\begin{equation}\label{eq:v85-overlap}
 \|(T_h)_*\mu-\mu\|_{\rm TV}
       =1-\int\min\{f(a),f(a+h)\}\,da
       =\mathbb E_\mu[1-\min\{1,q_h(a)\}].
\end{equation}
To check the sign of the translation,
the pushforward density is \(f(a-h)\);
change variables in its overlap with
\(f(a)\) to obtain the first integral.

On \(\mathcal A=\mathcal A_{E,\rho,\theta}\)
both configurations belong to the
same original chart. Their normalizers
cancel exactly and
\[
 \log q_h(a)=
 \log\frac{\det F(a+h)}{\det F(a)}
       -\beta\{S(a+h)-S(a)\}
       -\{\mathcal H(a+h)-\mathcal H(a)\}.
\]
For a real \(z\),
\(1-\min\{1,e^z\}\le|z|\).
On \(\mathcal A^c\) the rejection
is bounded by one, without asserting
that the output is admissible there.
Thus
\[
 \|(T_h)_*\mu-\mu\|_{\rm TV}
           \le\mu(\mathcal A^c)
                  +\mathbb E_\mu[\mathbf1_{\mathcal A}|\log q_h|].
\]
The guard probability is
\eqref{eq:v85-guard}. The determinant
term has expected upper bound
\(\eta|E|\widehat k/(1-\theta)\)
by \eqref{eq:v85-logdet} and
\eqref{eq:v85-local-trace}.
Use \eqref{eq:v85-action-mean} for
the Wilson term, which was bounded
without restricting to \(\mathcal A\),
and \eqref{eq:v85-Haar} for the
guarded Haar term. Their sum proves
\eqref{eq:v85-TV}.
No missing normalizer ratio or
boundary cancellation is assumed.
\end{proof}

The overlap identity also says that
the mean acceptance for this deterministic
translation with acceptance probability
\(\min\{1,q_h\}\) is exactly one minus
the total-variation distance.
This identity is about a proposed move,
not a proof of ergodicity or a spectral
gap for any update chain.
For a field-dependent move \(h(a)\),
neither this translation identity nor
the unit Jacobian can be imported
without a separate change-of-variables
argument.

\subsection{A concrete logarithmic-coupling consequence}

Let \(H_\beta=\log(2+\beta)\), and choose
the upper move size
\[
             \rho_\beta
               =\min\left\{\frac R{12},\
                           \frac1{\sqrt\beta\,H_\beta}\right\}.
\]
Every deterministic two-form with
supremum norm at most \(\rho_\beta\)
obeys the following conclusion.
In applying \eqref{eq:v85-TV}, the
upper bound \(\rho_\beta\) may replace
its actual norm throughout.

\begin{theorem}[Asymptotically faithful local coexact translations]
\label{thm:v85-joint}
For fixed \(R\), all sufficiently
large \(\beta\), and every deterministic
patch and move as just specified,
\begin{equation}\label{eq:v85-joint}
 \|(T_h)_*\mu_{0,\beta}-\mu_{0,\beta}\|_{\rm TV}
 \le
 C_R|E|\left\{
      H_\beta^{-1/2}
           +\frac{\sqrt\beta}{M^2H_\beta}\right\}.
\end{equation}
Thus the distance tends to zero for
fixed-size patches whenever
\begin{equation}\label{eq:v85-window}
 M\to\infty,\qquad \beta_M\to\infty,\qquad
                 \frac{\sqrt{\beta_M}}
                      {M^2\log(2+\beta_M)}\to0.
\end{equation}
In particular every
\(\beta_M=c\log M\), \(c>0\), is included.
The probability that the entire segment
\(a+sh\), \(0\le s\le1\), leaves the
positive chart also tends to zero.
\end{theorem}
\begin{proof}
Take \(\theta=1/2\) and use
\(R-6\rho_\beta\ge R/2\).
For sufficiently large \(\beta\),
\(\rho_\beta=\beta^{-1/2}H_\beta^{-1}\).
V84 gives
\[
 \widehat\tau\le C_R H_\beta/\beta,\qquad
 \widehat m\le C_R\sqrt{H_\beta/\beta}+C M^{-2},\qquad
 \widehat k\le12+C_R\beta\widehat m.
\]
Consequently
\[
 \begin{aligned}
 \rho_\beta\widehat k
   &\le C_R\left\{
       \frac1{\sqrt\beta H_\beta}
          +H_\beta^{-1/2}
          +\frac{\sqrt\beta}{M^2H_\beta}\right\},\\
 \beta\rho_\beta\sqrt{\widehat\tau}
   &\le C_R H_\beta^{-1/2},\qquad
 \beta\rho_\beta^2=H_\beta^{-2}.
 \end{aligned}
\]
The remaining radius and Haar terms
are bounded by the right side of
\eqref{eq:v85-joint}: eventually
\(\sqrt\beta/H_\beta\ge1\) and
\(\sqrt{H_\beta/\beta}\le H_\beta^{-1/2}\).
Insert these estimates into
\eqref{eq:v85-TV}.
The same estimates in
\eqref{eq:v85-guard}, combined with
\eqref{eq:v85-stability}, prove the
assertion about the whole segment.
\end{proof}

For a bounded measurable observable \(A\),
including a bounded gauge-invariant
local observable, the usual consequence is
\[
 |\mathbb E_\mu A(a+h)-\mathbb E_\mu A(a)|
       \le2\|A\|_\infty\|(T_h)_*\mu-\mu\|_{\rm TV}.
\]
The patch may also grow if its edge
count times the explicit error in
\eqref{eq:v85-joint} still tends to
zero. No estimate here permits an
extensive number of simultaneous
moves at fixed error.

\subsection{The remaining distinction from a physical mass gap}

The local trace result is an ordinary
inverse estimate for divergence sources.
The locality of the perturbation then
makes the relevant determinant and
resolvent comparison finite dimensional.
This bypasses a global minimum-eigenvalue
test for the stated local moves, but
does not improve the full vertex
inverse estimate of V83--V84.
The local matrix still depends on
the entire environment.

The total-variation statement is
stronger than a per-site entropy
comparison, but its comparison is
different: it compares one law with
its own small deterministic translation,
not the actual law with a Gaussian
or with the determinant-free reference.
The permitted moves are smaller than
the thermal scale by the displayed
logarithmic factor:
\(\sqrt\beta\,\|h\|_\infty
\le6/\log(2+\beta)\to0\).
No invariance under a fixed nonzero
thermal-scale shift is claimed.

In particular high acceptance for
vanishing local moves does not prove
spatial mixing, a scale-uniform
functional inequality, or physical
time contraction. Iterating these
moves for a macroscopic transport
would require controlling accumulated
errors and environment dependence.
The restricted chart also still
needs its legitimate connection to
the full compact theory and to
renormalized gauge-invariant
continuum observables.

The concrete new input is a
same-measure local transport estimate
with its inverse, action, Haar,
determinant and support costs paid.
The next step can test this input
against a specific local coarse
transport or conditional response.
It must not convert the finite
patch size, a fractional gauge
kernel, or a high acceptance
probability into a physical
Hamiltonian mass gap by definition.

\subsection{Exact checks and checkpoint preservation}

The certificate is
\begin{center}\small
\path{scripts/verify_ym_f14_v85_local_coexact_transport.py}.
\end{center}
Six exact stencil tests use one-
and two-plaquette two-forms on
actual tori of sides two, three
and four. They verify the
Coulomb and mean-zero constraints,
the symmetric edge factor \(H_B\),
its local support, and
\[
         D_E^*H_BD_E
                 =\tfrac12D^*\operatorname{ad}_{d^*B}S_{\rm end}.
\]
The two sides are assembled
independently. The two-plaquette
patch has seven edges and six
vertices; its compressed matrix
has dimension 21.

On the side-two free background,
a separate rational local-matrix
increment checks the determinant
identity and both resolvent formulas.
The full constant-regularized
matrix has dimension 48; the
constants cancel from all gradient
comparisons. Its compressed trace
is exactly \(315/64\), and the
computed guard bound is
\(656383/2048000<1/2\).
The rational diagonal added in
this algebra check is explicitly
a generic test increment, not
the full nonlinear cotangent
matrix at a Yang--Mills field.
The actual even increment has
the analytic derivative bound
used above.

Twelve products of exact rational
unit quaternions check the
plaquette squared-chord identity
and its Cauchy bound. Eleven
finite cyclic translation tests
check the total-variation/overlap
identity, with zero-density
boundary points retained.
These finite checks validate
algebra and conventions; the
actual probability, all-volume
and total-variation estimates
have the analytic proofs above.

The certificate contains 9,139 bytes and has SHA--256
\begin{center}\scriptsize\ttfamily
4128FB55C8A53C10C384303671E531DEDED5495146614BD07577BB6E561016BB.
\end{center}
The predecessor V84 contains 2,083,759 bytes and has SHA--256
\begin{center}\scriptsize\ttfamily
BEFEBF86CC82016BC00991923FD770412118917ED5564CFDE3841C3D33EACCA4.
\end{center}
Its entire body is preserved except
for the title version and this
terminal insertion. V85 replaces
V84 as the sole active main note.
Earlier certificates, supplied
sources, companions, monographs
and archived lineages are unchanged.

\begin{firewall}
The new theorem controls specified
small local translations of the
same restricted Wilson measure.
It does not construct the full
interacting continuum theory,
prove its nontriviality, or
establish a positive physical
Hamiltonian mass gap. The
original goal remains active.
\end{firewall}
\section{V86: actual weak-coupling plaquette nonconcentration and the thermal-scale gate}
\label{sec:v86}

\subsection{From local transport to a gauge-invariant fluctuation}

V85 was a progress turn: it controlled
specified local coexact translations in
total variation under the complete
restricted Wilson law. Its estimates
retain the chart boundary, gauge
determinant, Haar factor, Wilson action
and common normalization. Here we use
that result to obtain a lower fluctuation
bound for the actual plaquette trace.
The observable is not replaced by its
linear curvature approximation.

This does not repeat V48's compact
terminal-region variance bound or
V53's all-depth spatial Haar comparison.
Those concern their own compact
endpoint laws and blocking regimes.
The present measure is the fixed-radius
mean-zero Coulomb law
\(\mu_{0,\beta}\) of V80--V85, with
\(M\) and \(\beta\) tending to infinity
together. Its plaquettes approach
the identity in mean, so a positive
unscaled variance is neither proved
nor expected by these estimates.

The new bound is
\(\operatorname{Var}W_p\ge
c_R/(\beta\log(2+\beta))^2\)
in an explicit joint window. We also
prove a uniform nonconcentration
statement on that scale, which is
stronger than a variance lower bound
alone. The extra logarithm is retained.
A final implication identifies one
specific unproved improvement of the
same-measure free-energy estimate
that would remove it. Neither scale
has yet been identified with a
renormalized continuum field.
The V85 companion review remains
current; the next checkpoint is V90.

\subsection{A literal single-plaquette coexact direction}

Fix an oriented plaquette
\(p=(x;\mu,\nu)\), \(\mu<\nu\), and a
unit colour vector \(v\). Let \(B^p\)
be the two-form equal to \(v\) on
that plaquette and zero elsewhere.
Define
\[
                    h^p=d^*B^p,\qquad E_p=\partial p.
\]
For \(M\ge3\) the four edges of this
plaquette are distinct. Directly from
the codifferential convention,
\begin{equation}\label{eq:v86-direction}
 \begin{aligned}
 h^p_\mu(x)&=v,& h^p_\nu(x+e_\mu)&=v,\\
 h^p_\mu(x+e_\nu)&=-v,& h^p_\nu(x)&=-v .
 \end{aligned}
\end{equation}
All other edge values vanish. Thus
\[
 h^p\in\mathcal H_M,\qquad
 \|h^p\|_\infty=1,\qquad
                 (dh^p)_p=4v.
\]
These are properties of the full
periodic field, not merely of its
restriction to the plaquette.

Write the original gauge-invariant
Wilson observables as
\begin{equation}\label{eq:v86-observable}
 W_p(a)=\tfrac12\operatorname{Re}\operatorname{Tr}U_p(a),
                     \qquad s_p(a)=1-W_p(a)\in[0,2].
\end{equation}
They are well-defined through the
exponentials even outside the chart.
At zero input the four oriented
factors for \(t h^p\) are all
\(\exp(tv)\). Therefore
\begin{equation}\label{eq:v86-zero-curve}
        U_p(th^p)=\exp(4tv),\qquad
        s_p(th^p)=1-\cos(4t),\qquad
        D^2s_p(0)[h^p,h^p]=16.
\end{equation}
The same move also changes neighbouring
plaquettes. Their contributions to
the density are already included in
V85's full-action transport bound;
we will not replace that density by
a single-plaquette law.

\subsection{A uniform second difference of the nonlinear observable}

For a variation of the four plaquette
link arguments, let \(\|\cdot\|_{\infty,p}\)
be the maximum of the four colour norms.

\begin{lemma}[A quantitative third derivative and a convex pulse]
\label{lem:v86-convexity}
For all real plaquette link arguments
and all three variations,
\begin{equation}\label{eq:v86-third}
 |D^3s_p(a)[u_1,u_2,u_3]|
               \le64\prod_{j=1}^3\|u_j\|_{\infty,p}.
\end{equation}
Consequently, if
\[
        \max_{e\in E_p}|a_e|\le1/16,\qquad 0<b\le1/16,
\]
then
\begin{equation}\label{eq:v86-second-difference}
     s_p(a+b h^p)+s_p(a-b h^p)-2s_p(a)\ge8b^2.
\end{equation}
This assertion needs no positivity
of the gauge matrix at the shifted
inputs.
\end{lemma}
\begin{proof}
We give explicit constants in the
local derivative estimate underlying
V66's bounded-third-derivative statement.
In the two-dimensional matrix
representation, a Lie-algebra vector
\(u\) has operator norm \(|u|\), and
the exponential of a real such vector
is unitary. The \(k\)-th Fr\'echet
derivative of the exponential is a
sum over the \(k!\) orderings of
the inserted variations, integrated
over the ordered simplex of volume
\(1/k!\):
\[
 D^k e^A[X_1,\ldots,X_k]
   =\sum_{\pi\in S_k}\int_{\substack{t_j\ge0\\
                          t_0+\cdots+t_k=1}}
       e^{t_0A}X_{\pi(1)}e^{t_1A}\cdots
                X_{\pi(k)}e^{t_kA}\,dt_1\cdots dt_k .
\]
This formula follows by differentiating
the exponential differential equation,
or inductively from its first Duhamel
integral. Every exponential factor
in the integral is unitary. Hence
\[
                   \|D^k e^A[X_1,\ldots,X_k]\|
                                  \le\prod_j\|X_j\|.
\]
Inverse-oriented links are written
as \(\exp(-a_e)\), so the same bound
holds for them.

Differentiate the product of four
plaquette factors three times, with
three labelled variations. Each
variation may hit any of four
factors, giving \(4^3=64\) terms.
The derivative estimate just proved
bounds each term by the product
of the three maximum link norms.
The normalized absolute trace is
at most the operator norm, proving
\eqref{eq:v86-third}.

Along the segment \(t\mapsto a+t h^p\),
\(|t|\le b\), all four arguments have
norm at most \(1/8\).
Integrate \eqref{eq:v86-third} from
zero background to the input on
that segment and use
\eqref{eq:v86-zero-curve}:
\[
 \frac{d^2}{dt^2}s_p(a+t h^p)
       \ge16-64\max_{e\in E_p}|a_e+t h^p_e|
       \ge8.
\]
The exact symmetric second-difference
identity is
\[
 f(b)+f(-b)-2f(0)
           =\int_{-b}^{b}(b-|t|)f''(t)\,dt .
\]
The triangular weight has integral
\(b^2\). This proves
\eqref{eq:v86-second-difference}.
All derivatives concern the literal
nonlinear plaquette product.
\end{proof}

\subsection{A three-event nonconcentration principle}

The following elementary probability
argument is useful because it needs
total variation, not a density
derivative or an inverse second
moment at the gauge boundary.

\begin{lemma}[Two translations rule out concentration in a short interval]
\label{lem:v86-three-events}
Let \(\mu\) be a probability on a
Euclidean space, \(X\) a bounded
real observable, and \(h\) a
deterministic translation. Suppose
\[
 \|(T_{\pm h})_*\mu-\mu\|_{\rm TV}\le\epsilon_\pm,
                  \qquad \mu(G^c)\le\gamma,
\]
and
\[
             X(a+h)+X(a-h)-2X(a)\ge8b^2
                                      \quad\hbox{on }G.
\]
For every interval \(I\) of length
at most \(2b^2\),
\begin{equation}\label{eq:v86-interval}
        \mu\{X\in I\}
               \le\min\left\{1,\
                    \frac{2+\epsilon_++\epsilon_-+\gamma}{3}\right\}.
\end{equation}
Put
\(\delta=(1-\epsilon_+-\epsilon_--\gamma)_+/3\).
Then
\begin{equation}\label{eq:v86-variance}
 \inf_{c\in\mathbb R}\mathbb E_\mu|X-c|^q
                         \ge\delta b^{2q}\quad(q>0).
\end{equation}
In particular \(\operatorname{Var}_\mu X\ge\delta b^4\).
If \(X\ge0\), also
\begin{equation}\label{eq:v86-mean}
                         \mathbb E_\mu X\ge2\delta b^2.
\end{equation}
\end{lemma}
\begin{proof}
Let \(p_I=\mu\{X\in I\}\). The probabilities
of \(X(a+h)\in I\) and \(X(a-h)\in I\)
are at least \(p_I-\epsilon_+\) and
\(p_I-\epsilon_-\), respectively.
The probability that all three events
and \(G\) occur is therefore at least
\[
                    3p_I-2-\epsilon_+-\epsilon_--\gamma .
\]
But if all three values lie in the
same interval of length \(2b^2\),
their symmetric second difference
is at most \(4b^2\). This contradicts
the lower bound on \(G\). The preceding
probability lower bound must be
nonpositive, proving \eqref{eq:v86-interval}.

Apply it to \([c-b^2,c+b^2]\).
Its complement has probability at
least \(\delta\), and there
\(|X-c|>b^2\). This proves
\eqref{eq:v86-variance}. For \(q=2\)
minimization over \(c\) is the
variance. When \(X\ge0\), use instead
\([0,2b^2]\); on its complement
\(X>2b^2\), giving \eqref{eq:v86-mean}.
\end{proof}

\subsection{A finite-parameter lower bound under the unchanged YM law}

From now on all expectations are
under the actual \(\mu_{0,\beta}\).
Retain the V84 budgets with hats.
For \(0<b\le R/12\), define
\begin{equation}\label{eq:v86-transport-budget}
 \begin{aligned}
 \eta_b&=6(1+3R)b,\\
 \epsilon_b
 &=\min\Bigg\{1,\
 4\Bigg[\frac{\widehat m}{(R-6b)^2}
       +4\eta_b\widehat k
       +72\sqrt2\,\beta b\sqrt{\widehat\tau}
       +432\beta b^2+12Rb\Bigg]\Bigg\},\\
 \gamma&=\min\{1,1024\widehat m\},\qquad
              \delta_b=\frac{(1-2\epsilon_b-\gamma)_+}{3}.
 \end{aligned}
\end{equation}

\begin{theorem}[Actual nonlinear plaquette nonconcentration]
\label{thm:v86-finite}
For \(M\ge3\), \(0<R\le1/4\),
\(\beta\ge1\), and \(0<b\le R/12\),
every plaquette satisfies
\begin{equation}\label{eq:v86-finite}
 \begin{gathered}
 \sup_{I:\,|I|\le2b^2}\mu_{0,\beta}\{s_p\in I\}
                                          \le1-\delta_b,\\
 \mathbb E s_p\ge2\delta_b b^2,\qquad
 \operatorname{Var}W_p=\operatorname{Var}s_p
                                          \ge\delta_b b^4 .
 \end{gathered}
\end{equation}
The lower bound is informative
whenever \(2\epsilon_b+\gamma<1\).
\end{theorem}
\begin{proof}
Apply V85 with \(B=\pm bB^p\),
edge set \(E_p\) of size four,
and \(\theta=1/2\). Its theorem
gives
\[
        \|(T_{\pm b h^p})_*\mu_{0,\beta}
                              -\mu_{0,\beta}\|_{\rm TV}\le\epsilon_b.
\]
This uses the whole Wilson density,
including neighbouring plaquettes
touched by \(h^p\).
For
\[
                   G=\{\max_{e\in E_p}|a_e|\le1/16\},
\]
translation covariance and the
actual amplitude estimate give
\[
 \mu(G^c)\le\sum_{e\in E_p}
                         16^2\mathbb E|a_e|^2
                                      \le1024\widehat m .
\]
Since \(b\le R/12\le1/48<1/16\),
Lemma~\ref{lem:v86-convexity} supplies
the required second difference on
\(G\).
Apply Lemma~\ref{lem:v86-three-events}
with \(X=s_p\).
The variance of \(W_p=1-s_p\) is
identical to that of \(s_p\).
The translations used to compare
the laws are not asserted to be
gauge transformations; only the
tested observable is gauge invariant.
\end{proof}

\subsection{An explicit joint weak-coupling noncollapse scale}

Here are finite constants making the
asymptotic choice of \(b\) unambiguous.
Use V84's \(\kappa_R\) and
\(p_{R,13}>0\), and set
\begin{equation}\label{eq:v86-constants}
 \begin{gathered}
 H_\beta=\log(2+\beta),\qquad
 A_R=\frac92+3\log(\kappa_R^{-1})
            +\frac92\log(1+208eR^2)
            +18\log(p_{R,13}^{-1}),\\
 A_R^{\rm amp}=144\pi^2\sqrt{A_R},\qquad
 c_R^{\rm rad}=\frac{C_0R^2}{3},\\
 L_R=96(1+3R)c_R^{\rm rad}A_R^{\rm amp}
                                      +288\sqrt{2A_R},\qquad
 \sigma_R=\frac1{32(1+L_R)}.
 \end{gathered}
\end{equation}
All constants are positive and
independent of volume and coupling.
The fixed \(C_0\) is the same
proved Wilson derivative constant
as in V64 and V84.

For \(\beta\ge1\), \(b_{\rm norm}=12\beta+1\le13\beta\).
In the V84 entropy allowance,
\[
 \log(1+16eR^2b_{\rm norm})
          \le H_\beta+\log(1+208eR^2),\qquad H_\beta\ge1.
\]
Together with \(d_{\rm g}/V\le3\),
\(n/V\le9\), and
\(p_{R,b_{\rm norm}}\ge p_{R,13}\),
this proves the explicit estimate
\begin{equation}\label{eq:v86-explicit-amplitude}
 \widehat\tau\le A_R H_\beta/\beta,\qquad
 \widehat m\le A_R^{\rm amp}\sqrt{H_\beta/\beta}
                                      +72\pi^2/M^2
\end{equation}
for all sufficiently large \(\beta\)
in the second inequality, with
threshold depending only on \(R\).
This merely makes V84's constants
explicit; it is not a new
normalization comparison.

\begin{theorem}[A lower fluctuation bound in joint weak coupling]
\label{thm:v86-joint}
Fix \(0<R\le1/4\). Suppose
\begin{equation}\label{eq:v86-window}
 M\to\infty,\qquad \beta_M\to\infty,\qquad
                  \frac{\sqrt{\beta_M}}
                          {M^2\sqrt{\log(2+\beta_M)}}\to0 .
\end{equation}
For all sufficiently large \(M\),
put
\[
                     b_M=\frac{\sigma_R}
                            {\sqrt{\beta_M H_{\beta_M}}}.
\]
Then, uniformly over plaquettes,
\begin{equation}\label{eq:v86-joint}
 \begin{gathered}
 \sup_{I:\,|I|\le2b_M^2}\mu_{0,\beta_M}\{s_p\in I\}
                                                    \le\frac34,\\
 \mathbb E s_p\ge
           \frac{\sigma_R^2}{2\beta_M H_{\beta_M}},\qquad
 \operatorname{Var}W_p\ge
           \frac{\sigma_R^4}{4\beta_M^2H_{\beta_M}^2}.
 \end{gathered}
\end{equation}
Every sequence \(\beta_M=c\log M\)
with \(c>0\) is included.
\end{theorem}
\begin{proof}
Eventually \(b_M\le R/12\).
For \(b=\sigma/\sqrt{\beta H_\beta}\)
with fixed \(\sigma>0\), the untruncated
upper bound in
\eqref{eq:v86-transport-budget} obeys
\[
 \begin{aligned}
 \epsilon_b\le{}&
       \frac{16\widehat m}{R^2}
       +96(1+3R)b\widehat k\\
       &+288\sqrt2\,\beta b\sqrt{\widehat\tau}
                      +1728\beta b^2+48Rb.
 \end{aligned}
\]
The term \(b\widehat k\), using V84's
amplitude-paid inverse estimate and
\eqref{eq:v86-explicit-amplitude}, is
at most
\[
 \sigma\left\{
     \frac{12}{\sqrt{\beta H_\beta}}
       +c_R^{\rm rad}A_R^{\rm amp}
       +72\pi^2c_R^{\rm rad}
                   \frac{\sqrt\beta}{M^2\sqrt{H_\beta}}\right\}.
\]
Also
\[
 \beta b\sqrt{\widehat\tau}\le\sigma\sqrt{A_R},
            \qquad \beta b^2=\sigma^2/H_\beta.
\]
The radius and Haar terms vanish in
the joint limit, and the last term
in braces vanishes by
\eqref{eq:v86-window}. Thus
\[
                        \limsup\epsilon_b\le\sigma L_R.
\]
Choose \(\sigma=\sigma_R\).
Then \(\sigma_RL_R<1/32\), so
eventually \(\epsilon_{b_M}\le1/16\).
The amplitude estimate also gives
\(\gamma\le1/8\) eventually.
Hence
\[
          \delta_{b_M}\ge
          \frac{1-2(1/16)-1/8}{3}=\frac14 .
\]
Insert this in Theorem~\ref{thm:v86-finite}
to prove all claims.
\end{proof}

In particular the unscaled plaquette
variance is now bracketed, without
replacing the actual law, by
\[
 \frac{\sigma_R^4}{4\beta^2H_\beta^2}
        \le\operatorname{Var}W_p
        \le\mathbb E s_p^2
        \le2\mathbb E s_p
        \le2A_R H_\beta/\beta
\]
in the stated regime. The upper
bound is not of matching order.
Its role here is only to keep
the proved lower and upper claims
distinct.

\subsection{What this nonconcentration does not imply}

The concentration-function statement
is not merely a lower variance
produced by an uncontrolled rare
tail. For the scaled variable
\(Y_M=\beta_M H_{\beta_M}s_p\),
it excludes concentration of more
than \(3/4\) of the law in any
interval of length \(2\sigma_R^2\).
If this family has a weakly
convergent subsequence, its limit
cannot be a point mass: weak
convergence to a point mass would
put asymptotically all mass in
any fixed open interval about
that point, contrary to the
displayed bound.

Tightness of \(Y_M\) has not been
proved. Its available mean upper
bound is \(A_R H_{\beta_M}^2\),
which is not uniform. More
importantly \(Y_M\) is not the
same normalization as \(\beta_Ms_p\).
The following separate probability
example makes the logical gap
explicit. If \(X_H\) is uniform
on \(\{0,4/H,8/H\}\), then
\[
 \sup_{I:\,|I|\le2/H}\mathbb P\{X_H\in I\}=\frac13,\qquad
 \operatorname{Var}X_H=\frac{32}{3H^2},
 \qquad X_H\longrightarrow0 .
\]
Thus a nonconcentration bound at
a shrinking logarithmic scale
does not prevent collapse at
the thermal scale. This example
is not a model of the YM law.

No multiplication of the current
variance bound by a chosen field
normalization proves a nontrivial
continuum observable. The physical
lattice spacing, observable
renormalization, tightness and
cofinal limit remain separate
requirements.

\subsection{A precise upstream condition that would remove the logarithm}

The preceding proof also isolates
an actionable sufficient condition
for thermal-scale noncollapse.
It is stated conditionally rather
than inserted as a property of
the actual measure.

\begin{proposition}[Thermal noncollapse from an improved actual action bound]
\label{prop:v86-thermal}
Fix \(R\). Along a joint sequence
\(M,\beta_M\to\infty\), suppose that
there are constants \(A,C<\infty\)
such that
\begin{equation}\label{eq:v86-thermal-input}
       \frac1V\mathbb E_{\mu_{0,\beta_M}}S\le\frac A{\beta_M},
                         \qquad \frac{\beta_M}{M^4}\le C .
\end{equation}
Then there is \(c_{R,A,C}>0\) such
that, eventually and uniformly
over plaquettes,
\begin{equation}\label{eq:v86-thermal-output}
       \mathbb E s_p\ge\frac{c_{R,A,C}}{\beta_M},\qquad
       \operatorname{Var}W_p\ge\frac{c_{R,A,C}}{\beta_M^2}.
\end{equation}
The corresponding short-interval
nonconcentration also holds at
width proportional to \(\beta_M^{-1}\).
\end{proposition}
\begin{proof}
V81's deterministic action-to-amplitude
inequality can be integrated using
the actual first assumption in
\eqref{eq:v86-thermal-input}.
The same integer optimization gives
a new valid amplitude budget
\[
              m_A\le144\pi^2\sqrt A\,\beta^{-1/2}
                                  +72\pi^2M^{-2}
\]
eventually. It tends to zero.
The exact V79/V84 virial estimate
then gives a valid inverse budget
\[
 k_A\le12+c_R^{\rm rad}\beta m_A,\qquad
                \frac{k_A}{\sqrt\beta}\le K_{R,A,C}<\infty .
\]
For the last inequality use
\(\sqrt\beta/M^2\le\sqrt C\).

Revisit V85's proof with these
verified inputs, not its former
hatted upper bounds.
The actual local action expectation
is at most \(6|E_p|A/\beta\);
the local inverse trace is at
most \(|E_p|k_A\).
For \(b=\sigma/\sqrt\beta\),
its translation estimate therefore
gives
\[
 \limsup\epsilon_b
     \le96(1+3R)\sigma K_{R,A,C}
                       +288\sqrt{2A}\,\sigma+1728\sigma^2 .
\]
Choose a sufficiently small fixed
\(\sigma>0\) so this is below \(1/32\).
Eventually \(\epsilon_b\le1/16\)
and \(1024m_A\le1/8\).
The same second-difference and
three-event proof gives
\(\delta_b\ge1/4\).
Thus \(\mathbb E s_p\ge b^2/2\)
and \(\operatorname{Var}W_p\ge b^4/4\).
Taking the smaller of the two
positive constants proves
\eqref{eq:v86-thermal-output}.
\end{proof}

Only the improved action estimate
in \eqref{eq:v86-thermal-input}
is missing along logarithmic
coupling sequences; its volume
condition is then automatic.
This does not say that the
improved estimate is easy.
Here is a same-normalizer route
that would suffice to establish it.

\begin{lemma}[A dyadic free-energy gate]
\label{lem:v86-doubling}
For the original fixed chart and
geometric density, let \(Z_\beta\)
be its normalizer. At finite
volume,
\begin{equation}\label{eq:v86-doubling}
       \beta\,\mathbb E_{\mu_{0,\beta}}S
                  \le2\log\frac{Z_{\beta/2}}{Z_\beta}.
\end{equation}
Consequently a uniform bound
\[
                    \log\frac{Z_{\beta_M/2}}{Z_{\beta_M}}\le C_1V
\]
would supply the first assumption
of \eqref{eq:v86-thermal-input}
with \(A=2C_1\).
\end{lemma}
\begin{proof}
The support and geometric weight
are independent of the coupling,
and \(S\) is bounded at fixed
volume. Differentiation under
the integral gives
\[
 -\frac d{du}\log Z_u=\mathbb E_u S,\qquad
 \frac d{du}\mathbb E_uS=-\operatorname{Var}_uS\le0 .
\]
Hence
\[
 \log\frac{Z_{\beta/2}}{Z_\beta}
       =\int_{\beta/2}^{\beta}\mathbb E_uS\,du
       \ge\frac\beta2\,\mathbb E_\beta S .
\]
This proves the assertion without
a Gaussian replacement.
\end{proof}

V84 controls \(Z_0/Z_\beta\) by
an exponential whose logarithm
is of order \(V\log\beta\).
It has not proved the order-\(V\)
dyadic ratio needed here.
Applying its lower bound also
at \(\beta/2\) does not upper-bound
that ratio; an upper comparison
of the numerator, or another
action argument, is still required.
This is the upstream bottleneck
identified by the present
noncollapse calculation.
Even closing it would establish
thermal plaquette fluctuations,
not continuum existence or a
positive Hamiltonian mass gap.

\subsection{Exact checks and preservation}

The certificate is
\begin{center}\small
\path{scripts/verify_ym_f14_v86_plaquette_noncollapse.py}.
\end{center}
Six literal single-plaquette
directions, using all three
colours on tori of sides three
and four, verify the Coulomb
and mean-zero conditions, the
four-edge support, and
\((dh^p)_p=4v\).
Exact quaternion exponential
series reproduce the actual
single-pulse coefficients
\[
                 s_p(th^p)=8t^2-\frac{32}{3}t^4+O(t^6).
\]
Twelve noncommuting, mixed
third-derivative jets at zero
background are computed by
independent polynomial quaternion
multiplication; all twelve are
nonzero and satisfy the displayed
derivative bound. These finite
jets are regressions, not a
substitute for the global
unitary-simplex proof.

Further rational checks verify
the product-rule count 64, the
Hessian margin \(16-64/8=8\),
the radius-failure coefficient
1024, and the nonconcentration
margin \(1/4\).
Three separate finite cyclic
probability banks test 312
interval inequalities in the
abstract three-event argument,
retaining the locations where
the convexity condition fails.
Sixteen exact three-point
examples check the logarithmic
scale distinction. These
probability banks are explicitly
not substitutes for the
interacting Yang--Mills law.

The certificate contains 7,992 bytes and has SHA--256
\begin{center}\scriptsize\ttfamily
07D0D661827AA01DA29BB8F6D194A50EDA118F75607E5A848B80812B96DEFBDD.
\end{center}
The predecessor V85 contains 2,109,794 bytes and has SHA--256
\begin{center}\scriptsize\ttfamily
035E1E5ABE311D4A20FDC7AACA4334C126EC6D386F5B3226B85D88BE4A187FCB.
\end{center}
Its body is preserved except
for the title version and this
terminal insertion. V86 replaces
V85 as the sole active main note.
Earlier certificates, supplied
sources, companion notes,
monographs and archives are
unchanged.

\begin{firewall}
The proved lower bound concerns
a gauge-invariant plaquette
observable under the specified
restricted weak-coupling law.
It retains a logarithmic scale
loss and does not construct
a nontrivial interacting
continuum field or a positive
physical Hamiltonian mass gap.
The full objective remains
active and unproved.
\end{firewall}
\section{V87: selected dyadic free-energy bounds and nondegenerate thermal plaquette limits}
\label{sec:v87}

\subsection{Context and the change in direction}

V86 isolated an order-\(V\) dyadic
free-energy bound as a sufficient
upstream input for removing the
logarithmic loss in thermal
plaquette noncollapse. Its proof
did not establish that input
at a prescribed weak-coupling
sequence. This continuation
proves it at selected couplings
in an explicitly specified
window, using the actual
normalizer and its logarithmic
average. The selection is an
existence argument, not an
evaluation of the partition
function.

Combining this selection with
the already proved local
translation estimates gives
thermal-scale nonconcentration,
a uniform first-moment bound,
and a nondegenerate subsequential
law for the countable collection
of unit-lattice plaquette
energies. Unlike V86's
logarithmically rescaled family,
this family is proved tight.
No new Gaussian replacement,
radial homogeneity of the
Wilson action, or reflection
positivity is assumed.

The result remains about the
same restricted chart law.
It does not identify its limit
with continuum Yang--Mills or
construct a physical Hamiltonian.
The companion audit performed
at V85 remains applicable;
the next scheduled audit is V90.
The supplied source-selection,
stationarity and exact-action
notes remain separate inputs,
not assumptions that the
missing YM reconstruction
has been completed.

\subsection{Keeping the leading coefficient in the absolute pressure bound}

Fix \(0<R\le1/4\), \(M\ge3\),
and \(V=M^4\). In this section
write
\[
 Z_u=\int_{\Omega_R}e^{-uS(a)}
                      \det F(a)e^{-\mathcal H(a)}\,da,
             \qquad \mathbb E_u=\mathbb E_{\mu_{0,u}}.
\]
Thus \(Z_0\) means coupling zero
with the same geometric density;
it does not mean geometric
interpolation parameter zero.
The support, action and geometric
weight are independent of \(u\).
Set
\begin{equation}\label{eq:v87-pressure-constants}
 \begin{aligned}
 C_R^0&=3\log(\kappa_R^{-1})
       +\frac92\log(1+208eR^2)
       +18\log(p_{R,13}^{-1})
       =A_R-\frac92,\\
 C_R^1&=C_R^0+\frac92\log3 .
 \end{aligned}
\end{equation}
Here the constants on the right
are precisely those of V84 and
\eqref{eq:v86-constants}.

\begin{lemma}[Absolute pressure and its logarithmic average]
\label{lem:v87-pressure}
For \(u\ge1\),
\begin{equation}\label{eq:v87-pressure}
       0\le\frac1V\log\frac{Z_0}{Z_u}
           \le\frac92\log(2+u)+C_R^0.
\end{equation}
The continuous function
\[
                      q_M(u)=\frac{u}{V}\mathbb E_uS
\]
is nonnegative. For every \(B\ge2\),
\begin{equation}\label{eq:v87-average}
 \int_B^{B^2}q_M(u)\,\frac{du}{u}
        =\frac1V\log\frac{Z_B}{Z_{B^2}}
        \le9\log B+C_R^1.
\end{equation}
\end{lemma}
\begin{proof}
The two normalizer bounds of
Theorem~\ref{thm:v84-normalizer},
using \(D_M\ge1\), give
\[
 \log\frac{Z_0}{Z_u}
 \le d_{\rm g}\log(\kappa_R^{-1})
     +\frac n2\log(16eR^2(12u+1))
     +18V\log(p_{R,12u+1}^{-1}).
\]
Replace the logarithm in the
middle by the larger nonnegative
quantity
\(\log(1+16eR^2(12u+1))\).
For \(u\ge1\) it is at most
\(\log(2+u)+\log(1+208eR^2)\).
Now use \(d_{\rm g}/V\le3\),
\(n/(2V)\le9/2\), and
\(p_{R,12u+1}\ge p_{R,13}\).
This proves the upper bound
without absorbing its leading
logarithm into \(A_R\).
The lower bound follows from
\(S\ge0\).

At fixed volume the action is
bounded and \(Z_u>0\).
Differentiation under the integral
therefore gives
\[
 -(\log Z_u)'=\mathbb E_uS,\qquad
                (\mathbb E_uS)'=-\operatorname{Var}_uS\le0 .
\]
Integration proves the equality
in \eqref{eq:v87-average}.
Since \(Z_B\le Z_0\), its left
side is at most
\[
 \frac92\log(2+B^2)+C_R^0
       \le9\log B+\frac92\log3+C_R^0.
\]
This is the asserted bound.
\end{proof}

In particular,
\begin{equation}\label{eq:v87-minimum}
        \min_{u\in[B,B^2]}q_M(u)
                   \le9+\frac{C_R^1}{\log B}.
\end{equation}
There is also a quantitative
logarithmic-measure statement.
For \(A>0\),
\[
 \frac1{\log B}
 \int_{\{u\in[B,B^2]:q_M(u)>A\}}\frac{du}{u}
              \le\frac{9+C_R^1/\log B}{A}.
\]
This is just Markov's inequality
for the probability measure
\((\log B)^{-1}du/u\) on that
window. It is not a pointwise
bound at each coupling.

\subsection{Closing the dyadic gate at selected actual couplings}

\begin{theorem}[A selected order-volume dyadic cost]
\label{thm:v87-selection}
Let \(B\ge\max\{2,\exp(C_R^1)\}\).
Choose \(u_M(B)\) to be the
leftmost minimizer of \(q_M\)
on \([B,B^2]\), and put
\(\beta_M(B)=2u_M(B)\).
Then
\begin{equation}\label{eq:v87-selected-cost}
 \begin{gathered}
 2B\le\beta_M(B)\le2B^2,\\
 \log\frac{Z_{\beta_M(B)/2}}{Z_{\beta_M(B)}}\le10V,\qquad
       \frac{\beta_M(B)}{V}\mathbb E_{\beta_M(B)}S\le20 .
 \end{gathered}
\end{equation}
More generally \(q_M(t)\le20\)
throughout \(u_M(B)\le t\le2u_M(B)\).
\end{theorem}
\begin{proof}
Continuity on a compact interval
ensures a nonempty compact set
of minimizers, which has a
leftmost point. By
\eqref{eq:v87-minimum} it satisfies
\(q_M(u_M(B))\le10\).
Write this point as \(u\).
The already verified monotonicity
of \(\mathbb E_tS\) gives
\[
 \log\frac{Z_u}{Z_{2u}}
      =\int_u^{2u}\mathbb E_tS\,dt
      \le u\mathbb E_uS\le10V.
\]
It also gives, for \(u\le t\le2u\),
\[
             q_M(t)\le\frac{t}{u}q_M(u)\le20 .
\]
The bounds on \(\beta_M(B)=2u\)
follow from its definition.
\end{proof}

Take any \(B_M\to\infty\) with
\(B_M\le M^2\) eventually.
The theorem defines, for all
sufficiently large \(M\), an
actual coupling sequence
\begin{equation}\label{eq:v87-selected-window}
       2B_M\le\beta_M\le2B_M^2\le2M^4,\qquad
                \frac1V\mathbb E_{\beta_M}S\le\frac{20}{\beta_M}.
\end{equation}
These are precisely the
previously conditional inputs
of Proposition~\ref{prop:v86-thermal}
with \(A=20\) and \(C=2\).
For example, \(B_M=\log M\)
eventually yields
\[
                  2\log M\le\beta_M\le2(\log M)^2.
\]
The threshold can be large and
depends on \(R\). The construction
does not prove the same statement
at \(\beta_M=c\log M\) for a
prescribed \(c\), nor does it
produce an RG recursion or
compatibility between the
volume-dependent minimizers.

\subsection{Thermal nonconcentration with explicit uniform constants}

To record exactly what has now
become unconditional along the
selected sequences, define
\begin{equation}\label{eq:v87-thermal-constants}
 \begin{aligned}
 \chi_R&=c_R^{\rm rad}=C_0R^2/3,\\
 K_R^{\rm sel}
    &=12+\chi_R(144\pi^2\sqrt{20}+72\pi^2\sqrt2),\\
 L_R^{\rm sel}
    &=96(1+3R)K_R^{\rm sel}+288\sqrt{40},\\
 s_R^{\rm sel}
    &=\frac1{64(1+L_R^{\rm sel}+1728)} .
 \end{aligned}
\end{equation}
For brevity write \(s=s_R^{\rm sel}>0\)
in the remainder of this section.
These constants depend only on
the fixed chart radius and the
previously proved \(C_0\).

\begin{theorem}[Noncollapse at the thermal scale]
\label{thm:v87-thermal}
Along every selected sequence
\eqref{eq:v87-selected-window},
eventually and uniformly over
plaquettes,
\begin{equation}\label{eq:v87-thermal}
 \begin{gathered}
 \sup_{I:\,|I|\le2s^2/\beta_M}
          \mu_{0,\beta_M}\{s_p\in I\}\le\frac34,\\
 \frac{s^2}{2\beta_M}\le\mathbb E s_p
                                      \le\frac{20}{\beta_M},\qquad
        \operatorname{Var}W_p\ge\frac{s^4}{4\beta_M^2}.
 \end{gathered}
\end{equation}
Here \(s_p=1-\tfrac12\operatorname{ReTr}U_p\)
is the actual nonlinear plaquette
energy and \(W_p=1-s_p\).
\end{theorem}
\begin{proof}
For clarity the proof checks
the constants, rather than
merely invoking the conditional
proposition. Integrate V81's
deterministic amplitude estimate
using \(\mathbb E S/V\le20/\beta\).
Its integer optimization gives
the valid budget
\[
       m_{20}=144\pi^2\sqrt{20}\,\beta^{-1/2}
                                      +72\pi^2M^{-2}
\]
for \(\mathbb E Q/V\), eventually.
It tends to zero. V84's exact
virial bound then gives a valid
inverse budget
\[
 k_{20}=12+\chi_R\beta m_{20},\qquad
                 k_{20}/\sqrt\beta\le K_R^{\rm sel},
\]
because \(\beta\ge1\) and
\(\beta/M^4\le2\).

Set \(b=s/\sqrt\beta\).
Eventually \(b\le R/12\).
Use the proof of
\eqref{eq:v86-transport-budget}
with the verified budgets
\(m_{20},k_{20},20/\beta\).
The untruncated upper bound
on either translation error is
at most
\[
 \frac{16m_{20}}{R^2}
       +96(1+3R)b k_{20}
       +288\sqrt2\,\beta b\sqrt{20/\beta}
       +1728\beta b^2+48Rb.
\]
Consequently
\[
        \limsup\epsilon_b
             \le sL_R^{\rm sel}+1728s^2<\frac1{64}.
\]
Indeed \(s<1\) and the definition
in \eqref{eq:v87-thermal-constants}
makes the last strict inequality
immediate. Eventually
\(\epsilon_b\le1/16\) and
\(\gamma\le1024m_{20}\le1/8\).
The second-difference and
three-event argument in V86
therefore has
\[
              \delta_b\ge
                  \frac{1-2/16-1/8}{3}=\frac14 .
\]
It gives all the lower bounds
and the concentration-function
bound in \eqref{eq:v87-thermal}.
For the upper mean bound, fix
the orientation of \(p\).
Translation covariance makes
the expectation the same at
each of its \(V\) translates;
their nonnegative energy sum
is bounded by \(S\).
Thus \(\mathbb E s_p\le\mathbb E S/V\).
No rotation symmetry is needed.
\end{proof}

In particular \(X_{M,p}=\beta_Ms_p\)
satisfies
\begin{equation}\label{eq:v87-scaled}
       \frac{s^2}{2}\le\mathbb E X_{M,p}\le20,\qquad
       \sup_{I:\,|I|\le2s^2}\mathbb P\{X_{M,p}\in I\}
                                                        \le\frac34 .
\end{equation}
The finite-volume variance is
at least \(s^4/4\), but no
uniform second-moment upper
bound has been proved.
We will pass nonconcentration,
not an unjustified moment
convergence, to a limit.

\subsection{A countable, stationary, nondegenerate probe limit}

Let
\[
 \mathcal P_\infty
       =\mathbb Z^4\times\{(\mu,\nu):1\le\mu<\nu\le4\},
       \qquad \mathcal X=[0,\infty)^{\mathcal P_\infty}.
\]
Periodically extend the scalar
variables \(X_{M,p}\) from the
side-\(M\) torus to this index
set. Their simultaneous law
under the actual \(\mu_{0,\beta_M}\)
is a probability measure
\(\nu_M\) on \(\mathcal X\)
with its product topology.
This does not assume that
different plaquettes are
independent.

\begin{theorem}[Tightness and nondegenerate unit-lattice limits]
\label{thm:v87-probe-limit}
Along every sequence in
\eqref{eq:v87-selected-window},
the laws \(\nu_M\) are tight
after discarding finitely many
initial terms. Some subsequence
converges weakly to a probability
law \(\nu_\infty\) on \(\mathcal X\).
Every such limit is invariant
under \(\mathbb Z^4\) translations.
For each coordinate \(X_p\),
\begin{equation}\label{eq:v87-limit}
 \begin{gathered}
 \frac{s^2}{2}\le\mathbb E_{\nu_\infty}X_p\le20,\\
 \nu_\infty\{X_p\in J\}\le\frac34
       \quad\hbox{for every open interval \(J\) with \(|J|\le2s^2\)},\\
 \inf_{c\in\mathbb R}
      \mathbb E_{\nu_\infty}(X_p-c)^2\ge\frac{s^4}{4}.
 \end{gathered}
\end{equation}
The last expression is permitted
to be infinite. In particular
every coordinate marginal is
non-Dirac.
\end{theorem}
\begin{proof}
Enumerate \(\mathcal P_\infty\)
as \(p_1,p_2,\ldots\).
For \(\varepsilon>0\), set
\[
 r_j=\frac{20\,2^j}{\varepsilon},\qquad
                    K_\varepsilon=\prod_{j\ge1}[0,r_j].
\]
This product is compact in
\(\mathcal X\). For example,
one can prove sequential
compactness by the diagonal
subsequence argument in its
countable coordinates.
By Markov's inequality,
the union bound and
\eqref{eq:v87-scaled},
\[
 \nu_M(K_\varepsilon^c)
     \le\sum_{j\ge1}\nu_M\{X_{M,p_j}>r_j\}
     \le\sum_{j\ge1}\frac{20}{r_j}
     =\varepsilon .
\]
The space \(\mathcal X\) is
Polish: a complete compatible
metric is
\(\sum_{j\ge1}2^{-j}\min\{1,|x_{p_j}-y_{p_j}|\}\),
and finite rational coordinate
approximations are dense.
Prokhorov's theorem now gives
a weakly convergent subsequence.\footnote{For the
precise tightness-to-relative-compactness theorem,
see E. Sorensen, \emph{TA Session on tightness of
probability measures}, SLMath Stochastic Quantization
Summer School (1 July 2024), Theorem 1.1,
\href{https://www.mit.edu/~skycao/SLMath/Session_1_notes_(tightness).pdf}
{author's lecture notes}. The compact sets and the
hypotheses needed here were checked explicitly above.}

Each \(\nu_M\) is translation
invariant, because the original
chart law is translation invariant
and its periodic extension is
equivariant. A lattice shift
permutes the countably many
coordinates and is continuous.
Passing
\(\int f\circ T\,d\nu_M=\int f\,d\nu_M\)
to the limit for every bounded
continuous \(f\) proves
translation invariance of
\(\nu_\infty\).

For each \(L>0\), the function
\(\min\{X_p,L\}\) is bounded
continuous. Hence its limiting
expectation is at most 20.
Monotone convergence as
\(L\to\infty\) gives the stated
first-moment upper bound.

For an open interval \(J\)
as in the theorem, the cylinder
\(\{x:x_p\in J\}\) is open
in \(\mathcal X\).
The open-set direction of
Portmanteau and
\eqref{eq:v87-scaled} give
\[
 \nu_\infty\{X_p\in J\}
       \le\liminf_k\nu_{M_k}\{X_{M_k,p}\in J\}
       \le\frac34.
\]
This is why open intervals,
not the wrong closed-set
inequality, are used.
Choosing \(J=(c-s^2,c+s^2)\)
gives
\(\nu_\infty\{|X_p-c|\ge s^2\}\ge1/4\)
for every \(c\), and proves
the last inequality in
\eqref{eq:v87-limit}.
Finally for \(0<\eta<2s^2\),
use \(J=(-\eta,2s^2-\eta)\).
Since \(X_p\ge0\),
\[
 \mathbb E_{\nu_\infty}X_p
       \ge(2s^2-\eta)\,
                \nu_\infty\{X_p\ge2s^2-\eta\}
       \ge\frac{2s^2-\eta}{4}.
\]
Let \(\eta\downarrow0\).
This derives the positive
limiting mean without assuming
uniform integrability.
\end{proof}

This theorem constructs a
stationary scalar probe ensemble
on a unit lattice. Fixed
coordinate separations have
not been interpreted as fixed
physical distances. Neither
uniqueness of this limit nor
a local specification for it
has been proved. In particular,
stationarity and nondegeneracy
do not imply interacting
continuum YM, gauge-field
reconstruction, reflection
positivity, a physical transfer
operator, or a mass gap.

\subsection{Why the absolute bound does not close every prescribed coupling}

The selected-coupling restriction
is substantive. A separate
Laplace-transform example shows
that the present absolute
normalizer information alone
cannot remove it.

\begin{proposition}[A two-level obstruction to a pointwise pressure inference]
\label{prop:v87-obstruction}
For \(T\ge4\), \(p>0\), and
an integer \(V\ge1\), define
\[
 Z_{T,V}(u)=T^{-pV}
                 +\exp\{-u\,pV\log(T)/T\}.
\]
This is a positive two-level
Laplace transform with
nonnegative energies. It satisfies,
uniformly in \(T\), for \(u\ge1\),
\begin{equation}\label{eq:v87-toy-normalizer}
 Z_{T,V}(0)\le2,\qquad
                    Z_{T,V}(u)\ge e^{-pV}u^{-pV}.
\end{equation}
Nevertheless at \(u=T\),
\begin{equation}\label{eq:v87-toy-cost}
 \frac{T}{V}\mathbb E_T E=\frac p2\log T,\qquad
 \frac1V\log\frac{Z_{T,V}(T/2)}{Z_{T,V}(T)}
       \ge\frac p2\log T-\frac{\log2}{V}.
\end{equation}
Thus a Gaussian-power-type
absolute lower bound, even with
a volume-only exponential
constant, does not imply a
uniform order-\(V\) dyadic cost
at an arbitrary prescribed
coupling.
\end{proposition}
\begin{proof}
Give energy zero the unnormalized
weight \(T^{-pV}\), and energy
\(E_1=pV\log(T)/T\) weight one.
This yields the displayed
Laplace transform and its
upper bound at zero.
If \(u\ge T/2\), its ground
term alone gives
\[
 T^{-pV}\ge2^{-pV}u^{-pV}
                              \ge e^{-pV}u^{-pV}.
\]
For \(1\le u\le T/2\), consider
\(f(u)=\log u-u\log(T)/T\).
It is concave. Its values at
the two endpoints obey
\[
 f(1)=-\log(T)/T\ge-1,\qquad
 f(T/2)=\tfrac12\log T-\log2\ge0 .
\]
Concavity puts its minimum on
this interval at an endpoint.
Thus \(f(u)\ge-1\), which
rearranges to
\[
             e^{-u pV\log(T)/T}\ge e^{-pV}u^{-pV}.
\]
This proves \eqref{eq:v87-toy-normalizer}.

At \(u=T\) the two terms are
equal, so the two energy
probabilities are both \(1/2\).
This gives the first identity
in \eqref{eq:v87-toy-cost}.
Direct evaluation gives
\[
 \frac{Z_{T,V}(T/2)}{Z_{T,V}(T)}
                          =\frac{1+T^{pV/2}}2,
\]
and hence the second inequality.
Both normalized costs diverge
as \(T\to\infty\).
\end{proof}

One may take \(p=9/2\) to
match the power per volume
in \eqref{eq:v87-pressure}.
Equivalently this example
satisfies the uniform absolute
pressure upper bound
\[
 \frac1V\log\frac{Z_{T,V}(0)}{Z_{T,V}(u)}
                 \le p\log u+p+\frac{\log2}{V}.
\]
It is not a Yang--Mills measure
and does not show that the
actual law has these bad
couplings. It proves only
that such a conclusion cannot
be excluded by the absolute
pressure estimates alone.
The selection theorem uses
their valid averaged content.

\subsection{Checks, preservation, and the remaining upstream work}

The certificate is
\begin{center}\small
\path{scripts/verify_ym_f14_v87_selected_pressure_limits.py}.
\end{center}
It checks sixteen dyadic
integral brackets for four
finite nonnegative-energy
Laplace transforms, with exact
partition functions at multiples
of \(\log2\), and four exact
telescoping products. Logarithms
are enclosed by V82's certified
rational series intervals.
Twenty-four rational two-level
examples verify the obstruction's
dyadic ratio and equal-weight
calculation, with eight endpoint
checks for its uniform lower
bound. Fifteen compact-product
tail sums and the explicit
translation-error margins are
also checked.

These are algebraic diagnostics.
They neither sample the actual
YM normalizer nor replace the
analytic proofs of selection,
local transport and probability
compactness. No machine-checked
proof of the full YM programme
is being claimed.

The certificate contains 5,200 bytes and has SHA--256
\begin{center}\scriptsize\ttfamily
D36FDBF5E6A69BEE142C370006144617F1E05EA7F20808469BAF6F3F2AC01B08.
\end{center}
The predecessor V86 contains 2,131,895 bytes and has SHA--256
\begin{center}\scriptsize\ttfamily
FF31AB4FB878A8442FCEB4EA164DCAA4970A95A526835AE055E1DF527103CCAD.
\end{center}
Its body is preserved except
for the title version and this
terminal insertion. V87 replaces
V86 as the sole active main note.
Earlier certificates, supplied
sources, companion notes,
monographs and archives are
unchanged.

The next substantive problem
is not to repeat the compactness
argument. It is to obtain
additional control tied to
physical scale: a compatible
coupling trajectory, renormalized
observables at fixed physical
separation, and a reconstruction
that retains the interacting
theory. Even a pointwise
improvement of the dyadic bound
would not by itself supply
these requirements. The
two-level obstruction explains
why more structural information
than the present absolute
normalizer estimates is needed
for that particular improvement.

\begin{firewall}
The new unconditional result
is selected-coupling thermal
noncollapse and a tight,
nondegenerate unit-lattice
plaquette-energy limit for the
specified restricted law.
It is not a full continuum
Yang--Mills construction and
does not prove a positive
physical Hamiltonian mass gap.
The original full objective
remains active and unproved.
\end{firewall}
\section{V88: thermal curvature fields, exact cube transport, and the limiting Bianchi constraint}
\label{sec:v88}

\subsection{Context and the new upstream question}

V87 made concrete progress: at selected
actual couplings it removed the logarithmic
loss from plaquette noncollapse and
proved tightness of the scalar thermal
energies. A scalar ensemble, however,
does not yet retain the compatibility
relations of a field strength.
This continuation lifts those energies
to the actual Lie-algebra-valued
plaquette logarithms and proves a
quantitative Bianchi-defect estimate.
Every subsequential thermal curvature
limit is a nonzero closed lattice
two-form, not merely a collection of
nondegenerate scalar marginals.

The exact non-Abelian cube product is
used before taking any limit. Its
parallel transports are not dropped
formally. Their errors are paid by the
proved local amplitude expectation.
This is distinct from V13's inherited
classical minimum-branch convergence
and V59's fixed-volume perturbative
coefficients: the laws here are the
actual selected-coupling laws of V87.
The chart radius stays fixed.

We also quantify how far the transport
estimate extends along longer paths.
Finally two separate Gaussian examples
show why closedness and one-point
noncollapse do not yet control
physical-scale smearing. They identify
the missing correlation information;
they are not proposed replacements
for the interacting law.
The next companion checkpoint remains V90.

\subsection{Actual logarithmic curvature and its second-moment budget}

Fix \(0<R\le1/4\), and use a selected
sequence from \eqref{eq:v87-selected-window}.
Write \(\beta=\beta_M\), and retain
\(s=s_R^{\rm sel}>0\) from
\eqref{eq:v87-thermal-constants}.
Use the Lie-algebra identification
\(\iota:\mathbb R^3\to\mathfrak{su}(2)\)
in which
\[
 \iota(v)^2=-|v|^2I,\qquad
 \|\iota(v)\|_{\rm op}=|v|,\qquad
                       U_\mu(x)=e^{\iota(a_\mu(x))}.
\]
For an ordered pair \(i\ne j\), define
the based plaquette
\[
 P_{ij}(x)=U_i(x)U_j(x+e_i)
                    U_i(x+e_j)^{-1}U_j(x)^{-1}.
\]
In particular \(P_{ji}(x)=P_{ij}(x)^{-1}\).
All link radii are below \(R\).
The geodesic triangle inequality on
\(\mathrm{SU}(2)\), with the displayed
norm convention, puts the angle of
\(P_{ij}(x)\) below \(4R\le1<\pi\).
Equivalently concatenate the four
exponential paths; their total length
is at most the sum of the four link
norms. Thus its principal logarithm
is unambiguous:
\begin{equation}\label{eq:v88-curvature}
 P_{ij}(x)=e^{\iota(z_{ij}(x))},\qquad
 |z_{ij}(x)|\le1,\qquad
 z_{ji}(x)=-z_{ij}(x),\qquad
                   Y_{M,ij}(x)=\sqrt\beta\,z_{ij}(x).
\end{equation}
These are gauge-fixed colour variables.
Their squared norm is gauge invariant,
but the vector itself is not asserted
to be a gauge-invariant observable.

\begin{lemma}[Thermal curvature control]
\label{lem:v88-curvature}
Eventually, for every plaquette,
\begin{equation}\label{eq:v88-curvature-budget}
 \mathbb E|Y_{M,p}|^2\le80,\qquad
 0\le\frac12|Y_{M,p}|^2-X_{M,p}
                 \le\frac{|Y_{M,p}|^4}{24\beta},
               \quad X_{M,p}=\beta s_p .
\end{equation}
For each fixed \(p\), the difference
in the middle tends to zero in
probability. Also \(\mathbb E Y_{M,p}=0\).
\end{lemma}
\begin{proof}
For \(0\le r\le1\),
\[
 \frac{r^2}{4}\le
       \frac{r^2}{2}-\frac{r^4}{24}
       \le1-\cos r\le\frac{r^2}{2}.
\]
The middle inequalities follow from
Taylor's formula, or two successive
integrations of the corresponding
cosine bounds. Since
\(s_p=1-\cos|z_p|\), V87's
\(\mathbb E X_{M,p}\le20\) proves
the first inequality. The same
Taylor remainder gives the second.
For any \(K>0\), on \(|Y_{M,p}|\le K\)
the difference is at most
\(K^4/(24\beta)\), whereas
\[
                  \mathbb P\{|Y_{M,p}|>K\}\le80/K^2 .
\]
First let \(M\to\infty\) at fixed
\(K\), then \(K\to\infty\).
No fourth-moment bound is used.

The actual law is invariant under
global colour conjugation: the
Coulomb and mean-zero constraints,
link radii, Wilson action, Haar
density and determinant are all
unchanged by a common adjoint
rotation of the colour coordinates.
The Euclidean measure is unchanged
by this orthogonal transformation.
The principal logarithm transforms
by the same rotation. Its integrable
mean is therefore fixed by every
rotation in \(\mathrm{SO}(3)\),
and must be zero.
\end{proof}

In particular we have not replaced
\(\sqrt\beta\,z_p\) by
\(\sqrt\beta\,(da)_p\).
The available amplitude estimate
does not make all quadratic
commutator contributions vanish
at this normalization.

\subsection{The six-factor cube identity, with all connectors}

For three distinct directions \(i,j,k\),
let the path product from \(x\) to
\(x+e_i+e_j+e_k\) be
\[
 T_{ijk}(x)=U_i(x)U_j(x+e_i)
                              U_k(x+e_i+e_j).
\]
Consider the cycle of six orderings
\[
             ijk,\quad jik,\quad jki,\quad
                    kji,\quad kij,\quad ikj,\quad ijk.
\]
Consecutive path ratios telescope.
Writing them as based face holonomies
gives the following exact identity.
Parallel transport in non-Abelian
plaquette constraints is standard;
the particular ordering below is
derived directly, independently
of a plaquette change of variables.\footnote{
For context on connectors in non-Abelian
plaquette formulations, see O. Borisenko,
S. Voloshin and M. Faber,
\emph{Plaquette representation for 3D lattice
gauge models: I. Formulation and perturbation theory},
\href{https://arxiv.org/abs/hep-lat/0508003}
{arXiv:hep-lat/0508003}. No perturbative
uniformity claim from that work is used here.}

\begin{lemma}[Exact cube product and a logarithmic remainder]
\label{lem:v88-cube}
Define
\begin{equation}\label{eq:v88-cube-factors}
 \begin{aligned}
 g_1&=P_{ij}(x),&
 g_2&=U_j(x)P_{ik}(x+e_j)U_j(x)^{-1},\\
 g_3&=P_{jk}(x),&
 g_4&=U_k(x)P_{ji}(x+e_k)U_k(x)^{-1},\\
 g_5&=P_{ki}(x),&
 g_6&=U_i(x)P_{kj}(x+e_i)U_i(x)^{-1}.
 \end{aligned}
\end{equation}
Then \(g_1g_2g_3g_4g_5g_6=I\).
If \(g_r=e^{Z_r}\), where \(Z_r\)
are their principal skew-Hermitian
logarithms, then
\begin{equation}\label{eq:v88-log-remainder}
 \left\|\sum_{r=1}^6Z_r\right\|_{\rm op}
       \le\frac12\left(\sum_{r=1}^6\|Z_r\|_{\rm op}\right)^2
       \le3\sum_{r=1}^6\|Z_r\|_{\rm op}^2 .
\end{equation}
\end{lemma}
\begin{proof}
The first factor is \(T_{ijk}T_{jik}^{-1}\);
the second is \(T_{jik}T_{jki}^{-1}\);
and so on around the six-ordering
cycle. The common final link cancels
in the ratios exchanging the first
two directions. The common initial
link remains as a conjugation when
exchanging the last two directions.
This gives exactly
\eqref{eq:v88-cube-factors}.
Their product telescopes to \(I\).
It is not generally true after
removing the conjugations.

Put \(F(t)=e^{tZ_1}\cdots e^{tZ_6}\)
for \(0\le t\le1\). All factors
are unitary, and the product rule
gives
\[
 F'(0)=\sum_rZ_r,\qquad
             \|F''(t)\|_{\rm op}
                  \le\left(\sum_r\|Z_r\|_{\rm op}\right)^2 .
\]
The second derivative bound counts
both insertions into one factor
and the two orders of insertions
into different factors.
Since \(F(0)=F(1)=I\), Taylor's
formula with integral remainder
yields the first bound in
\eqref{eq:v88-log-remainder}.
Cauchy--Schwarz gives the second.
Commutativity of the \(Z_r\) is
neither needed nor assumed.
\end{proof}

\subsection{A quantitative actual-law Bianchi defect}

Use forward differences
\(D_i f(x)=f(x+e_i)-f(x)\), and set
\[
 (dY)_{ijk}(x)=D_iY_{jk}(x)
                     -D_jY_{ik}(x)+D_kY_{ij}(x)
                   \qquad(i<j<k).
\]
Recall V87's valid amplitude budget
\[
     m_{20}=144\pi^2\sqrt{20}\,\beta^{-1/2}
                                             +72\pi^2M^{-2}.
\]
In particular each link has
\(\mathbb E|a_e|^2\le m_{20}\),
by translation covariance and
the nonnegativity of \(Q\).

\begin{theorem}[Vanishing linear Bianchi defect]
\label{thm:v88-bianchi}
Along the selected actual laws,
eventually and uniformly in elementary
cubes,
\begin{equation}\label{eq:v88-bianchi}
 \mathbb E|(dY_M)_{ijk}(x)|
              \le\frac{1440}{\sqrt\beta}
                            +6\sqrt{80m_{20}}
                        \longrightarrow0 .
\end{equation}
\end{theorem}
\begin{proof}
Principal logarithms commute with
conjugation and inversion on this
angle domain. In vector coordinates,
the six logarithms in
\eqref{eq:v88-cube-factors} are
\[
 \begin{gathered}
 z_{ij}(x),\quad
 \operatorname{Ad}_{U_j(x)}z_{ik}(x+e_j),\quad
 z_{jk}(x),\\
 -\operatorname{Ad}_{U_k(x)}z_{ij}(x+e_k),\quad
 -z_{ik}(x),\quad
 -\operatorname{Ad}_{U_i(x)}z_{jk}(x+e_i).
 \end{gathered}
\]
On removing the three adjoint
rotations their sum is exactly
\(-(dz)_{ijk}(x)\). Keep the
errors of that removal.
For any \(a,z\in\mathbb R^3\),
\begin{equation}\label{eq:v88-adjoint}
             |\operatorname{Ad}_{e^{\iota(a)}}z-z|
                                     \le2|a|\,|z|.
\end{equation}
Indeed the adjoint generator has
norm at most \(2|a|\); integrate
its derivative along the unitary
conjugation path.
The cube lemma and this bound give
\[
 |(dY_M)_{ijk}(x)|
   \le3\sqrt\beta\sum_{\text{six faces }p}|z_p|^2
       +2\sqrt\beta
          \sum_{\text{three connectors }(e,p)}|a_e|\,|z_p|.
\]
The first sum has expectation
at most \(3\cdot6\cdot80/\sqrt\beta\)
by \eqref{eq:v88-curvature-budget}.
For each connector, Cauchy--Schwarz,
without an independence assumption,
gives
\[
 \sqrt\beta\,\mathbb E(|a_e|\,|z_p|)
      \le(\mathbb E|a_e|^2)^{1/2}
                   (\mathbb E|Y_{M,p}|^2)^{1/2}
      \le\sqrt{80m_{20}}.
\]
This proves the displayed estimate.
Both terms tend to zero in the
joint limit.
\end{proof}

This is a consequence of the
non-Abelian identity and proved
expectations, not an assumption
that the finite-volume curvature
is an abelian curl.

\subsection{Nonzero closed thermal curvature limits}

Periodically extend \(Y_{M,p}\) and
\(X_{M,p}\) to the countable
plaquette index set
\(\mathcal P_\infty\) of V87.
We retain their joint law, rather
than selecting unrelated energy
and colour subsequences.

\begin{theorem}[A compatible nondegenerate curvature ensemble]
\label{thm:v88-limit}
The joint laws of
\((Y_M,X_M)\) are tight in
\[
          (\mathbb R^3\times[0,\infty))^{\mathcal P_\infty}.
\]
Every subsequential weak limit
\((Y,X)\) is stationary under
\(\mathbb Z^4\) translations,
invariant under common
\(\mathrm{SO}(3)\) colour rotations,
and satisfies, almost surely,
\begin{equation}\label{eq:v88-limit-constraints}
            dY=0,\qquad X_p=\tfrac12|Y_p|^2
                           \quad(p\in\mathcal P_\infty).
\end{equation}
For each \(p\),
\begin{equation}\label{eq:v88-limit-moments}
       \mathbb E Y_p=0,\qquad
               s^2\le\mathbb E|Y_p|^2\le40 .
\end{equation}
Moreover for every open interval
\(J\) of length at most \(2s^2\),
\begin{equation}\label{eq:v88-limit-noncollapse}
                 \mathbb P\{\tfrac12|Y_p|^2\in J\}\le\frac34 .
\end{equation}
In particular the curvature
ensemble is not the zero field.
No assertion that this law is
Gaussian or unique is made.
\end{theorem}
\begin{proof}
Enumerate the plaquettes by \(p_j\).
For \(\varepsilon>0\) choose
\(K_j^2=80\,2^j/\varepsilon\).
The second-moment estimate and
the union bound put probability
at least \(1-\varepsilon\) in
\[
       \prod_{j\ge1}
             \left(\overline B_{\mathbb R^3}(0,K_j)
                                       \times[0,K_j^2/2]\right),
\]
because \(0\le X_{M,p}\le|Y_{M,p}|^2/2\).
This is a compact subset of the
countable product Polish space.
Apply the same tightness theorem
used in V87.
Stationarity and common colour
invariance pass to the limit
under their continuous coordinate
actions.

For each coordinate the discrepancy
\(\tfrac12|Y_{M,p}|^2-X_{M,p}\)
tends to zero in probability,
by Lemma~\ref{lem:v88-curvature}.
Its absolute value truncated at
one is a bounded continuous
function of the joint coordinate.
Its expectation tends to zero.
Weak convergence therefore proves
the second identity in
\eqref{eq:v88-limit-constraints}.
Similarly \(\min\{1,|(dY_M)_{ijk}(x)|\}\)
has expectation tending to zero
by Theorem~\ref{thm:v88-bianchi}.
It is a bounded continuous
function of six coordinates.
Thus \(dY=0\) on each cube,
almost surely. There are
countably many cubes and
coordinates, so all these
identities hold simultaneously.

The marginal law of \(X\) is
a V87 subsequential energy
limit. Its bounds
\(s^2/2\le\mathbb E X_p\le20\)
and open-interval nonconcentration
give \eqref{eq:v88-limit-moments}
and \eqref{eq:v88-limit-noncollapse},
using the pathwise identity.
The finite first moment of
\(Y_p\) follows from its finite
second moment. Common rotation
invariance then forces its mean
to vanish, as at finite volume.
This proof does not assume
convergence of second moments.
\end{proof}

The vector-valued compatibility
is new information beyond V87,
but \(dY=0\) is a kinematic
constraint. It is not the
Yang--Mills equation, an
integration-by-parts identity
for the limiting law, or
a physical spectral estimate.

\subsection{What the same estimates pay along longer paths}

\begin{proposition}[A mesoscopic transport range]
\label{prop:v88-path}
Let \(\gamma\) be any deterministic
oriented link path of length \(N\)
on the torus, and \(U_\gamma\)
its ordered holonomy. For any
plaquette \(p\), without requiring
disjointness or independence,
\begin{equation}\label{eq:v88-path}
 \mathbb E|\operatorname{Ad}_{U_\gamma}Y_{M,p}-Y_{M,p}|
       \le\min\{2\sqrt{80},\,2N\sqrt{80m_{20}}\}.
\end{equation}
In particular the error tends to
zero whenever
\[
                     N=o(\min\{\beta_M^{1/4},M\}).
\]
\end{proposition}
\begin{proof}
Write the path's adjoint action
as a product of orthogonal
rotations. Telescope that
product minus the identity.
Equation~\eqref{eq:v88-adjoint}
and invariance of the vector
norm give, pointwise,
\[
 |\operatorname{Ad}_{U_\gamma}Y-Y|
                       \le2|Y|\sum_{e\in\gamma}|a_e|.
\]
A reversed edge has the same
logarithmic norm; repeated edges
are counted with multiplicity.
Apply Cauchy--Schwarz to each
summand, with the same local
expectation bounds as above.
The alternative bound is
\(|\operatorname{Ad}_{U_\gamma}Y-Y|\le2|Y|\).
Finally
\(\sqrt{m_{20}}=O(\beta^{-1/4}+M^{-1})\).
\end{proof}

A path spanning a fixed physical
distance when the lattice spacing
tends to zero need not satisfy
this mesoscopic condition.
For instance the budget does not
vanish for paths of order \(M\)
across a torus with fixed
physical side length.
This is a limitation of the
estimate, not a proof that
actual long-path transport
fails to converge. The
one-cube argument cannot simply
be repeated over a physical
distance without paying its
accumulated error.

\subsection{Closedness and local noncollapse do not determine physical smearing}

Here is an explicit diagnostic
for the remaining covariance
problem. It concerns separate
Gaussian fields on the unit
lattice, not a claim about the
law in Theorem~\ref{thm:v88-limit}.
Let \(A_i(x)\), for all links,
be independent \(N(0,I_3)\)
vectors and define
\[
              Y^{\rm curl}_{ij}(x)
                        =D_iA_j(x)-D_jA_i(x).
\]
Alternatively take six independent
\(N(0,I_3)\) vectors \(G_{ij}\)
and put
\[
                        Y^{\rm const}_{ij}(x)=2G_{ij}
                    \quad\hbox{for every }x.
\]
Extend both antisymmetrically.

\begin{proposition}[Identical local curvature laws, incompatible scaling]
\label{prop:v88-scaling-diagnostic}
Both fields are stationary,
common-colour-rotation invariant
and closed under \(d\).
Each plaquette has precisely
the same law \(N(0,4I_3)\).
Its energy \(|Y_p|^2/2\) has
mean six, variance 24, and
satisfies the nonconcentration
bound in \eqref{eq:v88-limit-noncollapse}
for the value of \(s\) used here.

Nevertheless, fix one colour
component and the orientation
\((1,2)\), and for a real
\(f\in C_c^\infty(\mathbb R^4)\)
define the candidate dimension-two
smearing
\begin{equation}\label{eq:v88-smearing}
          \Phi_a(f)=a^2\sum_{x\in\mathbb Z^4}
                                     f(ax)Y_{12}(x),
                     \qquad a\downarrow0 .
\end{equation}
For \(Y^{\rm curl}\),
\(\mathbb E|\Phi_a(f)|^2=O_f(a^2)\),
so it collapses in \(L^2\).
For \(Y^{\rm const}\), if
\(\int f\ne0\), its variance
diverges like
\(4a^{-4}(\int f)^2\), and
the scalar laws are not tight.
\end{proposition}
\begin{proof}
The identity \(d^2=0\) proves
closedness of the curl field;
the constant field is closed
directly. Each curl plaquette
is a signed sum of four
independent standard Gaussian
vectors, proving the claimed
common marginal. All other
invariances follow from the
definitions.
The common energy has density
\[
                  h(t)=\frac{\sqrt t\,e^{-t/4}}{4\sqrt\pi},
                               \qquad t>0 .
\]
Its maximum is attained at
\(t=2\) and is less than one.
An interval of length \(2s^2\)
therefore has probability at
most \(2s^2<3/4\), since
\(s<1/2\). The elementary
chi-square mean and variance
give six and 24.

Write \(f_a(x)=f(ax)\).
Discrete summation by parts
for the curl field gives
\[
 \Phi_a(f)=a^2\sum_x
       \{[f_a(x-e_1)-f_a(x)]A_2(x)
            -[f_a(x-e_2)-f_a(x)]A_1(x)\}.
\]
Independence of the scalar link
components gives exactly
\begin{equation}\label{eq:v88-curl-smearing}
 \operatorname{Var}\Phi_a(f)
       =a^4\sum_x\sum_{i=1}^2
                         [f_a(x-e_i)-f_a(x)]^2 .
\end{equation}
There are \(O_f(a^{-4})\)
nonzero terms, each \(O_f(a^2)\),
so the result is \(O_f(a^2)\).
For the constant field,
\[
 \operatorname{Var}\Phi_a(f)
                 =4a^4\left(\sum_x f(ax)\right)^2
                 \sim4a^{-4}\left(\int f\right)^2
\]
by the ordinary Riemann-sum limit.
These variables are centered
Gaussian. When their variance
diverges, the probability of
any fixed bounded interval
tends to zero; hence they
are not tight.
\end{proof}

Thus the constraints now proved
for the actual thermal field
do not by themselves distinguish
a collapsing dimension-two
smearing from a non-tight one.
The two examples even share
their full one-plaquette laws,
not merely their moments.
Neither example identifies
the actual limiting field.
Nor does the diagnostic assert
that \eqref{eq:v88-smearing}
is already a renormalized,
gauge-invariant continuum YM
observable.

For any centered stationary
finite-second-moment scalar
component with covariance
\(C(x-y)\), the exact variance
of this smearing is
\begin{equation}\label{eq:v88-covariance-target}
 \operatorname{Var}\Phi_a(f)
       =a^4\sum_{x,y}f(ax)f(ay)C(x-y).
\end{equation}
Controlling this quadratic form
at shrinking physical lattice
spacing requires information
across growing separations.
The local Bianchi proof does
not supply it, and an ordinary
gauge-operator inverse bound
is not automatically a bound
on this physical curvature
covariance.

\subsection{Exact verification and preservation}

The certificate is
\begin{center}\small
\path{scripts/verify_ym_f14_v88_curvature_bianchi.py}.
\end{center}
It checks all six path-ratio
identities for each of the
four coordinate triples in
four dimensions, using free
word reduction. Twenty-four
rational, genuinely noncommuting
unit-quaternion cubes separately
verify the full product and
all path ratios. Dropping the
connectors fails in every one
of those twenty-four cases.
The four linearized cube
boundaries cancel exactly.

Seven tensor-product tent
smearings, at \(a=1/n\), verify
the diagnostic variances
\[
       \operatorname{Var}\Phi_{1/n}^{\rm curl}
                   =\frac{4(2n^2+1)^3}{27n^8}
                   \le\frac4{n^2},\qquad
       \operatorname{Var}\Phi_{1/n}^{\rm const}=4n^4 .
\]
Three are independently checked
by full four-dimensional finite
summation. These tests certify
the finite algebra, not a
numerically inferred YM limit.
The analytic expectation and
compactness proofs above remain
essential.

The certificate contains 6,768 bytes and has SHA--256
\begin{center}\scriptsize\ttfamily
86F801F10FF66BE89F51FA40B9AFD55ABB0C591D655A3119A7028881378E1633.
\end{center}
The predecessor V87 contains 2,151,034 bytes and has SHA--256
\begin{center}\scriptsize\ttfamily
8A79C730DD5292614DE6BC5D626F99D750EE0DCDD351FF99EE83AC10A4122EC9.
\end{center}
Its body is preserved except
for the title version and this
terminal insertion. V88 replaces
V87 as the sole active main note.
Earlier certificates, supplied
sources, companion notes,
monographs and archives remain
unchanged.

\begin{firewall}
The actual selected-coupling
curvature limits are now proved
nonzero, stationary and closed
as lattice two-forms, with a
quantified mesoscopic transport
range. These are field
compatibility and local
normalization results, not
an interacting continuum
reconstruction. Uniform
physical-scale covariance
control, the relevant
renormalized observables,
reflection positivity and
a positive physical Hamiltonian
mass gap remain unproved.
The full original objective
remains active.
\end{firewall}
\section{V89: an action-paid gauge barrier and exact Gaussian identities for local curvature probes}
\label{sec:v89}

\subsection{Context: beyond the kinematic constraint}

V88 established actual thermal curvature
limits satisfying a linear Bianchi
identity, but correctly left their
covariance undetermined. This section
adds a measure identity. It first
improves the adverse radial Wilson
cost by using the small action as
well as the link amplitude. The
resulting inverse-trace bound makes
the geometric density change vanish
under a fixed local coexact link
move of size \(\beta^{-1/2}\).
The Wilson density change itself
does not vanish: it converges to
the quadratic energy increment
of the thermal curvature.

An exact bounded-acceptance
change-of-variables identity then
passes to the limit. Dividing its
strictly positive limiting weights
gives the Gaussian exponential
translation formula. This proves
Gaussianity and an exact covariance
for all compact curvature probes
generated by \(dd^*\).
It is a property of limits of
the specified actual measures,
not of a substituted Gaussian
reference.

The proof does not infer
exponential integrability from
weak convergence. That
integrability is obtained only
after the limiting measure
identity has been proved.
Nor does the result identify
the whole limiting field or
its physical-scale correlations.
The companion review remains
scheduled for V90.

\subsection{A radial Wilson bound using the actual action}

Keep \(S=\sum_p s_p\), \(Q=\sum_e|a_e|^2\),
\(V=M^4\), and the actual chart law
\(\mu=\mu_{0,\beta}\).
The following deterministic estimate
is different from V64's Hessian
lower form.

\begin{lemma}[Action-amplitude control of radial work]
\label{lem:v89-radial}
For every field in the chart,
\begin{equation}\label{eq:v89-radial}
                  |DS(a)[a]|\le4\sqrt3\,\sqrt{S(a)Q(a)} .
\end{equation}
\end{lemma}
\begin{proof}
Represent a plaquette holonomy
along \(a\mapsto(1+t)a\) as a
unit quaternion \(q_p(t)\).
Write \(q_{p,0}\) for its scalar
coordinate. Since \(q_p'(0)\)
is tangent to the unit
three-sphere,
\[
 |q_{p,0}'(0)|
    \le\sqrt{1-q_{p,0}(0)^2}\,|q_p'(0)|
    \le\sqrt{2s_p(a)}\,|q_p'(0)|.
\]
The speed of each link
\(e^{(1+t)a_e}\), or its inverse,
is \(|a_e|\). Multiplication by
unit quaternions is isometric,
so the product rule gives
\[
                  |q_p'(0)|\le\sum_{e\in\partial p}|a_e|.
\]
Consequently
\[
 \begin{aligned}
 |DS(a)[a]|
 &\le\sqrt{2S(a)}
        \left\{\sum_p
                 \left(\sum_{e\in\partial p}|a_e|\right)^2
        \right\}^{1/2}\\
 &\le\sqrt{2S(a)}
        \left\{4\sum_p\sum_{e\in\partial p}|a_e|^2\right\}^{1/2}
   =\sqrt{48S(a)Q(a)} .
 \end{aligned}
\]
Each positively oriented edge
occurs in six positive plaquette
boundaries in four dimensions.
This proves the constant.
No radial monotonicity of
\(S\) is asserted.
\end{proof}

\begin{proposition}[A sharper actual inverse-trace budget]
\label{prop:v89-barrier}
Suppose \(\mathbb E S/V\le\tau\)
and \(\mathbb E Q/V\le m\).
Then
\begin{equation}\label{eq:v89-barrier}
 \mathbb E\operatorname{Tr}F^{-1}
        +\frac23\mathbb E Q
       \le n+d_{\rm g}
                         +4\sqrt3\,\beta\sqrt{\mathbb E S\,\mathbb E Q}.
\end{equation}
In particular
\begin{equation}\label{eq:v89-ksharp}
       \frac1V\mathbb E\operatorname{Tr}F^{-1}
             \le k^\sharp:=12+4\sqrt3\,\beta\sqrt{\tau m}.
\end{equation}
The original support and normalizer
are retained.
\end{proposition}
\begin{proof}
Use the same convex truncations
\(\Omega_{R,\epsilon}\) and radial
integration as in V75. At zero
homogeneous background its
determinant and Haar inequalities
give, for the full potential
\(U=\beta S+\mathcal H-\log\det F\),
\[
 DU(a)[a]\ge
      \operatorname{Tr}F^{-1}-d_{\rm g}
                  +\frac23 Q-4\sqrt3\,\beta\sqrt{SQ}.
\]
The outward radial flux is
nonnegative, so the integral
of \(DU[a]\) against the
unnormalized density is at
most \(n\) times its integral.
Move the nonnegative
\(\sqrt{SQ}\) term to the
right and let \(\epsilon\downarrow0\).
The inverse trace passes by
monotone convergence. All other
terms are bounded at fixed
volume and pass by dominated
convergence. Divide by the
unchanged \(Z_\beta\).
Finally
\(\mathbb E\sqrt{SQ}\le
\sqrt{\mathbb E S\,\mathbb E Q}\).
This proves \eqref{eq:v89-barrier}.
Use \((n+d_{\rm g})/V\le12\)
and drop the nonnegative Haar
term to get \eqref{eq:v89-ksharp}.
This is an alternative radial
bound; V79's kinetic term has
not been retained in it.
\end{proof}

Along the selected sequence of
V87, take \(\tau=20/\beta\) and
the valid V88 budget
\[
 m=m_{20}
      =144\pi^2\sqrt{20}\,\beta^{-1/2}
                                         +72\pi^2M^{-2}.
\]
Then
\begin{equation}\label{eq:v89-vanishing}
       \frac{k^\sharp}{\sqrt\beta}
              \le\frac{12}{\sqrt\beta}
                        +4\sqrt{60}\sqrt{m_{20}}
                         \longrightarrow0 .
\end{equation}
The old budget was only bounded,
not vanishing, after division by
\(\sqrt\beta\). The improvement
uses the action expectation,
not a smaller deterministic
chart radius. In particular
this step itself needs no
upper bound on \(\beta/M^4\)
once the action and amplitude
budgets have been supplied.

\subsection{Geometric density changes vanish at thermal move size}

Fix a real, finitely supported
colour two-form \(B\) on
\(\mathbb Z^4\). For sufficiently
large \(M\), embed it in the
torus without support collisions.
Set
\[
 h=d^*B,\qquad \ell=dh=dd^*B,\qquad
 h_\beta=\beta^{-1/2}h,\qquad
 E=\bigcup_{p\in\operatorname{supp}B}\partial p.
\]
Thus \(h_\beta\) preserves the
mean-zero Coulomb space, and
its Euclidean translation
Jacobian there is one.
Write \(b_0=\max_p|B_p|\),
\(\rho=b_0/\sqrt\beta\),
\(r=6\rho\), and
\(\eta=6(1+3R)\rho\).
The zero move is trivial.

\begin{proposition}[Negligible local geometric increment]
\label{prop:v89-geometry}
There are events \(G_M\) with
\(\mu(G_M^c)\to0\), on which
both segments \(a\pm th_\beta\),
\(0\le t\le1\), stay in
\(\Omega_R\). With
\(w(a)=\det F(a)e^{-\mathcal H(a)}\),
\begin{equation}\label{eq:v89-geometry}
  \log\frac{w(a\pm h_\beta)}{w(a)}
                       \longrightarrow0
                 \quad\hbox{in probability on }G_M.
\end{equation}
More explicitly, for all large \(M\),
\begin{equation}\label{eq:v89-geometry-budget}
 \begin{aligned}
 \mu(G_M^c)
    &\le |E|\left(\frac{4m_{20}}{R^2}
                                         +2\eta k^\sharp\right),\\
 \mathbb E\left[
  \mathbf1_{G_M}
       \left|\log\frac{w(a\pm h_\beta)}{w(a)}\right|\right]
    &\le2\eta|E|k^\sharp+12R\rho|E|.
 \end{aligned}
\end{equation}
\end{proposition}
\begin{proof}
Use V85's good event
\(\mathcal A_{E,\rho,1/2}\),
now with the proved inverse
budget \(k^\sharp\).
Its deterministic construction
works simultaneously for
\(B/\sqrt\beta\) and
\(-B/\sqrt\beta\), so call
that event \(G_M\).
Eventually \(r\le R/2\).
The local trace and guard
bounds in
\eqref{eq:v85-local-trace}
and \eqref{eq:v85-guard}
give the first inequality.
The determinant estimate
\eqref{eq:v85-logdet}, with
\(\theta=1/2\), gives
\(2\eta\operatorname{Tr}K_E\).
The unchanged local Haar bound
is \(12R\rho|E|\).
Taking expectations proves
the second inequality.
Both right sides vanish by
\eqref{eq:v89-vanishing} and
\(m_{20}\to0\).
Markov's inequality proves
\eqref{eq:v89-geometry}.
\end{proof}

This does not say that the
full Wilson law is invariant
under the move. The action
increment survives, as follows.

\subsection{The surviving quadratic action increment}

Use V88's actual logarithmic
curvature \(Y_M(a)=\sqrt\beta\,z(a)\).
The same local formula defines
\(Y_M(a\pm h_\beta)\) even when
the shifted field is outside
the gauge-positive chart:
for large \(M\) its link radii
are below \(1/3\), so every
plaquette angle is below
\(4/3<\pi\).
This is only an observable
extension; no density outside
the chart is introduced.
Pairings and norms of two-forms
sum over positive plaquettes
and the three real colours.

\begin{lemma}[Actual thermal translation and action expansion]
\label{lem:v89-action}
For each fixed plaquette,
\begin{equation}\label{eq:v89-curvature-shift}
  \mathbb E\left|
    Y_{M,p}(a\pm h_\beta)-Y_{M,p}(a)\mp\ell_p
             \right|^2\longrightarrow0 .
\end{equation}
Only the fixed finite set of
plaquettes touching \(E\)
can change. Moreover
\begin{equation}\label{eq:v89-action-shift}
 \beta\{S(a\pm h_\beta)-S(a)\}
        \mp\langle Y_M(a),\ell\rangle
                  -\tfrac12\|\ell\|_2^2
                      \longrightarrow0
                  \quad\hbox{in probability}.
\end{equation}
\end{lemma}
\begin{proof}
The map from the four link
logarithms to their plaquette
logarithm is smooth on the
compact product of radius-\(1/3\)
balls: the product stays
strictly away from the cut
at \(-I\). Its second derivative
is therefore uniformly bounded
there. At zero its derivative
is the signed link sum \(d\).
Taylor's formula in the moving
direction, with the difference
of first derivatives estimated
by this Hessian bound, gives
\[
 \left|Y_{M,p}(a\pm h_\beta)-Y_{M,p}(a)\mp(dh)_p\right|
       \le C_h\left(
               \Big(\sum_{e\in\partial p}|a_e|^2\Big)^{1/2}
                                     +\beta^{-1/2}\right).
\]
The constant depends on the
fixed move, not on volume.
The local amplitude expectation
proves \eqref{eq:v89-curvature-shift}.

V88 already proved
\[
        \beta s_p(a)-\tfrac12|Y_{M,p}(a)|^2
                       \longrightarrow0
                 \quad\hbox{in probability}.
\]
The same Taylor remainder
applies to the shifted field.
Its curvature coordinates are
tight by
\eqref{eq:v89-curvature-shift}
and V88's moment bound.
On every fixed bounded
curvature set the remainder
is \(O(\beta^{-1})\);
the complement has uniformly
small probability.
Thus its shifted remainder
also tends to zero in
probability.
Subtract the two quadratic
expressions, use
\eqref{eq:v89-curvature-shift},
and sum over the fixed finite
set of changed plaquettes.
The quadratic difference is
\(\pm\langle Y_M,\ell\rangle+
\|\ell\|_2^2/2\), up to a term
tending to zero in probability.
This proves
\eqref{eq:v89-action-shift}.
No convergence of fourth
moments is needed.
\end{proof}

\subsection{A bounded identity that survives weak convergence}

Extend the full normalized density
\[
           f_M(a)=Z_\beta^{-1}
                \mathbf1_{\Omega_R}(a)e^{-\beta S(a)}w(a)
\]
by zero to its finite-dimensional
mean-zero Coulomb space. For
\(\mu\)-almost every \(a\), let
\[
       q_\pm(a)=\frac{f_M(a\pm h_\beta)}{f_M(a)},
                      \qquad A_\pm(a)=\min\{1,q_\pm(a)\}.
\]
These definitions include all
possible chart-support loss.
The following equality is
exact at every finite volume:
for bounded measurable \(F\),
\begin{equation}\label{eq:v89-acceptance}
 \mathbb E\{A_+(a)F(Y_M(a))\}
     =\mathbb E\{A_-(a)F(Y_M(a-h_\beta))\}.
\end{equation}
Indeed the left side is
\(\int\min\{f_M(a),f_M(a+h_\beta)\}F(Y_M(a))\,da\);
substitute \(b=a+h_\beta\).
The right side follows with
no boundary omission or new
normalizer.

Along any V88 subsequence for
which \(Y_M\) converges weakly,
call the limit law \(\nu\).
For a finitely supported
\(\ell=dd^*B\), define
\begin{equation}\label{eq:v89-acceptance-limit}
 \begin{aligned}
 a_+(y)&=\min\{1,e^{-\langle y,\ell\rangle-\|\ell\|_2^2/2}\},\\
 a_-(y)&=\min\{1,e^{ \langle y,\ell\rangle-\|\ell\|_2^2/2}\}.
 \end{aligned}
\end{equation}
Both are bounded continuous
cylinder functions and strictly
positive at every field \(y\).

\begin{lemma}[The limiting balanced measure identity]
\label{lem:v89-balance}
For every bounded continuous
cylinder function \(F\),
\begin{equation}\label{eq:v89-balance}
        \int a_+(y)F(y)\,d\nu(y)
             =\int a_-(y)F(y-\ell)\,d\nu(y).
\end{equation}
\end{lemma}
\begin{proof}
On \(G_M\),
\[
 \log q_\pm(a)=
   -\beta\{S(a\pm h_\beta)-S(a)\}
                     +\log\frac{w(a\pm h_\beta)}{w(a)}.
\]
The geometry proposition and
action lemma show that this
differs in probability from
\(\mp\langle Y_M,\ell\rangle-\|\ell\|_2^2/2\)
by a vanishing term.
Since \(\mu(G_M^c)\to0\),
the bounded quantities
\(A_\pm-a_\pm(Y_M)\)
tend to zero in probability
on the whole law and hence
in \(L^1\).
This last conclusion concerns
only bounded acceptance
weights, not \(q_\pm\) themselves.

Equation~\eqref{eq:v89-curvature-shift}
also makes the shifted finite
coordinates equal to
\(Y_M-\ell+o_{\mathbb P}(1)\).
For a bounded continuous
cylinder function, tightness
of those coordinates and
uniform continuity on compact
sets imply
\[
          F(Y_M(a-h_\beta))-F(Y_M(a)-\ell)
                               \longrightarrow0
                         \quad\hbox{in }L^1 .
\]
Apply these facts to
\eqref{eq:v89-acceptance},
then pass its remaining
bounded continuous functions
through weak convergence.
This proves \eqref{eq:v89-balance}.
\end{proof}

\subsection{Exact exponential translation and Gaussian probe laws}

Let \(T_\ell y=y+\ell\).
The next result has the usual
Gaussian exponential translation
form, but is derived without
assuming that the full law
\(\nu\) is Gaussian.\footnote{
For the conventional Gaussian
quasi-invariance framework see
M. Hairer, \emph{An Introduction
to Stochastic PDEs}, Section 4.2,
\href{https://www.hairer.org/notes/SPDEs.pdf}
{author's lecture notes}. Here
the exponential identity is
deduced from the bounded
actual-law identity, not
imported for an unproved
Gaussian measure.}

\begin{theorem}[Thermal curvature translation formula]
\label{thm:v89-translation}
Every V88 limit law satisfies,
for every finitely supported
real colour two-form \(B\) and
\(\ell=dd^*B\),
\begin{equation}\label{eq:v89-RN}
     \frac{d(T_\ell)_*\nu}{d\nu}(y)
         =\exp\{\langle y,\ell\rangle-\tfrac12\|\ell\|_2^2\}.
\end{equation}
The translated and original
measures are mutually absolutely
continuous. Consequently
\begin{equation}\label{eq:v89-mgf}
       \mathbb E_\nu e^{t\langle Y,\ell\rangle}
                    =e^{t^2\|\ell\|_2^2/2}
                         \qquad(t\in\mathbb R).
\end{equation}
For any finite collection
\(\ell_j=dd^*B_j\), the real
variables \(\langle Y,\ell_j\rangle\)
are jointly centered Gaussian,
with
\begin{equation}\label{eq:v89-covariance}
 \operatorname{Cov}_\nu
       (\langle Y,\ell_i\rangle,\langle Y,\ell_j\rangle)
                                  =\langle\ell_i,\ell_j\rangle_{\ell^2}.
\end{equation}
\end{theorem}
\begin{proof}
The two sides of
\eqref{eq:v89-balance} are
finite measures tested on
bounded continuous cylinder
functions. These tests
determine finite Borel
measures on the countable
product space. Thus
\[
                a_+\nu=(T_{-\ell})_*(a_-\nu).
\]
Put \(\rho=(T_{-\ell})_*\nu\).
The right side is
\(a_-(y+\ell)\rho(dy)\).
The exact scalar relation
\[
 a_-(y+\ell)
      =e^{\langle y,\ell\rangle+\|\ell\|_2^2/2}a_+(y)
\]
and strict positivity of the
two weights give
\[
        \rho(dy)=
          e^{-\langle y,\ell\rangle-\|\ell\|_2^2/2}\nu(dy).
\]
Division is valid as an
identity of measures, even
though the positive weights
need not be bounded below.
For example divide first
on the sets where both
weights exceed \(1/n\)
and increase their union.
Replace \(\ell\) by \(-\ell\)
to obtain \eqref{eq:v89-RN}.
The positive density also
proves equivalence.
Its integrability is automatic
because the translated measure
is a probability measure.

Apply the same conclusion to
\(tB\), whose shift is \(t\ell\).
Integrating its density gives
\eqref{eq:v89-mgf}. The moment
generating function is exactly
that of \(N(0,\|\ell\|_2^2)\)
and is finite on the whole
real line, so it uniquely
determines that distribution.
Every real linear combination
of the \(\ell_j\) is again of
the form \(dd^*B\) with finite
support. Its Gaussian variance
is the squared norm of that
combination. This proves the
joint Gaussian statement and
the displayed covariance.
\end{proof}

The same formula gives, for
bounded \(C^1\) cylinder \(F\)
with bounded derivative,
\begin{equation}\label{eq:v89-Stein}
       \mathbb E_\nu D_\ell F(Y)
              =\mathbb E_\nu[F(Y)\langle Y,\ell\rangle].
\end{equation}
Differentiate the translated
expectation at \(t=0\).
The left derivative is
dominated by the bounded
directional derivative.
For the right derivative,
\eqref{eq:v89-mgf} supplies
integrable bounds of the form
\((1+|\langle Y,\ell\rangle|)
e^{|\langle Y,\ell\rangle|}\).
Thus differentiation is
justified after, not before,
the exponential formula.

These Gaussian statements
concern the limiting probes.
They do not by themselves
assert convergence of the
finite-volume second or
higher moments to their
Gaussian values.

\subsection{A determined derivative field and what remains invisible}

On lattice two-forms the
commuting forward differences
and their adjoints satisfy
the Hodge identity
\[
                 \Delta=d_1d_1^*+d_2^*d_2 .
\]
Here \(d_1\) maps links to
plaquettes, \(d_2\) maps
plaquettes to cubes, and
\(\Delta\) acts componentwise
as the positive scalar
Laplacian. The identity
follows by expanding the
alternating sums: mixed
differences cancel in pairs,
leaving \(\sum_iD_i^*D_i\)
on each component.
V88's almost-sure \(d_2Y=0\)
therefore gives, for every
finite two-form \(B\),
\begin{equation}\label{eq:v89-laplacian}
              \langle Y,\Delta B\rangle
                         =\langle Y,d_1d_1^*B\rangle .
\end{equation}
All sums are finite, so no
summability of \(Y\) is
required for summation by
parts. It follows that the
derivative field \(\Delta Y\)
has jointly Gaussian finite
coordinate laws, with
\begin{equation}\label{eq:v89-derivative-covariance}
 \mathbb E_\nu\langle\Delta Y,B\rangle
                         \langle\Delta Y,C\rangle
                    =\langle dd^*B,dd^*C\rangle .
\end{equation}
This is more information
than Bianchi compatibility.

For example take \(B\) equal
to one unit colour on a
single plaquette and zero
elsewhere. The field
\(\ell=dd^*B\) has value four
on that plaquette, values
of magnitude one on twenty
other plaquettes, and zero
elsewhere. Hence
\begin{equation}\label{eq:v89-single-probe}
          \|\ell\|_2^2=36,\qquad
                    \langle Y,\ell\rangle\sim N(0,36).
\end{equation}
Both separate examples from
V88 violate this identity:
the iid-link curl example
has probe variance
\(\|d^*\ell\|_2^2=352\),
while its constant-field
example has variance zero.
For the value 352, write the
two nonzero orientation
components of \(h=d^*B\) as
signed differences of two
neighbouring vertex deltas.
The scalar stencil has
\(\|\Delta\delta_0\|_2^2=8^2+8=72\)
and
\(\langle\Delta\delta_0,\Delta\delta_{e_i}\rangle=-16\):
the two adjacent centres
contribute \(-8\) each and
have no other common neighbours.
Each component of \(\Delta h\)
therefore has squared norm
\(2\cdot72-2(-16)=176\).
Since \(d^*\ell=\Delta h\),
the two components give 352.
Thus the new measure identity
rules out both former
diagnostics as descriptions
of the actual thermal limit.
The stencil and norm counts
are verified independently
by the exact certificate.

There remains a genuine
low-frequency question.
Every \(\ell=dd^*B\) has zero
sum in each orientation and
colour. A spatially constant
two-form \(H\) therefore obeys
\(\langle H,\ell\rangle=0\).
In particular the translation
identity alone is unchanged
by adding an independent
random constant two-form:
conditionally on \(H\), use
\eqref{eq:v89-RN} for the test
\(F(Y+H)\), and replace
\(\langle Y,\ell\rangle\) by
\(\langle Y+H,\ell\rangle\).
This observation does not
assert that such a component
is present in the actual law,
or that arbitrary added
components respect all its
other bounds. It shows what
the local translation identity
does not determine.

Likewise all moves and tests
in the limit argument have
fixed finite support.
The explicit error
\eqref{eq:v89-geometry-budget}
depends on \(|E|\) and
\(\|B\|_\infty\); the action
expansion also has move-dependent
constants. Replacing \(B\)
by a growing physical-scale
test is not justified by
the fixed-support proof.
Identifying and controlling
the remaining low-frequency
covariance, and then proving
the correct joint physical
limit, are upstream tasks
still required for the
mass-gap programme.

\subsection{Exact checks and preservation}

The certificate is
\begin{center}\small
\path{scripts/verify_ym_f14_v89_thermal_translation.py}.
\end{center}
Sixteen rational quaternion
first variations check the
action factor in the tangent
speed bound. The sixfold
edge incidence gives the
squared global constant 48.
Twenty compact colour
stencils check closedness,
zero total coefficients and
\(d^*dd^*B=\Delta d^*B\).
All eighteen single-plaquette
orientation-colour probes
have squared norm 36;
their iid-link-curl comparison
has squared source norm 352.

Sixty-four exact bounded
acceptance identities use
finite discrete probability
banks, retaining their
support-loss terms.
Sixty-eight scalar identities
check the positive-weight
division and its sign.
Thirty-six Gaussian
moment-generating coefficients
are checked rationally.
These finite probability
banks are diagnostics of
the algebra, not interacting
YM samples. The passage to
the limit and the density
formula are proved above,
not inferred numerically.

The certificate contains 7,450 bytes and has SHA--256
\begin{center}\scriptsize\ttfamily
5294DDF89BDE225DD21D74D22A9A68350D79AB21EFA48FF24E985D9B51C5E61D.
\end{center}
The predecessor V88 contains 2,171,921 bytes and has SHA--256
\begin{center}\scriptsize\ttfamily
BDC37DEA34744A940504FF6E994F05C58B65FAAC3DCB30A6D3CB81317B996EDE.
\end{center}
Its body is preserved except
for the title version and this
terminal insertion. V89 replaces
V88 as the sole active main note.
Earlier certificates, supplied
sources, companion notes,
monographs and archives remain
unchanged.

\begin{firewall}
The actual selected-coupling
thermal limits now satisfy
an exact exponential
translation identity and
Gaussian covariance formulas
for a specified family of
local curvature probes.
This is not a claim that
the full interacting
continuum theory has been
constructed, or that a
physical mass is zero or
positive. Physical-scale
correlations, reconstruction
and the positive Hamiltonian
mass gap remain unproved.
The original full objective
remains active.
\end{firewall}
\section{V90: classification of the thermal curvature law and a sharp selected normalization}
\label{sec:v90}

\subsection{Companion checkpoint and scope of the continuation}

The scheduled V90 checkpoint reopened
the current SM, NS, Shell and parallel
YM conclusions, and rechecked the
three supplied source-selection,
stationarity and exact-action notes.
All ten protected source hashes
match the preceding audit.
The active companion files are still:
\begin{center}\small
\begin{tabular}{lrr}
source&version&bytes\\\hline
SM--Einstein continuum&72&607,372\\
NS stability&26&278,283\\
YM Shell mixing&16&623,261\\
Parallel YM sector notes&5&54,174
\end{tabular}
\end{center}
SM V72's many-copy transport theorem
holds at a fixed number of physical
momentum modes. It is not a
cutoff-uniform physical YM theorem.
NS V26's rank-one kernel contractions
do not supply YM correlation control.
Shell V16 retains a complete flat
holonomy coefficient, including its
nonlocal low terms; its conclusion
does not identify a joint interacting
integral remainder. Parallel YM V5
continues to distinguish compact
sector coordinates from extensive
integer charge.

The supplied stationarity V21 also
retains its explicit real-phase
caustic obstruction and its separate
quadratic phase-space comparison.
Source-selection V17 does not identify
its comparison instrument with the
primitive physical source. The
exact-action bridge V16 does not
turn its comparison inverse into
a nonlinear continuum trajectory.
None supplies the missing physical
YM covariance estimate. These
documents are mathematical source
material, not instructions to
change the task or its goal.
The next regular checkpoint is V95.

V89 made genuine progress beyond
kinematics: it proved an exponential
translation formula for actual
thermal curvature limits. This
section determines the remaining
freedom in that formula. The
stationary law is a free lattice
Gaussian curvature plus an
independent random constant
two-form. A sharper logarithmic
coupling selection then removes
the constant part and proves
convergence of fixed-lattice
curvature covariances.

The qualifier \emph{fixed lattice}
is essential. The result concerns
the same restricted chart at
selected couplings. It is not
an interacting physical
continuum limit, and its Gaussian
character does not decide the
physical mass-gap question.

\subsection{The hypotheses and the deterministic curvature projector}

Let \(\nu\) be any thermal curvature
limit obtained in V88--V89.
We use exactly its proved properties:
stationarity, finite second moments,
common colour-rotation invariance,
\[
 d_2Y=0,\qquad
 \frac{d(T_\ell)_*\nu}{d\nu}(y)
       =e^{\langle y,\ell\rangle-\|\ell\|_2^2/2}
       \quad\hbox{for }\ell=d_1d_1^*B,\quad
                               B\hbox{ finitely supported}.
\]
Here \(d_1\) maps links to plaquettes,
\(d_2\) maps plaquettes to cubes,
and \(T_\ell y=y+\ell\).
The following arguments apply more
generally to any law with these
properties and a uniform finite
second-moment bound.

On the Fourier torus \(\mathbb T^4\),
with normalized measure
\(dk/(2\pi)^4\), put
\[
 q_i(k)=e^{ik_i}-1,\qquad
 \lambda(k)=\sum_i|q_i(k)|^2,\qquad
                   K(k)=d_1(k)d_1(k)^*.
\]
The Hodge identity and \(d_2d_1=0\)
give
\[
 K(k)^2=\lambda(k)K(k),\qquad
 K(k)+d_2(k)^*d_2(k)=\lambda(k)I .
\]
For \(k\ne0\), define
\begin{equation}\label{eq:v90-projector}
                P(k)=\frac{K(k)}{\lambda(k)}.
\end{equation}
It is the orthogonal projector
onto exact two-forms at that
momentum, has rank three, and
acts independently on the three
colours. Its value at \(k=0\)
may be set to zero in Lebesgue
integrals. We use \(P\) also
for the corresponding bounded
operator on square-summable
lattice two-forms.

For a deterministic absolutely
summable two-form \(C\), the
random pairing is well defined
in \(L^2(\nu)\):
\begin{equation}\label{eq:v90-l1-pairing}
 \|\langle Y,C\rangle\|_{L^2(\nu)}
       \le \sqrt{K_2}\sum_p|C_p|,
             \qquad K_2=\sup_p\mathbb E_\nu|Y_p|^2<\infty .
\end{equation}
This is Minkowski's inequality.
It will justify the passage
from finite tests to heat
kernels; no unproved global
\(\ell^2\) bound on a sample
of \(Y\) is used.

\subsection{Heat averaging isolates a spatially constant random part}

Let \(Q_t=e^{-t\Delta}\) be the
componentwise lattice heat
semigroup. Its convolution kernel
\(p_t\) is nonnegative, symmetric
and has total mass one. For
example the expansion
\[
 e^{-t\Delta}
   =e^{-8t}\sum_{n=0}^{\infty}
          \frac{t^n}{n!}
          \left(\sum_{i=1}^4(T_{e_i}+T_{-e_i})\right)^n
\]
proves these facts, with the
series converging in \(\ell^1\)
operator norm.

\begin{lemma}[The invariant heat limit]
\label{lem:v90-harmonic}
There is an \(L^2\) random colour
two-form \(H=(H_{ij})_{i<j}\),
independent of the spatial
coordinate, such that
\begin{equation}\label{eq:v90-heat-limit}
                  Q_tY_{ij}(x)\longrightarrow H_{ij}
                      \quad\hbox{in }L^2(\nu)
                          \quad(t\to\infty).
\end{equation}
It is a measurable function
of the field \(Y\), invariant
under spatial translations.
It can also be chosen so that
\begin{equation}\label{eq:v90-harmonic-shift}
                     H(Y+\ell)=H(Y)
\end{equation}
almost surely for every
finitely supported \(\ell\).
For the actual laws,
\(\mathbb E H=0\) and the
law of \(H\) is invariant
under common colour rotations.
\end{lemma}
\begin{proof}
Spatial translations define
commuting unitary operators
\(\mathcal U_i\) on \(L^2(\nu)\).
The bounded self-adjoint
nonnegative operator
\[
       \mathcal D=\sum_i(2I-\mathcal U_i-\mathcal U_i^*)
\]
satisfies
\(\langle f,\mathcal Df\rangle
=\sum_i\|(\mathcal U_i-I)f\|_2^2\).
Thus its kernel is exactly
the subspace of translation
invariant random variables.
Spectral calculus gives
\(e^{-t\mathcal D}\to\Pi_{\rm inv}\)
strongly: for every spectral
value \(u\ge0\),
\(e^{-tu}\to\mathbf1_{\{0\}}(u)\),
and dominated convergence
applies to each spectral
measure.

The same convergent heat
series shows that
\[
 e^{-t\mathcal D}Y_{ij}(x)
                  =\sum_zp_t(z)Y_{ij}(x+z)
\]
in \(L^2\), justified by
\eqref{eq:v90-l1-pairing}.
Set \(H_{ij}=\Pi_{\rm inv}Y_{ij}(0)\).
The projection is invariant
under translations, so its
value is the same at every
\(x\). This proves
\eqref{eq:v90-heat-limit}.

Choose a sequence \(t_n\to\infty\)
along which all eighteen
origin colour components
converge almost surely as
well as in \(L^2\).
Such a common subsequence
exists by summably small
\(L^2\) errors. The heat
series at each \(t_n\)
converges absolutely almost
surely, since its absolute
first moment is finite.
This constructs \(H\) as
a measurable function of
the coordinate field.

For a finitely supported
deterministic \(\ell\),
\(Q_t\ell(0)\to0\). Indeed
\[
 \sup_x|p_t(x)|
       \le\int_{\mathbb T^4}e^{-t\lambda(k)}
                                    \frac{dk}{(2\pi)^4}\longrightarrow0
\]
by dominated convergence,
since \(\lambda(k)>0\)
almost everywhere.
Adding \(\ell\) therefore
does not change the almost
sure heat limit on its
domain of definition. This
proves \eqref{eq:v90-harmonic-shift}.
The projection preserves
expectation, because constants
are in its range. Thus
\(\mathbb E H=\mathbb E Y(0)=0\).
Heat averaging commutes with
common colour rotations,
which proves the last claim.
\end{proof}

\subsection{Independence follows from the measure identity, not from zero covariance}

\begin{lemma}[Gaussian local scores independent of the heat limit]
\label{lem:v90-independence}
For every finite family
\(\ell_j=d_1d_1^*B_j\), the
Gaussian vector
\((\langle Y,\ell_j\rangle)_j\)
is independent of \(H\).
\end{lemma}
\begin{proof}
Let \(f\) be a bounded
nonnegative measurable
function of the finite
vector \(H\), and let
\(\ell=\sum_jt_j\ell_j\).
V89's density identity,
tested on \(f(H)\), and
\eqref{eq:v90-harmonic-shift}
give
\begin{equation}\label{eq:v90-weighted-mgf}
 \mathbb E_\nu\left[
   f(H)e^{\langle Y,\ell\rangle-\|\ell\|_2^2/2}\right]
                                      =\mathbb E_\nu f(H).
\end{equation}
The equality of the two
translated expectations is
legitimate for this Borel
test: the translation
formula is an identity of
measures, and translated
null sets remain null.

If \(\mathbb E f(H)>0\),
normalize the measure
weighted by \(f(H)\).
Equation~\eqref{eq:v90-weighted-mgf}
says that under every such
weight the score vector has
the same Gaussian moment
generating function, hence
the same distribution.
Taking \(f\) to be indicators
of Borel sets proves
independence. The case
\(\mathbb E f(H)=0\) is
immediate. This argument
is stronger than, and does
not rely on, vanishing
cross covariance.
\end{proof}

\subsection{The complete stationary thermal-law classification}

\begin{theorem}[Gaussian curvature plus an independent constant two-form]
\label{thm:v90-classification}
Under the stated actual-limit
hypotheses,
\begin{equation}\label{eq:v90-decomposition}
                          Y=G+H
\end{equation}
almost surely, where \(H\)
is the spatially constant
random two-form above,
and \(G\) is independent
of \(H\), centered Gaussian,
stationary, and has covariance
\begin{equation}\label{eq:v90-free-covariance}
 \mathbb E\langle G,B\rangle\langle G,C\rangle
          =\langle B,PC\rangle_{\ell^2}
          =\int_{\mathbb T^4}\widehat B(k)^*
                  P(k)\widehat C(k)\frac{dk}{(2\pi)^4}
\end{equation}
for finitely supported real
two-forms \(B,C\), with
\(P\otimes I_3\) understood
on colours.
In particular the finite
coordinate distributions
are exactly specified by
\begin{equation}\label{eq:v90-characteristic}
 \mathbb E e^{i\langle Y,B\rangle}
       =e^{-\langle B,PB\rangle/2}
         \mathbb E\exp\left\{
              i\sum_{i<j}\left(\sum_xB_{ij}(x)\right)\cdot H_{ij}
                        \right\}.
\end{equation}
No other stationary
finite-second-moment freedom
is left by the hypotheses.
\end{theorem}
\begin{proof}
Set \(G_t=(I-Q_t)Y\).
For a finite two-form \(B\),
put \(R_tB=\int_0^tQ_sB\,ds\).
This is in \(\ell^1\), with
norm at most \(t\|B\|_{\ell^1}\),
and \(\Delta R_tB=(I-Q_t)B\).
The Bianchi and Hodge
identities imply
\begin{equation}\label{eq:v90-filtered-test}
 \langle G_t,B\rangle
      =\langle Y,(I-Q_t)B\rangle
      =\langle Y,d_1d_1^*R_tB\rangle .
\end{equation}
To justify this beyond
finite support, approximate
\(R_tB\) by finite truncations
in \(\ell^1\).
The difference operators
are bounded on \(\ell^1\).
Equation~\eqref{eq:v90-l1-pairing}
passes all the pairings
in the finite-support
Bianchi identity to the
limit in \(L^2(\nu)\).

Let \(C_n\to R_tB\) in
\(\ell^1\), with finite
support, and
\(\ell_n=d_1d_1^*C_n\).
Then \(\ell_n\to d_1d_1^*R_tB\)
in both \(\ell^1\) and
\(\ell^2\), and the
corresponding random pairings
converge in \(L^2\).
By V89 and
Lemma~\ref{lem:v90-independence},
they are centered Gaussian
and independent of \(H\).
The same argument applies
jointly to finitely many
tests and times.
Passing to this \(L^2\)
limit proves those properties
for the filtered fields,
with covariance determined
by the limiting deterministic
\(\ell^2\) pairings.

In Fourier variables,
\[
 \widehat{d_1d_1^*R_tB}(k)
                  =(1-e^{-t\lambda(k)})P(k)\widehat B(k).
\]
Thus
\[
 \mathbb E\langle G_t,B\rangle\langle G_t,C\rangle
    =\int\widehat B^*
            (1-e^{-t\lambda})^2P\widehat C\,
                                         \frac{dk}{(2\pi)^4}.
\]
By Lemma~\ref{lem:v90-harmonic},
\(G_t\to G:=Y-H\) in
coordinatewise \(L^2\), and
therefore for every finite
collection of tested coordinates.
The integral converges to
\eqref{eq:v90-free-covariance}
by \(0\preceq P\preceq I\)
and dominated convergence.
Limits of these finite
Gaussian vectors are Gaussian.
Their joint independence
from the unchanged \(H\)
also passes to the limit.
This proves the decomposition,
covariance and characteristic
function. Stationarity follows
from that of the original
field and its invariant
heat projection.
\end{proof}

For the actual mean-zero
colour laws, the full covariance
is therefore the Fourier
kernel of \(P\), plus the
spatially constant matrix
\(\operatorname{Cov}(H)\).
Equivalently its spectral
measure consists of
\(P(k)\,dk/(2\pi)^4\)
and a possible atom at zero.
If the actual limiting law
were spatially ergodic, its
invariant \(H\) would be
deterministic and hence zero
by \(\mathbb E H=0\).
Ergodicity has not been
assumed or proved for the
original selected sequence.

The random \(H\) is a constant
\emph{thermal curvature}.
It has not been identified
with a compact flat-connection
holonomy or a topological
charge. Those are different
objects in the companion notes.

\subsection{The exact energy floor of the classified law}

For one colour and orientation
\((i,j)\), the diagonal symbol is
\[
                 P_{ij,ij}(k)
                   =\frac{|q_i(k)|^2+|q_j(k)|^2}{\lambda(k)} .
\]
Its integral is \(1/2\):
permutation of the four
coordinates makes the
integral of each
\(|q_i|^2/\lambda\) equal
to \(1/4\).
It follows that
\begin{equation}\label{eq:v90-energy-floor}
 \mathbb E|G_{ij}(0)|^2=\frac32,\qquad
 \frac12\sum_{i<j}\mathbb E|Y_{ij}(0)|^2
        =\frac92+\frac12\mathbb E\sum_{i<j}|H_{ij}|^2 .
\end{equation}
Independence and zero means
remove the cross terms.
For V87's original selected
couplings, the actual
expectation of
\(\sum_{i<j}\beta s_{ij}(0)\)
is at most 20.
Lower semicontinuity under
the joint curvature-energy
limit therefore gives
\[
                     \mathbb E\sum_{i<j}|H_{ij}|^2\le31 .
\]
This is a finite bound,
not elimination of \(H\).
The exact floor \(9/2\),
however, permits a sharper
selection to eliminate it.

\subsection{A sharper actual coupling selection}

Retain V87's
\(q_M(u)=u\mathbb E_uS/V\)
and \(C_R^1\) from
\eqref{eq:v87-pressure-constants}.
For sufficiently large \(M\),
define
\[
 U_M=\log M,\qquad r_M=\log U_M,\qquad
                         L_M=\exp(\sqrt{r_M}).
\]
Take \(r_M\ge4\), so
\(1<L_M<U_M\), and let
\(\beta_M^\circ\) be the
leftmost minimizer of
\(q_M\) on \([L_M,U_M]\).
This is an existence
selection using the actual
normalizer, not its numerical
evaluation.

\begin{theorem}[Sharp selected energy normalization]
\label{thm:v90-selection}
The selected couplings satisfy
\begin{equation}\label{eq:v90-selected-budget}
 \begin{gathered}
       e^{\sqrt{\log\log M}}\le\beta_M^\circ\le\log M,\\
       q_M(\beta_M^\circ)\le\frac92+\eta_R(r_M),\qquad
       \eta_R(r)=\frac{(9/2)\sqrt r+C_R^1}{r-\sqrt r}
                                           \longrightarrow0 .
 \end{gathered}
\end{equation}
Along these actual laws,
the joint periodic curvature-energy
fields converge weakly, in
the countable product topology,
to
\begin{equation}\label{eq:v90-selected-limit}
                  \left(G,\ \tfrac12|G_p|^2\right)_{p\in\mathcal P_\infty},
\end{equation}
where \(G\) has covariance
\eqref{eq:v90-free-covariance}.
In particular the limiting
constant component is zero
without an ergodicity assumption.
\end{theorem}
\begin{proof}
The same differentiability
and positivity used in V87
give
\[
 \begin{aligned}
 \int_{L_M}^{U_M}q_M(u)\,\frac{du}{u}
      &=\frac1V\log\frac{Z_{L_M}}{Z_{U_M}}\\
      &\le\frac92\log(2+U_M)+C_R^0
       \le\frac92r_M+C_R^1 .
 \end{aligned}
\]
Divide by the logarithmic
interval length
\(r_M-\sqrt{r_M}\).
The minimum is at most
this average, which is
exactly the right side of
\eqref{eq:v90-selected-budget}.
Continuity gives a compact
nonempty set of minimizers
and thus its leftmost point.
This proves the budget
and the coupling bounds.

Eventually \(q_M(\beta_M^\circ)\le20\),
\(\beta_M^\circ\to\infty\)
and \(\beta_M^\circ/M^4\le2\).
The proofs of V87's thermal
estimate and V88--V89 use
these actual bounds, not
the earlier convention
of doubling a particular
minimizer. In detail the
same deterministic
action-to-amplitude estimate
gives \(m_{20}\to0\);
V88 gives tight curvature,
the joint energy identity
and Bianchi closure;
V89's \(k^\sharp/\sqrt\beta\to0\)
gives the exponential
translation formula.
All limits on this newly
selected sequence therefore
satisfy the hypotheses of
Theorem~\ref{thm:v90-classification}.

For any such joint limit,
nonnegativity and weak
convergence give
\[
 \frac12\sum_{i<j}\mathbb E|Y_{ij}(0)|^2
     \le\liminf_M q_M(\beta_M^\circ)\le\frac92
\]
along its defining subsequence.
Equation~\eqref{eq:v90-energy-floor}
forces \(\mathbb E\sum_{i<j}|H_{ij}|^2=0\).
Hence \(H=0\) almost surely,
and the limit is
\eqref{eq:v90-selected-limit}.
Tightness and uniqueness of
all subsequential limits
give convergence of the
whole selected sequence.
\end{proof}

This selector has not proved
\(\beta_M^\circ\sim c\log M\)
for a prescribed \(c\).
It does not implement a
renormalization recursion,
choose a physical lattice
spacing, or establish
cofinal compatibility of
physical-scale observables.

\subsection{Uniform integrability and actual fixed-separation covariance convergence}

The sharp expectation budget
now supplies information that
V88--V89 explicitly left open.
It is important to derive it,
rather than infer moment
convergence merely from the
Gaussian limiting law.

\begin{theorem}[Convergence of thermal curvature second moments]
\label{thm:v90-moments}
Along \(\beta_M^\circ\),
\begin{equation}\label{eq:v90-moment-limits}
 \begin{gathered}
       q_M(\beta_M^\circ)\longrightarrow\frac92,\qquad
       \mathbb E[\beta_M^\circ s_p]\longrightarrow\frac34,\\
       \mathbb E|Y_{M,p}|^2\longrightarrow\frac32 .
 \end{gathered}
\end{equation}
For every pair of fixed
lattice plaquettes and
colour components, their
actual curvature covariance
converges to the corresponding
kernel entry of \(P\otimes I_3\).
The one-plaquette energy law
converges to
\(\tfrac14\chi_3^2\).
No finite-volume convergence
of energy variances or
fourth moments is asserted.
\end{theorem}
\begin{proof}
At the origin put
\[
                 Z_M=\sum_{i<j}\beta_M^\circ s_{ij}(0).
\]
Translation covariance gives
\(\mathbb E Z_M=q_M(\beta_M^\circ)\)
exactly. The joint weak
limit is
\(Z=\frac12\sum_{i<j}|G_{ij}(0)|^2\),
with \(\mathbb E Z=9/2\).
Lower semicontinuity and
\eqref{eq:v90-selected-budget}
therefore show
\(\mathbb E Z_M\to\mathbb E Z\).

Here is the required uniform
integrability argument.
For \(K>0\), bounded-continuous
truncation and convergence
of the means give
\[
 \mathbb E(Z_M-K)_+
       =\mathbb E Z_M-\mathbb E\min\{Z_M,K\}
                       \longrightarrow\mathbb E(Z-K)_+ .
\]
Use
\[
      z\mathbf1_{\{z>2K\}}\le2(z-K)_+
                                  \qquad(z\ge0)
\]
and then let \(K\to\infty\).
This proves uniform
integrability of \(Z_M\)
on the tail of the sequence;
finitely many initial
integrable variables do not
change the conclusion.
Each nonnegative origin
plaquette energy is bounded
by \(Z_M\), so it is
uniformly integrable too.
Stationarity gives the same
conclusion at any fixed
translate.

V88's pointwise comparison
\(|Y_{M,p}|^2\le4\beta_M^\circ s_p\)
then proves uniform
integrability of each
\(|Y_{M,p}|^2\).
Weak convergence of the
corresponding finite coordinates
now passes their expectations
to the Gaussian values.
For two fixed components,
their product is bounded
in absolute value by one
half the sum of the two
squared vector norms, so
the products are also
uniformly integrable.
Their expectations therefore
converge. Means are zero
by colour invariance, giving
the covariance statement.

Each limiting plaquette vector
has law \(N(0,\tfrac12I_3)\),
by the diagonal projector
calculation. Thus its energy
is \(\tfrac14\chi_3^2\), of
mean \(3/4\), and its
squared norm has mean \(3/2\).
This proves all claims.
It does not provide uniform
integrability of the squares
of the energy variables.
\end{proof}

\subsection{What remains between this result and the original goal}

The constant ambiguity in the
thermal law is now classified
and removed on a sharper
selected sequence. Fixed-lattice
curvature covariances are
identified under the actual
law, with their moment
passage justified. Repeating
local compactness or guessing
the same Gaussian covariance
is no longer the next task.

The unresolved requirement is
uniform control when spatial
separations and test supports
grow with the inverse physical
lattice spacing, together with
the appropriate renormalized
observables and physical
coupling trajectory.
A fixed-separation limit
does not give the covariance
at such moving separations,
and the compact-move estimates
have support-dependent errors.
The current Gaussian thermal
limit therefore neither
constructs the interacting
continuum theory nor proves
that its physical Hamiltonian
has zero or positive mass.
The restricted chart has also
not been proved to exhaust
the full compact theory.

\subsection{Exact verification and preservation}

The certificate is
\begin{center}\small
\path{scripts/verify_ym_f14_v90_thermal_classification.py}.
\end{center}
On the complete side-two and
side-four momentum grids it
checks the two-form Hodge
identity, the rank-three
Hermitian projector, Bianchi
closure, and the filtered
covariance identity over
\(\mathbb Q(i)\).
The 270 nonzero momenta are
all included. The one-colour
plaquette variances are
\[
                \frac{15}{32}\quad\hbox{and}\quad\frac{255}{512},
\]
and the three-colour total
energy per site is exactly
\((9/2)(1-M^{-4})\) in these
finite free-field diagnostics.
The continuum Fourier-torus
normalization in the proof
does not rely on inferring
a limit from those two grids.

Fifty-four compact-score
tests verify invisibility
of constant two-forms.
There are 252 exact
pressure-average algebra
checks and 324 nonnegative
tail inequalities.
These are finite algebraic
regressions. The infinite-volume
classification, independence,
selected actual-law convergence
and uniform integrability
are established by the
analytic arguments above.

The certificate contains 5,494 bytes and has SHA--256
\begin{center}\scriptsize\ttfamily
5226E8950827A0720878761AF09BCAD990FFBFA2C684FFA97186FE0D4F3E3BBD.
\end{center}
The predecessor V89 contains 2,194,170 bytes and has SHA--256
\begin{center}\scriptsize\ttfamily
0B13330A2B097E42E346C3C77F04FC179DB343AF6718C2ED04D680C0367AFCB5.
\end{center}
Its body is preserved except
for the title version and
this terminal insertion.
V90 replaces V89 as the
sole active main note.
Earlier certificates,
companion and supplied
sources, monographs and
archives remain unchanged.

\begin{firewall}
The new result classifies
and sharply normalizes
selected thermal lattice
curvature limits. It does
not replace the requested
full Yang--Mills mass-gap
problem by a Gaussian
lattice construction.
The interacting physical
continuum, its reconstruction
and a positive physical
Hamiltonian mass gap remain
unproved. The full original
objective remains active.
\end{firewall}
\section{V91: ordered Wilson rectangles, growing-loop errors, and the boundary-current coefficient}
\label{sec:v91}

\subsection{Context: pass from colour coordinates to actual loop traces}

V90 identifies selected thermal
curvature laws and fixed-lattice
second moments. The physical
goal still requires information
on observables whose support
grows. This continuation proves
an actual-law estimate for
rectangular Wilson loops,
retaining the order of all
plaquette factors and their
parallel transports.

There are two distinct results.
First, an explicit error
bound compares an area-normalized
Wilson chord with averaged
thermal curvature and remains
valid for growing rectangles.
It also gives a quantitative
comparison of bounded functions
of the gauge-invariant loop
trace. Second, for each fixed
rectangle, V90 identifies its
thermal trace law and first
moment. The resulting Gaussian
coefficient has an exact
boundary-current representation
and perimeter-order bounds.

The second result cannot
automatically be substituted
into the first for a growing
rectangle: uniform covariance
control is still missing.
In particular this is neither
an interacting Wilson area
law nor a mass-gap proof.
The companion audit was
completed at V90; the next
scheduled checkpoint is V95.

\subsection{An ordered rectangular plaquette factorization}

Fix directions \(\mu<\nu\),
a base vertex \(x_0\), and
positive integer side lengths
\(m,n<M\). The rectangle
\(\mathcal R\) consists of
the \(A=mn\) plaquettes with
bases
\[
            x_{ij}=x_0+i e_\mu+j e_\nu,\qquad
                   0\le i<m,\quad0\le j<n .
\]
It may be placed anywhere
on the torus; the indicated
grid is embedded without
identifying different faces.
Write
\[
                    P_C=2(m+n),\qquad D=m+n-2
\]
for its perimeter and the
largest connector length
used below.

Let \(\mathcal U_C\) be the
holonomy of the positively
oriented boundary: first
\(m\) steps in direction
\(\mu\), then \(n\) in
\(\nu\), then the reverse
top and left sides. Let
\(\Gamma_{ij}\) be the
holonomy of the path from
\(x_0\) taking \(j\) steps
in direction \(\nu\) and
then \(i\) steps in \(\mu\).
It has length \(i+j\le D\).
Let \(P_{\mu\nu}(x_{ij})\)
be the based plaquette of V88.

\begin{lemma}[Exact ordered rectangle]
\label{lem:v91-factorization}
For arbitrary compact link
variables,
\begin{equation}\label{eq:v91-factorization}
 \mathcal U_C
     =\prod_{j=0}^{n-1}
               \bigl(g_{m-1,j}g_{m-2,j}\cdots g_{0,j}\bigr),
       \qquad
       g_{ij}=\Gamma_{ij}P_{\mu\nu}(x_{ij})\Gamma_{ij}^{-1},
\end{equation}
where the smaller row index
is farther to the left
in the outer product.
No commutativity is assumed.
\end{lemma}
\begin{proof}
Compare the two paths
\(\mu^m\nu^n\) and
\(\nu^n\mu^m\), with the
same endpoints. Their
path-product ratio is
\(\mathcal U_C\).
Move the first \(\nu\)
in the first word to
the left across all \(m\)
copies of \(\mu\), then
do the same for the next
\(\nu\), and continue.
For row \(j\), the swaps
occur at columns
\(m-1,m-2,\ldots,0\).
Each swaps adjacent steps
\(\mu\nu\) to \(\nu\mu\),
with prefix \(\nu^j\mu^i\).
The ratio of the two
path products at that
swap is exactly \(g_{ij}\):
the common tail cancels,
and the common prefix
remains as a conjugation.
Multiplying all consecutive
ratios telescopes to the
initial-to-final ratio.
Every face occurs once,
in the order displayed.
\end{proof}

This is a direct finite-word
form of a lattice Stokes
calculation, not a replacement
of the non-Abelian product
by an unordered sum.\footnote{
For broader lattice Stokes
formulations see M. A. Zubkov,
\emph{Abelian representation for
nonabelian Wilson loops and the
Non-Abelian Stokes theorem on the lattice},
\href{https://arxiv.org/abs/hep-lat/0212001}
{arXiv:hep-lat/0212001}.
The identity used here is
proved by the explicit
path-word calculation above.}

Write a unit quaternion as
\(\mathcal U_C=u_0I+\iota(u)\),
with \(u_0^2+|u|^2=1\).
Its normalized trace is
\(W_C=u_0\), and its
nonnegative defect is
\(s_C=1-W_C\).
The operator norm of its
chord obeys
\[
                        \|\mathcal U_C-I\|_{\rm op}^2=2s_C .
\]

\begin{proposition}[A deterministic trace bound]
\label{prop:v91-trace-bound}
For every such rectangle,
\begin{equation}\label{eq:v91-trace-bound}
                 s_C\le A\sum_{p\in\mathcal R}s_p .
\end{equation}
Under any of the actual
sequences considered below
with \(\mathbb E s_p\le20/\beta\),
\begin{equation}\label{eq:v91-trace-mean}
       \mathbb E s_C\le\min\{2,20A^2/\beta\},\qquad
       \mu\{s_C>\delta\}\le
                      \min\{1,20A^2/(\beta\delta)\}
                            \quad(\delta>0).
\end{equation}
\end{proposition}
\begin{proof}
Telescope the product in
\eqref{eq:v91-factorization}
minus the identity.
All its factors are unitary,
so
\[
 \|\mathcal U_C-I\|_{\rm op}
       \le\sum_{p\in\mathcal R}\|P_p-I\|_{\rm op}
       =\sum_{p\in\mathcal R}\sqrt{2s_p}.
\]
Square and apply
Cauchy--Schwarz, then
divide by two.
Expectation and Markov's
inequality prove
\eqref{eq:v91-trace-mean}.
No independence or mixing
is used.
\end{proof}

\subsection{A quantitative chord-to-curvature comparison}

Use the actual laws of V90
at \(\beta=\beta_M^\circ\);
the same bounds below also
apply whenever the earlier
actual budgets
\[
 \mathbb E s_p\le20/\beta,\qquad
 \mathbb E|Y_{M,p}|^2\le80,\qquad
 \mathbb E|a_e|^2\le m_{20}
       =144\pi^2\sqrt{20}\,\beta^{-1/2}+72\pi^2M^{-2}
\]
are available. Define
\begin{equation}\label{eq:v91-normalized-chord}
 Z_C=\frac{\sqrt\beta}{A}(u_0-1,u)\in\mathbb R^4,\qquad
 \overline Y_C=\frac1A\sum_{p\in\mathcal R}Y_{M,p}\in\mathbb R^3 .
\end{equation}
The vector part of \(Z_C\)
uses the colour frame at
\(x_0\); its norm, and
\(s_C\), are gauge invariant.
The untransported average
\(\overline Y_C\) uses the
fixed chart trivialization
of V88--V90 and is not
invariant under arbitrary
site-dependent gauge changes.
The comparison below is
made in that chart, with
the omitted transports
bounded explicitly.
A logarithm of the whole
loop is not needed, so
there is no loop cut-locus
assumption.

\begin{theorem}[A growing-rectangle error bound under the actual law]
\label{thm:v91-growing}
For every admissible rectangle,
\begin{equation}\label{eq:v91-growing}
 \mathbb E|Z_C-(0,\overline Y_C)|
       \le\varepsilon_{M,C}
       :=\frac{40A}{\sqrt\beta}
                             +2D\sqrt{80m_{20}} .
\end{equation}
Also
\begin{equation}\label{eq:v91-second-budgets}
                  \mathbb E|Z_C|^2\le40,\qquad
                  \mathbb E|\overline Y_C|^2\le80 .
\end{equation}
The constants do not depend
on the number of faces.
\end{theorem}
\begin{proof}
The actual plaquette
logarithms \(z_p\) have
\(|z_p|\le1\), as in V88.
Conjugation in
\eqref{eq:v91-factorization}
gives
\(g_p=\exp(\iota(\zeta_p))\),
where
\(\zeta_p=\operatorname{Ad}_{\Gamma_p}z_p\)
and \(|\zeta_p|=|z_p|\).
For the product
\(F(t)=\prod_p\exp(t\iota(\zeta_p))\)
in the exact rectangle order,
the unitary product-rule
estimate used in V88 gives
\[
       \|F''(t)\|_{\rm op}
                    \le\left(\sum_p|z_p|\right)^2 .
\]
Here \(F(0)=I\) and
\(F'(0)=\iota(\sum_p\zeta_p)\).
Taylor's formula therefore
gives
\[
 \left\|\mathcal U_C-I-\iota\Big(\sum_p\zeta_p\Big)\right\|_{\rm op}
                    \le\frac12\left(\sum_p|z_p|\right)^2 .
\]
For quaternion matrices
this norm is exactly the
Euclidean norm of the
four real coefficients.
After multiplying by
\(\sqrt\beta/A\), its
expected upper bound is
\[
 \frac{\sqrt\beta}{2A}
          A\sum_{p\in\mathcal R}\mathbb E|z_p|^2
                   \le\frac{40A}{\sqrt\beta}.
\]
No fourth moment enters.

For every connector,
V88's path estimate gives
\[
 \mathbb E|\operatorname{Ad}_{\Gamma_p}Y_{M,p}-Y_{M,p}|
                       \le2\,\operatorname{length}(\Gamma_p)
                                                   \sqrt{80m_{20}}.
\]
The lengths are at most
\(D\). Sum the errors and
divide by \(A\) to get
the second term in
\eqref{eq:v91-growing}.
This retains correlations
between the connector and
its plaquette.

Finally
\(|Z_C|^2=2\beta s_C/A^2\),
so Proposition~\ref{prop:v91-trace-bound}
gives the first second-moment
bound. For the other,
the convexity of squared
norm gives
\[
              |\overline Y_C|^2
                      \le\frac1A\sum_{p\in\mathcal R}|Y_{M,p}|^2 .
\]
Taking expectation proves
the claim.
\end{proof}

\begin{proposition}[A bounded-observable comparison for the actual trace]
\label{prop:v91-bounded-trace}
If \(g:[0,\infty)\to\mathbb R\)
has \(\|g\|_\infty\le1\)
and Lipschitz constant at
most one, then
\begin{equation}\label{eq:v91-bounded-trace}
 \left|\mathbb E g\left(\frac{\beta s_C}{A^2}\right)
       -\mathbb E g\left(\tfrac12|\overline Y_C|^2\right)\right|
                \le\min\{2,481\varepsilon_{M,C}^{\,2/3}\}.
\end{equation}
\end{proposition}
\begin{proof}
Put \(Z=Z_C\) and
\(V=(0,\overline Y_C)\).
On \(|Z|+|V|\le K\),
the difference of the two
test values is at most
\(K|Z-V|/2\).
On the complement it is
at most two.
By \eqref{eq:v91-second-budgets},
\(\mathbb E(|Z|+|V|)^2\le240\).
Hence the expectation
difference is at most
\[
                   \frac K2\varepsilon_{M,C}+\frac{480}{K^2}.
\]
For positive \(\varepsilon_{M,C}\),
choose \(K=\varepsilon_{M,C}^{-1/3}\)
and round \(480+1/2\) upwards.
The zero-error case is
immediate; the trivial
bound is two.
\end{proof}

For rectangles with both
side lengths at most \(L_M\),
\begin{equation}\label{eq:v91-mesoscopic}
 \varepsilon_{M,C}
      =O\left(L_M^2\beta_M^{-1/2}
                 +L_M\beta_M^{-1/4}+L_M/M\right).
\end{equation}
Consequently the errors vanish
uniformly for
\[
                    L_M=o(\min\{\beta_M^{1/4},M\}).
\]
This is a quantitative
growing-support statement
for the actual measure.
The scalar trace itself
also tends to one in mean
if \(A_M=o(\sqrt{\beta_M})\),
by \eqref{eq:v91-trace-mean}.
None of these statements
replaces the random average
\(\overline Y_C\) by a
Gaussian average for an
arbitrary growing rectangle.
That would require the
unproved moving-separation
covariance control.

\subsection{The actual thermal Wilson law for each fixed rectangle}

Let \(b_C\) be the scalar
two-form equal to one on
the oriented rectangular
surface and zero elsewhere.
For each colour define
the flux of V90's Gaussian
curvature by
\[
                    F_C^\alpha=\sum_{p\in\mathcal R}G_p^\alpha .
\]
The three flux components
are independent centered
Gaussians with common
variance
\begin{equation}\label{eq:v91-flux-variance}
                         \sigma_C^2=\langle b_C,Pb_C\rangle .
\end{equation}

\begin{theorem}[Fixed-rectangle Wilson trace limit and first moment]
\label{thm:v91-fixed-loop}
For every fixed rectangle,
under the actual selected
couplings of V90,
\begin{equation}\label{eq:v91-fixed-loop}
 \beta_M^\circ(1-W_C)
           \ \Longrightarrow\ \frac{\sigma_C^2}{2}\chi_3^2,\qquad
 \beta_M^\circ\,\mathbb E(1-W_C)
                             \longrightarrow\frac32\sigma_C^2 .
\end{equation}
The joint distribution of
any finite family of fixed
rectangular trace defects
is obtained from the joint
Gaussian fluxes, whose
colour covariance is
\begin{equation}\label{eq:v91-joint-flux}
              \mathbb E F_C^\alpha F_{C'}^\gamma
                        =\delta_{\alpha\gamma}\langle b_C,Pb_{C'}\rangle .
\end{equation}
No finite-volume convergence
of trace-defect second
moments is asserted.
\end{theorem}
\begin{proof}
For fixed \(A,D\),
\(\varepsilon_{M,C}\to0\).
Theorem~\ref{thm:v91-growing}
and V90's joint curvature
limit show
\[
           \sqrt\beta\,(u_0-1,u)
                     \ \Longrightarrow\ (0,F_C).
\]
Since
\(\beta(1-W_C)=
\frac12|\sqrt\beta(u_0-1,u)|^2\),
the continuous mapping
theorem proves the
distributional assertion.
It works simultaneously
for finitely many fixed
rectangles, giving
\eqref{eq:v91-joint-flux}.

For the expectation, retain
the stronger pointwise
bound
\[
                   \beta s_C\le
                         A\sum_{p\in\mathcal R}\beta s_p .
\]
The finitely many local
thermal energies on its
right are uniformly
integrable by V90.
Thus \(\beta s_C\) is
uniformly integrable.
Its mean therefore converges
to the mean of the displayed
limit, \(3\sigma_C^2/2\).
This does not provide
uniform integrability of
its square.
\end{proof}

\subsection{The coefficient is a boundary-current energy}

Let \(j_C=d_1^*b_C\), with
the scalar colour suppressed.
It is the signed unit
current on the rectangle
boundary. It is divergence
free, and
\(\|j_C\|_{\ell^2}^2=P_C\).
The Fourier form of
\eqref{eq:v90-projector}
gives
\begin{equation}\label{eq:v91-current-energy}
 \sigma_C^2
     =\int_{\mathbb T^4}
          \frac{|\widehat j_C(k)|^2}{\lambda(k)}
                                             \frac{dk}{(2\pi)^4}
     =\langle j_C,\Delta^{-1}j_C\rangle .
\end{equation}
The last expression is
the quadratic-form pairing
defined by the integral.
It is finite in four
dimensions. In particular
the variance depends only
on the boundary current:
if a different finite
surface has the same
\(d_1^*b\), it has the
same Gaussian flux in
distribution and, under
the common \(G\), almost
surely. Indeed the variance
of the flux difference
is zero by the same formula.
This is a statement about
the identified limiting
Gaussian coefficient,
not an assumption of
surface-independent unordered
non-Abelian products.

\begin{proposition}[Perimeter-order bounds on the limiting coefficient]
\label{prop:v91-perimeter}
For every finite rectangular
boundary on \(\mathbb Z^4\),
\begin{equation}\label{eq:v91-perimeter}
       \frac{P_C}{16}\le\sigma_C^2
             \le g_3P_C,\qquad
       g_3=\int_{\mathbb T^3}
          \frac{dk}{(2\pi)^3\sum_{i=1}^3|e^{ik_i}-1|^2}
                           \le\frac{\sqrt3\,\pi}{8}<1 .
\end{equation}
Thus the fixed-loop
limiting mean coefficient
in \eqref{eq:v91-fixed-loop}
is strictly positive and
between \(3P_C/32\) and
\((3/2)g_3P_C\).
\end{proposition}
\begin{proof}
Since \(\lambda(k)\le16\),
Parseval and
\eqref{eq:v91-current-energy}
give the lower bound.

For the upper bound let
\[
 G_d(x)=\int_0^\infty p_t^{(d)}(x)\,dt
       =\int_{\mathbb T^d}
            \frac{e^{ik\cdot x}}{\lambda_d(k)}
                                             \frac{dk}{(2\pi)^d},
                               \qquad d=3,4 .
\]
These integrals are finite:
near zero
\(\lambda_d(k)\ge4|k|^2/\pi^2\),
so \(\lambda_d^{-1}\)
is integrable for \(d>2\).
Fubini is justified by
that integrability, and
the heat representation
shows \(G_d(x)\ge0\).
The heat kernel factors
by coordinate directions.
Consequently Tonelli's
theorem gives the exact
line-sum identity
\begin{equation}\label{eq:v91-line-sum}
              \sum_{r\in\mathbb Z}G_4(r e_\mu)
                    =\int_0^\infty p_t^{(3)}(0)\,dt=g_3 .
\end{equation}
Here the sum of the
one-dimensional heat kernel
over its coordinate is one.

The two \(\mu\)-directed
boundary segments have
opposite signs. Their
combined current energy
is twice the same-segment
energy, minus twice a
nonnegative cross energy.
Different orientation
components do not pair
in the componentwise
scalar \(\Delta^{-1}\).
Thus the two \(\mu\)
segments contribute at
most
\[
 2\sum_{r,s=0}^{m-1}G_4((r-s)e_\mu)
                            \le2m\,g_3 .
\]
The other two segments
contribute at most
\(2n\,g_3\).
This proves \(\sigma_C^2\le P_Cg_3\).

Finally on \([-\pi,\pi]^3\)
use \(\lambda_3(k)\ge4|k|^2/\pi^2\),
and enlarge the cube to
the ball of radius
\(\sqrt3\pi\). Then
\[
 g_3\le\frac{\pi^2}{4(2\pi)^3}
                      \int_{|k|\le\sqrt3\pi}|k|^{-2}\,dk
       =\frac{\sqrt3\,\pi}{8}<1 .
\]
The lattice symbol estimate
follows from
\(\sin t\ge2t/\pi\) for
\(0\le t\le\pi/2\).
\end{proof}

For a single plaquette,
\(\sigma_C^2=1/2\), so
\eqref{eq:v91-fixed-loop}
reduces to V90's
\(\tfrac14\chi_3^2\)
energy law and mean \(3/4\).
The bounds for larger
rectangles concern the
family of coefficients
obtained \emph{after}
taking each fixed-rectangle
thermal limit. They do
not interchange this limit
with a growing-loop or
fixed-physical-size limit.

\subsection{The remaining physical-scale covariance question}

The new actual-law growing
estimate is
\eqref{eq:v91-bounded-trace};
the new Gaussian coefficient
is \eqref{eq:v91-current-energy}.
The former still contains
the actual random flux
\(\overline Y_C\).
The latter is not yet
a uniform approximation
to that flux for
\(C=C_M\) growing with
the cutoff.

For example, with
\(\beta_M^\circ\le\log M\),
the proved geometric
range \(L_M=o(\beta_M^{1/4})\)
is much smaller than
rectangles of side of
order \(M\) in a fixed
physical box.
Even within the geometric
range, uniform distributional
comparison of the actual
growing flux to the
Gaussian flux remains
to be proved.
At fixed physical scale
one must also retain the
correct coupling trajectory,
observable normalization,
chart coverage and physical
reconstruction.
Calling the coefficient
bound a Yang--Mills
perimeter law or an
area-law exclusion would
therefore be unjustified.

\subsection{Exact verification and preservation}

The certificate is
\begin{center}\small
\path{scripts/verify_ym_f14_v91_wilson_rectangles.py}.
\end{center}
Twenty-four rectangle words
cover all six coordinate
planes and four side-length
pairs. Free-word reduction
checks their ordered
factorization, face count,
connector lengths and
boundary current length.
Seventy-two rational
noncommuting quaternion
configurations independently
check the product and
the pointwise trace bound.
Removing connectors fails
in all fifty-four cases
larger than one plaquette.

Five complete side-four
Fourier regressions compare
the surface-projector and
boundary-current variance
formulas exactly. Their
values for sides
\((1,1),(1,2),(2,2),(2,3),(3,3)\)
are, respectively,
\[
            \frac{255}{512},\quad
            \frac{381}{448},\quad
            \frac{21}{16},\quad
            \frac{3641}{2240},\quad
            \frac{927}{512}.
\]
These are finite-volume
Gaussian diagnostic values,
not the infinite-volume
\(\sigma_C^2\).
The finite periodic
four-to-three-dimensional
Green line sum is also
checked exactly, giving
\(1517/7680\).
Sixty-four rational checks
verify the bounded-test
constant, and four exponent
checks confirm examples
within the stated growth
range. None of these finite
tests certifies a physical
continuum limit by sampling.

The certificate contains 6,638 bytes and has SHA--256
\begin{center}\scriptsize\ttfamily
581FC9B2EDD4B9950E5CFB4261F3546E028ACD5FA5F6B5DA43700610F737B80A.
\end{center}
The predecessor V90 contains 2,216,933 bytes and has SHA--256
\begin{center}\scriptsize\ttfamily
F6ED8010EC46396E993F7297669A094F0EB735DE486D1BCD21DA14289A23693E.
\end{center}
Its body is preserved except
for the title version and
this terminal insertion.
V91 replaces V90 as the
sole active main note.
Earlier certificates,
companion and supplied
sources, monographs and
archives remain unchanged.

\begin{firewall}
The proved progress is an
actual growing-rectangle
geometric error estimate
and a fixed-rectangle
thermal Wilson-loop limit,
with its boundary-current
coefficient identified.
An interacting physical
Wilson law, full continuum
reconstruction and a
positive physical Hamiltonian
mass gap are still unproved.
The original full objective
remains active.
\end{firewall}
\section{V92: receding-window independence and uniform thermal covariance comparison}
\label{sec:v92}

\subsection{Context: remove a qualitative moving-separation bottleneck}

V91 leaves the actual curvature covariance at moving separations
uncontrolled. This section resolves its qualitative version for the
same selected, determinant-weighted small-chart laws as V90.
The additional observation is that V89's compact thermal moves can
be placed in several widely separated windows: the error bounds
depend on the number and shape of the moved plaquettes, not on the
distance between those windows. The joint limiting move identities,
together with V90's already identified marginal law, force the
different windows to become independent.

The consequences are uniform convergence of the full covariance
kernel to its periodic Gaussian comparator, a mean-square law of
large numbers for suitably normalized growing probes, and an
area-squared-normalized mean Wilson comparison. No numerical decay
rate is obtained. In particular, covariance errors multiplied by
physical fluctuation normalizations are still uncontrolled.
The full interacting continuum mass-gap objective is unchanged.
The companion audit was completed at V90; the next checkpoint is V95.

There is an older uniform covariance comparison in V89 of archive
\path{Renewal_Yang_Mills_Mass_Gap_Research_Notes_V100_f11.tex}.
That result uses a uniformly convex fluctuation-box Hessian and
exponentially decaying inverse columns. It concerns a different
local response problem. Those hypotheses have not been established
for the present determinant-weighted Coulomb chart; indeed the
Wilson-action convexity obstruction of the current V63 must still
be respected. The receding-window argument below does not import
those convexity or exponential-decay assumptions.

\subsection{Several thermal moves at mutually receding locations}

Throughout, \(\beta=\beta_M^\circ\) is V90's selected sequence,
\(\mu_M\) is the actual normalized chart law, and
\[
          Y_{M,p}=\sqrt{\beta_M^\circ}\,z_p,\qquad
          X_{M,p}=\beta_M^\circ s_p .
\]
Extend the torus fields periodically to \(\mathbb Z^4\).
All statements concern this sequence, not a new normalizer or an
unproved replacement by the full compact Wilson measure.
Write
\[
 d_M(x,y)=\min_{n\in\mathbb Z^4}|x-y+Mn|_\infty
\]
for the periodic lattice distance. A window based at \(c\) is the
coordinate field \(Y_M^{[c]}(x)=Y_M(c+x)\); each window keeps the
fixed chart colour trivialization.

\begin{lemma}[Uniformity of finitely many separated compact moves]
\label{lem:v92-moves}
Fix \(r<\infty\), finitely supported real colour two-forms
\(B_1,\ldots,B_r\), and centres \(c_{j,M}\) satisfying
\begin{equation}\label{eq:v92-separation}
            \min_{i\ne j}d_M(c_{i,M},c_{j,M})\longrightarrow\infty .
\end{equation}
Put \(h_j=d_1^*B_j\), \(\ell_j=d_1h_j\), and let
\[
 h_M=\beta^{-1/2}\sum_{j=1}^r \tau_{c_{j,M}}h_j ,
\]
where \(\tau_c\) places a finitely supported form at \(c\).
There are good events with probability tending to one on which
both segments \(a\pm t h_M\), \(0\le t\le1\), remain in the
original chart. On those events the logarithmic geometric weight
increments tend to zero in probability. Moreover
\begin{equation}\label{eq:v92-action}
 \beta\{S(a\pm h_M)-S(a)\}
       \mp\sum_{j=1}^r\langle Y_M^{[c_{j,M}]},\ell_j\rangle
       -\frac12\sum_{j=1}^r\|\ell_j\|_2^2
                     \longrightarrow0
                     \quad\hbox{in probability}.
\end{equation}
In every fixed finite set of coordinates of window \(j\), the
curvature increment is \(\pm\ell_j+o_{L^2}(1)\).
\end{lemma}
\begin{proof}
Let \(E_j\) be the set of edges bordering the support of \(B_j\),
and include the plaquettes touching those edges in a fixed enlarged
patch. By \eqref{eq:v92-separation}, their translated enlarged
patches are disjoint for all sufficiently large \(M\).
Set
\[
 N=\sum_j|E_j|,\qquad b=\max_{j,p}|(B_j)_p|,\qquad
 \rho=b/\sqrt\beta,\qquad \eta=6(1+3R)\rho .
\]
Zero forms can be omitted. V85's compressed edge-trace argument
applies to the union of the moved edges. It uses \(N\), \(b\) and
the local incidence bound, but never the diameter of that union.
Consequently V89's bounds hold in the form
\begin{align*}
 \mu_M(G_M^c)&\le N(4m_{20}/R^2+2\eta k^\sharp),\\
 \mathbb E\!\left[\mathbf1_{G_M}
           \left|\log\frac{w(a\pm h_M)}{w(a)}\right|\right]
       &\le 2\eta N k^\sharp+12R\rho N .
\end{align*}
Here \(m_{20}\to0\) and \(k^\sharp/\sqrt\beta\to0\) by V89.
Thus both bounds vanish uniformly in the centres. The construction
of the good event retains the radius guard, the gauge-positive
boundary and the original determinant.

The local plaquette-logarithm Hessian bound in
Lemma~\ref{lem:v89-action} is also independent of location.
Its curvature-increment error in each changed plaquette is bounded
in squared expectation by a constant, depending only on the fixed
moves, times \(m_{20}+\beta^{-1}\). There are only finitely many
changed plaquettes. Their unshifted curvature coordinates are
uniformly tight by \(\mathbb E|Y_{M,p}|^2\le80\), and the shifted
ones are tight by this increment bound. The same bounded-curvature
truncation as in V89 therefore controls the action's quartic Taylor
remainders in probability, without fourth-moment convergence.
The quadratic action increments add over the disjoint enlarged
patches. Their cross terms are zero, yielding \eqref{eq:v92-action}.

For any prescribed finite window coordinates, sufficiently distant
other moves change none of their incident links. Coordinates of
that window changed by its own move obey the stated local
\(L^2\) estimate; all remaining ones are unchanged. This proves
the final assertion.
\end{proof}

\begin{theorem}[Independent Gaussian limits in receding windows]
\label{thm:v92-windows}
For centres satisfying \eqref{eq:v92-separation},
\begin{equation}\label{eq:v92-window-limit}
       (Y_M^{[c_{1,M}]},\ldots,Y_M^{[c_{r,M}]})
                 \ \Longrightarrow\ (G^{(1)},\ldots,G^{(r)})
\end{equation}
on the product of \(r\) countable coordinate spaces.
The \(G^{(j)}\) are independent copies of V90's centered Gaussian
curvature field, with covariance \(P\otimes I_3\).
The assertion includes joint laws of arbitrary fixed finite patches
in those windows, not only the Laplacian-curvature scores.
\end{theorem}
\begin{proof}
Stationarity and V90 imply that each window separately converges to
the law of \(G\). The local second-moment bound gives joint
tightness on the countable product: choose successively larger
coordinate bounds with summably small tail probabilities.
Take any subsequential joint limit \(\nu\), with coordinate fields
\(Y^{(1)},\ldots,Y^{(r)}\). Its marginals are all exactly the
Gaussian law identified in V90.

For the move of Lemma~\ref{lem:v92-moves}, extend the actual
normalized density by zero, as in V89. Its bounded acceptance
identity is exact. The lemma allows passage to the limit for
bounded continuous cylinder functions of all \(r\) windows.
Writing
\[
 L=(\ell_1,\ldots,\ell_r),\quad
 s(y)=\sum_j\langle y^{(j)},\ell_j\rangle,\quad
 v=\sum_j\|\ell_j\|_2^2,\quad
 a_\pm(y)=\min\{1,e^{\mp s(y)-v/2}\},
\]
the result is \(a_+\nu=(T_{-L})_*(a_-\nu)\).
The acceptance weights are strictly positive. Dividing this
identity as in Theorem~\ref{thm:v89-translation}, and replacing
\(L\) by \(-L\), gives
\begin{equation}\label{eq:v92-joint-RN}
               \frac{d(T_L)_*\nu}{d\nu}(y)
                        =e^{s(y)-v/2}.
\end{equation}
There is no division of finite-volume small probabilities here:
the identity is first passed through bounded acceptance weights,
then divided as an identity of limiting measures.

Replace each \(B_j\) by \(t_jB_j\), for arbitrary real \(t_j\),
and integrate \eqref{eq:v92-joint-RN}. It follows that
\begin{equation}\label{eq:v92-score-mgf}
 \mathbb E_\nu\exp\left\{\sum_jt_j
                         \langle Y^{(j)},\ell_j\rangle\right\}
       =\exp\left\{\frac12\sum_j t_j^2\|\ell_j\|_2^2\right\}.
\end{equation}
Linear combinations of any finite collection of scores within a
window are still of the form \(d_1d_1^*B_j\).
Thus all such scores are jointly Gaussian, with independence
between different windows.

It remains to show that these scores determine each marginal
field; Gaussian one-window marginals alone would not imply joint
independence. For a finite two-form \(b\), put
\[
 R_tb=\int_0^t e^{-s\Delta}b\,ds .
\]
This form belongs to \(\ell^1\). Finite truncations \(C_{t,n}\)
converge to it in \(\ell^2\), and \(K=d_1d_1^*\) is a bounded
\(\ell^2\) operator. Its Fourier symbol gives
\[
 KR_tb=(I-e^{-t\Delta})Pb\longrightarrow Pb
                               \quad\hbox{in }\ell^2,\quad t\to\infty .
\]
The convergence follows from dominated convergence and
\(\lambda(k)>0\) almost everywhere. Choosing a diagonal sequence
of finite \(C_{t,n}\), obtain scores
\(\ell_n=KC_{t,n}\to Pb\) in \(\ell^2\).
In the already known marginal Gaussian law,
\[
 \mathbb E_\nu\big|
   \langle Y^{(j)},b\rangle-\langle Y^{(j)},\ell_n\rangle
                 \big|^2
      =\|Pb-\ell_n\|_2^2\longrightarrow0 .
\]
Here \(P\ell_n=\ell_n\). No sample of \(Y^{(j)}\) is assumed
square summable. A finite number of these approximations converge
jointly in probability, by their marginal \(L^2\) bounds.
The independence of the approximating score vectors from
\eqref{eq:v92-score-mgf} passes to their limit. Hence all finite
coordinate distributions factor into the required Gaussian
marginals. They determine \(\nu\), so every subsequential limit is
the same product law. This proves \eqref{eq:v92-window-limit}.
\end{proof}

For bounded continuous local functions \(\Phi,\Psi\), this implies
the qualitative clustering statement
\begin{equation}\label{eq:v92-local-clustering}
 \lim_{D\to\infty}\limsup_{M\to\infty}
       \sup_{d_M(x,y)\ge D}
 \left|\operatorname{Cov}_{\mu_M}
   \big(\Phi(Y_M^{[x]}),\Psi(Y_M^{[y]})\big)\right|=0 .
\end{equation}
Indeed failure would supply a sequence of centres with distance
tending to infinity and covariance bounded away from zero,
contrary to the theorem and bounded convergence of the relevant
test expectations. An empty supremum is taken to be zero.
This is not an exponential mixing estimate: neither a rate in
\(D\) nor a uniform finite-cutoff mass has been produced.

\subsection{A periodic comparator whose long-distance kernel is uniformly small}

Combine orientation and colour into an index \(a\) with eighteen
values, and put \(\mathsf P(k)=P(k)\otimes I_3\), with
\(\mathsf P(0)=0\). The finite torus operator \(\mathsf P_M\)
has this multiplier on the grid
\(\mathcal K_M=(2\pi/M)\mathbb Z_M^4\).
It is an orthogonal projection. Define its covariance kernel by
\begin{equation}\label{eq:v92-comparator}
 K_M^{ab}(z)=[\mathsf P_M]_{(0,a),(z,b)}
       =M^{-4}\sum_{k\in\mathcal K_M}
                             e^{-ik\cdot z}\mathsf P_{ab}(k).
\end{equation}
The kernels are real, since the multiplier at \(-k\) is the
complex conjugate of the multiplier at \(k\).
Let \(K_\infty\) denote the analogous Fourier integral on
\(\mathbb T^4\).

\begin{lemma}[Uniform tails of the discrete Fourier kernels]
\label{lem:v92-kernel}
For fixed \(z\in\mathbb Z^4\), \(K_M^{ab}(z)\to K_\infty^{ab}(z)\).
Also
\begin{equation}\label{eq:v92-kernel-tail}
 \lim_{D\to\infty}\limsup_{M\to\infty}
     \sup_{d_M(z,0)\ge D}|K_M^{ab}(z)|=0 .
\end{equation}
\end{lemma}
\begin{proof}
Every multiplier entry is bounded and continuous except possibly
at the single point \(k=0\). It is therefore Riemann integrable.
The same holds after multiplication by a fixed Fourier character,
which proves the fixed-\(z\) assertion by grid Riemann sums.

For the tail assertion, write \(f=\mathsf P_{ab}\). Given
\(\epsilon>0\), choose a trigonometric polynomial \(q\) with
\[
             \int_{\mathbb T^4}|f-q|\,\frac{dk}{(2\pi)^4}
                                                     <\epsilon .
\]
To see existence directly, multiply \(f\) by a continuous cutoff
vanishing near zero. The removed integral tends to zero with the
neighbourhood's volume. The resulting continuous periodic function
is uniformly approximable by trigonometric polynomials.
Since \(|f-q|\) is also Riemann integrable,
\[
              M^{-4}\sum_{k\in\mathcal K_M}|f(k)-q(k)|
                                                    <2\epsilon
\]
for all sufficiently large \(M\).
Let all Fourier frequencies of \(q\) have infinity norm at most
\(D_q\). Discrete character orthogonality implies that its sampled
Fourier coefficient at \(z\) is zero whenever
\(d_M(z,0)>D_q\), once \(M>2D_q\).
This statement includes periodic aliasing: such a \(z\) is not
congruent to any frequency of \(q\).
At these \(z\), \eqref{eq:v92-comparator} has absolute value at
most the last discrete \(\ell^1\) error, hence at most
\(2\epsilon\). Let \(\epsilon\downarrow0\).
\end{proof}

\subsection{Uniform covariance convergence, including moving separations}

The actual means \(\mathbb E_{\mu_M}Y_{M,a}(x)\) are zero by
common colour-rotation symmetry. Set
\begin{equation}\label{eq:v92-delta}
 \delta_M=
 \max_{a,b}\sup_{z\in\mathbb Z_M^4}
 \left|\mathbb E_{\mu_M}\{Y_{M,a}(0)Y_{M,b}(z)\}
                                      -K_M^{ab}(z)\right|.
\end{equation}

\begin{theorem}[Uniform entrywise covariance comparison]
\label{thm:v92-covariance}
Along V90's actual selected sequence,
\begin{equation}\label{eq:v92-covariance}
                         \delta_M\longrightarrow0 .
\end{equation}
Consequently, for every real deterministic torus two-form \(B\),
with the scalar component norms
\(\|B\|_1=\sum_{x,a}|B_a(x)|\) and
\(\|B\|_2^2=\sum_{x,a}|B_a(x)|^2\),
\begin{equation}\label{eq:v92-quadratic}
 \left|\mathbb E_{\mu_M}\langle Y_M,B\rangle^2
                    -\langle B,\mathsf P_MB\rangle\right|
                              \le\delta_M\|B\|_1^2 .
\end{equation}
In particular the variance is at most
\(\|B\|_2^2+\delta_M\|B\|_1^2\).
\end{theorem}
\begin{proof}
First, products of any two possibly moving components are uniformly
integrable. V90 proved uniform integrability of \(X_{M,p}\) and of
\(|Y_{M,p}|^2\); stationarity gives the same tail bounds at every
translated plaquette. There are only finitely many orientations
and colours. The inequality
\[
 |Y_{M,a}(0)Y_{M,b}(z)|
                 \le\tfrac12\{Y_{M,a}(0)^2+Y_{M,b}(z)^2\}
\]
therefore gives uniform integrability, independently of \(z\).
For completeness, the family of sums of two nonnegative uniformly
integrable families is uniformly integrable, since
\[
 (u+v)\mathbf1_{\{u+v>K\}}
       \le2u\mathbf1_{\{u>K/2\}}+2v\mathbf1_{\{v>K/2\}} .
\]

Suppose \eqref{eq:v92-covariance} fails. Extract a subsequence,
fixed indices \(a,b\), and separations \(z_M\) on which the
difference stays bounded away from zero.
If \(d_M(z_M,0)\) is bounded along a further subsequence, choose
bounded integer representatives and then a constant representative
\(z\). V90's fixed-coordinate convergence and the uniform
integrability just proved imply convergence of the actual
covariance to \(K_\infty^{ab}(z)\).
Lemma~\ref{lem:v92-kernel} gives the same limit for \(K_M^{ab}(z)\).

Otherwise take a further subsequence with
\(d_M(z_M,0)\to\infty\). Theorem~\ref{thm:v92-windows} with two
windows gives independent centered Gaussian limits of the two
coordinates. Uniform integrability makes their product expectation
tend to zero. Lemma~\ref{lem:v92-kernel} makes the comparator tend
to zero as well. Both cases contradict the selected positive
discrepancy. This proves \eqref{eq:v92-covariance}.

Expand the variance of \(\sum_{x,a}B_a(x)Y_{M,a}(x)\).
Stationarity and \eqref{eq:v92-delta} bound the absolute difference
of every covariance entry by \(\delta_M\); summing absolute
coefficients gives \eqref{eq:v92-quadratic}.
Finally \(0\preceq\mathsf P_M\preceq I\).
\end{proof}

This establishes a uniform conclusion, but the proof is a
compactness argument and gives no rate for \(\delta_M\).
It is stronger than fixed-separation convergence and weaker than
the summable or weighted operator estimates required below.

\begin{proposition}[Mean-square disappearance of spreading averages]
\label{prop:v92-averages}
If \(\sup_M\|B_M\|_1<\infty\) and \(\|B_M\|_2\to0\), then
\(\langle Y_M,B_M\rangle\to0\) in \(L^2(\mu_M)\).
In particular, for any rectangular surface with \(A\) faces, let
\(b_C\) be its scalar indicator and let \(P_M\) now denote the
one-colour projection. Then
\begin{equation}\label{eq:v92-average}
 \left|\mathbb E|\overline Y_C|^2
            -\frac{3}{A^2}\langle b_C,P_Mb_C\rangle\right|
                     \le3\delta_M,\qquad
 \mathbb E|\overline Y_C|^2\le\frac3A+3\delta_M .
\end{equation}
Thus \(\overline Y_{C_M}\to0\) in mean square whenever
\(A_M\to\infty\), without a restriction on the rectangle's diameter
other than its being an admissible torus rectangle.
\end{proposition}
\begin{proof}
The first assertion follows from the last variance bound in
Theorem~\ref{thm:v92-covariance}.
For each of the three colours of a rectangle, use \(B=b_C/A\)
in that colour, with \(\|B\|_1=1\).
The comparator variances add to the expression in
\eqref{eq:v92-average}, and the three absolute errors add to at
most \(3\delta_M\). Since \(P_M\) is a projection and
\(\|b_C\|_2^2=A\), its quadratic form is at most \(A\).
\end{proof}

\subsection{A growing-loop mean comparison, with its normalization retained}

V90's local energy uniform integrability supplies the finite
nonnegative tail modulus
\[
 \mathcal T(K)=\sup_{M,p}\mathbb E_{\mu_M}(X_{M,p}-K)_+,
                    \qquad \mathcal T(K)\longrightarrow0
                              \quad(K\to\infty).
\]
The supremum is over the eventual selected sequence; adding any
defined finite prefix does not change the limit.
For every collection of \(A\) faces, convexity gives
\[
 \mathbb E\left(\frac1A\sum_pX_{M,p}-K\right)_+
                                  \le\mathcal T(K).
\]
Hence these averages are uniformly integrable even when \(A\)
and the locations vary.

For V91's rectangle chord \(Z_C\), put
\(V_C=(0,\overline Y_C)\) and retain its proved bound
\[
 \mathbb E|Z_C-V_C|\le\varepsilon_{M,C}
       =40A/\sqrt\beta+2D\sqrt{80m_{20}} .
\]
Define a nonnegative modulus
\begin{equation}\label{eq:v92-modulus}
 \mathcal R(e)=\inf_{T>0}
       \{Te+24\mathcal T(T^2/24)\}.
\end{equation}
It satisfies \(\mathcal R(e)\to0\) as \(e\downarrow0\): first
choose \(T\) so that the second term is small, then decrease \(e\).

\begin{proposition}[Actual quadratic chord error and mean trace]
\label{prop:v92-wilson}
For every admissible rectangle,
\begin{equation}\label{eq:v92-L2-chord}
                  \mathbb E|Z_C-V_C|^2
                         \le\mathcal R(\varepsilon_{M,C})
\end{equation}
and
\begin{equation}\label{eq:v92-mean-wilson}
 \left|\frac{\beta\,\mathbb E(1-W_C)}{A^2}
                 -\frac{3}{2A^2}\langle b_C,P_Mb_C\rangle\right|
       \le\sqrt{60\mathcal R(\varepsilon_{M,C})}
                                     +\frac32\delta_M .
\end{equation}
Thus, if \(A_M\to\infty\) and \(\varepsilon_{M,C_M}\to0\),
\begin{equation}\label{eq:v92-vanishing-loop}
              \frac{\beta_M^\circ\,\mathbb E(1-W_{C_M})}{A_M^2}
                                                    \longrightarrow0 .
\end{equation}
\end{proposition}
\begin{proof}
Set
\[
                       U=A^{-1}\sum_{p\in C}X_{M,p}.
\]
The pointwise trace bound of V91 and the plaquette bound of V88 give
\[
       |Z_C|^2\le2U,\qquad
       |V_C|^2\le A^{-1}\sum_{p\in C}|Y_{M,p}|^2\le4U .
\]
Here the sum is over the rectangular surface faces.
Consequently \(d=|Z_C-V_C|\) satisfies \(d^2\le12U\).
For \(T>0\), split at \(d=T\). On \(d\le T\),
\(d^2\le Td\). On \(d>T\), \(12U>T^2\), so
\[
          d^2\le12U\le24(U-T^2/24)_+ .
\]
Adding the two bounds, taking expectation, and using convexity
of the positive part yields
\[
       \mathbb E d^2\le
          T\varepsilon_{M,C}+24\mathcal T(T^2/24).
\]
Taking the infimum proves \eqref{eq:v92-L2-chord}.

By V91, \(\mathbb E(|Z_C|+|V_C|)^2\le240\).
Cauchy--Schwarz therefore gives
\[
 \frac12\left|\mathbb E|Z_C|^2-\mathbb E|V_C|^2\right|
                   \le\sqrt{60\,\mathbb E|Z_C-V_C|^2}.
\]
Use \(\tfrac12|Z_C|^2=\beta(1-W_C)/A^2\) and
\eqref{eq:v92-average} to obtain \eqref{eq:v92-mean-wilson}.
Its comparator is at most \(3/(2A)\).
The assumptions and the vanishing of both moduli prove
\eqref{eq:v92-vanishing-loop}.
\end{proof}

For example, both side lengths may tend to infinity while staying
\(o(\min\{\beta_M^{1/4},M\})\), the geometric range of V91.
This is a mean statement for the actual growing Wilson trace,
not just a bounded-test comparison. However its absolute error is
measured \emph{after division by \(A^2\)}. It does not imply a
relative approximation to a perimeter-order coefficient, a central
limit theorem for a growing flux, or an area or perimeter law at
fixed physical size. The moduli \(\delta_M\) and \(\mathcal T\)
are proved to vanish in their respective limits, but have no
effective rate here.

\subsection{Why the remaining physical normalization cannot be dropped}

To make the missing rate explicit, set \(a=M^{-1}\) in a fixed
unit physical torus, and take a smooth real test function \(f\).
For a single curvature component the usual candidate pairing
has deterministic coefficients
\[
        B_M(x)=a^2 f(ax),\qquad
        \|B_M\|_2^2=O(1),\qquad \|B_M\|_1=O(M^2).
\]
Equation~\eqref{eq:v92-quadratic} then bounds the covariance error
only by \(O(\delta_M M^4)\).
The proved fact \(\delta_M\to0\) does not make this quantity
vanish. A rate \(\delta_M=o(M^{-4})\) would suffice for this
particular entrywise route, but is not claimed or shown necessary;
a suitable weighted covariance-operator estimate could be sharper.
Even such an estimate would still concern the present thermal
chart law, unless the physical coupling trajectory, chart coverage,
observable identification and reconstruction were separately proved.

\begin{proposition}[A vanishing common mode still defeats physical scaling]
\label{prop:v92-counterexample}
There are Gaussian comparison fields with the same local and
receding-window limits as above, with uniformly integrable local
thermal quadratic energies and uniform covariance error tending
to zero, whose constant-test physically normalized curvature
variance diverges.
\end{proposition}
\begin{proof}
Let \(\mathcal G_M\) be the finite Gaussian field of covariance
\(\mathsf P_M\), and independently let \(H\) be an eighteen-component
standard real Gaussian vector. Define the comparison field
\[
                  \widetilde Y_{M,a}(x)
                      =\mathcal G_{M,a}(x)+M^{-1}H_a .
\]
It is stationary and closed, and common colour rotations preserve
its law. Its covariance differs from \eqref{eq:v92-comparator} by
\(M^{-2}\delta_{ab}\) at every separation. Local variances are
bounded, so all local quadratic energies are uniformly integrable.
The fixed-coordinate limits follow from the Riemann-sum part of
Lemma~\ref{lem:v92-kernel}; the same lemma makes different
receding windows asymptotically independent, since the common-mode
covariance \(M^{-2}\) also vanishes. These are finite Gaussian
vectors, so covariance convergence proves both assertions.

In fact the comparison retains the exact finite coexact
translation identity: \(\ell=d_1d_1^*B\) belongs to
\(\operatorname{Ran}\mathsf P_M\) and has zero spatial sum.
On this range \(\mathcal G_M\) is a standard Gaussian, while \(H\)
is independent and unchanged by translation by \(\ell\).
Its translated density is therefore
\(e^{\langle\widetilde Y_M,\ell\rangle-\|\ell\|_2^2/2}\).
Thus local coexact identities alone cannot detect the added mode.

The comparator has no zero momentum, so
\(\sum_x\mathcal G_{M,a}(x)=0\) almost surely.
For the smooth constant test \(f=1\),
\[
       a^2\sum_x\widetilde Y_{M,a}(x)=M H_a,\qquad
       \operatorname{Var}\left(a^2\sum_x\widetilde Y_{M,a}(x)\right)
                                                       =M^2 .
\]
This diverges despite the uniform covariance error \(M^{-2}\to0\).
The mean of \(\tfrac12\sum_{a=1}^{18}\widetilde Y_{M,a}(0)^2\)
is \((9/2)(1-M^{-4})+9M^{-2}\), tending to V90's thermal energy
floor. To check the Gaussian part, the trace of \(P(k)\) is three
at every nonzero momentum, and symmetry divides its diagonal
mean equally among the six orientations.
\end{proof}

This is a comparison counterexample, not a claim that the actual
chart law contains this mode. Unlike V88's unscaled constant
diagnostic, its common mode disappears in every local and
receding-window limit; only the physical normalization reveals it.
It shows precisely why the new uniform qualitative covariance
theorem cannot be substituted for a physical fluctuation bound.

\subsection{Verification, preservation and next upstream estimate}

The exact diagnostic script is
\begin{center}\small
\path{scripts/verify_ym_f14_v92_receding_windows.py}.
\end{center}
Fifty-four periodic three-window checks cover all eighteen
orientation-colour scores and three torus sizes, including supports
crossing a periodic coordinate boundary. They check disjointness,
zero orientation-wise sums, closure, and the exact additive
score variance. A side-four Fourier calculation checks all
256 polynomial coefficients and 253 off-support kernel bounds,
with discrete symbol error \(3793/17920\).
Twenty-four signed quadratic-form checks verify the
entrywise-to-\(\ell^1\) bound; 384 rational tests check the tail
inequalities used for the quadratic chord upgrade.
Five exact common-mode scaling checks give covariance errors
\(M^{-2}\) but physical constant-test variances \(M^2\).
These tests check finite identities and constants. They do not
replace the compactness, uniform-integrability or limiting-measure
proofs, and do not sample the actual Yang--Mills law.

The certificate contains 6,404 bytes and has SHA--256
\begin{center}\scriptsize\ttfamily
D84E6A40340FDEF18A18A03D2487FC52E1466766FF5B3CACF5DAEB844D6BE792.
\end{center}
The predecessor V91 contains 2,236,230 bytes and has SHA--256
\begin{center}\scriptsize\ttfamily
2A39958DAC5007A22140222FC69306BD5A7FE9BF8CCF745547AA98A979705401.
\end{center}
Its body is preserved except for the title version and this
terminal insertion. V92 replaces V91 as the sole active main note.
All earlier certificates, supplied sources, companion notes,
monographs and archives are preserved.

The qualitative covariance question posed at the end of V91 has
now been answered for the selected thermal law; it must not be
listed again as wholly open. What remains upstream is effective
control of covariance in a norm compatible with physical test
normalization, including long-wavelength or near-zero modes.
The constants in the finite-window move bounds are explicit,
but the passage from finitely many scores to the whole curvature
field used nonquantitative approximation and compactness.
Quantifying that passage, while retaining the actual measure and
its chart boundary, is a concrete next target.

\begin{firewall}
V92 proves independent limits of receding finite windows and
uniform entrywise covariance comparison under the actual selected
thermal chart law. It also proves the stated normalized
growing-loop mean comparison. No summable decay rate, physical
Wilson law, full interacting continuum reconstruction, or positive
physical Hamiltonian mass gap has been proved. The original full
objective remains active.
\end{firewall}
\section{V93: quantitative thermal balance and Gaussian score estimates without fourth moments}
\label{sec:v93}

\subsection{Context: quantify the actual move before inverting its long wavelengths}

V92 proves uniform qualitative covariance comparison, but supplies
no rate compatible with physical test normalization. The first
upstream task is to quantify V89's actual change-of-measure
identity. Its action expansion was previously only in probability:
subtracting two quadratic approximations to the action left
quartic remainders without an effective expectation bound.

Here the action is differentiated along the thermal move itself.
Unitarity bounds its derivatives, and the already proved local
second moments control the surviving first derivative. This gives
an explicit \(L^1\) action-increment error, a quantitative bounded
balance identity, and a finite-volume Stein estimate. For every
fixed coexact score, its characteristic function approaches the
corresponding Gaussian characteristic function at rate
\(O(\beta^{-1/4}+M^{-1})\) on bounded frequency intervals.
The chart support loss is included, not discarded.

These are estimates for coexact curvature scores, not yet for
general physically normalized curvature tests. Passing from
scores to the long-wavelength field still needs uniform inverse
estimates and quantitative moment control. The full interacting
Yang--Mills mass-gap goal is unchanged. The next companion
checkpoint is V95.

\subsection{Notation and the actual budgets used}

Use V90's selected actual law at \(\beta=\beta_M^\circ\), with
\(M\ge4\) and the fixed radius \(0<R\le1/4\). All assertions below
are for sufficiently large members of that sequence. Retain
\[
 m=144\pi^2\sqrt{20}\,\beta^{-1/2}+72\pi^2M^{-2},
 \qquad k=12+4\sqrt{60}\,\sqrt{\beta m}.
\]
The already proved bounds are
\begin{equation}\label{eq:v93-budgets}
 \mathbb E|a_e|^2\le m,\qquad
 \mathbb E|Y_{M,p}|^2\le80,\qquad
 V^{-1}\mathbb E\operatorname{Tr}F^{-1}\le k .
\end{equation}
Thus \(m\to0\) and
\(k/\sqrt\beta\le12/\sqrt\beta+4\sqrt{60m}\).
No assumption about a fourth moment or global convexity is added.

Let \(B\) be a deterministic real colour two-form on the torus.
It may depend on \(M\). Define
\[
 h=d_1^*B,\qquad \ell=d_1h,\qquad
 E=\bigcup_{p\in\operatorname{supp}B}\partial p,\qquad
 N=|E|,\qquad b=\max_p|B_p|,
\]
and the link norms
\[
 H_1=\sum_e|h_e|,\qquad
 H_2=\left(\sum_e|h_e|^2\right)^{1/2},\qquad
 H_\infty=\max_e|h_e|.
\]
Here vector norms include the three colours. In particular,
\[
 H_\infty\le6b,\qquad H_1\le6bN,\qquad H_2\le6b\sqrt N.
\]
The zero move is immediate and can be omitted.
The amplitude \(\theta>0\) below means the translation
\(a\mapsto a\pm\theta h/\sqrt\beta\). Whenever curvature at shifted
points is used, impose
\begin{equation}\label{eq:v93-guard}
          \frac{6\theta b}{\sqrt\beta}
                        \le\min\{R/2,1/12\}.
\end{equation}
This leaves the geometric guard available and keeps every shifted
link logarithm within radius \(1/3\), even outside the chart.
It defines an observable extension only; the probability density
outside the chart remains zero.

\subsection{Unitary first variations control the action error}

For the four links of one plaquette set
\[
 A_p=\left(\sum_{e\in\partial p}|a_e|^2\right)^{1/2},
 \qquad u_p=\left(\sum_{e\in\partial p}|h_e|^2\right)^{1/2}.
\]
Then \(\mathbb E A_p^2\le4m\). Counting the six plaquettes incident
to a positive edge gives
\begin{equation}\label{eq:v93-incidence}
 \begin{gathered}
       \sum_pu_p\le6H_1,\qquad \sum_pu_p^2=6H_2^2,\\
       \sum_pu_p^3\le12H_\infty H_2^2,\qquad
       \sum_pu_p^4\le24H_\infty^2H_2^2 .
 \end{gathered}
\end{equation}
Indeed \(u_p\le\sum_{e\in\partial p}|h_e|\) and
\(u_p\le2H_\infty\); the stated sums follow.

\begin{lemma}[Plaquette body velocity and third derivative]
\label{lem:v93-velocity}
Write the plaquette product as
\(P_p=\cos r\,I+\iota(v)\), where
\(r=|z_p|\le1\) and \(v=(\sin r/r)z_p\).
Define its right body velocity by
\[
                 P_p^{-1}DP_p(a)[h]=\iota(w_p).
\]
Then
\begin{equation}\label{eq:v93-body}
 |w_p|\le2u_p,\qquad
 |w_p-(dh)_p|\le8A_pu_p,\qquad
                         Ds_p(a)[h]=v\cdot w_p .
\end{equation}
In the Euclidean space of four link logarithms, the Wilson
plaquette action has the global mixed derivative bound
\begin{equation}\label{eq:v93-third}
                     \|D^3s_p\|_{\rm op}\le8 .
\end{equation}
\end{lemma}
\begin{proof}
The differentiated exponential identity is
\[
 D e^X[H]=\int_0^1 e^{(1-t)X}H e^{tX}\,dt .
\]
For skew-Hermitian \(X,H\), multiplication by \(e^{-X}\) on the
left makes the body velocity an average of conjugates of \(H\).
For a factor \(e^{\iota(\epsilon a_e)}\), with
\(\epsilon=\pm1\), its body velocity therefore has norm at most
\(|h_e|\). Its difference from \(\epsilon h_e\) is at most
\[
        \int_0^1 2t|a_e|\,|h_e|\,dt=|a_e|\,|h_e|,
\]
using the adjoint bound of V88.

The body velocity of a product is the sum of the factor body
velocities conjugated by the inverse of their following suffix.
The suffix conjugation differs from the identity by operator norm
at most twice the sum of its link-logarithm lengths, by telescoping
orthogonal adjoint maps. Hence
\[
 |w_p-(dh)_p|
 \le\sum_{i=1}^4
       \left(|a_i|+2\sum_{j>i}|a_j|\right)|h_i|
 \le2\left(\sum_i|a_i|\right)\left(\sum_i|h_i|\right)
 \le8A_pu_p .
\]
The norm bound follows by summing \(|h_i|\le2u_p\).
The scalar part of \(P_p\iota(w_p)\) is \(-v\cdot w_p\);
since \(s_p=1-\tfrac12\operatorname{Tr}P_p\), this proves the
last identity in \eqref{eq:v93-body}.

Repeated differentiation of the exponential yields the sum over
ordered insertions of the direction matrices, integrated on an
ordered simplex. All intervening exponential factors are unitary.
The number of orders cancels the simplex volume, so the mixed
derivative norm is bounded by the product of the direction norms.
Applying the product rule to the four factors bounds a mixed
third derivative of \(P_p\), in directions \(u,v,w\), by
\[
       \left(\sum_e|u_e|\right)
       \left(\sum_e|v_e|\right)
       \left(\sum_e|w_e|\right)
                  \le8\|u\|_2\|v\|_2\|w\|_2 .
\]
The normalized trace is bounded by the operator norm, which proves
\eqref{eq:v93-third}. These arguments work at arbitrary real link
logarithms for the action. The radius restriction is needed for
the logarithmic observable, not for this third derivative.
\end{proof}

The exponential derivative above is also the standard
Fr\'echet-derivative formula; see A. H. Al-Mohy and N. J. Higham,
\emph{Computing the Fr\'echet Derivative of the Matrix Exponential,
with an Application to Condition Number Estimation},
\href{https://eprints.maths.manchester.ac.uk/1218/}
{Manchester author's repository}.
The unitary bounds required here were proved explicitly rather
than inferred from a numerical matrix-function algorithm.

\begin{theorem}[An effective \(L^1\) action-increment estimate]
\label{thm:v93-action}
Put \(S_\ell=\langle Y_M,\ell\rangle\) and \(v_\ell=\|\ell\|_2^2\).
Define
\begin{equation}\label{eq:v93-action-constants}
 \begin{aligned}
 D_1&=(96\sqrt{80m}+160/\sqrt\beta)H_1,\\
 D_2&=48\sqrt m\,H_2^2,\qquad
 D_3=16H_\infty H_2^2/\sqrt\beta .
 \end{aligned}
\end{equation}
For either sign and every \(\theta>0\),
\begin{equation}\label{eq:v93-action-error}
 \mathbb E\left|
 \beta\{S(a\pm\theta h/\sqrt\beta)-S(a)\}
          \mp\theta S_\ell-\tfrac12\theta^2v_\ell\right|
                          \le\theta D_1+\theta^2D_2+\theta^3D_3 .
\end{equation}
The action on the left is its ordinary link-exponential formula,
whether or not the shifted point lies in the chart.
\end{theorem}
\begin{proof}
By \eqref{eq:v93-body},
\[
 \sqrt\beta\,Ds_p(a)[h]
                   =Y_{M,p}\cdot\{(\sin r/r)w_p\}.
\]
Use \(|1-\sin r/r|\le r^2/6\), with its continuous value at zero,
and \(r\le1\). Since \(Y_{M,p}=\sqrt\beta z_p\),
\[
      |Y_{M,p}|r^2
          =\beta^{-1/2}|Y_{M,p}|^2r
          \le\beta^{-1/2}|Y_{M,p}|^2 .
\]
Consequently \eqref{eq:v93-budgets}, Cauchy--Schwarz and
\(\mathbb E A_p^2\le4m\) give
\begin{equation}\label{eq:v93-gradient}
 \mathbb E\left|
 \sqrt\beta\,Ds_p(a)[h]-Y_{M,p}\cdot(dh)_p\right|
       \le(16\sqrt{80m}+80/(3\sqrt\beta))u_p .
\end{equation}
The sinc correction was estimated by a second moment of \(Y\),
not a fourth moment.

At zero link logarithms,
\(D^2s_p(0)[h,h]=|(dh)_p|^2\): the linear term of the plaquette
is \(\iota(dh)\), and unitarity determines the scalar quadratic
term. By \eqref{eq:v93-third},
\[
 |D^2s_p(a+t\theta h/\sqrt\beta)[h,h]
                         -|(dh)_p|^2|
        \le8(A_p+t\theta u_p/\sqrt\beta)u_p^2 .
\]
Taylor's integral remainder along the move, followed by
\(\mathbb E A_p\le2\sqrt m\), thus bounds the expected quadratic
remainder by
\[
        8\theta^2\sqrt m\,u_p^2
                       +\frac{4\theta^3}{3\sqrt\beta}u_p^3 .
\]
For the negative move the first derivative changes sign and the
same absolute bounds apply. Sum this estimate and
\(\theta\) times \eqref{eq:v93-gradient} over plaquettes, and use
\eqref{eq:v93-incidence}. The constants become exactly
\eqref{eq:v93-action-constants}.
\end{proof}

\subsection{An explicit derivative bound for the logarithmic observable}

The following deliberately coarse constant is independent of
volume, support, and coupling:
\[
                        L_*=28\cdot9^5=1{,}653{,}372 .
\]

\begin{lemma}[Plaquette-logarithm Hessian and global move error]
\label{lem:v93-log}
On the product of four link-logarithm balls of radius \(1/3\),
the principal plaquette logarithm has Euclidean Hessian norm at
most \(L_*\). Under \eqref{eq:v93-guard}, put
\[
 \kappa_1=2\sqrt6\,L_*\sqrt m\,H_2,\qquad
 \kappa_2=\sqrt6\,L_*H_\infty H_2/\sqrt\beta .
\]
Then, in the full finite-dimensional curvature coordinate space,
\begin{equation}\label{eq:v93-curvature-error}
 \left(\mathbb E
   \|Y_M(a\pm\theta h/\sqrt\beta)-Y_M(a)\mp\theta\ell\|_2^2
                  \right)^{1/2}
                                  \le\theta\kappa_1+\theta^2\kappa_2 .
\end{equation}
\end{lemma}
\begin{proof}
The product angle is at most \(4/3<\pi/2\). Its scalar part is
therefore positive. If \(v\) is its quaternion vector part, its
principal logarithm is
\[
 f(v)=\frac{\arcsin|v|}{|v|}v
            =\int_0^1\frac{v}{\sqrt{1-t^2|v|^2}}\,dt .
\]
Set \(c=\cos(4/3)\). Then \(|v|\le\sin(4/3)\).
The radial and tangential eigenvalues of \(Df\) are at most
\(c^{-1}\). Differentiating the integral twice bounds the mixed
second derivative by
\[
                  \|D^2f\|\le3c^{-3}+3c^{-5}\le6c^{-5}.
\]
Indeed its three first terms have the form
\(t^2(1-t^2|v|^2)^{-3/2}(v\cdot w)u\), and its remaining
term has norm at most
\(3t^4(1-t^2|v|^2)^{-5/2}|v|^3\|u\|\|w\|\);
using \(|v|\le1\) and integrating over an interval of length one
gives the displayed bound.

The unitary product derivative argument gives \(\|Dv\|\le2\)
and \(\|D^2v\|\le4\). The chain rule therefore gives
\[
           \|D^2(f\circ v)\|
                   \le24c^{-5}+4c^{-1}\le28c^{-5}.
\]
Finally \(c\ge1-(4/3)^2/2=1/9\), so this is at most \(L_*\).
All estimates extend continuously to \(v=0\).

The derivative of the logarithm at zero is the signed link sum.
Integrating its Hessian estimate along the moving direction
therefore bounds the plaquette error by
\[
 L_*\left(\theta A_pu_p
                        +\frac{\theta^2}{2\sqrt\beta}u_p^2\right).
\]
Take the \(L^2\) norm over both probability and the finite
plaquette coordinate set. Minkowski's inequality,
\(\mathbb E A_p^2\le4m\), and
\eqref{eq:v93-incidence} give
\[
 L_*\left(2\theta\sqrt m\,\sqrt6 H_2
           +\frac{\theta^2}{2\sqrt\beta}\sqrt{24}H_\infty H_2\right),
\]
which is \eqref{eq:v93-curvature-error}.
\end{proof}

\subsection{Quantitative balance with every support-loss term retained}

For the original normalized density extended by zero, define
\[
 A_\pm^\theta(a)=
   \min\left\{1,\frac{f_M(a\pm\theta h/\sqrt\beta)}{f_M(a)}\right\},
 \qquad
 a_\pm^\theta(Y)=
   \min\{1,e^{\mp\theta S_\ell-\theta^2v_\ell/2}\}.
\]
These are acceptance weights, not additional probability measures.
Put
\begin{equation}\label{eq:v93-alpha-constants}
 \begin{aligned}
 A_0&=4Nm/R^2,\\
 A_1&=24(1+3R)bNk/\sqrt\beta+12RbN/\sqrt\beta+D_1,\\
 A_2&=D_2,\qquad A_3=D_3,\qquad
 \alpha(\theta)=A_0+\theta A_1+\theta^2A_2+\theta^3A_3 .
 \end{aligned}
\end{equation}
The symbols \(A_i\) here are deterministic bounds, not the random
acceptance weights \(A_\pm^\theta\).

\begin{theorem}[Finite-volume accepted balance with an explicit error]
\label{thm:v93-balance}
Under \eqref{eq:v93-guard},
\begin{equation}\label{eq:v93-weight-error}
       \mathbb E|A_\pm^\theta-a_\pm^\theta(Y_M)|
                                      \le\min\{1,\alpha(\theta)\}.
\end{equation}
For every bounded Lipschitz function \(F\) of the full curvature
coordinates, real or complex, with \(\|F\|_\infty\le1\),
\begin{equation}\label{eq:v93-balance}
 \begin{split}
 \big|\mathbb E\{a_+^\theta(Y_M)F(Y_M)\}
       -\mathbb E\{a_-^\theta(Y_M)F(Y_M-\theta\ell)\}\big|\\
       \le2\alpha(\theta)+
             \operatorname{Lip}(F)(\theta\kappa_1+\theta^2\kappa_2).
 \end{split}
\end{equation}
No constant depends on the diameter of the moved support.
\end{theorem}
\begin{proof}
Use V89's simultaneous good event, with
\(\rho=\theta b/\sqrt\beta\) and
\(\eta=6(1+3R)\rho\). The estimates are
\[
 \mu_M(G^c)\le N(4m/R^2+2\eta k),\qquad
 \mathbb E\left[\mathbf1_G
       \left|\log\frac{w(a\pm\theta h/\sqrt\beta)}{w(a)}\right|\right]
                                   \le2\eta Nk+12R\rho N .
\]
These bounds hold for arbitrary finite torus supports by the
compressed edge-trace proof; a support may grow with \(M\).
On \(G\), the log density ratio differs from
\(\mp\theta S_\ell-\theta^2v_\ell/2\) by the action error of
Theorem~\ref{thm:v93-action} plus this geometric log increment.
The function \(x\mapsto\min\{1,e^x\}\) is globally
one-Lipschitz. Off \(G\), the absolute difference of acceptance
weights is at most one. Adding the estimates gives
\[
 N(4m/R^2+4\eta k)+12R\rho N+
                       \theta D_1+\theta^2D_2+\theta^3D_3
                  =\alpha(\theta).
\]
This proves \eqref{eq:v93-weight-error}.

The density's change of variables on the mean-zero Coulomb space
has Jacobian one and gives the exact identity
\[
 \mathbb E\{A_+^\theta(a)F(Y_M(a))\}
   =\mathbb E\{A_-^\theta(a)
                           F(Y_M(a-\theta h/\sqrt\beta))\}.
\]
It includes every chart boundary because \(f_M\) is zero outside.
Replacing its two acceptance weights costs at most
\(2\alpha(\theta)\). Replacing its shifted observable by
\(F(Y_M-\theta\ell)\) costs at most its Lipschitz constant times
the \(L^1\) norm of the curvature increment error, bounded by
Lemma~\ref{lem:v93-log}. This proves \eqref{eq:v93-balance}.
\end{proof}

\subsection{An effective finite-volume Stein estimate}

Set
\[
          J=\sum_p|\ell_p|,\qquad q_1=\sqrt{80}\,J,\qquad q_2=80J^2 .
\]
Minkowski and \eqref{eq:v93-budgets} imply
\(\mathbb E|S_\ell|\le q_1\) and
\(\mathbb E S_\ell^2\le q_2\).
For a smooth test \(F\), write \(D_\ell F\) for its directional
derivative in the deterministic curvature direction \(\ell\).

\begin{proposition}[Stein identity with a controlled boundary cost]
\label{prop:v93-Stein}
Let \(F\) be bounded \(C^2\), with \(\|F\|_\infty\le1\), finite
Lipschitz constant and bounded \(D_\ell F,D_\ell^2F\).
For every \(\theta>0\) satisfying \eqref{eq:v93-guard},
\begin{equation}\label{eq:v93-Stein}
 \begin{aligned}
 \big|\mathbb E D_\ell F(Y_M)
                 -\mathbb E\{S_\ell F(Y_M)\}\big|
 &\le \frac{2\alpha(\theta)}{\theta}
          +\operatorname{Lip}(F)(\kappa_1+\theta\kappa_2)\\
 &\quad+\theta(q_2/2+v_\ell)\\
 &\quad+\|D_\ell F\|_\infty
                           (\theta q_1+\theta^2v_\ell/2)
          +\frac{\theta}{2}\|D_\ell^2F\|_\infty .
 \end{aligned}
\end{equation}
Complex-valued tests obey the same bound.
\end{proposition}
\begin{proof}
Put \(\varphi(x)=\min\{1,e^{-x}\}\).
For real \(u\), the exponential remainder at nonnegative
\(|u|\) gives
\[
        |\varphi(u)-\varphi(-u)+u|\le u^2/2 .
\]
The function \(\varphi\) is one-Lipschitz, so adding the same
nonnegative \(c\) to both its arguments changes the difference
by at most \(2c\). With \(u=\theta S_\ell\) and
\(c=\theta^2v_\ell/2\), this yields
\[
 |a_+^\theta-a_-^\theta+\theta S_\ell|
                      \le\theta^2(S_\ell^2/2+v_\ell),
 \qquad
 |1-a_-^\theta|\le\theta|S_\ell|+\theta^2v_\ell/2 .
\]
Taylor expansion of the test, not the density, gives
\[
 F(Y)-F(Y-\theta\ell)
          =\theta D_\ell F(Y)+R_F,\qquad
 |R_F|\le\theta^2\|D_\ell^2F\|_\infty/2 .
\]
Decompose the left difference in \eqref{eq:v93-balance} as
\[
 \mathbb E\{(a_+^\theta-a_-^\theta)F(Y)\}
  +\mathbb E\{a_-^\theta[F(Y)-F(Y-\theta\ell)]\}.
\]
Subtract
\(\theta\{\mathbb E D_\ell F-\mathbb E S_\ell F\}\),
use the last three inequalities, and divide by \(\theta\).
Taking the stated first and second score moment bounds proves
\eqref{eq:v93-Stein}.
No finite-volume integration by parts across a chart boundary was
performed. The term \(2A_0/\theta\) records its retained cost.
\end{proof}

\subsection{Quantitative Gaussian characteristic functions of compact scores}

Define the actual characteristic function
\(\phi_M(s)=\mathbb E e^{isS_\ell}\). For any permitted
\(\theta>0\), set
\begin{equation}\label{eq:v93-r-coefficients}
 \begin{aligned}
 r_0&=2A_0/\theta+2A_1+2\theta A_2+2\theta^2A_3
                              +\theta(q_2/2+v_\ell),\\
 r_1&=\sqrt{v_\ell}(\kappa_1+\theta\kappa_2)
                              +v_\ell(\theta q_1+\theta^2v_\ell/2),\\
 r_2&=\theta v_\ell^2/2 .
 \end{aligned}
\end{equation}

\begin{theorem}[Finite-volume Gaussian score rate]
\label{thm:v93-characteristic}
For every real \(s\),
\begin{equation}\label{eq:v93-characteristic}
     |\phi_M(s)-e^{-v_\ell s^2/2}|
             \le r_0|s|+\frac{r_1s^2}{2}+\frac{r_2|s|^3}{3}.
\end{equation}
For a fixed finitely supported \(B\) placed anywhere on the torus,
there is a constant depending on \(B,R\), not its location or
\(M\), such that, for all sufficiently large \(M\),
\begin{equation}\label{eq:v93-rate}
 |\phi_M(s)-e^{-v_\ell s^2/2}|
    \le C_{B,R}(\beta^{-1/4}+M^{-1})(|s|+s^2+|s|^3).
\end{equation}
The conclusion is a rate for characteristic functions, not a
claimed total-variation distance or a rate for second moments.
\end{theorem}
\begin{proof}
In \eqref{eq:v93-Stein} use
\(F(Y)=e^{is\langle Y,\ell\rangle}\).
Its Lipschitz constant is \(|s|\sqrt{v_\ell}\), and its first
two directional derivatives have supremum norms
\(|s|v_\ell\) and \(s^2v_\ell^2\).
Since \(\mathbb E|S_\ell|<\infty\),
\(\phi_M'(s)=i\mathbb E(S_\ell e^{isS_\ell})\).
Thus \eqref{eq:v93-Stein} becomes
\[
             |\phi_M'(s)+v_\ell s\phi_M(s)|
                                      \le r_0+r_1|s|+r_2s^2 .
\]
There is no assumption of a finite exponential moment.
Solve this first-order differential inequality using the
integrating factor \(e^{v_\ell s^2/2}\) and \(\phi_M(0)=1\).
For \(s\ge0\), the resulting error is bounded by
\[
 \int_0^s e^{-v_\ell(s^2-t^2)/2}
                      (r_0+r_1t+r_2t^2)\,dt
       \le r_0s+r_1s^2/2+r_2s^3/3 .
\]
For negative \(s\), use conjugate symmetry. The zero-score case
is immediate.

For fixed \(B\), choose \(\theta=\sqrt m\), which satisfies
\eqref{eq:v93-guard} eventually. All support and link-norm
parameters are then independent of \(M\). The displayed formulas
give
\[
 A_0=O_{B,R}(m),\quad
 A_1=O_{B,R}(\sqrt m+\beta^{-1/2}),\quad
 A_2=O_B(\sqrt m),\quad A_3=O_B(\beta^{-1/2}),
\]
and the same respective orders for \(\kappa_1,\kappa_2\).
It follows that every \(r_i\) is
\(O_{B,R}(\sqrt m+\beta^{-1/2})\).
Finally
\[
       \sqrt m\le12\pi\,20^{1/4}\beta^{-1/4}
                                          +6\sqrt2\pi\,M^{-1}.
\]
This proves \eqref{eq:v93-rate}.
\end{proof}

The same reasoning gives a joint characteristic-function rate for
any fixed number of compact scores. Apply the theorem to
\(B=\sum_j t_jB_j\); on each fixed compact set of coefficients
\((t_j)\), all constants can be bounded uniformly. The Gaussian
quadratic form is
\(\|\sum_jt_jd_1d_1^*B_j\|_2^2\).
If the patches are translated far enough apart that their score
supports are disjoint, its cross terms vanish exactly.
The rate is then uniform in their locations and separates the
Gaussian score factors. This quantitatively strengthens the
score part of V92; it does not quantify the subsequent
approximation of general curvature coordinates by such scores.

\subsection{Growing supports: what the explicit estimate does and does not pay}

For a sequence \(B_M\) with uniformly bounded \(b\), take
\(\theta=1\). Then \eqref{eq:v93-guard} eventually holds, and the
link-norm estimates above imply
\[
 \alpha(1)=O_{R,\sup b}\!\left(
                       N(\beta^{-1/4}+M^{-1})\right),\qquad
 \kappa_1+\kappa_2=O_{\sup b}\!\left(
                       \sqrt N(\beta^{-1/4}+M^{-1})\right).
\]
Consequently \eqref{eq:v93-balance} has a vanishing error, uniformly
for bounded tests with bounded Euclidean Lipschitz constant, if
\[
                           N=o(\min\{\beta^{1/4},M\}).
\]
Unlike the previous fixed-support passage, this assertion allows
the test and the number of moved edges to grow. It retains all
nonlinear action and geometric-density errors. A four-dimensional
block with \(N\) of order \(L^4\), for example, is covered only for
\(L=o(\min\{\beta^{1/16},M^{1/4}\})\) by this bounded-amplitude
sufficient condition. The sharper formulas remain available when
the move amplitude or link norms decrease with support size.

There are two explicit obstructions to treating these bounds as a
physical covariance estimate. First, the radius-guard cost
\(A_0=4Nm/R^2\) is a union bound independent of move amplitude.
For a move with extensive edge support it is not small under the
current budgets. Sending \(\theta\downarrow0\) alone makes
\(2A_0/\theta\) worse in the Stein estimate. A sharper near-boundary
estimate, not a declaration that the boundary is absent, is needed
for such a route.

Second, reconstructing a general curvature test from the scores
requires inverting \(K=d_1d_1^*\) on its exact-form range.
At nonzero Fourier momentum its nonzero eigenvalue is
\(\lambda(k)\). Thus the deterministic inverse costs
\(\lambda(k)^{-1}\), which is of order \(M^2\) near the smallest
nonzero torus momenta. V92 avoided this loss by taking a qualitative
heat approximation after the local limit; that order of limits
cannot be reused as a quantitative physical-scale estimate.
Moreover characteristic-function bounds cannot be differentiated
to obtain a second-moment error without quantitative tail control.
The present rate for compact scores therefore does not imply a
rate for V92's full covariance discrepancy \(\delta_M\).

\subsection{Exact checks and preservation}

The diagnostic script is
\begin{center}\small
\path{scripts/verify_ym_f14_v93_quantitative_balance.py}.
\end{center}
Twenty-four rational quaternion configurations check the exact
suffix-body-velocity formula and the action derivative
\(Ds=v\cdot w\). Eighteen compact coexact configurations check
the edge-incidence and fourth-power bounds used only for the
deterministic move, not for random curvature moments.
The script also checks 164 positive-part identities or inequalities,
48 exact bounded-acceptance balances with support loss,
48 unitary-product Leibniz coefficient identities, the integer
logarithm-Hessian constant, four admissible support-growth
exponents and eight characteristic-integral coefficient formulas.
These are finite algebra checks. The expectation bounds, Gaussian
score rate and support-control statements are the analytic results
proved above, not inferences from sampling.

The certificate contains 6,549 bytes and has SHA--256
\begin{center}\scriptsize\ttfamily
88E8DC7D456E1A824188925B3762FAF75D93ACA786559901588726F32CC23FC7.
\end{center}
The predecessor V92 contains 2,263,677 bytes and has SHA--256
\begin{center}\scriptsize\ttfamily
9F147190CE61F10A1FD4700C0FAA1E712D7FEAD7DAA457683B4D0D8722DC26C9.
\end{center}
Its contents are preserved except for the title version and this
terminal insertion. V93 replaces V92 as the sole active main note.
Earlier certificates, supplied sources, companion files,
monographs and archives remain unchanged.

\begin{firewall}
The quantitative bottleneck removed here is the actual action
increment and bounded thermal balance at finite volume. It yields
a proved compact-score characteristic-function rate using only
second moments. Physical long-wavelength covariance control,
full compact chart coverage, the physical coupling trajectory,
interacting continuum reconstruction and a positive physical
Hamiltonian mass gap remain unproved. The original full goal
remains active.
\end{firewall}
\section{V94: quantitative energy and curvature covariance}
\label{sec:v94}

\subsection{Context: use the sharp energy budget instead of assuming moment tails}

V93 gives a quantitative characteristic-function estimate for
compact coexact scores. It does not justify differentiating that
estimate to obtain moment convergence. This section supplies a
different route: use a bounded cosine to obtain a lower second
moment, then combine it with V90's sharp upper energy budget.
The key variable is the plaquette's quaternion vector part,
which obeys an exact energy identity.

The resulting energy deficit controls a finite-range
approximation to the full curvature, including its low-frequency
residual and nonlinear logarithm error. Applying the same argument
to sums and differences of two scores then gives an explicit
uniform covariance rate over all torus separations. This advances
the qualitative covariance conclusion of V92; it is not another
fixed-score Gaussian calculation.

The rates below are still insufficient after fixed-physical-scale
normalization, and apply only to the selected thermal chart law.
The original interacting continuum mass-gap objective is
unchanged. The next companion audit is V95.

\subsection{The vector part of a plaquette has an exact energy slack}

Continue with V90's selected law and V93's notation
\[
 \beta=\beta_M^\circ,\qquad
 m=144\pi^2\sqrt{20}\,\beta^{-1/2}+72\pi^2M^{-2},\qquad
 e_M=\sqrt m+\beta^{-1/2}.
\]
Write
\[
 \eta=\eta_R(\log\log M),\qquad
 q_M=\mathbb E\sum_{\mu<\nu}X_{M,\mu\nu}(0)\le9/2+\eta
\]
for the nonnegative excess bound proved in V90. Eventually
\(q_M\le20\), \(m\le1\), \(\beta\ge1\), and \(e_M\le1\).
All statements below are for this eventual selected sequence.

Write the plaquette as
\(P_p=u_{0,p}I+\iota(v_p)\), and define the three-component field
\[
 \mathcal W_{M,p}=\sqrt\beta\,v_p
          =\frac{\sin|z_p|}{|z_p|}Y_{M,p}.
\]
This vector field is not the scalar Wilson trace \(W_C\).
For a stationary field \(U\) with eighteen components per site,
put
\[
        \mathcal E_M(U)=M^{-4}\mathbb E\|U\|_2^2
                      =\mathbb E\sum_{a=1}^{18}|U_a(0)|^2 .
\]
Here \(a\) combines orientation and colour.

\begin{lemma}[Exact chord energy and logarithm comparison]
\label{lem:v94-chord}
Pointwise at every plaquette,
\begin{equation}\label{eq:v94-energy-identity}
             2X_{M,p}-|\mathcal W_{M,p}|^2
                                  =X_{M,p}^2/\beta\ge0 .
\end{equation}
Consequently \(\mathcal E_M(\mathcal W_M)\le9+2\eta\).
Also
\begin{equation}\label{eq:v94-log-comparison}
 \mathbb E|Y_{M,p}-\mathcal W_{M,p}|
                          \le\frac{40}{3\sqrt\beta},\qquad
 |Y_{M,p}-\mathcal W_{M,p}|^2
                          \le\frac4{9\beta}X_{M,p}^2 .
\end{equation}
\end{lemma}
\begin{proof}
The unit-quaternion identity gives
\[
 |\mathcal W_{M,p}|^2=\beta(1-u_{0,p}^2),\qquad
 X_{M,p}=\beta(1-u_{0,p}),
\]
which proves \eqref{eq:v94-energy-identity}.
Sum and use the actual bound on \(q_M\).
For \(r=|z_p|\le1\),
\[
 |Y-\mathcal W|\le r^2|Y|/6
              \le |Y|^2/(6\sqrt\beta),
\]
so V88's \(\mathbb E|Y_p|^2\le80\) proves the first inequality.
For the second use \(r^2\le4(1-\cos r)\), also proved in V88:
\[
 |Y-\mathcal W|^2\le\beta r^6/36
                 \le\beta r^4/36
                 \le4X^2/(9\beta).
\]
The last bound will be used only after its right side is controlled
by the energy identity; it assumes no fourth-moment estimate.
\end{proof}

\subsection{A finite-range coexact contraction}

Let \(\Delta\) be the scalar lattice Laplacian and
\(K=d_1d_1^*\), acting identically on the three colours.
Define the lazy walk operator and the polynomial
\[
 R_M=I-\Delta/16
      =\tfrac12I+\tfrac1{16}\sum_{i=1}^4(T_{e_i}+T_{-e_i}),
 \qquad q_n(s)=\frac1{16}\sum_{j=0}^{n-1}(1-s/16)^j ,
\]
for an integer \(n\ge1\). Set
\begin{equation}\label{eq:v94-filter}
        T_{n,M}=Kq_n(\Delta).
\end{equation}
If \(\rho(k)=1-\lambda(k)/16\), then
\[
   T_{n,M}(k)=(1-\rho(k)^n)\mathsf P(k),\qquad
                         0\preceq T_{n,M}\preceq I .
\]
At zero momentum the multiplier is zero. The operator is a
positive contraction, generally not an orthogonal projection.
Its columns are finite coexact scores.

Let
\[
 I_j=\int_{\mathbb T^4}\rho(k)^j\,\frac{dk}{(2\pi)^4}.
\]
Whenever \(M>4n+4\), the finite torus return probabilities up to
time \(2n\) equal these infinite-lattice probabilities: a walk of
that length cannot wind around the torus.

\begin{lemma}[Filter variance and its deficit from the limiting Gaussian energy]
\label{lem:v94-filter}
For \(M>4n+4\), every column of \(T_{n,M}\) has squared norm
\begin{equation}\label{eq:v94-filter-variance}
             v_n=\frac12(1-2I_n+I_{2n})\le\frac12 .
\end{equation}
The deficit relative to the limiting Gaussian energy \(9\) is
\[
       9-18v_n=9(2I_n-I_{2n}),\qquad
                         0\le I_{2n}\le I_n\le\pi^2/n^2 .
\]
\end{lemma}
\begin{proof}
The squared filter multiplier has trace
\(9(1-\rho^n)^2\) on all colours, including its zero value at
zero momentum. The six orientation diagonal means are equal:
the diagonal of \(P(k)\) at \(\mu\nu\) is
\((|q_\mu(k)|^2+|q_\nu(k)|^2)/\lambda(k)\), and the grid is
invariant under permutations of its four coordinates.
The colours contribute identically. Dividing the trace mean by
eighteen gives \eqref{eq:v94-filter-variance}. The no-winding
observation replaces its finite return probabilities by \(I_j\).

Since \(0\le\rho\le1\), \(I_{2n}\le I_n\).
On \([-\pi,\pi]^4\),
\(\lambda(k)\ge4|k|^2/\pi^2\), so
\[
 \rho(k)^n\le e^{-n\lambda(k)/16}
                  \le e^{-n|k|^2/(4\pi^2)} .
\]
Enlarging the integration region to \(\mathbb R^4\) gives
\[
 I_n\le(2\pi)^{-4}
          \int_{\mathbb R^4}e^{-n|k|^2/(4\pi^2)}\,dk
                                   =\pi^2/n^2 .
\]
\end{proof}

\subsection{Uniform score estimates for these growing filters}

Let \(\delta_{x,a}\) be a scalar unit coordinate in the real
curvature space. The sources
\[
 B_{n,x,a}=q_n(\Delta)\delta_{x,a}
\]
give \(\ell=KB_{n,x,a}=T_{n,M}\delta_{x,a}\).
We need both these sources and
\(B_{n,x,a}\pm B_{n,y,b}\), uniformly in \(x,y,a,b\).
All estimates below include coincident or overlapping supports.

For these one- or two-column sources, the parameters in V93 obey
\begin{equation}\label{eq:v94-source-bounds}
 \begin{gathered}
 b\le2,\quad N\le648n^4,\quad H_1\le n/2,\quad
 H_\infty\le12,\quad H_2^2\le6n,\\
 J=\sum_p|\ell_p|\le3n,\qquad v_\ell=\|\ell\|_2^2\le4 .
 \end{gathered}
\end{equation}
Here \(M>4n+4\) is understood.
Indeed the positive kernel \(q_n(\Delta)\delta_0\) has mass
\(n/16\) and support in the radius-\((n-1)\) lattice ball.
Its maximum is at most \(G_4(0)<1\): on \(\mathbb Z^4\),
\((1/16)\sum_{j\ge0}R^j=\Delta^{-1}\) at the kernel level,
and V91's positive Green line sum bounds \(G_4(0)\) by \(g_3<1\).
For a direct justification of the maximum, the multiplier of the
partial sum is nonnegative, so the absolute value of each Fourier
coefficient is at most its origin coefficient.
No periodic images overlap under the stated size condition.

Each source plaquette contributes at most four edges to \(E\),
giving \(N\le8(2n+1)^4\le648n^4\) for the two-source union.
The \(\ell^1\) norm of \(d_1^*\) is at most four, so
\(H_1\le n/2\); \(H_\infty\le6b\) and
\(H_2^2\le H_\infty H_1\) give the next bounds.
V89's \(K\)-column has absolute coefficient sum \(24\);
therefore \(J\le24(n/8)=3n\).
Finally \(T_{n,M}\) is a contraction, so its action on a sum or
difference of two unit coordinates has squared norm at most four.

For explicit, deliberately nonoptimized constants, put
\[
                 \mathfrak C_R=10^9(1+R^{-2}),\qquad
                 \epsilon_{n,M}=\mathfrak C_R n^4 e_M .
\]

\begin{lemma}[Characteristic functions and lower second moments]
\label{lem:v94-score}
For all sufficiently large \(M\), uniformly over the source family
above, the chord score \(Z=\langle\mathcal W_M,\ell\rangle\)
satisfies
\begin{equation}\label{eq:v94-score-characteristic}
 |\mathbb E e^{isZ}-e^{-v_\ell s^2/2}|
             \le\epsilon_{n,M}(|s|+s^2+|s|^3).
\end{equation}
If \(\epsilon_{n,M}\le1\), then
\begin{equation}\label{eq:v94-score-lower}
 \mathbb E Z^2\ge v_\ell-d_{n,M},\qquad
                         d_{n,M}=10\epsilon_{n,M}^{2/3}.
\end{equation}
The lower bound includes sums and differences of any two filtered
coordinates.
\end{lemma}
\begin{proof}
In V93 choose the move amplitude \(\theta=\sqrt m\).
The guard condition holds uniformly because \(b\le2\), independently
of \(n\). Use \eqref{eq:v94-source-bounds}, \(m\le1\),
\(e_M\le1\), and
\(k/\sqrt\beta\le44e_M\) in V93's explicit \(r_i\).
Writing \(L_*=1{,}653{,}372\), direct bounds give
\[
 \begin{aligned}
 r_0&\le(5184R^{-2}+5{,}486{,}348)n^4e_M,\\
 r_1&\le(168L_*+116)n^4e_M,\qquad
 r_2\le8n^4e_M .
 \end{aligned}
\]
For example \(A_0/\theta\le2592R^{-2}n^4e_M\),
\(A_1\le2{,}741{,}552n^4e_M\),
\(A_2\le288ne_M\), and \(A_3\le1152ne_M\).
Also \(\kappa_1\le12L_*\sqrt n\,e_M\) and
\(\kappa_2\le72L_*\sqrt n\,e_M\).
These bounds and \(q_1\le3\sqrt{80}\,n\),
\(q_2\le720n^2\) give the displayed \(r_i\).

Replacing the logarithmic score by the chord score costs at most
\[
 |s|\,\mathbb E|\langle Y_M-\mathcal W_M,\ell\rangle|
          \le\frac{40J}{3\sqrt\beta}|s|
          \le40n^4e_M|s|,
\]
by \eqref{eq:v94-log-comparison}.
The coefficient \(\mathfrak C_R\) dominates the resulting
coefficient of \(|s|\), the coefficient \(r_1/2\) of \(s^2\),
and \(r_2/3\) of \(|s|^3\). V93's characteristic estimate therefore
proves \eqref{eq:v94-score-characteristic}.

For a real variable with a finite second moment,
\(1-\cos(sZ)\le s^2Z^2/2\). Thus for \(s>0\),
\[
 \mathbb E Z^2
   \ge \frac{2(1-\mathbb E\cos(sZ))}{s^2}
   \ge v_\ell-\frac{v_\ell^2s^2}{4}
                       -2\epsilon_{n,M}(s^{-1}+1+s).
\]
Here \(1-e^{-u}\ge u-u^2/2\) was applied to
\(u=v_\ell s^2/2\). Since \(v_\ell\le4\), choose
\(s=\epsilon_{n,M}^{1/3}\le1\).
The total loss is at most
\[
 4\epsilon_{n,M}^{2/3}
   +2(\epsilon_{n,M}^{2/3}+\epsilon_{n,M}
                                     +\epsilon_{n,M}^{4/3})
                 \le10\epsilon_{n,M}^{2/3}.
\]
The actual finite second moment follows from
\(|\mathcal W_p|^2\le2X_p\) and the finite source.
This argument uses a bounded cosine; it does not differentiate a
characteristic-function error or assume fourth-moment tails.
\end{proof}

\subsection{Energy saturation controls the residual and the nonlinear slack}

For \(M>4n+4\) and \(\epsilon_{n,M}\le1\), define
\begin{equation}\label{eq:v94-deficit}
 \mathcal B_{n,M}
        =2\eta+9(2I_n-I_{2n})+18d_{n,M}.
\end{equation}
This is an explicit nonnegative bound, not a new measure or a
postulated spectral gap.

\begin{theorem}[Actual residual-energy and scaled fourth-moment bounds]
\label{thm:v94-energy}
Writing \(T=T_{n,M}\), one has
\begin{equation}\label{eq:v94-combined-energy}
 \mathcal E_M((I-T)\mathcal W_M)
       +\frac1\beta\mathbb E\sum_{\mu<\nu}X_{M,\mu\nu}(0)^2
                                      \le\mathcal B_{n,M}.
\end{equation}
Consequently,
\begin{equation}\label{eq:v94-Y-energy}
 \begin{gathered}
 \mathcal E_M(Y_M-\mathcal W_M)\le\tfrac49\mathcal B_{n,M},\qquad
 \mathcal E_M((I-T)Y_M)\le\tfrac{13}{9}\mathcal B_{n,M},\\
 \frac1\beta\mathbb E\sum_{\mu<\nu}|Y_{M,\mu\nu}(0)|^4
                                           \le16\mathcal B_{n,M}.
 \end{gathered}
\end{equation}
The last estimate is scaled by \(1/\beta\); it is not a uniform
unscaled fourth-moment bound.
\end{theorem}
\begin{proof}
The eighteen filtered coordinates at the origin are the single
scores in Lemma~\ref{lem:v94-score}. Therefore
\[
          \mathcal E_M(T\mathcal W_M)\ge18(v_n-d_{n,M}).
\]
Since \(0\preceq T\preceq I\),
\((I-T)^2\preceq I-T^2\), and hence
\[
 \mathcal E_M((I-T)\mathcal W_M)
                    \le\mathcal E_M(\mathcal W_M)
                                      -\mathcal E_M(T\mathcal W_M).
\]
Add the expectation of the exact slack identity
\eqref{eq:v94-energy-identity}. Its right side becomes
\[
          2q_M-\mathcal E_M(T\mathcal W_M)
                   \le9+2\eta-18(v_n-d_{n,M})
                   =\mathcal B_{n,M},
\]
proving \eqref{eq:v94-combined-energy}.

Let the two nonnegative terms on the left of
\eqref{eq:v94-combined-energy} be \(D\) and \(S\).
Equation~\eqref{eq:v94-log-comparison} gives
\(\mathcal E_M(Y_M-\mathcal W_M)\le4S/9\).
The contraction \(I-T\), Minkowski and Cauchy--Schwarz yield
\[
 \mathcal E_M((I-T)Y_M)^{1/2}
       \le\sqrt D+\tfrac23\sqrt S,\qquad
 (\sqrt D+\tfrac23\sqrt S)^2
                          \le\tfrac{13}{9}(D+S).
\]
This proves the first two statements in \eqref{eq:v94-Y-energy}.
Finally \(|Y_p|^2\le4X_p\) pointwise, so
\(|Y_p|^4/\beta\le16X_p^2/\beta\). Sum and use \(S\le\mathcal B_{n,M}\).
\end{proof}

There is also a quantitative lower bound for the actual thermal
energy itself:
\[
       9/2+\eta-\mathcal B_{n,M}/2
                              \le q_M\le9/2+\eta .
\]
Indeed \(2q_M\ge\mathcal E_M(\mathcal W_M)
\ge\mathcal E_M(T\mathcal W_M)\).
The appearance of \(\eta\) in the lower bound cancels its
appearance inside \(\mathcal B_{n,M}/2\).

\subsection{The uniform covariance rate}

Let \(\delta_M\) be exactly V92's discrepancy
\[
 \delta_M=\max_{a,b}\sup_z
 \left|\mathbb E\{Y_{M,a}(0)Y_{M,b}(z)\}
                         -[\mathsf P_M]_{(0,a),(z,b)}\right|.
\]
The means are zero by common colour-rotation symmetry.

\begin{theorem}[Quantitative covariance comparison at every separation]
\label{thm:v94-covariance}
Under the hypotheses of Theorem~\ref{thm:v94-energy}, if
\(\mathcal B_{n,M}\le1\), then
\begin{equation}\label{eq:v94-covariance}
                         \delta_M\le25\sqrt{\mathcal B_{n,M}} .
\end{equation}
There is no restriction on the separation in this supremum.
\end{theorem}
\begin{proof}
Put \(U=T\mathcal W_M\), \(d=d_{n,M}\) and
\(\mathcal B=\mathcal B_{n,M}\).
For any one coordinate, the lower moment bound gives
\(\mathbb E U_a(x)^2\ge v_n-d\).
The total eighteen-coordinate energy of \(U\) is at most
\(9+2\eta\). Subtract the lower bounds for the other seventeen
coordinates at the same site, obtaining
\[
 \mathbb E U_a(x)^2
       \le9+2\eta-17(v_n-d)
       \le v_n+\mathcal B .
\]
Stationarity makes these statements uniform in \(x\).

Apply \eqref{eq:v94-score-lower} to \(U_a(x)+U_b(y)\) and
\(U_a(x)-U_b(y)\). Their Gaussian comparison variances are
\(2v_n\pm2[T^2]_{(x,a),(y,b)}\).
Combining those two lower bounds with
\(\mathbb E U_a(x)^2+\mathbb E U_b(y)^2\le2v_n+2\mathcal B\)
gives
\begin{equation}\label{eq:v94-filtered-covariance}
 \left|\mathbb E\{U_a(x)U_b(y)\}
                    -[T^2]_{(x,a),(y,b)}\right|
                                      \le\mathcal B+d/2 .
\end{equation}
This is a second-moment comparison obtained by energy saturation
and polarization, not by differentiating characteristic functions.

Next compare \(Y_M\) with \(\mathcal W_M\).
Each coordinate of their difference has second moment at most
\(4\mathcal B/9\) by \eqref{eq:v94-Y-energy} and stationarity.
The coarse bounds
\(\mathbb E Y_{M,a}(x)^2\le80\) and
\(\mathbb E\mathcal W_{M,a}(x)^2\le40\) give
\[
 \left|\mathbb E(Y_a(x)Y_b(y))
           -\mathbb E(\mathcal W_a(x)\mathcal W_b(y))\right|
          \le\tfrac23(\sqrt{80}+\sqrt{40})\sqrt{\mathcal B}.
\]
Similarly, \eqref{eq:v94-combined-energy} and the bounds
\(\mathcal E_M(\mathcal W_M),\mathcal E_M(U)\le40\) imply
\[
 \left|\mathbb E(\mathcal W_a(x)\mathcal W_b(y))
                           -\mathbb E(U_a(x)U_b(y))\right|
                                  \le2\sqrt{40}\sqrt{\mathcal B}.
\]

Finally \(\mathsf P_M-T^2\) is positive semidefinite. Its common
diagonal is
\((2I_n-I_{2n}-M^{-4})/2\), because the finite comparator
omits the zero mode. This is nonnegative and at most
\(\sigma_n=(2I_n-I_{2n})/2\).
The positive-semidefinite Cauchy--Schwarz inequality bounds every
one of its entries in absolute value by \(\sigma_n\).
Since \(18d\le\mathcal B\) and
\(18\sigma_n\le\mathcal B\), adding all comparisons gives
\[
 \delta_M\le
   \left(\tfrac23\sqrt{80}+\tfrac83\sqrt{40}\right)\sqrt{\mathcal B}
                                           +\tfrac{13}{12}\mathcal B .
\]
For \(\mathcal B\le1\) the right side is at most
\(25\sqrt{\mathcal B}\); for example
\(\sqrt{80}<9\), \(\sqrt{40}<13/2\) suffice to verify this constant.
\end{proof}

\subsection{A concrete optimized rate and a low-frequency bound}

For sufficiently small \(e_M\), take
\[
                    n_M=\lfloor e_M^{-1/7}\rfloor .
\]
Eventually \(n_M\ge e_M^{-1/7}/2\) and \(M>4n_M+4\).
For the latter, the lower bound \(e_M\ge6\sqrt2\pi/M\)
already implies \(n_M=O(M^{1/7})\).
Also
\(\epsilon_{n_M,M}\le\mathfrak C_R e_M^{3/7}\to0\).
Consequently
\begin{equation}\label{eq:v94-optimized-energy}
 \mathcal B_{n_M,M}\le2\eta+C_R^\dagger e_M^{2/7},
       \qquad C_R^\dagger=72\pi^2+180\mathfrak C_R^{2/3}.
\end{equation}
Indeed the missing Gaussian energy is at most \(18\pi^2/n_M^2\),
and the score loss is
\[
               18d_{n_M,M}=180(\mathfrak C_R n_M^4e_M)^{2/3}.
\]
These are bounded respectively by
\(72\pi^2e_M^{2/7}\) and
\(180\mathfrak C_R^{2/3}e_M^{2/7}\).

Theorems~\ref{thm:v94-energy} and \ref{thm:v94-covariance} now give
the explicit eventual estimate
\begin{equation}\label{eq:v94-rate}
 \begin{split}
 \delta_M&\le25\sqrt{\,2\eta+C_R^\dagger e_M^{2/7}\,}\\
         &=O_R\!\left(
            \sqrt{\eta_R(\log\log M)}
                         +\beta_M^{-1/28}+M^{-1/7}\right).
 \end{split}
\end{equation}
Here \(\beta_M=\beta_M^\circ\).
In particular the energy-selection contribution is
\(O_R((\log\log M)^{-1/4})\).
The constants are deliberately large, but specified and independent
of volume and separation. V92's qualitative \(\delta_M\to0\)
is now quantitatively controlled for the same actual law.

For a direct long-wavelength interpretation, let
\(\Pi_{\le8/n}\) be the componentwise Fourier projection onto
\(\lambda(k)\le8/n\), including zero momentum.
On the exact-form range at such a momentum,
\[
          1-T_{n,M}=\rho(k)^n
                  \ge1-n\lambda(k)/16\ge1/2
\]
by Bernoulli's inequality. On the orthogonal range it equals one.
Therefore
\[
             \Pi_{\le8/n}\preceq4(I-T_{n,M})^2 .
\]
Equation~\eqref{eq:v94-Y-energy} proves
\begin{equation}\label{eq:v94-infrared}
           \mathcal E_M(\Pi_{\le8/n}Y_M)
                              \le\tfrac{52}{9}\mathcal B_{n,M}.
\end{equation}
This is an actual low-frequency energy estimate with its
per-lattice-volume normalization explicit, not a gap in the
physical Hamiltonian spectrum.

\subsection{The remaining physical-scale loss}

The new result does not leave the covariance rate wholly open:
\eqref{eq:v94-rate} is now proved. What remains is a rate or
operator estimate compatible with physical normalization.
For \(a=M^{-1}\) and \(B_M(x)=a^2f(ax)\) in one curvature
component, V92's entrywise route still costs
\[
       \delta_M\|B_M\|_1^2=O(\delta_M M^4).
\]
The rate \eqref{eq:v94-rate} does not show that this expression
tends to zero. The growth of this upper bound is not a claim
that the actual physical variance diverges.

Similarly, since \(T_{n,M}\) annihilates the constant mode,
\eqref{eq:v94-Y-energy} implies
\[
 \mathbb E\left|
       M^{-2}\sum_xY_{M,a}(x)\right|^2
                   \le\tfrac{13}{9}M^4\mathcal B_{n,M}.
\]
This is still too large to control the constant physical test.
The estimate \eqref{eq:v94-infrared} controls energy per lattice
volume; a single physical mode carries a different normalization.
One must not erase that factor of \(M^4\).

The sharp upper energy excess \(\eta\), the finite-range score
cost, and the need for a physically weighted covariance estimate
are thus explicit next targets. Full compact chart coverage,
the physical coupling trajectory, observables and reconstruction
remain additional obligations; the theorem does not replace them
by its selected weak-coupling Gaussian comparator.

\subsection{Exact verification and preservation}

The diagnostic script is
\begin{center}\small
\path{scripts/verify_ym_f14_v94_energy_covariance.py}.
\end{center}
Sixty-four rational quaternion checks verify the exact chord
energy identity. Fifty-four exact lazy-filter columns cover all
eighteen orientation-colour coordinates and degrees one through
three; they verify the source masses, incidence bounds and
Gaussian variances
\[
                       9/64,\qquad4501/16384,\qquad23301/65536.
\]
The return probabilities at those degrees are
\(1/2,9/32,11/64\).
There are also 528 scalar contraction checks, 66 low-frequency
checks, three explicit constant comparisons, 64 combined-energy
Cauchy identities, 32 lower-moment coefficient checks, and exact
checks of the optimized exponents and covariance constant.
Five finite Fourier checks retain the zero-mode subtraction in
the diagonal of \(\mathsf P_M-T^2\); the side-four, degree-one
value is \(183/512\).
These finite calculations do not establish the actual-law bounds
by sampling; the preceding proofs do that analytically.

The certificate contains 5,632 bytes and has SHA--256
\begin{center}\scriptsize\ttfamily
9387226214C3E2550D4CDAC25CA6B1D1C174AAA0F1A7944B9E12DD0E99F5B2EE.
\end{center}
The predecessor V93 contains 2,288,516 bytes and has SHA--256
\begin{center}\scriptsize\ttfamily
6D2777E85725E1C4C02DD114CFB8B80E6FB0D67BC324DB8995551A1C620FA7E8.
\end{center}
Its body is preserved except for the title version and this
terminal insertion. V94 replaces V93 as the sole active main note.
Earlier certificates, supplied notes, companion files, monographs
and archives remain unchanged.

\begin{firewall}
V94 proves a quantitative residual-energy bound, a scaled
fourth-moment consequence, and a uniform curvature-covariance rate
for the actual selected thermal chart law. None is a physical
mass-gap assertion. Physical-scale covariance control, full
interacting continuum construction and a positive physical
Hamiltonian mass gap remain unproved. The original goal is active.
\end{firewall}
\section{V95: full compact Wilson pressure and thermal noncollapse}
\label{sec:v95-full-compact}

\subsection{The scheduled companion audit and the upstream change}

This is the V95 five-version checkpoint, performed on 15 September
2026. The current SM continuum V72, NS stability V26, shell-mixing
V16, and parallel YM V5 files have unchanged hashes. The supplied
SM source-selection V17, regenerative stationarity V21, and exact
action Einstein bridge V16 are also unchanged, as are the supplied
suggestions, the closure monograph, and the preceding archive.
Their relevant end sections have been revisited, not treated as
additional instructions. The next scheduled companion audit is V100.

There are useful constraints on the next direction. SM V72's
many-copy comparison remains an auxiliary, fixed-mode and
fixed-time statement. Stationarity V21 retains a real-phase
caustic obstruction and proves a quadratic phase-space comparison,
not the nonlinear physical YM identification. Source-selection V17
does not identify its comparison port with the required primitive
physical instrument. Einstein bridge V16 distinguishes the full
flux/domain problem from a diagonal soft inverse. NS V26's
Oseen estimates do not supply a YM mixing constant. No such
constant or physical identification is imported here.

Shell-mixing V16 identifies a flat, one-holonomy one-loop
coefficient but does not supply a joint nonperturbative remainder.
The suggested noncommuting quartic calculation is not repeated:
the present lineage's V20--V30 already treats that direction.
The parallel YM V5 warning against a volume-uniform tail for an
extensive topological charge also remains relevant. The next
observable should be local rather than a total charge.

The internal audit is equally important. V48--V53 establish
compact strong-mixing and transfer interfaces in their stated
small-interaction regions; they do not prove that a weak-coupling
ultraviolet trajectory enters those regions. V80's deterministic
tree gauges compare orbit representatives. V90--V94 control a
selected, determinant-weighted mean-zero Coulomb-chart law.
The V94 covariance rate does not overcome physical-scale
normalization or prove chart coverage.

We therefore go upstream to the \emph{full compact periodic
Wilson measure} for the same \(\SU(2)\) subproblem. The results below
hold without assuming that a Coulomb chart carries almost all of
that measure. They establish a uniform first-moment bound and a
staple-uniform anticoncentration estimate for local thermally
rescaled plaquette energies. These are new interfaces in this
lineage, not a transfer of the earlier chart theorems.

Maximal-tree partition-function stability is a standard method,
not a claimed new invention here. Relevant primary sources are
P.~A.~Faria da Veiga and M.~O'Carroll,
\emph{On Thermodynamic and Ultraviolet Stability of Yang--Mills},
\href{https://arxiv.org/abs/1903.09829v2}{arXiv:1903.09829v2},
which treats free-boundary \(\mathrm U(N)\) and \(\SU(N)\) stability,
and their
\emph{On Yang--Mills Stability Bounds and Plaquette Field Generating
Function},
\href{https://arxiv.org/abs/2205.07376v2}{arXiv:2205.07376v2},
which discusses \(\mathrm U(N)\) free/periodic stability and
plaquette generating functions. We do not substitute a
\(\mathrm U(N)\) result for an \(\SU(2)\) proof. Every estimate used
below is proved directly, with the periodic boundary loss retained.

\subsection{A distinct full-measure normalizer}

Let \(M\ge3\), let \(\Lambda_M=(\mathbb Z/M\mathbb Z)^4\), and set
\(V=M^4\). There are \(4V\) positively directed edges and \(6V\)
positively oriented plaquettes. A positive edge \(e=(x,\mu)\)
carries \(U_e\in\SU(2)\); reversed edges carry \(U_e^{-1}\).
For \(\mu<\nu\), use
\[
 U_p=U_\mu(x)U_\nu(x+\hat\mu)
            U_\mu(x+\hat\nu)^{-1}U_\nu(x)^{-1},\qquad
 W_p=\tfrac12\Tr U_p,\qquad s_p=1-W_p.
\]
Here \(W_p\) is real and \(0\le s_p\le2\). With every Haar measure
normalized to mass one, define
\begin{equation}
 S(U)=\sum_p s_p(U),\qquad
 Z^{\mathrm W}_{M,\beta}
   =\int e^{-\beta S(U)}\prod_e dU_e,\qquad
 d\mu^{\mathrm W}_{M,\beta}
   =(Z^{\mathrm W}_{M,\beta})^{-1}
                e^{-\beta S(U)}\prod_e dU_e .
 \label{eq:v95-full-law}
\end{equation}
Expectations with this law are denoted
\(\mathbb E^{\mathrm W}_{M,\beta}\).

There is no chart restriction, cutoff, Faddeev--Popov determinant,
or extra renewal/profile factor in \eqref{eq:v95-full-law}.
In particular it is not the previously denoted chart law
\(\mu_M\). The superscript \(\mathrm W\) must not be dropped
when comparing the two programmes.

\subsection{An exact tree gauge retaining all wrap links}

Represent vertices by \(\{0,\ldots,M-1\}^4\).
Call \((x,\mu)\) open when \(x_\mu\le M-2\), and wrap otherwise.
Define
\begin{equation}
 T=\{(x,\mu):x_\mu\le M-2,\ 
                    x_\nu=0\text{ for every }\nu>\mu\}.
 \label{eq:v95-tree}
\end{equation}
The root is the origin.

\begin{lemma}[Global Haar-product tree coordinates]
\label{lem:v95-tree}
The set \(T\) is a spanning tree with \(V-1\) edges.
There is a bijective Haar-product change of variables
\[
 (U_e)_{e}\longleftrightarrow
       ((g_x)_{x\ne0},(k_e)_{e\notin T}),
 \qquad g_0=I,\qquad k_e=I\ (e\in T),
\]
such that
\[
                  k_{xy}=g_xU_{xy}g_y^{-1}.
\]
For every gauge-invariant integrable \(F\),
\begin{equation}
       \int F(U)\prod_e dU_e
             =\int F(k)\prod_{e\notin T}dk_e.
 \label{eq:v95-tree-haar}
\end{equation}
The number of remaining group variables is
\begin{equation}
                          D=3V+1.
 \label{eq:v95-D}
\end{equation}
\end{lemma}

\begin{proof}
A vertex is reached from the origin by increasing its first
coordinate, then its second, third and fourth. Every edge in that
path lies in \(T\). Moreover,
\[
             |T|=\sum_{\mu=1}^4(M-1)M^{\mu-1}=M^4-1.
\]
Connectedness and this count prove that \(T\) is a tree.
Starting at \(g_0=I\), recursively set \(g_y=g_xU_{xy}\) on
each oriented tree edge from its parent \(x\) to its child \(y\).
This is an invertible triangular change from tree links to
nonroot frames and preserves product Haar measure by successive
left invariance. At fixed frames, each off-tree change
\(U_{xy}\mapsto g_xU_{xy}g_y^{-1}\) preserves Haar measure by
left/right invariance. The inverse is
\(U_{xy}=g_x^{-1}k_{xy}g_y\), including tree edges.
Gauge invariance removes the frames from the integrand, and their
normalized integrals equal one. Finally
\(D=4V-(V-1)\).
\end{proof}

This is an exact compact change of variables, not a logarithmic
gauge slice. All wrap variables and their possible winding
holonomies remain. No assumption about a typical representative,
chart multiplicity or a small-field event entered the argument.

\subsection{A triangular family of reference plaquettes}

There are \(4(M-1)M^3\) open edges. If an open edge
\(e=(x,\mu)\) is not in \(T\), at least one coordinate above
\(\mu\) is positive. Let
\[
 \nu(e)=\max\{\nu>\mu:x_\nu>0\},\qquad y=x-\hat\nu(e),
 \qquad p(e)=(y;\mu,\nu(e)).
\]
The number of such edges is
\begin{equation}
 N=4(M-1)M^3-(V-1)=3V-4M^3+1,\qquad D-N=4M^3.
 \label{eq:v95-N}
\end{equation}

\begin{lemma}[Independent plaquettes under reference Haar measure]
\label{lem:v95-triangular}
The plaquettes \(p(e)\) are distinct. In the tree gauge their
products satisfy
\begin{equation}
        U_{p(e)}=k_\mu(y)k_\mu(x)^{-1}.
 \label{eq:v95-triangle}
\end{equation}
Replacing the \(N\) open off-tree link variables by these
plaquette products is an invertible Haar-product change of
variables. Consequently they are independent normalized Haar
variables under the reference measure on the right side of
\eqref{eq:v95-tree-haar}.
\end{lemma}

\begin{proof}
All coordinates above \(\nu(e)\) vanish. The two \(\nu(e)\)-edges
of \(p(e)\) are therefore in \(T\); they are open because their
\(\nu(e)\)-coordinate is \(x_{\nu(e)}-1\le M-2\).
This proves \eqref{eq:v95-triangle}. A selected oriented face
identifies its lower direction \(\mu\) and its top \(\mu\)-edge
\((x,\mu)\), so the selection is injective.

Order open off-tree \(\mu\)-edges by the grade
\(\sum_{\nu>\mu}x_\nu\). The predecessor \((y,\mu)\) has grade
one less and is either in \(T\) or already processed. If
\(Q_e=U_{p(e)}\), the inverse recursion is
\[
                        k_\mu(x)=Q_e^{-1}k_\mu(y).
\]
At fixed preceding variables, inversion and translation preserve
Haar measure. Induction proves the claimed product change.
The wrap variables are unchanged and independent in this
reference measure.
\end{proof}

The independence assertion is deliberately restricted to
\emph{reference Haar measure}. The omitted plaquette terms couple
these coordinates under \(\mu^{\mathrm W}_{M,\beta}\).

\subsection{Explicit periodic pressure bounds}

Identify \(\SU(2)\) with the unit quaternions
\(u=(u_0,\mathbf u)\in S^3\). The geodesic angle \(r\in[0,\pi]\)
has \(u_0=\cos r\) and normalized radial Haar density
\((2/\pi)\sin^2r\,dr\). Set
\begin{equation}
 z(\beta)=\int_{\SU(2)}e^{-\beta(1-u_0)}du,\qquad
 c_b=\frac{8}{3\pi^3},\qquad C_1=\left(\frac{\pi}{2}\right)^{5/2}.
 \label{eq:v95-single-constants}
\end{equation}

\begin{lemma}[One-group estimates]
\label{lem:v95-one-group}
For \(\beta\ge1\),
\begin{equation}
 z(\beta)\le C_1\beta^{-3/2}.
 \label{eq:v95-z-upper}
\end{equation}
For a geodesic identity ball \(B_r\), \(0<r\le1\),
\begin{equation}
                         |B_r|_{\mathrm{Haar}}\ge c_b r^3.
 \label{eq:v95-ball}
\end{equation}
\end{lemma}

\begin{proof}
On \([0,\pi]\), \(1-\cos r\ge2r^2/\pi^2\) and
\(\sin^2r\le r^2\). Thus
\[
 z(\beta)\le\frac2\pi\int_0^\infty
       r^2e^{-2\beta r^2/\pi^2}\,dr
           =C_1\beta^{-3/2}.
\]
For \(0\le t\le r\le1<\pi/2\), \(\sin t\ge2t/\pi\).
Integrating its square against \(2/\pi\) gives
\eqref{eq:v95-ball}.
\end{proof}

\begin{theorem}[Full compact pressure with an explicit boundary loss]
\label{thm:v95-pressure}
For \(M\ge3\) and \(\beta\ge1\),
\begin{equation}
 e^{-48V}c_b^D\beta^{-3D/2}
       \ \le Z^{\mathrm W}_{M,\beta}
       \ \le z(\beta)^N
       \ \le C_1^N\beta^{-3N/2},
 \label{eq:v95-pressure}
\end{equation}
where \(D=3V+1\) and \(N=3V-4M^3+1\).
\end{theorem}

\begin{proof}
For the upper bound retain only the selected \(N\) nonnegative
plaquette actions in the exponential. Apply
Lemma~\ref{lem:v95-triangular}; their integrals give \(z(\beta)^N\),
and the remaining normalized Haar integrals give one.

For the lower bound use the tree coordinates and restrict every
off-tree link to \(B_r\), where \(r=\beta^{-1/2}\le1\).
The Euclidean quaternion chord of such a link from \(I\) is at
most \(r\), and the chord of a product of unit quaternions is at
most the sum of its factor chords. This remains true for inverse
factors. Every plaquette therefore has chord at most \(4r\).
Since \(1-u_0=|u-I|^2/2\), it follows that
\[
          s_p\le8r^2,\qquad S\le48Vr^2,\qquad
                         e^{-\beta S}\ge e^{-48V}.
\]
The restricted product-Haar volume is at least
\((c_b r^3)^D\). Multiplication gives the lower bound.
\end{proof}

The difference \(D-N=4M^3\) must be retained. For example, the
free-energy bounds have leading coefficient \(9/2\) in
\(\log\beta\), but this elementary periodic proof has a
boundary error of order \((\log\beta)/M\). It does not give
a bound uniform in arbitrarily fast simultaneous growth of
\(\beta\) and \(M\).

\subsection{An energy bound at every coupling in its stated range}

Define
\begin{equation}
 C_P=48+4\log(c_b^{-1})+3\log C_1+\frac92\log2<65.
 \label{eq:v95-CP}
\end{equation}
One direct verification uses the following strict inequalities:
\[
 \pi<22/7,\qquad \log(3\pi^3/8)<5/2,\qquad
 \log(\pi/2)<23/50,\qquad \log2<7/10.
\]
The last three comparisons follow from positive Taylor partial
sums for the exponential. They give
\[
 C_P<48+4(5/2)+3(23/20)+(9/2)(7/10)
                      =323/5<65.
\]

\begin{theorem}[Pointwise full-Wilson thermal first moment]
\label{thm:v95-energy}
For every \(M\ge3\) and \(\beta\ge2\),
\begin{align}
 \frac{\beta}{V}\mathbb E^{\mathrm W}_{M,\beta} S
       &\le 2C_P+\frac{12\log\beta}{M}
          <130+\frac{12\log\beta}{M},                 \label{eq:v95-total-energy}\\
 \mathbb E^{\mathrm W}_{M,\beta}X_p
       &\le\frac{C_P}{3}+\frac{2\log\beta}{M}
          <\frac{65}{3}+\frac{2\log\beta}{M},
 \qquad X_p:=\beta s_p.                              \label{eq:v95-local-energy}
\end{align}
No choice of a favourable coupling from an interval is required.
\end{theorem}

\begin{proof}
Finite-volume compactness and \(0\le S\le12V\) justify all
derivatives under the integral. In particular
\[
 \frac{d}{dt}\log Z^{\mathrm W}_{M,t}
       =-\mathbb E^{\mathrm W}_{M,t} S,\qquad
 \frac{d}{dt}\mathbb E^{\mathrm W}_{M,t} S
       =-\Var_{\mu^{\mathrm W}_{M,t}}S\le0.
\]
Consequently
\[
 \frac{\beta}{2}\mathbb E^{\mathrm W}_{M,\beta}S
   \le \int_{\beta/2}^{\beta}\mathbb E^{\mathrm W}_{M,t}S\,dt
   =\log\frac{Z^{\mathrm W}_{M,\beta/2}}
                    {Z^{\mathrm W}_{M,\beta}}.
\]
Use the upper pressure bound at \(\beta/2\ge1\) and the lower
bound at \(\beta\). Their ratio obeys
\begin{align*}
 \frac1V\log\frac{Z^{\mathrm W}_{M,\beta/2}}
                       {Z^{\mathrm W}_{M,\beta}}
 &\le48+\frac{D}{V}\log(c_b^{-1})
       +\frac{N}{V}\log C_1+\frac{3N}{2V}\log2
       +\frac{3(D-N)}{2V}\log\beta\\
 &\le C_P+\frac{6\log\beta}{M},
\end{align*}
because \(D/V\le4\) and \(0<N/V<3\).
This proves \eqref{eq:v95-total-energy}. Translations and
coordinate permutations preserve the full periodic law.
They act transitively on unoriented plaquettes, and \(s_p\)
is unchanged by orientation reversal. Hence
\(\mathbb E S=6V\mathbb E s_p\), giving
\eqref{eq:v95-local-energy}.
\end{proof}

For any \(\beta_M\to\infty\) with
\begin{equation}
                L:=\sup_M\frac{\log\beta_M}{M}<\infty,
 \label{eq:v95-growth-window}
\end{equation}
the full-law \(X_p\) have uniformly bounded first moments.
This includes \(\beta_M\) proportional to \(\log M\), without
identifying that choice with a physical renormalized trajectory.
Writing \(K=65/3+2L\), Markov's inequality gives
\begin{equation}
 \mu^{\mathrm W}_{M,\beta_M}(s_p>\varepsilon)
                       \le\frac{K}{\beta_M\varepsilon}.
 \label{eq:v95-local-flat}
\end{equation}
Thus each fixed plaquette becomes flat in \(L^1\). A union bound
over all \(6M^4\) plaquettes does not show a globally small-field
configuration in the logarithmic coupling window. Local flatness
is not Coulomb-chart coverage.

\subsection{The exact six-staple conditional on the full law}

Fix a plaquette \(p\) and one of its edges \(e\).
Condition on all link variables except \(U_e=u\). Since \(M\ge3\),
exactly six distinct plaquettes contain this edge, once each.
For any fixed values of the other links, their normalized traces
are linear functions of the unit quaternion \(u\):
\begin{equation}
 W_{p_j}(u)=u\cdot a_j,\qquad |a_j|=1,\qquad
 A=\sum_{j=1}^6a_j,\quad |A|\le6.
 \label{eq:v95-staples}
\end{equation}
For example the scalar part of \(uB\) is
\(u\cdot(B_0,-\mathbf B)\), while that of \(u^{-1}B\)
is \(u\cdot(B_0,\mathbf B)\). Cyclic invariance of trace
reduces either position of \(u\) or \(u^{-1}\) to these forms.
Therefore the exact conditional probability has density
\begin{equation}
             \frac{e^{\beta u\cdot A}}
                    {\int_{S^3}e^{\beta v\cdot A}\,dv}\,du.
 \label{eq:v95-link-conditional}
\end{equation}
This formula is specific to \eqref{eq:v95-full-law}.
It is not the conditional density of a determinant-weighted
chart law or a Wilson law multiplied by an uncontrolled profile.

The archived f6 V95 already proves a different, \(\SU(3)\)
one-link concentration statement under a small-frustration
alignment hypothesis. We do not repeat that well analysis.
In \(\SU(2)\), the rotational symmetry of \(S^3\) instead gives
the following \emph{anticoncentration} estimate, uniformly for
all staples, aligned or frustrated.

\subsection{A cap lower bound uniform in the conditional field}

Rotate quaternion coordinates so that \(\beta A\) points along
the first axis. Set
\[
 \kappa=\beta|A|\in[0,6\beta],\qquad
 h=(1+\kappa)^{-1/2},\qquad
 C_*=\frac{\pi^{5/2}}2=2^{3/2}C_1.
\]
If \(A=0\), choose any rotation. Write \(\nu_\kappa\) for the
probability on \(S^3\) with density proportional to
\(e^{\kappa u_0}\).
For \(t\) in the three-dimensional ball
\(B=\{t\in\mathbb R^3:|t|\le1/2\}\), define
\[
             \Psi_h(t)=\big(\sqrt{1-h^2|t|^2},ht\big).
\]
Let \(\lambda_B\) be uniform probability on \(B\).

\begin{lemma}[A uniformly retained cap]
\label{lem:v95-cap}
For every \(\kappa\ge0\) and Borel set \(E\subset B\),
\begin{equation}
 \nu_\kappa(\Psi_h(E))\ge\gamma\,\lambda_B(E),
       \qquad \gamma=\frac{e^{-1/4}}{6\pi^{7/2}}.
 \label{eq:v95-cap}
\end{equation}
\end{lemma}

\begin{proof}
The normalizing integral is \(e^\kappa z(\kappa)\).
For \(\kappa\ge1\), Lemma~\ref{lem:v95-one-group} gives
\[
 z(\kappa)\le C_1\kappa^{-3/2}
                  \le C_*(1+\kappa)^{-3/2}.
\]
For \(0\le\kappa\le1\), use
\(z(\kappa)\le1\le2^{3/2}(1+\kappa)^{-3/2}
                 \le C_*(1+\kappa)^{-3/2}\).
Thus in every case the normalizer is at most
\(C_*e^\kappa h^3\).

The positive-hemisphere Haar density in
\(v\mapsto(\sqrt{1-|v|^2},v)\) is
\[
                      \frac{dv}{2\pi^2\sqrt{1-|v|^2}}.
 \]
On \(v=ht\), \(|t|\le1/2\), we have
\[
 \kappa\big(1-\sqrt{1-h^2|t|^2}\big)
        \le\kappa h^2|t|^2\le\frac14 .
\]
After division by the normalizer, the cap density with
respect to \(dt\) is consequently at least
\(e^{-1/4}/(2\pi^2C_*)\).
The Euclidean volume of \(B\) is \(\pi/6\).
Multiplication by that volume gives
\eqref{eq:v95-cap}. This is a measure lower bound on a cap,
not an assertion that the full conditional law is uniform there.
\end{proof}

\subsection{An interval-escape estimate for every plaquette direction}

For a unit vector \(a=(a_0,\mathbf a)\in\mathbb R^4\), the
plaquette trace on this cap has the form
\[
 F_a(t)=a_0\sqrt{1-h^2|t|^2}+h\,\mathbf a\cdot t.
\]
Both the radial and tangential cases matter: the radial case
has fluctuation scale \(h^2\), not \(h\).

\begin{lemma}[Uniform ball interval escape]
\label{lem:v95-escape}
For every \(0<h\le1\), every unit \(a\), and every real interval
\(I\) of length at most \(h^2/32\),
\begin{equation}
                 \lambda_B(F_a\notin I)\ge\frac1{16800}.
 \label{eq:v95-ball-escape}
\end{equation}
\end{lemma}

\begin{proof}
Elementary radial integration for uniform probability on the
radius-\(1/2\) ball in \(\mathbb R^3\) gives
\[
 \mathbb E_{\lambda_B}|t|^2=\frac3{20},\qquad
 \mathbb E_{\lambda_B}|t|^4=\frac3{112},\qquad
 \Var_{\lambda_B}|t|^2=\frac3{700},\qquad
 \mathbb E_{\lambda_B}(\mathbf a\cdot t)^2
                                      =\frac{|\mathbf a|^2}{20}.
\]
For \(g(z)=\sqrt{1-h^2z}\) on \(0\le z\le1/4\),
\(|g'(z)|\ge h^2/2\). The independent-copy formula for variance
therefore yields
\[
 \Var_{\lambda_B}\sqrt{1-h^2|t|^2}
        \ge\frac{h^4}{4}\Var_{\lambda_B}|t|^2
        =\frac{3h^4}{2800}.
\]
The covariance of the even radial part and the odd linear part
is zero. Defining
\[
 \sigma^2=a_0^2h^4+|\mathbf a|^2h^2\ge h^4,\qquad
                              c_0=\frac3{2800},
\]
we obtain
\begin{equation}
                        \Var_{\lambda_B}F_a\ge c_0\sigma^2.
 \label{eq:v95-cap-variance}
\end{equation}
On the same ball,
\[
 |\nabla F_a|
    \le |a_0|h^2/\sqrt3+h|\mathbf a|
    \le(2/\sqrt3)\sigma<2\sigma.
\]
Since its diameter is one, the range of \(F_a\) has length at
most \(2\sigma\).

It suffices to consider closed intervals of length \(h^2/32\),
since an interval of no greater length can be contained in one.
If such an \(I\) misses the range of \(F_a\), the conclusion is
immediate. Otherwise let \(\tau\) be its centre and
\(p_{\rm out}=\lambda_B(F_a\notin I)\). On \(I\),
\(|F_a-\tau|\le h^2/64\le\sigma/64\); everywhere on \(B\),
\(|F_a-\tau|\le2\sigma+\sigma/64<3\sigma\).
Using \eqref{eq:v95-cap-variance},
\[
 c_0\sigma^2
   \le\mathbb E_{\lambda_B}(F_a-\tau)^2
   \le \frac{\sigma^2}{4096}+9\sigma^2p_{\rm out}.
\]
Since \(1/4096\le c_0/2\), this implies
\(p_{\rm out}\ge c_0/18=1/16800\).
\end{proof}

\begin{theorem}[Full-Wilson thermal anticoncentration]
\label{thm:v95-anticoncentration}
Set
\begin{equation}
       w_*=\frac1{224},\qquad
       \delta_*=\frac{e^{-1/4}}{100800\pi^{7/2}}>0.
 \label{eq:v95-anticonstants}
\end{equation}
For every \(M\ge3\), \(\beta\ge1\), plaquette \(p\), edge
\(e\subset p\), and interval \(J\) of length at most \(w_*\),
\begin{equation}
 \mu^{\mathrm W}_{M,\beta}
   \big(X_p\in J\mid(U_f)_{f\ne e}\big)
                         \le1-\delta_* .
 \label{eq:v95-conditional-anti}
\end{equation}
The bound holds for the explicit conditional kernel at every
choice of exterior links, and hence also unconditionally.
In particular
\begin{align}
 \mathbb E^{\mathrm W}_{M,\beta}X_p
                              &\ge\delta_*w_*,
                                                       \label{eq:v95-mean-lower}\\
 \Var_{\mu^{\mathrm W}_{M,\beta}}X_p
                              &\ge\delta_*w_*^2/4 .
                                                       \label{eq:v95-var-lower}
\end{align}
\end{theorem}

\begin{proof}
In the conditional coordinates, \(W_p=u\cdot a\) for a unit
vector \(a\), so Lemmas~\ref{lem:v95-cap} and
\ref{lem:v95-escape} show that every trace interval of length
at most \(h^2/32\) has escape probability at least
\(\gamma/16800=\delta_*\).
Since \(\kappa\le6\beta\) and \(\beta\ge1\),
\[
           \beta h^2=\frac{\beta}{1+\kappa}
                    \ge\frac{\beta}{1+6\beta}\ge\frac17.
\]
The inverse image of \(J\) under \(W_p\mapsto\beta(1-W_p)\)
has length at most
\[
                     \frac1{224\beta}\le\frac{h^2}{32}.
\]
This proves the conditional bound. Integration over exterior
links proves the unconditional one.

Taking \(J=[0,w_*]\), and using \(X_p\ge0\), gives
\(\Pr(X_p>w_*)\ge\delta_*\) and
\eqref{eq:v95-mean-lower}. Taking the length-\(w_*\) closed
interval centred at \(\mathbb E X_p\) gives
\eqref{eq:v95-var-lower}. All moments at fixed \(M,\beta\)
exist because \(0\le X_p\le2\beta\).
\end{proof}

The conditioning in this theorem is on the \emph{other links},
not on all other plaquette energies. Those energies can also
depend on \(U_e\). Nor can one multiply the escape probabilities
of overlapping plaquettes as if they were independent.

\subsection{Tight noncollapsed local energy laws on the full measure}

\begin{theorem}[A full-compact microscopic subsequential law]
\label{thm:v95-local-limit}
Let \(M\to\infty\), \(\beta_M\to\infty\), and assume
\eqref{eq:v95-growth-window}. Periodically extend the field
\[
       X^{(M)}_{x,\mu\nu}
          =\beta_M\big(1-\tfrac12\Tr U_{(x;\mu,\nu)}\big)
 \]
to the countable index set
\(\mathbb Z^4\times\{(\mu,\nu):1\le\mu<\nu\le4\}\).
The laws of these nonnegative fields under
\(\mu^{\mathrm W}_{M,\beta_M}\) are tight in the product topology.
Every subsequential limit is translation invariant and inherits
the lattice coordinate symmetries. Every coordinate \(X_p^\infty\)
of every such limit is nonconstant in law. More quantitatively,
for every open interval \(J\) of length \(w_*\),
\begin{equation}
             \Pr(X_p^\infty\in J)\le1-\delta_*,
       \qquad \mathbb E X_p^\infty\le K=65/3+2L.
 \label{eq:v95-limit-anti}
\end{equation}
\end{theorem}

\begin{proof}
Discard finitely many terms so that \(\beta_M\ge2\).
Theorem~\ref{thm:v95-energy} bounds every coordinate mean by \(K\).
Enumerate the countable index set as \(p_1,p_2,\ldots\).
For any \(\varepsilon>0\), set
\(R_j=2^jK/\varepsilon\). Markov's inequality and the union
bound give
\[
 \Pr\big(\exists j:X^{(M)}_{p_j}>R_j\big)
                         \le\sum_{j\ge1}K/R_j=\varepsilon .
\]
The product \(\prod_j[0,R_j]\) is compact in this countable
product topology. The ambient countable product of
\([0,\infty)\) is Polish, so this proves tightness and
subsequential weak compactness.

Translation and coordinate symmetries hold for each periodic
law and act continuously on the product space; they pass to
weak limits. For an open real interval \(J\) of length \(w_*\),
the corresponding cylinder event is open, and Portmanteau gives
\[
 \Pr(X_p^\infty\in J)
       \le\liminf_M\Pr(X^{(M)}_p\in J)\le1-\delta_*.
\]
An interval of length \(w_*\) centred on a putative deterministic
value contradicts this bound, so the limiting coordinate is
not a point mass. Finally, apply weak convergence to the bounded
continuous truncations \(x\mapsto\min(x,R)\), then let
\(R\to\infty\), to obtain the stated upper first-moment bound.
\end{proof}

Anticoncentration is important here: a finite-volume variance
lower bound alone would not rule out weak collapse caused by
rare large values. The interval bound passes to the limit without
assuming uniform integrability of \(X_p^2\).
Neither the first-moment bound nor this argument proves
uniform integrability of \(X_p\), convergence of its means,
or the existence of limiting second moments. It also does not
identify the limiting law as Gaussian or chi-square.

There is one additional full-measure consequence requiring no
new group expansion. For a fixed contractible lattice rectangle
\(C\) tiled by \(A_C\) elementary plaquettes, V91's deterministic
ordered-product estimate holds for every compact link
configuration:
\[
          1-W_C\le A_C\sum_{p\text{ in its surface}}s_p.
\]
Consequently under the present full law,
\begin{equation}
       \mathbb E^{\mathrm W}_{M,\beta_M}
                    \big[\beta_M(1-W_C)\big]\le A_C^2K.
 \label{eq:v95-fixed-loop}
\end{equation}
Thus any fixed finite collection of such microscopic loop
defects is tight as well. We have used the deterministic V91
inequality, not transferred its chart-law limiting distribution.

\subsection{What this does and does not remove from the programme}

The result removes a specific chart-coverage premise from
one local task: thermal plaquette energies of the full compact
\(\SU(2)\) Wilson law now have a uniform first moment and
noncollapsed microscopic subsequential laws in the stated
growth window, at every coupling in that window.
The tree-gauge computation includes the wrap links exactly,
and the anticoncentration proof does not discard frustrated
staple configurations.

The constants are intentionally coarse. In particular
\eqref{eq:v95-total-energy} is not the sharp \(9/2\) saturation
used to eliminate the homogeneous chart mode in V90--V94.
There is no justified equality of the two local limiting laws.
A full-measure analogue of sharp energy saturation would need
additional work on the interacting pressure or an appropriate
local comparison, not a relabelling of the chart normalizer.

The index \(x\) in Theorem~\ref{thm:v95-local-limit} is fixed in
\emph{lattice units}. When physical spacing is \(a=M^{-1}\),
distinct fixed indices approach the same physical point.
This theorem does not construct a continuum random
distribution, prove a physically normalized covariance bound,
or control loops whose area grows with \(M\).
For example \eqref{eq:v95-fixed-loop} grows as \(A_C^2\);
it is not a physical area law or an exponential clustering
estimate. The \(\SU(2)\) argument also does not establish the
corresponding statements for every compact simple gauge group.

The upstream targets are now more concrete: control
source-local or block partition ratios beyond a bulk
first-moment estimate; obtain full-law higher-moment or
physically weighted correlation estimates; and identify a
nontrivial renormalized trajectory with the reconstruction
properties required by the original problem.
The existing strong-mixing basin is still only a conditional
endpoint, not a proven destination of that trajectory.

\subsection{Exact verification and preservation}

The certificate is
\begin{center}\small
\path{scripts/verify_ym_f14_v95_full_compact_energy.py}.
\end{center}
It checks the spanning-tree and triangular-plaquette counts on
seven tori, of sides three through nine. On the side-three
torus it verifies 136 exact quaternion plaquette-to-link
reconstructions, 80 nonroot-frame reconstructions and 486
gauge-covariance identities. It verifies six unit staple
coefficient vectors, 384 plaquette-linearity identities and
64 six-staple sum identities, covering both edge orientations.
Three rational exponential comparisons certify the logarithmic
majorants in \(C_P<65\); 27 further size checks retain the
periodic boundary coefficient. Exact uniform-ball moments and
500 rational scale/variance comparisons audit the
anticoncentration constants.

These are finite algebraic diagnostics. They do not prove
Haar invariance, estimates for every real parameter, tightness,
or a mass gap by sampling. Those mathematical statements, in
their precise scopes, are established by the analytic proofs above.

The certificate contains 8,807 bytes and has SHA--256
\begin{center}\scriptsize\ttfamily
B4AC549F633C6DF221018E49B406908BBA4B168864688E7F5471537D22EC27C1.
\end{center}
The predecessor V94 contains 2,310,510 bytes and has SHA--256
\begin{center}\scriptsize\ttfamily
14B85D1BF64EBD1EFC130F8C0D89D050057396A8FE332AA74CB795F5DD0F313D.
\end{center}
Its body is preserved except for the title version and this
terminal insertion. V95 replaces V94 as the sole active main
note; the predecessor is recoverable by reversing that insertion
and title change. Earlier certificates, companions, supplied
files, monographs and archives remain unchanged.

\begin{firewall}
V95 proves full compact periodic \(\SU(2)\) pressure bounds,
pointwise thermal first-moment control, staple-uniform local
anticoncentration, and tight noncollapsed microscopic energy
subsequential laws. These are not a construction of interacting
four-dimensional continuum Yang--Mills theory and not a
positive physical Hamiltonian mass gap. Both remain unproved.
The original goal remains active.
\end{firewall}
% V96 terminal insertion: allow three-digit section numbers in the contents.
\makeatletter
\addtocontents{toc}{\protect\renewcommand{\protect\l@section}{\protect\@dottedtocline{0}{0em}{2.5em}}}
\makeatother
\section{V96: full-measure staple alignment and a local Gaussian response}
\label{sec:v96-staple-response}

\subsection{The next step and its scope}

V95 established tight, noncollapsed microscopic plaquette-energy
laws for the full compact periodic \(\SU(2)\) Wilson measure.
This is genuine progress beyond a chart-restricted measure, but
the first-moment estimate alone did not identify a limiting local
law. Here we use that same estimate to control the frustration
of the six staples around one link, and then analyse the
\emph{exact} Wilson conditional distribution of that link.

The result is a quantitative Gaussian response in an
exterior-dependent orthogonal frame, followed by an explicit
mixture representation for the six adjoining thermal energies.
The exterior-dependent centres are not discarded. Their
distribution and correlations contain unresolved interaction
information. In particular, a local Gaussian response is not
a Gaussian law for the whole field or a mass-gap estimate.

The standard \(\SU(2)\) heat-bath conditional is not a new
construction. For historical context see A.~D.~Kennedy and
B.~J.~Pendleton, \emph{Improved heatbath method for Monte Carlo
calculations in lattice gauge theories}, Physics Letters B
156 (1985), 393--399,
\href{https://www.sciencedirect.com/science/article/pii/0370269385916326}
{publisher record}.
The arguments below are self-contained probability estimates
for that conditional, not claims about an algorithm's mixing time.

This also differs from the Gaussian shell innovations in
shell-mixing V16: those are innovations of a specified Gaussian
block covariance. Here the input is the full nonlinear compact
Wilson law and the Gaussian comparison is derived from its
one-link conditional. The last companion checkpoint was V95;
the next remains V100. No companion or archived source is edited.

\subsection{Exterior data, frames, and thermal frustration}

Fix \(M\ge3\), \(\beta\ge2\), and the full law
\(\mu^{\mathrm W}_{M,\beta}\) of \eqref{eq:v95-full-law}. Put
\begin{equation}
               K=K_{M,\beta}:=\frac{65}{3}+\frac{2\log\beta}{M}.
 \label{eq:v96-K}
\end{equation}
V95 gives \(\mathbb E^{\mathrm W}_{M,\beta}X_p\le K\) for every
plaquette, where \(X_p=\beta(1-\tfrac12\Tr U_p)\).
Fix one positive edge \(e\) and denote all other link variables
by \(\xi\). List the six plaquettes containing \(e\) as
\(p_1,\ldots,p_6\), and write their energies as \(X_1,\ldots,X_6\).

As in \eqref{eq:v95-staples}, with \(u=U_e\in S^3\),
\[
 W_{p_j}=u\cdot a_j(\xi),\qquad |a_j|=1,\qquad
 A=\sum_{j=1}^6a_j,\qquad \alpha=|A|\in[0,6].
\]
If \(\alpha>0\), set \(n=A/\alpha\). If \(\alpha=0\), set
\(n=(1,0,0,0)\). Choose a Borel orthogonal matrix \(R_n\)
with \(R_nn=(1,0,0,0)\). One concrete choice is the Householder
reflection with vector \(n-(1,0,0,0)\), taking the identity at
the exceptional point \(n=(1,0,0,0)\).
No continuous choice of frame or global gauge is required.

Write
\begin{equation}
 \begin{gathered}
 R_nu=(u_0,v),\qquad R_na_j=(c_j,\eta_j),\\
 b_j=\sqrt\beta\,\eta_j,\quad d_j=\beta(1-c_j),\quad
 Z=\sqrt\beta\,v,\quad Q=\beta(1-u_0).
 \end{gathered}
 \label{eq:v96-frame}
\end{equation}
Thus \(u_0,v\) here are the coordinates \emph{after} the
exterior-dependent rotation. Define the thermal staple
frustration by
\begin{equation}
                           F=\beta(6-\alpha).
 \label{eq:v96-frustration}
\end{equation}
All of \(F,\alpha,b_j,d_j\) are functions of \(\xi\) alone.
The bounds below include \(\alpha=0\).

\begin{lemma}[Exact star identities and full-measure alignment]
\label{lem:v96-alignment}
For every configuration,
\begin{align}
 \sum_{j=1}^6d_j&=F,&
 \sum_{j=1}^6b_j&=0,&
 |b_j|^2&=2d_j-d_j^2/\beta,                         \label{eq:v96-centres}\\
 \sum_{j=1}^6X_j&=F+\alpha Q,&
 0\le Q&\le2\beta,&
 \sum_{j=1}^6|b_j|^2&\le2F.                         \label{eq:v96-star}
\end{align}
With expectations under the full Wilson measure,
\begin{equation}
 \begin{gathered}
 \mathbb EF\le6K,\qquad
 \mathbb E\sum_j|b_j|^2\le12K,\\
 \Pr(\alpha<3)\le\frac{2K}{\beta},\qquad
 \mathbb EQ\le6K.
 \end{gathered}
 \label{eq:v96-first-budgets}
\end{equation}
In particular,
\begin{equation}
 \mathbb E\sum_{j=1}^6|a_j-n|^2
                         \le\frac{12K}{\beta}.
 \label{eq:v96-staple-alignment}
\end{equation}
\end{lemma}

\begin{proof}
Orthogonality preserves dot products and unit lengths.
The sum of the rotated \(a_j\)'s is \((\alpha,0)\), including
when \(A=0\). Therefore \(\sum_jc_j=\alpha\),
\(\sum_j\eta_j=0\), and
\(|b_j|^2=\beta(1-c_j^2)=2d_j-d_j^2/\beta\).
Moreover
\[
 \sum_jX_j=\beta(6-u\cdot A)
          =\beta(6-\alpha u_0)=F+\alpha Q\ge F.
\]
Taking expectations and using \(\mathbb EX_j\le K\)
gives the first two bounds in \eqref{eq:v96-first-budgets}.
The event \(\alpha<3\) implies \(F>3\beta\), giving the
probability bound by Markov's inequality.

On \(\{\alpha\ge3\}\), \(Q\le\sum_jX_j/3\), so
\(\mathbb E[Q;\alpha\ge3]\le2K\).
On its complement use \(Q\le2\beta\) and the preceding
probability bound to get \(\mathbb E[Q;\alpha<3]\le4K\).
Finally
\[
 \sum_j|a_j-n|^2
      =2\sum_j(1-a_j\cdot n)=2(6-\alpha)=2F/\beta,
\]
which proves \eqref{eq:v96-staple-alignment}.
\end{proof}

This is alignment in mean under the \emph{full} measure, not an
assumption that every exterior is aligned. A large-frustration
exterior still occurs in the space of configurations and must
be included in any unconditional estimate.

\subsection{A scalar spherical conditional and its Gaussian limit}

Let \(\nu_\kappa\) be normalized probability on \(S^3\) with
density proportional to \(e^{\kappa u_0}\), \(\kappa\ge0\).
We use the total-variation convention
\[
       d_{\mathrm{TV}}(\mu,\nu)
                  =\sup_E|\mu(E)-\nu(E)|
                  =\tfrac12\|p-q\|_1
\]
when densities exist. The letter \(I_3\) denotes the
three-dimensional identity matrix.

\begin{lemma}[Explicit spherical-to-Gaussian total variation]
\label{lem:v96-spherical-TV}
If \((u_0,v)\sim\nu_\kappa\) and \(\kappa\ge1\), then
\begin{equation}
 d_{\mathrm{TV}}\big(\mathcal L(\sqrt\kappa\,v),N(0,I_3)\big)
                              \le\min\{1,600/\kappa\}.
 \label{eq:v96-spherical-TV}
\end{equation}
Both hemispheres are included in this estimate.
\end{lemma}

\begin{proof}
Put \(y=\sqrt\kappa\,v\), \(r=|y|\), and
\(s=(1-r^2/\kappa)^{1/2}\) for \(r<\sqrt\kappa\).
The pushforward density is proportional to
\begin{align*}
 h_\kappa^+(y)&=s^{-1}e^{\kappa(s-1)}
                              1_{\{r^2<\kappa\}},\\
 h_\kappa^-(y)&=s^{-1}e^{-\kappa(s+1)}
                              1_{\{r^2<\kappa\}}.
\end{align*}
The common factor from Haar measure and rescaling cancels in
normalization. The two terms are the positive and negative
hemispheres, respectively. Let
\[
 g(y)=e^{-|y|^2/2},\qquad J_0=\int_{\mathbb R^3}g(y)\,dy
                                      =(2\pi)^{3/2}.
\]
On \(r^2\le\kappa/2\) we have \(s^{-1}\le2\),
\(s^{-1}-1\le2r^2/\kappa\), and
\[
 b:=\kappa(1-s)-r^2/2
      =\frac{r^4}{2\kappa(1+s)^2}
                  \in[0,r^4/(2\kappa)].
\]
Since \(1-e^{-b}\le b\),
\[
 |h_\kappa^+(y)-g(y)|
       \le \frac{g(y)}{\kappa}(2r^2+r^4).
\]
The Gaussian moments \(\mathbb E|G|^2=3\) and
\(\mathbb E|G|^4=15\), \(G\sim N(0,I_3)\), bound this
interior \(L^1\) error by \(21J_0/\kappa\).

For the exterior of that region,
\(s\le1/\sqrt2<3/4\), so the positive-hemisphere exponential
is at most \(e^{-\kappa/4}\).
The negative-hemisphere exponential is at most \(e^{-\kappa}\)
on its entire domain, and hence also at most \(e^{-\kappa/4}\).
The unweighted chart integral on one hemisphere is
\[
        \int_{r^2<\kappa}(1-r^2/\kappa)^{-1/2}\,dy
                                      =\pi^2\kappa^{3/2}.
\]
The positive exterior and the entire negative term thus have
combined mass at most \(2\pi^2\kappa^{3/2}e^{-\kappa/4}\).
Since \(e^{\kappa/4}\ge(\kappa/4)^3/6\) and \(\kappa\ge1\),
this is at most \(768\pi^2/\kappa\).
The Gaussian exterior mass is at most
\[
       \frac2\kappa\int r^2g(y)\,dy=6J_0/\kappa.
\]
Combining these bounds gives
\[
 \|h_\kappa^++h_\kappa^--g\|_1
       \le\frac{27J_0+768\pi^2}{\kappa}
       <\frac{600J_0}{\kappa}.
\]
For the last numerical comparison,
\(\pi^2/J_0=\sqrt{\pi/8}<2/3\), so the coefficient is
less than \(27+512=539<600\).

For any nonnegative integrable \(h,g\) with positive integrals,
\[
 d_{\mathrm{TV}}\left(\frac{h}{\int h},\frac{g}{\int g}\right)
                          \le\frac{\|h-g\|_1}{\int g}.
\]
Indeed, inserting \(h/\int g\) bounds the normalized \(L^1\)
distance by twice the right side. Apply this inequality and
the universal bound \(d_{\mathrm{TV}}\le1\).
\end{proof}

\begin{lemma}[Radial thermal moments]
\label{lem:v96-radial-Gamma}
Let
\[
                    m(\kappa)=\mathbb E_{\nu_\kappa}u_0.
\]
Then \(m(0)=0\), \(0\le m(\kappa)<1\), and
\(m'(\kappa)=\Var_{\nu_\kappa}u_0\ge0\).
As \(\kappa\to\infty\),
\begin{equation}
        T_\kappa:=\kappa(1-u_0)
                   \ \Longrightarrow\ \mathrm{Gamma}(3/2,1),
 \label{eq:v96-Gamma}
\end{equation}
with convergence of every nonnegative integer moment. In
particular
\begin{equation}
       \kappa(1-m(\kappa))\longrightarrow\frac32,\qquad
       \Var_{\nu_\kappa}T_\kappa\longrightarrow\frac32.
 \label{eq:v96-radial-moments}
\end{equation}
The second parameter of \(\mathrm{Gamma}(3/2,1)\) is its rate.
\end{lemma}

\begin{proof}
Differentiating the compact spherical integral gives
\(m'=\Var u_0\); reflection symmetry gives \(m(0)=0\), and
strictly positive smooth density on the sphere implies \(m<1\).
The first coordinate of normalized uniform \(S^3\) measure
has density \((2/\pi)\sqrt{1-u_0^2}\) on \([-1,1]\).
Consequently the density of \(T_\kappa=q\), after dropping
constants cancelled by normalization, is
\[
       e^{-q}q^{1/2}\sqrt{2-q/\kappa}\,
                        1_{\{0<q<2\kappa\}}.
\]
It converges pointwise on \((0,\infty)\) to
\(\sqrt2 e^{-q}q^{1/2}\).
For every integer \(r\ge0\), the \(q^r\)-weighted density is
dominated by \(\sqrt2e^{-q}q^{r+1/2}\), an integrable function.
Dominated convergence for the normalizer and all these weighted
integrals proves the assertion. The limiting mean and variance
are both \(3/2\).
\end{proof}

\subsection{A Gaussian fluctuation independent of the exterior in the limit}

Let \(\mathbb P_\xi\) be the exterior marginal of the full Wilson
law. Its exact conditional in the frame \eqref{eq:v96-frame}
is \(\nu_{\beta\alpha}\), by \eqref{eq:v95-link-conditional}.
Set
\begin{equation}
                \mathcal G=N(0,I_3/6),\qquad
                \epsilon_{M,\beta}
                      =\min\{1,(200+5K)/\beta\}.
 \label{eq:v96-epsilon}
\end{equation}

\begin{theorem}[An averaged conditional total-variation bound]
\label{thm:v96-innovation}
For every \(M\ge3\) and \(\beta\ge2\),
\begin{equation}
 \mathbb E_\xi\,d_{\mathrm{TV}}
          \big(\mathcal L(Z\mid\xi),\mathcal G\big)
                      \le\epsilon_{M,\beta}.
 \label{eq:v96-averaged-TV}
\end{equation}
In particular the actual joint law of \((\xi,Z)\) obeys
\begin{equation}
 d_{\mathrm{TV}}\big(\mathcal L(\xi,Z),
                    \mathbb P_\xi\otimes\mathcal G\big)
                      \le\epsilon_{M,\beta}.
 \label{eq:v96-joint-TV}
\end{equation}
The product on the right retains the actual exterior marginal.
\end{theorem}

\begin{proof}
On \(\{\alpha\ge3\}\), Lemma~\ref{lem:v96-spherical-TV},
followed by rescaling by \(1/\sqrt\alpha\), gives
\[
 d_{\mathrm{TV}}\big(\mathcal L(Z\mid\xi),
                                  N(0,I_3/\alpha)\big)
                           \le600/(\beta\alpha)\le200/\beta.
\]
For \(a\in[3,6]\), the density
\(\phi_a(z)=(a/(2\pi))^{3/2}e^{-a|z|^2/2}\) satisfies
\[
 \int|\partial_a\phi_a(z)|\,dz
    \le\frac{3}{2a}+\frac12\mathbb E_{\phi_a}|z|^2
    =\frac3a\le1.
\]
Integrating in \(a\), and retaining the factor \(1/2\)
in total variation, yields
\[
       d_{\mathrm{TV}}\big(N(0,I_3/\alpha),\mathcal G\big)
                                    \le(6-\alpha)/2=F/(2\beta).
\]
On \(\{\alpha<3\}\) use only \(d_{\mathrm{TV}}\le1\).
Lemma~\ref{lem:v96-alignment} then gives
\[
 \mathbb E_\xi d_{\mathrm{TV}}
     \le\frac{200}{\beta}+\frac{\mathbb EF}{2\beta}
                           +\Pr(\alpha<3)
     \le\frac{200+5K}{\beta}.
\]
The integrand is also at most one. Integrating conditional
probabilities over sections of any measurable joint event
proves \eqref{eq:v96-joint-TV}.
\end{proof}

The qualifier \emph{averaged} is essential. For a fixed exterior
with \(A=0\), the conditional is uniform on \(S^3\), not a
concentrated Gaussian. Such exteriors are controlled in
probability by the full-measure energy bound, not excluded
from the definition of the model.

\subsection{An exact quadratic energy plus a nonnegative remainder}

\begin{lemma}[Global two-component chord decomposition]
\label{lem:v96-energy-identity}
For every exterior and every link value, including the negative
hemisphere in \eqref{eq:v96-frame},
\begin{equation}
       X_j=\frac12|Z-b_j|^2+\frac{(d_j-Q)^2}{2\beta},
                          \qquad 1\le j\le6.
 \label{eq:v96-energy-identity}
\end{equation}
Thus, with
\(R_{\max}=\max_j|X_j-\tfrac12|Z-b_j|^2|\),
\begin{equation}
          0\le R_{\max}\le\frac{(F+Q)^2}{2\beta},
          \qquad \mathbb E(F+Q)\le12K.
 \label{eq:v96-remainder}
\end{equation}
\end{lemma}

\begin{proof}
For unit vectors \(u,a_j\), \(1-u\cdot a_j=|u-a_j|^2/2\).
In the rotated coordinates their transverse difference is
\((Z-b_j)/\sqrt\beta\), and their first-coordinate difference
is
\[
                   u_0-c_j=(d_j-Q)/\beta.
\]
Multiplication of the chord identity by \(\beta\) proves
\eqref{eq:v96-energy-identity}. Since \(0\le d_j\le F\),
the remainder bound follows; its first-moment budget is
Lemma~\ref{lem:v96-alignment}.
\end{proof}

No expansion of the exponential map or sign choice for
\(u_0=\pm\sqrt{1-|v|^2}\) is used in this identity.
There is also no unsupported fourth-moment estimate:
boundedness of \(\mathbb E(F+Q)\) implies smallness of the
remainder in probability, not a bound on \(\mathbb E(F+Q)^2\).

\begin{theorem}[Quantitative comparison of the six-energy experiment]
\label{thm:v96-energy-experiment}
Let \(\Phi(\xi,x_1,\ldots,x_6)\) be measurable, bounded in
absolute value by one, and, uniformly in \(\xi\), Lipschitz
with constant one in \(x\) for the \(\ell^\infty\) norm.
Let \(G\sim\mathcal G\) be independent of \(\xi\). Then
\begin{align}
 &\left|\mathbb E\Phi(\xi,X_1,\ldots,X_6)
       -\mathbb E_\xi\mathbb E_G
          \Phi\left(\xi,\left(\tfrac12|G-b_j(\xi)|^2\right)_{j=1}^6\right)
       \right| \notag\\
 &\hspace{3cm}\le\frac{12K}{\sqrt\beta}
                          +2\epsilon_{M,\beta}.
 \label{eq:v96-energy-experiment}
\end{align}
\end{theorem}

\begin{proof}
First keep the actual \((\xi,Z)\) and replace \(X_j\) by
\(\tfrac12|Z-b_j|^2\). The change of the test function is
bounded by
\[
 \min\{2,R_{\max}\}
    \le\min\{2,(F+Q)^2/(2\beta)\}
    \le(F+Q)/\sqrt\beta.
\]
The last inequality is \(\min\{2,t^2/2\}\le t\) for \(t\ge0\).
Expectation bounds this change by \(12K/\sqrt\beta\).
The remaining change replaces the joint law of \((\xi,Z)\)
by \(\mathbb P_\xi\otimes\mathcal G\). The composed test is
still bounded by one, so \eqref{eq:v96-joint-TV} bounds that
change by \(2\epsilon_{M,\beta}\).
\end{proof}

The comparison preserves \emph{all} exterior data in the
bounded test. It does not replace the exterior's distribution
by independent staples or by a free-field reference.

\subsection{The limiting mixture and what remains random}

Now let \(M\to\infty\), \(\beta_M\to\infty\), with
\(\sup_M(\log\beta_M)/M=L<\infty\), as in V95.
In the following limits one can use the fixed upper bound
\(K_*=65/3+2L\). Fix the edge \(e\) in lattice coordinates.

\begin{theorem}[A one-star limiting-law representation]
\label{thm:v96-mixture}
The following joint laws are tight:
\[
                 \mathcal L(b_1,\ldots,b_6,F,Z,X_1,\ldots,X_6).
\]
Along every subsequence there is a further subsequence on
which they converge in law to variables satisfying
\begin{equation}
 \begin{gathered}
 \sum_{j=1}^6B_j=0,\qquad H=\frac12\sum_{j=1}^6|B_j|^2,\\
 G\sim N(0,I_3/6),\qquad G\text{ independent of }(B,H),
 \end{gathered}
 \label{eq:v96-latent}
\end{equation}
and
\begin{equation}
             \mathcal X_j=\frac12|G-B_j|^2,\qquad
             \sum_{j=1}^6\mathcal X_j=H+3|G|^2.
 \label{eq:v96-mixture}
\end{equation}
Here \(b_j\to B_j\), \(F\to H\), \(Z\to G\), and
\(X_j\to\mathcal X_j\) along the further subsequence.
In particular, every subsequential six-energy law has such
a representation, with an a priori unspecified law for \(B\).
\end{theorem}

\begin{proof}
The first-moment bound for \(F\), the second-moment bound
for \(b\), and the first-moment bounds for \(X_j\) give
tightness of those finite-dimensional vectors. The
total-variation estimate gives tightness of \(Z\).
Thus the full finite-dimensional family is tight.

Because \((b,F)\) is exterior measurable,
\eqref{eq:v96-joint-TV} also bounds the distance between
the law of \((b,F,Z)\) and
\(\mathcal L(b,F)\otimes\mathcal G\) by
\(\epsilon_{M,\beta_M}\to0\).
Every joint subsequential limit therefore has the stated
Gaussian independence.

The exact relations \eqref{eq:v96-centres} give
\[
 \sum_jb_j=0,\qquad
 0\le F-\frac12\sum_j|b_j|^2
           =\frac1{2\beta_M}\sum_jd_j^2
           \le\frac{F^2}{2\beta_M}.
\]
Since \(F\) is tight and \(\beta_M\to\infty\), the last
quantity tends to zero in probability. This identifies \(H\)
and the zero-sum constraint. Similarly \(F+Q\) is tight by
\eqref{eq:v96-remainder}; hence the energy remainders tend to
zero in probability. The continuous mapping theorem applied
to \(\tfrac12|Z-b_j|^2\) proves the first identity in
\eqref{eq:v96-mixture}. Summing it and using \(\sum_jB_j=0\)
proves the second.
\end{proof}

The means \(\mathbb E|B_j|^2\) are finite by the lower
semicontinuity of second moments and \eqref{eq:v96-first-budgets}.
This is not a proof that the finite-volume second moments of
\(b_j\) converge. The law of \(B\) has not been classified.
The representation is also only for \emph{one chosen link's}
six-plaquette star at a time. It does not give independent
Gaussian variables for all links simultaneously.

\begin{proposition}[Explicit conditional laws and stronger local noncollapse]
\label{prop:v96-local-consequences}
Every one-star limit in Theorem~\ref{thm:v96-mixture} satisfies,
for \(t\ge0\),
\begin{equation}
 \mathbb E(e^{-t\mathcal X_j}\mid B)
    =(1+t/6)^{-3/2}
       \exp\left(-\frac{t|B_j|^2}{2(1+t/6)}\right).
 \label{eq:v96-Laplace}
\end{equation}
Consequently
\begin{align}
 \mathbb E(\mathcal X_j\mid B)
       &=\frac14+\frac12|B_j|^2,                       \label{eq:v96-limit-mean}\\
 \Var(\mathcal X_j\mid B)
       &=\frac1{24}+\frac16|B_j|^2,                    \label{eq:v96-limit-var}\\
 \operatorname{Cov}(\mathcal X_j,\mathcal X_k\mid B)
       &=\frac1{24}+\frac16 B_j\cdot B_k.              \label{eq:v96-limit-cov}
\end{align}
Every marginal \(\mathcal X_j\) is absolutely continuous on
\([0,\infty)\), has no atom at zero, and obeys
\begin{equation}
       \frac14\le\mathbb E\mathcal X_j\le K_*,
       \qquad \Var\mathcal X_j\ge\frac1{24},
 \label{eq:v96-noncollapse}
\end{equation}
where variance may be infinite.
\end{proposition}

\begin{proof}
Complete the square in the Gaussian integral for
\(\tfrac12|G-B_j|^2\) to obtain \eqref{eq:v96-Laplace}.
The mean, variance and covariance follow either by
differentiation or by expanding
\[
        \mathcal X_j=\tfrac12|G|^2-G\cdot B_j+\tfrac12|B_j|^2.
\]
The centred quadratic and linear Gaussian terms are
uncorrelated, \(\Var(|G|^2)=1/6\), and
\(\mathbb E(G\cdot b)(G\cdot c)=b\cdot c/6\).
These identities give \eqref{eq:v96-limit-mean}--%
\eqref{eq:v96-limit-cov}.

For any fixed \(b\), the pushforward
\(\tfrac12|G-b|^2\) has density on \(x>0\)
\[
 f_b(x)=\sqrt{2x}\int_{S^2}
          \left(\frac3\pi\right)^{3/2}
          e^{-3|b+\sqrt{2x}\,\omega|^2}\,d\omega,
\]
where \(d\omega\) is surface area measure. This follows
from polar coordinates centred at \(b\). The point \(G=b\)
has Gaussian probability zero. Averaging these densities
over the law of \(B_j\) and using Tonelli's theorem proves
absolute continuity and absence of atoms.

The upper first-moment bound passes to the weak limit by
bounded continuous truncation and monotone convergence,
as in V95. The lower bound follows from
\eqref{eq:v96-limit-mean}. The conditional variance identity
and the law of total variance imply the variance lower bound;
if the second moment is infinite the same inequality holds
in the extended sense.
\end{proof}

Thus the limiting law of one plaquette is a mixture of
noncentral squared-Gaussian laws, not necessarily a central
chi-square law. Conditional on \(B_j=b\), its distribution
is \(1/12\) times a noncentral chi-square distribution with
three degrees of freedom and noncentrality \(6|b|^2\).
This description gives a concrete unresolved variable rather
than silently setting the noncentrality to zero.

\subsection{Finite-volume lower moments without exchanging unbounded limits}

\begin{proposition}[A universal thermal mean floor]
\label{prop:v96-finite-mean}
For any \(M\ge3\), \(\beta>0\), plaquette \(p\), and an
edge \(e\subset p\), the full Wilson law satisfies
\begin{equation}
 \begin{gathered}
 \mathbb E(X_p\mid\xi)\ge
       \ell(\beta):=\beta\big(1-m(6\beta)\big)>0,\\
                         \ell(\beta)\longrightarrow\frac14
                         \quad(\beta\to\infty).
 \end{gathered}
 \label{eq:v96-finite-mean}
\end{equation}
Hence for arbitrary sequences \(\beta_M\to\infty\),
\[
       \liminf_M\mathbb EX_p\ge\frac14,\qquad
       \liminf_M\frac{\beta_M}{M^4}\mathbb ES\ge\frac32.
\]
The last statement does not require the growth condition
\eqref{eq:v95-growth-window}.
\end{proposition}

\begin{proof}
The conditional mean of \(u\cdot a_j\) is
\(c_jm(\beta\alpha)\). Since \(c_j\le1\), \(m\ge0\),
\(\alpha\le6\), and \(m\) is increasing,
\[
  \mathbb E(X_j\mid\xi)
      =\beta(1-c_jm(\beta\alpha))
      \ge\beta(1-m(\beta\alpha))
      \ge\beta(1-m(6\beta)).
\]
Strict positivity and the limit are
Lemma~\ref{lem:v96-radial-Gamma}, since
\(\beta(1-m(6\beta))=\tfrac16[6\beta(1-m(6\beta))]\).
Integrate the conditional inequality and sum over the
\(6M^4\) plaquettes.
\end{proof}

\begin{proposition}[A finite-volume variance lower limit]
\label{prop:v96-variance-liminf}
Under the growth condition used in
Theorem~\ref{thm:v96-mixture},
\begin{equation}
                  \liminf_M\Var X_p\ge\frac1{24}
 \label{eq:v96-variance-liminf}
\end{equation}
for every fixed plaquette \(p\).
\end{proposition}

\begin{proof}
If the left side is infinite, there is nothing to prove.
Otherwise choose a subsequence realizing its finite value.
The means lie in \([0,K_*]\), so pass to a further subsequence
where they converge to some \(c\). Choose an edge in \(p\)
and use Theorem~\ref{thm:v96-mixture} to extract a further
subsequence with \(X_p\Rightarrow\mathcal X_p\) of the
stated form. The nonnegative continuous function
\((x,a)\mapsto(x-a)^2\), applied by truncation or Portmanteau,
gives
\[
 \liminf_M\mathbb E(X_p-\mathbb EX_p)^2
        \ge\mathbb E(\mathcal X_p-c)^2
        \ge\Var\mathcal X_p\ge1/24.
\]
This argument does not assume that means or second moments
converge under weak convergence.
\end{proof}

These lower bounds sharpen the small universal constants
from V95. They do not give a uniform second-moment upper
bound for the finite-volume \(X_p\), nor do they prove
uniform integrability. The mean floor \(1/4\) per plaquette
is also below the \(3/4\) mean of the previously identified
selected chart limit. No equality of those laws follows.

\subsection{The remaining correlation problem}

The finite comparison \eqref{eq:v96-energy-experiment} makes
one part of the nonlinear local response explicit while keeping
the actual exterior. It would be incorrect to iterate the
comparison as if the independent Gaussian on its right could
be chosen independently for every overlapping star. Changing
a neighbouring link changes \(A\), the frame, and all relevant
centres \(b_j\). Compatibility of these responses across space
has not been proved.

In particular, the identities \(\sum_jB_j=0\) and
\(H=\tfrac12\sum_j|B_j|^2\) constrain the latent variables but
do not identify their spatial joint law. A one-link heat-bath
description, even with a Gaussian limit and local variance
lower bounds, supplies no volume-uniform mixing rate or
physical Hamiltonian gap.

The next upstream analytic objects are therefore the joint
staple/centre field and its response to neighbouring changes,
or source-local partition ratios that control those
correlations. V95's full pressure bound may also be used
to seek sharper energy information, but its coarse first
moment cannot by itself eliminate a latent collective mode.
Physical normalization, a nontrivial renormalized trajectory,
and reconstruction remain separate requirements. Every
limit here still concerns fixed microscopic lattice indices,
not a physically smeared continuum field.

\subsection{Exact checks and preservation}

The diagnostic certificate is
\begin{center}\small
\path{scripts/verify_ym_f14_v96_staple_innovations.py}.
\end{center}
It verifies 128 exact rational star configurations, including
both signs of the rotated first link coordinate, 768
energy-remainder identities, and 768 rotated dot-product
identities. It checks the noncentral Gaussian mean and
variance for 27 rational centres and all 729 corresponding
conditional covariance pairs. There are 164 exact density-chart
comparisons, the rational total-variation coefficient
\(539<600\), 129 bounded-Lipschitz truncation checks, and
108 exact Gaussian square completions for the Laplace transform.

The probability estimates and limit passages are proved above,
not inferred from these finite checks. In particular the
certificate is not a Monte Carlo simulation and is not evidence
of a physical gap.

The certificate contains 7,198 bytes and has SHA--256
\begin{center}\scriptsize\ttfamily
A57E3002179EC1FADA9707125ED0B3E7E1FD4349A35F897403EF95A9859225DF.
\end{center}
The predecessor V95 contains 2,340,397 bytes and has SHA--256
\begin{center}\scriptsize\ttfamily
13A64E274C828713799C8DCB663145EAD5330DEEFBE67C9D797C83E8B7D79423.
\end{center}
Its body is preserved apart from the title version and this
terminal insertion. V96 replaces V95 as the sole active main
note, with the predecessor recoverable by reversing those
two changes. Earlier certificates, all companion and supplied
files, monographs, and archives remain unchanged.

\begin{firewall}
V96 proves full-measure staple alignment, a quantitative
one-link Gaussian response, an exact energy-remainder identity,
and a local limiting mixture with explicit noncollapse.
It does not classify the interacting centre field, establish
physical-scale correlation decay, construct the interacting
four-dimensional continuum theory, or prove a positive
physical Hamiltonian mass gap. The original goal remains active.
\end{firewall}
\section{V97: full compact curvature, local identities, and sharp selection}
\label{sec:v97-full-curvature}

\subsection{Context and the precise upstream task}

V96 identifies the conditional fluctuations of a single full-Wilson
link, but does not make independent Gaussian choices on overlapping
stars. This section addresses compatibility directly. We transport
finitely many plaquette curvatures to a common root, prove that the
choice of fixed transport paths disappears in the microscopic
thermal limit, and use paths avoiding the varied edge to obtain
exact conditional Haar integration by parts.

The new input is a proof of the required local identities for
the \emph{full compact} periodic \(\SU(2)\) Wilson measure. We then
apply the abstract stationary-law classification already proved in
V90; we do not repeat its heat-semigroup and independence arguments.
A sharp selection using the full Wilson partition function removes
its possible constant curvature component. The selected couplings
and measure are distinct from the restricted-chart ones in V90.

The elementary use of Haar differentiation is standard. For
historical context see W.~Kerler, \emph{Schwinger--Dyson equations
and currents in lattice gauge theory}, Physics Letters B 100
(1981), 267--272,
\href{https://www.sciencedirect.com/science/article/abs/pii/0370269381903300}
{publisher record}; the corresponding
\href{https://www.slac.stanford.edu/pubs/slacpubs/2500/slac-pub-2646.pdf}
{SLAC-PUB-2646 preprint} dates from November 1980.
The transport estimates and limiting passages used here are
proved below rather than inferred from a formal continuum identity.

Throughout, \(M\to\infty\), \(\beta_M\to\infty\), and
\[
       \sup_M\frac{\log\beta_M}{M}=L<\infty.
\]
For a finite-volume estimate write
\[
       K=K_{M,\beta}=\frac{65}{3}+\frac{2\log\beta}{M},
       \qquad K_*=65/3+2L.
\]
We use \(\beta\ge2\), \(M\ge3\), and the full measure
\(\mu^{\mathrm W}_{M,\beta}\) from \eqref{eq:v95-full-law}.
Every fixed plaquette satisfies \(\mathbb EX_p\le K\).
No Faddeev--Popov factor, chart cutoff, or renewal profile is
inserted. The next companion checkpoint remains V100.

\subsection{Fixed loops and transport to one root}

We identify \(\SU(2)\) with unit quaternions, use their Euclidean
norm, and write \(\operatorname{Im}q\in\mathbb R^3\) for the
imaginary part. Products of unit quaternions preserve this norm.
Paths in \(\mathbb Z^4\) are projected to the periodic lattice;
all paths discussed at a given limiting step are fixed before
\(M,\beta\) are sent to infinity.

\begin{lemma}[A fixed-word filling bound]
\label{lem:v97-fixed-loop}
Let \(C\) be a fixed closed nearest-neighbour path in
\(\mathbb Z^4\). Choose a reduction of its word to the empty
word by square swaps and immediate backtrack cancellations,
and let \(A(C)\) count the square swaps. Then
\begin{equation}
        \mathbb E|U_C-I|^2\le\frac{2K A(C)^2}{\beta}.
 \label{eq:v97-loop-bound}
\end{equation}
Repeated plaquettes in the reduction are counted with
multiplicity. If there are no square swaps, \(U_C=I\).
\end{lemma}

\begin{proof}
Such a finite reduction exists: order the signed steps by
coordinate direction, exchanging adjacent steps in different
directions, then cancel the opposite steps within each
coordinate block. A closed word has zero exponent sum in
every direction. Each exchange is an elementary square,
possibly with reversed orientation; immediate backtracks
have holonomy \(I\).

As in the word proof of Lemma~\ref{lem:v91-factorization},
the ratio of successive path holonomies is a conjugate of
that square holonomy. Telescoping the successive ratios
writes \(U_C\) as an ordered product of \(A(C)\) conjugated
plaquette products or their inverses. Therefore
\[
 |U_C-I|
       \le\sum_{\ell=1}^{A(C)}|U_{p_\ell}-I|,\qquad
 |U_C-I|^2
       \le A(C)\sum_{\ell=1}^{A(C)}2s_{p_\ell}.
\]
Take expectations and use \(\mathbb Es_p\le K/\beta\).
The argument uses neither independence nor an embedded surface.
\end{proof}

Choose one finite path \(\gamma_x\) from the origin to each
\(x\in\mathbb Z^4\), with \(\gamma_0\) empty. These choices form
a countable deterministic family. Put
\begin{equation}
 \begin{gathered}
 g_x=U_{\gamma_x},\qquad
 k_{x,\mu}=g_xU_\mu(x)g_{x+\hat\mu}^{-1},\\
 \widetilde P_{\mu\nu}(x)=g_xP_{\mu\nu}(x)g_x^{-1},\qquad
 Y^\gamma_{\mu\nu}(x)=\sqrt\beta\,
                     \operatorname{Im}\widetilde P_{\mu\nu}(x).
 \end{gathered}
 \label{eq:v97-root-field}
\end{equation}
The \(g_x\)'s need not define a periodic frame on the covering
lattice. All identities in a fixed finite patch are consistent
because the same \(g_x\) is used at each lifted vertex.
The product of the four \(k\)'s around a plaquette equals
\(\widetilde P\) exactly.

Let \(A_e\) be a fixed filling count for
\(\gamma_x\,e\,\gamma_{x+\hat\mu}^{-1}\), where
\(e=(x,\mu)\). Lemma~\ref{lem:v97-fixed-loop} gives
\begin{equation}
       \mathbb E|k_e-I|^2\le\frac{2KA_e^2}{\beta}.
 \label{eq:v97-small-root-link}
\end{equation}
There is no claim that these constants are bounded uniformly
over all \(x\) or over physically growing paths.

\begin{lemma}[Path independence at fixed coordinates]
\label{lem:v97-path-independence}
The root-transported field satisfies
\begin{equation}
 |Y^\gamma_p|^2=2X_p-X_p^2/\beta\le2X_p,\qquad
                       \mathbb E|Y^\gamma_p|^2\le2K.
 \label{eq:v97-energy-chord}
\end{equation}
If \(\gamma_x\) is replaced by another fixed path
\(\widehat\gamma_x\), then
\begin{equation}
 \mathbb E|Y^{\widehat\gamma}_p-Y^\gamma_p|
           \le\frac{4K A(\widehat\gamma_x\gamma_x^{-1})}
                         {\sqrt\beta}
 \label{eq:v97-path-change}
\end{equation}
for a plaquette based at \(x\). Consequently fixed finite
collections have the same limiting laws for any such choices.
\end{lemma}

\begin{proof}
The first identity is
\(\beta(1-(\operatorname{Re}P)^2)
 =2\beta(1-\operatorname{Re}P)-\beta(1-\operatorname{Re}P)^2\).
Changing the path conjugates \(Y^\gamma_p\) by
\(H=g_x^{\,\widehat\gamma}(g_x^\gamma)^{-1}\).
For a unit quaternion \(H\) and imaginary vector \(v\),
\[
                     |\operatorname{Ad}_H v-v|
                                  \le2|H-I|\,|v|.
\]
Cauchy--Schwarz, \eqref{eq:v97-loop-bound}, and
\eqref{eq:v97-energy-chord} prove \eqref{eq:v97-path-change}.
Summing this bound over a fixed finite collection gives
convergence in probability of the difference of the vectors,
and hence equality of their subsequential limiting laws.
\end{proof}

\subsection{Closedness, colour invariance, and stationarity}

Let \(d_1\) be the usual oriented lattice curl from edge
cochains to plaquette cochains, and \(d_2\) the coboundary
from plaquettes to cubes. Thus \(d_2d_1=0\), with all
orientations and colours as in V89--V90.

\begin{lemma}[A quantitative local linearization]
\label{lem:v97-linearization}
Set \(a_e=\sqrt\beta\,\operatorname{Im}k_e\).
For a plaquette with its four boundary edges
\(e_1,\ldots,e_4\),
\begin{equation}
 \mathbb E|Y^\gamma_p-(d_1a)_p|
       \le\frac{2K}{\sqrt\beta}
                  \sum_{1\le i<j\le4}A_{e_i}A_{e_j}.
 \label{eq:v97-linearization}
\end{equation}
In particular \((d_2Y^\gamma)_c\to0\) in \(L^1\) on every
fixed cube \(c\).
\end{lemma}

\begin{proof}
Write the oriented factors in the plaquette product as
\(q_1,\ldots,q_4\), with inverse links where required.
Set \(\delta_i=q_i-I\). An exact telescoping identity gives
\[
 \prod_{i=1}^4q_i-I-\sum_{i=1}^4\delta_i
       =\sum_{j=2}^4(q_1\cdots q_{j-1}-I)\delta_j.
\]
The norm of the right side is at most
\(\sum_{i<j}|\delta_i||\delta_j|\).
For an inverse unit quaternion, the imaginary part changes
sign and its chord norm is unchanged. Taking imaginary
parts, multiplying by \(\sqrt\beta\), and applying
Cauchy--Schwarz with \eqref{eq:v97-small-root-link}
proves \eqref{eq:v97-linearization}.
Apply \(d_2\), sum the six fixed face errors, and use
\(d_2d_1a=0\). No logarithmic branch or global gauge slice
is involved.
\end{proof}

\begin{theorem}[Compatible full-measure curvature limits]
\label{thm:v97-compatible}
The joint laws of \((Y^\gamma,X)\) on the countable
plaquette index set are tight in the product topology.
Every subsequential limit \((Y,X)\) is translation invariant,
invariant under simultaneous \(\mathrm{SO}(3)\) colour rotations,
and satisfies almost surely
\begin{equation}
       d_2Y=0,\qquad X_p=\tfrac12|Y_p|^2,\qquad
       \mathbb EY_p=0,\qquad \sup_p\mathbb E|Y_p|^2\le2K_*.
 \label{eq:v97-compatible}
\end{equation}
\end{theorem}

\begin{proof}
The bounds \(\mathbb EX_p\le K_*\) and \(|Y^\gamma_p|^2\le2X_p\)
give countable-product tightness by the compact-box argument
in Theorem~\ref{thm:v95-local-limit}. At each coordinate,
\[
             X_p-\tfrac12|Y^\gamma_p|^2=X_p^2/(2\beta_M)
                                                  \longrightarrow0
\]
in probability because \(X_p\) is tight. The local linearization
gives \(d_2Y^\gamma\to0\) in probability. Bounded continuous
truncations of these discrepancies pass both identities to
every joint weak limit. Lower semicontinuity gives the
second-moment bound.

A gauge transformation at the root conjugates every based
loop by the same group element. Gauge invariance of the full
Wilson measure gives exact common colour-rotation invariance
at finite \(M\), and this passes to the limit. Its finite
first moment then implies \(\mathbb EY_p=0\).

Stationarity needs a separate argument: changing the root
introduces a \emph{random} colour rotation, which cannot be
dropped merely by citing deterministic colour invariance.
Fix \(z\ne0\) in \(\mathbb Z^4\), and take \(M\) large enough
that \(z\) and \(0\) are distinct torus vertices.
Let \(Z_x\) be the field rooted at \(z\), using the translated
paths \(z+\gamma_x\). Its joint law with the corresponding
energies is the law of \((Y^\gamma_x,X_x)\), by translation
invariance of the original link measure.
Let \(g_z\) be the holonomy of \(\gamma_z\).
The field \(\operatorname{Ad}_{g_z}Z_x\) uses root-zero paths
\(\gamma_z\) followed by \(z+\gamma_x\).
Lemma~\ref{lem:v97-path-independence} compares it in \(L^1\),
on each fixed patch, with \(Y^\gamma_{z+x}\).

For a bounded measurable test \(f\) of the \(Z\)'s and their
energies, make a gauge transformation \(h\) at vertex zero
alone. It leaves every loop rooted at \(z\) unchanged, while
sending \(g_z\) to \(hg_z\). Invariance of the measure and
Haar averaging therefore give
\[
 \begin{split}
 \mathbb E f(\operatorname{Ad}_{g_z}Z,X)
   &=\mathbb E\int f(\operatorname{Ad}_{hg_z}Z,X)\,dh\\
   &=\mathbb E\int f(\operatorname{Ad}_hZ,X)\,dh
    =\mathbb E f(Z,X).
 \end{split}
\]
The last equality uses common colour invariance at root \(z\);
the energies are gauge invariant throughout.
Thus the translated field differs, on fixed patches, from
a field with the original law by an error tending to zero.
Translation invariance follows in the limit. The case \(z=0\)
is immediate. Countably many plaquettes, cubes and lattice
translations give the stated simultaneous conclusions.
\end{proof}

\subsection{Edge-avoiding transports and exact Haar differentiation}

Fix one positive edge \(e\), a finite collection of tested
plaquettes, and the six plaquettes containing \(e\). The graph
\(\mathbb Z^4\) remains connected after that unoriented edge is
removed. We may therefore choose all root paths needed for
these plaquettes to avoid both orientations of \(e\).
They are finite paths. For all sufficiently large \(M\),
their projections also avoid that torus edge; no translated
copy can alias it in this fixed finite collection.

The resulting \(g_x\)'s are functions of the exterior links
alone. With \(e=(x,\mu)\) and \(h\in\mathbb R^3\), define the
conditional flow
\begin{equation}
 U_e(t)=g_x^{-1}\exp\big(t\,\iota(h)/\sqrt\beta\big)g_xU_e,
 \label{eq:v97-Haar-flow}
\end{equation}
where \(\iota(h)\) is the imaginary quaternion with vector \(h\).
Conditionally on all other links this is left Haar translation.
In the transported variables,
\[
             k_e(t)=\exp\big(t\,\iota(h)/\sqrt\beta\big)k_e,
\]
while every other relevant \(k_f\) stays fixed.
Write \(D_e\) for the derivative at \(t=0\).

Let \(\sigma_{p,e}\in\{-1,0,1\}\) be the boundary incidence.
For \(p\) containing \(e\), let \(C_{p,e}\) be its ordered
transported boundary prefix \emph{before} \(k_e\) if the
incidence is positive, and the prefix \emph{through}
\(k_e^{-1}\) if it is negative. Direct differentiation gives
\begin{equation}
 \begin{gathered}
 v_{p,e}=\sigma_{p,e}\operatorname{Ad}_{C_{p,e}}h,\qquad
 D_e\widetilde P_p=\iota(v_{p,e})\widetilde P_p/\sqrt\beta,\\
 D_e(\beta s_p)=v_{p,e}\cdot Y^\gamma_p,\qquad
 D_eY^\gamma_p=\operatorname{Im}
                            \big(\iota(v_{p,e})\widetilde P_p\big).
 \end{gathered}
 \label{eq:v97-exact-variation}
\end{equation}
All these derivatives vanish on a tested plaquette not
containing \(e\). For negative incidence, the formula uses
\(D_e(k_e^{-1})=-k_e^{-1}\iota(h)/\sqrt\beta\);
omitting the inverse factor from the prefix would be incorrect.

Put
\[
        A_\star=\sum_{p\ni e}\ \sum_{f\in\partial p}A_f,
\]
with edge multiplicities included. This is a fixed finite
constant for the chosen adapted paths.

\begin{proposition}[A quantitative full-measure local identity]
\label{prop:v97-finite-Stein}
Let \(f\) be a bounded \(C^1\) function of the tested curvature
coordinates with continuous bounded gradient. Set
\[
 B_f=\|f\|_\infty,\qquad
 L_f=\max_p\sup|\nabla_{Y_p}f|.
\]
Let \(h_e\) denote the edge cochain equal to \(h\) at \(e\)
and zero elsewhere, and put \(\ell=d_1h_e\). Then
\begin{equation}
 \left|\mathbb E D_\ell f(Y^\gamma)
         -\mathbb E f(Y^\gamma)\langle Y^\gamma,\ell\rangle\right|
 \le\frac{|h|A_\star}{\sqrt\beta}
                     \big(3L_f\sqrt{2K}+4B_fK\big).
 \label{eq:v97-finite-Stein}
\end{equation}
Here \(D_\ell\) is the ordinary directional derivative in
the finitely many tested curvature coordinates.
\end{proposition}

\begin{proof}
Conditioning on the exterior, Haar invariance of
\eqref{eq:v97-Haar-flow} gives the exact identity
\[
                  \mathbb ED_e f
                           =\mathbb E f\,D_e(\beta S).
\]
There is no divergence term: the \(g_x\) in this flow is
independent of the varied link.
The prefix \(C_{p,e}\) contains at most the four oriented
boundary links. Telescoping its chord and using the adjoint
bound from Lemma~\ref{lem:v97-path-independence} gives
\[
 |v_{p,e}-\sigma_{p,e}h|
              \le2|h|\sum_{f\in\partial p}|k_f-I|.
\]
Since \(|v_{p,e}|=|h|\) and
\(|\widetilde P_p-I|\le\sum_{f\in\partial p}|k_f-I|\),
\eqref{eq:v97-exact-variation} also gives
\[
 |D_eY^\gamma_p-\sigma_{p,e}h|
              \le3|h|\sum_{f\in\partial p}|k_f-I|.
\]
Thus \eqref{eq:v97-small-root-link} bounds the expectation
of the derivative error by
\[
               \frac{3|h|L_f\sqrt{2K}}{\sqrt\beta}A_\star.
\]
For the action error use Cauchy--Schwarz:
\[
 \mathbb E\big(|k_f-I|\,|Y^\gamma_p|\big)
             \le\frac{2KA_f}{\sqrt\beta}.
\]
The difference between \(D_e(\beta S)\) and
\(\langle Y^\gamma,\ell\rangle\) therefore has first moment
at most \(4K|h|A_\star/\sqrt\beta\).
Insert these two estimates into the exact integration-by-parts
identity and multiply the second one by \(B_f\).
\end{proof}

The path choice in this proposition is adapted to one edge.
One must not differentiate a frame that depends on that same
link and then omit its derivative. In the next step, path
independence is used \emph{after} taking the finite identity;
it is not used to justify differentiating a small error.

\subsection{The limiting identity and its measure consequence}

\begin{theorem}[Full-curvature Stein and translation identities]
\label{thm:v97-Stein-CM}
Every limit law \(\nu\) in Theorem~\ref{thm:v97-compatible}
satisfies
\begin{equation}
       \mathbb E_\nu D_{d_1h}f(Y)
              =\mathbb E_\nu f(Y)\langle Y,d_1h\rangle
 \label{eq:v97-Stein}
\end{equation}
for every finitely supported edge cochain \(h\) and bounded
\(C^1\) cylinder function with continuous bounded gradient.
For every \(\ell=d_1h\) and bounded Borel function \(F\) on
the countable field space, it also satisfies
\begin{equation}
 \mathbb E_\nu F(Y+\ell)
    =\mathbb E_\nu F(Y)
                \exp\big(\langle Y,\ell\rangle-\|\ell\|_2^2/2\big).
 \label{eq:v97-CM}
\end{equation}
\end{theorem}

\begin{proof}
First take \(h\) supported at one edge. The adapted path
fields have the same finite-coordinate weak limits as the
original path family, by \eqref{eq:v97-path-change}.
The right side of \eqref{eq:v97-finite-Stein} tends to zero.
The directional derivative of \(f\) is bounded continuous.
The other integrand is bounded by a fixed constant times
a finite sum of \(|Y_p|\). Uniform second moments make
these linear-growth integrands uniformly integrable.
Both expectations therefore pass to the weak limit,
proving \eqref{eq:v97-Stein} for one edge. Linearity gives
it for every finitely supported \(h\). Different single-edge
path choices have not changed the limiting law.

For the measure statement let \(\ell\ne0\) and put
\(\varphi(B)=\mathbb E_\nu e^{i\langle Y,B\rangle}\)
for finite real plaquette cochains \(B\).
First moments justify differentiation in the following
real parameter \(t\). Applying \eqref{eq:v97-Stein} to
the real and imaginary parts of the exponential yields
\[
 \frac{d}{dt}\varphi(B+t\ell)
                  =-\langle B+t\ell,\ell\rangle\varphi(B+t\ell).
\]
Solving this scalar linear equation gives
\begin{equation}
 \varphi(B+t\ell)
       =e^{-t\langle B,\ell\rangle-t^2\|\ell\|_2^2/2}\varphi(B).
 \label{eq:v97-characteristic-line}
\end{equation}
In particular \(S=\langle Y,\ell\rangle\) is
\(N(0,\|\ell\|_2^2)\). For any finite family \(B_j\), define
\[
 B_j^\perp=B_j-\frac{\langle B_j,\ell\rangle}{\|\ell\|_2^2}\ell.
\]
Equation~\eqref{eq:v97-characteristic-line}, applied to
all real combinations of the \(B_j^\perp\), proves that
\(S\) is independent of their joint score vector.
Multiplication by \(e^{S-\|\ell\|_2^2/2}\) shifts this
one-dimensional Gaussian by \(\|\ell\|_2^2\) and leaves
the perpendicular vector unchanged. It thus shifts
each \(\langle Y,B_j\rangle\) by \(\langle\ell,B_j\rangle\).
This proves \eqref{eq:v97-CM} on cylinder events.
The weight has expectation one, so both sides define
probability measures; the cylinder generating class and
the monotone-class theorem extend the identity to all
bounded Borel tests. The case \(\ell=0\) is immediate.
\end{proof}

\subsection{Applying the established classification to the full measure}

The hypotheses of Theorem~\ref{thm:v90-classification} were
stated there for any stationary law with uniform finite
second moments, common colour invariance, closedness, and
the translation identity for
\(\ell=d_1d_1^*B\), \(B\) finitely supported.
Theorems~\ref{thm:v97-compatible} and
\ref{thm:v97-Stein-CM} now establish all of them for the
full Wilson curvature limits. In the last condition take
\(h=d_1^*B\). No hypothesis about the earlier chart law
is being imported.

For clarity, on the Fourier torus let
\[
 q_i(k)=e^{ik_i}-1,\qquad
 \lambda(k)=\sum_{i=1}^4|q_i(k)|^2,\qquad
 P(k)=\frac{d_1(k)d_1(k)^*}{\lambda(k)}
                         \quad(k\ne0).
\]
The value at \(k=0\) is immaterial to Lebesgue integrals.
The operator \(P\) is the rank-three exact-two-form
projector, acting independently in the three colour coordinates.

\begin{theorem}[Full-measure microscopic classification]
\label{thm:v97-classification}
Every joint limit in Theorem~\ref{thm:v97-compatible}
has a decomposition
\begin{equation}
        Y=G+H,\qquad G\text{ independent of }H,\qquad
                         X_p=\tfrac12|G_p+H_p|^2.
 \label{eq:v97-classification}
\end{equation}
Here \(G\) is centred stationary Gaussian with covariance
\(P\otimes I_3\), and \(H=(H_{\mu\nu})_{\mu<\nu}\) is an
almost surely spatially constant random colour two-form
with finite second moments. Its mean is zero. In particular
\begin{equation}
 \frac12\sum_{\mu<\nu}\mathbb E|Y_{\mu\nu}(0)|^2
       =\frac92+\frac12\mathbb E
                             \sum_{\mu<\nu}|H_{\mu\nu}|^2 .
 \label{eq:v97-energy-floor}
\end{equation}
\end{theorem}

\begin{proof}
Apply Theorem~\ref{thm:v90-classification} with the newly
verified full-measure hypotheses. Its proof constructs \(H\)
as the invariant heat average of \(Y\), proves independence
using the measure translation identity, and identifies the
Gaussian covariance using the lattice Hodge identity.
Those arguments require no further measure-specific input
and are not repeated here.

The relation for \(X\) is \eqref{eq:v97-compatible}.
Common colour invariance gives zero mean for \(H\).
For one colour and orientation \(\mu\nu\),
\[
 P_{\mu\nu,\mu\nu}(k)
      =\frac{|q_\mu(k)|^2+|q_\nu(k)|^2}{\lambda(k)}.
\]
Permutation symmetry of the four momentum coordinates
makes its integral \(1/2\). Thus each plaquette vector of
\(G\) has mean-square norm \(3/2\). Independence and
\(\mathbb EG=0\) remove the cross term, giving
\eqref{eq:v97-energy-floor}.
\end{proof}

For any sequence in the stated growth window, this also gives
\begin{equation}
        \liminf_M\frac{\beta_M}{M^4}
                  \mathbb E^{\mathrm W}_{M,\beta_M}S\ge\frac92.
 \label{eq:v97-full-liminf}
\end{equation}
Indeed, take a subsequence realizing the liminf and a further
joint curvature-energy limit. The sum of the six origin
energies is nonnegative and has expectation
\(\beta_M\mathbb ES/M^4\). Lower semicontinuity and
\eqref{eq:v97-energy-floor} prove the inequality.
This improves the \(3/2\) lower bound from one-link conditioning
under the present additional joint-limit hypotheses.
It does not yet show that \(H=0\) for every such sequence.

\subsection{A sharp selection for the full Wilson normalizer}

Write
\[
 f_M(u)=-M^{-4}\log Z^{\mathrm W}_{M,u},\qquad
 q_M^{\mathrm W}(u)=u\,\mathbb E^{\mathrm W}_{M,u}S/M^4.
\]
These are functions of the full compact normalizer, not
the chart normalizer used in V90.
Set \(V=M^4\), \(D=3V+1\), \(N=3V-4M^3+1\), and retain
the constants \(c_b,C_1\) from \eqref{eq:v95-single-constants}.
The V95 pressure bounds imply, for \(U\ge L\ge1\),
\begin{equation}
 f_M(U)-f_M(L)
  \le C_S+\frac{3D}{2V}\log U-\frac{3N}{2V}\log L,
 \qquad C_S=48+4\log(c_b^{-1})+3\log C_1<62.
 \label{eq:v97-pressure-difference}
\end{equation}
The strict numerical bound follows from the V95 logarithmic
majorants:
\[
             C_S<48+4(5/2)+3(23/20)=1229/20<62.
\]

For \(M\) sufficiently large that \(r_M=\log\log M\ge4\), put
\begin{equation}
       L_M=e^{\sqrt{r_M}},\qquad U_M=\log M=e^{r_M}.
 \label{eq:v97-selection-window}
\end{equation}
Let \(\beta_M^{\mathrm W,*}\) be the leftmost minimizer of
the continuous function \(q_M^{\mathrm W}\) on this compact
interval. This is a well-defined existence selection from
the exact finite-volume model, not a numerical evaluation
of that normalizer.

\begin{theorem}[Sharp full-measure thermal normalization]
\label{thm:v97-selection}
The selected couplings satisfy
\begin{equation}
 \begin{gathered}
       e^{\sqrt{\log\log M}}\le\beta_M^{\mathrm W,*}\le\log M,\\
 q_M^{\mathrm W}(\beta_M^{\mathrm W,*})
                 \le\frac92+\eta_M,\qquad
 \eta_M=\frac{3}{2M^4}
       +\frac{62+6\sqrt{r_M}/M}{r_M-\sqrt{r_M}}
                                      \longrightarrow0.
 \end{gathered}
 \label{eq:v97-sharp-budget}
\end{equation}
Along the corresponding full compact measures, the joint
countable-coordinate laws converge to
\begin{equation}
                  (Y^\gamma,X)\Longrightarrow
                              \big(G,\tfrac12|G|^2\big),
       \qquad \operatorname{Cov}G=P\otimes I_3 .
 \label{eq:v97-selected-limit}
\end{equation}
The limit is independent of every fixed choice of root paths.
\end{theorem}

\begin{proof}
Differentiating the finite compact integral gives
\(u f_M'(u)=q_M^{\mathrm W}(u)\). Hence
\[
 \min_{[L_M,U_M]}q_M^{\mathrm W}
       \le\frac{f_M(U_M)-f_M(L_M)}
                          {\log U_M-\log L_M}.
\]
Use \eqref{eq:v97-pressure-difference}, noting that
\[
 \frac{3D}{2V}=\frac92+\frac{3}{2M^4},\qquad
 \frac{3N}{2V}=\frac92-\frac6M+\frac{3}{2M^4}.
\]
The numerator is bounded above by
\[
 \frac92(r_M-\sqrt{r_M})
 +\frac{3}{2M^4}(r_M-\sqrt{r_M})
 +\frac{6\sqrt{r_M}}{M}+62.
\]
Division proves \eqref{eq:v97-sharp-budget}; the periodic
boundary term \(6\sqrt{r_M}/M\) has not been omitted.

The selected couplings tend to infinity and satisfy the
growth window for all preceding full-measure results.
For any joint subsequential limit, nonnegativity gives
\[
 \frac12\sum_{\mu<\nu}\mathbb E|Y_{\mu\nu}(0)|^2
       \le\liminf_M q_M^{\mathrm W}(\beta_M^{\mathrm W,*})
       \le\frac92.
\]
The exact floor \eqref{eq:v97-energy-floor} forces
\(\mathbb E\sum_{\mu<\nu}|H_{\mu\nu}|^2=0\), so \(H=0\)
almost surely. Every subsequential joint limit is therefore
the law in \eqref{eq:v97-selected-limit}. Tightness and
uniqueness of all subsequential limits give convergence
of the whole sequence. Path independence was proved in
Lemma~\ref{lem:v97-path-independence}.
\end{proof}

\begin{proposition}[Full-measure first moments and fixed-coordinate covariance]
\label{prop:v97-moments}
Along this selected full-measure sequence,
\begin{equation}
 \begin{gathered}
 q_M^{\mathrm W}(\beta_M^{\mathrm W,*})\longrightarrow\frac92,
 \qquad \mathbb EX_p\longrightarrow\frac34,\qquad
 \mathbb E|Y^\gamma_p|^2\longrightarrow\frac32,\\
                  X_p\Longrightarrow\tfrac14\chi_3^2 .
 \end{gathered}
 \label{eq:v97-selected-moments}
\end{equation}
The families \(X_p\) and \(|Y^\gamma_p|^2\) are uniformly
integrable for each fixed \(p\). Every fixed pair of
colour-curvature coordinates has covariance converging to
the corresponding entry of \(P\otimes I_3\). Also
\begin{equation}
              \frac{\mathbb EX_p^2}{\beta_M^{\mathrm W,*}}
                                      \longrightarrow0 .
 \label{eq:v97-scaled-second}
\end{equation}
No uniform bound on \(\mathbb EX_p^2\) is asserted.
\end{proposition}

\begin{proof}
At the origin let \(Z_M=\sum_{\mu<\nu}X_{\mu\nu}(0)\).
Its expectation equals \(q_M^{\mathrm W}\).
Its weak limit has expectation \(9/2\), while
\eqref{eq:v97-sharp-budget} bounds its limsup by \(9/2\).
Thus \(\mathbb EZ_M\to9/2\).
The nonnegative truncation argument proved in
Theorem~\ref{thm:v90-moments} applies directly:
convergence of the means and of the bounded truncations
\(\min(Z_M,R)\) implies uniform integrability of \(Z_M\).
Every origin energy is bounded by \(Z_M\).
Exact translation and coordinate-permutation invariance of
the scalar full-Wilson energies extend this to any fixed
plaquette and identify its limiting mean as \(3/4\).

The pointwise inequality \(|Y^\gamma_p|^2\le2X_p\)
gives uniform integrability of squared curvature norms.
The joint weak convergence then gives their Gaussian
moments and all fixed-coordinate cross moments, since
a product is bounded by half the sum of the two squared
norms. Common colour invariance makes all the means zero.
The projector diagonal is \(I_3/2\), giving the stated
chi-square law. Finally, take expectations in the exact
identity
\[
                  X_p^2/\beta_M^{\mathrm W,*}
                                  =2X_p-|Y^\gamma_p|^2.
\]
The two terms on the right have the proved limits \(3/2\)
and \(3/2\). This yields \eqref{eq:v97-scaled-second},
not an unscaled second-moment bound.
\end{proof}

\subsection{What has been removed and what remains}

The gap between a chart-restricted local limit and a full
compact local limit has been removed for this selected
microscopic regime. All compact link configurations remained
in the measure. The proof did not assume that the torus
admits a typical globally small Coulomb representative,
and did not drop winding variables. Fixed closed paths
were controlled by finite fillings; this is not a claim
about noncontractible torus loops whose lengths grow with \(M\).

V96's local energy-mixture freedom is correspondingly
identified, at the level of the joint scalar energies and
the consistently transported curvature, along this selected
sequence. This does not require independently resampling
overlapping stars. For arbitrary couplings in the larger
growth window, the present classification still permits
the independent spatially constant \(H\).

The selection is not a renormalization recursion and does
not prove \(\beta_M^{\mathrm W,*}\sim c\log M\) for a
prescribed physical constant. Its fixed-coordinate limit
is a Gaussian lattice curvature field. It does not
construct the interacting continuum Yang--Mills theory.
In particular, the local error constants contain the
filling counts \(A_e\) and \(A_\star\); these grow when
transport paths and supports grow. The proof contains
no uniform estimate at fixed nonzero physical separation.

Even the fixed-coordinate covariance convergence in
Proposition~\ref{prop:v97-moments} gives no rate compatible
with the \(M^4\)-type amplification of dimension-two
physical smearings discussed in V92--V94. There is still
no justified passage from this microscopic Gaussian law
to a nontrivial interacting physical field, no proved
entry into the existing strong-mixing basin, and no
physical Hamiltonian mass gap.

The next work should therefore retain the full measure
and move to estimates with growing supports or fixed
physical normalization, including the quantitative
coupling/scale relation. Repeating the fixed-coordinate
Gaussian classification would not address that remaining
bottleneck.

\subsection{Exact diagnostics and preservation}

The certificate is
\begin{center}\small
\path{scripts/verify_ym_f14_v97_full_curvature.py}.
\end{center}
It constructs 625 finite root paths avoiding one specified
unoriented edge and checks 96 adapted-patch gauge covariance
identities. Sixteen closed-word reductions verify 65 exact
signed-square swap factors. There are 128 exact signed
left-variation identities, including the negative-incidence
prefix through the inverse link, and 32 quadratic word
remainder checks. Twelve closed edge-source stencils and
144 Gaussian perpendicular-decomposition identities audit
the linear algebra. Finally 147 rational coefficient
identities verify the sharp pressure-selection formula
with its periodic boundary term.

These are finite algebraic checks, not samples from the
full Wilson measure and not a machine proof of the limit.
The Haar identities, transport estimates, tightness,
stationarity, translation formula and limit passages are
proved analytically above. The already established
abstract classification is explicitly reused, not claimed
as a new independent calculation.

The certificate contains 8,953 bytes and has SHA--256
\begin{center}\scriptsize\ttfamily
1692F89534C6C38C761592A6F01260B7A4C957E9F25C94163BA904A1F294C6B4.
\end{center}
The predecessor V96 contains 2,367,700 bytes and has SHA--256
\begin{center}\scriptsize\ttfamily
6AF817D587F66013D32A9161E7B19B641DBE5E042E9F851A5B526A59547AEA47.
\end{center}
Its body is preserved apart from the title version and this
terminal insertion. V97 replaces V96 as the sole active
main note; reversing those two changes recovers the
predecessor. Earlier certificates, companion and supplied
files, monographs, and archives remain unchanged.

\begin{firewall}
V97 proves compatible full compact microscopic curvature
limits, their local translation identities and stationary
classification, and a sharp selected full-measure Gaussian
normalization. The interacting four-dimensional continuum
construction and a positive physical Hamiltonian mass gap
remain unproved. The original goal remains active.
\end{firewall}
\section{V98: thermal moments and polynomial Haar identities}
\label{sec:v98-thermal-exponential}

\subsection{The missing moment interface and the nonduplication audit}

V97 identifies a selected microscopic Gaussian curvature limit
under the full compact Wilson measure. Its final moment statement
does not bound the unscaled second moment of \(X_p=\beta s_p\),
let alone higher nonlinear insertions. This section supplies
exponential integrability and then tracks how the resulting
polynomial identity bounds depend on a growing support. It
does not repeat the Gaussian classification or turn that
microscopic limit into an interacting continuum theory.

The reflection method is inherited, not a new invention.
The earlier archives f6 V28/V41 and f12 V17--V19 already
establish bare-Wilson chessboard defect estimates; f7 V69
and f6 V86 also develop reflection-positive defect bounds.
Those estimates retain a polynomial factor in \(\beta\)
when the disseminated numerator is bounded by an absolute
fixed-defect action cost. Repeating that calculation does
not give the uniform thermal-scale exponential bound below.
The new use of V95 is to compare two partition functions
with matching extensive powers of \(\beta\). Only their
explicit periodic boundary loss remains.

For literature context, the multireflection and generating-function
method is also used by P.~A.~Faria da Veiga and M.~O'Carroll in
\emph{On Yang--Mills Stability Bounds and Plaquette Field Generating
Function}, \href{https://arxiv.org/abs/2205.07376v2}
{arXiv:2205.07376v2}, Sections V--VI. Their displayed linear
trace observable involves \(\operatorname{Im}\Tr(U_p-I)\);
that observable vanishes identically for \(\SU(2)\).
Their Remark 16 also discusses an action-type observable.
We do not import either statement as a proof for the present
normalization: the following \(\SU(2)\) energy estimate is
proved directly. No novelty relative to that literature is claimed.

The source audit remains the V95 checkpoint. The next
scheduled companion-file audit is V100. The sources continue
to impose the distinction between auxiliary/local bounds
and a physically normalized interacting construction.

\subsection{The pressure ratio and the exact site-reflection input}

Throughout this section the measure is precisely
\(\mu^{\mathrm W}_{M,\beta}\) in \eqref{eq:v95-full-law}.
All link configurations are retained. For reflection replication
we restrict to the cofinal family
\[
                         M=2^n,\qquad n\ge2.
\]
Put \(V=M^4\). We use \(N_0=3V-4M^3+1\) for the number
of triangular plaquettes in V95, to distinguish it from
the number of reflection copies below.

\begin{lemma}[Thermal pressure-ratio cancellation]
\label{lem:v98-pressure-ratio}
If \(0<t<1\) and \((1-t)\beta\ge1\), then
\begin{equation}
 \frac1V\log\frac{Z^{\mathrm W}_{M,(1-t)\beta}}
                       {Z^{\mathrm W}_{M,\beta}}
 \le62+\frac92\log\frac1{1-t}+\frac{6\log\beta}{M}.
 \label{eq:v98-pressure-ratio}
\end{equation}
This estimate itself holds for every integer \(M\ge3\).
\end{lemma}

\begin{proof}
The upper and lower estimates of Theorem~\ref{thm:v95-pressure}
give, with \(D=3V+1\),
\begin{align*}
 \frac1V\log\frac{Z^{\mathrm W}_{M,(1-t)\beta}}
                       {Z^{\mathrm W}_{M,\beta}}
 &\le48+\frac{D}{V}\log(c_b^{-1})
       +\frac{N_0}{V}\log C_1
       +\frac{3N_0}{2V}\log\frac1{1-t}
       +\frac{3(D-N_0)}{2V}\log\beta .
\end{align*}
Here \(D/V\le4\), \(0<N_0/V<3\), and
\(3(D-N_0)/(2V)=6/M\). The elementary logarithm bounds
already verified in V95 give
\[
 48+4\log(c_b^{-1})+3\log C_1
 <48+4(5/2)+3(23/20)=1229/20<62 .
\]
This proves the stated inequality. In particular the
extensive \(\log\beta\) terms cancel; dropping the
coupling dependence of the numerator would lose this fact.
\end{proof}

For completeness, only the following elementary site-reflection
case of the archived reflection framework is needed. Fix a
coordinate reflection through the two site planes
\(x_i=c\) and \(x_i=c+M/2\), and denote the induced pullback
on link functions by \(\theta\). Reversed oriented links
are inverted. Let \(E_0\) consist of tangent links in
the two planes and let \(E_+\), \(E_-\) be the remaining
links in the two closed halves. These three sets are disjoint.
Each elementary plaquette has vertices in one closed half
or entirely in a boundary plane: a unit plaquette cannot
meet both open halves. Consequently
\[
                        S=S_++\theta S_++S_0.
\]
Boundary links are fixed by reflection, while reflection
is a Haar-preserving bijection between the other two sets.
For a complex function \(F\) of \(E_+\cup E_0\), conditional
factorization therefore yields
\begin{equation}
 \mathbb E^{\mathrm W}_{M,\beta}[F\,\theta\overline F]
 =\frac1{Z^{\mathrm W}_{M,\beta}}
    \int e^{-\beta S_0}
       \left|\int e^{-\beta S_+}F\prod_{e\in E_+}dU_e\right|^2
                              \prod_{e\in E_0}dU_e\ge0 .
 \label{eq:v98-site-RP}
\end{equation}
This argument uses site planes, not a claim that crossing
bond-plane plaquettes factor without a character expansion.
The Cauchy--Schwarz inequality for this positive
sesquilinear form, applied to \(F\) and \(1\), gives
\begin{equation}
                  |\mathbb EF|^2\le\mathbb E[F\,\theta\overline F].
 \label{eq:v98-one-reflection}
\end{equation}
All functions used below are bounded at each finite \(M,\beta\);
no integrability extension is needed at this stage.

\subsection{A full-measure thermal exponential bound}

\begin{theorem}[Uniform exponential energy moments on dyadic tori]
\label{thm:v98-exponential}
For \(M=2^n\ge4\), \(0<t<1\), and \((1-t)\beta\ge1\),
every plaquette obeys
\begin{equation}
 \mathbb E^{\mathrm W}_{M,\beta}e^{tX_p}
 \le
 \exp\left\{256\left[
       62+\frac92\log\frac1{1-t}+\frac{6\log\beta}{M}\right]\right\}.
 \label{eq:v98-general-exponential}
\end{equation}
In particular, for \(\beta\ge2\) and
\((\log\beta)/M\le L<\infty\), define
\begin{equation}
 C_L:=\exp\left\{256\left[62+\frac92\log2+6L\right]\right\}.
 \label{eq:v98-CL}
\end{equation}
Then, independently of the plaquette,
\begin{equation}
                   \mathbb E e^{X_p/2}\le C_L .
 \label{eq:v98-half-exponential}
\end{equation}
The large constant is not optimized.
\end{theorem}

\begin{proof}
Translation invariance lets us place the plaquette at
base \((1,1,1,1)\) in the open cube \((0,4)^4\), with
any fixed orientation. Put \(F=e^{tX_p}\).
For the first coordinate, reflect successively through
the planes with positions \(4,8,\ldots,M/2\), each time
including the reflected copy of the entire pattern already
constructed. Before the reflection at \(a\), that pattern
has first-coordinate support strictly between \(0\) and \(a\).
It lies in the closed half \([a-M/2,a]\); its reflection
lies strictly between \(a\) and \(2a\). Thus
\eqref{eq:v98-one-reflection} applies at every step and
duplicates the number of factors without overlaps.
Do the same in the other three coordinates.

There are
\[
                         N_{\rm ref}=(M/4)^4
\]
final copies, exactly one in each open side-four block.
Their plaquettes are distinct. Trace and hence \(s_p\)
are unchanged by orientation reversal or by a change
of the starting corner of a plaquette. Repeated square
roots in \eqref{eq:v98-one-reflection} give
\[
 \mathbb Ee^{tX_p}
 \le\left(\mathbb E
          e^{t\beta\sum_{q\in{\cal D}(p)}s_q}\right)^{1/N_{\rm ref}},
\]
where \({\cal D}(p)\) is the set of reflected plaquettes.
For \(M=4\) this is simply the identity before the next
inequality, with \(N_{\rm ref}=1\). Since all actions
are nonnegative and the selected plaquettes are distinct,
\[
 \mathbb E e^{t\beta\sum_{{\cal D}(p)}s_q}
 \le \mathbb Ee^{t\beta S}
 =\frac{Z^{\mathrm W}_{M,(1-t)\beta}}
             {Z^{\mathrm W}_{M,\beta}}.
\]
Apply Lemma~\ref{lem:v98-pressure-ratio} and
\(V/N_{\rm ref}=256\). Choosing \(t=1/2\) proves the
uniform statement. No independence of plaquette
energies was asserted or used.
\end{proof}

As immediate consequences, for every integer \(r\ge0\),
\begin{equation}
 \mathbb EX_p^r\le C_L\,2^r r!,\qquad
 \mathbb P(X_p>R)\le C_Le^{-R/2}\quad(R\ge0).
 \label{eq:v98-energy-moments}
\end{equation}
The first inequality follows from
\((X_p/2)^r/r!\le e^{X_p/2}\). In the transported
curvature notation of V97,
\[
 |Y^\gamma_p|^2=2X_p-X_p^2/\beta\le2X_p,
\]
so, for any choice of root paths,
\begin{equation}
 \mathbb E e^{|Y^\gamma_p|^2/4}\le C_L,\qquad
 \mathbb E|Y^\gamma_p|^{2r}\le C_L\,4^r r! .
 \label{eq:v98-curvature-moments}
\end{equation}
These are full compact-measure bounds, not estimates
conditioned on a small-field event.

\subsection{Support-uniform source bounds and their actual normalization}

\begin{proposition}[An energy \(\ell^1\) source ball and arbitrary linear sources]
\label{prop:v98-sources}
Take any finite set of distinct plaquettes \(p_1,\ldots,p_m\).
Let \(z_j\in\mathbb C\), \(J_j\in\mathbb C^3\), and set
\[
 q=\sum_{j=1}^m|z_j|<\tfrac12,\qquad
 A=\sum_{j=1}^m|J_j|.
\]
For arbitrary root paths, define the finite-volume function
\begin{equation}
 {\cal Z}_{M,\beta}(z,J)
  =\mathbb E\exp\left(\sum_jz_jX_{p_j}
                         +\sum_jJ_j\cdot Y^\gamma_{p_j}\right).
 \label{eq:v98-local-source}
\end{equation}
Here the dot product is bilinear, not Hermitian; absolute
values of complex vectors use the Euclidean norm.
Under the hypotheses of \eqref{eq:v98-half-exponential},
\begin{equation}
             |{\cal Z}_{M,\beta}(z,J)|
               \le C_L\exp\left(\frac{A^2}{1-2q}\right).
 \label{eq:v98-source-bound}
\end{equation}
The bound is independent of \(m\), the plaquette positions,
and the root paths when \(q,A\) are fixed.
\end{proposition}

\begin{proof}
First, for nonnegative \(a_j\) with \(\sum_ja_j\le1/2\),
convexity of the exponential gives pointwise
\[
 e^{\sum_ja_jX_{p_j}}
 \le(1-2\sum_ja_j)+\sum_j2a_j e^{X_{p_j}/2}.
\]
Its expectation is at most \(C_L\), since \(C_L\ge1\).
This also allows arbitrary repetitions of an energy
after combining its coefficients.

If \(A=0\), apply this observation to \(a_j=|z_j|\).
Otherwise put \(\epsilon=1/2-q>0\) and
\(u_j=\epsilon |J_j|/A\). For \(J_j\ne0\), the identity
\(|Y_{p_j}^\gamma|\le\sqrt{2X_{p_j}}\) and a scalar
square completion give
\[
 |J_j||Y^\gamma_{p_j}|
 \le u_jX_{p_j}+\frac{|J_j|^2}{2u_j}.
\]
Terms with \(J_j=0\) are omitted from this inequality.
The total energy coefficient is \(q+\sum_ju_j=1/2\),
and the sum of the constant terms is
\(A^2/(2\epsilon)\). Taking absolute values inside the
expectation proves \eqref{eq:v98-source-bound}.
\end{proof}

There is also a volume-independent complex zero-free
neighborhood for energy sources, although its radius here
is very small. If \(J=0\) and \(q\le1/(8C_L)\), then
\begin{equation}
                    |{\cal Z}_{M,\beta}(z,0)-1|\le\tfrac13.
 \label{eq:v98-zero-free}
\end{equation}
Indeed, for \(q>0\), put
\(T=\sum_j(|z_j|/q)X_{p_j}\), so
\(\mathbb Ee^{T/2}\le C_L\). With
\(\delta=1/2-q\), use
\[
 |e^{\sum_jz_jX_{p_j}}-1|
 \le qT e^{qT},\qquad T\le\delta^{-1}e^{\delta T}.
\]
The expectation is at most \(qC_L/(1/2-q)\le1/3\).
The case \(q=0\) is immediate. Thus the logarithm
normalized to zero at \(z=0\) is well defined in this
source ball. This is a local generating-function statement,
not reflection positivity of a complex-tilted measure.

Every fixed polynomial in the \(X_{p_j},Y^\gamma_{p_j}\)
has uniformly bounded moments of every fixed order:
use \(\mathbb E\exp[(2m)^{-1}\sum_jX_{p_j}]\le C_L\)
and \(|Y_{p_j}|^2\le2X_{p_j}\). This proves the
higher uniform integrability missing from V97.
For supports whose size grows, the explicitly controlled
source quantities are \(q\) and \(A\); the result does
not replace an \(\ell^1\) norm by an \(\ell^2\) norm.

\subsection{All fixed moments and the joint analytic Gaussian limit}

Now restrict the selected sequence
\(\beta_M^{\mathrm W,*}\) of Theorem~\ref{thm:v97-selection}
to dyadic \(M\). Its defining window has
\(\beta_M^{\mathrm W,*}\le\log M\), so \(L=1\)
is a valid common bound for all its sufficiently large terms.
The selection and the full-measure convergence of V97
remain valid on this subsequence.

\begin{theorem}[Polynomial and locally uniform analytic convergence]
\label{thm:v98-analytic-limit}
For every fixed finite plaquette set and fixed finite
root paths, along the selected dyadic sequence,
\[
       (Y^\gamma_{p_j},X_{p_j})_{j=1}^m
          \ \Longrightarrow\ (G_{p_j},|G_{p_j}|^2/2)_{j=1}^m
\]
with convergence of every mixed polynomial moment.
Moreover \({\cal Z}_{M,\beta_M^{\mathrm W,*}}(z,J)\)
converges locally uniformly in
\begin{equation}
 {\cal U}_m=\{(z,J)\in\mathbb C^m\times\mathbb C^{3m}:
                                  \sum_j|z_j|<1/2\}.
 \label{eq:v98-analytic-domain}
\end{equation}
All its complex derivatives converge locally uniformly there.

Here \(G\) is the V97 Gaussian curvature with covariance
\(P\otimes I_3\). On the chosen coordinates let \(C\)
be its \(3m\)-dimensional covariance matrix, put
\[
 D_z=\operatorname{diag}(z_1I_3,\ldots,z_mI_3),\qquad
 B_z=C^{1/2}D_zC^{1/2},\qquad
 K_z=C^{1/2}(I-B_z)^{-1}C^{1/2}.
\]
The limiting generating function is
\begin{equation}
 {\cal Z}_\infty(z,J)
   =\det(I-B_z)^{-1/2}
                      \exp\big(\tfrac12J^{\mathsf T}K_zJ\big),
 \label{eq:v98-determinant-limit}
\end{equation}
with the determinant branch equal to one at \(z=0\).
Degeneracy of \(C\) causes no difficulty.
\end{theorem}

\begin{proof}
V97 gives the displayed weak convergence.
The exponential bounds just proved give uniform
integrability of every fixed polynomial and its square;
truncation therefore upgrades weak convergence to moment
convergence.

At any fixed point of \({\cal U}_m\), choose \(a>1\)
so close to one that \(a\sum_j|z_j|<1/2\).
Proposition~\ref{prop:v98-sources}, applied to absolute
values with strengths multiplied by \(a\), bounds the
\(a\)-th absolute moment of the exponential integrand
uniformly in \(M\). It is therefore uniformly integrable.
Its expectation converges by weak convergence and truncation.
Each finite-volume generating function is entire, since
its random variables are bounded at fixed \(M,\beta\).
The source bound is locally uniform on \({\cal U}_m\).
Cauchy estimates on slightly larger compact polydiscs
give equicontinuity and normality there. Every subsequential
analytic limit equals the already identified pointwise
limit. This proves local uniform convergence of the full
sequence; the Cauchy formula also gives convergence of
all derivatives.

For the formula, \(0\le C\le I\), because \(C\) is a
finite-coordinate compression of the orthogonal covariance
projection \(P\otimes I_3\). Thus
\(\|B_z\|\le\max_j|z_j|<1/2\). Write \(G=C^{1/2}\xi\),
with \(\xi\) a standard real \(3m\)-dimensional Gaussian.
For real sources, completion of the square in its
ordinary Gaussian integral gives
\eqref{eq:v98-determinant-limit}. Analytic continuation
gives the same identity for the complex sources in
\({\cal U}_m\). The analytic logarithm of \(I-B_z\),
defined by its norm-convergent power series, fixes the
determinant branch. The representation uses no inverse of
\(C\), so it also proves the degenerate case.
\end{proof}

Some newly justified unscaled limits are
\begin{align}
 \mathbb EX_p^2&\longrightarrow15/16,&
 \Var(X_p)&\longrightarrow3/8,                         \label{eq:v98-second-energy}\\
 \operatorname{Cov}(X_p,X_q)
     &\longrightarrow\tfrac32 P_{pq}^{\,2}.           \label{eq:v98-energy-covariance}
\end{align}
Indeed \(G_p\) has three independent components of variance
\(1/2\), so \(X_p\) tends to \(\chi_3^2/4\); Wick's formula
gives the cross covariance. More generally, the limiting
one-plaquette energy cumulant of order \(r\ge1\) is
\begin{equation}
                         \kappa_r(X_p)
                         \longrightarrow
                      \frac{3(r-1)!}{2^{r+1}}.
 \label{eq:v98-energy-cumulants}
\end{equation}
The zero-free estimate \eqref{eq:v98-zero-free} and local
uniform convergence justify differentiating the corresponding
logarithms, not merely the moment generating functions.
Thus all fixed-support connected polynomial coefficients
are controlled. This is not a decay estimate as the
separation or number of insertions grows.

\subsection{Polynomial Haar identities with an explicit growing-support cost}

The following consequence addresses a limited growing-patch
interface without assuming the selected Gaussian limit.
It holds for every dyadic \(M,\beta\) in the exponential
window of Theorem~\ref{thm:v98-exponential}.

Use V97's root-path construction adapted to a single
varied unoriented edge \(e\): every root path needed for
the tested plaquettes and the star of \(e\) avoids that
edge. For each boundary edge \(f\) of that star, let \(A_f\)
be the number of faces, counted with multiplicity, in
a specified finite filling of its rooted closed word.
The chosen paths and fillings must be valid on the
regulator; fixed lattice fillings embed for sufficiently
large \(M\). For growing paths this is a hypothesis to
check, not a consequence of their lengths alone. Put
\begin{equation}
 T_f=\sqrt\beta\,|k_f-I|,\qquad
 T_\star=\sum_{p\ni e}\sum_{f\in\partial p}T_f,\qquad
 A_\star=\sum_{p\ni e}\sum_{f\in\partial p}A_f .
 \label{eq:v98-transport-cost}
\end{equation}

The deterministic filling inequality used in V97 gives
\[
 T_f^2\le2A_f\sum_{a=1}^{A_f}X_{p_a}.
\]
Repeated faces are allowed. For \(A_f=0\) the rooted
word freely reduces to the identity and \(T_f=0\).
Convexity and \eqref{eq:v98-energy-moments} show, for
every integer \(r\ge1\),
\begin{equation}
 \mathbb ET_f^{2r}\le C_L\,4^r r!\,A_f^{2r},\qquad
                  \|T_\star\|_2\le2\sqrt{C_L}\,A_\star .
 \label{eq:v98-transport-moments}
\end{equation}
For example, the first inequality follows by applying
\((\sum_{a=1}^{A_f}x_a)^r\le
A_f^{r-1}\sum_ax_a^r\), and the second follows by
the triangle inequality in \(L^2\).

Let \({\cal I}\) be the tested plaquette set together with
the at most six plaquettes in the star, let \(m=|{\cal I}|\),
and put \(Q=1+\sum_{p\in{\cal I}}|Y^\gamma_p|\).
Define
\begin{equation}
 B_0(m)=1,\qquad
 B_k(m)=\left[1+m(C_L4^k k!)^{1/(2k)}\right]^k
                                    \quad(k\ge1).
 \label{eq:v98-Bk}
\end{equation}
Minkowski's inequality and \eqref{eq:v98-curvature-moments}
give \(\|Q^k\|_2\le B_k(m)\).

\begin{proposition}[Polynomial full-measure Stein error]
\label{prop:v98-polynomial-Stein}
Let \(f\) be a polynomial of total degree at most \(d\)
in the real tested curvature coordinates. Let \(a_f\)
be the sum of the absolute values of its monomial
coefficients, including its constant coefficient.
For \(\ell=d_1h_e\), with \(h_e=h\) on \(e\) and zero
elsewhere, the adapted root fields satisfy
\begin{align}
 &\left|\mathbb E D_\ell f(Y^\gamma)
      -\mathbb E f(Y^\gamma)\langle Y^\gamma,\ell\rangle\right|
 \nonumber\\
 &\hspace{8mm}\le
 \frac{2\sqrt{C_L}\,|h|\,a_f A_\star}{\sqrt\beta}
           \left[3d B_{d-1}(m)+2B_{d+1}(m)\right].
 \label{eq:v98-polynomial-Stein}
\end{align}
For \(d=0\) the term \(dB_{d-1}\) is defined to be zero.
The constants do not depend on a selected coupling.
\end{proposition}

\begin{proof}
At each finite regulator the polynomial and its derivatives
on the compact link space are bounded. The exact
conditional Haar flow of \eqref{eq:v97-Haar-flow}
therefore gives
\[
                         \mathbb ED_e f
                              =\mathbb E fD_e(\beta S).
\]
Avoidance of the varied edge by every relevant root path
is essential for this identity; no frame derivative has
been omitted.

The pointwise estimates in the proof of
Proposition~\ref{prop:v97-finite-Stein} are
\[
 |D_eY^\gamma_p-\sigma_{p,e}h|
       \le\frac{3|h|}{\sqrt\beta}
                                \sum_{f\in\partial p}T_f,\qquad
 |v_{p,e}-\sigma_{p,e}h|
       \le\frac{2|h|}{\sqrt\beta}
                                \sum_{f\in\partial p}T_f .
\]
A monomial estimate gives
\[
 |f|\le a_fQ^d,\qquad
 |\nabla_{Y_p}f|\le d a_fQ^{d-1}.
\]
For the latter, bound the Euclidean gradient by its
coordinate \(\ell^1\) norm and sum the exponents of
each monomial. Subtracting the desired ordinary
directional derivative and linear score from the exact
Haar identity now bounds its error by
\[
 \frac{|h|a_f}{\sqrt\beta}\,
       \mathbb E T_\star\left(3dQ^{d-1}+2Q^{d+1}\right).
\]
The action term uses \(|Y_p|\le Q\), including every
plaquette of the star even when it is not tested by \(f\).
Cauchy--Schwarz, \eqref{eq:v98-transport-moments}, and
\eqref{eq:v98-Bk} prove the proposition.
\end{proof}

For fixed \(d,L\), bounded \(a_f,|h|\), and a family of
valid adapted patches, this error tends to zero if
\begin{equation}
             \frac{A_\star(1+m)^{d+1}}{\sqrt\beta}
                              \longrightarrow0 .
 \label{eq:v98-growing-patch-condition}
\end{equation}
This is an explicit, albeit restrictive, growing-support
identity criterion. It does not by itself prove Gaussian
approximation of all observables on the growing patch:
that would additionally require quantitative identification
and control of the harmonic component. It also does not
justify replacing these adapted fields by a different
growing family of root frames without estimating that
replacement.

\subsection{The remaining physical-scale obstruction}

The exponential moment estimate controls rare large
thermal energies without a chart cutoff. The analytic
limit closes the fixed-degree nonlinear moment interface,
and \eqref{eq:v98-growing-patch-condition} records the
cost of some increasing supports. These are useful
upstream inputs, not a closure of the RG programme.

There are several precise limits on the conclusions.
First, \eqref{eq:v98-source-bound} depends on the absolute
source sums. Dimension-two physical smearings can have
an \(\ell^1\) source norm of order \(M^2\); this bound then
allows an exponent of order \(M^4\). It does not exploit
the cancellations required for the physical normalization.
Second, \eqref{eq:v98-growing-patch-condition} retains
the transport filling cost and is far short of uniform
fixed-physical-separation estimates when
\(\beta\) is at most logarithmic in \(M\).

Third, all reflection estimates here apply to the bare
Wilson law. They do not establish reflection positivity,
stability of conditional source derivatives, or a mixing
constant for a response-generated blocked law. Fourth,
exponential control of a quadratic curvature energy does
not permit a positive uncapped quartic Gaussian
exponential. The obstruction already recorded in V25
must not be bypassed by calling polynomial moment bounds
an exponential quartic bound.

Finally, a union bound gives only
\[
 \mathbb P\!\left(\max_pX_p>R\right)
                               \le6M^4C_Le^{-R/2}.
\]
At \(R\) proportional to \(\log M\), this does not make
\(\max_ps_p= \max_pX_p/\beta\) tend to zero along V97's
selected window. Thus the new estimates do not establish
global small-chart coverage either.

The next substantive target is an estimate exploiting
source cancellations or a genuinely local response norm
on growing supports, with the coupling/physical-scale
relation retained. Repeating another fixed-coordinate
moment calculation would not address that target.

\subsection{Exact diagnostics and preservation}

The certificate is
\begin{center}\small
\path{scripts/verify_ym_f14_v98_thermal_exponential.py}.
\end{center}
It checks 208,896 plaquette/site-reflection classifications
on periodic regulators, including reflection involution
and the exchange of the two open halves. There are
144 dissemination rounds producing 26,214 distinct final
faces across all six orientations and four dyadic sizes.
Each face lies strictly inside its own side-four block;
the copy count and the factor \(V/N_{\rm ref}=256\) are
checked exactly. The script also checks 97 periodic
pressure coefficient identities, seven exact source
allocations, twelve marginal moment/cumulant orders,
five Gaussian energy cross covariances, and 45 rational
determinants.

These finite checks are not Wilson-measure simulations
or a substitute for the analytic proofs. In particular
they do not prove reflection positivity or a thermodynamic
limit numerically. Those arguments are given above.

The certificate contains 7,137 bytes and has SHA--256
\begin{center}\scriptsize\ttfamily
AF4DCA25F435C8D5B9B0554052B04DDE222BDB32C0A78F484E1F7C726D6F28B4.
\end{center}
The predecessor V97 contains 2,399,314 bytes and has SHA--256
\begin{center}\scriptsize\ttfamily
9025CE717026115D30CAC81DC4612D8F52BFEE09F6C69B68590232F52B5DE8FE.
\end{center}
Its body is preserved apart from the title version and
this terminal insertion. V98 replaces V97 as the sole
active main note; reversing those two changes recovers
the predecessor. Earlier certificates and all protected
companion, supplied, monograph, and archive files are
unchanged. The next scheduled companion audit is V100.

\begin{firewall}
V98 proves full compact thermal exponential moments,
fixed-support polynomial and analytic convergence, and
explicit polynomial Haar-identity errors for qualifying
growing patches. It proves neither a physically normalized
interacting four-dimensional continuum Yang--Mills
construction nor a positive physical Hamiltonian mass gap.
The original goal remains active.
\end{firewall}
\section{V99: a sharp action-density susceptibility and response norm}
\label{sec:v99-susceptibility}

\subsection{A source-cancellation route that does not require root paths}

V98 supplies exponential moments but its general source
bound uses absolute source sums. We now obtain an
\(\ell^2\) covariance bound for one specific gauge-invariant
observable: the symmetrically distributed Wilson action
in a unit four-dimensional cell. The observable is chosen
so that reflection dissemination gives exactly the full
action, including all shared faces. This avoids growing
transport paths, but it does not estimate every individual
or colour-valued curvature source.

The previous chessboard defect bounds and V98's thermal
exponential argument are not repeated as new achievements.
The new steps are the exact action partition, its
quadratic source consequence, a joint energy/susceptibility
selection, and a sharp infinite-volume sum rule.
In particular, selecting a minimum of the mean energy
alone, as in V97, does not automatically control the
second derivative of the pressure. A different, explicitly
defined selection will be needed here.

The scalar chessboard argument below is a specialization
of the reflection framework already used in the archives.
A standard formulation allowing shared boundaries in
site-reflection blocks is M.~Biskup,
\emph{Reflection Positivity and Phase Transitions in Lattice
Spin Models}, Section 4.3, Theorem 4.8 in the
\href{https://www.math.ucla.edu/~biskup/PDFs/papers/Prague-School.pdf}
{author's lecture manuscript}.
We give the scalar-label verification below to account
explicitly for the present link variables and shared
plaquette faces. No spin-model infrared bound is imported
into Yang--Mills.

All expectations in this section use the full compact
\(\SU(2)\) law \(\mu^{\mathrm W}_{M,\beta}\), with
dyadic \(M\ge4\) and \(\beta>0\), unless a smaller range
is stated. Set \(V=M^4\). There is no chart restriction,
auxiliary instrument, or response-generated factor.

\subsection{A reflection-covariant partition of the Wilson action}

For \(\mu<\nu\), write
\[
 {\cal C}_{\mu\nu}
   =\{\epsilon\in\{0,1\}^4:\epsilon_\mu=\epsilon_\nu=0\}.
\]
The unit cell based at \(x\) is \(x+[0,1]^4\).
Define its symmetrically assigned action and thermal energy by
\begin{equation}
 E_x=\frac14\sum_{\mu<\nu}\ \sum_{\epsilon\in{\cal C}_{\mu\nu}}
                   s_{(x+\epsilon;\mu,\nu)},\qquad
 T_x=\beta E_x
     =\frac14\sum_{\mu<\nu}\ \sum_{\epsilon\in{\cal C}_{\mu\nu}}
                   X_{(x+\epsilon;\mu,\nu)}.
 \label{eq:v99-cell-energy}
\end{equation}
There are 24 face terms, each with weight \(1/4\).
They are not 24 independent random variables.

\begin{lemma}[Exact action partition and reflection covariance]
\label{lem:v99-action-partition}
Every positively oriented plaquette belongs to exactly
four unit cells, and
\begin{equation}
                  \sum_xE_x=S,\qquad \sum_xT_x=\beta S,
       \qquad \mathbb ET_x=\frac{\beta}{V}\mathbb ES.
 \label{eq:v99-action-partition}
\end{equation}
Reflection through an integer coordinate plane takes
\(E_x\) to \(E_{x'}\), where \(x'\) is the base of the
reflected unit cell. Its pullback is exactly the ordinary
translated cell-energy observable at that reflected cell.
\end{lemma}

\begin{proof}
A plaquette based at \(y\) in directions \(\mu,\nu\)
is a face of the cells based at
\(x=y-\epsilon\), \(\epsilon\in{\cal C}_{\mu\nu}\).
These four cells are distinct for \(M\ge4\).
Thus its total assigned weight is one, proving the first
two identities. Translation invariance gives the third.

A coordinate reflection permutes the 24 geometric faces
of a unit cell onto those of the reflected cell.
For a face whose orientation changes, its ordered
holonomy is inverted and possibly conjugated by a
change of starting corner. The scalar
\(s_p=1-\Tr(U_p)/2\) is unchanged under either operation.
All face weights are \(1/4\), so the entire function
transforms as claimed. Although adjacent cells share
links and faces, each \(E_x\) is supported in its
closed unit cell.
\end{proof}

The use of closed cells is intentional. At a site plane,
cell supports on opposite sides may share boundary links.
The conditional boundary-Haar reflection form
\eqref{eq:v98-site-RP} remains valid for these functions.
There is no step in which shared faces are treated as
independent or counted with full weight in every cell.

\subsection{The exact chessboard source bound and its Hessian}

\begin{proposition}[Scalar cell sources]
\label{prop:v99-chessboard-source}
For any real source array \(a=(a_x)_{x\in\Lambda_M}\),
\begin{equation}
 \log\mathbb E\exp\left(\sum_xa_xT_x\right)
 \le\frac1V\sum_x
        \log\frac{Z^{\mathrm W}_{M,(1-a_x)\beta}}
                         {Z^{\mathrm W}_{M,\beta}}.
 \label{eq:v99-source-chessboard}
\end{equation}
The compact integral defining \(Z^{\mathrm W}_{M,u}\)
is finite and positive for every real \(u\), so the
formula is meaningful even when some \(a_x\ge1\).
Reflection positivity is used only for the base law at
\(\beta>0\), not for those other couplings.
\end{proposition}

\begin{proof}
For a real scalar \(b\), set
\[
 c(b)=\left(\mathbb Ee^{b\sum_xT_x}\right)^{1/V}>0 .
\]
For a scalar array \(b\), put
\[
                    {\cal G}(b)
                    =\frac{\mathbb Ee^{\sum_xb_xT_x}}
                           {\prod_xc(b_x)}.
\]
Split the unit cells into the two halves determined
by an integer site plane and its opposite plane.
There is no cell with interior on both sides. Reflection
positivity, applied to the product of the source factors
in each half, gives
\[
                     {\cal G}(b)^2
                       \le{\cal G}(b^+){\cal G}(b^-).
\]
Here \(b^+\) duplicates the positive half by reflection,
and \(b^-\) duplicates the negative half. Lemma
\ref{lem:v99-action-partition} identifies the reflected
factors with the corresponding cell sources. The
normalizing factors cancel in this inequality because
the two duplicated arrays together contain twice
each label multiplicity of the original array.
Shared boundary links cause no additional factor:
they are exactly the variables retained in the
conditional reflection form of V98.

Let the labels range over the finite set of values
appearing in the given array \(a\). Among all arrays
with these labels, choose one maximizing the positive
number \({\cal G}\); among maximizers choose one with
the fewest disagreeing nearest-neighbor cell pairs.
If it is not constant, choose a site plane separating
at least one such pair. Let \(n_+,n_-\) count the
disagreeing pairs internal to the two halves and
\(n_\times>0\) those crossing their boundaries.
The two reflected arrays have respectively \(2n_+\)
and \(2n_-\) disagreements: all pairs across the
reflection boundaries now agree. At least one of
these counts is strictly smaller than
\(n_++n_-+n_\times\).

The reflection inequality and maximality imply that
both duplicated arrays are also maximizers.
This contradicts the choice of the disagreement count.
The maximizing array must therefore be constant.
For a constant array with label \(b\), the definition
of \(c(b)\) gives \({\cal G}=1\). Consequently
\({\cal G}(a)\le1\). Finally,
\[
 c(b)^V=\mathbb Ee^{b\beta S}
                    =Z^{\mathrm W}_{M,(1-b)\beta}/
                                      Z^{\mathrm W}_{M,\beta}.
\]
Taking logarithms proves the proposition.
\end{proof}

Define the thermal susceptibility and centered cell energy by
\begin{equation}
 \chi_{M,\beta}=\frac{\beta^2}{V}\Var(S),\qquad
                  \widetilde T_x=T_x-\mathbb ET_x.
 \label{eq:v99-susceptibility}
\end{equation}

\begin{theorem}[Sharp covariance norm at every finite coupling]
\label{thm:v99-covariance-norm}
For every real or complex source array \(a\),
\begin{equation}
 \mathbb E\left|\sum_xa_x\widetilde T_x\right|^2
                         \le\chi_{M,\beta}\sum_x|a_x|^2 .
 \label{eq:v99-l2-bound}
\end{equation}
In fact, if \(\Gamma_{M,\beta}\) denotes the covariance
operator of the cell field on \(\ell^2(\Lambda_M)\), then
\begin{equation}
       0\le\Gamma_{M,\beta}\le\chi_{M,\beta}I,\qquad
       \Gamma_{M,\beta}{\bf1}=\chi_{M,\beta}{\bf1},\qquad
       \|\Gamma_{M,\beta}\|=\chi_{M,\beta}.
 \label{eq:v99-exact-norm}
\end{equation}
\end{theorem}

\begin{proof}
Set
\[
 g(t)=\log\frac{Z^{\mathrm W}_{M,(1-t)\beta}}
                       {Z^{\mathrm W}_{M,\beta}}
                               -t\beta\mathbb ES .
\]
Finite-volume compactness permits differentiation, giving
\[
                 g(0)=g'(0)=0,\qquad
                 g''(0)=\beta^2\Var(S)=V\chi_{M,\beta}.
\]
Center \eqref{eq:v99-source-chessboard}, replace \(a\)
by \(ta\), and differentiate its second-order inequality
at \(t=0\). Both sides and their first derivatives
vanish there, so the nonnegative difference of the
two sides has nonnegative second derivative. This gives
\eqref{eq:v99-l2-bound} for real \(a\).
For complex \(a\), apply the real result to its real
and imaginary parts and add.

Positivity of a covariance operator is automatic.
The inequality gives its upper bound. Its row sums
are constant by translation invariance, and
\[
 \frac1V\sum_{x,y}\operatorname{Cov}(T_x,T_y)
      =\frac1V\Var\!\left(\sum_xT_x\right)
      =\chi_{M,\beta}.
\]
Each row sum is therefore \(\chi_{M,\beta}\).
The constant vector is an eigenvector attaining the
upper bound, proving the exact norm.
\end{proof}

Equivalently, every discrete Fourier structure factor
of the cell energy lies in \([0,\chi_{M,\beta}]\), with
equality at zero momentum. This does not claim that
every individual off-diagonal covariance is nonnegative.
For the source-tilted law proportional to
\(e^{\lambda\sum_ya_yT_y}d\mu^{\mathrm W}_{M,\beta}\),
ordinary differentiation gives
\begin{equation}
 \left.\frac{d}{d\lambda}\mathbb E_\lambda T_x\right|_{\lambda=0}
                =\sum_y\operatorname{Cov}(T_x,T_y)a_y .
 \label{eq:v99-response}
\end{equation}
Thus \eqref{eq:v99-exact-norm} is an exact
\(\ell^2\)-to-\(\ell^2\) norm for the linear response
at the bare Wilson law, not just a one-point moment bound.

\subsection{A joint selection of mean energy and susceptibility}

Write \(s=\log\beta\) and, only in this subsection, set
\[
 q_M(s)=\frac{e^s}{V}\mathbb E^{\mathrm W}_{M,e^s}S,
                  \qquad \chi_M(s)=\chi_{M,e^s}.
\]
Differentiating the same finite compact integral yields
the exact identity
\begin{equation}
                             q_M'(s)=q_M(s)-\chi_M(s).
 \label{eq:v99-log-derivative}
\end{equation}
Use V97's logarithmic window
\[
 r_M=\log\log M\ge4,\qquad
        a_M=\sqrt{r_M},\quad b_M=r_M,\quad
        \tau_M=b_M-a_M.
\]
Define averages with respect to \(ds/\tau_M\) on this
interval, denoted \(\bar q_M,\bar\chi_M\), and put
\begin{equation}
 m_M=\min_{[a_M,b_M]}q_M,\quad
 \delta_M=\bar q_M-m_M,\quad
 \rho_M=\sqrt{\delta_M}+\tau_M^{-1/2}.
 \label{eq:v99-selection-parameters}
\end{equation}
Let \(s_M^\sharp\) be the leftmost minimizer of
\(q_M+\rho_M\chi_M\) on \([a_M,b_M]\), and set
\(\beta_M^\sharp=e^{s_M^\sharp}\). All these quantities
are defined from the exact normalizer. No numerical
evaluation of that normalizer is assumed.

\begin{theorem}[Simultaneously sharp energy and bounded susceptibility]
\label{thm:v99-joint-selection}
This selection is well defined and satisfies
\begin{equation}
 q_M(s_M^\sharp)\longrightarrow\frac92,\qquad
                         \limsup_M\chi_M(s_M^\sharp)\le\frac92.
 \label{eq:v99-selection-limit}
\end{equation}
More explicitly, with the \(\eta_M\) of
\eqref{eq:v97-sharp-budget},
\begin{equation}
 \chi_M(s_M^\sharp)\le
                 \frac92+\eta_M+\frac{142}{\tau_M}
                                           +\sqrt{\delta_M},
 \qquad \delta_M\longrightarrow0.
 \label{eq:v99-finite-susceptibility}
\end{equation}
The sequence \(\beta_M^\sharp\) need not coincide with
V97's \(\beta_M^{\mathrm W,*}\).
\end{theorem}

\begin{proof}
All functions are continuous on the compact interval,
and \(\rho_M>0\); hence all minima and the leftmost
minimizers exist. Integrating
\eqref{eq:v99-log-derivative} gives
\begin{equation}
 \bar\chi_M=\bar q_M
                      -\frac{q_M(b_M)-q_M(a_M)}{\tau_M}.
 \label{eq:v99-average-susceptibility}
\end{equation}
The pressure calculation in V97, before taking the
minimum, proves \(\bar q_M\le9/2+\eta_M\).
The V95 energy bound gives throughout this interval
\[
 0\le q_M(s)<130+\frac{12s}{M}\le142.
\]
Here \(s\le r_M<M\). Consequently
\[
                 0\le\bar\chi_M
                     \le\frac92+\eta_M+\frac{142}{\tau_M}.
\]

We also have \(m_M\to9/2\). For the lower limit,
choose a minimizer of \(q_M\) in each interval and apply
\eqref{eq:v97-full-liminf} to those couplings: their
lower bound \(e^{\sqrt{r_M}}\) tends to infinity and
their growth satisfies the required window. For the
upper limit use \(m_M\le\bar q_M\le9/2+\eta_M\).
It follows that \(\bar q_M\to9/2\), that
\(\delta_M=\bar q_M-m_M\to0\), and that \(\rho_M\to0\).

The minimum of a function is no greater than its average,
so
\[
 q_M(s_M^\sharp)+\rho_M\chi_M(s_M^\sharp)
                           \le\bar q_M+\rho_M\bar\chi_M.
\]
Using \(\chi_M\ge0\) bounds the first term above by
\(\bar q_M+\rho_M\bar\chi_M\); its lower bound is \(m_M\).
This proves the first limit in \eqref{eq:v99-selection-limit}.
Using instead \(q_M(s_M^\sharp)\ge m_M\) gives
\[
 \chi_M(s_M^\sharp)\le\bar\chi_M+\frac{\delta_M}{\rho_M}
                           \le\bar\chi_M+\sqrt{\delta_M}.
\]
The final inequality also holds when \(\delta_M=0\).
This proves the claimed finite bound and the upper limit.
\end{proof}

This is an existence selection in the full compact theory,
not a proposed physical running coupling. Its definition
uses a susceptibility penalty, not the unsupported
assumption that a minimum of \(q_M\) lies in the interior
or has derivative zero.

\subsection{The Gaussian cell-energy sum rule}

Along \(\beta_M^\sharp\), the joint local
curvature-energy limit is again
\[
                     (Y^\gamma,X)
                          \Longrightarrow(G,|G|^2/2).
\]
Indeed, Theorem~\ref{thm:v99-joint-selection} saturates
the energy floor in the full-law classification
of Theorem~\ref{thm:v97-classification}, forcing its
constant harmonic component to vanish. That argument
uses energy saturation, not the particular definition
of V97's minimizer. The exponential bounds of V98
apply with \(L=1\) on the entire present window.
They therefore give all fixed mixed polynomial moments
for this new selection as well.

Let \(w_x(p)\) be the face weight in
\eqref{eq:v99-cell-energy}, now on the infinite lattice.
Thus each \(w_x\) has 24 entries of value \(1/4\),
\(\sum_pw_x(p)=6\), and \(\sum_xw_x(p)=1\).
The limiting centered cell field is
\begin{equation}
 {\cal T}_x=\frac12\sum_pw_x(p)
                           \big(|G_p|^2-\mathbb E|G_p|^2\big).
 \label{eq:v99-Gaussian-cell}
\end{equation}
It is a quadratic Gaussian observable, not a Gaussian
random variable in general.

\begin{lemma}[An absolutely summable susceptibility identity]
\label{lem:v99-Gaussian-sum-rule}
The covariance
\(\Gamma_\infty(x)=\mathbb E({\cal T}_0{\cal T}_x)\)
satisfies
\begin{equation}
 \Gamma_\infty(x)=\frac32
            \sum_{p,q}w_0(p)w_x(q)P_{pq}^{\,2}\ge0,\qquad
                       \sum_{x\in\mathbb Z^4}\Gamma_\infty(x)=\frac92 .
 \label{eq:v99-Gaussian-sum-rule}
\end{equation}
In particular \(\Gamma_\infty\in\ell^1(\mathbb Z^4)\).
Its convolution operator on \(\ell^2(\mathbb Z^4)\)
has norm exactly \(9/2\).
\end{lemma}

\begin{proof}
The three colours are independent copies of a real
Gaussian field with covariance \(P\). Wick's formula
gives
\[
 \operatorname{Cov}(|G_p|^2/2,|G_q|^2/2)
                                                  =\tfrac32P_{pq}^{\,2}.
\]
This proves the first formula and its nonnegativity.
The infinite-lattice operator \(P\) is an orthogonal
projection, with \(P_{pp}=1/2\), as established in
V90 and reused in V97. Therefore
\[
               \sum_qP_{pq}^{\,2}
                =\|P\delta_p\|_2^2=P_{pp}=\tfrac12.
\]
All terms in the covariance sum are nonnegative.
Tonelli's theorem and the action partition give
\[
 \sum_x\Gamma_\infty(x)
  =\frac32\sum_pw_0(p)\sum_qP_{pq}^{\,2}
                            \sum_xw_x(q)
  =\frac34\sum_pw_0(p)=\frac92.
\]
This is an exact projection sum rule, not an assumption
about pointwise decay of the interacting covariance.

The \(\ell^1\) sum bounds the convolution norm above.
For the reverse inequality let \(D_R=[-R,R]^4\cap\mathbb Z^4\).
Then
\[
 \frac1{|D_R|}\mathbb E\left(\sum_{x\in D_R}{\cal T}_x\right)^2
   =\sum_z\Gamma_\infty(z)
                   \frac{|D_R\cap(D_R-z)|}{|D_R|}
                       \longrightarrow\frac92 .
\]
For each fixed \(z\), the ratio tends to one and is
bounded by one; dominated convergence applies to the
summable kernel. These are Rayleigh quotients, so they
give the matching lower bound on the operator norm.
\end{proof}

\subsection{Sharp susceptibility and arbitrary-support linear response}

\begin{theorem}[Full compact susceptibility saturation]
\label{thm:v99-sharp-susceptibility}
Along the jointly selected dyadic regulators,
\begin{equation}
                   \frac{(\beta_M^\sharp)^2}{M^4}
                          \Var^{\mathrm W}_{M,\beta_M^\sharp}(S)
                                      \longrightarrow\frac92 .
 \label{eq:v99-sharp-susceptibility}
\end{equation}
Consequently there are nonnegative \(\varepsilon_M\to0\)
such that, for every real or complex array on the entire
regulator, with no restriction on its support size,
\begin{equation}
 \mathbb E^{\mathrm W}_{M,\beta_M^\sharp}
       \left|\sum_xa_x\widetilde T_x\right|^2
                    \le\left(\frac92+\varepsilon_M\right)
                                                     \sum_x|a_x|^2.
 \label{eq:v99-selected-response}
\end{equation}
The susceptibility is the exact operator norm, not merely
one possible upper bound. The finite choice
\[
       \varepsilon_M=\eta_M+142/\tau_M+\sqrt{\delta_M}
\]
is valid for the displayed inequality.
\end{theorem}

\begin{proof}
Only the lower limit in
\eqref{eq:v99-sharp-susceptibility} remains to be shown.
Fix \(R\), embed \(D_R\) in the torus for sufficiently
large \(M\), and apply Theorem~\ref{thm:v99-covariance-norm}
to its indicator:
\[
 \chi_{M,\beta_M^\sharp}\ge
   \frac1{|D_R|}\mathbb E
                      \left(\sum_{x\in D_R}\widetilde T_x\right)^2.
\]
There are only finitely many plaquettes in this expression.
The local limit and the higher uniform integrability
proved above show that its right-hand side converges
to the Gaussian cell-energy Rayleigh quotient in
Lemma~\ref{lem:v99-Gaussian-sum-rule}.
Taking a lower limit first in \(M\), and then letting
\(R\) tend to infinity, gives
\(\liminf_M\chi_{M,\beta_M^\sharp}\ge9/2\).
The upper limit from Theorem~\ref{thm:v99-joint-selection}
proves equality. The response bound follows from the
exact finite-volume covariance norm and
\eqref{eq:v99-finite-susceptibility}.
\end{proof}

For example, for every nonempty subset \(D\subset\Lambda_M\),
the variance of its centered cell-energy average is at most
\begin{equation}
 \mathbb E\left|\frac1{|D|}\sum_{x\in D}\widetilde T_x\right|^2
                        \le\frac{9/2+\varepsilon_M}{|D|}.
 \label{eq:v99-average-variance}
\end{equation}
No independence, sign assumption on finite-volume
off-diagonal covariances, or bound on a transport filling
is used. The improvement over the V98 absolute-source
bound is genuinely at the level of arbitrary supports.
It concerns this particular scalar action-density
combination.

\begin{proposition}[Strong \(\ell^2\) approximation of source observables]
\label{prop:v99-l2-source-limit}
Represent the tori by centered fundamental cubes and extend
their coefficient arrays by zero to \(\mathbb Z^4\).
Suppose real arrays \(a^{(M)}\) converge strongly in
\(\ell^2(\mathbb Z^4)\) to \(a\). Then
\begin{equation}
 \sum_xa_x^{(M)}\widetilde T_x
             \ \Longrightarrow\ {\cal T}(a)
                         :=\sum_{x\in\mathbb Z^4}a_x{\cal T}_x ,
 \label{eq:v99-l2-source-limit}
\end{equation}
and their second moments converge. The series on the right
is defined by convergence in \(L^2\) of the limiting
Gaussian probability space, independently of the finite
exhaustion used. The assertion also holds jointly for
finitely many such coefficient sequences.
\end{proposition}

\begin{proof}
Lemma~\ref{lem:v99-Gaussian-sum-rule} bounds the variance
of finite sums by \((9/2)\|a\|_2^2\). It therefore defines
the right side as an \(L^2\) completion of finitely supported
coefficients. This asserts no coupling in \(L^2\) between
different finite-volume probability spaces.

For a fixed finite box, truncated source observables
converge in law and second moment by the established
local polynomial limit; coefficient convergence on that
box does not change this conclusion. The second moment
of the finite-volume remainder is bounded by
\((9/2+\varepsilon_M)\) times the squared \(\ell^2\)
tail norm, using \eqref{eq:v99-selected-response}.
Strong \(\ell^2\) convergence makes these tails uniformly
small after enlarging the fixed box. The limiting tail
has the analogous bound. Truncation now proves convergence
in distribution and of the \(L^2\) norms, hence of the
second moments. Applying the same argument to every
finite linear combination proves the joint assertion.
\end{proof}

\subsection{What this norm does not yet prove}

The result removes an absolute-source loss for the
linear response of a symmetrically assigned bare action
density. It does not prove the full Yang--Mills mass gap.
The following distinctions are essential.

The \(\ell^2\) estimate is a second-derivative statement
at zero source. There is no established uniform
neighborhood in which all the source-tilted covariance
operators satisfy the same bound. The chessboard formula
contains pressure differences at shifted couplings,
and the new selection only controls susceptibility at
the selected coupling. In particular the present proof
does not produce a Gaussian exponential bound, a
logarithmic Sobolev inequality, or a central limit theorem
for the total action.

The observable \(T_x\) sums the 24 faces of a unit cell
with fixed weights. The theorem does not resolve the
colour-valued, path-transported curvature source problem.
It also cannot be applied to an arbitrary individual
orientation by removing the other faces: covariance
inequalities do not follow from the pointwise order of
uncentered nonnegative variables.

For a macroscopic profile, the coefficients
\(a_x^{(M)}=M^{-2}f(x/M)\) can have bounded \(\ell^2\)
norm and hence bounded variance by
\eqref{eq:v99-selected-response}. For a nonzero profile
they need not converge strongly in \(\ell^2\);
their mass spreads over an increasing number of sites.
Proposition~\ref{prop:v99-l2-source-limit} therefore
does not identify their limiting law. Bounded variance
is not proof of a white-noise or interacting continuum
limit. Nor is this central-fluctuation normalization
the unrenormalized dimension-four action-density
smearing: coefficients of order one over a physical
region have squared \(\ell^2\) norm of order \(M^4\),
so the present variance bound still diverges.

The finite structure-factor bound gives no positive
physical correlation length or exponential decay.
The limit sum rule arises from the Gaussian projection,
not a nonzero physical mass. Moreover the selection
\(\beta_M^\sharp\) is not an identified renormalized
Yang--Mills trajectory. No response-generated blocked
law has been shown to retain this covariance norm.

The next question is whether this bare, gauge-invariant
response estimate can be made stable under the source
and scale transformations actually required by the RG
programme, or can be coupled to estimates for physically
normalized nonlocal sources. Repeating the mean-energy
selection alone would not answer that question.

\subsection{Exact diagnostics and preservation}

The certificate is
\begin{center}\small
\path{scripts/verify_ym_f14_v99_energy_susceptibility.py}.
\end{center}
It checks 104,448 cell-face occurrences, the multiplicity
four of every plaquette, and 17,408 reflected cell-face
sets on finite periodic lattices. There are 192
half-duplication label checks, including 136 strict
decreases of the disagreement count. These audit the
combinatorial step in the normalized source argument.
Forty exact finite-ensemble derivative calculations
check the logarithmic energy/susceptibility identity,
and 320 rational finite-array tests check the weighted
selection inequalities.

For a separate finite Gaussian reference, 16 rational
Fourier projections, six projection row-sum identities,
16 energy structure factors, and 17 signed-source
variance inequalities are checked. Its finite
susceptibility is \(135/32=(9/2)(1-1/16)\), retaining
the absent zero mode. This binary Fourier reference is
not a side-two Wilson regulator and is not being used
in place of the infinite-lattice projection proof.

The certificate has 8,937 bytes and SHA--256
\begin{center}\scriptsize\ttfamily
1564AA567CD03031D0DF37319FA103ED917792CF94C3EFA50AE86DD2FF21C563.
\end{center}
These are finite algebraic diagnostics, not samples from
the Wilson law or a machine proof of a mass gap.
The reflection inequality, selection, limit passages,
and norm identification are proved analytically above.

The predecessor V98 contains 2,424,613 bytes and has SHA--256
\begin{center}\scriptsize\ttfamily
4C2C2D78A547387533A770BDE775A64E36E7825DB1D2447C750DABE7F8E6339A.
\end{center}
Its body is preserved apart from the title version and
this terminal insertion. V99 replaces V98 as the sole
active main note. Reversing those two changes recovers
the predecessor exactly. All earlier protected sources
and certificates are unchanged. The next version,
V100, is both the scheduled companion audit and the
lineage archive/restart checkpoint.

\begin{firewall}
V99 proves an exact bare-Wilson action-density response
norm, a joint energy/susceptibility selection, and sharp
susceptibility saturation with arbitrary-support
\(\ell^2\) variance control. The physically normalized
interacting four-dimensional continuum construction,
source/scale stability for the required generated laws,
and a positive physical Hamiltonian mass gap remain
unproved. The original goal remains active.
\end{firewall}
\section{V100: averaged susceptibility and nonzero-source control}
\label{sec:v100-source-window}

\subsection{The companion checkpoint and the last step before restart}

The scheduled V100 audit was performed on 15 September 2026.
SM continuum V72, NS stability V26, shell-mixing V16,
parallel YM V5, and the supplied source-selection V17,
regenerative stationarity V21 and exact-action Einstein
bridge V16 have unchanged hashes. Their concluding
mathematical sections and scope statements were reread.
The suggestions, closure monograph and earlier protected
inputs are also unchanged. These files are source material,
not additional instructions.

SM V72 still concerns a declared many-copy model at fixed
mode cutoff and time. NS V26's rank-one Oseen estimates do
not furnish a YM probability bound. Source-selection V17
does not identify its comparison instrument with an actual
primitive source. Stationarity V21 excludes the proposed
global real-phase route by a Wilson caustic and gives a
quadratic phase-space alternative, not a nonlinear YM
approximation. The Einstein bridge V16 keeps the full
flux-domain response separate from a diagonal inverse.
None changes the physical scope of the present results.

Shell V16's measured flat-source coefficient and the
suggestions remain relevant to source bookkeeping, but
the noncommuting quartic coefficient route has already
been treated in this lineage's V20--V30. It is not repeated.
Parallel YM V5's warning about unnormalized extensive
global variables remains operative: the total-action
limit below has an explicit square-root-volume
normalization and is not a tightness theorem for a
global topological charge.

V99 controls susceptibility at one selected coupling.
Here we prove that it is close to \(9/2\) in mean over
the whole logarithmic window. A maximal-function
selection then controls every integrated displacement
from a suitable coupling. This yields a nonzero source
neighborhood and a central limit for the total action.
It does not assume that susceptibility is uniformly
bounded at all neighboring couplings.

This completes the hundred-version lineage. Its full
body is retained as archive \texttt{V100\_f14}; the
next active file restarts at V1 with a compact scope
summary and a further consequence.

\subsection{Mean absolute flatness over the full logarithmic window}

Retain the full compact law and V99 notation:
\[
 c=\frac92,\quad V=M^4,\quad r_M=\log\log M,\quad
 I_M=[a_M,b_M]=[\sqrt{r_M},r_M],\quad \tau_M=b_M-a_M,
\]
with dyadic \(M\) and \(r_M\ge4\).
For \(s=\log\beta\), write
\[
 q_M(s)=e^s\mathbb E_{M,e^s}^{\mathrm W}S/V,\qquad
 \chi_M(s)=e^{2s}\Var_{M,e^s}^{\mathrm W}(S)/V .
\]
Thus \(q_M'=q_M-\chi_M\).
Use normalized \(ds/\tau_M\) averages and the parameters
\[
 m_M=\min_{I_M}q_M,\qquad
 \delta_M=\bar q_M-m_M,\qquad
 \kappa_M=\sqrt{\delta_M}+\tau_M^{-1}.
\]
V99 proves \(m_M,\bar q_M\to c\), \(\delta_M\to0\),
and
\[
 \bar\chi_M\le c+\eta_M+142/\tau_M,\qquad 0\le q_M\le142,
\]
where \(\eta_M\) is the explicit pressure error in V97.

\begin{theorem}[Window-averaged susceptibility]
\label{thm:v100-L1-flatness}
One has
\begin{equation}
       e_M:=\frac1{\tau_M}\int_{I_M}|\chi_M(s)-c|\,ds
                                                    \longrightarrow0 .
 \label{eq:v100-L1-flatness}
\end{equation}
\end{theorem}

\begin{proof}
Let
\[
 G_M=\{s\in I_M:q_M(s)\le m_M+\kappa_M\}.
\]
This compact set is nonempty, and Markov's inequality
for \(q_M-m_M\ge0\) gives
\[
                  |I_M\setminus G_M|/\tau_M
                                          \le\delta_M/\kappa_M\to0.
\]
We claim that
\[
             \zeta_M:=\max\{0,c-\min_{G_M}\chi_M\}\to0.
\]
Otherwise, along a subsequence there would be points
\(s_M\in G_M\) with \(\chi_M(s_M)\le c-\epsilon\)
for some fixed \(\epsilon>0\). But
\(m_M\le q_M(s_M)\le m_M+\kappa_M\) tends to \(c\).
Energy saturation therefore gives the full local
Gaussian curvature limit by V97's classification,
with all local polynomial moments by V98. V99's
finite-box covariance bound then implies
\[
 \liminf_M\chi_M(s_M)\ge
     |D_R|^{-1}\mathbb E\left(\sum_{x\in D_R}{\cal T}_x\right)^2
\]
for every fixed \(R\). Letting \(R\) grow and using
the Gaussian cell-energy sum rule gives a lower
bound \(c\), a contradiction. This reasoning applies
to any sequence in \(G_M\), not just the minimizer
selected in V99.

Since \(\chi_M\ge0\) everywhere,
\[
 \frac1{\tau_M}\int_{I_M}(c-\chi_M)_+\,ds
                                \le\zeta_M+c\delta_M/\kappa_M.
\]
For a real number \(x\), \(|x-c|=x-c+2(c-x)_+\).
After averaging, this gives
\[
 e_M\le\eta_M+142/\tau_M
                      +2\zeta_M+2c\delta_M/\kappa_M\longrightarrow0.
\]
No pointwise upper bound for \(\chi_M\) across the
window is inferred from this \(L^1\) statement.
\end{proof}

\subsection{A coupling with controlled integrated displacements}

For nonnegative integrable \(f\) on \(\mathbb R\), define
the uncentered interval maximal function
\[
 {\cal M}f(s)=\sup_{\substack{J\ {\rm bounded\ open\ interval}\\s\in J}}
                            \frac1{|J|}\int_Jf.
\]
We use the elementary bound
\begin{equation}
               |\{{\cal M}f>\lambda\}|
                            \le3\|f\|_1/\lambda,\qquad\lambda>0 .
 \label{eq:v100-maximal-weak}
\end{equation}
For completeness, cover a compact subset of the set on
the left by finitely many intervals of average exceeding
\(\lambda\). Select disjoint intervals in decreasing
order of length, discarding those meeting an earlier
selection. Each discarded interval lies in the triple
dilate of its selected partner. Thus the compact
subset has measure at most three times the sum of
the selected lengths, which is at most
\(3\|f\|_1/\lambda\). Inner regularity gives the result.
This is the standard covering argument; see, for context,
T.~Tao, \href{https://www.math.ucla.edu/~tao/247a.1.06f/notes3.pdf}
{Lecture Notes 3 for 247A, Proposition 1.1 and Lemma 1.3}.
The one-dimensional uncentered form needed here has
just been proved explicitly.

Extend \(f_M=|\chi_M-c|\mathbf1_{I_M}\) by zero outside
\(I_M\), and set
\begin{equation}
 \lambda_M=\sqrt{e_M}+\tau_M^{-1},\qquad
 J_M=[a_M+\tau_M/4,b_M-\tau_M/4].
 \label{eq:v100-maximal-parameters}
\end{equation}
For sufficiently large \(M\), choose \(s_M^\diamond\)
to be the leftmost point of
\[
                 J_M\cap\{{\cal M}f_M\le\lambda_M\},
 \qquad \beta_M^\diamond=e^{s_M^\diamond}.
\]
This is a new selection, not an assertion about the
previous \(\beta_M^{\mathrm W,*}\) or \(\beta_M^\sharp\).

\begin{lemma}[Existence and a cumulative error budget]
\label{lem:v100-maximal-selection}
The selection is well defined for all sufficiently
large dyadic \(M\). Put
\[
 d_M=\min\{s_M^\diamond-a_M,b_M-s_M^\diamond\}\ge\tau_M/4 .
\]
For either sign and \(0\le h\le d_M\),
\begin{equation}
       \int_0^h|\chi_M(s_M^\diamond\pm v)-c|\,dv
                                                \le\lambda_M h.
 \label{eq:v100-cumulative}
\end{equation}
Furthermore
\begin{equation}
 |\chi_M(s_M^\diamond)-c|\le\lambda_M,\qquad
 |q_M(s_M^\diamond)-c|\le142e^{-d_M}+\lambda_M .
 \label{eq:v100-point-control}
\end{equation}
\end{lemma}

\begin{proof}
The maximal superlevel set is open: it is a union
of the intervals witnessing averages above the
threshold. Its complement is closed.
By \eqref{eq:v100-maximal-weak}, the excluded measure
is at most \(3e_M\tau_M/\lambda_M=o(\tau_M)\),
whereas \(J_M\) has length \(\tau_M/2\).
Thus the indicated closed subset of \(J_M\) is
nonempty and compact, and has a leftmost point.

Intervals starting at that point, with length \(h\)
and lying in \(I_M\), have average at most \(\lambda_M\).
Endpoints may be included by approximation with open
intervals. This proves \eqref{eq:v100-cumulative}.
Continuity of \(\chi_M\) at the selected interior
point gives the first inequality of
\eqref{eq:v100-point-control}.

For the second, solve \(q_M'=q_M-\chi_M\) backwards
from \(b_M\):
\[
 q_M(s)-c=e^{s-b_M}(q_M(b_M)-c)
                +\int_s^{b_M}e^{s-u}(\chi_M(u)-c)\,du .
\]
The same maximal-function bound applies to every
interval from \(s=s_M^\diamond\) up to \(b_M\),
not just those of length at most \(d_M\).
If \(F(h)=\int_0^hf_M(s+v)\,dv\), then
\(F(h)\le\lambda_Mh\) over this entire range.
Integration by parts gives
\[
 \int_0^He^{-v}\,dF(v)
     =e^{-H}F(H)+\int_0^He^{-v}F(v)\,dv
                         \le\lambda_M(1-e^{-H})\le\lambda_M.
\]
Use \(|q_M(b_M)-c|\le142\) and \(b_M-s\ge d_M\).
\end{proof}

In particular this selection retains the energy
saturation and all fixed local polynomial limits
of V97--V99. More importantly, it controls integrated
susceptibility errors down to arbitrarily short
coupling intervals. It has not bounded
\(\sup_{|v|\le d_M}|\chi_M(s_M^\diamond+v)-c|\).

\subsection{A two-sided centered pressure estimate}

Freeze \(M\), write \(\beta=\beta_M^\diamond\), and set
\begin{equation}
 \Psi_M(t)=\frac1V\log\mathbb E^{\mathrm W}_{M,\beta}
                  e^{t\beta(S-\mathbb ES)},\qquad
                       J(t)=-\log(1-t)-t\quad(t<1).
 \label{eq:v100-Psi}
\end{equation}

\begin{theorem}[A source window with sharp quadratic coefficient]
\label{thm:v100-pressure-window}
If \(t<1\) and \(|\log(1-t)|\le d_M\), then
\begin{equation}
                  |\Psi_M(t)-cJ(t)|\le\lambda_MJ(t).
 \label{eq:v100-pressure-window}
\end{equation}
Here \(J(t)\ge0\), with equality only at zero.
In particular the interval \(|t|\le R_M\), where
\begin{equation}
                            R_M=1-e^{-d_M}\longrightarrow1,
 \label{eq:v100-source-radius}
\end{equation}
is included in this source window.
\end{theorem}

\begin{proof}
Finite compact integrals are smooth, and the pressure
ratio formula gives
\[
 \Psi_M(0)=\Psi_M'(0)=0,\qquad
 \Psi_M''(u)=
       \frac{\chi_M(s_M^\diamond+\log(1-u))}{(1-u)^2}.
\]
For \(t>0\), let \(h=-\log(1-t)\) and
make the change \(u=1-e^{-v}\) in the twice-integrated
derivative. The result is
\[
 \Psi_M(t)=\int_0^h
       (1-e^{v-h})\,\chi_M(s_M^\diamond-v)\,dv .
\]
The kernel \(k(v)=1-e^{v-h}\) is nonnegative,
decreasing, and zero at \(h\). Its integral is
\(h-t=J(t)\).

For \(t<0\), put \(h=\log(1-t)\) and \(u=1-e^v\).
The corresponding formula is
\[
 \Psi_M(t)=\int_0^h
       (e^{h-v}-1)\,\chi_M(s_M^\diamond+v)\,dv .
\]
This kernel has the same three properties, and its
integral is \(-t-h=J(t)\).

For either sign, write \(F(v)=\int_0^v
|\chi_M(s_M^\diamond\pm z)-c|\,dz\le\lambda_Mv\).
Integration by parts, using \(k(h)=0\), yields
\[
 \left|\int_0^hk(v)(\chi_M-c)\,dv\right|
       \le\int_0^hk(v)\,dF(v)
       =-\int_0^hk'(v)F(v)\,dv
       \le\lambda_M\int_0^hk(v)\,dv .
\]
This proves \eqref{eq:v100-pressure-window}; at zero
the assertion is immediate.
For \(|t|\le R_M\), the positive endpoint obeys
\(-\log(1-R_M)=d_M\), while
\(\log(1+R_M)\le d_M\) follows from
\(2-e^{-d_M}\le e^{d_M}\).
\end{proof}

The cancellation of the centered linear term is exact.
The proof uses cumulative susceptibility control,
not differentiation of an \(L^1\) convergence statement
and not a bound on the largest nearby susceptibility.

\subsection{Nonzero cell sources and arbitrary-support tail bounds}

Let \(\widetilde T_x\) be the centered cell field
of V99 at \(\beta_M^\diamond\), and put \(k_M=c+\lambda_M\).
The exact chessboard inequality of
Proposition~\ref{prop:v99-chessboard-source} now gives,
for every real array \(b\) with \(\|b\|_\infty\le R_M\),
\begin{equation}
 0\le\log\mathbb E e^{\sum_xb_x\widetilde T_x}
                       \le k_M\sum_xJ(b_x).
 \label{eq:v100-nonzero-source}
\end{equation}
The lower bound is Jensen's inequality.
The upper bound follows by centering the chessboard
formula and using \(\Psi_M(b_x)\le k_MJ(b_x)\)
for each site. In particular, the admissible source
amplitude does not decrease when the support grows.

For a real array \(a\), write
\[
 X(a)=\sum_xa_x\widetilde T_x,\qquad
 v(a)=k_M\|a\|_2^2,\qquad B(a)=\|a\|_\infty/R_M .
\]
The power series of \(J\) gives
\[
 J(u)\le\frac{u^2}{2(1-|u|)}\quad(|u|<1).
\]
Since \(R_M<1\), \eqref{eq:v100-nonzero-source} implies
\begin{equation}
 \log\mathbb E e^{\pm tX(a)}
            \le\frac{v(a)t^2}{2(1-B(a)t)}
                  \quad(0\le t<1/B(a)),
 \label{eq:v100-subgamma}
\end{equation}
with the usual interpretation when \(a=0\).

\begin{proposition}[A Bernstein-type bound without a support-size loss]
\label{prop:v100-Bernstein}
For every \(u>0\),
\begin{equation}
 \mathbb P\!\left\{|X(a)|>
                \sqrt{2v(a)u}+B(a)u\right\}\le2e^{-u}.
 \label{eq:v100-Bernstein}
\end{equation}
\end{proposition}

\begin{proof}
For \(v(a)>0\), apply exponential Markov to
\eqref{eq:v100-subgamma}. Choose
\[
 t=\frac{\sqrt{2u/v(a)}}{1+B(a)\sqrt{2u/v(a)}}.
\]
The exponent at the displayed threshold is \(-u\).
Apply the same argument to \(-X(a)\) and add the
two probabilities. If \(v(a)=0\), then \(a=0\)
and the result is immediate.
\end{proof}

This is a nonzero-source improvement over V99's
zero-source Hessian bound. It does not establish a
uniform bound for the Hessian at every tilted law:
an upper bound for a cumulant generating function
cannot simply be differentiated twice as an inequality
away from a point of equality.

\subsection{A central limit for the total centered action}

\begin{theorem}[Selected full-Wilson total-action fluctuations]
\label{thm:v100-total-CLT}
Along the new dyadic selection,
\begin{equation}
 W_M=\frac{\beta_M^\diamond}{\sqrt{M^4}}
          \left(S-\mathbb E^{\mathrm W}_{M,\beta_M^\diamond}S\right)
                     \ \Longrightarrow\ N(0,9/2).
 \label{eq:v100-total-CLT}
\end{equation}
All fixed moments converge to the corresponding
Gaussian moments.
\end{theorem}

\begin{proof}
For each fixed real \(t\) and sufficiently large \(M\),
\[
 \log\mathbb E e^{tW_M}
                         =V\Psi_M(t/\sqrt V).
\]
Since \(VJ(t/\sqrt V)\to t^2/2\),
\eqref{eq:v100-pressure-window} and \(\lambda_M\to0\)
give
\[
                   \mathbb E e^{tW_M}\longrightarrow e^{ct^2/2}.
\]
The susceptibility control already gives bounded
second moments and hence tightness. For every fixed
\(A>0\), the same estimate at \(t=\pm2A\) gives
uniform bounds on \(\mathbb E e^{2A|W_M|}\), using
\(e^{2A|w|}\le e^{2Aw}+e^{-2Aw}\).
Every subsequential limit therefore has moment
generating function \(e^{ct^2/2}\), by uniform
integrability and weak convergence. This identifies
the limit uniquely as \(N(0,c)\). The same exponential
bound gives uniform integrability of every polynomial,
and hence convergence of all moments.
\end{proof}

This global theorem is not obtained by replacing
all cells by independent copies of their local limits.
It follows from the controlled finite-volume pressure
and the complete compact measure. It concerns a
single centered extensive observable with
square-root-volume normalization.

\subsection{Remaining source and physical-scale distinctions}

The present argument controls a source neighborhood
for the symmetrically assigned bare action density,
with no restriction on support size. Its lower
pressure estimate and central limit use uniform
sources. For nonuniform macroscopic arrays the
chessboard bound is only an upper estimate; it does
not determine their limiting covariance or law.

Neither \(\beta_M^\diamond\) nor either earlier selected
sequence has been identified with an interacting
renormalized physical trajectory. The source window
belongs to the bare Wilson law. Reflection positivity,
conditional-source Hessian bounds, and locality for
the required generated laws have not been established
by this argument. In particular no nonlinear
observation-budget hypothesis of V70 has been silently
supplied.

The cell observable is a scalar quadratic-energy
combination. These results do not give the
physically normalized colour-curvature distribution.
Unrenormalized dimension-four action-density
smearings still carry their ultraviolet volume
factors. A Gaussian limit for \(W_M\) is neither
a construction of an interacting continuum field
nor evidence of a positive physical Hamiltonian gap.

The next useful interface is to quantify how much
these source bounds survive a specified change of
law, with its actual density, normalization and
relative entropy retained. This is consistent with
the companion notes' bookkeeping requirements,
but proving the required entropy or source costs
for an actual RG step remains additional work.

\subsection{Verification and archive record}

The certificate is
\begin{center}\small
\path{scripts/verify_ym_f14_v100_source_window.py}.
\end{center}
It checks 100 interval families with 453 disjoint
selections, 192 discrete cumulative-kernel inequalities,
88 signed source-kernel transformations, 180
window-\(L^1\) ledger identities, twenty formal
Gaussian-limit moment orders, and 27 Chernoff
optimizer identities. The Gaussian reference is a
formal centered gamma calculation used only to
audit the Taylor coefficients; it is not a claim
that the finite Wilson total action has a gamma law.

The certificate contains 6,290 bytes and has SHA--256
\begin{center}\scriptsize\ttfamily
4F1B683033421579ED8737A3F626B504B62F42B24719EB38679A9885001F1E0F.
\end{center}
The maximal inequality, susceptibility convergence,
coupling selection and probability limits are proved
analytically above, not inferred from these finite tests.

The predecessor V99 contains 2,451,009 bytes and has SHA--256
\begin{center}\scriptsize\ttfamily
FBC5984CB133163E665950F54C164054A8B4CAFCC434ADB3CBCB3631CE955E2D.
\end{center}
Its body is unchanged apart from the title version and
this terminal insertion. The resulting complete
hundred-version file is retained as
\begin{center}\small
\path{Renewal_Yang_Mills_Mass_Gap_Research_Notes_V100_f14.tex}.
\end{center}
The next active note is a compact V1, not a deletion
of this accumulated mathematical record. Companion
files, supplied sources, previous archives and earlier
certificates remain unchanged.

\begin{firewall}
V100 proves mean susceptibility flatness, a selected
nonzero bare-action source window, arbitrary-support
Bernstein bounds, and a normalized total-action
central limit. The full interacting four-dimensional
continuum Yang--Mills construction and a positive
physical Hamiltonian mass gap remain unproved.
The original goal remains active through the restart.
\end{firewall}
\end{document}
